% Generated by matins/tools/build_latex.py from data/corpus.json — do not edit.
\documentclass[10pt,twoside,openright]{book}
\usepackage{matins}
\begin{document}
\BodySize
\frontmatter
\thispagestyle{empty}
\vspace*{0.25\textheight}
{\centering
{\Huge\rubr{Lectiones ad Matutinum}}\par\vspace{0.4em}
{\large\itshape ex Breviario Romano anni MCMLIV}\par\vspace{2em}
{\LARGE The Lessons of Matins}\par\vspace{0.4em}
{\large\itshape from the Roman Breviary of 1954, in Latin and English}\par\vspace{3em}
{\Large\rubr{Pars Verna}}\par\vspace{0.3em}{\large\itshape Spring}\par\vspace{2em}
{\small }\par}
\clearpage
\thispagestyle{empty}
{\footnotesize\setlength{\parindent}{0pt}\setlength{\parskip}{0.35em}\raggedright
\par\addvspace{0.6em}\textsc{Divinum Officium}\par
The texts, and the program that arranges them according to the rubrics, are the work of the Divinum Officium Project (divinumofficium.com), begun by the late Laszlo Kiss and continued by the Project's volunteers. Its code and data are available at github.com/DivinumOfficium/divinum-officium under the MIT License.\par
\par\addvspace{0.6em}\textsc{Latin}\par
The Latin text derives from the Ratisbon 1888 edition of the Breviarium Romanum, scanned by the University of St Michael's College, Toronto, amended from the 1906 edition, and corrected by Laszlo Kiss and the Project team against later printings: Pustet (Regensburg, 1943), Benziger Brothers (New York, 1945), Burns Oates \& Washbourne (London, 1946), La Presse Catholique Panaméricaine (Montréal, 1943), and Desclée and Mame (1962).\par
The accented text of the Sixto-Clementine Vulgate was supplied by Frère Romain-Marie of the Abbaye Saint-Joseph de Clairval, Flavigny-sur-Ozerain.\par
\par\addvspace{0.6em}\textsc{English}\par
Holy Scripture is given in the Douay-Rheims translation.\par
The other lessons and the responsories are from the translation of the Roman Breviary by John, Marquess of Bute (edition of 1908), made available by the University of St Michael's College, Toronto.\par
\par\addvspace{0.6em}\textsc{This edition}\par
The arrangement of this edition follows the Divinum Officium program's choice of lessons for each day under the rubrics of 1954. Any errors in the arrangement are the editor's, not the Project's.\par
\par\addvspace{0.6em}\textsc{Typeface}\par
Set in EB Garamond, designed by Georg Duffner and Octavio Pardo, under the SIL Open Font License 1.1.\par
}
\mainmatter
\PartTitle{Proprium de Tempore}{Proper of Time}

\addcontentsline{toc}{chapter}{Proprium de Tempore\enspace\textit{Proper of Time}}

\SeasonTitle{Tempus Quadragesimæ}{Lent}

\addcontentsline{toc}{section}{Tempus Quadragesimæ\enspace\textit{Lent}}

\Office{Dominica I in Quadragesima}{First Sunday of Lent}{Semiduplex I classis}
\addcontentsline{toc}{subsection}{Dominica I in Quadragesima\enspace\textit{First Sunday of Lent}}
\begin{paracol}{2}
\LA
\LessonHead{Lectio I}
\switchcolumn
\EN
\LessonHead{Lesson I}
\switchcolumn*

\LA
\LessonTitle{De Epístola secúnda beáti Pauli Apóstoli ad Corínthios}
\switchcolumn
\EN
\LessonTitle{Lesson from the second letter of St. Paul the Apostle to the Corinthians}
\switchcolumn*

\LA
\Cite{2 Cor 6:1-10}
\switchcolumn
\EN
\Cite{2 Cor 6:1-10}
\switchcolumn*

\LA
\noindent Adjuvántes autem exhortámur ne in vácuum grátiam Dei recipiátis. Ait enim: Témpore accépto exaudívi te, et in die salútis adjúvi te. Ecce nunc tempus acceptábile, ecce nunc dies salútis. Némini dantes ullam offensiónem, ut non vituperétur ministérium nostrum: sed in ómnibus exhibeámus nosmetípsos sicut Dei minístros in multa patiéntia, in tribulatiónibus, in necessitátibus, in angústiis, in plagis, in carcéribus, in seditiónibus, in labóribus, in vigíliis, in jejúniis, in castitáte, in sciéntia, in longanimitáte, in suavitáte, in Spíritu Sancto, in caritáte non ficta, in verbo veritátis, in virtúte Dei, per arma justítiæ a dextris et a sinístris, per glóriam, et ignobilitátem, per infámiam, et bonam famam: ut seductóres, et veráces, sicut qui ignóti, et cógniti: quasi moriéntes, et ecce vívimus: ut castigáti, et non mortificáti: quasi tristes, semper autem gaudéntes: sicut egéntes, multos autem locupletántes: tamquam nihil habéntes, et ómnia possidéntes.\par
\switchcolumn
\EN
\noindent And we helping do exhort you, that you receive not the grace of God in vain. For he saith: In an accepted time have I heard thee; and in the day of salvation have I helped thee. Behold, now is the acceptable time; behold, now is the day of salvation. Giving no offence to any man, that our ministry be not blamed: But in all things let us exhibit ourselves as the ministers of God, in much patience, in tribulation, in necessities, in distresses, In stripes, in prisons, in seditions, in labours, in watchings, in fastings, In chastity, in knowledge, in longsuffering, in sweetness, in the Holy Ghost, in charity unfeigned, In the word of truth, in the power of God; by the armour of justice on the right hand and on the left; By honour and dishonour, by evil report and good report; as deceivers, and yet true; as unknown, and yet known; As dying, and behold we live; as chastised, and not killed; As sorrowful, yet always rejoicing; as needy, yet enriching many; as having nothing, and possessing all things.\par
\switchcolumn*

\LA
\begin{Resp}\Rsign Ecce nunc tempus acceptábile, ecce nunc dies salútis: commendémus nosmetípsos in multa patiéntia, in jejúniis multis, \Star{} Per arma justítiæ virtútis Dei.\end{Resp}
\begin{Resp}\Vsign In ómnibus exhibeámus nosmetípsos sicut Dei minístros, in multa patiéntia, in jejúniis multis.\end{Resp}
\begin{Resp}\Rsign Per arma justítiæ virtútis Dei.\end{Resp}
\switchcolumn
\EN
\begin{Resp}\Rsign Behold, now is the acceptable time; behold, now is the day of salvation: let us approve ourselves in much patience, in much fasting; \Star{} In the power of God, by the armour of righteousness.\end{Resp}
\begin{Resp}\Vsign In all things let us approve ourselves as the ministers of God, in much patience, in much fasting.\end{Resp}
\begin{Resp}\Rsign In the power of God, by the armour of righteousness.\end{Resp}
\switchcolumn*

\LA
\LessonHead{Lectio II}
\switchcolumn
\EN
\LessonHead{Lesson II}
\switchcolumn*

\LA
\Cite{2 Cor 6:11-16}
\switchcolumn
\EN
\Cite{2 Cor 6:11-16}
\switchcolumn*

\LA
\noindent Os nostrum patet ad vos, o Corínthii; cor nostrum dilatátum est. Non angustiámini in nobis: angustiámini autem in viscéribus vestris: eándem autem habéntes remuneratiónem, tamquam fíliis dico, dilatámini et vos. Nolíte jugum dúcere cum infidélibus. Quæ enim participátio justítiæ cum iniquitáte? aut quæ socíetas luci ad ténebras? quæ autem convéntio Christi ad Bélial? aut quæ pars fidéli cum infidéli? qui autem consénsus templo Dei cum idólis? vos enim estis templum Dei vivi, sicut dicit Deus: Quóniam inhabitábo in illis, et inambulábo inter eos, et ero illórum Deus, et ipsi erunt mihi pópulus.\par
\switchcolumn
\EN
\noindent Our mouth is open to you, O ye Corinthians, our heart is enlarged. You are not straitened in us, but in your own bowels you are straitened. But having the same recompense, (I speak as to my children,) be you also enlarged. Bear not the yoke with unbelievers. For what participation hath justice with injustice? Or what fellowship hath light with darkness? And what concord hath Christ with Belial? Or what part hath the faithful with the unbeliever? And what agreement hath the temple of God with idols? For you are the temple of the living God; as God saith: I will dwell in them, and walk among them; and I will be their God, and they shall be my people.\par
\switchcolumn*

\LA
\begin{Resp}\Rsign In ómnibus exhibeámus nosmetípsos sicut Dei minístros in multa patiéntia: \Star{} Ut non vituperétur ministérium nostrum.\end{Resp}
\begin{Resp}\Vsign Ecce nunc tempus acceptábile, ecce nunc dies salútis: commendémus nosmetípsos in multa patiéntia.\end{Resp}
\begin{Resp}\Rsign Ut non vituperétur ministérium nostrum.\end{Resp}
\switchcolumn
\EN
\begin{Resp}\Rsign In all things let us approve ourselves as the ministers of God, in much patience; \Star{} That our ministry be not blamed.\end{Resp}
\begin{Resp}\Vsign Behold, now is the acceptable time; behold, now is the day of salvation: let us approve ourselves in much patience.\end{Resp}
\begin{Resp}\Rsign That our ministry be not blamed.\end{Resp}
\switchcolumn*

\LA
\LessonHead{Lectio III}
\switchcolumn
\EN
\LessonHead{Lesson III}
\switchcolumn*

\LA
\Cite{2 Cor 7:4-9}
\switchcolumn
\EN
\Cite{2 Cor 7:4-9}
\switchcolumn*

\LA
\noindent Replétus sum consolatióne; superabúndo gáudio in omni tribulatióne nostra. Nam et cum venissémus in Macedóniam, nullam réquiem hábuit caro nostra, sed omnem tribulatiónem passi sumus: foris pugnæ, intus timóres. Sed qui consolátur húmiles, consolátus est nos Deus in advéntu Titi. Non solum autem in advéntu ejus, sed étiam in consolatióne, qua consolátus est in vobis, réferens nobis vestrum desidérium, vestrum fletum, vestram æmulatiónem pro me, ita ut magis gaudérem. Quóniam etsi contristávi vos in epístola, non me pœ́nitet: etsi pœnitéret, videns quod epístola illa (etsi ad horam) vos contristávit, nunc gáudeo: non quia contristáti estis, sed quia contristáti estis ad pœniténtiam.\par
\switchcolumn
\EN
\noindent I am filled with comfort: I exceedingly abound with joy in all our tribulation. For also when we were come into Macedonia, our flesh had no rest, but we suffered all tribulation; combats without, fears within. But God, who comforteth the humble, comforted us by the coming of Titus. And not by his coming only, but also by the consolation, wherewith he was comforted in you, relating to us your desire, your mourning, your zeal for me, so that I rejoiced the more. For although I made you sorrowful by my epistle, I do not repent; and if I did repent, seeing that the same epistle (although but for a time) did make you sorrowful; Now I am glad: not because you were made sorrowful; but because you were made sorrowful unto penance.\par
\switchcolumn*

\LA
\begin{Resp}\Rsign In jejúnio et fletu orábunt sacerdótes, dicéntes: \Star{} Parce, Dómine, parce pópulo tuo; et ne des hereditátem tuam in perditiónem.\end{Resp}
\begin{Resp}\Vsign Inter vestíbulum et altáre plorábunt sacerdótes, dicéntes.\end{Resp}
\begin{Resp}\Rsign Parce, Dómine, parce pópulo tuo; et ne des hereditátem tuam in perditiónem.\end{Resp}
\begin{Resp}\Vsign Glória Patri, et Fílio, \Star{} et Spirítui Sancto.\end{Resp}
\begin{Resp}\Rsign Parce, Dómine, parce pópulo tuo; et ne des hereditátem tuam in perditiónem.\end{Resp}
\switchcolumn
\EN
\begin{Resp}\Rsign The Priests shall pray, with fasting and with weeping, and shall say: \Star{} Spare, O Lord, spare thy people; and give not thine heritage to destruction.\end{Resp}
\begin{Resp}\Vsign The Priests shall weep between the porch and the altar, and shall say:\end{Resp}
\begin{Resp}\Rsign Spare, O Lord, spare thy people; and give not thine heritage to destruction.\end{Resp}
\begin{Resp}\Vsign Glory be to the Father, and to the Son, \Star{} and to the Holy Ghost.\end{Resp}
\begin{Resp}\Rsign Spare, O Lord, spare thy people; and give not thine heritage to destruction.\end{Resp}
\switchcolumn*

\LA
\LessonHead{Lectio IV}
\switchcolumn
\EN
\LessonHead{Lesson IV}
\switchcolumn*

\LA
\LessonTitle{Sermo sancti Leónis Papæ}
\switchcolumn
\EN
\LessonTitle{From the Sermons of Pope St. Leo (the Great.)}
\switchcolumn*

\LA
\Source{Sermo 4 de Quadragesima}
\switchcolumn
\EN
\Source{4th on Lent}
\switchcolumn*

\LA
\noindent Prædicatúrus vobis, dilectíssimi, sacratíssimum maximúmque jejúnium, quo áptius utar exórdio, quam ut verbis Apóstoli, in quo Christus loquebátur, incípiam, dicámque quod lectum est: Ecce nunc tempus acceptábile: ecce nunc dies salútis? Quamvis enim nulla sint témpora, quæ divínis non sint plena munéribus, et semper nobis ad misericórdiam Dei per ipsíus grátiam præstétur accéssus: nunc tamen ómnium mentes majóri stúdio ad spiritáles proféctus movéri, et amplióri fidúcia opórtet animári, quando ad univérsa pietátis offícia, illíus nos diéi, in quo redémpti sumus, recúrsus invítat: ut excéllens super ómnia passiónis Domínicæ sacraméntum, purificátis et corpóribus et ánimis celebrémus.\par
\switchcolumn
\EN
\noindent Dearly beloved brethren, I am to preach to you the holiest and the greatest of Fasts; and with what words can I more fitly begin than with those words of the Apostle, in whom Christ spake, which have just been read? Behold, now is the acceptable time! Behold, now is the day of salvation! It is true that there are no times which are not rich with God's gifts; His grace doth ever give us an entry unto His mercy; nevertheless, more especially at this time doth it behove that the minds of all men be earnestly stirred up to make progress in things spiritual, and to be nerved by a trust in God stronger than ever; for now the anniversary of that day on which we were redeemed is drawing near, and thereby moving us to work all godliness, to the end that we may be able to celebrate, with clean minds and bodies, that mystery which exceedeth all others, the mystery of the Lord's sufferings.\par
\switchcolumn*

\LA
\begin{Resp}\Rsign Emendémus in mélius, quæ ignoránter peccávimus: ne súbito præoccupáti die mortis, quærámus spátium pœniténtiæ, et inveníre non possímus: \Star{} Atténde, Dómine, et miserére, quia peccávimus tibi.\end{Resp}
\begin{Resp}\Vsign Adjuva nos, Deus salutáris noster, et propter honórem nóminis tui, Dómine, líbera nos.\end{Resp}
\begin{Resp}\Rsign Atténde, Dómine, et miserére, quia peccávimus tibi.\end{Resp}
\switchcolumn
\EN
\begin{Resp}\Rsign Let us amend for the better in that wherein we have sinned unknowingly, or ever the day of death suddenly prevent us, and we seek a place of repentance, and find none. \Star{} Give heed, O Lord, and have mercy upon us, for we have sinned against thee.\end{Resp}
\begin{Resp}\Vsign Help us, O God of our salvation, and for the glory of thy Name deliver us, O Lord.\end{Resp}
\begin{Resp}\Rsign Give heed, O Lord, and have mercy upon us, for we have sinned against thee.\end{Resp}
\switchcolumn*

\LA
\LessonHead{Lectio V}
\switchcolumn
\EN
\LessonHead{Lesson V}
\switchcolumn*

\LA
\noindent Debebátur quidem tantis mystériis ita incessábilis devótio, et continuáta reveréntia, ut tales permanerémus in conspéctu Dei, quales nos in ipso Pascháli festo dignum est inveníri. Sed quia hæc fortitúdo paucórum est: et dum carnis fragilitáte austérior observántia relaxátur, dumque per várias actiónes vitæ hujus sollicitúdo disténditur, necésse est de mundáno púlvere étiam religiósa corda sordéscere: magna divínæ institutiónis salubritáte provísum est, ut ad reparándam méntium puritátem quadragínta nobis diérum exercitátio mederétur, in quibus aliórum témporum culpas, et pia ópera redímerent, et jejúnia casta decóquerent.\par
\switchcolumn
\EN
\noindent Mysteries so great demand sustained earnestness, and continuous worship, if we would ever abide in the sight of God, such as it is meet that He should find us on the Feast of the Passover. But since few have the strength to do thus, and the frailty of the body rebelleth against such hardness, while the diverse actions of this life distract us with their cares, it necessarily befalleth that the dust of earth befouleth the hearts even of the godly. To meet this befoulment therefore, and to restore the cleanness of our souls, it is provided by the healthful institution of God, that we should be purged by an exercise of forty days, wherein godly works may redeem the mis-spending of our other time, and purifying fasts rid us of the same.\par
\switchcolumn*

\LA
\begin{Resp}\Rsign Derelínquat ímpius viam suam, et vir iníquus cogitatiónes suas, et revertátur ad Dóminum, et miserébitur ejus: \Star{} Quia benígnus et miséricors est, et præstábilis super malítia Dóminus Deus noster.\end{Resp}
\begin{Resp}\Vsign Non vult Dóminus mortem peccatóris, sed ut convertátur et vivat.\end{Resp}
\begin{Resp}\Rsign Quia benígnus et miséricors est, et præstábilis super malítia Dóminus Deus noster.\end{Resp}
\switchcolumn
\EN
\begin{Resp}\Rsign Let the wicked forsake his way, and the unrighteous man his thoughts, and let him return unto the Lord, and He will have mercy upon him; \Star{} For the Lord our God is gracious, and merciful, and repenteth Him of the evil.\end{Resp}
\begin{Resp}\Vsign The Lord hath no pleasure in the death of the wicked; but that he turn from his way and live.\end{Resp}
\begin{Resp}\Rsign For the Lord our God is gracious, and merciful, and repenteth Him of the evil.\end{Resp}
\switchcolumn*

\LA
\LessonHead{Lectio VI}
\switchcolumn
\EN
\LessonHead{Lesson VI}
\switchcolumn*

\LA
\noindent Ingressúri ígitur, dilectíssimi, dies mýsticos, et purificándis ánimis atque corpóribus sacrátius institútos, præcéptis apostólicis obedíre curémus, emundántes nos ab omni inquinaménto carnis ac spíritus: ut castigátis colluctatiónibus, quæ sunt inter utrámque substántiam, ánimus, quem sub Dei gubernáculis constitútum, córporis sui decet esse rectórem, dominatiónis suæ obtíneat dignitátem: ut némini dantes ullam offensiónem, vituperatiónibus obloquéntium non simus obnóxii. Digna enim ab infidélibus reprehensióne carpémur, et nostro vítio linguæ ímpiæ in injúriam se religiónis armábunt, si jejunántium mores a puritáte perféctæ continéntiæ discrepárint. Non enim in sola abstinéntia cibi stat nostri summa jejúnii: aut fructuóse córpori esca subtráhitur, nisi mens ab iniquitáte revocétur.\par
\switchcolumn
\EN
\noindent Therefore, dearly beloved brethren, as we are now about to enter upon these mystic days, the end of whose most holy ordinance is the cleansing both of our souls and bodies, let us take heed that we be obedient unto the command of the Apostle, putting far away from us every defilement of flesh and spirit, ordering the strife which there is between the two substances whereof we are compounded; that the soul, which is ordained under the rule of God, and which it beseemeth under His rule to rule the body, may enjoy the fulness of her lordship; giving no offence to any so that we may give no cause to such as revile us. For if our ways during the Fast agree not with the purity of perfect temperance, the reproaches of the unbelievers will be just, and our sins will arm the tongues of the ungodly to the harming of our religion. The sum of our Fast standeth not only in abstaining from meats; neither is it profitable to deny food to the body, if the mind be not bridled from iniquity.\par
\switchcolumn*

\LA
\begin{Resp}\Rsign Paradísi portas apéruit nobis jejúnii tempus: suscipiámus illud orántes, et deprecántes: \Star{} Ut in die resurrectiónis cum Dómino gloriémur.\end{Resp}
\begin{Resp}\Vsign In ómnibus exhibeámus nosmetípsos sicut Dei minístros in multa patiéntia.\end{Resp}
\begin{Resp}\Rsign Ut in die resurrectiónis cum Dómino gloriémur.\end{Resp}
\begin{Resp}\Vsign Glória Patri, et Fílio, \Star{} et Spirítui Sancto.\end{Resp}
\begin{Resp}\Rsign Ut in die resurrectiónis cum Dómino gloriémur.\end{Resp}
\switchcolumn
\EN
\begin{Resp}\Rsign The season of the Fast openeth unto us the gates of heaven; let us enter thereon in prayer and supplication, \Star{} That on the day when the Lord riseth again we may rejoice with Him.\end{Resp}
\begin{Resp}\Vsign In all things let us approve ourselves the ministers of God, in much patience.\end{Resp}
\begin{Resp}\Rsign That on the day when the Lord riseth again we may rejoice with Him.\end{Resp}
\begin{Resp}\Vsign Glory be to the Father, and to the Son, \Star{} and to the Holy Ghost.\end{Resp}
\begin{Resp}\Rsign That on the day when the Lord riseth again we may rejoice with Him.\end{Resp}
\switchcolumn*

\LA
\LessonHead{Lectio VII}
\switchcolumn
\EN
\LessonHead{Lesson VII}
\switchcolumn*

\LA
\LessonTitle{Léctio sancti Evangélii secúndum Matthǽum}
\switchcolumn
\EN
\LessonTitle{From the Holy Gospel according to Matthew}
\switchcolumn*

\LA
\Cite{Matt 4:1-11}
\switchcolumn
\EN
\Cite{Matt 4:1-11}
\switchcolumn*

\LA
\noindent In illo témpore: Ductus est Jesus in desértum a Spíritu, ut tentarétur a diábolo. Et cum jejunásset quadragínta diébus et quadragínta nóctibus, póstea esúriit. Et réliqua.\par
\switchcolumn
\EN
\noindent In that time, Jesus was led by the spirit into the desert, to be tempted by the devil. And when he had fasted forty days and forty nights, afterwards he was hungry. And so on.\par
\switchcolumn*

\LA
\LessonTitle{Homilía sancti Gregórii Papæ}
\switchcolumn
\EN
\LessonTitle{Homily by Pope St. Gregory (the Great.)}
\switchcolumn*

\LA
\Source{Homilia 16 in Evangelia}
\switchcolumn
\EN
\Source{16th on the Gospels}
\switchcolumn*

\LA
\noindent Dubitári a quibúsdam solet, a quo spíritu sit Jesus ductus in desértum, propter hoc quod súbditur: Assúmpsit eum diábolus in sanctam civitátem: et rursum: Assúmpsit eum in montem excélsum valde. Sed vere et absque ulla quæstióne conveniénter accípitur, ut a Sancto Spíritu in desértum ductus credátur: ut illuc eum suus Spíritus dúceret, ubi hunc ad tentándum malígnus spíritus inveníret. Sed ecce cum dícitur Deus homo vel in excélsum montem, vel in sanctam civitátem a diábolo assúmptus, mens réfugit crédere, humánæ hoc audíre aures expavéscunt. Qui tamen non esse incredibília ista cognóscimus, si in illo et ália facta pensámus.\par
\switchcolumn
\EN
\noindent Some persons are accustomed to question what Spirit it was of which Jesus was led up into the wilderness, on account of the words a little farther on Then the devil taketh Him up into the holy city and again The devil taketh Him up into an exceeding high mountain. But in truth, and without any searching, we may very fitly take it that we are to believe it was the Holy Ghost Who led Him up into the wilderness; His own Spirit led Him where the evil spirit found Him to tempt Him. When however it is said that He, God and man, was taken up by the devil either into an exceeding high mountain or into the holy city, the mind shrinketh from believing, and the ears of man tingle to hear it. Yet these things we know not to be incredible, when we consider certain other things concerning Him.\par
\switchcolumn*

\LA
\begin{Resp}\Rsign Scíndite corda vestra, et non vestiménta vestra: et convertímini ad Dóminum Deum vestrum: \Star{} Quia benígnus et miséricors est.\end{Resp}
\begin{Resp}\Vsign Derelínquat ímpius viam suam, et vir iníquus cogitatiónes suas, et revertátur ad Dóminum, et miserébitur ejus.\end{Resp}
\begin{Resp}\Rsign Quia benígnus et miséricors est.\end{Resp}
\switchcolumn
\EN
\begin{Resp}\Rsign Rend your hearts and not your garments, and turn unto the Lord your God; \Star{} For He is gracious and merciful.\end{Resp}
\begin{Resp}\Vsign Let the wicked forsake his way, and the unrighteous man his thoughts, and let him return unto the Lord, and He will have mercy upon him.\end{Resp}
\begin{Resp}\Rsign For He is gracious and merciful.\end{Resp}
\switchcolumn*

\LA
\LessonHead{Lectio VIII}
\switchcolumn
\EN
\LessonHead{Lesson VIII}
\switchcolumn*

\LA
\noindent Certe iniquórum ómnium caput diábolus est: et hujus cápitis membra sunt omnes iníqui. An non diáboli membrum fuit Pilátus? an non diáboli membra Judǽi persequéntes, et mílites crucifigéntes Christum fuérunt? Quid ergo mirum, si se ab illo permísit in montem duci, qui se pértulit étiam a membris illíus crucifígi? Non est ergo indígnum Redemptóri nostro quod tentári vóluit, qui vénerat occídi. Justum quippe erat, ut sic tentatiónes nostras suis tentatiónibus vínceret, sicut mortem nostram vénerat sua morte superáre.\par
\switchcolumn
\EN
\noindent In truth, the devil is the head of all the wicked, and every wicked man is a member of this body, of which the devil is the head. Was not Pilate a limb of Satan? Were not the Jews that persecuted, and the soldiers that crucified Christ, likewise limbs of Satan? Is it then strange that He should allow Himself to be led up into a mountain by the head, Who allowed Himself to be crucified by the members? Therefore it is not unworthy of our Redeemer, Who came to be slain, that He was willing to be tempted. It was meet that He should thus overcome our temptations by His own, even as He came to conquer our death by His own.\par
\switchcolumn*

\LA
\begin{Resp}\Rsign Frange esuriénti panem tuum, et egénos vagósque induc in domum tuam: \Star{} Tunc erúmpet quasi mane lumen tuum, et anteíbit fáciem tuam justítia tua.\end{Resp}
\begin{Resp}\Vsign Cum víderis nudum, óperi eum, et carnem tuam ne despéxeris.\end{Resp}
\begin{Resp}\Rsign Tunc erúmpet quasi mane lumen tuum, et anteíbit fáciem tuam justítia tua.\end{Resp}
\switchcolumn
\EN
\begin{Resp}\Rsign Deal thy bread to the hungry, and bring the poor and the wanderer to thine house. \Star{} Then shall thy light break forth as the morning, and thy righteousness shall go before thee.\end{Resp}
\begin{Resp}\Vsign When thou seest the naked, cover him; and hide not thyself from thine own flesh.\end{Resp}
\begin{Resp}\Rsign Then shall thy light break forth as the morning, and thy righteousness shall go before thee.\end{Resp}
\switchcolumn*

\LA
\LessonHead{Lectio IX}
\switchcolumn
\EN
\LessonHead{Lesson IX}
\switchcolumn*

\LA
\noindent Sed sciéndum nobis est, quia tribus modis tentátio ágitur: suggestióne, delectatióne et consénsu. Et nos cum tentámur, plerúmque in delectatiónem, aut étiam in consénsum lábimur: quia de carnis peccáto propagáti, in nobis ipsis étiam gérimus, unde certámina tolerámus. Deus vero, qui in útero Vírginis incarnátus, in mundum sine peccáto vénerat, nihil contradictiónis in semetípso tolerábat. Tentári ergo per suggestiónem pótuit: sed ejus mentem peccáti delectátio non momórdit. Atque ídeo omnis diabólica illa tentátio foris, non intus fuit.\par
\switchcolumn
\EN
\noindent We ought to know that temptation worketh through three forms. There is, first, the suggestion; then the delectation; lastly, the consent. When we are tempted, it often happeneth that we fall into delectation, and even into consent, because in the sinful flesh of which we are begotten, we carry in ourselves matter to favor the attack. But God, when He took Flesh in the womb of the Virgin, and came into the world without sin, did so without having in Himself anything of this lusting of the flesh against the spirit. It was possible therefore for Him to be tempted in the first stage, namely suggestion; but there was nothing in His Mind in which delectation could fix its teeth. And thus all the temptation which He endured from the devil was without, and none within Him.\par
\switchcolumn*

\LA
\begin{Resp}\Rsign Angelis suis Deus mandávit de te, ut custódiant te in ómnibus viis tuis: \Star{} In mánibus portábunt te, ne umquam offéndas ad lápidem pedem tuum.\end{Resp}
\begin{Resp}\Vsign Super áspidem et basilíscum ambulábis, et conculcábis leónem et dracónem.\end{Resp}
\begin{Resp}\Rsign In mánibus portábunt te, ne umquam offéndas ad lápidem pedem tuum.\end{Resp}
\begin{Resp}\Vsign Glória Patri, et Fílio, \Star{} et Spirítui Sancto.\end{Resp}
\begin{Resp}\Rsign In mánibus portábunt te, ne umquam offéndas ad lápidem pedem tuum.\end{Resp}
\switchcolumn
\EN
\begin{Resp}\Rsign God hath given His Angels charge over thee, to keep thee in all thy ways. \Star{} They shall bear thee up in their hands, lest haply thou dash thy foot against a stone.\end{Resp}
\begin{Resp}\Vsign Thou shalt tread upon the adder and the cockatrice, the lion also, and the dragon shalt thou trample under feet.\end{Resp}
\begin{Resp}\Rsign They shall bear thee up in their hands, lest haply thou dash thy foot against a stone.\end{Resp}
\begin{Resp}\Vsign Glory be to the Father, and to the Son, \Star{} and to the Holy Ghost.\end{Resp}
\begin{Resp}\Rsign They shall bear thee up in their hands, lest haply thou dash thy foot against a stone.\end{Resp}
\switchcolumn*

\end{paracol}

\Office{Feria Secunda infra Hebdomadam I in Quadragesima}{Monday of the First Week of Lent}{Feria major}
\addcontentsline{toc}{subsection}{Feria Secunda infra Hebdomadam I in Quadragesima\enspace\textit{Monday of the First Week of Lent}}
\begin{paracol}{2}
\LA
\LessonHead{Lectio I}
\switchcolumn
\EN
\LessonHead{Lesson I}
\switchcolumn*

\LA
\LessonTitle{Léctio sancti Evangélii secúndum Matthǽum}
\switchcolumn
\EN
\LessonTitle{Continuation of the Holy Gospel according to Matthew}
\switchcolumn*

\LA
\Cite{Matt 25:31-46}
\switchcolumn
\EN
\Cite{Matt 25:31-46}
\switchcolumn*

\LA
\noindent In illo témpore: Dixit Jesus discípulis suis: Cum vénerit Fílius hóminis in majestáte sua, et omnes Angeli cum eo, tunc sedébit super sedem majestátis suæ: et congregabúntur ante eum omnes gentes. Et réliqua.\par
\switchcolumn
\EN
\noindent In that time, Jesus said to his disciples: When the Son of man shall come in his majesty, and all the angels with him, then shall he sit upon the seat of his majesty. And all nations shall be gathered together before him. And so on.\par
\switchcolumn*

\LA
\LessonTitle{Homilía sancti Augustíni Epíscopi}
\switchcolumn
\EN
\LessonTitle{Homily by St. Augustine, Bishop (of Hippo.)}
\switchcolumn*

\LA
\Source{Lib. de fide et operibus. cap. 15 tom. 4, circa medium}
\switchcolumn
\EN
\Source{On Faith and Works, xv, 4}
\switchcolumn*

\LA
\noindent Si mandátis non servátis, ad vitam veníri potest per solam fidem, quæ sine opéribus mórtua est: illud deínde quómodo verum erit, quod eis, quos ad sinístram positúrus est, dicet: Ite in ignem ætérnum, qui parátus est diábolo, et ángelis ejus: nec íncrepat, quia in eum non credidérunt; sed quia bona ópera non fecérunt? Nam profécto, ne sibi quisquam de fide, quæ sine opéribus mórtua est, promíttat ætérnam vitam; proptérea omnes gentes segregatúrum se dixit, quæ permíxtæ eisdem páscuis utebántur: ut appáreat, eos illi dictúros: Dómine, quando te vídimus illa et illa patiéntem, et non ministrávimus tibi? qui in eum credíderant, sed bona operári non curáverant, tamquam de ipsa fide mórtua ad vitam pervenirétur ætérnam.\par
\switchcolumn
\EN
\noindent If, without keeping the commandments, it be possible to attain unto life by faith only, (and faith, if it hath not works, is dead, James ii. 17,) how can it be true that the Lord will say to such as He shall have set on His left hand Depart from Me, ye cursed, into everlasting fire, prepared for the devil and his angels? He rebuketh them, not because they have not believed in Him, but because they have not wrought good works. Yea, lest any man should promise himself life eternal by faith only, (and faith, if it hath not works, is dead,) the Lord saith that He will gather together all nations, nations who have lived mingled together in the same countries, that we may seem to hear them which have believed indeed in Him, but have not wrought good works, (as though that their dead faith could, being alone, lead them into life eternal,) that we may seem to hear such crying unto Him, Lord, when saw we thee suffering such and such things, and did not minister unto thee?\par
\switchcolumn*

\LA
\begin{Resp}\Rsign Ecce nunc tempus acceptábile, ecce nunc dies salútis: commendémus nosmetípsos in multa patiéntia, in jejúniis multis, \Star{} Per arma justítiæ virtútis Dei.\end{Resp}
\begin{Resp}\Vsign In ómnibus exhibeámus nosmetípsos sicut Dei minístros, in multa patiéntia, in jejúniis multis.\end{Resp}
\begin{Resp}\Rsign Per arma justítiæ virtútis Dei.\end{Resp}
\switchcolumn
\EN
\begin{Resp}\Rsign Behold, now is the acceptable time; behold, now is the day of salvation: let us approve ourselves in much patience, in much fasting; \Star{} In the power of God, by the armour of righteousness.\end{Resp}
\begin{Resp}\Vsign In all things let us approve ourselves as the ministers of God, in much patience, in much fasting.\end{Resp}
\begin{Resp}\Rsign In the power of God, by the armour of righteousness.\end{Resp}
\switchcolumn*

\LA
\LessonHead{Lectio II}
\switchcolumn
\EN
\LessonHead{Lesson II}
\switchcolumn*

\LA
\noindent An forte ibunt in ignem ætérnum, qui ópera misericórdiæ non fecérunt: et non ibunt, qui aliéna rapuérunt? vel corrumpéndo in se templum Dei, in seípsos immisericórdes fuérunt: quasi ópera misericórdiæ prosint áliquid sine dilectióne, dicénte Apóstolo: Si distríbuam ómnia mea paupéribus, caritátem autem non hábeam, nihil mihi prodest? Aut díligat quisquam próximum sicut seípsum, qui non díligit seípsum? Qui enim díligit iniquitátem, odit ánimam suam.\par
\switchcolumn
\EN
\noindent If they shall go into everlasting fire who have not done works of mercy, shall not they go who have taken their neighbour's goods? Or shall not they go who have outraged the temple of God in their own selves, and so been merciless to themselves? As if works of mercy could avail anything without love, contrary to the words of the Apostle Though I bestow all my goods to feed the poor, and have not charity, it profiteth me nothing. (1 Cor. xiii. 3.) And what manner of love to his neighbour hath he who loveth him as himself and loveth not himself? remembering that he that loveth iniquity hateth his own soul. (Ps. x. 6.)\par
\switchcolumn*

\LA
\begin{Resp}\Rsign In ómnibus exhibeámus nosmetípsos sicut Dei minístros in multa patiéntia: \Star{} Ut non vituperétur ministérium nostrum.\end{Resp}
\begin{Resp}\Vsign Ecce nunc tempus acceptábile, ecce nunc dies salútis: commendémus nosmetípsos in multa patiéntia.\end{Resp}
\begin{Resp}\Rsign Ut non vituperétur ministérium nostrum.\end{Resp}
\switchcolumn
\EN
\begin{Resp}\Rsign In all things let us approve ourselves as the ministers of God, in much patience; \Star{} That our ministry be not blamed.\end{Resp}
\begin{Resp}\Vsign Behold, now is the acceptable time; behold, now is the day of salvation: let us approve ourselves in much patience.\end{Resp}
\begin{Resp}\Rsign That our ministry be not blamed.\end{Resp}
\switchcolumn*

\LA
\LessonHead{Lectio III}
\switchcolumn
\EN
\LessonHead{Lesson III}
\switchcolumn*

\LA
\noindent Neque illud dici hic póterit, in quo nonnúlli seípsos sedúcunt, ignem ætérnum dictum, non ipsam combustiónem ætérnam. Per ignem quippe, qui ætérnus erit, transitúros arbitrántur eos, quibus propter fidem mórtuam per ignem promíttunt salútem: ut vidélicet ipse ignis ætérnus sit, combústio vero eórum, hoc est, operátio ignis, non sit in eos ætérna: cum et hoc prǽvidens Dóminus senténtiam suam conclúsit ita dicens: Sic ibunt illi in combustiónem ætérnam, justi autem in vitam ætérnam. Erit ergo ætérna combústio, sicut ignis: et eos in illam itúros Véritas dicit, quorum non fidem, sed bona ópera defuísse declarávit.\par
\switchcolumn
\EN
\noindent Neither dare we say here that by which some delude themselves, namely, that the fire indeed is everlasting, but that they will not burn therein everlastingly. Such men say that they whose faith is dead, will pass through that everlasting fire, and that they are they to whom it is promised that they themselves shall be saved, yet so as by fire. (1 Cor. iii. 15.) So that, though the fire itself be everlasting, the burning of the damned therein, that is, the work of the fire upon them, will not be everlasting. As though the Lord were answering this beforehand, the last words of His Sermon are And these shall go away into everlasting punishment, but the righteous into life eternal. As the fire, so shall the burning be; and the Truth biddeth us know that they shall burn therein, who have lacked, not faith, but good works.\par
\switchcolumn*

\LA
\begin{Resp}\Rsign In jejúnio et fletu orábunt sacerdótes, dicéntes: \Star{} Parce, Dómine, parce pópulo tuo; et ne des hereditátem tuam in perditiónem.\end{Resp}
\begin{Resp}\Vsign Inter vestíbulum et altáre plorábunt sacerdótes, dicéntes.\end{Resp}
\begin{Resp}\Rsign Parce, Dómine, parce pópulo tuo; et ne des hereditátem tuam in perditiónem.\end{Resp}
\begin{Resp}\Vsign Glória Patri, et Fílio, \Star{} et Spirítui Sancto.\end{Resp}
\begin{Resp}\Rsign Parce, Dómine, parce pópulo tuo; et ne des hereditátem tuam in perditiónem.\end{Resp}
\switchcolumn
\EN
\begin{Resp}\Rsign The Priests shall pray, with fasting and with weeping, and shall say: \Star{} Spare, O Lord, spare thy people; and give not thine heritage to destruction.\end{Resp}
\begin{Resp}\Vsign The Priests shall weep between the porch and the altar, and shall say:\end{Resp}
\begin{Resp}\Rsign Spare, O Lord, spare thy people; and give not thine heritage to destruction.\end{Resp}
\begin{Resp}\Vsign Glory be to the Father, and to the Son, \Star{} and to the Holy Ghost.\end{Resp}
\begin{Resp}\Rsign Spare, O Lord, spare thy people; and give not thine heritage to destruction.\end{Resp}
\switchcolumn*

\end{paracol}

\Office{Feria Tertia infra Hebdomadam I in Quadragesima}{Tuesday of the First Week of Lent}{Feria major}
\addcontentsline{toc}{subsection}{Feria Tertia infra Hebdomadam I in Quadragesima\enspace\textit{Tuesday of the First Week of Lent}}
\begin{paracol}{2}
\LA
\LessonHead{Lectio I}
\switchcolumn
\EN
\LessonHead{Lesson I}
\switchcolumn*

\LA
\LessonTitle{Léctio sancti Evangélii secúndum Matthǽum}
\switchcolumn
\EN
\LessonTitle{Continuation of the Holy Gospel according to Matthew}
\switchcolumn*

\LA
\Cite{Matt 21:10-17}
\switchcolumn
\EN
\Cite{Matt 21:10-17}
\switchcolumn*

\LA
\noindent In illo témpore: Cum intrásset Jesus Jerosólymam, commóta est univérsa cívitas, dicens: Quis est hic? Et réliqua.\par
\switchcolumn
\EN
\noindent In that time, when Jesus was come into Jerusalem, the whole city was moved, saying: Who is this? And so on.\par
\switchcolumn*

\LA
\LessonTitle{Homilía sancti Bedæ Venerábilis Presbýteri}
\switchcolumn
\EN
\LessonTitle{Homily by the Venerable Bede, Priest (at Jarrow.)}
\switchcolumn*

\LA
\Source{Homil. 7 in Quadrag. tom. 7}
\switchcolumn
\EN
\Source{11th for Lent, Tom. vii}
\switchcolumn*

\LA
\noindent Quod maledicéndo ficum infructuósam per figúram fecit Dóminus, hoc idem mox apértius osténdit, eiciéndo ímprobos e templo. Neque enim áliquid peccávit arbor, quod esuriénte Dómino poma non hábuit, quorum necdum tempus advénerat: sed peccavére sacerdótes, qui in domo Dómini negótia sæculária gerébant, et fructum pietátis, quem debúerant, quemque in eis Dóminus esuriébat, ferre superséderant. Arefécit Dóminus árborem maledícto, ut hómines hæc vidéntes, sive audiéntes, multo magis intellégerent sese divíno condemnándos esse judício, si absque óperum fructu, de plausu tantum sibi religiósi sermónis, velut de sónitu et teguménto blandiréntur viridántium foliórum.\par
\switchcolumn
\EN
\noindent The same thing which the Lord showed in a figure by cursing the barren fig-tree, He afterwards more plainly put before us by casting the desecrators out of the temple. The tree herself had not sinned by bearing no fruit when the Lord was hungry, for the time of figs was not yet come, but those Priests had sinned who were carrying on worldly business in the Lord's house, and who neglected to bring forth that fruit of godliness which they owed, and which the Lord was hungry to find in them. The Lord made the fig-tree to wither away under His curse, that all men who saw it, and all men who hear of it, might know that they will be condemned by the judgment of God, if they content themselves with the talk of godliness, without the solid fruit of good works, even as that barren figtree was clothed only with a rustling garb of green leaves.\par
\switchcolumn*

\LA
\begin{Resp}\Rsign Emendémus in mélius, quæ ignoránter peccávimus: ne súbito præoccupáti die mortis, quærámus spátium pœniténtiæ, et inveníre non possímus: \Star{} Atténde, Dómine, et miserére, quia peccávimus tibi.\end{Resp}
\begin{Resp}\Vsign Adjuva nos, Deus salutáris noster, et propter honórem nóminis tui, Dómine, líbera nos.\end{Resp}
\begin{Resp}\Rsign Atténde, Dómine, et miserére, quia peccávimus tibi.\end{Resp}
\switchcolumn
\EN
\begin{Resp}\Rsign Let us amend for the better in that wherein we have sinned unknowingly, or ever the day of death suddenly prevent us, and we seek a place of repentance, and find none. \Star{} Give heed, O Lord, and have mercy upon us, for we have sinned against thee.\end{Resp}
\begin{Resp}\Vsign Help us, O God of our salvation, and for the glory of thy Name deliver us, O Lord.\end{Resp}
\begin{Resp}\Rsign Give heed, O Lord, and have mercy upon us, for we have sinned against thee.\end{Resp}
\switchcolumn*

\LA
\LessonHead{Lectio II}
\switchcolumn
\EN
\LessonHead{Lesson II}
\switchcolumn*

\LA
\noindent Verum quia non intellexérunt, in ipsos consequénter districtiónem méritæ ultiónis exércuit: et ejécit commércia rerum humanárum de domo illa, in qua divínas tantum res agi, hóstias et oratiónes Dei offérri, verbum Dei legi, audíri, et decantári præcéptum erat. Et quidem credéndum est, quia ea tantum vendi vel emi repérerit in templo quæ ad ministérium necessária essent ejúsdem templi, juxta hoc quod álias factum légimus, cum idem templum ingrédiens, invénit in eo vendéntes et eméntes oves, et boves, et colúmbas: quia nimírum hæc ómnia non nisi ut offerréntur in domo Dómini, eos qui de longe vénerant, ab indígenis comparáre credéndum est.\par
\switchcolumn
\EN
\noindent But because the buyers and sellers understood not the parable of the barren fig-tree, the Lord brought upon them the stroke of the punishment that they had deserved, and cast out the traffickers in earthly things, from that house, wherein it had been commanded that nothing should be done save the work of God, sacrifices and prayers offered up to Him, and His word read, taught, and sung. And yet it may be believed that nothing was being sold or bought in the temple save such things as were needful for the service thereof, as we read in another place, (John ii. 14,) that when Jesus went into the temple He found those that sold oxen and sheep and doves, and all these things were doubtless there for no other end but to be offered to God in that His holy house, and were sold by the natives to those worshippers who came from a distance, to be so used.\par
\switchcolumn*

\LA
\begin{Resp}\Rsign Derelínquat ímpius viam suam, et vir iníquus cogitatiónes suas, et revertátur ad Dóminum, et miserébitur ejus: \Star{} Quia benígnus et miséricors est, et præstábilis super malítia Dóminus Deus noster.\end{Resp}
\begin{Resp}\Vsign Non vult Dóminus mortem peccatóris, sed ut convertátur et vivat.\end{Resp}
\begin{Resp}\Rsign Quia benígnus et miséricors est, et præstábilis super malítia Dóminus Deus noster.\end{Resp}
\switchcolumn
\EN
\begin{Resp}\Rsign Let the wicked forsake his way, and the unrighteous man his thoughts, and let him return unto the Lord, and He will have mercy upon him; \Star{} For the Lord our God is gracious, and merciful, and repenteth Him of the evil.\end{Resp}
\begin{Resp}\Vsign The Lord hath no pleasure in the death of the wicked; but that he turn from his way and live.\end{Resp}
\begin{Resp}\Rsign For the Lord our God is gracious, and merciful, and repenteth Him of the evil.\end{Resp}
\switchcolumn*

\LA
\LessonHead{Lectio III}
\switchcolumn
\EN
\LessonHead{Lesson III}
\switchcolumn*

\LA
\noindent Si ergo Dóminus nec ea volébat venúmdari in templo, quæ in templo volébat offérri, vidélicet propter stúdium avarítiæ, sive fraudis, quod próprium solet esse negotiántium fácinus: quanta putas animadversióne puníret, si invenísset ibi áliquos rísui vel vanilóquio vacántes aut álii cuílibet vítio mancipátos? Si enim ea quæ álibi líbere geri póterant, Dóminus in domo sua temporália negótia geri non pátitur: quanto magis ea quæ nusquam fíeri licet, plus cæléstis iræ meréntur, si in ǽdibus Deo sacrátis agúntur? Verum quia Spíritus Sanctus in colúmba super Dóminum appáruit, recte per colúmbas Sancti Spíritus charísmata signántur. Qui autem sunt in templo Dei hódie, qui colúmbas vendunt, nisi qui in Ecclésia prétium de impositióne manus accípiunt, per quam vidélicet impositiónem Spíritus Sanctus cǽlitus datur?\par
\switchcolumn
\EN
\noindent If, therefore, the Lord would not have to be sold in the temple, even such things as He willed should be offered therein, (On account, that is, of the greed or dishonesty which is often the stain of such transactions,) with what anger, suppose ye, would He visit such as He might find laughing or gossiping there, or yielding to any other sin. If the Lord suffer not to be carried on in His house such worldly business as may be freely done elsewhere, how much more shall such things as ought never to be done anywhere, draw down the anger of God if they be done in His own holy house. Lastly the Holy Ghost came down upon the Lord in the shape of a dove, and by doves therefore may be signified the gifts of that Holy Spirit. They, then, to this day sell doves in the temple of God, who take money in the Church for the laying on of their hands, whereby the Holy Ghost is given from heaven.\par
\switchcolumn*

\LA
\begin{Resp}\Rsign Paradísi portas apéruit nobis jejúnii tempus: suscipiámus illud orántes, et deprecántes: \Star{} Ut in die resurrectiónis cum Dómino gloriémur.\end{Resp}
\begin{Resp}\Vsign In ómnibus exhibeámus nosmetípsos sicut Dei minístros in multa patiéntia.\end{Resp}
\begin{Resp}\Rsign Ut in die resurrectiónis cum Dómino gloriémur.\end{Resp}
\begin{Resp}\Vsign Glória Patri, et Fílio, \Star{} et Spirítui Sancto.\end{Resp}
\begin{Resp}\Rsign Ut in die resurrectiónis cum Dómino gloriémur.\end{Resp}
\switchcolumn
\EN
\begin{Resp}\Rsign The season of the Fast openeth unto us the gates of heaven; let us enter thereon in prayer and supplication, \Star{} That on the day when the Lord riseth again we may rejoice with Him.\end{Resp}
\begin{Resp}\Vsign In all things let us approve ourselves the ministers of God, in much patience.\end{Resp}
\begin{Resp}\Rsign That on the day when the Lord riseth again we may rejoice with Him.\end{Resp}
\begin{Resp}\Vsign Glory be to the Father, and to the Son, \Star{} and to the Holy Ghost.\end{Resp}
\begin{Resp}\Rsign That on the day when the Lord riseth again we may rejoice with Him.\end{Resp}
\switchcolumn*

\end{paracol}

\Office{Feria Quarta Quattuor Temporum Quadragesimæ}{Ember Wednesday of Lent}{Feria major}
\addcontentsline{toc}{subsection}{Feria Quarta Quattuor Temporum Quadragesimæ\enspace\textit{Ember Wednesday of Lent}}
\begin{paracol}{2}
\LA
\LessonHead{Lectio I}
\switchcolumn
\EN
\LessonHead{Lesson I}
\switchcolumn*

\LA
\LessonTitle{Léctio sancti Evangélii secúndum Matthǽum}
\switchcolumn
\EN
\LessonTitle{Continuation of the Holy Gospel according to Matthew}
\switchcolumn*

\LA
\Cite{Matt 12:38-50}
\switchcolumn
\EN
\Cite{Matt 12:38-50}
\switchcolumn*

\LA
\noindent In illo témpore: Respondérunt Jesu quidam de scribis et pharisǽis, dicéntes: Magíster, vólumus a te signum vidére. Et réliqua.\par
\switchcolumn
\EN
\noindent In that time, some of the scribes and Pharisees answered him, saying: Master, we would see a sign from thee. Who answering said to them: An evil and adulterous generation seeketh a sign. And so on.\par
\switchcolumn*

\LA
\LessonTitle{Homilía sancti Ambrósii Epíscopi}
\switchcolumn
\EN
\LessonTitle{Homily by St. Ambrose, Bishop (of Milan.)}
\switchcolumn*

\LA
\Source{Lib. 7 in Lucæ cap. 11}
\switchcolumn
\EN
\Source{7th Bk. on Luke, ch. xii}
\switchcolumn*

\LA
\noindent Judæórum plebe damnáta, Ecclésiæ mystérium evidénter exprímitur, quæ in Ninivítis per pœniténtiam, et in regína Austri per stúdium percipiéndæ sapiéntiæ, de totíus orbis fínibus congregátur, ut pacífici Salomónis verba cognóscat. Regína plane, cujus regnum est indivísum, de divérsis et distántibus pópulis in unum corpus assúrgens.\par
\switchcolumn
\EN
\noindent After the condemnation of the Jewish people, the mystery of the Church is plainly declared in the figures of the repentant Ninevites, and of the Queen of the South. Like that Queen, the Church cometh from the uttermost parts of the earth, to hear the wisdom of the true Solomon, the Prince of Peace. A Queen she is, and a Queen of one indivisible realm, wrought into one body out of all nations, however diverse and distant.\par
\switchcolumn*

\LA
\begin{Resp}\Rsign Scíndite corda vestra, et non vestiménta vestra: et convertímini ad Dóminum Deum vestrum: \Star{} Quia benígnus et miséricors est.\end{Resp}
\begin{Resp}\Vsign Derelínquat ímpius viam suam, et vir iníquus cogitatiónes suas, et revertátur ad Dóminum, et miserébitur ejus.\end{Resp}
\begin{Resp}\Rsign Quia benígnus et miséricors est.\end{Resp}
\switchcolumn
\EN
\begin{Resp}\Rsign Rend your hearts and not your garments, and turn unto the Lord your God; \Star{} For He is gracious and merciful.\end{Resp}
\begin{Resp}\Vsign Let the wicked forsake his way, and the unrighteous man his thoughts, and let him return unto the Lord, and He will have mercy upon him.\end{Resp}
\begin{Resp}\Rsign For He is gracious and merciful.\end{Resp}
\switchcolumn*

\LA
\LessonHead{Lectio II}
\switchcolumn
\EN
\LessonHead{Lesson II}
\switchcolumn*

\LA
\noindent Itaque sacraméntum illud magnum est de Christo et Ecclésia. Sed tamen hoc majus est, quia illud in figúra ante præcéssit, nunc autem plenum in veritáte mystérium est. Illic enim Sálomon typus, hic autem Christus in suo córpore est. Ex duóbus ígitur constat Ecclésia: ut aut peccáre nésciat, aut peccáre désinat. Pœniténtia enim delíctum ábolet, sapiéntia cavet.\par
\switchcolumn
\EN
\noindent And thus cometh that great mystery of Christ and the Church, a mystery more excellent now in the fulness of truth, than in the ancient type. For there they had in Solomon only a type of that which Christ is now in His own Person. And the Church is of two classes, whereof the one knoweth not how to sin, and the other sinneth no more. To wash away sin is the work of repentance, to eschew it that of wisdom.\par
\switchcolumn*

\LA
\begin{Resp}\Rsign Frange esuriénti panem tuum, et egénos vagósque induc in domum tuam: \Star{} Tunc erúmpet quasi mane lumen tuum, et anteíbit fáciem tuam justítia tua.\end{Resp}
\begin{Resp}\Vsign Cum víderis nudum, óperi eum, et carnem tuam ne despéxeris.\end{Resp}
\begin{Resp}\Rsign Tunc erúmpet quasi mane lumen tuum, et anteíbit fáciem tuam justítia tua.\end{Resp}
\switchcolumn
\EN
\begin{Resp}\Rsign Deal thy bread to the hungry, and bring the poor and the wanderer to thine house. \Star{} Then shall thy light break forth as the morning, and thy righteousness shall go before thee.\end{Resp}
\begin{Resp}\Vsign When thou seest the naked, cover him; and hide not thyself from thine own flesh.\end{Resp}
\begin{Resp}\Rsign Then shall thy light break forth as the morning, and thy righteousness shall go before thee.\end{Resp}
\switchcolumn*

\LA
\LessonHead{Lectio III}
\switchcolumn
\EN
\LessonHead{Lesson III}
\switchcolumn*

\LA
\noindent Céterum Jonæ signum, ut typus Domínicæ passiónis, ita étiam grávium, quæ Judǽi commíserint, testificátio peccatórum est. Simul advértere licet et majestátis oráculum, et pietátis indícium. Namque Ninivitárum exémplo et denuntiátur supplícium, et remédium demonstrátur. Unde étiam Judǽi debent non desperáre indulgéntiam, si velint ágere pœniténtiam.\par
\switchcolumn
\EN
\noindent Lastly, the sign of the Prophet Jonas, as it was a figure of the Lord's sufferings, was also a witness to the gravity of those sins which the Jews committed. At the same time, we see in these words of the Lord a declaration at once of His power, and of His love; for, by turning our eyes on the Ninevites, He showeth us a way of escape, while He setteth before us the horror of what will otherwise be our punishment. Even the Jews need not cease to hope for pardon, if only they would repent.\par
\switchcolumn*

\LA
\begin{Resp}\Rsign Abscóndite eleemósynam in sinu páuperum, et ipsa orábit pro vobis ad Dóminum: \Star{} Quia sicut aqua extínguit ignem, ita eleemósyna extínguit peccátum.\end{Resp}
\begin{Resp}\Vsign Date eleemósynam, et ecce ómnia munda sunt vobis.\end{Resp}
\begin{Resp}\Rsign Quia sicut aqua extínguit ignem, ita eleemósyna extínguit peccátum.\end{Resp}
\begin{Resp}\Vsign Glória Patri, et Fílio, \Star{} et Spirítui Sancto.\end{Resp}
\begin{Resp}\Rsign Quia sicut aqua extínguit ignem, ita eleemósyna extínguit peccátum.\end{Resp}
\switchcolumn
\EN
\begin{Resp}\Rsign Shut up alms in the breast of the poor, and it shall plead for you with the Lord. \Star{} For, as water will quench fire, so alms maketh an atonement for sins.\end{Resp}
\begin{Resp}\Vsign Give alms, and, behold, all things are clean unto you.\end{Resp}
\begin{Resp}\Rsign For, as water will quench fire, so alms maketh an atonement for sins.\end{Resp}
\begin{Resp}\Vsign Glory be to the Father, and to the Son, \Star{} and to the Holy Ghost.\end{Resp}
\begin{Resp}\Rsign For, as water will quench fire, so alms maketh an atonement for sins.\end{Resp}
\switchcolumn*

\end{paracol}

\Office{Feria Quinta infra Hebdomadam I in Quadragesima}{Thursday of the First Week of Lent}{Feria major}
\addcontentsline{toc}{subsection}{Feria Quinta infra Hebdomadam I in Quadragesima\enspace\textit{Thursday of the First Week of Lent}}
\begin{paracol}{2}
\LA
\LessonHead{Lectio I}
\switchcolumn
\EN
\LessonHead{Lesson I}
\switchcolumn*

\LA
\LessonTitle{Léctio sancti Evangélii secúndum Matthǽum}
\switchcolumn
\EN
\LessonTitle{Continuation of the Holy Gospel according to Matthew}
\switchcolumn*

\LA
\Cite{Matt 15:21-28}
\switchcolumn
\EN
\Cite{Matt 15:21-28}
\switchcolumn*

\LA
\noindent In illo témpore: Egréssus Jesus secéssit in partes Tyri et Sidónis. Et réliqua.\par
\switchcolumn
\EN
\noindent In that time, Jesus went from thence, and retired into the coasts of Tyre and Sidon. And so on.\par
\switchcolumn*

\LA
\LessonTitle{Homilía sancti Hierónymi Presbýteri}
\switchcolumn
\EN
\LessonTitle{Homily by St. Jerome, Priest (at Bethlehem.)}
\switchcolumn*

\LA
\Source{Liber 2 Comm. in cap. 15. Matthæi}
\switchcolumn
\EN
\Source{Bk. ii Comm. on Matth. xv}
\switchcolumn*

\LA
\noindent Scribis et pharisǽis calumniatóribus derelíctis, transgréditur in partes Tyri et Sidónis, ut Týrios Sidoniósque curáret. Múlier autem Chananǽa egréditur de fínibus prístinis, ut clamans fíliæ ímpetret sanitátem. Obsérva quod in quintodécimo loco fília Chananǽæ sanétur. Miserére mei, Dómine, Fili David. Inde novit vocáre Fílium David, quia egréssa jam fúerat de fínibus suis, et errórem Tyriórum ac Sidoniórum loci et fídei commutatióne dimíserat.\par
\switchcolumn
\EN
\noindent Christ leaveth the Scribes and Pharisees who had spoken falsely against Him, and goeth into the coasts of Tyre and Sidon, that He may heal the Tyrians and Sidonians. But a woman of Canaan cometh to Him out of the land He had left, and crieth to Him to give health to her daughter. Remark that the case of the daughter of this woman of Canaan is the fifteenth case of healing. Have mercy on me, O Lord, Thou Son of David! She knew that He was to be called Son of David because she was come out of His own country, and had left the errors of the Tyrians and Sidonians when she changed her home and her faith.\par
\switchcolumn*

\LA
\begin{Resp}\Rsign Tribulárer, si nescírem misericórdias tuas, Dómine; tu dixísti: Nolo mortem peccatóris, sed ut magis convertátur et vivat: \Star{} Qui Chananǽam et publicánum vocásti ad pœniténtiam.\end{Resp}
\begin{Resp}\Vsign Secúndum multitúdinem dolórum meórum in corde meo, consolatiónes tuæ lætificavérunt ánimam meam.\end{Resp}
\begin{Resp}\Rsign Qui Chananǽam et publicánum vocásti ad pœniténtiam.\end{Resp}
\switchcolumn
\EN
\begin{Resp}\Rsign I had been troubled, but that I knew thy mercy, O Lord; Thou hast said: I have no pleasure in the death of the wicked, but that he turn from his way and live. \Star{} O Thou, Who didst call the Canaanite woman and the Publican unto repentance!\end{Resp}
\begin{Resp}\Vsign In the multitude of the sorrows within my heart, thy comforts delight my soul.\end{Resp}
\begin{Resp}\Rsign O Thou Who didst call the Canaanite woman and the Publican unto repentance!\end{Resp}
\switchcolumn*

\LA
\LessonHead{Lectio II}
\switchcolumn
\EN
\LessonHead{Lesson II}
\switchcolumn*

\LA
\noindent Fília mea male a dæmónio vexátur. Ego fíliam Chananǽæ puto ánimas esse credéntium, quæ male a dæmónio vexabántur, ignorántes Creatórem, et adorántes lápidem. Qui non respóndit ei verbum: non de supérbia pharisáica, nec de scribárum supercílio: sed ne ipse senténtiæ suæ viderétur esse contrárius, per quam jússerat: In viam géntium ne abiéritis, et in civitátes Samaritanórum ne intravéritis. Nolébat enim occasiónem calumniatóribus dare: perfectámque salútem géntium passiónis et resurrectiónis témpori reservábat.\par
\switchcolumn
\EN
\noindent My daughter is grievously vexed with a devil. I think that the daughter of this woman of Canaan, (whom the Lord at length delivered,) was a figure of the souls of such as now believe, but were once grievously vexed by the devil, knowing not Him Who made them, and bowing down to stocks and stones. But He answered not a word not because He was puffed up with the pride of the Pharisees, or shared the high looks of the Scribes, but that He might fulfil His own word that He had spoken, saying: Go not into the way of the Gentiles, and into any city of the Samaritans enter ye not. (Matth. x. 5.) He would not give an occasion to such as spoke falsely against Him, and He kept back perfect salvation from the Gentiles until such time as He should have suffered and risen again.\par
\switchcolumn*

\LA
\begin{Resp}\Rsign In ómnibus exhibeámus nosmetípsos sicut Dei minístros in multa patiéntia: \Star{} Ut non vituperétur ministérium nostrum.\end{Resp}
\begin{Resp}\Vsign Ecce nunc tempus acceptábile, ecce nunc dies salútis: commendémus nosmetípsos in multa patiéntia.\end{Resp}
\begin{Resp}\Rsign Ut non vituperétur ministérium nostrum.\end{Resp}
\switchcolumn
\EN
\begin{Resp}\Rsign In all things let us approve ourselves as the ministers of God, in much patience; \Star{} That our ministry be not blamed.\end{Resp}
\begin{Resp}\Vsign Behold, now is the acceptable time; behold, now is the day of salvation: let us approve ourselves in much patience.\end{Resp}
\begin{Resp}\Rsign That our ministry be not blamed.\end{Resp}
\switchcolumn*

\LA
\LessonHead{Lectio III}
\switchcolumn
\EN
\LessonHead{Lesson III}
\switchcolumn*

\LA
\noindent Et accedéntes discípuli ejus, rogábant eum, dicéntes: Dimítte eam, quia clamat post nos. Discípuli illo adhuc témpore mystéria Dómini nesciéntes, vel misericórdia commóti, rogábant pro Chananǽa mulíere, quam alter Evangelísta Syrophœníssam appéllat: vel importunitáte ejus carére cupiéntes, quia non ut cleméntem, sed ut durum médicum crébrius inclamáret. Ipse autem respóndens ait: Non sum missus, nisi ad oves quæ periérunt domus Israël. Non quo et ad gentes non missus sit, sed quo primum missus sit ad Israël: ut illis non recipiéntibus Evangélium, justa fíeret ad gentes transmigrátio.\par
\switchcolumn
\EN
\noindent And His disciples came and besought Him, saying: Send her away; for she crieth after us. The disciples, knowing not as yet the mysterious things of the Lord, said this, either because they were moved with compassion and so interceded for this Canaanite woman, whom another Evangelist calleth a Syrophoenician, (Mark vii. 26,) or because she was crying out that the Lord was an hard, instead of a merciful physician, and they desired to be rid of her clamour. But He answered and said: I am not sent but unto the lost sheep of the house of Israel, not that He was not sent unto the Gentiles, but because it was to Israel in the first instance that He was sent, whom refusing the Gospel, He might justly pass away from, and go to the Gentiles.\par
\switchcolumn*

\LA
\begin{Resp}\Rsign In jejúnio et fletu orábunt sacerdótes, dicéntes: \Star{} Parce, Dómine, parce pópulo tuo; et ne des hereditátem tuam in perditiónem.\end{Resp}
\begin{Resp}\Vsign Inter vestíbulum et altáre plorábunt sacerdótes, dicéntes.\end{Resp}
\begin{Resp}\Rsign Parce, Dómine, parce pópulo tuo; et ne des hereditátem tuam in perditiónem.\end{Resp}
\begin{Resp}\Vsign Glória Patri, et Fílio, \Star{} et Spirítui Sancto.\end{Resp}
\begin{Resp}\Rsign Parce, Dómine, parce pópulo tuo; et ne des hereditátem tuam in perditiónem.\end{Resp}
\switchcolumn
\EN
\begin{Resp}\Rsign The Priests shall pray, with fasting and with weeping, and shall say: \Star{} Spare, O Lord, spare thy people; and give not thine heritage to destruction.\end{Resp}
\begin{Resp}\Vsign The Priests shall weep between the porch and the altar, and shall say:\end{Resp}
\begin{Resp}\Rsign Spare, O Lord, spare thy people; and give not thine heritage to destruction.\end{Resp}
\begin{Resp}\Vsign Glory be to the Father, and to the Son, \Star{} and to the Holy Ghost.\end{Resp}
\begin{Resp}\Rsign Spare, O Lord, spare thy people; and give not thine heritage to destruction.\end{Resp}
\switchcolumn*

\end{paracol}

\Office{Feria Sexta Quattuor Temporum Quadragesimæ}{Ember Friday of Lent}{Feria major}
\addcontentsline{toc}{subsection}{Feria Sexta Quattuor Temporum Quadragesimæ\enspace\textit{Ember Friday of Lent}}
\begin{paracol}{2}
\LA
\LessonHead{Lectio I}
\switchcolumn
\EN
\LessonHead{Lesson I}
\switchcolumn*

\LA
\LessonTitle{Léctio sancti Evangélii secúndum Joánnem}
\switchcolumn
\EN
\LessonTitle{Continuation of the Holy Gospel according to John}
\switchcolumn*

\LA
\Cite{Joannes 5:1-15}
\switchcolumn
\EN
\Cite{John 5:1-15}
\switchcolumn*

\LA
\noindent In illo témpore: Erat dies festus Judæórum, et ascéndit Jesus Jerosólymam. Et réliqua.\par
\switchcolumn
\EN
\noindent In that time was a festival day of the Jews, and Jesus went up to Jerusalem. And so on.\par
\switchcolumn*

\LA
\LessonTitle{Homilía sancti Augustíni Epíscopi}
\switchcolumn
\EN
\LessonTitle{Homily of St. Augustine, Bishop of Hippo}
\switchcolumn*

\LA
\Source{Tractatus 17 in Joannem, post initium}
\switchcolumn
\EN
\Source{17th Tract on John}
\switchcolumn*

\LA
\noindent Videámus quid volúerit significáre in illo uno, quem étiam ipse servans unitátis mystérium, de tot languéntibus unum sanáre dignátus est. Invénit in annis ejus númerum quemdam languóris: trigínta et octo annos habébat in infirmitáte. Hic númerus quómodo magis ad languórem pertíneat, quam ad sanitátem, paulo diligéntius exponéndum est. Inténtos vos volo: áderit Dóminus, ut cóngrue loquar, et sufficiénter audiátis. Quadragenárius númerus sacrátus nobis in quadam perfectióne commendátur; notum esse árbitror caritáti vestræ: testántur sæpíssime divínæ Scriptúræ: jejúnium hoc número consecrátum esse, bene nostis. Nam et Móyses quadragínta diébus jejunávit, et Elías tótidem: et ipse Dóminus noster et Salvátor Jesus Christus hunc jejúnii númerum implévit. Per Móysen significátur Lex, per Elíam significántur Prophétæ, per Dóminum significátur Evangélium. Ideo in illo monte tres apparuérunt, ubi se discípulis osténdit in claritáte vultus et vestis suæ: appáruit enim médius inter Móysen et Elíam, tamquam Evangélium testimónium habéret a Lege et Prophétis.\par
\switchcolumn
\EN
\noindent Let us see what is mystically signified by that one infirm man whom alone the Lord, keeping to a mysterious unity, chose out of so many sufferers, to be the subject of His healing power. He found in him a certain number of years of sickness. He had had an infirmity thirty and eight years. How this number is proper rather to weakness than to health, will now be the subject of a few careful remarks. I bespeak your attention; the Lord will be present, that I may speak fitly, and you may understand. The number forty is put before us as hallowed, and, in a way, perfect. I think that your love knoweth this God's Scriptures often and; often witness it. Ye well know that a Fast of this number of days is hallowed. Moses fasted forty days. Elias did the same. And our Lord and Saviour Jesus Christ Himself fasted this number of days complete. Moses representeth the Law, Elias the Prophets, and the Lord the Gospel. And therefore these three appeared on the Mount of the Transfiguration. There the Lord showed Himself to His disciples with His Face shining as the sun, and His raiment glistering; and He stood between Moses and Elias; as it were, the Gospel receiving testimony, on the one hand from the Law, and, on the other, from the Prophets.\par
\switchcolumn*

\LA
\begin{Resp}\Rsign Emendémus in mélius, quæ ignoránter peccávimus: ne súbito præoccupáti die mortis, quærámus spátium pœniténtiæ, et inveníre non possímus: \Star{} Atténde, Dómine, et miserére, quia peccávimus tibi.\end{Resp}
\begin{Resp}\Vsign Adjuva nos, Deus salutáris noster, et propter honórem nóminis tui, Dómine, líbera nos.\end{Resp}
\begin{Resp}\Rsign Atténde, Dómine, et miserére, quia peccávimus tibi.\end{Resp}
\switchcolumn
\EN
\begin{Resp}\Rsign Let us amend for the better in that wherein we have sinned unknowingly, or ever the day of death suddenly prevent us, and we seek a place of repentance, and find none. \Star{} Give heed, O Lord, and have mercy upon us, for we have sinned against thee.\end{Resp}
\begin{Resp}\Vsign Help us, O God of our salvation, and for the glory of thy Name deliver us, O Lord.\end{Resp}
\begin{Resp}\Rsign Give heed, O Lord, and have mercy upon us, for we have sinned against thee.\end{Resp}
\switchcolumn*

\LA
\LessonHead{Lectio II}
\switchcolumn
\EN
\LessonHead{Lesson II}
\switchcolumn*

\LA
\noindent Sive ergo in Lege, sive in Prophétis, sive in Evangélio, quadragenárius númerus nobis in jejúnio commendátur. Jejúnium autem magnum et generále est, abstinére ab iniquitátibus et illícitis voluptátibus sǽculi, quod est perféctum jejúnium: Ut abnegántes impietátem et sæculáres cupiditátes, temperánter et juste et pie vivámus in hoc sǽculo. Huic jejúnio quam mercédem addit Apóstolus? Séquitur et dicit: Exspectántes illam beátam spem, et manifestatiónem glóriæ beáti Dei et Salvatóris nostri Jesu Christi. In hoc ergo sǽculo quasi quadragésimam abstinéntiæ celebrámus, cum bene vívimus, cum ab iniquitátibus et ab illícitis voluptátibus abstinémus: sed quia hæc abstinéntia sine mercéde non erit, exspectámus beátam illam spem, et revelatiónem glóriæ magni Dei et Salvatóris nostri Jesu Christi. In illa spe, cum fúerit de spe facta res, acceptúri sumus mercédem denárium. Ipsa enim merces rédditur operáriis in vínea laborántibus, secúndum Evangélium, quod vos credo reminísci: neque enim ómnia commemoránda sunt tamquam rúdibus et imperítis. Denárius ergo, qui accépit nomen a número decem, rédditur, et conjúnctus quadragenário fit quinquagenárius: unde cum labóre celebrámus Quadragésimam ante Pascha; cum lætítia vero, tamquam accépta mercéde, Quinquagésimam post Pascha.\par
\switchcolumn
\EN
\noindent Whether, therefore, it be in the Law, or in the Prophets, or in the Gospel, the number of forty is recommended to us for Fast-days. The great and general Fast is this to abstain from the iniquity of the world, and her forbidden pleasures. This is the perfect Fast, that, denying ungodliness, and worldly lusts, we should live soberly, righteously, and godly in this present world. After such a Fast, what is the Feast that followeth? Hear what the Apostle saith in continuation Looking for that blessed hope, and the glorious appearing of our great God and Saviour Jesus Christ. (Titus ii. 12, 13.) We, then, make our pilgrimage in this world a Lent, by living good lives, and abstaining from her iniquities and her forbidden pleasures. But at the end of this life-long Lent there will be an Easter indeed. We look for that blessed hope, and the glorious appearing of our great God and Saviour Jesus Christ When that hope is realised, when that faith is swallowed up in knowledge, then indeed shall we receive every man a penny. In good sooth, it is true that every labourer in the vineyard will get his wages witness that Gospel which I believe ye have not forgotten, (Matth. xx. 116) and which it is not my business to quote again as if ye were ignorant children. Now, the word used in the original for this penny which the labourers received is denarion. And the derivation of the word denarion is the numeral decem, ten. There are forty days in Lent, and if we add ten, we get fifty. So do we toil in fasting for the forty days of Lent before Easter, and, then, when we have, as it were, received our reward, we keep holiday for the fifty days of Eastertide.\par
\switchcolumn*

\LA
\begin{Resp}\Rsign Derelínquat ímpius viam suam, et vir iníquus cogitatiónes suas, et revertátur ad Dóminum, et miserébitur ejus: \Star{} Quia benígnus et miséricors est, et præstábilis super malítia Dóminus Deus noster.\end{Resp}
\begin{Resp}\Vsign Non vult Dóminus mortem peccatóris, sed ut convertátur et vivat.\end{Resp}
\begin{Resp}\Rsign Quia benígnus et miséricors est, et præstábilis super malítia Dóminus Deus noster.\end{Resp}
\switchcolumn
\EN
\begin{Resp}\Rsign Let the wicked forsake his way, and the unrighteous man his thoughts, and let him return unto the Lord, and He will have mercy upon him; \Star{} For the Lord our God is gracious, and merciful, and repenteth Him of the evil.\end{Resp}
\begin{Resp}\Vsign The Lord hath no pleasure in the death of the wicked; but that he turn from his way and live.\end{Resp}
\begin{Resp}\Rsign For the Lord our God is gracious, and merciful, and repenteth Him of the evil.\end{Resp}
\switchcolumn*

\LA
\LessonHead{Lectio III}
\switchcolumn
\EN
\LessonHead{Lesson III}
\switchcolumn*

\LA
\noindent Mementóte quod proposúerim númerum trigínta octo annórum in illo lánguido. Volo expónere, quare númerus ille trigésimus et octávus, languóris sit pótius quam sanitátis. Ergo, ut dicébam, cáritas implet Legem: ad plenitúdinem Legis in ómnibus opéribus pértinet quadragenárius númerus. In caritáte autem duo præcépta nobis commendántur: Díliges Dóminum Deum tuum ex toto corde tuo, et ex tota ánima tua, et ex tota mente tua: et díliges próximum tuum sicut teípsum. In his duóbus præcéptis tota Lex pendet, et Prophétæ. Mérito et illa vídua omnes facultátes suas, duo minúta misit in dona Dei: mérito et pro illo lánguido a latrónibus sauciáto stabulárius duos nummos accépit, unde sanarétur: mérito apud Samaritános bíduum fecit Jesus, ut eos caritáte firmáret. Binário ergo isto número cum áliquid boni significátur, máxime bipertíta cáritas commendátur. Si ergo quadragenárius númerus habet perfectiónem Legis, et Lex non implétur nisi in gémino præcépto caritátis: quid miráris, quia languébat, qui ad quadragínta, duo minus habébat?\par
\switchcolumn
\EN
\noindent Remember how I remarked, that the man healed by our Lord at the pool of Bethesda had had an infirmity thirty and eight years. I wish to explain why this number of thirty-eight is proper rather to weakness than to health. Love is the fulfilling of the law (Rom. xiii. 10;) to the fulfilling of the law belongeth in every work the number forty. But in love we have given us two precepts Thou shalt love the Lord thy God with all thy heart, and with all thy soul, and with all thy mind. This is the first and great commandment. And the second is like unto it Thou shalt love thy neighbour as thyself. On these two commandments hang all the law and the prophets. (Matth. xxii. 37-40.) When the widow gave all she had for an offering to God she gave two mites (Mark xii. 42;) the inn-keeper received two pence wherewith to cure him that had fallen among thieves (Luke x. 35;) Jesus abode for two days among the Samaritans (John iv. 40,) that He might establish them in love. When, then, anything good is spoken of as two, the two great divisions of love are the chief mystic interpretation. If, then, the law is fulfilled in the number forty, and it is not fulfilled if there be lacking the two precepts of love, what wonder is it that he was infirm who lacked two of forty?\par
\switchcolumn*

\LA
\begin{Resp}\Rsign Paradísi portas apéruit nobis jejúnii tempus: suscipiámus illud orántes, et deprecántes: \Star{} Ut in die resurrectiónis cum Dómino gloriémur.\end{Resp}
\begin{Resp}\Vsign In ómnibus exhibeámus nosmetípsos sicut Dei minístros in multa patiéntia.\end{Resp}
\begin{Resp}\Rsign Ut in die resurrectiónis cum Dómino gloriémur.\end{Resp}
\begin{Resp}\Vsign Glória Patri, et Fílio, \Star{} et Spirítui Sancto.\end{Resp}
\begin{Resp}\Rsign Ut in die resurrectiónis cum Dómino gloriémur.\end{Resp}
\switchcolumn
\EN
\begin{Resp}\Rsign The season of the Fast openeth unto us the gates of heaven; let us enter thereon in prayer and supplication, \Star{} That on the day when the Lord riseth again we may rejoice with Him.\end{Resp}
\begin{Resp}\Vsign In all things let us approve ourselves the ministers of God, in much patience.\end{Resp}
\begin{Resp}\Rsign That on the day when the Lord riseth again we may rejoice with Him.\end{Resp}
\begin{Resp}\Vsign Glory be to the Father, and to the Son, \Star{} and to the Holy Ghost.\end{Resp}
\begin{Resp}\Rsign That on the day when the Lord riseth again we may rejoice with Him.\end{Resp}
\switchcolumn*

\end{paracol}

\Office{Sabbato Quattuor Temporum Quadragesimæ}{Ember Saturday of Lent}{Feria major}
\addcontentsline{toc}{subsection}{Sabbato Quattuor Temporum Quadragesimæ\enspace\textit{Ember Saturday of Lent}}
\begin{paracol}{2}
\LA
\LessonHead{Lectio I}
\switchcolumn
\EN
\LessonHead{Lesson I}
\switchcolumn*

\LA
\LessonTitle{Léctio sancti Evangélii secúndum Matthǽum}
\switchcolumn
\EN
\LessonTitle{Continuation of the Holy Gospel according to Matthew}
\switchcolumn*

\LA
\Cite{Matt 17:1-9}
\switchcolumn
\EN
\Cite{Matt 17:1-9}
\switchcolumn*

\LA
\noindent In illo témpore: Assúmpsit Jesus Petrum, et Jacóbum, et Joánnem fratrem ejus, et duxit illos in montem excélsum seórsum: et transfigurátus est ante eos. Et réliqua.\par
\switchcolumn
\EN
\noindent In that time, Jesus taketh unto him Peter and James, and John his brother, and bringeth them up into a high mountain apart: And he was transfigured before them. And so on.\par
\switchcolumn*

\LA
\LessonTitle{Homilía sancti Leónis Papæ}
\switchcolumn
\EN
\LessonTitle{Homily by Pope St. Leo (the Great.)}
\switchcolumn*

\LA
\Source{Homilia de Transfiguratione Domini}
\switchcolumn
\EN
\Source{On the Transfiguration of the Lord}
\switchcolumn*

\LA
\noindent Evangélica lectio, dilectíssimi, quæ per aures córporis interiórem méntium nostrárum pulsávit audítum, ad magni sacraménti nos intellegéntiam vocat: quam, aspiránte grátia Dei, facílius assequémur, si consideratiónem nostram ad ea, quæ paulo supérius sunt narráta, referámus. Salvátor enim humáni géneris Jesus Christus, condens eam fidem, quæ et ímpios ad justítiam, et mórtuos révocat ad vitam, ad hoc discípulos suos doctrínæ mónitis, et óperum miráculis imbuébat, ut idem Christus et Unigénitus Dei, et hóminis Fílius crederétur. Nam unum horum sine áltero non próderat ad salútem: et æquális erat perículi, Dóminum Jesum Christum aut Deum tantúmmodo sine hómine, aut sine Deo solum hóminem credidísse: cum utrúmque esset páriter confiténdum: quia sicut Deo vera humánitas, ita hómini ínerat vera divínitas.\par
\switchcolumn
\EN
\noindent Dearly beloved brethren, the Lesson from the Holy Gospel which, entering in by our bodily ears, hath knocked at the door of our inner mind, calleth us to understand a great mystery. This, by the grace of God, we shall the more readily do, if we return to consider what hath been told us just before. The Saviour of mankind, even Jesus Christ, laying the foundations of that faith whereby the ungodly are called to righteousness and the dead to life, instilled into the minds of His disciples, both by the voice of His teaching and the wonder of His works, that they should believe Him, the one Christ, to be both the Only-begotten Son of God and the Son of man. Had they believed Him one of these and not the other, it had availed them nothing to salvation; and the danger was equally great, of holding the Lord Jesus Christ to be God without the Manhood, or Man only without the Godhead, since we are constrained to acknowledge that He is perfect God and perfect Man, and that as there is in the Godhead perfect Manhood, so there is in the Manhood perfect Godhead.\par
\switchcolumn*

\LA
\begin{Resp}\Rsign Scíndite corda vestra, et non vestiménta vestra: et convertímini ad Dóminum Deum vestrum: \Star{} Quia benígnus et miséricors est.\end{Resp}
\begin{Resp}\Vsign Derelínquat ímpius viam suam, et vir iníquus cogitatiónes suas, et revertátur ad Dóminum, et miserébitur ejus.\end{Resp}
\begin{Resp}\Rsign Quia benígnus et miséricors est.\end{Resp}
\switchcolumn
\EN
\begin{Resp}\Rsign Rend your hearts and not your garments, and turn unto the Lord your God; \Star{} For He is gracious and merciful.\end{Resp}
\begin{Resp}\Vsign Let the wicked forsake his way, and the unrighteous man his thoughts, and let him return unto the Lord, and He will have mercy upon him.\end{Resp}
\begin{Resp}\Rsign For He is gracious and merciful.\end{Resp}
\switchcolumn*

\LA
\LessonHead{Lectio II}
\switchcolumn
\EN
\LessonHead{Lesson II}
\switchcolumn*

\LA
\noindent Ad confírmandam ergo hujus fídei salubérrimam cognitiónem interrogáverat discípulos suos Dóminus, inter divérsas aliórum opiniónes quid ipsi de eo créderent, quidve sentírent. Ubi Petrus Apóstolus, per revelatiónem summi Patris corpórea súperans, et humána transcéndens, vidit mentis óculis Fílium Dei vivi, et conféssus est glóriam Deitátis; quia non ad solam respéxit substántiam carnis et sánguinis: tantúmque in hac fídei sublimitáte complácuit, ut beatitúdinis felicitáte donátus, sacram inviolábilis petræ accíperet firmitátem: super quam fundáta Ecclésia, portis ínferi et mortis légibus prævaléret: nec in solvéndis aut ligándis quorumcúmque causis áliud ratum esset in cælis, quam quod Petri sedísset arbítrio.\par
\switchcolumn
\EN
\noindent To strengthen, therefore, the saving knowledge of this faith, the Lord had asked His disciples what, among the differing opinions of men, it was their own belief and judgment as to Who He was. Then did the Apostle Peter, by the revelation of That Father Who is above all, rising above fleshly things, yea, outstripping the thoughts of men, then did he fix the eyes of his mind upon the Son of the living God, and confess the glory of the Godhead, for he looked not on the substance of the flesh and blood only. And in all the exaltation of this faith so well did he please God, that he was gifted with that joyous blessing, the hallowed establishment of that impregnable rock, whereon the Church being founded, should prevail against the gates of hell and the laws of death; neither, when anything is to be bound or loosed, is any bound or loosed in heaven, otherwise than as the judgment of Peter hath bound or loosed it upon earth.\par
\switchcolumn*

\LA
\begin{Resp}\Rsign Frange esuriénti panem tuum, et egénos vagósque induc in domum tuam: \Star{} Tunc erúmpet quasi mane lumen tuum, et anteíbit fáciem tuam justítia tua.\end{Resp}
\begin{Resp}\Vsign Cum víderis nudum, óperi eum, et carnem tuam ne despéxeris.\end{Resp}
\begin{Resp}\Rsign Tunc erúmpet quasi mane lumen tuum, et anteíbit fáciem tuam justítia tua.\end{Resp}
\switchcolumn
\EN
\begin{Resp}\Rsign Deal thy bread to the hungry, and bring the poor and the wanderer to thine house. \Star{} Then shall thy light break forth as the morning, and thy righteousness shall go before thee.\end{Resp}
\begin{Resp}\Vsign When thou seest the naked, cover him; and hide not thyself from thine own flesh.\end{Resp}
\begin{Resp}\Rsign Then shall thy light break forth as the morning, and thy righteousness shall go before thee.\end{Resp}
\switchcolumn*

\LA
\LessonHead{Lectio III}
\switchcolumn
\EN
\LessonHead{Lesson III}
\switchcolumn*

\LA
\noindent Hæc autem, dilectíssimi, laudátæ intellegéntiæ celsitúdo instruénda erat de inferióris substántiæ sacraménto: ne apostólica fides ad glóriam confiténdæ in Christo Deitátis evécta, infirmitátis nostræ receptiónem indígnam impassíbili Deo atque incóngruam judicáret: et ita jam in Christo humánam créderet glorificátam esse natúram, ut nec supplício posset áffici, nec morte dissólvi. Et ídeo dicénte Dómino, quod oportéret eum ire Jerosólymam, et multa pati a senióribus et scribis, ac princípibus sacerdótum, et occídi, et tértia die resúrgere: cum beátus Petrus, qui supérno illustrátus lúmine, de ardentíssima Fílii Dei confessióne fervébat, contumélias illusiónum et crudelíssimæ mortis oppróbrium religióso, ut putábat, et líbero fastídio respuísset; benígna a Jesu increpatióne corréptus, et ad cupiditátem participándæ cum eo passiónis animátus est.\par
\switchcolumn
\EN
\noindent But, dearly beloved brethren, it behoved that the height of this understanding, which the Lord praised, should rest upon a foundation, and that foundation, the mystery of the lower nature, lest the faith of the Apostle, carried away by the glorious acknowledgment of the Godhead in Christ, should deem it unworthy and unnatural for the impassible God to take into Himself the frailty of our nature; and should thus believe that in Christ the Manhood had been so glorified as to be no longer able to suffer pain, or be dissolved in death. And therefore it was that, when the Lord said how that He must go up unto Jerusalem, and suffer many things of the elders and chief priests, and scribes, and be killed, and rise again the third day, and the blessed Peter, bright with heavenly illumination, and still glowing from the passionate acknowledgment of the Divine Sonship, by a natural, and, as seemed to him, a godly shrinking, could not bear the mention of mockery and insult and a cruel death, he was corrected by the merciful rebuke of Jesus, and moved rather to desire to be a partaker in the sufferings of his Master.\par
\switchcolumn*

\LA
\begin{Resp}\Rsign Abscóndite eleemósynam in sinu páuperum, et ipsa orábit pro vobis ad Dóminum: \Star{} Quia sicut aqua extínguit ignem, ita eleemósyna extínguit peccátum.\end{Resp}
\begin{Resp}\Vsign Date eleemósynam, et ecce ómnia munda sunt vobis.\end{Resp}
\begin{Resp}\Rsign Quia sicut aqua extínguit ignem, ita eleemósyna extínguit peccátum.\end{Resp}
\begin{Resp}\Vsign Glória Patri, et Fílio, \Star{} et Spirítui Sancto.\end{Resp}
\begin{Resp}\Rsign Quia sicut aqua extínguit ignem, ita eleemósyna extínguit peccátum.\end{Resp}
\switchcolumn
\EN
\begin{Resp}\Rsign Shut up alms in the breast of the poor, and it shall plead for you with the Lord. \Star{} For, as water will quench fire, so alms maketh an atonement for sins.\end{Resp}
\begin{Resp}\Vsign Give alms, and, behold, all things are clean unto you.\end{Resp}
\begin{Resp}\Rsign For, as water will quench fire, so alms maketh an atonement for sins.\end{Resp}
\begin{Resp}\Vsign Glory be to the Father, and to the Son, \Star{} and to the Holy Ghost.\end{Resp}
\begin{Resp}\Rsign For, as water will quench fire, so alms maketh an atonement for sins.\end{Resp}
\switchcolumn*

\end{paracol}

\Office{Dominica II in Quadragesima}{Second Sunday of Lent}{Semiduplex I classis}
\addcontentsline{toc}{subsection}{Dominica II in Quadragesima\enspace\textit{Second Sunday of Lent}}
\begin{paracol}{2}
\LA
\LessonHead{Lectio I}
\switchcolumn
\EN
\LessonHead{Lesson I}
\switchcolumn*

\LA
\LessonTitle{De libro Génesis}
\switchcolumn
\EN
\LessonTitle{Lesson from the book of Genesis}
\switchcolumn*

\LA
\Cite{Gen 27:1-10}
\switchcolumn
\EN
\Cite{Gen 27:1-10}
\switchcolumn*

\LA
\noindent Sénuit autem Isaac, et calligavérunt óculi ejus, et vidére non póterat: vocavítque Esau fílium suum majórem, et dixit ei: Fili mi! Qui respóndit: Adsum. Cui pater: Vides, inquit, quod senúerim et ignórem diem mortis meæ. Sume arma tua, pháretram, et arcum, et egrédere foras: cumque venátu áliquid apprehénderis, fac mihi inde pulméntum sicut velle me nosti, et affer ut cómedam: et benedícat tibi ánima mea ántequam móriar. Quod cum audísset Rebécca, et ille abiísset in agrum ut jussiónem patris impléret, dixit fílio suo Jacob: Audívi patrem tuum loquéntem cum Esau fratre tuo, et dicéntem ei: Affer mihi de venatióne tua, et fac cibos ut cómedam, et benedícam tibi coram Dómino ántequam móriar. Nunc ergo fili mi, acquiésce consíliis meis; et pergens ad gregem, affer mihi duos hædos óptimos, ut fáciam ex eis escas patri tuo, quibus libénter véscitur: quas cum intúleris, et coméderit, benedícat tibi priúsquam moriátur.\par
\switchcolumn
\EN
\noindent Now Isaac was old, and his eyes were dim, and he could not see: and he called Esau, his elder son, and said to him: My son? And he answered: Here I am. And his father said to him: Thou seest that I am old, and know not the day of my death. Take thy arms, thy quiver, and bow, and go abroad: and when thou hast taken some thing by hunting, Make me savoury meat thereof, as thou knowest I like, and bring it, that I may eat: and my soul may bless thee before I die. And when Rebecca had heard this, and he was gone into the field to fulfill his father's commandment, She said to her son Jacob: I heard thy father talking with Esau thy brother, and saying to him: Bring me of thy hunting, and make me meats that I may eat, and bless thee in the sight of the Lord, before I die. Now, therefore, my son, follow my counsel: And go thy way to the flock, bring me two kids of the best, that I may make of them meat for thy father, such as he gladly eateth: Which when thou hast brought in, and he hath eaten, he may bless thee before he die.\par
\switchcolumn*

\LA
\begin{Resp}\Rsign Tolle arma tua, pháretram et arcum, et affer de venatióne tua, ut cómedam: \Star{} Et benedícat tibi ánima mea.\end{Resp}
\begin{Resp}\Vsign Cumque venátu áliquid attúleris, fac mihi inde pulméntum, ut cómedam.\end{Resp}
\begin{Resp}\Rsign Et benedícat tibi ánima mea.\end{Resp}
\switchcolumn
\EN
\begin{Resp}\Rsign Take thy weapons, thy quiver and thy bow, and bring me some of thy venison, that I may eat, \Star{} and my soul may bless thee.\end{Resp}
\begin{Resp}\Vsign And when thou hast taken somewhat, make me thereof savoury meat, that I may eat.\end{Resp}
\begin{Resp}\Rsign And my soul may bless thee.\end{Resp}
\switchcolumn*

\LA
\LessonHead{Lectio II}
\switchcolumn
\EN
\LessonHead{Lesson II}
\switchcolumn*

\LA
\Cite{Gen 27:11-20}
\switchcolumn
\EN
\Cite{Gen 27:11-20}
\switchcolumn*

\LA
\noindent Cui ille respóndit: Nosti quod Esau frater meus homo pilósus sit, et ego lenis: si attractáverit me pater meus, et sénserit, tímeo ne putet sibi voluísse illúdere, et indúcam super me maledictiónem pro benedictióne. Ad quem mater: In me sit, ait, ista maledíctio, fili mi: tantum audi vocem meam, et pergens, affer quæ dixi. Abiit, et áttulit, dedítque matri. Parávit illa cibos, sicut velle nóverat patrem illíus. Et véstibus Esau valde bonis, quas apud se habébat domi, índuit eum: pelliculásque hædórum circúmdedit mánibus, et colli nuda protéxit. Dedítque pulméntum, et panes, quos cóxerat, trádidit. Quibus illátis, dixit: Pater mi! At ille respóndit: Audio. Quis es tu, fili mi? Dixítque Jacob: Ego sum primogénitus tuus Esau: feci sicut præcepísti mihi: surge, sede, et cómede de venatióne mea, ut benedícat mihi ánima tua. Rursúmque Isaac ad fílium suum: Quómodo, inquit, tam cito inveníre potuísti, fili mi? Qui respóndit: Volúntas Dei fuit ut cito occúrreret mihi quod volébam.\par
\switchcolumn
\EN
\noindent And he answered her: Thou knowest that Esau my brother is a hairy man, and I am smooth. If my father shall feel me, and perceive it, I fear lest he will think I would have mocked him, and I shall bring upon me a curse instead of a blessing. And his mother said to him: Upon me be this curse, my son: only hear thou my voice, and go, fetch me the things which I have said. He went, and brought, and gave them to his mother. She dressed meats, such as she knew his father liked. And she put on him very good garments of Esau, which she had at home with her: And the little skins of the kids she put about his hands, and covered the bare of his neck. And she gave him the savoury meat, and delivered him bread that she had baked. Which when he had carried in, he said: My father? But he answered: I hear. Who art thou, my son? And Jacob said: I am Esau thy firstborn: I have done as thou didst command me: arise, sit, and eat of my venison, that thy soul may bless me. And Isaac said to his son: How couldst thou find it so quickly, my son? He answered: It was the will of God, that what I sought came quickly in my way.\par
\switchcolumn*

\LA
\begin{Resp}\Rsign Ecce odor fílii mei sicut odor agri pleni, cui benedíxit Dóminus: créscere te fáciat Deus meus sicut arénam maris: \Star{} Et donet tibi de rore cæli benedictiónem.\end{Resp}
\begin{Resp}\Vsign Deus autem omnípotens benedícat tibi, atque multíplicet.\end{Resp}
\begin{Resp}\Rsign Et donet tibi de rore cæli benedictiónem.\end{Resp}
\switchcolumn
\EN
\begin{Resp}\Rsign See the smell of my son is as the smell of a field which the Lord hath blessed. May my God multiply thee as the sand of the sea. \Star{} And give thee a blessing of the dew of heaven.\end{Resp}
\begin{Resp}\Vsign And God Almighty bless thee, and multiply thee.\end{Resp}
\begin{Resp}\Rsign And give thee a blessing of the dew of heaven.\end{Resp}
\switchcolumn*

\LA
\LessonHead{Lectio III}
\switchcolumn
\EN
\LessonHead{Lesson III}
\switchcolumn*

\LA
\Cite{Gen 27:21-29}
\switchcolumn
\EN
\Cite{Gen 27:21-29}
\switchcolumn*

\LA
\noindent Dixítque Isaac: Accéde huc ut tangam te, fili mi, et probem utrum tu sis fílius meus Esau, an non. Accéssit ille ad patrem, et, palpáto eo, dixit Isaac: Vox quidem, vox Jacob est: sed manus, manus sunt Esau. Et non cognóvit eum, quia pilósæ manus similitúdinem majóris exprésserant. Benedícens ergo illi, ait: Tu es fílius meus Esau? Respóndit: Ego sum. At ille: Affer mihi, inquit, cibos de venatióne tua, fili mi, ut benedícat tibi ánima mea. Quos cum oblátos comedísset, óbtulit ei étiam vinum. Quo hausto, dixit ad eum: Accéde ad me, et da mihi ósculum, fili mi. Accéssit, et osculátus est eum. Statímque ut sensit vestimentórum illíus flagrántiam, benedícens illi, ait: Ecce odor fílii mei sicut odor agri pleni, cui benedíxit Dóminus. Det tibi Deus de rore cæli, et de pinguédine terræ abundántiam fruménti et vini. Et sérviant tibi pópuli, et adórent te tribus: esto dóminus fratrum tuórum, et incurvéntur ante te fílii matris tuæ. Qui maledíxerit tibi, sit ille maledíctus: et qui benedíxerit tibi, benedictiónibus repleátur.\par
\switchcolumn
\EN
\noindent And Isaac said: Come hither, that I may feel thee, my son, and may prove whether thou be my son Esau, or not. He came near to his father, and when he had felt him, Isaac said: The voice indeed is the voice of Jacob; but the hands are the hands of Esau. He said: Art thou my son Esau? He answered: I am. Then he said: Bring me the meats of thy hunting, my son, that my soul may bless thee. And when they were brought, and he had eaten, he offered him wine also, which after he had drunk, He said to him: Come near me, and give me a kiss, my son. He came near, and kissed him. And immediately as he smelled the fragrant smell of his garments, blessing him, he said: Behold the smell of my son is as the smell of a plentiful field, which Lord hath blessed. God give thee the dew of heaven, and of the fatness of the earth, abundance of corn and wine. And let peoples serve thee, and tribes worship thee: be thou lord of thy brethren, and let they mother's children bow down before thee. Cursed be he that curseth thee: and let him that blesseth thee be filled with blessings.\par
\switchcolumn*

\LA
\begin{Resp}\Rsign Det tibi Deus de rore cæli et de pinguédine terræ abundántiam: sérviant tibi tribus et pópuli: \Star{} Esto dóminus fratrum tuórum.\end{Resp}
\begin{Resp}\Vsign Et incurvéntur ante te fílii matris tuæ.\end{Resp}
\begin{Resp}\Rsign Esto dóminus fratrum tuórum.\end{Resp}
\begin{Resp}\Vsign Glória Patri, et Fílio, \Star{} et Spirítui Sancto.\end{Resp}
\begin{Resp}\Rsign Esto dóminus fratrum tuorum.\end{Resp}
\switchcolumn
\EN
\begin{Resp}\Rsign God give thee of the dew of heaven and the fatness of the earth. Let people and nations serve thee. \Star{} Be lord over thy brethren.\end{Resp}
\begin{Resp}\Vsign And let thy mother's sons bow down to thee.\end{Resp}
\begin{Resp}\Rsign Be lord over thy brethren.\end{Resp}
\begin{Resp}\Vsign Glory be to the Father, and to the Son, \Star{} and to the Holy Ghost.\end{Resp}
\begin{Resp}\Rsign Be lord over thy brethren.\end{Resp}
\switchcolumn*

\LA
\LessonHead{Lectio IV}
\switchcolumn
\EN
\LessonHead{Lesson IV}
\switchcolumn*

\LA
\LessonTitle{Ex libro sancti Augustíni Epíscopi contra mendacium}
\switchcolumn
\EN
\LessonTitle{From the Book against Lying written by St. Augustine, Bishop (of Hippo.)}
\switchcolumn*

\LA
\Source{Cap. 10, tom. 4; post init.}
\switchcolumn
\EN
\Source{Ch. ix, tom. 4}
\switchcolumn*

\LA
\noindent Jacob quod matre fecit auctóre, ut patrem fállere viderétur, si diligénter et fidéliter attendátur, non est mendácium, sed mystérium. Quæ si mendácia dixérimus, omnes étiam parábolæ ac figúræ significandárum quarumcúmque rerum, quæ non ad proprietátem accipiéndæ sunt, sed in eis áliud ex álio est intelligéndum, dicéntur esse mendácia: quod absit omníno. Nam qui hoc putat, trópicis étiam tam multis locutiónibus ómnibus potest hanc importáre calúmniam; ita ut et hæc ipsa, quæ appellátur metáphora, hoc est, de re própria ad rem non própriam verbi alicújus usurpáta translátio, possit ista ratióne mendácium nuncupári.\par
\switchcolumn
\EN
\noindent If we consider faithfully and carefully what it was that Jacob did by the advice of his mother, and wherein he seemeth to have deceived his father, it will appear that (it hath an aspect in which) it is not a lie, but an allegory. If we denounce this (its mystic sense) as a lie, then must we also give the name of lies to even all parable, and to every figure devised to set forth the nature of anything, which is not to be taken in its literal sense, but in which one thing is to be understood under the name of another. And this be far from us. Whoso should do this, would bring the charge of falsehood against very many figures of speech, including that one called metaphor (in which a word is transferred from that meaning which belongeth to it, to some other) to which would, by such reasoning, be given the name of a lie.\par
\switchcolumn*

\LA
\begin{Resp}\Rsign Dum exíret Jacob de terra sua, vidit glóriam Dei, et ait: Quam terríbilis est locus iste! \Star{} Non est hic áliud, nisi domus Dei, et porta cæli.\end{Resp}
\begin{Resp}\Vsign Vere Deus est in loco isto, et ego nesciébam.\end{Resp}
\begin{Resp}\Rsign Non est hic áliud, nisi domus Dei, et porta cæli.\end{Resp}
\switchcolumn
\EN
\begin{Resp}\Rsign As Jacob went out from his own land, he saw the glory of God, and said: How dreadful is this place. \Star{} This is none other but the house of God and this is the gate of heaven.\end{Resp}
\begin{Resp}\Vsign Surely God is in this place, and I knew it not.\end{Resp}
\begin{Resp}\Rsign This is none other but the house of God and this is the gate of heaven.\end{Resp}
\switchcolumn*

\LA
\LessonHead{Lectio V}
\switchcolumn
\EN
\LessonHead{Lesson V}
\switchcolumn*

\LA
\noindent Quæ significántur enim, útique ipsa dicúntur: putántur autem mendácia, quóniam non ea quæ vere significántur, dicta intelligúntur; sed ea, quæ falsa sunt, dicta esse credúntur. Hoc ut exémplis fiat plánius, idípsum quod Jacob fecit, atténde. Hædínis certe péllibus membra contéxit. Si causam próximam requirámus, mentítum putábimus: hoc enim fecit, ut putarétur esse qui non erat. Si autem hoc factum ad illud, propter quod significándum revéra factum est, referátur: per hædínas pelles, peccáta; per eum vero, qui eis se opéruit, ille significátus est, qui non sua, sed aliéna peccáta portávit.\par
\switchcolumn
\EN
\noindent The deep meaning is given; but what is considered is the lie because men do not understand the way in which that signification, which is a truth, is set forth but the falsehood is plainly expressed, and believed. That we may understand this more plainly by taking some points in illustration, consider with me what Jacob did. It is certain that he covered his limbs with the skins of goats. If we consider his object in point of fact, we shall find that it was to lie, because he did this that he might be thought to be he who he was not. But if we consider this his deed in that deep typical sense which it undoubtedly possesseth, we find that by the goat-skins are represented sins, and by him who covered himself therewith Him Who bore not His own sins, but the sins of others.\par
\switchcolumn*

\LA
\begin{Resp}\Rsign Si Dóminus Deus meus fúerit mecum in via ista, per quam ego ámbulo, et custodíerit me, et déderit mihi panem ad edéndum, et vestiméntum quo opériar, et revocáverit me cum salúte: \Star{} Erit mihi Dóminus in refúgium, et lapis iste in signum.\end{Resp}
\begin{Resp}\Vsign Surgens ergo mane Jacob, tulit lápidem quem supposúerat cápiti suo, et eréxit in títulum, fundénsque óleum désuper, dixit.\end{Resp}
\begin{Resp}\Rsign Erit mihi Dóminus in refúgium, et lapis iste in signum.\end{Resp}
\switchcolumn
\EN
\begin{Resp}\Rsign If the Lord my God will be with me, in this way that I go, and will keep me, and will give me bread to eat, and raiment to put on, and will bring me again safely. \Star{} The Lord shall be my refuge, and this stone shall be a sign.\end{Resp}
\begin{Resp}\Vsign So Jacob rose up early in the morning, and took the stone that he had put for his pillow, and set it up for a pillar, and poured oil upon the top of it, and said:\end{Resp}
\begin{Resp}\Rsign The Lord shall be my refuge, and this shall be a sign.\end{Resp}
\switchcolumn*

\LA
\LessonHead{Lectio VI}
\switchcolumn
\EN
\LessonHead{Lesson VI}
\switchcolumn*

\LA
\noindent Verax ergo significátio nullo modo mendácium recte dici potest: ut autem in facto, ita et in verbo. Nam cum ei pater dixísset: Quis es tu, fili? ille respóndit: Ego sum Esau primogénitus tuus. Hoc si referátur ad duos illos géminos, mendácium vidébitur: si autem ad illud, propter quod significándum ista gesta dictáque conscrípta sunt; ille est hic intelligéndus in córpore suo, quod est ejus Ecclésia, qui de hac re loquens, ait: Cum vidéritis Abraham, et Isaac et Jacob et omnes Prophétas in regno Dei, vos autem expélli foras. Et, Vénient ab Oriénte et Occidénte, et Aquilóne et Austro, et accúmbent in regno Dei. Et, Ecce sunt novíssimi qui erant primi: et sunt primi, qui erant novíssimi. Sic enim quodámmodo minor majóris primátum frater ábstulit, atque in se tránstulit fratris.\par
\switchcolumn
\EN
\noindent It is impossible to apply the term a “lie” to that mystic aspect of this transaction in which it was true and such an aspect there is, not only in the acts, but in the words. When Isaac said to Jacob: “Who art thou, my son” and Jacob answered: “I am Esau, thy first-born”, if we take this in its sense relative to the two brothers, it will be apparent that it was a lie. If, however, we look at it relatively to that for the sake of which these words and deeds were written down, we shall see that Christ is here signified in His mystic body, the Church. Concerning her, (the younger covenant,) He saith (to them of the older covenant): “Ye shall see Abraham, and Isaac, and Jacob, and all the Prophets in the kingdom of God, and you yourselves thrust out. And they shall come from the east, and from the west, and from the north, and from the south, and shall sit down in the kingdom of God. And, behold, there are last which shall be first, and there are first which shall be last.” (Luke xiii. 28-30.) Thus did the younger take away the title and inheritance from the elder, and acquire it to himself.\par
\switchcolumn*

\LA
\begin{Resp}\Rsign Erit mihi Dóminus in Deum, et lapis iste quem eréxi in títulum, vocábitur domus Dei: et de univérsis quæ déderis mihi, \Star{} Décimas et hóstias pacíficas ófferam tibi.\end{Resp}
\begin{Resp}\Vsign Si revérsus fúero próspere ad domum patris mei.\end{Resp}
\begin{Resp}\Rsign Décimas et hóstias pacíficas ófferam tibi.\end{Resp}
\begin{Resp}\Vsign Glória Patri, et Fílio, \Star{} et Spirítui Sancto.\end{Resp}
\begin{Resp}\Rsign Décimas et hóstias pacíficas ófferam tibi.\end{Resp}
\switchcolumn
\EN
\begin{Resp}\Rsign The Lord shall be my God, and this stone, which I have set for a pillar, shall be called God's house, and of all that Thou shalt give me; \Star{} I will offer tithes and peace-offerings to thee.\end{Resp}
\begin{Resp}\Vsign If I come again to my father's house in peace.\end{Resp}
\begin{Resp}\Rsign I will offer tithes and peace-offerings unto thee.\end{Resp}
\begin{Resp}\Vsign Glory be to the Father, and to the Son, \Star{} and to the Holy Ghost.\end{Resp}
\begin{Resp}\Rsign I will offer tithes and peace-offerings unto thee.\end{Resp}
\switchcolumn*

\LA
\LessonHead{Lectio VII}
\switchcolumn
\EN
\LessonHead{Lesson VII}
\switchcolumn*

\LA
\LessonTitle{Léctio sancti Evangélii secúndum Matthǽum}
\switchcolumn
\EN
\LessonTitle{From the Holy Gospel according to Matthew}
\switchcolumn*

\LA
\Cite{Matt 17:1-9}
\switchcolumn
\EN
\Cite{Matt 17:1-9}
\switchcolumn*

\LA
\noindent In illo témpore: Assúmpsit Jesus Petrum, et Jacóbum, et Joánnem fratrem ejus, et duxit illos in montem excélsum seórsum: et transfigurátus est ante eos. Et réliqua.\par
\switchcolumn
\EN
\noindent At that time, Jesus taketh unto him Peter and James, and John his brother, and bringeth them up into a high mountain apart: And he was transfigured before them. And so on.\par
\switchcolumn*

\LA
\LessonTitle{De Homilía sancti Leónis Papæ}
\switchcolumn
\EN
\LessonTitle{Homily by Pope St. Leo (the Great.)}
\switchcolumn*

\LA
\Source{Ex Homil. de Transfiguratione Domini}
\switchcolumn
\EN
\Source{From a homily on the Transfiguration of the Lord}
\switchcolumn*

\LA
\noindent Assúmpsit Jesus Petrum, et Jacóbum, et fratrem ejus Joánnem, et conscénso cum eis seórsum monte præcélso, claritátem suæ glóriæ demonstrávit: quia licet intellexíssent in eo majestátem Dei, ipsíus tamen córporis, quo divínitas tegebátur, poténtiam nesciébant. Et ídeo próprie signantérque promíserat, quosdam de astántibus discípulis non prius gustáre mortem, quam vidérent Fílium hóminis veniéntem in regno suo, id est, in régia claritáte, quam spiritáliter ad natúram suscépti hóminis pertinéntem, his tribus viris vóluit esse conspícuam. Nam illam ipsíus Deitátis ineffábilem et inaccessíbilem visiónem, quæ in ætérnam vitam mundis corde servátur, nullo modo mortáli adhuc carne circúmdati intuéri póterant et vidére.\par
\switchcolumn
\EN
\noindent Jesus took Peter, and James, and John his brother, and brought them up into an exceeding high mountain apart, and manifested forth the brightness of His glory. Hitherto, though they understood that there was in Him the Majesty of God, they knew not the power of that Body which veiled the Godhead. And therefore He had individually and markedly promised to some of the disciples that had stood by Him (Matth. xvi. 28) that they should “not taste of death till they had seen the Son of Man coming in His kingdom”, that is, in the kingly splendour, which is the right of the Manhood taken into God, and which He willed to make visible to those three men. This it was that they saw, for the unspeakable and unapproachable vision of the Godhead Himself which will be the everlasting life of the pure in heart, (Matth. v. 8) can no man, who is still burdened with a dying body, see and live.\par
\switchcolumn*

\LA
\begin{Resp}\Rsign Dixit Angelus ad Jacob: \Star{} Dimítte me, auróra est. Respóndit ei: Non dimíttam te, nisi benedíxeris mihi. Et benedíxit ei in eódem loco.\end{Resp}
\begin{Resp}\Vsign Cumque surrexísset Jacob, ecce vir luctabátur cum eo usque mane: et cum vidéret quod eum superáre non posset, dixit ad eum.\end{Resp}
\begin{Resp}\Rsign Dimítte me, auróra est. Respóndit ei: Non dimíttam te, nisi benedíxeris mihi. Et benedíxit ei in eódem loco.\end{Resp}
\switchcolumn
\EN
\begin{Resp}\Rsign The Angel said unto Jacob \Star{} Let me go, for the day breaketh. And he said: I will not let thee go, except thou bless me. And he blessed him there.\end{Resp}
\begin{Resp}\Vsign And when Jacob arose, behold there wrestled a man with him, until the breaking of the day and, when he saw that he prevailed not, he said unto him\end{Resp}
\begin{Resp}\Rsign Let me go, for the day breaketh. And he said: I will not let thee go, except thou bless me. And he blessed him there.\end{Resp}
\switchcolumn*

\LA
\LessonHead{Lectio VIII}
\switchcolumn
\EN
\LessonHead{Lesson VIII}
\switchcolumn*

\LA
\noindent Dicénte Patre: Hic est Fílius meus diléctus, in quo mihi bene complácui, ipsum audíte: nonne evidénter audítum est: Hic est Fílius meus, cui ex me et mecum esse sine témpore est? quia nec génitor génito prior, nec génitus est genitóre postérior. Hic est Fílius meus, quem a me non séparat Déitas, non dívidit potéstas, non discérnit ætérnitas. Hic est Fílius meus, non adoptívus, sed próprius: non aliúnde creátus, sed ex me génitus: nec de ália natúra mihi factus comparábilis, sed de mea esséntia mihi natus æquális.\par
\switchcolumn
\EN
\noindent When the Father saith: “This is My beloved Son, in Whom I am well pleased; hear ye Him” did they not plainly hear Him say: “This is My Son, Whose it is to be of Me and with Me without all time. For neither is He That begetteth, before Him That is begotten, neither He That is begotten, after Him That begetteth Him.” This is My Son between Whom and Me, to be God is not a point of difference to be Almighty, a point of separation nor to be Eternal, a point of distinction. “This is My Son not by adoption, but My very Own; not created from, or of another substance, or out of nothing, but begotten of Me not of another nature, and made like unto Me, but of Mine own Being, born of Me, equal unto Me.”\par
\switchcolumn*

\LA
\begin{Resp}\Rsign Vidi Dóminum fácie ad fáciem: \Star{} Et salva facta est ánima mea.\end{Resp}
\begin{Resp}\Vsign Et dixit mihi: Nequáquam vocáberis Jacob, sed Israël erit nomen tuum.\end{Resp}
\begin{Resp}\Rsign Et salva facta est ánima mea.\end{Resp}
\switchcolumn
\EN
\begin{Resp}\Rsign I have seen God face to face; \Star{} And my life is preserved.\end{Resp}
\begin{Resp}\Vsign And he said unto me: thy name shall be called no more Jacob, but Israel shall be thy name.\end{Resp}
\begin{Resp}\Rsign And my life is preserved.\end{Resp}
\switchcolumn*

\LA
\LessonHead{Lectio IX}
\switchcolumn
\EN
\LessonHead{Lesson IX}
\switchcolumn*

\LA

\switchcolumn
\EN
\LessonTitle{“This is My Son by Whom all things were made, and without Whom was not anything made that was made, (John i. 3) Who maketh likewise all things whatsoever I make and what things soever I do He doeth likewise, (v. 19), inseparably and indifferently.” This is My Son Who thought it not robbery, nor hath taken it by violence, to be equal with Me, but, abiding still in the form of My glory, that He may fulfill Our common decree for the restoration of mankind, hath bowed the unchangeable Godhead even to the form of a servant. (Phil, ii. 6, 7.) Him therefore in Whom I am in all things well pleased, by Whose preaching I am manifested, and by Whose lowliness I am glorified, Him instantly hear ye. For He is the Truth and the Life, (John xiv. 6,) My Power, and My Wisdom. (1 Cor. i. 24.)}
\switchcolumn*

\LA
\noindent Hic est Fílius meus, per quem ómnia facta sunt, et sine quo factum est nihil: qui ómnia quæ fácio, simíliter facit; et quidquid óperor, inseparabíliter mecum atque indifferénter operátur. Hic est Fílius meus, qui eam, quam mecum habet æqualitátem, non rapína appétiit, nec usurpatióne præsúmpsit: sed manens in forma glóriæ meæ, ut ad reparándum genus humánum exsequerétur commúne consílium, usque ad formam servílem inclinávit incommutábilem Deitátem. Hunc ergo, in quo mihi per ómnia bene compláceo, et cujus prædicatióne maniféstor, cujus humilitáte claríficor, incunctánter audíte: quia ipse est véritas et vita, ipse virtus mea atque sapiéntia.\par
\switchcolumn
\EN

\switchcolumn*

\LA
\begin{Resp}\Rsign Cum audísset Jacob quod Esau veníret contra eum, divísit fílios suos et uxóres, dicens: Si percússerit Esau unam turmam, salvábitur áltera. \Star{} Líbera me, Dómine, qui dixísti mihi: * Multiplicábo semen tuum sicut stellas cæli, et sicut arénam maris, quæ præ multitúdine numerári non potest.\end{Resp}
\begin{Resp}\Vsign Dómine, qui dixísti mihi, Revértere in terram nativitátis tuæ: Dómine, qui pascis me a juventúte mea.\end{Resp}
\begin{Resp}\Rsign Líbera me, Dómine, qui dixísti mihi.\end{Resp}
\begin{Resp}\Vsign Glória Patri, et Fílio, \Star{} et Spirítui Sancto.\end{Resp}
\begin{Resp}\Rsign Multiplicábo semen tuum sicut stellas cæli, et sicut arénam maris, quæ præ multitúdine numerári non potest.\end{Resp}
\switchcolumn
\EN
\begin{Resp}\Rsign When Jacob heard that Esau came to meet him, he divided his sons and his wives, saying: If Esau smite the one company, then the other shall escape. \Star{} Deliver me, O Lord, Who saidst unto me: I will multiply thy seed as the stars of heaven, and as the sand of the sea, which cannot be numbered for multitude.\end{Resp}
\begin{Resp}\Vsign O Lord, Who saidst unto me: Return unto thy country O Lord, Which feedest me still from my youth up.\end{Resp}
\begin{Resp}\Rsign Deliver me, O Lord, Who saidst unto me:\end{Resp}
\begin{Resp}\Vsign Glory be to the Father, and to the Son, \Star{} and to the Holy Ghost.\end{Resp}
\begin{Resp}\Rsign I will multiply thy seed as the stars of heaven, and as the sand of the sea, which cannot be numbered for multitude.\end{Resp}
\switchcolumn*

\end{paracol}

\Office{Feria Secunda infra Hebdomadam II in Quadragesima}{Monday of the Second Week of Lent}{Feria major}
\addcontentsline{toc}{subsection}{Feria Secunda infra Hebdomadam II in Quadragesima\enspace\textit{Monday of the Second Week of Lent}}
\begin{paracol}{2}
\LA
\LessonHead{Lectio I}
\switchcolumn
\EN
\LessonHead{Lesson I}
\switchcolumn*

\LA
\LessonTitle{Léctio sancti Evangélii secúndum Joánnem}
\switchcolumn
\EN
\LessonTitle{Continuation of the Holy Gospel according to John}
\switchcolumn*

\LA
\Cite{Joannes 8:21-29}
\switchcolumn
\EN
\Cite{John 8:21-29}
\switchcolumn*

\LA
\noindent In illo témpore: Dixit Jesus turbis Judæórum: Ego vado, et quærétis me, et in peccáto vestro moriémini. Et réliqua.\par
\switchcolumn
\EN
\noindent In that time, Jesus said to the multitude of Jews: I go, and you shall seek me, and you shall die in your sin. And so on.\par
\switchcolumn*

\LA
\LessonTitle{Homilía sancti Augustíni Epíscopi}
\switchcolumn
\EN
\LessonTitle{Homily on this passage by St. Augustine, Bishop (of Hippo.)}
\switchcolumn*

\LA
\Source{Tract. 38 in Joann. post init.}
\switchcolumn
\EN
\Source{Tract 38 on John}
\switchcolumn*

\LA
\noindent Locútus est Dóminus Judǽis, dicens: Ego vado. Christo enim Dómino mors proféctio fuit illo, unde vénerat, et unde non discésserat. Ego, inquit, vado, et quærétis me, non desidério, sed ódio. Nam illum póstea quam abscéssit ab óculis hóminum, inquisiérunt et qui óderant, et qui amábant: illi persequéndo, isti habére cupiéndo. In Psalmis ait ipse Dóminus per Prophétam: Périit fuga a me, et non est qui requírat ánimam meam. Et íterum ait álio loco in Psalmo: Confundántur et revereántur requiréntes ánimam meam.\par
\switchcolumn
\EN
\noindent The Lord spake unto the Jews, saying: I go My way for, to the Lord Christ, death was a departure to that place whence He had come, and whence He had never departed. I go My way, saith He, and ye shall seek Me not from love, but from hatred. Yea after He had withdrawn Himself from the sight of men, two classes sought Him, even they that loved, and they that hated Him; the one because they longed for His presence, the other because they were fain to hunt Him down. In the Psalms the Lord Himself saith by His Prophet: Refuge failed me, and no man cared for my soul. (Ps. cxli. 5.) And again He said in another Psalm: Let them be confounded and put to shame that seek after my soul. (Ps. xxxiv.)\par
\switchcolumn*

\LA
\begin{Resp}\Rsign Dum iret Jacob de Bersabée, et pérgeret Haram, locútus est ei Dóminus, dicens: \Star{} Terram, in qua dormis, tibi dabo, et sémini tuo.\end{Resp}
\begin{Resp}\Vsign Ædificávit ex lapídibus altáre in honórem Dómini, fundens óleum désuper: et benedíxit eum Deus, dicens.\end{Resp}
\begin{Resp}\Rsign Terram, in qua dormis, tibi dabo, et sémini tuo.\end{Resp}
\switchcolumn
\EN
\begin{Resp}\Rsign While as Jacob went from Beersheba, and hasted unto Haran, the Lord spake unto him, saying \Star{} The land whereon thou sleepest, to thee will I give it, and to thy seed.\end{Resp}
\begin{Resp}\Vsign He built an altar of stones unto the Name of the Lord, and poured oil upon the top of it; and God blessed him and said,:\end{Resp}
\begin{Resp}\Rsign The land whereon thou sleepest, to thee will I give it, and to thy seed.\end{Resp}
\switchcolumn*

\LA
\LessonHead{Lectio II}
\switchcolumn
\EN
\LessonHead{Lesson II}
\switchcolumn*

\LA
\noindent Culpávit, qui non requírerent: damnávit requiréntes. Bonum est enim quǽrere ánimam Christi, sed quo modo eam quæsiérunt discípuli: et malum est quǽrere ánimam Christi, sed quo modo eam Judǽi quæsiérunt: illi enim ut habérent, isti ut pérderent. Dénique istis, quia sic quærébant more malo, corde pervérso, quid secutus adjúnxit? Quærétis me; et ne putétis, quia bene me quærétis, in peccáto vestro moriémini. Hoc est Christum male quǽrere, in peccáto suo mori: hoc est illum odísse, per quem possit solum salvus esse.\par
\switchcolumn
\EN
\noindent Thus doth He blame them that seek not, and condemn such as seek. Yea, it is a good thing to seek the soul of Christ, as the disciples sought it; and an evil thing to seek it, as the Jews sought it; the first sought it to possess, the second to destroy it. What then doth He bid us know will be the reward of such as seek it evilly in a perverse heart? Ye shall seek Me, and lest ye think that ye shall do well so to seek Me, I tell you that ye shall die in your sins. To seek Christ with bad intent, is as much as to die in sin, for it is to hate Him through Whom alone we can be saved.\par
\switchcolumn*

\LA
\begin{Resp}\Rsign Appáruit Deus Jacob, et benedíxit eum, et dixit: Ego sum Deus Bethel, ubi unxísti lápidem, et votum vovísti mihi: \Star{} Créscere te fáciam, et multiplicábo te.\end{Resp}
\begin{Resp}\Vsign Vere Dóminus est in loco isto, et ego nesciébam.\end{Resp}
\begin{Resp}\Rsign Créscere te fáciam, et multiplicábo te.\end{Resp}
\switchcolumn
\EN
\begin{Resp}\Rsign God appeared unto Jacob, and blessed him, and said I am the God of Bethel, where thou anointedst the pillar, and where thou vowedst a vow unto Me. \Star{} I will make thee fruitful, and multiply thee.\end{Resp}
\begin{Resp}\Vsign Surely the Lord is in this place, and I knew it not.\end{Resp}
\begin{Resp}\Rsign I will make thee fruitful, and multiply thee.\end{Resp}
\switchcolumn*

\LA
\LessonHead{Lectio III}
\switchcolumn
\EN
\LessonHead{Lesson III}
\switchcolumn*

\LA
\noindent Cum enim hómines, quorum spes in Deo est, non débeant mala réddere nec pro malis; reddébant isti mala pro bonis. Prænuntiávit ergo illis Dóminus, dixítque senténtiam prǽscius, quod in suo peccáto moreréntur. Deínde adjúnxit: Quo ego vado, vos non potéstis veníre. Hoc et discípulis suis álio loco dixit: nec tamen eis dixit: In peccáto vestro moriémini. Quid autem dixit? quod et istis: Quo ego vado, vos non potéstis veníre. Non ábstulit spem, sed prædíxit dilatiónem. Quando enim hoc discípulis Dóminus loquebátur, tunc non póterant veníre, quo ille ibat, sed póstea ventúri erant: isti autem nunquam, quibus prǽscius dixit, In peccáto vestro moriémini.\par
\switchcolumn
\EN
\noindent Whereas men whose hope is in God ought to return good even for evil, those men returned evil for good. The Lord therefore told them beforehand, and, because He knew it, He let them know their coming end, how that they should die in their sins. Then He said farther Whither I go, ye cannot come. This He said in another place (xiii. 33) to His disciples, but He never said to them: Ye shall die in your sins. What said He? The same words as to the Jews: Whither I go, ye cannot come. Yet, to the disciples, these words only deferred, they cut not away hope for they, though for a little while they could not come whither He was to go, were yet in the end to go there. Not so they, to whom He foretold and said: Ye shall die in your sins.\par
\switchcolumn*

\LA
\begin{Resp}\Rsign Det tibi Deus de rore cæli et de pinguédine terræ abundántiam: sérviant tibi tribus et pópuli: \Star{} Esto dóminus fratrum tuórum.\end{Resp}
\begin{Resp}\Vsign Et incurvéntur ante te fílii matris tuæ.\end{Resp}
\begin{Resp}\Rsign Esto dóminus fratrum tuórum.\end{Resp}
\begin{Resp}\Vsign Glória Patri, et Fílio, \Star{} et Spirítui Sancto.\end{Resp}
\begin{Resp}\Rsign Esto dóminus fratrum tuorum.\end{Resp}
\switchcolumn
\EN
\begin{Resp}\Rsign God give thee of the dew of heaven and the fatness of the earth. Let people and nations serve thee. \Star{} Be lord over thy brethren.\end{Resp}
\begin{Resp}\Vsign And let thy mother's sons bow down to thee.\end{Resp}
\begin{Resp}\Rsign Be lord over thy brethren.\end{Resp}
\begin{Resp}\Vsign Glory be to the Father, and to the Son, \Star{} and to the Holy Ghost.\end{Resp}
\begin{Resp}\Rsign Be lord over thy brethren.\end{Resp}
\switchcolumn*

\end{paracol}

\Office{Feria Tertia infra Hebdomadam II in Quadragesima}{Tuesday of the Second Week of Lent}{Feria major}
\addcontentsline{toc}{subsection}{Feria Tertia infra Hebdomadam II in Quadragesima\enspace\textit{Tuesday of the Second Week of Lent}}
\begin{paracol}{2}
\LA
\LessonHead{Lectio I}
\switchcolumn
\EN
\LessonHead{Lesson I}
\switchcolumn*

\LA
\LessonTitle{Léctio sancti Evangélii secúndum Matthǽum}
\switchcolumn
\EN
\LessonTitle{Continuation of the Holy Gospel according to Matthew}
\switchcolumn*

\LA
\Cite{Matt 23:1-12}
\switchcolumn
\EN
\Cite{Matt 23:1-12}
\switchcolumn*

\LA
\noindent In illo témpore: Locútus est Jesus ad turbas, et ad discípulos suos, dicens: Super cáthedram Móysi sedérunt scribæ et pharisǽi. Omnia ergo quæcúmque díxerint vobis, serváte, et fácite: secúndum ópera vero eórum nolíte fácere. Et réliqua.\par
\switchcolumn
\EN
\noindent At that time, Jesus spoke to the multitudes and to his disciples, saying: The scribes and the Pharisees have sitten on the chair of Moses. All things therefore whatsoever they shall say to you, observe and do: but according to their works do ye not; for they say, and do not. And so on.\par
\switchcolumn*

\LA
\LessonTitle{Homilía sancti Hierónymi Presbýteri}
\switchcolumn
\EN
\LessonTitle{Homily by St. Jerome, Priest (at Bethlehem.)}
\switchcolumn*

\LA
\Source{Lib. 4 Comment. in cap. 23 Matthæi}
\switchcolumn
\EN
\Source{Bk. iv Comm. on Matth. xxiii}
\switchcolumn*

\LA
\noindent Quid mansuétius, quid benígnius Dómino? Tentátur a pharisǽis, confringúntur insídiæ eórum, et secúndum Psalmístam: Sagíttæ parvulórum factæ sunt plagæ eórum: et nihilóminus propter sacerdótii et nóminis dignitátem hortátur pópulos, ut subiciántur eis, non ópera, sed doctrínam considerántes. Quod autem ait, Super cáthedram Móysi sedérunt scribæ et pharisǽi: per cáthedram, doctrínam Legis osténdit. Ergo et illud quod dícitur in Psalmo: In cáthedra pestiléntiæ non sedit: et, Cáthedras vendéntium colúmbas evértit: doctrínam debémus accípere.\par
\switchcolumn
\EN
\noindent Was there ever man gentler and kinder than the Lord? The Pharisees tempted Him; their craft was confounded, and, in the words of the Psalmist, The arrows of babes have pierced them, (Ps. lxiii. 8,) and nevertheless, because of the dignity of their priesthood and name, He exhorteth the people to be subject to them, by doing according to their words, though not according to their works. By the words “Moses' seat” we are to understand the teaching of the law. Thus also must we mystically take: “Sitteth in the seat of the scornful”, (Ps. i. 1,) and likewise, “overthrew the seats of them that sold doves”, (Matth. xxi. 12,) to describe doctrine.\par
\switchcolumn*

\LA
\begin{Resp}\Rsign Dum exíret Jacob de terra sua, vidit glóriam Dei, et ait: Quam terríbilis est locus iste! \Star{} Non est hic áliud, nisi domus Dei, et porta cæli.\end{Resp}
\begin{Resp}\Vsign Vere Deus est in loco isto, et ego nesciébam.\end{Resp}
\begin{Resp}\Rsign Non est hic áliud, nisi domus Dei, et porta cæli.\end{Resp}
\switchcolumn
\EN
\begin{Resp}\Rsign As Jacob went out from his own land, he saw the glory of God, and said: How dreadful is this place. \Star{} This is none other but the house of God and this is the gate of heaven.\end{Resp}
\begin{Resp}\Vsign Surely God is in this place, and I knew it not.\end{Resp}
\begin{Resp}\Rsign This is none other but the house of God and this is the gate of heaven.\end{Resp}
\switchcolumn*

\LA
\LessonHead{Lectio II}
\switchcolumn
\EN
\LessonHead{Lesson II}
\switchcolumn*

\LA
\noindent Alligant enim ónera grávia et importabília, et impónunt in húmeros hóminum, dígito autem suo nolunt ea movére. Hoc generáliter advérsus omnes magístros, qui grávia jubent, et minóra non fáciunt. Notándum autem, quod et húmeri, et dígitus, et ónera, et víncula quibus alligántur ónera, spirituáliter intellegénda sunt. Omnia vero ópera sua fáciunt, ut videántur ab homínibus. Quicúmque ígitur ita facit quódlibet, ut videátur ab homínibus, scriba et pharisǽus est.\par
\switchcolumn
\EN
\noindent How they bind heavy burdens, and grievous to be borne, and lay them on men's shoulders, but they themselves will not move them with one of their fingers. This is generally directed against all teachers who command things hard, and themselves do not even things easy. But it is to be remarked that the shoulders, the fingers, and the binding of the burdens, have a spiritual interpretation. But all their works they do for to be seen of men. Whosoever therefore doth anything for to be seen of men, the same is, so far, a Scribe and a Pharisee.\par
\switchcolumn*

\LA
\begin{Resp}\Rsign Si Dóminus Deus meus fúerit mecum in via ista, per quam ego ámbulo, et custodíerit me, et déderit mihi panem ad edéndum, et vestiméntum quo opériar, et revocáverit me cum salúte: \Star{} Erit mihi Dóminus in refúgium, et lapis iste in signum.\end{Resp}
\begin{Resp}\Vsign Surgens ergo mane Jacob, tulit lápidem quem supposúerat cápiti suo, et eréxit in títulum, fundénsque óleum désuper, dixit.\end{Resp}
\begin{Resp}\Rsign Erit mihi Dóminus in refúgium, et lapis iste in signum.\end{Resp}
\switchcolumn
\EN
\begin{Resp}\Rsign If the Lord my God will be with me, in this way that I go, and will keep me, and will give me bread to eat, and raiment to put on, and will bring me again safely. \Star{} The Lord shall be my refuge, and this stone shall be a sign.\end{Resp}
\begin{Resp}\Vsign So Jacob rose up early in the morning, and took the stone that he had put for his pillow, and set it up for a pillar, and poured oil upon the top of it, and said:\end{Resp}
\begin{Resp}\Rsign The Lord shall be my refuge, and this shall be a sign.\end{Resp}
\switchcolumn*

\LA
\LessonHead{Lectio III}
\switchcolumn
\EN
\LessonHead{Lesson III}
\switchcolumn*

\LA
\noindent Dilátant enim phylactéria sua, et magníficant fímbrias. Amant quoque primos recúbitus in cenis, et primas cáthedras in synagógis, et salutatiónes in foro, et vocári ab homínibus Rabbi. Væ nobis míseris, ad quos pharisæórum vítia transiérunt. Dóminus cum dedísset mandáta Legis per Móysen, ad extrémum íntulit: Ligábis ea in manu tua, et erunt immóta ante óculos tuos. Et est sensus: Præcépta mea sint in manu tua, ut ópere compleántur: sint ante óculos tuos, ut die ac nocte meditéris in eis. Hoc pharisǽi male interpretántes, scribébant in membránis decálogum Móysi, id est, decem verba Legis, complicántes ea et ligántes in fronte, et quasi corónam cápiti faciéntes: ut semper ante óculos moveréntur.\par
\switchcolumn
\EN
\noindent They make broad their phylacteries, and enlarge the borders of their garments. And love the uppermost rooms at feasts, and the chief seats in the synagogues, and greetings in the markets, and to be called of men, Rabbi. Woe to us miserable sinners who have inherited the vices of the Pharisees! When the Lord had given the commandments of the law to Moses He added afterwards Thou shalt bind them for a sign upon thine hand, and they shall be as frontlets between thine eyes, (Deut. vi. 8.) The sense of these words is: My Law shall be in thine hand to order whatsoever thou doest, and ever before thine eyes that thou mayest meditate therein day and night. But the Pharisees, by a bad interpretation, were accustomed to write on pieces of parchment the Decalogue of Moses, that is, the Ten Words of the Law, and to tie these pieces of parchment, plaited in a peculiar manner, on their foreheads, so as to make a sort of crown round their heads, which projected in front of their eyes, and always moved before them.\par
\switchcolumn*

\LA
\begin{Resp}\Rsign Erit mihi Dóminus in Deum, et lapis iste quem eréxi in títulum, vocábitur domus Dei: et de univérsis quæ déderis mihi, \Star{} Décimas et hóstias pacíficas ófferam tibi.\end{Resp}
\begin{Resp}\Vsign Si revérsus fúero próspere ad domum patris mei.\end{Resp}
\begin{Resp}\Rsign Décimas et hóstias pacíficas ófferam tibi.\end{Resp}
\begin{Resp}\Vsign Glória Patri, et Fílio, \Star{} et Spirítui Sancto.\end{Resp}
\begin{Resp}\Rsign Décimas et hóstias pacíficas ófferam tibi.\end{Resp}
\switchcolumn
\EN
\begin{Resp}\Rsign The Lord shall be my God, and this stone, which I have set for a pillar, shall be called God's house, and of all that Thou shalt give me; \Star{} I will offer tithes and peace-offerings to thee.\end{Resp}
\begin{Resp}\Vsign If I come again to my father's house in peace.\end{Resp}
\begin{Resp}\Rsign I will offer tithes and peace-offerings unto thee.\end{Resp}
\begin{Resp}\Vsign Glory be to the Father, and to the Son, \Star{} and to the Holy Ghost.\end{Resp}
\begin{Resp}\Rsign I will offer tithes and peace-offerings unto thee.\end{Resp}
\switchcolumn*

\end{paracol}

\Office{Feria Quarta infra Hebdomadam II in Quadragesima}{Wednesday of the Second Week of Lent}{Feria major}
\addcontentsline{toc}{subsection}{Feria Quarta infra Hebdomadam II in Quadragesima\enspace\textit{Wednesday of the Second Week of Lent}}
\begin{paracol}{2}
\LA
\LessonHead{Lectio I}
\switchcolumn
\EN
\LessonHead{Lesson I}
\switchcolumn*

\LA
\LessonTitle{Léctio sancti Evangélii secúndum Matthǽum}
\switchcolumn
\EN
\LessonTitle{Continuation of the Holy Gospel according to Matthew}
\switchcolumn*

\LA
\Cite{Matt 20:17-28}
\switchcolumn
\EN
\Cite{Matt 20:17-28}
\switchcolumn*

\LA
\noindent In illo témpore: Ascéndens Jesus Jerosólymam, assúmpsit duódecim discípulos secréto, et ait illis: Ecce ascéndimus Jerosólymam, et Fílius hóminis tradétur princípibus sacerdótum, et scribis, et condemnábunt eum morte. Et réliqua.\par
\switchcolumn
\EN
\noindent In that time, Jesus, going up to Jerusalem, took the twelve disciples apart, and said to them: Behold we go up to Jerusalem, and the Son of man shall be betrayed to the chief priests and the scribes, and they shall condemn him to death. And so on.\par
\switchcolumn*

\LA
\LessonTitle{Homilía sancti Ambrósii Epíscopi}
\switchcolumn
\EN
\LessonTitle{Homily by St. Ambrose, Bishop (of Milan.)}
\switchcolumn*

\LA
\Source{Lib. 5 de Fide ad Gratianum, cap. 2, post initium}
\switchcolumn
\EN
\Source{Bk. v to Gratian, on Faith, c. ii}
\switchcolumn*

\LA
\noindent Consideráte, quæ mater filiórum Zebedǽi cum fíliis et pro fíliis petat: mater est útique, cui pro filiórum honóre sollícitæ, immoderátior quidem, sed tamen ignoscénda mensúra votórum est. Atque mater ætáte longǽva, stúdio religiósa, solátio destitúta, quæ tunc témporis, quando vel juvánda, vel alénda foret válidæ prolis auxílio, abésse sibi líberos patiebátur, et voluptáti suæ mercédem sequéntium Christum prætúlerat filiórum. Qui prima voce vocáti a Dómino (ut légimus) relíctis rétibus et patre, secúti sunt eum.\par
\switchcolumn
\EN
\noindent Consider what it was that the mother of Zebedee's children came to Christ desiring, with, and for her sons. She was a mother, who, longing for the honour of her sons, preferred a request immoderate, and yet pardonable. She was a mother who, albeit stricken in years and comfortless, at an age when she had sore need of the strength of her offspring to help and keep her, was yet so earnest in godliness and motherly love, that she had liefer suffer the loss of her sons, that they might gain the reward of following Christ still, as we read they had already done, when, at the first call of the Lord, they left their nets and their father, (iv. 21, 22.)\par
\switchcolumn*

\LA
\begin{Resp}\Rsign Dixit Angelus ad Jacob: \Star{} Dimítte me, auróra est. Respóndit ei: Non dimíttam te, nisi benedíxeris mihi. Et benedíxit ei in eódem loco.\end{Resp}
\begin{Resp}\Vsign Cumque surrexísset Jacob, ecce vir luctabátur cum eo usque mane: et cum vidéret quod eum superáre non posset, dixit ad eum.\end{Resp}
\begin{Resp}\Rsign Dimítte me, auróra est. Respóndit ei: Non dimíttam te, nisi benedíxeris mihi. Et benedíxit ei in eódem loco.\end{Resp}
\switchcolumn
\EN
\begin{Resp}\Rsign The Angel said unto Jacob \Star{} Let me go, for the day breaketh. And he said: I will not let thee go, except thou bless me. And he blessed him there.\end{Resp}
\begin{Resp}\Vsign And when Jacob arose, behold there wrestled a man with him, until the breaking of the day and, when he saw that he prevailed not, he said unto him\end{Resp}
\begin{Resp}\Rsign Let me go, for the day breaketh. And he said: I will not let thee go, except thou bless me. And he blessed him there.\end{Resp}
\switchcolumn*

\LA
\LessonHead{Lectio II}
\switchcolumn
\EN
\LessonHead{Lesson II}
\switchcolumn*

\LA
\noindent Hæc ígitur, stúdio matérnæ sedulitátis indulgéntior, obsecrábat Salvatórem, dicens: Ut sédeant hi duo fílii mei, unus ad déxteram tuam, et alter ad sinístram in regno tuo. Etsi error, pietátis tamen error est. Nésciunt enim matérna víscera patiéntiam: etsi voti avára, tamen veniábilis cupíditas, quæ non pecúniæ est ávida, sed grátiæ. Nec inverecúnda petítio, quæ non sibi, sed líberis consulébat. Matrem consideráte, matrem cogitáte.\par
\switchcolumn
\EN
\noindent She, then, yielding to the intensity of her motherly love, besought the Saviour, saying, Grant that these my two sons may sit, the one at thy right hand and the other at thy left hand, in thy kingdom. Although it was a mistake, it was a mistake of love. For a mother's love knoweth no moderation. Yet, although it was a greedy prayer, that was a pardonable greed, which hungered, not for riches, but for grace. Neither was that request shameless which sought, not her own good, but her children's. Remember that she was a mother. Think how that she was a mother.\par
\switchcolumn*

\LA
\begin{Resp}\Rsign Vidi Dóminum fácie ad fáciem: \Star{} Et salva facta est ánima mea.\end{Resp}
\begin{Resp}\Vsign Et dixit mihi: Nequáquam vocáberis Jacob, sed Israël erit nomen tuum.\end{Resp}
\begin{Resp}\Rsign Et salva facta est ánima mea.\end{Resp}
\switchcolumn
\EN
\begin{Resp}\Rsign I have seen God face to face; \Star{} And my life is preserved.\end{Resp}
\begin{Resp}\Vsign And he said unto me: thy name shall be called no more Jacob, but Israel shall be thy name.\end{Resp}
\begin{Resp}\Rsign And my life is preserved.\end{Resp}
\switchcolumn*

\LA
\LessonHead{Lectio III}
\switchcolumn
\EN
\LessonHead{Lesson III}
\switchcolumn*

\LA
\noindent Considerábat Christus matris dilectiónem, quæ filiórum mercéde grandǽvam solabátur senéctam: et desidériis licet fessa matérnis, carissimórum pignórum tolerábat abséntiam. Consideráte étiam féminam, hoc est, sexum fragiliórem, quem Dóminus própria nondum confirmáverat passióne. Consideráte, inquam, Hevæ illíus primæ mulíeris herédem, transfúsa in omnes immoderátæ cupiditátis successióne labéntem: quam Dóminus adhuc próprio sánguine non redémerat, nondum inólitam afféctibus ómnium immódici contra fas honóris appeténtiam suo Christus cruóre dilúerat. Hereditário ígitur múlier delinquébat erróre.\par
\switchcolumn
\EN
\noindent Christ took into His consideration that mother's love of hers, which made her sons' reward the comfort of her own old age, and which could bear the loss of her loved ones, broken as she was by a mother's yearnings. Consider also that she was a woman, that is, of the weaker sex, to which the Lord had not yet given strength by His Passion. Consider, I say, that she was an heiress of Eve, and weakened by that transmission of the unbridled covetousness of the first woman, which the Lord had not yet disarmed by His Blood, even that craving for undue dignity, wherewith all our natures are imbued, and which Christ's Bloodshedding had not yet washed away. She erred indeed, but the mistake was an inherited weakness.\par
\switchcolumn*

\LA
\begin{Resp}\Rsign Cum audísset Jacob quod Esau veníret contra eum, divísit fílios suos et uxóres, dicens: Si percússerit Esau unam turmam, salvábitur áltera. \Star{} Líbera me, Dómine, qui dixísti mihi: * Multiplicábo semen tuum sicut stellas cæli, et sicut arénam maris, quæ præ multitúdine numerári non potest.\end{Resp}
\begin{Resp}\Vsign Dómine, qui dixísti mihi, Revértere in terram nativitátis tuæ: Dómine, qui pascis me a juventúte mea.\end{Resp}
\begin{Resp}\Rsign Líbera me, Dómine, qui dixísti mihi.\end{Resp}
\begin{Resp}\Vsign Glória Patri, et Fílio, \Star{} et Spirítui Sancto.\end{Resp}
\begin{Resp}\Rsign Multiplicábo semen tuum sicut stellas cæli, et sicut arénam maris, quæ præ multitúdine numerári non potest.\end{Resp}
\switchcolumn
\EN
\begin{Resp}\Rsign When Jacob heard that Esau came to meet him, he divided his sons and his wives, saying: If Esau smite the one company, then the other shall escape. \Star{} Deliver me, O Lord, Who saidst unto me: I will multiply thy seed as the stars of heaven, and as the sand of the sea, which cannot be numbered for multitude.\end{Resp}
\begin{Resp}\Vsign O Lord, Who saidst unto me: Return unto thy country O Lord, Which feedest me still from my youth up.\end{Resp}
\begin{Resp}\Rsign Deliver me, O Lord, Who saidst unto me:\end{Resp}
\begin{Resp}\Vsign Glory be to the Father, and to the Son, \Star{} and to the Holy Ghost.\end{Resp}
\begin{Resp}\Rsign I will multiply thy seed as the stars of heaven, and as the sand of the sea, which cannot be numbered for multitude.\end{Resp}
\switchcolumn*

\end{paracol}

\Office{Feria Quinta infra Hebdomadam II in Quadragesima}{Thursday of the Second Week of Lent}{Feria major}
\addcontentsline{toc}{subsection}{Feria Quinta infra Hebdomadam II in Quadragesima\enspace\textit{Thursday of the Second Week of Lent}}
\begin{paracol}{2}
\LA
\LessonHead{Lectio I}
\switchcolumn
\EN
\LessonHead{Lesson I}
\switchcolumn*

\LA
\LessonTitle{Léctio sancti Evangélii secúndum Lucam}
\switchcolumn
\EN
\LessonTitle{Continuation of the Holy Gospel according to Luke}
\switchcolumn*

\LA
\Cite{Luc 16:19-31}
\switchcolumn
\EN
\Cite{Luke 16:19-31}
\switchcolumn*

\LA
\noindent In illo témpore: Dixit Jesus pharisǽis: Homo quidam erat dives, qui induebátur púrpura et bysso: et epulabátur cotídie spléndide. Et réliqua.\par
\switchcolumn
\EN
\noindent In that time, Jesus said to the pharisees: There was a certain rich man, who was clothed in purple and fine linen; and feasted sumptuously every day. And so on.\par
\switchcolumn*

\LA
\LessonTitle{Homilía sancti Gregórii Papæ}
\switchcolumn
\EN
\LessonTitle{Homily by Pope St. Gregory (the Great.)}
\switchcolumn*

\LA
\Source{Homilia 40 in Evangelia}
\switchcolumn
\EN
\Source{40th on the Gospels}
\switchcolumn*

\LA
\noindent Quem, fratres caríssimi, quem dives iste, qui induebátur púrpura et bysso, et epulabátur cotídie spléndide, nisi Judáicum pópulum signat: qui cultum vitæ extérius hábuit, qui accéptæ legis delíciis ad nitórem usus est, non ad utilitátem? Quem vero Lázarus ulcéribus plenus, nisi gentílem pópulum figuráliter éxprimit? Qui dum convérsus ad Deum peccáta confitéri sua non erúbuit, huic vulnus in cute fuit. In cutis quippe vúlnere virus a viscéribus tráhitur, et foras erúmpit.\par
\switchcolumn
\EN
\noindent Whom, dearly beloved brethren, whom are we to understand as signified by that rich man which was clothed in purple and fine linen, and fared sumptuously every day, whom, I ask, are we to understand, but the Jewish people, who had all the outward life of religious ordinances, and who turned the treasure of the law they had received to show and not to use? What but the herd of the Gentiles is figured in Lazarus, full of sores? Whosoever turneth himself to God and is not ashamed to confess his sin, hath his sores on the skin, for in a sore on the skin breaketh out the corruption, which is drawn from within.\par
\switchcolumn*

\LA
\begin{Resp}\Rsign Tolle arma tua, pháretram et arcum, et affer de venatióne tua, ut cómedam: \Star{} Et benedícat tibi ánima mea.\end{Resp}
\begin{Resp}\Vsign Cumque venátu áliquid attúleris, fac mihi inde pulméntum, ut cómedam.\end{Resp}
\begin{Resp}\Rsign Et benedícat tibi ánima mea.\end{Resp}
\switchcolumn
\EN
\begin{Resp}\Rsign Take thy weapons, thy quiver and thy bow, and bring me some of thy venison, that I may eat, \Star{} and my soul may bless thee.\end{Resp}
\begin{Resp}\Vsign And when thou hast taken somewhat, make me thereof savoury meat, that I may eat.\end{Resp}
\begin{Resp}\Rsign And my soul may bless thee.\end{Resp}
\switchcolumn*

\LA
\LessonHead{Lectio II}
\switchcolumn
\EN
\LessonHead{Lesson II}
\switchcolumn*

\LA
\noindent Quid est ergo peccatórum conféssio, nisi quædam vúlnerum rúptio? Quia peccáti virus salúbriter aperítur in confessióne, quod pestífere latébat in mente. Vúlnera étenim cutis in superfíciem trahunt humórem putrédinis. Et confiténdo peccáta, quid áliud ágimus, nisi malum, quod in nobis latébat, aperímus? Sed Lázarus vulnerátus cupiébat saturári de micis, quæ cadébant de mensa dívitis, et nemo illi dabat: quia gentílium quemque ad cognitiónem legis admíttere supérbus ille pópulus despiciébat.\par
\switchcolumn
\EN
\noindent What is, then, the confession of our sins but the breaking out of our sores? The corrupt matter of sin is healthily opened in confession, instead of remaining in the mind to rot it. Open sores on the skin bring the poisonous matter to the surface, and when we confess our sins, what do we but open up the evil that there is lurking in us? But Lazarus desired to be fed with the crumbs which fell from the rich man's table, and no man gave unto him; even so did that proud people scorn to admit a Gentile to the knowledge of their law.\par
\switchcolumn*

\LA
\begin{Resp}\Rsign Ecce odor fílii mei sicut odor agri pleni, cui benedíxit Dóminus: créscere te fáciat Deus meus sicut arénam maris: \Star{} Et donet tibi de rore cæli benedictiónem.\end{Resp}
\begin{Resp}\Vsign Deus autem omnípotens benedícat tibi, atque multíplicet.\end{Resp}
\begin{Resp}\Rsign Et donet tibi de rore cæli benedictiónem.\end{Resp}
\switchcolumn
\EN
\begin{Resp}\Rsign See the smell of my son is as the smell of a field which the Lord hath blessed. May my God multiply thee as the sand of the sea. \Star{} And give thee a blessing of the dew of heaven.\end{Resp}
\begin{Resp}\Vsign And God Almighty bless thee, and multiply thee.\end{Resp}
\begin{Resp}\Rsign And give thee a blessing of the dew of heaven.\end{Resp}
\switchcolumn*

\LA
\LessonHead{Lectio III}
\switchcolumn
\EN
\LessonHead{Lesson III}
\switchcolumn*

\LA
\noindent Qui dum doctrínam legis non ad caritátem hábuit, sed ad elatiónem, quasi de accéptis ópibus túmuit: et quia ei verba defluébant de sciéntia, quasi micæ cadébant de mensa. At contra, jacéntis páuperis vúlnera lingébant canes. Nonnúnquam solent in sacro elóquio, per canes prædicatóres intéllegi. Canum étenim lingua, vulnus dum lingit, curat: quia et doctóres sancti, dum in confessióne peccáti nostri nos ínstruunt, quasi vulnus mentis per linguam tangunt.\par
\switchcolumn
\EN
\noindent The teaching of the law moved them to pride, and not to love, as though they swelled with selfimportance at the thought of their riches, and the words which some Gentiles caught of their knowledge were as crumbs falling from their sumptuous table. On the other hand, the dogs came and licked the sores of the beggar that was laid at their gate. Sometimes in Holy Writ, under the figure of dogs, preachers are understood. A dog's tongue healeth the sore which it licketh, and so do holy teachers, when we confess our sins, and they speak to us, mollify by their tongues the sores of our souls.\par
\switchcolumn*

\LA
\begin{Resp}\Rsign Det tibi Deus de rore cæli et de pinguédine terræ abundántiam: sérviant tibi tribus et pópuli: \Star{} Esto dóminus fratrum tuórum.\end{Resp}
\begin{Resp}\Vsign Et incurvéntur ante te fílii matris tuæ.\end{Resp}
\begin{Resp}\Rsign Esto dóminus fratrum tuórum.\end{Resp}
\begin{Resp}\Vsign Glória Patri, et Fílio, \Star{} et Spirítui Sancto.\end{Resp}
\begin{Resp}\Rsign Esto dóminus fratrum tuorum.\end{Resp}
\switchcolumn
\EN
\begin{Resp}\Rsign God give thee of the dew of heaven and the fatness of the earth. Let people and nations serve thee. \Star{} Be lord over thy brethren.\end{Resp}
\begin{Resp}\Vsign And let thy mother's sons bow down to thee.\end{Resp}
\begin{Resp}\Rsign Be lord over thy brethren.\end{Resp}
\begin{Resp}\Vsign Glory be to the Father, and to the Son, \Star{} and to the Holy Ghost.\end{Resp}
\begin{Resp}\Rsign Be lord over thy brethren.\end{Resp}
\switchcolumn*

\end{paracol}

\Office{Feria Sexta infra Hebdomadam II in Quadragesima}{Friday of the Second Week of Lent}{Feria major}
\addcontentsline{toc}{subsection}{Feria Sexta infra Hebdomadam II in Quadragesima\enspace\textit{Friday of the Second Week of Lent}}
\begin{paracol}{2}
\LA
\LessonHead{Lectio I}
\switchcolumn
\EN
\LessonHead{Lesson I}
\switchcolumn*

\LA
\LessonTitle{Léctio sancti Evangélii secúndum Matthǽum}
\switchcolumn
\EN
\LessonTitle{Continuation of the Holy Gospel according to Matthew}
\switchcolumn*

\LA
\Cite{Matt 21:33-46}
\switchcolumn
\EN
\Cite{Matt 21:33-46}
\switchcolumn*

\LA
\noindent In illo témpore: Dixit Jesus turbis Judæórum, et princípibus sacerdótum parábolam hanc: Homo erat paterfamílias, qui plantávit víneam, et sepem circúmdedit ei. Et réliqua.\par
\switchcolumn
\EN
\noindent In that time, Jesus said to the multitude of Jews and the chief priests: Hear ye another parable. There was a man an householder, who planted a vineyard, and made a hedge round about it. And so on.\par
\switchcolumn*

\LA
\LessonTitle{Homilía sancti Ambrósii Epíscopi}
\switchcolumn
\EN
\LessonTitle{Homily by St. Ambrose, Bishop (of Milan.)}
\switchcolumn*

\LA
\Source{Lib. 9 in cap. 20 Lucæ}
\switchcolumn
\EN
\Source{Bk. ix on Luke xx}
\switchcolumn*

\LA
\noindent Pleríque várias significatiónes de víneæ appellatióne derívant: sed evidénter Isaías víneam Dómini Sábaoth, domum Israël esse memorávit. Hanc víneam quis álius, nisi Deus, cóndidit? Hic est ergo qui eam locávit colónis, et ipse péregre fuit: non quia ex loco ad locum proféctus est Dóminus, qui ubíque semper præsens est: sed quia est præséntior diligéntibus, negligéntibus abest. Multis tempóribus ábfuit, ne præprópera viderétur exáctio. Nam quo indulgéntior liberálitas, eo inexcusabílior pervicácia.\par
\switchcolumn
\EN
\noindent Many derive diverse spiritual meanings from the term vineyard, but Isaias giveth us to know that the vineyard of the Lord of Sabaoth is the house of Israel. (v. 7.) Who but God planted that vineyard? He it was that let it out to husbandmen, and went into a far country; not that the Lord, Who is everywhere present, moveth from place to place; but because He is nigh unto them that seek Him, and from such as regard Him not He standeth afar off. For a long time He tarried away, lest He might seem to ask too early for the fruits of His vineyard. For where kindness is greatest, there ingratitude is worst.\par
\switchcolumn*

\LA
\begin{Resp}\Rsign Dum exíret Jacob de terra sua, vidit glóriam Dei, et ait: Quam terríbilis est locus iste! \Star{} Non est hic áliud, nisi domus Dei, et porta cæli.\end{Resp}
\begin{Resp}\Vsign Vere Deus est in loco isto, et ego nesciébam.\end{Resp}
\begin{Resp}\Rsign Non est hic áliud, nisi domus Dei, et porta cæli.\end{Resp}
\switchcolumn
\EN
\begin{Resp}\Rsign As Jacob went out from his own land, he saw the glory of God, and said: How dreadful is this place. \Star{} This is none other but the house of God and this is the gate of heaven.\end{Resp}
\begin{Resp}\Vsign Surely God is in this place, and I knew it not.\end{Resp}
\begin{Resp}\Rsign This is none other but the house of God and this is the gate of heaven.\end{Resp}
\switchcolumn*

\LA
\LessonHead{Lectio II}
\switchcolumn
\EN
\LessonHead{Lesson II}
\switchcolumn*

\LA
\noindent Unde bene secúndum Matthǽum habes, quia et sepem circúmdedit: hoc est, divínæ custódiæ munitióne vallávit, ne fácile spiritálium patéret incúrsibus bestiárum. Et fodit in ea tórcular. Quómodo intellégimus quid sit tórcular, nisi forte quia Psalmi pro torculáribus inscribúntur; eo quod in his mystéria Domínicæ passiónis, modo musti Sancto fervéntis Spíritu, redundántius æstuáverint? Unde ébrii putabántur, quibus Spíritus Sanctus inundábat. Ergo et hic fodit tórcular, in quod uvæ rationábilis fructus intérior spiritáli infusióne deflúeret.\par
\switchcolumn
\EN
\noindent Therefore it is well written in Matthew, for our instruction, that He hedged it round about, that is, He girded it with the fortifications of His own Divine protection, that it might not easily lie open to the ravages of spiritual wild beasts. And digged a wine-press in it. What sense are we to put upon the wine-press, unless it be that the Psalms are here described under that title, because in them the mysteries of the Lord's Passion flow over like new wine, working under the power of the Holy Ghost? Whence also, they upon whom the Holy Ghost was outpoured were deemed to be drunken (Acts ii. 13.). God therefore digged a wine-press, whereinto the reasonable grapes of inward fruitfulness poured their spiritual richness.\par
\switchcolumn*

\LA
\begin{Resp}\Rsign Si Dóminus Deus meus fúerit mecum in via ista, per quam ego ámbulo, et custodíerit me, et déderit mihi panem ad edéndum, et vestiméntum quo opériar, et revocáverit me cum salúte: \Star{} Erit mihi Dóminus in refúgium, et lapis iste in signum.\end{Resp}
\begin{Resp}\Vsign Surgens ergo mane Jacob, tulit lápidem quem supposúerat cápiti suo, et eréxit in títulum, fundénsque óleum désuper, dixit.\end{Resp}
\begin{Resp}\Rsign Erit mihi Dóminus in refúgium, et lapis iste in signum.\end{Resp}
\switchcolumn
\EN
\begin{Resp}\Rsign If the Lord my God will be with me, in this way that I go, and will keep me, and will give me bread to eat, and raiment to put on, and will bring me again safely. \Star{} The Lord shall be my refuge, and this stone shall be a sign.\end{Resp}
\begin{Resp}\Vsign So Jacob rose up early in the morning, and took the stone that he had put for his pillow, and set it up for a pillar, and poured oil upon the top of it, and said:\end{Resp}
\begin{Resp}\Rsign The Lord shall be my refuge, and this shall be a sign.\end{Resp}
\switchcolumn*

\LA
\LessonHead{Lectio III}
\switchcolumn
\EN
\LessonHead{Lesson III}
\switchcolumn*

\LA
\noindent Ædificávit turrim, vérticem scílicet legis attóllens: atque ita hanc víneam munítam, instrúctam, ornátam, locávit Judǽis. Et témpore frúctuum sérvulos misit. Bene tempus frúctuum pósuit, non provéntuum. Nullus enim fructus éxstitit Judæórum, nullus víneæ hujus provéntus, de qua Dóminus ait: Exspectávi ut fáceret uvas, fecit autem spinas. Itaque non lætítiæ vino, non spiritáli musto, sed cruénto Prophetárum sánguine torculária redundárunt.\par
\switchcolumn
\EN
\noindent And built a tower that is, He raised up the goodly structure of the Law. And so this His vineyard, thus fortified, furnished, and garnished, He gave over to the Jews. And when the time of the fruit drew near, He sent His servants to the husbandmen. Well doth He call it the time of the fruit, not the time of the in-gathering. For the Jews yielded Him no fruit; the Lord had no ingathering from that vineyard of which He said: When I looked that it should bring forth grapes, it brought forth wild grapes. (Isa. v. 4.) Not that wine that maketh glad the heart of man, not with the new wine of the spirit, reeked that wine-press, but with the blood of the Prophets, brutally shed.\par
\switchcolumn*

\LA
\begin{Resp}\Rsign Erit mihi Dóminus in Deum, et lapis iste quem eréxi in títulum, vocábitur domus Dei: et de univérsis quæ déderis mihi, \Star{} Décimas et hóstias pacíficas ófferam tibi.\end{Resp}
\begin{Resp}\Vsign Si revérsus fúero próspere ad domum patris mei.\end{Resp}
\begin{Resp}\Rsign Décimas et hóstias pacíficas ófferam tibi.\end{Resp}
\begin{Resp}\Vsign Glória Patri, et Fílio, \Star{} et Spirítui Sancto.\end{Resp}
\begin{Resp}\Rsign Décimas et hóstias pacíficas ófferam tibi.\end{Resp}
\switchcolumn
\EN
\begin{Resp}\Rsign The Lord shall be my God, and this stone, which I have set for a pillar, shall be called God's house, and of all that Thou shalt give me; \Star{} I will offer tithes and peace-offerings to thee.\end{Resp}
\begin{Resp}\Vsign If I come again to my father's house in peace.\end{Resp}
\begin{Resp}\Rsign I will offer tithes and peace-offerings unto thee.\end{Resp}
\begin{Resp}\Vsign Glory be to the Father, and to the Son, \Star{} and to the Holy Ghost.\end{Resp}
\begin{Resp}\Rsign I will offer tithes and peace-offerings unto thee.\end{Resp}
\switchcolumn*

\end{paracol}

\Office{Sabbato infra Hebdomadam II in Quadragesima}{Saturday of the Second Week of Lent}{Feria major}
\addcontentsline{toc}{subsection}{Sabbato infra Hebdomadam II in Quadragesima\enspace\textit{Saturday of the Second Week of Lent}}
\begin{paracol}{2}
\LA
\LessonHead{Lectio I}
\switchcolumn
\EN
\LessonHead{Lesson I}
\switchcolumn*

\LA
\LessonTitle{Léctio sancti Evangélii secúndum Lucam}
\switchcolumn
\EN
\LessonTitle{Continuation of the Holy Gospel according to Luke}
\switchcolumn*

\LA
\Cite{Luc 15:11-32}
\switchcolumn
\EN
\Cite{Luke 15:11-32}
\switchcolumn*

\LA
\noindent In illo témpore: Dixit Jesus pharisǽis et scribis parábolam istam: Homo quidam hábuit duos fílios: et dixit adolescéntior ex illis patri: Pater, da mihi portiónem substántiæ, quæ me contíngit. Et réliqua.\par
\switchcolumn
\EN
\noindent In that time, Jesus said to the pharisees and scribes in parable: A certain man had two sons: And the younger of them said to his father: Father, give me the portion of substance that falleth to me. And so on.\par
\switchcolumn*

\LA
\LessonTitle{Homilía sancti Ambrósii Epíscopi}
\switchcolumn
\EN
\LessonTitle{Homily by St. Ambrose, Bishop (of Milan.)}
\switchcolumn*

\LA
\Source{Lib. 8 Comment. in cap. 15 Lucæ, post initium}
\switchcolumn
\EN
\Source{Bk. vii Comment. on Luke xv}
\switchcolumn*

\LA
\noindent Vides, quod divínum patrimónium peténtibus datur. Nec putes culpam patris, quod adolescentióri dedit. Nulla Dei regno infírma ætas: nec fides gravátur annis. Ipse certe se judicávit idóneum, qui popóscit. Atque útinam non recessísset a patre, impediméntum nescísset ætátis. Sed posteáquam domum pátriam derelínquens péregre proféctus est, cœpit egére. Mérito ergo prodégit patrimónium, qui recéssit ab Ecclésia.\par
\switchcolumn
\EN
\noindent Thou seest how that the heavenly goods are given to such as seek them. Neither oughtest thou to think the father to blame, because he gave to his younger son. In the kingdom of God there is no age of weakness, neither doth faith wax infirm with years. He, surely, who asked, deemed himself of sufficient age. And would that he had not left his father! then had he been ignorant of the obstacle of his age! But after that he had left his father's house, and had gone into a far country, he began to be in want. Well is he said to have wasted his substance, who hath cut himself off from the Church!\par
\switchcolumn*

\LA
\begin{Resp}\Rsign Pater, peccávi in cælum, et coram te: jam non sum dignus vocári fílius tuus: \Star{} Fac me sicut unum ex mercenáriis tuis.\end{Resp}
\begin{Resp}\Vsign Quanti mercenárii in domo patris mei abúndant pánibus, ego autem hic fame péreo! Surgam, et ibo ad patrem meum, et dicam ei.\end{Resp}
\begin{Resp}\Rsign Fac me sicut unum ex mercenáriis tuis.\end{Resp}
\switchcolumn
\EN
\begin{Resp}\Rsign Father, I have sinned against heaven and before thee, and am no more worthy to be called thy son. \Star{} Make me as one of thine hired servants.\end{Resp}
\begin{Resp}\Vsign How many hired servants of my father's have bread enough and to spare, and I perish with hunger! I will arise and go to my father and will say unto him\end{Resp}
\begin{Resp}\Rsign Make me as one of thine hired servants.\end{Resp}
\switchcolumn*

\LA
\LessonHead{Lectio II}
\switchcolumn
\EN
\LessonHead{Lesson II}
\switchcolumn*

\LA
\noindent Péregre proféctus est in regiónem longínquam. Quid longínquius, quam a se recédere: nec regiónibus, sed móribus separári: stúdiis discrétum esse, non terris; et quasi interfúso luxúriæ sæculáris æstu, divórtia habére sanctórum? Etenim qui se a Christo séparat, exsul est pátriæ, civis est mundi. Sed nos non sumus ádvenæ atque peregríni, sed cives sumus Sanctórum, et doméstici Dei. Qui enim erámus longe, facti sumus prope in sánguine Christi. Non invideámus de longínqua regióne remeántibus: quia et nos fúimus in regióne longínqua, sicut Isaías docet. Sic enim habes: Qui sedébant in regióne umbræ mortis, lux orta est illis. Régio ergo longínqua, umbra est mortis.\par
\switchcolumn
\EN
\noindent He took his journey into a far country. No man can go farther than to abandon his own better self, to leave, not his country, but his morals, and, as it were, in an hideous fever of lust after the world, to divorce himself from the ties that bind him to holy things. Yea, he that turneth his back on Christ, banisheth himself from his Fatherland, and becometh a citizen of the world. But we are no more strangers and foreigners, but fellowcitizens with the saints, and of the household of God, since we who sometimes were afar off, are made nigh by the Blood of Christ. (Eph. ii. 19, 13.) Let us not envy the pleasures of them who remain in the far country. We too have once been there, but, as saith Isaiah, they that dwelt in the land of the shadow of death, upon them hath the light shined. (ix. 2.) And that far country is the land of the shadow of death.\par
\switchcolumn*

\LA
\begin{Resp}\Rsign Vidi Dóminum fácie ad fáciem: \Star{} Et salva facta est ánima mea.\end{Resp}
\begin{Resp}\Vsign Et dixit mihi: Nequáquam vocáberis Jacob, sed Israël erit nomen tuum.\end{Resp}
\begin{Resp}\Rsign Et salva facta est ánima mea.\end{Resp}
\switchcolumn
\EN
\begin{Resp}\Rsign I have seen God face to face; \Star{} And my life is preserved.\end{Resp}
\begin{Resp}\Vsign And he said unto me: thy name shall be called no more Jacob, but Israel shall be thy name.\end{Resp}
\begin{Resp}\Rsign And my life is preserved.\end{Resp}
\switchcolumn*

\LA
\LessonHead{Lectio III}
\switchcolumn
\EN
\LessonHead{Lesson III}
\switchcolumn*

\LA
\noindent Nos autem, quibus spíritus ante fáciem Christus est Dóminus, in umbra vívimus Christi. Et ídeo dicit Ecclésia: In umbra ejus concupívi, et sedi. Ille ígitur vivéndo luxurióse, consúmpsit ómnia ornaménta natúræ. Unde tu, qui accepísti imáginem Dei, qui habes similitúdinem ejus, noli eam irrationábili fœditáte consúmere. Opus Dei es: noli ligno dícere, Pater meus es tu: ne accípias similitúdinem ligni, quia scriptum est: Símiles illis fiant qui fáciunt ea.\par
\switchcolumn
\EN
\noindent But we to whom the Lord Christ is the breath of life, are alive under the shadow of Christ. And therefore it is that the Church saith: I sat down under His shadow with great delight. (Cant. ii. 3.) The prodigal son by riotous living wasted all the gifts of nature. Take warning, O thou who art made in the image and likeness of God, lest thou waste the same by brutish wallowing. Thou art the work of God; say not to a stock: Thou art my father, (Jer. ii. 27,) lest thou grow into the likeness of a stock, as it is written: They that make them are like unto them. (Ps. cxiii. 16.)\par
\switchcolumn*

\LA
\begin{Resp}\Rsign Cum audísset Jacob quod Esau veníret contra eum, divísit fílios suos et uxóres, dicens: Si percússerit Esau unam turmam, salvábitur áltera. \Star{} Líbera me, Dómine, qui dixísti mihi: * Multiplicábo semen tuum sicut stellas cæli, et sicut arénam maris, quæ præ multitúdine numerári non potest.\end{Resp}
\begin{Resp}\Vsign Dómine, qui dixísti mihi, Revértere in terram nativitátis tuæ: Dómine, qui pascis me a juventúte mea.\end{Resp}
\begin{Resp}\Rsign Líbera me, Dómine, qui dixísti mihi.\end{Resp}
\begin{Resp}\Vsign Glória Patri, et Fílio, \Star{} et Spirítui Sancto.\end{Resp}
\begin{Resp}\Rsign Multiplicábo semen tuum sicut stellas cæli, et sicut arénam maris, quæ præ multitúdine numerári non potest.\end{Resp}
\switchcolumn
\EN
\begin{Resp}\Rsign When Jacob heard that Esau came to meet him, he divided his sons and his wives, saying: If Esau smite the one company, then the other shall escape. \Star{} Deliver me, O Lord, Who saidst unto me: I will multiply thy seed as the stars of heaven, and as the sand of the sea, which cannot be numbered for multitude.\end{Resp}
\begin{Resp}\Vsign O Lord, Who saidst unto me: Return unto thy country O Lord, Which feedest me still from my youth up.\end{Resp}
\begin{Resp}\Rsign Deliver me, O Lord, Who saidst unto me:\end{Resp}
\begin{Resp}\Vsign Glory be to the Father, and to the Son, \Star{} and to the Holy Ghost.\end{Resp}
\begin{Resp}\Rsign I will multiply thy seed as the stars of heaven, and as the sand of the sea, which cannot be numbered for multitude.\end{Resp}
\switchcolumn*

\end{paracol}

\Office{Dominica III in Quadragesima}{Third Sunday of Lent}{Semiduplex I classis}
\addcontentsline{toc}{subsection}{Dominica III in Quadragesima\enspace\textit{Third Sunday of Lent}}
\begin{paracol}{2}
\LA
\LessonHead{Lectio I}
\switchcolumn
\EN
\LessonHead{Lesson I}
\switchcolumn*

\LA
\LessonTitle{De libro Génesis}
\switchcolumn
\EN
\LessonTitle{Lesson from the book of Genesis}
\switchcolumn*

\LA
\Cite{Gen 37:2-10}
\switchcolumn
\EN
\Cite{Gen 37:2-10}
\switchcolumn*

\LA
\noindent Joseph, cum sédecim esset annórum, pascébat gregem cum frátribus suis adhuc puer: et erat cum fíliis Balæ et Zelphæ uxórum patris sui: accusavítque fratres suos apud patrem crímine péssimo. Israël autem diligébat Joseph super omnes fílios suos, eo quod in senectúte genuísset eum: fecítque ei túnicam polýmitam. Vidéntes autem fratres ejus quod a patre plus cunctis fíliis amarétur, óderant eum, nec póterant ei quidquam pacífice loqui. Accidit quoque ut visum sómnium reférret frátribus suis: quæ causa majóris ódii seminárium fuit. Dixítque ad eos: Audíte sómnium meum quod vidi: putábam nos ligáre manípulos in agro: et quasi consúrgere manípulum meum, et stare, vestrósque manípulos circumstántes adoráre manípulum meum. Respondérunt fratres ejus: Numquid rex noster eris? aut subiciémur ditióni tuæ? Hæc ergo causa somniórum atque sermónum, invídiæ et ódii fómitem ministrávit. Aliud quoque vidit sómnium, quod narrans frátribus, ait: Vidi per sómnium, quasi solem, et lunam, et stellas úndecim adoráre me. Quod cum patri suo, et frátribus retulísset, increpávit eum pater suus, et dixit: Quid sibi vult hoc sómnium quod vidísti? num ego et mater tua, et fratres tui adorábimus te super terram?\par
\switchcolumn
\EN
\noindent And these are his generations: Joseph, when he was sixteen years old, was feeding the flock with his brethren, being but a boy: and he was with the sons of Bala and of Zelpha his father's wives: and he accused his brethren to his father of a most wicked crime. Now Israel loved Joseph above all his sons, because he had him in his old age: and he made him a coat of diverse colours. And his brethren seeing that he was loved by his father, more than all his sons, hated him, and could not speak peaceably to him. Now it fell out also that he told his brethren a dream, that he had dreamed: which occasioned them to hate him the more. And he said to them: Hear my dream which I dreamed. I thought we were binding sheaves in the field: and my sheaf arose as it were, and stood, and your sheaves standing about, bowed down before my sheaf. His brethren answered Shalt thou be our king? or shall we be subject to thy dominion? Therefore this matter of his dreams and words ministered nourishment to their envy and hatred. He dreamed also another dream, which he told his brethren, saying: I saw in a dream, as it were the sun, and the moon, and eleven stars worshipping me. And when he had told this to his father and brethren, his father rebuked him, and said: What meaneth this dream that thou hast dreamed? shall I and thy mother, and thy brethren worship thee upon the earth?\par
\switchcolumn*

\LA
\begin{Resp}\Rsign Vidéntes Joseph a longe, loquebántur mútuo fratres, dicéntes: Ecce somniátor venit: \Star{} Veníte, occidámus eum, et videámus si prosint illi sómnia sua.\end{Resp}
\begin{Resp}\Vsign Cumque vidíssent Joseph fratres sui, quod a patre cunctis frátribus plus amarétur, óderant eum, nec póterant ei quidquam pacífice loqui, unde et dicébant.\end{Resp}
\begin{Resp}\Rsign Veníte, occidámus eum, et videámus si prosint illi sómnia sua.\end{Resp}
\switchcolumn
\EN
\begin{Resp}\Rsign And when his brethren saw Joseph afar off, they said one to another: Behold, this dreamer cometh. \Star{} Come, let us slay him; and we shall see what will become of his dreams.\end{Resp}
\begin{Resp}\Vsign And when his brethren saw that their father loved Joseph more than all his brethren, they hated him, and could not speak peaceably unto him; therefore they said:\end{Resp}
\begin{Resp}\Rsign Come, let us slay him; and we shall see what will become of his dreams.\end{Resp}
\switchcolumn*

\LA
\LessonHead{Lectio II}
\switchcolumn
\EN
\LessonHead{Lesson II}
\switchcolumn*

\LA
\Cite{Gen 37:11-20}
\switchcolumn
\EN
\Cite{Gen 37:11-20}
\switchcolumn*

\LA
\noindent Invidébant ei ígitur fratres sui: pater vero rem tácitus considerábat. Cumque fratres illíus in pascéndis grégibus patris moraréntur in Sichem, dixit ad eum Israël: Fratres tui pascunt oves in Síchimis: veni, mittam te ad eos. Quo respondénte, Præsto sum, ait ei: Vade, et vide si cuncta próspera sint erga fratres tuos, et pécora: et renúntia mihi quid agátur. Missus de valle Hebron, venit in Sichem: invenítque eum vir errántem in agro, et interrogávit quid quǽreret. At ille respóndit: Fratres meos quæro: índica mihi ubi pascant greges. Dixítque ei vir: Recessérunt de loco isto: audívi autem eos dicéntes: Eámus in Dóthain. Perréxit ergo Joseph post fratres suos, et invénit eos in Dóthain. Qui cum vidíssent eum procul, ántequam accéderet ad eos, cogitavérunt illum occídere: et mútuo loquebántur: Ecce somniátor venit: veníte, occidámus eum, et mittámus in cistérnam véterem: dicemúsque: Fera péssima devorávit eum: et tunc apparébit quid illi prosint sómnia sua.\par
\switchcolumn
\EN
\noindent His brethren therefore envied him: but his father considered the thing with himself. And when his brethren abode in Sichem feeding their father's flocks, Israel said to him: thy brethren feed the sheep in Sichem: come, I will send thee to them. And when he answered: I am ready: he said to him: Go, and see if all things be well with thy brethren, and the cattle: and bring me word again what is doing. So being sent from the vale of Hebron, he came to Sichem: And a man found him there wandering in the field, and asked what he sought. But he answered: I seek my brethren; tell me where they feed the flocks. And the man said to him: They are departed from this place: for I heard them say: Let us go to Dothain. And Joseph went forward after his brethren, and found them in Dothain. And when they saw him afar off, before he came nigh them, they thought to kill him. And said one to another: Behold the dreamer cometh. Come, let us kill him, and cast him into some old pit and we will say: Some evil beast hath devoured him: and then it shall appear what his dreams avail him.\par
\switchcolumn*

\LA
\begin{Resp}\Rsign Dixit Judas frátribus suis: Ecce Ismaëlítæ tránseunt; veníte, venumdétur, et manus nostræ non polluántur: \Star{} Caro enim et frater noster est.\end{Resp}
\begin{Resp}\Vsign Quid enim prodest, si occidérimus fratrem nostrum, et celavérimus sánguinem ipsíus? mélius est ut venumdétur.\end{Resp}
\begin{Resp}\Rsign Caro enim et frater noster est.\end{Resp}
\switchcolumn
\EN
\begin{Resp}\Rsign Judah said unto his brethren: Behold, the Ismaelites pass by; come, let us sell him, and let not our hands be defiled. \Star{} For he is our flesh, and our brother.\end{Resp}
\begin{Resp}\Vsign What profit is it if we slay our brother, and conceal his blood? It is better to sell him.\end{Resp}
\begin{Resp}\Rsign For he is our flesh, and our brother.\end{Resp}
\switchcolumn*

\LA
\LessonHead{Lectio III}
\switchcolumn
\EN
\LessonHead{Lesson III}
\switchcolumn*

\LA
\Cite{Gen 37:21-28}
\switchcolumn
\EN
\Cite{Gen 37:21-28}
\switchcolumn*

\LA
\noindent Audiens autem hoc Ruben, nitebátur liberáre eum de mánibus eórum, et dicébat: Non interficiátis ánimam ejus, nec effundátis sánguinem: sed proícite eum in cistérnam hanc, quæ est in solitúdine, manúsque vestras serváte innóxias: hoc autem dicébat, volens erípere eum de mánibus eórum, et réddere patri suo. Conféstim ígitur ut pervénit ad fratres suos, nudavérunt eum túnica talári et polýmita: miserúntque eum in cistérnam véterem, quæ non habébat aquam. Et sedéntes ut coméderent panem, vidérunt Ismaëlítas viatóres veníre de Gálaad, et camélos eórum portántes arómata, et resínam, et stacten in Ægýptum. Dixit ergo Judas frátribus suis: Quid nobis prodest si occidérimus fratrem nostrum, et celavérimus sánguinem ipsíus? mélius est ut venundétur Ismaëlítis, et manus nostræ non polluántur: frater enim et caro nostra est. Acquievérunt fratres sermónibus illíus. Et prætereúntibus Madianítis negotiatóribus, extrahéntes eum de cistérna, vendidérunt eum Ismaëlítis, vigínti argénteis: qui duxérunt eum in Ægýptum.\par
\switchcolumn
\EN
\noindent And Ruben hearing this, endeavoured to deliver him out of their hands, and said: Do not take away his life, nor shed his blood: but cast him into this pit, that is in the wilderness, and keep your hands harmless: now he said this, being desirous to deliver him out of their hands and to restore him to his father. And as soon as he came to his brethren, they forthwith stript him of his outside coat, that was of diverse colours: And cast him into an old pit, where there was no water. And sitting down to eat bread, they saw some Ismaelites on their way coming from Calaad, with their camels, carrying spices, and balm, and myrrh to Egypt. And Juda said to his brethren: What will it profit us to kill our brother, and conceal his blood? It is better that he be sold to the Ismaelites, and that our hands be not defiled: for he is our brother and our flesh. His brethren agreed to his words. And when the Madianite merchants passed by, they drew him out of the pit, and sold him to the Ismaelites, for twenty pieces of silver: and they led him into Egypt.\par
\switchcolumn*

\LA
\begin{Resp}\Rsign Extrahéntes Joseph de lacu, vendidérunt Ismaëlítæ vigínti argénteis: \Star{} Reversúsque Ruben ad púteum, cum non invenísset eum, scidit vestiménta sua cum fletu, et dixit: * Puer non compáret, et ego quo ibo?\end{Resp}
\begin{Resp}\Vsign At illi, intíncta túnica Joseph in sánguine hædi, misérunt qui ferret eam ad patrem, et díceret: Vide, si túnica fílii tui sit, an non.\end{Resp}
\begin{Resp}\Rsign Reversúsque Ruben ad púteum, cum non invenísset eum, scidit vestiménta sua cum fletu, et dixit.\end{Resp}
\begin{Resp}\Vsign Glória Patri, et Fílio, \Star{} et Spirítui Sancto.\end{Resp}
\begin{Resp}\Rsign Puer non compáret, et ego quo ibo?\end{Resp}
\switchcolumn
\EN
\begin{Resp}\Rsign They drew up Joseph out of the pit, and sold him to the Ishmaelites for twenty pieces of silver. \Star{} And Reuben returned unto the pit, and when he found not Joseph, he rent his clothes, and wept, and said: * The child is not, and I, whither shall I go?\end{Resp}
\begin{Resp}\Vsign And they took Joseph's coat, and dipped it in the blood of a kid of the goats, and they sent one that brought the coat unto their father, and said: See now whether this be thy son's coat or no.\end{Resp}
\begin{Resp}\Rsign And Reuben returned unto the pit, and when he found not Joseph, he rent his clothes, and wept.\end{Resp}
\begin{Resp}\Vsign Glory be to the Father, and to the Son, \Star{} and to the Holy Ghost.\end{Resp}
\begin{Resp}\Rsign And Reuben said: The child is not, and I, whither shall I go?\end{Resp}
\switchcolumn*

\LA
\LessonHead{Lectio IV}
\switchcolumn
\EN
\LessonHead{Lesson IV}
\switchcolumn*

\LA
\LessonTitle{Ex libro sancti Ambrósii Epíscopi de sancto Joseph}
\switchcolumn
\EN
\LessonTitle{From the Book upon holy Joseph written by St. Ambrose, Bishop (of Milan.)}
\switchcolumn*

\LA
\Source{Cap. 1.}
\switchcolumn
\EN
\Source{Ch. 1}
\switchcolumn*

\LA
\noindent Sanctórum vita céteris norma vivéndi est. Ideóque digéstam plénius accépimus sériem Scripturárum; ut dum Abraham, Isaac, et Jacob, ceterósque justos legéndo cognóscimus, velut quemdam nobis innocéntiæ trámitem, virtúte eórum reserátum imitántibus vestígiis persequámur. De quibus mihi cum frequens tractátus fúerit, hódie sancti Joseph história occúrrit: in quo cum plúrima fúerint génera virtútum, præcípue tamen insígne effúlsit castimóniæ. Justum est ígitur, ut cum in Abraham didicéritis ímpigram fídei devotiónem, in Isaac sincéræ mentis puritátem, in Jacob singulárem ánimi laborúmque patiéntiam: ex illa generalitáte virtútum in ipsas spécies disciplinárum intendátis ánimum.\par
\switchcolumn
\EN
\noindent The lives of the saints are the models for the lives of others. This is one of the reasons why we have been given the wise tale of the Scriptures, that while, by reading therein, we come to know Abraham, and Isaac, and Jacob, and others of the righteous, we may follow them in that path of innocency which is opened to us for our imitation by the record of their godly conversation. Of them I have often treated, and today the story of the holy Joseph cometh before me. In that story there are patterns of many virtues, but chiefly is he glorious on account of his clean living. Right is it then that ye who have learnt in Abraham the devotedness of a faith that nothing could daunt, in Isaac the transparency of an upright soul, in Jacob a wonderful patience of spirit in great travails, should now turn from their worthy deeds, to see the bright example of Joseph's self-control.\par
\switchcolumn*

\LA
\begin{Resp}\Rsign Videns Jacob vestiménta Joseph, scidit vestiménta sua cum fletu, et dixit: \Star{} Fera péssima devorávit fílium meum Joseph.\end{Resp}
\begin{Resp}\Vsign Tulérunt autem fratres ejus túnicam illíus, mitténtes ad patrem: quam cum cognovísset pater, ait.\end{Resp}
\begin{Resp}\Rsign Fera péssima devorávit fílium meum Joseph.\end{Resp}
\switchcolumn
\EN
\begin{Resp}\Rsign When Jacob saw Joseph's coat he rent his clothes, and mourned; and he said: \Star{} An evil beast hath devoured my son Joseph.\end{Resp}
\begin{Resp}\Vsign And his brethren took his coat, and sent it to his father and he knew it, and said:\end{Resp}
\begin{Resp}\Rsign An evil beast hath devoured my son Joseph.\end{Resp}
\switchcolumn*

\LA
\LessonHead{Lectio V}
\switchcolumn
\EN
\LessonHead{Lesson V}
\switchcolumn*

\LA
\noindent Sit ígitur nobis propósitus sanctus Joseph tamquam spéculum castitátis. In ejus enim móribus, in ejus áctibus lucet pudicítia, et quidam splendet castimóniæ comes, nitor grátiæ. Unde étiam a paréntibus plus quam céteri fílii diligebátur. Sed ea res invídiæ fuit: quod siléntio prætereúndum non fuit: hinc enim arguméntum totíus históriæ procéssit: simul ut cognoscámus, perféctum virum non movéri ulciscéndi dolóris invídia, nec malórum repéndere vicem. Unde et David ait: Si réddidi retribuéntibus mihi mala.\par
\switchcolumn
\EN
\noindent The holy Joseph is put before us as a pattern of chastity. Modesty shineth in his manners and in his deeds, and a certain loveliness, which is found with chastity, shineth there also. Hence his parents loved him more than their other children. But this love caused him to be the object of an envy, which we must needs not pass by, and upon this the whole story turneth. Yet, at the same time, we learn how that just man was not swayed by any desire to avenge his own sufferings, neither repaid evil for evil. Whence also David saith: If I have rewarded evil (Ps. vii. 5.)\par
\switchcolumn*

\LA
\begin{Resp}\Rsign Joseph dum intráret in terram Ægýpti, linguam quam non nóverat, audívit: manus ejus in labóribus serviérunt: \Star{} Et lingua ejus inter príncipes loquebátur sapiéntiam.\end{Resp}
\begin{Resp}\Vsign Humiliavérunt in compédibus pedes ejus: ferrum petránsiit ánimam ejus, donec veníret verbum ejus.\end{Resp}
\begin{Resp}\Rsign Et lingua ejus inter príncipes loquebátur sapiéntiam.\end{Resp}
\switchcolumn
\EN
\begin{Resp}\Rsign When Joseph came into the land of Egypt, he heard a language that he understood not; his hands were burdened with labour; \Star{} And his tongue spake wisdom among princes.\end{Resp}
\begin{Resp}\Vsign Whose feet they hurt with fetters; the iron entered into his soul, until the time that his word came\end{Resp}
\begin{Resp}\Rsign And his tongue spake wisdom among princes.\end{Resp}
\switchcolumn*

\LA
\LessonHead{Lectio VI}
\switchcolumn
\EN
\LessonHead{Lesson VI}
\switchcolumn*

\LA
\noindent Quid autem esset, quod præférri Joseph mererétur céteris, si aut lædéntes læsísset, aut diligéntes dilexísset? Hoc enim pleríque fáciunt. Sed illud mirábile, si díligas inimícum tuum: quod Salvátor docet. Jure ergo mirándus, qui hoc fecit ante Evangélium, ut læsus párceret, appetítus ignósceret, vénditus non reférret injúriam, sed grátiam pro contumélia sólveret: quod post Evangélium omnes didícimus, et serváre non póssumus. Discámus ergo et Sanctórum invídiam, ut imitémur patiéntiam: et cognoscámus, illos non natúræ præstantióris fuísse, sed observantióris: nec vítia nescísse, sed emendásse. Quod si invídia étiam Sanctos adússit, quanto magis cavéndum est, ne inflámmet peccatóres?\par
\switchcolumn
\EN
\noindent In what would Joseph have been worthy to be chosen before others, if he had harmed them which harmed him, and loved them which loved him? For this do many do. But it is a wonder if one do that which the Saviour teacheth, and love his enemy. Well, then, may we wonder at him who did this before the Gospel came; who, being injured, spared; being assailed, forgave; being sold, returned no evil; but repaid insult with favour. We, from the Gospel, have been taught to do all this, and we cannot. Let us also, then, learn how that there was envy even among some of the holy (Patriarchs), that we may follow the example of the patience (wherewith others of them bore it;) and let us feel that they were not men of another and higher nature than ours, but only more heedful; that they were not sinless, but that they repented. But if the passion of envy scorched even some of the holy race, how much more need is there for the sinful to take heed lest it set fire to them?\par
\switchcolumn*

\LA
\begin{Resp}\Rsign Meménto mei, dum bene tibi fúerit: \Star{} Ut súggeras pharaóni, ut edúcat me de isto cárcere: * Quia furtim sublátus sum, et hic ínnocens in lacum missus sum.\end{Resp}
\begin{Resp}\Vsign Tres enim adhuc dies sunt, post quos recordábitur phárao ministérii tui, et restítuet te in gradum prístinum: tunc meménto mei.\end{Resp}
\begin{Resp}\Rsign Ut súggeras pharaóni, ut edúcat me de isto cárcere.\end{Resp}
\begin{Resp}\Vsign Glória Patri, et Fílio, \Star{} et Spirítui Sancto.\end{Resp}
\begin{Resp}\Rsign Quia furtim sublátus sum, et hic ínnocens in lacum missus sum.\end{Resp}
\switchcolumn
\EN
\begin{Resp}\Rsign Think on me when it shall be well with thee; \Star{} And make mention of me unto Pharaoh, that he may bring me out of this prison. * For I was stolen away; and here have I done nothing, that they should put me into the dungeon.\end{Resp}
\begin{Resp}\Vsign For yet three days, and then Pharaoh shall remember thy service, and restore thee unto thy place; then think of me.\end{Resp}
\begin{Resp}\Rsign And make mention of me unto Pharaoh, that he may bring me out of this prison.\end{Resp}
\begin{Resp}\Vsign Glory be to the Father, and to the Son, \Star{} and to the Holy Ghost.\end{Resp}
\begin{Resp}\Rsign For I was stolen away; and here have I done nothing, that they should put me into the dungeon.\end{Resp}
\switchcolumn*

\LA
\LessonHead{Lectio VII}
\switchcolumn
\EN
\LessonHead{Lesson VII}
\switchcolumn*

\LA
\LessonTitle{Léctio sancti Evangélii secúndum Lucam}
\switchcolumn
\EN
\LessonTitle{From the Holy Gospel according to Luke}
\switchcolumn*

\LA
\Cite{Luc 11:14-28}
\switchcolumn
\EN
\Cite{Luke 11:14-28}
\switchcolumn*

\LA
\noindent In illo témpore: Erat Jesus eíciens dæmónium, et illud erat mutum. Et cum ejecísset dæmónium, locútus est mutus, et admirátæ sunt turbæ. Et réliqua.\par
\switchcolumn
\EN
\noindent At that time, Jesus was casting out a devil, and the same was dumb: and when he had cast out the devil, the dumb spoke: and the multitudes were in admiration at it. And so on.\par
\switchcolumn*

\LA
\LessonTitle{Homilía sancti Bedæ Venerábilis Presbýteri}
\switchcolumn
\EN
\LessonTitle{Homily by the Venerable Bede, Priest (at Jarrow.)}
\switchcolumn*

\LA
\Source{Lib. 4, cap. 48, in cap. 11 Lucæ}
\switchcolumn
\EN
\Source{Bk. iv, 48 on Luke xi}
\switchcolumn*

\LA
\noindent Dæmoníacus iste apud Matthǽum non solum mutus, sed et cæcus fuísse narrátur: curatúsque dícitur a Dómino, ita ut loquerétur, et vidéret. Tria ergo signa simul in uno hómine perpetráta sunt: cæcus videt, mutus lóquitur, posséssus a dǽmone liberátur. Quod et tunc quidem carnáliter factum est, sed et cotídie complétur in conversióne credéntium: ut, expúlso primum dǽmone, fídei lumen aspíciant; deínde ad laudes Dei tacéntia prius ora laxéntur. Quidam autem ex eis dixérunt: in Beélzebub príncipe dæmoniórum éicit dæmónia. Non hæc áliqui de turba, sed pharisǽi calumniabántur, et scribæ, sicut álii Evangelístæ testántur.\par
\switchcolumn
\EN
\noindent We read in Matthew (xii. 22) that the devil, by which this poor creature was possessed, was not only dumb, but also blind; and that, when he was healed by the Lord, he saw as well as spake. Three miracles, therefore, were performed on this one man; the blind saw, the dumb spake, and the possessed was delivered. This mighty work was then indeed wrought carnally, but it is still wrought spiritually in the conversion of believers, when the devil is cast out of them, so that their eyes see the light of faith, and the lips, that before were dumb, are opened that their mouth may show forth the praise of God. But some of them said: He casteth out devils through Beelzebub, the chief of the devils. These some were not of the multitude, but liars among the Pharisees and Scribes, as we are told by the other Evangelist.\par
\switchcolumn*

\LA
\begin{Resp}\Rsign Mérito hæc pátimur, quia peccávimus in fratrem nostrum, vidéntes angústias ánimæ ejus, dum deprecarétur nos, et non audívimus: \Star{} Idcírco venit super nos tribulátio.\end{Resp}
\begin{Resp}\Vsign Dixit Ruben frátribus suis: Numquid non dixi vobis, Nolíte peccáre in púerum; et non audístis me?\end{Resp}
\begin{Resp}\Rsign Idcírco venit super nos tribulatio.\end{Resp}
\switchcolumn
\EN
\begin{Resp}\Rsign We are verily guilty concerning our brother, in that we saw the anguish of his soul, when he besought us, and we would not hear. \Star{} Therefore is this distress come upon us.\end{Resp}
\begin{Resp}\Vsign And Reuben answered his brethren, saying Spake I not unto you, saying: Do not sin against the child; and ye would not hear?\end{Resp}
\begin{Resp}\Rsign Therefore is this distress come upon us.\end{Resp}
\switchcolumn*

\LA
\LessonHead{Lectio VIII}
\switchcolumn
\EN
\LessonHead{Lesson VIII}
\switchcolumn*

\LA
\noindent Turbis quippe, quæ minus erudítæ videbántur, Dómini semper facta mirántibus; illi contra, vel negáre hæc, vel quæ negáre nequíverant, sinístra interpretatióne pervértere laborábant: quasi non hæc divinitátis, sed immúndi spíritus ópera fuíssent. Et álii tentántes, signum de cælo quærébant ab eo. Vel in morem Elíæ ignem de sublími veníre cupiébant; vel in similitúdinem Samuélis témpore æstívo mugíre tonítrua, coruscáre fúlgura, imbres rúere: quasi non possent et illa calumniári, et dícere, ex occúltis et váriis aëris passiónibus accidísse. At tu, qui calumniáris ea, quæ óculis vides, manu tenes, utilitáte sentis; quid féceris de iis, quæ de cælo vénerint? Utique respondébis, et magos in Ægýpto multa signa fecísse de cælo.\par
\switchcolumn
\EN
\noindent While the multitude, who were less instructed, wondered ever at the works of the Lord, the Pharisees and Scribes, on the other hand, denied the facts when they could, and when they were not able, twisted them by an evil interpretation, and asserted that the works of God were the works of an unclean spirit. And others, tempting Him, sought of Him a sign from heaven. They would have had Christ either to call down fire from heaven like Elias, (4 Kings i. 10,) or, like Samuel, (1 Kings vii. 10,) to have made thunder roll, and lightning flash, and rain fall at midsummer. And yet and if he had so done, they had been still able to explain away these signs also, as being the natural result of some unusual, though, till that moment, unremarked state of the atmosphere. O thou, who stubbornly deniest that which thine eye seeth, thine hand holdeth, and thy sense perceiveth, what wilt thou say to a sign from heaven? In good sooth, thou wilt say that the magicians in Egypt also wrought diverse signs from heaven. (Ex. vii., viii.)\par
\switchcolumn*

\LA
\begin{Resp}\Rsign Dixit Ruben frátribus suis: Numquid non dixi vobis, Nolíte peccáre in púerum, et non audístis me? \Star{} En sanguis ejus exquíritur.\end{Resp}
\begin{Resp}\Vsign Mérito hæc pátimur, quia peccávimus in fratrem nostrum, vidéntes angústias ánimæ ejus, dum deprecarétur nos, et non audívimus.\end{Resp}
\begin{Resp}\Rsign En sanguis ejus exquíritur.\end{Resp}
\switchcolumn
\EN
\begin{Resp}\Rsign And Reuben answered his brethren, saying: Spake I not unto you, saying: Do not sin against the child; and ye would not hear? \Star{} Behold, his blood is required.\end{Resp}
\begin{Resp}\Vsign We are verily guilty concerning our brother, in that we saw the anguish of his soul, when he besought us, and we would not hear.\end{Resp}
\begin{Resp}\Rsign Behold, his blood is required.\end{Resp}
\switchcolumn*

\LA
\LessonHead{Lectio IX}
\switchcolumn
\EN
\LessonHead{Lesson IX}
\switchcolumn*

\LA
\noindent Ipse autem ut vidit cogitatiónes eórum, dixit eis: Omne regnum in seípsum divísum desolábitur, et domus supra domum cadet. Non ad dicta, sed ad cogitáta respóndit: ut vel sic compelleréntur crédere poténtiæ ejus, qui cordis vidébat occúlta. Si autem omne regnum in seípsum divísum desolátur; ergo Patris et Fílii et Spíritus Sancti regnum non est divísum; quod sine ulla contradictióne, non áliquo impúlsu desolándum, sed ætérna est stabilitáte mansúrum. Si autem sátanas in seípsum divísus est: quómodo stabit regnum ipsíus, quia dícitis, in Beélzebub ejícere me dæmónia? Hoc dicens, ex ipsórum confessióne volébat intélligi, quod in eum non credéndo, in regno diáboli esse elegíssent, quod útique advérsum se divísum stare non posset.\par
\switchcolumn
\EN
\noindent But He, knowing their thoughts, said unto them: Every kingdom divided against itself is brought to desolation, and an house divided against an house falleth. He answered not their words, but their thoughts; as though He would compel them to believe in the power of Him Who seeth the secrets of the heart. But if every kingdom divided against itself is brought to desolation, then have not the Father, the Son, and the Holy Ghost a divided kingdom, since His is a kingdom that, without all contradiction, shall never be brought to desolation by any shock, but abideth unchanged and unchangeable for ever. If Satan also be divided against himself, how shall his kingdom stand? Because ye say that I cast out devils by Beelzebub. In saying this, He sought to draw from their own mouth a confession that they had chosen for themselves to be part of that devil's kingdom, which, if it be divided against itself, cannot stand.\par
\switchcolumn*

\LA
\begin{Resp}\Rsign Lamentabátur Jacob de duóbus fíliis suis: Heu me, dolens sum de Joseph pérdito, et tristis nimis de Bénjamin ducto pro alimóniis: \Star{} Precor cæléstem Regem, ut me doléntem nímium fáciat eos cérnere.\end{Resp}
\begin{Resp}\Vsign Prostérnens se Jacob veheménter cum lácrimis pronus in terram, et adórans ait.\end{Resp}
\begin{Resp}\Rsign Precor cæléstem Regem, ut me doléntem nímium fáciat eos cérnere.\end{Resp}
\begin{Resp}\Vsign Glória Patri, et Fílio, \Star{} et Spirítui Sancto.\end{Resp}
\begin{Resp}\Rsign Precor cæléstem Regem, ut me doléntem nímium fáciat eos cérnere.\end{Resp}
\switchcolumn
\EN
\begin{Resp}\Rsign Jacob lamented for his two sons, saying: Woe is me; I am bereaved of Joseph, for he is not; and afflicted because of Benjamin, because he is taken away for bread. \Star{} I pray the King of heaven in my distress, that He may make me to see them yet again.\end{Resp}
\begin{Resp}\Vsign And Jacob cast him down upon his face upon the ground, and wept sore; and he prayed, saying:\end{Resp}
\begin{Resp}\Rsign I pray the King of heaven in my distress, that He may make me to see them yet again.\end{Resp}
\begin{Resp}\Vsign Glory be to the Father, and to the Son, \Star{} and to the Holy Ghost.\end{Resp}
\begin{Resp}\Rsign I pray the King of heaven in my distress, that He may make me to see them yet again.\end{Resp}
\switchcolumn*

\end{paracol}

\Office{Feria Secunda infra Hebdomadam III in Quadragesima}{Monday of the Third Week of Lent}{Feria major}
\addcontentsline{toc}{subsection}{Feria Secunda infra Hebdomadam III in Quadragesima\enspace\textit{Monday of the Third Week of Lent}}
\begin{paracol}{2}
\LA
\LessonHead{Lectio I}
\switchcolumn
\EN
\LessonHead{Lesson I}
\switchcolumn*

\LA
\LessonTitle{Léctio sancti Evangélii secúndum Lucam}
\switchcolumn
\EN
\LessonTitle{Continuation of the Holy Gospel according to Luke}
\switchcolumn*

\LA
\Cite{Luc 4:23-30}
\switchcolumn
\EN
\Cite{Luke 4:23-30}
\switchcolumn*

\LA
\noindent In illo témpore: Dixit Jesus pharisǽis: Utique dicétis mihi hanc similitúdinem: Médice, cura te ipsum: quanta audívimus facta in Caphárnaum, fac et hic in pátria tua. Et réliqua.\par
\switchcolumn
\EN
\noindent At that time, Jesus said to the pharisees: Doubtless you will say to me this similitude: Physician, heal thyself: as great things as we have heard done in Capharnaum, do also here in thy own country. And so on.\par
\switchcolumn*

\LA
\LessonTitle{Homilía sancti Ambrósii Epíscopi}
\switchcolumn
\EN
\LessonTitle{Homily by St. Ambrose, Bishop (of Milan.)}
\switchcolumn*

\LA
\Source{Lib. 4 in c. 4 Lucæ, post med.}
\switchcolumn
\EN
\Source{Bk. iv on Luke iv}
\switchcolumn*

\LA
\noindent Non medíocris invídia próditur, quæ cívicæ caritátis oblíta, in acérba ódia causas amóris infléctit. Simul hoc exémplo páriter et oráculo declarátur, quod frustra opem misericórdiæ cæléstis exspéctes, si aliénæ frúctibus virtútis invídeas. Aspernátor enim Dóminus invidórum est: et ab iis qui divína benefícia in áliis persequúntur, mirácula suæ potestátis avértit. Domínicæ quippe carnis actus, divinitátis exémplum est: et invisibília nobis ejus, per ea quæ sunt visibília, demonstrántur.\par
\switchcolumn
\EN
\noindent Here we have a display of a spite not very common. Their hatred of Christ, and their desire to find grounds for that hatred in what in Him appealed for their love, had made them forget their local friendliness to a fellow-citizen. By this example as well as by God's declaration, thou mayest learn that thou wilt wait in vain to be helped of His mercy, whilst thou art envious of the spiritual good of thy neighbour. Yea, the Lord turneth Him away from the envious, and will not show the mighty works of His power to such as are bitter against His gifts to others. The example of Himself which God hath been pleased to set before us is that of His doings in the Flesh, and it is by these His doings which He suffered to be seen, that we are taught touching those which are unseen.\par
\switchcolumn*

\LA
\begin{Resp}\Rsign Tóllite hinc vobíscum múnera, et ite ad dóminum terræ: et cum invenéritis, adoráte eum super terram: \Star{} Deus autem meus fáciat eum vobis placábilem: et remíttat et hunc fratrem vestrum vobíscum, et eum quem tenet in vínculis.\end{Resp}
\begin{Resp}\Vsign Súmite de óptimis terræ frúgibus in vasis vestris, et deférte viro múnera.\end{Resp}
\begin{Resp}\Rsign Deus autem meus fáciat eum vobis placábilem: et remíttat et hunc fratrem vestrum vobíscum, et eum quem tenet in vínculis.\end{Resp}
\switchcolumn
\EN
\begin{Resp}\Rsign Take hence presents with you, and go unto the lord of the land, and when ye be come into his presence, bow yourselves to him to the earth. \Star{} And my God give you mercy before the man, that he may send away again this your brother, and him which he keepeth in ward.\end{Resp}
\begin{Resp}\Vsign Take of the best fruits of the land in your vessels, and carry down the man a present.\end{Resp}
\begin{Resp}\Rsign And my God give you mercy before the man, that he may send away again this your brother, and him which he keepeth in ward.\end{Resp}
\switchcolumn*

\LA
\LessonHead{Lectio II}
\switchcolumn
\EN
\LessonHead{Lesson II}
\switchcolumn*

\LA
\noindent Non otióse ítaque Salvátor excúsat, quod nulla in pátria sua mirácula virtútis operátus sit: ne fortássis áliquis viliórem pátriæ nobis esse debére putáret afféctum. Neque enim cives póterat non amáre, qui amáret omnes: sed ipsi se caritáte pátriæ, dum ínvident, abdicárunt. In veritáte dico vobis: multæ víduæ fuérunt in diébus Elíæ. Non quia Elíæ dies fuérunt, sed in quibus Elías operátus est: aut quia Elías dies faciébat illis, qui in ejus opéribus lucem vidébant grátiæ spiritális, et convertebántur ad Dóminum. Et ídeo aperiebátur cælum vidéntibus ætérna et divína mystéria: claudebátur, et fames erat, quando nulla erat cognoscéndæ divinitátis ubértas. Sed de hoc plénius díximus, cum de víduis scriberémus.\par
\switchcolumn
\EN
\noindent The Saviour then doth not lightly excuse Himself that He had wrought none of His mighty works in His own country, lest perchance any should thence learn to think lightly of our duty to love our Fatherland. Neither was it possible that He Who loved all, should not love His own countrymen; they it was who failed in that love because of their very envy. I tell you of a truth, many widows were in Israel in the days of Elias. The days of Elias not that the said days belonged to Elias, but either because those were the days when Elias lived and worked; or, else, this is a mystic phrase, meaning that Elias by his works made many souls to awake spiritually from the night of sin to the day of grace, and turn to the Lord. In this latter sense that holy Prophet was a mean whereby heaven was opened to such as looked to the eternal and mysterious things of God, and again was shut, (and there was a famine,) when there were no means of knowing God through outward ordinances. This subject, however, I have treated before at full length, when I was writing on the subject of widows.\par
\switchcolumn*

\LA
\begin{Resp}\Rsign Iste est frater vester mínimus, de quo dixerátis mihi? Deus misereátur tibi, fili mi. \Star{} Festinavítque in domum, et plorávit: quia erumpébant lácrimæ, et non póterat se continére.\end{Resp}
\begin{Resp}\Vsign Attóllens autem Joseph óculos, vidit Bénjamin stantem: et commóta sunt ómnia víscera ejus super fratre suo.\end{Resp}
\begin{Resp}\Rsign Festinavítque in domum, et plorávit: quia erumpébant lácrimæ, et non póterat se continére.\end{Resp}
\switchcolumn
\EN
\begin{Resp}\Rsign Is this your younger brother, of whom ye spake unto me? God be gracious unto thee, my son. \Star{} And he made haste, and entered into the house, and wept there, for his tears brake forth, and he could not refrain himself.\end{Resp}
\begin{Resp}\Vsign And Joseph lifted up his eyes, and saw his brother Benjamin, and his bowels yearned upon his brother.\end{Resp}
\begin{Resp}\Rsign And he made haste, and entered into the house, and wept there, for his tears brake forth, and he could not refrain himself.\end{Resp}
\switchcolumn*

\LA
\LessonHead{Lectio III}
\switchcolumn
\EN
\LessonHead{Lesson III}
\switchcolumn*

\LA
\noindent Et multi leprósi erant in Judǽa tempóribus Eliséi prophétæ: et nemo eórum mundátus est, nisi Náaman Syrus. Evidénter hic sermo nos Dómini salutáris infórmat, et ad stúdium venerándæ divinitátis hortátur: quod nemo sanátus osténditur, et maculósi morbo córporis absolútus, nisi qui religióso offício stúduit sanitáti. Non enim dormiéntibus divína benefícia, sed observántibus deferúntur. Díximus in libro álio, in vídua illa, ad quam Elías diréctus est, typum Ecclésiæ præmíssum. Pópulus Ecclésiam congregávit, ut sequátur pópulus ille ex alienígenis congregátus. Pópulus ille ante leprósus, pópulus ille ante maculósus, priúsquam mýstico baptizarétur in flúmine: idem post sacraménta baptísmatis máculis córporis et mentis ablútus, jam non lepra, sed immaculáta virgo cœpit esse sine ruga.\par
\switchcolumn
\EN
\noindent And many lepers were in Israel in the days of Eliseus the Prophet, and none of them was cleansed, saving Naaman the Syrian. By these words of the Lord our great Physician, we are plainly taught and urged to put our trust in the Adorable God, since we see that none was healed, or cleansed from bodily plague-spots, save him who took a religious means to regain health. For the blessings of God are not given to them who close their eyes in sleep, but to them that look to Him. We have remarked in our other book, (alluded to above,) that the widow to whom Elias was sent was a type of the Church. And next after (the mention of the type of) the Church cometh meetly the (mention of him who was a type of the Gentile) people, (her converts.) Yea, the Gentiles were a people foreigners by birth, leprous, and covered with plague-spots, till they were baptized in the stream of the mystic Jordan; but from the sacramental waters they rise, lepers no more, but cleansed in body and soul, a glorious virgin Church, not having spot, or wrinkle, or any such thing. (Eph. v. 27.)\par
\switchcolumn*

\LA
\begin{Resp}\Rsign Dixit Joseph úndecim frátribus suis: Ego sum Joseph, quem vendidístis in Ægýptum: adhuc vivit pater noster sénior, de quo dixerátis mihi? \Star{} Ite, addúcite eum ad me, ut possit vívere.\end{Resp}
\begin{Resp}\Vsign Biénnium enim est, quod cœpit esse fames in terra: et adhuc restant anni quinque, quibus nec arári póterit, nec meti.\end{Resp}
\begin{Resp}\Rsign Ite, addúcite eum ad me, ut possit vívere.\end{Resp}
\begin{Resp}\Vsign Glória Patri, et Fílio, \Star{} et Spirítui Sancto.\end{Resp}
\begin{Resp}\Rsign Ite, addúcite eum ad me, ut possit vívere.\end{Resp}
\switchcolumn
\EN
\begin{Resp}\Rsign Joseph said unto his eleven brethren: I am Joseph, whom ye sold into Egypt; is our father yet alive, the old man of whom ye spake unto me? \Star{} Go, bring him down unto me, that he may live.\end{Resp}
\begin{Resp}\Vsign For these two years hath the famine been in the land; and yet there are five years, in the which there shall neither be earing nor harvest.\end{Resp}
\begin{Resp}\Rsign Go, bring him down unto me, that he may live.\end{Resp}
\begin{Resp}\Vsign Glory be to the Father, and to the Son, \Star{} and to the Holy Ghost.\end{Resp}
\begin{Resp}\Rsign Go, bring him down unto me, that he may live.\end{Resp}
\switchcolumn*

\end{paracol}

\Office{Feria Tertia infra Hebdomadam III in Quadragesima}{Tuesday of the Third Week of Lent}{Feria major}
\addcontentsline{toc}{subsection}{Feria Tertia infra Hebdomadam III in Quadragesima\enspace\textit{Tuesday of the Third Week of Lent}}
\begin{paracol}{2}
\LA
\LessonHead{Lectio I}
\switchcolumn
\EN
\LessonHead{Lesson I}
\switchcolumn*

\LA
\LessonTitle{Léctio sancti Evangélii secúndum Matthǽum}
\switchcolumn
\EN
\LessonTitle{Continuation of the Holy Gospel according to Matthew}
\switchcolumn*

\LA
\Cite{Matt 18:15-22}
\switchcolumn
\EN
\Cite{Matt 18:15-22}
\switchcolumn*

\LA
\noindent In illo témpore: Dixit Jesus discípulis suis: Si peccáverit in te frater tuus, vade, et córripe eum inter te et ipsum solum. Et réliqua.\par
\switchcolumn
\EN
\noindent At that time, Jesus saith to his disciples: But if thy brother shall offend against thee, go, and rebuke him between thee and him alone. And so on.\par
\switchcolumn*

\LA
\LessonTitle{Homilía sancti Augustíni Epíscopi}
\switchcolumn
\EN
\LessonTitle{Homily by St. Augustine, Bishop (of Hippo.)}
\switchcolumn*

\LA
\Source{Sermo 16 de Verbis Domini, tomus 10, post initium}
\switchcolumn
\EN
\Source{16th Sermon on the Words of the Lord, vol. x}
\switchcolumn*

\LA
\noindent Quare illum córripis? Quia tu doles, quod peccáverit in te? Absit. Si amóre tui id facis, nihil facis: si amóre illíus facis, óptime facis. Dénique in ipsis verbis atténde, cujus amóre id fácere débeas, utrum tui, an illíus. Si te audíerit, inquit, lucrátus es fratrem tuum. Ergo propter illum fac, ut lucréris illum. Sic faciéndo lucráris: nisi fecísses, períerat. Quid est ergo, quod pleríque hómines ista peccáta contémnunt, et dicunt: Quid magnum feci? In hóminem peccávi. Noli contémnere: in hóminem peccásti.\par
\switchcolumn
\EN
\noindent Why tell him his fault? Because he hath made thee smart by trespassing against thee? God forbid. If thou tell him his fault because thou lovest thyself, thou dost nothing. But if thou tell it him because thou lovest him, then dost thou do exceeding well. Hear now, in the words of the Gospel itself, for love of whom, thou oughtest to do it, of thyself, or of him. The Lord saith: If he shall hear thee, thou hast gained thy brother. Therefore it behoveth thee to do it for his sake, that thou mayest gain him; since, if thou so do, haply thou mayest gain him; whereas, if thou do it not, he may haply perish. Why then are there so many who reckon lightly of a trespass against their brother, and say I have done no great offence, for I have trespassed only against my fellow man? Deem it not light; thou hast trespassed, though it be against thy fellow man.\par
\switchcolumn*

\LA
\begin{Resp}\Rsign Nuntiavérunt Jacob dicéntes: Joseph fílius tuus vivit, et ipse dominátur in tota terra Ægýpti: quo audíto revíxit spíritus ejus, et dixit: \Star{} Súfficit mihi, vadam et vidébo eum ántequam móriar.\end{Resp}
\begin{Resp}\Vsign Cumque audísset Jacob quod fílius ejus víveret, quasi de gravi somno evígilans, ait.\end{Resp}
\begin{Resp}\Rsign Súfficit mihi, vadam et vidébo eum ántequam móriar.\end{Resp}
\switchcolumn
\EN
\begin{Resp}\Rsign They told Jacob, saying: thy son Joseph is yet alive, and he is governor over all the land of Egypt; and, when he heard it, his spirit revived, and he said \Star{} It is enough; I will go, and see him before I die.\end{Resp}
\begin{Resp}\Vsign And when Jacob heard that his son yet lived, he was as one that awakeneth from a deep sleep, and said:\end{Resp}
\begin{Resp}\Rsign It is enough; I will go, and see him before I die.\end{Resp}
\switchcolumn*

\LA
\LessonHead{Lectio II}
\switchcolumn
\EN
\LessonHead{Lesson II}
\switchcolumn*

\LA
\noindent Vis nosse, quia in hóminem peccándo perísti? Si te ille, in quem peccásti, corripúerit inter te et ipsum solum, et audíeris illum, lucrátus est te. Quid est, Lucrátus est te; nisi quia períeras, si non lucrarétur te? Nam si non períeras, quómodo te lucrátus est? Nemo ergo contémnat, quando peccat in fratrem. Ait enim quodam loco Apóstolus: Sic autem peccántes in fratres, et percutiéntes consciéntiam eórum infírmam, in Christum peccátis: ídeo quia membra Christi omnes facti sumus. Quómodo non peccas in Christum, qui peccas in membrum Christi?\par
\switchcolumn
\EN
\noindent Wouldest thou know that thy trespass against thy brother hath destroyed thee? If he against whom thou hast trespassed tell thee thy fault between himself and thee alone, and thou hear him, he hath gained thee. Gained thee! And what signify those words, if it be not that thou, if thou be not gained, shalt perish? For if thou shouldest not otherwise perish, in what sense can he be said to gain thee? Therefore let no man deem it a light thing when he trespasseth against his brother. For the Apostle Paul saith in a certain place: When ye sin so against the brethren, and wound their weak conscience, ye sin against Christ. (1 Cor. viii. 12.) We are all members of Christ. How dost thou not trespass against Christ, which trespassest against one of His members?\par
\switchcolumn*

\LA
\begin{Resp}\Rsign Joseph dum intráret in terram Ægýpti, linguam quam non nóverat, audívit: manus ejus in labóribus serviérunt: \Star{} Et lingua ejus inter príncipes loquebátur sapiéntiam.\end{Resp}
\begin{Resp}\Vsign Humiliavérunt in compédibus pedes ejus: ferrum petránsiit ánimam ejus, donec veníret verbum ejus.\end{Resp}
\begin{Resp}\Rsign Et lingua ejus inter príncipes loquebátur sapiéntiam.\end{Resp}
\switchcolumn
\EN
\begin{Resp}\Rsign When Joseph came into the land of Egypt, he heard a language that he understood not; his hands were burdened with labour; \Star{} And his tongue spake wisdom among princes.\end{Resp}
\begin{Resp}\Vsign Whose feet they hurt with fetters; the iron entered into his soul, until the time that his word came\end{Resp}
\begin{Resp}\Rsign And his tongue spake wisdom among princes.\end{Resp}
\switchcolumn*

\LA
\LessonHead{Lectio III}
\switchcolumn
\EN
\LessonHead{Lesson III}
\switchcolumn*

\LA
\noindent Nemo ergo dicat, quia non peccávi in Deum, sed peccávi in fratrem: in hóminem peccávi, leve, vel nullum peccátum est. Forte inde dicis: Leve est, quia cito curátur. Peccásti in fratrem: fac satis, et sanátus es. Cito fecísti mortíferam rem, sed remédium cito invenísti. Quis nostrum speret regnum cælórum, fratres mei, quando dicit Evangélium: Qui díxerit fratri suo, Fátue: reus erit gehénnæ ignis? Magnus terror: sed vide ibi remédium. Si obtúleris munus tuum ad altáre, et ibi recordátus fúeris, quia frater tuus habet áliquid advérsum te, relínque ibi munus tuum ante altáre. Non iráscitur Deus, quia differs impónere munus tuum: te quærit Deus magis, quam munus tuum.\par
\switchcolumn
\EN
\noindent Let no man therefore say I have not trespassed against God, but only against my brother; that is, I have trespassed against my fellowman; and so the sin is light, if any at all. And perchance thou wilt argue that it is light, because it is quickly mended; thou hast trespassed against thy brother, but thou canst make satisfaction, and be right again; thou hast done the deadly thing quickly, and quickly canst thou find a remedy. O my brethren, which of us can hope for the kingdom of heaven, when we remember that the Gospel saith: “Whosoever shall say to his brother: Thou fool, shall be in danger of hell fire?” (Matth. v. 22.) It is a thought full of dread; but, lo! the remedy: If thou bring thy gift to the altar, and there rememberest that thy brother hath aught against thee, leave there thy gift before the altar, and go thy way; first be reconciled to thy brother, and then come and offer thy gift. God is not wroth that thou tarry or ever thou offer thy gift; for God seeketh thyself more than thy gift.\par
\switchcolumn*

\LA
\begin{Resp}\Rsign Meménto mei, dum bene tibi fúerit: \Star{} Ut súggeras pharaóni, ut edúcat me de isto cárcere: * Quia furtim sublátus sum, et hic ínnocens in lacum missus sum.\end{Resp}
\begin{Resp}\Vsign Tres enim adhuc dies sunt, post quos recordábitur phárao ministérii tui, et restítuet te in gradum prístinum: tunc meménto mei.\end{Resp}
\begin{Resp}\Rsign Ut súggeras pharaóni, ut edúcat me de isto cárcere.\end{Resp}
\begin{Resp}\Vsign Glória Patri, et Fílio, \Star{} et Spirítui Sancto.\end{Resp}
\begin{Resp}\Rsign Quia furtim sublátus sum, et hic ínnocens in lacum missus sum.\end{Resp}
\switchcolumn
\EN
\begin{Resp}\Rsign Think on me when it shall be well with thee; \Star{} And make mention of me unto Pharaoh, that he may bring me out of this prison. * For I was stolen away; and here have I done nothing, that they should put me into the dungeon.\end{Resp}
\begin{Resp}\Vsign For yet three days, and then Pharaoh shall remember thy service, and restore thee unto thy place; then think of me.\end{Resp}
\begin{Resp}\Rsign And make mention of me unto Pharaoh, that he may bring me out of this prison.\end{Resp}
\begin{Resp}\Vsign Glory be to the Father, and to the Son, \Star{} and to the Holy Ghost.\end{Resp}
\begin{Resp}\Rsign For I was stolen away; and here have I done nothing, that they should put me into the dungeon.\end{Resp}
\switchcolumn*

\end{paracol}

\Office{Feria Quarta infra Hebdomadam III in Quadragesima}{Wednesday of the Third Week of Lent}{Feria major}
\addcontentsline{toc}{subsection}{Feria Quarta infra Hebdomadam III in Quadragesima\enspace\textit{Wednesday of the Third Week of Lent}}
\begin{paracol}{2}
\LA
\LessonHead{Lectio I}
\switchcolumn
\EN
\LessonHead{Lesson I}
\switchcolumn*

\LA
\LessonTitle{Léctio sancti Evangélii secúndum Matthǽum}
\switchcolumn
\EN
\LessonTitle{Continuation of the Holy Gospel according to Matthew}
\switchcolumn*

\LA
\Cite{Matt 15:1-20}
\switchcolumn
\EN
\Cite{Matt 15:1-20}
\switchcolumn*

\LA
\noindent In illo témpore: Accessérunt ad Jesum ab Jerosólymis scribæ et pharisǽi, dicéntes: Quare discípuli tui transgrediúntur traditiónem seniórum? Et réliqua.\par
\switchcolumn
\EN
\noindent At that time, scribes and Pharisees came to Jesus from Jerusalem, saying: Why do thy disciples trangress the tradition of the ancients? And so on.\par
\switchcolumn*

\LA
\LessonTitle{Homilía sancti Hierónymi Presbýteri}
\switchcolumn
\EN
\LessonTitle{Homily by St. Jerome, Priest (at Bethlehem.)}
\switchcolumn*

\LA
\Source{Lib. 2 Comment. in cap. 15 Matthæi}
\switchcolumn
\EN
\Source{Bk. ii Comm. on Matth. xv}
\switchcolumn*

\LA
\noindent Mira pharisæórum scribarúmque stultítia. Dei Fílium árguunt, quare hóminum traditiónes et præcépta non servet: Non enim lavant manus suas, cum panem mandúcant. Manus, id est ópera, non córporis útique, sed ánimæ lavándæ sunt, ut fiat in illis verbum Dei. Ipse autem respóndens ait illis: Quare et vos transgredímini mandátum Dei propter traditiónem vestram? Falsam calúmniam vera responsióne confútat. Cum, inquit, vos propter traditiónem hóminum præcépta Dómini negligátis: quare discípulos meos arguéndos putátis, quod seniórum jussa parvipéndant, ut Dei scita custódiant?\par
\switchcolumn
\EN
\noindent The stupidity of the Pharisees and Scribes is something extraordinary. They rebuke the Son of God because He doth not observe the traditions and commandments of men for they wash not their hands when they eat bread. It behoveth us to cleanse, not the hands of the body, but the hands of the soul, namely, our works, that we may do the commandments of God. But He answered and said unto them Why do ye also transgress the commandment of God by your tradition? He meeteth here their false accusation by a true. “How,” saith He, “do ye, who pass over the commandments of God, in order to keep to the traditions of men, hold that My disciples are to be rebuked, because they deem the tradition of the elders of little moment in comparison with the doing of what they know to be the Laws of God?”\par
\switchcolumn*

\LA
\begin{Resp}\Rsign Mérito hæc pátimur, quia peccávimus in fratrem nostrum, vidéntes angústias ánimæ ejus, dum deprecarétur nos, et non audívimus: \Star{} Idcírco venit super nos tribulátio.\end{Resp}
\begin{Resp}\Vsign Dixit Ruben frátribus suis: Numquid non dixi vobis, Nolíte peccáre in púerum; et non audístis me?\end{Resp}
\begin{Resp}\Rsign Idcírco venit super nos tribulatio.\end{Resp}
\switchcolumn
\EN
\begin{Resp}\Rsign We are verily guilty concerning our brother, in that we saw the anguish of his soul, when he besought us, and we would not hear. \Star{} Therefore is this distress come upon us.\end{Resp}
\begin{Resp}\Vsign And Reuben answered his brethren, saying Spake I not unto you, saying: Do not sin against the child; and ye would not hear?\end{Resp}
\begin{Resp}\Rsign Therefore is this distress come upon us.\end{Resp}
\switchcolumn*

\LA
\LessonHead{Lectio II}
\switchcolumn
\EN
\LessonHead{Lesson II}
\switchcolumn*

\LA
\noindent Nam Deus dixit: Honóra patrem et matrem; et, Qui maledíxerit patri, vel matri, morte moriátur. Vos autem dícitis: Quicúmque díxerit patri, vel matri: Munus quodcúmque est ex me, tibi próderit: et non honorificábit patrem suum, aut matrem suam. Honor in Scriptúris non tantum in salutatiónibus et offíciis deferéndis, quantum in eleemósynis, ac múnerum oblatióne sentítur. Honóra, inquit Apóstolus, víduas, quæ vere víduæ sunt. Hic honor donum intellégitur. Et in álio loco: Presbýteri dúplici honóre honorándi sunt, máxime qui labórant in verbo et doctrína Dei. Et per hoc mandátum jubémur bovi trituránti os non cláudere: et dignus sit operárius mercéde sua.\par
\switchcolumn
\EN
\noindent Now God commanded, saying: Honour thy father and mother; and: He that curseth father or mother, let him die the death. But ye say: Whosoever shall say to his father or his mother: It is a gift, by whatsoever thou mightest be profited by me; and honour not his father or his mother, he shall be free. The word ‘honour’ is used in Scripture, not so much in the sense of paying salutations and services, as in that of giving alms and gifts. Honour widows, saith the Apostle, which are widows indeed. (1 Tim. v. 3.) And here ‘honour’ signifieth support. So again, (17, 18): Let the Priests that rule well be counted worthy of double honour, especially they who labour in the word and doctrine. For the Scripture saith: “Thou shalt not muzzle the ox that treadeth out the corn” and “The labourer is worthy of his reward.”\par
\switchcolumn*

\LA
\begin{Resp}\Rsign Dixit Ruben frátribus suis: Numquid non dixi vobis, Nolíte peccáre in púerum, et non audístis me? \Star{} En sanguis ejus exquíritur.\end{Resp}
\begin{Resp}\Vsign Mérito hæc pátimur, quia peccávimus in fratrem nostrum, vidéntes angústias ánimæ ejus, dum deprecarétur nos, et non audívimus.\end{Resp}
\begin{Resp}\Rsign En sanguis ejus exquíritur.\end{Resp}
\switchcolumn
\EN
\begin{Resp}\Rsign And Reuben answered his brethren, saying: Spake I not unto you, saying: Do not sin against the child; and ye would not hear? \Star{} Behold, his blood is required.\end{Resp}
\begin{Resp}\Vsign We are verily guilty concerning our brother, in that we saw the anguish of his soul, when he besought us, and we would not hear.\end{Resp}
\begin{Resp}\Rsign Behold, his blood is required.\end{Resp}
\switchcolumn*

\LA
\LessonHead{Lectio III}
\switchcolumn
\EN
\LessonHead{Lesson III}
\switchcolumn*

\LA
\noindent Præcéperat Dóminus, vel imbecillitátes, vel ætátes, vel penúrias paréntum consíderans, ut fílii honorárent, étiam in vitæ necessáriis ministrándis, paréntes suos. Hanc providentíssimam Dei legem voléntes scribæ et pharisǽi subvértere, ut impietátem sub nómine pietátis indúcerent, docuérunt péssimos fílios, ut si quis ea, quæ paréntibus offerénda sunt, Deo vovére volúerit, qui verus est pater, oblátio Dómini præponátur paréntum munéribus: vel certe ipsi paréntes, quæ Deo consecráta cernébant, ne sacrilégii crimen incúrrerent, declinántes, egestáte conficiebántur. Atque ita fiébat, ut oblátio liberórum sub occasióne templi et Dei, in sacerdótum lucra céderet.\par
\switchcolumn
\EN
\noindent The Lord being mindful of the helplessness, or age, or poverty of parents, had commanded their children to honour them even by giving them the necessaries of life. The Scribes and Pharisees, scrupling not to make of none effect this most benign law, and bringing in ungodliness under the very form of godliness, taught, for the benefit of unnatural children, that if any one vowed to God, Who is our very Father in heaven, whatsoever he was bound to give to his parents, the duty of discharging his debt to his heavenly Father ought to come before that which he owed to his earthly father; or, at least, that parents in such case incurred the guilt of sacrilege by taking for themselves what they knew had been made a gift to God. And so parents were left unsuccoured, and the offerings of such children, under pretence of being given to God and His temple, became the gain of the Priests.\par
\switchcolumn*

\LA
\begin{Resp}\Rsign Lamentabátur Jacob de duóbus fíliis suis: Heu me, dolens sum de Joseph pérdito, et tristis nimis de Bénjamin ducto pro alimóniis: \Star{} Precor cæléstem Regem, ut me doléntem nímium fáciat eos cérnere.\end{Resp}
\begin{Resp}\Vsign Prostérnens se Jacob veheménter cum lácrimis pronus in terram, et adórans ait.\end{Resp}
\begin{Resp}\Rsign Precor cæléstem Regem, ut me doléntem nímium fáciat eos cérnere.\end{Resp}
\begin{Resp}\Vsign Glória Patri, et Fílio, \Star{} et Spirítui Sancto.\end{Resp}
\begin{Resp}\Rsign Precor cæléstem Regem, ut me doléntem nímium fáciat eos cérnere.\end{Resp}
\switchcolumn
\EN
\begin{Resp}\Rsign Jacob lamented for his two sons, saying: Woe is me; I am bereaved of Joseph, for he is not; and afflicted because of Benjamin, because he is taken away for bread. \Star{} I pray the King of heaven in my distress, that He may make me to see them yet again.\end{Resp}
\begin{Resp}\Vsign And Jacob cast him down upon his face upon the ground, and wept sore; and he prayed, saying:\end{Resp}
\begin{Resp}\Rsign I pray the King of heaven in my distress, that He may make me to see them yet again.\end{Resp}
\begin{Resp}\Vsign Glory be to the Father, and to the Son, \Star{} and to the Holy Ghost.\end{Resp}
\begin{Resp}\Rsign I pray the King of heaven in my distress, that He may make me to see them yet again.\end{Resp}
\switchcolumn*

\end{paracol}

\Office{Feria Quinta infra Hebdomadam III in Quadragesima}{Thursday of the Third Week of Lent}{Feria major}
\addcontentsline{toc}{subsection}{Feria Quinta infra Hebdomadam III in Quadragesima\enspace\textit{Thursday of the Third Week of Lent}}
\begin{paracol}{2}
\LA
\LessonHead{Lectio I}
\switchcolumn
\EN
\LessonHead{Lesson I}
\switchcolumn*

\LA
\LessonTitle{Léctio sancti Evangélii secúndum Lucam}
\switchcolumn
\EN
\LessonTitle{Continuation of the Holy Gospel according to Luke}
\switchcolumn*

\LA
\Cite{Luc 4:38-44}
\switchcolumn
\EN
\Cite{Luke 4:38}
\switchcolumn*

\LA
\noindent In illo témpore: Surgens Jesus de synagóga, introívit in domum Simónis. Socrus autem Simónis tenebátur magnis fébribus. Et réliqua.\par
\switchcolumn
\EN
\noindent At that time Jesus arose out of the synagogue, and entered into Simon's house. And Simon's wife's mother was taken with a great fever. And so on.\par
\switchcolumn*

\LA
\LessonTitle{Homilía sancti Ambrósii Epíscopi}
\switchcolumn
\EN
\LessonTitle{Homily by St. Ambrose, Bishop of Milan}
\switchcolumn*

\LA
\Source{Liber 4 in cap. 4 Lucæ, circa finem}
\switchcolumn
\EN
\Source{Bk. 4 on Luke iv}
\switchcolumn*

\LA
\noindent Vide cleméntiam Dómini Salvatóris: nec indignatióne commótus nec scélere offénsus, nec injúria violátus Judǽam déserit: quin étiam ímmemor injúriæ, memor cleméntiæ, nunc docéndo, nunc liberándo, nunc sanándo, infídæ plebis corda demúlcet. Et bene sanctus Lucas virum ab spíritu nequítiæ liberátum ante præmísit, et subdit féminæ sanitátem. Utrúmque enim sexum Dóminus curatúrus advénerat: sed prior sanári débuit, qui prior creátus est; nec prætermítti illa, quæ mobilitáte magis ánimi, quam pravitáte peccáverat.\par
\switchcolumn
\EN
\noindent Behold here how long-suffering is the Lord our Redeemer! Neither moved to anger against them, nor sickened at their guilt, nor outraged by their attacks, did He leave the Jews' country. Nay, forgetting their iniquity, and mindful only of His mercy, He strove to soften their hard and unbelieving hearts, sometimes by His teaching, and sometimes by freeing some of them, and sometimes by healing them. St. Luke doth well to tell us first of the man who was delivered from an unclean spirit, and then of the healing of a woman. The Lord indeed came to heal both sexes, but that must be healed first which was created first, and then must not she be passed by whose first sin arose rather from fickleness of heart than from depraved will.\par
\switchcolumn*

\LA
\begin{Resp}\Rsign Vidéntes Joseph a longe, loquebántur mútuo fratres, dicéntes: Ecce somniátor venit: \Star{} Veníte, occidámus eum, et videámus si prosint illi sómnia sua.\end{Resp}
\begin{Resp}\Vsign Cumque vidíssent Joseph fratres sui, quod a patre cunctis frátribus plus amarétur, óderant eum, nec póterant ei quidquam pacífice loqui, unde et dicébant.\end{Resp}
\begin{Resp}\Rsign Veníte, occidámus eum, et videámus si prosint illi sómnia sua.\end{Resp}
\switchcolumn
\EN
\begin{Resp}\Rsign And when his brethren saw Joseph afar off, they said one to another: Behold, this dreamer cometh. \Star{} Come, let us slay him; and we shall see what will become of his dreams.\end{Resp}
\begin{Resp}\Vsign And when his brethren saw that their father loved Joseph more than all his brethren, they hated him, and could not speak peaceably unto him; therefore they said:\end{Resp}
\begin{Resp}\Rsign Come, let us slay him; and we shall see what will become of his dreams.\end{Resp}
\switchcolumn*

\LA
\LessonHead{Lectio II}
\switchcolumn
\EN
\LessonHead{Lesson II}
\switchcolumn*

\LA
\noindent Sábbato medicínæ Domínicæ ópera cœpta signíficat, ut inde nova creatúra cœ́perit, ubi vetus creatúra ante desívit: nec sub lege esse Dei Fílium, sed supra legem in ipso princípio designáret: nec solvi legem, sed impléri. Neque enim per legem, sed verbo factus est mundus, sicut légimus: Verbo Dómini cæli firmáti sunt. Non sólvitur ergo lex, sed implétur: ut fiat renovátio hóminis jam labéntis. Unde et Apóstolus ait: Exspoliántes vos véterem hóminem, indúite novum, qui secúndum Deum creátus est.\par
\switchcolumn
\EN
\noindent That the Lord began to heal on the Sabbath-day showeth in a figure how that the new creation beginneth where the old creation ended. It showeth, moreover, that the Son of God, Who is come not to destroy the law but to fulfill the law, (Matth. v. 17,) is not under the law, but above the law. Neither was it by the law, but by the Word, that the world was created, as it is written “By the Word of the Lord were the heavens made.” (Ps. xxxii. 6.) The law, then, is not destroyed, but fulfilled, in the Redemption of fallen man. Whence also the Apostle saith: “Put off, concerning the former conversation, the old man, which is corrupt according to the deceitful lusts and be renewed in the spirit of your mind and put on the new man, which after God is created in righteousness and true holiness.” (Eph. iv. 22.)\par
\switchcolumn*

\LA
\begin{Resp}\Rsign Dixit Judas frátribus suis: Ecce Ismaëlítæ tránseunt; veníte, venumdétur, et manus nostræ non polluántur: \Star{} Caro enim et frater noster est.\end{Resp}
\begin{Resp}\Vsign Quid enim prodest, si occidérimus fratrem nostrum, et celavérimus sánguinem ipsíus? mélius est ut venumdétur.\end{Resp}
\begin{Resp}\Rsign Caro enim et frater noster est.\end{Resp}
\switchcolumn
\EN
\begin{Resp}\Rsign Judah said unto his brethren: Behold, the Ismaelites pass by; come, let us sell him, and let not our hands be defiled. \Star{} For he is our flesh, and our brother.\end{Resp}
\begin{Resp}\Vsign What profit is it if we slay our brother, and conceal his blood? It is better to sell him.\end{Resp}
\begin{Resp}\Rsign For he is our flesh, and our brother.\end{Resp}
\switchcolumn*

\LA
\LessonHead{Lectio III}
\switchcolumn
\EN
\LessonHead{Lesson III}
\switchcolumn*

\LA
\noindent Et bene sábbato cœpit, ut ipsum se osténderet Creatórem, qui ópera opéribus intéxeret, et prosequerétur opus, quod ipse jam cœ́perat: ut si domum faber renováre dispónat, non a fundaméntis, sed a culmínibus íncipit sólvere vetustátem. Itaque ibi prius manum ádmovet, ubi ante desíerat: deínde a minóribus íncipit, ut ad majóra pervéniat. Liberáre a dǽmone et hómines, sed in verbo Dei possunt: resurrectiónem mórtuis imperáre, divínæ solíus est potestátis. Fortássis étiam in typo mulíeris illíus socrus Simónis et Andréæ, váriis críminum fébribus caro nostra languébat, et diversárum cupiditátum immódicis æstuábat illécebris. Nec minórem febrem amóris esse díxerim, quam calóris. Itaque illa ánimum, hæc corpus inflámmat. Febris enim nostra, avarítia est: febris nostra, libído est: febris nostra, luxúria est: febris nostra, ambítio est: febris nostra, iracúndia est.\par
\switchcolumn
\EN
\noindent It was well that He began to heal on the Sabbath, that He might show Himself to be the Creator, weaving in one with another of His works, and continuing that which He had already begun, even as a workman, being to repair an house, beginneth not to take down that which is old from the foundations, but from the roof. Thus doth the Lord begin to lay to His hand again, in that place whence last He hath lifted it; then He beginneth which after God is created in righteousness and true holiness: (Eph. iv. 22.) with things lesser, that He may go on to things greater. Even men are able to deliver other men from evil spirits, albeit with the word of God: to command the dead to rise again is for God's power alone. Perchance, also, this woman, the mother-in-law of Simon and Andrew, was a type of our nature, stricken down with the great fever of sin, and burning with unlawful lusts after diverse objects. Nor would I say that the passion which rageth in the mind is a lesser fire than that fever which burneth the body. Covetousness, and lust, and uncleanness, and vain desires, and strivings, and anger, these be our fevers.\par
\switchcolumn*

\LA
\begin{Resp}\Rsign Extrahéntes Joseph de lacu, vendidérunt Ismaëlítæ vigínti argénteis: \Star{} Reversúsque Ruben ad púteum, cum non invenísset eum, scidit vestiménta sua cum fletu, et dixit: * Puer non compáret, et ego quo ibo?\end{Resp}
\begin{Resp}\Vsign At illi, intíncta túnica Joseph in sánguine hædi, misérunt qui ferret eam ad patrem, et díceret: Vide, si túnica fílii tui sit, an non.\end{Resp}
\begin{Resp}\Rsign Reversúsque Ruben ad púteum, cum non invenísset eum, scidit vestiménta sua cum fletu, et dixit.\end{Resp}
\begin{Resp}\Vsign Glória Patri, et Fílio, \Star{} et Spirítui Sancto.\end{Resp}
\begin{Resp}\Rsign Puer non compáret, et ego quo ibo?\end{Resp}
\switchcolumn
\EN
\begin{Resp}\Rsign They drew up Joseph out of the pit, and sold him to the Ishmaelites for twenty pieces of silver. \Star{} And Reuben returned unto the pit, and when he found not Joseph, he rent his clothes, and wept, and said: * The child is not, and I, whither shall I go?\end{Resp}
\begin{Resp}\Vsign And they took Joseph's coat, and dipped it in the blood of a kid of the goats, and they sent one that brought the coat unto their father, and said: See now whether this be thy son's coat or no.\end{Resp}
\begin{Resp}\Rsign And Reuben returned unto the pit, and when he found not Joseph, he rent his clothes, and wept.\end{Resp}
\begin{Resp}\Vsign Glory be to the Father, and to the Son, \Star{} and to the Holy Ghost.\end{Resp}
\begin{Resp}\Rsign And Reuben said: The child is not, and I, whither shall I go?\end{Resp}
\switchcolumn*

\end{paracol}

\Office{Feria Sexta infra Hebdomadam III in Quadragesima}{Friday of the Third Week of Lent}{Feria major}
\addcontentsline{toc}{subsection}{Feria Sexta infra Hebdomadam III in Quadragesima\enspace\textit{Friday of the Third Week of Lent}}
\begin{paracol}{2}
\LA
\LessonHead{Lectio I}
\switchcolumn
\EN
\LessonHead{Lesson I}
\switchcolumn*

\LA
\LessonTitle{Léctio sancti Evangélii secúndum Joánnem}
\switchcolumn
\EN
\LessonTitle{Continuation of the Holy Gospel according to John}
\switchcolumn*

\LA
\Cite{Joannes 4:5-42}
\switchcolumn
\EN
\Cite{John 4:5-42}
\switchcolumn*

\LA
\noindent In illo témpore: Venit Jesus in civitátem Samaríæ, quæ dícitur Sichar: juxta prǽdium, quod dedit Jacob Joseph fílio suo. Et réliqua.\par
\switchcolumn
\EN
\noindent At that time, Jesus cometh therefore to a city of Samaria, which is called Sichar, near the land which Jacob gave to his son Joseph. And so on.\par
\switchcolumn*

\LA
\LessonTitle{Homilía sancti Augustíni Epíscopi}
\switchcolumn
\EN
\LessonTitle{Homily by St. Augustine, Bishop (of Hippo.)}
\switchcolumn*

\LA
\Source{Tract. 15 in Joánnem, post init.}
\switchcolumn
\EN
\Source{15th Tract on John}
\switchcolumn*

\LA
\noindent Jam incípiunt mystéria. Non enim frustra fatigátur Jesus: non enim frustra fatigátur virtus Dei: non enim frustra fatigátur, per quem fatigáti recreántur: non enim frustra fatigátur, quo deserénte fatigámur, quo præsénte firmámur. Fatigátur tamen Jesus, et fatigátur ab itínere, et sedet, et juxta púteum sedet, et hora sexta fatigátus sedet. Omnia ista ínnuunt áliquid, indicáre volunt áliquid: ut pulsémus, hortántur. Ipse ergo apériat et nobis et vobis, qui dignátus est ita hortári, ut díceret: Pulsáte, et aperiétur vobis.\par
\switchcolumn
\EN
\noindent Jesus, wearied with His journey, the mysteries are beginning now. It is not for nothing that Jesus is wearied. It is not for nothing that the Mighty One of God is wearied. It is not for nothing that He is wearied Who Himself giveth Rest to all them that are weary and heavy-laden. It is not for nothing that He is wearied Whose absence prostrateth us, and Whose presence maketh us to be strong.\par
\switchcolumn*

\LA
\begin{Resp}\Rsign Videns Jacob vestiménta Joseph, scidit vestiménta sua cum fletu, et dixit: \Star{} Fera péssima devorávit fílium meum Joseph.\end{Resp}
\begin{Resp}\Vsign Tulérunt autem fratres ejus túnicam illíus, mitténtes ad patrem: quam cum cognovísset pater, ait.\end{Resp}
\begin{Resp}\Rsign Fera péssima devorávit fílium meum Joseph.\end{Resp}
\switchcolumn
\EN
\begin{Resp}\Rsign When Jacob saw Joseph's coat he rent his clothes, and mourned; and he said: \Star{} An evil beast hath devoured my son Joseph.\end{Resp}
\begin{Resp}\Vsign And his brethren took his coat, and sent it to his father and he knew it, and said:\end{Resp}
\begin{Resp}\Rsign An evil beast hath devoured my son Joseph.\end{Resp}
\switchcolumn*

\LA
\LessonHead{Lectio II}
\switchcolumn
\EN
\LessonHead{Lesson II}
\switchcolumn*

\LA
\noindent Tibi fatigátus est ab itínere Jesus. Invenímus virtútem Jesum, et invenímus infírmum Jesum: fortem, et infírmum. Fortem, quia in princípio erat Verbum, et Verbum erat apud Deum, et Deus erat Verbum: hoc erat in princípio apud Deum. Vis vidére quam iste Fílius Dei fortis sit? Omnia per ipsum facta sunt, et sine ipso factum est nihil: et sine labóre facta sunt. Quid ergo illo fórtius, per quem sine labóre facta sunt ómnia? Infírmum vis nosse? Verbum caro factum est, et habitávit in nobis. Fortitúdo Christi te creávit: infírmitas Christi te recreávit. Fortitúdo Christi fecit, ut quod non erat, esset: infírmitas Christi fecit, ut quod erat, non períret. Cóndidit nos fortitúdine sua, quæsívit nos infirmitáte sua.\par
\switchcolumn
\EN
\noindent Jesus, therefore, being wearied with His journey, sat thus on the well, and it was about the sixth hour. There is a depth in all these details they all have something to say for us to learn. Upon them we gaze. Knock, saith the Lord, and it shall be opened unto you. Let us knock then and, O, may He open to me and to you, even He Who hath spoken to us those words: Knock, and it shall be opened unto you. (Matth. vii. 7.) It is for thy sake that Jesus was wearied with His journey. We find the strength of Jesus, and we find Jesus weak; yea, strong and weak. Strong, for In the beginning was the Word, and the Word was with God, and the Word was God: the Same was in the beginning with God. Wouldest thou know again how that the Son of God is strong? All things were made by Him, and without Him was not anything made that was made without effort. (John i. 1-3.) What then is stronger than He by Whom all things were made without effort? Wouldest thou know His weakness? The Word was made Flesh and dwelt among us. Christ, strong, made thee; Christ, weak, redeemed thee. Christ, strong, made all things out of nothing; Christ, weak, so wrought that that was made perished not. His strength hath made us, and His weakness saved us.\par
\switchcolumn*

\LA
\begin{Resp}\Rsign Joseph dum intráret in terram Ægýpti, linguam quam non nóverat, audívit: manus ejus in labóribus serviérunt: \Star{} Et lingua ejus inter príncipes loquebátur sapiéntiam.\end{Resp}
\begin{Resp}\Vsign Humiliavérunt in compédibus pedes ejus: ferrum petránsiit ánimam ejus, donec veníret verbum ejus.\end{Resp}
\begin{Resp}\Rsign Et lingua ejus inter príncipes loquebátur sapiéntiam.\end{Resp}
\switchcolumn
\EN
\begin{Resp}\Rsign When Joseph came into the land of Egypt, he heard a language that he understood not; his hands were burdened with labour; \Star{} And his tongue spake wisdom among princes.\end{Resp}
\begin{Resp}\Vsign Whose feet they hurt with fetters; the iron entered into his soul, until the time that his word came\end{Resp}
\begin{Resp}\Rsign And his tongue spake wisdom among princes.\end{Resp}
\switchcolumn*

\LA
\LessonHead{Lectio III}
\switchcolumn
\EN
\LessonHead{Lesson III}
\switchcolumn*

\LA
\noindent Nutrit ergo ipse infírmus infírmos, tamquam gallína pullos suos: huic enim se símilem fecit. Quóties vólui, inquit ad Jerúsalem, congregáre fílios tuos sub alas tamquam gallína pullos suos, et noluísti? Vidétis autem, fratres, quemádmodum gallína infirmétur cum pullis suis. Nulla enim ália avis, quod sit mater, agnóscitur. Vidémus nidificáre pásseres quóslibet ante óculos nostros: hirúndines, cicónias, colúmbas cotídie vidémus nidificáre: quos, nisi quando in nidis vidémus, paréntes esse non agnóscimus. Gallína vero sic infirmátur in pullis suis, ut étiam si ipsi pulli non sequántur, fílios non vídeas, matrem tamen intéllegas.\par
\switchcolumn
\EN
\noindent He then, being Himself made weak, is strength to all such as are weak, gathering them together, to use His own figure, even as an hen gathereth her chickens under her wings. O Jerusalem, Jerusalem! how often would I have gathered thy children together, even as an hen gathereth her chickens under her wings, and ye would not! (Matth. xxiii. 37.) Consider now, my brethren, in what bondage is an hen to her chickens. There is no other bird in whom motherhood is unmistakeable. We watch the sparrows building their nests under our eyes; we see swallows, and storks, and pigeons building theirs every day. But, unless we actually see them in their nests, we know not if they have little ones, or no. But the hen's motherhood is so much a part of herself, that even if at the minute we see not her children the chickens following after her, nevertheless we see by her ways if she be a mother.\par
\switchcolumn*

\LA
\begin{Resp}\Rsign Meménto mei, dum bene tibi fúerit: \Star{} Ut súggeras pharaóni, ut edúcat me de isto cárcere: * Quia furtim sublátus sum, et hic ínnocens in lacum missus sum.\end{Resp}
\begin{Resp}\Vsign Tres enim adhuc dies sunt, post quos recordábitur phárao ministérii tui, et restítuet te in gradum prístinum: tunc meménto mei.\end{Resp}
\begin{Resp}\Rsign Ut súggeras pharaóni, ut edúcat me de isto cárcere.\end{Resp}
\begin{Resp}\Vsign Glória Patri, et Fílio, \Star{} et Spirítui Sancto.\end{Resp}
\begin{Resp}\Rsign Quia furtim sublátus sum, et hic ínnocens in lacum missus sum.\end{Resp}
\switchcolumn
\EN
\begin{Resp}\Rsign Think on me when it shall be well with thee; \Star{} And make mention of me unto Pharaoh, that he may bring me out of this prison. * For I was stolen away; and here have I done nothing, that they should put me into the dungeon.\end{Resp}
\begin{Resp}\Vsign For yet three days, and then Pharaoh shall remember thy service, and restore thee unto thy place; then think of me.\end{Resp}
\begin{Resp}\Rsign And make mention of me unto Pharaoh, that he may bring me out of this prison.\end{Resp}
\begin{Resp}\Vsign Glory be to the Father, and to the Son, \Star{} and to the Holy Ghost.\end{Resp}
\begin{Resp}\Rsign For I was stolen away; and here have I done nothing, that they should put me into the dungeon.\end{Resp}
\switchcolumn*

\end{paracol}

\Office{Sabbato infra Hebdomadam III in Quadragesima}{Saturday of the Third Week of Lent}{Feria major}
\addcontentsline{toc}{subsection}{Sabbato infra Hebdomadam III in Quadragesima\enspace\textit{Saturday of the Third Week of Lent}}
\begin{paracol}{2}
\LA
\LessonHead{Lectio I}
\switchcolumn
\EN
\LessonHead{Lesson I}
\switchcolumn*

\LA
\LessonTitle{Léctio sancti Evangélii secúndum Joánnem}
\switchcolumn
\EN
\LessonTitle{Continuation of the Holy Gospel according to John}
\switchcolumn*

\LA
\Cite{Joannes 8:1-11}
\switchcolumn
\EN
\Cite{John 8:1-11}
\switchcolumn*

\LA
\noindent In illo témpore: Perréxit Jesus in montem Olivéti, et dilúculo íterum venit in templum. Et réliqua.\par
\switchcolumn
\EN
\noindent In that time, Jesus went unto Mount Olivet. And early in the morning he came again into the temple. And so on.\par
\switchcolumn*

\LA
\LessonTitle{Homilía sancti Augustíni Epíscopi}
\switchcolumn
\EN
\LessonTitle{Homily by St. Augustine, Bishop (of Hippo.)}
\switchcolumn*

\LA
\Source{Tract. 33 in Joannem, post init.}
\switchcolumn
\EN
\Source{33rd Tract on John.}
\switchcolumn*

\LA
\noindent Jesus perréxit in montem Olivéti, in montem fructuósum, in montem unguénti, in montem chrísmatis. Ubi enim decébat docére Christum, nisi in monte Olivéti? Christi enim nomen a chrísmate dictum est: chrisma autem Græce, Latíne únctio nominátur. Ideo autem nos unxit, quia luctatóres contra diábolum fecit. Et dilúculo íterum venit in templum, et omnis pópulus venit ad eum: et sedens docébat eos, et non tenebátur, quia nondum pati dignabátur. Nunc jam atténdite, ubi ab inimícis tentáta sit Dómini mansuetúdo.\par
\switchcolumn
\EN
\noindent Jesus went unto the Mount of Olives, even unto that fruitful Mount, that anointing Mount, that Mount of Chrism. Where else became it Christ to teach, if not on the Mount of Olives? For the word Christ is derived from Chrisma, and Chrisma is the Greek for ointment. He hath anointed us that we may be able to wrestle with the devil. And, early in the morning, He came again into the temple; and all the people came unto Him; and He sat down, and taught them and no man laid hands on Him, because He was not yet pleased to suffer. And now listen how His enemies tried the Lord's meekness.\par
\switchcolumn*

\LA
\begin{Resp}\Rsign Mérito hæc pátimur, quia peccávimus in fratrem nostrum, vidéntes angústias ánimæ ejus, dum deprecarétur nos, et non audívimus: \Star{} Idcírco venit super nos tribulátio.\end{Resp}
\begin{Resp}\Vsign Dixit Ruben frátribus suis: Numquid non dixi vobis, Nolíte peccáre in púerum; et non audístis me?\end{Resp}
\begin{Resp}\Rsign Idcírco venit super nos tribulatio.\end{Resp}
\switchcolumn
\EN
\begin{Resp}\Rsign We are verily guilty concerning our brother, in that we saw the anguish of his soul, when he besought us, and we would not hear. \Star{} Therefore is this distress come upon us.\end{Resp}
\begin{Resp}\Vsign And Reuben answered his brethren, saying Spake I not unto you, saying: Do not sin against the child; and ye would not hear?\end{Resp}
\begin{Resp}\Rsign Therefore is this distress come upon us.\end{Resp}
\switchcolumn*

\LA
\LessonHead{Lectio II}
\switchcolumn
\EN
\LessonHead{Lesson II}
\switchcolumn*

\LA
\noindent Addúcunt autem illi scribæ et pharisǽi mulíerem in adultério deprehénsam, et statuérunt eam in médio, et dixérunt ei: Magíster, hæc múlier modo deprehénsa est in adultério: in lege autem Móyses mandávit nobis hujúsmodi lapidáre: tu ergo quid dicis? Hæc autem dicébant tentántes eum: ut possent accusáre eum. Unde accusáre? Numquid ipsum in áliquo facínore deprehénderant, aut illa múlier ad eum áliquo modo pertinuísse dicebátur?\par
\switchcolumn
\EN
\noindent And the Scribes and Pharisees brought unto Him a woman taken in adultery; and when they had set her in the midst, they say unto Him: Master, this woman was taken in adultery, in the very act. Now, Moses in the law commanded that such should be stoned; but what sayest Thou? This they said, tempting Him, that they might have to accuse Him. Whereof to accuse Him? Had they taken Him in any sin? Or was the woman said to have anything to do with Him?\par
\switchcolumn*

\LA
\begin{Resp}\Rsign Dixit Ruben frátribus suis: Numquid non dixi vobis, Nolíte peccáre in púerum, et non audístis me? \Star{} En sanguis ejus exquíritur.\end{Resp}
\begin{Resp}\Vsign Mérito hæc pátimur, quia peccávimus in fratrem nostrum, vidéntes angústias ánimæ ejus, dum deprecarétur nos, et non audívimus.\end{Resp}
\begin{Resp}\Rsign En sanguis ejus exquíritur.\end{Resp}
\switchcolumn
\EN
\begin{Resp}\Rsign And Reuben answered his brethren, saying: Spake I not unto you, saying: Do not sin against the child; and ye would not hear? \Star{} Behold, his blood is required.\end{Resp}
\begin{Resp}\Vsign We are verily guilty concerning our brother, in that we saw the anguish of his soul, when he besought us, and we would not hear.\end{Resp}
\begin{Resp}\Rsign Behold, his blood is required.\end{Resp}
\switchcolumn*

\LA
\LessonHead{Lectio III}
\switchcolumn
\EN
\LessonHead{Lesson III}
\switchcolumn*

\LA
\noindent Intellegámus, fratres, admirábilem mansuetúdinem in Dómino fuísse. Animadvertérunt eum nímium esse mitem, nímium esse mansuétum. De illo quippe fúerat ante prædíctum: Accíngere gládio tuo circa femur tuum, potentíssime. Spécie tua et pulchritúdine tua inténde, próspere procéde, et regna: propter veritátem, et mansuetúdinem, et justítiam. Ergo áttulit veritátem ut doctor, mansuetúdinem ut liberátor, justítiam ut cógnitor. Propter hæc eum esse regnatúrum in Spíritu Sancto prophéta prædíxerat. Cum loquerétur, véritas agnoscebátur: cum advérsus inimícos non moverétur, mansuetúdo laudabátur. Cum ergo de duóbus istis, id est, de veritáte et mansuetúdine ejus, inimíci livóre et invídia torqueréntur; in tértio, id est justítia, scándalum posuérunt.\par
\switchcolumn
\EN
\noindent We must understand, my brethren, that there was a wonderful gentleness in the Lord. They knew that He was most mild and most gentle. Of Him indeed it had been said of old time: Gird thy sword upon thy thigh, O most Mighty! In thy comeliness and thy beauty go forward, fare prosperously, and reign, because of truth, and meekness, and righteousness. (Ps. xliv. 4, 5.) And He came bringing truth as one that teacheth, meekness as one that delivereth, and righteousness as one that knoweth. Because of these it was that the Prophet declared, in the Holy Ghost, that He was to reign. Whenever He spake, truth shone forth whenever He spared His enemies, meekness was made glorious. And His enemies, racked with envy and hatred by His truth and His meekness, laid a stumbling-block for His righteousness.\par
\switchcolumn*

\LA
\begin{Resp}\Rsign Lamentabátur Jacob de duóbus fíliis suis: Heu me, dolens sum de Joseph pérdito, et tristis nimis de Bénjamin ducto pro alimóniis: \Star{} Precor cæléstem Regem, ut me doléntem nímium fáciat eos cérnere.\end{Resp}
\begin{Resp}\Vsign Prostérnens se Jacob veheménter cum lácrimis pronus in terram, et adórans ait.\end{Resp}
\begin{Resp}\Rsign Precor cæléstem Regem, ut me doléntem nímium fáciat eos cérnere.\end{Resp}
\begin{Resp}\Vsign Glória Patri, et Fílio, \Star{} et Spirítui Sancto.\end{Resp}
\begin{Resp}\Rsign Precor cæléstem Regem, ut me doléntem nímium fáciat eos cérnere.\end{Resp}
\switchcolumn
\EN
\begin{Resp}\Rsign Jacob lamented for his two sons, saying: Woe is me; I am bereaved of Joseph, for he is not; and afflicted because of Benjamin, because he is taken away for bread. \Star{} I pray the King of heaven in my distress, that He may make me to see them yet again.\end{Resp}
\begin{Resp}\Vsign And Jacob cast him down upon his face upon the ground, and wept sore; and he prayed, saying:\end{Resp}
\begin{Resp}\Rsign I pray the King of heaven in my distress, that He may make me to see them yet again.\end{Resp}
\begin{Resp}\Vsign Glory be to the Father, and to the Son, \Star{} and to the Holy Ghost.\end{Resp}
\begin{Resp}\Rsign I pray the King of heaven in my distress, that He may make me to see them yet again.\end{Resp}
\switchcolumn*

\end{paracol}

\Office{Dominica IV in Quadragesima}{Fourth Sunday of Lent}{Semiduplex I classis}
\addcontentsline{toc}{subsection}{Dominica IV in Quadragesima\enspace\textit{Fourth Sunday of Lent}}
\begin{paracol}{2}
\LA
\LessonHead{Lectio I}
\switchcolumn
\EN
\LessonHead{Lesson I}
\switchcolumn*

\LA
\LessonTitle{De libro Exodi}
\switchcolumn
\EN
\LessonTitle{Lesson from the book of Exodus}
\switchcolumn*

\LA
\Cite{Exod 3:1-6}
\switchcolumn
\EN
\Cite{Exod 3:1-6}
\switchcolumn*

\LA
\noindent Móyses autem pascébat oves Jethro sóceri sui sacerdótis Mádian; cumque minásset gregem ad interióra desérti, venit ad montem Dei Horeb. Apparuítque ei Dóminus in flamma ignis de médio rubi; et vidébat quod rubus ardéret et non comburerétur. Dixit ergo Móyses: Vadam, et vidébo visiónem hanc magnam, quare non comburátur rubus. Cernens autem Dóminus quod pérgeret ad vidéndum, vocávit eum de médio rubi, et ait: Móyses, Móyses! Qui respóndit: Adsum. At ille: Ne apprópies, inquit, huc: solve calceaméntum de pédibus tuis; locus enim, in quo stas, terra sancta est. Et ait: Ego sum Deus patris tui, Deus Abraham, Deus Isaac, et Deus Jacob. Abscóndit Móyses fáciem suam: non enim audébat aspícere contra Deum.\par
\switchcolumn
\EN
\noindent Now Moses fed the sheep of Jethro his father in law, the priest of Madian: and he drove the flock to the inner parts of the desert, and came to the mountain of God, Horeb. And the Lord appeared to him in a flame of fire out of the midst of a bush: and he saw that the bush was on fire and was not burnt. And Moses said: I will go and see this great sight, why the bush is not burnt. And when the Lord saw that he went forward to see, he called to him out of the midst of the bush, and said: Moses, Moses. And he answered: Here I am. And he said: Come not nigh hither, put off the shoes from thy feet: for the place whereon thou standest is holy ground. And he said: I am the God of thy father, the God of Abraham, the God of Isaac, and the God of Jacob. Moses hid his face: for he durst not look at God.\par
\switchcolumn*

\LA
\begin{Resp}\Rsign Locútus est Dóminus ad Móysen, dicens: Descénde in Ægýptum, et dic Pharaóni: \Star{} Ut dimíttat pópulum meum: indurátum est cor Pharaónis: non vult dimíttere pópulum meum, nisi in manu forti.\end{Resp}
\begin{Resp}\Vsign Clamor filiórum Israël venit ad me, vidíque afflictiónem eórum: sed veni, mittam te ad Pharaónem.\end{Resp}
\begin{Resp}\Rsign Ut dimíttat pópulum meum: indurátum est cor Pharaónis: non vult dimíttere pópulum meum, nisi in manu forti.\end{Resp}
\switchcolumn
\EN
\begin{Resp}\Rsign The Lord spake unto Moses, saying: Go down now into Egypt, and say unto Pharaoh: \Star{} Let My people go. And the heart of Pharaoh shall be hardened, that he will not let My people go but by a mighty hand.\end{Resp}
\begin{Resp}\Vsign The cry of the children of Israel is come unto Me, and I have seen their affliction: come now, therefore, and I will send thee unto Pharaoh, and thou shalt say unto him:\end{Resp}
\begin{Resp}\Rsign Let My people go. And the heart of Pharaoh shall be hardened, that he will not let My people go but by a mighty hand.\end{Resp}
\switchcolumn*

\LA
\LessonHead{Lectio II}
\switchcolumn
\EN
\LessonHead{Lesson II}
\switchcolumn*

\LA
\Cite{Exod 3:7-10}
\switchcolumn
\EN
\Cite{Exod 3:7-10}
\switchcolumn*

\LA
\noindent Cui ait Dóminus: Vidi afflictiónem pópuli mei in Ægýpto, et clamórem ejus audívi propter durítiam eórum qui præsunt opéribus: et sciens dolórem ejus, descéndi ut líberem eum de mánibus Ægyptiórum, et edúcam de terra illa in terram bonam et spatiósam, in terram quæ fluit lacte et melle, ad loca Chananǽi, et Hethǽi, et Amorrhǽi, et Pherezǽi, et Hevǽi, et Jebusǽi. Clamor ergo filiórum Israël venit ad me: vidíque afflictiónem eórum, qua ab Ægýptiis opprimúntur. Sed veni, et mittam te ad Pharaónem, ut edúcas pópulum meum, fílios Israël de Ægýpto.\par
\switchcolumn
\EN
\noindent And the Lord said to him: I have seen the affliction of my people in Egypt, and I have heard their cry because of the rigour of them that are over the works: And knowing their sorrow, I am come down to deliver them out of the hands of the Egyptians, and to bring them out of that land into a good and spacious land, into a land that floweth with milk and honey, to the places of the Chanaanite, and Hethite, and Amorrhite, and Pherezite, and Hevite, and Jebusite. For the cry of the children of Israel is come unto me: and I have seen their affliction, wherewith they are oppressed by the Egyptians. But come, and I will send thee to Pharao, that thou mayst bring forth my people, the children of Israel out of Egypt.\par
\switchcolumn*

\LA
\begin{Resp}\Rsign Stetit Móyses coram Pharaóne, et dixit: Hæc dicit Dóminus: \Star{} Dimítte pópulum meum, ut sacríficet mihi in desérto.\end{Resp}
\begin{Resp}\Vsign Dóminus Deus Hebræórum misit me ad te, dicens.\end{Resp}
\begin{Resp}\Rsign Dimítte pópulum meum, ut sacríficet mihi in desérto.\end{Resp}
\switchcolumn
\EN
\begin{Resp}\Rsign Moses stood before Pharaoh, and said: Thus saith the Lord: \Star{} Let My people go, that they may hold a feast unto Me in the wilderness.\end{Resp}
\begin{Resp}\Vsign The Lord God of the Hebrews hath sent me unto thee, saying:\end{Resp}
\begin{Resp}\Rsign Let My people go, that they may hold a feast unto Me in the wilderness.\end{Resp}
\switchcolumn*

\LA
\LessonHead{Lectio III}
\switchcolumn
\EN
\LessonHead{Lesson III}
\switchcolumn*

\LA
\Cite{Exod 3:11-15}
\switchcolumn
\EN
\Cite{Exod 3:11-15}
\switchcolumn*

\LA
\noindent Dixítque Móyses ad Deum: Quis sum ego, ut vadam ad Pharaónem, et edúcam fílios Israël de Ægýpto? Qui dixit ei: Ego ero tecum: et hoc habébis signum, quod míserim te: Cum edúxeris pópulum de Ægýpto, immolábis Deo super montem istum. Ait Móyses ad Deum: Ecce, ego vadam ad fílios Israël, et dicam eis: Deus patrum vestrórum misit me ad vos. Si díxerint mihi: Quod est nomen ejus? quid dicam eis? Dixit Deus ad Móysen: Ego sum qui sum. Ait: Sic dices fíliis Israël: Qui est, misit me ad vos. Dixítque íterum Deus ad Móysen: Hæc dices fíliis Israël: Dóminus Deus patrum vestrórum, Deus Abraham, Deus Isaac, et Deus Jacob, misit me ad vos; hoc nomen mihi est in ætérnum, et hoc memoriále meum in generatiónem et generatiónem.\par
\switchcolumn
\EN
\noindent And Moses said to God: Who am I that I should go to Pharao, and should bring forth the children of Israel out of Egypt? And he said to him: I will be with thee: and this thou shalt have for a sign, that I have sent thee: When thou shalt have brought my people out of Egypt, thou shalt offer sacrifice to God upon this mountain. Moses said to God: Lo, I shall go to the children of Israel, and say to them: The God of your fathers hath sent me to you. If they should say to me: What is his name? what shall I say to them? God said to Moses: I AM WHO AM. He said: Thus shalt thou say to the children of Israel: HE WHO IS, hath sent me to you. And God said again to Moses: Thus shalt thou say to the children of Israel: The Lord God of your fathers, the God of Abraham, the God of Isaac, and the God of Jacob, hath sent me to you: This is my name for ever, and this is my memorial unto all generations.\par
\switchcolumn*

\LA
\begin{Resp}\Rsign Cantémus Dómino: glorióse enim honorificátus est, equum et ascensórem projécit in mare: \Star{} Adjútor et protéctor factus est mihi Dóminus in salútem.\end{Resp}
\begin{Resp}\Vsign Dóminus quasi vir pugnátor, Omnípotens nomen ejus.\end{Resp}
\begin{Resp}\Rsign Adjútor et protéctor factus est mihi Dóminus in salútem.\end{Resp}
\begin{Resp}\Vsign Glória Patri, et Fílio, \Star{} et Spirítui Sancto.\end{Resp}
\begin{Resp}\Rsign Adjútor et protéctor factus est mihi Dóminus in salútem.\end{Resp}
\switchcolumn
\EN
\begin{Resp}\Rsign Let us sing unto the Lord, for He hath triumphed gloriously; the horse and his rider hath He thrown into the sea. \Star{} The Lord is my strength and song, and He is become my salvation.\end{Resp}
\begin{Resp}\Vsign The Lord is a man of war; Almighty is His Name.\end{Resp}
\begin{Resp}\Rsign The Lord is my strength and song, and He is become my salvation.\end{Resp}
\begin{Resp}\Vsign Glory be to the Father, and to the Son, \Star{} and to the Holy Ghost.\end{Resp}
\begin{Resp}\Rsign The Lord is my strength, and my song, and He is become my salvation.\end{Resp}
\switchcolumn*

\LA
\LessonHead{Lectio IV}
\switchcolumn
\EN
\LessonHead{Lesson IV}
\switchcolumn*

\LA
\LessonTitle{Sermo sancti Basilíi Magni}
\switchcolumn
\EN
\LessonTitle{From the Sermons of St. Basil the Great, Archbishop (of Caesarea-in-Pontus.)}
\switchcolumn*

\LA
\Source{Homilia 1 de jejúnio, ante med.}
\switchcolumn
\EN
\Source{1st on Fasting}
\switchcolumn*

\LA
\noindent Móysen per jejúnium nóvimus in montem ascendísse: neque enim áliter ausus esset vérticem fumántem adíre, atque in calíginem íngredi, nisi jejúnio munítus. Per jejúnium mandáta dígito Dei in tábulis conscrípta suscépit. Item supra montem jejúnium legis latæ conciliátor fuit: inférius vero, gula ad idololatríam pópulum dedúxit, ac contaminávit. Sedit, inquit, pópulus manducáre et bíbere, et surrexérunt lúdere. Quadragínta diérum labórem ac perseverántiam, Dei servo contínuo jejunánte ac oránte, una tantum pópuli ebríetas cassam irritámque réddidit. Quas enim tábulas Dei dígito conscríptas jejúnium accépit, has ebríetas contrívit: Prophéta sanctíssimo indígnum existimánte, vinoléntum pópulum a Deo legem accípere.\par
\switchcolumn
\EN
\noindent We know that it was with and by fasting that Moses went up into the Mount, for he had not dared to go up to that smoking summit, nor to have entered that darkness, except he had been made strong by a Fast. It was with fasting that he received the commandments, written by the finger of God upon tables of stone. Upon the mountain, that Fast made interest with Him Whose law was given unto it; but, below, gluttony was leading the people to the worship of idols and polluting them. It is written The people sat down to eat and to drink, and rose up to play. (Ex. xxxii. 6.) That one fit of drunken frenzy, on the part of the people, made void and of none effect all the toil and patience of the forty days, during the which the servant of God had fasted and prayed unceasingly. To the Fast had been given those tables of stone written on with the finger of God; the Feast's work was to break them, by the hand of the most holy prophet, who deemed a nation of drunkards a nation unmeet to receive law from God.\par
\switchcolumn*

\LA
\begin{Resp}\Rsign In mare viæ tuæ, et sémitæ tuæ in aquis multis: \Star{} Deduxísti sicut oves pópulum tuum in manu Móysi et Aaron.\end{Resp}
\begin{Resp}\Vsign Transtulísti illos per mare Rubrum, et transvexísti eos per aquam nímiam.\end{Resp}
\begin{Resp}\Rsign Deduxísti sicut oves pópulum tuum in manu Móysi et Aaron.\end{Resp}
\switchcolumn
\EN
\begin{Resp}\Rsign Thy way is in the sea, and thy paths in the great waters. \Star{} Thou led thy people like a flock, by the hand of Moses and Aaron.\end{Resp}
\begin{Resp}\Vsign Thou broughtest them through the Red Sea, and led them through much water.\end{Resp}
\begin{Resp}\Rsign Thou leddest thy people like a flock by the hand of Moses and Aaron.\end{Resp}
\switchcolumn*

\LA
\LessonHead{Lectio V}
\switchcolumn
\EN
\LessonHead{Lesson V}
\switchcolumn*

\LA
\noindent Uno témporis moménto ob gulam pópulus ille per máxima prodígia Dei cultum edóctus, in Ægyptíacam idololatríam turpíssime devolútus est. Ex quo si utrúmque simul cónferas, vidére licet, jejúnium ad Deum dúcere, delícias vero salútem pérdere. Quid Esau inquinávit, servúmque fratris réddidit? nonne esca una, propter quam primogénita véndidit? Samuélem vero nonne per jejúnium orátio largíta est matri? Quid fortíssimum Samsónem inexpugnábilem réddidit? nonne jejúnium, cum quo in matris ventre concéptus est? Jejúnium concépit, jejúnium nutrívit, jejúnium virum effécit. Quod sane Angelus matri præcépit, monens quæcúmque ex vite procéderent, ne attíngeret, non vinum, non síceram bíberet. Jejúnium prophétas génuit, poténtes confírmat atque róborat.\par
\switchcolumn
\EN
\noindent In a moment of time, that people, who had by great wonders been taught to worship God, were, by gluttony, dropped back into the cesspool of Egyptian idolatry. The which things if thou wilt consider, thou shalt see that the tendency of fasting is to God-ward, and that that of feasting is to hell-ward. What was it that degraded Esau, and made him a slave to his brother? Was it not that one dish of pottage for which he sold his birthright? (Gen. xxv. 29-34.) Was it not prayer when joined to fasting that gave Samuel to his mother? (1 Kings (Sam.) i. 7, 19.) What made the mighty Samson invincible? Was it not the fast during the which he was conceived in his mother's womb? The fast it was which made him to be conceived; the fast, which fed him; the fast, which made a man of him, even as the Angel of the Lord commanded his mother, saying: She may not eat of anything that cometh of the vine, neither let her drink wine or strong drink. (Judges xiii. 14.) Fasting is the mother of prophets, the strength and stay of mighty men.\par
\switchcolumn*

\LA
\begin{Resp}\Rsign Qui persequebántur pópulum tuum, Dómine, demersísti eos in profúndum: \Star{} Et in colúmna nubis ductor eórum fuísti.\end{Resp}
\begin{Resp}\Vsign Deduxísti sicut oves pópulum tuum in manu Móysi et Aaron.\end{Resp}
\begin{Resp}\Rsign Et in colúmna nubis ductor eórum fuísti.\end{Resp}
\switchcolumn
\EN
\begin{Resp}\Rsign O Lord, Thou hast overwhelmed in the deep of the sea them which persecuted thy people; \Star{} Even thy people which Thou led in the pillar of the cloud.\end{Resp}
\begin{Resp}\Vsign Thou led thy people like a flock, by the hand of Moses and Aaron.\end{Resp}
\begin{Resp}\Rsign Even thy people, which Thou led in the pillar of the cloud.\end{Resp}
\switchcolumn*

\LA
\LessonHead{Lectio VI}
\switchcolumn
\EN
\LessonHead{Lesson VI}
\switchcolumn*

\LA
\noindent Jejúnium legislatóres sapiéntes facit: ánimæ óptima custódia, córporis sócius secúrus, fórtibus viris muniméntum et arma, athlétis et certántibus exercitátio. Hoc prætérea tentatiónes propúlsat, ad pietátem armat, cum sobrietáte hábitat, temperántiæ ópifex est: in bellis fortitúdinem affert, in pace quiétem docet: nazarǽum sanctíficat, sacerdótem pérficit: neque enim fas est sine jejúnio sacrifícium attíngere, non solum in mýstica nunc et vera Dei adoratióne, sed nec in illa, in qua sacrifícium secúndum legem in figúra offerebátur. Jejúnium Elíam magnæ visiónis spectatórem fecit: quadragínta namque diérum jejúnio cum ánimam purgásset, in spelúnca méruit, quantum fas est hómini, Deum vidére. Móyses íterum legem accípiens, íterum jejúnia secútus est. Ninivítæ, nisi cum illis et bruta jejunássent, ruínæ minas nequáquam evasíssent. In desérto autem quorúmnam membra cecidérunt? nonne illórum, qui carnes appetivére?\par
\switchcolumn
\EN
\noindent It is fasting which giveth wisdom to lawgivers; fasting which is the trustiest keeper of the soul, and the safest companion for the body. It is fasting which is strength and armour to mighty men; fasting which maketh supple them which run and which wrestle. It is fasting which maketh a man strong to strive against temptation, and which is to godliness as a fenced city; even fasting, whose fellow is soberness, and her work temperance. It is fasting which maketh men to I wax valiant in fight; fasting which teacheth to rest in time of peace. Fasting maketh a Nazarite to be holy, and a priest perfect. Without a fast it is unlawful to touch the Sacrifice, not only in that mystic and true worship of God which now is, but also according to the law, in those sacrifices which were offered of old time as figures of the true. It was fasting which opened the eyes of Elias to look upon the visions of God, even as it is written, that when he had fasted forty days and forty nights he was in the mount of God, even Horeb, and he was made able, so far as man may be made able, to see God. (3 Kings xix. seq.) Even so also was Moses in that Mount forty days and forty nights, fasting, at what time he again received the Law. (Ex. xxxiv. 28.) Unless the Ninevites had fasted, both man and beast, herd and flock, they had not escaped from the ruin that hung over them. (Jonah iii. 710.) In the wilderness fell some and who were they? Yea, they were such as lusted after flesh meat. (Num. xi. 33.)\par
\switchcolumn*

\LA
\begin{Resp}\Rsign Móyses fámulus Dei jejunávit quadragínta diébus et quadragínta nóctibus: \Star{} Ut legem Dómini mererétur accípere.\end{Resp}
\begin{Resp}\Vsign Ascéndens Móyses in montem Sínai ad Dóminum, fuit ibi quadragínta diébus et quadragínta nóctibus.\end{Resp}
\begin{Resp}\Rsign Ut legem Dómini mererétur accípere.\end{Resp}
\begin{Resp}\Vsign Glória Patri, et Fílio, \Star{} et Spirítui Sancto.\end{Resp}
\begin{Resp}\Rsign Ut legem Dómini mererétur accípere.\end{Resp}
\switchcolumn
\EN
\begin{Resp}\Rsign Moses, the servant of God, fasted forty days and forty nights \Star{} To make him meet to receive the Law of God.\end{Resp}
\begin{Resp}\Vsign Moses went up unto the Lord into Mount Sinai, and he was in the Mount forty days and forty nights.\end{Resp}
\begin{Resp}\Rsign To make him meet to receive the Law of God.\end{Resp}
\begin{Resp}\Vsign Glory be to the Father, and to the Son, \Star{} and to the Holy Ghost.\end{Resp}
\begin{Resp}\Rsign To make him meet to receive the Law of God.\end{Resp}
\switchcolumn*

\LA
\LessonHead{Lectio VII}
\switchcolumn
\EN
\LessonHead{Lesson VII}
\switchcolumn*

\LA
\LessonTitle{Léctio sancti Evangélii secúndum Joánnem}
\switchcolumn
\EN
\LessonTitle{From the Holy Gospel according to John}
\switchcolumn*

\LA
\Cite{Joannes 6:1-15}
\switchcolumn
\EN
\Cite{John 6:1-15}
\switchcolumn*

\LA
\noindent In illo témpore: Abiit Jesus trans mare Galilǽæ, quod est Tiberíadis: et sequebátur eum multitúdo magna, quia vidébant signa, quæ faciébat super his qui infirmabántur. Et réliqua.\par
\switchcolumn
\EN
\noindent In that time, Jesus went over the sea of Galilee, which is that of Tiberias. And a great multitude followed him, because they saw the miracles which he did on them that were diseased. And so on.\par
\switchcolumn*

\LA
\LessonTitle{Homilía sancti Augustíni Epíscopi}
\switchcolumn
\EN
\LessonTitle{Homily by St. Augustine, Bishop (of Hippo.)}
\switchcolumn*

\LA
\Source{Tract. 24 in Joannem}
\switchcolumn
\EN
\Source{24th Tract on John}
\switchcolumn*

\LA
\noindent Mirácula, quæ fecit Dóminus noster Jesus Christus, sunt quidem divína ópera, et ad intellegéndum Deum de visibílibus ádmonent humánam mentem. Quia enim ille non est talis substántia, quæ vidéri óculis possit; et mirácula ejus, quibus totum mundum regit, universámque creatúram adminístrat, assiduitáte viluérunt, ita ut pene nemo dignétur atténdere ópera Dei mira et stupénda in quólibet séminis grano: secúndum ipsam suam misericórdiam, servávit sibi quædam, quæ fáceret opportúno témpore præter usitátum cursum ordinémque natúræ; ut non majóra, sed insólita vidéndo stupérent, quibus cotidiána vilúerant.\par
\switchcolumn
\EN
\noindent The miracles which our Lord Jesus Christ did were the very works of God, and they enlighten the mind of man by mean of things which are seen, that he may know more of God. God is Himself of such a Substance as eye cannot see, and the miracles, by the which He ruleth the whole world continually, and satisfieth the need of everything that He hath made, are by use become so common, that scarce any will vouchsafe to see that there are wonderful and amazing works of God in every grain of seed of grass. According to His mercy He kept some works to be done in their due season, but out of the common course and order of nature, that men might see them and be astonished, not because they are greater, but because they are rarer than those which they lightly esteem, since they see them day by day.\par
\switchcolumn*

\LA
\begin{Resp}\Rsign Spléndida facta est fácies Móysi, dum respíceret in eum Dóminus: \Star{} Vidéntes senióres claritátem vultus ejus, admirántes timuérunt valde.\end{Resp}
\begin{Resp}\Vsign Cumque descendísset de monte Sínai, portábat duas tábulas testimónii, ignórans quod cornúta esset fácies ejus ex consórtio sermónis Dei.\end{Resp}
\begin{Resp}\Rsign Vidéntes senióres claritátem vultus ejus, admirántes timuérunt valde.\end{Resp}
\switchcolumn
\EN
\begin{Resp}\Rsign After that the Lord had looked upon him, the face of Moses shone. \Star{} And when the elders saw that his face shone, they marvelled and were sore afraid.\end{Resp}
\begin{Resp}\Vsign And when he came down from Mount Sinai with the two tables of testimony in his hand, he wist not that his face shone, because that God had spoken with him.\end{Resp}
\begin{Resp}\Rsign And when the elders saw that his face shone, they marvelled and were sore afraid.\end{Resp}
\switchcolumn*

\LA
\LessonHead{Lectio VIII}
\switchcolumn
\EN
\LessonHead{Lesson VIII}
\switchcolumn*

\LA
\noindent Majus enim miráculum est gubernátio totíus mundi, quam saturátio quinque míllium hóminum de quinque pánibus. Et tamen hoc nemo mirátur: illud mirántur hómines, non quia majus est, sed quia rarum est. Quis enim et nunc pascit univérsum mundum, nisi ille, qui de paucis granis ségetes creat? Fecit ergo quo modo Deus. Unde enim multíplicat de paucis granis ségetes, inde in mánibus suis multiplicávit quinque panes: potéstas enim erat in mánibus Christi. Panes autem illi quinque, quasi sémina erant, non quidem terræ mandáta, sed ab eo, qui terram fecit, multiplicáta.\par
\switchcolumn
\EN
\noindent Or it is a greater miracle to govern the whole universe, than to satisfy five thousand men with five loaves of bread; and yet no man marvelleth at it. At the feeding of the five thousand, men marvel, not because it is a greater miracle than the other, but because it is rarer. For Who is He Who now feedeth the whole world, but He Who, from a little grain that is sown, maketh the fulness of the harvest? God worketh in both cases in one and the same manner. He Who of the sowing maketh to come the harvest, is He Who of the five barley loaves in His Hands made bread to feed five thousand men; for Christ's are the Hands which are able to do both the one and the other. He Who multiplieth the grains of corn multiplied the loaves, only not by committing them to the earth whereof He is the Maker.\par
\switchcolumn*

\LA
\begin{Resp}\Rsign Ecce mitto Angelum meum, qui præcédat te, et custódiat semper: \Star{} Obsérva et audi vocem meam, et inimícus ero inimícis tuis, et affligéntes te afflígam: et præcédet te Angelus meus.\end{Resp}
\begin{Resp}\Vsign Israël, si me audíeris, non erit in te deus recens, neque adorábis deum aliénum: ego enim Dóminus.\end{Resp}
\begin{Resp}\Rsign Obsérva et audi vocem meam, et inimícus ero inimícis tuis, et affligéntes te afflígam: et præcédet te Angelus meus.\end{Resp}
\switchcolumn
\EN
\begin{Resp}\Rsign Behold, I send My Angel before thee, to keep thee. \Star{} Beware, and obey My voice; then I will be an enemy unto thine enemies, and an adversary unto thine adversaries; for Mine Angel shall go before thee.\end{Resp}
\begin{Resp}\Vsign O Israel, if thou wilt hearken unto Me, there shall no strange god be in thee, neither shalt thou worship any strange god, for I am the Lord.\end{Resp}
\begin{Resp}\Rsign Beware, and obey My voice; then I will be an enemy unto thine enemies, and an adversary unto thine adversaries; for Mine Angel shall go before thee.\end{Resp}
\switchcolumn*

\LA
\LessonHead{Lectio IX}
\switchcolumn
\EN
\LessonHead{Lesson IX}
\switchcolumn*

\LA
\noindent Hoc ergo admótum est sénsibus, quo erigerétur mens: et exhíbitum óculis, ubi exercerétur intelléctus: ut invisíbilem Deum per visibília ópera mirarémur, et erécti ad fidem, et purgáti per fidem, étiam ipsum invisíbilem vidére cuperémus, quem de rebus visibílibus invisíbilem noscerémus. Nec tamen súfficit hæc intuéri in miráculis Christi. Interrogémus ipsa mirácula, quid nobis loquántur de Christo: habent enim, si intelligántur, linguam suam. Nam quia ipse Christus Verbum Dei est: étiam factum Verbi, verbum nobis est.\par
\switchcolumn
\EN
\noindent This miracle, then, is brought to bear upon our bodies, that our souls may thereby be quickened; shown to our eyes, to give food to our understanding; that, through His works which we see, we may marvel at that God Whom we cannot see, and, being roused up to believe, and purified by believing, we may long to see Him, yea, may know by things which are seen Him Who is Unseen. Nor yet sufficeth it for us to see only this meaning in Christ's miracles. Let us ask of the miracles themselves what they have to tell us concerning Christ for, soothly, they have a tongue of their own, if only we will understand it. For, because Christ is the Word of God, therefore the work of the Word is a Word for us.\par
\switchcolumn*

\LA
\begin{Resp}\Rsign Atténdite, pópule meus, legem meam: \Star{} Inclináte aurem vestram in verba oris mei.\end{Resp}
\begin{Resp}\Vsign Apériam in parábolis os meum: loquar propositiónes ab inítio sǽculi.\end{Resp}
\begin{Resp}\Rsign Inclináte aurem vestram in verba oris mei.\end{Resp}
\begin{Resp}\Vsign Glória Patri, et Fílio, \Star{} et Spirítui Sancto.\end{Resp}
\begin{Resp}\Rsign Inclináte aurem vestram in verba oris mei.\end{Resp}
\switchcolumn
\EN
\begin{Resp}\Rsign Give ear, O My people, to My law; \Star{} Incline your ears to the words of My mouth.\end{Resp}
\begin{Resp}\Vsign I will open My mouth in parables; I will utter dark sayings of old.\end{Resp}
\begin{Resp}\Rsign Incline your ears to the words of My mouth.\end{Resp}
\begin{Resp}\Vsign Glory be to the Father, and to the Son, \Star{} and to the Holy Ghost.\end{Resp}
\begin{Resp}\Rsign Incline your ears to the words of My mouth.\end{Resp}
\switchcolumn*

\end{paracol}

\Office{Feria Secunda infra Hebdomadam IV in Quadragesima}{Monday of the Fourth Week of Lent}{Feria major}
\addcontentsline{toc}{subsection}{Feria Secunda infra Hebdomadam IV in Quadragesima\enspace\textit{Monday of the Fourth Week of Lent}}
\begin{paracol}{2}
\LA
\LessonHead{Lectio I}
\switchcolumn
\EN
\LessonHead{Lesson I}
\switchcolumn*

\LA
\LessonTitle{Léctio sancti Evangélii secúndum Joánnem}
\switchcolumn
\EN
\LessonTitle{Continuation of the Holy Gospel according to John}
\switchcolumn*

\LA
\Cite{Joannes 2:13-25}
\switchcolumn
\EN
\Cite{John 2:13-25}
\switchcolumn*

\LA
\noindent In illo témpore: Prope erat Pascha Judæórum, et ascéndit Jesus Jerosólymam: et invénit in templo vendéntes boves, et oves, et colúmbas. Et réliqua.\par
\switchcolumn
\EN
\noindent In that time, the pasch of the Jews was at hand, and Jesus went up to Jerusalem. And he found in the temple them that sold oxen and sheep and doves. And so on.\par
\switchcolumn*

\LA
\LessonTitle{Homilía sancti Augustíni Epíscopi}
\switchcolumn
\EN
\LessonTitle{Homily by St. Augustine, Bishop (of Hippo.)}
\switchcolumn*

\LA
\Source{Tract. 10 in Joannem, post init.}
\switchcolumn
\EN
\Source{10th Tract on John.}
\switchcolumn*

\LA
\noindent Quid audívimus, fratres? Ecce templum illud figúra adhuc erat, et ejécit inde Dóminus omnes qui sua quærébant, qui ad núndinas vénerant. Et quæ ibi vendébant illi? Quæ opus habébant hómines in sacrifíciis illíus témporis. Novit enim cáritas vestra, quod sacrifícia illi pópulo pro ejus carnalitáte, et corde adhuc lapídeo, tália data sunt, quibus tenerétur, ne in idóla deflúeret: et immolábant ibi sacrifícia, boves, oves et colúmbas. Nostis, quia legístis.\par
\switchcolumn
\EN
\noindent What hear we now, my brethren? Behold, that temple was still but a figure, and the Lord drove out therefrom all them that sought their own, even them that were come to deal in merchandise. And what was it that they sold there? Only such things as were needful to men for the sacrifices that then were. For your love knoweth that, because of that people's carnal-mindedness and the stoniness of their heart, there were commanded unto them such sacrifices as these, thereby to hold them back from idolatry and there, according, they offered up oxen, and sheep, and doves. This ye have read, and know.\par
\switchcolumn*

\LA
\begin{Resp}\Rsign Vos, qui transitúri estis Jordánem, ædificáte altáre Dómino \Star{} De lapídibus, quos ferrum non tétigit: et offérte super illud holocáusta, et hóstias pacíficas Deo vestro.\end{Resp}
\begin{Resp}\Vsign Cumque intravéritis terram, quam Dóminus datúrus est vobis, ædificáte ibi altáre Dómino.\end{Resp}
\begin{Resp}\Rsign De lapídibus, quos ferrum non tétigit: et offérte super illud holocáusta, et hóstias pacíficas Deo vestro.\end{Resp}
\switchcolumn
\EN
\begin{Resp}\Rsign When ye be gone over Jordan, there shall ye build an altar unto the Lord \Star{} Of whole stones; ye shall not lift up any iron tool upon them; and ye shall offer burnt-offerings thereon, and peace-offerings, unto your God.\end{Resp}
\begin{Resp}\Vsign When ye shall pass over Jordan unto the land which the Lord giveth you, there shall ye build an altar unto the Lord.\end{Resp}
\begin{Resp}\Rsign Of whole stones; ye shall not lift up any iron tool upon them; and ye shall offer burnt-offerings thereon, and peace-offerings, unto your God.\end{Resp}
\switchcolumn*

\LA
\LessonHead{Lectio II}
\switchcolumn
\EN
\LessonHead{Lesson II}
\switchcolumn*

\LA
\noindent Non ergo magnum peccátum, si hoc vendébant in templo, quod emebátur ut offerétur in templo: et tamen ejécit inde illos. Quid si ibi ebriósos inveníret, quid fáceret Dóminus, si vendéntes ea quæ lícita sunt, et contra justítiam non sunt (quæ enim honéste emúntur, non illícite vendúntur) éxpulit tamen, et non est passus domum oratiónis fíeri domum negotiatiónis?\par
\switchcolumn
\EN
\noindent It was no great sin, therefore, if they sold in the temple that which was bought to be offered in the temple and yet He drove them out. If, then, the Lord drove out of His temple them which sold such things as are lawful and right (for to buy and sell is lawful, if only it be done honestly,) and suffered not the house of prayer to be made an house of merchandise, what would He have done if He had found there men drunken?\par
\switchcolumn*

\LA
\begin{Resp}\Rsign Audi, Israël, præcépta Dómini, et ea in corde tuo quasi in libro scribe: \Star{} Et dabo tibi terram fluéntem lac et mel.\end{Resp}
\begin{Resp}\Vsign Obsérva ígitur, et audi vocem meam: et inimícus ero inimícis tuis.\end{Resp}
\begin{Resp}\Rsign Et dabo tibi terram fluéntem lac et mel.\end{Resp}
\switchcolumn
\EN
\begin{Resp}\Rsign Hear, O Israel, the law of the Lord, and write it in thine heart as in a book \Star{} And I will give unto thee a land flowing with milk and honey.\end{Resp}
\begin{Resp}\Vsign Take heed therefore, and hearken unto My voice and I will be an enemy unto thine enemies.\end{Resp}
\begin{Resp}\Rsign And I will give unto thee a land flowing with milk and honey.\end{Resp}
\switchcolumn*

\LA
\LessonHead{Lectio III}
\switchcolumn
\EN
\LessonHead{Lesson III}
\switchcolumn*

\LA
\noindent Si negotiatiónis domus non debet fíeri domus Dei, potatiónis debet fíeri? Nos autem quando ista dícimus, strident déntibus suis advérsus nos: et consolátur nos Psalmus, quem audístis: Stridérunt in me déntibus suis. Nóvimus et nos audíre unde curémur: etsi ingeminántur flagélla Christo, quia flagellátur sermo ipsíus. Congregáta sunt, inquit, in me flagélla, et nesciébant. Flagellátus est flagéllis Judæórum: flagellátur blasphémiis falsórum Christianórum: multíplicant flagélla Dómino Deo suo, et nésciunt. Faciámus nos, quantum ipse ádjuvat. Ego autem, cum mihi molésti essent, induébam me cilício, et humiliábam in jejúnio ánimam meam.\par
\switchcolumn
\EN
\noindent If the house of God must not be a house of merchandise, must it be an house to drink in? And yet, when we say this, men gnash upon us with their teeth. But we find consolation in remembering that so far we are even as the Psalmist, who saith: They gnashed upon me with their teeth. (Ps. xxxiv. 16.) Yea, we have also learnt to listen to words that heal us, though, of a verity, the lashes that are made at His word are really made at Christ. Lashes, saith He, were heaped upon Me; and they knew not what they did. He was lashed by the scourges of the Jews, and He is lashed still by the blasphemies of false Christians; they heap lashes upon the Lord their God; and know not what they do. As for us, we will do that which He hath holpen us to do; But as for me, when they troubled me, my clothing was sackcloth, and I humbled my soul with fasting\par
\switchcolumn*

\LA
\begin{Resp}\Rsign Sicut fui cum Moyse ita ero tecum, dicit Dóminus: \Star{} Confortáre, et esto robústus: introdúces pópulum meum ad terram lacte et melle manántem.\end{Resp}
\begin{Resp}\Vsign Noli timére, quóniam tecum sum: ad quæcúmque perréxeris, non dimíttam te, neque derelínquam.\end{Resp}
\begin{Resp}\Rsign Confortáre, et esto robústus: introdúces pópulum meum ad terram lacte et melle manántem.\end{Resp}
\begin{Resp}\Vsign Glória Patri, et Fílio, \Star{} et Spirítui Sancto.\end{Resp}
\begin{Resp}\Rsign Confortáre, et esto robústus: introdúces pópulum meum ad terram lacte et melle manántem.\end{Resp}
\switchcolumn
\EN
\begin{Resp}\Rsign As I was with Moses, so I will be with thee, saith the Lord. \Star{} Be strong and of a good courage, and thou shalt bring My people into a land flowing with milk and honey.\end{Resp}
\begin{Resp}\Vsign Fear not, for I am with thee whithersoever thou goest I will not fail thee, nor forsake thee.\end{Resp}
\begin{Resp}\Rsign Be strong and of a good courage, and thou shalt bring My people into a land flowing with milk and honey.\end{Resp}
\begin{Resp}\Vsign Glory be to the Father, and to the Son, \Star{} and to the Holy Ghost.\end{Resp}
\begin{Resp}\Rsign Be strong and of a good courage, and thou shalt bring My people into a land flowing with milk and honey.\end{Resp}
\switchcolumn*

\end{paracol}

\Office{Feria Tertia infra Hebdomadam IV in Quadragesima}{Tuesday of the Fourth Week of Lent}{Feria major}
\addcontentsline{toc}{subsection}{Feria Tertia infra Hebdomadam IV in Quadragesima\enspace\textit{Tuesday of the Fourth Week of Lent}}
\begin{paracol}{2}
\LA
\LessonHead{Lectio I}
\switchcolumn
\EN
\LessonHead{Lesson I}
\switchcolumn*

\LA
\LessonTitle{Léctio sancti Evangélii secúndum Joánnem}
\switchcolumn
\EN
\LessonTitle{Continuation of the Holy Gospel according to John}
\switchcolumn*

\LA
\Cite{Joannes 7:14-31}
\switchcolumn
\EN
\Cite{John 7:14-31}
\switchcolumn*

\LA
\noindent In illo témpore: Jam die festo mediánte, ascéndit Jesus in templum, et docébat. Et mirabántur Judǽi. Et réliqua.\par
\switchcolumn
\EN
\noindent In that time, about the midst of the feast, Jesus went up into the temple, and taught. And the Jews wondered. And so on.\par
\switchcolumn*

\LA
\LessonTitle{Homilía sancti Augustíni Epíscopi}
\switchcolumn
\EN
\LessonTitle{Homily by St. Augustine, Bishop (of Hippo.)}
\switchcolumn*

\LA
\Source{Tract. 29 in Joannem, sub init.}
\switchcolumn
\EN
\Source{29th Tract on John}
\switchcolumn*

\LA
\noindent Ille qui latébat, docébat, et palam loquebátur, et non tenebátur. Illud enim ut latéret, erat causa exémpli, hoc potestátis. Sed cum docéret, mirabántur Judǽi. Omnes quidem, quantum árbitror, mirabántur, sed non omnes convertebántur. Et unde admirátio? Quia multi nóverant ubi natus, quemádmodum fúerit educátus. Nunquam eum víderant lítteras discéntem: audiébant autem de lege disputántem, legis testimónia proferéntem, quæ nemo posset proférre, nisi legísset, nemo légeret, nisi lítteras didicísset: et ídeo mirabántur. Eórum autem admirátio, magístro facta est insinuándæ áltius veritátis occásio.\par
\switchcolumn
\EN
\noindent He Who had gone up unto the Feast, not openly, but as it were in secret, the Same taught, and spake openly, and no man laid hands upon Him. That He had hid Himself, was for example's sake; that He manifested Himself, was to show His power. And when He taught, the Jews marvelled. As seemeth to my mind, they all marvelled, but were not all converted. And wherefore marvelled they? Because many of them knew where He was born, and how He had been brought up. They had never seen Him learn letters; but they heard Him dispute concerning the law, and allege the testimony of the same, as no man could do who had not read it; and no man can read unless he learn; and therefore they marvelled. But their marvelling was unto the Teacher an occasion for the revealing of higher truth.\par
\switchcolumn*

\LA
\begin{Resp}\Rsign Quid me quǽritis interfícere, hóminem qui vera locútus sum vobis? \Star{} Si male locútus sum, testimónium pérhibe de malo: si autem bene, cur me cædis?\end{Resp}
\begin{Resp}\Vsign Multa bona ópera operátus sum vobis: propter quod opus vultis me occídere?\end{Resp}
\begin{Resp}\Rsign Si male locútus sum, testimónium pérhibe de malo: si autem bene, cur me cædis?\end{Resp}
\switchcolumn
\EN
\begin{Resp}\Rsign Why go ye about to kill Me, a Man That hath told you the truth? \Star{} If I have spoken evil, bear witness of the evil; but if well, why smitest thou Me?\end{Resp}
\begin{Resp}\Vsign Many good works have I wrought among you; for which of those works go ye about to kill Me?\end{Resp}
\begin{Resp}\Rsign If I have spoken evil, bear witness of the evil; but if well, why smitest thou Me?\end{Resp}
\switchcolumn*

\LA
\LessonHead{Lectio II}
\switchcolumn
\EN
\LessonHead{Lesson II}
\switchcolumn*

\LA
\noindent Ex eórum quippe admiratióne et verbis, dixit Dóminus profúndum áliquid, et diligéntius ínspici et discúti dignum. Quid ergo Dóminus respóndit eis, admirántibus quómodo sciret lítteras, quas non didícerat? Mea, inquit, doctrína non est mea, sed ejus qui misit me. Hæc est profúnditas prima: vidétur enim paucis verbis quasi contrária locútus. Non enim ait: Ista doctrína non est mea: sed, Mea doctrína non est mea. Si non tua, quómodo tua? si tua, quómodo non tua? Tu enim dicis utrúmque: et mea doctrína, et non mea.\par
\switchcolumn
\EN
\noindent For when they marvelled and whispered, the Lord said a certain deep thing, yea, a thing worthy of very careful thought and discussion. And what was this thing which the Lord gave for an answer to such as marvelled that He knew letters, having never learned? Jesus answered them and said: My doctrine is not Mine, but His That sent Me. Here is the first depth, for He seemeth in these few words to enunciate a contradiction. He saith not: This doctrine is not Mine but: My doctrine is not Mine. O how is it thine? If it be thine, wherefore sayest Thou that it is not thine? For Thou sayest: My doctrine is not Mine.\par
\switchcolumn*

\LA
\begin{Resp}\Rsign Addúxi vos per desértum quadragínta annis ego Dóminus, et non sunt attríta vestiménta vestra: \Star{} Manna de cælo plui vobis, et oblíti estis me, dicit Dóminus.\end{Resp}
\begin{Resp}\Vsign Ego addúxi vos de terra Ægýpti, et de domo servitútis liberávi vos.\end{Resp}
\begin{Resp}\Rsign Manna de cælo plui vobis, et oblíti estis me, dicit Dóminus.\end{Resp}
\switchcolumn
\EN
\begin{Resp}\Rsign I, even I, the Lord, have led you forty years in the wilderness, and your clothes are not waxen old upon you. \Star{} I rained down manna upon you from heaven, and ye have forgotten Me, saith the Lord.\end{Resp}
\begin{Resp}\Vsign I led you forth out of the land of Egypt, and delivered you from the house of bondage.\end{Resp}
\begin{Resp}\Rsign I rained down manna upon you from heaven, and ye have forgotten Me, saith the Lord.\end{Resp}
\switchcolumn*

\LA
\LessonHead{Lectio III}
\switchcolumn
\EN
\LessonHead{Lesson III}
\switchcolumn*

\LA
\noindent Si ergo intueámur diligénter quod ipse in exórdio dicit sanctus Evangelísta: In princípio erat Verbum, et Verbum erat apud Deum, et Deus erat Verbum: inde pendet hujus solútio quæstiónis. Quæ est doctrína Patris, nisi Verbum Patris? Ipse ergo Christus doctrína Patris, si Verbum Patris. Sed quia Verbum, non potest esse nullíus, sed alicújus: et suam doctrínam dixit seípsum, et non suam, quia Patris est Verbum. Quid enim tam tuum quam tu? Et quid tam non tuum quam tu, si alicújus est, quod es?\par
\switchcolumn
\EN
\noindent Let us then carefully regard what this same holy Evangelist saith in the beginning of his Gospel, and we shall find there wherewith to loose the knot of this difficulty. There it is written: In the beginning was the Word, and the Word was with God, and the Word was God. (i. 1.) What is the doctrine of the Father but the Word of the Father? If Christ therefore be the Word of the Father, He is the doctrine of the Father. But a Word cannot be of no one, but must needs, if it be a Word, have some one whose word it is. Christ therefore saith that His doctrine is Himself, and therefore not His, forasmuch as He is the Word of the Father. And what hast thou that is so much thine own as thy self? Or what is there that is so little thine own as thyself, if that which thou art is another's?\par
\switchcolumn*

\LA
\begin{Resp}\Rsign Móyses fámulus Dei jejunávit quadragínta diébus et quadragínta nóctibus: \Star{} Ut legem Dómini mererétur accípere.\end{Resp}
\begin{Resp}\Vsign Ascéndens Móyses in montem Sínai ad Dóminum, fuit ibi quadragínta diébus et quadragínta nóctibus.\end{Resp}
\begin{Resp}\Rsign Ut legem Dómini mererétur accípere.\end{Resp}
\begin{Resp}\Vsign Glória Patri, et Fílio, \Star{} et Spirítui Sancto.\end{Resp}
\begin{Resp}\Rsign Ut legem Dómini mererétur accípere.\end{Resp}
\switchcolumn
\EN
\begin{Resp}\Rsign Moses, the servant of God, fasted forty days and forty nights \Star{} To make him meet to receive the Law of God.\end{Resp}
\begin{Resp}\Vsign Moses went up unto the Lord into Mount Sinai, and he was in the Mount forty days and forty nights.\end{Resp}
\begin{Resp}\Rsign To make him meet to receive the Law of God.\end{Resp}
\begin{Resp}\Vsign Glory be to the Father, and to the Son, \Star{} and to the Holy Ghost.\end{Resp}
\begin{Resp}\Rsign To make him meet to receive the Law of God.\end{Resp}
\switchcolumn*

\end{paracol}

\Office{Feria Quarta infra Hebdomadam IV in Quadragesima}{Wednesday of the Fourth Week of Lent}{Feria major}
\addcontentsline{toc}{subsection}{Feria Quarta infra Hebdomadam IV in Quadragesima\enspace\textit{Wednesday of the Fourth Week of Lent}}
\begin{paracol}{2}
\LA
\LessonHead{Lectio I}
\switchcolumn
\EN
\LessonHead{Lesson I}
\switchcolumn*

\LA
\LessonTitle{Léctio sancti Evangélii secúndum Joánnem}
\switchcolumn
\EN
\LessonTitle{Continuation of the Holy Gospel according to John}
\switchcolumn*

\LA
\Cite{Joannes 9:1-38}
\switchcolumn
\EN
\Cite{John 9:1-38}
\switchcolumn*

\LA
\noindent In illo témpore: Prætériens Jesus, vidit hóminem cæcum a nativitáte: et interrogavérunt eum discípuli ejus: Rabbi, quis peccávit, hic, aut paréntes ejus, ut cæcus nascerétur? Et réliqua.\par
\switchcolumn
\EN
\noindent At that time, Jesus, passing by, saw a man, who was blind from his birth. And so on.\par
\switchcolumn*

\LA
\LessonTitle{Homilía sancti Augustíni Epíscopi}
\switchcolumn
\EN
\LessonTitle{Homily by St. Augustine, Bishop (of Hippo.)}
\switchcolumn*

\LA
\Source{Tract. 44 in Joannem, circa initium}
\switchcolumn
\EN
\Source{44th Tract on John}
\switchcolumn*

\LA
\noindent Ea quæ fecit Dóminus noster Jesus Christus, stupénda atque miránda, et ópera, et verba sunt: ópera, quia facta sunt: verba, quia signa sunt. Si ergo quid signíficet hoc quod factum est, cogitémus: genus humánum est iste cæcus. Hæc enim cǽcitas cóntigit in primo hómine per peccátum, de quo omnes oríginem dúximus, non solum mortis, sed étiam iniquitátis. Si enim cǽcitas est infidélitas, et illuminátio fides: quem fidélem, quando venit Christus, invénit? Quandóquidem Apóstolus natus in gente prophetárum dicit: Fúimus et nos aliquándo natúra fílii iræ, sicut et céteri. Si fílii iræ, fílii vindíctæ, fílii pœnæ, fílii gehénnæ: quómodo natúra, nisi quia peccánte primo hómine vítium pro natúra inolévit? Si vítium pro natúra inolévit, secúndum mentem omnis homo cæcus natus est.\par
\switchcolumn
\EN
\noindent Dread and wondrous are all the things which our Lord Jesus Christ did, both His works and His words; the works, because He wrought them; the words, because they are deep. If, therefore, we consider the meaning of this work of His, we see that that man which was blind from his birth was a figure of mankind. This spiritual blindness was the consequence of the sin of the first man, from whom we all inherit by birth, not death only, but depravity also. For if blindness be unbelief, and faith, light, whom, when Christ came, did He find faithful? May, the Apostle who had himself been born of the race of which the Prophets came, saith: We also were by nature children of wrath, even as others. (Eph. ii. 3.) And if children of wrath, then children also of vengeance, children of damnation, children of hell. And wherefore so by nature, unless it were that the sin of the first man had made all his descendants to be born in sin, in that they partook of his nature? If, then, our nature bring sin with it, all men, according to the spirit, are born blind.\par
\switchcolumn*

\LA
\begin{Resp}\Rsign Spléndida facta est fácies Móysi, dum respíceret in eum Dóminus: \Star{} Vidéntes senióres claritátem vultus ejus, admirántes timuérunt valde.\end{Resp}
\begin{Resp}\Vsign Cumque descendísset de monte Sínai, portábat duas tábulas testimónii, ignórans quod cornúta esset fácies ejus ex consórtio sermónis Dei.\end{Resp}
\begin{Resp}\Rsign Vidéntes senióres claritátem vultus ejus, admirántes timuérunt valde.\end{Resp}
\switchcolumn
\EN
\begin{Resp}\Rsign After that the Lord had looked upon him, the face of Moses shone. \Star{} And when the elders saw that his face shone, they marvelled and were sore afraid.\end{Resp}
\begin{Resp}\Vsign And when he came down from Mount Sinai with the two tables of testimony in his hand, he wist not that his face shone, because that God had spoken with him.\end{Resp}
\begin{Resp}\Rsign And when the elders saw that his face shone, they marvelled and were sore afraid.\end{Resp}
\switchcolumn*

\LA
\LessonHead{Lectio II}
\switchcolumn
\EN
\LessonHead{Lesson II}
\switchcolumn*

\LA
\noindent Venit Dóminus: quid fecit? Magnum mystérium commendávit. Exspuit in terram, de salíva sua lutum fecit: quia Verbum caro factum est, et inúnxit óculos cæci. Inúnctus erat, et nondum vidébat. Misit illum ad piscínam, quæ vocátur Síloë. Pertínuit autem ad Evangelístam commendáre nobis nomen hujus piscínæ, et ait: Quod interpretátur Missus. Jam quis sit missus agnóscitis. Nisi enim ille fuísset missus, nemo nostrum esset ab iniquitáte dimíssus. Lavit ergo óculos in ea piscína, quæ interpretátur Missus; baptizátus est in Christo. Si ergo quando eum in seípso quodámmodo baptizávit, tunc illuminávit: quando inúnxit, fortásse catechúmenum fecit.\par
\switchcolumn
\EN
\noindent The Lord came; and what did He? He set before us a great mystery. Jesus spat on the ground, and made clay of the Spittle for the Word was made flesh. And He anointed the eyes of the blind man with the clay but yet that man saw not. He was anointed, indeed, but yet still he saw not. And He said unto him: Go, wash in the Pool of Siloam. Now, it was the duty of the Evangelist to impress upon us the name of this Pool, and therefore he saith Siloam, which is, by interpretation, Sent. Ye, my brethren, know Who is signified where it is written: (The sceptre shall not depart from Judah, nor a law-giver from his loins, until) He that shall be Sent cometh. (Gen. xlix. 10.) Yea, He it is, Who, if He had not been sent, we had never been sent loose out of the prison-house of sin. The blind man went his way therefore, and washed his eyes in that Pool, which is, by interpretation, Sent in other words, he was baptized in Christ. When, therefore, he had figuratively been baptized in Him Whom the Father hath Sent into the world, he came seeing. When he was anointed, he was perchance made a figure of a Catechumen.\par
\switchcolumn*

\LA
\begin{Resp}\Rsign Ecce mitto Angelum meum, qui præcédat te, et custódiat semper: \Star{} Obsérva et audi vocem meam, et inimícus ero inimícis tuis, et affligéntes te afflígam: et præcédet te Angelus meus.\end{Resp}
\begin{Resp}\Vsign Israël, si me audíeris, non erit in te deus recens, neque adorábis deum aliénum: ego enim Dóminus.\end{Resp}
\begin{Resp}\Rsign Obsérva et audi vocem meam, et inimícus ero inimícis tuis, et affligéntes te afflígam: et præcédet te Angelus meus.\end{Resp}
\switchcolumn
\EN
\begin{Resp}\Rsign Behold, I send My Angel before thee, to keep thee. \Star{} Beware, and obey My voice; then I will be an enemy unto thine enemies, and an adversary unto thine adversaries; for Mine Angel shall go before thee.\end{Resp}
\begin{Resp}\Vsign O Israel, if thou wilt hearken unto Me, there shall no strange god be in thee, neither shalt thou worship any strange god, for I am the Lord.\end{Resp}
\begin{Resp}\Rsign Beware, and obey My voice; then I will be an enemy unto thine enemies, and an adversary unto thine adversaries; for Mine Angel shall go before thee.\end{Resp}
\switchcolumn*

\LA
\LessonHead{Lectio III}
\switchcolumn
\EN
\LessonHead{Lesson III}
\switchcolumn*

\LA
\noindent Audístis grande mystérium. Intérroga hóminem: Christiánus es? Respóndet tibi: Non sum. Si pagánus es, aut Judǽus? Si autem díxerit, Non sum: adhuc quæris ab eo, Catechúmenus, an fidélis? Si respónderit tibi, Catechúmenus: inúnctus est, nondum lotus. Sed unde inúnctus? Quære, et respóndet. Quære ab illo, in quem credat? Eo ipso quo catechúmenus est, dicit: In Christum. Ecce modo loquor et fidélibus et catechúmenis. Quid dixi de sputo et luto? Quia Verbum caro factum est; hoc catechúmeni áudiunt: sed non eis súfficit ad quod inúncti sunt: festínent ad lavácrum, si lumen inquírunt.\par
\switchcolumn
\EN
\noindent We have heard this great mystery. Ask of a man: Art thou a Christian? He answereth thee: I am not. Then, if thou ask him: Art thou a pagan then, or a Jew? And he still saith unto thee Nay and thou say Art thou then a Catechumen, though not yet one of the faithful? and he saith Yea, a Catechumen then there thou seest a man anointed, but not yet washed. With what hath he been anointed? Ask of him, and he will tell thee. Ask of him in Whom he believeth, and, being a Catechumen, he will say: In Christ. But, behold, I speak before both Faithful and Catechumens. What said I touching the. Spittle and the clay? I said for 'the Word was made flesh.' This the Catechumens hear, but it is not enough for them to be anointed; they must make haste to the washing, if they would have their eyes opened.\par
\switchcolumn*

\LA
\begin{Resp}\Rsign Atténdite, pópule meus, legem meam: \Star{} Inclináte aurem vestram in verba oris mei.\end{Resp}
\begin{Resp}\Vsign Apériam in parábolis os meum: loquar propositiónes ab inítio sǽculi.\end{Resp}
\begin{Resp}\Rsign Inclináte aurem vestram in verba oris mei.\end{Resp}
\begin{Resp}\Vsign Glória Patri, et Fílio, \Star{} et Spirítui Sancto.\end{Resp}
\begin{Resp}\Rsign Inclináte aurem vestram in verba oris mei.\end{Resp}
\switchcolumn
\EN
\begin{Resp}\Rsign Give ear, O My people, to My law; \Star{} Incline your ears to the words of My mouth.\end{Resp}
\begin{Resp}\Vsign I will open My mouth in parables; I will utter dark sayings of old.\end{Resp}
\begin{Resp}\Rsign Incline your ears to the words of My mouth.\end{Resp}
\begin{Resp}\Vsign Glory be to the Father, and to the Son, \Star{} and to the Holy Ghost.\end{Resp}
\begin{Resp}\Rsign Incline your ears to the words of My mouth.\end{Resp}
\switchcolumn*

\end{paracol}

\Office{Feria Quinta infra Hebdomadam IV in Quadragesima}{Thursday of the Fourth Week of Lent}{Feria major}
\addcontentsline{toc}{subsection}{Feria Quinta infra Hebdomadam IV in Quadragesima\enspace\textit{Thursday of the Fourth Week of Lent}}
\begin{paracol}{2}
\LA
\LessonHead{Lectio I}
\switchcolumn
\EN
\LessonHead{Lesson I}
\switchcolumn*

\LA
\LessonTitle{Léctio sancti Evangélii secúndum Lucam}
\switchcolumn
\EN
\LessonTitle{Continuation of the Holy Gospel according to Luke}
\switchcolumn*

\LA
\Cite{Luc 7:11-16}
\switchcolumn
\EN
\Cite{Luke 7:11-16}
\switchcolumn*

\LA
\noindent In illo témpore: Ibat Jesus in civitátem, quæ vocátur Naim: et ibant cum eo discípuli ejus, et turba copiósa. Et réliqua.\par
\switchcolumn
\EN
\noindent At that time, Jesus went into a city that is called Naim; and there went with him his disciples, and a great multitude. And so on.\par
\switchcolumn*

\LA
\LessonTitle{Homilía sancti Ambrósii Epíscopi}
\switchcolumn
\EN
\LessonTitle{Homily by St. Ambrose, Bishop (of Milan.)}
\switchcolumn*

\LA
\Source{Lib. 5 Comment. in Lucæ cap. 7, post initium}
\switchcolumn
\EN
\Source{Bk. v Comm. on Luke vii}
\switchcolumn*

\LA
\noindent Et hic locus ad utrámque redúndat grátiam; et ut cito flecti divínam misericórdiam matris víduæ lamentatióne credámus, ejus præcípue, quæ únici fílii vel labóre, vel morte frangátur; cui tamen víduæ gravitátis méritum exsequiárum turba concíliet: et ut hanc víduam populórum turba septam, plus vidéri esse quam féminam, quæ resurrectiónem únici et adolescéntis fílii suis lácrimis merúerit impetráre: eo quod sancta Ecclésia pópulum juniórem a pompa fúneris atque a suprémis sepúlcri, suárum révocet ad vitam contemplatióne lacrimárum: quæ flere prohibétur eum, cui resurréctio debebátur.\par
\switchcolumn
\EN
\noindent The history which we here read in the Holy Gospel hath for us specially two gracious lessons, the one from the literal, the other from the mystic interpretation thereof. According to the letter then, we see how quickly the compassion of God was aroused by the sorrow of this mother, who was a widow, a widow broken down by nursing her only son, or by the bitterness of her grief for his death. She was a widow also whose worshipful conversation is borne witness to by this, that, much people of the city was with her. Mystically however, this widow encompassed by the multitude was something more than a poor woman whose tears won from the Lord the resurrection of her young and only son; for she is a type of our holy Mother the Church, who calleth back her young children to life from the pursuit of deathly vanities, and soul-slaying honours, by bidding them look on those tears which she sheddeth for such as they, and which it is unlawful for her to shed for them of whom she knoweth that they will rise again.\par
\switchcolumn*

\LA
\begin{Resp}\Rsign Locútus est Dóminus ad Móysen, dicens: Descénde in Ægýptum, et dic Pharaóni: \Star{} Ut dimíttat pópulum meum: indurátum est cor Pharaónis: non vult dimíttere pópulum meum, nisi in manu forti.\end{Resp}
\begin{Resp}\Vsign Clamor filiórum Israël venit ad me, vidíque afflictiónem eórum: sed veni, mittam te ad Pharaónem.\end{Resp}
\begin{Resp}\Rsign Ut dimíttat pópulum meum: indurátum est cor Pharaónis: non vult dimíttere pópulum meum, nisi in manu forti.\end{Resp}
\switchcolumn
\EN
\begin{Resp}\Rsign The Lord spake unto Moses, saying: Go down now into Egypt, and say unto Pharaoh: \Star{} Let My people go. And the heart of Pharaoh shall be hardened, that he will not let My people go but by a mighty hand.\end{Resp}
\begin{Resp}\Vsign The cry of the children of Israel is come unto Me, and I have seen their affliction: come now, therefore, and I will send thee unto Pharaoh, and thou shalt say unto him:\end{Resp}
\begin{Resp}\Rsign Let My people go. And the heart of Pharaoh shall be hardened, that he will not let My people go but by a mighty hand.\end{Resp}
\switchcolumn*

\LA
\LessonHead{Lectio II}
\switchcolumn
\EN
\LessonHead{Lesson II}
\switchcolumn*

\LA
\noindent Qui quidem mórtuus in lóculo materiálibus quátuor ad sepúlcrum ferebátur eleméntis, sed spem resurgéndi habébat, quia ferebátur in ligno. Quod etsi nobis ante non próderat, tamen posteáquam Jesus id tétigit, profícere cœpit ad vitam: ut esset indício, salútem pópulo per crucis patíbulum refundéndam. Audíto ígitur Dei verbo, stetérunt acérbi illi fúneris portitóres, qui corpus humánum letháli fluxu natúræ materiális urgébant. Quid enim áliud, nisi quasi in quodam féretro, hoc est, suprémi fúneris instruménto, jacémus exánimes, cum vel ignis immódicæ cupiditátis exǽstuat, vel frígidus humor exúndat, vel pigra quadam terréni córporis habitúdine vigor hebetátur animórum; vel concréta noster spíritus labe, puræ lucis vácuus mentem alit? Hi sunt nostri fúneris portitóres.\par
\switchcolumn
\EN
\noindent This man, then, being dead, was carried out on a bier to the grave by four bearers, even as the sinner is borne to destruction by the four elements of which he is composed. But there was hope in his latter end, from this, that that whereon he was carried was of wood, and wood, albeit it had profited us little before, is become everything to us now since Jesus touched it, being a figure of that gibbet, the Cross, which was made thereof, and wherefrom salvation floweth unto all people. When, therefore, the horrid bearers of the corpse heard the commandment of God, they stood still, and carried no farther him who was dead through the fatal course of a material nature. And is not our case even as that of the widow's son, when we lie, as it were, lifeless, in our spiritual coffin, that is, in the last bed of our soul's death, consumed by the fever of unbridled lust, or frozen by cold-heartedness, or with our whole manliness sapped by some degrading habit of this earthly body, or starved by a spiritual lockjaw that shutteth our mouth to the bright food of our soul? These, and such as these, are they which carry us out to burial.\par
\switchcolumn*

\LA
\begin{Resp}\Rsign Stetit Móyses coram Pharaóne, et dixit: Hæc dicit Dóminus: \Star{} Dimítte pópulum meum, ut sacríficet mihi in desérto.\end{Resp}
\begin{Resp}\Vsign Dóminus Deus Hebræórum misit me ad te, dicens.\end{Resp}
\begin{Resp}\Rsign Dimítte pópulum meum, ut sacríficet mihi in desérto.\end{Resp}
\switchcolumn
\EN
\begin{Resp}\Rsign Moses stood before Pharaoh, and said: Thus saith the Lord: \Star{} Let My people go, that they may hold a feast unto Me in the wilderness.\end{Resp}
\begin{Resp}\Vsign The Lord God of the Hebrews hath sent me unto thee, saying:\end{Resp}
\begin{Resp}\Rsign Let My people go, that they may hold a feast unto Me in the wilderness.\end{Resp}
\switchcolumn*

\LA
\LessonHead{Lectio III}
\switchcolumn
\EN
\LessonHead{Lesson III}
\switchcolumn*

\LA
\noindent Sed quamvis supréma mortis spem vitæ omnis aboléverint, et túmulo próxima córpora jáceant defunctórum: verbo tamen Dei jam mórtua resúrgunt cadávera: vox redit, rédditur fílius matri, revocátur a túmulo, erípitur a sepúlcro. Quis iste est túmulus tuus, nisi mali mores? Túmulus tuus perfídia est: sepúlcrum tuum guttur est. Sepúlcrum enim patens, est guttur eórum, unde verba mórtua proferúntur. Ab hoc sepúlcro te líberat Christus: ab hoc túmulo surges, si áudias verbum Dei. Et si grave peccátum est, quod pœniténtiæ lácrimis ipse laváre non possis; fleat pro te mater Ecclésia, quæ pro síngulis tamquam pro únicis fíliis vídua mater intérvenit. Compátitur enim quodam spiritáli dolóre natúræ, cum suos líberos lethálibus vítiis ad mortem cernit urgéri.\par
\switchcolumn
\EN
\noindent But even at the last hour, when the hope of life hath been utterly extinguished, and the bodies of the dead are lying by the side of the grave, by the word of God those carcases live again, yea, arise and speak. Then doth Jesus deliver the son to his mother, for Jesus calleth him out of the grave, and delivereth him from death. O, what is the grave of the soul but a bad life? Sinner! thy grave is unbelief, and thy throat is a sepulchre! Even so is it written: Their throat is an open sepulchre, (Ps. v. ii,) whereout breathe their pestilential words. Lo! Christ maketh thee free from that grave! If only thou wilt hear the word of God, thou shalt yet arise from that sepulchre! Yea, though thy sin be exceeding weighty, so that the tears of thine own sorrow cannot wash it away, let thy Mother the Church weep for thee, that longing Mother who weepeth for every one of her children as though he were the only son of his mother, and she was a widow. Believe me, her spiritual anguish is keen like the anguish of nature, when she seeth her children dead in sin, and carried out to be buried for ever.\par
\switchcolumn*

\LA
\begin{Resp}\Rsign Cantémus Dómino: glorióse enim honorificátus est, equum et ascensórem projécit in mare: \Star{} Adjútor et protéctor factus est mihi Dóminus in salútem.\end{Resp}
\begin{Resp}\Vsign Dóminus quasi vir pugnátor, Omnípotens nomen ejus.\end{Resp}
\begin{Resp}\Rsign Adjútor et protéctor factus est mihi Dóminus in salútem.\end{Resp}
\begin{Resp}\Vsign Glória Patri, et Fílio, \Star{} et Spirítui Sancto.\end{Resp}
\begin{Resp}\Rsign Adjútor et protéctor factus est mihi Dóminus in salútem.\end{Resp}
\switchcolumn
\EN
\begin{Resp}\Rsign Let us sing unto the Lord, for He hath triumphed gloriously; the horse and his rider hath He thrown into the sea. \Star{} The Lord is my strength and song, and He is become my salvation.\end{Resp}
\begin{Resp}\Vsign The Lord is a man of war; Almighty is His Name.\end{Resp}
\begin{Resp}\Rsign The Lord is my strength and song, and He is become my salvation.\end{Resp}
\begin{Resp}\Vsign Glory be to the Father, and to the Son, \Star{} and to the Holy Ghost.\end{Resp}
\begin{Resp}\Rsign The Lord is my strength, and my song, and He is become my salvation.\end{Resp}
\switchcolumn*

\end{paracol}

\Office{Feria Sexta infra Hebdomadam IV in Quadragesima}{Friday of the Fourth Week of Lent}{Feria major}
\addcontentsline{toc}{subsection}{Feria Sexta infra Hebdomadam IV in Quadragesima\enspace\textit{Friday of the Fourth Week of Lent}}
\begin{paracol}{2}
\LA
\LessonHead{Lectio I}
\switchcolumn
\EN
\LessonHead{Lesson I}
\switchcolumn*

\LA
\LessonTitle{Léctio sancti Evangélii secúndum Joánnem}
\switchcolumn
\EN
\LessonTitle{Continuation of the Holy Gospel according to John}
\switchcolumn*

\LA
\Cite{Joannes 11:1-45}
\switchcolumn
\EN
\Cite{John 11:1-45}
\switchcolumn*

\LA
\noindent In illo témpore: Erat quidam languens Lázarus a Bethánia, de castéllo Maríæ et Marthæ soróris ejus. Et réliqua.\par
\switchcolumn
\EN
\noindent At that time, there was a certain man sick, named Lazarus, of Bethania, of the town of Mary and Martha her sister. And so on.\par
\switchcolumn*

\LA
\LessonTitle{Homilía sancti Augustíni Epíscopi}
\switchcolumn
\EN
\LessonTitle{Homily by St. Augustine, Bishop (of Hippo.)}
\switchcolumn*

\LA
\Source{Tract. 49 in Joannem, post init.}
\switchcolumn
\EN
\Source{49th Tract on John}
\switchcolumn*

\LA
\noindent In superióri lectióne meminístis, quod Dóminus éxiit de mánibus eórum, qui lapidáre illum volúerant, et discéssit trans Jordánem, ubi Joánnes baptizábat. Ibi ergo Dómino constitúto, infirmabátur in Bethánia Lázarus: quod castéllum erat próximum Jerosólymis. María autem erat, quæ unxit Dóminum unguénto, et extérsit pedes ejus capíllis suis, cujus frater Lázarus infirmabátur. Misérunt ergo soróres ejus ad eum. Jam intellégimus quo misérunt, ubi erat Jesus: quóniam absens erat, trans Jordánem scílicet. Misérunt ad Dóminum, nuntiántes quod ægrotáret frater eárum, ut si dignarétur, veníret, et eum ab ægritúdine liberáret. Ille dístulit sanáre, ut posset resuscitáre.\par
\switchcolumn
\EN
\noindent Ye remember that in our last reading we learnt how that the Lord escaped out of the hands of them which took up stones to stone Him, and went away again beyond Jordan, into the place where John at first baptized. (John x. 31, 39, 40-) While, then, the Lord still tarried there, Lazarus was sick at Bethany, which was a town near to Jerusalem. It was that Mary which anointed the Lord with ointment, and wiped His Feet with her hair, whose brother Lazarus was sick. Therefore his sisters sent unto Him. We know already whither it was that they sent, for we know where Jesus was: He was gone away again beyond Jordan. His sisters sent unto Him, saying: Lord, behold, he whom Thou lovest is sick, in order that, if He so pleased, He might come and free him from his sickness. But Jesus healed not, that He might afterward quicken.\par
\switchcolumn*

\LA
\begin{Resp}\Rsign In mare viæ tuæ, et sémitæ tuæ in aquis multis: \Star{} Deduxísti sicut oves pópulum tuum in manu Móysi et Aaron.\end{Resp}
\begin{Resp}\Vsign Transtulísti illos per mare Rubrum, et transvexísti eos per aquam nímiam.\end{Resp}
\begin{Resp}\Rsign Deduxísti sicut oves pópulum tuum in manu Móysi et Aaron.\end{Resp}
\switchcolumn
\EN
\begin{Resp}\Rsign Thy way is in the sea, and thy paths in the great waters. \Star{} Thou led thy people like a flock, by the hand of Moses and Aaron.\end{Resp}
\begin{Resp}\Vsign Thou broughtest them through the Red Sea, and led them through much water.\end{Resp}
\begin{Resp}\Rsign Thou leddest thy people like a flock by the hand of Moses and Aaron.\end{Resp}
\switchcolumn*

\LA
\LessonHead{Lectio II}
\switchcolumn
\EN
\LessonHead{Lesson II}
\switchcolumn*

\LA
\noindent Quid ergo nuntiavérunt soróres ejus? Dómine, ecce quem amas, infirmátur. Non dixérunt: Veni: amánti enim tantúmmodo nuntiándum fuit. Non ausæ sunt dícere: Veni, et sana. Non ausæ sunt dícere: Ibi jube, et hic fiet. Cur enim non et istæ, si fides illíus centuriónis inde laudátur? Ait enim: Non sum dignus ut intres sub tectum meum; sed tantum dic verbo, et sanábitur puer meus. Nihil horum istæ, sed tantúmmodo: Dómine, ecce quem amas, infirmátur. Súfficit ut nóveris: non enim amas, et déseris.\par
\switchcolumn
\EN
\noindent What therefore sent his sisters to say? Lord, behold, he whom Thou lovest is sick and no more. They said not: Come, for Jesus loved him; and to tell Him that he was sick was enough. They dared not to say: Come, and heal him, they dared not to say: Speak the word where Thou art, and it shall be done here. And wherefore should they not have said this if they had the faith which won the Centurion so much praise? He had said: Lord, I am not worthy that Thou shouldest come under my roof; but speak the word only, and my servant shall be healed. (Matth. viii. 8.) But they said none of these things, only: Lord, behold, he whom Thou lovest is sick. It is enough that Thou shouldest know it. Thou art not one that lovest and leavest.\par
\switchcolumn*

\LA
\begin{Resp}\Rsign Qui persequebántur pópulum tuum, Dómine, demersísti eos in profúndum: \Star{} Et in colúmna nubis ductor eórum fuísti.\end{Resp}
\begin{Resp}\Vsign Deduxísti sicut oves pópulum tuum in manu Móysi et Aaron.\end{Resp}
\begin{Resp}\Rsign Et in colúmna nubis ductor eórum fuísti.\end{Resp}
\switchcolumn
\EN
\begin{Resp}\Rsign O Lord, Thou hast overwhelmed in the deep of the sea them which persecuted thy people; \Star{} Even thy people which Thou led in the pillar of the cloud.\end{Resp}
\begin{Resp}\Vsign Thou led thy people like a flock, by the hand of Moses and Aaron.\end{Resp}
\begin{Resp}\Rsign Even thy people, which Thou led in the pillar of the cloud.\end{Resp}
\switchcolumn*

\LA
\LessonHead{Lectio III}
\switchcolumn
\EN
\LessonHead{Lesson III}
\switchcolumn*

\LA
\noindent Dicit áliquis: Quómodo per Lázarum peccátor significabátur, et a Dómino sic amabátur? Audiat enim dicéntem: Non veni vocáre justos, sed peccatóres. Si enim peccatóres Deus non amáret, de cælo ad terram non descénderet. Audiens autem Jesus, dixit eis: Infírmitas hæc non est ad mortem, sed pro glória Dei, ut glorificétur Fílius Dei. Talis glorificátio ipsíus non ipsum auxit, sed nobis prófuit. Hoc est ergo quod ait: Non est ad mortem, sed pótius ad miráculum: quo facto créderent hómines in Christum, et vitárent veram mortem. Sane vidéte quemádmodum tamquam ex oblíquo Dóminus Deum se dixit: propter quosdam qui negant Fílium Dei Deum esse.\par
\switchcolumn
\EN
\noindent But some man will say: How shall Lazarus be a type of the sinner, and yet the Lord so love him? Let such an one hear the words of the same Lord, which He said: I am not come to call the righteous, but sinners. (Matth. ix. 13.) For if God had not loved sinners, He had not come down from heaven to earth. When Jesus heard that, He said: This sickness is not unto death, but for the glory of God, that the Son of God might be glorified thereby. Such a glorification is no increase of majesty for Him, but of profit for us. He therefore meaneth to say: This sickness is not unto death, but for the working of a miracle, the which being wrought, if men will thereby believe in Christ, they shall escape the real death. Note especially how the Lord doth in this place declare Himself to be God, as it were by implication, for the sake of some which say that He is not the Son of God.\par
\switchcolumn*

\LA
\begin{Resp}\Rsign Móyses fámulus Dei jejunávit quadragínta diébus et quadragínta nóctibus: \Star{} Ut legem Dómini mererétur accípere.\end{Resp}
\begin{Resp}\Vsign Ascéndens Móyses in montem Sínai ad Dóminum, fuit ibi quadragínta diébus et quadragínta nóctibus.\end{Resp}
\begin{Resp}\Rsign Ut legem Dómini mererétur accípere.\end{Resp}
\begin{Resp}\Vsign Glória Patri, et Fílio, \Star{} et Spirítui Sancto.\end{Resp}
\begin{Resp}\Rsign Ut legem Dómini mererétur accípere.\end{Resp}
\switchcolumn
\EN
\begin{Resp}\Rsign Moses, the servant of God, fasted forty days and forty nights \Star{} To make him meet to receive the Law of God.\end{Resp}
\begin{Resp}\Vsign Moses went up unto the Lord into Mount Sinai, and he was in the Mount forty days and forty nights.\end{Resp}
\begin{Resp}\Rsign To make him meet to receive the Law of God.\end{Resp}
\begin{Resp}\Vsign Glory be to the Father, and to the Son, \Star{} and to the Holy Ghost.\end{Resp}
\begin{Resp}\Rsign To make him meet to receive the Law of God.\end{Resp}
\switchcolumn*

\end{paracol}

\Office{Sabbato infra Hebdomadam IV in Quadragesima}{Saturday of the Fourth Week of Lent}{Feria major}
\addcontentsline{toc}{subsection}{Sabbato infra Hebdomadam IV in Quadragesima\enspace\textit{Saturday of the Fourth Week of Lent}}
\begin{paracol}{2}
\LA
\LessonHead{Lectio I}
\switchcolumn
\EN
\LessonHead{Lesson I}
\switchcolumn*

\LA
\LessonTitle{Léctio sancti Evangélii secúndum Joánnem}
\switchcolumn
\EN
\LessonTitle{Continuation of the Holy Gospel according to John}
\switchcolumn*

\LA
\Cite{Joannes 8:12-20}
\switchcolumn
\EN
\Cite{John 8:12-20}
\switchcolumn*

\LA
\noindent In illo témpore: Locútus est Jesus turbis Judæórum, dicens: Ego sum lux mundi: qui séquitur me, non ámbulat in ténebris, sed habébit lumen vitæ. Et réliqua.\par
\switchcolumn
\EN
\noindent At that time, Jesus spoke to the multitude of Jews saying: I am the light of the world: he that followeth me, walketh not in darkness, but shall have the light of life. And so on.\par
\switchcolumn*

\LA
\LessonTitle{Homilía sancti Augustíni Epíscopi}
\switchcolumn
\EN
\LessonTitle{Homily on this passage by St. Augustine, Bishop (of Hippo.)}
\switchcolumn*

\LA
\Source{Tract. 34 in Joannem, post init.}
\switchcolumn
\EN
\Source{34th Tract on John}
\switchcolumn*

\LA
\noindent Quod ait Dóminus: Ego sum lux mundi: clarum puto esse eis, qui habent óculos, unde hujus lucis partícipes fiant: qui autem non habent óculos, nisi in sola carne, mirántur quod dictum est a Dómino Jesu Christo: Ego sum lux mundi. Et forte non desit qui dicat apud semetípsum: Numquid forte Dóminus Christus est sol iste, qui ortu et occásu péragit diem? Non enim defuérunt hærétici, qui ista sensérunt. Manichǽi solem istum óculis cárneis visíbilem, expósitum et públicum non tantum homínibus, sed étiam pecóribus ad vidéndum, Christum Dóminum esse putavérunt.\par
\switchcolumn
\EN
\noindent I take it that these words of the Lord: “I am the Light of the world” are sufficiently clear to all men who have eyes which see that Light. At the same time, such men as have no eyes except those which are in their bodies, are surprised to find our Lord Jesus Christ saying, “I am the Light of the world”. And that we might not want somebody to say, “Is our Lord Jesus Christ, then, the same sun that riseth and setteth every day?” There have actually been heretics who did say it. The Manicheans believed that that sun which we see with our bodily eyes, and to see which is plain and common to beasts as well as men, was the Lord Christ.\par
\switchcolumn*

\LA
\begin{Resp}\Rsign Spléndida facta est fácies Móysi, dum respíceret in eum Dóminus: \Star{} Vidéntes senióres claritátem vultus ejus, admirántes timuérunt valde.\end{Resp}
\begin{Resp}\Vsign Cumque descendísset de monte Sínai, portábat duas tábulas testimónii, ignórans quod cornúta esset fácies ejus ex consórtio sermónis Dei.\end{Resp}
\begin{Resp}\Rsign Vidéntes senióres claritátem vultus ejus, admirántes timuérunt valde.\end{Resp}
\switchcolumn
\EN
\begin{Resp}\Rsign After that the Lord had looked upon him, the face of Moses shone. \Star{} And when the elders saw that his face shone, they marvelled and were sore afraid.\end{Resp}
\begin{Resp}\Vsign And when he came down from Mount Sinai with the two tables of testimony in his hand, he wist not that his face shone, because that God had spoken with him.\end{Resp}
\begin{Resp}\Rsign And when the elders saw that his face shone, they marvelled and were sore afraid.\end{Resp}
\switchcolumn*

\LA
\LessonHead{Lectio II}
\switchcolumn
\EN
\LessonHead{Lesson II}
\switchcolumn*

\LA
\noindent Sed cathólicæ Ecclésiæ recta fides ímprobat tale comméntum, et diabólicam doctrínam esse cognóscit: nec solum agnóscit credéndo, sed in quibus potest convíncit étiam disputándo. Improbémus ítaque hujúsmodi errórem, quem sancta ab inítio anathematizávit Ecclésia. Non arbitrémur Dóminum Jesum Christum hunc esse solem, quem vidémus oríri ab Oriénte, occídere in Occidénte: cujus cúrsui nox succédit, cujus rádii nube obumbrántur: qui certa de loco in locum motióne cómmigrat. Non est hoc Dóminus Christus. Non est Dóminus Christus sol factus, sed per quem sol factus est. Omnia enim per ipsum facta sunt, et sine ipso factum est nihil.\par
\switchcolumn
\EN
\noindent But the right faith of the Catholic Church damneth such comment, and recogniseth in it a doctrine of devils. And as it is her practice not only to brand errors by the difference of her own Creed, but also to remove them, if possible, by dint of argument, let us take up arms against this falsehood, which hath from the very beginning been the object of the curse of the Holy Church. God forbid that we should believe that our Lord Jesus Christ is this sun whose apparent movement is to rise every day in the East, and set every day in the West; which when we see no more, night cometh over us; and whose rays are sometimes intercepted by clouds and which hath some law of motion of its own whereby it describeth an orbit. The planet is not the same thing as our Lord Jesus Christ. Our Lord Jesus Christ is not that created sun, but He by Whom that sun was created; for all things were made by Him, and without Him was not anything made that was made. (John i. 30)\par
\switchcolumn*

\LA
\begin{Resp}\Rsign Ecce mitto Angelum meum, qui præcédat te, et custódiat semper: \Star{} Obsérva et audi vocem meam, et inimícus ero inimícis tuis, et affligéntes te afflígam: et præcédet te Angelus meus.\end{Resp}
\begin{Resp}\Vsign Israël, si me audíeris, non erit in te deus recens, neque adorábis deum aliénum: ego enim Dóminus.\end{Resp}
\begin{Resp}\Rsign Obsérva et audi vocem meam, et inimícus ero inimícis tuis, et affligéntes te afflígam: et præcédet te Angelus meus.\end{Resp}
\switchcolumn
\EN
\begin{Resp}\Rsign Behold, I send My Angel before thee, to keep thee. \Star{} Beware, and obey My voice; then I will be an enemy unto thine enemies, and an adversary unto thine adversaries; for Mine Angel shall go before thee.\end{Resp}
\begin{Resp}\Vsign O Israel, if thou wilt hearken unto Me, there shall no strange god be in thee, neither shalt thou worship any strange god, for I am the Lord.\end{Resp}
\begin{Resp}\Rsign Beware, and obey My voice; then I will be an enemy unto thine enemies, and an adversary unto thine adversaries; for Mine Angel shall go before thee.\end{Resp}
\switchcolumn*

\LA
\LessonHead{Lectio III}
\switchcolumn
\EN
\LessonHead{Lesson III}
\switchcolumn*

\LA
\noindent Est ergo lux, quæ fecit hanc lucem. Hanc amémus, hanc intellégere cupiámus, ipsam sitíamus, ut ad ipsam duce ipsa aliquándo veniámus: et in illa ita vivámus, ut nunquam omníno moriámur. Ista enim lux est, de qua prophetía olim præmíssa ita in Psalmo cécinit: Quóniam apud te est fons vitæ, et in lúmine tuo vidébimus lumen. Advértite quid de tali luce antíquus sanctórum hóminum Dei sermo præmíserit. Hómines, inquit, et juménta salvos fácies, Dómine: sicut multiplicáta est misericórdia tua, Deus.\par
\switchcolumn
\EN
\noindent He is therefore the Light by Whom the material light was made. Him may we love, Him may we long to know, Him may we thirst after; to Him may His own beams one day lead us, and in Him may we so live that we shall never die! For He, even He, and none other, He is that Light, of Whom the Prophet that was given of old time sang in the Psalms, when he said: “For with thee is the fountain of life, and in thy Light shall we see light.” (Ps. xxxv. 10.) Remember ye likewise what the word of God's ancient saints saith of such Light: “O Lord, Thou preservest man and beast! How excellent is thy loving-kindness, O God!” (7, 8.)\par
\switchcolumn*

\LA
\begin{Resp}\Rsign Atténdite, pópule meus, legem meam: \Star{} Inclináte aurem vestram in verba oris mei.\end{Resp}
\begin{Resp}\Vsign Apériam in parábolis os meum: loquar propositiónes ab inítio sǽculi.\end{Resp}
\begin{Resp}\Rsign Inclináte aurem vestram in verba oris mei.\end{Resp}
\begin{Resp}\Vsign Glória Patri, et Fílio, \Star{} et Spirítui Sancto.\end{Resp}
\begin{Resp}\Rsign Inclináte aurem vestram in verba oris mei.\end{Resp}
\switchcolumn
\EN
\begin{Resp}\Rsign Give ear, O My people, to My law; \Star{} Incline your ears to the words of My mouth.\end{Resp}
\begin{Resp}\Vsign I will open My mouth in parables; I will utter dark sayings of old.\end{Resp}
\begin{Resp}\Rsign Incline your ears to the words of My mouth.\end{Resp}
\begin{Resp}\Vsign Glory be to the Father, and to the Son, \Star{} and to the Holy Ghost.\end{Resp}
\begin{Resp}\Rsign Incline your ears to the words of My mouth.\end{Resp}
\switchcolumn*

\end{paracol}

\SeasonTitle{Tempus Passionis}{Passiontide}

\addcontentsline{toc}{section}{Tempus Passionis\enspace\textit{Passiontide}}

\Office{Dominica de Passione}{Passion Sunday}{Semiduplex I classis}
\addcontentsline{toc}{subsection}{Dominica de Passione\enspace\textit{Passion Sunday}}
\begin{paracol}{2}
\LA
\LessonHead{Lectio I}
\switchcolumn
\EN
\LessonHead{Lesson I}
\switchcolumn*

\LA
\LessonTitle{Incipit liber Jeremíæ Prophétæ}
\switchcolumn
\EN
\LessonTitle{Lesson from the book of Jeremias}
\switchcolumn*

\LA
\Cite{Jer 1:1-6}
\switchcolumn
\EN
\Cite{Jer 1:1-6}
\switchcolumn*

\LA
\noindent Verba Jeremíæ fílii Helcíæ, de sacerdótibus, qui fuérunt in Anathoth, in terra Bénjamin. Quod factum est verbum Dómini ad eum in diébus Josíæ fílii Amon regis Juda, in tertiodécimo anno regni ejus. Et factum est in diébus Joakim fílii Josíæ regis Juda, usque ad consummatiónem undécimi anni Sedecíæ fílii Josíæ regis Juda, usque ad transmigratiónem Jerúsalem, in mense quinto. Et factum est verbum Dómini ad me, dicens: Priúsquam te formárem in útero, novi te: et ántequam exíres de vulva, sanctificávi te, et prophétam in géntibus dedi te. Et dixi: A, a, a, Dómine Deus: ecce néscio loqui, quia puer ego sum.\par
\switchcolumn
\EN
\noindent The words of Jeremias the son of Helcias, of the priests that were in Anathoth, in the land of Benjamin. The word of the Lord which came to him in the days of Josias the son of Amon king of Juda, in the thirteenth year of his reign. And which came to him in the days of Joakim the son of Josias king of Juda, unto the end of the eleventh year of Sedecias the son of Josias king of Juda, even unto the carrying away of Jerusalem captive, in the fifth month. And the word of the Lord came to me, saying: Before I formed thee in the bowels of thy mother, I knew thee: and before thou camest forth out of the womb, I sanctified thee, and made thee a prophet unto the nations. And I said: Ah, ah, ah, Lord God: behold, I cannot speak, for I am a child.\par
\switchcolumn*

\LA
\begin{Resp}\Rsign Isti sunt dies, quos observáre debétis tempóribus suis: \Star{} Quartadécima die ad vésperum Pascha Dómini est: et in quintadécima solemnitátem celebrábitis altíssimo Dómino.\end{Resp}
\begin{Resp}\Vsign Locútus est Dóminus ad Móysen, dicens: Lóquere fíliis Israël, et dices ad eos.\end{Resp}
\begin{Resp}\Rsign Quartadécima die ad vésperum Pascha Dómini est: et in quintadécima solemnitátem celebrábitis altíssimo Dómino.\end{Resp}
\switchcolumn
\EN
\begin{Resp}\Rsign These are the days to be observed of you in their seasons. \Star{} In the fourteenth day at even is the Lord's Passover, and on the fifteenth day ye shall keep a Feast unto the Lord, the Most High.\end{Resp}
\begin{Resp}\Vsign The Lord spoke unto Moses, saying: Speak unto the children of Israel, and say unto them\end{Resp}
\begin{Resp}\Rsign In the fourteenth day at even is the Lord's Passover, and on the fifteenth day ye shall keep a Feast unto the Lord, the Most High.\end{Resp}
\switchcolumn*

\LA
\LessonHead{Lectio II}
\switchcolumn
\EN
\LessonHead{Lesson II}
\switchcolumn*

\LA
\Cite{Jer 1:7-13}
\switchcolumn
\EN
\Cite{Jer 1:7-13}
\switchcolumn*

\LA
\noindent Et dixit Dóminus ad me: Noli dícere: Puer sum: quóniam ad ómnia, quæ mittam te, ibis: et univérsa, quæcúmque mandávero tibi, loquéris. Ne tímeas a fácie eórum: quia tecum ego sum, ut éruam te, dicit Dóminus. Et misit Dóminus manum suam, et tétigit os meum: et dixit Dóminus ad me: Ecce dedi verba mea in ore tuo: ecce constítui te hódie super gentes, et super regna, ut evéllas, et déstruas, et dispérdas, et díssipes, et ædífices, et plantes. Et factum est verbum Dómini ad me, dicens: Quid tu vides, Jeremía? Et dixi: Virgam vigilántem ego vídeo. Et dixit Dóminus ad me: Bene vidísti, quia vigilábo ego super verbo meo, ut fáciam illud. Et factum est verbum Dómini secúndo ad me, dicens: Quid tu vides? Et dixi: Ollam succénsam ego vídeo, et fáciem ejus a fácie Aquilónis.\par
\switchcolumn
\EN
\noindent And the Lord said to me: Say not: I am a child: for thou shalt go to all that I shall send thee: and whatsoever I shall command thee, thou shalt speak. Be not afraid at their presence: for I am with thee to deliver thee, saith the Lord. And the Lord put forth his hand, and touched my mouth: and the Lord said to me: Behold I have given my words in thy mouth: Lo, I have set thee this day over the nations, and over the kingdoms, to root up, and pull down, and to waste, and to destroy, and to build, and to plant. And the word of the Lord came to me, saying: What seest thou, Jeremias? And I said: I see a rod watching. And the Lord said to me: Thou hast seen well: for I will watch over my word to perform it. And the word of the Lord came to me a second time, saying: What seest thou? I see a boiling caldron, and the face thereof from the face of the north.\par
\switchcolumn*

\LA
\begin{Resp}\Rsign Multiplicáti sunt qui tríbulant me, et dicunt: Non est salus illi in Deo ejus: \Star{} Exsúrge, Dómine, salvum me fac, Deus meus.\end{Resp}
\begin{Resp}\Vsign Nequándo dicat inimícus meus: Præválui advérsus eum.\end{Resp}
\begin{Resp}\Rsign Exsúrge, Dómine, salvum me fac, Deus meus.\end{Resp}
\switchcolumn
\EN
\begin{Resp}\Rsign They be increased that trouble me, and that say There is no help for him in his God. \Star{} Arise, O Lord! Save me, O my God!\end{Resp}
\begin{Resp}\Vsign Lest mine enemy say I have prevailed against him.\end{Resp}
\begin{Resp}\Rsign Arise, O Lord! Save me, O my God!\end{Resp}
\switchcolumn*

\LA
\LessonHead{Lectio III}
\switchcolumn
\EN
\LessonHead{Lesson III}
\switchcolumn*

\LA
\Cite{Jer 1:14-19}
\switchcolumn
\EN
\Cite{Jer 1:14-19}
\switchcolumn*

\LA
\noindent Et dixit Dóminus ad me: Ab Aquilóne pandétur malum super omnes habitatóres terræ. Quia ecce ego convocábo omnes cognatiónes regnórum Aquilónis, ait Dóminus: et vénient et ponent unusquísque sólium suum in intróitu portárum Jerúsalem, et super omnes muros ejus in circúitu, et super univérsas urbes Juda. Et loquar judícia mea cum eis super omnem malítiam eórum, qui dereliquérunt me, et libavérunt diis aliénis, et adoravérunt opus mánuum suárum. Tu ergo, accínge lumbos tuos, et surge, et lóquere ad eos ómnia quæ ego præcípio tibi. Ne formídes a fácie eórum: nec enim timére te fáciam vultum eórum. Ego quippe dedi te hódie in civitátem munítam, et in colúmnam férream, et in murum ǽreum, super omnem terram, régibus Juda, princípibus ejus, et sacerdótibus et pópulo terræ. Et bellábunt advérsum te, et non prævalébunt: quia ego tecum sum, ait Dóminus, ut líberem te.\par
\switchcolumn
\EN
\noindent And the Lord said to me: from the north shall an evil break forth upon all the inhabitants of the land. For behold I will call together all the families of the kingdoms of the north: saith the Lord: and they shall come, and shall set every one his throne in the entrance of the gates of Jerusalem, and upon all the walls thereof round about, and upon all the cities of Juda, And I will pronounce my judgements against them, touching all their wickedness, who have forsaken me, and have sacrificed to strange gods, and have adored the work of their own hands. Thou therefore gird up thy loins, and arise, and speak to them all that I command thee. Be not afraid at their presence for I will make thee not to fear their countenance. For behold I have made thee this day a fortified city, and a pillar of iron, and a wall of brass, over all the land, to the kings of Juda, to the princes thereof, and to the priests, and to the people of the land. And they shall fight against thee, and shall not prevail: for I am with thee, saith the Lord, to deliver thee.\par
\switchcolumn*

\LA
\begin{Resp}\Rsign Usquequo exaltábitur inimícus meus super me? \Star{} Réspice, et exáudi me, Dómine, Deus meus.\end{Resp}
\begin{Resp}\Vsign Qui tríbulant me, exsultábunt si motus fúero: ego autem in misericórdia tua sperábo.\end{Resp}
\begin{Resp}\Rsign Réspice, et exáudi me, Dómine, Deus meus.\end{Resp}
\begin{Resp}\Vsign Glória Patri, et Fílio, \Star{} et Spirítui Sancto.\end{Resp}
\begin{Resp}\Rsign Usquequo exaltábitur inimícus meus super me? * Réspice, et exáudi me, Dómine, Deus meus.\end{Resp}
\switchcolumn
\EN
\begin{Resp}\Rsign How long shall mine enemy be exalted over me? \Star{} Consider, and hear me, O Lord my God!\end{Resp}
\begin{Resp}\Vsign Those that trouble me will rejoice when I am moved but I have trusted in thy mercy.\end{Resp}
\begin{Resp}\Rsign Consider and hear me, O Lord my God!\end{Resp}
\begin{Resp}\Vsign Glory be to the Father, and to the Son, \Star{} and to the Holy Ghost.\end{Resp}
\begin{Resp}\Rsign How long shall mine enemy be exalted over me? Consider and hear me, O Lord my God!\end{Resp}
\switchcolumn*

\LA
\LessonHead{Lectio IV}
\switchcolumn
\EN
\LessonHead{Lesson IV}
\switchcolumn*

\LA
\LessonTitle{Sermo sancti Leónis Papæ}
\switchcolumn
\EN
\LessonTitle{From the Sermons of Pope St. Leo (the Great.)}
\switchcolumn*

\LA
\Source{Sermo 9 de Quadragesima}
\switchcolumn
\EN
\Source{9th for Lent}
\switchcolumn*

\LA
\noindent In ómnibus, dilectíssimi, solemnitátibus christiánis non ignorámus paschále sacraméntum esse præcípuum: cui condígne et cóngrue suscipiéndo, totíus quidem nos témporis institúta refórmant: sed devotiónem nostram præséntes vel máxime dies éxigunt, quos illi sublimíssimo divínæ misericórdiæ sacraménto scimus esse contíguos. In quibus mérito a sanctis Apóstolis per doctrínam Spíritus Sancti majóra sunt ordináta jejúnia: ut per commúne consórtium crucis Christi, étiam nos áliquid in eo quod propter nos gessit, agerémus, sicut Apóstolus ait: Si compátimur, et conglorificábimur. Certa atque secúra est exspectátio promíssæ beatitúdinis, ubi est participátio Domínicæ passiónis.\par
\switchcolumn
\EN
\noindent Dearly beloved brethren, we know that of all the solemn Feasts which are kept by Christians the Passover is the chief. The ordinances of the whole rest of the year are ordered to the end of preparing us to come to this one in worthy and meet manner. But these days, which now are, are they which ought most especially to stir up a godly mind in us, seeing that they are they which are nearest to that most glorious mystery of God's mercy. In these days the holy Apostles, taught by the Holy Ghost, ordered the chiefest store of Fasting, that we, sharing His Cross with Christ, might, albeit we are what we are, in Him, do some of the same things which He did for our sakes, and so realise the saying of the Apostle: “If we suffer with Him, we shall be also glorified together.” (Rom. viii. 17.) He that is partaker of the sufferings (2 Cor. i. 7) of the Lord hath a sure and certain hope of that blessedness which He hath promised unto us.\par
\switchcolumn*

\LA
\begin{Resp}\Rsign Deus meus es tu, ne discédas a me: \Star{} Quóniam tribulátio próxima est, et non est qui ádjuvet.\end{Resp}
\begin{Resp}\Vsign Tu autem, Dómine, ne elongáveris auxílium tuum a me: ad defensiónem meam áspice.\end{Resp}
\begin{Resp}\Rsign Quóniam tribulátio próxima est, et non est qui ádjuvet.\end{Resp}
\switchcolumn
\EN
\begin{Resp}\Rsign Thou art my God; be not far from me, for trouble is near; \Star{} For there is none to help.\end{Resp}
\begin{Resp}\Vsign But be not thy strength far from me; O Lord, haste thee to help me.\end{Resp}
\begin{Resp}\Rsign For trouble is near; for there is none to help.\end{Resp}
\switchcolumn*

\LA
\LessonHead{Lectio V}
\switchcolumn
\EN
\LessonHead{Lesson V}
\switchcolumn*

\LA
\noindent Nemo est, dilectíssimi, cui per conditiónem témporis socíetas hujus glóriæ denegétur, tamquam tranquíllitas pacis vácua sit occasióne virtútis. Apóstolus enim prǽdicat, dicens: Omnes qui pie volunt vívere in Christo, persecutiónem patiéntur: et ídeo numquam deest tribulátio persecutiónis, si numquam desit observántia pietátis. Dóminus enim in exhortatiónibus suis dicit: Qui non áccipit crucem suam, et séquitur me, non est me dignus. Nec dubitáre debémus, hanc vocem non solum ad discípulos Christi, sed ad cunctos fidéles, totámque Ecclésiam pertinére, quæ salutáre suum in his qui áderant, universáliter audiébat.\par
\switchcolumn
\EN
\noindent Dearly beloved brethren, there is no man to whom the state of the age in which he liveth denieth a share in this glory of partaking, first the sufferings, and then the triumph and joy, of Christ. It is not as though this time of peace were barren in occasions of valour. The Apostle giveth us this warning All that will live godly in Christ Jesus shall suffer persecution. (2 Tim. iii. 12.) And therefore, as long as godliness is watchful, persecution will never be asleep. The Lord Himself saith in one of His own exhortations He that taketh not his cross, and followeth after Me, is not worthy of Me. (Matth. x. 38.) And we must not doubt that these words of Christ apply not only to His immediate disciples, to whom He spoke them, but belong to all the faithful and to the whole Church, who, whosoever be the believers of whom she is for the time composed on earth, heareth in these words the way to be saved which her Lord hath appointed for them.\par
\switchcolumn*

\LA
\begin{Resp}\Rsign In te jactátus sum ex útero, de ventre matris meæ Deus meus es tu, ne discédas a me: \Star{} Quóniam tribulátio próxima est, et non est qui ádjuvet.\end{Resp}
\begin{Resp}\Vsign Salva me ex ore leónis, et a córnibus unicórnium humilitátem meam.\end{Resp}
\begin{Resp}\Rsign Quóniam tribulátio próxima est, et non est qui ádjuvet.\end{Resp}
\switchcolumn
\EN
\begin{Resp}\Rsign I was cast upon thee from the womb; Thou art my God from my mother's belly; be not far from me. \Star{} For trouble is near, and there is none to help.\end{Resp}
\begin{Resp}\Vsign Save me from the lion's mouth, and mine affliction from the horns of the unicorns.\end{Resp}
\begin{Resp}\Rsign For trouble is near, and there is none to help.\end{Resp}
\switchcolumn*

\LA
\LessonHead{Lectio VI}
\switchcolumn
\EN
\LessonHead{Lesson VI}
\switchcolumn*

\LA
\noindent Sicut ergo totíus est córporis pie vívere, ita totíus est témporis crucem ferre: quæ mérito ferri unicuíque suadétur, quia própriis modis atque mensúris ab unoquóque tolerátur. Unum nomen est persecutiónis, sed non una est causa certáminis: et plus plerúmque perículi est in insidiatóre occúlto, quam in hoste manifésto. Beátus Job alternántibus bonis ac malis mundi hujus erudítus, pie veracitérque dicébat: Nonne tentátio est vita hóminis super terram? Quóniam non solis dolóribus córporis atque supplíciis ánima fidélis impétitur, verum étiam, salva incolumitáte membrórum, gravi morbo urgétur, si carnis voluptáte mollítur. Sed cum caro concupíscit advérsus spíritum, spíritus autem advérsus carnem; præsídio crucis Christi mens rationális instrúitur, nec cupiditátibus nóxiis illécta conséntit, quóniam continéntiæ clavis et Dei timóre transfígitur.\par
\switchcolumn
\EN
\noindent As then, it is the duty of the whole body of the Church to live godly, so is it her right at all times to be a-bearing of her Master's Cross, and that not only in her general body, but individually in the person of each one of her members, who differ every one from another in the way in which they have to carry it, and the shape in which it is laid upon them. The one common name for all their carrying of the Cross is persecution, but the manner of his wrestling is special to each; and there is often more danger in the ambush than in the pitched field of battle. Blessed Job, who had tried both the goods and the ills of this world, said: Is not the life of man upon earth a warfare? (vii. 1.) The attack upon the faithful soul arrayeth itself not alone in bodily torture and punishment; yea, when the limbs are sound enough, fearful is the ravage that threateneth us when the lusts of the flesh unman us. But when the flesh lusteth against the spirit, and the spirit against the flesh (Gal. iv. 7) the reasonable mind findeth her reinforcement in the helpful Cross of Christ, and though she be lured by foul cravings, yet refuseth to give her consent, for God maketh her pure thoughts to tremble for fear of Him. (Ps. cxviii. 120.)\par
\switchcolumn*

\LA
\begin{Resp}\Rsign In próximo est tribulátio mea, Dómine, et non est qui ádjuvet; ut fódiant manus meas et pedes meos: líbera me de ore leónis, \Star{} Ut enárrem nomen tuum frátribus meis.\end{Resp}
\begin{Resp}\Vsign Erue a frámea, Deus, ánimam meam, et de manu canis únicam meam.\end{Resp}
\begin{Resp}\Rsign Ut enárrem nomen tuum frátribus meis.\end{Resp}
\begin{Resp}\Vsign Glória Patri, et Fílio, \Star{} et Spirítui Sancto.\end{Resp}
\begin{Resp}\Rsign In próximo est tribulátio mea, Dómine, et non est qui ádjuvet; ut fódiant manus meas et pedes meos: líbera me de ore leónis, * Ut enárrem nomen tuum frátribus meis.\end{Resp}
\switchcolumn
\EN
\begin{Resp}\Rsign O Lord, my trouble is near, and there is none to help me; or ever they pierce my hands and my feet, save me from the lion's mouth! \Star{} That I may declare thy Name unto my brethren.\end{Resp}
\begin{Resp}\Vsign O God, deliver my soul from the sword, and my darling from the power of the dog.\end{Resp}
\begin{Resp}\Rsign That I may declare thy Name unto my brethren.\end{Resp}
\begin{Resp}\Vsign Glory be to the Father, and to the Son, \Star{} and to the Holy Ghost.\end{Resp}
\begin{Resp}\Rsign O Lord, my trouble is near, and there is none to help me; or ever they pierce my hands and my feet, save me from the lion's mouth! that I may declare thy Name unto my brethren.\end{Resp}
\switchcolumn*

\LA
\LessonHead{Lectio VII}
\switchcolumn
\EN
\LessonHead{Lesson VII}
\switchcolumn*

\LA
\LessonTitle{Léctio sancti Evangélii secúndum Joánnem}
\switchcolumn
\EN
\LessonTitle{From the Holy Gospel according to John}
\switchcolumn*

\LA
\Cite{Joannes 8:46-59}
\switchcolumn
\EN
\Cite{John 8:46-59}
\switchcolumn*

\LA
\noindent In illo témpore: Dicébat Jesus turbis Judæórum: Quis ex vobis árguet me de peccáto? Si veritátem dico vobis, quare non créditis mihi? Et réliqua.\par
\switchcolumn
\EN
\noindent At that time, Jesus said to the multitude of Jews: Which of you shall convince me of sin? If I say the truth to you, why do you not believe me? And so on.\par
\switchcolumn*

\LA
\LessonTitle{Homilía sancti Gregórii Papæ}
\switchcolumn
\EN
\LessonTitle{Homily by Pope St. Gregory (the Great)}
\switchcolumn*

\LA
\Source{Homilia 18 in Evangelia}
\switchcolumn
\EN
\Source{18th on the Gospels}
\switchcolumn*

\LA
\noindent Pensáte, fratres caríssimi, mansuetúdinem Dei. Relaxáre peccáta vénerat, et dicébat: Quis ex vobis árguet me de peccáto? Non dedignátur ex ratióne osténdere se peccatórem non esse, qui ex virtúte divinitátis póterat peccatóres justificáre. Sed terríbile est valde, quod súbditur: Qui ex Deo est, verba Dei audit: proptérea vos non audítis, quia ex Deo non estis. Si enim ipse verba Dei audit qui ex Deo est, et audíre verba ejus non potest quisquis ex illo non est: intérroget se unusquísque, si verba Dei in aure cordis pércipit; et intélleget unde sit. Cæléstem pátriam desideráre Véritas jubet, carnis desidéria cónteri, mundi glóriam declináre, aliéna non appétere, própria largíri.\par
\switchcolumn
\EN
\noindent Dearly beloved brethren, consider the gentleness of God. He came to take away sins, and He saith Which of you convinceth Me of sin? He Who, through the might of His Godhead, was able to justify sinners, was contented to show by argument that He was not Himself a sinner. But exceeding dread is that which followeth. He that is of God heareth God's words; ye, therefore, hear them not, because ye are not of God. If, then, whosoever is of God heareth God's words, and whosoever is not of Him cannot hear His words, let each one ask himself if he, in the ear of his heart, heareth God's words, and understandeth Whose words they are? The Truth commandeth us to long for a Fatherland in heaven, to bridle the lusts of the flesh, to turn away from the glory of the world, to seek no man's goods, and to give away our own.\par
\switchcolumn*

\LA
\begin{Resp}\Rsign Tota die contristátus ingrediébar, Dómine: quóniam ánima mea compléta est illusiónibus: \Star{} Et vim faciébant, qui quærébant ánimam meam.\end{Resp}
\begin{Resp}\Vsign Amíci mei et próximi mei advérsum me appropinquavérunt et stetérunt: et qui juxta me erant, de longe stetérunt.\end{Resp}
\begin{Resp}\Rsign Et vim faciébant, qui quærébant ánimam meam.\end{Resp}
\switchcolumn
\EN
\begin{Resp}\Rsign O Lord, I go mourning all the day long, for my soul is filled with a loathsome disease \Star{} They also that sought after my life have used violence against me.\end{Resp}
\begin{Resp}\Vsign My friends and my neighbours draw near, and stand over against me; and they that are nearest to me stand afar off.\end{Resp}
\begin{Resp}\Rsign They also that sought after my life have used violence against me.\end{Resp}
\switchcolumn*

\LA
\LessonHead{Lectio VIII}
\switchcolumn
\EN
\LessonHead{Lesson VIII}
\switchcolumn*

\LA
\noindent Penset ergo apud se unusquísque vestrum, si hæc vox Dei in cordis ejus aure conváluit, et quia jam ex Deo sit, agnóscit. Nam sunt nonnúlli, qui præcépta Dei nec aure córporis percípere dignántur. Et sunt nonnúlli, qui hæc quidem córporis aure percípiunt, sed nullo ea mentis desidério complectúntur. Et sunt nonnúlli, qui libénter verba Dei suscípiunt, ita ut étiam in flétibus compungántur, sed post lacrimárum tempus ad iniquitátem rédeunt. Hi profécto verba Dei non áudiunt, qui hæc exercére in ópere contémnunt. Vitam ergo vestram, fratres caríssimi, ante mentis óculos revocáte, et alta consideratióne pertiméscite hoc quod ex ore Veritátis sonat: Proptérea vos non audítis, quia ex Deo non estis.\par
\switchcolumn
\EN
\noindent Let each of you, therefore, think within himself if this voice of God is heard in the ear of his heart, and if he knoweth already if he is of God. For some there be, whom it pleaseth not to hear the commandments of God even with their bodily ears. And some there be, who receive the same with their bodily ears, but whose heart is far from them. And some also there be, who hear the words of God with joy, so that they are moved thereby even to tears; but when their fit of weeping is past they turn again to iniquity. They hear not the words of God, who despise to do them. Therefore, dearly beloved brethren, call up your own life before your mind's eye, and then ponder with trembling those awful words which the mouth of the Truth spoke. Ye therefore hear them not, because ye are not of God.\par
\switchcolumn*

\LA
\begin{Resp}\Rsign Ne avértas fáciem tuam a púero tuo, Dómine: \Star{} Quóniam tríbulor, velóciter exáudi me.\end{Resp}
\begin{Resp}\Vsign Inténde ánimæ meæ, et líbera eam: propter inimícos meos éripe me.\end{Resp}
\begin{Resp}\Rsign Quóniam tríbulor, velóciter exáudi me.\end{Resp}
\switchcolumn
\EN
\begin{Resp}\Rsign O Lord, hide not thy face from thy servant \Star{} For I am in trouble; hear me speedily.\end{Resp}
\begin{Resp}\Vsign Draw nigh unto my soul, and redeem it; deliver me, because of mine enemies.\end{Resp}
\begin{Resp}\Rsign For I am in trouble; hear me speedily.\end{Resp}
\switchcolumn*

\LA
\LessonHead{Lectio IX}
\switchcolumn
\EN
\LessonHead{Lesson IX}
\switchcolumn*

\LA
\noindent Sed hoc quod de réprobis Véritas lóquitur, ipsi hoc de semetípsis réprobi iníquis suis opéribus osténdunt: nam séquitur: Respondérunt ígitur Judǽi et dixérunt ei: Nonne bene dícimus nos, quia Samaritánus es tu, et dæmónium habes? Accépta autem tanta contumélia, quid Dóminus respóndeat, audiámus: Ego dæmónium non hábeo, sed honorífico Patrem meum, et vos inhonorástis me. Quia enim Samaritánus interpretátur custos: et ipse veráciter custos est, de quo Psalmísta ait: Nisi Dóminus custodíerit civitátem, in vanum vígilant qui custódiunt eam: et cui per Isaíam dícitur: Custos quid de nocte? custos quid de nocte? respondére nóluit Dóminus, Samaritánus non sum; sed, Ego dæmónium non hábeo. Duo quippe ei illáta fuérunt: unum negávit, áliud tacéndo consénsit.\par
\switchcolumn
\EN
\noindent The Truth speaketh these words concerning the reprobate; but the reprobate make manifest the same thing concerning themselves, by their evil works. Thus immediately followeth: Then answered the Jews, and said unto Him: Say we not well that Thou art a Samaritan, and hast a devil? But let us hear what the Lord said to this insult. I have not a devil, but I honour My Father, and ye do dishonour Me. The Lord said: I have not a devil, but He did not say: I am not a Samaritan, for in a sense a Samaritan He was indeed, since the word Samaritan, in the Hebrew tongue, signifieth, being interpreted, a Watcher, and the Lord is that Watcher, of Whom the Psalmist saith (cxxviii. 2) that unless He keep the city, other watchman waketh but in vain. He also is that Watchman unto Whom crieth Isaiah (xxi. 11): Watchman, what of the night? Watchman, what of the night? Therefore the Lord said: I have not a devil, but not I am not a Samaritan. Of the two things brought against Him He denied one; but by His silence, admitted the other.\par
\switchcolumn*

\LA
\begin{Resp}\Rsign Quis dabit cápiti meo aquam, et óculis meis fontem lacrimárum, et plorábo die ac nocte? quia frater propínquus supplantávit me, \Star{} Et omnis amícus fraudulénter incéssit in me.\end{Resp}
\begin{Resp}\Vsign Fiant viæ eórum ténebræ et lúbricum: et Angelus Dómini pérsequens eos.\end{Resp}
\begin{Resp}\Rsign Et omnis amícus fraudulénter incéssit in me.\end{Resp}
\begin{Resp}\Vsign Glória Patri, et Fílio, \Star{} et Spirítui Sancto.\end{Resp}
\begin{Resp}\Rsign Quis dabit cápiti meo aquam, et óculis meis fontem lacrimárum, et plorábo die ac nocte? quia frater propínquus supplantávit me, et omnis amícus fraudulénter incéssit in me.\end{Resp}
\switchcolumn
\EN
\begin{Resp}\Rsign O that my head were waters, and mine eyes a fountain of tears, that I might weep day and night! for my nearest brother hath supplanted me, \Star{} And my neighbour hath walked with slanders against me.\end{Resp}
\begin{Resp}\Vsign Let their way be dark and slippery, and let the Angel of the Lord persecute them.\end{Resp}
\begin{Resp}\Rsign And my neighbour hath walked with slanders against me.\end{Resp}
\begin{Resp}\Vsign Glory be to the Father, and to the Son, \Star{} and to the Holy Ghost.\end{Resp}
\begin{Resp}\Rsign O that my head were waters, and mine eyes a fountain of tears, that I might weep day and night! for my nearest brother hath supplanted me, and my neighbour hath walked with slanders against me.\end{Resp}
\switchcolumn*

\end{paracol}

\Office{Feria Secunda infra Hebdomadam Passionis}{Monday in Passion Week}{Feria major}
\addcontentsline{toc}{subsection}{Feria Secunda infra Hebdomadam Passionis\enspace\textit{Monday in Passion Week}}
\begin{paracol}{2}
\LA
\LessonHead{Lectio I}
\switchcolumn
\EN
\LessonHead{Lesson I}
\switchcolumn*

\LA
\LessonTitle{Léctio sancti Evangélii secúndum Joánnem}
\switchcolumn
\EN
\LessonTitle{Continuation of the Holy Gospel according to John}
\switchcolumn*

\LA
\Cite{Joannes 7:32-39}
\switchcolumn
\EN
\Cite{John 7:32-39}
\switchcolumn*

\LA
\noindent In illo témpore: Misérunt príncipes et pharisǽi minístros, ut apprehénderent Jesum. Et réliqua.\par
\switchcolumn
\EN
\noindent At that time, the rulers and Pharisees sent ministers to apprehend him. And so on.\par
\switchcolumn*

\LA
\LessonTitle{Homilía sancti Augustíni Epíscopi}
\switchcolumn
\EN
\LessonTitle{Homily by St. Augustine, Bishop (of Hippo.)}
\switchcolumn*

\LA
\Source{Tract. 31 in Joannem, circa medium}
\switchcolumn
\EN
\Source{31th Tract on John}
\switchcolumn*

\LA
\noindent Quómodo apprehénderent adhuc noléntem? Quia ergo non póterant apprehéndere noléntem, missi sunt, ut audírent docéntem. Quid docéntem? Dicit ergo Jesus: Adhuc módicum tempus vobíscum sum. Quod modo vultis fácere, factúri estis; sed non modo, quia modo nolo. Quare adhuc modo nolo? Quia adhuc módicum tempus vobíscum sum, et tunc vado ad eum qui me misit. Implére débeo dispensatiónem meam, et sic perveníre ad passiónem meam.\par
\switchcolumn
\EN
\noindent How could they take Him until such time as He willed to be taken? If, then, they could not take Him until He willed to be taken, were they sent to watch His teaching? Then said Jesus unto them Yet a little while am I with you what ye now seek to do, ye shall do; but not yet, for I will not so yet. And why will I not so yet? Because yet a little while am I with you, and then I go unto Him that sent Me I must fulfill that which I am sent to do, and so go to suffer.\par
\switchcolumn*

\LA
\begin{Resp}\Rsign Deus meus, éripe me de manu peccatóris: et de manu contra legem agéntis, et iníqui: \Star{} Quóniam tu es patiéntia mea.\end{Resp}
\begin{Resp}\Vsign Deus meus, ne elongéris a me: Deus meus, in auxílium meum réspice.\end{Resp}
\begin{Resp}\Rsign Quóniam tu es patiéntia mea.\end{Resp}
\switchcolumn
\EN
\begin{Resp}\Rsign Deliver me, O my God, out of the hand of the wicked, and out of the hand of the unrighteous and cruel man. \Star{} For Thou art my hope.\end{Resp}
\begin{Resp}\Vsign O my God, be not far from me O my God, make haste for my help.\end{Resp}
\begin{Resp}\Rsign For Thou art my hope.\end{Resp}
\switchcolumn*

\LA
\LessonHead{Lectio II}
\switchcolumn
\EN
\LessonHead{Lesson II}
\switchcolumn*

\LA
\noindent Quærétis me, et non inveniétis, et ubi sum ego, vos non potéstis veníre. Hic jam resurrectiónem suam prædíxit: noluérunt enim agnóscere præséntem, et póstea quæsiérunt, cum vidérent in eum multitúdinem jam credéntem. Magna enim signa facta sunt étiam cum Dóminus resurréxit, et ascéndit in cælum. Tunc per discípulos facta sunt magna: sed ille per illos, qui et per seípsum: ipse quippe illis díxerat: Sine me nihil potéstis fácere. Quando claudus ille, qui sedébat ad portam, ad vocem Petri surréxit, et suis pédibus ambulávit, ita ut hómines miraréntur, sic eos allocútus est Petrus, quia non in sua potestáte ista fecit, sed in virtúte illíus, quem ipsi occidérunt. Multi compúncti dixérunt: Quid faciémus?\par
\switchcolumn
\EN
\noindent You shall seek Me, and shall not find Me, and where I am thither ye cannot come. In these words He foretold already His rising again from the dead. While He was with them they would not know Him; and afterwards they sought Him, when they saw that a multitude already believed in Him. For great signs were wrought also when the Lord rose again, and ascended up into heaven. Then were great signs again wrought through the Disciples, that is, through them by Him Who worketh the same directly also by Himself, according as He had said unto them: Without Me ye can do nothing. (John xv. 5.) When that lame man that was laid daily at the Beautiful Gate of the Temple stood up at the voice of Peter (Acts iii.) and walked, and all the people were filled with wonder, Peter bade them know that it was not by his own power that he had made him to walk, but by the power of Him Whom they had killed. And when they heard this, many were pricked in their heart, and said: What shall we do? (Acts ii. 37.)\par
\switchcolumn*

\LA
\begin{Resp}\Rsign Qui custodiébant ánimam meam, consílium fecérunt in unum, dicéntes: Deus derelíquit eum, \Star{} Persequímini et comprehéndite eum: quia non est qui líberet eum: Deus meus, ne elongéris a me: Deus meus, in adjutórium meum inténde.\end{Resp}
\begin{Resp}\Vsign Omnes inimíci advérsum me cogitábant mala mihi: verbum iníquum mandavérunt advérsum me, dicéntes.\end{Resp}
\begin{Resp}\Rsign Persequímini et comprehéndite eum: quia non est qui líberet eum: Deus meus, ne elongéris a me: Deus meus, in adjutórium meum inténde.\end{Resp}
\switchcolumn
\EN
\begin{Resp}\Rsign They that lay wait for my soul take counsel together, saying God hath forsaken him; \Star{} Persecute and take him, for there is none to deliver him. O my God, be not far from me; O my God, make haste for my help.\end{Resp}
\begin{Resp}\Vsign All that hate me whispered together against me; against me did they devise my hurt, saying:\end{Resp}
\begin{Resp}\Rsign Persecute and take him, for there is none to deliver him. O my God, be not far from me; O my God, make haste for my help.\end{Resp}
\switchcolumn*

\LA
\LessonHead{Lectio III}
\switchcolumn
\EN
\LessonHead{Lesson III}
\switchcolumn*

\LA
\noindent Vidérunt enim se ingénti crímine impietátis adstríctos, quando illum occidérunt, quem venerári et adoráre debuérunt: et hoc putábant esse inexpiábile. Magnum enim fácinus erat, cujus considerátio illos fáceret desperáre: sed non debébant desperáre, pro quibus in cruce pendens Dóminus est dignátus oráre. Díxerat enim: Pater, ignósce illis, quia nésciunt quid fáciunt. Vidébat quosdam suos inter multos aliénos: illis jam petébat véniam, a quibus adhuc accipiébat injúriam. Non enim attendébat quod ab ipsis moriebátur, sed quia pro ipsis moriebátur.\par
\switchcolumn
\EN
\noindent Nor they saw that they were burdened with the guilt of an exceeding great sin, in that they had killed Him, Whom it was their duty to worship and adore and for that guilt they knew of no propitiation. Yea, their sin was indeed exceeding great; and the consideration of it made them to despair for whom the Lord, when He hung upon the Cross, had been willing to pray, as it is written. Then said Jesus: Father, forgive them, for they know not what they do. (Luke xxiii. 34.) At that hour He had seen among many aliens some that were His Own; for them He asked forgiveness, while yet He suffered at their hand, nor considered that they were putting Him to death, but only that He was dying for them.\par
\switchcolumn*

\LA
\begin{Resp}\Rsign Pacífice loquebántur mihi inimíci mei, et in ira molésti erant mihi: \Star{} Vidísti, Dómine, ne síleas, ne discédas a me.\end{Resp}
\begin{Resp}\Vsign Ego autem cum mihi molésti essent, induébam me cilício, et humiliábam in jejúnio ánimam meam.\end{Resp}
\begin{Resp}\Rsign Vidísti, Dómine, ne síleas, ne discédas a me.\end{Resp}
\begin{Resp}\Vsign Glória Patri, et Fílio, \Star{} et Spirítui Sancto.\end{Resp}
\begin{Resp}\Rsign Pacífice loquebántur mihi inimíci mei, et in ira molésti erant mihi: * Vidísti, Dómine, ne síleas, ne discédas a me.\end{Resp}
\switchcolumn
\EN
\begin{Resp}\Rsign Mine enemies spoke to me peaceably, but in wrath they troubled me. \Star{} This Thou hast seen, O Lord; keep not silence be not far from me.\end{Resp}
\begin{Resp}\Vsign But as for me, when they troubled me, my clothing was sackcloth, and I humbled my soul with fasting.\end{Resp}
\begin{Resp}\Rsign This Thou hast seen, O Lord; keep not silence, be not far from me.\end{Resp}
\begin{Resp}\Vsign Glory be to the Father, and to the Son, \Star{} and to the Holy Ghost.\end{Resp}
\begin{Resp}\Rsign Mine enemies spoke to me peaceably, but in wrath they troubled me. * This Thou hast seen, O Lord; keep not silence, be not far from me.\end{Resp}
\switchcolumn*

\end{paracol}

\Office{Feria Tertia infra Hebdomadam Passionis}{Tuesday in Passion Week}{Feria major}
\addcontentsline{toc}{subsection}{Feria Tertia infra Hebdomadam Passionis\enspace\textit{Tuesday in Passion Week}}
\begin{paracol}{2}
\LA
\LessonHead{Lectio I}
\switchcolumn
\EN
\LessonHead{Lesson I}
\switchcolumn*

\LA
\LessonTitle{Léctio sancti Evangélii secúndum Joánnem}
\switchcolumn
\EN
\LessonTitle{Continuation of the Holy Gospel according to John}
\switchcolumn*

\LA
\Cite{Joannes 7:1-13}
\switchcolumn
\EN
\Cite{John 7:1-13}
\switchcolumn*

\LA
\noindent In illo témpore: Ambulábat Jesus in Galilǽam: non enim volébat in Judǽam ambuláre, quia quærébant eum Judǽi interfícere. Et réliqua.\par
\switchcolumn
\EN
\noindent At that time: Jesus walked in Galilee; for he would not walk in Judea, because the Jews sought to kill him. And so on.\par
\switchcolumn*

\LA
\LessonTitle{Homilía sancti Augustíni Epíscopi}
\switchcolumn
\EN
\LessonTitle{Homily by St. Augustine, Bishop (of Hippo.)}
\switchcolumn*

\LA
\Source{Tractatus 28 in Joannem}
\switchcolumn
\EN
\Source{28th Tract on John}
\switchcolumn*

\LA
\noindent In isto Evangélii capítulo, fratres, Dóminus noster Jesus Christus secúndum hóminem se plúrimum commendávit fídei nostræ. Etenim semper hoc egit dictis et factis suis, ut Deus credátur et homo: Deus qui nos fecit, homo qui nos quæsívit: Deus cum Patre semper, homo nobíscum ex témpore. Non enim quǽreret quem fécerat, nisi fíeret ipse quod fécerat. Verum hoc mementóte, et de córdibus vestris nolíte dimíttere: sic esse Christum hóminem factum, ut non destíterit Deus esse. Manens Deus accépit hóminem, qui fecit hóminem.\par
\switchcolumn
\EN
\noindent In this chapter of the Gospel, my brethren, our Lord Jesus Christ hath much commended Himself unto our faith, as touching His Manhood. At the same time, His words and works were alway such as to give us to believe that He is both God and Man, yea, that God Who made us, and that Man Who hath sought us, yea, God the Son, Who, as touching His Godhead, is alway with the Father, (John i. 18; iii. 13,) and, as touching His Manhood, hath been with us in time. (Matth. i. 23.) For He had not sought the work of His hands unless He had been made His own work. (John i. 14.) Keep this well in mind, and let your hearts never forget it, namely, that Christ was not made Man so as to cease to be God. He, Who made the Manhood, took It into that Godhead Which is His from everlasting to everlasting.\par
\switchcolumn*

\LA
\begin{Resp}\Rsign Adjútor et suscéptor meus es tu, Dómine: et in verbum tuum sperávi: \Star{} Declináte a me, malígni: et scrutábor mandáta Dei mei.\end{Resp}
\begin{Resp}\Vsign Iníquos ódio hábui: et legem tuam diléxi.\end{Resp}
\begin{Resp}\Rsign Declináte a me, malígni: et scrutábor mandáta Dei mei.\end{Resp}
\switchcolumn
\EN
\begin{Resp}\Rsign Thou art my Helper and my Protector, O Lord, and in thy word do I hope. \Star{} Depart from me, ye evil doers, for I will keep the commandments of my God.\end{Resp}
\begin{Resp}\Vsign I hate the unrighteous, but thy law do I love.\end{Resp}
\begin{Resp}\Rsign Depart from me, ye evil doers, for I will keep the commandments of my God.\end{Resp}
\switchcolumn*

\LA
\LessonHead{Lectio II}
\switchcolumn
\EN
\LessonHead{Lesson II}
\switchcolumn*

\LA
\noindent Quando ergo látuit ut homo, non poténtiam perdidísse putándus est, sed exémplum infirmitáti præbuísse. Ille enim quando vóluit, deténtus est: quando vóluit, occísus est. Sed quóniam futúra erant membra ejus, id est, fidéles ejus, qui non habérent illam potestátem, quam habébat ipse Deus noster: quod latébat, quod se tamquam ne occiderétur, occultábat, hoc indicábat factúra esse membra sua, in quibus útique membris suis ipse erat.\par
\switchcolumn
\EN
\noindent While therefore He lay hid in the Manhood, we must not think that He had suffered any lessening of power, but that He was giving example to our weakness. When He willed it, He was taken; when He willed it, He was put to death. (John x. 18.) But, since He was to have members, that is, His faithful people, who would not have that power over their lives which He, our God, had over His, He hid Himself, He concealed Himself, as if it were to escape being put to death, to show what should be done by those His members in whom He should dwell.\par
\switchcolumn*

\LA
\begin{Resp}\Rsign Docébo iníquos vias tuas: et ímpii ad te converténtur: \Star{} Líbera me de sanguínibus, Deus, Deus salútis meæ.\end{Resp}
\begin{Resp}\Vsign Dómine, lábia mea apéries: et os meum annuntiábit laudem tuam.\end{Resp}
\begin{Resp}\Rsign Líbera me de sanguínibus, Deus, Deus salútis meæ.\end{Resp}
\switchcolumn
\EN
\begin{Resp}\Rsign I will teach transgressors thy ways, and sinners shall be converted unto thee. \Star{} Deliver me from blood-guiltiness, O God, Thou God of my salvation.\end{Resp}
\begin{Resp}\Vsign O Lord, open Thou my lips, and my mouth shall show forth thy praise.\end{Resp}
\begin{Resp}\Rsign Deliver me from blood-guiltiness, O God, Thou God of my salvation.\end{Resp}
\switchcolumn*

\LA
\LessonHead{Lectio III}
\switchcolumn
\EN
\LessonHead{Lesson III}
\switchcolumn*

\LA
\noindent Non enim Christus in cápite, et non in córpore: sed Christus totus in cápite, et in córpore. Quod ergo membra ejus, ipse: quod autem ipse, non contínuo membra ejus. Nam si non ipsi essent membra ejus, non díceret Saulo: Quid me perséqueris? Non enim Saulus ipsum, sed membra ejus, id est, fidéles ejus, in terra persequebátur. Nóluit tamen dícere sanctos meos, servos meos, postrémo honorabílius, fratres meos: sed me, hoc est membra mea, quibus ego sum caput.\par
\switchcolumn
\EN
\noindent Christ is not the Head of His Church in such sense that He is not in her Body; but the whole Christ is in the Head, and the whole Christ is in the Body. That, then, which His members are is Himself, though That Which He is, That are not therefore His members. For if His members had not been indeed His Own, how had He said unto Saul, (Acts ix. 4): Why persecutest thou Me? since Saul was not persecuting Him in Himself, but in His members, that is, in His faithful ones which were upon earth. He said not: Why persecutest thou My holy ones, nor: My servants, no, nor yet called He them by that more honourable name “My brethren”, but, “Why persecutest thou Me?”; that is, the members of My Body, whose Head I am.\par
\switchcolumn*

\LA
\begin{Resp}\Rsign Ne perdas cum ímpiis, Deus, ánimam meam, et cum viris sánguinum vitam meam: \Star{} Rédime me, Dómine.\end{Resp}
\begin{Resp}\Vsign Eripe me, Dómine, ab hómine malo, a viro iníquo líbera me.\end{Resp}
\begin{Resp}\Rsign Rédime me, Dómine.\end{Resp}
\begin{Resp}\Vsign Glória Patri, et Fílio, \Star{} et Spirítui Sancto.\end{Resp}
\begin{Resp}\Rsign Ne perdas cum ímpiis, Deus, ánimam meam, et cum viris sánguinum vitam meam: * Rédime me, Dómine.\end{Resp}
\switchcolumn
\EN
\begin{Resp}\Rsign Make not my soul to perish with sinners, O God, nor my life with bloody men. \Star{} Redeem me, O Lord!\end{Resp}
\begin{Resp}\Vsign Deliver me, O Lord, from the evil man, preserve me from the wicked man.\end{Resp}
\begin{Resp}\Rsign Redeem me, O Lord!\end{Resp}
\begin{Resp}\Vsign Glory be to the Father, and to the Son, \Star{} and to the Holy Ghost.\end{Resp}
\begin{Resp}\Rsign Make not my soul to perish with sinners, O God, nor my life with bloody men. Redeem me, O Lord!\end{Resp}
\switchcolumn*

\end{paracol}

\Office{Feria Quarta infra Hebdomadam Passionis}{Wednesday in Passion Week}{Feria major}
\addcontentsline{toc}{subsection}{Feria Quarta infra Hebdomadam Passionis\enspace\textit{Wednesday in Passion Week}}
\begin{paracol}{2}
\LA
\LessonHead{Lectio I}
\switchcolumn
\EN
\LessonHead{Lesson I}
\switchcolumn*

\LA
\LessonTitle{Léctio sancti Evangélii secúndum Joánnem}
\switchcolumn
\EN
\LessonTitle{Continuation of the Holy Gospel according to John}
\switchcolumn*

\LA
\Cite{Joannes 10:22-38}
\switchcolumn
\EN
\Cite{John 10:22-38}
\switchcolumn*

\LA
\noindent In illo témpore: Facta sunt encǽnia in Jerosólymis: et hiems erat. Et ambulábat Jesus in templo, in pórticu Salomónis. Et réliqua.\par
\switchcolumn
\EN
\noindent In that time, it was the feast of the dedication at Jerusalem: and it was winter. And Jesus walked in the temple, in Solomon's porch. And so on.\par
\switchcolumn*

\LA
\LessonTitle{Homilía sancti Augustíni Epíscopi}
\switchcolumn
\EN
\LessonTitle{Homily by St. Augustine, Bishop (of Hippo.)}
\switchcolumn*

\LA
\Source{Tractatus 48 in Joannem, circa init.}
\switchcolumn
\EN
\Source{42th Tract on John}
\switchcolumn*

\LA
\noindent Encǽnia festívitas erat dedicatiónis templi. Græce enim cænon dícitur novum. Quandocúmque novum áliquid fúerit dedicátum, encǽnia vocántur. Jam et usus habet hoc verbum. Si quis nova túnica induátur, encæniáre dícitur. Illum enim diem, quo templum dedicátum est, Judǽi solémniter celebrábant: ipse dies festus agebátur, cum ea quæ lecta sunt, locútus est Dóminus.\par
\switchcolumn
\EN
\noindent The Greek word Enkainia, used by the Evangelist, signifieth the Feast of the Dedication of the Temple. The derivation thereof is kainon, which is, being interpreted, new; and the Dedication of anything new is thence called Enkainia. The use of this word is still preserved among ourselves; if any man put on his new coat for the first time we say that he enkainiateth. It was the use of the Jews to keep solemn holiday upon the Anniversary of the Dedication of the Temple, and this was the Feast-day which was being observed when the Lord spake the words which have been read.\par
\switchcolumn*

\LA
\begin{Resp}\Rsign Tota die contristátus ingrediébar, Dómine: quóniam ánima mea compléta est illusiónibus: \Star{} Et vim faciébant, qui quærébant ánimam meam.\end{Resp}
\begin{Resp}\Vsign Amíci mei et próximi mei advérsum me appropinquavérunt et stetérunt: et qui juxta me erant, de longe stetérunt.\end{Resp}
\begin{Resp}\Rsign Et vim faciébant, qui quærébant ánimam meam.\end{Resp}
\switchcolumn
\EN
\begin{Resp}\Rsign O Lord, I go mourning all the day long, for my soul is filled with a loathsome disease \Star{} They also that sought after my life have used violence against me.\end{Resp}
\begin{Resp}\Vsign My friends and my neighbours draw near, and stand over against me; and they that are nearest to me stand afar off.\end{Resp}
\begin{Resp}\Rsign They also that sought after my life have used violence against me.\end{Resp}
\switchcolumn*

\LA
\LessonHead{Lectio II}
\switchcolumn
\EN
\LessonHead{Lesson II}
\switchcolumn*

\LA
\noindent Hiems erat, et ambulábat Jesus in templo, in pórticu Salomónis. Circumdedérunt ergo eum Judǽi, et dicébant ei: Quoúsque ánimam nostram tollis? Si tu es Christus, dic nobis palam. Non veritátem desiderábant, sed calúmniam præparábant. Hiems erat, et frígidi erant: ad illum enim divínum ignem accédere pigri erant. Si accédere est crédere: qui credit, accédit: qui negat, recédit. Non movétur ánima pédibus, sed afféctibus.\par
\switchcolumn
\EN
\noindent It was winter. And Jesus walked in the Temple in Solomon's Porch. Then came the Jews round about Him, and said unto Him: How long dost Thou make us to doubt? If Thou be the Christ, tell us plainly. They sought not to know the truth, but to have whereof to accuse Him. It was winter, and they were cold; for they were slow to draw near to God's fire. If to believe is to draw near thereto, then he which believeth draweth near thereto and he which denieth, goeth away therefrom. The soul is not moved by the feet but by the affections.\par
\switchcolumn*

\LA
\begin{Resp}\Rsign Ne avértas fáciem tuam a púero tuo, Dómine: \Star{} Quóniam tríbulor, velóciter exáudi me.\end{Resp}
\begin{Resp}\Vsign Inténde ánimæ meæ, et líbera eam: propter inimícos meos éripe me.\end{Resp}
\begin{Resp}\Rsign Quóniam tríbulor, velóciter exáudi me.\end{Resp}
\switchcolumn
\EN
\begin{Resp}\Rsign O Lord, hide not thy face from thy servant \Star{} For I am in trouble; hear me speedily.\end{Resp}
\begin{Resp}\Vsign Draw nigh unto my soul, and redeem it; deliver me, because of mine enemies.\end{Resp}
\begin{Resp}\Rsign For I am in trouble; hear me speedily.\end{Resp}
\switchcolumn*

\LA
\LessonHead{Lectio III}
\switchcolumn
\EN
\LessonHead{Lesson III}
\switchcolumn*

\LA
\noindent Frigúerant diligéndi caritáte, et ardébant nocéndi cupiditáte. Longe áberant, et ibi erant: non accedébant credéndo, et premébant persequéndo. Quærébant audíre a Dómino, Ego sum Christus: et fortásse de Christo secúndum hóminem sapiébant. Prædicavérunt enim prophétæ Christum: sed divinitátem Christi et in prophétis et in ipso Evangélio nec hærétici intéllegunt: quanto minus Judǽi, quámdiu velámen est super cor eórum?\par
\switchcolumn
\EN
\noindent They were frozen with want of love, and at the same time on fire with thirst to do injury. They stood afar off, and yet came near; for though they drew not near by faith, they were eager to persecute. They sought to hear the Lord say I am the Christ; and perchance they knew somewhat concerning Christ, as touching His Manhood, for the Prophets had prophesied of Christ. But the Godhead of Christ even some heretics do not see witnessed either in the Prophets or in the Gospel; how much less the Jews, as long as the veil is upon their heart. (2 Cor. iii. 15.)\par
\switchcolumn*

\LA
\begin{Resp}\Rsign Quis dabit cápiti meo aquam, et óculis meis fontem lacrimárum, et plorábo die ac nocte? quia frater propínquus supplantávit me, \Star{} Et omnis amícus fraudulénter incéssit in me.\end{Resp}
\begin{Resp}\Vsign Fiant viæ eórum ténebræ et lúbricum: et Angelus Dómini pérsequens eos.\end{Resp}
\begin{Resp}\Rsign Et omnis amícus fraudulénter incéssit in me.\end{Resp}
\begin{Resp}\Vsign Glória Patri, et Fílio, \Star{} et Spirítui Sancto.\end{Resp}
\begin{Resp}\Rsign Quis dabit cápiti meo aquam, et óculis meis fontem lacrimárum, et plorábo die ac nocte? quia frater propínquus supplantávit me, et omnis amícus fraudulénter incéssit in me.\end{Resp}
\switchcolumn
\EN
\begin{Resp}\Rsign O that my head were waters, and mine eyes a fountain of tears, that I might weep day and night! for my nearest brother hath supplanted me, \Star{} And my neighbour hath walked with slanders against me.\end{Resp}
\begin{Resp}\Vsign Let their way be dark and slippery, and let the Angel of the Lord persecute them.\end{Resp}
\begin{Resp}\Rsign And my neighbour hath walked with slanders against me.\end{Resp}
\begin{Resp}\Vsign Glory be to the Father, and to the Son, \Star{} and to the Holy Ghost.\end{Resp}
\begin{Resp}\Rsign O that my head were waters, and mine eyes a fountain of tears, that I might weep day and night! for my nearest brother hath supplanted me, and my neighbour hath walked with slanders against me.\end{Resp}
\switchcolumn*

\end{paracol}

\Office{Feria Quinta infra Hebdomadam Passionis}{Thursday in Passion Week}{Feria major}
\addcontentsline{toc}{subsection}{Feria Quinta infra Hebdomadam Passionis\enspace\textit{Thursday in Passion Week}}
\begin{paracol}{2}
\LA
\LessonHead{Lectio I}
\switchcolumn
\EN
\LessonHead{Lesson I}
\switchcolumn*

\LA
\LessonTitle{Léctio sancti Evangélii secúndum Lucam}
\switchcolumn
\EN
\LessonTitle{Continuation of the Holy Gospel according to Luke}
\switchcolumn*

\LA
\Cite{Luc 7:36-50}
\switchcolumn
\EN
\Cite{Luke 7:36-50}
\switchcolumn*

\LA
\noindent In illo témpore: Rogábat Jesum quidam de pharisǽis, ut manducáret cum illo. Et ingréssus domum pharisǽi discúbuit. Et réliqua.\par
\switchcolumn
\EN
\noindent In that time, one of the Pharisees desired him to eat with him. And he went into the house of the Pharisee, And so on.\par
\switchcolumn*

\LA
\LessonTitle{Homilía sancti Gregórii Papæ}
\switchcolumn
\EN
\LessonTitle{Homily by Pope St. Gregory (the Great.)}
\switchcolumn*

\LA
\Source{Homilia 33 in Evangelia}
\switchcolumn
\EN
\Source{33rd on the Gospels}
\switchcolumn*

\LA
\noindent Cogitánti mihi de Maríæ Magdalénæ pœniténtia, flere magis libet, quam áliquid dícere. Cujus enim vel sáxeum pectus illæ hujus peccatrícis lácrimæ ad exémplum pœniténdi non emólliant? Considerávit namque quid fecit, et nóluit moderári quid fáceret. Super convivántes ingréssa est, non jussa venit, inter épulas lácrimas óbtulit. Díscite, quo dolóre ardet, quæ flere et inter épulas non erubéscit.\par
\switchcolumn
\EN
\noindent When I think of the repentance of Mary Magdalene I feel nigher to weep than to say ought. Is there indeed any man, however stony his heart, who is not somewhat moved to follow the example of her repentance by the tears of that poor sinful woman? She weighed what she did, and would not that what she did should be niggardly. She came unbidden among the guests, and obtruded her tears upon the banquet. Ye may hence gather her sorrow, that she was content to weep at a feast.\par
\switchcolumn*

\LA
\begin{Resp}\Rsign Deus meus, éripe me de manu peccatóris: et de manu contra legem agéntis, et iníqui: \Star{} Quóniam tu es patiéntia mea.\end{Resp}
\begin{Resp}\Vsign Deus meus, ne elongéris a me: Deus meus, in auxílium meum réspice.\end{Resp}
\begin{Resp}\Rsign Quóniam tu es patiéntia mea.\end{Resp}
\switchcolumn
\EN
\begin{Resp}\Rsign Deliver me, O my God, out of the hand of the wicked, and out of the hand of the unrighteous and cruel man. \Star{} For Thou art my hope.\end{Resp}
\begin{Resp}\Vsign O my God, be not far from me O my God, make haste for my help.\end{Resp}
\begin{Resp}\Rsign For Thou art my hope.\end{Resp}
\switchcolumn*

\LA
\LessonHead{Lectio II}
\switchcolumn
\EN
\LessonHead{Lesson II}
\switchcolumn*

\LA
\noindent Hanc vero, quam Lucas peccatrícem mulíerem, Joánnes Maríam nóminat, illam esse Maríam crédimus, de qua Marcus septem dæmónia ejécta fuísse testátur. Et quid per septem dæmónia, nisi univérsa vítia designántur? Quia enim septem diébus omne tempus comprehénditur, recte septenário número univérsitas figurátur. Septem ergo dæmónia María hábuit, quæ univérsis vítiis plena fuit.\par
\switchcolumn
\EN
\noindent We believe that this woman, of whom Luke saith that she was a woman in the city, which was a sinner, and whom John nameth Mary, (xi. 2,) was the same as she of whom it is written in Mark (xvi. 9) that the Lord had cast out of her seven devils. And what signify seven devils but all manner of sin? For even as seven days do represent all time, so doth the number seven stand for all. Therefore is it said that Mary had seven devils, because she was full of all sin.\par
\switchcolumn*

\LA
\begin{Resp}\Rsign Multiplicáti sunt qui tríbulant me, et dicunt: Non est salus illi in Deo ejus: \Star{} Exsúrge, Dómine, salvum me fac, Deus meus.\end{Resp}
\begin{Resp}\Vsign Nequándo dicat inimícus meus: Præválui advérsus eum.\end{Resp}
\begin{Resp}\Rsign Exsúrge, Dómine, salvum me fac, Deus meus.\end{Resp}
\switchcolumn
\EN
\begin{Resp}\Rsign They be increased that trouble me, and that say There is no help for him in his God. \Star{} Arise, O Lord! Save me, O my God!\end{Resp}
\begin{Resp}\Vsign Lest mine enemy say I have prevailed against him.\end{Resp}
\begin{Resp}\Rsign Arise, O Lord! Save me, O my God!\end{Resp}
\switchcolumn*

\LA
\LessonHead{Lectio III}
\switchcolumn
\EN
\LessonHead{Lesson III}
\switchcolumn*

\LA
\noindent Sed ecce quia turpitúdinis suæ máculas aspéxit, lavánda ad fontem misericórdiæ cucúrrit, convivántes non erúbuit. Nam quia semetípsam gráviter erubescébat intus, nihil esse crédidit, quod verecundarétur foris. Quid ergo mirámur, fratres? Maríam veniéntem, an Dóminum suscipiéntem? Suscipiéntem dicam, an trahéntem? Sed mélius trahéntem dicam, et suscipiéntem: quia nimírum ipse eam per misericórdiam traxit intus, qui per mansuetúdinem suscépit foris.\par
\switchcolumn
\EN
\noindent But see how she realized the depth of her own filthiness, and came to be washed to the Well of Mercy, before all them which were bidden to the feast. The bitterness of her inward shame made her esteem it a light thing to be despised outwardly. At what then do we marvel, my brethren? That she came, or that the Lord welcomed her? Or would it be truer for me to say that He drew her to Him and welcomed her when she came? for His mercy inwardly drew her, and, when she came, His gentleness openly welcomed her.\par
\switchcolumn*

\LA
\begin{Resp}\Rsign Usquequo exaltábitur inimícus meus super me? \Star{} Réspice, et exáudi me, Dómine, Deus meus.\end{Resp}
\begin{Resp}\Vsign Qui tríbulant me, exsultábunt si motus fúero: ego autem in misericórdia tua sperábo.\end{Resp}
\begin{Resp}\Rsign Réspice, et exáudi me, Dómine, Deus meus.\end{Resp}
\begin{Resp}\Vsign Glória Patri, et Fílio, \Star{} et Spirítui Sancto.\end{Resp}
\begin{Resp}\Rsign Usquequo exaltábitur inimícus meus super me? * Réspice, et exáudi me, Dómine, Deus meus.\end{Resp}
\switchcolumn
\EN
\begin{Resp}\Rsign How long shall mine enemy be exalted over me? \Star{} Consider, and hear me, O Lord my God!\end{Resp}
\begin{Resp}\Vsign Those that trouble me will rejoice when I am moved but I have trusted in thy mercy.\end{Resp}
\begin{Resp}\Rsign Consider and hear me, O Lord my God!\end{Resp}
\begin{Resp}\Vsign Glory be to the Father, and to the Son, \Star{} and to the Holy Ghost.\end{Resp}
\begin{Resp}\Rsign How long shall mine enemy be exalted over me? Consider and hear me, O Lord my God!\end{Resp}
\switchcolumn*

\end{paracol}

\Office{Septem Dolorum Beatæ Mariæ Virginis}{Septem Dolorum Beatae Mariae Virginis}{Duplex majus}
\addcontentsline{toc}{subsection}{Septem Dolorum Beatæ Mariæ Virginis\enspace\textit{Septem Dolorum Beatae Mariae Virginis}}
\begin{paracol}{2}
\LA
\LessonHead{Lectio I}
\switchcolumn
\EN
\LessonHead{Lesson I}
\switchcolumn*

\LA
\LessonTitle{De Isaía Prophéta}
\switchcolumn
\EN
\LessonTitle{Lesson from the book of Isaias}
\switchcolumn*

\LA
\Cite{Isa 53:1-5}
\switchcolumn
\EN
\Cite{Isa 53:1-5}
\switchcolumn*

\LA
\noindent Quis crédidit audítui nostro? et bráchium Dómini cui revelátum est? Et ascéndet sicut virgúltum coram eo, et sicut radix de terra sitiénti. Non est spécies ei, neque decor, et vídimus eum, et non erat aspéctus, et desiderávimus eum: despéctum, et novíssimum virórum, virum dolórum, et sciéntem infirmitátem, et quasi abscónditus vultus ejus et despéctus, unde nec reputávimus eum. Vere languóres nostros ipse tulit, et dolóres nostros ipse portávit; et nos putávimus eum quasi leprósum, et percússum a Deo, et humiliátum. Ipse autem vulnerátus est propter iniquitátes nostras; attrítus est propter scélera nostra: disciplína pacis nostræ super eum, et livóre ejus sanáti sumus.\par
\switchcolumn
\EN
\noindent WHO hath believed our report? and to whom is the arm of the Lord revealed? And he shall grow up as a tender plant before him, and as a root out of a thirsty ground: there is no beauty in him, nor comeliness: and we have seen him, and there was no sightliness, that we should be desirous of him: Despised, and the most abject of men, a man of sorrows, and acquainted with infirmity: and his look was as it were hidden and despised, whereupon we esteemed him not. Surely he hath borne our infirmities and carried our sorrows: and we have thought him as it were a leper, and as one struck by God and afflicted. But he was wounded for our iniquities, he was bruised for our sins: the chastisement of our peace was upon him, and by his bruises we are healed.\par
\switchcolumn*

\LA
\begin{Resp}\Rsign Diléctus meus cándidus, et rubicúndus, et totus desiderábilis: \Star{} Omnis enim figúra ejus amórem spirat, et ad redamándum próvocat caput inclinátum, manus expánsæ, pectus apértum.\end{Resp}
\begin{Resp}\Vsign Piis, o Virgo, spectas eum óculis, contémplans in eo non tam vúlnerum livórem, quam mundi salútem.\end{Resp}
\begin{Resp}\Rsign Omnis enim figúra ejus amórem spirat, et ad redamándum próvocat caput inclinátum, manus expánsæ, pectus apértum.\end{Resp}
\switchcolumn
\EN
\begin{Resp}\Rsign My Beloved is white and ruddy, yea, He is altogether lovely; \Star{} For the sight of Him doth altogether breathe of love, and stirreth up to love in return; His Head is bowed down, His Hands are stretched out, and His Side is opened.\end{Resp}
\begin{Resp}\Vsign Maiden and Mother, thou didst look upon Him with eyes full of tenderness, and there thou sawest not only that thy Son was smitten, but that the world was saved.\end{Resp}
\begin{Resp}\Rsign For the sight of Him doth altogether breathe of love, and stirreth up to love in return; His Head is bowed down, His Hands are stretched out, and His Side is opened.\end{Resp}
\switchcolumn*

\LA
\LessonHead{Lectio II}
\switchcolumn
\EN
\LessonHead{Lesson II}
\switchcolumn*

\LA
\Cite{Isa 53:6-9}
\switchcolumn
\EN
\Cite{Isa 53:6-9}
\switchcolumn*

\LA
\noindent Omnes nos quasi oves errávimus, unusquísque in viam suam declinávit: et pósuit Dóminus in eo iniquitátem ómnium nostrum. Oblátus est quia ipse vóluit, et non apéruit os suum; sicut ovis ad occisiónem ducétur, et quasi agnus coram tondénte se obmutéscet, et non apériet os suum. De angústia, et de judício sublátus est. Generatiónem ejus quis enarrábit? quia abscíssus est de terra vivéntium: propter scelus pópuli mei percússi eum. Et dabit ímpios pro sepultúra, et dívitem pro morte sua, eo quod iniquitátem non fécerit, neque dolus fúerit in ore ejus.\par
\switchcolumn
\EN
\noindent All we like sheep have gone astray, every one hath turned aside into his own way: and the Lord hath laid on him the iniquity of us all. He was offered because it was his own will, and he opened not his mouth: he shall be led as a sheep to the slaughter, and shall be dumb as a lamb before his shearer, and he shall not open his mouth. He was taken away from distress, and from judgment: who shall declare his generation? because he is cut off out of the land of the living: for the wickedness of my people have I struck him. And he shall give the ungodly for his burial, and the rich for his death: because he hath done no iniquity, neither was there deceit in his mouth.\par
\switchcolumn*

\LA
\begin{Resp}\Rsign Manus ejus tornátiles, clavórum cúspide terebrátæ, \Star{} Humánæ salútis prétio quasi hyacínthis refértæ.\end{Resp}
\begin{Resp}\Vsign Córnua in mánibus ejus: ibi abscóndita est fortitúdo ejus: sunt enim manus ejus.\end{Resp}
\begin{Resp}\Rsign Humánæ salútis prétio quasi hyacínthis refértæ.\end{Resp}
\switchcolumn
\EN
\begin{Resp}\Rsign His hands are like rings, pierced with the points of the nails; \Star{} Set with price of man's salvation, as it were with jacinths.\end{Resp}
\begin{Resp}\Vsign He had horns coming out of His hands there was the hiding of His power for His Hands are\end{Resp}
\begin{Resp}\Rsign Set with the price of man's salvation, as it were with jacinths.\end{Resp}
\switchcolumn*

\LA
\LessonHead{Lectio III}
\switchcolumn
\EN
\LessonHead{Lesson III}
\switchcolumn*

\LA
\Cite{Isa 53:10-12}
\switchcolumn
\EN
\Cite{Isa 53:10-12}
\switchcolumn*

\LA
\noindent Et Dóminus vóluit contérere eum in infirmitáte. Si posúerit pro peccáto ánimam suam, vidébit semen longǽvum, et volúntas Dómini in manu ejus dirigétur. Pro eo quod laborávit ánima ejus, vidébit et saturábitur. In sciéntia sua justificábit ipse justus servus meus multos, et iniquitátes eórum ipse portábit. Ideo dispértiam ei plúrimos, et fórtium dívidet spólia, pro eo quod trádidit in mortem ánimam suam, et cum scelerátis reputátus est, et ipse peccáta multórum tulit, et pro transgressóribus rogávit.\par
\switchcolumn
\EN
\noindent And the Lord was pleased to bruise him in infirmity: if he shall lay down his life for sin, he shall see a long-lived seed, and the will of the Lord shall be prosperous in his hand. Because his soul hath laboured, he shall see and be filled: by his knowledge shall this my just servant justify many, and he shall bear their iniquities. Therefore will I distribute to him very many, and he shall divide the spoils of the strong, because he hath delivered his soul unto death, and was reputed with the wicked: and he hath borne the sins of many, and hath prayed for the transgressors.\par
\switchcolumn*

\LA
\begin{Resp}\Rsign Diligébat Jesus Joánnem, quóniam speciális prærogatíva castitátis amplióri dilectióne fécerat dignum: \Star{} Quia virgo eléctus ab ipso, virgo in ævum permánsit.\end{Resp}
\begin{Resp}\Vsign In cruce dénique moritúrus huic Matrem suam vírginem vírgini commendávit.\end{Resp}
\begin{Resp}\Rsign Quia virgo eléctus ab ipso, virgo in ævum permánsit.\end{Resp}
\begin{Resp}\Vsign Glória Patri, et Fílio, \Star{} et Spirítui Sancto.\end{Resp}
\begin{Resp}\Rsign Quia virgo eléctus ab ipso, virgo in ævum permánsit.\end{Resp}
\switchcolumn
\EN
\begin{Resp}\Rsign Jesus loved John because his singular gift of purity made him more worthy of love. \Star{} He chose him for a virgin unto Himself, and he remaineth a virgin for ever.\end{Resp}
\begin{Resp}\Vsign At the end, when He was dying upon the Cross, to him did He commit His mother, maiden to maiden.\end{Resp}
\begin{Resp}\Rsign He chose him for a virgin unto Himself, and he remaineth a virgin for ever.\end{Resp}
\begin{Resp}\Vsign Glory be to the Father, and to the Son, \Star{} and to the Holy Ghost.\end{Resp}
\begin{Resp}\Rsign He chose him for a virgin unto Himself, and he remaineth a virgin for ever.\end{Resp}
\switchcolumn*

\LA
\LessonHead{Lectio IV}
\switchcolumn
\EN
\LessonHead{Lesson IV}
\switchcolumn*

\LA
\LessonTitle{Sermo sancti Bernárdi Abbátis}
\switchcolumn
\EN
\LessonTitle{From the Sermons of St. Bernard, Abbot (of Clairvaux.)}
\switchcolumn*

\LA
\Source{Sermo de duodecim stellis}
\switchcolumn
\EN
\Source{On the twelve stars}
\switchcolumn*

\LA
\noindent Martýrium Vírginis tam in Simeónis prophetía, quam in ipsa Domínicæ passiónis história commendátur. Pósitus est hic (ait sanctus senex de púero Jesu) in signum cui contradicétur; et tuam ipsíus ánimam (ad Maríam autem dicébat) pertransíbit gládius. Vere tuam, o beáta Mater, ánimam pertransívit. Alióquin non nisi eam pertránsiens, carnem Fílii tui penetráret. Et quidem posteáquam emísit spíritum tuus ille Jesus, ipsíus plane non áttigit ánimam crudélis láncea, quæ ipsíus apéruit latus, sed tuam útique ánimam pertransívit. Ipsíus nimírum ánima jam ibi non erat, sed tua plane inde nequíbat avélli.\par
\switchcolumn
\EN
\noindent The Martyrdom of the Virgin is set before us, not only in the prophecy of Simeon, but also in the story itself of the Lord's Passion. The holy old man said of the Child Jesus (Luke ii. 34,) Behold, this Child is set for the fall and the rising again of many in Israel; and for a sign which shall be spoken against; yea, said he unto Mary, a sword shall pierce through thine own soul also Even so, O Blessed Mother! The sword did indeed pierce through thy soul! for nought could pierce the Body of thy Son, nor pierce thy soul likewise. Yea, and when this Jesus of thine had given up the ghost, and the bloody spear could torture Him no more, thy soul winced as it pierced His dead Side His Own Soul might leave Him, but thine could not.\par
\switchcolumn*

\LA
\begin{Resp}\Rsign Ténebræ factæ sunt, dum crucifixíssent Jesum Judǽi, et circa horam nonam exclamávit Jesus voce magna: Deus meus, ut quid dereliquísti me? \Star{} Et inclináto cápite emísit spíritum.\end{Resp}
\begin{Resp}\Vsign Quis tibi nunc sensus, dum cernis tália, Virgo?\end{Resp}
\begin{Resp}\Rsign Et inclináto cápite emísit spíritum.\end{Resp}
\switchcolumn
\EN
\begin{Resp}\Rsign The Jews crucified Jesus; and there was darkness; and about the ninth hour Jesus cried with a loud voice My God, why hast Thou forsaken Me? \Star{} And He bowed His Head, and gave up the Ghost.\end{Resp}
\begin{Resp}\Vsign O what a sickening at heart was thine at that moment, O Mother!\end{Resp}
\begin{Resp}\Rsign And He bowed His Head, and gave up the Ghost.\end{Resp}
\switchcolumn*

\LA
\LessonHead{Lectio V}
\switchcolumn
\EN
\LessonHead{Lesson V}
\switchcolumn*

\LA
\noindent Tuam ergo pertransívit ánimam vis dolóris, ut plusquam Mártyrem non immérito prædicémus, in qua nimírum corpóreæ sensum passiónis excésserit compassiónis afféctus. An non tibi plusquam gládius fuit sermo ille, revéra pertránsiens ánimam, et pertíngens usque ad divisiónem ánimæ et spíritus: Múlier, ecce fílius tuus? O commutatiónem! Joánnes tibi pro Jesu tráditur, servus pro Dómino, discípulus pro Magístro, fílius Zebedǽi pro Fílio Dei, homo purus pro Deo vero. Quómodo non tuam affectuosíssimam ánimam pertransíret hæc audítio, quando et nostra, licet sáxea, licet férrea péctora, sola recordátio scindit?\par
\switchcolumn
\EN
\noindent The sword of sorrow pierced through thy soul, so that we may truly call thee more than martyr, in whom the love, that made thee suffer along with thy Son, wrung thy heart more bitterly than any pang of bodily pain could do. Did not that word of His indeed pierce through thy soul, sharper than any two-edged sword, even to the dividing asunder of soul and spirit, (Heb. iv. 12,) Woman, behold thy son! (John xix. 26.) O what a change to thee! Thou art given John for Jesus, the servant for his Lord, the disciple for his master, the son of Zebedee for the Son of God, a mere man for Very God. O how keenly must the hearing of those words have pierced through thy most loving soul, when even our hearts, stony, iron, as they are, are wrung at the memory thereof only!\par
\switchcolumn*

\LA
\begin{Resp}\Rsign Pássio Dómini, \Star{} Ipsam ejus Matrem, carnáli orbitáte gráviter percússam, vehementíssime contristávit.\end{Resp}
\begin{Resp}\Vsign Ferrum lánceæ militáris, latus quidem Salvatóris, ánimam vero transívit Vírginis Matris.\end{Resp}
\begin{Resp}\Rsign Ipsam ejus Matrem, carnáli orbitáte gráviter percússam, vehementíssime contristávit.\end{Resp}
\switchcolumn
\EN
\begin{Resp}\Rsign The suffering of the Lord was \Star{} Asorrow exceeding sorrowful to her, His Mother, crushed by a natural bereavement.\end{Resp}
\begin{Resp}\Vsign The iron of the soldier's lance pierced through the Side of the Redeemer, and through the soul of the Virgin Mother.\end{Resp}
\begin{Resp}\Rsign A sorrow exceeding sorrowful to her, His Mother, crushed by a natural bereavement.\end{Resp}
\switchcolumn*

\LA
\LessonHead{Lectio VI}
\switchcolumn
\EN
\LessonHead{Lesson VI}
\switchcolumn*

\LA
\noindent Non mirémini, fratres, quod María Martyr in ánima fuísse dicátur. Mirétur qui non memínerit se audivísse Paulum inter máxima géntium crímina memorántem, quod sine affectióne fuíssent. Longe id fuit a Maríæ viscéribus, longe sit a sérvulis ejus. Sed forte quis dicat: Numquid non eum præscíerat moritúrum? Et indubitánter. Numquid non sperábat contínuo resurrectúrum? Et fidéliter. Super hæc dóluit crucifíxum? Et veheménter. Alióquin quisnam tu, frater, aut unde tibi hæc sapiéntia, ut miréris plus Maríæ Fílium patiéntem? Ille étiam mori córpore pótuit; ista cómmori corde non pótuit? Fecit illud cáritas, qua majórem nemo hábuit; fecit et hoc cáritas, cui post illam símilis áltera non fuit.\par
\switchcolumn
\EN
\noindent Marvel not, my brethren, that Mary should be called a Martyr in spirit. He indeed may marvel who remembereth not what Paul saith, naming the greater sins of the Gentiles, that they were without natural affection, (Rom. i. 31.) Far other were the bowels of Mary, and far other may those of her servants be! But some man perchance will say Did she not know that He was to die? Yea, without doubt, she knew it. Did she not hope that He was soon to rise again? Yea, she most faithfully hoped it. And did she still mourn because He was crucified? Yea, bitterly. But who art thou, my brother, or whence hast thou such wisdom, to marvel less that the Son of Mary suffered than that Mary suffered with Him? He could die in the Body, and could not she die with Him in her heart? His was the deed of that Love, greater than which hath no man (John xv. 13;) hers, of a love, like to which hath no man, save He.\par
\switchcolumn*

\LA
\begin{Resp}\Rsign Quis mihi det te fratrem meum sugéntem úbera matris meæ, et inhæréndo láteri tuo, ut sanguis tuus sánguinem meum tangat, et tergat: \Star{} Ut fons aquæ tuæ de scaturígine recti cordis, per venas boni óperis, in finem ætérnæ felicitátis exsíliat?\end{Resp}
\begin{Resp}\Vsign Fílii tui de longe vénient, et fíliæ tuæ de látere surgent.\end{Resp}
\begin{Resp}\Rsign Ut fons aquæ tuæ de scaturígine recti cordis, per venas boni óperis, in finem ætérnæ felicitátis exsíliat?\end{Resp}
\begin{Resp}\Vsign Glória Patri, et Fílio, \Star{} et Spirítui Sancto.\end{Resp}
\begin{Resp}\Rsign Ut fons aquæ tuæ de scaturígine recti cordis, per venas boni óperis, in finem ætérnæ felicitátis exsíliat?\end{Resp}
\switchcolumn
\EN
\begin{Resp}\Rsign O that Thou wert my brother, that sucked the breasts of my mother, that I might cleave unto thy Side, till thy Blood touched my blood, and cleansed it! \Star{} O that the Fountain of Water Which floweth from the Well-head of thy Righteous Heart, through thy Veins, Who hast done all things well, may at the last spring up for us into everlasting blessedness!\end{Resp}
\begin{Resp}\Vsign thy sons shall come from far, and thy daughters shall be nursed at thy Side.\end{Resp}
\begin{Resp}\Rsign O that the Fountain of Water Which floweth from the Wellhead of thy Righteous Heart, through thy Veins, Who hast done all things well, may at the last spring up for us into everlasting blessedness!\end{Resp}
\begin{Resp}\Vsign Glory be to the Father, and to the Son, \Star{} and to the Holy Ghost.\end{Resp}
\begin{Resp}\Rsign O that the Fountain of Water Which floweth from the Wellhead of thy Righteous Heart, through thy Veins, Who hast done all things well, may at the last spring up for us into everlasting blessedness!\end{Resp}
\switchcolumn*

\LA
\LessonHead{Lectio VII}
\switchcolumn
\EN
\LessonHead{Lesson VII}
\switchcolumn*

\LA
\LessonTitle{Léctio sancti Evangélii secúndum Joánnem.}
\switchcolumn
\EN
\LessonTitle{From the Holy Gospel according to John}
\switchcolumn*

\LA
\Cite{Joannes 19:25-27}
\switchcolumn
\EN
\Cite{John 19:25-27}
\switchcolumn*

\LA
\noindent In illo témpore: Stabant juxta crucem Jesu Mater ejus, et soror Matris ejus Maria Cleophæ, et Maria Magdalene. Et réliqua.\par
\switchcolumn
\EN
\noindent In that time stood by the cross of Jesus, his mother, and his mother's sister, Mary of Cleophas, and Mary Magdalen. And so on.\par
\switchcolumn*

\LA
\LessonTitle{Homilía sancti Augustíni Epíscopi.}
\switchcolumn
\EN
\LessonTitle{Homily by St. Augustine, Bishop (of Hippo.)}
\switchcolumn*

\LA
\Source{Tract. 119 in Joánnem.}
\switchcolumn
\EN
\Source{11th Tract on John.}
\switchcolumn*

\LA
\noindent Hæc nimírum est illa hora, de qua Jesus, aquam conversúrus in vinum, díxerat Matri: Quid mihi et tibi est múlier? nondum venit hora mea. Hanc ítaque horam prædíxerat, quæ tunc nondum vénerat, in qua debéret agnóscere moritúrus, de qua fúerat mortáliter natus. Tunc ergo divína factúrus, non divinitátis, sed infirmitátis matrem velut incógnitam repellébat. nunc autem humána jam pátiens, ex qua fúerat factus homo, afféctu commendábat humáno. Morális ígitur insinuátur locus. Facit quod faciéndum ádmonet, et exémplo suo suos minístros instrúxit præcéptor bonus, ut a fíliis piis impendátur cura paréntibus: tamquam lignum illud, ubi fixa erant membra moriéntis, étiam cáthedra fúerit magístri docéntis.\par
\switchcolumn
\EN
\noindent This is that hour whereof Jesus, when He was about to turn water into wine, had said unto His Mother Woman, what have I to do with thee? Mine hour is not yet come. (John ii. 4.) He had spoken of this hour, which then was not yet come, wherein, being about to die, it should be His duty to acknowledge her of whom He had been born in a dying Body. Then, since He was about to work the works of God, He thrust from Him, as though He knew her not, her who was His Mother, not in that nature as touching which He is equal to the Father, but in that as touching which He is inferior to the Father. But now, since He is suffering the pains of Man, He careth, with a Man's love, for her of whom He hath been made Man. And herein He giveth us a lesson. He doth that which He would have us to do. The Good Master, by His Own example, commandeth that among His disciples, dutiful children should succour their parents, as though even that Tree whereupon His dying Limbs were nailed, even that Tree were to be a pulpit for His teaching.\par
\switchcolumn*

\LA
\begin{Resp}\Rsign Dóleo super te fili mi Jesu, decórus nimis, et amábilis super amórem mulíerum! \Star{} Sicut enim mater únicum díligit fílium, ita ego te diligébam.\end{Resp}
\begin{Resp}\Vsign Defécit in dolóre vita mea, et anni mei in gemítibus.\end{Resp}
\begin{Resp}\Rsign Sicut enim mater únicum díligit fílium, ita ego te diligébam.\end{Resp}
\switchcolumn
\EN
\begin{Resp}\Rsign I am distressed for thee, my Son Jesus, very pleasant hast Thou been unto me; thy love to me was wonderful, passing the love of women; \Star{} For even as a mother loveth her only Son, so loved I thee.\end{Resp}
\begin{Resp}\Vsign My life is spent with grief, and my years with sighing.\end{Resp}
\begin{Resp}\Rsign For even as a mother loveth her only Son, so loved I thee.\end{Resp}
\switchcolumn*

\LA
\LessonHead{Lectio VIII}
\switchcolumn
\EN
\LessonHead{Lesson VIII}
\switchcolumn*

\LA
\noindent Ex hac sana doctrína didícerat Paulus Apóstolus quod docébat quando dicébat: Si quis autem suis, et máxime domésticis non próvidet, fidem negávit, et est infidéli detérior, Quid autem tam cuíque domésticum, quam paréntes fíliis, aut paréntibus fílii? Hujus ítaque salubérrimi præcépti ipse Magíster Sanctórum de seípso constituébat exémplum: quando non ut fámulæ Deus, quam creáverat, et regébat; sed ut matri homo, de qua creátus fúerat, et quam relinquébat, álterum pro se quodámmodo fílium providébat.\par
\switchcolumn
\EN
\noindent And of this teaching by Jesus Crucified cometh that which the Apostle Paul commandeth, where he saith, (1 Tim. v. 8) If any provide not for his own, and specially for those of his own house, he hath denied the faith, and is worse than an infidel. But what is so much of a man's own house, as children are of their parents'? and parents of their children's? Of this most healthy law the Master of the Saints was pleased Himself to give an example, when, being God, He treated not as His handmaid her of whom He was the Maker and the Lord, but, being also Man, gave another to be as a son in His stead, to her of whom as Man He had been made, and whom He was leaving.\par
\switchcolumn*

\LA
\begin{Resp}\Rsign Eja, Mater fons amóris, fac nos sentíre vim dolóris, ut tecum lugeámus: \Star{} Et Domínicæ passiónis fructum sentiámus.\end{Resp}
\begin{Resp}\Vsign Ut sicut Fílius tuus Jesus pro nobis mórtuus est, et resurréxit; ita et nos commórtui cum eódem resurgámus.\end{Resp}
\begin{Resp}\Rsign Et Domínicæ passiónis fructum sentiámus.\end{Resp}
\switchcolumn
\EN
\begin{Resp}\Rsign Fount of love and holy sorrow, Mother! may my spirit borrow Somewhat of thy woe profound; \Star{} Unto Christ, with pure emotion, Raise my contrite heart's devotion, Love to read in every Wound.\end{Resp}
\begin{Resp}\Vsign That as thy Son Jesus for our sakes died and rose again, so we also who have died with Him may rise again with Him.\end{Resp}
\begin{Resp}\Rsign Unto Christ, with pure emotion, Raise my contrite heart's devotion, Love to read in every Wound.\end{Resp}
\switchcolumn*

\LA
\LessonHead{Lectio IX}
\switchcolumn
\EN
\LessonHead{Lesson IX}
\switchcolumn*

\LA
\Cite{Joannes 11:47-54}
\switchcolumn
\EN
\Cite{John 11:47-54}
\switchcolumn*

\LA
\Source{de Homilía Feriæ.}
\switchcolumn
\EN
\Source{Homily on the Feria}
\switchcolumn*

\LA
\noindent Léctio sancti Evangélii secúndum Joánnem\par
\switchcolumn
\EN
\noindent From the holy Gospel according to John.\par
\switchcolumn*

\LA
In illo témpore: Collegérunt pontífices et pharisǽi concílium advérsus Jesum, et dicébant: Quid fácimus, quia hic homo multa signa facit? Et réliqua.\par
\switchcolumn
\EN
The chief priests therefore, and the Pharisees, gathered a council, and said: What do we, for this man doth many miracles? If we let him alone so, all will believe in him; and the Romans will come, and take away our place and nation. And so on.\par
\switchcolumn*

\LA
\LessonTitle{Homilía sancti Augustíni Epíscopi}
\switchcolumn
\EN
\LessonTitle{Homily by St. Augustine, Bishop (of Hippo.)}
\switchcolumn*

\LA
\Source{Tract. 49 in Joannem, sub finem}
\switchcolumn
\EN
\Source{49th Tract on John.}
\switchcolumn*

\LA
\noindent Pontífices et pharisǽi sibi consulébant: nec tamen dicébant: Credámus. Plus enim pérditi hómines cogitábant, quómodo nocérent, ut pérderent, quam quómodo sibi consúlerent, ne perírent: et tamen timébant, et quasi consulébant. Dicébant enim: Quid fácimus, quia hic homo multa signa facit? Si dimíttimus eum sic, omnes credent in eum: et vénient Románi, et tollent nostrum locum et gentem. Temporália pérdere timuérunt, et vitam ætérnam non cogitavérunt, ac sic utrúmque amisérunt.\par
\switchcolumn
\EN
\noindent The chief Priests and the Pharisees took counsel together, but Let us believe in Him was not one of the suggestions offered. Those lost creatures thought much more how they might hurt and undo Him, than how they might save themselves from perishing. And yet they were afraid, and took counsel together, and said What do we? For this Man doeth many miracles. If we let Him thus alone, all men will believe on Him; and the Romans shall come and take away both our place and our nation. They were afraid of losing temporal things, but they gave no thought to eternal life, and so they lost both.\par
\switchcolumn*

\end{paracol}

\Office{Sabbato infra Hebdomadam Passionis}{Saturday in Passion Week}{Feria major}
\addcontentsline{toc}{subsection}{Sabbato infra Hebdomadam Passionis\enspace\textit{Saturday in Passion Week}}
\begin{paracol}{2}
\LA
\LessonHead{Lectio I}
\switchcolumn
\EN
\LessonHead{Lesson I}
\switchcolumn*

\LA
\LessonTitle{Léctio sancti Evangélii secúndum Joánnem}
\switchcolumn
\EN
\LessonTitle{Continuation of the Holy Gospel according to John}
\switchcolumn*

\LA
\Cite{Joannes 12:10-36}
\switchcolumn
\EN
\Cite{John 12:10-36}
\switchcolumn*

\LA
\noindent In illo témpore: Cogitavérunt príncipes sacerdótum ut et Lázarum interfícerent: quia multi propter illum abíbant ex Judǽis, et credébant in Jesum. Et réliqua.\par
\switchcolumn
\EN
\noindent In that time, the chief priests thought to kill Lazarus also: Because many of the Jews, by reason of him, went away, and believed in Jesus. And so on.\par
\switchcolumn*

\LA
\LessonTitle{Homilía sancti Augustíni Epíscopi}
\switchcolumn
\EN
\LessonTitle{Homily by St. Augustine, Bishop (of Hippo)}
\switchcolumn*

\LA
\Source{Tractatus 50 in Joannem, in fine}
\switchcolumn
\EN
\Source{50th Tract on John}
\switchcolumn*

\LA
\noindent Viso Lázaro resuscitáto, quia tantum miráculum Dómini tanta erat evidéntia diffamátum, tanta manifestatióne declarátum, ut non possent vel occultáre quod factum est, vel negáre: quid invenérunt, vidéte. Cogitavérunt autem príncipes sacerdótum ut et Lázarum interfícerent. O stulta cogitátio, et cæca sævítia! Dóminus Christus, qui suscitáre pótuit mórtuum, non posset occísum! Quando Lázaro inferebátis necem, numquid auferebátis Dómino potestátem? Si áliud vobis vidétur mórtuus, áliud occísus: ecce Dóminus utrúmque fecit, et Lázarum mórtuum, et seípsum suscitávit occísum.\par
\switchcolumn
\EN
\noindent When they saw Lazarus who had been raised from the dead, and knew that the miracle which the Lord had worked was so great, spread about by so many witnesses, and so plain and manifest that it could neither be concealed nor denied, they invented an expedient; and see here what it was: “But the chief Priests consulted that they might put Lazarus also to death.” What stupidity of thought, what blindness of cruelty is here! If the Lord Christ had raised up again a man who had died a natural death, could He not also raise up one that had died by violence? Would killing Lazarus paralyse the Lord? But if ye consider that there is a difference between a man dead of disease, and a man killed, behold, the Lord hath raised up both for He first raised up Lazarus, who had died a natural death, and then Himself, after a violent one.\par
\switchcolumn*

\LA
\begin{Resp}\Rsign Tota die contristátus ingrediébar, Dómine: quóniam ánima mea compléta est illusiónibus: \Star{} Et vim faciébant, qui quærébant ánimam meam.\end{Resp}
\begin{Resp}\Vsign Amíci mei et próximi mei advérsum me appropinquavérunt et stetérunt: et qui juxta me erant, de longe stetérunt.\end{Resp}
\begin{Resp}\Rsign Et vim faciébant, qui quærébant ánimam meam.\end{Resp}
\switchcolumn
\EN
\begin{Resp}\Rsign O Lord, I go mourning all the day long, for my soul is filled with a loathsome disease \Star{} They also that sought after my life have used violence against me.\end{Resp}
\begin{Resp}\Vsign My friends and my neighbours draw near, and stand over against me; and they that are nearest to me stand afar off.\end{Resp}
\begin{Resp}\Rsign They also that sought after my life have used violence against me.\end{Resp}
\switchcolumn*

\LA
\LessonHead{Lectio II}
\switchcolumn
\EN
\LessonHead{Lesson II}
\switchcolumn*

\LA

\switchcolumn
\EN
\LessonTitle{“On the next day much people that were come to the feast, when they heard that Jesus was coming to Jerusalem, took branches of palm-trees, and went forth to meet Him, and cried Hosanna! Blessed is the King of Israel That cometh in the Name of the Lord!” Palm branches are glorious boughs which tell of victory; yea, the Lord was now ready by His Own Death to trample down death, and to carry the victorious banner of His Cross in triumph over the devil, the prince of death. The cry with which He was greeted, namely “Hosanna,” hath not, as we are assured by some who are acquainted with the Hebrew language, any meaning in particular, but is a shout after the manner of interjections, as they are called, just as in Latin when we lament we say “Heu,” or when we are pleased, “Vah.”}
\switchcolumn*

\LA
\noindent In crástinum autem turba multa, quæ vénerat ad diem festum, cum audíssent quia venit Jesus Jerosólymam: accepérunt ramos palmárum, et processérunt óbviam ei, et clamábant: Hosánna, benedíctus qui venit in nómine Dómini, Rex Israël. Rami palmárum laudes sunt, significántes victóriam: quia erat Dóminus mortem moriéndo superatúrus, et trophǽo crucis de diábolo mortis príncipe triumphatúrus. Vox autem obsecrántis est Hosánna, sicut nonnúlli dicunt, qui Hebrǽam linguam novérunt, magis afféctum índicans, quam rem áliquam signíficans, sicut sunt in lingua Latína, quas interjectiónes vocant: velut cum doléntes dícimus, heu; vel cum delectámur, vah dícimus.\par
\switchcolumn
\EN

\switchcolumn*

\LA
\begin{Resp}\Rsign Ne avértas fáciem tuam a púero tuo, Dómine: \Star{} Quóniam tríbulor, velóciter exáudi me.\end{Resp}
\begin{Resp}\Vsign Inténde ánimæ meæ, et líbera eam: propter inimícos meos éripe me.\end{Resp}
\begin{Resp}\Rsign Quóniam tríbulor, velóciter exáudi me.\end{Resp}
\switchcolumn
\EN
\begin{Resp}\Rsign O Lord, hide not thy face from thy servant \Star{} For I am in trouble; hear me speedily.\end{Resp}
\begin{Resp}\Vsign Draw nigh unto my soul, and redeem it; deliver me, because of mine enemies.\end{Resp}
\begin{Resp}\Rsign For I am in trouble; hear me speedily.\end{Resp}
\switchcolumn*

\LA
\LessonHead{Lectio III}
\switchcolumn
\EN
\LessonHead{Lesson III}
\switchcolumn*

\LA
\noindent Has ei laudes turba dicébat: Hosánna, benedíctus, qui venit in nómine Dómini, Rex Israël. Quam crucem mentis invidéntia príncipum Judæórum pérpeti potúerat, quando Regem suum Christum tanta multitúdo clamábat? Sed quid fuit Dómino Regem esse Israël? Quid magnum fuit Regi sæculórum, Regem fíeri hóminum? Non enim Rex Israël Christus ad exigéndum tribútum, vel exércitum ferro armándum, hostésque visibíliter debellándos: sed Rex Israël, quod mentes regat, quod in ætérnum cónsulat, quod in regnum cælórum credéntes, sperántes, amantésque perdúcat.\par
\switchcolumn
\EN
\noindent These were the shouts of applause with which the crowd greeted Him, “Hosanna! Blessed is the King of Israel That cometh in the Name of the Lord!” What inward torture must the jealousy of the Jewish leaders have caused them, when they heard that great multitude hailing Christ as their King! But, for the Lord, what was it to be King of Israel? To the Eternal King what mattered it to become a King of men? And Christ is not King of Israel in the sense of monarchs who exact tribute, or arm hosts with steel to conquer enemies that are seen. But King of Israel He is, as He Who is Lord of our intellect, a Ruler Whose power shall never wane, and Who openeth a Kingdom in heaven to all such as centre in Him their faith, their hope, and their love.\par
\switchcolumn*

\LA
\begin{Resp}\Rsign Quis dabit cápiti meo aquam, et óculis meis fontem lacrimárum, et plorábo die ac nocte? quia frater propínquus supplantávit me, \Star{} Et omnis amícus fraudulénter incéssit in me.\end{Resp}
\begin{Resp}\Vsign Fiant viæ eórum ténebræ et lúbricum: et Angelus Dómini pérsequens eos.\end{Resp}
\begin{Resp}\Rsign Et omnis amícus fraudulénter incéssit in me.\end{Resp}
\begin{Resp}\Vsign Glória Patri, et Fílio, \Star{} et Spirítui Sancto.\end{Resp}
\begin{Resp}\Rsign Quis dabit cápiti meo aquam, et óculis meis fontem lacrimárum, et plorábo die ac nocte? quia frater propínquus supplantávit me, et omnis amícus fraudulénter incéssit in me.\end{Resp}
\switchcolumn
\EN
\begin{Resp}\Rsign O that my head were waters, and mine eyes a fountain of tears, that I might weep day and night! for my nearest brother hath supplanted me, \Star{} And my neighbour hath walked with slanders against me.\end{Resp}
\begin{Resp}\Vsign Let their way be dark and slippery, and let the Angel of the Lord persecute them.\end{Resp}
\begin{Resp}\Rsign And my neighbour hath walked with slanders against me.\end{Resp}
\begin{Resp}\Vsign Glory be to the Father, and to the Son, \Star{} and to the Holy Ghost.\end{Resp}
\begin{Resp}\Rsign O that my head were waters, and mine eyes a fountain of tears, that I might weep day and night! for my nearest brother hath supplanted me, and my neighbour hath walked with slanders against me.\end{Resp}
\switchcolumn*

\end{paracol}

\Office{Dominica in Palmis}{Palm Sunday}{Semiduplex I classis}
\addcontentsline{toc}{subsection}{Dominica in Palmis\enspace\textit{Palm Sunday}}
\begin{paracol}{2}
\LA
\LessonHead{Lectio I}
\switchcolumn
\EN
\LessonHead{Lesson I}
\switchcolumn*

\LA
\LessonTitle{De Jeremía Prophéta}
\switchcolumn
\EN
\LessonTitle{Lesson from the book of Jeremias}
\switchcolumn*

\LA
\Cite{Jer 2:12-17}
\switchcolumn
\EN
\Cite{Jer 2:12-17}
\switchcolumn*

\LA
\noindent Obstupéscite, cæli, super hoc, et, portæ ejus, desolámini veheménter, dicit Dóminus. Duo enim mala fecit pópulus meus: Me dereliquérunt fontem aquæ vivæ, et fodérunt sibi cistérnas, cistérnas dissipátas, quæ continére non valent aquas. Numquid servus est Israël, aut vernáculus? Quare ergo factus est in prædam? Super eum rugiérunt leónes, et dedérunt vocem suam, posuérunt terram ejus in solitúdinem: civitátes ejus exústæ sunt, et non est qui hábitet in eis. Fílii quoque Mémpheos et Taphnes constupravérunt te usque ad vérticem. Numquid non istud factum est tibi, quia dereliquísti Dóminum Deum tuum eo témpore, quo ducébat te per viam?\par
\switchcolumn
\EN
\noindent Be astonished, O ye heavens, at this, and ye gates thereof, be very desolate, saith the Lord. For my people have done two evils. They have forsaken me, the fountain of living water, and have digged to themselves cisterns, broken cisterns, that can hold no water. Is Israel a bondman, or a homeborn slave? why then is he become prey? The lions have roared upon him, and have made a noise, they have made his land a wilderness: his cities are burnt down and there is none to dwell in them. The children also of Memphis, and of Taphnes have deflowered thee, even to the crown of the head. Hath not this been done to thee, because thou hast forsaken the Lord thy God at that time, when he led thee by the way?\par
\switchcolumn*

\LA
\begin{Resp}\Rsign In die qua invocávi te, Dómine, dixísti: Noli timére: \Star{} Judicásti causam meam, et liberásti me, Dómine, Deus meus.\end{Resp}
\begin{Resp}\Vsign In die tribulatiónis meæ clamávi ad te, quia exaudísti me.\end{Resp}
\begin{Resp}\Rsign Judicásti causam meam, et liberásti me, Dómine, Deus meus.\end{Resp}
\switchcolumn
\EN
\begin{Resp}\Rsign O Lord, in the day that I called upon thee, Thou saidst: Fear not. \Star{} Thou hast pleaded my cause, and hast redeemed me, O Lord my God.\end{Resp}
\begin{Resp}\Vsign In the day of my trouble I called upon thee, for Thou hast heard me.\end{Resp}
\begin{Resp}\Rsign Thou hast pleaded my cause, and hast redeemed me, O Lord my God.\end{Resp}
\switchcolumn*

\LA
\LessonHead{Lectio II}
\switchcolumn
\EN
\LessonHead{Lesson II}
\switchcolumn*

\LA
\Cite{Jer 2:18-22}
\switchcolumn
\EN
\Cite{Jer 2:18-22}
\switchcolumn*

\LA
\noindent Et nunc quid tibi vis in via Ægýpti, ut bibas aquam túrbidam? et quid tibi cum via Assyriórum, ut bibas aquam flúminis? Arguet te malítia tua, et avérsio tua increpábit te. Scito, et vide quia malum et amárum est reliquísse te Dóminum Deum tuum, et non esse timórem mei apud te, dicit Dóminus Deus exercítuum. A sǽculo confregísti jugum meum, rupísti víncula mea, et dixísti: Non sérviam. In omni enim colle sublími, et sub omni ligno frondóso, tu prosternebáris méretrix. Ego autem plantávi te víneam eléctam, omne semen verum: quómodo ergo convérsa es mihi in pravum, vínea aliéna? Si láveris te nitro, et multiplicáveris tibi herbam borith, maculáta es in iniquitáte tua coram me, dicit Dóminus Deus.\par
\switchcolumn
\EN
\noindent And now what hast thou to do in the way of Egypt, to drink the troubled water? And what hast thou to do with the way of the Assyrians, to drink the water of the river? thy own wickedness shall reprove thee, and thy apostasy shall rebuke thee. Know thou, and see that it is an evil and a bitter thing for thee, to have left the Lord thy God, and that my fear is not with thee, saith the Lord the God of hosts. Of old time thou hast broken my yoke, thou hast burst my bands, and thou saidst: I will not serve. For on every high hill, and under every green tree thou didst prostitute thyself. Yet I planted thee a chosen vineyard, all true seed: how then art thou turned unto me into that which is good for nothing, O strange vineyard? Though thou wash thyself with nitre, and multiply to thyself the herb borith, thou art stained in thy iniquity before me, saith the Lord God.\par
\switchcolumn*

\LA
\begin{Resp}\Rsign Fratres mei elongavérunt se a me: et noti mei \Star{} Quasi aliéni recessérunt a me.\end{Resp}
\begin{Resp}\Vsign Dereliquérunt me próximi mei, et qui me novérunt.\end{Resp}
\begin{Resp}\Rsign Quasi aliéni recessérunt a me.\end{Resp}
\switchcolumn
\EN
\begin{Resp}\Rsign My brethren stand afar off from me, and they which have known me \Star{} Make themselves strange unto me, and leave me.\end{Resp}
\begin{Resp}\Vsign My neighbours forsake me, and mine acquaintance\end{Resp}
\begin{Resp}\Rsign Make themselves strange unto me, and leave me me.\end{Resp}
\switchcolumn*

\LA
\LessonHead{Lectio III}
\switchcolumn
\EN
\LessonHead{Lesson III}
\switchcolumn*

\LA
\Cite{Jer 2:29-32}
\switchcolumn
\EN
\Cite{Jer 2:29-32}
\switchcolumn*

\LA
\noindent Quid vultis mecum judício conténdere? Omnes dereliquístis me, dicit Dóminus. Frustra percússi fílios vestros, disciplínam non recepérunt: devorávit gládius vester prophétas vestros, quasi leo vastátor generátio vestra. Vidéte verbum Dómini: Numquid solitúdo factus sum Israéli, aut terra serótina? Quare ergo dixit pópulus meus: Recéssimus, non veniémus ultra ad te? Numquid obliviscétur virgo ornaménti sui, aut sponsa fásciæ pectorális suæ? pópulus vero meus oblítus est mei diébus innúmeris.\par
\switchcolumn
\EN
\noindent Why will you contend with me in judgement? you have all forsaken me, saith the Lord. In vain have I struck your children, they have not received correction: your sword hath devoured your prophets, your generation is like a ravaging lion. See ye the word of the Lord: Am I become a wilderness to Israel, or a lateward springing land? why then have my people said: We are revolted, we will come to thee no more. Will a virgin forget her ornament, or a bride her stomacher? but my people hath forgotten me days without number.\par
\switchcolumn*

\LA
\begin{Resp}\Rsign Atténde, Dómine, ad me, et audi voces adversariórum meórum: \Star{} Numquid rédditur pro bono malum, quia fodérunt fóveam ánimæ meæ?\end{Resp}
\begin{Resp}\Vsign Recordáre quod stéterim in conspéctu tuo, ut lóquerer pro eis bonum, et avérterem indignatiónem tuam ab eis.\end{Resp}
\begin{Resp}\Rsign Numquid rédditur pro bono malum, quia fodérunt fóveam ánimæ meæ?\end{Resp}
\begin{Resp}\Vsign Glória Patri, et Fílio, \Star{} et Spirítui Sancto.\end{Resp}
\begin{Resp}\Rsign Atténde, Dómine, ad me, et audi voces adversariórum meórum: * Numquid rédditur pro bono malum, quia fodérunt fóveam ánimæ meæ?\end{Resp}
\switchcolumn
\EN
\begin{Resp}\Rsign Give heed to me, O Lord, and hearken to the voice of them that contend with me. \Star{} Shall evil be recompensed for good? for they have digged a pit for my soul.\end{Resp}
\begin{Resp}\Vsign Remember that I stood before thee to speak good for them, and to turn away thy wrath from them.\end{Resp}
\begin{Resp}\Rsign Shall evil be recompensed for good? for they have digged a pit for my soul.\end{Resp}
\begin{Resp}\Vsign Glory be to the Father, and to the Son, \Star{} and to the Holy Ghost.\end{Resp}
\begin{Resp}\Rsign Give heed to me, O Lord, and hearken to the voice of them that contend with me. * Shall evil be recompensed for good? for they have digged a pit for my soul.\end{Resp}
\switchcolumn*

\LA
\LessonHead{Lectio IV}
\switchcolumn
\EN
\LessonHead{Lesson IV}
\switchcolumn*

\LA
\LessonTitle{Sermo sancti Leónis Papæ}
\switchcolumn
\EN
\LessonTitle{From the Sermons of Pope St. Leo (the Great.)}
\switchcolumn*

\LA
\Source{Sermo 11 de Passióne Dómini}
\switchcolumn
\EN
\Source{Second on the Passion of the Lord.}
\switchcolumn*

\LA
\noindent Desideráta nobis, dilectíssimi, et univérso optábilis mundo adest festívitas Domínicæ passiónis, quæ nos inter exsultatiónes spirituálium gaudiórum silére non pátitur. Quia etsi diffícile est, de eádem solemnitáte sǽpius digne aptéque dissérere: non est tamen líberum sacerdóti in tanto divínæ misericórdiæ sacraménto fidélibus pópulis subtráhere sermónis offícium: cum ipsa matéria ex eo quod est ineffábilis, fandi tríbuat facultátem: nec possit defícere quod dicátur, dum numquam potest satis esse quod dícitur. Succúmbat ergo humána infírmitas glóriæ Dei, et in explicándis opéribus misericórdiæ ejus, ímparem se semper invéniat. Laborémus sensu, hæreámus ingénio, deficiámus elóquio: bonum est ut nobis parum sit, quod étiam recte de Dómini majestáte sentímus.\par
\switchcolumn
\EN
\noindent Dearly beloved brethren, the jubilant and triumphal day which ushereth in the commemoration of the Lord's Passion is come; even that day for which we have longed so much, and for whose yearly coming the whole world may well look. Shouts of spiritual exultation are ringing, and suffer not that we should be silent. It is indeed hard to preach often on the same Festival, and that always meetly and rightly, but a Priest is not free, when we celebrate so great and mysterious an out-pouring of God's mercy, to leave his faithful people without the service of a discourse. Nay, that his subject-matter is unspeakable should in itself make him eloquent, since where enough can never be said, there must needs ever be somewhat to say. Let man's weakness, then, fall down before the glory of God, and acknowledge herself ever too feeble to unfold all the works of His mercy. We may jade our emotions, break down in our understanding, and fail in our speech it is good for us, that even what we truly feel in presence of the Divine Majesty is but little, (compared to the vastness of the subject.)\par
\switchcolumn*

\LA
\begin{Resp}\Rsign Conclúsit vias meas inimícus, insidiátor factus est mihi sicut leo in abscóndito, replévit et inebriávit me amaritúdine: deduxérunt in lacum mortis vitam meam, et posuérunt lápidem contra me. \Star{} Vide, Dómine, iniquitátes illórum: et júdica causam ánimæ meæ, defénsor vitæ meæ.\end{Resp}
\begin{Resp}\Vsign Factus sum in derísum omni pópulo meo, cánticum eórum tota die.\end{Resp}
\begin{Resp}\Rsign Vide, Dómine, iniquitátes illórum: et júdica causam ánimæ meæ, defénsor vitæ meæ.\end{Resp}
\switchcolumn
\EN
\begin{Resp}\Rsign The enemy hath enclosed my ways; he lay in wait for me as a lion; in secret places he hath filled me, and made me drunken with bitterness; they have cut off my life in the dungeon, and cast a stone upon me. \Star{} O Lord, behold all their iniquity, and plead the cause of my soul, Thou That art the Redeemer of my life!\end{Resp}
\begin{Resp}\Vsign I was a derision to all my people, and their song all the day.\end{Resp}
\begin{Resp}\Rsign O Lord, behold all their iniquity, and plead the cause of my soul, Thou That art the Redeemer of my life!\end{Resp}
\switchcolumn*

\LA
\LessonHead{Lectio V}
\switchcolumn
\EN
\LessonHead{Lesson V}
\switchcolumn*

\LA
\noindent Dicénte enim prophéta: Quǽrite Dóminum, et confirmámini, quǽrite fáciem ejus semper: némini præsuméndum est, quod totum quod quærit, invénerit, ne désinat propinquáre, qui cessárit accédere. Quid autem inter ómnia ópera Dei, in quibus humánæ admiratiónis fatigátur inténtio, ita contemplatiónem mentis nostræ et obléctat et súperat, sicut pássio Salvatóris? Qui ut humánum genus vínculis mortíferæ prævaricatiónis absólveret, et sæviénti diábolo poténtiam suæ majestátis occúluit, et infirmitátem nostræ humilitátis objécit. Si enim crudélis et supérbus inimícus consílium misericórdiæ Dei nosse potuísset, Judæórum ánimos mansuetúdine pótius temperáre, quam injústis ódiis studuísset accéndere: ne ómnium captivórum amítteret servitútem, dum nihil sibi debéntis perséquitur libertátem.\par
\switchcolumn
\EN
\noindent For when the Prophet saith: Seek the Lord and be strong; seek His face evermore, \textasciitilde{}(Ps. civ. 4,) let no man thence conclude that he will ever have found all that he seeketh, lest he which hath ceased to come near should cease to be near. But among all the works of God which foil and weary the steadfast gaze of man's wonder, what is there that doth at once so ravish and so exceed the power of our mind's eye as do the sufferings of the Saviour? He it was Who, to loose man from the bands wherewith he had bound himself by the first death-dealing transgression, spared to bring against the rage of the devil the power of the Divine Majesty, and met him with the weakness of our lowly nature. For if our proud and cruel enemy had been able to know the counsel of God's mercy, it had been his task rather to have softened the minds of the Jews into gentleness, than to have inflamed them with unrighteous hatred; and so lost the service of all his slaves, by pursuing for his Debtor One That owed him nothing.\par
\switchcolumn*

\LA
\begin{Resp}\Rsign Salvum me fac, Deus, quóniam intravérunt aquæ usque ad ánimam meam: ne avértas fáciem tuam a me: \Star{} Quóniam tríbulor, exáudi me, Dómine, Deus meus.\end{Resp}
\begin{Resp}\Vsign Inténde ánimæ meæ, et líbera eam: propter inimícos meos éripe me.\end{Resp}
\begin{Resp}\Rsign Quóniam tríbulor, exáudi me, Dómine, Deus meus.\end{Resp}
\switchcolumn
\EN
\begin{Resp}\Rsign Save me, O God, for the waters are come in unto my soul; hide not thy face from me; \Star{} For I am in trouble. Hear me speedily, O Lord my God.\end{Resp}
\begin{Resp}\Vsign Draw nigh unto my soul, and redeem it deliver me because of mine enemies.\end{Resp}
\begin{Resp}\Rsign For I am in trouble. Hear me speedily, O Lord my God.\end{Resp}
\switchcolumn*

\LA
\LessonHead{Lectio VI}
\switchcolumn
\EN
\LessonHead{Lesson VI}
\switchcolumn*

\LA
\noindent Feféllit ergo illum malígnitas sua, íntulit supplícium Fílio Dei, quod cunctis fíliis hóminum in remédium verterétur. Fudit sánguinem justum, qui reconciliándo mundo et prétium esset, et póculum. Suscépit Dóminus, quod secúndum propósitum suæ voluntátis elégit. Admísit in se ímpias manus furéntium: quæ dum próprio incúmbunt scéleri, famulátæ sunt Redemptóri. Cujus étiam circa interfectóres suos tanta erat pietátis afféctio, ut de cruce súpplicans Patri, non se vindicári, sed illis postuláret ignósci.\par
\switchcolumn
\EN
\noindent But his own hate dug a pit-fall for him he brought upon the Son of God that death which is become life to all the sons of men. He shed that innocent Blood, Which hath reconciled the world unto God, and become at once the price of our redemption and the cup of our salvation. The Lord hath received that which according to the purpose of His Own good pleasure He hath chosen. He hath let fall on Him the hands of bloody men, but while they were bent only on their own sin, they were servants ministering to the Redeemer's work. And such was His tenderness even for His murderers that His prayer to His Father from the Cross, as touching them, was, not that He might be avenged upon them, but that they might be forgiven.\par
\switchcolumn*

\LA
\begin{Resp}\Rsign Noli esse mihi, Dómine, aliénus: parce mihi in die mala: confundántur omnes qui me persequúntur, \Star{} Et non confúndar ego.\end{Resp}
\begin{Resp}\Vsign Confundántur omnes inimíci mei, qui quærunt ánimam meam.\end{Resp}
\begin{Resp}\Rsign Et non confúndar ego.\end{Resp}
\begin{Resp}\Vsign Glória Patri, et Fílio, \Star{} et Spirítui Sancto.\end{Resp}
\begin{Resp}\Rsign Noli esse mihi, Dómine, aliénus: parce mihi in die mala: confundántur omnes qui me persequúntur, * Et non confúndar ego.\end{Resp}
\switchcolumn
\EN
\begin{Resp}\Rsign O Lord, be not Thou far from me; spare me in the day of evil, let them be confounded that persecute me; \Star{} But let not me be confounded.\end{Resp}
\begin{Resp}\Vsign Let all mine enemies which seek after my soul be confounded.\end{Resp}
\begin{Resp}\Rsign But let not me be confounded.\end{Resp}
\begin{Resp}\Vsign Glory be to the Father, and to the Son, \Star{} and to the Holy Ghost.\end{Resp}
\begin{Resp}\Rsign O Lord, be not Thou far from me spare me in the evil day let them be confounded that persecute me; * but let not me be confounded.\end{Resp}
\switchcolumn*

\LA
\LessonHead{Lectio VII}
\switchcolumn
\EN
\LessonHead{Lesson VII}
\switchcolumn*

\LA
\LessonTitle{Léctio sancti Evangélii secúndum Matthǽum}
\switchcolumn
\EN
\LessonTitle{From the Holy Gospel according to Matthew}
\switchcolumn*

\LA
\Cite{Matt 21:1-9}
\switchcolumn
\EN
\Cite{Matt 21:1-9}
\switchcolumn*

\LA
\noindent In illo témpore: Cum appropinquásset Jesus Jerosólymis, et venísset Béthphage ad montem Olivéti: tunc misit duos discípulos, dicens eis. Et réliqua.\par
\switchcolumn
\EN
\noindent At that time, when Jesus drew nigh to Jerusalem, and were come to Bethphage, unto mount Olivet, then Jesus sent two disciples, Saying to them: And so on.\par
\switchcolumn*

\LA
\LessonTitle{Homilía sancti Ambrósii Epíscopi}
\switchcolumn
\EN
\LessonTitle{Homily by St. Ambrose, Bishop (of Milan.)}
\switchcolumn*

\LA
\Source{Liber 9 in Lucam}
\switchcolumn
\EN
\Source{9th Book on Luke}
\switchcolumn*

\LA
\noindent Pulchre relíctis Judǽis, habitatúrus in afféctibus géntium, templum Dóminus ascéndit. Hoc enim templum est verum, in quo non in líttera, sed in spíritu Dóminus adorátur. Hoc Dei templum est, quod fídei séries, non lápidum structúra fundávit. Deserúntur ergo qui óderant: eligúntur qui amatúri erant. Et ídeo ad montem venit Olivéti, ut novéllas óleas in sublími virtúte plantáret, quarum mater est illa, quæ sursum est, Jerúsalem. In hoc monte est ille cæléstis agrícola: ut plantáti omnes in domo Dei, possint virítim dícere: Ego autem sicut olíva fructífera in domo Dómini.\par
\switchcolumn
\EN
\noindent Beautiful is the type, when the Lord, about to leave the Jews, and to take up His abode in the hearts of the Gentiles, goeth up into the Temple; a figure of His going to the true Temple wherein He is worshipped, not in the deadness of the letter, but in spirit and in truth, even that Temple of God whereof the foundations are laid, not in buildings of stone, but in faith. He leaveth behind Him such as hate Him, and getteth Him to such as will love Him. And therefore cometh He unto the Mount of Olives that He may plant upon the heights of grace those young olive-branches, whose Mother is the Jerusalem which is above. Upon this mountain standeth He, the Heavenly Husbandman, that all they which be planted in the House of the Lord may be able each one to say: “But I am like a fruitful olive-tree in the House of God.” (Ps. li. 10.)\par
\switchcolumn*

\LA
\begin{Resp}\Rsign Dóminus mecum est tamquam bellátor fortis: proptérea persecúti sunt me, et intellégere non potuérunt: Dómine, probas renes et corda: \Star{} Tibi revelávi causam meam.\end{Resp}
\begin{Resp}\Vsign Vidísti, Dómine, iniquitátes eórum advérsum me: júdica judícium meum.\end{Resp}
\begin{Resp}\Rsign Tibi revelávi causam meam.\end{Resp}
\switchcolumn
\EN
\begin{Resp}\Rsign The Lord is with me as a Mighty Terrible One; therefore have they persecuted me, and have not been able to understand. O Lord, Thou triest the reins and the heart \Star{} Unto thee have I opened my cause.\end{Resp}
\begin{Resp}\Vsign O Lord, Thou hast seen my wrong that they do me; judge Thou my cause.\end{Resp}
\begin{Resp}\Rsign Unto thee have I opened my cause.\end{Resp}
\switchcolumn*

\LA
\LessonHead{Lectio VIII}
\switchcolumn
\EN
\LessonHead{Lesson VIII}
\switchcolumn*

\LA
\noindent Et fortásse ipse mons Christus est. Quis enim álius tales fructus ferret oleárum, non curvescéntium ubertáte baccárum, sed spíritus plenitúdine géntium fœcundárum? Ipse est per quem ascéndimus, et ad quem ascéndimus. Ipse est jánua, ipse est via, qui aperítur, et qui áperit: qui pulsátur ab ingrediéntibus, et ab eméritis adorátur. Ergo in castéllo erat, et ligátus erat pullus cum ásina: non póterat solvi nisi jussu Dómini. Solvit eum manus apostólica. Talis actus, talis vita, talis grátia. Esto talis et tu, ut possis ligátos sólvere.\par
\switchcolumn
\EN
\noindent And perchance that mountain doth signify Christ Himself. For what other is there that beareth such fruit of olives as He doth, not rich with store of loaded branches, but spiritually fruitful with the fulness of the Gentiles? He also it is on Whom we go up, and unto Whom we go up; He is the Door; He is the Way; He is He Which is opened and Which openeth; He is He upon Whom knocketh whosoever entereth in, and to Whom they that have entered in, do worship. A figure also was it that the disciples went into a village, and that there they found an ass tied and a colt with her neither could they be loosed, save at the command of the Lord. It was the hand of His Apostles which loosed them. He whose work and life are like theirs will have such grace as was theirs. Be thou also such as they, if thou wouldest loose them that are bound.\par
\switchcolumn*

\LA
\begin{Resp}\Rsign Dixérunt ímpii apud se, non recte cogitántes: Circumveniámus justum, quóniam contrárius est opéribus nostris: promíttit se sciéntiam Dei habére, Fílium Dei se nóminat, et gloriátur patrem se habére Deum: \Star{} Videámus si sermónes illíus veri sunt: et si est vere Fílius Dei, líberet eum de mánibus nostris: morte turpíssima condemnémus eum.\end{Resp}
\begin{Resp}\Vsign Tamquam nugáces æstimáti sumus ab illo, et ábstinet se a viis nostris tamquam ab immundítiis: et præfert novíssima justórum.\end{Resp}
\begin{Resp}\Rsign Videámus si sermónes illíus veri sunt: et si est vere Fílius Dei, líberet eum de mánibus nostris: morte turpíssima condemnémus eum.\end{Resp}
\switchcolumn
\EN
\begin{Resp}\Rsign The ungodly said, reasoning with themselves, but not aright; Let us lie in wait for the righteous, because he is clean contrary to our doings he professeth to have the knowledge of God, he calleth himself the Son of God, and boasteth that he hath God to his Father. \Star{} Let us see if his words be true; and, if he be indeed the Son of God, let Him deliver him from our hand; let us condemn him with a shameful death.\end{Resp}
\begin{Resp}\Vsign We are esteemed of him as counterfeits, and he abstaineth from our ways as from filthiness, and commendeth the end of the just.\end{Resp}
\begin{Resp}\Rsign Let us see if his words be true; and, if he be indeed the Son of God, let Him deliver him from our hand; let us condemn him with a shameful death.\end{Resp}
\switchcolumn*

\LA
\LessonHead{Lectio IX}
\switchcolumn
\EN
\LessonHead{Lesson IX}
\switchcolumn*

\LA
\noindent Nunc considerémus qui fúerint illi, qui erróre detécto, de paradíso ejécti, in castéllum sint relegáti. Et vides, quemádmodum quos mors expúlerat, vita revocáverit. Et ídeo secúndum Matthǽum, et ásinam et pullum légimus: ut quia in duóbus homínibus utérque fúerat sexus expúlsus, in duóbus animálibus sexus utérque revocétur. Ergo illic in ásina matre quasi Hevam figurávit erróris: hic autem in pullo generalitátem pópuli Gentílis expréssit: et ídeo pullo sedétur ásinæ. Et bene, in quo nemo sedit: quia nullus, ántequam Christus, natiónum pópulos vocávit ad Ecclésiam. Dénique secúndum Marcum sic habes: Quem nemo adhuc sedit hóminum.\par
\switchcolumn
\EN
\noindent Now, let us consider who they were, who, being convicted of transgression, were banished from their home in the Garden of Eden into a village, and in this thou wilt see how Life called back again them whom death had cast out. For this reason, we read in Matthew that there were tied both an ass and her colt; thus, as man was banished from Eden in a member of either sex, so is it in animals of both sexes that his re-call is figured. The she-ass is a type of our sinful Mother Eve, and the colt of the multitude of the Gentiles; and it was upon the colt that Christ took His seat. And thus it is well written of the colt, (Luke xix. 30,) that thereon never yet had man sat, for no man before Christ ever called the Gentiles into the Church which statement thou hast in Mark also (xi. 2): Whereon never man sat.\par
\switchcolumn*

\LA
\begin{Resp}\Rsign Circumdedérunt me viri mendáces: sine causa flagéllis cecidérunt me: \Star{} Sed tu, Dómine defénsor, víndica me.\end{Resp}
\begin{Resp}\Vsign Quóniam tribulátio próxima est, et non est qui ádjuvet.\end{Resp}
\begin{Resp}\Rsign Sed tu, Dómine defénsor, víndica me.\end{Resp}
\begin{Resp}\Vsign Glória Patri, et Fílio, \Star{} et Spirítui Sancto.\end{Resp}
\begin{Resp}\Rsign Circumdedérunt me viri mendáces: sine causa flagéllis cecidérunt me: * Sed tu, Dómine defénsor, víndica me.\end{Resp}
\switchcolumn
\EN
\begin{Resp}\Rsign Liars are come round about me, they have fallen upon me with scourges without a cause. \Star{} But do Thou, O Lord my Redeemer, avenge me!\end{Resp}
\begin{Resp}\Vsign For trouble is near, and there is none to help.\end{Resp}
\begin{Resp}\Rsign But do Thou, O Lord my Redeemer, avenge me!\end{Resp}
\begin{Resp}\Vsign Glory be to the Father, and to the Son, \Star{} and to the Holy Ghost.\end{Resp}
\begin{Resp}\Rsign Liars are come round about me, they have fallen upon me with scourges without a cause. * But do Thou, O Lord my Redeemer, avenge me!\end{Resp}
\switchcolumn*

\end{paracol}

\Office{Feria Secunda Majoris Hebdomadæ}{Monday in Holy Week}{Feria privilegiata}
\addcontentsline{toc}{subsection}{Feria Secunda Majoris Hebdomadæ\enspace\textit{Monday in Holy Week}}
\begin{paracol}{2}
\LA
\LessonHead{Lectio I}
\switchcolumn
\EN
\LessonHead{Lesson I}
\switchcolumn*

\LA
\LessonTitle{Léctio sancti Evangélii secúndum Joánnem}
\switchcolumn
\EN
\LessonTitle{Continuation of the Holy Gospel according to John}
\switchcolumn*

\LA
\Cite{Joannes 12:1-9}
\switchcolumn
\EN
\Cite{John 12:1-9}
\switchcolumn*

\LA
\noindent Ante sex dies Paschæ venit Jesus Bethániam, ubi Lázarus fúerat mórtuus, quem suscitávit Jesus. Et réliqua.\par
\switchcolumn
\EN
\noindent Six days before the pasch, Jesus came to Bethania, where Lazarus had been dead, whom Jesus raised to life.\par
\switchcolumn*

\LA
\LessonTitle{Homilía sancti Augustíni Epíscopi}
\switchcolumn
\EN
\LessonTitle{Homily by St. Augustine, Bishop (of Hippo.)}
\switchcolumn*

\LA
\Source{Tractus 50 in Joann., post initium}
\switchcolumn
\EN
\Source{50th Tract on John}
\switchcolumn*

\LA
\noindent Ne putárent hómines phantásma esse factum, quia mórtuus resurréxit, Lázarus unus erat ex recumbéntibus: vivébat, loquebátur, epulabátur, véritas ostendebátur, infidélitas Judæórum confundebátur. Discumbébat ergo Jesus cum Lázaro, et céteris: ministrábat Martha, una ex soróribus Lázari. María vero, áltera soror Lázari, accépit libram unguénti nardi pístici pretiósi, et unxit pedes Jesu, et extérsit capíllis suis pedes ejus, et domus impléta est ex odóre unguénti. Factum audívimus: mystérium requirámus.\par
\switchcolumn
\EN
\noindent There they made Him a supper and Lazarus was one of them that sat at the table lest men should deem that it was but by an ocular delusion that they had seen him arise from the dead. He lived therefore, spake, and ate; to the manifestation of the truth, and the confusion of the unbelieving Jews. Jesus, then, sat down to meat with Lazarus and others, and Martha, being one of Lazarus' sisters, served. But Mary, Lazarus' other sister, took a pound of ointment of spikenard, very costly, and anointed the Feet of Jesus, and wiped His Feet with her hair; and the house was filled with the odour of the ointment. We have now heard that which was done; let us search out the mystic meaning thereof.\par
\switchcolumn*

\LA
\begin{Resp}\Rsign Viri ímpii dixérunt: Opprimámus virum justum injúste, et deglutiámus eum tamquam inférnus vivum: auferámus memóriam illíus de terra: et de spóliis ejus sortem mittámus inter nos: ipsi enim homicídæ thesaurizavérunt sibi mala. \Star{} Insipiéntes et malígni odérunt sapiéntiam: et rei facti sunt in cogitatiónibus suis.\end{Resp}
\begin{Resp}\Vsign Hæc cogitavérunt, et erravérunt: et excæcávit illos malítia eórum.\end{Resp}
\begin{Resp}\Rsign Insipiéntes et malígni odérunt sapiéntiam: et rei facti sunt in cogitatiónibus suis.\end{Resp}
\switchcolumn
\EN
\begin{Resp}\Rsign The ungodly said: Let us oppress the righteous man without cause, and swallow him up alive, as the grave; let us make his memorial to perish from the earth, and cast lots among us for his spoils and those murderers laid by store for themselves, but of evil. \Star{} Fools and haters loathe wisdom, and are guilty in their thoughts.\end{Resp}
\begin{Resp}\Vsign Such things they did imagine, and were deceived, for their own wickedness blinded them.\end{Resp}
\begin{Resp}\Rsign Fools and haters loathe wisdom, and are guilty in their thoughts.\end{Resp}
\switchcolumn*

\LA
\LessonHead{Lectio II}
\switchcolumn
\EN
\LessonHead{Lesson II}
\switchcolumn*

\LA
\noindent Quæcúmque ánima fidélis vis esse, cum María unge pedes Dómini pretióso unguénto. Unguéntum illud justítia fuit, ídeo libra fuit: erat autem unguéntum nardi pístici pretiósi. Quod ait, pístici, locum áliquem crédere debémus, unde hoc erat unguéntum pretiósum: nec tamen hoc vacat, et sacraménto óptime cónsonat. Pistis Græce, fides Latíne dícitur. Quærébas operári justítiam. Justus ex fide vivit. Unge pedes Jesu bene vivéndo: Domínica sectáre vestígia. Capíllis terge: si habes supérflua, da paupéribus, et Dómini pedes tersísti: capílli enim supérflua córporis vidéntur. Habes quod agas de supérfluis tuis: tibi supérflua sunt, sed Dómini pédibus necessária sunt. Forte in terra Dómini pedes índigent.\par
\switchcolumn
\EN
\noindent Whatsover thou art that wilt be a faithful soul, seek with Mary to anoint the Feet of the Lord with costly ointment. This ointment was a figure of justice, and therefore is there said to have been a pound thereof, a pound being a weight used in scales. The word pistikes used by the Evangelist as the name of this ointment, we must believe to be that of some place, from which this costly perfume was imported. Neither is this name meaningless for us, but agreeth well with our mystic interpretation, since Pistis is the Greek word which signifieth Faith, and whosoever will do justice must know that: The just shall live by faith. (Rom. i. 17; Hab. ii. 4.) Anoint therefore the Feet of Jesus by thy good life, following in the marks which those Feet of the Lord have traced. Wipe His Feet likewise with thy hair; that is, if thou have aught which is not needful to thee, give it to the poor; and then thou hast wiped the Feet of Jesus with thy hair, that is, with that which thou needest not, and which is therefore to thee as is hair, being a needless out-growth to the body. Here thou hast what to do with that which thou needest not. To thee it is needless, but the Lord's Feet have need of it; yea, the Feet which the Lord hath on earth are sorely needy.\par
\switchcolumn*

\LA
\begin{Resp}\Rsign Oppróbrium factus sum nimis inimícis meis: vidérunt me, et movérunt cápita sua: \Star{} Adjuva me, Dómine, Deus meus.\end{Resp}
\begin{Resp}\Vsign Locúti sunt advérsum me lingua dolósa, et sermónibus ódii circumdedérunt me.\end{Resp}
\begin{Resp}\Rsign Adjuva me, Dómine, Deus meus.\end{Resp}
\switchcolumn
\EN
\begin{Resp}\Rsign I became a reproach unto mine enemies they looked upon me and shaked their heads. \Star{} Help me, O Lord my God!\end{Resp}
\begin{Resp}\Vsign They have spoken against me with a lying tongue they compassed me about also with words of hatred.\end{Resp}
\begin{Resp}\Rsign Help me, O Lord my God!\end{Resp}
\switchcolumn*

\LA
\LessonHead{Lectio III}
\switchcolumn
\EN
\LessonHead{Lesson III}
\switchcolumn*

\LA
\noindent De quibus enim, nisi de membris suis in fine dictúrus est: Cum uni ex mínimis meis fecístis, mihi fecístis? Supérflua vestra impendístis: sed pédibus meis obsecúti estis. Domus autem impléta est odóre: mundus implétus est fama bona: nam odor bonus, fama bona est. Qui male vivunt, et Christiáni vocántur, injúriam Christo fáciunt: de quálibus dictum est, quod per eos nomen Dómini blasphemátur. Si per tales nomen Dei blasphemátur, per bonos nomen Dómini laudátur. Audi Apóstolum: Christi bonus odor sumus, inquit, in omni loco.\par
\switchcolumn
\EN
\noindent For of whom save of His members, will He say at the latter day: Inasmuch as ye have done it unto one of the least of these My brethren, ye have done it unto Me. (Matth. xxv. 40.) That is, ye have spent nothing save that which ye needed not, but ye have ministered unto My Feet. And the house was filled with the odour of the ointment. That is, the fragrance of your good example filleth the world; for this odour is a figure of reputation. They which are called Christians, and yet live bad lives, cast a slur on Christ and it is even such as they unto whom it is said: The Name of God is blasphemed among the Gentiles through you. (Rom. ii. 24; Ezek. xxxvi. 20, 23.) But if, through such, the Name of God be blasphemed, through the godly is praise ascribed to the Same His Holy Name, as the Apostle doth likewise say: In every place we are unto God a sweet savour of Christ, in them that are saved, and in them that perish. (2 Cor. ii. 14, 15.)\par
\switchcolumn*

\LA
\begin{Resp}\Rsign Insurrexérunt in me viri iníqui absque misericórdia, quæsiérunt me interfícere: et non pepercérunt in fáciem meam spúere, et lánceis suis vulneravérunt me: et concússa sunt ómnia ossa mea: \Star{} Ego autem existimábam me tamquam mórtuum super terram.\end{Resp}
\begin{Resp}\Vsign Effudérunt furórem suum in me: fremuérunt contra me déntibus suis.\end{Resp}
\begin{Resp}\Rsign Ego autem existimábam me tamquam mórtuum super terram.\end{Resp}
\begin{Resp}\Vsign Glória Patri, et Fílio, \Star{} et Spirítui Sancto.\end{Resp}
\begin{Resp}\Rsign Insurrexérunt in me viri iníqui absque misericórdia, quæsiérunt me interfícere: et non pepercérunt in fáciem meam spúere, et lánceis suis vulneravérunt me: et concússa sunt ómnia ossa mea: * Ego autem existimábam me tamquam mórtuum super terram.\end{Resp}
\switchcolumn
\EN
\begin{Resp}\Rsign False witnesses are risen up against me, and such as breathe out cruelty they have gone about to kill me, neither spared they to spit in my face; their spears have wounded me, and all my bones are out of joint. \Star{} But as for me, I counted myself as one that is dead upon the earth.\end{Resp}
\begin{Resp}\Vsign They poured forth their fury upon me, they gnashed upon me with their teeth.\end{Resp}
\begin{Resp}\Rsign But as for me, I counted myself as one that is dead upon the earth.\end{Resp}
\begin{Resp}\Vsign Glory be to the Father, and to the Son, \Star{} and to the Holy Ghost.\end{Resp}
\begin{Resp}\Rsign False witnesses are risen up against me, and such as breathe out cruelty; they have gone about to kill me, neither spared they to spit in my face; their spears have wounded me, all my bones are out of joint. But as for me, I counted myself as one that is dead upon the earth.\end{Resp}
\switchcolumn*

\end{paracol}

\Office{Feria Tertia Majoris Hebdomadæ}{Tuesday in Holy Week}{Feria privilegiata}
\addcontentsline{toc}{subsection}{Feria Tertia Majoris Hebdomadæ\enspace\textit{Tuesday in Holy Week}}
\begin{paracol}{2}
\LA
\LessonHead{Lectio I}
\switchcolumn
\EN
\LessonHead{Lesson I}
\switchcolumn*

\LA
\LessonTitle{De Jeremía Prophéta}
\switchcolumn
\EN
\LessonTitle{Lesson from the book of Jeremias}
\switchcolumn*

\LA
\Cite{Jer 11:15-20}
\switchcolumn
\EN
\Cite{Jer 11:15-20}
\switchcolumn*

\LA
\noindent Quid est quod diléctus meus in domo mea fecit scélera multa? Numquid carnes sanctæ áuferent a te malítias tuas, in quibus gloriáta es? Olívam úberem, pulchram, fructíferam, speciósam vocávit Dóminus nomen tuum: ad vocem loquélæ, grandis exársit ignis in ea, et combústa sunt frutéta ejus. Et Dóminus exercítuum, qui plantávit te, locútus est super te malum: pro malis domus Israël et domus Juda, quæ fecérunt sibi ad irritándum me, libántes Báalim. Tu autem, Dómine, demonstrásti mihi, et cognóvi: tunc ostendísti mihi stúdia eórum. Et ego quasi agnus mansuétus, qui portátur ad víctimam: et non cognóvi quia cogitavérunt super me consília, dicéntes: Mittámus lignum in panem ejus, et eradámus eum de terra vivéntium, et nomen ejus non memorétur ámplius. Tu autem, Dómine Sábaoth, qui júdicas juste, et probas renes et corda, vídeam ultiónem tuam ex eis: tibi enim revelávi causam meam.\par
\switchcolumn
\EN
\noindent What is the meaning that my beloved hath wrought much wickedness in my house? shall the holy flesh take away from thee thy crimes, in which thou hast boasted? The Lord called thy name, a plentiful olive tree, fair, fruitful, and beautiful: at the noise of a word, a great fire was kindled in it and the branches thereof are burnt. And the Lord of hosts that planted thee, hath pronounced evil against thee: for the evils of the house of Israel, and the house of Juda, which they have done to themselves, to provoke me, offering sacrifice to Baalim. But thou, O Lord, hast shewn me, and I have known: then thou shewedst me their doings. And I was as a meek lamb, that is carried to be a victim: and I knew not that they had devised counsels against me, saying: Let us put wood on his bread, and cut him off from the land of the living, and let his name be remembered no more. But thou, O Lord of Sabaoth, who judgest justly, and triest the reins and hearts, let me see thy revenge on them: for to thee I have revealed my cause.\par
\switchcolumn*

\LA
\begin{Resp}\Rsign Contumélias et terróres passus sum ab eis, qui erant pacífici mei, et custodiéntes latus meum, dicéntes: Decipiámus eum, et prævaleámus illi: sed tu, Dómine, mecum es tamquam bellátor fortis. \Star{} Cadant in oppróbrium sempitérnum, ut vídeam vindíctam in eis, quia tibi revelávi causam meam.\end{Resp}
\begin{Resp}\Vsign Júdica, Dómine, causam ánimæ meæ, defénsor vitæ meæ.\end{Resp}
\begin{Resp}\Rsign Cadant in oppróbrium sempitérnum, ut vídeam vindíctam in eis, quia tibi revelávi causam meam.\end{Resp}
\switchcolumn
\EN
\begin{Resp}\Rsign I have suffered defaming and fear from them that were my familiars they watched for my halting, saying Let us entice him, and prevail against him. But Thou, O Lord, art with me, as a Mighty Terrible One. \Star{} Let them stumble into everlasting confusion, that I may see thy vengeance upon them, for unto thee have I opened my cause.\end{Resp}
\begin{Resp}\Vsign O Lord, plead Thou the cause of my soul, Thou That art the Redeemer of my life.\end{Resp}
\begin{Resp}\Rsign Let them stumble into everlasting confusion, that I may see thy vengeance upon them, for unto thee have I opened my cause.\end{Resp}
\switchcolumn*

\LA
\LessonHead{Lectio II}
\switchcolumn
\EN
\LessonHead{Lesson II}
\switchcolumn*

\LA
\Cite{Jer 12:1-4}
\switchcolumn
\EN
\Cite{Jer 12:1-4}
\switchcolumn*

\LA
\noindent Justus quidem tu es, Dómine, si dispútem tecum: verúmtamen justa loquar ad te: Quare via impiórum prosperátur: bene est ómnibus, qui prævaricántur, et iníque agunt? Plantásti eos, et radícem misérunt: profíciunt et fáciunt fructum: prope es tu ori eórum, et longe a rénibus eórum. Et tu, Dómine, nosti me, vidísti me, et probásti cor meum tecum: cóngrega eos quasi gregem ad víctimam, et sanctífica eos in die occisiónis. Usquequo lugébit terra, et herba omnis regiónis siccábitur propter malítiam habitántium in ea? Consúmptum est ánimal et vólucre, quóniam dixérunt: Non vidébit novíssima nostra.\par
\switchcolumn
\EN
\noindent Thou indeed, O Lord, art just, if I plead with thee, but yet I will speak what is just to thee: Why doth the way of the wicked prosper: why is it well with all them that transgress, and do wickedly? Thou hast planted them, and they have taken root: they prosper and bring forth fruit: thou art near in their mouth, and far from their reins. And thou, O Lord, hast known me, thou hast seen me, and proved my heart with thee: gather them together as sheep for a sacrifice, and prepare them for the day of slaughter. How long shall the land mourn, and the herb of every field wither for the wickedness of them that dwell therein? The beasts and the birds are consumed: because they have said: He shall not see our last end.\par
\switchcolumn*

\LA
\begin{Resp}\Rsign Deus Israël, propter te sustínui impropérium, opéruit reveréntia fáciem meam, extráneus factus sum frátribus meis, et hospes fíliis matris meæ: \Star{} Quóniam zelus domus tuæ comédit me.\end{Resp}
\begin{Resp}\Vsign Inténde ánimæ meæ, et líbera eam, propter inimícos meos éripe me.\end{Resp}
\begin{Resp}\Rsign Quóniam zelus domus tuæ comédit me.\end{Resp}
\switchcolumn
\EN
\begin{Resp}\Rsign For thy sake, O God of Israel, I have borne reproach; shame hath covered my face; I am become a stranger unto my brethren, and an alien unto my mother's children. \Star{} For the zeal of thine house hath eaten me up.\end{Resp}
\begin{Resp}\Vsign Draw nigh unto my soul, and redeem it; deliver me, because of mine enemies.\end{Resp}
\begin{Resp}\Rsign For the zeal of thine house hath eaten me up.\end{Resp}
\switchcolumn*

\LA
\LessonHead{Lectio III}
\switchcolumn
\EN
\LessonHead{Lesson III}
\switchcolumn*

\LA
\Cite{Jer 12:7-11}
\switchcolumn
\EN
\Cite{Jer 12:7-11}
\switchcolumn*

\LA
\noindent Relíqui domum meam, dimísi hereditátem meam: dedi diléctam ánimam meam in manu inimicórum ejus. Facta est mihi heréditas mea quasi leo in silva: dedit contra me vocem, ídeo odívi eam. Numquid avis díscolor heréditas mea mihi? numquid avis tincta per totum? Veníte, congregámini, omnes béstiæ terræ, properáte ad devorándum. Pastóres multi demolíti sunt víneam meam, conculcavérunt partem meam: dedérunt portiónem meam desiderábilem in desértum solitúdinis. Posuérunt eam in dissipatiónem, luxítque super me: desolatióne desoláta est omnis terra, quia nullus est qui recógitet corde.\par
\switchcolumn
\EN
\noindent I have forsaken my house, I have left my inheritance: I have given my dear soul into the land of her enemies. My inheritance is become to me as a lion in the wood: it hath cried out against me, therefore have I hated it. Is my inheritance to me as a speckled bird? Is it as a bird died throughout? come ye, assemble yourselves, all the beasts of the earth, make haste to devour. Many pastors have destroyed my vineyard, they have trodden my portion under foot: they have changed my delightful portion into a desolate wilderness. They have laid it waste, and it hath mourned for me. With desolation is all the land made desolate; because there is none that considereth in the heart.\par
\switchcolumn*

\LA
\begin{Resp}\Rsign Synagóga populórum circumdedérunt me: et non réddidi retribuéntibus mihi mala. \Star{} Consumétur, Dómine, nequítia peccatórum, et díriges justum.\end{Resp}
\begin{Resp}\Vsign Júdica me, Dómine, secúndum justítiam meam, et secúndum innocéntiam meam super me.\end{Resp}
\begin{Resp}\Rsign Consumétur, Dómine, nequítia peccatórum, et díriges justum.\end{Resp}
\begin{Resp}\Vsign Glória Patri, et Fílio, \Star{} et Spirítui Sancto.\end{Resp}
\begin{Resp}\Rsign Synagóga populórum circumdedérunt me: et non réddidi retribuéntibus mihi mala. * Consumétur, Dómine, nequítia peccatórum, et díriges justum.\end{Resp}
\switchcolumn
\EN
\begin{Resp}\Rsign The congregation of the people hath compassed me about, but I rewarded no evil unto him that rewarded evil unto me. \Star{} O Lord, let the wickedness of the wicked come to an end, but establish the just.\end{Resp}
\begin{Resp}\Vsign Judge me, O Lord, according to my righteousness, and according to mine integrity that is in me.\end{Resp}
\begin{Resp}\Rsign O Lord, let the wickedness of the wicked come to an end, but establish the just.\end{Resp}
\begin{Resp}\Vsign Glory be to the Father, and to the Son, \Star{} and to the Holy Ghost.\end{Resp}
\begin{Resp}\Rsign The congregation of the people hath compassed me about, but I rewarded no evil unto him that rewarded evil unto me. O Lord, let the wickedness of the wicked come to an end, but establish the just.\end{Resp}
\switchcolumn*

\end{paracol}

\Office{Feria Quarta Majoris Hebdomadæ}{Wednesday in Holy Week}{Feria privilegiata}
\addcontentsline{toc}{subsection}{Feria Quarta Majoris Hebdomadæ\enspace\textit{Wednesday in Holy Week}}
\begin{paracol}{2}
\LA
\LessonHead{Lectio I}
\switchcolumn
\EN
\LessonHead{Lesson I}
\switchcolumn*

\LA
\LessonTitle{De Jeremía Prophéta}
\switchcolumn
\EN
\LessonTitle{Lesson from the book of Jeremias}
\switchcolumn*

\LA
\Cite{Jer 17:13-18}
\switchcolumn
\EN
\Cite{Jer 17:13-18}
\switchcolumn*

\LA
\noindent Exspectátio Israël, Dómine: omnes, qui te derelínquunt, confundéntur: recedéntes a te, in terra scribéntur: quóniam dereliquérunt venam aquárum vivéntium Dóminum. Sana me, Dómine, et sanábor: salvum me fac, et salvus ero: quóniam laus mea tu es. Ecce ipsi dicunt ad me: Ubi est verbum Dómini? véniat. Et ego non sum turbátus, te pastórem sequens: et diem hóminis non desiderávi, tu scis. Quod egréssum est de lábiis meis, rectum in conspéctu tuo fuit. Non sis tu mihi formídini, spes mea tu in die afflictiónis. Confundántur qui me persequúntur, et non confúndar ego: páveant illi, et non páveam ego: induc super eos diem afflictiónis, et dúplici contritióne cóntere eos.\par
\switchcolumn
\EN
\noindent O Lord, the hope of Israel: all that forsake thee shall be confounded: they that depart from thee, shall be written in the earth: because they have forsaken the Lord, the vein of living waters. Heal me, O Lord, and I shall be healed: save me, and I shall be saved, for thou art my praise. Behold they say to me: Where is the word of the Lord? let it come. And I am not troubled, following thee for my pastor, and I have not desired the day of man, thou knowest. That which went out of my lips, hath been right in thy sight. Be not thou a terror unto me, thou art my hope in the day of affliction. Let them be confounded that persecute me, and let not me be confounded: let them be afraid, and let not me be afraid: bring upon them the day of affliction, and with a double destruction, destroy them.\par
\switchcolumn*

\LA
\begin{Resp}\Rsign Locúti sunt advérsum me lingua dolósa, et sermónibus ódii circumdedérunt me: pro eo ut me dilígerent, detrahébant mihi: \Star{} Ego autem orábam, et exaudísti me, Dómine, Deus meus.\end{Resp}
\begin{Resp}\Vsign Et posuérunt advérsum me mala pro bonis, et ódium pro dilectióne mea.\end{Resp}
\begin{Resp}\Rsign Ego autem orábam, et exaudísti me, Dómine, Deus meus.\end{Resp}
\switchcolumn
\EN
\begin{Resp}\Rsign They have spoken against me with a lying tongue; they compassed me about also with words of hatred: in return for my love they were my adversaries: \Star{} But I gave myself unto prayer; and Thou hast heard me, Lord my God!\end{Resp}
\begin{Resp}\Vsign And they have rewarded me evil for good, and hatred for my love.\end{Resp}
\begin{Resp}\Rsign But I gave myself unto prayer; and Thou hast heard me, Lord my God!\end{Resp}
\switchcolumn*

\LA
\LessonHead{Lectio II}
\switchcolumn
\EN
\LessonHead{Lesson II}
\switchcolumn*

\LA
\Cite{Jer 18:13-18}
\switchcolumn
\EN
\Cite{Jer 18:13-18}
\switchcolumn*

\LA
\noindent Quis audívit tália horribília, quæ fecit nimis virgo Israël? Numquid defíciet de petra agri nix Líbani? aut evélli possunt aquæ erumpéntes frígidæ, et defluéntes? Quia oblítus est mei pópulus meus, frustra libántes, et impingéntes in viis suis, in sémitis sǽculi, ut ambulárent per eas in itínere non trito: ut fíeret terra eórum in desolatiónem, et in síbilum sempitérnum: omnis qui præteríerit per eam obstupéscet, et movébit caput suum. Sicut ventus urens dispérgam eos coram inimíco: dorsum, et non fáciem osténdam eis in die perditiónis eórum. Et dixérunt: Veníte et cogitémus contra Jeremíam cogitatiónes: non enim períbit lex a sacerdóte, neque consílium a sapiénte, nec sermo a prophéta: veníte, et percutiámus eum lingua, et non attendámus ad univérsos sermónes ejus.\par
\switchcolumn
\EN
\noindent Therefore thus saith the Lord: Ask among the nations: Who hath heard such horrible things, as the virgin of Israel hath done to excess? Shall the snow of Libanus fail from the rock of the field? or can the cold waters that gush out and run down, be taken away? Because my people have forgotten me, sacrificing in vain, and stumbling in their ways, in ancient paths, to walk by them in a way not trodden: That their land might be given up to desolation, and to a perpetual hissing: every one that shall pass by it, shall be astonished, and wag his head. As a burning will I scatter them before the enemy: I will shew them the back, and not the face, in the day of their destruction. And they said: Come, and let us invent devices against Jeremias: for the law shall not perish from the priest, nor counsel from the wise, nor the word from the prophet: come, and let us strike him with the tongue, and let us give no heed to all his words.\par
\switchcolumn*

\LA
\begin{Resp}\Rsign Dixérunt ímpii apud se, non recte cogitántes: Circumveniámus justum, quóniam contrárius est opéribus nostris: promíttit se sciéntiam Dei habére, Fílium Dei se nóminat, et gloriátur patrem se habére Deum: \Star{} Videámus si sermónes illíus veri sunt: et si est vere Fílius Dei, líberet eum de mánibus nostris: morte turpíssima condemnémus eum.\end{Resp}
\begin{Resp}\Vsign Tamquam nugáces æstimáti sumus ab illo, et ábstinet se a viis nostris tamquam ab immundítiis: et præfert novíssima justórum.\end{Resp}
\begin{Resp}\Rsign Videámus si sermónes illíus veri sunt: et si est vere Fílius Dei, líberet eum de mánibus nostris: morte turpíssima condemnémus eum.\end{Resp}
\switchcolumn
\EN
\begin{Resp}\Rsign The ungodly said, reasoning with themselves, but not aright; Let us lie in wait for the righteous, because he is clean contrary to our doings he professeth to have the knowledge of God, he calleth himself the Son of God, and boasteth that he hath God to his Father. \Star{} Let us see if his words be true; and, if he be indeed the Son of God, let Him deliver him from our hand; let us condemn him with a shameful death.\end{Resp}
\begin{Resp}\Vsign We are esteemed of him as counterfeits, and he abstaineth from our ways as from filthiness, and commendeth the end of the just.\end{Resp}
\begin{Resp}\Rsign Let us see if his words be true; and, if he be indeed the Son of God, let Him deliver him from our hand; let us condemn him with a shameful death.\end{Resp}
\switchcolumn*

\LA
\LessonHead{Lectio III}
\switchcolumn
\EN
\LessonHead{Lesson III}
\switchcolumn*

\LA
\Cite{Jer 18:19-23}
\switchcolumn
\EN
\Cite{Jer 18:19-23}
\switchcolumn*

\LA
\noindent Atténde, Dómine, ad me, et audi vocem adversariórum meórum. Numquid rédditur pro bono malum, quia fodérunt fóveam ánimæ meæ? Recordáre quod stéterim in conspéctu tuo, ut lóquerer pro eis bonum, et avérterem indignatiónem tuam ab eis. Proptérea da fílios eórum in famem, et deduc eos in manus gládii: fiant uxóres eórum absque líberis, et víduæ: et viri eárum interficiántur morte: júvenes eórum confodiántur gládio in prǽlio. Audiátur clamor de dómibus eórum: addúces enim super eos latrónem repénte: quia fodérunt fóveam ut cáperent me, et láqueos abscondérunt pédibus meis. Tu autem, Dómine, scis omne consílium eórum advérsum me in mortem: ne propitiéris iniquitáti eórum, et peccátum eórum a fácie tua non deleátur: fiant corruéntes in conspéctu tuo, in témpore furóris tui abútere eis.\par
\switchcolumn
\EN
\noindent Give heed to me, O Lord, and hear the voice of my adversaries. Shall evil be rendered for good, because they have digged a pit for my soul? Remember that I have stood in thy sight, to speak good for them, and turn away thy indignation from them. Therefore deliver up their children to famine, and bring them into the hands of the sword: let their wives be bereaved of children and widows: and let their husbands be slain by death: let their young men be stabbed with the sword in battle. Let a cry be heard out of their houses: for thou shalt bring the robber upon them suddenly: because they have digged a pit to take me, and have hid snares for my feet. But thou, O Lord, knowest all their counsel against me unto death: forgive not their iniquity, and let not their sin be blotted out from thy sight: let them be overthrown before thy eyes, in the time of thy wrath do thou destroy them.\par
\switchcolumn*

\LA
\begin{Resp}\Rsign Circumdedérunt me viri mendáces: sine causa flagéllis cecidérunt me: \Star{} Sed tu, Dómine defénsor, víndica me.\end{Resp}
\begin{Resp}\Vsign Quóniam tribulátio próxima est, et non est qui ádjuvet.\end{Resp}
\begin{Resp}\Rsign Sed tu, Dómine defénsor, víndica me.\end{Resp}
\begin{Resp}\Vsign Glória Patri, et Fílio, \Star{} et Spirítui Sancto.\end{Resp}
\begin{Resp}\Rsign Circumdedérunt me viri mendáces: sine causa flagéllis cecidérunt me: * Sed tu, Dómine defénsor, víndica me.\end{Resp}
\switchcolumn
\EN
\begin{Resp}\Rsign Liars are come round about me, they have fallen upon me with scourges without a cause. \Star{} But do Thou, O Lord my Redeemer, avenge me!\end{Resp}
\begin{Resp}\Vsign For trouble is near, and there is none to help.\end{Resp}
\begin{Resp}\Rsign But do Thou, O Lord my Redeemer, avenge me!\end{Resp}
\begin{Resp}\Vsign Glory be to the Father, and to the Son, \Star{} and to the Holy Ghost.\end{Resp}
\begin{Resp}\Rsign Liars are come round about me, they have fallen upon me with scourges without a cause. * But do Thou, O Lord my Redeemer, avenge me!\end{Resp}
\switchcolumn*

\end{paracol}

\Office{Feria Quinta in Cena Domini}{Maundy Thursday}{Feria privilegiata Duplex I classis}
\addcontentsline{toc}{subsection}{Feria Quinta in Cena Domini\enspace\textit{Maundy Thursday}}
\begin{paracol}{2}
\LA
\LessonHead{Lectio I}
\switchcolumn
\EN
\LessonHead{Lesson I}
\switchcolumn*

\LA
\LessonTitle{Incipit Lamentátio Jeremíæ Prophétæ}
\switchcolumn
\EN
\LessonTitle{Lesson from the book of Lamentations}
\switchcolumn*

\LA
\Cite{Lam 1:1-5}
\switchcolumn
\EN
\Cite{Lam 1:1-5}
\switchcolumn*

\LA
\noindent Aleph. Quómodo sedet sola cívitas plena pópulo: facta est quasi vídua dómina géntium: princeps provinciárum facta est sub tribúto.\par
\switchcolumn
\EN
\noindent Aleph. How doth the city sit solitary that was full of people! how is the mistress of the Gentiles become as a widow: the princes of provinces made tributary!\par
\switchcolumn*

\LA
Beth. Plorans plorávit in nocte, et lácrimæ ejus in maxíllis ejus: non est qui consolétur eam ex ómnibus caris ejus: omnes amíci ejus sprevérunt eam, et facti sunt ei inimíci.\par
\switchcolumn
\EN
Beth. Weeping she hath wept in the night, and her tears are on her cheeks: there is none to comfort her among all them that were dear to her: all her friends have despised her, and are become her enemies.\par
\switchcolumn*

\LA
Ghimel. Migrávit Judas propter afflictiónem, et multitúdinem servitútis: habitávit inter gentes, nec invénit réquiem: omnes persecutóres ejus apprehendérunt eam inter angústias.\par
\switchcolumn
\EN
Ghimel. Juda hath removed her dwelling place because of her affliction, and the greatness of her bondage: she hath dwelt among the nations, and she hath found no rest: all her persecutors have taken her in the midst of straits.\par
\switchcolumn*

\LA
Daleth. Viæ Sion lugent eo quod non sint qui véniant ad solemnitátem: omnes portæ ejus destrúctæ: sacerdótes ejus geméntes: vírgines ejus squálidæ, et ipsa oppréssa amaritúdine.\par
\switchcolumn
\EN
Daleth. The ways of Sion mourn, because there are none that come to the solemn feast: all her gates are broken down: her priests sigh: her virgins are in affliction, and she is oppressed with bitterness.\par
\switchcolumn*

\LA
He. Facti sunt hostes ejus in cápite, inimíci ejus locupletáti sunt: quia Dóminus locútus est super eam propter multitúdinem iniquitátum ejus: párvuli ejus ducti sunt in captivitátem, ante fáciem tribulántis.\par
\switchcolumn
\EN
He. Her adversaries are become her lords, her enemies are enriched: because the Lord hath spoken against her for the multitude of her iniquities: her children are led into captivity: before the face of the oppressor.\par
\switchcolumn*

\LA
Jerúsalem, Jerúsalem, convértere ad Dóminum Deum tuum.\par
\switchcolumn
\EN
Jerusalem, Jerusalem, return to the Lord thy God.\par
\switchcolumn*

\LA
\begin{Resp}\Rsign In monte Olivéti orávit ad Patrem: Pater, si fíeri potest, tránseat a me calix iste: \Star{} Spíritus quidem promptus est, caro autem infírma.\end{Resp}
\begin{Resp}\Vsign Vigiláte, et oráte, ut non intrétis in tentatiónem.\end{Resp}
\begin{Resp}\Rsign Spíritus quidem promptus est, caro autem infírma.\end{Resp}
\switchcolumn
\EN
\begin{Resp}\Rsign At the Mount of Olives He prayed unto the Father: O My Father, if it be possible, let this cup pass from Me! \Star{} The spirit indeed is willing, but the flesh is weak.\end{Resp}
\begin{Resp}\Vsign Watch and pray, that ye enter not into temptation.\end{Resp}
\begin{Resp}\Rsign The spirit indeed is willing, but the flesh is weak.\end{Resp}
\switchcolumn*

\LA
\LessonHead{Lectio II}
\switchcolumn
\EN
\LessonHead{Lesson II}
\switchcolumn*

\LA
\Cite{Lam 1:6-9}
\switchcolumn
\EN
\Cite{Lam 1:6-9}
\switchcolumn*

\LA
\noindent Vau. Et egréssus est a fília Sion omnis decor ejus: facti sunt príncipes ejus velut aríetes non inveniéntes páscua: et abiérunt absque fortitúdine ante fáciem subsequéntis.\par
\switchcolumn
\EN
\noindent Vau. And from the daughter of Sion all her beauty is departed: her princes are become like rams that find no pastures: and they are gone away without strength before the face of the pursuer.\par
\switchcolumn*

\LA
Zain. Recordáta est Jerúsalem diérum afflictiónis suæ, et prævaricatiónis ómnium desiderabílium suórum, quæ habúerat a diébus antíquis, cum cáderet pópulus ejus in manu hostíli, et non esset auxiliátor: vidérunt eam hostes, et derisérunt sábbata ejus.\par
\switchcolumn
\EN
Zain. Jerusalem hath remembered the days of her affliction, and prevarication of all her desirable things which she had from the days of old, when her people fell in the enemy's hand, and there was no helper: the enemies have seen her, and have mocked at her sabbaths.\par
\switchcolumn*

\LA
Heth. Peccátum peccávit Jerúsalem, proptérea instábilis facta est: omnes, qui glorificábant eam, sprevérunt illam, quia vidérunt ignomíniam ejus: ipsa autem gemens convérsa est retrórsum.\par
\switchcolumn
\EN
Heth. Jerusalem hath grievously sinned, therefore is she become unstable: all that honoured her have despised her, because they have seen her shame: but she sighed and turned backward.\par
\switchcolumn*

\LA
Teth. Sordes ejus in pédibus ejus, nec recordáta est finis sui: depósita est veheménter, non habens consolatórem: vide, Dómine, afflictiónem meam, quóniam eréctus est inimícus.\par
\switchcolumn
\EN
Teth. Her filthiness is on her feet, and she hath not remembered her end: she is wonderfully cast down, not having a comforter: behold, O Lord, my affliction, because the enemy is lifted up.\par
\switchcolumn*

\LA
Jerúsalem, Jerúsalem, convértere ad Dóminum Deum tuum.\par
\switchcolumn
\EN
Jerusalem, Jerusalem, return to the Lord thy God.\par
\switchcolumn*

\LA
\begin{Resp}\Rsign Tristis est ánima mea usque ad mortem: sustinéte hic, et vigiláte mecum: nunc vidébitis turbam, quæ circúmdabit me: \Star{} Vos fugam capiétis, et ego vadam immolári pro vobis.\end{Resp}
\begin{Resp}\Vsign Ecce appropínquat hora, et Fílius hóminis tradétur in manus peccatórum.\end{Resp}
\begin{Resp}\Rsign Vos fugam capiétis, et ego vadam immolári pro vobis.\end{Resp}
\switchcolumn
\EN
\begin{Resp}\Rsign My Soul is exceeding sorrowful, even unto death: tarry ye here and watch with Me \Star{} Yet a little while, and ye shall see the multitude close Me in. Ye shall flee; and I will go to be offered a sacrifice for you.\end{Resp}
\begin{Resp}\Vsign Behold, the hour is at hand, and the Son of man is betrayed into the hands of sinners.\end{Resp}
\begin{Resp}\Rsign Ye shall flee; and I will go to be offered a sacrifice for you.\end{Resp}
\switchcolumn*

\LA
\LessonHead{Lectio III}
\switchcolumn
\EN
\LessonHead{Lesson III}
\switchcolumn*

\LA
\Cite{Lam 1:10-14}
\switchcolumn
\EN
\Cite{Lam 1:10-14}
\switchcolumn*

\LA
\noindent Jod. Manum suam misit hostis ad ómnia desiderabília ejus: quia vidit gentes ingréssas sanctuárium suum, de quibus præcéperas ne intrárent in ecclésiam tuam.\par
\switchcolumn
\EN
\noindent Jod. The enemy hath put out his hand to all her desirable things: for she hath seen the Gentiles enter into her sanctuary, of whom thou gavest commandment that they should not enter into thy church.\par
\switchcolumn*

\LA
Caph. Omnis pópulus ejus gemens, et quærens panem: dedérunt pretiósa quæque pro cibo ad refocillándam ánimam. Vide, Dómine, et consídera, quóniam facta sum vilis.\par
\switchcolumn
\EN
Caph. All her people sigh, they seek bread: they have given all their precious things for food to relieve the soul: see, O Lord, and consider, for I am become vile.\par
\switchcolumn*

\LA
Lamed. O vos omnes, qui transítis per viam, atténdite, et vidéte, si est dolor sicut dolor meus: quóniam vindemiávit me, ut locútus est Dóminus in die iræ furóris sui.\par
\switchcolumn
\EN
Lamed. O all ye that pass by the way, attend, and see if there be any sorrow like to my sorrow: for he hath made a vintage of me, as the Lord spoke in the day of his fierce anger.\par
\switchcolumn*

\LA
Mem. De excélso misit ignem in óssibus meis, et erudívit me: expándit rete pédibus meis, convértit me retrórsum: pósuit me desolátam, tota die mæróre conféctam.\par
\switchcolumn
\EN
Mem. From above he hath sent fire into my bones, and hath chastised me: he hath spread a net for my feet, he hath turned me back: he hath made me desolate, wasted with sorrow all the day long.\par
\switchcolumn*

\LA
Nun. Vigilávit jugum iniquitátum meárum: in manu ejus convolútæ sunt, et impósitæ collo meo: infirmáta est virtus mea: dedit me Dóminus in manu, de qua non pótero súrgere.\par
\switchcolumn
\EN
Nun. The yoke of my iniquities hath watched: they are folded together in his hand, and put upon my neck: my strength is weakened: the Lord hath delivered me into a hand out of which I am not able to rise.\par
\switchcolumn*

\LA
Jerúsalem, Jerúsalem, convértere ad Dóminum Deum tuum.\par
\switchcolumn
\EN
Jerusalem, Jerusalem, return to the Lord thy God.\par
\switchcolumn*

\LA
\begin{Resp}\Rsign Ecce vídimus eum non habéntem spéciem, neque decórem: aspéctus ejus in eo non est: hic peccáta nostra portávit, et pro nobis dolet: ipse autem vulnerátus est propter iniquitátes nostras: \Star{} Cujus livóre sanáti sumus.\end{Resp}
\begin{Resp}\Vsign Vere languóres nostros ipse tulit, et dolóres nostros ipse portávit.\end{Resp}
\begin{Resp}\Rsign Cujus livóre sanáti sumus.\end{Resp}
\begin{Resp}\Vsign Glória Patri, et Fílio, \Star{} et Spirítui Sancto.\end{Resp}
\begin{Resp}\Rsign Ecce vídimus eum non habéntem spéciem, neque decórem: aspéctus ejus in eo non est: hic peccáta nostra portávit, et pro nobis dolet: ipse autem vulnerátus est propter iniquitátes nostras: * Cujus livóre sanáti sumus.\end{Resp}
\switchcolumn
\EN
\begin{Resp}\Rsign Behold, when we shall see Him, He hath no form nor comeliness: there is no beauty in Him: this is He Which hath borne our griefs and carried our sorrows; but He was wounded for our transgressions \Star{} And with His stripes we are healed.\end{Resp}
\begin{Resp}\Vsign Surely He hath borne our griefs and carried our sorrows.\end{Resp}
\begin{Resp}\Rsign And with His stripes we are healed.\end{Resp}
\begin{Resp}\Vsign Glory be to the Father, and to the Son, \Star{} and to the Holy Ghost.\end{Resp}
\begin{Resp}\Rsign Behold, when we shall see Him, He hath no form nor comeliness: there is no beauty in Him; this is He Which hath borne our sins and carried our sorrows: but He was wounded for our transgressions, and with His stripes we are healed.\end{Resp}
\switchcolumn*

\LA
\LessonHead{Lectio IV}
\switchcolumn
\EN
\LessonHead{Lesson IV}
\switchcolumn*

\LA
\LessonTitle{Ex tractátu sancti Augustíni Epíscopi super Psalmos}
\switchcolumn
\EN
\LessonTitle{From the Treatise of St. Augustine, Bishop (of Hippo) Upon the Psalms}
\switchcolumn*

\LA
\Source{In Psalmum 54 ad 1 versum}
\switchcolumn
\EN
\Source{On Psalm liv, 1}
\switchcolumn*

\LA
\noindent Exáudi, Deus, oratiónem meam, et ne despéxeris deprecatiónem meam: inténde mihi, et exáudi me. Satagéntis, sollíciti, in tribulatióne pósiti, verba sunt ista. Orat multa pátiens, de malo liberári desíderans. Súperest ut videámus in quo malo sit: et cum dícere cœ́perit, agnoscámus ibi nos esse: ut communicáta tribulatióne, conjungámus oratiónem. Contristátus sum, inquit, in exercitatióne mea, et conturbátus sum. Ubi contristátus? ubi conturbátus? In exercitatióne mea, inquit. Hómines malos, quos pátitur, commemorátus est: eandémque passiónem malórum hóminum exercitatiónem suam dixit. Ne putétis gratis esse malos in hoc mundo, et nihil boni de illis ágere Deum. Omnis malus aut ídeo vivit, ut corrigátur; aut ídeo vivit, ut per illum bonus exerceátur.\par
\switchcolumn
\EN
\noindent Give ear to my prayer, O God, and despise not my supplication: attend unto me and hear me. These are the words of a man travailing, anxious, and troubled. He prayeth in the midst of much suffering, longing to be rid of his affliction. Our part is to see what that his affliction was, and when he hath told us, to acknowledge that we also suffer therefrom; that so, partaking in his trouble, we may take part also in his exercise, and am troubled. Wherein mourned he? Wherein was he troubled? He saith: In my exercise. In the next words he giveth us to know that his affliction was the oppression of the wicked, because of the voice of the enemy, and because of the oppression of the wicked, and this suffering which came upon him at the hands of wicked men, he hath called his exercise. Think not that wicked men are in this world for nothing, or that God doth no good with them. Every wicked man liveth, either to repent, or to exercise the righteous.\par
\switchcolumn*

\LA
\begin{Resp}\Rsign Amicus meus ósculi me trádidit signo: Quem osculátus fúero, ipse est, tenéte eum: hoc malum fecit signum, qui per ósculum adimplévit homicídium. \Star{} Infélix prætermísit prétium sánguinis, et in fine láqueo se suspéndit.\end{Resp}
\begin{Resp}\Vsign Bonum erat ei, si natus non fuísset homo ille.\end{Resp}
\begin{Resp}\Rsign Infélix prætermísit prétium sánguinis, et in fine láqueo se suspéndit.\end{Resp}
\switchcolumn
\EN
\begin{Resp}\Rsign Mine own friend hath betrayed Me by the sign of a kiss: Whomsoever I shall kiss, That Same is He; hold Him fast. This was the traitorous sign which he gave, even he who murdered with a kiss. \Star{} Woe unto that man! He cast down the price of blood, and went, and hanged himself.\end{Resp}
\begin{Resp}\Vsign It had been good for that man if he had not been born.\end{Resp}
\begin{Resp}\Rsign Woe unto that man! He cast down the price of blood, and went, and hanged himself.\end{Resp}
\switchcolumn*

\LA
\LessonHead{Lectio V}
\switchcolumn
\EN
\LessonHead{Lesson V}
\switchcolumn*

\LA
\noindent Utinam ergo qui nos modo exércent, convertántur, et nobíscum exerceántur: tamen quámdiu ita sunt ut exérceant, non eos odérimus: quia in eo quod malus est quis eórum, utrum usque in finem perseveratúrus sit, ignorámus. Et plerúmque cum tibi vidéris odísse inimícum, fratrem odísti, et nescis. Diábolus, et ángeli ejus in Scriptúris sanctis manifestáti sunt nobis, quod ad ignem ætérnum sint destináti. Ipsórum tantum desperánda est corréctio, contra quos habémus occúltam luctam: ad quam luctam nos armat Apóstolus, dicens: Non est nobis colluctátio advérsus carnem et sánguinem: id est, non advérsus hómines, quos vidétis, sed advérsus príncipes, et potestátes, et rectóres mundi, tenebrárum harum. Ne forte cum dixísset, mundi, intellégeres dǽmones esse rectóres cæli et terræ. Mundi dixit, tenebrárum harum: mundi dixit, amatórum mundi: mundi dixit, impiórum et iniquórum: mundi dixit, de quo dicit Evangélium: Et mundus eum non cognóvit.\par
\switchcolumn
\EN
\noindent Would to God that they which now exercise us were converted and exercised with us! Yet, while they are as they are, and exercise us, we will not hate them: for we know not of any one of them whether he will endure to the end in his sin. Yea, oftentimes, when thou deemest that thou hatest thine enemy, he whom thou hatest is thy brother, and thou knowest it not. The Holy Scriptures show us that the devil and his angels are already damned unto everlasting fire, and therefore of their repentance it behoveth us to despair; but of theirs only. These are they against whom we wrestle within; to the which wrestling the Apostle stirreth us up where he saith: We wrestle not against flesh and blood, (that is, not against men whom we see,) but against principalities, against powers, against the rulers of the darkness of this world. (Eph. vi. 12.) He saith not the rulers of this world, lest perchance thou shouldest deem that devils are the lords of heaven and earth; what he doth say is, rulers of the darkness of this world, of that world which they love who love the world, of that world wherein the ungodly and unrighteous do prosper, of that world, in fine, of which the Gospel saith: And the world knew Him not.\par
\switchcolumn*

\LA
\begin{Resp}\Rsign Judas mercátor péssimus ósculo pétiit Dóminum: ille ut agnus ínnocens non negávit Judæ ósculum: \Star{} Denariórum número Christum Judǽis trádidit.\end{Resp}
\begin{Resp}\Vsign Mélius illi erat, si natus non fuísset.\end{Resp}
\begin{Resp}\Rsign Denariórum número Christum Judǽis trádidit.\end{Resp}
\switchcolumn
\EN
\begin{Resp}\Rsign The vile trader Judas came to the Lord to kiss Him, and He, as a guileless Lamb, refused not a kiss to Judas, \Star{} Who, for a certain number of pence, betrayed Christ to the Jews.\end{Resp}
\begin{Resp}\Vsign It had been good for that man if he had not been born.\end{Resp}
\begin{Resp}\Rsign Who, for a certain number of pence, betrayed Christ to the Jews.\end{Resp}
\switchcolumn*

\LA
\LessonHead{Lectio VI}
\switchcolumn
\EN
\LessonHead{Lesson VI}
\switchcolumn*

\LA
\noindent Quóniam vidi iniquitátem, et contradictiónem in civitáte. Atténde glóriam crucis ipsíus. Jam in fronte regum crux illa fixa est, cui inimíci insultavérunt. Efféctus probávit virtútem: dómuit orbem non ferro, sed ligno. Lignum crucis contuméliis dignum visum est inimícis, et ante ipsum lignum stantes caput agitábant, et dicébant: Si Fílius Dei est, descéndat de cruce. Extendébat ille manus suas ad pópulum non credéntem, et contradicéntem. Si enim justus est, qui ex fide vivit; iníquus est, qui non habet fidem. Quod ergo hic ait, iniquitátem: perfídiam intéllege. Vidébat ergo Dóminus in civitáte iniquitátem et contradictiónem, et extendébat manus suas ad pópulum non credéntem et contradicéntem: et tamen et ipsos exspéctans dicébat: Pater, ignósce illis, quia nésciunt quid fáciunt.\par
\switchcolumn
\EN
\noindent We have seen iniquity and strife in the city. (v. 10.) Behold, the glory of the Cross. That Cross which was the object of the insults of God's enemies, is established now above the brows of kings. The end hath shown the measure of its power: it hath conquered the world, not by the sword, but by its wood. The enemies of God thought the Cross a meet object of insult and ridicule, yea, they stood before it, wagging their heads and saying: If He be the Son of God, let Him come down from the Cross! (Matth. xxvii. 39, 40.) And He stretched forth His Hands unto a disobedient and gainsaying people. (Rom. x. 21.) If he is just which liveth by faith, (Rom. i. 17; Hab. ii. 4,) he is unjust that hath not faith. Therefore where is written iniquity we may understand unbelief. The Lord therefore saith that He saw iniquity and strife in the city, and that He stretched forth His Hands unto that disobedient and gainsaying people, and, disobedient and gainsaying as they were, He was hungry for their salvation, and said: Father, forgive them, for they know not what they do. (Luke xxiii. 34.)\par
\switchcolumn*

\LA
\begin{Resp}\Rsign Unus ex discípulis meis tradet me hódie: Væ illi per quem tradar ego: \Star{} Mélius illi erat, si natus non fuísset.\end{Resp}
\begin{Resp}\Vsign Qui intíngit mecum manum in parópside, hic me traditúrus est in manus peccatórum.\end{Resp}
\begin{Resp}\Rsign Mélius illi erat, si natus non fuísset.\end{Resp}
\begin{Resp}\Vsign Glória Patri, et Fílio, \Star{} et Spirítui Sancto.\end{Resp}
\begin{Resp}\Rsign Unus ex discípulis meis tradet me hódie: Væ illi per quem tradar ego: * Mélius illi erat, si natus non fuísset.\end{Resp}
\switchcolumn
\EN
\begin{Resp}\Rsign One of My disciples shall betray Me this night. Woe unto that man by whom I am betrayed! \Star{} It had been good for that man if he had not been born.\end{Resp}
\begin{Resp}\Vsign He that dippeth his hand with Me in the dish, the same shall betray Me into the hands of sinners.\end{Resp}
\begin{Resp}\Rsign It had been good for that man if he had not been born.\end{Resp}
\begin{Resp}\Vsign Glory be to the Father, and to the Son, \Star{} and to the Holy Ghost.\end{Resp}
\begin{Resp}\Rsign One of My disciples shall betray Me this night. Woe unto that man by whom I am betrayed. * It had been good for that man if he had not been born.\end{Resp}
\switchcolumn*

\LA
\LessonHead{Lectio VII}
\switchcolumn
\EN
\LessonHead{Lesson VII}
\switchcolumn*

\LA
\LessonTitle{De Epístola prima beáti Pauli Apóstoli ad Corínthios}
\switchcolumn
\EN
\LessonTitle{From the first letter of blessed Apostle Paul to Corinthians}
\switchcolumn*

\LA
\Cite{1 Cor 11:17-22}
\switchcolumn
\EN
\Cite{1 Cor 11:17-22}
\switchcolumn*

\LA
\noindent Hoc autem præcípio: non laudans quod non in mélius, sed in detérius convenítis. Primum quidem conveniéntibus vobis in Ecclésiam, áudio scissúras esse inter vos, et ex parte credo. Nam opórtet et hǽreses esse, ut et qui probáti sunt, manifésti fiant in vobis. Conveniéntibus ergo vobis in unum, jam non est Domínicam cenam manducáre. Unusquísque enim suam cenam præsúmit ad manducándum. Et álius quidem ésurit, álius autem ébrius est. Numquid domos non habétis ad manducándum et bibéndum? aut Ecclésiam Dei contémnitis, et confúnditis eos, qui non habent? Quid dicam vobis? Laudo vos? In hoc non laudo.\par
\switchcolumn
\EN
\noindent Now this I ordain: not praising you, that you come together not for the better, but for the worse. For first of all I hear that when you come together in the church, there are schisms among you; and in part I believe it. For there must be also heresies: that they also, who are approved, may be made manifest among you. When you come therefore together into one place, it is not now to eat the Lord's supper. For every one taketh before his own supper to eat. And one indeed is hungry and another is drunk. What, have you not houses to eat and to drink in? Or despise ye the church of God; and put them to shame that have not? What shall I say to you? Do I praise you? In this I praise you not.\par
\switchcolumn*

\LA
\begin{Resp}\Rsign Eram quasi agnus ínnocens: ductus sum ad immolándum, et nesciébam: consílium fecérunt inimíci mei advérsum me, dicéntes: \Star{} Veníte, mittámus lignum in panem ejus, et eradámus eum de terra vivéntium.\end{Resp}
\begin{Resp}\Vsign Omnes inimíci mei advérsum me cogitábant mala mihi: verbum iníquum mandavérunt advérsum me, dicéntes.\end{Resp}
\begin{Resp}\Rsign Veníte, mittámus lignum in panem ejus, et eradámus eum de terra vivéntium.\end{Resp}
\switchcolumn
\EN
\begin{Resp}\Rsign I was like a gentle lamb that is brought to the slaughter, and I knew not that mine enemies had devised devices against me, saying: \Star{} Come, let us put (poison of a deadly) tree into his bread, and let us cut him off from the land of the living.\end{Resp}
\begin{Resp}\Vsign All they that hate me devised my hurt against me: they plotted together to do me evil, saying:\end{Resp}
\begin{Resp}\Rsign Come, let us put (poison of a deadly) tree into his bread, and let us cut him off from the land of the living.\end{Resp}
\switchcolumn*

\LA
\LessonHead{Lectio VIII}
\switchcolumn
\EN
\LessonHead{Lesson VIII}
\switchcolumn*

\LA
\Cite{1 Cor 11:23-26}
\switchcolumn
\EN
\Cite{1 Cor 11:23-26}
\switchcolumn*

\LA
\noindent Ego enim accépi a Dómino quod et trádidi vobis, quóniam Dóminus Jesus, in qua nocte tradebátur, accépit panem, et grátias agens fregit, et dixit: Accípite, et manducáte: hoc est corpus meum, quod pro vobis tradétur: hoc fácite in meam commemoratiónem. Simíliter et cálicem, postquam cœnávit, dicens: Hic calix novum testaméntum est in meo sánguine: hoc fácite, quotiescúmque bibétis, in meam commemoratiónem. Quotiescúmque enim manducábitis panem hunc, et cálicem bibétis, mortem Dómini annuntiábitis donec véniat.\par
\switchcolumn
\EN
\noindent For I have received of the Lord that which also I delivered unto you, that the Lord Jesus, the same night in which he was betrayed, took bread. And giving thanks, broke, and said: Take ye, and eat: this is my body, which shall be delivered for you: this do for the commemoration of me. In like manner also the chalice, after he had supped, saying: This chalice is the new testament in my blood: this do ye, as often as you shall drink, for the commemoration of me. For as often as you shall eat this bread, and drink the chalice, you shall show the death of the Lord, until he come.\par
\switchcolumn*

\LA
\begin{Resp}\Rsign Una hora non potuístis vigiláre mecum, qui exhortabámini mori pro me? \Star{} Vel Judam non vidétis, quómodo non dormit, sed festínat trádere me Judǽis.\end{Resp}
\begin{Resp}\Vsign Quid dormítis? súrgite, et oráte, ne intrétis in tentatiónem.\end{Resp}
\begin{Resp}\Rsign Vel Judam non vidétis, quómodo non dormit, sed festínat trádere me Judǽis.\end{Resp}
\switchcolumn
\EN
\begin{Resp}\Rsign Could ye not watch with Me one hour, ye that called one on the other to die for Me? \Star{} Or see ye not Judas, how that he sleepeth not, but maketh haste to betray Me to the Jews?\end{Resp}
\begin{Resp}\Vsign Why sleep ye? Rise, and pray, lest ye enter into temptation.\end{Resp}
\begin{Resp}\Rsign Or see ye not Judas, how that he sleepeth not, but maketh haste to betray Me to the Jews?\end{Resp}
\switchcolumn*

\LA
\LessonHead{Lectio IX}
\switchcolumn
\EN
\LessonHead{Lesson IX}
\switchcolumn*

\LA
\Cite{1 Cor 11:27-34}
\switchcolumn
\EN
\Cite{1 Cor 11:27-34}
\switchcolumn*

\LA
\noindent Itaque quicúmque manducáverit panem hunc, vel bíberit cálicem Dómini indígne, reus erit córporis et sánguinis Dómini. Probet autem seípsum homo: et sic de pane illo edat, et de cálice bibat. Qui enim mandúcat et bibit indígne, judícium sibi mandúcat et bibit, non dijúdicans corpus Dómini. Ideo inter vos multi infírmi et imbecílles, et dórmiunt multi. Quod, si nosmetípsos dijudicarémus, non útique judicarémur. Dum judicámur autem, a Dómino corrípimur, ut non cum hoc mundo damnémur. Itaque, fratres mei, cum convenítis ad manducándum, ínvicem exspectáte. Si quis ésurit, domi mandúcet: ut non in judícium conveniátis. Cétera autem, cum vénero, dispónam.\par
\switchcolumn
\EN
\noindent Therefore whosoever shall eat this bread, or drink the chalice of the Lord unworthily, shall be guilty of the body and of the blood of the Lord. But let a man prove himself: and so let him eat of that bread, and drink of the chalice. For he that eateth and drinketh unworthily, eateth and drinketh judgment to himself, not discerning the body of the Lord. Therefore are there many infirm and weak among you, and many sleep. But if we would judge ourselves, we should not be judged. But whilst we are judged, we are chastised by the Lord, that we be not condemned with this world. Wherefore, my brethren, when you come together to eat, wait for one another. If any man be hungry, let him eat at home; that you come not together unto judgment. And the rest I will set in order, when I come.\par
\switchcolumn*

\LA
\begin{Resp}\Rsign Senióres pópuli consílium fecérunt, \Star{} Ut Jesum dolo tenérent, et occíderent: cum gládiis et fústibus exiérunt tamquam ad latrónem.\end{Resp}
\begin{Resp}\Vsign Collegérunt pontífices et pharisǽi concílium.\end{Resp}
\begin{Resp}\Rsign Ut Jesum dolo tenérent, et occíderent: cum gládiis et fústibus exiérunt tamquam ad latrónem.\end{Resp}
\begin{Resp}\Vsign Glória Patri, et Fílio, \Star{} et Spirítui Sancto.\end{Resp}
\begin{Resp}\Rsign Senióres pópuli consílium fecérunt, * Ut Jesum dolo tenérent, et occíderent: cum gládiis et fústibus exiérunt tamquam ad latrónem.\end{Resp}
\switchcolumn
\EN
\begin{Resp}\Rsign The elders of the people consulted \Star{} That they might take Jesus by subtilty, and kill Him they came out, as against a thief, with swords and staves.\end{Resp}
\begin{Resp}\Vsign The chief Priests and the Pharisees gathered a council.\end{Resp}
\begin{Resp}\Rsign That they might take Jesus by subtilty, and kill Him: they came out, as against a thief, with swords and staves.\end{Resp}
\begin{Resp}\Vsign Glory be to the Father, and to the Son, \Star{} and to the Holy Ghost.\end{Resp}
\begin{Resp}\Rsign The elders of the people consulted * That they might take Jesus by subtilty, and kill Him: they came out, as against a thief, with swords and staves.\end{Resp}
\switchcolumn*

\end{paracol}

\Office{Feria Sexta in Parasceve}{Good Friday}{Feria privilegiata Duplex I classis}
\addcontentsline{toc}{subsection}{Feria Sexta in Parasceve\enspace\textit{Good Friday}}
\begin{paracol}{2}
\LA
\LessonHead{Lectio I}
\switchcolumn
\EN
\LessonHead{Lesson I}
\switchcolumn*

\LA
\LessonTitle{De Lamentatióne Jeremíæ Prophétæ}
\switchcolumn
\EN
\LessonTitle{Lesson from the book of Lamentations}
\switchcolumn*

\LA
\Cite{Lam 2:8-11}
\switchcolumn
\EN
\Cite{Lam 2:8-11}
\switchcolumn*

\LA
\noindent Heth. Cogitávit Dóminus dissipáre murum fíliæ Sion: teténdit funículum suum, et non avértit manum suam a perditióne: luxítque antemurále, et murus páriter dissipátus est.\par
\switchcolumn
\EN
\noindent Heth. The Lord hath purposed to destroy the wall of the daughter of Sion: he hath stretched out his line, and hath not withdrawn his hand from destroying: and the bulwark hath mourned, and the wall hath been destroyed together.\par
\switchcolumn*

\LA
Teth. Defíxæ sunt in terra portæ ejus: pérdidit, et contrívit vectes ejus: regem ejus et príncipes ejus in géntibus: non est lex, et prophétæ ejus non invenérunt visiónem a Dómino.\par
\switchcolumn
\EN
Teth. Her gates are sunk into the ground: he hath destroyed, and broken her bars: her king and her princes are among the Gentiles: the law is no more, and her prophets have found no vision from the Lord.\par
\switchcolumn*

\LA
Jod. Sedérunt in terra, conticuérunt senes fíliæ Sion: conspersérunt cínere cápita sua, accíncti sunt cilíciis, abjecérunt in terram cápita sua vírgines Jerúsalem.\par
\switchcolumn
\EN
Jod. The ancients of the daughter of Sion sit upon the ground, they have held their peace: they have sprinkled their heads with dust, they are girded with haircloth, the virgins of Jerusalem hang down their heads to the ground.\par
\switchcolumn*

\LA
Caph. Defecérunt præ lácrimis óculi mei, conturbáta sunt víscera mea: effúsum est in terra jecur meum super contritióne fíliæ pópuli mei, cum defíceret párvulus et lactens in platéis óppidi.\par
\switchcolumn
\EN
Caph. My eyes have failed with weeping, my bowels are troubled: my liver is poured out upon the earth, for the destruction of the daughter of my people, when the children, and the sucklings, fainted away in the streets of the city.\par
\switchcolumn*

\LA
Jerúsalem, Jerúsalem, convértere ad Dóminum Deum tuum.\par
\switchcolumn
\EN
Jerusalem! Jerusalem! Return unto the Lord thy God.\par
\switchcolumn*

\LA
\begin{Resp}\Rsign Omnes amíci mei dereliquérunt me, et prævaluérunt insidiántes mihi: trádidit me quem diligébam: \Star{} Et terribílibus óculis plaga crudéli percutiéntes, acéto potábant me.\end{Resp}
\begin{Resp}\Vsign Inter iníquos projecérunt me, et non pepercérunt ánimæ meæ.\end{Resp}
\begin{Resp}\Rsign Et terribílibus óculis plaga crudéli percutiéntes, acéto potábant me.\end{Resp}
\switchcolumn
\EN
\begin{Resp}\Rsign All my friends have forsaken me, and mine enemies have prevailed against me; he whom I loved hath betrayed me. \Star{} Mine enemy sharpeneth his eyes upon me; he breaketh me with breach upon breach: and in my thirst they gave me vinegar to drink.\end{Resp}
\begin{Resp}\Vsign I am numbered with the transgressors; and my life is not spared.\end{Resp}
\begin{Resp}\Rsign Mine enemy sharpeneth his eyes upon me; he breaketh me with breach upon breach; and (in my thirst) they gave me vinegar to drink.\end{Resp}
\switchcolumn*

\LA
\LessonHead{Lectio II}
\switchcolumn
\EN
\LessonHead{Lesson II}
\switchcolumn*

\LA
\Cite{Lam 2:12-15}
\switchcolumn
\EN
\Cite{Lam 2:12-15}
\switchcolumn*

\LA
\noindent Lamed. Mátribus suis dixérunt: Ubi est tríticum et vinum? cum defícerent quasi vulneráti in platéis civitátis: cum exhalárent ánimas suas in sinu matrum suárum.\par
\switchcolumn
\EN
\noindent Lamed. They said to their mothers: Where is corn and wine? when they fainted away as the wounded in the streets of the city: when they breathed out their souls in the bosoms of their mothers.\par
\switchcolumn*

\LA
Mem. Cui comparábo te? vel cui assimilábo te, fília Jerúsalem? cui exæquábo te, et consolábor te, virgo fília Sion? Magna est enim velut mare contrítio tua: quis medébitur tui?\par
\switchcolumn
\EN
Mem. To what shall I compare thee? or to what shall I liken thee, O daughter of Jerusalem? to what shall I equal thee, that I may comfort thee, O virgin daughter of Sion? for great as the sea is thy destruction: who shall heal thee?\par
\switchcolumn*

\LA
Nun. Prophétæ tui vidérunt tibi falsa et stulta, nec aperiébant iniquitátem tuam, ut te ad pœniténtiam provocárent: vidérunt autem tibi assumptiónes falsas, et ejectiónes.\par
\switchcolumn
\EN
Nun. thy prophets have seen false and foolish things for thee: and they have not laid open thy iniquity, to excite thee to penance: but they have seen for thee false revelations and banishments.\par
\switchcolumn*

\LA
Samech. Plausérunt super te mánibus omnes transeúntes per viam: sibilavérunt, et movérunt caput suum super fíliam Jerúsalem: Hǽccine est urbs, dicéntes, perfécti decóris, gáudium univérsæ terræ?\par
\switchcolumn
\EN
Samech. All they that passed by the way have clapped their hands at thee: they have hissed, and wagged their heads at the daughter of Jerusalem, saying: Is this the city of perfect beauty, the joy of all the earth?\par
\switchcolumn*

\LA
Jerúsalem, Jerúsalem, convértere ad Dóminum Deum tuum.\par
\switchcolumn
\EN
Jerusalem! Jerusalem! Return unto the Lord thy God.\par
\switchcolumn*

\LA
\begin{Resp}\Rsign Velum templi scissum est, \Star{} Et omnis terra trémuit: latro de cruce clamábat, dicens: Meménto mei, Dómine, dum véneris in regnum tuum.\end{Resp}
\begin{Resp}\Vsign Petræ scissæ sunt, et monuménta apérta sunt, et multa córpora sanctórum, qui dormíerant, surrexérunt.\end{Resp}
\begin{Resp}\Rsign Et omnis terra trémuit: latro de cruce clamábat, dicens: Meménto mei, Dómine, dum véneris in regnum tuum.\end{Resp}
\switchcolumn
\EN
\begin{Resp}\Rsign The veil of the Temple was rent in twain, from the top to the bottom, \Star{} And all the earth did quake: the thief on the cross cried, saying: Lord, remember me when Thou comest into thy kingdom!\end{Resp}
\begin{Resp}\Vsign The rocks rent, and the graves were opened, and many bodies of the saints, which slept, arose.\end{Resp}
\begin{Resp}\Rsign And all the earth did quake: the thief on the cross cried, saying: Lord, remember me when thou comest into thy kingdom.\end{Resp}
\switchcolumn*

\LA
\LessonHead{Lectio III}
\switchcolumn
\EN
\LessonHead{Lesson III}
\switchcolumn*

\LA
\Cite{Lam 3:1-9}
\switchcolumn
\EN
\Cite{Lam 3:1-9}
\switchcolumn*

\LA
\noindent Aleph. Ego vir videns paupertátem meam in virga indignatiónis ejus.\par
\switchcolumn
\EN
\noindent Aleph. I am the man that see my poverty by the rod of his indignation.\par
\switchcolumn*

\LA
Aleph. Me minávit, et addúxit in ténebras, et non in lucem.\par
\switchcolumn
\EN
Aleph. He hath led me, and brought me into darkness, and not into light.\par
\switchcolumn*

\LA
Aleph. Tantum in me vertit, et convértit manum suam tota die.\par
\switchcolumn
\EN
Aleph. Only against me he hath turned, and turned again his hand all the day.\par
\switchcolumn*

\LA
Beth. Vetústam fecit pellem meam, et carnem meam, contrívit ossa mea.\par
\switchcolumn
\EN
Beth. My skin and my flesh he hath made old, he hath broken my bones.\par
\switchcolumn*

\LA
Beth. Ædificávit in gyro meo, et circúmdedit me felle et labóre.\par
\switchcolumn
\EN
Beth. He hath built round about me, and he hath compassed me with gall and labour.\par
\switchcolumn*

\LA
Beth. In tenebrósis collocávit me, quasi mórtuos sempitérnos.\par
\switchcolumn
\EN
Beth. He hath set me in dark places as those that are dead for ever.\par
\switchcolumn*

\LA
Ghimel. Circumædificávit advérsum me, ut non egrédiar: aggravávit cómpedem meum.\par
\switchcolumn
\EN
Ghimel. He hath built against me round about, that I may not get out: he hath made my fetters heavy.\par
\switchcolumn*

\LA
Ghimel. Sed et, cum clamávero et rogávero, exclúsit oratiónem meam.\par
\switchcolumn
\EN
Ghimel. Yea, and when I cry, and entreat, he hath shut out my prayer.\par
\switchcolumn*

\LA
Ghimel. Conclúsit vias meas lapídibus quadris, sémitas meas subvértit.\par
\switchcolumn
\EN
Ghimel. He hath shut up my ways with square stones, he hath turned my paths upside down\par
\switchcolumn*

\LA
Jerúsalem, Jerúsalem, convértere ad Dóminum Deum tuum.\par
\switchcolumn
\EN
Jerusalem! Jerusalem! Return unto the Lord thy God.\par
\switchcolumn*

\LA
\begin{Resp}\Rsign Vínea mea elécta, ego te plantávi: \Star{} Quómodo convérsa es in amaritúdinem, ut me crucifígeres et Barábbam dimítteres?\end{Resp}
\begin{Resp}\Vsign Sepívi te, et lápides elégi ex te, et ædificávi turrim.\end{Resp}
\begin{Resp}\Rsign Quómodo convérsa es in amaritúdinem, ut me crucifígeres et Barábbam dimítteres?\end{Resp}
\begin{Resp}\Vsign Glória Patri, et Fílio, \Star{} et Spirítui Sancto.\end{Resp}
\begin{Resp}\Rsign Vínea mea elécta, ego te plantávi: * Quómodo convérsa es in amaritúdinem, ut me crucifígeres et Barábbam dimítteres?\end{Resp}
\switchcolumn
\EN
\begin{Resp}\Rsign I had planted thee a noble vineyard; \Star{} How then art thou turned into a degenerate plant, which willest that Barabbas should be released unto thee, and that I should be crucified.\end{Resp}
\begin{Resp}\Vsign I fenced thee, and gathered out the stones from thee, and built a tower in the midst of the land.\end{Resp}
\begin{Resp}\Rsign How then art thou turned into a degenerate plant, which willest that Barabbas should be released unto thee, and that I should be crucified.\end{Resp}
\begin{Resp}\Vsign Glory be to the Father, and to the Son, \Star{} and to the Holy Ghost.\end{Resp}
\begin{Resp}\Rsign I had planted thee a noble vineyard; * How then art thou turned into a degenerate plant, which willest that Barabbas should be released unto thee, and that I should be crucified.\end{Resp}
\switchcolumn*

\LA
\LessonHead{Lectio IV}
\switchcolumn
\EN
\LessonHead{Lesson IV}
\switchcolumn*

\LA
\LessonTitle{Ex tractátu sancti Augustíni Epíscopi super Psalmos}
\switchcolumn
\EN
\LessonTitle{From the Treatise of St. Augustine, Bishop (of Hippo,) Upon the Psalms}
\switchcolumn*

\LA
\Source{In Psalm 63 ad versum 2}
\switchcolumn
\EN
\Source{On Psalm lxiii, 2}
\switchcolumn*

\LA
\noindent Protexísti me, Deus, a convéntu malignántium, a multitúdine operántium iniquitátem. Jam ipsum caput nostrum intueámur. Multi Mártyres tália passi sunt, sed nihil sic elúcet, quómodo caput Mártyrum: ibi mélius intuémur, quod illi expérti sunt. Protéctus est a multitúdine malignántium, protegénte se Deo, protegénte carnem suam ipso Fílio, et hómine, quem gerébat: quia fílius hóminis est, et Fílius Dei est. Fílius Dei, propter formam Dei: fílius hóminis, propter formam servi, habens in potestáte pónere ánimam suam, et recípere eam. Quid ei potuérunt fácere inimíci? Occidérunt corpus, ánimam non occidérunt. Inténdite. Parum ergo erat, Dóminum hortári Mártyres verbo, nisi firmáret exémplo.\par
\switchcolumn
\EN
\noindent Thou hast hidden me from the secret counsel of the wicked, from the insurrection of the workers of iniquity. Now let us fix our eyes upon our Head. Many martyrs have suffered such things as He suffered, but God's hiding of His suffering servants is not so well seen in the Martyrs, as it is in the Captain of the Martyrs. And it is in Him that we best see how it fared with them. He was hidden from the secret counsel of the wicked; hidden by God, being Himself God; hidden, as touching the Manhood, by God the Son, and the very Manhood, Which is taken into God the Son; because He is the Son of man, and He is the Son of God; Son of God, as being in the form of God; Son of man, as having taken upon Him the form of a servant, (Phil. ii. 6, 7,) Whose life no man taketh from Him, but Who layeth it down of Himself. He hath power to lay it down, and He hath power to take it again, (John x. 18.) What then was all that they which hated Him could do? They could kill the Body, but they were not able to kill the Soul. Consider this very earnestly. It had been a small thing for the Lord to preach to the Martyrs by His word, if He had not also nerved them by His example.\par
\switchcolumn*

\LA
\begin{Resp}\Rsign Tamquam ad latrónem exístis cum gládiis et fústibus comprehéndere me: \Star{} Cotídie apud vos eram in templo docens, et non me tenuístis: et ecce flagellátum dúcitis ad crucifigéndum.\end{Resp}
\begin{Resp}\Vsign Cumque injecíssent manus in Jesum, et tenuíssent eum, dixit ad eos.\end{Resp}
\begin{Resp}\Rsign Cotídie apud vos eram in templo docens, et non me tenuístis: et ecce flagellátum dúcitis ad crucifigéndum.\end{Resp}
\switchcolumn
\EN
\begin{Resp}\Rsign Are ye come out, as against a thief, with swords and staves, for to take Me? \Star{} I sat daily with you, teaching in the Temple, and ye laid no hold on Me; and, now when ye have scourged Me, ye lead Me away to crucify Me!\end{Resp}
\begin{Resp}\Vsign And when they had laid hands on Jesus, and taken Him, He said unto them:\end{Resp}
\begin{Resp}\Rsign I sat daily with you, teaching in the Temple, and ye laid no hold on Me; and now, when ye have scourged Me, ye lead Me away to crucify Me!\end{Resp}
\switchcolumn*

\LA
\LessonHead{Lectio V}
\switchcolumn
\EN
\LessonHead{Lesson V}
\switchcolumn*

\LA
\noindent Nostis qui convéntus erat malignántium Judæórum, et quæ multitúdo erat operántium iniquitátem. Quam iniquitátem? Quia voluérunt occídere Dóminum Jesum Christum. Tanta ópera bona, inquit, osténdi vobis: propter quod horum me vultis occídere? Pértulit omnes infírmos eórum, curávit omnes lánguidos eórum, prædicávit regnum cælórum, non tácuit vítia eórum, ut ipsa pótius eis displicérent, non médicus, a quo sanabántur. His ómnibus curatiónibus ejus ingráti, tamquam multa febre phrenétici, insaniéntes in médicum, qui vénerat curáre eos, excogitavérunt consílium perdéndi eum: tamquam ibi voléntes probáre, utrum vere homo sit, qui mori possit, an áliquid super hómines sit, et mori se non permíttat. Verbum ipsórum agnóscimus in Sapiéntia Salomónis: Morte turpíssima, ínquiunt, condemnémus eum. Interrogémus eum: erit enim respéctus in sermónibus illíus. Si enim vere Fílius Dei est, líberet eum.\par
\switchcolumn
\EN
\noindent We know what secret counsel was that of the wicked Jews, and what insurrection was that of the workers of iniquity. Of what iniquity were they the workers? The murder of our Lord Jesus Christ. Many good works, saith He, have I showed you for which of those works go ye about to kill Me? He had borne with all their weaknesses: He had healed all their diseases: He had preached unto them the kingdom of heaven: He had discovered to them their iniquities, that they might rather hate them, than the Physician That came to cure them. And now at last, without gratitude for all the tenderness of His healing love, like men raging in an high delirium, throwing themselves madly on the Physician, Who had come to cure them, they took counsel together how they might kill Him, as if to see if He were a Man and could die, or Something more than a man, and That would not let Himself die. In the Wisdom of Solomon we recognize their words, (ii. 18, 19, 20,) Let us condemn Him with a shameful death. Let us examine Him; for, by His own saying, He shall be respected. If He be the Son of God, let Him help Him.\par
\switchcolumn*

\LA
\begin{Resp}\Rsign Ténebræ factæ sunt, dum crucifixíssent Jesum Judǽi: et circa horam nonam exclamávit Jesus voce magna: Deus meus, ut quid me dereliquísti? \Star{} Et inclináto cápite, emísit spíritum.\end{Resp}
\begin{Resp}\Vsign Exclámans Jesus voce magna, ait: Pater, in manus tuas comméndo spíritum meum.\end{Resp}
\begin{Resp}\Rsign Et inclináto cápite, emísit spíritum.\end{Resp}
\switchcolumn
\EN
\begin{Resp}\Rsign The Jews crucified Jesus: and there was darkness (over all the land, unto the ninth hour): and about the ninth hour Jesus cried with a loud voice, (saying): My God, (My God,) why hast Thou forsaken Me? \Star{} And He bowed His Head, and gave up the Ghost.\end{Resp}
\begin{Resp}\Vsign When Jesus had cried with a loud voice, He said: Father, into thy hands I commend My Spirit.\end{Resp}
\begin{Resp}\Rsign And He bowed His Head, and gave up the Ghost.\end{Resp}
\switchcolumn*

\LA
\LessonHead{Lectio VI}
\switchcolumn
\EN
\LessonHead{Lesson VI}
\switchcolumn*

\LA
\noindent Exacuérunt tamquam gládium linguas suas. Non dicant Judǽi: Non occídimus Christum. Etenim proptérea eum dedérunt júdici Piláto, ut quasi ipsi a morte ejus videréntur immúnes. Nam cum dixísset eis Pilátus: Vos eum occídite: respondérunt, Nobis non licet occídere quemquam. Iniquitátem facínoris sui in júdicem hóminem refúndere volébant: sed numquid Deum júdicem fallébant? Quod fecit Pilátus, in eo ipso quod fecit, aliquántum párticeps fuit: sed in comparatióne illórum multo ipse innocéntior. Institit enim quantum pótuit, ut illum ex eórum mánibus liberáret: nam proptérea flagellátum prodúxit ad eos. Non persequéndo Dóminum flagellávit, sed eórum furóri satisfácere volens: ut vel sic jam mitéscerent, et desínerent velle occídere, cum flagellátum vidérent. Fecit et hoc. At ubi perseveravérunt, nostis illum lavísse manus, et dixísse, quod ipse non fecísset, mundum se esse a morte illíus. Fecit tamen. Sed si reus, quia fecit vel invítus: illi innocéntes, qui coëgérunt ut fáceret? Nullo modo. Sed ille dixit in eum senténtiam, et jussit eum crucifígi, et quasi ipse occídit: et vos, o Judǽi, occidístis. Unde occidístis? Gládio linguæ: acuístis enim linguas vestras. Et quando percussístis, nisi quando clamástis: Crucifíge, crucifíge?\par
\switchcolumn
\EN
\noindent They whetted their tongue like a sword. The Jews cannot say: We did not murder Christ, albeit they gave Him over to Pilate His judge, that they themselves might seem free of His death. For when Pilate said unto them, Take ye Him: and kill Him, they answered, It is not lawful for us to put any man to death. They could throw the blame of their sin upon a human judge: but did they deceive God, the Great Judge? In that which Pilate did, he was their accomplice, but in comparison with them, he had far the lesser sin. (John xix. 11.) Pilate strove as far as he could, to deliver Him out of their hands; for the which reason also he scourged Him, (John xix. 1,) and brought Him forth to them. He scourged not the Lord for cruelty's sake, but in the hope that; he might so slake their wild thirst for blood: that, perchance, even they might be touched with compassion, and cease to lust for His death, when they saw What He was after the flagellation. Even this effort he made! But when Pilate saw that he could not prevail, but that rather a tumult was made, (Matth. xxvii. 24,) ye know how that he took water, and washed his hands before the multitude, saying: I am innocent of the blood of this Just Person. And yet he delivered Him to be crucified! But if he were guilty who did it against his will, were they innocent; who goaded him on to it? No. Pilate gave sentence against Him. and commanded Him to be crucified. But ye, O ye Jews, ye also are His murderers! Wherewith? With your tongue, whetted like a sword. And when? But when ye cried, Crucify Him! Crucify Him! (Mark xv.13-14)\par
\switchcolumn*

\LA
\begin{Resp}\Rsign Animam meam diléctam trádidi in manus iniquórum, et facta est mihi heréditas mea sicut leo in silva: dedit contra me voces adversárius, dicens: Congregámini, et properáte ad devorándum illum: posuérunt me in desérto solitúdinis, et luxit super me omnis terra: \Star{} Quia non est invéntus qui me agnósceret, et fáceret bene.\end{Resp}
\begin{Resp}\Vsign Insurrexérunt in me viri absque misericórdia, et non pepercérunt ánimæ meæ.\end{Resp}
\begin{Resp}\Rsign Quia non est invéntus qui me agnósceret, et fáceret bene.\end{Resp}
\begin{Resp}\Vsign Glória Patri, et Fílio, \Star{} et Spirítui Sancto.\end{Resp}
\begin{Resp}\Rsign Animam meam diléctam trádidi in manus iniquórum, et facta est mihi heréditas mea sicut leo in silva: dedit contra me voces adversárius, dicens: Congregámini, et properáte ad devorándum illum: posuérunt me in desérto solitúdinis, et luxit super me omnis terra: * Quia non est invéntus qui me agnósceret, et fáceret bene.\end{Resp}
\switchcolumn
\EN
\begin{Resp}\Rsign I have given the dearly-beloved of My soul into the hand of her enemies and Mine heritage is become unto Me as a lion in the forest; the enemy crieth out against Me, saying: Assemble yourselves together, hasten to devour Him: they have made My portion a desolate wilderness, and the whole land mourneth unto Me: \Star{} Because there is none found that will know Me, nor do well.\end{Resp}
\begin{Resp}\Vsign There be risen up against me such as breathe out cruelty, and they have not spared my soul.\end{Resp}
\begin{Resp}\Rsign Because there is none found that will know Me, nor do well.\end{Resp}
\begin{Resp}\Vsign Glory be to the Father, and to the Son, \Star{} and to the Holy Ghost.\end{Resp}
\begin{Resp}\Rsign I have given the dearly-beloved of My soul into the hand of her enemies, and Mine heritage is become unto Me as a lion in the forest: the enemy crieth out against Me, saying: Assemble yourselves together, hasten to devour Him: they have made My portion a desolate wilderness, and the whole land mourneth unto me: * Because there is none found that will know Me, nor do well.\end{Resp}
\switchcolumn*

\LA
\LessonHead{Lectio VII}
\switchcolumn
\EN
\LessonHead{Lesson VII}
\switchcolumn*

\LA
\LessonTitle{De Epístola beáti Pauli Apóstoli ad Hebrǽos}
\switchcolumn
\EN
\LessonTitle{From the letter of the blessed Apostle Paul to the Hebrews}
\switchcolumn*

\LA
\Cite{Heb 4:11-15}
\switchcolumn
\EN
\Cite{Heb 4:11-15}
\switchcolumn*

\LA
\noindent Festinémus íngredi in illam réquiem: ut ne in idípsum quis íncidat incredulitátis exémplum. Vivus est enim sermo Dei, et éfficax et penetrabílior omni gládio ancípiti: et pertíngens usque ad divisiónem ánimæ ac spíritus, compágum quoque ac medullárum, et discrétor cogitatiónum et intentiónum cordis. Et non est ulla creatúra invisíbilis in conspéctu ejus: ómnia autem nuda et apérta sunt óculis ejus, ad quem nobis sermo. Habéntes ergo Pontíficem magnum, qui penetrávit cælos, Jesum Fílium Dei: teneámus confessiónem. Non enim habémus Pontíficem, qui non possit cómpati infirmitátibus nostris: tentátum autem per ómnia pro similitúdine absque peccáto.\par
\switchcolumn
\EN
\noindent Let us hasten therefore to enter into that rest; lest any man fall into the same example of unbelief. For the word of God is living and effectual, and more piercing than any two edged sword; and reaching unto the division of the soul and the spirit, of the joints also and the marrow, and is a discerner of the thoughts and intents of the heart. Neither is there any creature invisible in his sight: but all things are naked and open to his eyes, to whom our speech is. Having therefore a great high priest that hath passed into the heavens, Jesus the Son of God: let us hold fast our confession. For we have not a high priest, who can not have compassion on our infirmities: but one tempted in all things like as we are, without sin.\par
\switchcolumn*

\LA
\begin{Resp}\Rsign Tradidérunt me in manus impiórum, et inter iníquos projecérunt me, et non pepercérunt ánimæ meæ: congregáti sunt advérsum me fortes: \Star{} Et sicut gigántes stetérunt contra me.\end{Resp}
\begin{Resp}\Vsign Aliéni insurrexérunt advérsum me, et fortes quæsiérunt ánimam meam.\end{Resp}
\begin{Resp}\Rsign Et sicut gigántes stetérunt contra me.\end{Resp}
\switchcolumn
\EN
\begin{Resp}\Rsign They have turned me over into the hands of the wicked: they also have numbered me with the transgressors, neither have they spared my life: the mighty are gathered together against me; \Star{} And stand up against me like giants.\end{Resp}
\begin{Resp}\Vsign Strangers are risen up against me, and oppressors seek after my soul.\end{Resp}
\begin{Resp}\Rsign And stand up against me like giants.\end{Resp}
\switchcolumn*

\LA
\LessonHead{Lectio VIII}
\switchcolumn
\EN
\LessonHead{Lesson VIII}
\switchcolumn*

\LA
\Cite{Heb 4:16; 5:1-3}
\switchcolumn
\EN
\Cite{Heb 4:16; 5:1-3}
\switchcolumn*

\LA
\noindent Adeámus ergo cum fidúcia ad thronum grátiæ: ut misericórdiam consequámur, et grátiam inveniámus in auxílio opportúno. Omnis namque Póntifex ex homínibus assúmptus, pro homínibus constitúitur in iis, quæ sunt ad Deum, ut ófferat dona, et sacrifícia pro peccátis: qui condolére possit iis, qui ignórant et errant: quóniam et ipse circúmdatus est infirmitáte: et proptérea debet quemádmodum pro pópulo, ita étiam pro semetípso offérre pro peccátis.\par
\switchcolumn
\EN
\noindent Let us go therefore with confidence to the throne of grace: that we may obtain mercy, and find grace in seasonable aid. For every high priest taken from among men, is ordained for men in the things that appertain to God, that he may offer up gifts and sacrifices for sins: Who can have compassion on them that are ignorant and that err: because he himself also is compassed with infirmity. And therefore he ought, as for the people, so also for himself, to offer for sins.\par
\switchcolumn*

\LA
\begin{Resp}\Rsign Jesum trádidit ímpius summis princípibus sacerdótum, et senióribus pópuli: \Star{} Petrus autem sequebátur eum a longe, ut vidéret finem.\end{Resp}
\begin{Resp}\Vsign Adduxérunt autem eum ad Cáipham príncipem sacerdótum, ubi scribæ et pharisǽi convénerant.\end{Resp}
\begin{Resp}\Rsign Petrus autem sequebátur eum a longe, ut vidéret finem.\end{Resp}
\switchcolumn
\EN
\begin{Resp}\Rsign That wicked one betrayed Jesus to the chief-priests and elders of the people \Star{} But Peter followed Him afar off, to see the end.\end{Resp}
\begin{Resp}\Vsign And they led Him away to Caiaphas the High Priest, where the Scribes and Pharisees were assembled.\end{Resp}
\begin{Resp}\Rsign But Peter followed Him afar off, to see the end.\end{Resp}
\switchcolumn*

\LA
\LessonHead{Lectio IX}
\switchcolumn
\EN
\LessonHead{Lesson IX}
\switchcolumn*

\LA
\Cite{Heb 5:4-10}
\switchcolumn
\EN
\Cite{Heb 5:4-10}
\switchcolumn*

\LA
\noindent Nec quisquam sumit sibi honórem, sed qui vocátur a Deo, tamquam Aaron. Sic et Christus non semetípsum clarificávit ut Póntifex fíeret, sed qui locútus est ad eum: Fílius meus es tu, ego hódie génui te. Quemádmodum et in álio loco dicit: Tu es sacérdos in ætérnum, secúndum órdinem Melchísedech. Qui in diébus carnis suæ preces, supplicationésque ad eum, qui possit illum salvum fácere a morte, cum clamóre válido et lácrimis ófferens, exaudítus est pro sua reveréntia. Et quidem cum esset Fílius Dei, dídicit ex iis, quæ passus est, obediéntiam: et consummátus, factus est ómnibus obtemperántibus sibi causa salútis ætérnæ, appellátus a Deo Póntifex juxta órdinem Melchísedech.\par
\switchcolumn
\EN
\noindent Neither doth any man take the honour to himself, but he that is called by God, as Aaron was. So Christ also did not glorify himself, that he might be made a high priest: but he that said unto him: Thou art my Son, this day have I begotten thee. As he saith also in another place: Thou art a priest for ever, according to the order of Melchisedech. Who in the days of his flesh, with a strong cry and tears, offering up prayers and supplications to him that was able to save him from death, was heard for his reverence. And whereas indeed he was the Son of God, he learned obedience by the things which he suffered: And being consummated, he became, to all that obey him, the cause of eternal salvation. Called by God a high priest according to the order of Melchisedech.\par
\switchcolumn*

\LA
\begin{Resp}\Rsign Caligavérunt óculi mei a fletu meo: quia elongátus est a me, qui consolabátur me: Vidéte omnes pópuli, \Star{} Si est dolor símilis sicut dolor meus.\end{Resp}
\begin{Resp}\Vsign O vos omnes, qui transítis per viam, atténdite et vidéte.\end{Resp}
\begin{Resp}\Rsign Si est dolor símilis sicut dolor meus.\end{Resp}
\begin{Resp}\Vsign Glória Patri, et Fílio, \Star{} et Spirítui Sancto.\end{Resp}
\begin{Resp}\Rsign Caligavérunt óculi mei a fletu meo: quia elongátus est a me, qui consolabátur me: Vidéte omnes pópuli, * Si est dolor símilis sicut dolor meus.\end{Resp}
\switchcolumn
\EN
\begin{Resp}\Rsign Mine eyes do fail with tears, because the Comforter that should relieve me is far from me. Behold, O all ye nations \Star{} If there be any sorrow like unto my sorrow.\end{Resp}
\begin{Resp}\Vsign O all ye that pass by, behold, and see\end{Resp}
\begin{Resp}\Rsign If there be any sorrow like unto my sorrow.\end{Resp}
\begin{Resp}\Vsign Glory be to the Father, and to the Son, \Star{} and to the Holy Ghost.\end{Resp}
\begin{Resp}\Rsign Mine eyes do fail with tears, because the Comforter that should relieve me is far from me. Behold, O all ye nations, * If there be any sorrow like unto my sorrow.\end{Resp}
\switchcolumn*

\end{paracol}

\Office{Sabbato Sancto}{Holy Saturday}{Feria privilegiata Duplex I classis}
\addcontentsline{toc}{subsection}{Sabbato Sancto\enspace\textit{Holy Saturday}}
\begin{paracol}{2}
\LA
\LessonHead{Lectio I}
\switchcolumn
\EN
\LessonHead{Lesson I}
\switchcolumn*

\LA
\LessonTitle{De Lamentatióne Jeremíæ Prophétæ}
\switchcolumn
\EN
\LessonTitle{Lesson from the book of Lamentations}
\switchcolumn*

\LA
\Cite{Lam 3:22-30}
\switchcolumn
\EN
\Cite{Lam 3:22-30}
\switchcolumn*

\LA
\noindent Heth. Misericórdiæ Dómini quia non sumus consúmpti: quia non defecérunt miseratiónes ejus.\par
\switchcolumn
\EN
\noindent Heth. The mercies of the Lord that we are not consumed: because his commiserations have not failed.\par
\switchcolumn*

\LA
Heth. Novi dilúculo, multa est fides tua.\par
\switchcolumn
\EN
Heth. They are new every morning, great is thy faithfulness.\par
\switchcolumn*

\LA
Heth. Pars mea Dóminus, dixit ánima mea: proptérea exspectábo eum.\par
\switchcolumn
\EN
Heth. The Lord is my portion, said my soul: therefore will I wait for him.\par
\switchcolumn*

\LA
Teth. Bonus est Dóminus sperántibus in eum, ánimæ quærénti illum.\par
\switchcolumn
\EN
Teth. The Lord is good to them that hope in him, to the soul that seeketh him.\par
\switchcolumn*

\LA
Teth. Bonum est præstolári cum siléntio salutáre Dei.\par
\switchcolumn
\EN
Teth. It is good to wait with silence for the salvation of God.\par
\switchcolumn*

\LA
Teth. Bonum est viro cum portáverit jugum ab adulescéntia sua.\par
\switchcolumn
\EN
Teth. It is good for a man, when he hath borne the yoke from his youth.\par
\switchcolumn*

\LA
Jod. Sedébit solitárius, et tacébit: quia levávit super se.\par
\switchcolumn
\EN
Jod. He shall sit solitary, and hold his peace: because he hath taken it up upon himself.\par
\switchcolumn*

\LA
Jod. Ponet in púlvere os suum, si forte sit spes.\par
\switchcolumn
\EN
Jod. He shall put his mouth in the dust, if so be there may be hope.\par
\switchcolumn*

\LA
Jod. Dabit percutiénti se maxíllam, saturábitur oppróbriis.\par
\switchcolumn
\EN
Jod. He shall give his cheek to him that striketh him, he shall be filled with reproaches.\par
\switchcolumn*

\LA
Jerúsalem, Jerúsalem, convértere ad Dóminum Deum tuum.\par
\switchcolumn
\EN
Jerusalem! Jerusalem! Return unto the Lord thy God.\par
\switchcolumn*

\LA
\begin{Resp}\Rsign Sicut ovis ad occisiónem ductus est, et dum male tractarétur, non apéruit os suum: tráditus est ad mortem, \Star{} Ut vivificáret pópulum suum.\end{Resp}
\begin{Resp}\Vsign Trádidit in mortem ánimam suam, et inter scelerátos reputátus est.\end{Resp}
\begin{Resp}\Rsign Ut vivificáret pópulum suum.\end{Resp}
\switchcolumn
\EN
\begin{Resp}\Rsign He hath been brought as a lamb to the slaughter, and while he was evil entreated he opened not his mouth: he was delivered up to death \Star{} That he might quicken his people.\end{Resp}
\begin{Resp}\Vsign He hath poured out his soul unto death, and he was numbered with the transgressors.\end{Resp}
\begin{Resp}\Rsign That he might quicken his people.\end{Resp}
\switchcolumn*

\LA
\LessonHead{Lectio II}
\switchcolumn
\EN
\LessonHead{Lesson II}
\switchcolumn*

\LA
\Cite{Lam 4:1-6}
\switchcolumn
\EN
\Cite{Lam 4:1-8}
\switchcolumn*

\LA
\noindent Aleph. Quómodo obscurátum est aurum, mutátus est color óptimus, dispérsi sunt lápides sanctuárii in cápite ómnium plateárum?\par
\switchcolumn
\EN
\noindent Aleph. How is the gold become dim, the finest colour is changed, the stones of the sanctuary are scattered in the top of every street?\par
\switchcolumn*

\LA
Beth. Fílii Sion íncliti, et amícti auro primo: quómodo reputáti sunt in vasa téstea, opus mánuum fíguli?\par
\switchcolumn
\EN
Beth. The noble sons of Sion, and they that were clothed with the best gold: how are they esteemed as earthen vessels, the work of the potter's hands?\par
\switchcolumn*

\LA
Ghimel. Sed et lámiæ nudavérunt mammam, lactavérunt cátulos suos: fília pópuli mei crudélis, quasi strúthio in desérto.\par
\switchcolumn
\EN
Ghimel. Even the sea monsters have drawn out the breast, they have given suck to their young: the daughter of my people is cruel, like the ostrich in the desert.\par
\switchcolumn*

\LA
Daleth. Adhǽsit lingua lacténtis ad palátum ejus in siti: párvuli petiérunt panem, et non erat qui frángeret eis.\par
\switchcolumn
\EN
Daleth. The tongue of the sucking child hath stuck to the roof of his mouth for thirst: the little ones have asked for bread, and there was none to break it unto them.\par
\switchcolumn*

\LA
He. Qui vescebántur voluptuóse, interiérunt in viis: qui nutriebántur in cróceis, amplexáti sunt stércora.\par
\switchcolumn
\EN
He. They that were fed delicately have died in the streets; they that were brought up in scarlet have embraced the dung.\par
\switchcolumn*

\LA
Vau. Et major effécta est iníquitas fíliæ pópuli mei peccáto Sodomórum, quæ subvérsa est in moménto, et non cepérunt in ea manus.\par
\switchcolumn
\EN
Vau. And the iniquity of the daughter of my people is made greater than the sin of Sodom, which was overthrown in a moment, and hands took nothing in her.\par
\switchcolumn*

\LA
Jerúsalem, Jerúsalem, convértere ad Dóminum Deum tuum.\par
\switchcolumn
\EN
Jerusalem! Jerusalem! Return unto the Lord thy God.\par
\switchcolumn*

\LA
\begin{Resp}\Rsign Jerúsalem, surge, et éxue te véstibus jucunditátis: indúere cínere et cilício, \Star{} Quia in te occísus est Salvátor Israël.\end{Resp}
\begin{Resp}\Vsign Deduc quasi torréntem lácrimas per diem et noctem, et non táceat pupílla óculi tui.\end{Resp}
\begin{Resp}\Rsign Quia in te occísus est Salvátor Israël.\end{Resp}
\switchcolumn
\EN
\begin{Resp}\Rsign Arise, O Jerusalem, and put off thy garments of rejoicing: cover thee with sackcloth and ashes \Star{} For the Saviour of Israel hath been slain in the midst of thee.\end{Resp}
\begin{Resp}\Vsign Let thy tears run down like a river day and night, and let not the apple of thine eye cease.\end{Resp}
\begin{Resp}\Rsign For the Saviour of Israel hath been slain in the midst of thee.\end{Resp}
\switchcolumn*

\LA
\LessonHead{Lectio III}
\switchcolumn
\EN
\LessonHead{Lesson III}
\switchcolumn*

\LA
\LessonTitle{Incipit Orátio Jeremíæ Prophétæ}
\switchcolumn
\EN

\switchcolumn*

\LA
\Cite{Lam 5:1-11}
\switchcolumn
\EN
\Cite{Lam 5:1-11}
\switchcolumn*

\LA
\noindent Recordáre, Dómine, quid accíderit nobis: intuére, et réspice oppróbrium nostrum. Heréditas nostra versa est ad aliénos: domus nostræ ad extráneos. Pupílli facti sumus absque patre, matres nostræ quasi víduæ. Aquam nostram pecúnia bíbimus: ligna nostra prétio comparávimus. Cervícibus nostris minabámur, lassis non dabátur réquies. Ægýpto dédimus manum, et Assýriis, ut saturarémur pane. Patres nostri peccavérunt, et non sunt: et nos iniquitátes eórum portávimus. Servi domináti sunt nostri: non fuit qui redímeret de manu eórum. In animábus nostris afferebámus panem nobis, a fácie gládii in desérto. Pellis nostra quasi clíbanus exústa est a fácie tempestátum famis. Mulíeres in Sion humiliavérunt, et vírgines in civitátibus Juda.\par
\switchcolumn
\EN
\noindent Remember, O Lord, what is come upon us: consider and behold our reproach. Our inheritance is turned to aliens: our houses to strangers. We are become orphans without a father: our mothers are as widows. We have drunk our water for money: we have bought our wood. We were dragged by the necks, we were weary and no rest was given us. We have given our hand to Egypt, and to the Assyrians, that we might be satisfied with bread. Our fathers have sinned, and are not: and we have borne their iniquities. Servants have ruled over us: there was none to redeem us out of their hand. We fetched our bread at the peril of our lives, because of the sword in the desert. Our skin was burnt as an oven, by reason of the violence of the famine. They oppressed the women in Sion, and the virgins in the cities of Juda.\par
\switchcolumn*

\LA
Jerúsalem, Jerúsalem, convértere ad Dóminum Deum tuum.\par
\switchcolumn
\EN
Jerusalem! Jerusalem! Return unto the Lord thy God.\par
\switchcolumn*

\LA
\begin{Resp}\Rsign Plange quasi virgo, plebs mea: ululáte, pastóres, in cínere et cilício: \Star{} Quia venit dies Dómini magna, et amára valde.\end{Resp}
\begin{Resp}\Vsign Accíngite vos, sacerdótes, et plángite, minístri altáris, aspérgite vos cínere.\end{Resp}
\begin{Resp}\Rsign Quia venit dies Dómini magna, et amára valde.\end{Resp}
\begin{Resp}\Vsign Glória Patri, et Fílio, \Star{} et Spirítui Sancto.\end{Resp}
\begin{Resp}\Rsign Plange quasi virgo, plebs mea: ululáte, pastóres, in cínere et cilício: * Quia venit dies Dómini magna, et amára valde.\end{Resp}
\switchcolumn
\EN
\begin{Resp}\Rsign O my people! lament, like a virgin girded with sack-cloth for the husband of her youth, howl, ye shepherds, in sack-cloth and ashes \Star{} For the day of the Lord is at hand, and it is great and very terrible.\end{Resp}
\begin{Resp}\Vsign Gird yourselves, ye Priests, and howl, ye ministers of the altar: cast up ashes upon you.\end{Resp}
\begin{Resp}\Rsign For the day of the Lord is at hand, and it is great and very terrible.\end{Resp}
\begin{Resp}\Vsign Glory be to the Father, and to the Son, \Star{} and to the Holy Ghost.\end{Resp}
\begin{Resp}\Rsign O my people! lament, like a virgin, girded with sack-cloth for the husband of her youth, howl, ye shepherds, in sack-cloth and ashes * For the day of the Lord is at hand, and it is great and very terrible.\end{Resp}
\switchcolumn*

\LA
\LessonHead{Lectio IV}
\switchcolumn
\EN
\LessonHead{Lesson IV}
\switchcolumn*

\LA
\LessonTitle{Ex Tractátu sancti Augustíni Epíscopi super Psalmos}
\switchcolumn
\EN
\LessonTitle{From the Treatise of St. Augustine, Bishop (of Hippo,) Upon the Psalms}
\switchcolumn*

\LA
\Source{In Psalmum 63 versum 7}
\switchcolumn
\EN
\Source{On Psalm lxiii, 7}
\switchcolumn*

\LA
\noindent Accédet homo ad cor altum, et exaltábitur Deus. Illi dixérunt: Quis nos vidébit? Defecérunt scrutántes scrutatiónes, consília mala. Accéssit homo ad ipsa consília, passus est se tenéri ut homo. Non enim tenerétur nisi homo, aut viderétur nisi homo, aut cæderétur nisi homo, aut crucifigerétur, aut morerétur nisi homo. Accéssit ergo homo ad illas omnes passiónes, quæ in illo nihil valérent, nisi esset homo. Sed si ille non esset homo, non liberarétur homo. Accéssit homo ad cor altum, id est, cor secrétum, obíciens aspéctibus humánis hóminem, servans intus Deum: celans formam Dei, in qua æquális est Patri, et ófferens formam servi, qua minor est Patre.\par
\switchcolumn
\EN
\noindent We shall attain to thoughts that are very deep: but God shall still be exalted. The enemies of our Lord had communed of laying snares privily; they had said, Who shall see them? They had searched out iniquities; they had accomplished a diligent search. And Man attained even unto (the realization of) their counsels, for the Lord, as Man, suffered Himself to be taken. For He had not been taken at all, unless He had been a Man, or seen, unless He had been a Man, or smitten, unless He had been a Man, or crucified, unless He had been a Man, or have died, unless He had been a Man. Man therefore, He attained unto all those sufferings, which had had nothing in Him, unless He had been a Man. But if He had not been Man, man had not been redeemed. And the Lord as Man attained to thoughts that were very deep, yea, secret; showing the Manhood to the eyes of men, and keeping the Godhead within Him; veiling the form of God, as touching Which, He is Equal to the Father, and manifesting the form of a servant, as touching which, He is inferior to the Father.\par
\switchcolumn*

\LA
\begin{Resp}\Rsign Recéssit pastor noster, fons aquæ vivæ, ad cujus tránsitum sol obscurátus est: \Star{} Nam et ille captus est, qui captívum tenébat primum hóminem: hódie portas mortis et seras páriter Salvátor noster disrúpit.\end{Resp}
\begin{Resp}\Vsign Destrúxit quidem claustra inférni, et subvértit poténtias diáboli.\end{Resp}
\begin{Resp}\Rsign Nam et ille captus est, qui captívum tenébat primum hóminem: hódie portas mortis et seras páriter Salvátor noster disrúpit.\end{Resp}
\switchcolumn
\EN
\begin{Resp}\Rsign Our Shepherd, even the Fountain of living waters, is gone from us; He passed away, and the sun was darkened. \Star{} For now hath our Saviour bound him captive, which bound the first man captive; this day hath He burst the gates and bars of death.\end{Resp}
\begin{Resp}\Vsign The bands of hell He hath utterly abolished, and hath done away the power of the devil.\end{Resp}
\begin{Resp}\Rsign For now hath our Saviour bound him captive, which bound the first man captive; this day hath He burst the gates and bars of death.\end{Resp}
\switchcolumn*

\LA
\LessonHead{Lectio V}
\switchcolumn
\EN
\LessonHead{Lesson V}
\switchcolumn*

\LA
\noindent Quo perduxérunt illas scrutatiónes suas, quas perscrutántes defecérunt, ut étiam mórtuo Dómino et sepúlto, custódes pónerent ad sepúlcrum? Dixérunt enim Piláto: Sedúctor ille: hoc appellabátur nómine Dóminus Jesus Christus, ad solátium servórum suórum, quando dicúntur seductóres: ergo illi Piláto: Sedúctor ille, ínquiunt, dixit adhuc vivens: Post tres dies resúrgam. Jube ítaque custodíri sepúlcrum usque in diem tértium, ne forte véniant discípuli ejus, et furéntur eum, et dicant plebi: Surréxit a mórtuis: et erit novíssimus error pejor prióre. Ait illis Pilátus: Habétis custódiam, ite, custodíte sicut scitis. Illi autem abeúntes, muniérunt sepúlcrum, signántes lápidem cum custódibus.\par
\switchcolumn
\EN
\noindent How far did the accomplishment of their diligent search reach? Even to the setting a watch of soldiers at the sepulchre, to guard the Lord, even after He was dead and buried. For they said unto Pilate: Sir, we remember that that deceiver (Matth. xxvii. 63.) This was the term by which they designated the Lord Jesus Christ, and the remembrance that He was so named is a sweet consolation to us His servants, when we are called impostors. So they said unto Pilate, that deceiver said, while He was yet alive: After three days I will rise again. Command therefore that the sepulchre be made sure until the third day, lest His disciples come and steal Him away, and say unto the people: He is risen again from the dead: so the last error shall be worse than the first. Pilate said unto them: Ye have a watch; go your way; make it as sure as ye can. So they went and made the sepulchre sure, sealing the stone, and setting a watch.\par
\switchcolumn*

\LA
\begin{Resp}\Rsign O vos omnes, qui transítis per viam, atténdite, et vidéte, \Star{} Si est dolor símilis sicut dolor meus.\end{Resp}
\begin{Resp}\Vsign Atténdite, univérsi pópuli, et vidéte dolórem meum.\end{Resp}
\begin{Resp}\Rsign Si est dolor símilis sicut dolor meus.\end{Resp}
\switchcolumn
\EN
\begin{Resp}\Rsign O all ye that pass by, behold and see; \Star{} If there be any sorrow like unto my sorrow.\end{Resp}
\begin{Resp}\Vsign O all ye nations, behold, and see my sorrow,\end{Resp}
\begin{Resp}\Rsign If there be any sorrow like unto my sorrow.\end{Resp}
\switchcolumn*

\LA
\LessonHead{Lectio VI}
\switchcolumn
\EN
\LessonHead{Lesson VI}
\switchcolumn*

\LA
\noindent Posuérunt custódes mílites ad sepúlcrum. Concússa terra Dóminus resurréxit: mirácula facta sunt tália circa sepúlcrum, ut et ipsi mílites, qui custódes advénerant, testes fíerent, si vellent vera nuntiáre. Sed avarítia illa, quæ captivávit discípulum cómitem Christi, captivávit et mílitem custódem sepúlcri. Damus, ínquiunt, vobis pecúniam: et dícite, quia vobis dormiéntibus venérunt discípuli ejus, et abstulérunt eum. Vere defecérunt scrutántes scrutatiónes. Quid est quod dixísti, o infélix astútia? Tantúmne déseris lucem consílii pietátis, et in profúnda versútiæ demérgeris, ut hoc dicas: Dícite quia vobis dormiéntibus venérunt discípuli ejus, et abstulérunt eum? Dormiéntes testes ádhibes: vere tu ipse obdormísti, qui scrutándo tália defecísti.\par
\switchcolumn
\EN
\noindent So they went, and made the sepulchre sure, sealing the stone, and setting a watch and anon, behold, there was a great earthquake, and the Lord arose. So great wonders were wrought about the sepulchre that the very soldiers, which were put to guard it, were witnesses thereto, if only they would have told the truth. But the same love of money which had made a slave of that disciple who was a companion of Christ, made slaves also of the soldiers that were put to watch His sepulchre. Some of the watch came into the city, and showed unto the chief-priests all the things that were done: and when they were assembled with the elders, and had taken counsel, they gave large money unto the soldiers, saying: Say ye, His disciples came by night and stole Him away while we slept. In good sooth, their diligent search had been accomplished and ended before this. What didst thou say, O stupid cunning? Wast thou indeed so utterly void of the light of godly wisdom, and confounded in the bottomless pit of thine own falsehood as to tell them to say: His disciples came by night, and stole Him away while we slept? Part of the testimony of thine eye-witnesses was that they were asleep at the time: thou thyself wast asleep not to be able to see that on their own testimony, their testimony must have been worthless.\par
\switchcolumn*

\LA
\begin{Resp}\Rsign Ecce quómodo móritur justus, et nemo pércipit corde: et viri justi tollúntur, et nemo consíderat: a fácie iniquitátis sublátus est justus: \Star{} Et erit in pace memória ejus.\end{Resp}
\begin{Resp}\Vsign Tamquam agnus coram tondénte se obmútuit, et non apéruit os suum: de angústia et de judício sublátus est.\end{Resp}
\begin{Resp}\Rsign Et erit in pace memória ejus.\end{Resp}
\begin{Resp}\Vsign Glória Patri, et Fílio, \Star{} et Spirítui Sancto.\end{Resp}
\begin{Resp}\Rsign Ecce quómodo móritur justus, et nemo pércipit corde: et viri justi tollúntur, et nemo consíderat: a fácie iniquitátis sublátus est justus: * Et erit in pace memória ejus.\end{Resp}
\switchcolumn
\EN
\begin{Resp}\Rsign Behold how the righteous dieth, and no man taketh it to heart; and the just are taken away, and none considereth. From the midst of sinners is the righteous translated; \Star{} And his memory is in peace.\end{Resp}
\begin{Resp}\Vsign As a lamb before his shearers is dumb, so He opened not His mouth; He was taken from prison and from judgment.\end{Resp}
\begin{Resp}\Rsign And his memory is in peace.\end{Resp}
\begin{Resp}\Vsign Glory be to the Father, and to the Son, \Star{} and to the Holy Ghost.\end{Resp}
\begin{Resp}\Rsign Behold how the righteous dieth, and no man taketh it to heart; and the just are taken away, and none considereth. From the midst of sinners is the righteous translated; * And his memory is in peace.\end{Resp}
\switchcolumn*

\LA
\LessonHead{Lectio VII}
\switchcolumn
\EN
\LessonHead{Lesson VII}
\switchcolumn*

\LA
\LessonTitle{De Epístola beáti Pauli Apóstoli ad Hebrǽos}
\switchcolumn
\EN
\LessonTitle{From the letter of blessed Apostle Paul to the Hebrews}
\switchcolumn*

\LA
\Cite{Heb 9:11-14}
\switchcolumn
\EN
\Cite{Heb 9:11-14}
\switchcolumn*

\LA
\noindent Christus assístens Póntifex futurórum bonórum, per ámplius et perféctius tabernáculum non manufáctum, id est, non hujus creatiónis: neque per sánguinem hircórum, aut vitulórum, sed per próprium sánguinem introívit semel in Sancta, ætérna redemptióne invénta. Si enim sanguis hircórum, et taurórum, et cinis vítulæ aspérsus inquinátos sanctíficat ad emundatiónem carnis: quanto magis sanguis Christi, qui per Spíritum Sanctum semetípsum óbtulit immaculátum Deo, emundábit consciéntiam nostram ab opéribus mórtuis, ad serviéndum Deo vivénti?\par
\switchcolumn
\EN
\noindent But Christ, being come an high priest of the good things to come, by a greater and more perfect tabernacle not made with hand, that is, not of this creation: Neither by the blood of goats, or of calves, but by his own blood, entered once into the holies, having obtained eternal redemption. For if the blood of goats and of oxen, and the ashes of an heifer being sprinkled, sanctify such as are defiled, to the cleansing of the flesh: How much more shall the blood of Christ, who by the Holy Ghost offered himself unspotted unto God, cleanse our conscience from dead works, to serve the living God?\par
\switchcolumn*

\LA
\begin{Resp}\Rsign Astitérunt reges terræ, et príncipes convenérunt in unum, \Star{} Advérsus Dóminum, et advérsus Christum ejus.\end{Resp}
\begin{Resp}\Vsign Quare fremuérunt gentes, et pópuli meditáti sunt inánia?\end{Resp}
\begin{Resp}\Rsign Advérsus Dóminum, et advérsus Christum ejus.\end{Resp}
\switchcolumn
\EN
\begin{Resp}\Rsign The kings of the earth set themselves, and the rulers take counsel together \Star{} Against the Lord, and against His Anointed.\end{Resp}
\begin{Resp}\Vsign Why do the heathen rage? and the people imagine a vain thing,\end{Resp}
\begin{Resp}\Rsign Against the Lord, and against His Anointed?\end{Resp}
\switchcolumn*

\LA
\LessonHead{Lectio VIII}
\switchcolumn
\EN
\LessonHead{Lesson VIII}
\switchcolumn*

\LA
\Cite{Heb 9:15-18}
\switchcolumn
\EN
\Cite{Heb 9:15-18}
\switchcolumn*

\LA
\noindent Et ídeo novi testaménti mediátor est: ut, morte intercedénte, in redemptiónem eárum prævaricatiónum, quæ erant sub prióri testaménto, repromissiónem accípiant, qui vocáti sunt ætérnæ hereditátis. Ubi enim testaméntum est: mors necésse est intercédat testatóris. Testaméntum enim in mórtuis confirmátum est: alióquin nondum valet, dum vivit qui testátus est. Unde nec primum quidem sine sánguine dedicátum est.\par
\switchcolumn
\EN
\noindent And therefore he is the mediator of the new testament: that by means of his death, for the redemption of those transgressions, which were under the former testament, they that are called may receive the promise of eternal inheritance. For where there is a testament, the death of the testator must of necessity come in. For a testament is of force, after men are dead: otherwise it is as yet of no strength, whilst the testator liveth. Whereupon neither was the first indeed dedicated without blood.\par
\switchcolumn*

\LA
\begin{Resp}\Rsign Æstimátus sum cum descendéntibus in lacum: \Star{} Factus sum sicut homo sine adjutório, inter mórtuos liber.\end{Resp}
\begin{Resp}\Vsign Posuérunt me in lacu inferióri, in tenebrósis, et in umbra mortis.\end{Resp}
\begin{Resp}\Rsign Factus sum sicut homo sine adjutório, inter mórtuos liber.\end{Resp}
\switchcolumn
\EN
\begin{Resp}\Rsign I am counted with them that go down into the pit. \Star{} I am as a man that hath no strength, lying nerveless among the dead.\end{Resp}
\begin{Resp}\Vsign They have laid me in the lowest pit, in darkness, and in the shadow of death.\end{Resp}
\begin{Resp}\Rsign I am as a man that hath no strength, lying nerveless among the dead.\end{Resp}
\switchcolumn*

\LA
\LessonHead{Lectio IX}
\switchcolumn
\EN
\LessonHead{Lesson IX}
\switchcolumn*

\LA
\Cite{Heb 9:19-22}
\switchcolumn
\EN
\Cite{Heb 9:19-22}
\switchcolumn*

\LA
\noindent Lecto enim omni mandáto legis a Móyse univérso pópulo: accípiens sánguinem vitulórum, et hircórum cum aqua et lana coccínea, et hyssópo: ipsum quoque librum, et omnem pópulum aspérsit, dicens: Hic sanguis testaménti, quod mandávit ad vos Deus. Etiam tabernáculum, et ómnia vasa ministérii sánguine simíliter aspérsit: et ómnia pene in sánguine secúndum legem mundántur: et sine sánguinis effusióne non fit remíssio.\par
\switchcolumn
\EN
\noindent For when every commandment of the law had been read by Moses to all the people, he took the blood of calves and goats, with water, and scarlet wool and hyssop, and sprinkled both the book itself and all the people, Saying: This is the blood of the testament, which God hath enjoined unto you. The tabernacle also and all the vessels of the ministry, in like manner, he sprinkled with blood. And almost all things, according to the law, are cleansed with blood: and without shedding of blood there is no remission.\par
\switchcolumn*

\LA
\begin{Resp}\Rsign Sepúlto Dómino, signátum est monuméntum, volvéntes lápidem ad óstium monuménti: \Star{} Ponéntes mílites, qui custodírent illum.\end{Resp}
\begin{Resp}\Vsign Accedéntes príncipes sacerdótum ad Pilátum, petiérunt illum.\end{Resp}
\begin{Resp}\Rsign Ponéntes mílites, qui custodírent illum.\end{Resp}
\begin{Resp}\Vsign Glória Patri, et Fílio, \Star{} et Spirítui Sancto.\end{Resp}
\begin{Resp}\Rsign Sepúlto Dómino, signátum est monuméntum, volvéntes lápidem ad óstium monuménti: * Ponéntes mílites, qui custodírent illum.\end{Resp}
\switchcolumn
\EN
\begin{Resp}\Rsign After that the Lord was buried, they sealed the sepulchre, rolling a stone to the door of the sepulchre \Star{} Setting a watch to keep Him.\end{Resp}
\begin{Resp}\Vsign The chief priests came together unto Pilate, and made that request unto him.\end{Resp}
\begin{Resp}\Rsign Setting a watch to keep Him.\end{Resp}
\begin{Resp}\Vsign Glory be to the Father, and to the Son, \Star{} and to the Holy Ghost.\end{Resp}
\begin{Resp}\Rsign After that the Lord was buried, they sealed the sepulchre, rolling a stone to the door of the sepulchre, * Setting a watch to keep Him.\end{Resp}
\switchcolumn*

\end{paracol}

\SeasonTitle{Tempus Paschale}{Paschaltide}

\addcontentsline{toc}{section}{Tempus Paschale\enspace\textit{Paschaltide}}

\Office{Dominica Resurrectionis}{Easter Sunday}{Duplex I classis cum Octava privilegiata I}
\addcontentsline{toc}{subsection}{Dominica Resurrectionis\enspace\textit{Easter Sunday}}
\begin{paracol}{2}
\LA
\LessonHead{Lectio I}
\switchcolumn
\EN
\LessonHead{Lesson I}
\switchcolumn*

\LA
\LessonTitle{Léctio sancti Evangélii secúndum Marcum}
\switchcolumn
\EN
\LessonTitle{Continuation of the Holy Gospel according to Mark}
\switchcolumn*

\LA
\Cite{Marc 16:1-7}
\switchcolumn
\EN
\Cite{Mark 16:1-7}
\switchcolumn*

\LA
\noindent In illo témpore: María Magdaléne, et María Jacóbi, et Salóme emérunt arómata, ut veniéntes úngerent Jesum. Et réliqua.\par
\switchcolumn
\EN
\noindent And when the sabbath was past, Mary Magdalen, and Mary the mother of James, and Salome, bought sweet spices, that coming, they might anoint Jesus. And what follows.\par
\switchcolumn*

\LA
\LessonTitle{Homilía sancti Gregórii Papæ}
\switchcolumn
\EN
\LessonTitle{Homily of St. Gregory, Pope}
\switchcolumn*

\LA
\Source{Homilia 21 in Evangelia}
\switchcolumn
\EN
\Source{21st Homily on the Gospels}
\switchcolumn*

\LA
\noindent Audístis, fratres caríssimi, quod sanctæ mulíeres, quæ Dóminum fúerant secútæ, cum aromátibus ad monuméntum venérunt, et ei, quem vivéntem diléxerant, étiam mórtuo, stúdio humanitátis obsequúntur. Sed res gesta áliquid in sancta Ecclésia signat geréndum. Sic quippe necésse est ut audiámus quæ facta sunt, quátenus cogitémus étiam quæ nobis sint ex eórum imitatióne faciénda. Et nos ergo in eum, qui est mórtuus, credéntes, si odóre virtútum reférti, cum opinióne bonórum óperum Dóminum quǽrimus, ad monuméntum profécto illíus cum aromátibus venímus. Illæ autem mulíeres Angelos vident, quæ cum aromátibus venérunt: quia vidélicet illæ mentes supérnos cives aspíciunt, quæ cum virtútum odóribus ad Dóminum per sancta desidéria proficiscúntur.\par
\switchcolumn
\EN
\noindent Dearly beloved brethren, ye have heard the deed of the holy women which had followed the Lord; how that they brought sweet spices to His sepulchre, and, now that He was dead, having loved Him while He was yet alive, they followed Him with careful tenderness still. But the deed of these holy women doth point to somewhat which must needs be done in the holy Church. And it behoveth us well to give ear to what they did, that we may afterward consider with ourselves what we must do likewise after their example. We also, who believe in Him That was dead, do come to His sepulchre bearing sweet spices, when we seek the Lord with the savour of good living, and the fragrant report of good works. Those women, when they brought their spices, saw a vision of Angels, and, in sooth, those souls whose godly desires do move them to seek the Lord with the savour of good lives, do see the countrymen of our Fatherland which is above.\par
\switchcolumn*

\LA
\begin{Resp}\Rsign Angelus Dómini descéndit de cælo, et accédens revólvit lápidem, et super eum sedit, et dixit muliéribus: \Star{} Nolíte timére: scio enim quia crucifíxum quǽritis: jam surréxit: veníte, et vidéte locum, ubi pósitus erat Dóminus, allelúja.\end{Resp}
\begin{Resp}\Vsign Et introëúntes in monuméntum, vidérunt júvenem sedéntem in dextris, coopértum stola cándida, et obstupuérunt: qui dixit illis.\end{Resp}
\begin{Resp}\Rsign Nolíte timére: scio enim quia crucifíxum quǽritis: jam surréxit: veníte, et vidéte locum, ubi pósitus erat Dóminus, allelúja.\end{Resp}
\begin{Resp}\Vsign Glória Patri, et Fílio, \Star{} et Spirítui Sancto.\end{Resp}
\begin{Resp}\Rsign Angelus Dómini descéndit de cælo, et accédens revólvit lápidem, et super eum sedit, et dixit muliéribus: * Nolíte timére: scio enim quia crucifíxum quǽritis: jam surréxit: veníte, et vidéte locum, ubi pósitus erat Dóminus, allelúja.\end{Resp}
\switchcolumn
\EN
\begin{Resp}\Rsign The Angel of the Lord descended from heaven, and came and rolled back the stone, and sat upon it, and said unto the women: \Star{} Fear not ye; for I know that ye seek Him That was crucified: He is risen already. Come, see the place where the Lord was laid, alleluia.\end{Resp}
\begin{Resp}\Vsign And entering into the sepulchre, they saw a young man sitting on the right side, clothed in a long white garment, and they were affrighted; and he saith unto them:\end{Resp}
\begin{Resp}\Rsign Fear not ye: for I know that ye seek Him That was crucified: He is risen already; come, see the place where the Lord was laid, alleluia.\end{Resp}
\begin{Resp}\Vsign Glory be to the Father, and to the Son, \Star{} and to the Holy Ghost.\end{Resp}
\begin{Resp}\Rsign The Angel of the Lord descended from heaven, and came and rolled back the stone and sat upon it, and said unto the women: Fear not ye: for I know that ye seek Him That was crucified: He is risen already: Come, see the place where the Lord was laid, alleluia.\end{Resp}
\switchcolumn*

\LA
\LessonHead{Lectio II}
\switchcolumn
\EN
\LessonHead{Lesson II}
\switchcolumn*

\LA
\noindent Notándum vero nobis est, quidnam sit, quod in dextris sedére Angelus cérnitur. Quid namque per sinístram, nisi vita præsens: quid vero per déxteram, nisi perpétua vita designátur? Unde in Cánticis canticórum scriptum est: Læva ejus sub cápite meo, et déxtera illíus amplexábitur me. Quia ergo Redémptor noster jam præséntis vitæ corruptiónem transíerat, recte Angelus qui nuntiáre perénnem ejus vitam vénerat, in déxtera sedébat. Qui stola cándida coopértus appáruit: quia festivitátis nostræ gáudia nuntiávit. Candor étenim vestis, splendórem nostræ denúntiat solemnitátis. Nostræ dicámus, an suæ? Sed ut fateámur vérius, et suæ dicámus, et nostræ. Illa quippe Redemptóris nostri resurréctio et nostra festívitas fuit, quia nos ad immortalitátem redúxit: et Angelórum festívitas éxstitit, quia nos revocándo ad cæléstia, eórum númerum implévit.\par
\switchcolumn
\EN
\noindent It behoveth us to mark what this meaneth, that they saw the Angel sitting on the right side. For what signifieth the left, but this life which now is? or the right, but life everlasting? Whence also it is written in the Song of Songs (ii. 6): His left hand is under my head, and His right hand doth embrace me. Since, therefore, our Redeemer had passed from the corruption of this life which now is, the Angel which told that His undying life was come, sat, as became him, on the right side. They saw him clothed in a white garment, for he was herald of the joy of this our great solemnity, and the glistering whiteness of his raiment told of the brightness of this holy Festival of ours. Of ours, said I? or of his? But if we will speak the truth, we must acknowledge that it is both his and ours. The Again-rising of our Redeemer is a Festival of gladness for us, for us it biddeth know that we shall not die for ever; and for Angels also it is a festival of gladness, for it biddeth them know that we are called to fulfill their number in heaven.\par
\switchcolumn*

\LA
\begin{Resp}\Rsign Cum transísset sábbatum, María Magdaléne, et María Jacóbi, et Salóme emérunt arómata, \Star{} Ut veniéntes úngerent Jesum, allelúja, allelúja.\end{Resp}
\begin{Resp}\Vsign Et valde mane una sabbatórum, véniunt ad monuméntum, orto jam sole.\end{Resp}
\begin{Resp}\Rsign Ut veniéntes úngerent Jesum, allelúja, allelúja.\end{Resp}
\begin{Resp}\Vsign Glória Patri, et Fílio, \Star{} et Spirítui Sancto.\end{Resp}
\begin{Resp}\Rsign Ut veniéntes úngerent Jesum, allelúja, allelúja.\end{Resp}
\switchcolumn
\EN
\begin{Resp}\Rsign When the Sabbath was passed, Mary Magdalene, and Mary the mother of James, and Salome, had bought sweet spices, \Star{} That they might come and anoint Jesus, alleluia, alleluia.\end{Resp}
\begin{Resp}\Vsign And very early in the morning, the first day of the week, they came unto the sepulchre, at the rising of the sun.\end{Resp}
\begin{Resp}\Rsign That they might come and anoint Jesus, alleluia, alleluia.\end{Resp}
\begin{Resp}\Vsign Glory be to the Father, and to the Son, \Star{} and to the Holy Ghost.\end{Resp}
\begin{Resp}\Rsign That they might come and anoint Jesus, alleluia, alleluia.\end{Resp}
\switchcolumn*

\LA
\LessonHead{Lectio III}
\switchcolumn
\EN
\LessonHead{Lesson III}
\switchcolumn*

\LA
\noindent In sua ergo ac nostra festivitáte Angelus in albis véstibus appáruit: quia dum nos per resurrectiónem Domínicam ad supérna redúcimur, cæléstis pátriæ damna reparántur. Sed quid adveniéntes féminas affátur, audiámus: Nolíte expavéscere. Ac si apérte dicat, Páveant illi, qui non amant advéntum supernórum cívium: pertiméscant, qui carnálibus desidériis pressi, ad eórum se societátem pertíngere posse despérant. Vos autem, cur pertiméscitis, quæ vestros concíves vidétis? Unde et Matthǽus Angelum apparuísse descríbens, ait: Erat aspéctus ejus sicut fulgur, et vestiménta ejus sicut nix. In fúlgure étenim terror timóris est, in nive autem blandiméntum candóris.\par
\switchcolumn
\EN
\noindent See this glad Festival then, which is both his and ours, the Angel appeared in white raiment. For as the Lord, rising again from the dead, leadeth us unto the mansions above, He repaireth the breaches of the heavenly Fatherland. But what meaneth this, that the Angel said unto the women which came to the sepulchre: Fear not? Is it not as though he had said openly: Let them fear which love not the coming of the heavenly countrymen; let them be afraid who are so laden by fleshly lusts, that they have lost all hope ever to be joined to their company. But as for you, why fear ye, who, when ye see us, see but your fellow countrymen? Hence also Matthew, writing of the guise of the Angel, saith (xxviii. 3): His countenance was like lightning, and His raiment white as snow. The lightning speaketh of fear and great dread, the snow of the soft brilliancy of rejoicing.\par
\switchcolumn*

\LA
\TeDeum{Te Deum}
\switchcolumn
\EN
\TeDeum{Te Deum}
\switchcolumn*

\end{paracol}

\Office{Die II infra octavam Paschæ}{Monday within the Octave of Easter}{Duplex I classis}
\addcontentsline{toc}{subsection}{Die II infra octavam Paschæ\enspace\textit{Monday within the Octave of Easter}}
\begin{paracol}{2}
\LA
\LessonHead{Lectio I}
\switchcolumn
\EN
\LessonHead{Lesson I}
\switchcolumn*

\LA
\LessonTitle{Léctio sancti Evangélii secúndum Lucam}
\switchcolumn
\EN
\LessonTitle{Continuation of the Holy Gospel according to Luke}
\switchcolumn*

\LA
\Cite{Luc 24:13-35}
\switchcolumn
\EN
\Cite{Luke 24:13-35}
\switchcolumn*

\LA
\noindent In illo témpore: Duo ex discípulis Jesu ibant ipsa die in castéllum, quod erat in spátio stadiórum sexagínta ab Jerúsalem, nómine Emmaus. Et réliqua.\par
\switchcolumn
\EN
\noindent And behold, two of them went, the same day, to a town which was sixty furlongs from Jerusalem, named Emmaus. And so forth.\par
\switchcolumn*

\LA
\LessonTitle{Homilía sancti Gregórii Papæ}
\switchcolumn
\EN
\LessonTitle{Homily by Pope St. Gregory (the Great)}
\switchcolumn*

\LA
\Source{Homilia 23 in Evangelia}
\switchcolumn
\EN
\Source{23rd on the Gospels}
\switchcolumn*

\LA
\noindent Audístis, fratres caríssimi, quia duóbus discípulis ambulántibus in via, non quidem credéntibus, sed tamen de se loquéntibus, Dóminus appáruit: sed eis spéciem, quam recognóscerent, non osténdit. Hoc ergo egit foris Dóminus in óculis córporis, quod apud ipsos agebátur intus in óculis cordis. Ipsi namque apud semetípsos intus et amábant, et dubitábant: eis autem Dóminus foris et præsens áderat, et quis esset non ostendébat. De se ergo loquéntibus præséntiam exhíbuit: sed de se dubitántibus cognitiónis suæ spéciem abscóndit.\par
\switchcolumn
\EN
\noindent Dearly beloved brethren, ye hear, how that while two of His disciples walked together in the way, not believing in His Resurrection, but talking, together concerning Him, the Lord manifested Himself unto them, but yet held their eyes that they should not know Him. This holding of the eyes of their body, wrought by the Lord, was a figure of the spiritual veil which was yet upon the eyes of their heart. For in their heart they loved and yet doubted: even as the Lord drew near to them outwardly, but showed not Who He was. To them that talked together of Him, He revealed His immediate presence; but hid, from them that doubted, the knowledge of His Person.\par
\switchcolumn*

\LA
\begin{Resp}\Rsign María Magdaléne, et áltera María ibant dilúculo ad monuméntum: \Star{} Jesum quem quǽritis, non est hic, surréxit sicut locútus est, præcédet vos in Galilǽam, ibi eum vidébitis, allelúja, allelúja.\end{Resp}
\begin{Resp}\Vsign Et valde mane una sabbatórum véniunt ad monuméntum, orto jam sole: et introëúntes vidérunt júvenem sedéntem in dextris, qui dixit illis.\end{Resp}
\begin{Resp}\Rsign Jesum quem quǽritis, non est hic, surréxit sicut locútus est, præcédet vos in Galilǽam, ibi eum vidébitis, allelúja, allelúja.\end{Resp}
\switchcolumn
\EN
\begin{Resp}\Rsign Mary Magdalene and the other Mary went very early to the sepulchre. \Star{} That Jesus whom ye seek, is not here: for He is risen, as He said: He goeth before you into Galilee; there shall ye see Him, alleluia, alleluia.\end{Resp}
\begin{Resp}\Vsign And very early in the morning, the first day of the week, they came unto the sepulchre, at the rising of the sun; and, entering into the sepulchre, they saw a young man sitting upon the right side, who saith unto them:\end{Resp}
\begin{Resp}\Rsign That Jesus Whom ye seek is not here: for He is risen, as He said: He goeth before you into Galilee: there shall ye see Him, alleluia, alleluia.\end{Resp}
\switchcolumn*

\LA
\LessonHead{Lectio II}
\switchcolumn
\EN
\LessonHead{Lesson II}
\switchcolumn*

\LA
\noindent Verba quidem cóntulit, durítiam intelléctus increpávit, sacræ Scriptúræ mystéria, quæ de seípso erant, apéruit: et tamen quia adhuc in eórum córdibus peregrínus erat a fide, se ire lóngius finxit. Fíngere namque, compónere dícimus: unde et compositóres luti, fígulos vocámus. Nihil ergo simplex Véritas per duplicitátem fecit: sed talem se eis exhíbuit in córpore, qualis apud illos erat in mente. Probándi autem erant, si hi, qui eum etsi necdum ut Deum dilígerent, saltem ut peregrínum amáre potuíssent.\par
\switchcolumn
\EN
\noindent He spoke to them; He rebuked the hardness of their heart; He expounded unto them in all the Scriptures the things concerning Himself: and, nevertheless, seeing that He was yet a stranger to faith in their hearts, He made as though He would have gone further. These words He made as though would here seem to mean He feigned, but He Who is simple Truth doth nothing with feigning: He only showed Himself to them in bodily manners, as He was towards them spiritually; but they were put to the proof whether, though they loved Him not yet as their God, they could love Him at least as a wayfarer.\par
\switchcolumn*

\LA
\begin{Resp}\Rsign Surréxit pastor bonus, qui ánimam suam pósuit pro óvibus suis, et pro grege suo mori dignátus est: \Star{} Allelúja, allelúja, allelúja.\end{Resp}
\begin{Resp}\Vsign Etenim Pascha nostrum immolátus est Christus.\end{Resp}
\begin{Resp}\Rsign Allelúja, allelúja, allelúja.\end{Resp}
\begin{Resp}\Vsign Glória Patri, et Fílio, \Star{} et Spirítui Sancto.\end{Resp}
\begin{Resp}\Rsign Allelúja, allelúja, allelúja.\end{Resp}
\switchcolumn
\EN
\begin{Resp}\Rsign The Good Shepherd, Who laid down His life for the sheep, yea, Who was contented even to die for His flock, the Good Shepherd is risen again. \Star{} Alleluia, alleluia, alleluia.\end{Resp}
\begin{Resp}\Vsign For even Christ our Passover is sacrificed for us.\end{Resp}
\begin{Resp}\Rsign Alleluia, alleluia, alleluia.\end{Resp}
\begin{Resp}\Vsign Glory be to the Father, and to the Son, \Star{} and to the Holy Ghost.\end{Resp}
\begin{Resp}\Rsign Alleluia, alleluia, alleluia.\end{Resp}
\switchcolumn*

\LA
\LessonHead{Lectio III}
\switchcolumn
\EN
\LessonHead{Lesson III}
\switchcolumn*

\LA
\noindent Sed quia esse extránei a caritáte non póterant hi, cum quibus Véritas gradiebátur: eum ad hospítium quasi peregrínum vocant. Cur autem dícimus, vocant, cum illic scriptum sit: Et coëgérunt eum? Ex quo nimírum exémplo collígitur, quia peregríni ad hospítium non solum invitándi sunt, sed étiam trahéndi. Mensam ígitur ponunt, panes cibósque ófferunt: et Deum, quem in Scriptúræ sacræ expositióne non cognóverant, in panis fractióne cognóscunt. Audiéndo ergo præcépta Dei illumináti non sunt, faciéndo illumináti sunt: quia scriptum est: Non auditóres legis justi sunt apud Deum, sed factóres legis justificabúntur. Quisquis ergo vult audíta intellégere, festínet ea, quæ jam audíre pótuit, ópere implére. Ecce Dóminus non est cógnitus dum loquerétur, et dignátus est cognósci, dum páscitur.\par
\switchcolumn
\EN
\noindent But since it was impossible, that they with whom Truth walked, should be loveless, they asked Him as a wayfarer to take of their hospitality. But why say we that they asked Him, when it is written: And they constrained Him? From their example we learn that we ought not only to bid, but also to urge, wayfarers to our hospitable entertainment. They laid a table therefore, and set before Him bread and meat; and that God Whom they had not known in the expounding of the Holy Scripture, they knew in the breaking of bread. In hearing the commandments of God they were not enlightened, but they were enlightened in the doing of them: as it is written: Not the hearers of the law are just before God, but the doers of the law shall be justified. (Rom. ii. 13.) Whosoever therefore will understand that which he heareth, let him make haste to practice in his works that which he hath already been able to hear. Behold, the Lord was not known while He spake, but He was contented to be known when He broke bread.\par
\switchcolumn*

\LA
\TeDeum{Te Deum}
\switchcolumn
\EN
\TeDeum{Te Deum}
\switchcolumn*

\end{paracol}

\Office{Die III infra octavam Paschæ}{Tuesday within the Octave of Easter}{Duplex I classis}
\addcontentsline{toc}{subsection}{Die III infra octavam Paschæ\enspace\textit{Tuesday within the Octave of Easter}}
\begin{paracol}{2}
\LA
\LessonHead{Lectio I}
\switchcolumn
\EN
\LessonHead{Lesson I}
\switchcolumn*

\LA
\LessonTitle{Léctio sancti Evangélii secúndum Lucam}
\switchcolumn
\EN
\LessonTitle{Continuation of the Holy Gospel according to Luke}
\switchcolumn*

\LA
\Cite{Luc 24:36-47}
\switchcolumn
\EN
\Cite{Luke 24:36-47}
\switchcolumn*

\LA
\noindent In illo témpore: Stetit Jesus in médio discipulórum, et dicit eis: Pax vobis: ego sum, nolíte timére. Et réliqua.\par
\switchcolumn
\EN
\noindent At that time: Now whilst they were speaking these things, Jesus stood in the midst of them, and saith to them: Peace be to you; it is I, fear not. And so forth.\par
\switchcolumn*

\LA
\LessonTitle{Homilía sancti Ambrósii Epíscopi}
\switchcolumn
\EN
\LessonTitle{Homily by St. Ambrose, Bishop (of Milan.)}
\switchcolumn*

\LA
\Source{Liber 10. Comment. in Lucam, cap. 24, ante finem}
\switchcolumn
\EN
\Source{Bk. x Comm. on Luke xxiv}
\switchcolumn*

\LA
\noindent Mirum, quo modo se natúra corpórea per impenetrábile corpus infúderit invisíbili áditu, visíbili conspéctu: tangi fácilis, diffícilis æstimári. Dénique conturbáti discípuli æstimábant se spíritum vidére. Et ídeo Dóminus, ut spéciem nobis resurrectiónis osténderet: Palpáte, inquit, et vidéte, quia spíritus carnem et ossa non habet, sicut me vidétis habére. Non ergo per incorpóream natúram, sed per resurrectiónis qualitátem, impérvia usu clausa penetrávit. Nam quod tángitur, corpus est: quod palpátur, corpus est.\par
\switchcolumn
\EN
\noindent We see here the marvelous nature of the Lord's glorified Body. It could enter unseen, and then become seen. It could easily be touched, but Its nature is hard to understand. The disciples were affrighted, and supposed that they had seen a spirit. And therefore the Lord, that He might show us the evidence of His Resurrection, said: Handle Me, and see; for a spirit hath not flesh and bones as ye see Me have. Therefore it was not by being in a disembodied state, but by the peculiar qualities of the risen and glorified Body that He had passed through closed doors. (John xx. 19.) For that which is touched or handled is a body.\par
\switchcolumn*

\LA
\begin{Resp}\Rsign Virtúte magna reddébant Apóstoli, \Star{} Testimónium resurrectiónis Jesu Christi Dómini nostri, allelúja, allelúja.\end{Resp}
\begin{Resp}\Vsign Repléti quidem Spíritu Sancto loquebántur cum fidúcia verbum Dei.\end{Resp}
\begin{Resp}\Rsign Testimónium resurrectiónis Jesu Christi Dómini nostri, allelúja, allelúja.\end{Resp}
\switchcolumn
\EN
\begin{Resp}\Rsign With great power gave the Apostles \Star{} Witness of the Resurrection of our Lord Jesus Christ, alleluia, alleluia.\end{Resp}
\begin{Resp}\Vsign They were all filled with the Holy Ghost, and they spoke the Word of God with boldness.\end{Resp}
\begin{Resp}\Rsign Witness of the Resurrection of our Lord Jesus Christ, alleluia, alleluia.\end{Resp}
\switchcolumn*

\LA
\LessonHead{Lectio II}
\switchcolumn
\EN
\LessonHead{Lesson II}
\switchcolumn*

\LA
\noindent In córpore autem resurgémus. Seminátur enim corpus animále, surgit corpus spiritále: sed illud subtílius, hoc crássius, útpote adhuc terrénæ labis qualitáte concrétum. Nam quómodo non corpus, in quo manébant insígnia vúlnerum, vestígia cicatrícum, quæ Dóminus palpánda óbtulit? In quo non solum fidem firmat, sed étiam devotiónem ácuit, quod vúlnera suscépta pro nobis cælo inférre máluit, abolére nóluit: ut Deo Patri nostræ prétia libertátis osténderet. Talem sibi Pater ad déxteram locat, trophǽum nostræ salútis ampléctens: tales illic Mártyres nobis cicatrícis suæ coróna monstrávit.\par
\switchcolumn
\EN
\noindent We shall all rise again with our bodies. But it is sown a natural body; it is raised a spiritual body. (1. Cor. xv. 44.) The spiritual body is the finer, and the natural body is the grosser, besodden as yet by the corruption of earth. Was not That a real Body, wherein remained those marks of His Wounds, those holes of the nail-prints, which the Lord bade His disciples to handle? Hereby, also, He hath not only strengthened our faith, but also quickened our love, since we know that it has been His will to carry to heaven those Wounds which He bore for our sake, and wherewith He would not make away; but plainly showeth to His Eternal Father the price of our freedom. It is as marked with these Wounds and embracing the trophy of our salvation that the Father hath said to Him, Sit Thou at My right Hand: and it is, like Him, marked with their wounds, that He hath shown us that the Martyrs, whose Crown He is, are, and will be with Him there.\par
\switchcolumn*

\LA
\begin{Resp}\Rsign De ore prudéntis procédit mel, allelúja: dulcédo mellis est sub língua ejus, allelúja: \Star{} Favus distíllans lábia ejus, allelúja, allelúja.\end{Resp}
\begin{Resp}\Vsign Sapiéntia requiéscit in corde ejus, et prudéntia in sermóne oris illíus.\end{Resp}
\begin{Resp}\Rsign Favus distíllans lábia ejus, allelúja, allelúja.\end{Resp}
\begin{Resp}\Vsign Glória Patri, et Fílio, \Star{} et Spirítui Sancto.\end{Resp}
\begin{Resp}\Rsign Favus distíllans lábia ejus, allelúja, allelúja.\end{Resp}
\switchcolumn
\EN
\begin{Resp}\Rsign From the mouth of the wise doth proceed honey, alleluia: the sweetness of honey is under his tongue, alleluia. \Star{} His lips drop as the honey-comb, alleluia, alleluia.\end{Resp}
\begin{Resp}\Vsign Wisdom doth abide in his heart, and out of his mouth cometh understanding.\end{Resp}
\begin{Resp}\Rsign His lips drop as the honey-comb, alleluia, alleluia.\end{Resp}
\begin{Resp}\Vsign Glory be to the Father, and to the Son, \Star{} and to the Holy Ghost.\end{Resp}
\begin{Resp}\Rsign His lips drop as the honey-comb, alleluia, alleluia.\end{Resp}
\switchcolumn*

\LA
\LessonHead{Lectio III}
\switchcolumn
\EN
\LessonHead{Lesson III}
\switchcolumn*

\LA
\noindent Et quóniam sermo huc noster evásit, considerémus qua grátia secúndum Joánnem credíderint Apóstoli, qui gavísi sunt; secúndum Lucam quasi incréduli redarguántur: ibi Spíritum Sanctum accéperint, hic sedére in civitáte jubeántur, quoadúsque induántur virtúte ex alto. Et vidétur mihi ille quasi Apóstolus majóra et altióra tetigísse, hic sequéntia et humánis próxima: hic histórico usus circúitu, ille compéndio: quia et de illo dubitári non potest, qui testimónium pérhibet de iis, quibus ipse intérfuit, et verum est testimónium ejus: et ab hoc quoque, qui Evangelísta esse méruit, vel negligéntiæ, vel mendácii suspiciónem æquum est propulsári. Et ídeo verum putámus utrúmque, non sententiárum varietáte, nec personárum diversitáte distínctum. Nam etsi primo Lucas eos non credidísse dicat, póstea tamen credidísse demónstrat: et si prima considerémus, contrária sunt: si sequéntia, certum est conveníre.\par
\switchcolumn
\EN
\noindent And now, since our Lesson from Luke here faileth, let us have recourse to John, and consider how that, according to him, (xx. 20,) then were the disciples glad when they saw the Lord, and received the grace of faith. According to Luke, He upbraided them with their unbelief, but according to John He said also, Receive ye the Holy Ghost. Luke, not John, hath, Tarry ye in the city of Jerusalem, until ye be endued with power from on high. Indeed, to me it seemeth as though the one Evangelist had busied himself with the greater and higher matters, and the other with the narrative, and such things as are more human: the one with the course, the other with the essence, of history. For as it is impossible to doubt the word of him who testifieth of these things, (John xxi. 24,) and who saw these things, and concerning whom we know that his testimony is true, (xxi. 24,) so is it sinful to think of negligence or falsehood as attaching to the other, even Luke, who earned to himself to be an Evangelist, albeit he was not an Apostle, and therefore we hold that both are truthful, neither are they at variance one with the other, either in the difference of the words they use, or in the sacredness of their characters as Evangelists. For though Luke saith that at the first the Apostles believed not, yet he showeth that afterward they believed: and although, if we regard only the first fact, the Evangelists seem divergent one from the other, yet, when we consider what cometh afterward, we see that they are at one.\par
\switchcolumn*

\LA
\TeDeum{Te Deum}
\switchcolumn
\EN
\TeDeum{Te Deum}
\switchcolumn*

\end{paracol}

\Office{Die IV infra octavam Paschæ}{Wednesday within the Octave of Easter}{Semiduplex I classis}
\addcontentsline{toc}{subsection}{Die IV infra octavam Paschæ\enspace\textit{Wednesday within the Octave of Easter}}
\begin{paracol}{2}
\LA
\LessonHead{Lectio I}
\switchcolumn
\EN
\LessonHead{Lesson I}
\switchcolumn*

\LA
\LessonTitle{Léctio sancti Evangélii secúndum Joánnem}
\switchcolumn
\EN
\LessonTitle{Continuation of the Holy Gospel according to John}
\switchcolumn*

\LA
\Cite{Joannes 21:1-14}
\switchcolumn
\EN
\Cite{John 21:1-14}
\switchcolumn*

\LA
\noindent In illo témpore: Manifestávit se íterum Jesus discípulis ad mare Tiberíadis. Manifestávit autem sic: Erant simul Simon Petrus, et Thomas, qui dícitur Dídymus. Et réliqua.\par
\switchcolumn
\EN
\noindent At that time: After this, Jesus showed himself again to the disciples at the sea of Tiberias. And he showed himself after this manner. There were together Simon Peter, and Thomas, who is called Didymus, And so forth.\par
\switchcolumn*

\LA
\LessonTitle{Homilía sancti Gregórii Papæ}
\switchcolumn
\EN
\LessonTitle{Homily by Pope St. Gregory (the Great.)}
\switchcolumn*

\LA
\Source{Homilia 24 in Evangelia}
\switchcolumn
\EN
\Source{24th on the Gospels}
\switchcolumn*

\LA
\noindent Léctio sancti Evangélii, quæ modo in áuribus vestris lecta est, fratres mei, quæstióne ánimum pulsat, sed pulsatióne sua vim discretiónis índicat. Quæri étenim potest, cur Petrus, qui piscátor ante conversiónem fuit, post conversiónem ad piscatiónem rédiit: et cum Véritas dicat: Nemo mittens manum suam ad arátrum, et aspíciens retro, aptus est regno Dei: cur repétiit quod derelíquit? Sed si virtus discretiónis inspícitur, cítius vidétur: quia nimírum negótium, quod ante conversiónem sine peccáto éxstitit, hoc étiam post conversiónem repétere culpa non fuit.\par
\switchcolumn
\EN
\noindent Dearly beloved brethren, the portion of the Holy Gospel which hath but now been read in your ears, knocketh loudly at the door of your heart, with a certain question, the answer whereto calleth for thought. This same question is: Wherefore did Peter, who had before his conversion been a fisher, wherefore did he, after his conversion, again go a-fishing? since the Truth hath said: No man, having put his hand to the plough, and looking back, is fit for the kingdom of God? (Luke ix. 62.) Wherefore did Peter return to that which he had left? But we thought we see the answer to his question. The trade which was harmless before his conversion, did not become harmful because he had been converted.\par
\switchcolumn*

\LA
\begin{Resp}\Rsign Ecce vicit leo de tribu Juda, radix David, aperíre librum, et sólvere septem signácula ejus: \Star{} Allelúja, allelúja, allelúja.\end{Resp}
\begin{Resp}\Vsign Dignus est Agnus, qui occísus est, accípere virtútem, et divinitátem, et sapiéntiam, et fortitúdinem, et honórem, et glóriam, et benedictiónem.\end{Resp}
\begin{Resp}\Rsign Allelúja, allelúja, allelúja.\end{Resp}
\switchcolumn
\EN
\begin{Resp}\Rsign Behold, the Lion of the tribe of Judah, the Root of David, hath prevailed to open the Book, and to loose the seven seals thereof. \Star{} Alleluia, alleluia, alleluia.\end{Resp}
\begin{Resp}\Vsign Worthy is the Lamb That was slain to receive power, and riches, in wisdom, and strength, and honour, and glory, and blessing.\end{Resp}
\begin{Resp}\Rsign Alleluia, alleluia, alleluia.\end{Resp}
\switchcolumn*

\LA
\LessonHead{Lectio II}
\switchcolumn
\EN
\LessonHead{Lesson II}
\switchcolumn*

\LA
\noindent Nam piscatórem Petrum, Matthǽum vero teloneárium scimus: et post conversiónem suam ad piscatiónem Petrus rédiit, Matthǽus vero ad telónei negótium non resédit: quia áliud est victum per piscatiónem quǽrere, áliud autem telónei lucris pecúnias augére. Sunt enim pléraque negótia, quæ sine peccátis exhibéri aut vix, aut nullátenus possunt. Quæ ergo ad peccátum ímplicant, ad hæc necésse est, ut post conversiónem ánimus non recúrrat.\par
\switchcolumn
\EN
\noindent We know that Peter had been a fisherman, and Matthew a publican, and that Peter after his conversion went back to his fishing, but Matthew did not return to the receipt of custom. It is one thing to seek a livelihood by fishing, and another to amass money by farming of taxes. There are many kinds of business in which it is difficult or impossible to be engaged without committing sin, and to such kinds of business as these, he which hath once been converted must not again betake himself.\par
\switchcolumn*

\LA
\begin{Resp}\Rsign Ego sum vitis vera, et vos pálmites: \Star{} Qui manet in me, et ego in eo, hic fert fructum multum, allelúja, allelúja.\end{Resp}
\begin{Resp}\Vsign Sicut diléxit me Pater, et ego diléxi vos.\end{Resp}
\begin{Resp}\Rsign Qui manet in me, et ego in eo, hic fert fructum multum, allelúja, allelúja.\end{Resp}
\begin{Resp}\Vsign Glória Patri, et Fílio, \Star{} et Spirítui Sancto.\end{Resp}
\begin{Resp}\Rsign Qui manet in me, et ego in eo, hic fert fructum multum, allelúja, allelúja.\end{Resp}
\switchcolumn
\EN
\begin{Resp}\Rsign I am the True Vine, and ye are the branches; \Star{} He that abideth in Me, and I in him, the same bringeth forth much fruit, alleluia, alleluia.\end{Resp}
\begin{Resp}\Vsign As the Father hath loved Me, so have I loved you.\end{Resp}
\begin{Resp}\Rsign He that abideth in Me, and I in him, the same bringeth forth much fruit, alleluia, alleluia.\end{Resp}
\begin{Resp}\Vsign Glory be to the Father, and to the Son, \Star{} and to the Holy Ghost.\end{Resp}
\begin{Resp}\Rsign He that abideth in Me, and I in him, the same bringeth forth much fruit, alleluia, alleluia.\end{Resp}
\switchcolumn*

\LA
\LessonHead{Lectio III}
\switchcolumn
\EN
\LessonHead{Lesson III}
\switchcolumn*

\LA
\noindent Quæri étiam potest, cur discípulis in mari laborántibus, post resurrectiónem suam Dóminus in líttore stetit, qui ante resurrectiónem suam coram discípulis in flúctibus maris ambulávit. Cujus rei rátio festíne cognóscitur, si ipsa, quæ tunc ínerat, causa pensétur. Quid enim mare, nisi præsens sǽculum signat, quod se cásuum tumúltibus, et undis vitæ corruptíbilis illídit? Quid per soliditátem líttoris, nisi illa perpetúitas quiétis ætérnæ figurátur? Quia ergo discípuli adhuc flúctibus mortális vitæ ínerant, in mari laborábant: quia autem Redémptor noster jam corruptiónem carnis excésserat, post resurrectiónem suam in líttore stabat.\par
\switchcolumn
\EN
\noindent It may likewise be asked why, when the disciples were toiling in the sea, the Lord, after His Resurrection, stood on the shore, whereas, before His Resurrection, He had walked on the waves before them all. The reason of this is quickly known if we will think of the end which it then served. The sea is a figure of this present world, tossed to and fro by changing fortune, and continually ebbing and flowing with the diverse tides of life. The stableness of the shore is an image of the never-ending rest of the eternal home. The disciples therefore, for that they were yet tossed to and fro upon the waves of a dying life, were toiling in the sea, but He our Redeemer, Who had already laid aside that which in this body is subject to corruption, and had risen again from the dead, He stood upon the shore.\par
\switchcolumn*

\LA
\TeDeum{Te Deum}
\switchcolumn
\EN
\TeDeum{Te Deum}
\switchcolumn*

\end{paracol}

\Office{Die V infra octavam Paschæ}{Thursday within the Octave of Easter}{Semiduplex I classis}
\addcontentsline{toc}{subsection}{Die V infra octavam Paschæ\enspace\textit{Thursday within the Octave of Easter}}
\begin{paracol}{2}
\LA
\LessonHead{Lectio I}
\switchcolumn
\EN
\LessonHead{Lesson I}
\switchcolumn*

\LA
\LessonTitle{Léctio sancti Evangélii secúndum Joánnem}
\switchcolumn
\EN
\LessonTitle{Continuation of the Holy Gospel according to John}
\switchcolumn*

\LA
\Cite{Joannes 20:11-18}
\switchcolumn
\EN
\Cite{John 20:11-18}
\switchcolumn*

\LA
\noindent In illo témpore: María stabat ad monuméntum foris, plorans. Dum ergo fleret, inclinávit se, et prospéxit in monuméntum: et vidit duos Angelos in albis, sedéntes. Et réliqua.\par
\switchcolumn
\EN
\noindent At that time: But Mary stood at the sepulchre without, weeping. Now as she was weeping, she stooped down, and looked into the sepulchre. And so forth.\par
\switchcolumn*

\LA
\LessonTitle{Homilía sancti Gregórii Papæ}
\switchcolumn
\EN
\LessonTitle{Homily by Pope St. Gregory (the Great.)}
\switchcolumn*

\LA
\Source{Homilia 25 in Evangelia}
\switchcolumn
\EN
\Source{25th on the Gospels}
\switchcolumn*

\LA
\noindent María Magdaléne, quæ fúerat in civitáte peccátrix, amándo veritátem, lavit lácrimis máculas críminis: et vox Veritátis implétur, qua dícitur: Dimíssa sunt ei peccáta multa, quia diléxit multum. Quæ enim prius frígida peccándo remánserat, póstmodum amándo fórtiter ardébat. Nam postquam venit ad monuméntum, ibíque corpus Domínicum non invénit, sublátum crédidit, atque discípulis nuntiávit: qui veniéntes vidérunt, atque ita esse, ut múlier díxerat, credidérunt. Et de eis prótinus scriptum est: Abiérunt ergo discípuli ad semetípsos: ac deínde subjúngitur: María autem stabat ad monuméntum foris, plorans.\par
\switchcolumn
\EN
\noindent Mary Magdalene, a woman in the city, who was a sinner, through love of the truth washed away by her tears the befoulment of her sin, and the word of the Truth was fulfilled which He spake: Her sins, which are many, are forgiven: for she loved much. (Luke vii. 47.) She that had remained cold while she sinned, became burning when she loved. For after that she had been to the Sepulchre, and had not found there the Body of the Lord, and had believed that It had been taken away, and had told His disciples, they came and saw, and thought it was even as the woman had said, and it is written: Then the disciples went away again unto their own home but Mary stood without at the sepulchre, weeping.\par
\switchcolumn*

\LA
\begin{Resp}\Rsign Tulérunt Dóminum meum, et néscio ubi posuérunt eum. Dicunt ei Angeli: Múlier, quid ploras? surréxit sicut dixit: \Star{} Præcédet vos in Galilǽam: ibi eum vidébitis, allelúja, allelúja.\end{Resp}
\begin{Resp}\Vsign Cum ergo fleret, inclinávit se, et prospéxit in monuméntum: et vidit duos Angelos in albis, sedéntes, qui dicunt ei.\end{Resp}
\begin{Resp}\Rsign Præcédet vos in Galilǽam: ibi eum vidébitis, allelúja, allelúja.\end{Resp}
\switchcolumn
\EN
\begin{Resp}\Rsign They have taken away my Lord, and I know not where they have laid Him. The Angels say unto her: Woman, why weepest thou? He is risen, as He said. \Star{} He goeth before you into Galilee; there shall ye see Him, alleluia, alleluia.\end{Resp}
\begin{Resp}\Vsign And as she wept, she stooped down and looked into the Sepulchre, and saw two Angels in white, sitting; and they say unto her:\end{Resp}
\begin{Resp}\Rsign He goeth before you into Galilee; there shall ye see Him, alleluia, alleluia.\end{Resp}
\switchcolumn*

\LA
\LessonHead{Lectio II}
\switchcolumn
\EN
\LessonHead{Lesson II}
\switchcolumn*

\LA
\noindent Qua in re pensándum est, hujus mulíeris mentem quanta vis amóris accénderat, quæ a monuménto Dómini, étiam discípulis recedéntibus, non recedébat. Exquirébat quem non invénerat: flebat inquírendo, et amóris sui igne succénsa, ejus, quem ablátum crédidit, ardébat desidério. Unde cóntigit, ut eum sola tunc vidéret, quæ remánserat ut quǽreret: quia nimírum virtus boni óperis, perseverántia est: et voce Veritátis dícitur: Qui autem perseveráverit usque in finem, hic salvus erit.\par
\switchcolumn
\EN
\noindent In connection with this matter, we ought to ponder what great store of love there was in that woman's heart, who, when even His disciples were gone away, could not tear herself from the grave of the Lord. She sought Him Whom she had not found there, and as she sought, she wept, and the fire of love in her heart yearned after Him, Who she believed had been taken away. And so it came to pass that she, who had lingered to seek Him, was the only one who then saw Him, since the back-bone of a good work is endurance, and the voice of the Truth Himself hath said: He that endureth to the end shall be saved. (Matth. x. 22; xxiv. 13.)\par
\switchcolumn*

\LA
\begin{Resp}\Rsign Congratulámini mihi omnes qui dilígitis Dóminum, quia quem quærébam, appáruit mihi: \Star{} Et dum flerem ad monuméntum, vidi Dóminum, allelúja, allelúja.\end{Resp}
\begin{Resp}\Vsign Recedéntibus discípulis, non recedébam, et amóris ejus igne succénsa, ardébam desidério.\end{Resp}
\begin{Resp}\Rsign Et dum flerem ad monuméntum, vidi Dóminum, allelúja, allelúja.\end{Resp}
\begin{Resp}\Vsign Glória Patri, et Fílio, \Star{} et Spirítui Sancto.\end{Resp}
\begin{Resp}\Rsign Et dum flerem ad monuméntum, vidi Dóminum, allelúja, allelúja.\end{Resp}
\switchcolumn
\EN
\begin{Resp}\Rsign Rejoice with me, all ye that love the Lord: for I sought Him and He hath appeared unto me; \Star{} And while as I was weeping at the Sepulchre, I saw the Lord, alleluia, alleluia.\end{Resp}
\begin{Resp}\Vsign When His disciples were gone away, I tarried still; and the fire of love in mine heart glowed for Him.\end{Resp}
\begin{Resp}\Rsign And while as I was weeping at the Sepulchre, I saw the Lord, alleluia, alleluia.\end{Resp}
\begin{Resp}\Vsign Glory be to the Father, and to the Son, \Star{} and to the Holy Ghost.\end{Resp}
\begin{Resp}\Rsign And while as I was weeping at the Sepulchre, I saw the Lord, alleluia, alleluia.\end{Resp}
\switchcolumn*

\LA
\LessonHead{Lectio III}
\switchcolumn
\EN
\LessonHead{Lesson III}
\switchcolumn*

\LA
\noindent María ergo cum fleret, inclinávit se, et prospéxit in monuméntum. Certe jam monuméntum vácuum víderat, jam sublátum Dóminum nuntiáverat: quid est, quod se íterum inclínat, íterum vidére desíderat? Sed amánti semel aspexísse non súfficit: quia vis amóris intentiónem multíplicat inquisitiónis. Quæsívit ergo prius, et mínime invénit: perseverávit ut quǽreret, unde et cóntigit, ut inveníret: actúmque est, ut desidéria diláta créscerent, et crescéntia cáperent quod inveníssent.\par
\switchcolumn
\EN
\noindent As Mary wept there, she stooped down and looked into the Sepulchre. It was but a little while and she had seen how the Sepulchre was empty, and had told that the Lord was taken away. Why then should she stoop down and look in again? But she loved Him so well, that one look was not enough; the energy of her affection constrained her to search again and again. She began by searching and not finding; but she endured in her search, and, behold, it came to pass that she found. And this was done that our own longings for Christ's presence might be taught to expand, and know that as they expand they will meet with Him to Whom they aspire.\par
\switchcolumn*

\LA
\TeDeum{Te Deum}
\switchcolumn
\EN
\TeDeum{Te Deum}
\switchcolumn*

\end{paracol}

\Office{Die VI infra octavam Paschæ}{Friday within the Octave of Easter}{Semiduplex I classis}
\addcontentsline{toc}{subsection}{Die VI infra octavam Paschæ\enspace\textit{Friday within the Octave of Easter}}
\begin{paracol}{2}
\LA
\LessonHead{Lectio I}
\switchcolumn
\EN
\LessonHead{Lesson I}
\switchcolumn*

\LA
\LessonTitle{Léctio sancti Evangélii secúndum Matthǽum}
\switchcolumn
\EN
\LessonTitle{Continuation of the Holy Gospel according to Matthew}
\switchcolumn*

\LA
\Cite{Matt 28:16-20}
\switchcolumn
\EN
\Cite{Matt 28:16-20}
\switchcolumn*

\LA
\noindent In illo témpore: Undecim discípuli abiérunt in Galilǽam in montem, ubi constitúerat illis Jesus. Et réliqua.\par
\switchcolumn
\EN
\noindent At that time: And the eleven disciples went into Galilee, unto the mountain where Jesus had appointed them. And so forth.\par
\switchcolumn*

\LA
\LessonTitle{Homilía sancti Hierónymi Presbýteri}
\switchcolumn
\EN
\LessonTitle{Homily by St. Jerome, Priest (at Bethlehem.)}
\switchcolumn*

\LA
\Source{Lib. 4 Comment. in Matth., in fine}
\switchcolumn
\EN
\Source{Bk. iv Comm. on the end of Matth.}
\switchcolumn*

\LA
\noindent Post resurrectiónem Jesus in monte Galilǽæ conspícitur, ibíque adorátur; licet quidam dúbitent, et dubitátio eórum nostram áugeat fidem. Tunc maniféstius osténditur Thomæ, et latus láncea vulnerátum, et manus fixas demónstrat clavis. Accédens Jesus locútus est eis, dicens: Data est mihi omnis potéstas in cælo et in terra. Illi potéstas data est, qui paulo ante crucifíxus, qui sepúltus in túmulo, qui mórtuus jacúerat, qui póstea resurréxit. In cælo autem et in terra potéstas data est: ut qui ante regnábat in cælo, per fidem credéntium regnet et in terris.\par
\switchcolumn
\EN
\noindent After His Resurrection Jesus was seen on a mountain in Galilee, and there He was worshipped; and, albeit some doubted, their doubts have led to a further establishing of our faith. Then He showed Himself more openly unto Thomas, and made him handle the Side That was pierced with the spear, and the Hands wherein were the holes of the nails. And Jesus came and spake unto them, saying: All power is given unto Me in heaven and in earth. Yea, all power is given unto Him Who but a little while before had been crucified, and buried in the grave, and had lain among the dead, but Who also had risen again. Power is given unto Him in heaven and in earth, that He Who of everlasting had been King of heaven, might have a Monarchy on earth also, through the faith of them which believe in Him.\par
\switchcolumn*

\LA
\begin{Resp}\Rsign Surgens Jesus Dóminus noster, stans in médio discipulórum suórum, dixit: \Star{} Pax vobis, allelúja: gavísi sunt discípuli viso Dómino, allelúja.\end{Resp}
\begin{Resp}\Vsign Una ergo sabbatórum, cum fores essent clausæ, ubi erant discípuli congregáti, venit Jesus, et stetit in medio eórum, et dixit eis.\end{Resp}
\begin{Resp}\Rsign Pax vobis, allelúja: gavísi sunt discípuli viso Dómino, allelúja.\end{Resp}
\switchcolumn
\EN
\begin{Resp}\Rsign After that our Lord Jesus was risen again, He came and stood in the midst of His disciples, and said unto them: \Star{} Peace be unto you, alleluia. Then were the disciples glad, when they saw the Lord, alleluia.\end{Resp}
\begin{Resp}\Vsign The first day of the week, when the doors were shut where the disciples were assembled, came Jesus, and stood in the midst, and said unto them:\end{Resp}
\begin{Resp}\Rsign Peace be unto you, alleluia. Then were the disciples glad, when they saw the Lord, alleluia.\end{Resp}
\switchcolumn*

\LA
\LessonHead{Lectio II}
\switchcolumn
\EN
\LessonHead{Lesson II}
\switchcolumn*

\LA
\noindent Eúntes autem docéte omnes gentes, baptizántes eos in nómine Patris, et Fílii, et Spíritus Sancti. Primum docent omnes gentes, deínde doctas intíngunt aqua. Non enim potest fíeri, ut corpus baptísmi recípiat sacraméntum, nisi ante ánima fídei suscéperit veritátem. Baptizántur autem in nómine Patris, et Fílii, et Spíritus Sancti: ut quorum una est divínitas, una sit largítio: noménque Trinitátis, unus Deus est.\par
\switchcolumn
\EN
\noindent Go ye therefore and teach all nations, baptizing them in the Name of the Father, and of the Son, and of the Holy Ghost. First, they teach all nations; then, they wash with water them whom they have taught. For it is impossible for the body to receive the Sacrament of Baptism, unless the mind first receive the truth of the faith. And they are baptized In the Name of the Father, and of the Son, and of the Holy Ghost for, even as the Godhead of the Father, and of the Son, and of the Holy Ghost, is all One, so is the one grace of Baptism the gift of all the Three Divine Persons: and the Name of the Trinity is the Name of One God.\par
\switchcolumn*

\LA
\begin{Resp}\Rsign Expurgáte vetus ferméntum, ut sitis nova conspérsio: étenim Pascha nostrum immolátus est Christus: \Star{} Itaque epulémur in Dómino, allelúja.\end{Resp}
\begin{Resp}\Vsign Mórtuus est propter delícta nostra, et resurréxit propter justificatiónem nostram.\end{Resp}
\begin{Resp}\Rsign Itaque epulémur in Dómino, allelúja.\end{Resp}
\begin{Resp}\Vsign Glória Patri, et Fílio, \Star{} et Spirítui Sancto.\end{Resp}
\begin{Resp}\Rsign Itaque epulémur in Dómino, allelúja.\end{Resp}
\switchcolumn
\EN
\begin{Resp}\Rsign Purge out the old leaven, that ye may be new dough: for even Christ our Passover is sacrificed for us: \Star{} Therefore let us keep the Feast, in the Lord, alleluia.\end{Resp}
\begin{Resp}\Vsign He died for our offences, and rose again for our justification.\end{Resp}
\begin{Resp}\Rsign Therefore let us keep the Feast, in the Lord, alleluia.\end{Resp}
\begin{Resp}\Vsign Glory be to the Father, and to the Son, \Star{} and to the Holy Ghost.\end{Resp}
\begin{Resp}\Rsign Therefore let us keep the Feast, in the Lord, alleluia.\end{Resp}
\switchcolumn*

\LA
\LessonHead{Lectio III}
\switchcolumn
\EN
\LessonHead{Lesson III}
\switchcolumn*

\LA
\noindent Docéntes eos serváre ómnia, quæcúmque mandávi vobis. Ordo præcípuus: jussit Apóstolis, ut primum docérent univérsas gentes, deínde fídei intíngerent sacraménto, et post fidem ac baptísma, quæ essent observánda præcíperent. Ac ne putémus lévia esse, quæ jussa sunt, et pauca, áddidit: Omnia quæcúmque mandávi vobis: ut quicúmque credíderint, qui in Trinitáte fúerint baptizáti, ómnia fáciant, quæ præcépta sunt. Et ecce ego vobíscum sum usque ad consummatiónem sǽculi. Qui usque ad consummatiónem sǽculi cum discípulis se futúrum esse promíttit, et illos osténdit semper esse victúros, et se nunquam a credéntibus recessúrum.\par
\switchcolumn
\EN
\noindent Preaching them to observe all things whatsoever I have commanded you. The order of the Lord's commands to the Apostles is markedly this. First, to teach all nations; secondly, to make them partake in the Sacrament of the faith; thirdly, when they had believed and been baptized, to teach them what to observe. And lest we should think that He commanded things light and few, He hath said: All things whatsoever I have commanded you, so that all, who have believed and been baptized in the Name of the Trinity, are bound to observe all things whatsoever He hath commanded. And, lo, I am with you alway, even unto the end of the world. He Who promiseth that He will be with His disciples even unto the end of the world, doth give them thereby to know that they will be always conquerors, and that He will never fail any who believe in Him.\par
\switchcolumn*

\LA
\TeDeum{Te Deum}
\switchcolumn
\EN
\TeDeum{Te Deum}
\switchcolumn*

\end{paracol}

\Office{Sabbato in Albis}{Saturday in Easter Week}{Semiduplex I classis}
\addcontentsline{toc}{subsection}{Sabbato in Albis\enspace\textit{Saturday in Easter Week}}
\begin{paracol}{2}
\LA
\LessonHead{Lectio I}
\switchcolumn
\EN
\LessonHead{Lesson I}
\switchcolumn*

\LA
\LessonTitle{Léctio sancti Evangélii secúndum Joánnem}
\switchcolumn
\EN
\LessonTitle{Continuation of the Holy Gospel according to John}
\switchcolumn*

\LA
\Cite{Joannes 20:1-9}
\switchcolumn
\EN
\Cite{John 20:1-9}
\switchcolumn*

\LA
\noindent In illo témpore: Una sábbati María Magdaléne venit mane, cum adhuc ténebræ essent, ad monuméntum. Et réliqua.\par
\switchcolumn
\EN
\noindent And on the first day of the week, Mary Magdalen cometh early, when it was yet dark, unto the sepulchre. And so forth.\par
\switchcolumn*

\LA
\LessonTitle{Homilía sancti Gregórii Papæ}
\switchcolumn
\EN
\LessonTitle{Homily by Pope St. Gregory (the Great.)}
\switchcolumn*

\LA
\Source{Homilia 22 in Evangelia}
\switchcolumn
\EN
\Source{22nd on the Gospels}
\switchcolumn*

\LA
\noindent Léctio sancti Evangélii, quam modo, fratres, audístis, valde in superfície histórica est apérta: sed ejus nobis sunt mystéria sub brevitáte requirénda. María Magdaléne, cum adhuc ténebræ essent, venit ad monuméntum. Juxta históriam notátur hora: juxta intelléctum vero mýsticum, requiréntis signátur intellegéntia. María étenim auctórem ómnium, quem in carne víderat mórtuum, quærébat in monuménto; et quia hunc mínime invénit, furátum crédidit. Adhuc ergo erant ténebræ, cum venit ad monuméntum. Cucúrrit cítius, discípulis nuntiávit: sed illi præ céteris cucurrérunt, qui præ céteris amavérunt, vidélicet Petrus et Joánnes.\par
\switchcolumn
\EN
\noindent Dearly beloved brethren, the portion of the Holy Gospel which hath just now been read in your ears, is exceeding simple on the face of it, which is its historical sense; but the mystic sense, which underlieth that other, requireth from us a little searching. Mary Magdalene came unto the Sepulchre when it was yet dark. The historic sense telleth us what was the hour of day; the mystic sense, the state of her understanding who sought. Mary Magdalene sought for Him, by Whom all things were made, and Whom she had seen die, as concerning the flesh; she sought for Him, I say, in the grave, and finding Him not, she believed that He had been stolen away. Yea, it was yet dark, when she came unto the sepulchre. Then she ran and told the disciples, but they who had loved Him most, namely Peter and John, did outrun the others.\par
\switchcolumn*

\LA
\begin{Resp}\Rsign Christus resúrgens ex mórtuis, jam non móritur, mors illi ultra non dominábitur: quod enim mórtuus est peccáto, mórtuus est semel: \Star{} Quod autem vivit, vivit Deo, allelúja, allelúja.\end{Resp}
\begin{Resp}\Vsign Mórtuus est semel propter delícta nostra, et resurréxit propter justificatiónem nostram.\end{Resp}
\begin{Resp}\Rsign Quod autem vivit, vivit Deo, allelúja, allelúja.\end{Resp}
\switchcolumn
\EN
\begin{Resp}\Rsign Knowing that Christ rising again from the dead, dieth now no more, death shall no more have dominion over him. For in that he died to sin, he died once; \Star{} In that he liveth, he liveth unto God, alleluia, alleluia.\end{Resp}
\begin{Resp}\Vsign He died once for the sins, and he rose for our justification.\end{Resp}
\begin{Resp}\Rsign In that he liveth, he liveth unto God, alleluia, alleluia.\end{Resp}
\switchcolumn*

\LA
\LessonHead{Lectio II}
\switchcolumn
\EN
\LessonHead{Lesson II}
\switchcolumn*

\LA
\noindent Currébant autem duo simul: sed Joánnes præcucúrrit cítius Petro. Venit prior ad monuméntum, et íngredi non præsúmpsit. Venit ergo postérior Petrus, et intrávit. Quid, fratres, quid cursus signíficat? Numquid hæc tam subtílis Evangelístæ descríptio a mystériis vacáre credénda est? Mínime. Neque enim se Joánnes et præísse, et non intrásse díceret, si in ipsa sui trepidatióne mystérium defuísse credidísset. Quid ergo per Joánnem, nisi synagóga: quid per Petrum, nisi Ecclésia designátur?\par
\switchcolumn
\EN
\noindent So they ran both together, but John did outrun Peter, and came first to the Sepulchre, but yet took he not upon himself to go in first. Then cometh Peter following him, and went in. What, my brethren, what did the racing of these Apostles signify? Can we believe that the description given by the deepest of the Evangelists is without a mystic interpretation? By no means. John had never told how that he did outrun Peter, and yet went not into the Sepulchre, if he had not believed that his hesitation veiled some mystery. What signifieth John but the Synagogue? or Peter, but the Church?\par
\switchcolumn*

\LA
\begin{Resp}\Rsign Isti sunt agni novélli, qui annuntiavérunt, allelúja: modo venérunt ad fontes, \Star{} Repléti sunt claritáte, allelúja, allelúja.\end{Resp}
\begin{Resp}\Vsign In conspéctu Agni amícti sunt stolis albis, et palmæ in mánibus eórum.\end{Resp}
\begin{Resp}\Rsign Repléti sunt claritáte, allelúja, allelúja.\end{Resp}
\begin{Resp}\Vsign Glória Patri, et Fílio, \Star{} et Spirítui Sancto.\end{Resp}
\begin{Resp}\Rsign Repléti sunt claritáte, allelúja, allelúja.\end{Resp}
\switchcolumn
\EN
\begin{Resp}\Rsign These are the new lambs, who have proclaimed, alleluia: they came but just now to the well \Star{} They are all filled with light, alleluia, alleluia.\end{Resp}
\begin{Resp}\Vsign In the presence of the Lamb they are clothed with white robes, and hold palms in their hands.\end{Resp}
\begin{Resp}\Rsign They are all filled with light, alleluia, alleluia.\end{Resp}
\begin{Resp}\Vsign Glory be to the Father, and to the Son, \Star{} and to the Holy Ghost.\end{Resp}
\begin{Resp}\Rsign They are all filled with light, alleluia, alleluia.\end{Resp}
\switchcolumn*

\LA
\LessonHead{Lectio III}
\switchcolumn
\EN
\LessonHead{Lesson III}
\switchcolumn*

\LA
\noindent Nec mirum esse videátur, quod per juniórem synagóga, per seniórem vero Ecclésia signári perhibétur: quia etsi ad Dei cultum prior est synagóga, quam Ecclésia géntium, ad usum tamen sǽculi prior est multitúdo géntium, quam synagóga, Paulo attestánte, qui ait: Quia non prius quod spiritále est, sed quod animále. Per seniórem ergo Petrum significátur Ecclésia géntium: per juniórem vero Joánnem synagóga Judæórum. Currunt ambo simul: quia ab ortus sui témpore usque ad occásum, pari et commúni via, etsi non pari et commúni sensu, gentílitas cum synagóga cucúrrit. Venit synagóga prior ad monuméntum, sed mínime intrávit: quia legis quidem mandáta percépit, prophetías de incarnatióne ac passióne Domínica audívit, sed crédere in mórtuum nóluit.\par
\switchcolumn
\EN
\noindent Neither must ye take it as strange that the elder Apostle should represent the Church, and the younger the Synagogue: for although the Synagogue was first to worship God, yet the herd of Gentiles is in the world older than the Synagogue, as witnesseth Paul where he saith: That was not first which is spiritual, but that which is natural. (i Cor. xv. 46.) By Peter, then, who was the elder, is signified the Church of the Gentiles; and by John, who was the younger, the Synagogue of the Jews. They run both of them together, for from the time of her birth until now, and so will it be until the end, the Church of the Gentiles hath run in a parallel road and manywise a common road with the Synagogue, albeit not with equal understandings. The Synagogue came first to the Sepulchre, but she hath not yet entered in; for, though she hath received the commandments of the law, and hath heard the Prophets tell of the Incarnation and Passion of the Lord, she will not believe in Him Who died for her.\par
\switchcolumn*

\LA
\TeDeum{Te Deum}
\switchcolumn
\EN
\TeDeum{Te Deum}
\switchcolumn*

\end{paracol}

\Office{Dominica in Albis in Octava Paschæ}{Low Sunday (Octave of Easter)}{Duplex majus I classis}
\addcontentsline{toc}{subsection}{Dominica in Albis in Octava Paschæ\enspace\textit{Low Sunday (Octave of Easter)}}
\begin{paracol}{2}
\LA
\LessonHead{Lectio I}
\switchcolumn
\EN
\LessonHead{Lesson I}
\switchcolumn*

\LA
\LessonTitle{De Epístola beáti Pauli Apóstoli ad Colossénses}
\switchcolumn
\EN
\LessonTitle{Lesson from the letter of St. Paul the Apostle to the Colossians}
\switchcolumn*

\LA
\Cite{Col 3:1-7}
\switchcolumn
\EN
\Cite{Col 3:1-7}
\switchcolumn*

\LA
\noindent Si consurrexístis cum Christo: quæ sursum sunt quǽrite, ubi Christus est in déxtera Dei sedens: quæ sursum sunt sápite, non quæ super terram. Mórtui enim estis, et vita vestra est abscóndita cum Christo in Deo. Cum Christus apparúerit, vita vestra: tunc et vos apparébitis cum ipso in glória. Mortificáte ergo membra vestra, quæ sunt super terram: fornicatiónem, immundítiam, libídinem, concupiscéntiam malam, et avarítiam, quæ est simulacrórum sérvitus: propter quæ venit ira Dei super fílios incredulitátis: in quibus et vos ambulástis aliquándo, cum viverétis in illis.\par
\switchcolumn
\EN
\noindent Therefore, if you be risen with Christ, seek the things that are above; where Christ is sitting at the right hand of God: Mind the things that are above, not the things that are upon the earth. For you are dead; and your life is hid with Christ in God. When Christ shall appear, who is your life, then you also shall appear with him in glory. Mortify therefore your members which are upon the earth; fornication, uncleanness, lust, evil concupiscence, and covetousness, which is the service of idols. For which things the wrath of God cometh upon the children of unbelief, In which you also walked some time, when you lived in them.\par
\switchcolumn*

\LA
\begin{Resp}\Rsign Angelus Dómini descéndit de cælo, et accédens revólvit lápidem, et super eum sedit, et dixit muliéribus: \Star{} Nolíte timére: scio enim quia crucifíxum quǽritis: jam surréxit: veníte, et vidéte locum, ubi pósitus erat Dóminus, allelúja.\end{Resp}
\begin{Resp}\Vsign Et introëúntes in monuméntum, vidérunt júvenem sedéntem in dextris, coopértum stola cándida, et obstupuérunt: qui dixit illis.\end{Resp}
\begin{Resp}\Rsign Nolíte timére: scio enim quia crucifíxum quǽritis: jam surréxit: veníte, et vidéte locum, ubi pósitus erat Dóminus, allelúja.\end{Resp}
\switchcolumn
\EN
\begin{Resp}\Rsign The Angel of the Lord descended from heaven, and came and rolled back the stone, and sat upon it, and said unto the women: \Star{} Fear not ye; for I know that ye seek Him That was crucified: He is risen already. Come, see the place where the Lord was laid, alleluia.\end{Resp}
\begin{Resp}\Vsign And entering into the sepulchre, they saw a young man sitting on the right side, clothed in a long white garment, and they were affrighted; and he saith unto them:\end{Resp}
\begin{Resp}\Rsign Fear not ye: for I know that ye seek Him That was crucified: He is risen already; come, see the place where the Lord was laid, alleluia.\end{Resp}
\switchcolumn*

\LA
\LessonHead{Lectio II}
\switchcolumn
\EN
\LessonHead{Lesson II}
\switchcolumn*

\LA
\Cite{Col 3:8-13}
\switchcolumn
\EN
\Cite{Col 3:8-13}
\switchcolumn*

\LA
\noindent Nunc autem depónite et vos ómnia: iram, indignatiónem, malítiam, blasphémiam, turpem sermónem de ore vestro. Nolíte mentíri ínvicem, expoliántes vos véterem hóminem cum áctibus suis, et induéntes novum eum, qui renovátur in agnitiónem secúndum imáginem ejus qui creávit illum: ubi non est gentílis et Judǽus, circumcísio et præpútium, Bárbarus et Scytha, servus et liber: sed ómnia, et in ómnibus Christus. Indúite vos ergo, sicut elécti Dei, sancti, et dilécti, víscera misericórdiæ, benignitátem, humilitátem, modéstiam, patiéntiam: supportántes ínvicem, et donántes vobismetípsis si quis advérsus áliquem habet querélam: sicut et Dóminus donávit vobis, ita et vos.\par
\switchcolumn
\EN
\noindent But now put you also all away: anger, indignation, malice, blasphemy, filthy speech out of your mouth. Lie not one to another: stripping yourselves of the old man with his deeds, And putting on the new, him who is renewed unto knowledge, according to the image of him that created him. Where there is neither Gentile nor Jew, circumcision nor uncircumcision, Barbarian nor Scythian, bond nor free. But Christ is all, and in all. Put ye on therefore, as the elect of God, holy, and beloved, the bowels of mercy, benignity, humility, modesty, patience: Bearing with one another, and forgiving one another, if any have a complaint against another: even as the Lord hath forgiven you, so do you also\par
\switchcolumn*

\LA
\begin{Resp}\Rsign Angelus Dómini locútus est muliéribus, dicens: Quem quǽritis? an Jesum quǽritis? Jam surréxit: \Star{} Veníte, et vidéte, allelúja, allelúja.\end{Resp}
\begin{Resp}\Vsign Jesum quǽritis Nazarénum crucifíxum? Surréxit, non est hic.\end{Resp}
\begin{Resp}\Rsign Veníte, et vidéte, allelúja, allelúja.\end{Resp}
\switchcolumn
\EN
\begin{Resp}\Rsign The Angel of the Lord spake unto the woman, saying: Whom seek ye? Seek ye Jesus? He is risen now: \Star{} Come and see, alleluia, alleluia.\end{Resp}
\begin{Resp}\Vsign Seek ye Jesus of Nazareth, Which was crucified? He is risen, He is not here.\end{Resp}
\begin{Resp}\Rsign Come and see, alleluia, alleluia.\end{Resp}
\switchcolumn*

\LA
\LessonHead{Lectio III}
\switchcolumn
\EN
\LessonHead{Lesson III}
\switchcolumn*

\LA
\Cite{Col 3:14-17}
\switchcolumn
\EN
\Cite{Col 3:14-17}
\switchcolumn*

\LA
\noindent Super ómnia autem hæc, caritátem habéte, quod est vínculum perfectiónis: et pax Christi exsúltet in córdibus vestris, in qua et vocáti estis in uno córpore: et grati estóte. Verbum Christi hábitet in vobis abundánter, in omni sapiéntia, docéntes, et commonéntes vosmetípsos, psalmis, hymnis, et cánticis spirituálibus, in grátia cantántes in córdibus vestris Deo. Omne, quodcúmque fácitis in verbo aut in ópere, ómnia in nómine Dómini Jesu Christi, grátias agéntes Deo et Patri per ipsum.\par
\switchcolumn
\EN
\noindent But above all these things have charity, which is the bond of perfection: And let the peace of Christ rejoice in your hearts, wherein also you are called in one body: and be ye thankful. Let the word of Christ dwell in you abundantly, in all wisdom: teaching and admonishing one another in psalms, hymns, and spiritual canticles, singing in grace in your hearts to God. All whatsoever you do in word or in work, do all in the name of the Lord Jesus Christ, giving thanks to God and the Father by him\par
\switchcolumn*

\LA
\begin{Resp}\Rsign Cum transísset sábbatum, María Magdaléne, et María Jacóbi, et Salóme emérunt arómata, \Star{} Ut veniéntes úngerent Jesum, allelúja, allelúja.\end{Resp}
\begin{Resp}\Vsign Et valde mane una sabbatórum, véniunt ad monuméntum, orto jam sole.\end{Resp}
\begin{Resp}\Rsign Ut veniéntes úngerent Jesum, allelúja, allelúja.\end{Resp}
\begin{Resp}\Vsign Glória Patri, et Fílio, \Star{} et Spirítui Sancto.\end{Resp}
\begin{Resp}\Rsign Ut veniéntes úngerent Jesum, allelúja, allelúja.\end{Resp}
\switchcolumn
\EN
\begin{Resp}\Rsign When the Sabbath was passed, Mary Magdalene, and Mary the mother of James, and Salome, had bought sweet spices, \Star{} That they might come and anoint Jesus, alleluia, alleluia.\end{Resp}
\begin{Resp}\Vsign And very early in the morning, the first day of the week, they came unto the sepulchre, at the rising of the sun.\end{Resp}
\begin{Resp}\Rsign That they might come and anoint Jesus, alleluia, alleluia.\end{Resp}
\begin{Resp}\Vsign Glory be to the Father, and to the Son, \Star{} and to the Holy Ghost.\end{Resp}
\begin{Resp}\Rsign That they might come and anoint Jesus, alleluia, alleluia.\end{Resp}
\switchcolumn*

\LA
\LessonHead{Lectio IV}
\switchcolumn
\EN
\LessonHead{Lesson IV}
\switchcolumn*

\LA
\LessonTitle{Sermo sancti Augustíni Epíscopi}
\switchcolumn
\EN
\LessonTitle{From the Sermons of St. Augustine, Bishop (of Hippo.)}
\switchcolumn*

\LA
\Source{Sermo 1 in Octava Paschæ, qui est 157 de Tempore}
\switchcolumn
\EN
\Source{1st Sermon for the Octave of the Passover, being the 157th for the Seasons.}
\switchcolumn*

\LA
\noindent Paschális solémnitas hodiérna festivitáte conclúditur, et ídeo hódie Neophytórum hábitus commutátur: ita tamen, ut candor, qui de hábitu depónitur, semper in corde teneátur. In qua quidem primum nobis agéndum est, ut quia pascháles dies sunt, id est, indulgéntiæ ac remissiónis, ita a nobis sanctórum diérum festívitas agátur, ut relaxatióne córporum púritas non obfuscétur: sed pótius abstinéntes ab omni luxu, ebrietáte, lascívia, demus óperam sóbriæ remissióni, ac sanctæ sinceritáti: ut, quidquid modo corporáli abstinéntia non acquírimus, méntium puritáte quærámus.\par
\switchcolumn
\EN
\noindent The Feast of this day is the end of the Paschal solemnity, and therefore it is today that the Newly-Baptized put off their white garments: but, though they lay aside the outward mark of washing in their raiment, the mark of that washing in their souls remaineth to eternity. Now are the days of the Pass-over, that is, of God's Passing-over our iniquity by His pardon and remission; and therefore our first duty is so to sanctify the mirth of these holy days, that our bodily recreation may be taken without defilement to our spiritual cleanness. Let us strive that our relaxation may be sober and our freedom holy, holding ourselves carefully aloof from anything like excess, drunkenness or lechery. Let us try so to keep in our souls their Lenten cleansing, that if our Fasting hath left us aught yet unwon, we may still be able to seek it.\par
\switchcolumn*

\LA
\begin{Resp}\Rsign María Magdaléne, et áltera María ibant dilúculo ad monuméntum: \Star{} Jesum quem quǽritis, non est hic, surréxit sicut locútus est, præcédet vos in Galilǽam, ibi eum vidébitis, allelúja, allelúja.\end{Resp}
\begin{Resp}\Vsign Et valde mane una sabbatórum véniunt ad monuméntum, orto jam sole: et introëúntes vidérunt júvenem sedéntem in dextris, qui dixit illis.\end{Resp}
\begin{Resp}\Rsign Jesum quem quǽritis, non est hic, surréxit sicut locútus est, præcédet vos in Galilǽam, ibi eum vidébitis, allelúja, allelúja.\end{Resp}
\switchcolumn
\EN
\begin{Resp}\Rsign Mary Magdalene and the other Mary went very early to the sepulchre. \Star{} That Jesus whom ye seek, is not here: for He is risen, as He said: He goeth before you into Galilee; there shall ye see Him, alleluia, alleluia.\end{Resp}
\begin{Resp}\Vsign And very early in the morning, the first day of the week, they came unto the sepulchre, at the rising of the sun; and, entering into the sepulchre, they saw a young man sitting upon the right side, who saith unto them:\end{Resp}
\begin{Resp}\Rsign That Jesus Whom ye seek is not here: for He is risen, as He said: He goeth before you into Galilee: there shall ye see Him, alleluia, alleluia.\end{Resp}
\switchcolumn*

\LA
\LessonHead{Lectio V}
\switchcolumn
\EN
\LessonHead{Lesson V}
\switchcolumn*

\LA
\noindent Ad omnes quidem pértinet sermo, quos cura nostra compléctitur: verúmtamen hódie termináta sacramentórum solemnitáte, vos allóquimur, novélla gérmina sanctitátis, regeneráta ex aqua et Spíritu Sancto: germen pium, exámen novéllum, flos nostri honóris, et fructus labóris, gáudium et coróna mea, omnes qui statis in Dómino. Apostólicis verbis vos álloquor: Ecce nox præcéssit, dies autem appropinquávit: abjícite ópera tenebrárum, et indúite vos arma lucis. Sicut in die honéste ambulémus: non in comessatiónibus et ebrietátibus, non in cubílibus et impudicítiis, non in contentióne et æmulatióne: sed induímini Dóminum Jesum Christum.\par
\switchcolumn
\EN
\noindent At discourse concerneth all them which are committed unto my spiritual charge; but, nevertheless, since the first happy week of your Sacramental life draweth this day to a close, I address myself in especial to you who are the new olive-plants of holiness round about the Table of the Lord, (Ps. cxxvii. 4,) to you, who have but a little while been born again of water and the Holy Ghost, (John iii. 5,) to you, O holy generation (1 Pet. ii. 9) to you, O new creation, (Gal. vi. 15,) to you, the excellency of my dignity, (Gen. xlix. 3,) and the fruit of my labour, my brethren dearly beloved and longed for, my joy and my crown, all ye who now stand so fast in the Lord. (Phil. iv. i.) To you I address the words of the Apostle (Rom. xiii. 12.) Behold! the night is past! the day is come! Cast off therefore the works of darkness, and put on the armour of light. Let us walk honestly, as in the day; not in rioting and drunkenness, not in chambering and wantonness, not in strife and envying: but put ye on the Lord Jesus Christ.\par
\switchcolumn*

\LA
\begin{Resp}\Rsign Surréxit pastor bonus, qui ánimam suam pósuit pro óvibus suis, et pro grege suo mori dignátus est: \Star{} Allelúja, allelúja, allelúja.\end{Resp}
\begin{Resp}\Vsign Etenim Pascha nostrum immolátus est Christus.\end{Resp}
\begin{Resp}\Rsign Allelúja, allelúja, allelúja.\end{Resp}
\switchcolumn
\EN
\begin{Resp}\Rsign The Good Shepherd, Who laid down His life for the sheep, yea, Who was contented even to die for His flock, the Good Shepherd is risen again. \Star{} Alleluia, alleluia, alleluia.\end{Resp}
\begin{Resp}\Vsign For even Christ our Passover is sacrificed for us.\end{Resp}
\begin{Resp}\Rsign Alleluia, alleluia, alleluia.\end{Resp}
\switchcolumn*

\LA
\LessonHead{Lectio VI}
\switchcolumn
\EN
\LessonHead{Lesson VI}
\switchcolumn*

\LA
\noindent Habémus, inquit, certiórem prophéticum sermónem: cui bene fácitis intendéntes tamquam lucérnæ in obscúro loco, donec dies lucéscat, et lúcifer oriátur in córdibus vestris. Sint ergo lumbi vestri accíncti, et lucérnæ ardéntes in mánibus vestris: et vos símiles homínibus exspectántibus dóminum suum, quando revertátur a núptiis. Ecce dies advéniunt, in quibus Dóminus dicit: Pusíllum, inquit, et non vidébitis me, et íterum pusíllum, et vidébitis me. Hæc est hora, de qua dixit: Vos tristes éritis, sǽculum autem gaudébit: id est, vita ista tentatiónibus plena, in qua peregrinámur ab eo. Sed íterum, inquit, vidébo vos, et gaudébit cor vestrum, et gáudium vestrum nemo tollet a vobis.\par
\switchcolumn
\EN
\noindent We have, saith Peter, (2 Peter i. 19,) a more sure word of Prophecy, whereunto ye do well that ye take heed, as unto a light that shineth in a dark place, until the day dawn, and the day-star arise in your hearts. Let your loins therefore be girded about, and your lights burning in your hands, and ye yourselves like unto men that wait for their lord, when he will return from the wedding. (Luke xii. 36.) Behold, the days come, whereof the Lord saith, (John xvi. 16, 17, 19,) A little while, and ye shall not see Me, and again a little while and ye shall see Me. Now is the hour whereof He said (20), Ye shall weep and lament, but the work shall rejoice that is to say, this present life, wherein we walk as strangers and pilgrims, (1 Pet. ii. 11,) far away from Him Who is our Home, this present life is very full of trials. But, saith Jesus, but I will see you again, and your heart shall rejoice, and your joy no man taketh from you. (22.)\par
\switchcolumn*

\LA
\begin{Resp}\Rsign Virtúte magna reddébant Apóstoli, \Star{} Testimónium resurrectiónis Jesu Christi Dómini nostri, allelúja, allelúja.\end{Resp}
\begin{Resp}\Vsign Repléti quidem Spíritu Sancto loquebántur cum fidúcia verbum Dei.\end{Resp}
\begin{Resp}\Rsign Testimónium resurrectiónis Jesu Christi Dómini nostri, allelúja, allelúja.\end{Resp}
\begin{Resp}\Vsign Glória Patri, et Fílio, \Star{} et Spirítui Sancto.\end{Resp}
\begin{Resp}\Rsign Testimónium resurrectiónis Jesu Christi Dómini nostri, allelúja, allelúja.\end{Resp}
\switchcolumn
\EN
\begin{Resp}\Rsign With great power gave the Apostles \Star{} Witness of the Resurrection of our Lord Jesus Christ, alleluia, alleluia.\end{Resp}
\begin{Resp}\Vsign They were all filled with the Holy Ghost, and they spoke the Word of God with boldness.\end{Resp}
\begin{Resp}\Rsign Witness of the Resurrection of our Lord Jesus Christ, alleluia, alleluia.\end{Resp}
\begin{Resp}\Vsign Glory be to the Father, and to the Son, \Star{} and to the Holy Ghost.\end{Resp}
\begin{Resp}\Rsign Witness of the Resurrection of our Lord Jesus Christ, alleluia, alleluia.\end{Resp}
\switchcolumn*

\LA
\LessonHead{Lectio VII}
\switchcolumn
\EN
\LessonHead{Lesson VII}
\switchcolumn*

\LA
\LessonTitle{Léctio sancti Evangélii secúndum Joánnem}
\switchcolumn
\EN
\LessonTitle{From the Holy Gospel according to John}
\switchcolumn*

\LA
\Cite{Joannes 20:19-31}
\switchcolumn
\EN
\Cite{John 20:19-31}
\switchcolumn*

\LA
\noindent In illo témpore: Cum sero esset die illo, una sabbatórum, et fores essent clausæ, ubi erant discípuli congregáti propter metum Judæórum: venit Jesus, et stetit in médio, et dixit eis: Pax vobis. Et réliqua.\par
\switchcolumn
\EN
\noindent At that time: Now when it was late that same day, the first of the week, and the doors were shut, where the disciples were gathered together, for fear of the Jews, Jesus came and stood in the midst, and said to them: Peace be to you. And so forth.\par
\switchcolumn*

\LA
\LessonTitle{Homilía sancti Gregórii Papæ}
\switchcolumn
\EN
\LessonTitle{Homily by Pope St. Gregory (the Great.)}
\switchcolumn*

\LA
\Source{Homilia 26 in Evangelia}
\switchcolumn
\EN
\Source{26th on the Gospels}
\switchcolumn*

\LA
\noindent Prima lectiónis hujus evangélicæ quǽstio ánimum pulsat: quómodo post resurrectiónem corpus Domínicum verum fuit, quod clausis jánuis ad discípulos íngredi pótuit? Sed sciéndum nobis est, quod divína operátio, si ratióne comprehénditur, non est admirábilis: nec fides habet méritum, cui humána rátio præbet experiméntum. Sed hæc ipsa nostri Redemptóris ópera, quæ ex semetípsis comprehéndi nequáquam possunt, ex ália ejus operatióne pensánda sunt: ut rebus mirabílibus fidem prǽbeant facta mirabilióra. Illud enim corpus Dómini intrávit ad discípulos jánuis clausis, quod vidélicet ad humános óculos per nativitátem suam clauso exívit útero Vírginis. Quid ergo mirum, si clausis jánuis post resurrectiónem suam in ætérnum jam victúrus intrávit, qui moritúrus véniens, non apérto útero Vírginis exívit?\par
\switchcolumn
\EN
\noindent When we hear this passage of the Gospel read, a question straightway knocketh at the door of our mind. How was it that the Body of the Risen Lord was a real Body, if It was able to pass through closed doors into the assembly of His disciples? But we ought to know that the works of God are no more wonderful when they can be understood by man's reason, and faith has lost her worth when her subject-matter is the subject-matter of human demonstration. Nevertheless, those very works of our Redeemer which are in themselves impossible to be understood, must be thought over in connection with other of His works, that we may be led to believe in things wonderful, by mean of things more wonderful still. That Body of the Lord, Which came into the assembly of the disciples through closed doors, was the Same, Which at Its birth, had become manifest to the eyes of men by passing out of the cloister of the Virgin's womb without breaking the seal thereof. What wonder is it if that Body Which had come out of the Virgin's womb, without opening the matrix, albeit It was then on Its way to die, now that It was risen again from the dead and instinct for ever with undying life, what wonder is it, I say, if that Body passed through closed doors?\par
\switchcolumn*

\LA
\begin{Resp}\Rsign De ore prudéntis procédit mel, allelúja: dulcédo mellis est sub língua ejus, allelúja: \Star{} Favus distíllans lábia ejus, allelúja, allelúja.\end{Resp}
\begin{Resp}\Vsign Sapiéntia requiéscit in corde ejus, et prudéntia in sermóne oris illíus.\end{Resp}
\begin{Resp}\Rsign Favus distíllans lábia ejus, allelúja, allelúja.\end{Resp}
\switchcolumn
\EN
\begin{Resp}\Rsign From the mouth of the wise doth proceed honey, alleluia: the sweetness of honey is under his tongue, alleluia. \Star{} His lips drop as the honey-comb, alleluia, alleluia.\end{Resp}
\begin{Resp}\Vsign Wisdom doth abide in his heart, and out of his mouth cometh understanding.\end{Resp}
\begin{Resp}\Rsign His lips drop as the honey-comb, alleluia, alleluia.\end{Resp}
\switchcolumn*

\LA
\LessonHead{Lectio VIII}
\switchcolumn
\EN
\LessonHead{Lesson VIII}
\switchcolumn*

\LA
\noindent Sed quia ad illud corpus, quod vidéri póterat, fides intuéntium dubitábat: osténdit eis prótinus manus et latus, palpándam carnem prǽbuit, quam clausis jánuis introdúxit. Qua in re duo mira, et juxta humánam ratiónem sibi valde contrária osténdit: dum post resurrectiónem suam corpus suum incorruptíbile, et tamen palpábile demonstrávit. Nam et corrúmpi necésse est quod palpátur: et palpári non potest quod non corrúmpitur. Sed miro modo atque inæstimábili Redémptor noster et incorruptíbile post resurrectiónem, et palpábile corpus exhíbuit: ut monstrándo incorruptíbile, invitáret ad prǽmium, et præbéndo palpábile, firmáret ad fidem. Et incorruptíbilem se ergo, et palpábilem demonstrávit: ut profécto esse post resurrectiónem osténderet corpus suum et ejúsdem natúræ, et altérius glóriæ.\par
\switchcolumn
\EN
\noindent But since the beholders doubted of the reality of that Body Which they saw, He showed unto them His Hands and His Side, and allowed them to handle that Same Flesh Which had just passed through the closed doors. (Luke xxiv. 39.) In this there were two strange things manifested; yea, things which, according to our understanding, are contrary the one to the other. His Risen Body was incorruptible and yet palpable. For whatever can be touched, must needs be subject to corruption; and whatever is not subject to corruption, cannot be touched. But, in a way altogether wonderful and incomprehensible, our Redeemer after His Resurrection revealed Himself in a Body at once palpable and incorruptible: revealed Himself in an incorruptible Body, that we might learn to seek a like glorification; and in a palpable Body, for the strengthening of our faith. He revealed Himself in a Body at once incorruptible and palpable, that He might thereby make manifest the fact that His Risen Body was unaltered in nature, albeit transfigured in glory.\par
\switchcolumn*

\LA
\begin{Resp}\Rsign Surgens Jesus Dóminus noster, stans in médio discipulórum suórum, dixit: \Star{} Pax vobis, allelúja: gavísi sunt discípuli viso Dómino, allelúja.\end{Resp}
\begin{Resp}\Vsign Una ergo sabbatórum, cum fores essent clausæ, ubi erant discípuli congregáti, venit Jesus, et stetit in medio eórum, et dixit eis.\end{Resp}
\begin{Resp}\Rsign Pax vobis, allelúja: gavísi sunt discípuli viso Dómino, allelúja.\end{Resp}
\begin{Resp}\Vsign Glória Patri, et Fílio, \Star{} et Spirítui Sancto.\end{Resp}
\begin{Resp}\Rsign Pax vobis, allelúja: gavísi sunt discípuli viso Dómino, allelúja.\end{Resp}
\switchcolumn
\EN
\begin{Resp}\Rsign After that our Lord Jesus was risen again, He came and stood in the midst of His disciples, and said unto them: \Star{} Peace be unto you, alleluia. Then were the disciples glad, when they saw the Lord, alleluia.\end{Resp}
\begin{Resp}\Vsign The first day of the week, when the doors were shut where the disciples were assembled, came Jesus, and stood in the midst, and said unto them:\end{Resp}
\begin{Resp}\Rsign Peace be unto you, alleluia. Then were the disciples glad, when they saw the Lord, alleluia.\end{Resp}
\begin{Resp}\Vsign Glory be to the Father, and to the Son, \Star{} and to the Holy Ghost.\end{Resp}
\begin{Resp}\Rsign Peace be unto you, alleluia. Then were the disciples glad, when they saw the Lord, alleluia.\end{Resp}
\switchcolumn*

\LA
\LessonHead{Lectio IX}
\switchcolumn
\EN
\LessonHead{Lesson IX}
\switchcolumn*

\LA
\noindent Dixit eis: Pax vobis. Sicut misit me Pater, et ego mitto vos: id est, sicut misit me Pater Deus Deum, et ego mitto vos homo hómines. Pater Fílium misit, qui hunc pro redemptióne géneris humáni incarnári constítuit. Quem vidélicet in mundum veníre ad passiónem vóluit: sed tamen amávit Fílium, quem ad passiónem misit. Eléctos vero Apóstolos Dóminus non ad mundi gáudia, sed sicut ipse missus est, ad passiónes in mundum mittit. Quia ergo et Fílius amátur a Patre, et tamen ad passiónem míttitur: ita et discípuli a Dómino amántur, qui tamen ad passiónem mittúntur in mundum. Itaque recte dícitur: Sicut misit me Pater, et ego mitto vos: id est, ea vos caritáte díligo, cum inter scándala persecutórum mitto, qua me caritáte Pater díligit, quem veníre ad tolerándas passiónes fecit.\par
\switchcolumn
\EN
\noindent Then said Jesus to them again: Peace be unto you. As My Father hath sent Me, even so send I you that is, as My Father, Who is God, hath sent Me, Who am God, even so do I, Who am Man, send you, who are men. The Father sent the Son, Whom He appointed to be made Man for the redemption of man. Him He willed to send into the world to suffer, albeit He Whom He sent to suffer was the Son of His love. The Lord sendeth His chosen Apostles into the world, not to be happy in the world, but, as He had been Himself sent, to suffer. As the Father loveth the Son and yet sendeth Him to suffer, even so doth the Lord love His disciples, albeit He sendeth them into the world, to suffer therein; and therefore it is well said: “As My Father hath sent Me, even so send I you”; that is, while I send you into the wild storm of persecution, I love you all the same, I love you, yea, I love you with a love like that wherewith the Father loveth Me, Who sent Me into the world to bear agony therein.\par
\switchcolumn*

\LA
\TeDeum{Te Deum}
\switchcolumn
\EN
\TeDeum{Te Deum}
\switchcolumn*

\end{paracol}

\Office{Feria Secunda infra Hebdomadam I post Octavam Paschæ}{Monday of the First Week after the Octave of Easter}{Feria}
\addcontentsline{toc}{subsection}{Feria Secunda infra Hebdomadam I post Octavam Paschæ\enspace\textit{Monday of the First Week after the Octave of Easter}}
\begin{paracol}{2}
\LA
\LessonHead{Lectio I}
\switchcolumn
\EN
\LessonHead{Lesson I}
\switchcolumn*

\LA
\LessonTitle{Incipit liber Actuum Apostolórum}
\switchcolumn
\EN
\LessonTitle{Lesson from the Acts of the Apostles}
\switchcolumn*

\LA
\Cite{Act. 1:1-8}
\switchcolumn
\EN
\Cite{Acts 1:1-8}
\switchcolumn*

\LA
\noindent Primum quidem sermónem feci de ómnibus, o Theóphile, quæ cœpit Jesus fácere et docére usque in diem qua præcípiens Apóstolis per Spíritum Sanctum, quos elégit, assúmptus est: quibus et prǽbuit seípsum vivum post passiónem suam in multis arguméntis, per dies quadragínta appárens eis, et loquens de regno Dei. Et convéscens, præcépit eis ab Jerosólymis ne discéderent, sed exspectárent promissiónem Patris, quam audístis (inquit) per os meum: quia Joánnes quidem baptizávit aqua, vos autem baptizabímini Spíritu Sancto non post multos hos dies. Igitur qui convénerant, interrogábant eum, dicéntes: Dómine, si in témpore hoc restítues regnum Israël? Dixit autem eis: Non est vestrum nosse témpora vel moménta quæ Pater pósuit in sua potestáte: sed accipiétis virtútem superveniéntis Spíritus Sancti in vos, et éritis mihi testes in Jerúsalem, et in omni Judǽa, et Samaría, et usque ad últimum terræ.\par
\switchcolumn
\EN
\noindent The former treatise I made, O Theophilus, of all things which Jesus began to do and to teach, Until the day on which, giving commandments by the Holy Ghost to the apostles whom he had chosen, he was taken up. To whom also he showed himself alive after his passion, by many proofs, for forty days appearing to them, and speaking of the kingdom of God. And eating together with them, he commanded them, that they should not depart from Jerusalem, but should wait for the promise of the Father, which you have heard (saith he) by my mouth. For John indeed baptized with water, but you shall be baptized with the Holy Ghost, not many days hence. They therefore who were come together, asked him, saying: Lord, wilt thou at this time restore again the kingdom to Israel? But he said to them: It is not for you to know the times or moments, which the Father hath put in his own power: But you shall receive the power of the Holy Ghost coming upon you, and you shall be witnesses unto me in Jerusalem, and in all Judea, and Samaria, and even to the uttermost part of the earth.\par
\switchcolumn*

\LA
\begin{Resp}\Rsign Virtúte magna reddébant Apóstoli, \Star{} Testimónium resurrectiónis Jesu Christi Dómini nostri, allelúja, allelúja.\end{Resp}
\begin{Resp}\Vsign Repléti quidem Spíritu Sancto loquebántur cum fidúcia verbum Dei.\end{Resp}
\begin{Resp}\Rsign Testimónium resurrectiónis Jesu Christi Dómini nostri, allelúja, allelúja.\end{Resp}
\switchcolumn
\EN
\begin{Resp}\Rsign With great power gave the Apostles \Star{} Witness of the Resurrection of our Lord Jesus Christ, alleluia, alleluia.\end{Resp}
\begin{Resp}\Vsign They were all filled with the Holy Ghost, and they spoke the Word of God with boldness.\end{Resp}
\begin{Resp}\Rsign Witness of the Resurrection of our Lord Jesus Christ, alleluia, alleluia.\end{Resp}
\switchcolumn*

\LA
\LessonHead{Lectio II}
\switchcolumn
\EN
\LessonHead{Lesson II}
\switchcolumn*

\LA
\Cite{Act. 1:9-14}
\switchcolumn
\EN
\Cite{Acts 1:9-14}
\switchcolumn*

\LA
\noindent Et cum hæc dixísset, vidéntibus illis, elevátus est: et nubes suscépit eum ab óculis eórum. Cumque intueréntur in cælum eúntem illum, ecce duo viri astitérunt juxta illos in véstibus albis, qui et dixérunt: Viri Galilǽi, quid statis aspiciéntes in cælum? Hic Jesus, qui assúmptus est a vobis in cælum, sic véniet quemádmodum vidístis eum eúntem in cælum. Tunc revérsi sunt Jerosólymam a monte qui vocátur Olivéti, qui est juxta Jerúsalem, sábbati habens iter. Et cum introíssent in cœnáculum, ascendérunt ubi manébant Petrus, et Joánnes, Jacóbus, et Andréas, Philíppus, et Thomas, Bartholomǽus, et Matthǽus, Jacóbus Alphǽi, et Simon Zelótes, et Judas Jacóbi. Hi omnes erant perseverántes unanímiter in oratióne cum muliéribus, et María matre Jesu, et frátribus ejus.\par
\switchcolumn
\EN
\noindent And when he had said these things, while they looked on, he was raised up: and a cloud received him out of their sight. And while they were beholding him going up to heaven, behold two men stood by them in white garments. Who also said: Ye men of Galilee, why stand you looking up to heaven? This Jesus who is taken up from you into heaven, shall so come, as you have seen him going into heaven. Then they returned to Jerusalem from the mount that is called Olivet, which is nigh Jerusalem, within a sabbath day's journey. And when they were come in, they went up into an upper room, where abode Peter and John, James and Andrew, Philip and Thomas, Bartholomew and Matthew, James of Alpheus, and Simon Zelotes, and Jude the brother of James. All these were persevering with one mind in prayer with the women, and Mary the mother of Jesus, and with his brethren\par
\switchcolumn*

\LA
\begin{Resp}\Rsign De ore prudéntis procédit mel, allelúja: dulcédo mellis est sub língua ejus, allelúja: \Star{} Favus distíllans lábia ejus, allelúja, allelúja.\end{Resp}
\begin{Resp}\Vsign Sapiéntia requiéscit in corde ejus, et prudéntia in sermóne oris illíus.\end{Resp}
\begin{Resp}\Rsign Favus distíllans lábia ejus, allelúja, allelúja.\end{Resp}
\begin{Resp}\Vsign Glória Patri, et Fílio, \Star{} et Spirítui Sancto.\end{Resp}
\begin{Resp}\Rsign Favus distíllans lábia ejus, allelúja, allelúja.\end{Resp}
\switchcolumn
\EN
\begin{Resp}\Rsign From the mouth of the wise doth proceed honey, alleluia: the sweetness of honey is under his tongue, alleluia. \Star{} His lips drop as the honey-comb, alleluia, alleluia.\end{Resp}
\begin{Resp}\Vsign Wisdom doth abide in his heart, and out of his mouth cometh understanding.\end{Resp}
\begin{Resp}\Rsign His lips drop as the honey-comb, alleluia, alleluia.\end{Resp}
\begin{Resp}\Vsign Glory be to the Father, and to the Son, \Star{} and to the Holy Ghost.\end{Resp}
\begin{Resp}\Rsign His lips drop as the honey-comb, alleluia, alleluia.\end{Resp}
\switchcolumn*

\LA
\LessonHead{Lectio III}
\switchcolumn
\EN
\LessonHead{Lesson III}
\switchcolumn*

\LA
\Cite{Act. 1:15-26}
\switchcolumn
\EN
\Cite{Acts 1:15-26}
\switchcolumn*

\LA
\noindent In diébus illis, exsúrgens Petrus in médio fratrum, dixit (erat autem turba hóminum simul, fere centum vigínti): Viri fratres, opórtet impléri Scriptúram quam prædíxit Spíritus Sanctus per os David de Juda, qui fuit dux eórum qui comprehendérunt Jesum: qui connumerátus erat in nobis, et sortítus est sortem ministérii hujus. Et hic quidem possédit agrum de mercéde iniquitátis, et suspénsus crépuit médius: et diffúsa sunt ómnia víscera ejus. Et notum factum est ómnibus habitántibus Jerúsalem, ita ut appellarétur ager ille, lingua eórum, Hacéldama, hoc est, ager sánguinis. Scriptum est enim in libro Psalmórum: Fiat commorátio eórum desérta, et non sit qui inhábitet in ea: et episcopátum ejus accípiat alter. Opórtet ergo ex his viris qui nobíscum sunt congregáti in omni témpore quo intrávit et exívit inter nos Dóminus Jesus, incípiens a baptísmate Joánnis usque in diem qua assúmptus est a nobis, testem resurrectiónis ejus nobíscum fíeri unum ex istis. Et statuérunt duos, Joseph, qui vocabátur Bársabas, qui cognominátus est Justus, et Matthíam. Et orántes dixérunt: Tu Dómine, qui corda nosti ómnium, osténde quem elégeris ex his duóbus unum, accípere locum ministérii hujus et apostolátus, de quo prævaricátus est Judas ut abíret in locum suum. Et dedérunt sortes eis, et cécidit sors super Matthíam: et annumerátus est cum úndecim Apóstolis.\par
\switchcolumn
\EN
\noindent In those days Peter rising up in the midst of the brethren, said: (now the number of persons together was about an hundred and twenty:) Men, brethren, the scripture must needs be fulfilled, which the Holy Ghost spoke before by the mouth of David concerning Judas, who was the leader of them that apprehended Jesus: Who was numbered with us, and had obtained part of this ministry. And he indeed hath possessed a field of the reward of iniquity, and being hanged, burst asunder in the midst: and all his bowels gushed out. And it became known to all the inhabitants of Jerusalem: so that the same field was called in their tongue, Haceldama, that is to say, The field of blood. For it is written in the book of Psalms: Let their habitation become desolate, and let there be none to dwell therein. And his bishopric let another take. Wherefore of these men who have companied with us all the time that the Lord Jesus came in and went out among us, Beginning from the baptism of John, until the day wherein he was taken up from us, one of these must be made a witness with us of his resurrection. And they appointed two, Joseph, called Barsabas, who was surnamed Justus, and Matthias. And praying, they said: Thou, Lord, who knowest the hearts of all men, show whether of these two thou hast chosen, To take the place of this ministry and apostleship, from which Judas hath by transgression fallen, that he might go to his own place. And they gave them lots, and the lot fell upon Matthias, and he was numbered with the eleven apostles.\par
\switchcolumn*

\LA
\TeDeum{Te Deum}
\switchcolumn
\EN
\TeDeum{Te Deum}
\switchcolumn*

\end{paracol}

\Office{Feria Tertia infra Hebdomadam I post Octavam Paschæ}{Tuesday of the First Week after the Octave of Easter}{Feria}
\addcontentsline{toc}{subsection}{Feria Tertia infra Hebdomadam I post Octavam Paschæ\enspace\textit{Tuesday of the First Week after the Octave of Easter}}
\begin{paracol}{2}
\LA
\LessonHead{Lectio I}
\switchcolumn
\EN
\LessonHead{Lesson I}
\switchcolumn*

\LA
\LessonTitle{De Actibus Apostolórum}
\switchcolumn
\EN
\LessonTitle{Lesson from the Acts of the Apostles}
\switchcolumn*

\LA
\Cite{Act. 2:1-8}
\switchcolumn
\EN
\Cite{Acts 2:1-8}
\switchcolumn*

\LA
\noindent Et cum compleréntur dies Pentecóstes, erant omnes páriter in eódem loco: et factus est repénte de cælo sonus, tamquam adveniéntis spíritus veheméntis, et replévit totam domum ubi erant sedéntes. Et apparuérunt illis dispertítæ linguæ tamquam ignis, sedítque supra síngulos eórum: et repléti sunt omnes Spíritu Sancto, et cœpérunt loqui váriis linguis, prout Spíritus Sanctus dabat éloqui illis. Erant autem in Jerúsalem habitántes Judǽi, viri religiósi ex omni natióne quæ sub cælo est. Facta autem hac voce, convénit multitúdo, et mente confúsa est, quóniam audiébat unusquísque lingua sua illos loquéntes. Stupébant autem omnes, et mirabántur, dicéntes: Nonne ecce omnes isti qui loquúntur, Galilǽi sunt? et quómodo nos audívimus unusquísque linguam nostram in qua nati sumus?\par
\switchcolumn
\EN
\noindent And when the days of the Pentecost were accomplished, they were all together in one place: And suddenly there came a sound from heaven, as of a mighty wind coming, and it filled the whole house where they were sitting. And there appeared to them parted tongues as it were of fire, and it sat upon every one of them: And they were all filled with the Holy Ghost, and they began to speak with diverse tongues, according as the Holy Ghost gave them to speak. Now there were dwelling at Jerusalem, Jews, devout men, out of every nation under heaven. And when this was noised abroad, the multitude came together, and were confounded in mind, because that every man heard them speak in his own tongue. And they were all amazed, and wondered, saying: Behold, are not all these, that speak, Galileans? And how have we heard, every man our own tongue wherein we were born?\par
\switchcolumn*

\LA
\begin{Resp}\Rsign Ego sum vitis vera, et vos pálmites: \Star{} Qui manet in me, et ego in eo, hic fert fructum multum, allelúja, allelúja.\end{Resp}
\begin{Resp}\Vsign Sicut diléxit me Pater, et ego diléxi vos.\end{Resp}
\begin{Resp}\Rsign Qui manet in me, et ego in eo, hic fert fructum multum, allelúja, allelúja.\end{Resp}
\switchcolumn
\EN
\begin{Resp}\Rsign I am the True Vine, and ye are the branches; \Star{} He that abideth in Me, and I in him, the same bringeth forth much fruit, alleluia, alleluia.\end{Resp}
\begin{Resp}\Vsign As the Father hath loved Me, so have I loved you.\end{Resp}
\begin{Resp}\Rsign He that abideth in Me, and I in him, the same bringeth forth much fruit, alleluia, alleluia.\end{Resp}
\switchcolumn*

\LA
\LessonHead{Lectio II}
\switchcolumn
\EN
\LessonHead{Lesson II}
\switchcolumn*

\LA
\Cite{Act. 2:14-21}
\switchcolumn
\EN
\Cite{Acts 2:14-21}
\switchcolumn*

\LA
\noindent Stans autem Petrus cum úndecim, levávit vocem suam, et locútus est eis: Viri Judǽi, et qui habitátis Jerúsalem univérsi, hoc vobis notum sit, et áuribus percípite verba mea. Non enim, sicut vos æstimátis, hi ébrii sunt, cum sit hora diéi tértia: sed hoc est quod dictum est per prophétam Joël: Et erit in novíssimis diébus, dicit Dóminus, effúndam de Spíritu meo super omnem carnem: et prophetábunt fílii vestri et fíliæ vestræ, et júvenes vestri visiónes vidébunt, et senióres vestri sómnia somniábunt. Et quidem super servos meos, et super ancíllas meas, in diébus illis effúndam de Spíritu meo, et prophetábunt: et dabo prodígia in cælo sursum, et signa in terra deórsum, sánguinem, et ignem, et vapórem fumi: sol convertétur in ténebras, et luna in sánguinem, ántequam véniat dies Dómini magnus et maniféstus. Et erit: omnis quicúmque invocáverit nomen Dómini, salvus erit.\par
\switchcolumn
\EN
\noindent But Peter standing up with the eleven, lifted up his voice, and spoke to them: Ye men of Judea, and all you that dwell in Jerusalem, be this known to you, and with your ears receive my words. For these are not drunk, as you suppose, seeing it is but the third hour of the day: But this is that which was spoken of by the prophet Joel: And it shall come to pass, in the last days, (saith the Lord,) I will pour out of my Spirit upon all flesh: and your sons and your daughters shall prophesy, and your young men shall see visions, and your old men shall dream dreams. And upon my servants indeed, and upon my handmaids will I pour out in those days of my spirit, and they shall prophesy. And I will show wonders in the heaven above, and signs on the earth beneath: blood and fire, and vapour of smoke. The sun shall be turned into darkness, and the moon into blood, before the great and manifest day of the Lord come. And it shall come to pass, that whosoever shall call upon the name of the Lord, shall be saved.\par
\switchcolumn*

\LA
\begin{Resp}\Rsign Surgens Jesus Dóminus noster, stans in médio discipulórum suórum, dixit: \Star{} Pax vobis, allelúja: gavísi sunt discípuli viso Dómino, allelúja.\end{Resp}
\begin{Resp}\Vsign Una ergo sabbatórum, cum fores essent clausæ, ubi erant discípuli congregáti, venit Jesus, et stetit in medio eórum, et dixit eis.\end{Resp}
\begin{Resp}\Rsign Pax vobis, allelúja: gavísi sunt discípuli viso Dómino, allelúja.\end{Resp}
\begin{Resp}\Vsign Glória Patri, et Fílio, \Star{} et Spirítui Sancto.\end{Resp}
\begin{Resp}\Rsign Pax vobis, allelúja: gavísi sunt discípuli viso Dómino, allelúja.\end{Resp}
\switchcolumn
\EN
\begin{Resp}\Rsign After that our Lord Jesus was risen again, He came and stood in the midst of His disciples, and said unto them: \Star{} Peace be unto you, alleluia. Then were the disciples glad, when they saw the Lord, alleluia.\end{Resp}
\begin{Resp}\Vsign The first day of the week, when the doors were shut where the disciples were assembled, came Jesus, and stood in the midst, and said unto them:\end{Resp}
\begin{Resp}\Rsign Peace be unto you, alleluia. Then were the disciples glad, when they saw the Lord, alleluia.\end{Resp}
\begin{Resp}\Vsign Glory be to the Father, and to the Son, \Star{} and to the Holy Ghost.\end{Resp}
\begin{Resp}\Rsign Peace be unto you, alleluia. Then were the disciples glad, when they saw the Lord, alleluia.\end{Resp}
\switchcolumn*

\LA
\LessonHead{Lectio III}
\switchcolumn
\EN
\LessonHead{Lesson III}
\switchcolumn*

\LA
\Cite{Act. 2:22-27}
\switchcolumn
\EN
\Cite{Acts 2:22-27}
\switchcolumn*

\LA
\noindent Viri Israëlítæ, audíte verba hæc: Jesum Nazarénum, virum approbátum a Deo in vobis, virtútibus, et prodígiis, et signis, quæ fecit Deus per illum in médio vestri, sicut et vos scitis: hunc, definíto consílio et præsciéntia Dei tráditum, per manus iniquórum affligéntes interemístis: quem Deus suscitávit, solútis dolóribus inférni, juxta quod impossíbile erat tenéri illum ab eo. David enim dicit in eum: Providébam Dóminum in conspéctu meo semper: quóniam a dextris est mihi, ne commóvear: propter hoc lætátum est cor meum, et exsultávit lingua mea, ínsuper et caro mea requiéscet in spe: quóniam non derelínques ánimam meam in inférno, nec dabis sanctum tuum vidére corruptiónem.\par
\switchcolumn
\EN
\noindent Ye men of Israel, hear these words: Jesus of Nazareth, a man approved of God among you, by miracles, and wonders, and signs, which God did by him, in the midst of you, as you also know: This same being delivered up, by the determinate counsel and foreknowledge of God, you by the hands of wicked men have crucified and slain. Whom God hath raised up, having loosed the sorrows of hell, as it was impossible that he should be holden by it. For David saith concerning him: I foresaw the Lord before my face: because he is at my right hand, that I may not be moved. For this my heart hath been glad, and any tongue hath rejoiced: moreover my flesh also shall rest in hope. Because thou wilt not leave my soul in hell, nor suffer thy Holy One to see corruption.\par
\switchcolumn*

\LA
\TeDeum{Te Deum}
\switchcolumn
\EN
\TeDeum{Te Deum}
\switchcolumn*

\end{paracol}

\Office{Feria Quarta infra Hebdomadam I post Octavam Paschæ}{Wednesday of the First Week after the Octave of Easter}{Feria}
\addcontentsline{toc}{subsection}{Feria Quarta infra Hebdomadam I post Octavam Paschæ\enspace\textit{Wednesday of the First Week after the Octave of Easter}}
\begin{paracol}{2}
\LA
\LessonHead{Lectio I}
\switchcolumn
\EN
\LessonHead{Lesson I}
\switchcolumn*

\LA
\LessonTitle{De Actibus Apostolórum}
\switchcolumn
\EN
\LessonTitle{Lesson from the Acts of the Apostles}
\switchcolumn*

\LA
\Cite{Act. 3:1-6}
\switchcolumn
\EN
\Cite{Acts 3:1-6}
\switchcolumn*

\LA
\noindent Petrus autem et Joánnes ascendébant in templum ad horam oratiónis nonam. Et quidam vir, qui erat claudus ex útero matris suæ, bajulabátur; quem ponébant cotídie ad portam templi, quæ dícitur Speciósa, ut péteret eleemósynam ab introëúntibus in templum. Is, cum vidísset Petrum et Joánnem incipiéntes introíre in templum, rogábat ut eleemósynam accíperet. Intuens autem in eum Petrus cum Joánne, dixit: Réspice in nos. At ille intendébat in eos, sperans se áliquid acceptúrum ab eis. Petrus autem dixit: Argéntum et aurum non est mihi: quod autem hábeo, hoc tibi do: In nómine Jesu Christi Nazaréni surge, et ámbula.\par
\switchcolumn
\EN
\noindent Now Peter and John went up into the temple at the ninth hour of prayer. And a certain man who was lame from his mother's womb, was carried: whom they laid every day at the gate of the temple, which is called Beautiful, that he might ask alms of them that went into the temple. He, when he had seen Peter and John about to go into the temple, asked to receive an alms. But Peter with John fastening his eyes upon him, said: Look upon us. But he looked earnestly upon them, hoping that he should receive something of them. But Peter said: Silver and gold I have none; but what I have, I give thee: In the name of Jesus Christ of Nazareth, arise, and walk.\par
\switchcolumn*

\LA
\begin{Resp}\Rsign Christus resúrgens ex mórtuis, jam non móritur, mors illi ultra non dominábitur: quod enim mórtuus est peccáto, mórtuus est semel: \Star{} Quod autem vivit, vivit Deo, allelúja, allelúja.\end{Resp}
\begin{Resp}\Vsign Mórtuus est semel propter delícta nostra, et resurréxit propter justificatiónem nostram.\end{Resp}
\begin{Resp}\Rsign Quod autem vivit, vivit Deo, allelúja, allelúja.\end{Resp}
\switchcolumn
\EN
\begin{Resp}\Rsign Knowing that Christ rising again from the dead, dieth now no more, death shall no more have dominion over him. For in that he died to sin, he died once; \Star{} In that he liveth, he liveth unto God, alleluia, alleluia.\end{Resp}
\begin{Resp}\Vsign He died once for the sins, and he rose for our justification.\end{Resp}
\begin{Resp}\Rsign In that he liveth, he liveth unto God, alleluia, alleluia.\end{Resp}
\switchcolumn*

\LA
\LessonHead{Lectio II}
\switchcolumn
\EN
\LessonHead{Lesson II}
\switchcolumn*

\LA
\Cite{Act. 3:7-11}
\switchcolumn
\EN
\Cite{Acts 3:7-11}
\switchcolumn*

\LA
\noindent Et, apprehénsa manu ejus déxtera, allevávit eum, et prótinus consolidátæ sunt bases ejus, et plantæ. Et exsíliens stetit, et ambulábat; et intrávit cum illis in templum ámbulans, et exsíliens, et laudans Deum. Et vidit omnis pópulus eum ambulántem, et laudántem Deum. Cognoscébant autem illum, quod ipse erat, qui ad eleemósynam sedébat ad Speciósam portam templi; et impléti sunt stupóre et éxtasi in eo, quod contígerat illi. Cum tenéret autem Petrum et Joánnem, cucúrrit omnis pópulus ad eos ad pórticum, quæ appellátur Salomónis, stupéntes.\par
\switchcolumn
\EN
\noindent And taking him by the right hand, he lifted him up, and forthwith his feet and soles received strength. And he leaping up, stood, and walked, and went in with them into the temple, walking, and leaping, and praising God. And all the people saw him walking and praising God. And they knew him, that it was he who sat begging alms at the Beautiful gate of the temple: and they were filled with wonder and amazement at that which had happened to him. And as he held Peter and John, all the people ran to them to the porch which is called Solomon's, greatly wondering\par
\switchcolumn*

\LA
\begin{Resp}\Rsign Surréxit pastor bonus, qui ánimam suam pósuit pro óvibus suis, et pro grege suo mori dignátus est: \Star{} Allelúja, allelúja, allelúja.\end{Resp}
\begin{Resp}\Vsign Etenim Pascha nostrum immolátus est Christus.\end{Resp}
\begin{Resp}\Rsign Allelúja, allelúja, allelúja.\end{Resp}
\begin{Resp}\Vsign Glória Patri, et Fílio, \Star{} et Spirítui Sancto.\end{Resp}
\begin{Resp}\Rsign Allelúja, allelúja, allelúja.\end{Resp}
\switchcolumn
\EN
\begin{Resp}\Rsign The Good Shepherd, Who laid down His life for the sheep, yea, Who was contented even to die for His flock, the Good Shepherd is risen again. \Star{} Alleluia, alleluia, alleluia.\end{Resp}
\begin{Resp}\Vsign For even Christ our Passover is sacrificed for us.\end{Resp}
\begin{Resp}\Rsign Alleluia, alleluia, alleluia.\end{Resp}
\begin{Resp}\Vsign Glory be to the Father, and to the Son, \Star{} and to the Holy Ghost.\end{Resp}
\begin{Resp}\Rsign Alleluia, alleluia, alleluia.\end{Resp}
\switchcolumn*

\LA
\LessonHead{Lectio III}
\switchcolumn
\EN
\LessonHead{Lesson III}
\switchcolumn*

\LA
\Cite{Act. 3:12-16}
\switchcolumn
\EN
\Cite{Acts 3:12-16}
\switchcolumn*

\LA
\noindent Videns autem Petrus, respóndit ad pópulum: Viri Israëlítæ, quid mirámini in hoc, aut nos quid intuémini, quasi nostra virtúte aut potestáte fecérimus hunc ambuláre? Deus Abraham, et Deus Isaac, et Deus Jacob, Deus patrum nostrórum glorificávit Fílium suum Jesum, quem vos quidem tradidístis et negástis ante fáciem Piláti, judicánte illo dimítti. Vos autem Sanctum et Justum negástis, et petístis virum homicídam donári vobis: Auctórem vero vitæ interfecístis, quem Deus suscitávit a mórtuis, cujus nos testes sumus. Et in fide nóminis ejus, hunc, quem vos vidístis et nostis, confirmávit nomen ejus: et fides, quæ per eum est, dedit íntegram sanitátem istam in conspéctu ómnium vestrum.\par
\switchcolumn
\EN
\noindent But Peter seeing, made answer to the people: Ye men of Israel, why wonder you at this? or why look you upon us, as if by our strength or power we had made this man to walk? The God of Abraham, and the God of Isaac, and the God of Jacob, the God of our fathers, hath glorified his Son Jesus, whom you indeed delivered up and denied before the face of Pilate, when he judged he should be released. But you denied the Holy One and the Just, and desired a murderer to be granted unto you. But the author of life you killed, whom God hath raised from the dead, of which we are witnesses. And in the faith of his name, this man, whom you have seen and known, hath his name strengthened; and the faith which is by him, hath given this perfect soundness in the sight of you all.\par
\switchcolumn*

\LA
\TeDeum{Te Deum}
\switchcolumn
\EN
\TeDeum{Te Deum}
\switchcolumn*

\end{paracol}

\Office{Feria Quinta infra Hebdomadam I post Octavam Paschæ}{Thursday of the First Week after the Octave of Easter}{Feria}
\addcontentsline{toc}{subsection}{Feria Quinta infra Hebdomadam I post Octavam Paschæ\enspace\textit{Thursday of the First Week after the Octave of Easter}}
\begin{paracol}{2}
\LA
\LessonHead{Lectio I}
\switchcolumn
\EN
\LessonHead{Lesson I}
\switchcolumn*

\LA
\LessonTitle{De Actibus Apostolórum}
\switchcolumn
\EN
\LessonTitle{Lesson from the Acts of the Apostles}
\switchcolumn*

\LA
\Cite{Act. 5:1-6}
\switchcolumn
\EN
\Cite{Acts 5:1-6}
\switchcolumn*

\LA
\noindent Vir autem quidam nómine Ananías, cum Saphíra uxóre suo véndidit agrum, et fraudávit de prétio agri, cónscia uxóre sua: et áfferens partem quamdam, ad pedes Apostolórum pósuit. Dixit autem Petrus: Ananía, cur tentávit Sátanas cor tuum, mentíri te Spirítui Sancto, et fraudáre de prétio agri? nonne manens tibi manébat, et venúndatum in tua erat potestáte? quare posuísti in corde tuo hanc rem? non es mentítus homínibus, sed Deo. Audiens autem Ananías hæc verba, cécidit, et expirávit. Et factus est timor magnus super omnes qui audiérunt. Surgéntes autem júvenes amovérunt eum, et efferéntes sepeliérunt.\par
\switchcolumn
\EN
\noindent But a certain man named Ananias, with Saphira his wife, sold a piece of land, And by fraud kept back part of the price of the land, his wife being privy thereunto: and bringing a certain part of it, laid it at the feet of the apostles. But Peter said: Ananias, why hath Satan tempted thy heart, that thou shouldst lie to the Holy Ghost, and by fraud keep part of the price of the land? Whilst it remained, did it not remain to thee? and after it was sold, was it not in thy power? Why hast thou conceived this thing in thy heart? Thou hast not lied to men, but to God. And Ananias hearing these words, fell down, and gave up the ghost. And there came great fear upon all that heard it. And the young men rising up, removed him, and carrying him out, buried him.\par
\switchcolumn*

\LA
\begin{Resp}\Rsign Virtúte magna reddébant Apóstoli, \Star{} Testimónium resurrectiónis Jesu Christi Dómini nostri, allelúja, allelúja.\end{Resp}
\begin{Resp}\Vsign Repléti quidem Spíritu Sancto loquebántur cum fidúcia verbum Dei.\end{Resp}
\begin{Resp}\Rsign Testimónium resurrectiónis Jesu Christi Dómini nostri, allelúja, allelúja.\end{Resp}
\switchcolumn
\EN
\begin{Resp}\Rsign With great power gave the Apostles \Star{} Witness of the Resurrection of our Lord Jesus Christ, alleluia, alleluia.\end{Resp}
\begin{Resp}\Vsign They were all filled with the Holy Ghost, and they spoke the Word of God with boldness.\end{Resp}
\begin{Resp}\Rsign Witness of the Resurrection of our Lord Jesus Christ, alleluia, alleluia.\end{Resp}
\switchcolumn*

\LA
\LessonHead{Lectio II}
\switchcolumn
\EN
\LessonHead{Lesson II}
\switchcolumn*

\LA
\Cite{Act. 5:7-11}
\switchcolumn
\EN
\Cite{Acts 5:7-11}
\switchcolumn*

\LA
\noindent Factum est autem quasi horárum trium spátium, et uxor ipsíus, nésciens quod factum fúerat, introívit. Dixit autem ei Petrus: Dic mihi múlier, si tanti agrum vendidístis? At illa dixit: Etiam tanti. Petrus autem ad eam: Quid útique convénit vobis tentáre Spíritum Dómini? Ecce pedes eórum qui sepeliérunt virum tuum ad óstium, et éfferent te. Conféstim cécidit ante pedes ejus, et expirávit. Intrántes autem júvenes invenérunt illam mórtuam: et extulérunt, et sepeliérunt ad virum suum. Et factus est timor magnus in univérsa ecclésia, et in omnes qui audiérunt hæc.\par
\switchcolumn
\EN
\noindent And it was about the space of three hours after, when his wife, not knowing what had happened, came in. And Peter said to her: Tell me, woman, whether you sold the land for so much? And she said: Yea, for so much. And Peter said unto her: Why have you agreed together to tempt the Spirit of the Lord? Behold the feet of them who have buried thy husband are at the door, and they shall carry thee out. Immediately she fell down before his feet, and gave up the ghost. And the young men coming in, found her dead: and carried her out, and buried her by her husband. And there came great fear upon the whole church, and upon all that heard these things.\_\par
\switchcolumn*

\LA
\begin{Resp}\Rsign De ore prudéntis procédit mel, allelúja: dulcédo mellis est sub língua ejus, allelúja: \Star{} Favus distíllans lábia ejus, allelúja, allelúja.\end{Resp}
\begin{Resp}\Vsign Sapiéntia requiéscit in corde ejus, et prudéntia in sermóne oris illíus.\end{Resp}
\begin{Resp}\Rsign Favus distíllans lábia ejus, allelúja, allelúja.\end{Resp}
\begin{Resp}\Vsign Glória Patri, et Fílio, \Star{} et Spirítui Sancto.\end{Resp}
\begin{Resp}\Rsign Favus distíllans lábia ejus, allelúja, allelúja.\end{Resp}
\switchcolumn
\EN
\begin{Resp}\Rsign From the mouth of the wise doth proceed honey, alleluia: the sweetness of honey is under his tongue, alleluia. \Star{} His lips drop as the honey-comb, alleluia, alleluia.\end{Resp}
\begin{Resp}\Vsign Wisdom doth abide in his heart, and out of his mouth cometh understanding.\end{Resp}
\begin{Resp}\Rsign His lips drop as the honey-comb, alleluia, alleluia.\end{Resp}
\begin{Resp}\Vsign Glory be to the Father, and to the Son, \Star{} and to the Holy Ghost.\end{Resp}
\begin{Resp}\Rsign His lips drop as the honey-comb, alleluia, alleluia.\end{Resp}
\switchcolumn*

\LA
\LessonHead{Lectio III}
\switchcolumn
\EN
\LessonHead{Lesson III}
\switchcolumn*

\LA
\Cite{Act. 5:12-16}
\switchcolumn
\EN
\Cite{Acts 5:12-16}
\switchcolumn*

\LA
\noindent Per manus autem Apostolórum fiébant signa et prodígia multa in plebe. Et erant unanímiter omnes in pórticu Salomónis. Ceterórum autem nemo audébat se conjúngere illis: sed magnificábat eos pópulus. Magis autem augebátur credéntium in Dómino multitúdo virórum ac mulíerum, ita ut in platéas eícerent infírmos, et pónerent in léctulis et grabátis, ut, veniénte Petro, saltem umbra illíus obumbráret quemquam illórum, et liberaréntur ab infirmitátibus suis. Concurrébat autem et multitúdo vicinárum civitátum Jerúsalem, afferéntes ægros, et vexátos a spirítibus immúndis: qui curabántur omnes.\par
\switchcolumn
\EN
\noindent And by the hands of the apostles were many signs and wonders wrought among the people. And they were all with one accord in Solomon's porch. But of the rest no man durst join himself unto them; but the people magnified them. And the multitude of men and women who believed in the Lord, was more increased: Insomuch that they brought forth the sick into the streets, and laid them on beds and couches, that when Peter came, his shadow at the least, might overshadow any of them, and they might be delivered from their infirmities. And there came also together to Jerusalem a multitude out of the neighboring cities, bringing sick persons, and such as were troubled with unclean spirits; who were all healed.\par
\switchcolumn*

\LA
\TeDeum{Te Deum}
\switchcolumn
\EN
\TeDeum{Te Deum}
\switchcolumn*

\end{paracol}

\Office{Feria Sexta infra Hebdomadam I post Octavam Paschæ}{Friday of the First Week after the Octave of Easter}{Feria}
\addcontentsline{toc}{subsection}{Feria Sexta infra Hebdomadam I post Octavam Paschæ\enspace\textit{Friday of the First Week after the Octave of Easter}}
\begin{paracol}{2}
\LA
\LessonHead{Lectio I}
\switchcolumn
\EN
\LessonHead{Lesson I}
\switchcolumn*

\LA
\LessonTitle{De Actibus Apostolórum}
\switchcolumn
\EN
\LessonTitle{Lesson from the Acts of the Apostles}
\switchcolumn*

\LA
\Cite{Act. 8:9-13}
\switchcolumn
\EN
\Cite{Acts 8:9-13}
\switchcolumn*

\LA
\noindent Vir autem quidam nómine Simon, qui ante fúerat in civitáte magus, sedúcens gentem Samaríæ, dicens se esse áliquem magnum: cui auscultábant omnes a mínimo usque ad máximum, dicéntes: Hic est virtus Dei, quæ vocátur magna. Attendébant autem eum: propter quod multo témpore magíis suis dementásset eos. Cum vero credidíssent Philíppo evangelizánti de regno Dei, in nómine Jesu Christi baptizabántur viri ac mulíeres. Tunc Simon et ipse crédidit: et cum baptizátus esset, adhærébat Philíppo. Videns étiam signa et virtútes máximas fíeri, stupens admirabátur.\par
\switchcolumn
\EN
\noindent There was therefore great joy in that city. Now there was a certain man named Simon, who before had been a magician in that city, seducing the people of Samaria, giving out that he was some great one: To whom they all gave ear, from the least to the greatest, saying: This man is the power of God, which is called great. And they were attentive to him, because, for a long time, he had bewitched them with his magical practices. But when they had believed Philip preaching of the kingdom of God, in the name of Jesus Christ, they were baptized, both men and women. Then Simon himself believed also; and being baptized, he adhered to Philip. And being astonished, wondered to see the signs and exceeding great miracles which were done.\par
\switchcolumn*

\LA
\begin{Resp}\Rsign Ego sum vitis vera, et vos pálmites: \Star{} Qui manet in me, et ego in eo, hic fert fructum multum, allelúja, allelúja.\end{Resp}
\begin{Resp}\Vsign Sicut diléxit me Pater, et ego diléxi vos.\end{Resp}
\begin{Resp}\Rsign Qui manet in me, et ego in eo, hic fert fructum multum, allelúja, allelúja.\end{Resp}
\switchcolumn
\EN
\begin{Resp}\Rsign I am the True Vine, and ye are the branches; \Star{} He that abideth in Me, and I in him, the same bringeth forth much fruit, alleluia, alleluia.\end{Resp}
\begin{Resp}\Vsign As the Father hath loved Me, so have I loved you.\end{Resp}
\begin{Resp}\Rsign He that abideth in Me, and I in him, the same bringeth forth much fruit, alleluia, alleluia.\end{Resp}
\switchcolumn*

\LA
\LessonHead{Lectio II}
\switchcolumn
\EN
\LessonHead{Lesson II}
\switchcolumn*

\LA
\Cite{Act. 8:14-19}
\switchcolumn
\EN
\Cite{Acts 8:14-19}
\switchcolumn*

\LA
\noindent Cum autem audíssent Apóstoli qui erant Jerosólymis, quod recepísset Samaría verbum Dei, misérunt ad eos Petrum et Joánnem. Qui cum veníssent, oravérunt pro ipsis ut accíperent Spíritum Sanctum: nondum enim in quemquam illórum vénerat, sed baptizáti tantum erant in nómine Dómini Jesu. Tunc imponébant manus super illos, et accipiébant Spíritum Sanctum. Cum vidísset autem Simon quia per impositiónem manus Apostolórum darétur Spíritus Sanctus, óbtulit eis pecúniam, dicens: Date et mihi hanc potestátem, ut cuicúmque imposúero manus, accípiat Spíritum Sanctum.\par
\switchcolumn
\EN
\noindent Now when the apostles, who were in Jerusalem, had heard that Samaria had received the word of God, they sent unto them Peter and John. Who, when they were come, prayed for them, that they might receive the Holy Ghost. For he was not as yet come upon any of them; but they were only baptized in the name of the Lord Jesus. Then they laid their hands upon them, and they received the Holy Ghost. And when Simon saw, that by the imposition of the hands of the apostles, the Holy Ghost was given, he offered them money, Saying: Give me also this power, that on whomsoever I shall lay my hands, he may receive the Holy Ghost.\par
\switchcolumn*

\LA
\begin{Resp}\Rsign Surgens Jesus Dóminus noster, stans in médio discipulórum suórum, dixit: \Star{} Pax vobis, allelúja: gavísi sunt discípuli viso Dómino, allelúja.\end{Resp}
\begin{Resp}\Vsign Una ergo sabbatórum, cum fores essent clausæ, ubi erant discípuli congregáti, venit Jesus, et stetit in medio eórum, et dixit eis.\end{Resp}
\begin{Resp}\Rsign Pax vobis, allelúja: gavísi sunt discípuli viso Dómino, allelúja.\end{Resp}
\begin{Resp}\Vsign Glória Patri, et Fílio, \Star{} et Spirítui Sancto.\end{Resp}
\begin{Resp}\Rsign Pax vobis, allelúja: gavísi sunt discípuli viso Dómino, allelúja.\end{Resp}
\switchcolumn
\EN
\begin{Resp}\Rsign After that our Lord Jesus was risen again, He came and stood in the midst of His disciples, and said unto them: \Star{} Peace be unto you, alleluia. Then were the disciples glad, when they saw the Lord, alleluia.\end{Resp}
\begin{Resp}\Vsign The first day of the week, when the doors were shut where the disciples were assembled, came Jesus, and stood in the midst, and said unto them:\end{Resp}
\begin{Resp}\Rsign Peace be unto you, alleluia. Then were the disciples glad, when they saw the Lord, alleluia.\end{Resp}
\begin{Resp}\Vsign Glory be to the Father, and to the Son, \Star{} and to the Holy Ghost.\end{Resp}
\begin{Resp}\Rsign Peace be unto you, alleluia. Then were the disciples glad, when they saw the Lord, alleluia.\end{Resp}
\switchcolumn*

\LA
\LessonHead{Lectio III}
\switchcolumn
\EN
\LessonHead{Lesson III}
\switchcolumn*

\LA
\Cite{Act. 8:19-24}
\switchcolumn
\EN
\Cite{Acts 8:19-24}
\switchcolumn*

\LA
\noindent Petrus autem dixit ad eum: Pecúnia tua tecum sit in perditiónem: quóniam donum Dei existimásti pecúnia possidéri. Non est tibi pars neque sors in sermóne isto: cor enim tuum non est rectum coram Deo. Pœniténtiam ítaque age ab hac nequítia tua: et roga Deum, si forte remittátur tibi hæc cogitátio cordis tui. In felle enim amaritúdinis, et obligatióne iniquitátis, vídeo te esse. Respóndens autem Simon, dixit: Precámini vos pro me ad Dóminum, ut nihil véniat super me horum quæ dixístis.\par
\switchcolumn
\EN
\noindent But Peter said to him: Keep thy money to thyself, to perish with thee, because thou hast thought that the gift of God may be purchased with money. Thou hast no part nor lot in this matter. For thy heart is not right in the sight of God. Do penance therefore for this thy wickedness; and pray to God, that perhaps this thought of thy heart may be forgiven thee. For I see thou art in the gall of bitterness, and in the bonds of iniquity. Then Simon answering, said: Pray you for me to the Lord, that none of these things which you have spoken may come upon me.\par
\switchcolumn*

\LA
\TeDeum{Te Deum}
\switchcolumn
\EN
\TeDeum{Te Deum}
\switchcolumn*

\end{paracol}

\Office{Sabbato infra Hebdomadam I post Octavam Paschæ}{Saturday of the First Week after the Octave of Easter}{Feria}
\addcontentsline{toc}{subsection}{Sabbato infra Hebdomadam I post Octavam Paschæ\enspace\textit{Saturday of the First Week after the Octave of Easter}}
\begin{paracol}{2}
\LA
\LessonHead{Lectio I}
\switchcolumn
\EN
\LessonHead{Lesson I}
\switchcolumn*

\LA
\LessonTitle{De Actibus Apostolórum}
\switchcolumn
\EN
\LessonTitle{Lesson from the Acts of the Apostles}
\switchcolumn*

\LA
\Cite{Act. 10:1-8}
\switchcolumn
\EN
\Cite{Acts 10:1-8}
\switchcolumn*

\LA
\noindent Vir autem quidam erat in Cæsaréa, nómine Cornélius, centúrio cohórtis quæ dícitur Itálica, religiósus, ac timens Deum cum omni domo sua, fáciens eleemósynas multas plebi, et déprecans Deum semper. Is vidit in visu maniféste, quasi hora diéi nona, ángelum Dei introëúntem ad se, et dicéntem sibi: Cornéli. At ille íntuens eum, timóre corréptus, dixit: Quid est, dómine? Dixit autem illi: Oratiónes tuæ et eleemósynæ tuæ ascendérunt in memóriam in conspéctu Dei. Et nunc mitte viros in Joppen, et accérsi Simónem quemdam, qui cognominátur Petrus: hic hospitátur apud Simónem quemdam coriárium, cujus est domus juxta mare: hic dicet tibi quid te opórteat fácere. Et cum discessísset ángelus qui loquebátur illi, vocávit duos domésticos suos, et mílitem metuéntem Dóminum ex his qui illi parébant. Quibus cum narrásset ómnia, misit illos in Joppen.\par
\switchcolumn
\EN
\noindent And there was a certain man in Caesarea, named Cornelius, a centurion of that which is called the Italian band; A religious man, and fearing God with all his house, giving much alms to the people, and always praying to God. This man saw in a vision manifestly, about the ninth hour of the day, an angel of God coming in unto him, and saying to him: Cornelius. And he, beholding him, being seized with fear, said: What is it, Lord? And he said to him: thy prayers and thy alms are ascended for a memorial in the sight of God. And now send men to Joppe, and call hither one Simon, who is surnamed Peter: He lodgeth with one Simon a tanner, whose house is by the sea side. He will tell thee what thou must do. And when the angel who spoke to him was departed, he called two of his household servants, and a soldier who feared the Lord, of them that were under him. To whom when he had related all, he sent them to Joppe.\par
\switchcolumn*

\LA
\begin{Resp}\Rsign Christus resúrgens ex mórtuis, jam non móritur, mors illi ultra non dominábitur: quod enim mórtuus est peccáto, mórtuus est semel: \Star{} Quod autem vivit, vivit Deo, allelúja, allelúja.\end{Resp}
\begin{Resp}\Vsign Mórtuus est semel propter delícta nostra, et resurréxit propter justificatiónem nostram.\end{Resp}
\begin{Resp}\Rsign Quod autem vivit, vivit Deo, allelúja, allelúja.\end{Resp}
\switchcolumn
\EN
\begin{Resp}\Rsign Knowing that Christ rising again from the dead, dieth now no more, death shall no more have dominion over him. For in that he died to sin, he died once; \Star{} In that he liveth, he liveth unto God, alleluia, alleluia.\end{Resp}
\begin{Resp}\Vsign He died once for the sins, and he rose for our justification.\end{Resp}
\begin{Resp}\Rsign In that he liveth, he liveth unto God, alleluia, alleluia.\end{Resp}
\switchcolumn*

\LA
\LessonHead{Lectio II}
\switchcolumn
\EN
\LessonHead{Lesson II}
\switchcolumn*

\LA
\Cite{Act. 10:9-17}
\switchcolumn
\EN
\Cite{Acts 10:9-17}
\switchcolumn*

\LA
\noindent Póstera autem die, iter illis faciéntibus, et appropinquántibus civitáti, ascéndit Petrus in superióra ut oráret circa horam sextam. Et cum esuríret, vóluit gustáre. Parántibus autem illis, cécidit super eum mentis excéssus: et vidit cælum apértum, et descéndens vas quoddam, velut línteum magnum, quátuor inítiis submítti de cælo in terram, in quo erant ómnia quadrupédia, et serpéntia terræ, et volatília cæli. Et facta est vox ad eum: Surge, Petre: occíde, et mandúca. Ait autem Petrus: Absit Dómine, quia nunquam manducávi omne commúne et immúndum. Et vox íterum secúndo ad eum: Quod Deus purificávit, tu commúne ne díxeris. Hoc autem factum est per ter: et statim recéptum est vas in cælum. Et dum intra se hæsitáret Petrus quidnam esset vísio quam vidísset, ecce viri qui missi erant a Cornélio, inquiréntes domum Simónis astitérunt ad jánuam.\par
\switchcolumn
\EN
\noindent And on the next day, whilst they were going on their journey, and drawing nigh to the city, Peter went up to the higher parts of the house to pray, about the sixth hour. And being hungry, he was desirous to taste somewhat. And as they were preparing, there came upon him an ecstasy of mind. And he saw the heaven opened, and a certain vessel descending, as it were a great linen sheet let down by the four corners from heaven to the earth: Wherein were all manner of fourfooted beasts, and creeping things of the earth, and fowls of the air. And there came a voice to him: Arise, Peter; kill and eat. But Peter said: Far be it from me; for I never did eat any thing that is common and unclean. And the voice spoke to him again the second time: That which God hath cleansed, do not thou call common. And this was done thrice; and presently the vessel was taken up into heaven. Now, whilst Peter was doubting within himself, what the vision that he had seen should mean, behold the men who were sent from Cornelius, inquiring for Simon's house, stood at the gate.\par
\switchcolumn*

\LA
\begin{Resp}\Rsign Surréxit pastor bonus, qui ánimam suam pósuit pro óvibus suis, et pro grege suo mori dignátus est: \Star{} Allelúja, allelúja, allelúja.\end{Resp}
\begin{Resp}\Vsign Etenim Pascha nostrum immolátus est Christus.\end{Resp}
\begin{Resp}\Rsign Allelúja, allelúja, allelúja.\end{Resp}
\switchcolumn
\EN
\begin{Resp}\Rsign The Good Shepherd, Who laid down His life for the sheep, yea, Who was contented even to die for His flock, the Good Shepherd is risen again. \Star{} Alleluia, alleluia, alleluia.\end{Resp}
\begin{Resp}\Vsign For even Christ our Passover is sacrificed for us.\end{Resp}
\begin{Resp}\Rsign Alleluia, alleluia, alleluia.\end{Resp}
\switchcolumn*

\LA
\LessonHead{Lectio III}
\switchcolumn
\EN
\LessonHead{Lesson III}
\switchcolumn*

\LA
\Cite{Act. 10:34-41}
\switchcolumn
\EN
\Cite{Acts 10:34-41}
\switchcolumn*

\LA
\noindent Apériens autem Petrus os suum, dixit: In veritáte cómperi quia non est personárum accéptor Deus; sed in omni gente qui timet eum, et operátur justítiam, accéptus est illi. Verbum misit Deus fíliis Israël, annúntians pacem per Jesum Christum (hic est ómnium Dóminus). Vos scitis quod factum est verbum per univérsam Judǽam: incípiens enim a Galilǽa post baptísmum quod prædicávit Joánnes, Jesum a Názareth: quómodo unxit eum Deus Spíritu Sancto, et virtúte, qui pertránsiit benefaciéndo, et sanándo omnes oppréssos a diábolo, quóniam Deus erat cum illo. Et nos testes sumus ómnium quæ fecit in regióne Judæórum, et Jerúsalem, quem occidérunt suspendéntes in ligno. Hunc Deus suscitávit tértia die, et dedit eum maniféstum fíeri, non omni pópulo, sed téstibus præordinátis a Deo: nobis, qui manducávimus et bíbimus cum illo postquam resurréxit a mórtuis.\par
\switchcolumn
\EN
\noindent And Peter opening his mouth, said: In very deed I perceive, that God is not a respecter of persons. But in every nation, he that feareth him, and worketh justice, is acceptable to him. God sent the word to the children of Israel, preaching peace by Jesus Christ: he is Lord of all. You know the word which hath been published through all Judea: for it began from Galilee, after the baptism which John preached, Jesus of Nazareth: how God anointed him with the Holy Ghost, and with power, who went about doing good, and healing all that were oppressed by the devil, for God was with him. And we are witnesses of all things that he did in the land of the Jews and in Jerusalem, whom they killed, hanging him upon a tree. Him God raised up the third day, and gave him to be made manifest, Not to all the people, but to witnesses preordained by God, even to us, who did eat and drink with him after he arose again from the dead;\par
\switchcolumn*

\LA
\begin{Resp}\Rsign Ecce vicit leo de tribu Juda, radix David, aperíre librum, et sólvere septem signácula ejus: \Star{} Allelúja, allelúja, allelúja.\end{Resp}
\begin{Resp}\Vsign Dignus est Agnus, qui occísus est, accípere virtútem, et divinitátem, et sapiéntiam, et fortitúdinem, et honórem, et glóriam, et benedictiónem.\end{Resp}
\begin{Resp}\Rsign Allelúja, allelúja, allelúja.\end{Resp}
\begin{Resp}\Vsign Glória Patri, et Fílio, \Star{} et Spirítui Sancto.\end{Resp}
\begin{Resp}\Rsign Allelúja, allelúja, allelúja.\end{Resp}
\switchcolumn
\EN
\begin{Resp}\Rsign Behold, the Lion of the tribe of Judah, the Root of David, hath prevailed to open the Book, and to loose the seven seals thereof. \Star{} Alleluia, alleluia, alleluia.\end{Resp}
\begin{Resp}\Vsign Worthy is the Lamb That was slain to receive power, and riches, in wisdom, and strength, and honour, and glory, and blessing.\end{Resp}
\begin{Resp}\Rsign Alleluia, alleluia, alleluia.\end{Resp}
\begin{Resp}\Vsign Glory be to the Father, and to the Son, \Star{} and to the Holy Ghost.\end{Resp}
\begin{Resp}\Rsign Alleluia, alleluia, alleluia.\end{Resp}
\switchcolumn*

\end{paracol}

\Office{Dominica II Post Pascha}{Second Sunday after Easter}{Semiduplex}
\addcontentsline{toc}{subsection}{Dominica II Post Pascha\enspace\textit{Second Sunday after Easter}}
\begin{paracol}{2}
\LA
\LessonHead{Lectio I}
\switchcolumn
\EN
\LessonHead{Lesson I}
\switchcolumn*

\LA
\LessonTitle{De Actibus Apostolórum}
\switchcolumn
\EN
\LessonTitle{Lesson from the Acts of the Apostles}
\switchcolumn*

\LA
\Cite{Act. 13:13-20}
\switchcolumn
\EN
\Cite{Acts 13:13-20}
\switchcolumn*

\LA
\noindent Cum a Papho navigássent Paulus et qui cum eo erant, venérunt Pergen Pamphýliæ. Joánnes autem discédens ab eis, revérsus est Jerosólymam. Illi vero pertranseúntes Pergen, venérunt Antiochíam Pisídiæ: et ingréssi synagógam die sabbatórum, sedérunt. Post lectiónem autem legis et prophetárum, misérunt príncipes synagógæ ad eos, dicéntes: Viri fratres, si quis est in vobis sermo exhortatiónis ad plebem, dícite. Surgens autem Paulus, et manu siléntium indícens, ait: Viri Israëlítæ, et qui timétis Deum, audíte: Deus plebis Israël elégit patres nostros, et plebem exaltávit cum essent íncolæ in terra Ægýpti, et in brácchio excélso edúxit eos ex ea, et per quadragínta annórum tempus mores eórum sustínuit in desérto. Et déstruens gentes septem in terra Chánaan, sorte distríbuit eis terram eórum, quasi post quadringéntos et quinquagínta annos: et post hæc dedit júdices, usque ad Sámuel prophétam.\par
\switchcolumn
\EN
\noindent Now when Paul and they that were with him had sailed from Paphos, they came to Perge in Pamphylia. And John departing from them, returned to Jerusalem. But they passing through Perge, came to Antioch in Pisidia: and entering into the synagogue on the sabbath day, they sat down. And after the reading of the law and the prophets, the rulers of the synagogue sent to them, saying: Ye men, brethren, if you have any word of exhortation to make to the people, speak. Then Paul rising up, and with his hand bespeaking silence, said: Ye men of Israel, and you that fear God, give ear. The God of the people of Israel chose our fathers, and exalted the people when they were sojourners in the land of Egypt, and with an high arm brought them out from thence, And for the space of forty years endured their manners in the desert. And destroying seven nations in the land of Chanaan, divided their land among them, by lot, As it were, after four hundred and fifty years: and after these things, he gave unto them judges, until Samuel the prophet.\par
\switchcolumn*

\LA
\begin{Resp}\Rsign Virtúte magna reddébant Apóstoli, \Star{} Testimónium resurrectiónis Jesu Christi Dómini nostri, allelúja, allelúja.\end{Resp}
\begin{Resp}\Vsign Repléti quidem Spíritu Sancto loquebántur cum fidúcia verbum Dei.\end{Resp}
\begin{Resp}\Rsign Testimónium resurrectiónis Jesu Christi Dómini nostri, allelúja, allelúja.\end{Resp}
\switchcolumn
\EN
\begin{Resp}\Rsign With great power gave the Apostles \Star{} Witness of the Resurrection of our Lord Jesus Christ, alleluia, alleluia.\end{Resp}
\begin{Resp}\Vsign They were all filled with the Holy Ghost, and they spoke the Word of God with boldness.\end{Resp}
\begin{Resp}\Rsign Witness of the Resurrection of our Lord Jesus Christ, alleluia, alleluia.\end{Resp}
\switchcolumn*

\LA
\LessonHead{Lectio II}
\switchcolumn
\EN
\LessonHead{Lesson II}
\switchcolumn*

\LA
\Cite{Act. 13:21-25}
\switchcolumn
\EN
\Cite{Acts 13:21-25}
\switchcolumn*

\LA
\noindent Et exínde postulavérunt regem: et dedit illis Deus Saul fílium Cis, virum de tribu Bénjamin, annis quadragínta: et amóto illo, suscitávit illis David regem: cui testimónium pérhibens, dixit: Invéni David fílium Jesse, virum secúndum cor meum, qui fáciet omnes voluntátes meas. Hujus Deus ex sémine secúndum promissiónem edúxit Israël salvatórem Jesum, prædicánte Joánne ante fáciem advéntus ejus baptísmum pœniténtiæ omni pópulo Israël. Cum impléret autem Joánnes cursum suum, dicébat: Quem me arbitrámini esse, non sum ego: sed ecce venit post me, cujus non sum dignus calceaménta pedum sólvere.\par
\switchcolumn
\EN
\noindent And after that they desired a king: and God gave them Saul the son of Cis, a man of the tribe of Benjamin, forty years. And when he had removed him, he raised them up David to be king: to whom giving testimony, he said: I have found David, the son of Jesse, a man according to my own heart, who shall do all my wills. Of this man's seed God according to his promise, hath raised up to Israel a Saviour, Jesus: John first preaching, before his coming, the baptism of penance to all the people of Israel. And when John was fulfilling his course, he said: I am not he, whom you think me to be: but behold, there cometh one after me, whose shoes of his feet I am not worthy to loose.\par
\switchcolumn*

\LA
\begin{Resp}\Rsign De ore prudéntis procédit mel, allelúja: dulcédo mellis est sub língua ejus, allelúja: \Star{} Favus distíllans lábia ejus, allelúja, allelúja.\end{Resp}
\begin{Resp}\Vsign Sapiéntia requiéscit in corde ejus, et prudéntia in sermóne oris illíus.\end{Resp}
\begin{Resp}\Rsign Favus distíllans lábia ejus, allelúja, allelúja.\end{Resp}
\switchcolumn
\EN
\begin{Resp}\Rsign From the mouth of the wise doth proceed honey, alleluia: the sweetness of honey is under his tongue, alleluia. \Star{} His lips drop as the honey-comb, alleluia, alleluia.\end{Resp}
\begin{Resp}\Vsign Wisdom doth abide in his heart, and out of his mouth cometh understanding.\end{Resp}
\begin{Resp}\Rsign His lips drop as the honey-comb, alleluia, alleluia.\end{Resp}
\switchcolumn*

\LA
\LessonHead{Lectio III}
\switchcolumn
\EN
\LessonHead{Lesson III}
\switchcolumn*

\LA
\Cite{Act. 13:26-33}
\switchcolumn
\EN
\Cite{Acts 13:26-33}
\switchcolumn*

\LA
\noindent Viri fratres, fílii géneris Abraham, et qui in vobis timent Deum, vobis verbum salútis hujus missum est. Qui enim habitábant Jerúsalem, et príncipes ejus hunc ignorántes, et voces prophetárum quæ per omne sábbatum legúntur, judicántes implevérunt, et nullam causam mortis inveniéntes in eo, petiérunt a Piláto ut interfícerent eum. Cumque consummássent ómnia quæ de eo scripta erant, deponéntes eum de ligno, posuérunt eum in monuménto. Deus vero suscitávit eum a mórtuis tértia die: qui visus est per dies multos his qui simul ascénderant cum eo de Galilǽa in Jerúsalem: qui usque nunc sunt testes ejus ad plebem. Et nos vobis annuntiámus eam, quæ ad patres nostros repromíssio facta est: quóniam hanc Deus adimplévit fíliis nostris resúscitans Jesum, sicut et in psalmo secúndo scriptum est: Fílius meus es tu, ego hódie génui te.\par
\switchcolumn
\EN
\noindent Men, brethren, children of the stock of Abraham, and whosoever among you fear God, to you the word of this salvation is sent. For they that inhabited Jerusalem, and the rulers thereof, not knowing him, nor the voices of the prophets, which are read every sabbath, judging him have fulfilled them. And finding no cause of death in him, they desired of Pilate, that they might kill him. And when they had fulfilled all things that were written of him, taking him down from the tree, they laid him in a sepulchre. But God raised him up from the dead the third day: Who was seen for many days, by them who came up with him from Galilee to Jerusalem, who to this present are his witnesses to the people. And we declare unto you, that the promise which was made to our fathers, This same God hath fulfilled to our children, raising up Jesus, as in the second psalm also is written: Thou art my Son, this day have I begotten thee.\par
\switchcolumn*

\LA
\begin{Resp}\Rsign Ecce vicit leo de tribu Juda, radix David, aperíre librum, et sólvere septem signácula ejus: \Star{} Allelúja, allelúja, allelúja.\end{Resp}
\begin{Resp}\Vsign Dignus est Agnus, qui occísus est, accípere virtútem, et divinitátem, et sapiéntiam, et fortitúdinem, et honórem, et glóriam, et benedictiónem.\end{Resp}
\begin{Resp}\Rsign Allelúja, allelúja, allelúja.\end{Resp}
\begin{Resp}\Vsign Glória Patri, et Fílio, \Star{} et Spirítui Sancto.\end{Resp}
\begin{Resp}\Rsign Allelúja, allelúja, allelúja.\end{Resp}
\switchcolumn
\EN
\begin{Resp}\Rsign Behold, the Lion of the tribe of Judah, the Root of David, hath prevailed to open the Book, and to loose the seven seals thereof. \Star{} Alleluia, alleluia, alleluia.\end{Resp}
\begin{Resp}\Vsign Worthy is the Lamb That was slain to receive power, and riches, in wisdom, and strength, and honour, and glory, and blessing.\end{Resp}
\begin{Resp}\Rsign Alleluia, alleluia, alleluia.\end{Resp}
\begin{Resp}\Vsign Glory be to the Father, and to the Son, \Star{} and to the Holy Ghost.\end{Resp}
\begin{Resp}\Rsign Alleluia, alleluia, alleluia.\end{Resp}
\switchcolumn*

\LA
\LessonHead{Lectio IV}
\switchcolumn
\EN
\LessonHead{Lesson IV}
\switchcolumn*

\LA
\LessonTitle{Sermo sancti Leónis Papæ}
\switchcolumn
\EN
\LessonTitle{From the Sermons of Pope St. Leo (the Great)}
\switchcolumn*

\LA
\Source{Sermo 1 de Ascensione Domini, post initium}
\switchcolumn
\EN
\Source{1st for the Lord's Ascension}
\switchcolumn*

\LA
\noindent Hi dies, dilectíssimi, qui inter resurrectiónem Dómini ascensionémque fluxérunt, non otióso transiére decúrsu, sed magna in eis confirmáta sacraménta, magna sunt reveláta mystéria. In iis metus diræ mortis aufértur, et non solum ánimæ, sed étiam carnis immortálitas declarátur. In iis per insufflatiónem Dómini infúnditur Apóstolis ómnibus Spíritus Sanctus: et beáto Apóstolo Petro supra céteros, post regni claves, ovílis Domínici cura mandátur.\par
\switchcolumn
\EN
\noindent Dearly beloved brethren, the days which passed between the Resurrection and the Ascension of the Lord, wore not idly by, but in them were established great Sacraments, and great Mysteries were revealed. In them was abolished the terror of that fearful death, and it was shown that not the soul only, but the body also, will not die eternally. In them the breathing of the Lord on His Apostles shed upon them the Holy Ghost, and the Blessed Apostle Peter, being given the keys of the kingdom of heaven, was chosen out of the rest to receive the chief care of the Lord's fold.\par
\switchcolumn*

\LA
\begin{Resp}\Rsign Ego sum vitis vera, et vos pálmites: \Star{} Qui manet in me, et ego in eo, hic fert fructum multum, allelúja, allelúja.\end{Resp}
\begin{Resp}\Vsign Sicut diléxit me Pater, et ego diléxi vos.\end{Resp}
\begin{Resp}\Rsign Qui manet in me, et ego in eo, hic fert fructum multum, allelúja, allelúja.\end{Resp}
\switchcolumn
\EN
\begin{Resp}\Rsign I am the True Vine, and ye are the branches; \Star{} He that abideth in Me, and I in him, the same bringeth forth much fruit, alleluia, alleluia.\end{Resp}
\begin{Resp}\Vsign As the Father hath loved Me, so have I loved you.\end{Resp}
\begin{Resp}\Rsign He that abideth in Me, and I in him, the same bringeth forth much fruit, alleluia, alleluia.\end{Resp}
\switchcolumn*

\LA
\LessonHead{Lectio V}
\switchcolumn
\EN
\LessonHead{Lesson V}
\switchcolumn*

\LA
\noindent In iis diébus, duóbus discípulis tértius in via Dóminus comes júngitur, et ad omnem nostræ ambiguitátis calíginem detergéndam, pavéntium ac trepidántium tárditas increpátur. Flammam fídei illumináta corda concípiunt: et quæ erant tépida, reseránte Scriptúras Dómino, efficiúntur ardéntia. In fractióne quoque panis, convescéntium aperiúntur obtútus: multo felícius eórum óculis patefáctis, quibus natúræ suæ manifestáta est glorificátio, quam illórum géneris nostri príncipum, quibus prævaricatiónis suæ est ingésta confúsio.\par
\switchcolumn
\EN
\noindent It was during those days, that as two of His disciples were walking together, the Lord Himself joined them, and made Himself One of three companions. Then that, to clear away all shadow of doubt from our mind, He rebuked the slowness of such as still feared and trembled. Their hearts enlightened by faith, caught the flame; and, whereas they had afore been cold, they glowed again as the Lord opened to them the Scriptures. In the breaking of bread their eyes were opened, and they knew Him. And, O, how much happier were they with their eyes opened, and gazing upon the glorification of our nature in His Person, than were the first father and mother of our race, upon whom their own transgression had brought shame!\par
\switchcolumn*

\LA
\begin{Resp}\Rsign Surgens Jesus Dóminus noster, stans in médio discipulórum suórum, dixit: \Star{} Pax vobis, allelúja: gavísi sunt discípuli viso Dómino, allelúja.\end{Resp}
\begin{Resp}\Vsign Una ergo sabbatórum, cum fores essent clausæ, ubi erant discípuli congregáti, venit Jesus, et stetit in medio eórum, et dixit eis.\end{Resp}
\begin{Resp}\Rsign Pax vobis, allelúja: gavísi sunt discípuli viso Dómino, allelúja.\end{Resp}
\switchcolumn
\EN
\begin{Resp}\Rsign After that our Lord Jesus was risen again, He came and stood in the midst of His disciples, and said unto them: \Star{} Peace be unto you, alleluia. Then were the disciples glad, when they saw the Lord, alleluia.\end{Resp}
\begin{Resp}\Vsign The first day of the week, when the doors were shut where the disciples were assembled, came Jesus, and stood in the midst, and said unto them:\end{Resp}
\begin{Resp}\Rsign Peace be unto you, alleluia. Then were the disciples glad, when they saw the Lord, alleluia.\end{Resp}
\switchcolumn*

\LA
\LessonHead{Lectio VI}
\switchcolumn
\EN
\LessonHead{Lesson VI}
\switchcolumn*

\LA
\noindent Inter hæc autem, alíaque mirácula, cum discípuli trépidis cogitatiónibus æstuárent, et apparuísset in médio eórum Dóminus, dixissétque, Pax vobis: ne hoc remanéret in eórum opiniónibus, quod volvebátur in córdibus (putábant enim se spíritum vidére, non carnem) redárguit cogitatiónes a veritáte discórdes: íngerit dubitántium óculis manéntia in mánibus suis et pédibus crucis signa; et ut diligéntius pertractétur, invítat. Quia ad sanánda infidélium cordium vúlnera, clavórum et lánceæ erant serváta vestígia: ut non dúbia fide, sed constantíssima sciéntia tenerétur, eam natúram in Dei Patris consessúram throno, quæ jacúerat in sepúlcro.\par
\switchcolumn
\EN
\noindent Amid these and other miracles, while the disciples were still troubled with fearful thoughts, the Lord manifested Himself in the midst of them, and said: Peace be unto you. And lest their reason should be deceived by the vain imaginations which lurked in their hearts, (for they thought that What they saw was a spirit, and not Flesh,) He rebuked thoughts so inconsistent with the truth; and pointed out to the eyes of the doubters the marks of crucifixion which still remained in His Hands and His Feet, and bade them handle Him more closely. Those open Wounds made by the nails and spear in His Body remain ever open to close the wounds in unbelievers' hearts: that we may hold, not with doubtful faith, but with most firm and absolute knowledge, that the Manhood Which lay in the grave is the Same Which now sitteth at the right hand of God the Father.\par
\switchcolumn*

\LA
\begin{Resp}\Rsign Expurgáte vetus ferméntum, ut sitis nova conspérsio: étenim Pascha nostrum immolátus est Christus: \Star{} Itaque epulémur in Dómino, allelúja.\end{Resp}
\begin{Resp}\Vsign Mórtuus est propter delícta nostra, et resurréxit propter justificatiónem nostram.\end{Resp}
\begin{Resp}\Rsign Itaque epulémur in Dómino, allelúja.\end{Resp}
\begin{Resp}\Vsign Glória Patri, et Fílio, \Star{} et Spirítui Sancto.\end{Resp}
\begin{Resp}\Rsign Itaque epulémur in Dómino, allelúja.\end{Resp}
\switchcolumn
\EN
\begin{Resp}\Rsign Purge out the old leaven, that ye may be new dough: for even Christ our Passover is sacrificed for us: \Star{} Therefore let us keep the Feast, in the Lord, alleluia.\end{Resp}
\begin{Resp}\Vsign He died for our offences, and rose again for our justification.\end{Resp}
\begin{Resp}\Rsign Therefore let us keep the Feast, in the Lord, alleluia.\end{Resp}
\begin{Resp}\Vsign Glory be to the Father, and to the Son, \Star{} and to the Holy Ghost.\end{Resp}
\begin{Resp}\Rsign Therefore let us keep the Feast, in the Lord, alleluia.\end{Resp}
\switchcolumn*

\LA
\LessonHead{Lectio VII}
\switchcolumn
\EN
\LessonHead{Lesson VII}
\switchcolumn*

\LA
\LessonTitle{Léctio sancti Evangélii secúndum Joánnem}
\switchcolumn
\EN
\LessonTitle{From the Holy Gospel according to John}
\switchcolumn*

\LA
\Cite{Joannes 10:11-16}
\switchcolumn
\EN
\Cite{John 10:11-16}
\switchcolumn*

\LA
\noindent In illo témpore: Dixit Jesus pharisǽis: Ego sum pastor bonus. Bonus pastor ánimam suam dat pro óvibus suis. Et réliqua.\par
\switchcolumn
\EN
\noindent At that time, Jesus said unto the Pharisees: I am the Good Shepherd. The Good Shepherd giveth His life for His sheep. And so on.\par
\switchcolumn*

\LA
\LessonTitle{Homilía sancti Gregórii Papæ}
\switchcolumn
\EN
\LessonTitle{Homily by Pope St. Gregory (the Great.)}
\switchcolumn*

\LA
\Source{Homilia 14 in Evangelia}
\switchcolumn
\EN
\Source{14th on the Gospels}
\switchcolumn*

\LA
\noindent Audístis, fratres caríssimi, ex lectióne evangélica eruditiónem vestram: audístis et perículum nostrum. Ecce enim is, qui non ex accidénti dono, sed essentiáliter bonus est, dicit: Ego sum pastor bonus. Atque ejúsdem bonitátis formam, quam nos imitémur, adjúngit, dicens: Bonus pastor ánimam suam ponit pro óvibus suis. Fecit quod mónuit: osténdit quod jussit. Bonus pastor pro óvibus suis ánimam suam pósuit, ut in sacraménto nostro corpus suum et sánguinem vérteret, et oves quas redémerat, carnis suæ aliménto satiáret.\par
\switchcolumn
\EN
\noindent Dearly beloved brethren, ye have heard from the Holy Gospel what is at once your instruction, and our danger. Behold, how He Who, not by the varying gifts of nature, but of the very essence of His being, is Good, behold how He saith: I am the Good Shepherd. And then He saith what is the character of His goodness, even of that goodness of His which we must strive to copy: The Good Shepherd giveth His life for the Sheep. As He had foretold, even so did He; as He had commanded, so gave He example. The Good Shepherd gave His life for the sheep, and made His Own Body and His Own Blood to be our Sacramental Food, pasturing upon His Own Flesh the sheep whom He had bought.\par
\switchcolumn*

\LA
\begin{Resp}\Rsign Christus resúrgens ex mórtuis, jam non móritur, mors illi ultra non dominábitur: quod enim mórtuus est peccáto, mórtuus est semel: \Star{} Quod autem vivit, vivit Deo, allelúja, allelúja.\end{Resp}
\begin{Resp}\Vsign Mórtuus est semel propter delícta nostra, et resurréxit propter justificatiónem nostram.\end{Resp}
\begin{Resp}\Rsign Quod autem vivit, vivit Deo, allelúja, allelúja.\end{Resp}
\switchcolumn
\EN
\begin{Resp}\Rsign Knowing that Christ rising again from the dead, dieth now no more, death shall no more have dominion over him. For in that he died to sin, he died once; \Star{} In that he liveth, he liveth unto God, alleluia, alleluia.\end{Resp}
\begin{Resp}\Vsign He died once for the sins, and he rose for our justification.\end{Resp}
\begin{Resp}\Rsign In that he liveth, he liveth unto God, alleluia, alleluia.\end{Resp}
\switchcolumn*

\LA
\LessonHead{Lectio VIII}
\switchcolumn
\EN
\LessonHead{Lesson VIII}
\switchcolumn*

\LA
\noindent Osténsa nobis est de contémptu mortis via, quam sequámur: appósita est forma, cui imprimámur. Primum nobis est, exterióra nostra misericórditer óvibus ejus impéndere: postrémum vero, si necésse sit, étiam mortem nostram pro eísdem óvibus ministráre. A primo autem hoc mínimo pervenítur ad postrémum majus. Sed cum incomparabíliter longe sit mélior ánima, qua vívimus, quam terréna substántia, quam extérius possidémus: qui non dat pro óvibus substántiam suam, quando pro his datúrus est ánimam suam?\par
\switchcolumn
\EN
\noindent He, by despising death, hath shown us how to do the like; He hath set before us the mould wherein it behoveth us to be cast. Our first duty is, freely and tenderly to spend our outward things for His sheep, but lastly, if need be, to serve the same by our death also. From the light offering of the first, we go on to the stern offering of the last, and, if we be ready to give our life for the sheep, why should we scruple to give our substance, seeing how much more is the life than meat? (Matth. vi. 25.)\par
\switchcolumn*

\LA
\begin{Resp}\Rsign Surréxit pastor bonus, qui ánimam suam pósuit pro óvibus suis, et pro grege suo mori dignátus est: \Star{} Allelúja, allelúja, allelúja.\end{Resp}
\begin{Resp}\Vsign Etenim Pascha nostrum immolátus est Christus.\end{Resp}
\begin{Resp}\Rsign Allelúja, allelúja, allelúja.\end{Resp}
\begin{Resp}\Vsign Glória Patri, et Fílio, \Star{} et Spirítui Sancto.\end{Resp}
\begin{Resp}\Rsign Allelúja, allelúja, allelúja.\end{Resp}
\switchcolumn
\EN
\begin{Resp}\Rsign The Good Shepherd, Who laid down His life for the sheep, yea, Who was contented even to die for His flock, the Good Shepherd is risen again. \Star{} Alleluia, alleluia, alleluia.\end{Resp}
\begin{Resp}\Vsign For even Christ our Passover is sacrificed for us.\end{Resp}
\begin{Resp}\Rsign Alleluia, alleluia, alleluia.\end{Resp}
\begin{Resp}\Vsign Glory be to the Father, and to the Son, \Star{} and to the Holy Ghost.\end{Resp}
\begin{Resp}\Rsign Alleluia, alleluia, alleluia.\end{Resp}
\switchcolumn*

\LA
\LessonHead{Lectio IX}
\switchcolumn
\EN
\LessonHead{Lesson IX}
\switchcolumn*

\LA
\noindent Et sunt nonnúlli, qui dum plus terrénam substántiam quam oves díligunt, mérito nomen pastóris perdunt: de quibus prótinus súbditur: Mercenárius autem, et qui non est pastor, cujus non sunt oves própriæ, videt lupum veniéntem, et dimíttet oves, et fugit. Non pastor, sed mercenárius vocátur, qui non pro amóre íntimo oves Domínicas, sed ad temporáles mercédes pascit. Mercenárius quippe est, qui locum quidem pastóris tenet, sed lucra animárum non quærit: terrénis cómmodis ínhiat, honóre prælatiónis gaudet, temporálibus lucris páscitur, impénsa sibi ab homínibus reveréntia lætátur.\par
\switchcolumn
\EN
\noindent And some there be which love the things of this world better than they love the sheep; and such as they deserve no longer to be called shepherds. These are they of whom it is written: But he that is an hireling, and not the shepherd, whose own the sheep are not, seeth the wolf coming, and leaveth the sheep, and fleeth (12.) He is not a shepherd but an hireling which feedeth the Lord's sheep, not because he loveth their souls, but because he doth gain earthly wealth thereby. He that taketh a shepherd's place, but seeketh not gain of souls, that same is but an hireling; such an one is ever ready for creature comforts, he loveth his pre-eminence, he groweth sleek upon his income, and he liketh well to see men bow down to him.\par
\switchcolumn*

\LA
\TeDeum{Te Deum}
\switchcolumn
\EN
\TeDeum{Te Deum}
\switchcolumn*

\end{paracol}

\Office{Feria Secunda infra Hebdomadam II post Octavam Paschæ}{Monday of the Second Week after the Octave of Easter}{Feria}
\addcontentsline{toc}{subsection}{Feria Secunda infra Hebdomadam II post Octavam Paschæ\enspace\textit{Monday of the Second Week after the Octave of Easter}}
\begin{paracol}{2}
\LA
\LessonHead{Lectio I}
\switchcolumn
\EN
\LessonHead{Lesson I}
\switchcolumn*

\LA
\LessonTitle{De Actibus Apostolórum}
\switchcolumn
\EN
\LessonTitle{Lesson from the Acts of the Apostles}
\switchcolumn*

\LA
\Cite{Act. 15:5-12}
\switchcolumn
\EN
\Cite{Acts 15:5-12}
\switchcolumn*

\LA
\noindent Surrexérunt autem quidam de hǽresi Pharisæórum, qui credidérunt, dicéntes quia opórtet circumcídi eos, præcípere quoque serváre legem Móysi. Convenerúntque Apóstoli et senióres vidére de verbo hoc. Cum autem magna conquisítio fíeret, surgens Petrus dixit ad eos: Viri fratres, vos scitis quóniam ab antíquis diébus Deus in nobis elégit, per os meum audíre gentes verbum Evangélii et crédere. Et qui novit corda Deus, testimónium perhíbuit, dans illis Spíritum Sanctum, sicut et nobis, et nihil discrévit inter nos et illos, fide puríficans corda eórum. Nunc ergo quid tentátis Deum, impónere jugum super cervíces discipulórum quod neque patres nostri, neque nos portáre potúimus? sed per grátiam Dómini Jesu Christi crédimus salvári, quemádmodum et illi. Tácuit autem omnis multitúdo: et audiébant Bárnabam et Paulum narrántes quanta Deus fecísset signa et prodígia in géntibus per eos.\par
\switchcolumn
\EN
\noindent But there arose some of the sect of the Pharisees that believed, saying: They must be circumcised, and be commanded to observe the law of Moses. And the apostles and ancients assembled to consider of this matter. And when there had been much disputing, Peter, rising up, said to them: Men, brethren, you know, that in former days God made choice among us, that by my mouth the Gentiles should hear the word of the gospel, and believe. And God, who knoweth the hearts, gave testimony, giving unto them the Holy Ghost, as well as to us; And put no difference between us and them, purifying their hearts by faith. Now therefore, why tempt you God to put a yoke upon the necks of the disciples, which neither our fathers nor we have been able to bear? But by the grace of the Lord Jesus Christ, we believe to be saved, in like manner as they also. And all the multitude held their peace; and they heard Barnabas and Paul telling what great signs and wonders God had wrought among the Gentiles by them\par
\switchcolumn*

\LA
\begin{Resp}\Rsign Virtúte magna reddébant Apóstoli, \Star{} Testimónium resurrectiónis Jesu Christi Dómini nostri, allelúja, allelúja.\end{Resp}
\begin{Resp}\Vsign Repléti quidem Spíritu Sancto loquebántur cum fidúcia verbum Dei.\end{Resp}
\begin{Resp}\Rsign Testimónium resurrectiónis Jesu Christi Dómini nostri, allelúja, allelúja.\end{Resp}
\switchcolumn
\EN
\begin{Resp}\Rsign With great power gave the Apostles \Star{} Witness of the Resurrection of our Lord Jesus Christ, alleluia, alleluia.\end{Resp}
\begin{Resp}\Vsign They were all filled with the Holy Ghost, and they spoke the Word of God with boldness.\end{Resp}
\begin{Resp}\Rsign Witness of the Resurrection of our Lord Jesus Christ, alleluia, alleluia.\end{Resp}
\switchcolumn*

\LA
\LessonHead{Lectio II}
\switchcolumn
\EN
\LessonHead{Lesson II}
\switchcolumn*

\LA
\Cite{Act. 15:13-21}
\switchcolumn
\EN
\Cite{Acts 15:13-21}
\switchcolumn*

\LA
\noindent Et postquam tacuérunt, respóndit Jacóbus, dicens: Viri fratres, audíte me. Simon narrávit quemádmodum primum Deus visitávit súmere ex géntibus pópulum nómini suo. Et huic concórdant verba prophetárum: sicut scriptum est: Post hæc revértar, et reædificábo tabernáculum David quod décidit: et díruta ejus reædificábo, et érigam illud: ut requírant céteri hóminum Dóminum, et omnes gentes super quas invocátum est nomen meum, dicit Dóminus fáciens hæc. Notum a sǽculo est Dómino opus suum. Propter quod ego júdico non inquietári eos qui ex géntibus convertúntur ad Deum, sed scríbere ad eos ut abstíneant se a contaminatiónibus simulacrórum, et fornicatióne, et suffocátis, et sánguine. Móyses enim a tempóribus antíquis habet in síngulis civitátibus qui eum prǽdicent in synagógis, ubi per omne sábbatum légitur.\par
\switchcolumn
\EN
\noindent And after they had held their peace, James answered, saying: Men, brethren, hear me. Simon hath related how God first visited to take of the Gentiles a people to his name. And to this agree the words of the prophets, as it is written: After these things I will return, and will rebuild the tabernacle of David, which is fallen down; and the ruins thereof I will rebuild, and I will set it up: That the residue of men may seek after the Lord, and all nations upon whom my name is invoked, saith the Lord, who doth these things. To the Lord was his own work known from the beginning of the world. For which cause I judge that they, who from among the Gentiles are converted to God, are not to be disquieted. But that we write unto them, that they refrain themselves from the pollutions of idols, and from fornication, and from things strangled, and from blood. For Moses of old time hath in every city them that preach him in the synagogues, where he is read every sabbath.\par
\switchcolumn*

\LA
\begin{Resp}\Rsign De ore prudéntis procédit mel, allelúja: dulcédo mellis est sub língua ejus, allelúja: \Star{} Favus distíllans lábia ejus, allelúja, allelúja.\end{Resp}
\begin{Resp}\Vsign Sapiéntia requiéscit in corde ejus, et prudéntia in sermóne oris illíus.\end{Resp}
\begin{Resp}\Rsign Favus distíllans lábia ejus, allelúja, allelúja.\end{Resp}
\begin{Resp}\Vsign Glória Patri, et Fílio, \Star{} et Spirítui Sancto.\end{Resp}
\begin{Resp}\Rsign Favus distíllans lábia ejus, allelúja, allelúja.\end{Resp}
\switchcolumn
\EN
\begin{Resp}\Rsign From the mouth of the wise doth proceed honey, alleluia: the sweetness of honey is under his tongue, alleluia. \Star{} His lips drop as the honey-comb, alleluia, alleluia.\end{Resp}
\begin{Resp}\Vsign Wisdom doth abide in his heart, and out of his mouth cometh understanding.\end{Resp}
\begin{Resp}\Rsign His lips drop as the honey-comb, alleluia, alleluia.\end{Resp}
\begin{Resp}\Vsign Glory be to the Father, and to the Son, \Star{} and to the Holy Ghost.\end{Resp}
\begin{Resp}\Rsign His lips drop as the honey-comb, alleluia, alleluia.\end{Resp}
\switchcolumn*

\LA
\LessonHead{Lectio III}
\switchcolumn
\EN
\LessonHead{Lesson III}
\switchcolumn*

\LA
\Cite{Act. 15:22-29}
\switchcolumn
\EN
\Cite{Acts 15:22-29}
\switchcolumn*

\LA
\noindent Tunc plácuit Apóstolis et senióribus cum omni ecclésia elígere viros ex eis, et míttere Antiochíam cum Paulo et Bárnaba: Judam, qui cognominabátur Bársabas, et Silam, viros primos in frátribus: scribéntes per manus eórum: Apóstoli et senióres fratres, his qui sunt Antiochíæ, et Sýriæ, et Cilíciæ, frátribus ex géntibus, salútem. Quóniam audívimus quia quidam ex nobis exeúntes, turbavérunt vos verbis, everténtes ánimas vestras, quibus non mandávimus, plácuit nobis colléctis in unum elígere viros, et míttere ad vos cum caríssimis nostris Bárnaba et Paulo, homínibus qui tradidérunt ánimas suas pro nómine Dómini nostri Jesu Christi. Mísimus ergo Judam et Silam, qui et ipsi vobis verbis réferent éadem. Visum est enim Spirítui Sancto et nobis nihil ultra impónere vobis óneris quam hæc necessária: ut abstineátis vos ab immolátis simulacrórum, et sánguine, et suffocáto, et fornicatióne: a quibus custodiéntes vos, bene agétis. Valéte.\par
\switchcolumn
\EN
\noindent Then it pleased the apostles and ancients, with the whole church, to choose men of their own company, and to send to Antioch, with Paul and Barnabas, namely, Judas, who was surnamed Barsabas, and Silas, chief men among the brethren. Writing by their hands: The apostles and ancients, brethren, to the brethren of the Gentiles that are at Antioch, and in Syria and Cilicia, greeting. Forasmuch as we have heard, that some going out from us have troubled you with words, subverting your souls; to whom we gave no commandment: It hath seemed good to us, being assembled together, to choose out men, and to send them unto you, with our well beloved Barnabas and Paul: Men that have given their lives for the name of our Lord Jesus Christ. We have sent therefore Judas and Silas, who themselves also will, by word of mouth, tell you the same things. For it hath seemed good to the Holy Ghost and to us, to lay no further burden upon you than these necessary things: That you abstain from things sacrificed to idols, and from blood, and from things strangled, and from fornication; from which things keeping yourselves, you shall do well. Fare ye well.\par
\switchcolumn*

\LA
\TeDeum{Te Deum}
\switchcolumn
\EN
\TeDeum{Te Deum}
\switchcolumn*

\end{paracol}

\Office{Feria Tertia infra Hebdomadam II post Octavam Paschæ}{Tuesday of the Second Week after the Octave of Easter}{Feria}
\addcontentsline{toc}{subsection}{Feria Tertia infra Hebdomadam II post Octavam Paschæ\enspace\textit{Tuesday of the Second Week after the Octave of Easter}}
\begin{paracol}{2}
\LA
\LessonHead{Lectio I}
\switchcolumn
\EN
\LessonHead{Lesson I}
\switchcolumn*

\LA
\LessonTitle{De Actibus Apostolórum}
\switchcolumn
\EN
\LessonTitle{Lesson from the Acts of the Apostles}
\switchcolumn*

\LA
\Cite{Act. 17:22-27}
\switchcolumn
\EN
\Cite{Acts 17:22-27}
\switchcolumn*

\LA
\noindent Stans autem Paulus in médio Areópagi, ait: Viri Atheniénses, per ómnia quasi superstitiosióres vos vídeo. Prætériens enim, et videns simulácra vestra, invéni et aram in qua scriptum erat: Ignóto Deo. Quod ergo ignorántes cólitis, hoc ego annúntio vobis. Deus, qui fecit mundum, et ómnia quæ in eo sunt, hic cæli et terræ cum sit Dóminus, non in manufáctis templis hábitat, nec mánibus humánis cólitur índigens áliquo, cum ipse det ómnibus vitam, et inspiratiónem, et ómnia: fecítque ex uno omne genus hóminum inhabitáre super univérsam fáciem terræ, defíniens statúta témpora, et términos habitatiónis eórum, quǽrere Deum si forte attréctent eum, aut invéniant, quamvis non longe sit ab unoquóque nostrum.\par
\switchcolumn
\EN
\noindent But Paul standing in the midst of the Areopagus, said: Ye men of Athens, I perceive that in all things you are too superstitious. For passing by, and seeing your idols, I found an altar also, on which was written: To the unknown God. What therefore you worship, without knowing it, that I preach to you: God, who made the world, and all things therein; he, being Lord of heaven and earth, dwelleth not in temples made with hands; Neither is he served with men's hands, as though he needed any thing; seeing it is he who giveth to all life, and breath, and all things: And hath made of one, all mankind, to dwell upon the whole face of the earth, determining appointed times, and the limits of their habitation. That they should seek God, if happily they may feel after him or find him, although he be not far from every one of us:\par
\switchcolumn*

\LA
\begin{Resp}\Rsign Ego sum vitis vera, et vos pálmites: \Star{} Qui manet in me, et ego in eo, hic fert fructum multum, allelúja, allelúja.\end{Resp}
\begin{Resp}\Vsign Sicut diléxit me Pater, et ego diléxi vos.\end{Resp}
\begin{Resp}\Rsign Qui manet in me, et ego in eo, hic fert fructum multum, allelúja, allelúja.\end{Resp}
\switchcolumn
\EN
\begin{Resp}\Rsign I am the True Vine, and ye are the branches; \Star{} He that abideth in Me, and I in him, the same bringeth forth much fruit, alleluia, alleluia.\end{Resp}
\begin{Resp}\Vsign As the Father hath loved Me, so have I loved you.\end{Resp}
\begin{Resp}\Rsign He that abideth in Me, and I in him, the same bringeth forth much fruit, alleluia, alleluia.\end{Resp}
\switchcolumn*

\LA
\LessonHead{Lectio II}
\switchcolumn
\EN
\LessonHead{Lesson II}
\switchcolumn*

\LA
\Cite{Act. 17:28-33}
\switchcolumn
\EN
\Cite{Acts 17:28-33}
\switchcolumn*

\LA
\noindent In ipso enim vívimus, et movémur, et sumus: sicut et quidam vestrórum poëtárum dixérunt: Ipsíus enim et genus sumus. Genus ergo cum simus Dei, non debémus æstimáre auro, aut argénto, aut lápidi, sculptúræ artis, et cogitatiónis hóminis, divínum esse símile. Et témpora quidem hujus ignorántiæ despíciens Deus, nunc annúntiat homínibus ut omnes ubíque pœniténtiam agant, eo quod státuit diem in quo judicatúrus est orbem in æquitáte, in viro in quo státuit, fidem præbens ómnibus, súscitans eum a mórtuis. Cum audíssent autem resurrectiónem mortuórum, quidam quidem irridébant, quidam vero dixérunt: Audiémus te de hoc íterum. Sic Paulus exívit de médio eórum.\par
\switchcolumn
\EN
\noindent For in him we live, and move, and are; as some also of your own poets said: For we are also his offspring. Being therefore the offspring of God, we must not suppose the divinity to be like unto gold, or silver, or stone, the graving of art, and device of man. And God indeed having winked at the times of this ignorance, now declareth unto men, that all should every where do penance. Because he hath appointed a day wherein he will judge the world in equity, by the man whom he hath appointed; giving faith to all, by raising him up from the dead. And when they had heard of the resurrection of the dead, some indeed mocked, but others said: We will hear thee again concerning this matter. So Paul went out from among them.\par
\switchcolumn*

\LA
\begin{Resp}\Rsign Surgens Jesus Dóminus noster, stans in médio discipulórum suórum, dixit: \Star{} Pax vobis, allelúja: gavísi sunt discípuli viso Dómino, allelúja.\end{Resp}
\begin{Resp}\Vsign Una ergo sabbatórum, cum fores essent clausæ, ubi erant discípuli congregáti, venit Jesus, et stetit in medio eórum, et dixit eis.\end{Resp}
\begin{Resp}\Rsign Pax vobis, allelúja: gavísi sunt discípuli viso Dómino, allelúja.\end{Resp}
\begin{Resp}\Vsign Glória Patri, et Fílio, \Star{} et Spirítui Sancto.\end{Resp}
\begin{Resp}\Rsign Pax vobis, allelúja: gavísi sunt discípuli viso Dómino, allelúja.\end{Resp}
\switchcolumn
\EN
\begin{Resp}\Rsign After that our Lord Jesus was risen again, He came and stood in the midst of His disciples, and said unto them: \Star{} Peace be unto you, alleluia. Then were the disciples glad, when they saw the Lord, alleluia.\end{Resp}
\begin{Resp}\Vsign The first day of the week, when the doors were shut where the disciples were assembled, came Jesus, and stood in the midst, and said unto them:\end{Resp}
\begin{Resp}\Rsign Peace be unto you, alleluia. Then were the disciples glad, when they saw the Lord, alleluia.\end{Resp}
\begin{Resp}\Vsign Glory be to the Father, and to the Son, \Star{} and to the Holy Ghost.\end{Resp}
\begin{Resp}\Rsign Peace be unto you, alleluia. Then were the disciples glad, when they saw the Lord, alleluia.\end{Resp}
\switchcolumn*

\LA
\LessonHead{Lectio III}
\switchcolumn
\EN
\LessonHead{Lesson III}
\switchcolumn*

\LA
\Cite{Act. 17:34; 18:1-4}
\switchcolumn
\EN
\Cite{Acts 17:34, 18:1-4}
\switchcolumn*

\LA
\noindent Quidam vero viri adhæréntes ei, credidérunt: in quibus et Dionýsius Areopagíta, et múlier nómine Dámaris, et álii cum eis. Post hæc egréssus ab Athénis, venit Corínthum: et invéniens quemdam Judǽum nómine Aquilam, Pónticum génere, qui nuper vénerat ab Itália, et Priscíllam uxórem ejus (eo quod præcepísset Cláudius discédere omnes Judǽos a Roma), accéssit ad eos. Et quia ejúsdem erat artis, manébat apud eos, et operabátur. (Erant autem scenofactóriæ artis.) Et disputábat in synagóga per omne sábbatum, interpónens nomen Dómini Jesu: suadebátque Judǽis et Græcis.\par
\switchcolumn
\EN
\noindent But certain men adhering to him, did believe; among whom was also Dionysius, the Areopagite, and a woman named Damaris, and others with them. After these things, departing from Athens, he came to Corinth. And finding a certain Jew, named Aquila, born in Pontus, lately come from Italy, with Priscilla his wife, (because that Claudius had commanded all Jews to depart from Rome,) he came to them. And because he was of the same trade, he remained with them, and wrought; (now they were tentmakers by trade.) And he reasoned in the synagogue every sabbath, bringing in the name of the Lord Jesus; and he persuaded the Jews and the Greeks.\par
\switchcolumn*

\LA
\TeDeum{Te Deum}
\switchcolumn
\EN
\TeDeum{Te Deum}
\switchcolumn*

\end{paracol}

\Office{Patrocinii St. Joseph Confessoris Sponsi B.M.V. Confessoris}{Patrocinii St. Joseph Confessoris Sponsi B.M.V. confessoris}{Duplex I classis}
\addcontentsline{toc}{subsection}{Patrocinii St. Joseph Confessoris Sponsi B.M.V. Confessoris\enspace\textit{Patrocinii St. Joseph Confessoris Sponsi B.M.V. confessoris}}
\begin{paracol}{2}
\LA
\LessonHead{Lectio I}
\switchcolumn
\EN
\LessonHead{Lesson I}
\switchcolumn*

\LA
\LessonTitle{De libro Génesis}
\switchcolumn
\EN
\LessonTitle{Lesson from the book of Genesis}
\switchcolumn*

\LA
\Cite{Gen 39:1-6}
\switchcolumn
\EN
\Cite{Gen 39:1-6}
\switchcolumn*

\LA
\noindent Ígitur Joseph ductus est in Ægýptum, emítque eum Putíphar eunúchus Pharaónis, princeps exércitus, vir Ægýptius, de manu Ismaelitárum, a quibus perdúctus erat. Fuítque Dóminus cum eo, et erat vir in cunctis próspere agens: habitavítque in domo dómini sui, qui óptime nóverat Dóminum esse cum eo, et ómnia, quæ gerébat, ab eo dírigi in manu illíus. Invenítque Joseph grátiam coram dómino suo, et ministrábat ei: a quo præpósitus ómnibus gubernábat créditam sibi domum, et univérsa quæ ei trádita fúerant: benedixítque Dóminus dómui Ægýptii propter Joseph, et multiplicávit tam in ǽdibus quam in agris cunctam ejus substántiam: nec quidquam áliud nóverat, nisi panem quo vescebátur. Erat autem Joseph pulchra fácie, et decórus aspéctu.\par
\switchcolumn
\EN
\noindent And Joseph was brought into Egypt, and Putiphar an eunuch of Pharao, chief captain of the army, an Egyptian, bought him of the Ismaelites, by whom he was brought. And the Lord was with him, and he was a prosperous man in all things: and he dwelt in his master's house, Who knew very well that the Lord was with him, and made all that he did to prosper in his hand. And Joseph found favour in the sight of his master, and ministered to him: and being set over all by him, he governed the house committed to him, and all things that were delivered to him: And the Lord blessed the house of the Egyptian for Joseph's sake, and multiplied all his substance, both at home, and in the fields. Neither knew he any other thing, but the bread which he ate. And Joseph was of a beautiful countenance, and comely to behold.\par
\switchcolumn*

\LA
\begin{Resp}\Rsign Clamávit pópulus ad regem, aliménta petens: \Star{} Quibus ille respóndit: Ite ad Joseph, allelúja.\end{Resp}
\begin{Resp}\Vsign Salus nostra in manu tua est: réspice nos tantum, et læti serviémus regi.\end{Resp}
\begin{Resp}\Rsign Quibus ille respóndit: Ite ad Joseph, allelúja.\end{Resp}
\switchcolumn
\EN
\begin{Resp}\Rsign The people cried to Pharaoh for bread \Star{} And he answered them Go unto Joseph. Alleluia.\end{Resp}
\begin{Resp}\Vsign The saving of our lives is in thy hand; only let us find grace in thy sight, and we will gladly be Pharaoh's servants.\end{Resp}
\begin{Resp}\Rsign And he answered them Go unto Joseph. Alleluia.\end{Resp}
\switchcolumn*

\LA
\LessonHead{Lectio II}
\switchcolumn
\EN
\LessonHead{Lesson II}
\switchcolumn*

\LA
\Cite{Gen 41:37-43}
\switchcolumn
\EN
\Cite{Gen 41:37-43}
\switchcolumn*

\LA
\noindent Plácuit Pharaóni consílium et cunctis minístris ejus: locutúsque est ad eos: Num inveníre potérimus talem virum, qui Spíritu Dei plenus sit? Dixit ergo ad Joseph: Quia osténdit tibi Deus ómnia quæ locútus es, numquid sapientiórem et consímilem tui inveníre pótero? Tu eris super domum meam, et ad tui oris impérium cunctus pópulus obédiet: uno tantum regni sólio te præcédam. Dixítque rursus Phárao ad Joseph: Ecce, constítui te super univérsam terram Ægýpti. Tulítque ánnulum de manu sua, et dedit eum in manu ejus: vestivítque eum stola býssina, et collo torquem áuream circumpósuit. Fecítque eum ascéndere super currum suum secúndum, clamánte præcóne, ut omnes coram eo genu flécterent, et præpósitum esse scirent univérsæ terræ Ægýpti.\par
\switchcolumn
\EN
\noindent The counsel pleased Pharao and all his servants. And he said to them: Can we find such another man, that is full of the spirit of God? He said therefore to Joseph: Seeing God hath shown thee all that thou hast said, can I find one wiser and one like unto thee? Thou shalt be over my house, and at the commandment of thy mouth all the people shall obey: only in the kingly throne will I be above thee. And again Pharao said to Joseph: Behold, I have appointed thee over the whole land of Egypt. And he took his ring from his own hand, and gave it into his hand: and he put upon him a robe of silk, and put a chain of gold about his neck. And he made him go up into his second chariot, the crier proclaiming that all should bow their knee before him, and that they should know he was made governor over the whole land of Egypt.\par
\switchcolumn*

\LA
\begin{Resp}\Rsign Fecit me Deus quasi patrem regis, et dóminum univérsæ domus ejus: \Star{} Exaltávit me, ut salvos fáceret multos pópulos, allelúja.\end{Resp}
\begin{Resp}\Vsign Veníte ad me, et ego dabo vobis ómnia bona Ægýpti, ut comedátis medúllam terræ.\end{Resp}
\begin{Resp}\Rsign Exaltávit me, ut salvos fáceret multos pópulos, allelúja.\end{Resp}
\switchcolumn
\EN
\begin{Resp}\Rsign God hath made me as a father to Pharaoh, and lord of all his house. \Star{} He hath made me great, to save much people alive. Alleluia.\end{Resp}
\begin{Resp}\Vsign Come unto me, and I wili give you all the good of the land of Egypt, and ye shall eat the fat of the land.\end{Resp}
\begin{Resp}\Rsign He hath made me great, to save much people alive. Alleluia.\end{Resp}
\switchcolumn*

\LA
\LessonHead{Lectio III}
\switchcolumn
\EN
\LessonHead{Lesson III}
\switchcolumn*

\LA
\Cite{Gen 41:44-49}
\switchcolumn
\EN
\Cite{Gen 41:44-49}
\switchcolumn*

\LA
\noindent Dixit quoque rex ad Joseph: Ego sum Phárao: absque tuo império non movébit quisquam manum aut pedem in omni terra Ægýpti. Vertítque nomen ejus, et vocávit eum, lingua Ægyptíaca, Salvatórem mundi. Dedítque illi uxórem Áseneth fíliam Putíphare sacerdótis Heliopóleos. Egréssus est ítaque Joseph ad terram Ægýpti (trigínta autem annórum erat quando stetit in conspéctu regis Pharaónis), et circuívit omnes regiónes Ægýpti. Venítque fertílitas septem annórum: et in manípulos redáctæ ségetes congregátæ sunt in hórrea Ægýpti. Omnis étiam frugum abundántia in síngulis úrbibus cóndita est. Tantáque fuit abundántia trítici, ut arénæ maris coæquarétur, et cópia mensúram excéderet.\par
\switchcolumn
\EN
\noindent And the king said to Joseph: I am Pharao; without thy commandment no man shall move hand or foot in all the land of Egypt. And he turned his name, and called him in the Egyptian tongue, The saviour of the world. And he gave him to wife Aseneth the daughter of Putiphare priest of Heliopolis. Then Joseph went out to the land of Egypt: (Now he was thirty years old when he stood before king Pharao) and he went round all the countries of Egypt. And the fruitfulness of the seven years came: and the corn being bound up into sheaves was gathered together into the barns of Egypt. And all the abundance of grain was laid up in every city. And there was so great abundance of wheat, that it was equal to the sand of the sea, and the plenty exceeded measure.\par
\switchcolumn*

\LA
\begin{Resp}\Rsign Jam lætus móriar, quia vidi fáciem tuam, et supérstitem te relínquo. Non sum fraudátus aspéctu tuo: \Star{} Ínsuper osténdit mihi Dóminus semen tuum, allelúja.\end{Resp}
\begin{Resp}\Vsign Qui pascit me ab adolescéntia mea, benedícat púeris istis, et invocétur super eos nomen meum.\end{Resp}
\begin{Resp}\Rsign Ínsuper osténdit mihi Dóminus semen tuum, allelúja.\end{Resp}
\begin{Resp}\Vsign Glória Patri, et Fílio, \Star{} et Spirítui Sancto.\end{Resp}
\begin{Resp}\Rsign Ínsuper osténdit mihi Dóminus semen tuum, allelúja.\end{Resp}
\switchcolumn
\EN
\begin{Resp}\Rsign Now shall I die happy, since I have seen thy face, and do leave thee behind me. I am not disappointed of seeing thee. \Star{} The Lord hath showed me also thy seed. Alleluia.\end{Resp}
\begin{Resp}\Vsign He That hath fed me from my youth up, bless the lads, and let my name be named on them.\end{Resp}
\begin{Resp}\Rsign The Lord hath showed me also thy seed. Alleluia.\end{Resp}
\begin{Resp}\Vsign Glory be to the Father, and to the Son, \Star{} and to the Holy Ghost.\end{Resp}
\begin{Resp}\Rsign The Lord hath showed me also thy seed. Alleluia.\end{Resp}
\switchcolumn*

\LA
\LessonHead{Lectio IV}
\switchcolumn
\EN
\LessonHead{Lesson IV}
\switchcolumn*

\LA
\LessonTitle{Sermo sancti Bernardíni Senénsis.}
\switchcolumn
\EN
\LessonTitle{From the Sermons of St. Bernardine of Siena.}
\switchcolumn*

\LA
\Source{Sermo 1 de S. Joseph.}
\switchcolumn
\EN
\Source{1st of St. Joseph.)}
\switchcolumn*

\LA
\noindent Ómnium singulárium gratiárum alícui rationábili creatúræ communicatárum generális régula est, quod, quandocúmque divína grátia éligit áliquem ad áliquam grátiam sublímem statum, ómnia charísmata donet, quæ illi persónæ sic eléctæ et ejus offício necessária sunt atque illam copióse décorant. Quod máxime verificátum est in sancto Joseph, putatívo patre Dómini nostri Jesu Christi, et vero sponso Regínæ mundi et Dóminæ Angelórum; qui ab ætérno Patre eléctus est fidélis nutrítius atque custos principálium thesaurórum suórum, scílicet Fílii ejus et Sponsæ suæ: quod offícium fidelíssime prosecútus est. Cui proptérea Dóminus ait: Serve bone et fidélis, intra in gáudium Dómini tui.\par
\switchcolumn
\EN
\noindent When any special favours are conferred upon a reasonable being, it is the common rule that whenever the grace of God electeth such and such an one for such and such a grace, or for such and such an high post of duty, the person so elected receiveth all the gifts of grace which be needful for him in that state of life whereunto he is called, and receiveth them abundantly. Of this there is an excellent instance in the case of the holy Joseph, the socalled father of our Lord Jesus Christ, and the real husband of her, who is Queen of the world, and Lady of Angels. He had been elected by the Eternal Father to be the faithful nurse and warder of His two chief treasures, that is, His Son, and Joseph's own Wife. This duty Joseph faithfully discharged, and consequently the Lord hath said to him: Well done, thou good and faithful servant enter thou into the joy of thy Lord. (Matth. xxv. 21.)\par
\switchcolumn*

\LA
\begin{Resp}\Rsign Dedísti mihi protectiónem salútis tuæ, et déxtera tua suscépit me: \Star{} Protéctor meus, et cornu salútis meæ, et suscéptor meus, allelúja.\end{Resp}
\begin{Resp}\Vsign Ego protéctor tuus sum, et merces tua magna nimis.\end{Resp}
\begin{Resp}\Rsign Protéctor meus, et cornu salútis meæ, et suscéptor meus, allelúja.\end{Resp}
\switchcolumn
\EN
\begin{Resp}\Rsign Thou hast given me the shield of thy salvation, and thy right hand hath holden me up. \Star{} My buckler, and the horn of my salvation, and my refuge. Alleluia.\end{Resp}
\begin{Resp}\Vsign I am thy shield and thy exceeding great reward.\end{Resp}
\begin{Resp}\Rsign My buckler, and the horn of my salvation, and my refuge. Alleluia.\end{Resp}
\switchcolumn*

\LA
\LessonHead{Lectio V}
\switchcolumn
\EN
\LessonHead{Lesson V}
\switchcolumn*

\LA
\noindent Si cómpares eum ad totam Ecclésiam Christi, nonne iste est homo eléctus et speciális, per quem et sub quo Christus est ordináte et honéste introdúctus in mundum? Si ergo Vírgini Matri tota Ecclésia sancta débitrix est, quia per eam Christum suscípere digna facta est: sic profécto post eam huic debet grátiam et reveréntiam singulárem. Ipse enim est clavis véteris testaménti, in qua patriarchális et prophetális dígnitas promíssum conséquitur fructum. Porro hic est solus, qui corporáliter possédit, quod eis divína dignátio repromísit. Mérito ígitur figurátur per illum patriárcham Joseph, qui pópulis fruménta servávit. Sed et hic illum præcéllit, quia non solum Ægýptiis panem corporális vitæ, sed ómnibus eléctis panem de cælo, qui cæléstem vitam tríbuit, cum multa sollértia enutrívit.\par
\switchcolumn
\EN
\noindent This man Joseph, if we compare him with the Universal Church of Christ, is he not that elect and chosen one, through whom, and under whom, Christ is orderly and honestly brought into the world? If, then, the Holy Universal Church be under a debt to the Virgin Mother, because it is through her that she hath been made to receive Christ, next to Mary she oweth love and worship to Joseph. Joseph is the key of the (Church of the Saints) which were under the Old Testament, in whose person the noble structure of Patriarchs and Prophets reacheth her completion and realiseth her promises. He is the only one of them who actually enjoyed in full fruition what God had been pleased to promise before to them. It is, therefore, with good reason that we see a type of him in that Patriarch Joseph who stored up corn for the people. But the second Joseph hath a more excellent dignity than the first, seeing that the first only gave to the Egyptians bread for the body, but the second was the watchful guardian for all the elect of that Living Bread Which came down from heaven, of Which whosoever eateth will never die.\par
\switchcolumn*

\LA
\begin{Resp}\Rsign Státuet fílios suos sub tégmine illíus, et sub ramis ejus morábitur: protegétur sub tégmine illíus a fervóre: \Star{} Et in glória ejus requiéscet, allelúja.\end{Resp}
\begin{Resp}\Vsign Speráte in eo, omnis congregátio pópuli, effúndite coram illo corda vestra.\end{Resp}
\begin{Resp}\Rsign Et in glória ejus requiéscet, allelúja.\end{Resp}
\switchcolumn
\EN
\begin{Resp}\Rsign He shall set his children under her shelter, and shall lodge under her branches by her shall he be covered from heat, \Star{} And in her glory shall he dwell. Alleluia.\end{Resp}
\begin{Resp}\Vsign Trust in Him, ye congregation of the people, pour out your heart before Him.\end{Resp}
\begin{Resp}\Rsign And in her glory shall he dwell. Alleluia.\end{Resp}
\switchcolumn*

\LA
\LessonHead{Lectio VI}
\switchcolumn
\EN
\LessonHead{Lesson VI}
\switchcolumn*

\LA
\noindent Profécto dubitándum non est, quod Christus familiaritátem, reveréntiam atque sublimíssimam dignitátem, quam illi exhíbuit, dum ágeret in humánis, tamquam fílius patri suo, in cælis útique non negávit, quin pótius complévit et consummávit. Unde non immérito in verbo propósito a Dómino subinfértur: Intra in gáudium Dómini tui. Unde, licet gáudium ætérnæ beatitúdinis in cor hóminis intret, máluit tamen Dóminus ei dícere, Intra in gáudium; ut mýstice innuátur, quod gáudium illud non solum in eo sit intra, sed úndique illum circúmdans et absórbens, et ipsum velut abýssus infiníta submérgens. Meménto ígitur nostri, beáte Joseph, et tuæ oratiónis suffrágio apud tuum putatívum Fílium intercéde; sed et beatíssimam Vírginem sponsam tuam nobis propítiam redde, quæ Mater est ejus, qui cum Patre et Spíritu Sancto vivit et regnat per infiníta sǽcula sæculórum. Amen.\par
\switchcolumn
\EN
\noindent There can be no doubt that Christ still treateth Joseph in heaven with that familiarity, honour, and most high condescension which He paid him, like a Son to a father, while He walked among men; nay, rather, that He hath now crowned and completed those habits. We may very reasonably suspect that it was with a peculiar meaning that Christ said (to him) Enter thou into the joy of thy Lord. The joy of being blessed for ever entereth into the heart of man, but when the Lord said (to Joseph), Enter thou into joy, He probably meant mystically to bid him realise a joy which should not be within him only, but outside him also, above him, and below him, and all round about him, and overflowing him as it were a great bottomless pit of joy to swallow him up altogether. Therefore, O thou blessed Joseph! remember us! In thy helpful prayers, make intercession for us with Him Who vouchsafed to be supposed thy Son! Likewise, obtain some pity for us from that most blessed Maiden who was thy wife, and the Mother of Him, Who, with the Father and the Holy Ghost, liveth and reigneth, God, world without end. Amen.\par
\switchcolumn*

\LA
\begin{Resp}\Rsign Si consístant advérsum me castra, non timébit cor meum: \Star{} Si exsúrgat advérsum me prǽlium, in hoc ego sperábo, allelúja.\end{Resp}
\begin{Resp}\Vsign In te cantátio mea semper, quóniam tu adjútor fortis.\end{Resp}
\begin{Resp}\Rsign Si exsúrgat advérsum me prǽlium, in hoc ego sperábo, allelúja.\end{Resp}
\begin{Resp}\Vsign Glória Patri, et Fílio, \Star{} et Spirítui Sancto.\end{Resp}
\begin{Resp}\Rsign Si exsúrgat advérsum me prǽlium, in hoc ego sperábo, allelúja.\end{Resp}
\switchcolumn
\EN
\begin{Resp}\Rsign Though an host should encamp against me, my heart shall not fear. \Star{} Though war should rise against me, in this will I be confident. Alleluia.\end{Resp}
\begin{Resp}\Vsign My praise shall be continually of thee, for Thou art my strong refuge. R.* Though war should rise against me, in this will I be confident. Alleluia.\end{Resp}
\begin{Resp}\Vsign Glory be to the Father, and to the Son, \Star{} and to the Holy Ghost.\end{Resp}
\begin{Resp}\Rsign Though war should rise against me, in this will I be confident. Alleluia.\end{Resp}
\switchcolumn*

\LA
\LessonHead{Lectio VII}
\switchcolumn
\EN
\LessonHead{Lesson VII}
\switchcolumn*

\LA
\LessonTitle{Léctio sancti Evangélii secúndum Lucam}
\switchcolumn
\EN
\LessonTitle{From the Holy Gospel according to Luke}
\switchcolumn*

\LA
\Cite{Luc 3:21-23}
\switchcolumn
\EN
\Cite{Luke 3:21}
\switchcolumn*

\LA
\noindent In illo témpore: Factum est autem cum baptizarétur omnis pópulus, et Jesu baptizáto et oránte, apértum est cælum. Et réliqua.\par
\switchcolumn
\EN
\noindent At that time: When all the people were baptized, it came to pass, that Jesus also being baptized and praying, the heaven was opened. And so on.\par
\switchcolumn*

\LA
\LessonTitle{Homilía sancti Augustíni Epíscopi}
\switchcolumn
\EN
\LessonTitle{Homily by St. Augustine, Bishop (of Hippo.)}
\switchcolumn*

\LA
\Source{Liber 2 de Consensu Evang.}
\switchcolumn
\EN
\Source{Book ii., on the Harmony of the Evangelists.}
\switchcolumn*

\LA
\noindent Maniféstum est illud quod ait, Ut putabátur fílius Joseph, propter illos dixísse, qui eum ex Joseph, sicut álii hómines nascúntur, natum arbitrántur. Quos autem movet, quod álios progenitóres Matthǽus enúmerat, descéndens a David usque ad Joseph, álios autem Lucas, ascéndens a Joseph usque ad David; fácile est, ut advértant duos patres habére potuísse Joseph: unum, a quo génitus, álterum, a quo fúerit adoptátus. Antíqua est enim consuetúdo adoptándi étiam in illo pópulo Dei, ut sibi fílios fácerent, quos non ipsi genuíssent. Unde intellígitur Lucas patrem Joseph, non a quo génitus, sed a quo fúerat adoptátus, suscepísse in Evangélio suo, cujus progenitóres sursum versus commémorat, donec exíret ad David.\par
\switchcolumn
\EN
\noindent And Jesus Himself began to be about thirty years of age, being (as was supposed) the Son of Joseph. These words, as was supposed, were evidently here written for the correction of such as might think that the Lord was the Son of Joseph, in the same sense as other men are called the children of their fathers. Those who find any trouble in the fact that the ancestors reckoned downward by Matthew from David to Joseph, are other than those reckoned upward by Luke from Joseph to David, such, I say, as are troubled by this, may get over it by supposing that Joseph had two fathers; one, that is, who begat him, and another who adopted him. The custom of adopting children, whereby those who have none of their own surround themselves with a family, is very ancient, even among the people of God. Hence, Luke is understood to have included in his Gospel, under the name of father of Joseph, that, not of the father by whom he was begotten, but of him by whom he was adopted, and it is the ancestors of this adoptive father who are reckoned up as far as David.\par
\switchcolumn*

\LA
\begin{Resp}\Rsign Joseph, fili David, noli timére accípere Maríam cónjugem tuam; quod enim in ea natum est, de Spíritu Sancto est: páriet autem fílium, \Star{} Et vocábis nomen ejus Jesum, allelúja.\end{Resp}
\begin{Resp}\Vsign Ipse enim salvum fáciet pópulum suum a peccátis eórum.\end{Resp}
\begin{Resp}\Rsign Et vocábis nomen ejus Jesum, allelúja.\end{Resp}
\switchcolumn
\EN
\begin{Resp}\Rsign Joseph, thou Son of David, fear not to take unto thee Mary thy wife; for That Which is conceived in her is of the Holy Ghost and she shall bring forth a Son; \Star{} And thou shalt call His Name Jesus. Alleluia.\end{Resp}
\begin{Resp}\Vsign For He shall save His people from their sins.\end{Resp}
\begin{Resp}\Rsign And thou shalt call His Name Jesus. Alleluia.\end{Resp}
\switchcolumn*

\LA
\LessonHead{Lectio VIII}
\switchcolumn
\EN
\LessonHead{Lesson VIII}
\switchcolumn*

\LA
\noindent Cum enim necésse sit, utróque Evangelísta vera narránte, et Matthǽo scílicet et Luca, ut unus eórum ejus patris oríginem tenéret qui genúerat, alter ejus qui adoptáverat Joseph; quem probabílius intellígimus adoptántis oríginem tenuísse quam eum, qui nóluit Joseph génitum dícere ab illo cujus eum fílium esse narrábat? Matthǽus autem dicens, Ábraham génuit Isaac, Isaac génuit Jacob, atque ita in hoc verbo, quod est, Génuit, persevérans, donec in último díceret, Jacob autem génuit Joseph; satis expréssit ad eum patrem se perduxísse oríginem generátium, a quo Joseph non adoptátus, sed génitus erat.\par
\switchcolumn
\EN
\noindent Thus since we are behoven to believe that what each of the Evangelists said was true, Matthew as well as Luke; and therefore that one of them nameth the father who begat, and the other, the father who adopted, Joseph; we naturally suppose that the Evangelist, who nameth the adoptive father, was he who abstaineth from using the term beget.Matthew beginneth (i. 2) Abraham begat Isaac; and Isaac begat Jacob, and so on, always with the use of this word begat, till he cometh to: and Jacob begat Joseph. By the word which he useth he doth sufficiently indicate that the genealogy which he is giving is that of him who begat.\par
\switchcolumn*

\LA
\begin{Resp}\Rsign Surge, et áccipe Púerum et matrem ejus, et fuge in Ægýptum: \Star{} Et esto ibi usque dum dicam tibi, allelúja.\end{Resp}
\begin{Resp}\Vsign Ut adimplerétur quod dictum est a Dómino per Prophétam dicéntem: Ex Ægýpto vocávi Fílium meum.\end{Resp}
\begin{Resp}\Rsign Et esto ibi usque dum dicam tibi, allelúja.\end{Resp}
\begin{Resp}\Vsign Glória Patri, et Fílio, \Star{} et Spirítui Sancto.\end{Resp}
\begin{Resp}\Rsign Et esto ibi usque dum dicam tibi, allelúja.\end{Resp}
\switchcolumn
\EN
\begin{Resp}\Rsign Arise, and take the young Child, and His Mother, and flee into Egypt; \Star{} And be thou there until I bring thee word. Alleluia.\end{Resp}
\begin{Resp}\Vsign That it might be fulfilled which was spoken of the Lord by the Prophets, saying Out of Egypt have I called My Son.\end{Resp}
\begin{Resp}\Rsign And be thou there until I bring thee word. Alleluia.\end{Resp}
\begin{Resp}\Vsign Glory be to the Father, and to the Son, \Star{} and to the Holy Ghost.\end{Resp}
\begin{Resp}\Rsign And be thou there until I bring thee word. Alleluia.\end{Resp}
\switchcolumn*

\LA
\LessonHead{Lectio IX}
\switchcolumn
\EN
\LessonHead{Lesson IX}
\switchcolumn*

\LA
\noindent Quamquam si étiam Lucas génitum díceret Joseph ab Heli, nec sic nos hoc verbum perturbáre debéret, ut áliud crederémus quam ab uno Evangelísta gignéntem, ab áltero adoptántem patrem fuísse commemorátum. Neque enim absúrde quisque dícitur non carne sed caritáte genuísse, quem fílium sibi adoptáverit. At vero étiam nos, quibus dedit Deus potestátem fílios ejus fíeri, de natúra atque substántia sua non nos génuit, sicut únicum Fílium, sed útique dilectióne adoptávit.\par
\switchcolumn
\EN
\noindent Luke saith Joseph was the son of Heli, not Joseph was begotten of Heli; but even if he had said the latter, it would not have troubled this interpretation of ours, that one Evangelist nameth the natural, and the other the adoptive father of Joseph. It is not an outrageous thing to say that one who adopteth another hath begotten him, albeit he hath done it, not carnally, but by love. Even so hath God given to us the power to become His sons, albeit He hath not begotten us of His Own Nature and Substance, as He hath His Only - Begotten Son, but only reckoneth us, in His love, among His children.\par
\switchcolumn*

\LA
\TeDeum{Te Deum}
\switchcolumn
\EN
\TeDeum{Te Deum}
\switchcolumn*

\end{paracol}

\Office{De II die infra Octavam S. Joseph}{Second Day within the Octave of St. Joseph}{Semiduplex}
\addcontentsline{toc}{subsection}{De II die infra Octavam S. Joseph\enspace\textit{Second Day within the Octave of St. Joseph}}
\begin{paracol}{2}
\LA
\LessonHead{Lectio I}
\switchcolumn
\EN
\LessonHead{Lesson I}
\switchcolumn*

\LA
\LessonTitle{De Actibus Apostolórum}
\switchcolumn
\EN
\LessonTitle{Lesson from the Acts of the Apostles}
\switchcolumn*

\LA
\Cite{Act. 24:10-16}
\switchcolumn
\EN
\Cite{Acts 24:10-16}
\switchcolumn*

\LA
\noindent Respóndit autem Paulus (annuénte sibi prǽside dícere): Ex multis annis te esse júdicem genti huic sciens, bono ánimo pro me satisfáciam. Potes enim cognóscere quia non plus sunt mihi dies quam duódecim, ex quo ascéndi adoráre in Jerúsalem: et neque in templo invenérunt me cum áliquo disputántem, aut concúrsum faciéntem turbæ, neque in synagógis, neque in civitáte: neque probáre possunt tibi de quibus nunc me accúsant. Confíteor autem hoc tibi, quod secúndum sectam quam dicunt hǽresim, sic desérvio Patri et Deo meo, credens ómnibus quæ in lege et prophétis scripta sunt: spem habens in Deum, quam et hi ipsi exspéctant, resurrectiónem futúram justórum et iniquórum. In hoc et ipse stúdeo sine offendículo consciéntiam habére ad Deum et ad hómines semper.\par
\switchcolumn
\EN
\noindent Then Paul answered, (the governor making a sign to him to speak:) Knowing that for many years thou hast been judge over this nation, I will with good courage answer for myself. For thou mayest understand, that there are yet but twelve days, since I went up to adore in Jerusalem: And neither in the temple did they find me disputing with any man, or causing any concourse of the people, neither in the synagogues, nor in the city: Neither can they prove unto thee the things whereof they now accuse me. But this I confess to thee, that according to the way, which they call a heresy, so do I serve the Father and my God, believing all things which are written in the law and the prophets: Having hope in God, which these also themselves look for, that there shall be a resurrection of the just and unjust. And herein do I endeavour to have always a conscience without offence toward God, and towards men.\par
\switchcolumn*

\LA
\begin{Resp}\Rsign Virtúte magna reddébant Apóstoli, \Star{} Testimónium resurrectiónis Jesu Christi Dómini nostri, allelúja, allelúja.\end{Resp}
\begin{Resp}\Vsign Repléti quidem Spíritu Sancto loquebántur cum fidúcia verbum Dei.\end{Resp}
\begin{Resp}\Rsign Testimónium resurrectiónis Jesu Christi Dómini nostri, allelúja, allelúja.\end{Resp}
\switchcolumn
\EN
\begin{Resp}\Rsign With great power gave the Apostles \Star{} Witness of the Resurrection of our Lord Jesus Christ, alleluia, alleluia.\end{Resp}
\begin{Resp}\Vsign They were all filled with the Holy Ghost, and they spoke the Word of God with boldness.\end{Resp}
\begin{Resp}\Rsign Witness of the Resurrection of our Lord Jesus Christ, alleluia, alleluia.\end{Resp}
\switchcolumn*

\LA
\LessonHead{Lectio II}
\switchcolumn
\EN
\LessonHead{Lesson II}
\switchcolumn*

\LA
\Cite{Act. 24:17-21}
\switchcolumn
\EN
\Cite{Acts 24:17-21}
\switchcolumn*

\LA
\noindent Post annos autem plures eleemósynas factúrus in gentem meam, veni, et oblatiónes, et vota, in quibus invenérunt me purificátum in templo: non cum turba, neque cum tumúltu. Quidam autem ex Asia Judǽi, quos oportébat apud te præsto esse, et accusáre si quid habérent advérsum me: aut hi ipsi dicant si quid invenérunt in me iniquitátis cum stem in concílio, nisi de una hac solúmmodo voce qua clamávi inter eos stans: Quóniam de resurrectióne mortuórum ego júdicor hódie a vobis.\par
\switchcolumn
\EN
\noindent Now after many years, I came to bring alms to my nation, and offerings, and vows. In which I was found purified in the temple: neither with multitude, nor with tumult. But certain Jews of Asia, who ought to be present before thee, and to accuse, if they had any thing against me: Or let these men themselves say, if they found in me any iniquity, when standing before the council, Except it be for this one voice only that I cried, standing among them, Concerning the resurrection of the dead am I judged this day by you.\par
\switchcolumn*

\LA
\begin{Resp}\Rsign De ore prudéntis procédit mel, allelúja: dulcédo mellis est sub língua ejus, allelúja: \Star{} Favus distíllans lábia ejus, allelúja, allelúja.\end{Resp}
\begin{Resp}\Vsign Sapiéntia requiéscit in corde ejus, et prudéntia in sermóne oris illíus.\end{Resp}
\begin{Resp}\Rsign Favus distíllans lábia ejus, allelúja, allelúja.\end{Resp}
\switchcolumn
\EN
\begin{Resp}\Rsign From the mouth of the wise doth proceed honey, alleluia: the sweetness of honey is under his tongue, alleluia. \Star{} His lips drop as the honey-comb, alleluia, alleluia.\end{Resp}
\begin{Resp}\Vsign Wisdom doth abide in his heart, and out of his mouth cometh understanding.\end{Resp}
\begin{Resp}\Rsign His lips drop as the honey-comb, alleluia, alleluia.\end{Resp}
\switchcolumn*

\LA
\LessonHead{Lectio III}
\switchcolumn
\EN
\LessonHead{Lesson III}
\switchcolumn*

\LA
\Cite{Act. 24:22-27}
\switchcolumn
\EN
\Cite{Acts 24:22-27}
\switchcolumn*

\LA
\noindent Dístulit autem illos Felix, certíssime sciens de via hac, dicens: Cum tribúnus Lýsias descénderit, áudiam vos. Jussítque centurióni custodíre eum, et habére réquiem, nec quemquam de suis prohibére ministráre ei. Post áliquot autem dies véniens Felix cum Drusílla uxóre sua, quæ erat Judǽa, vocávit Paulum, et audívit ab eo fidem quæ est in Christum Jesum. Disputánte autem illo de justítia, et castitáte, et de judício futúro, tremefáctus Felix, respóndit: Quod nunc áttinet, vade: témpore autem opportúno accérsam te: simul et sperans quod pecúnia ei darétur a Paulo, propter quod et frequénter accérsens eum, loquebátur cum eo. Biénnio autem expléto, accépit successórem Felix Pórtium Festum. Volens autem grátiam præstáre Judǽis Felix, relíquit Paulum vinctum.\par
\switchcolumn
\EN
\noindent And Felix put them off, having most certain knowledge of this way, saying: When Lysias the tribune shall come down, I will hear you. And he commanded a centurion to keep him, and that he should be easy, and that he should not prohibit any of his friends to minister unto him. And after some days, Felix, coming with Drusilla his wife, who was a Jew, sent for Paul, and heard of him the faith, that is in Christ Jesus. And as he treated of justice, and chastity, and of the judgment to come, Felix being terrified, answered: For this time, go thy way: but when I have a convenient time, I will send for thee. Hoping also withal, that money should be given him by Paul; for which cause also oftentimes sending for him, he spoke with him. But when two years were ended, Felix had for successor Portius Festus. And Felix being willing to show the Jews a pleasure, left Paul bound.\par
\switchcolumn*

\LA
\begin{Resp}\Rsign Ecce vicit leo de tribu Juda, radix David, aperíre librum, et sólvere septem signácula ejus: \Star{} Allelúja, allelúja, allelúja.\end{Resp}
\begin{Resp}\Vsign Dignus est Agnus, qui occísus est, accípere virtútem, et divinitátem, et sapiéntiam, et fortitúdinem, et honórem, et glóriam, et benedictiónem.\end{Resp}
\begin{Resp}\Rsign Allelúja, allelúja, allelúja.\end{Resp}
\begin{Resp}\Vsign Glória Patri, et Fílio, \Star{} et Spirítui Sancto.\end{Resp}
\begin{Resp}\Rsign Allelúja, allelúja, allelúja.\end{Resp}
\switchcolumn
\EN
\begin{Resp}\Rsign Behold, the Lion of the tribe of Judah, the Root of David, hath prevailed to open the Book, and to loose the seven seals thereof. \Star{} Alleluia, alleluia, alleluia.\end{Resp}
\begin{Resp}\Vsign Worthy is the Lamb That was slain to receive power, and riches, in wisdom, and strength, and honour, and glory, and blessing.\end{Resp}
\begin{Resp}\Rsign Alleluia, alleluia, alleluia.\end{Resp}
\begin{Resp}\Vsign Glory be to the Father, and to the Son, \Star{} and to the Holy Ghost.\end{Resp}
\begin{Resp}\Rsign Alleluia, alleluia, alleluia.\end{Resp}
\switchcolumn*

\LA
\LessonHead{Lectio IV}
\switchcolumn
\EN
\LessonHead{Lesson IV}
\switchcolumn*

\LA
\LessonTitle{Sermo sancti Bernardíni Senénsis}
\switchcolumn
\EN
\LessonTitle{From a Sermon by St. Bernardine of Siena}
\switchcolumn*

\LA
\Source{Sermo de S. Joseph}
\switchcolumn
\EN
\Source{Sermo de S. Joseph}
\switchcolumn*

\LA
\noindent Cum inter Maríam et Joseph fúerit veríssimum matrimónium per divínam inspiratiónem contráctum, et in matrimónio fiat únio animórum in tantum quod una dicúntur persóna sponsus et sponsa, ut possit dici únitas quasi summa; quómodo cogitáre potest mens discréta quod Spíritus Sanctus tanta unióne uníret menti tantæ Vírginis áliquam ánimam, nisi ei virtútum operatióne simíllimam? Unde credo, istum virum sanctum Joseph fuísse mundíssimum in virginitáte, profundíssimum in humilitáte, ardentíssimum in Dei amóre et caritáte, altíssimum in contemplatióne. Et quia nóverat Virgo istum sibi a Spíritu Sancto datum in sponsum, et in suæ virginitátis fidum custódem, et ad participándum secum in caritátis amóre et obsequiósa sollicitúdine erga diviníssimam Prolem Dei; ídeo credo, quod totíus cordis afféctu hunc, sanctum Joseph, sinceríssime diligébat.\par
\switchcolumn
\EN
\noindent The marriage between Mary and Joseph was a real marriage, for it was contracted under divine inspiration. Now in marriage there is so close a union of souls that the bridegroom and the bride are said to be one person, for which reason marriage is like unto the very perfection of unity. Hence how can any discerning mind think that the Holy Spirit would unite, in a union of this intimacy, a mind such as the Virgin's, with the soul of a man who had not within him the operation of a godliness like unto hers? Wherefore I believe that this Joseph was holy, the chastest of men and a virgin, completely humble, burning with a passion of charity towards God, and full of the highest graces of contemplation. And since the Virgin knew that he was given her by the Holy Spirit to be her spouse, and the faithful guardian of her virginity, and to share besides in devoted love and affectionate care towards that One who was in the divinest fashion the very offspring of God; therefore I believe that she sincerely loved Saint Joseph with all the affection of her heart.\par
\switchcolumn*

\LA
\begin{Resp}\Rsign Dedísti mihi protectiónem salútis tuæ, et déxtera tua suscépit me: \Star{} Protéctor meus, et cornu salútis meæ, et suscéptor meus, allelúja.\end{Resp}
\begin{Resp}\Vsign Ego protéctor tuus sum, et merces tua magna nimis.\end{Resp}
\begin{Resp}\Rsign Protéctor meus, et cornu salútis meæ, et suscéptor meus, allelúja.\end{Resp}
\switchcolumn
\EN
\begin{Resp}\Rsign Thou hast given me the shield of thy salvation, and thy right hand hath holden me up. \Star{} My buckler, and the horn of my salvation, and my refuge. Alleluia.\end{Resp}
\begin{Resp}\Vsign I am thy shield and thy exceeding great reward.\end{Resp}
\begin{Resp}\Rsign My buckler, and the horn of my salvation, and my refuge. Alleluia.\end{Resp}
\switchcolumn*

\LA
\LessonHead{Lectio V}
\switchcolumn
\EN
\LessonHead{Lesson V}
\switchcolumn*

\LA
\noindent Hábuit Joseph erga Christum ardentíssimam caritátem. Quis déneget, óbsecro, quod ipsi tenénti Christum in brácchiis aut confabulánti cum ipso, Christus, sive infans sive adúltus, ingéreret et imprímeret ineffábiles sensus atque jucunditátes de semetípso, et hoc cooperánte extérius grátia Christi cum filiáli aspéctu, affátu atque compléxu? O quanta dúlcia óscula ab ipso recépit! o quanta dulcédine audiébat balbutiéntem Párvulum se patrem vocáre, et quanta suavitáte sentiébat se dúlciter amplexári! Consídera étiam, cum quanta compassióne in itinéribus, quæ fecérunt, párvulum Jesum ex labóre lassum, cum grandiúsculus esset, in suo grémio requiéscere faciébat: quia omni amóre transformatívo ferebátur in eum, ut in dulcíssimum Fílium sibi in cónjuge sua Vírgine per Spíritum Sanctum datum.\par
\switchcolumn
\EN
\noindent Now Joseph was most ardent in his love for Christ. For who, pray, would deny that Christ, whether as a child or as a grown man, would most deeply inspire ineffable affection, and the peculiar joys which he alone could give? And what would be the effect on one who held him in his arms, and conversed at will with him? And besides all this, who can reckon the bliss of receiving from the Christ Child those gazes of filial love? or his words spoken as a devoted son? or the giving of his trustful embraces? O how sweet were the kisses that Joseph received from him! O how sweet to hear little One lisp the name of father, and how delightful to feel his gentle carresses! Think again how often (when the little Jesus was growing bigger, and was wearied with much walking on the journeys which they made) Joseph must have been filled with compassion, and so carried him at rest in his bosom. For Joseph bore towards Jesus all the fulness of an adoptive love, as to a most dear son, given to him by the Holy Ghost in his Virgin bride.\par
\switchcolumn*

\LA
\begin{Resp}\Rsign Státuet fílios suos sub tégmine illíus, et sub ramis ejus morábitur: protegétur sub tégmine illíus a fervóre: \Star{} Et in glória ejus requiéscet, allelúja.\end{Resp}
\begin{Resp}\Vsign Speráte in eo, omnis congregátio pópuli, effúndite coram illo corda vestra.\end{Resp}
\begin{Resp}\Rsign Et in glória ejus requiéscet, allelúja.\end{Resp}
\switchcolumn
\EN
\begin{Resp}\Rsign He shall set his children under her shelter, and shall lodge under her branches by her shall he be covered from heat, \Star{} And in her glory shall he dwell. Alleluia.\end{Resp}
\begin{Resp}\Vsign Trust in Him, ye congregation of the people, pour out your heart before Him.\end{Resp}
\begin{Resp}\Rsign And in her glory shall he dwell. Alleluia.\end{Resp}
\switchcolumn*

\LA
\LessonHead{Lectio VI}
\switchcolumn
\EN
\LessonHead{Lesson VI}
\switchcolumn*

\LA
\noindent Ideo prudentíssima Mater, quæ expérta fúerat ejus afféctum, ad Fílium suum Jesum in templo reinvéntum ait: Fili, quid fecísti nobis sic? Ecce pater tuus et ego doléntes quærebámus te. Ad hujus verbi intelléctum notándum est, quod duo sapórum génera cóntinet in se Christus, dulcóris et dolóris; et quia sanctíssimus Joseph horum duórum gústuum mirabíliter párticeps fuit, ídeo beáta Virgo vocat eum singuláriter patrem Christi. Hic solum légitur Vírginem Joseph appellásse patrem Jesu: quia sensus dolóris, quem hábuit de Jesu pérdito, verum in eo monstrávit patris afféctum. Si enim secúndum humánas leges divínitus approbátas potest extráneus áliquem adoptáre in fílium, multo magis Dei Fílius datus ipsi Joseph in sua sanctíssima Sponsa sub virginális matrimónii admirábili sacraménto, debet ejus fílius appellári; et étiam credi quod in eo fúerit gustus paternális amóris atque dolóris respéctu dilécti Jesu.\par
\switchcolumn
\EN
\noindent Hence it was that a most prudent Mother, who knew the devotion of Joseph to Jesus, said to her Son, when she found him in the temple: Son, why hast thou thus dealt with us? behold, thy father and I have sought thee sorrowing. In order to understand this, we must note that Christ hath within himself, as it were, two savours, sweetness and bitterness. And since the most holy Joseph was in a wonderful manner (as we shall see) a partaker of these two savours, therefore the Blessed Virgin doth bestow upon him in a special sense the title of Father of Christ. This is the only place where we read that she did call Saint Joseph the father of Jesus, doubtless because the bitterness of sorrow which he felt at the loss of Jesus showed the fatherly affection which he bore him. For if according to human laws, which are approved by God, a man can adopt as his son the child of another family, how much more truly ought the Son of God to be called the Son of Joseph. For he was given to this Joseph by his most holy Spouse, in the wonderful mystery of a virginal marriage. And so it is also to be believed that in Joseph there were the two savours of Jesus, sweetness and bitterness, which were manifested as the sweetness of paternal love, and the bitterness of his compassion, towards his beloved Jesus.\par
\switchcolumn*

\LA
\begin{Resp}\Rsign Si consístant advérsum me castra, non timébit cor meum: \Star{} Si exsúrgat advérsum me prǽlium, in hoc ego sperábo, allelúja.\end{Resp}
\begin{Resp}\Vsign In te cantátio mea semper, quóniam tu adjútor fortis.\end{Resp}
\begin{Resp}\Rsign Si exsúrgat advérsum me prǽlium, in hoc ego sperábo, allelúja.\end{Resp}
\begin{Resp}\Vsign Glória Patri, et Fílio, \Star{} et Spirítui Sancto.\end{Resp}
\begin{Resp}\Rsign Si exsúrgat advérsum me prǽlium, in hoc ego sperábo, allelúja.\end{Resp}
\switchcolumn
\EN
\begin{Resp}\Rsign Though an host should encamp against me, my heart shall not fear. \Star{} Though war should rise against me, in this will I be confident. Alleluia.\end{Resp}
\begin{Resp}\Vsign My praise shall be continually of thee, for Thou art my strong refuge. R.* Though war should rise against me, in this will I be confident. Alleluia.\end{Resp}
\begin{Resp}\Vsign Glory be to the Father, and to the Son, \Star{} and to the Holy Ghost.\end{Resp}
\begin{Resp}\Rsign Though war should rise against me, in this will I be confident. Alleluia.\end{Resp}
\switchcolumn*

\LA
\LessonHead{Lectio VII}
\switchcolumn
\EN
\LessonHead{Lesson VII}
\switchcolumn*

\LA
\LessonTitle{Léctio sancti Evangélii secúndum Lucam}
\switchcolumn
\EN
\LessonTitle{From the Holy Gospel according to Luke}
\switchcolumn*

\LA
\Cite{Luc 3:21-23}
\switchcolumn
\EN
\Cite{Luke 3:21-23}
\switchcolumn*

\LA
\noindent In illo témpore: Factum est autem cum baptizarétur omnis pópulus, et Jesu baptizáto et oránte, apértum est cælum. Et réliqua.\par
\switchcolumn
\EN
\noindent At that time: When all the people were baptized, it came to pass, that Jesus also being baptized and praying, the heaven was opened. And so on, and that which followeth.\par
\switchcolumn*

\LA
\LessonTitle{De Homilía sancti Augustíni Epíscopi}
\switchcolumn
\EN
\LessonTitle{A Homily by St. Augustine the Bishop}
\switchcolumn*

\LA
\Source{Liber 2 de Consensu Evang.}
\switchcolumn
\EN
\Source{Liber 2 de Consensu Evang.}
\switchcolumn*

\LA
\noindent Neque proptérea non erat appellándus Joseph pater Christi, quia non eum concumbéndo genúerat; quandóquidem pater esset étiam ejus, quem non ex sua cónjuge procreátum aliúnde adoptásset. Putabátur quidem Christus étiam áliter fílius Joseph, tamquam ex ejus omníno carne progénitus; sed ab eis hoc putabátur, quos Maríæ latébat virgínitas: nam Lucas ait: Et ipse Jesus erat incípiens quasi annórum trigínta, ut putabátur, fílius Joseph. Qui tamen Lucas non ejus paréntem solam Maríam, sed ambos paréntes ejus appelláre mínime dubitávit, ubi ait: Puer autem crescébat et confortabátur plenus sapiéntia, et grátia Dei erat in illo; et ibant paréntes ejus per omnes annos in Jerúsalem, in die solémni Paschæ.\par
\switchcolumn
\EN
\noindent Joseph cannot be denied the name of father of Christ, merely because he did not beget him by coition. For he would have been called the father of any child whom he adopted, even if the child were not the issue of his wife, but from another family. It is true that Christ was supposed to be the son of Joseph in another sense; namely, in that of having been actually begotten by Joseph according to the flesh. But this supposition was made only by those from whom Mary's virginity was concealed. It is of this that Luke saith: And Jesus himself began to be about thirty years of age, being (as was supposed) the son of Joseph. However, Luke sheweth no hesitation in giving the name of parent, not to Mary only, but also to Joseph, when in another place he saith: And the child grew, and waxed strong in spirit, filled with wisdom: and the grace of God was upon him: Now his parents went every year to Jerusalem, at the Feast of Passover.\par
\switchcolumn*

\LA
\begin{Resp}\Rsign Joseph, fili David, noli timére accípere Maríam cónjugem tuam; quod enim in ea natum est, de Spíritu Sancto est: páriet autem fílium, \Star{} Et vocábis nomen ejus Jesum, allelúja.\end{Resp}
\begin{Resp}\Vsign Ipse enim salvum fáciet pópulum suum a peccátis eórum.\end{Resp}
\begin{Resp}\Rsign Et vocábis nomen ejus Jesum, allelúja.\end{Resp}
\switchcolumn
\EN
\begin{Resp}\Rsign Joseph, thou Son of David, fear not to take unto thee Mary thy wife; for That Which is conceived in her is of the Holy Ghost and she shall bring forth a Son; \Star{} And thou shalt call His Name Jesus. Alleluia.\end{Resp}
\begin{Resp}\Vsign For He shall save His people from their sins.\end{Resp}
\begin{Resp}\Rsign And thou shalt call His Name Jesus. Alleluia.\end{Resp}
\switchcolumn*

\LA
\LessonHead{Lectio VIII}
\switchcolumn
\EN
\LessonHead{Lesson VIII}
\switchcolumn*

\LA
\noindent Sed, ne quisnam hic paréntes consanguíneos pótius Maríæ cum ipsa Matre ejus intelligéndos putet, quid ad illud respondébit quod ipse item Lucas supérius dixit, Et erant Pater ejus et Mater mirántes super his, quæ dicebántur de illo? Cum ígitur ipse narret, non ex concúbitu Joseph, sed ex María Vírgine natum Christum; unde eum patrem ejus appéllat, nisi quia et virum Maríæ recte intellígimus sine commixtióne carnis, ipsa copulatióne conjúgii; et ob hoc étiam Christi patrem multo conjúnctius, qui ex ejus cónjuge natus sit, quam si esset aliúnde adoptátus? Ac per hoc, si demonstráre áliquis posset Maríam ex David nullam consanguinitátis oríginem dúcere, sat erat secúndum istam ratiónem, accípere Christum Fílium David; qua ratióne étiam Joseph pater ejus appellátus est.\par
\switchcolumn
\EN
\noindent But lest anyone should think that by the word Parents there is here to understood Mary and her forbears only, we must take into account what Luke recordeth earlier, And Joseph and his mother marvelled at those things which were spoken of him. Since, therefore, Luke witnesseth that Christ was born, not by the begetting of Joseph, but of the Virgin Mary, how can he call Joseph the father of Jesus, except in the sense that Joseph was a real husband to Mary by virtue of the true bond of marriage, saving only that there never was any carnal intercourse between them? And yet, on account of this bond of marriage, Joseph was the father of Jesus in a much closer sense (seeing that the Christ Child was born of his wife) than if Joseph had adopted Jesus from another family. Hence, also, if anyone could prove that Mary did not trace her origin from David, the same reasoning by which Joseph is called the father of Christ would be sufficient reason for giving Christ the name, Son of David.\par
\switchcolumn*

\LA
\begin{Resp}\Rsign Surge, et áccipe Púerum et matrem ejus, et fuge in Ægýptum: \Star{} Et esto ibi usque dum dicam tibi, allelúja.\end{Resp}
\begin{Resp}\Vsign Ut adimplerétur quod dictum est a Dómino per Prophétam dicéntem: Ex Ægýpto vocávi Fílium meum.\end{Resp}
\begin{Resp}\Rsign Et esto ibi usque dum dicam tibi, allelúja.\end{Resp}
\begin{Resp}\Vsign Glória Patri, et Fílio, \Star{} et Spirítui Sancto.\end{Resp}
\begin{Resp}\Rsign Et esto ibi usque dum dicam tibi, allelúja.\end{Resp}
\switchcolumn
\EN
\begin{Resp}\Rsign Arise, and take the young Child, and His Mother, and flee into Egypt; \Star{} And be thou there until I bring thee word. Alleluia.\end{Resp}
\begin{Resp}\Vsign That it might be fulfilled which was spoken of the Lord by the Prophets, saying Out of Egypt have I called My Son.\end{Resp}
\begin{Resp}\Rsign And be thou there until I bring thee word. Alleluia.\end{Resp}
\begin{Resp}\Vsign Glory be to the Father, and to the Son, \Star{} and to the Holy Ghost.\end{Resp}
\begin{Resp}\Rsign And be thou there until I bring thee word. Alleluia.\end{Resp}
\switchcolumn*

\LA
\LessonHead{Lectio IX}
\switchcolumn
\EN
\LessonHead{Lesson IX}
\switchcolumn*

\LA
\noindent Lucas autem non ab inítio Evangélii sui, sed a baptísmo Christi generatiónes enárrat, nec descendéndo, sed ascendéndo, tamquam sacerdótem in expiándis peccátis magis assígnans; ubi eum vox de cælo declarávit, ubi testimónium Joánnes ipsi perhíbuit dicens: Ecce, qui tollit peccáta mundi. Ascendéndo autem transit Abraham et pérvenit ad Deum, cui mundáti et expiáti reconciliámur. Mérito et adoptiónis oríginem ipse suscépit, quia per adoptiónem effícimur fílii Dei, credéndo in Fílium Dei. Satis autem demonstrávit, non se ídeo dixísse. Joseph fílium Heli, quod de illo génitus, sed quod ab illo fúerat adoptátus; cum et ipsum Adam fílium Dei dixit, cum sit factus a Deo, sed per grátiam, quam póstea peccándo amísit, tamquam fílius in paradíso constitútus sit.\par
\switchcolumn
\EN
\noindent Luke giveth the genealogy, not at the beginning of his Gospel, but after the account of the Baptism of Christ. And he giveth it, not in the descending order, but in the ascending, more as if he were pointing to Christ as Priest, making atonement for sins. This was the occasion when the voice spake in testimony from heaven. And also at this time John himself gave testimony, saying: Behold the Lamb of God, which taketh away the sin of the world. By thus beginning with Jesus and tracing back, Luke passeth up through Abraham and eventually cometh to God, to whom we are reconciled after purification and atonement. Rightly then doth Luke give the origin by adoption, for through adoption and faith in the Son of God we become God's sons. In this fashion, Luke sheweth clearly enough why he nameth Joseph as the son of Heli; that is, not because Joseph was begotten by Heli, but because he was adopted by him. For Luke calleth Adam himself the son of God, and this because he was made by God, being set in the paradise of Eden as a son by virtue of the grace which afterwards he lost in sinning.\par
\switchcolumn*

\LA
\TeDeum{Te Deum}
\switchcolumn
\EN
\TeDeum{Te Deum}
\switchcolumn*

\end{paracol}

\Office{De III die infra Octavam S. Joseph}{Third Day within the Octave of St. Joseph}{Semiduplex}
\addcontentsline{toc}{subsection}{De III die infra Octavam S. Joseph\enspace\textit{Third Day within the Octave of St. Joseph}}
\begin{paracol}{2}
\LA
\LessonHead{Lectio I}
\switchcolumn
\EN
\LessonHead{Lesson I}
\switchcolumn*

\LA
\LessonTitle{De Actibus Apostolórum}
\switchcolumn
\EN
\LessonTitle{Lesson from the Acts of the Apostles}
\switchcolumn*

\LA
\Cite{Act. 25:1-5}
\switchcolumn
\EN
\Cite{Acts 25:1-5}
\switchcolumn*

\LA
\noindent Festus ergo cum venísset in provínciam, post tríduum ascéndit Jerosólymam a Cæsaréa. Adierúntque eum príncipes sacerdótum et primi Judæórum advérsus Paulum: et rogábant eum, postulántes grátiam advérsus eum, ut jubéret perdúci eum in Jerúsalem, insídias tendéntes ut interfícerent eum in via. Festus autem respóndit servári Paulum in Cæsaréa: se autem matúrius profectúrum. Qui ergo in vobis, ait, poténtes sunt, descendéntes simul, si quod est in viro crimen, accúsent eum.\par
\switchcolumn
\EN
\noindent Now when Festus was come into the province, after three days, he went up to Jerusalem from Caesarea. And the chief priests, and principal men of the Jews, went unto him against Paul: and they besought him, Requesting favour against him, that he would command him to be brought to Jerusalem, laying wait to kill him in the way. But Festus answered: That Paul was kept in Caesarea, and that he himself would very shortly depart thither. Let them, therefore, saith he, among you that are able, go down with me, and accuse him, if there be any crime in the man.\par
\switchcolumn*

\LA
\begin{Resp}\Rsign Ego sum vitis vera, et vos pálmites: \Star{} Qui manet in me, et ego in eo, hic fert fructum multum, allelúja, allelúja.\end{Resp}
\begin{Resp}\Vsign Sicut diléxit me Pater, et ego diléxi vos.\end{Resp}
\begin{Resp}\Rsign Qui manet in me, et ego in eo, hic fert fructum multum, allelúja, allelúja.\end{Resp}
\switchcolumn
\EN
\begin{Resp}\Rsign I am the True Vine, and ye are the branches; \Star{} He that abideth in Me, and I in him, the same bringeth forth much fruit, alleluia, alleluia.\end{Resp}
\begin{Resp}\Vsign As the Father hath loved Me, so have I loved you.\end{Resp}
\begin{Resp}\Rsign He that abideth in Me, and I in him, the same bringeth forth much fruit, alleluia, alleluia.\end{Resp}
\switchcolumn*

\LA
\LessonHead{Lectio II}
\switchcolumn
\EN
\LessonHead{Lesson II}
\switchcolumn*

\LA
\Cite{Act. 25:6-8}
\switchcolumn
\EN
\Cite{Acts 25:6-8}
\switchcolumn*

\LA
\noindent Demorátus autem inter eos dies non ámplius quam octo aut decem, descéndit Cæsaréam, et áltera die sedit pro tribunáli, et jussit Paulum addúci. Qui cum perdúctus esset, circumstetérunt eum, qui ab Jerosólyma descénderant Judǽi, multas et graves causas obiciéntes, quas non póterant probáre: Paulo ratiónem reddénte: Quóniam neque in legem Judæórum, neque in templum, neque in Cǽsarem quidquam peccávi.\par
\switchcolumn
\EN
\noindent And having tarried among them no more than eight or ten days, he went down to Caesarea, and the next day he sat in the judgment seat; and commanded Paul to be brought. Who being brought, the Jews stood about him, who were come down from Jerusalem, objecting many and grievous causes, which they could not prove; Paul making answer for himself: Neither against the law of the Jews, nor against the temple, nor against Caesar, have I offended in any thing.\par
\switchcolumn*

\LA
\begin{Resp}\Rsign Surgens Jesus Dóminus noster, stans in médio discipulórum suórum, dixit: \Star{} Pax vobis, allelúja: gavísi sunt discípuli viso Dómino, allelúja.\end{Resp}
\begin{Resp}\Vsign Una ergo sabbatórum, cum fores essent clausæ, ubi erant discípuli congregáti, venit Jesus, et stetit in medio eórum, et dixit eis.\end{Resp}
\begin{Resp}\Rsign Pax vobis, allelúja: gavísi sunt discípuli viso Dómino, allelúja.\end{Resp}
\switchcolumn
\EN
\begin{Resp}\Rsign After that our Lord Jesus was risen again, He came and stood in the midst of His disciples, and said unto them: \Star{} Peace be unto you, alleluia. Then were the disciples glad, when they saw the Lord, alleluia.\end{Resp}
\begin{Resp}\Vsign The first day of the week, when the doors were shut where the disciples were assembled, came Jesus, and stood in the midst, and said unto them:\end{Resp}
\begin{Resp}\Rsign Peace be unto you, alleluia. Then were the disciples glad, when they saw the Lord, alleluia.\end{Resp}
\switchcolumn*

\LA
\LessonHead{Lectio III}
\switchcolumn
\EN
\LessonHead{Lesson III}
\switchcolumn*

\LA
\Cite{Act. 25:9-12}
\switchcolumn
\EN
\Cite{Acts 25:9-12}
\switchcolumn*

\LA
\noindent Festus autem volens grátiam præstáre Judǽis, respóndens Paulo, dixit: Vis Jerosólymam ascéndere, et ibi de his judicári apud me? Dixit autem Paulus: Ad tribúnal Cǽsaris sto: ibi me opórtet judicári: Judǽis non nócui, sicut tu mélius nosti. Si enim nócui, aut dignum morte áliquid feci, non recúso mori: si vero nihil est eórum quæ hi accúsant me, nemo potest me illis donáre. Cǽsarem appéllo. Tunc Festus cum concílio locútus, respóndit: Cǽsarem appellásti? ad Cǽsarem ibis.\par
\switchcolumn
\EN
\noindent But Festus, willing to show the Jews a pleasure, answering Paul, said: Wilt thou go up to Jerusalem, and there be judged of these things before me? Then Paul said: I stand at Caesar's judgment seat, where I ought to be judged. To the Jews I have done no injury, as thou very well knowest. For if I have injured them, or have committed any thing worthy of death, I refuse not to die. But if there be none of these things whereof they accuse me, no man may deliver me to them: I appeal to Caesar. Then Festus having conferred with the council, answered: Hast thou appealed to Caesar? To Caesar shalt thou go.\par
\switchcolumn*

\LA
\begin{Resp}\Rsign Expurgáte vetus ferméntum, ut sitis nova conspérsio: étenim Pascha nostrum immolátus est Christus: \Star{} Itaque epulémur in Dómino, allelúja.\end{Resp}
\begin{Resp}\Vsign Mórtuus est propter delícta nostra, et resurréxit propter justificatiónem nostram.\end{Resp}
\begin{Resp}\Rsign Itaque epulémur in Dómino, allelúja.\end{Resp}
\begin{Resp}\Vsign Glória Patri, et Fílio, \Star{} et Spirítui Sancto.\end{Resp}
\begin{Resp}\Rsign Itaque epulémur in Dómino, allelúja.\end{Resp}
\switchcolumn
\EN
\begin{Resp}\Rsign Purge out the old leaven, that ye may be new dough: for even Christ our Passover is sacrificed for us: \Star{} Therefore let us keep the Feast, in the Lord, alleluia.\end{Resp}
\begin{Resp}\Vsign He died for our offences, and rose again for our justification.\end{Resp}
\begin{Resp}\Rsign Therefore let us keep the Feast, in the Lord, alleluia.\end{Resp}
\begin{Resp}\Vsign Glory be to the Father, and to the Son, \Star{} and to the Holy Ghost.\end{Resp}
\begin{Resp}\Rsign Therefore let us keep the Feast, in the Lord, alleluia.\end{Resp}
\switchcolumn*

\LA
\LessonHead{Lectio IV}
\switchcolumn
\EN
\LessonHead{Lesson IV}
\switchcolumn*

\LA
\LessonTitle{Sermo sancti Joánnis Chrysóstomi}
\switchcolumn
\EN

\switchcolumn*

\LA
\Source{Homilia 4 in Matth.}
\switchcolumn
\EN

\switchcolumn*

\LA
\noindent Hunc morem plerúmque tenébat antíquitas, ut sponsæ in sponsórum dómibus haberéntur: sic habitábat étiam María cum Sponso. Et cujus tandem rei grátia, non ántequam desponderétur, Virgo concépit? Ut vidélicet mystérium ínterim quasi obumbrátum latéret, et ut Virgo omnem prorsus occasiónem malígnæ suspiciónis effúgeret. Quando enim ille, qui præcípuo zelo posset ardére, cérnitur non solum non abjícere Sponsam, nec eam ignomínia notáre, sed étiam recípere in consórtium et inservíre post conceptiónem; profécto maniféstum est, quod, nisi apérte nosset ex operatióne Sancti Spíritum illum exstitísse concéptum, numquam vel apud se illam retinuísset, vel ei in ómnibus, quorum indíguit, ministrásset.\par
\switchcolumn
\EN
\noindent The Lesson is taken from a Sermon by St. John Chrysostom\par
Homilia 4 in Matth.\par
It was the custom in ancient times for betrothed brides to dwell in the houses of their bridegrooms. And it would seem that Mary thus dwelt with her Spouse. And herein we find answer to the question: Why did not the virgin conception take place before Mary was wedded? In order that the mystery might be hid in the meantime, and that the Virgin might escape all danger of evil suspicion. For Joseph had the best right to be moved by jealousy. Yet we see that he not only refrained from sending away his Espoused, or branding her with infamy, but that he received her as his own, and did cherish her after she conceived. And verily, it is evident that he would never have kept her in his house, or ministered to all her needs, if he had not clearly come to know that she had conceived by the operation of the Holy Ghost.\par
\switchcolumn*

\LA
\begin{Resp}\Rsign Dedísti mihi protectiónem salútis tuæ, et déxtera tua suscépit me: \Star{} Protéctor meus, et cornu salútis meæ, et suscéptor meus, allelúja.\end{Resp}
\begin{Resp}\Vsign Ego protéctor tuus sum, et merces tua magna nimis.\end{Resp}
\begin{Resp}\Rsign Protéctor meus, et cornu salútis meæ, et suscéptor meus, allelúja.\end{Resp}
\switchcolumn
\EN
\begin{Resp}\Rsign Thou hast given me the shield of thy salvation, and thy right hand hath holden me up. \Star{} My buckler, and the horn of my salvation, and my refuge. Alleluia.\end{Resp}
\begin{Resp}\Vsign I am thy shield and thy exceeding great reward.\end{Resp}
\begin{Resp}\Rsign My buckler, and the horn of my salvation, and my refuge. Alleluia.\end{Resp}
\switchcolumn*

\LA
\LessonHead{Lectio V}
\switchcolumn
\EN
\LessonHead{Lesson V}
\switchcolumn*

\LA
\noindent Joseph autem, cum esset justus et nollet eam tradúcere, vóluit occúlte dimíttere eam. Postquam dixit, quod esset ex Spíritu Sancto et absque ulla mixtióne séxuum, sermónem suum étiam aliúnde confírmat. Ne enim áliquis díceret: Et unde hoc potest esse maniféstum? quis vidit? quis audívit aliquándo tale áliquid contigísse? neque putáres discípulum, quasi gratificátum magístro, ista confíngere; introdúcit Joseph ex his, quæ passus est, fidem dictis per cuncta faciéntem, ut plane hinc dícere Evangelísta videátur: Si non credis mihi et testimónium tibi meum forte suspéctum est, crede ígitur vel maríto.\par
\switchcolumn
\EN
\noindent Joseph being a just man, and not willing to make her a publick example, was minded to put her away privily. After the Evangelist hath told us that she was with child, of the Holy Ghost, and not of any sexual commerce, Luke bringeth testimony from another source to confirm the statement. For lest anyone should say: And how can this be proved? who saw it? who ever heard of any such thing having happened? and lest you might think that a disciple had invented this tale to please his Master, the Evangelist thus bringeth forward the grief of Joseph, and what he did thereafter, to confirm the story, as if to say: If ye will not believe me, or if ye hold my testimony in suspicion, at least believe the husband.\par
\switchcolumn*

\LA
\begin{Resp}\Rsign Státuet fílios suos sub tégmine illíus, et sub ramis ejus morábitur: protegétur sub tégmine illíus a fervóre: \Star{} Et in glória ejus requiéscet, allelúja.\end{Resp}
\begin{Resp}\Vsign Speráte in eo, omnis congregátio pópuli, effúndite coram illo corda vestra.\end{Resp}
\begin{Resp}\Rsign Et in glória ejus requiéscet, allelúja.\end{Resp}
\switchcolumn
\EN
\begin{Resp}\Rsign He shall set his children under her shelter, and shall lodge under her branches by her shall he be covered from heat, \Star{} And in her glory shall he dwell. Alleluia.\end{Resp}
\begin{Resp}\Vsign Trust in Him, ye congregation of the people, pour out your heart before Him.\end{Resp}
\begin{Resp}\Rsign And in her glory shall he dwell. Alleluia.\end{Resp}
\switchcolumn*

\LA
\LessonHead{Lectio VI}
\switchcolumn
\EN
\LessonHead{Lesson VI}
\switchcolumn*

\LA
\noindent Joseph enim vir ejus, inquit, cum esset justus. Justum hic, in omni virtúte dicit esse perféctum. Cum ígitur esset justus, hoc est, frugi bonúsque vir, vóluit occúlte dimíttere eam. Proptérea vero Evangelísta dixit quid justo illi accíderit ante notítiam, ut nequáquam de his, quæ post notítiam rei sunt facta, dubitáres. Et certe, si María talis fuísset qualem illam suspício fingébat, non modo publicári merúerat, verum étiam ex legis auctoritáte puníri: sed Joseph non tantum eam damnáre nóluit, sed nec publicáre quidem. Vidísti nempe virum sublímiter philosophántem et tyránnicæ illíus passiónis immúnem. Quamquam hic quæ tandem dicerétur esse suspício, ubi ipse úteri tumor videbátur factum argúere? Sed tamen ita erat ille vir ab hujúsmodi passióne mundus ac liber, ut ne in mínimis quidem Vírgini vellet inférre mæstítiam; et adhuc sub lege vivens, supra legem philosophátur: síquidem adventánte jam grátia, multo sublimióris disciplínæ documénta fulgére oportébat.\par
\switchcolumn
\EN
\noindent He saith: Joseph her husband, being a just man. To be just, as the word is here used, implieth that full growth of righteousness which cometh from the habitual service of God. Being therefore a just man (that is, a worthy and good man), he was minded to put her away privily. Thus the Evangelist recordeth the grief of this just man before he knew the secret of the virgin conception, lest thou shouldst have doubts concerning what happened after he knew that secret. And certainly if Mary had been such as suspicion would make her out to be, she would have deserved not only to be denounced, but to be punished by the authority of the Law. But Joseph was not only unwilling to condemn her, but even to denounce her. Herein thou dost see an instance of a man full of spiritual understanding, and free from the tyranny of suspicion. But was this a matter of mere suspicion, seeing that the very swelling of her body seemed to prove a fact? Nonetheless, this man was so pure, and free from that kind of jealousy, that he would not cause the Virgin even the slightest grief. And although he lived under the Law, his spiritual understanding was above the Law. For now that the reign of grace was approaching, it was fitting that there should be a shining example of a more sublime spirituality than was common under the Law.\par
\switchcolumn*

\LA
\begin{Resp}\Rsign Si consístant advérsum me castra, non timébit cor meum: \Star{} Si exsúrgat advérsum me prǽlium, in hoc ego sperábo, allelúja.\end{Resp}
\begin{Resp}\Vsign In te cantátio mea semper, quóniam tu adjútor fortis.\end{Resp}
\begin{Resp}\Rsign Si exsúrgat advérsum me prǽlium, in hoc ego sperábo, allelúja.\end{Resp}
\begin{Resp}\Vsign Glória Patri, et Fílio, \Star{} et Spirítui Sancto.\end{Resp}
\begin{Resp}\Rsign Si exsúrgat advérsum me prǽlium, in hoc ego sperábo, allelúja.\end{Resp}
\switchcolumn
\EN
\begin{Resp}\Rsign Though an host should encamp against me, my heart shall not fear. \Star{} Though war should rise against me, in this will I be confident. Alleluia.\end{Resp}
\begin{Resp}\Vsign My praise shall be continually of thee, for Thou art my strong refuge. R.* Though war should rise against me, in this will I be confident. Alleluia.\end{Resp}
\begin{Resp}\Vsign Glory be to the Father, and to the Son, \Star{} and to the Holy Ghost.\end{Resp}
\begin{Resp}\Rsign Though war should rise against me, in this will I be confident. Alleluia.\end{Resp}
\switchcolumn*

\LA
\LessonHead{Lectio VII}
\switchcolumn
\EN
\LessonHead{Lesson VII}
\switchcolumn*

\LA
\LessonTitle{Léctio sancti Evangélii secúndum Lucam}
\switchcolumn
\EN
\LessonTitle{From the Holy Gospel according to Luke}
\switchcolumn*

\LA
\Cite{Luc 3:21-23}
\switchcolumn
\EN
\Cite{Luke 3:21-23}
\switchcolumn*

\LA
\noindent In illo témpore: Factum est autem cum baptizarétur omnis pópulus, et Jesu baptizáto et oránte, apértum est cælum. Et réliqua.\par
\switchcolumn
\EN
\noindent At that time: When all the people were baptized, it came to pass, that Jesus also being baptized and praying, the heaven was opened. And so on, and that which followeth.\par
\switchcolumn*

\LA
\LessonTitle{De Homilía sancti Augustíni Epíscopi}
\switchcolumn
\EN
\LessonTitle{A Homily by St. Augustine the Bishop}
\switchcolumn*

\LA
\Source{Liber 23 contra Faust. cap. 7-8}
\switchcolumn
\EN
\Source{Liber 23 contra Faust. cap. 7-8}
\switchcolumn*

\LA
\noindent Sic de cælo dictum est super aquam Jordánis: Hic est Fílius meus diléctus, in quo mihi bene complácui; quemádmodum dictum est et in monte. Neque enim quia et ibi de cælo vox ipsa sónuit, Fílius Dei ante non fuit; quandóquidem ex útero Vírginis ille accépit formam servi, qui cum in forma Dei esset, non rapínam arbitrátus est esse æquális Deo. Dénique idem Apóstolus Paulus álio loco apertíssime dicit: Cum autem venit plenitúdo témporis, misit Deus Fílium suum factum ex mulíere, factum sub lege, ut eos, qui sub lege erant, redímeret, ut filiórum adoptiónem reciperémus. Ipse ergo est Fílius Dei, qui et Dóminus David secúndum divinitátem, et idem ipse fílius David ex sémine David secúndum carnem.\par
\switchcolumn
\EN
\noindent The words uttered from heaven over the river Jordan: This is my beloved Son, in whom I am well pleased: were said also on the Mount of the Transfigurátion. Now these words should not be understood to imply that he was not the Son of God before the voice from heaven was heard. For we know that he who received the form of a servant from the Virgin's womb, being in the form of God, thought it not robbery to be equal with God. Indeed, the same Apostle Paul clearly declareth elsewhere: But when the fulness of the time was come, God sent forth his Son, made of a Woman, made under the Law, that he might redeem them that were under the Law, that we might receive the adoption of sons. He therefore is the Son of God, who is according to his divinity the Lord of David, and at the same time according to the flesh the Son of David, and of the seed of David.\par
\switchcolumn*

\LA
\begin{Resp}\Rsign Joseph, fili David, noli timére accípere Maríam cónjugem tuam; quod enim in ea natum est, de Spíritu Sancto est: páriet autem fílium, \Star{} Et vocábis nomen ejus Jesum, allelúja.\end{Resp}
\begin{Resp}\Vsign Ipse enim salvum fáciet pópulum suum a peccátis eórum.\end{Resp}
\begin{Resp}\Rsign Et vocábis nomen ejus Jesum, allelúja.\end{Resp}
\switchcolumn
\EN
\begin{Resp}\Rsign Joseph, thou Son of David, fear not to take unto thee Mary thy wife; for That Which is conceived in her is of the Holy Ghost and she shall bring forth a Son; \Star{} And thou shalt call His Name Jesus. Alleluia.\end{Resp}
\begin{Resp}\Vsign For He shall save His people from their sins.\end{Resp}
\begin{Resp}\Rsign And thou shalt call His Name Jesus. Alleluia.\end{Resp}
\switchcolumn*

\LA
\LessonHead{Lectio VIII}
\switchcolumn
\EN
\LessonHead{Lesson VIII}
\switchcolumn*

\LA
\noindent Quod si nobis crédere non prodésset, non hoc tam atténte idem Apóstolus Timótheo commendáret dicens: Memor esto Christum Jesum resurrexísse a mórtuis, ex sémine David, secúndum Evangélium meum. Quid ergo jam móveat sancti Evangélii sectatórem, quod sine concúbitu Joseph Christus natus ex Vírgine, fílius tamen David appellátur, cum generatiónum sériem non usque ad Maríam sed usque ad Joseph Matthæus Evangelísta perdúcat? Primo quia maríti ejus fúerat propter virílem sexum pótius honoránda persóna; neque enim quia concúbitu non permíxtus, ídeo non marítus, cum ipse Matthæus narret ab Angelo Maríam cónjugem ipsíus appellátam, qui narrat, quod de Spíritu Sancto concéperat.\par
\switchcolumn
\EN
\noindent Except it were profitable to believe this, the same Apostle would not have exhorted Timothy so earnestly, saying: Remember that Jesus Christ, of the seed of David, is risen again from the dead, according to my Gospel. Why then should any follower of the holy Gospel be troubled concerning this; to wit, that Christ, who was born of the Virgin without cohabitation with Joseph, hath his line of descent traced by the Evangelist Matthew through Joseph rather than through Mary; or again, that because of Joseph's descent from David, the Evangelist should call Christ the Son of David? For we know good reasons for this. The first is that the genealogy of her husband would be preferred to hers as a matter of honour to the male sex. For even though he was not joined to her in cohabitation, he was not on that account any the less her husband, since Matthew himself witnesseth that Mary was called the wife of Joseph by the Angel, and yet he also saith that she had conceived by the Holy Ghost.\par
\switchcolumn*

\LA
\begin{Resp}\Rsign Surge, et áccipe Púerum et matrem ejus, et fuge in Ægýptum: \Star{} Et esto ibi usque dum dicam tibi, allelúja.\end{Resp}
\begin{Resp}\Vsign Ut adimplerétur quod dictum est a Dómino per Prophétam dicéntem: Ex Ægýpto vocávi Fílium meum.\end{Resp}
\begin{Resp}\Rsign Et esto ibi usque dum dicam tibi, allelúja.\end{Resp}
\begin{Resp}\Vsign Glória Patri, et Fílio, \Star{} et Spirítui Sancto.\end{Resp}
\begin{Resp}\Rsign Et esto ibi usque dum dicam tibi, allelúja.\end{Resp}
\switchcolumn
\EN
\begin{Resp}\Rsign Arise, and take the young Child, and His Mother, and flee into Egypt; \Star{} And be thou there until I bring thee word. Alleluia.\end{Resp}
\begin{Resp}\Vsign That it might be fulfilled which was spoken of the Lord by the Prophets, saying Out of Egypt have I called My Son.\end{Resp}
\begin{Resp}\Rsign And be thou there until I bring thee word. Alleluia.\end{Resp}
\begin{Resp}\Vsign Glory be to the Father, and to the Son, \Star{} and to the Holy Ghost.\end{Resp}
\begin{Resp}\Rsign And be thou there until I bring thee word. Alleluia.\end{Resp}
\switchcolumn*

\LA
\LessonHead{Lectio IX}
\switchcolumn
\EN
\LessonHead{Lesson IX}
\switchcolumn*

\LA
\noindent Cum vero unus idémque narrátor utrúmque dicat, utrúmque comméndet, et virum Maríæ Joseph et Christi Vírginem Matrem, et Christum ex sémine David et Joseph in série progeneratórum Christi ex David; quid restat, nisi et Maríam non fuísse extráneam a cognatióne David, et eam Joseph cónjugem non frustra appellátam propter órdinem sexus et animórum confœderatiónem; et Joseph pótius propter dignitátem virílem ab órdine generatiónum illárum non fuísse separándum, ne hoc ipso viderétur ab illa fémina separátus, cui eum conjungébat mentis afféctus?\par
\switchcolumn
\EN
\noindent Note that one and the same narrátor both maketh and approveth all these statements: to wit, that Joseph was the husband of Mary; that the Mother of Christ was a Virgin; that Christ was of the seed of David; and that Joseph was in the line of Christ's ancestors from David. From all this we perceive several other reasons for the giving of Joseph's genealogy: to wit, that Mary was not without blood-relationship to David; and that in consideration of their union as souls, according to the due order of sex, she was not falsely given the title of Joseph's wife; and that Joseph, chiefly on account of his dignity as a man, was not to be separated from the line of their common genealogy, lest he should appear as separated from that Woman, to whom the affection of his soul bound him.\par
\switchcolumn*

\LA
\TeDeum{Te Deum}
\switchcolumn
\EN
\TeDeum{Te Deum}
\switchcolumn*

\end{paracol}

\Office{De IV die infra Octavam S. Joseph}{Fourth Day within the Octave of St. Joseph}{Semiduplex}
\addcontentsline{toc}{subsection}{De IV die infra Octavam S. Joseph\enspace\textit{Fourth Day within the Octave of St. Joseph}}
\begin{paracol}{2}
\LA
\LessonHead{Lectio I}
\switchcolumn
\EN
\LessonHead{Lesson I}
\switchcolumn*

\LA
\LessonTitle{De Actibus Apostolórum}
\switchcolumn
\EN
\LessonTitle{Lesson from the Acts of the Apostles}
\switchcolumn*

\LA
\Cite{Act. 28:16-20}
\switchcolumn
\EN
\Cite{Acts 28:16-20}
\switchcolumn*

\LA
\noindent Cum autem venissémus Romam, permíssum est Paulo manére síbimet cum custodiénte se mílite. Post tértium autem diem convocávit primos Judæórum. Cumque conveníssent, dicébat eis: Ego, viri fratres, nihil advérsus plebem fáciens, aut morem patérnum, vinctus ab Jerosólymis tráditus sum in manus Romanórum, qui cum interrogatiónem de me habuíssent, voluérunt me dimíttere, eo quod nulla esset causa mortis in me. Contradicéntibus autem Judǽis, coáctus sum appelláre Cǽsarem, non quasi gentem meam habens áliquid accusáre. Propter hanc ígitur causam rogávi vos vidére, et álloqui. Propter spem enim Israël caténa hac circúmdatus sum.\par
\switchcolumn
\EN
\noindent And when we were come to Rome, Paul was suffered to dwell by himself, with a soldier that kept him. And after the third day, he called together the chief of the Jews. And when they were assembled, he said to them: Men, brethren, I, having done nothing against the people, or the custom of our fathers, was delivered prisoner from Jerusalem into the hands of the Romans; Who, when they had examined me, would have released me, for that there was no cause of death in me; But the Jews contradicting it, I was constrained to appeal unto Caesar; not that I had any thing to accuse my nation of. For this cause therefore I desired to see you, and to speak to you. Because that for the hope of Israel, I am bound with this chain.\par
\switchcolumn*

\LA
\begin{Resp}\Rsign Christus resúrgens ex mórtuis, jam non móritur, mors illi ultra non dominábitur: quod enim mórtuus est peccáto, mórtuus est semel: \Star{} Quod autem vivit, vivit Deo, allelúja, allelúja.\end{Resp}
\begin{Resp}\Vsign Mórtuus est semel propter delícta nostra, et resurréxit propter justificatiónem nostram.\end{Resp}
\begin{Resp}\Rsign Quod autem vivit, vivit Deo, allelúja, allelúja.\end{Resp}
\switchcolumn
\EN
\begin{Resp}\Rsign Knowing that Christ rising again from the dead, dieth now no more, death shall no more have dominion over him. For in that he died to sin, he died once; \Star{} In that he liveth, he liveth unto God, alleluia, alleluia.\end{Resp}
\begin{Resp}\Vsign He died once for the sins, and he rose for our justification.\end{Resp}
\begin{Resp}\Rsign In that he liveth, he liveth unto God, alleluia, alleluia.\end{Resp}
\switchcolumn*

\LA
\LessonHead{Lectio II}
\switchcolumn
\EN
\LessonHead{Lesson II}
\switchcolumn*

\LA
\Cite{Act. 28:21-24}
\switchcolumn
\EN
\Cite{Acts 28:21-24}
\switchcolumn*

\LA
\noindent At illi dixérunt ad eum: Nos neque lítteras accépimus de te a Judǽa, neque advéniens áliquis fratrum nuntiávit, aut locútus est quid de te malum. Rogámus autem a te audíre quæ sentis: nam de secta hac notum est nobis quia ubíque ei contradícitur. Cum constituíssent autem illi diem, venérunt ad eum in hospítium plúrimi, quibus exponébat testíficans regnum Dei, suadénsque eis de Jesu ex lege Móysi et prophétis a mane usque ad vésperam. Et quidam credébant his quæ dicebántur: quidam vero non credébant.\par
\switchcolumn
\EN
\noindent But they said to him: We neither received letters concerning thee from Judea, neither did any of the brethren that came hither, relate or speak any evil of thee. But we desire to hear of thee what thou thinkest; for as concerning this sect, we know that it is every where contradicted. And when they had appointed him a day, there came very many to him unto his lodgings; to whom he expounded, testifying the kingdom of God, and persuading them concerning Jesus, out of the law of Moses and the prophets, from morning until evening. And some believed the things that were said; but some believed not.\par
\switchcolumn*

\LA
\begin{Resp}\Rsign Surréxit pastor bonus, qui ánimam suam pósuit pro óvibus suis, et pro grege suo mori dignátus est: \Star{} Allelúja, allelúja, allelúja.\end{Resp}
\begin{Resp}\Vsign Etenim Pascha nostrum immolátus est Christus.\end{Resp}
\begin{Resp}\Rsign Allelúja, allelúja, allelúja.\end{Resp}
\switchcolumn
\EN
\begin{Resp}\Rsign The Good Shepherd, Who laid down His life for the sheep, yea, Who was contented even to die for His flock, the Good Shepherd is risen again. \Star{} Alleluia, alleluia, alleluia.\end{Resp}
\begin{Resp}\Vsign For even Christ our Passover is sacrificed for us.\end{Resp}
\begin{Resp}\Rsign Alleluia, alleluia, alleluia.\end{Resp}
\switchcolumn*

\LA
\LessonHead{Lectio III}
\switchcolumn
\EN
\LessonHead{Lesson III}
\switchcolumn*

\LA
\Cite{Act. 28:25-31}
\switchcolumn
\EN
\Cite{Acts 28:25-31}
\switchcolumn*

\LA
\noindent Cumque ínvicem non essent consentiéntes, discedébant, dicénte Paulo unum verbum: Quia bene Spíritus Sanctus locútus est per Isaíam prophétam ad patres nostros, dicens: Vade ad pópulum istum, et dic ad eos: Aure audiétis, et non intellegétis, et vidéntes vidébitis, et non perspiciétis. Incrassátum est enim cor pópuli hujus, et áuribus gráviter audiérunt, et óculos suos compressérunt: ne forte vídeant óculis, et áuribus áudiant, et corde intéllegant, et convertántur, et sanem eos. Notum ergo sit vobis, quóniam géntibus missum est hoc salutáre Dei, et ipsi áudient. Et cum hæc dixísset, exiérunt ab eo Judǽi, multam habéntes inter se quæstiónem. Mansit autem biénnio toto in suo condúcto: et suscipiébat omnes qui ingrediebántur ad eum, prǽdicans regnum Dei, et docens quæ sunt de Dómino Jesu Christo cum omni fidúcia, sine prohibitióne.\par
\switchcolumn
\EN
\noindent And when they agreed not among themselves, they departed, Paul speaking this one word: Well did the Holy Ghost speak to our fathers by Isaias the prophet, Saying: Go to this people, and say to them: With the ear you shall hear, and shall not understand; and seeing you shall see, and shall not perceive. For the heart of this people is grown gross, and with their ears have they heard heavily, and their eyes they have shut; lest perhaps they should see with their eyes, and hear with their ears, and understand with their heart, and should be converted, and I should heal them. Be it known therefore to you, that this salvation of God is sent to the Gentiles, and they will hear it. And when he had said these things, the Jews went out from him, having much reasoning among themselves. And he remained two whole years in his own hired lodging; and he received all that came in to him, Preaching the kingdom of God, and teaching the things which concern the Lord Jesus Christ, with all confidence, without prohibition.\par
\switchcolumn*

\LA
\begin{Resp}\Rsign Ecce vicit leo de tribu Juda, radix David, aperíre librum, et sólvere septem signácula ejus: \Star{} Allelúja, allelúja, allelúja.\end{Resp}
\begin{Resp}\Vsign Dignus est Agnus, qui occísus est, accípere virtútem, et divinitátem, et sapiéntiam, et fortitúdinem, et honórem, et glóriam, et benedictiónem.\end{Resp}
\begin{Resp}\Rsign Allelúja, allelúja, allelúja.\end{Resp}
\begin{Resp}\Vsign Glória Patri, et Fílio, \Star{} et Spirítui Sancto.\end{Resp}
\begin{Resp}\Rsign Allelúja, allelúja, allelúja.\end{Resp}
\switchcolumn
\EN
\begin{Resp}\Rsign Behold, the Lion of the tribe of Judah, the Root of David, hath prevailed to open the Book, and to loose the seven seals thereof. \Star{} Alleluia, alleluia, alleluia.\end{Resp}
\begin{Resp}\Vsign Worthy is the Lamb That was slain to receive power, and riches, in wisdom, and strength, and honour, and glory, and blessing.\end{Resp}
\begin{Resp}\Rsign Alleluia, alleluia, alleluia.\end{Resp}
\begin{Resp}\Vsign Glory be to the Father, and to the Son, \Star{} and to the Holy Ghost.\end{Resp}
\begin{Resp}\Rsign Alleluia, alleluia, alleluia.\end{Resp}
\switchcolumn*

\LA
\LessonHead{Lectio IV}
\switchcolumn
\EN
\LessonHead{Lesson IV}
\switchcolumn*

\LA
\LessonTitle{Sermo sancti Joánnis Chrysóstomi}
\switchcolumn
\EN
\LessonTitle{From a Sermon by St. John Chrysostom}
\switchcolumn*

\LA
\Source{Homilia 4 in Matth.}
\switchcolumn
\EN
\Source{Homilia 4 in Matth.}
\switchcolumn*

\LA
\noindent Joseph, fili David, noli timére accípere Maríam cónjugem tuam. Quid autem est accípere? Domi profécto retinére; jam enim illam mente dimíserat: sed dimíssam, inquit, retíneas, quam Deus tibi cópulat, non paréntes; cópulat vero, non in fœdus solémne conjúgii, sed in consórtium commúnis habitáculi, et cópulat per meæ vocis offícium. Sicut enim illam póstea comméndat Christus ipse discípulo, ita étiam nunc Angelus Sponso; solátium tantúmmodo ejus habitúram absque fœdere nuptiárum. Deínde honéstius multóque dígnius causa partus expósita, suspiciónem quoque prorsus restínxit. Non modo, inquit, illícito non est violáta compléxu, verum étiam supra natúram morémque fœcúnda est. Noli ígitur de tam felíci partu Sponsæ attráhere mærórem, immo vero in majorem prorúmpe lætítiam; quod enim in ea natum est, de Spíritu Sancto est.\par
\switchcolumn
\EN
\noindent Joseph, thou son of David, fear not to take unto thee Mary thy wife. But what is: To take? Undoubtedly, to maintain, and that in his own house. For he had already sent her away in his mind. But now the Angel commandeth: Her whom thou wouldst send away, maintain; her, do thou, and not her parents, maintain, for God joineth her to thee; her, God verily joineth to thee, not in the sacred commerce of marriage, but in the fellowship of a common home; and her, God joineth to thee through the ministry of my words. Just as Christ himself later entrusted her to the care of his disciple, so now the Angel giveth her to her spouse; in such manner that she may have the consolation of his company without other conjugal rights. By this means her confinement would be explained in a worthier and more honourable way, and suspicion would be allayed. It is as though the Angel said: Not only was she not dishonoured by an unlawful embrace, but indeed she is fruitful in a manner above nature and usage; therefore grieve not at the happy confinement of thy Bride, but break forth into greater joy! For that which is conceived in her is of the Holy Ghost.\par
\switchcolumn*

\LA
\begin{Resp}\Rsign Dedísti mihi protectiónem salútis tuæ, et déxtera tua suscépit me: \Star{} Protéctor meus, et cornu salútis meæ, et suscéptor meus, allelúja.\end{Resp}
\begin{Resp}\Vsign Ego protéctor tuus sum, et merces tua magna nimis.\end{Resp}
\begin{Resp}\Rsign Protéctor meus, et cornu salútis meæ, et suscéptor meus, allelúja.\end{Resp}
\switchcolumn
\EN
\begin{Resp}\Rsign Thou hast given me the shield of thy salvation, and thy right hand hath holden me up. \Star{} My buckler, and the horn of my salvation, and my refuge. Alleluia.\end{Resp}
\begin{Resp}\Vsign I am thy shield and thy exceeding great reward.\end{Resp}
\begin{Resp}\Rsign My buckler, and the horn of my salvation, and my refuge. Alleluia.\end{Resp}
\switchcolumn*

\LA
\LessonHead{Lectio V}
\switchcolumn
\EN
\LessonHead{Lesson V}
\switchcolumn*

\LA
\noindent Páriet autem Fílium, et vocábis nomen ejus Jesum; non enim quia ex Spíritu Sancto est, idcírco te a ministério tantæ exístimes dispensatiónis extráneum. Nam etsi nihil hábeas in hac generatióne commúne (Virgo quippe permánsit intácta), tamen quod est próprium patris quodque nihil offúscat Vírginis dignitátem, hoc tibi fácile concédo, ut scílicet Nato nomen impónas; tu enim illum primum vocábis. Quamquam enim non sit fílius tuus iste, qui náscitur, tu tamen curam erga illum et sollicitúdinem osténdes paréntis; et proptérea te illi ab ipsa statim nóminis impositióne conjúngo. Deínde, ne quis illum ex hoc patrem esse suspicarétur, Páriet, inquit, Fílium. Non dixit, Páriet tibi, sed pósuit illud indefinítum et suspénsum; non enim illi, sed univérso prorsus orbi péperit Christum.\par
\switchcolumn
\EN
\noindent And she shall bring forth a Son, and thou shalt call his Name Jesus. That is: Think not that the ministry of this great dispensation, because it is of the Holy Ghost, is a thing apart from thee. For even thought thou hast no part in his generation, since the Virgin remaineth ínviolate, yet do I readily grant thee this, namely; that thine are all the rights of a father, in so far as they obscure not the dignity of the Virgin; thou shalt certainly give the new-born his Name; thou shalt be the first to call him by his Name. For even though he who is born is not thy son, nonetheless thou shalt shew him the care and solicitude of a parent; and therefore I unite thee to him by this immediate giving of the Name. But, lest anyone might think from this that Joseph was the begetter of Christ, the Angel was first careful to say: She shall bring forth a Son. He doth not say: She shall bear thee a son: but maketh his statement in an undetermined and indefinite way. For Mary did not bear a son to Joseph, but brought forth Christ to the whole world.\par
\switchcolumn*

\LA
\begin{Resp}\Rsign Státuet fílios suos sub tégmine illíus, et sub ramis ejus morábitur: protegétur sub tégmine illíus a fervóre: \Star{} Et in glória ejus requiéscet, allelúja.\end{Resp}
\begin{Resp}\Vsign Speráte in eo, omnis congregátio pópuli, effúndite coram illo corda vestra.\end{Resp}
\begin{Resp}\Rsign Et in glória ejus requiéscet, allelúja.\end{Resp}
\switchcolumn
\EN
\begin{Resp}\Rsign He shall set his children under her shelter, and shall lodge under her branches by her shall he be covered from heat, \Star{} And in her glory shall he dwell. Alleluia.\end{Resp}
\begin{Resp}\Vsign Trust in Him, ye congregation of the people, pour out your heart before Him.\end{Resp}
\begin{Resp}\Rsign And in her glory shall he dwell. Alleluia.\end{Resp}
\switchcolumn*

\LA
\LessonHead{Lectio VI}
\switchcolumn
\EN
\LessonHead{Lesson VI}
\switchcolumn*

\LA
\noindent Proptérea et nomen ejus de cælo Angelum detulísse Evangelísta memorávit, ut hinc quoque osténderet illum mirábilem esse partum, quo ejus nomen ad Joseph et per Angelum et a Deo missum docéret. Nam et ipsum vocábulum non inániter pósitum est, quod certe mille cóntinet thesáuros bonórum. Propter quod illud étiam Angelus interpretátur bonis mæréntem spebus animándo; et hoc quoque illum modo ad credéndum quod lóquitur, invítat. Fácile namque sollicitámur ad próspera, et prómptius fidem accommodámus secúndis. Ipse enim, inquit, salvum fáciet pópulum suum a peccátis eórum. Hinc quoque benefícii nóvitas indicátur. Non enim a bello visíbili neque a gládio barbarórum, sed, quod his longe majus est, a peccáto suo pópulum suum núntiat liberándum: quod præstáre nulli fuit hóminum aliquándo possíbile.\par
\switchcolumn
\EN
\noindent Therefore the Evangelist relateth that the Angel brought his Name from heaven, so that thus might be shewn how wonderful was his birth, seeing that he himself taught his Name to Joseph by an Angel sent from God. For this Name, which verily containeth a thousand treasures of good, was not given without meaning. Therefore the Angel doth himself interpret it, thereby consoling Joseph's grief with good hopes; and thus also inviting him to believe these words. For we are easily summoned to that which is pleasant, and give prompt credence unto good tidings. Wherefore the Angel said: He shall save his people from their sins. This also sheweth the novelty of the gift. For he announceth that this people are to be saved, not indeed from external wars, nor from the swords of barbarians, but from what is far greater than these: From their sins. And no mere man could ever accomplish this.\par
\switchcolumn*

\LA
\begin{Resp}\Rsign Si consístant advérsum me castra, non timébit cor meum: \Star{} Si exsúrgat advérsum me prǽlium, in hoc ego sperábo, allelúja.\end{Resp}
\begin{Resp}\Vsign In te cantátio mea semper, quóniam tu adjútor fortis.\end{Resp}
\begin{Resp}\Rsign Si exsúrgat advérsum me prǽlium, in hoc ego sperábo, allelúja.\end{Resp}
\begin{Resp}\Vsign Glória Patri, et Fílio, \Star{} et Spirítui Sancto.\end{Resp}
\begin{Resp}\Rsign Si exsúrgat advérsum me prǽlium, in hoc ego sperábo, allelúja.\end{Resp}
\switchcolumn
\EN
\begin{Resp}\Rsign Though an host should encamp against me, my heart shall not fear. \Star{} Though war should rise against me, in this will I be confident. Alleluia.\end{Resp}
\begin{Resp}\Vsign My praise shall be continually of thee, for Thou art my strong refuge. R.* Though war should rise against me, in this will I be confident. Alleluia.\end{Resp}
\begin{Resp}\Vsign Glory be to the Father, and to the Son, \Star{} and to the Holy Ghost.\end{Resp}
\begin{Resp}\Rsign Though war should rise against me, in this will I be confident. Alleluia.\end{Resp}
\switchcolumn*

\LA
\LessonHead{Lectio VII}
\switchcolumn
\EN
\LessonHead{Lesson VII}
\switchcolumn*

\LA
\LessonTitle{Léctio sancti Evangélii secúndum Lucam}
\switchcolumn
\EN
\LessonTitle{From the Holy Gospel according to Luke}
\switchcolumn*

\LA
\Cite{Luc 3:21-23}
\switchcolumn
\EN
\Cite{Luke 3:21-23}
\switchcolumn*

\LA
\noindent In illo témpore: Factum est autem cum baptizarétur omnis pópulus, et Jesu baptizáto et oránte, apértum est cælum. Et réliqua.\par
\switchcolumn
\EN
\noindent At that time: When all the people were baptized, it came to pass, that Jesus also being baptized and praying, the heaven was opened. And so on, and that which followeth.\par
\switchcolumn*

\LA
\LessonTitle{De Homilía sancti Ambrósii Epíscopi}
\switchcolumn
\EN
\LessonTitle{A Homily by St. Ambrose the Bishop}
\switchcolumn*

\LA
\Source{Expositio in Luc. lib. 3}
\switchcolumn
\EN
\Source{Expositio in Luc. lib. 3}
\switchcolumn*

\LA
\noindent Néminem movére debet, quod ita scriptum est: Qui putabátur fílius Joseph. Bene enim putabátur, quia natúra non erat; sed ídeo putabátur, quia eum María, quæ Joseph viro suo erat desponsáta genúerat. Sic enim habes: Nonne hic est fílius Joseph fabri? Díximus supra qua ratióne per Vírginem, díximus étiam, qua ratióne per desponsátam, et quare census témpore nasci volúerit Dóminus salutáris; non aliénum étiam vidétur, ut qua ratióne fabrum patrem habúerit, declarémus. Hoc enim typo eum patrem sibi esse demónstrat, qui fabricátor ómnium cóndidit mundum. Nam etsi humána non sunt comparánda divínis, typus tamen ínteger est, quod Pater Christi igne operátur et spíritu, et tamquam bonus ánimæ faber vítia nostra circúmdolat; cito secúrim ádmovens arbóribus infœcúndis, secáre doctus exígua, culmínibus serváre sublímia, rígida méntium spíritus igne mollíre, et in vários usus omne humánum genus divérsa ministeriórum qualitáte formáre.\par
\switchcolumn
\EN
\noindent No one should be troubled at the words: As was supposed, the son of Joseph. For it was no more than a supposition, seeing that Christ was not the son of Joseph by nature. Albeit, it was so supposed because Mary who was espoused to her husband Joseph, gave Christ birth. And so, referring to Joseph as father, it is written: Is not this the carpenter's son? We have already discussed why the Lord of salvation chose to be born of a Virgin. We have also discussed why she was an espoused Virgin when her conception took place; and why it took place at the time of the enrollment for taxing. Hence it is not fitting to explain why Christ had a working-man for his father. For thereby is figured Christ's divine Father, who as Maker of all things, framed the world. Even though human and divine matters be not equal to each other, yet is this figure a complete one. Christ's Father worketh by fire and by breathing on things. Yea, like a good carpenter of the soul, he chippeth away our defects. Promptly doth he lay his axe to the barren trees and hew them down. Skilful is he in correcting whatever is built scantily, and in buttressing whatever is to be built magnificently. He tempereth the hardness of hearts as with fire, and with his gentle Breath. And by his diverse workings he formeth the quality of the human race.\par
\switchcolumn*

\LA
\begin{Resp}\Rsign Joseph, fili David, noli timére accípere Maríam cónjugem tuam; quod enim in ea natum est, de Spíritu Sancto est: páriet autem fílium, \Star{} Et vocábis nomen ejus Jesum, allelúja.\end{Resp}
\begin{Resp}\Vsign Ipse enim salvum fáciet pópulum suum a peccátis eórum.\end{Resp}
\begin{Resp}\Rsign Et vocábis nomen ejus Jesum, allelúja.\end{Resp}
\switchcolumn
\EN
\begin{Resp}\Rsign Joseph, thou Son of David, fear not to take unto thee Mary thy wife; for That Which is conceived in her is of the Holy Ghost and she shall bring forth a Son; \Star{} And thou shalt call His Name Jesus. Alleluia.\end{Resp}
\begin{Resp}\Vsign For He shall save His people from their sins.\end{Resp}
\begin{Resp}\Rsign And thou shalt call His Name Jesus. Alleluia.\end{Resp}
\switchcolumn*

\LA
\LessonHead{Lectio VIII}
\switchcolumn
\EN
\LessonHead{Lesson VIII}
\switchcolumn*

\LA
\noindent Cur autem Joseph magis quam Maríæ generátio describátur, cum María de Spíritu Sancto generávit Christum, et Joseph a generatióne Dómini videátur aliénus, dubitáre possémus, nisi consuetúdo nos instrúeret Scripturárum, quæ semper viri oríginem quærit. Viri enim persóna quæritur, qui étiam in senátu et réliquis cúriis civitátum géneris assérvat dignitátem. Quam defórme autem, si relícta viri orígine, orígo féminæ quærerétur, ut viderétur patrem non habuísse ille totíus mundi pópulis prædicándus! Vides ubíque famíliam per virórum generatiónes esse decúrsam. Noli mirári quod Joseph orígo descrípta est. Etenim secúndum carnem natus, usum débuit sequi carnis, et qui in sæculum venit, sæculi débuit more descríbi; máxime cum in Joseph orígine étiam orígo sit Maríæ.\par
\switchcolumn
\EN
\noindent We might wonder why the genealogy of Joseph, rather than of Mary, is given (since Mary conceived Christ by the Holy Ghost, and Joseph had not part in the Lord's conception), were it not that the Holy Scripture teacheth us how it was the custom to trace descent on the male side. For in this fashion the person of the man is set forth as pre-eminent, and his dignity maintained, even as it is wont to be done in the Senate and the other high places of the commonwealth. And how unseemly it would have been to have passed over the lineage of the father, and to have given that of the mother, since the so doing would have appeared to proclaim to all the people in the world that Christ had not father. It is a world-wide custom to trace the genealogy of a family in the male line. Therefore, be not perplexed that the lineage of Joseph is given. Forasmuch as Christ was born in the flesh, he was bound to follow this custom of the flesh. And he who came into the world had to be enrolled for taxing in the worldly manner, the more so that Joseph's descent was the same as Mary's.\par
\switchcolumn*

\LA
\begin{Resp}\Rsign Surge, et áccipe Púerum et matrem ejus, et fuge in Ægýptum: \Star{} Et esto ibi usque dum dicam tibi, allelúja.\end{Resp}
\begin{Resp}\Vsign Ut adimplerétur quod dictum est a Dómino per Prophétam dicéntem: Ex Ægýpto vocávi Fílium meum.\end{Resp}
\begin{Resp}\Rsign Et esto ibi usque dum dicam tibi, allelúja.\end{Resp}
\begin{Resp}\Vsign Glória Patri, et Fílio, \Star{} et Spirítui Sancto.\end{Resp}
\begin{Resp}\Rsign Et esto ibi usque dum dicam tibi, allelúja.\end{Resp}
\switchcolumn
\EN
\begin{Resp}\Rsign Arise, and take the young Child, and His Mother, and flee into Egypt; \Star{} And be thou there until I bring thee word. Alleluia.\end{Resp}
\begin{Resp}\Vsign That it might be fulfilled which was spoken of the Lord by the Prophets, saying Out of Egypt have I called My Son.\end{Resp}
\begin{Resp}\Rsign And be thou there until I bring thee word. Alleluia.\end{Resp}
\begin{Resp}\Vsign Glory be to the Father, and to the Son, \Star{} and to the Holy Ghost.\end{Resp}
\begin{Resp}\Rsign And be thou there until I bring thee word. Alleluia.\end{Resp}
\switchcolumn*

\LA
\LessonHead{Lectio IX}
\switchcolumn
\EN
\LessonHead{Lesson IX}
\switchcolumn*

\LA
\noindent Cur autem sanctus Matthæus ab Abraham generatiónem enumeráre cœperit Christi, sanctus vero Lucas a Christo usque ad Deum perdúxerit, explanándum vidétur. Lucas ad Deum putávit oríginem ejus esse referéndam, quod verus Christi generátor Deus sit, vel secúndum veram generatiónem Pater, vel secúndum lavácri regeneratiónem mystici múneris auctor. Et ídeo non a primo generatiónem ejus cœpit descríbere; sed posteáquam baptísmum ejus explícuit, Auctórem ómnium Deum per baptísmum cúpiens demonstráre. Christum quoque a Deo órdine manásse successiónis asséruit, univérsa contéxens, ut et secúndum natúram et secúndum grátiam et secúndum carnem Dei Fílium demonstráret. Quod autem evidéntius divínæ generatiónis indícium, quam quod de generatióne dictúrus, ipsum Patrem præmísit loquéntem: Hic est Fílius meus diléctus, in quo complácui?\par
\switchcolumn
\EN
\noindent But some explanation is required as to why Saint Matthew reckoneth Christ's descent from Abraham forward, whilst Saint Luke traceth the same from Christ backward to the creation of Adam by God. By this, Luke would have us understand that Christ's lineage should be traced to God, because God was Christ's true Progenitor, both as his Father whereof he was begotten, and as the Author, in the laver of baptism, of the mystical gift of the Spirit. Wherefore Luke doth not begin his Gospel with the reckoning of Christ's lineage, but recordeth it after the account of the baptism, thereby shewing forth in baptism the working of God, the Author of all things. Thus also Luke asserteth that Christ came forth from God according to a certain rule of orderliness. For he weaveth all things together to prove that Christ is by nature, by grace, and in the flesh, the Son of God. But what more evident proof of Christ's divine descent could we have than what Luke giveth? For before the Evangelist reckoneth the genealogy of Christ, he giveth the words of the Father himself: This is my beloved Son, in whom I am well pleased.\par
\switchcolumn*

\LA
\TeDeum{Te Deum}
\switchcolumn
\EN
\TeDeum{Te Deum}
\switchcolumn*

\end{paracol}

\Office{Dominica III Post Pascha}{Third Sunday after Easter}{Semiduplex}
\addcontentsline{toc}{subsection}{Dominica III Post Pascha\enspace\textit{Third Sunday after Easter}}
\begin{paracol}{2}
\LA
\LessonHead{Lectio I}
\switchcolumn
\EN
\LessonHead{Lesson I}
\switchcolumn*

\LA
\LessonTitle{Incipit liber Apocalýpsis beáti Joánnis Apóstoli}
\switchcolumn
\EN
\LessonTitle{Lesson from the book of Revelation}
\switchcolumn*

\LA
\Cite{Apo 1:1-6}
\switchcolumn
\EN
\Cite{Rev 1:1-6}
\switchcolumn*

\LA
\noindent Apocalýpsis Jesu Christi, quam dedit illi Deus palam fácere servis suis, quæ opórtet fíeri cito: et significávit, mittens per ángelum suum servo suo Joánni, qui testimónium perhíbuit verbo Dei, et testimónium Jesu Christi, quæcúmque vidit. Beátus qui legit, et audit verba prophetíæ hujus, et servat ea, quæ in ea scripta sunt: tempus enim prope est. Joánnes septem ecclésiis, quæ sunt in Asia. Grátia vobis, et pax ab eo, qui est, et qui erat, et qui ventúrus est: et a septem spirítibus qui in conspéctu throni ejus sunt: et a Jesu Christo, qui est testis fidélis, primogénitus mortuórum, et princeps regum terræ, qui diléxit nos, et lavit nos a peccátis nostris in sánguine suo, et fecit nos regnum, et sacerdótes Deo et Patri suo: ipsi glória et impérium in sǽcula sæculórum. Amen.\par
\switchcolumn
\EN
\noindent The Revelation of Jesus Christ, which God gave unto him, to make known to his servants the things which must shortly come to pass: and signified, sending by his angel to his servant John, who hath given testimony to the word of God, and the testimony of Jesus Christ, what things soever he hath seen. Blessed is he, that readeth and heareth the words of this prophecy; and keepeth those things which are written in it; for the time is at hand. John to the seven churches which are in Asia. Grace be unto you and peace from him that is, and that was, and that is to come, and from the seven spirits which are before his throne, And from Jesus Christ, who is the faithful witness, the first begotten of the dead, and the prince of the kings of the earth, who hath loved us, and washed us from our sins in his own blood, And hath made us a kingdom, and priests to God and his Father, to him be glory and empire for ever and ever. Amen.\par
\switchcolumn*

\LA
\begin{Resp}\Rsign Dignus es, Dómine, accípere librum, et aperíre signácula ejus, allelúja: quóniam occísus es, et redemísti nos Deo \Star{} In sánguine tuo, allelúja.\end{Resp}
\begin{Resp}\Vsign Fecísti enim nos Deo nostro regnum et sacerdótium.\end{Resp}
\begin{Resp}\Rsign In sánguine tuo, allelúja.\end{Resp}
\switchcolumn
\EN
\begin{Resp}\Rsign Thou art worthy, O Lord, to take the book, and to open the seals thereof; alleluia: because thou wast slain, and hast redeemed us to God, \Star{} In thy blood, alleluia.\end{Resp}
\begin{Resp}\Vsign And hast made us to our God a kingdom and priests.\end{Resp}
\begin{Resp}\Rsign In thy blood, alleluia.\end{Resp}
\switchcolumn*

\LA
\LessonHead{Lectio II}
\switchcolumn
\EN
\LessonHead{Lesson II}
\switchcolumn*

\LA
\Cite{Apo 1:7-11}
\switchcolumn
\EN
\Cite{Rev 1:7-11}
\switchcolumn*

\LA
\noindent Ecce venit cum núbibus, et vidébit eum omnis óculus, et qui eum pupugérunt. Et plangent se super eum omnes tribus terræ. Etiam: amen. Ego sum alpha et ómega, princípium et finis, dicit Dóminus Deus: qui est, et qui erat, et qui ventúrus est, omnípotens. Ego Joánnes frater vester, et párticeps in tribulatióne, et regno, et patiéntia in Christo Jesu: fui in ínsula, quæ appellátur Patmos, propter verbum Dei, et testimónium Jesu: fui in spíritu in domínica die, et audívi post me vocem magnam tamquam tubæ, dicéntis: Quod vides, scribe in libro: et mitte septem ecclésiis, quæ sunt in Asia, Epheso, et Smyrnæ, et Pérgamo, et Thyatíræ, et Sardis, et Philadelphíæ, et Laodicíæ.\par
\switchcolumn
\EN
\noindent Behold, he cometh with the clouds, and every eye shall see him, and they also that pierced him. And all the tribes of the earth shall bewail themselves because of him. Even so. Amen. I am Alpha and Omega, the beginning and the end, saith the Lord God, who is, and who was, and who is to come, the Almighty. I John, your brother and your partner in tribulation, and in the kingdom, and patience in Christ Jesus, was in the island, which is called Patmos, for the word of God, and for the testimony of Jesus. I was in the spirit on the Lord's day, and heard behind me a great voice, as of a trumpet, Saying: What thou seest, write in a book, and send to the seven churches which are in Asia, to Ephesus, and to Smyrna, and to Pergamus, and to Thyatira, and to Sardis, and to Philadelphia, and to Laodicea.\par
\switchcolumn*

\LA
\begin{Resp}\Rsign Ego sicut vitis fructificávi suavitátem odóris, allelúja: \Star{} Transíte ad me, omnes qui concupíscitis me, et a generatiónibus meis adimplémini, allelúja, allelúja.\end{Resp}
\begin{Resp}\Vsign In me omnis grátia viæ et veritátis: in me omnes spes vitæ et virtútis.\end{Resp}
\begin{Resp}\Rsign Transíte ad me, omnes qui concupíscitis me, et a generatiónibus meis adimplémini, allelúja, allelúja.\end{Resp}
\switchcolumn
\EN
\begin{Resp}\Rsign As the vine I have brought forth a pleasant odour, alleluia. \Star{} Come over to me, all ye that desire me, and be filled with my fruits, alleluia, alleluia.\end{Resp}
\begin{Resp}\Vsign in me is all hope of life and of virtue.\end{Resp}
\begin{Resp}\Rsign Come over to me, all ye that desire me, and be filled with my fruits, alleluia, alleluia.\end{Resp}
\switchcolumn*

\LA
\LessonHead{Lectio III}
\switchcolumn
\EN
\LessonHead{Lesson III}
\switchcolumn*

\LA
\Cite{Apo 1:12-19}
\switchcolumn
\EN
\Cite{Rev 1:12-19}
\switchcolumn*

\LA
\noindent Et convérsus sum ut vidérem vocem, quæ loquebátur mecum: et convérsus vidi septem candelábra áurea: et in médio septem candelabrórum aureórum, símilem Fílio hóminis vestítum podére, et præcínctum ad mamíllas zona áurea: caput autem ejus, et capílli erant cándidi tamquam lana alba, et tamquam nix, et óculi ejus tamquam flamma ignis: et pedes ejus símiles aurichálco, sicut in camíno ardénti, et vox illíus tamquam vox aquárum multárum: et habébat in déxtera sua stellas septem: et de ore ejus gládius utráque parte acútus exíbat: et fácies ejus sicut sol lucet in virtúte sua. Et cum vidíssem eum, cécidi ad pedes ejus tamquam mórtuus. Et pósuit déxteram suam super me, dicens: Noli timére: ego sum primus, et novíssimus, et vivus, et fui mórtuus, et ecce sum vivens in sǽcula sæculórum: et hábeo claves mortis, et inférni. Scribe ergo quæ vidísti, et quæ sunt, et quæ opórtet fíeri post hæc.\par
\switchcolumn
\EN
\noindent And I turned to see the voice that spoke with me. And being turned, I saw seven golden candlesticks: And in the midst of the seven golden candlesticks, one like to the Son of man, clothed with a garment down to the feet, and girt about the paps with a golden girdle. And his head and his hairs were white, as white wool, and as snow, and his eyes were as a flame of fire, And his feet like unto fine brass, as in a burning furnace. And his voice as the sound of many waters. And he had in his right hand seven stars. And from his mouth came out a sharp two edged sword: and his face was as the sun shineth in his power. And when I had seen him, I fell at his feet as dead. And he laid his right hand upon me, saying: Fear not. I am the First and the Last, And alive, and was dead, and behold I am living for ever and ever, and have the keys of death and of hell. Write therefore the things which thou hast seen, and which are, and which must be done hereafter.\par
\switchcolumn*

\LA
\begin{Resp}\Rsign Audívi vocem de cælo, tamquam vocem tonítrui magni, allelúja: Regnábit Deus noster in ætérnum, allelúja: \Star{} Quia facta est salus, et virtus, et potéstas Christi ejus, allelúja, allelúja.\end{Resp}
\begin{Resp}\Vsign Et vox de throno exívit, dicens: Laudem dícite Deo nostro, omnes Sancti ejus, et qui timétis Deum, pusílli et magni.\end{Resp}
\begin{Resp}\Rsign Quia facta est salus, et virtus, et potéstas Christi ejus, allelúja, allelúja.\end{Resp}
\begin{Resp}\Vsign Glória Patri, et Fílio, \Star{} et Spirítui Sancto.\end{Resp}
\begin{Resp}\Rsign Quia facta est salus, et virtus, et potéstas Christi ejus, allelúja, allelúja.\end{Resp}
\switchcolumn
\EN
\begin{Resp}\Rsign And I heard a voice from heaven, the voice of great thunder, alleluia: the Lord our God shall reign for ever, alleluia: \Star{} Now is come salvation, and strength, and the power of his Christ, alleluia, alleluia.\end{Resp}
\begin{Resp}\Vsign And a voice came out from the throne, saying: Give praise to our God, all saints of Him that fear him, little and great.\end{Resp}
\begin{Resp}\Rsign Now is come salvation, and strength, and the power of his Christ, alleluia, alleluia.\end{Resp}
\begin{Resp}\Vsign Glory be to the Father, and to the Son, \Star{} and to the Holy Ghost.\end{Resp}
\begin{Resp}\Rsign Now is come salvation, and strength, and the power of his Christ, alleluia, alleluia.\end{Resp}
\switchcolumn*

\LA
\LessonHead{Lectio IV}
\switchcolumn
\EN
\LessonHead{Lesson IV}
\switchcolumn*

\LA
\LessonTitle{Sermo sancti Augustíni Epíscopi}
\switchcolumn
\EN
\LessonTitle{Sermon from St. Augustine, Bishop (of Hippo)}
\switchcolumn*

\LA
\Source{Sermo 147 de Tempore}
\switchcolumn
\EN
\Source{Sermo 147 de Tempore}
\switchcolumn*

\LA
\noindent Diébus his sanctis resurrectióni Dómini dedicátis, quantum donánte ipso póssumus, de carnis resurrectióne tractémus. Hæc enim est fides nostra: hoc donum in Dómini nostri Jesu Christi nobis carne promíssum est, et in ipso præcéssit exémplum. Vóluit enim nobis, quod promísit in fine, non solum prænuntiáre, sed étiam demonstráre. Illi quidem qui tunc fuérunt, cum illum vidérent, et cum expavéscerent, et spíritum se vidére créderent, soliditátem córporis tenuérunt. Locútus est enim non solum verbis ad aures eórum, sed étiam spécie ad óculos eórum: parúmque erat se præbére cernéndum, nisi étiam offérret pertractándum atque palpándum.\par
\switchcolumn
\EN
\noindent During these Holy Days in commemoration of the Lord's resurrection, we purpose to preach, so far as he will empower us, the doctrine of the resurrection of the body. For this is the Faith, to wit: the gift of resurrection, which was bestowed upon the flesh of our Lord Jesus Christ, is what is promised to us, for it was first made manifest in him that we might know what to hope for ourselves. What he hath thus promised would come to us at the last, he willed not only to foretell but to demonstrate. Those who were present at that time (even though they were terrified and affrighted, and supposed they had seen a spirit), handled him, and saw that a spirit would not have flesh and bones, such as they saw him to have. Thus he spake to them, not only in words which they could hear, but in a body which they could see; as if it had not been enough to shew himself to their sight, but must needs even offer himself to be touched and handled.\par
\switchcolumn*

\LA
\begin{Resp}\Rsign Locútus est ad me unus ex septem Angelis, dicens: Veni, osténdam tibi novam nuptam, sponsam Agni: \Star{} Et vidi Jerúsalem descendéntem de cælo, ornátam monílibus suis, allelúja, allelúja, allelúja.\end{Resp}
\begin{Resp}\Vsign Et sústulit me in spíritu in montem magnum et altum.\end{Resp}
\begin{Resp}\Rsign Et vidi Jerúsalem descendéntem de cælo, ornátam monílibus suis, allelúja, allelúja, allelúja.\end{Resp}
\switchcolumn
\EN
\begin{Resp}\Rsign One of the seven Angels talked with me, saying: Come hither, I will show thee the bride, the Lamb's wife. \Star{} And I saw Jerusalem descending out of heaven, adorned with her jewels. Alleluia, Alleluia, Alleluia.\end{Resp}
\begin{Resp}\Vsign And he carried me away in the Spirit to a great and high mountain\end{Resp}
\begin{Resp}\Rsign And I saw Jerusalem descending out of heaven, adorned with her jewels. Alleluia, Alleluia, Alleluia.\end{Resp}
\switchcolumn*

\LA
\LessonHead{Lectio V}
\switchcolumn
\EN
\LessonHead{Lesson V}
\switchcolumn*

\LA
\noindent Ait enim: Quid turbáti estis, et cogitatiónes ascéndunt in cor vestrum? Putavérunt enim se spíritum vidére. Quid turbáti estis, inquit, et cogitatiónes ascéndunt in cor vestrum? Vidéte manus meas, et pedes meos: palpáte, et vidéte: quia spíritus ossa et carnem non habet, sicut me vidétis habére. Contra istam evidéntiam disputábant hómines. Quid enim áliud fácerent hómines, qui ea, quæ sunt hóminum, sápiunt, quam sic disputáre de Deo contra Deum? Ille enim Deus est, isti hómines sunt. Sed Deus novit cogitatiónes hóminum, quóniam vanæ sunt.\par
\switchcolumn
\EN
\noindent For he said: Why are ye troubled? and why do thoughts arise in your hearts? For they supposed that they had seen a spirit. Therefore he added: Behold my hands and my feet, that it is I myself; handle me and see; for a spirit hath not flesh and bones, as ye see me have. Of course, men have disputed this evidence; for what else could men do, seeing that they are wise only according to man's wisdom, which thus permitteth them to dispute concerning God in spite of what God hath shewn them of himself. He is God, they are men. But God knoweth the thoughts of man, that they are but vain.\par
\switchcolumn*

\LA
\begin{Resp}\Rsign Audívi vocem in cælo Angelórum multórum dicéntium: \Star{} Timéte Dóminum, et date claritátem illi, et adoráte eum, qui fecit cælum et terram, mare et fontes aquárum, allelúja, allelúja.\end{Resp}
\begin{Resp}\Vsign Vidi Angelum Dei fortem, volántem per médium cæli, voce magna clamántem et dicéntem.\end{Resp}
\begin{Resp}\Rsign Timéte Dóminum, et date claritátem illi, et adoráte eum, qui fecit cælum et terram, mare et fontes aquárum, allelúja, allelúja.\end{Resp}
\switchcolumn
\EN
\begin{Resp}\Rsign I heard in heaven the voice of many Angels, saying \Star{} Fear the Lord, and give glory to Him, and worship Him That made heaven and earth, the sea, and the fountains of waters. Alleluia, Alleluia.\end{Resp}
\begin{Resp}\Vsign I saw a strong Angel of God fly through the midst of heaven, crying with a loud voice and saying\end{Resp}
\begin{Resp}\Rsign Fear the Lord, and give glory to Him, and worship Him That made heaven and earth, the sea, and the fountains of waters. Alleluia, Alleluia.\end{Resp}
\switchcolumn*

\LA
\LessonHead{Lectio VI}
\switchcolumn
\EN
\LessonHead{Lesson VI}
\switchcolumn*

\LA
\noindent In hómine carnáli tota régula intelligéndi est consuetúdo cernéndi. Quod solent vidére, credunt: quod non solent, non credunt. Præter consuetúdinem facit Deus mirácula, quia Deus est. Majóra quidem mirácula sunt, tot quotídie hómines nasci, qui non erant, quam paucos resurrexísse, qui erant: et tamen ista mirácula non consideratióne comprehénsa sunt, sed assiduitáte viluérunt. Resurréxit Christus: absolúta est res. Corpus erat, caro erat: pepéndit in cruce, emísit ánimam, pósita est caro in sepúlcro. Exhíbuit illam vivam, qui vivébat in illa. Quare mirámur? quare non crédimus? Deus est, qui fecit.\par
\switchcolumn
\EN
\noindent To carnal men, the one rule of understanding is his ordinary experience; seeing is believing. What men are accustomed to see, that they credit; what they are not accustomed to see, that they deem incredible. But God often worketh wonders (that is, things contrary to what we are accustomed), because he is God. Every day many men are born that previously had no existence at all; and this is a greater miracle than that a few, who did exist, have been raised from the dead. Yet this wonder is not recognized as such; on the contrary, it is disregarded because man is accustomed to it. Christ rose again from the dead; that is a fact. He had a body: he took flesh, he hung upon the cross, he gave up the ghost; his flesh was laid in the tomb. After that, he shewed his flesh as alive again, he lived again the flesh. Why wonder, why deny it? God wrought this.\par
\switchcolumn*

\LA
\begin{Resp}\Rsign Véniens a Líbano quam pulchra facta est, allelúja: \Star{} Et odor vestimentórum ejus super ómnia arómata, allelúja, allelúja.\end{Resp}
\begin{Resp}\Vsign Favus distíllans lábia ejus, mel et lac sub lingua ejus.\end{Resp}
\begin{Resp}\Rsign Et odor vestimentórum ejus super ómnia arómata, allelúja, allelúja.\end{Resp}
\begin{Resp}\Vsign Glória Patri, et Fílio, \Star{} et Spirítui Sancto.\end{Resp}
\begin{Resp}\Rsign Et odor vestimentórum ejus super ómnia arómata, allelúja, allelúja.\end{Resp}
\switchcolumn
\EN
\begin{Resp}\Rsign Coming from Lebanon how sweet she is, alleluia: \Star{} And the sweet smell of her garments, above all aromatical spices, alleluia, alleluia.\end{Resp}
\begin{Resp}\Vsign Thy lips are as a dropping honeycomb, honey and milk are under thy tongue.\end{Resp}
\begin{Resp}\Rsign And the sweet smell of her garments, above all aromatical spices, alleluia, alleluia.\end{Resp}
\begin{Resp}\Vsign Glory be to the Father, and to the Son, \Star{} and to the Holy Ghost.\end{Resp}
\begin{Resp}\Rsign And the sweet smell of her garments, above all aromatical spices, alleluia, alleluia.\end{Resp}
\switchcolumn*

\LA
\LessonHead{Lectio VII}
\switchcolumn
\EN
\LessonHead{Lesson VII}
\switchcolumn*

\LA
\LessonTitle{Léctio sancti Evangélii secúndum Joánnem}
\switchcolumn
\EN
\LessonTitle{From the Holy Gospel according to John}
\switchcolumn*

\LA
\Cite{Joannes 16:16-22}
\switchcolumn
\EN
\Cite{John 16:16-22}
\switchcolumn*

\LA
\noindent In illo témpore: Dixit Jesus discípulis suis: Módicum, et jam non vidébitis me; et íterum módicum, et vidébitis me: quia vado ad Patrem. Et réliqua.\par
\switchcolumn
\EN
\noindent At that time, Jesus said to his disciples: A little while, and now you shall not see me; and again a little while, and you shall see me: because I go to the Father. And so on.\par
\switchcolumn*

\LA
\LessonTitle{Homilía sancti Augustíni Epíscopi}
\switchcolumn
\EN
\LessonTitle{Homily by St. Augustine, Bishop (of Hippo.)}
\switchcolumn*

\LA
\Source{Tr. 101 in Joann., sub fine}
\switchcolumn
\EN
\Source{Tract 101 of John}
\switchcolumn*

\LA
\noindent Módicum est hoc totum spátium, quo præsens pérvolat sǽculum. Unde dicit idem ipse Evangelísta in Epístola sua: Novíssima hora est. Ideo namque áddidit: Quia vado ad Patrem: quod ad priórem senténtiam referéndum est, ubi ait: Módicum et jam non vidébitis me: non ad posteriórem, ubi ait: Et íterum módicum, et vidébitis me. Eúndo quippe ad Patrem, factúrus erat ut eum non vidérent. Ac per hoc non ídeo dictum est, quia fúerat moritúrus, et donec resúrgeret, ab eórum aspéctibus recessúrus: sed quod esset itúrus ad Patrem, quod fecit posteáquam resurréxit, et cum eis per quadragínta dies conversátus, ascéndit in cælum.\par
\switchcolumn
\EN
\noindent This little while is the whole duration of this present world. In the same sense this same Evangelist saith in his Epistle (ii. 18): It is the last time. The words, because I go to the Father, refer to the first clause of the text, namely, A little while and ye shall not see Me, and not to the latter clause, that is, and again a little while, and ye shall see Me. By His going to the Father He was about to bring it to pass that they should see Him no more. And thus it was that He said, not that He was about to die, and that after His death they should not see Him until He rose again, but that He was going to the Father, which He did when, after that He was risen again and had manifested Himself to them for forty days, He ascended up into heaven.\par
\switchcolumn*

\LA
\begin{Resp}\Rsign Decantábat pópulus Israël, allelúja: et univérsa multitúdo Jacob canébat legítime: \Star{} Et David cum cantóribus cítharam percutiébat in domo Dómini, et laudes Deo canébat, allelúja, allelúja.\end{Resp}
\begin{Resp}\Vsign Sanctificáti sunt ergo sacerdótes et levítæ: et univérsus Israël deducébat arcam fœ́deris Dómini in júbilo.\end{Resp}
\begin{Resp}\Rsign Et David cum cantóribus cítharam percutiébat in domo Dómini, et laudes Deo canébat, allelúja, allelúja.\end{Resp}
\switchcolumn
\EN
\begin{Resp}\Rsign The people of Israel sung: Alleluia, and all the multitude of Jacob sung in measure. \Star{} And David was with the singers, (and) played upon an harp in the house of the Lord, and sung praises unto God. Alleluia, Alleluia.\end{Resp}
\begin{Resp}\Vsign So the Priests and the Levites were sanctified, and all Israel brought up the ark of the covenant of the Lord with shouting.\end{Resp}
\begin{Resp}\Rsign And David was with the singers, (and) played upon an harp in the house of the Lord, and sung praises unto God. Alleluia, Alleluia.\end{Resp}
\switchcolumn*

\LA
\LessonHead{Lectio VIII}
\switchcolumn
\EN
\LessonHead{Lesson VIII}
\switchcolumn*

\LA
\noindent Illis ergo ait: Módicum, et jam non vidébitis me; qui eum corporáliter tunc vidébant: quia itúrus erat ad Patrem, et eum deínceps mortálem visúri non erant, qualem, cum ista loquerétur, vidébant. Quod vero áddidit: Et íterum módicum, et vidébitis me: univérsæ promísit Ecclésiæ, sicut univérsæ promísit: Ecce ego vobíscum sum usque ad consummatiónem sǽculi. Non tardat Dóminus promíssum. Módicum et vidébimus eum: ubi jam nihil rogémus, nihil interrogémus, quia nihil desiderándum remanébit, nihil quæréndum latébit.\par
\switchcolumn
\EN
\noindent But now, to them which were looking on Him in the Body, He saith: A little while, and ye shall not see Me, a little while, and they who now saw Him clad in a dying nature, should see Him so no more, because He was about to go to the Father. But He saith: And again a little while, and ye shall see Me, and these words are a promise to the Universal Church, just as are those others: Lo, I am with you alway, even unto the end of the world (Matth. xxviii. 20.) Our Lord delayeth not His promised coming. Again a little while, and we shall see Him. We shall see Him. And, O, when we shall see Him, then we shall beg, we shall ask no more; for no desire will be unsatisfied, and no riddle unsolved.\par
\switchcolumn*

\LA
\begin{Resp}\Rsign Tristítia vestra, allelúja, \Star{} Convertétur in gáudium, allelúja, allelúja.\end{Resp}
\begin{Resp}\Vsign Mundus autem gaudébit, vos vero contristabímini, sed tristítia vestra.\end{Resp}
\begin{Resp}\Rsign Convertétur in gáudium, allelúja, allelúja.\end{Resp}
\begin{Resp}\Vsign Glória Patri, et Fílio, \Star{} et Spirítui Sancto.\end{Resp}
\begin{Resp}\Rsign Convertétur in gáudium, allelúja, allelúja.\end{Resp}
\switchcolumn
\EN
\begin{Resp}\Rsign Your sorrow, alleluia, \Star{} Shall be turned into joy. Alleluia.\end{Resp}
\begin{Resp}\Vsign The world shall rejoice; and you shall be made sorrowful: but your sorrow\end{Resp}
\begin{Resp}\Rsign Shall be turned into joy. Alleluia.\end{Resp}
\begin{Resp}\Vsign Glory be to the Father, and to the Son, \Star{} and to the Holy Ghost.\end{Resp}
\begin{Resp}\Rsign Shall be turned into joy. Alleluia.\end{Resp}
\switchcolumn*

\LA
\LessonHead{Lectio IX}
\switchcolumn
\EN
\LessonHead{Lesson IX}
\switchcolumn*

\LA
\noindent Hoc módicum longum nobis vidétur, quóniam adhuc ágitur; cum finítum fúerit, tunc sentiémus quam módicum fúerit. Non ergo sit gáudium nostrum quale habet mundus, de quo dictum est: Mundus autem gaudébit. Nec tamen in hujus desidérii parturitióne sine gáudio tristes simus: sed, sicut ait Apóstolus: Spe gaudéntes: In tribulatióne patiéntes: quia et ipsa partúriens, cui comparáti sumus, plus gaudet de mox futúra prole, quam tristis est de præsénti dolóre. Sed hujus sermónis iste sit finis: habent enim quæstiónem molestíssimam, quæ sequúntur: nec brevitáte coarctánda sunt, ut possint commódius, si Dóminus volúerit, explicári.\par
\switchcolumn
\EN
\noindent This little while seemeth a very long while to us now, while as it is still going on, but when it is over we shall feel indeed how truly it is but a little while. Therefore, may our rejoicing never be like the rejoicing of that world whereof it is said: The world shall rejoice. A woman when she is in travail hath sorrow, and yet, while hitherto our gladness is still coming to the birth through throes of sorrow, let us not be altogether sorrowful, but, as the Apostle hath it (Rom. xii. 12): Rejoicing in hope: patient in tribulation. A woman, when she is in travail hath sorrow, because her hour is come: but as soon as she is delivered of the child, she remembereth no more the anguish, for joy that a man is born into the world. And so will it be with us. And with that let me end my discourse. The next passage is one of extreme difficulty; nor is it possible to treat it briefly, if, (with the will of God,) it is to be treated satisfactorily.\par
\switchcolumn*

\LA
\TeDeum{Te Deum}
\switchcolumn
\EN
\TeDeum{Te Deum}
\switchcolumn*

\end{paracol}

\Office{De VI die infra Octavam S. Joseph}{Sixth Day within the Octave of St. Joseph}{Semiduplex}
\addcontentsline{toc}{subsection}{De VI die infra Octavam S. Joseph\enspace\textit{Sixth Day within the Octave of St. Joseph}}
\begin{paracol}{2}
\LA
\LessonHead{Lectio I}
\switchcolumn
\EN
\LessonHead{Lesson I}
\switchcolumn*

\LA
\LessonTitle{De libro Apocalýpsis beáti Joánnis Apóstoli}
\switchcolumn
\EN
\LessonTitle{Lesson from the book of Revelation}
\switchcolumn*

\LA
\Cite{Apo 2:1-7}
\switchcolumn
\EN
\Cite{Rev 2:1-7}
\switchcolumn*

\LA
\noindent Angelo Ephesi ecclésiæ scribe: Hæc dicit, qui tenet septem stellas in déxtera sua, qui ámbulat in médio septem candelabrórum aureórum: Scio ópera tua, et labórem, et patiéntiam tuam, et quia non potes sustinére malos: et tentásti eos, qui se dicunt apóstolos esse, et non sunt: et invenísti eos mendáces: et patiéntiam habes, et sustinuísti propter nomen meum, et non defecísti. Sed hábeo advérsum te, quod caritátem tuam primam reliquísti. Memor esto ítaque unde excíderis: et age pœniténtiam, et prima ópera fac: sin autem, vénio tibi, et movébo candelábrum tuum de loco suo, nisi pœniténtiam égeris. Sed hoc habes, quia odísti facta Nicolaitárum, quæ et ego odi. Qui habet aurem, áudiat quid Spíritus dicat ecclésiis: Vincénti dabo édere de ligno vitæ, quod est in paradíso Dei mei.\par
\switchcolumn
\EN
\noindent Unto the angel of the church of Ephesus write: These things saith he, who holdeth the seven stars in his right hand, who walketh in the midst of the seven golden candlesticks: I know thy works, and thy labour, and thy patience, and how thou canst not bear them that are evil, and thou hast tried them, who say they are apostles, and are not, and hast found them liars: And thou hast patience, and hast endured for my name, and hast not fainted. But I have somewhat against thee, because thou hast left thy first charity. Be mindful therefore from whence thou art fallen: and do penance, and do the first works. Or else I come to thee, and will move thy candlestick out of its place, except thou do penance. But this thou hast, that thou hatest the deeds of the Nicolaites, which I also hate. He, that hath an ear, let him hear what the Spirit saith to the churches: To him, that overcometh, I will give to eat of the tree of life, which is in the paradise of my God.\par
\switchcolumn*

\LA
\begin{Resp}\Rsign Vidi portam civitátis ad Oriéntem pósitam, et Apostolórum nómina et Agni super eam scripta: \Star{} Et super muros ejus Angelórum custódiam, allelúja.\end{Resp}
\begin{Resp}\Vsign Vidi cælum novum, et terram novam, et civitátem novam descendéntem de cælo.\end{Resp}
\begin{Resp}\Rsign Et super muros ejus Angelórum custódiam, allelúja.\end{Resp}
\switchcolumn
\EN
\begin{Resp}\Rsign I saw the gate of the city, which looketh toward the East, and written thereon the names of the (Twelve) Apostles and of the Lamb. \Star{} And upon the walls thereof a guard of Angels. alleluia.\end{Resp}
\begin{Resp}\Vsign I saw a new heaven and a new earth coming down out of heaven.\end{Resp}
\begin{Resp}\Rsign And upon the walls thereof a guard of Angels. Alleluia.\end{Resp}
\switchcolumn*

\LA
\LessonHead{Lectio II}
\switchcolumn
\EN
\LessonHead{Lesson II}
\switchcolumn*

\LA
\Cite{Apo 2:8-11}
\switchcolumn
\EN
\Cite{Rev 2:8-11}
\switchcolumn*

\LA
\noindent Et ángelo Smyrnæ ecclésiæ scribe: Hæc dicit primus, et novíssimus, qui fuit mórtuus, et vivit: Scio tribulatiónem tuam, et paupertátem tuam, sed dives es: et blasphemáris ab his, qui se dicunt Judǽos esse, et non sunt, sed sunt synagóga Sátanæ. Nihil horum tímeas quæ passúrus es. Ecce missúrus est diábolus áliquos ex vobis in cárcerem ut tentémini: et habébitis tribulatiónem diébus decem. Esto fidélis usque ad mortem, et dabo tibi corónam vitæ. Qui habet aurem, áudiat quid Spíritus dicat ecclésiis: Qui vícerit, non lædétur a morte secúnda.\par
\switchcolumn
\EN
\noindent And to the angel of the church of Smyrna write: These things saith the First and the Last, who was dead, and is alive: I know thy tribulation and thy poverty, but thou art rich: and thou art blasphemed by them that say they are Jews and are not, but are the synagogue of Satan. Fear none of those things which thou shalt suffer. Behold, the devil will cast some of you into prison that you may be tried: and you shall have tribulation ten days. Be thou faithful until death: and I will give thee the crown of life. He, that hath an ear, let him hear what the Spirit saith to the churches: He that shall overcome, shall not be hurt by the second death.\par
\switchcolumn*

\LA
\begin{Resp}\Rsign Osténdit mihi Angelus fontem aquæ vivæ, et dixit ad me, allelúja: \Star{} Hic Deum adóra, allelúja, allelúja, allelúja.\end{Resp}
\begin{Resp}\Vsign Postquam audíssem et vidíssem, cécidi ut adorárem ante pedes Angeli, qui mihi hæc ostendébat, et dixit mihi.\end{Resp}
\begin{Resp}\Rsign Hic Deum adóra, allelúja, allelúja, allelúja.\end{Resp}
\switchcolumn
\EN
\begin{Resp}\Rsign The Angel showed me the fountain of the water of life and he said unto me, Alleluia. \Star{} Here worship God. Alleluia, Alleluia, Alleluia.\end{Resp}
\begin{Resp}\Vsign When I had heard and seen, I fell down to worship before the feet of the Angel, which showed me these things, and he said unto me\end{Resp}
\begin{Resp}\Rsign Here worship God. Alleluia, Alleluia, Alleluia.\end{Resp}
\switchcolumn*

\LA
\LessonHead{Lectio III}
\switchcolumn
\EN
\LessonHead{Lesson III}
\switchcolumn*

\LA
\Cite{Apo 2:12-17}
\switchcolumn
\EN
\Cite{Rev 2:12-17}
\switchcolumn*

\LA
\noindent Et ángelo Pérgami ecclésiæ scribe: Hæc dicit qui habet rhomphǽam utráque parte acútam: Scio ubi hábitas, ubi sedes est Sátanæ: et tenes nomen meum, et non negásti fidem meam. Et in diébus illis Antipas testis meus fidélis, qui occísus est apud vos ubi Sátanas hábitat. Sed hábeo advérsus te pauca: quia habes illic tenéntes doctrínam Bálaam, qui docébat Balac míttere scándalum coram fíliis Israël, édere, et fornicári: ita habes et tu tenéntes doctrínam Nicolaitárum. Simíliter pœniténtiam age: si quóminus véniam tibi cito, et pugnábo cum illis in gládio oris mei. Qui habet aurem, áudiat quid Spíritus dicat ecclésiis: Vincénti dabo manna abscónditum, et dabo illi cálculum cándidum: et in cálculo nomen novum scriptum, quod nemo scit, nisi qui áccipit.\par
\switchcolumn
\EN
\noindent And to the angel of the church of Pergamus write: These things, saith he, that hath the sharp two edged sword: I know where thou dwellest, where the seat of Satan is: and thou holdest fast my name, and hast not denied my faith. Even in those days when Antipas was my faithful witness, who was slain among you, where Satan dwelleth. But I have against thee a few things: because thou hast there them that hold the doctrine of Balaam, who taught Balac to cast a stumblingblock before the children of Israel, to eat, and to commit fornication: So hast thou also them that hold the doctrine of the Nicolaites. In like manner do penance: if not, I will come to thee quickly, and will fight against them with the sword of my mouth. He, that hath an ear, let him hear what the Spirit saith to the churches: To him that overcometh, I will give the hidden manna, and will give him a white counter, and in the counter, a new name written, which no man knoweth, but he that receiveth it.\par
\switchcolumn*

\LA
\begin{Resp}\Rsign Audívi vocem de cælo, tamquam vocem tonítrui magni, allelúja: Regnábit Deus noster in ætérnum, allelúja: \Star{} Quia facta est salus, et virtus, et potéstas Christi ejus, allelúja, allelúja.\end{Resp}
\begin{Resp}\Vsign Et vox de throno exívit, dicens: Laudem dícite Deo nostro, omnes Sancti ejus, et qui timétis Deum, pusílli et magni.\end{Resp}
\begin{Resp}\Rsign Quia facta est salus, et virtus, et potéstas Christi ejus, allelúja, allelúja.\end{Resp}
\begin{Resp}\Vsign Glória Patri, et Fílio, \Star{} et Spirítui Sancto.\end{Resp}
\begin{Resp}\Rsign Quia facta est salus, et virtus, et potéstas Christi ejus, allelúja, allelúja.\end{Resp}
\switchcolumn
\EN
\begin{Resp}\Rsign And I heard a voice from heaven, the voice of great thunder, alleluia: the Lord our God shall reign for ever, alleluia: \Star{} Now is come salvation, and strength, and the power of his Christ, alleluia, alleluia.\end{Resp}
\begin{Resp}\Vsign And a voice came out from the throne, saying: Give praise to our God, all saints of Him that fear him, little and great.\end{Resp}
\begin{Resp}\Rsign Now is come salvation, and strength, and the power of his Christ, alleluia, alleluia.\end{Resp}
\begin{Resp}\Vsign Glory be to the Father, and to the Son, \Star{} and to the Holy Ghost.\end{Resp}
\begin{Resp}\Rsign Now is come salvation, and strength, and the power of his Christ, alleluia, alleluia.\end{Resp}
\switchcolumn*

\LA
\LessonHead{Lectio IV}
\switchcolumn
\EN
\LessonHead{Lesson IV}
\switchcolumn*

\LA
\LessonTitle{Sermo sancti Bernardi Abbátis}
\switchcolumn
\EN
\LessonTitle{From a Sermon by St. Bernard the Abbot}
\switchcolumn*

\LA
\Source{Homilia 2 super Missus est}
\switchcolumn
\EN
\Source{Homilía 2 super Missus est}
\switchcolumn*

\LA
\noindent Desponsáta est María Joseph, vel pótius (sicut ponit Evangelísta) viro, cui nomen erat Joseph. Virum nóminat, non quia marítus, sed quod homo virtútis erat; vel pótius quia juxta álium Evangelístam non vir simplíciter, sed vir ejus dictus est, mérito appellátur quod necessário putátur. Débuit ergo vir ejus appellári, quia necésse fuit et putári; sicut et pater Salvatóris non quidem esse, sed dici méruit, ut putarétur esse, dicénte hoc ipso Evangelísta: Et ipse Jesus erat incípiens quasi annórum trigínta, ut putabátur, fílius Joseph.\par
\switchcolumn
\EN
\noindent Mary was espoused to Joseph, or rather, as saith the Evangelist Luke: To a man whose name was Joseph. He is called a man, not because he was her husband, but because he was a person of manliness. And again, the same is said by the Evangelist Matthew, to wit: Joseph the husband of Mary: and: Joseph her man: for he rightly calleth Joseph by this title of manliness, for so Joseph was expected to be, that his virtuous manhood might be given in marriage to Mary. And we must conclude that he is here called what he was, a man; and further, that he was called her man because it was necessary that he should be publickly accepted as her man. And likewise, he was found worthy to be called the father of the Saviour, not that he was, but that he was publickly accepted as such, as the Evangelist himself saith: And Jesus himself began to be about thirty years of age, being (as it was supposed) the son of Joseph.\par
\switchcolumn*

\LA
\begin{Resp}\Rsign Dedísti mihi protectiónem salútis tuæ, et déxtera tua suscépit me: \Star{} Protéctor meus, et cornu salútis meæ, et suscéptor meus, allelúja.\end{Resp}
\begin{Resp}\Vsign Ego protéctor tuus sum, et merces tua magna nimis.\end{Resp}
\begin{Resp}\Rsign Protéctor meus, et cornu salútis meæ, et suscéptor meus, allelúja.\end{Resp}
\switchcolumn
\EN
\begin{Resp}\Rsign Thou hast given me the shield of thy salvation, and thy right hand hath holden me up. \Star{} My buckler, and the horn of my salvation, and my refuge. Alleluia.\end{Resp}
\begin{Resp}\Vsign I am thy shield and thy exceeding great reward.\end{Resp}
\begin{Resp}\Rsign My buckler, and the horn of my salvation, and my refuge. Alleluia.\end{Resp}
\switchcolumn*

\LA
\LessonHead{Lectio V}
\switchcolumn
\EN
\LessonHead{Lesson V}
\switchcolumn*

\LA
\noindent Non est dúbium, quin bonus et fidélis homo fúerit iste Joseph, cui Mater desponsáta est Salvatóris. Fidélis, inquam, servus et prudens, quem constítuit Dóminus suæ Matris solátium, suæ carnis nutrítium, solum dénique in terris magni consílii coadjutórem fidelíssimum. Huc accédit, quod dícitur fuísse de domo David. Vere enim de domo David, vere de régia stirpe descéndit vir iste Joseph, nóbilis génere, mente nobílior. Plane fílius David, non degénerans a patre suo David; prorsus, inquam, fílius David, non tantum carne, sed fide, sed sanctitáte, sed devotióne, quem tamquam álterum David Dóminus invénit secúndum cor suum; cui tuto commítteret sacratíssimum atque secretíssimum sui cordis arcánum; cui, tamquam álteri David, incérta et occúlta sapiéntiæ suæ manifestávit, et dedit illi non ignárum esse mystérii, quod nemo príncipum hujus sǽculi agnóvit.\par
\switchcolumn
\EN
\noindent Without doubt, good and faithful was this Joseph who espoused the Mother of the Saviour. Yea, I say unto you, he is that faithful and wise servant whom the Lord hath made ruler over his household. For the Lord appointed him to be the comfort of his Mother, the keeper of his own body, and, in a word, his chief and most trusty helper on earth in the carrying out the eternal counsels. Add to this that he is said to have been of the house of David, as he verily was. For this Joseph was a true son of a race of kings, noble in descent, nobler yet in mind. A true son of David, not so much according to the flesh as in faith, holiness, and devotion. Whom, like another David, the Lord found to be a man after his own heart, to whom he therefore safely entrusted the most hold and hidden secret of his heart. To whom also, like another David, he shewed the uncertain and hidden things of his wisdom, and granted that he should not be ignorant of a mystery which was known to none of the princes of this world.\par
\switchcolumn*

\LA
\begin{Resp}\Rsign Státuet fílios suos sub tégmine illíus, et sub ramis ejus morábitur: protegétur sub tégmine illíus a fervóre: \Star{} Et in glória ejus requiéscet, allelúja.\end{Resp}
\begin{Resp}\Vsign Speráte in eo, omnis congregátio pópuli, effúndite coram illo corda vestra.\end{Resp}
\begin{Resp}\Rsign Et in glória ejus requiéscet, allelúja.\end{Resp}
\switchcolumn
\EN
\begin{Resp}\Rsign He shall set his children under her shelter, and shall lodge under her branches by her shall he be covered from heat, \Star{} And in her glory shall he dwell. Alleluia.\end{Resp}
\begin{Resp}\Vsign Trust in Him, ye congregation of the people, pour out your heart before Him.\end{Resp}
\begin{Resp}\Rsign And in her glory shall he dwell. Alleluia.\end{Resp}
\switchcolumn*

\LA
\LessonHead{Lectio VI}
\switchcolumn
\EN
\LessonHead{Lesson VI}
\switchcolumn*

\LA
\noindent Cui dénique datum est, quod multi reges et prophétæ, cum vellent vidére, non vidérunt, audíre et non audiérunt; non solum vidére et audíre, sed étiam portáre, dedúcere, amplécti, deosculári, et nutríre et custodíre. Non tantum autem Joseph, sed et María descendísse credénda est de domo David; alióquin non esset desponsáta viro de domo David, si non esset et ipsa de domo David. Ambo ígitur erant de domo David; sed in áltera compléta est véritas, quam jurávit Dóminus David, áltero tamen cónscio et teste adimplétæ promissiónis.\par
\switchcolumn
\EN
\noindent Lastly, there was given to him not only to see and hear, him whom many kings desired to see yet saw not, and to hear yet heard not, but even to carry him in his arms, to kiss him with his lips, to clothe him and to guard him. We must believe that Mary too, like Joseph, was descended from the house of David. For she would not have been espoused to a man of the house of David, if she had not herself been of the house of David. Both therefore were of the house of David. But in Mary the truth, which the Lord had sworn to David, was fulfilled; whereas to Joseph it was therefore given to know and bear witness unto the fulfilment of the promise.\par
\switchcolumn*

\LA
\begin{Resp}\Rsign Si consístant advérsum me castra, non timébit cor meum: \Star{} Si exsúrgat advérsum me prǽlium, in hoc ego sperábo, allelúja.\end{Resp}
\begin{Resp}\Vsign In te cantátio mea semper, quóniam tu adjútor fortis.\end{Resp}
\begin{Resp}\Rsign Si exsúrgat advérsum me prǽlium, in hoc ego sperábo, allelúja.\end{Resp}
\begin{Resp}\Vsign Glória Patri, et Fílio, \Star{} et Spirítui Sancto.\end{Resp}
\begin{Resp}\Rsign Si exsúrgat advérsum me prǽlium, in hoc ego sperábo, allelúja.\end{Resp}
\switchcolumn
\EN
\begin{Resp}\Rsign Though an host should encamp against me, my heart shall not fear. \Star{} Though war should rise against me, in this will I be confident. Alleluia.\end{Resp}
\begin{Resp}\Vsign My praise shall be continually of thee, for Thou art my strong refuge. R.* Though war should rise against me, in this will I be confident. Alleluia.\end{Resp}
\begin{Resp}\Vsign Glory be to the Father, and to the Son, \Star{} and to the Holy Ghost.\end{Resp}
\begin{Resp}\Rsign Though war should rise against me, in this will I be confident. Alleluia.\end{Resp}
\switchcolumn*

\LA
\LessonHead{Lectio VII}
\switchcolumn
\EN
\LessonHead{Lesson VII}
\switchcolumn*

\LA
\LessonTitle{Léctio sancti Evangélii secúndum Lucam}
\switchcolumn
\EN
\LessonTitle{From the Holy Gospel according to Luke}
\switchcolumn*

\LA
\Cite{Luc 3:21-23}
\switchcolumn
\EN
\Cite{Luke 3:21-23}
\switchcolumn*

\LA
\noindent In illo témpore: Factum est autem cum baptizarétur omnis pópulus, et Jesu baptizáto et oránte, apértum est cælum. Et réliqua.\par
\switchcolumn
\EN
\noindent At that time: When all the people were baptized, it came to pass, that Jesus also being baptized and praying, the heaven was opened. And so on, and that which followeth.\par
\switchcolumn*

\LA
\LessonTitle{De Homilía sancti Ambrósii Epíscopi}
\switchcolumn
\EN
\LessonTitle{A Homily by St. Ambrose the Bishop}
\switchcolumn*

\LA
\Source{Expositio in Luc. lib. 3}
\switchcolumn
\EN
\Source{Expositio in Luc. lib. 3}
\switchcolumn*

\LA
\noindent Quod per Salomónem Matthǽus generatiónem derivándum putávit, Lucas vero per Nathan, álteram regálem, álteram sacerdotálem Christi famíliam vidétur osténdere. Quod non ita accípere debémus, quod álterum áltero vérius; sed alter álteri pari fide et veritáte concórdet. Fuit enim vere et secúndum carnem regális et sacerdotális famíliæ, Rex ex régibus, Sacérdos ex sacerdótibus; licet oráculum non de carnálibus, sed de cæléstibus exprimátur: quóniam et Rex in Dei virtúte lætátur, cui judícium a Patre Rege defértur, et Sacérdos est in ætérnum, secúndum quod scriptum est: Tu es Sacérdos in ætérnum secúndum órdinem Melchísedech.\par
\switchcolumn
\EN
\noindent That Matthew should trace the lineage of Christ through Solomon, and Luke through Nathan, would seem to indicate that the one desireth to shew the royalty of Christ's descent, and the other the priestliness thereof. We need not infer from this that one is more accurate than the other. On the contrary, each agreeth with the other, with an equal good faith and veracity. For he was indeed, according to the flesh, of both a royal and a priestly family, a King sprung from kings, a Priest from priests. But the voice from heaven is speaking of divine things rather than of human. So then, as it is written: The King shall rejoice in God: that is, in God's strength, from which come to him the judgments of his royal Father; and likewise, he is that Priest of whom it is written: Thou art a Priest for ever after the order of Melchisedech.\par
\switchcolumn*

\LA
\begin{Resp}\Rsign Joseph, fili David, noli timére accípere Maríam cónjugem tuam; quod enim in ea natum est, de Spíritu Sancto est: páriet autem fílium, \Star{} Et vocábis nomen ejus Jesum, allelúja.\end{Resp}
\begin{Resp}\Vsign Ipse enim salvum fáciet pópulum suum a peccátis eórum.\end{Resp}
\begin{Resp}\Rsign Et vocábis nomen ejus Jesum, allelúja.\end{Resp}
\switchcolumn
\EN
\begin{Resp}\Rsign Joseph, thou Son of David, fear not to take unto thee Mary thy wife; for That Which is conceived in her is of the Holy Ghost and she shall bring forth a Son; \Star{} And thou shalt call His Name Jesus. Alleluia.\end{Resp}
\begin{Resp}\Vsign For He shall save His people from their sins.\end{Resp}
\begin{Resp}\Rsign And thou shalt call His Name Jesus. Alleluia.\end{Resp}
\switchcolumn*

\LA
\LessonHead{Lectio VIII}
\switchcolumn
\EN
\LessonHead{Lesson VIII}
\switchcolumn*

\LA
\noindent Bene ígitur utérque ténuit fidem, ut Matthǽus per reges ductam oríginem comprobáret, et Lucas per sacerdótes a Deo transmíssam in Christum sériem géneris deducéndo, sanctiórem ipsam oríginem declaráret. Simul in hoc quoque vítuli figúra signátur, quod ubíque sacerdotále mystérium putat esse servándum. Nec miréris, si ab Abraham plures secúndum Lucam successiónes usque ad Christum sunt, paucióres secúndum Matthǽum, cum per álias persónas generatiónem fateáris esse decúrsam. Potest enim fíeri, ut álii longǽvam transégerint vitam, altérius vero generatiónis viri immatúra ætáte decésserint; cum videámus complúres senes cum suis nepótibus vívere, álios vero viros statim fíliis obíre suscéptis\par
\switchcolumn
\EN
\noindent Therefore both Evangelists keep well within the truth. For Matthew doth establish descent through the kings; whereas Luke, by tracing through the priests the lineage transmitted to Christ from God, manifesteth his more sacred origin. And from this we perceive the significance of the symbol used for this Evangelist; namely, the sacrifícial calf , for everywhere he bringeth forward the mystery of the sacrifícial priesthood. Nor need it surprise us that Luke giveth many more generations from Abraham to Christ than doth Matthew, since we can recognize that the line of descent is led through different persons. It may be that some lived long lives, whilst persons of the other line died young. For we are used to seeing many old men living with their grandchildren, and others dying soon after the birth of their children.\par
\switchcolumn*

\LA
\begin{Resp}\Rsign Surge, et áccipe Púerum et matrem ejus, et fuge in Ægýptum: \Star{} Et esto ibi usque dum dicam tibi, allelúja.\end{Resp}
\begin{Resp}\Vsign Ut adimplerétur quod dictum est a Dómino per Prophétam dicéntem: Ex Ægýpto vocávi Fílium meum.\end{Resp}
\begin{Resp}\Rsign Et esto ibi usque dum dicam tibi, allelúja.\end{Resp}
\begin{Resp}\Vsign Glória Patri, et Fílio, \Star{} et Spirítui Sancto.\end{Resp}
\begin{Resp}\Rsign Et esto ibi usque dum dicam tibi, allelúja.\end{Resp}
\switchcolumn
\EN
\begin{Resp}\Rsign Arise, and take the young Child, and His Mother, and flee into Egypt; \Star{} And be thou there until I bring thee word. Alleluia.\end{Resp}
\begin{Resp}\Vsign That it might be fulfilled which was spoken of the Lord by the Prophets, saying Out of Egypt have I called My Son.\end{Resp}
\begin{Resp}\Rsign And be thou there until I bring thee word. Alleluia.\end{Resp}
\begin{Resp}\Vsign Glory be to the Father, and to the Son, \Star{} and to the Holy Ghost.\end{Resp}
\begin{Resp}\Rsign And be thou there until I bring thee word. Alleluia.\end{Resp}
\switchcolumn*

\LA
\LessonHead{Lectio IX}
\switchcolumn
\EN
\LessonHead{Lesson IX}
\switchcolumn*

\LA
\noindent Illud quoque advértimus, quod sanctus Matthǽus, Jacob, qui fuit pater Joseph, fílium Mathat esse memoráverit; Lucas vero Joseph, cui desponsáta erat María, fílium Heli, Heli autem fílium Melchi esse descrípserit. Quómodo uníus duo patres, id est, Heli et Jacob? Duórum fílius dictus est, quia altérius secúndum generatiónem, altérius secúndum legem factus est fílius. In quo præscrípto legis futúram perpetuitátem defunctórum séminis nobis esse promíssam non intelléxit pópulus Judæórum, sed secúndum lítteram accípiens, grátiam corrúpit oráculi. Alius enim erat frater, qui defunctórum fratrum semen suscitáret, non frater secúndum carnis germanitátem, sed secúndum grátiæ puritátem. Et ídeo fortásse Frater non redémit, redémit homo; quia non germánus frater ille, sed mediátor Dei et hóminum, homo Christus Jesus, resurrectiónis grátiam propagávit.\par
\switchcolumn
\EN
\noindent We notice also a further difference. Saint Matthew saith that Jacob, the son of Matthan, was the father of Joseph. Whereas Luke saith that Joseph, to whom Mary was espoused, was the son of Heli, and that Heli was the son of Matthat. How then could Joseph have had two fathers, that is Heli and Jacob? Perchance he is called the son of two men, because one was his father according to nature, whereas the other became his father according to the Law. The particulars of the Law regarding the raising up of seed to a dead brother were not understood by the Jewish people as a promise to us that the seed of the dead should be perpetuated for ever. But in so far as they read it only according to the letter, they failed to grasp its revelation of spiritual truth. For the living brother that raised up seed unto his dead brother, is not to be considered a brother after the flesh, but only according to the purity of his motives. And on that account, perchance we read: But no man may deliver his brother, nor make agreement unto God for him (for it cost more to redeem their souls, so that he must let that alone for ever): yea, though he live long, and see not the grave. For the man Christ Jesus was not our natural brother, but the Mediator between God and man, whereby he hath engendered in us the grace of the resurrection unto perpetual life.\par
\switchcolumn*

\LA
\TeDeum{Te Deum}
\switchcolumn
\EN
\TeDeum{Te Deum}
\switchcolumn*

\end{paracol}

\Office{De VII die infra Octavam S. Joseph}{Seventh Day within the Octave of St. Joseph}{Semiduplex}
\addcontentsline{toc}{subsection}{De VII die infra Octavam S. Joseph\enspace\textit{Seventh Day within the Octave of St. Joseph}}
\begin{paracol}{2}
\LA
\LessonHead{Lectio I}
\switchcolumn
\EN
\LessonHead{Lesson I}
\switchcolumn*

\LA
\LessonTitle{De libro Apocalýpsis beáti Joánnis Apóstoli}
\switchcolumn
\EN
\LessonTitle{Lesson from the book of Revelation}
\switchcolumn*

\LA
\Cite{Apo 4:1-5}
\switchcolumn
\EN
\Cite{Rev 4:1-5}
\switchcolumn*

\LA
\noindent Post hæc vidi: et ecce óstium apértum in cælo, et vox prima, quam audívi tamquam tubæ loquéntis mecum, dicens: Ascénde huc, et osténdam tibi quæ opórtet fíeri post hæc. Et statim fui in spíritu: et ecce sedes pósita erat in cælo, et supra sedem sedens. Et qui sedébat símilis erat aspéctui lápidis jáspidis, et sárdinis: et iris erat in circúitu sedis símilis visióni smarágdinæ. Et in circúitu sedis sedília vigínti quátuor: et super thronos vigínti quátuor senióres sedéntes, circumamícti vestiméntis albis, et in capítibus eórum corónæ áureæ. Et de throno procedébant fúlgura, et voces, et tonítrua: et septem lámpades ardéntes ante thronum, qui sunt septem spíritus Dei.\par
\switchcolumn
\EN
\noindent After these things I looked, and behold a door was opened in heaven, and the first voice which I heard, as it were, of a trumpet speaking with me, said: Come up hither, and I will shew thee the things which must be done hereafter. And immediately I was in the spirit: and behold there was a throne set in heaven, and upon the throne one sitting. And he that sat, was to the sight like the jasper and the sardine stone; and there was a rainbow round about the throne, in sight like unto an emerald. And round about the throne were four and twenty seats; and upon the seats, four and twenty ancients sitting, clothed in white garments, and on their heads were crowns of gold. And from the throne proceeded lightnings, and voices, and thunders; and there were seven lamps burning before the throne, which are the seven spirits of God.\par
\switchcolumn*

\LA
\begin{Resp}\Rsign Vidi Jerúsalem descendéntem de cælo, ornátam auro mundo, et lapídibus pretiósis intéxtam: \Star{} Allelúja, allelúja.\end{Resp}
\begin{Resp}\Vsign Et erat structúra muri ejus ex lápide jáspide; ipsa vero aurum mundum, símile vitro mundo.\end{Resp}
\begin{Resp}\Rsign Allelúja, allelúja.\end{Resp}
\switchcolumn
\EN
\begin{Resp}\Rsign I saw Jerusalem coming down out of heaven, adorned with pure gold, and garnished with precious stones. \Star{} Alleluia, Alleluia.\end{Resp}
\begin{Resp}\Vsign The building of the wall of it was of jasper and the city was pure gold, like unto clear glass.\end{Resp}
\begin{Resp}\Rsign Alleluia, Alleluia.\end{Resp}
\switchcolumn*

\LA
\LessonHead{Lectio II}
\switchcolumn
\EN
\LessonHead{Lesson II}
\switchcolumn*

\LA
\Cite{Apo 4:6-8}
\switchcolumn
\EN
\Cite{Rev 4:6-8}
\switchcolumn*

\LA
\noindent Et in conspéctu sedis tamquam mare vítreum símile crystállo: et in médio sedis, et in circúitu sedis quátuor animália plena óculis ante et retro. Et ánimal primum símile leóni, et secúndum ánimal símile vítulo, et tértium ánimal habens fáciem quasi hóminis, et quartum ánimal símile áquilæ volánti. Et quátuor animália, síngula eórum habébant alas senas: et in circúitu, et intus plena sunt óculis: et réquiem non habébant die ac nocte, dicéntia: Sanctus, Sanctus, Sanctus Dóminus Deus omnípotens, qui erat, et qui est, et qui ventúrus est.\par
\switchcolumn
\EN
\noindent And in the sight of the throne was, as it were, a sea of glass like to crystal; and in the midst of the throne, and round about the throne, were four living creatures, full of eyes before and behind. And the first living creature was like a lion: and the second living creature like a calf: and the third living creature, having the face, as it were, of a man: and the fourth living creature was like an eagle flying. And the four living creatures had each of them six wings; and round about and within they are full of eyes. And they rested not day and night, saying: Holy, holy, holy, Lord God Almighty, who was, and who is, and who is to come.\par
\switchcolumn*

\LA
\begin{Resp}\Rsign In diadémate cápitis Aaron magnificéntia Dómini sculpta erat: \Star{} Dum perficerétur opus Dei, allelúja, allelúja, allelúja.\end{Resp}
\begin{Resp}\Vsign In veste enim podéris quam habébat, totus erat orbis terrárum, et paréntum magnália in quátuor ordínibus lápidum sculpta erant.\end{Resp}
\begin{Resp}\Rsign Dum perficerétur opus Dei, allelúja, allelúja, allelúja.\end{Resp}
\switchcolumn
\EN
\begin{Resp}\Rsign Upon the diadem of Aaron's head was graven the Majesty of the Lord, \Star{} While as the work of God was in doing. Alleluia, Alleluia, Alleluia.\end{Resp}
\begin{Resp}\Vsign For in the long garment which he had, was the whole world, and in the four rows of the stones was the glory of the fathers graven.\end{Resp}
\begin{Resp}\Rsign While as the work of God was in doing. Alleluia, Alleluia, Alleluia.\end{Resp}
\switchcolumn*

\LA
\LessonHead{Lectio III}
\switchcolumn
\EN
\LessonHead{Lesson III}
\switchcolumn*

\LA
\Cite{Apo 4:9-11}
\switchcolumn
\EN
\Cite{Rev 4:9-11}
\switchcolumn*

\LA
\noindent Et cum darent illa animália glóriam, et honórem, et benedictiónem sedénti super thronum, vivénti in sǽcula sæculórum, procidébant vigínti quátuor senióres ante sedéntem in throno, et adorábant vivéntem in sǽcula sæculórum, et mittébant corónas suas ante thronum, dicéntes: Dignus es, Dómine, Deus noster, accípere glóriam, et honórem, et virtútem: quia tu creásti ómnia, et propter voluntátem tuam erant, et creáta sunt.\par
\switchcolumn
\EN
\noindent And when those living creatures gave glory, and honour, and benediction to him that sitteth on the throne, who liveth for ever and ever; The four and twenty ancients fell down before him that sitteth on the throne, and adored him that liveth for ever and ever, and cast their crowns before the throne, saying: Thou art worthy, O Lord our God, to receive glory, and honour, and power: because thou hast created all things; and for thy will they were, and have been created.\par
\switchcolumn*

\LA
\begin{Resp}\Rsign Véniens a Líbano quam pulchra facta est, allelúja: \Star{} Et odor vestimentórum ejus super ómnia arómata, allelúja, allelúja.\end{Resp}
\begin{Resp}\Vsign Favus distíllans lábia ejus, mel et lac sub lingua ejus.\end{Resp}
\begin{Resp}\Rsign Et odor vestimentórum ejus super ómnia arómata, allelúja, allelúja.\end{Resp}
\begin{Resp}\Vsign Glória Patri, et Fílio, \Star{} et Spirítui Sancto.\end{Resp}
\begin{Resp}\Rsign Et odor vestimentórum ejus super ómnia arómata, allelúja, allelúja.\end{Resp}
\switchcolumn
\EN
\begin{Resp}\Rsign Coming from Lebanon how sweet she is, alleluia: \Star{} And the sweet smell of her garments, above all aromatical spices, alleluia, alleluia.\end{Resp}
\begin{Resp}\Vsign Thy lips are as a dropping honeycomb, honey and milk are under thy tongue.\end{Resp}
\begin{Resp}\Rsign And the sweet smell of her garments, above all aromatical spices, alleluia, alleluia.\end{Resp}
\begin{Resp}\Vsign Glory be to the Father, and to the Son, \Star{} and to the Holy Ghost.\end{Resp}
\begin{Resp}\Rsign And the sweet smell of her garments, above all aromatical spices, alleluia, alleluia.\end{Resp}
\switchcolumn*

\LA
\LessonHead{Lectio IV}
\switchcolumn
\EN
\LessonHead{Lesson IV}
\switchcolumn*

\LA
\LessonTitle{De Sermóne sancti Bernárdi Abbátis}
\switchcolumn
\EN
\LessonTitle{The Lesson is taken from a Sermon by St. Bernard the Abbot}
\switchcolumn*

\LA
\Source{Homilia 2 super Missus est}
\switchcolumn
\EN
\Source{Homilia 2 super Missus est}
\switchcolumn*

\LA
\noindent Scriptum est: Joseph autem vir ejus, cum esset justus et nollet eam tradúcere, vóluit occúlte dimíttere eam. Bene, cum esset justus, nóluit eam tradúcere; quia, sicut nequáquam justus esset, si cógnitam ream consensísset; sic nihilóminus justus non esset, si probátam innóxiam condemnásset. Cum ergo justus esset et nollet eam tradúcere, vóluit occúlte dimíttere eam. Quare vóluit dimíttere? Accipe et in hoc, non meam, sed Patrum senténtiam. Propter hoc Joseph vóluit dimíttere eam, propter quod et Petrus Dóminum a se repellébat, dicens: Exi a me, Dómine, quia homo peccátor sum; propter quod et centúrio a domo sua eum prohibébat, cum díceret: Dómine, non sum dignus, ut intres sub tectum meum.\par
\switchcolumn
\EN
\noindent It is written: Joseph, her husband, being a just man, and not willing to make her a publick example, was minded to put her away privily. Being a just man, he was rightly unwilling to expose her. For as he would not have been a just man if he had connived at known guilt, so he would have been even less just if he had condemned proven innocence. Being a just man, therefore, and not willing to make her a public example, he was minded to put her away privily. Why did he wish to put her away? On this point hear, not my opinion, but that of the Fathers. Perchance Joseph wished to put her away for the same reason of reverence that made Peter seek to put away the Lord, when he said: Depart from me, for I am a sinful man, O Lord: just as the centurion also sought to keep the Lord away from his house, when he said: Lord, I am not worthy that thou shouldst come under my roof.\par
\switchcolumn*

\LA
\begin{Resp}\Rsign Dedísti mihi protectiónem salútis tuæ, et déxtera tua suscépit me: \Star{} Protéctor meus, et cornu salútis meæ, et suscéptor meus, allelúja.\end{Resp}
\begin{Resp}\Vsign Ego protéctor tuus sum, et merces tua magna nimis.\end{Resp}
\begin{Resp}\Rsign Protéctor meus, et cornu salútis meæ, et suscéptor meus, allelúja.\end{Resp}
\switchcolumn
\EN
\begin{Resp}\Rsign Thou hast given me the shield of thy salvation, and thy right hand hath holden me up. \Star{} My buckler, and the horn of my salvation, and my refuge. Alleluia.\end{Resp}
\begin{Resp}\Vsign I am thy shield and thy exceeding great reward.\end{Resp}
\begin{Resp}\Rsign My buckler, and the horn of my salvation, and my refuge. Alleluia.\end{Resp}
\switchcolumn*

\LA
\LessonHead{Lectio V}
\switchcolumn
\EN
\LessonHead{Lesson V}
\switchcolumn*

\LA
\noindent Ita ergo et Joseph, indígnum et peccatórem se réputans, dicébat intra se, a tali et tanta non debére sibi ultra familiáre præstári contubérnium, cujus supra se mirábilem expavescébat dignitátem. Vidébat et horrébat divínæ præséntiæ certíssimum gestántem insígne; et quia mystérium penetráre non póterat, volébat dimíttere eam. Expávit Petrus poténtiæ magnitúdinem: expávit centúrio præséntiæ majestátem: exhórruit nimírum et Joseph, sicut homo, hujus tanti miráculi novitátem. Miráris quod Joseph prægnántis se consórtio Vírginis judicábat indígnum, cum áudias et sanctam Elísabeth ejus non posse ferre præséntiam, nisi cum tremóre quidem et reveréntia? Ait namque: Unde hoc mihi, ut véniat Mater Dómini mei ad me?\par
\switchcolumn
\EN
\noindent In like manner Joseph may have held himself to be sinful and unworthy, in such wise that he thought he ought no longer to enjoy the familiar companionship of her whose marvellous dignity filled him with awe. Perchance he saw and trembled at the unmistakeable signs of the divine presence; and, since he could not fathom the mystery, he was minded to put her away. Peter trembled at the greatness of the divine power. The Centurion trembled at the presence of the divine Majesty. Joseph too, being but a man, was filled with awe at the strangeness of this mystery. Dost thou wonder that Joseph judged himself unworthy of the companionship of this pregnant Virgin, when thou hearest that Saint Elizabeth too was filled with reverence and trembling at her presence? For she said: Whence is this to me, that the Mother of my Lord should come to me?\par
\switchcolumn*

\LA
\begin{Resp}\Rsign Státuet fílios suos sub tégmine illíus, et sub ramis ejus morábitur: protegétur sub tégmine illíus a fervóre: \Star{} Et in glória ejus requiéscet, allelúja.\end{Resp}
\begin{Resp}\Vsign Speráte in eo, omnis congregátio pópuli, effúndite coram illo corda vestra.\end{Resp}
\begin{Resp}\Rsign Et in glória ejus requiéscet, allelúja.\end{Resp}
\switchcolumn
\EN
\begin{Resp}\Rsign He shall set his children under her shelter, and shall lodge under her branches by her shall he be covered from heat, \Star{} And in her glory shall he dwell. Alleluia.\end{Resp}
\begin{Resp}\Vsign Trust in Him, ye congregation of the people, pour out your heart before Him.\end{Resp}
\begin{Resp}\Rsign And in her glory shall he dwell. Alleluia.\end{Resp}
\switchcolumn*

\LA
\LessonHead{Lectio VI}
\switchcolumn
\EN
\LessonHead{Lesson VI}
\switchcolumn*

\LA
\noindent Ideo ítaque Joseph vóluit dimíttere eam. Sed quare occúlte, et non palam? Ne vidélicet divórtii causa inquirerétur, exigerétur rátio. Quid enim vir justus respónderet pópulo duræ cervícis, pópulo non credénti et contradicénti? Si díceret quod sentiébat, quod de illíus puritáte comprobáverat, nonne mox incréduli et crudéles Judæi subsannárent illum, lapidárent illam? Quómodo namque Veritáti créderent tacénti in útero, quam póstea contempsérunt clamántem in templo? quid fácerent necdum apparénti, cui póstmodum ímpias manus injecérunt étiam miráculis coruscánti? Mérito ergo Vir justus, ne aut mentíri aut diffamáre cogerétur innóxiam, vóluit occúlte dimíttere eam.\par
\switchcolumn
\EN
\noindent And so Joseph was minded to put her away. But why privily, and not publickly? Lest perhaps enquiry should be made about this separation, and he should be asked for reasons. What should a just man reply to a stiffnecked people, a faithless and perverse generation? If he were to have said what he thought, and what he had proved, concerning her purity, would not all the cruel and unbelieving amongst the Jews have soon laughed him to scorn, and stoned her to death? How would they have believed in the Truth lying silent in her womb, when they afterwards despised the Truth preaching in the temple? What would they have done to him before his appearance in the flesh, when afterwards they laid impious hands on him in spite of his signs and wonders? The just man, therefore, was right in wishing to put her away privily, lest he should be thought to lie, or to defame an innocent woman.\par
\switchcolumn*

\LA
\begin{Resp}\Rsign Si consístant advérsum me castra, non timébit cor meum: \Star{} Si exsúrgat advérsum me prǽlium, in hoc ego sperábo, allelúja.\end{Resp}
\begin{Resp}\Vsign In te cantátio mea semper, quóniam tu adjútor fortis.\end{Resp}
\begin{Resp}\Rsign Si exsúrgat advérsum me prǽlium, in hoc ego sperábo, allelúja.\end{Resp}
\begin{Resp}\Vsign Glória Patri, et Fílio, \Star{} et Spirítui Sancto.\end{Resp}
\begin{Resp}\Rsign Si exsúrgat advérsum me prǽlium, in hoc ego sperábo, allelúja.\end{Resp}
\switchcolumn
\EN
\begin{Resp}\Rsign Though an host should encamp against me, my heart shall not fear. \Star{} Though war should rise against me, in this will I be confident. Alleluia.\end{Resp}
\begin{Resp}\Vsign My praise shall be continually of thee, for Thou art my strong refuge. R.* Though war should rise against me, in this will I be confident. Alleluia.\end{Resp}
\begin{Resp}\Vsign Glory be to the Father, and to the Son, \Star{} and to the Holy Ghost.\end{Resp}
\begin{Resp}\Rsign Though war should rise against me, in this will I be confident. Alleluia.\end{Resp}
\switchcolumn*

\LA
\LessonHead{Lectio VII}
\switchcolumn
\EN
\LessonHead{Lesson VII}
\switchcolumn*

\LA
\LessonTitle{Léctio sancti Evangélii secúndum Lucam}
\switchcolumn
\EN
\LessonTitle{From the Holy Gospel according to Luke}
\switchcolumn*

\LA
\Cite{Luc 3:21-23}
\switchcolumn
\EN
\Cite{Luke 3:21-23}
\switchcolumn*

\LA
\noindent In illo témpore: Factum est autem cum baptizarétur omnis pópulus, et Jesu baptizáto et oránte, apértum est cælum. Et réliqua.\par
\switchcolumn
\EN
\noindent At that time: When all the people were baptized, it came to pass, that Jesus also being baptized and praying, the heaven was opened. And so on, and that which followeth.\par
\switchcolumn*

\LA
\LessonTitle{Homilía sancti Joánnis Damascéni}
\switchcolumn
\EN
\LessonTitle{A Homily by St. John of Damascus}
\switchcolumn*

\LA
\Source{Oratio 3 de B.M.V. Nativ.}
\switchcolumn
\EN
\Source{Oratio 3 de B.M.V. Nativ.}
\switchcolumn*

\LA
\noindent Matthæus opus suum ínchoans, Liber generatiónis, inquit, Jesu Christi fílii David, fílii Abraham: verum hic non subsístit; étenim ipsíus sermo étiam usque ad Vírginis Sponsum progréssus est. Lucas autem post Salvatóris in baptísmo declaratiónem, oratiónem suam non nihil derívans, ad hunc modum scribit: Ipse Jesus erat incípiens quasi annórum trigínta, ut putabátur, fílius Joseph, qui fuit Heli, qui fuit Mathat; et sic deínceps in altum ascendéndo usque ad Seth, qui fuit Adæ, qui fuit Dei. Proínde cum Joséphi genus ad hunc modum recenseátur, certe Virgo quoque ipsa ac Dei Génitrix María ejúsdem cum eo tribus esse simul demonstrátur. Síquidem Móysis lege non licébat ulli tríbui cum áltera tribu permiscéri, ne géneris heréditas ab una tribu ad áliam laberétur.\par
\switchcolumn
\EN
\noindent Matthew beginneth his Gospel with the words: The book of the generation of Jesus Christ, the son of David, the son of Abraham. But he doth not stop here. In fact, he continueth his genealogy down to the very Spouse of the Virgin. Luke, on the other hand, after relating the manifestation of the Saviour at his baptism, maketh a digression in his account, and writeth thus: And Jesus himself began to be about thirty years of age, being (as was supposed) the son of Joseph, which was the son of Heli, which was the son of Matthat: and so on, in an ascending line, going up even to Seth: which was the son of Adam, which was the Son of God. Thus after reckoning up Joseph's genealogy in this fashion, we are shewn clearly at the same time how Mary, the Virgin Mother of God, was herself also of the same lineage as Joseph. For the Mosaic Law straitly forbade marriages between the different tribes, in order that the hereditary rights of one tribe might not pass into another.\par
\switchcolumn*

\LA
\begin{Resp}\Rsign Joseph, fili David, noli timére accípere Maríam cónjugem tuam; quod enim in ea natum est, de Spíritu Sancto est: páriet autem fílium, \Star{} Et vocábis nomen ejus Jesum, allelúja.\end{Resp}
\begin{Resp}\Vsign Ipse enim salvum fáciet pópulum suum a peccátis eórum.\end{Resp}
\begin{Resp}\Rsign Et vocábis nomen ejus Jesum, allelúja.\end{Resp}
\switchcolumn
\EN
\begin{Resp}\Rsign Joseph, thou Son of David, fear not to take unto thee Mary thy wife; for That Which is conceived in her is of the Holy Ghost and she shall bring forth a Son; \Star{} And thou shalt call His Name Jesus. Alleluia.\end{Resp}
\begin{Resp}\Vsign For He shall save His people from their sins.\end{Resp}
\begin{Resp}\Rsign And thou shalt call His Name Jesus. Alleluia.\end{Resp}
\switchcolumn*

\LA
\LessonHead{Lectio VIII}
\switchcolumn
\EN
\LessonHead{Lesson VIII}
\switchcolumn*

\LA
\noindent Non abs re Christi ex Spíritu Sancto natívitas apud vulgus silebátur, ac Joseph loco patris assumebátur; atque inde, ut vero consentáneum est, tamquam Púeri pater recensebátur. Nam, nisi hoc exstitísset, patre carére Puer existimátus fuísset, quod a patérno látere genus ipsíus mínime recenserétur. Quam ob rem ab exímiis Evangelístis necessário tunc factum est, ut Joséphi genus recensérent. Nam si, hoc prætermísso, a matérno látere géneris ipsíus sériem texuíssent, prætérquam quod istud indecórum fuísset, a divinárum quoque Scripturárum consuetúdine abhorruísset. Cómmode ígitur Joséphi genus a Davíde, ob eam quam attúlimus causam, ducéntes, simul quoque Vírginem Maríam ex Davíde ortam esse confírmant, per Sponsum scílicet uxóris genus una inferéntes.\par
\switchcolumn
\EN
\noindent Note that there was good reason for these following things: namely, that the birth of Christ by the power of the Holy Ghost was kept secret from the people; and that Joseph stood to Jesus in the place of a father; and that on that account, as was truly fitting, he was counted the father of the Child. Otherwise, it would have seemed that the Child had no father, because he had no recorded descent from his father's side. Wherefore it was of the utmost importance that the Evangelists should record Joseph's lineage. Had they not done so, but had given the Child's lineage on his mother's side, they would have done unseemly, and gone contrary to the usage of divine Scripture. It was therefore fitting that they should give the lineage of Joseph from David, for the reason which we have already given of the kinship between her and her husband. They thereby attest that the Virgin Mary was of the lineage of David.\par
\switchcolumn*

\LA
\begin{Resp}\Rsign Surge, et áccipe Púerum et matrem ejus, et fuge in Ægýptum: \Star{} Et esto ibi usque dum dicam tibi, allelúja.\end{Resp}
\begin{Resp}\Vsign Ut adimplerétur quod dictum est a Dómino per Prophétam dicéntem: Ex Ægýpto vocávi Fílium meum.\end{Resp}
\begin{Resp}\Rsign Et esto ibi usque dum dicam tibi, allelúja.\end{Resp}
\begin{Resp}\Vsign Glória Patri, et Fílio, \Star{} et Spirítui Sancto.\end{Resp}
\begin{Resp}\Rsign Et esto ibi usque dum dicam tibi, allelúja.\end{Resp}
\switchcolumn
\EN
\begin{Resp}\Rsign Arise, and take the young Child, and His Mother, and flee into Egypt; \Star{} And be thou there until I bring thee word. Alleluia.\end{Resp}
\begin{Resp}\Vsign That it might be fulfilled which was spoken of the Lord by the Prophets, saying Out of Egypt have I called My Son.\end{Resp}
\begin{Resp}\Rsign And be thou there until I bring thee word. Alleluia.\end{Resp}
\begin{Resp}\Vsign Glory be to the Father, and to the Son, \Star{} and to the Holy Ghost.\end{Resp}
\begin{Resp}\Rsign And be thou there until I bring thee word. Alleluia.\end{Resp}
\switchcolumn*

\LA
\LessonHead{Lectio IX}
\switchcolumn
\EN
\LessonHead{Lesson IX}
\switchcolumn*

\LA
\noindent Quod quidem justítia præditus Joseph esset, ac vitam legi consentáneam dúceret, némini obscúrum est. Ex legis porro præscrípto vivens, profécto non aliúnde quam ex sua tribu ortam uxórem despondébat. Quocírca, si Joseph ex tribu Juda atque ex Davídis sorte ac família erat, an non consentáneum est Maríam quoque ab iísdem proficísci? Ex quo factum est, ut Sponsi genus recenserétur. Nam, cum de Apóstoli senténtia caput mulíeris vir sit; quid tandem afférri potest, quin, cum cápitis genus recensétur, corpus quoque una cum cápite recenséri consequátur? Maniféste ígitur osténsum esse árbitror, Joséphi genus apud Evangelístas non frustra recenséri, ex quo necessário Virgo quoque a Davíde oriúnda esse per consecutiónem intellígitur, et qui præcellénti miráculo ex ipsa génitus est Christus, ante sæcula Dei Fílius.\par
\switchcolumn
\EN
\noindent It is indeed clear to all that Joseph was endued with righteousness, and led a life in accordance with the Law. Therefore, living by what the Law prescribed, he certainly would not marry a wife sprung from any other but his own tribe. If, then, Joseph belonged to the tribe of Judah, and came of the seed and family of David, is it not a matter of course that Mary should come from the same? Whence it is that her husband's descent is recorded. For if, according to the Apostle's saying, the head of the woman is the man, doth it not follow in consequence that when the descent of the head is registered, that of the body in included in that of the head? I think it is therefore clearly shewn that the Evangelists purposely chronicle Joseph's genealogy, so that, in consequence, it would be understood that the Virgin was also sprung of the family of David; thereby implying the surpassing wonder that it was the Christ, before all ages the Son of God, who was born of her.\par
\switchcolumn*

\LA
\TeDeum{Te Deum}
\switchcolumn
\EN
\TeDeum{Te Deum}
\switchcolumn*

\end{paracol}

\Office{In Octava Patrocinii S. Joseph}{Octave Day of St. Joseph}{Duplex majus}
\addcontentsline{toc}{subsection}{In Octava Patrocinii S. Joseph\enspace\textit{Octave Day of St. Joseph}}
\begin{paracol}{2}
\LA
\LessonHead{Lectio I}
\switchcolumn
\EN
\LessonHead{Lesson I}
\switchcolumn*

\LA
\LessonTitle{De libro Apocalýpsis beáti Joánnis Apóstoli}
\switchcolumn
\EN
\LessonTitle{Lesson from the book of Revelation}
\switchcolumn*

\LA
\Cite{Apo 5:1-7}
\switchcolumn
\EN
\Cite{Rev 5:1-7}
\switchcolumn*

\LA
\noindent Et vidi in déxtera sedéntis supra thronum, librum scriptum intus et foris, signátum sigíllis septem. Et vidi ángelum fortem, prædicántem voce magna: Quis est dignus aperíre librum, et sólvere signácula ejus? Et nemo póterat neque in cælo, neque in terra, neque subtus terram aperíre librum, neque respícere illum. Et ego flebam multum, quóniam nemo dignus invéntus est aperíre librum, nec vidére eum. Et unus de senióribus dixit mihi: Ne fléveris: ecce vicit leo de tribu Juda, radix David, aperíre librum, et sólvere septem signácula ejus. Et vidi: et ecce in médio throni et quátuor animálium, et in médio seniórum, Agnum stantem tamquam occísum, habéntem córnua septem, et óculos septem: qui sunt septem spíritus Dei, missi in omnem terram. Et venit: et accépit de déxtera sedéntis in throno librum.\par
\switchcolumn
\EN
\noindent And I saw in the right hand of him that sat on the throne, a book written within and without, sealed with seven seals. And I saw a strong angel, proclaiming with a loud voice: Who is worthy to open the book, and to loose the seals thereof? And no man was able, neither in heaven, nor on earth, nor under the earth, to open the book, nor to look on it. And I wept much, because no man was found worthy to open the book, nor to see it. And one of the ancients said to me: Weep not; behold the lion of the tribe of Juda, the root of David, hath prevailed to open the book, and to loose the seven seals thereof. And I saw: and behold in the midst of the throne and of the four living creatures, and in the midst of the ancients, a Lamb standing as it were slain, having seven horns and seven eyes: which are the seven Spirits of God, sent forth into all the earth. And he came and took the book out of the right hand of him that sat on the throne.\par
\switchcolumn*

\LA
\begin{Resp}\Rsign Platéæ tuæ, Jerúsalem, sternéntur auro mundo, allelúja: et cantábitur in te cánticum lætítiæ, allelúja: \Star{} Et per omnes vicos tuos ab univérsis dicétur, allelúja, allelúja.\end{Resp}
\begin{Resp}\Vsign Luce spléndida fulgébis, et omnes fines terræ adorábunt te.\end{Resp}
\begin{Resp}\Rsign Et per omnes vicos tuos ab univérsis dicétur, allelúja, allelúja.\end{Resp}
\switchcolumn
\EN
\begin{Resp}\Rsign Thy streets, O Jerusalem, shall bel paved with pure gold, Alleluia, and the song of joy shall be sung in thee. Alleluia. \Star{} And all that pass through all thy streets shall say Alleluia, Alleluia.\end{Resp}
\begin{Resp}\Vsign Thy light shall be exceeding glorious, and all the ends of the earth shall worship thee.\end{Resp}
\begin{Resp}\Rsign And all that pass through all thy streets shall say Alleluia, Alleluia.\end{Resp}
\switchcolumn*

\LA
\LessonHead{Lectio II}
\switchcolumn
\EN
\LessonHead{Lesson II}
\switchcolumn*

\LA
\Cite{Apo 5:8-10}
\switchcolumn
\EN
\Cite{Rev 5:8-10}
\switchcolumn*

\LA
\noindent Et cum aperuísset librum, quátuor animália, et vigínti quátuor senióres cecidérunt coram Agno, habéntes sínguli cítharas, et phíalas áureas plenas odoramentórum, quæ sunt oratiónes sanctórum: et cantábant cánticum novum, dicéntes: Dignus es, Dómine, accípere librum, et aperíre signácula ejus: quóniam occísus es, et redemísti nos Deo in sánguine tuo ex omni tribu, et lingua, et pópulo, et natióne: et fecísti nos Deo nostro regnum, et sacerdótes: et regnábimus super terram.\par
\switchcolumn
\EN
\noindent And when he had opened the book, the four living creatures, and the four and twenty ancients fell down before the Lamb, having every one of them harps, and golden vials full of odours, which are the prayers of saints: And they sung a new canticle, saying: Thou art worthy, O Lord, to take the book, and to open the seals thereof; because thou wast slain, and hast redeemed us to God, in thy blood, out of every tribe, and tongue, and people, and nation. And hast made us to our God a kingdom and priests, and we shall reign on the earth.\par
\switchcolumn*

\LA
\begin{Resp}\Rsign Decantábat pópulus Israël, allelúja: et univérsa multitúdo Jacob canébat legítime: \Star{} Et David cum cantóribus cítharam percutiébat in domo Dómini, et laudes Deo canébat, allelúja, allelúja.\end{Resp}
\begin{Resp}\Vsign Sanctificáti sunt ergo sacerdótes et levítæ: et univérsus Israël deducébat arcam fœ́deris Dómini in júbilo.\end{Resp}
\begin{Resp}\Rsign Et David cum cantóribus cítharam percutiébat in domo Dómini, et laudes Deo canébat, allelúja, allelúja.\end{Resp}
\switchcolumn
\EN
\begin{Resp}\Rsign The people of Israel sung: Alleluia, and all the multitude of Jacob sung in measure. \Star{} And David was with the singers, (and) played upon an harp in the house of the Lord, and sung praises unto God. Alleluia, Alleluia.\end{Resp}
\begin{Resp}\Vsign So the Priests and the Levites were sanctified, and all Israel brought up the ark of the covenant of the Lord with shouting.\end{Resp}
\begin{Resp}\Rsign And David was with the singers, (and) played upon an harp in the house of the Lord, and sung praises unto God. Alleluia, Alleluia.\end{Resp}
\switchcolumn*

\LA
\LessonHead{Lectio III}
\switchcolumn
\EN
\LessonHead{Lesson III}
\switchcolumn*

\LA
\Cite{Apo 5:11-14}
\switchcolumn
\EN
\Cite{Rev 5:11-14}
\switchcolumn*

\LA
\noindent Et vidi, et audívi vocem angelórum multórum in circúitu throni, et animálium, et seniórum: et erat númerus eórum míllia míllium, dicéntium voce magna: Dignus est Agnus, qui occísus est, accípere virtútem, et divinitátem, et sapiéntiam, et fortitúdinem, et honórem, et glóriam, et benedictiónem. Et omnem creatúram, quæ in cælo est, et super terram, et sub terra, et quæ sunt in mari, et quæ in eo: omnes audívi dicéntes: Sedénti in throno, et Agno, benedíctio et honor, et glória, et potéstas in sǽcula sæculórum. Et quátuor animália dicébant: Amen. Et vigínti quátuor senióres cecidérunt in fácies suas: et adoravérunt vivéntem in sǽcula sæculórum.\par
\switchcolumn
\EN
\noindent And I beheld, and I heard the voice of many angels round about the throne, and the living creatures, and the ancients; and the number of them was thousands of thousands, Saying with a loud voice: The Lamb that was slain is worthy to receive power, and divinity, and wisdom, and strength, and honour, and glory, and benediction. And every creature, which is in heaven, and on the earth, and under the earth, and such as are in the sea, and all that are in them: I heard all saying: To him that sitteth on the throne, and to the Lamb, benediction, and honour, and glory, and power, for ever and ever. And the four living creatures said: Amen. And the four and twenty ancients fell down on their faces, and adored him that liveth for ever and ever.\par
\switchcolumn*

\LA
\begin{Resp}\Rsign Vidi portam civitátis ad Oriéntem pósitam, et Apostolórum nómina et Agni super eam scripta: \Star{} Et super muros ejus Angelórum custódiam, allelúja.\end{Resp}
\begin{Resp}\Vsign Vidi cælum novum, et terram novam, et civitátem novam descendéntem de cælo.\end{Resp}
\begin{Resp}\Rsign Et super muros ejus Angelórum custódiam, allelúja.\end{Resp}
\begin{Resp}\Vsign Glória Patri, et Fílio, \Star{} et Spirítui Sancto.\end{Resp}
\begin{Resp}\Rsign Et super muros ejus Angelórum custódiam, allelúja.\end{Resp}
\switchcolumn
\EN
\begin{Resp}\Rsign I saw the gate of the city, which looketh toward the East, and written thereon the names of the (Twelve) Apostles and of the Lamb. \Star{} And upon the walls thereof a guard of Angels. alleluia.\end{Resp}
\begin{Resp}\Vsign I saw a new heaven and a new earth coming down out of heaven.\end{Resp}
\begin{Resp}\Rsign And upon the walls thereof a guard of Angels. Alleluia.\end{Resp}
\begin{Resp}\Vsign Glory be to the Father, and to the Son, \Star{} and to the Holy Ghost.\end{Resp}
\begin{Resp}\Rsign And upon the walls thereof a guard of Angels. Alleluia.\end{Resp}
\switchcolumn*

\LA
\LessonHead{Lectio IV}
\switchcolumn
\EN
\LessonHead{Lesson IV}
\switchcolumn*

\LA
\LessonTitle{Sermo sancti Augustíni Epíscopi}
\switchcolumn
\EN
\LessonTitle{From a Sermon by St. Augustine the Bishop}
\switchcolumn*

\LA
\Source{Liber 1 de Nupt. et Concup., c. 11}
\switchcolumn
\EN
\Source{Liber 1 de Nupt. et Concup., c. 11}
\switchcolumn*

\LA
\noindent Non falláciter ab Angelo dictum est ad Joseph: Noli timére accípere Maríam cónjugem tuam. Conjux vocátur ex prima fide desponsatiónis, quam concúbitu nec cognóverat nec fúerat cognitúrus; nec períerat nec mendax mánserat cónjugis appellátio, ubi nec fúerat nec futúra erat ulla carnis commíxtio. Erat quippe illa Virgo; ídeo et sánctius et mirabílius jucúnda suo viro, quia étiam fœcúnda sine viro, prole dispar, fide compar. Propter quod fidéle conjúgium paréntes Christi vocári ambo meruérunt, et non solum illa mater, verum étiam ille pater ejus, sicut conjux matris ejus, utrúmque mente, non carne. Sive autem ille pater sola mente, sive illa mater et carne, paréntes tamen ambo humilitátis ejus, non sublimitátis; infirmitátis, non divinitátis.\par
\switchcolumn
\EN
\noindent The Angel did not speak falsely when he said to Joseph: Fear not to take unto thee Mary thy wife. She is called wife because of the mutual confidence established between them at the time of her espousal, although he had not known her carnally, nor was he ever so to do. And the name of wife was not lost or rendered untrue because there had not been any carnal intercourse, and would not be in the future. She was, in fact, The Virgin; and therefore she was holier and a more wonderful source of joy to her husband just because she became a mother without a man's intervention. Thus he knew her to be like unto himself in faithfulness, unlike him as regards her offspring. On account of his faithful union, both of them merited the name of Christ's parents. And not only is she called his Mother, but he also is called his father, as being the husband of his Mother, not according to the flesh, but according to the spirit. But even though he was a father only in spirit, whilst she was Mother according to the flesh, yet they both were the parents of his humility, not of his glory; of his infirmity, not of his divinity.\par
\switchcolumn*

\LA
\begin{Resp}\Rsign Dedísti mihi protectiónem salútis tuæ, et déxtera tua suscépit me: \Star{} Protéctor meus, et cornu salútis meæ, et suscéptor meus, allelúja.\end{Resp}
\begin{Resp}\Vsign Ego protéctor tuus sum, et merces tua magna nimis.\end{Resp}
\begin{Resp}\Rsign Protéctor meus, et cornu salútis meæ, et suscéptor meus, allelúja.\end{Resp}
\switchcolumn
\EN
\begin{Resp}\Rsign Thou hast given me the shield of thy salvation, and thy right hand hath holden me up. \Star{} My buckler, and the horn of my salvation, and my refuge. Alleluia.\end{Resp}
\begin{Resp}\Vsign I am thy shield and thy exceeding great reward.\end{Resp}
\begin{Resp}\Rsign My buckler, and the horn of my salvation, and my refuge. Alleluia.\end{Resp}
\switchcolumn*

\LA
\LessonHead{Lectio V}
\switchcolumn
\EN
\LessonHead{Lesson V}
\switchcolumn*

\LA
\noindent Neque enim mentítur Evangélium, ubi légitur: Et erat Pater ejus et Mater mirántes super his, quæ dicebántur de illo. Et álio loco: Ibant paréntes ejus per omnes annos in Jerúsalem. Item paulo post: Et dixit Mater ejus ad illum: Fili, quid fecísti nobis sic? Ecce pater tuus et ego doléntes quærebámus te. At ille, ut osténderet habére se præter illos Patrem, qui eum génuit præter matrem, respóndit eis: Quid est, quod me quærebátis? Nesciebátis, quia in his, quæ Patris mei sunt, opórtet me esse? Et rursum, ne hoc dicto paréntes illos negásse putarétur, Evangelísta secútus adjúnxit: Et ipsi non intellexérunt verbum quod locútus est ad illos; et descéndit cum eis, et venit Názareth, et erat súbditus illis. Quibus súbditus nisi paréntibus? Quis autem súbditus, nisi Jesus Christus, qui, cum in forma Dei esset, non rapínam arbitrátus est, esse se æquálem Deo?\par
\switchcolumn
\EN
\noindent For the Gospel doth not lie, when it saith: And Joseph and his Mother marvelled at those things which were spoken of him. And in another place: His parents went to Jerusalem every year. And a little further on: And his Mother said unto him; Son, why hast thou thus dealt with us? behold thy father and I have sought thee sorrowing. But, to shew that, apart from them, he had a Father who begat him without a mother, he answered them: How is it that ye sought me? wist ye not that I must be about my Father's business? And, as a set-off to this, lest anyone might think that by these words he denied his parents, the Evangelist immediately addeth: And they understood not the saying which he spoke unto them; and he went down with them, and came to Nazareth, and was subject unto them. To whom was he subject but to his parents? And who was thus subject but Jesus Christ, who, being in the form of God, thought it not robbery to be equal to God?\par
\switchcolumn*

\LA
\begin{Resp}\Rsign Státuet fílios suos sub tégmine illíus, et sub ramis ejus morábitur: protegétur sub tégmine illíus a fervóre: \Star{} Et in glória ejus requiéscet, allelúja.\end{Resp}
\begin{Resp}\Vsign Speráte in eo, omnis congregátio pópuli, effúndite coram illo corda vestra.\end{Resp}
\begin{Resp}\Rsign Et in glória ejus requiéscet, allelúja.\end{Resp}
\switchcolumn
\EN
\begin{Resp}\Rsign He shall set his children under her shelter, and shall lodge under her branches by her shall he be covered from heat, \Star{} And in her glory shall he dwell. Alleluia.\end{Resp}
\begin{Resp}\Vsign Trust in Him, ye congregation of the people, pour out your heart before Him.\end{Resp}
\begin{Resp}\Rsign And in her glory shall he dwell. Alleluia.\end{Resp}
\switchcolumn*

\LA
\LessonHead{Lectio VI}
\switchcolumn
\EN
\LessonHead{Lesson VI}
\switchcolumn*

\LA
\noindent Cur ergo illis súbditus, qui longe infra formam Dei erant, nisi quia semetípsum exinanívit formam servi accípiens, cujus formæ paréntes erant? Sed profécto nec ipsíus formæ servi paréntes ambo essent, nisi inter se étiam sine carnis commixtióne cónjuges essent. Unde et séries generatiónum, cum paréntes Christi connexióne successiónis commemorántur, usque ad Joseph pótius, sicut factum est, fúerat perducénda; ne in illo conjúgio viríli séxui, útique potióri, fíeret injúria, cum veritáti nihil períret, quia ex sémine David, ex quo ventúrus prædíctus est Christus, et Joseph erat et María. Omne ítaque nuptiárum bonum implétum est in illis paréntibus Christi: proles, fides sacraméntum. Prolem cognóscimus ipsum Dóminum Jesum; fidem, quia nullum adultérium; sacraméntum, quia nullum divórtium.\par
\switchcolumn
\EN
\noindent Why therefore was he subject to them who were so far below the form of God, except that he humbled himself, taking upon himself the form of a servant, of which form they were the parents? But truly, neither of them would have attained unto the parenthood of this form of a servant, except they had become respectively husband and wife, albeit without any carnal intercourse. And hence, when the ancestors of Christ are recounted in direct line of succession, the genealogy was fittingly traced down to Joseph. Otherwise, it would have been a slur upon the male sex, which is wont to be accorded the greater dignity. At the same time the truth did not suffer, for both Joseph and Mary were of the seed of David, from which it was prophesied that Christ should come. Note how thus all the good things of marriage are found in these parents of Christ: offspring, fidelity, the marriage bond. The offspring we know, was the Lord Jesus himself; their fidelity is proved because there was no adultery; the marriage bond, because there was no divorce.\par
\switchcolumn*

\LA
\begin{Resp}\Rsign Si consístant advérsum me castra, non timébit cor meum: \Star{} Si exsúrgat advérsum me prǽlium, in hoc ego sperábo, allelúja.\end{Resp}
\begin{Resp}\Vsign In te cantátio mea semper, quóniam tu adjútor fortis.\end{Resp}
\begin{Resp}\Rsign Si exsúrgat advérsum me prǽlium, in hoc ego sperábo, allelúja.\end{Resp}
\begin{Resp}\Vsign Glória Patri, et Fílio, \Star{} et Spirítui Sancto.\end{Resp}
\begin{Resp}\Rsign Si exsúrgat advérsum me prǽlium, in hoc ego sperábo, allelúja.\end{Resp}
\switchcolumn
\EN
\begin{Resp}\Rsign Though an host should encamp against me, my heart shall not fear. \Star{} Though war should rise against me, in this will I be confident. Alleluia.\end{Resp}
\begin{Resp}\Vsign My praise shall be continually of thee, for Thou art my strong refuge. R.* Though war should rise against me, in this will I be confident. Alleluia.\end{Resp}
\begin{Resp}\Vsign Glory be to the Father, and to the Son, \Star{} and to the Holy Ghost.\end{Resp}
\begin{Resp}\Rsign Though war should rise against me, in this will I be confident. Alleluia.\end{Resp}
\switchcolumn*

\LA
\LessonHead{Lectio VII}
\switchcolumn
\EN
\LessonHead{Lesson VII}
\switchcolumn*

\LA
\LessonTitle{Léctio sancti Evangélii secúndum Lucam}
\switchcolumn
\EN
\LessonTitle{From the Holy Gospel according to Luke}
\switchcolumn*

\LA
\Cite{Luc 3:21-23}
\switchcolumn
\EN
\Cite{Luke 3:21-23}
\switchcolumn*

\LA
\noindent In illo témpore: Factum est autem cum baptizarétur omnis pópulus, et Jesu baptizáto et oránte, apértum est cælum. Et réliqua.\par
\switchcolumn
\EN
\noindent At that time: When all the people were baptized, it came to pass, that Jesus also being baptized and praying, the heaven was opened. And so on, and that which followeth.\par
\switchcolumn*

\LA
\LessonTitle{Homilía sancti Augustíni Epíscopi}
\switchcolumn
\EN
\LessonTitle{A Homily by St. Augustine the Bishop}
\switchcolumn*

\LA
\Source{Serm. 36 de Temp. Bapt. Christi}
\switchcolumn
\EN
\Source{Serm. 36 de Temp. Bapt. Christi}
\switchcolumn*

\LA
\noindent Natális hódie alter est quodam modo Salvatóris. Nam eísdem eum signis, eísdem miráculis cognóscimus génitum, sed nunc majóri mystério baptizátum. Ait enim Deus: Hic est Fílius meus diléctus, in quo mihi complácui. Præclárior plane est secúnda, quam prima natívitas. Illa enim sine teste siléntio Christum génuit; ista cum divinitátis professióne Dóminum baptizávit. Ab illa se Joseph, qui pater putabátur, excúsat; in hac se Pater, qui non credebátur, insínuat. Ibi labórat suspiciónibus Mater, quia professióni déerat pater; hic honorátur Génitrix, quia Divínitas Fílium protestátur.\par
\switchcolumn
\EN
\noindent The day of his baptism is, as it were, a second birthday of the Saviour. For we know that he was born with signs and wonders like to those of his baptism, and that in the latter is a great mystery like to his birth. For God saith: This is my beloved Son, in whom I am well pleased. This second birth is indeed more glorious than the first. For then, he was born in silence, and without witnesses. Now, the Lord is baptized with a proclamation of his divinity. Then, Joseph, who was supposed to be his father, denied that he was. Now, his true Father, who was not believed to be so, proclaimeth himself so to be. Then, the Mother was enduring suspicion, because no father was acknowledged. Now she that bore him is honoured because the Divinity maketh him known as his Son.\par
\switchcolumn*

\LA
\begin{Resp}\Rsign Joseph, fili David, noli timére accípere Maríam cónjugem tuam; quod enim in ea natum est, de Spíritu Sancto est: páriet autem fílium, \Star{} Et vocábis nomen ejus Jesum, allelúja.\end{Resp}
\begin{Resp}\Vsign Ipse enim salvum fáciet pópulum suum a peccátis eórum.\end{Resp}
\begin{Resp}\Rsign Et vocábis nomen ejus Jesum, allelúja.\end{Resp}
\switchcolumn
\EN
\begin{Resp}\Rsign Joseph, thou Son of David, fear not to take unto thee Mary thy wife; for That Which is conceived in her is of the Holy Ghost and she shall bring forth a Son; \Star{} And thou shalt call His Name Jesus. Alleluia.\end{Resp}
\begin{Resp}\Vsign For He shall save His people from their sins.\end{Resp}
\begin{Resp}\Rsign And thou shalt call His Name Jesus. Alleluia.\end{Resp}
\switchcolumn*

\LA
\LessonHead{Lectio VIII}
\switchcolumn
\EN
\LessonHead{Lesson VIII}
\switchcolumn*

\LA
\noindent Honorátior, inquam, secúnda, quam prima natívitas. Síquidem Pater hic Deus majestátis inscríbitur; illic Joseph ártifex æstimátur. Et licet in utráque Dóminus per Spíritum Sanctum et natus sit et baptizátus, tamen honorátior est qui de cælis clamat, quam qui in terris labórat. Joseph ergo faber in terris pater putabátur esse Dómini Salvatóris, nec ab hoc ópere Deus, qui vere est Pater Dómini nostri Jesu Christi, exclúditur; nam est et ipse faber.\par
\switchcolumn
\EN
\noindent I say that the second birth was more glorious than the first. For now, the God of majesty proclaimeth himself as his father. Then, the carpenter Joseph was so accounted. And although it was the Holy Ghost through whom the Lord was born and baptized, yet the Father, whose voice was heard from heaven, is greater than the father who laboured on earth. Therefore Joseph the workman on earth was supposed to be the father of the Lord and Saviour. But God, the true Father of our Lord Jesus Christ, is also a workman, and cannot be excluded from those who work at the carpenter's trade.\par
\switchcolumn*

\LA
\begin{Resp}\Rsign Surge, et áccipe Púerum et matrem ejus, et fuge in Ægýptum: \Star{} Et esto ibi usque dum dicam tibi, allelúja.\end{Resp}
\begin{Resp}\Vsign Ut adimplerétur quod dictum est a Dómino per Prophétam dicéntem: Ex Ægýpto vocávi Fílium meum.\end{Resp}
\begin{Resp}\Rsign Et esto ibi usque dum dicam tibi, allelúja.\end{Resp}
\begin{Resp}\Vsign Glória Patri, et Fílio, \Star{} et Spirítui Sancto.\end{Resp}
\begin{Resp}\Rsign Et esto ibi usque dum dicam tibi, allelúja.\end{Resp}
\switchcolumn
\EN
\begin{Resp}\Rsign Arise, and take the young Child, and His Mother, and flee into Egypt; \Star{} And be thou there until I bring thee word. Alleluia.\end{Resp}
\begin{Resp}\Vsign That it might be fulfilled which was spoken of the Lord by the Prophets, saying Out of Egypt have I called My Son.\end{Resp}
\begin{Resp}\Rsign And be thou there until I bring thee word. Alleluia.\end{Resp}
\begin{Resp}\Vsign Glory be to the Father, and to the Son, \Star{} and to the Holy Ghost.\end{Resp}
\begin{Resp}\Rsign And be thou there until I bring thee word. Alleluia.\end{Resp}
\switchcolumn*

\LA
\LessonHead{Lectio IX}
\switchcolumn
\EN
\LessonHead{Lesson IX}
\switchcolumn*

\LA
\noindent Ipse enim est ártifex, qui hujus mundi máchinam non solum mirábili, sed étiam ineffábili poténtia fabricávit; tamquam sápiens architéctus cælum sublimitáte suspéndit, terram mole fundávit, mária cálculis alligávit. Ipse est ártifex, qui ad mensúram quamdam supérbiæ depónit fastígia, humilitátis extréma sublímat. Ipse est ártifex, qui in nostris móribus præcídit supérflua ópera, utília quæque consérvat. Ipse est ártifex, cujus secúrim ad radícem nostram pósitam Joánnes Baptísta comminátur, ut omnis arbor, quæ normam justæ discretiónis excésserit, excísa radícitus tradátur incéndio; quæ autem mensúram veritátis habúerit, cælésti fábricæ deputétur.\par
\switchcolumn
\EN
\noindent For he is the artificer who hath wrought the fabric of this world with power not only wondrous but ineffable. Like a wise architect hath he erected the heavens on high; he hath laid the foundations of the earth; he hath constrained the sea within its beaches. He is the artificer who, in due measure, lowereth the pinnacles of pride and bringeth to the surface the bedrock of humility. He is the artificer who doth chip off the unnecessary substance in our behaviour, and preserveth whatever is useful. He is the artificer whose axe, as John the Baptist warneth us, is laid to the root of our tree. So every tree, which measureth not up to his standard of due growth, is cut down at the roots, and used as fuel for the fire. But that which measureth up rightly, according to his rule and standard, is squared and fitted by his divine workmanship.\par
\switchcolumn*

\LA
\TeDeum{Te Deum}
\switchcolumn
\EN
\TeDeum{Te Deum}
\switchcolumn*

\end{paracol}

\Office{Feria Quinta infra Hebdomadam III post Octavam Paschæ}{Thursday of the Third Week after the Octave of Easter}{Feria}
\addcontentsline{toc}{subsection}{Feria Quinta infra Hebdomadam III post Octavam Paschæ\enspace\textit{Thursday of the Third Week after the Octave of Easter}}
\begin{paracol}{2}
\LA
\LessonHead{Lectio I}
\switchcolumn
\EN
\LessonHead{Lesson I}
\switchcolumn*

\LA
\LessonTitle{De libro Apocalýpsis beáti Joánnis Apóstoli}
\switchcolumn
\EN
\LessonTitle{Lesson from the book of Revelation}
\switchcolumn*

\LA
\Cite{Apo 15:1-4}
\switchcolumn
\EN
\Cite{Rev 15:1-4}
\switchcolumn*

\LA
\noindent Et vidi áliud signum in cælo magnum et mirábile, ángelos septem, habéntes plagas septem novíssimas: quóniam in illis consummáta est ira Dei. Et vidi tamquam mare vítreum mistum igne, et eos, qui vicérunt béstiam, et imáginem ejus, et númerum nóminis ejus, stantes super mare vítreum, habéntes cítharas Dei: et cantántes cánticum Móysi servi Dei, et cánticum Agni, dicéntes: Magna et mirabília sunt ópera tua, Dómine Deus omnípotens: justæ et veræ sunt viæ tuæ, Rex sæculórum. Quis non timébit te, Dómine, et magnificábit nomen tuum? quia solus pius es: quóniam omnes gentes vénient, et adorábunt in conspéctu tuo, quóniam judícia tua manifésta sunt.\par
\switchcolumn
\EN
\noindent And I saw another sign in heaven, great and wonderful: seven angels having the seven last plagues. For in them is filled up the wrath of God. And I saw as it were a sea of glass mingled with fire, and them that had overcome the beast, and his image, and the number of his name, standing on the sea of glass, having the harps of God: And singing the canticle of Moses, the servant of God, and the canticle of the Lamb, saying: Great and wonderful are thy works, O Lord God Almighty; just and true are thy ways, O King of ages. Who shall not fear thee, O Lord, and magnify thy name? For thou only art holy: for all nations shall come, and shall adore in thy sight, because thy judgments are manifest.\par
\switchcolumn*

\LA
\begin{Resp}\Rsign Dignus es, Dómine, accípere librum, et aperíre signácula ejus, allelúja: quóniam occísus es, et redemísti nos Deo \Star{} In sánguine tuo, allelúja.\end{Resp}
\begin{Resp}\Vsign Fecísti enim nos Deo nostro regnum et sacerdótium.\end{Resp}
\begin{Resp}\Rsign In sánguine tuo, allelúja.\end{Resp}
\switchcolumn
\EN
\begin{Resp}\Rsign Thou art worthy, O Lord, to take the book, and to open the seals thereof; alleluia: because thou wast slain, and hast redeemed us to God, \Star{} In thy blood, alleluia.\end{Resp}
\begin{Resp}\Vsign And hast made us to our God a kingdom and priests.\end{Resp}
\begin{Resp}\Rsign In thy blood, alleluia.\end{Resp}
\switchcolumn*

\LA
\LessonHead{Lectio II}
\switchcolumn
\EN
\LessonHead{Lesson II}
\switchcolumn*

\LA
\Cite{Apo 15:5-8}
\switchcolumn
\EN
\Cite{Rev 15:5-8}
\switchcolumn*

\LA
\noindent Et post hæc vidi: et ecce apértum est templum tabernáculi testimónii in cælo, et exiérunt septem ángeli habéntes septem plagas de templo, vestíti lino mundo et cándido, et præcíncti circa péctora zonis áureis. Et unum de quátuor animálibus dedit septem ángelis septem phíalas áureas, plenas iracúndiæ Dei vivéntis in sǽcula sæculórum. Et implétum est templum fumo a majestáte Dei, et de virtúte ejus: et nemo póterat introíre in templum, donec consummaréntur septem plagæ septem angelórum.\par
\switchcolumn
\EN
\noindent And after these things I looked; and behold, the temple of the tabernacle of the testimony in heaven was opened: And the seven angels came out of the temple, having the seven plagues, clothed with clean and white linen, and girt about the breasts with golden girdles. And one of the four living creatures gave to the seven angels seven golden vials, full of the wrath of God, who liveth for ever and ever. And the temple was filled with smoke from the majesty of God, and from his power; and no man was able to enter into the temple, till the seven plagues of the seven angels were fulfilled.\par
\switchcolumn*

\LA
\begin{Resp}\Rsign Ego sicut vitis fructificávi suavitátem odóris, allelúja: \Star{} Transíte ad me, omnes qui concupíscitis me, et a generatiónibus meis adimplémini, allelúja, allelúja.\end{Resp}
\begin{Resp}\Vsign In me omnis grátia viæ et veritátis: in me omnes spes vitæ et virtútis.\end{Resp}
\begin{Resp}\Rsign Transíte ad me, omnes qui concupíscitis me, et a generatiónibus meis adimplémini, allelúja, allelúja.\end{Resp}
\begin{Resp}\Vsign Glória Patri, et Fílio, \Star{} et Spirítui Sancto.\end{Resp}
\begin{Resp}\Rsign Transíte ad me, omnes qui concupíscitis me, et a generatiónibus meis adimplémini, allelúja, allelúja.\end{Resp}
\switchcolumn
\EN
\begin{Resp}\Rsign As the vine brought I forth pleasant savour, Alleluia. \Star{} Come unto me, all ye that be desirous of me, and fill yourselves with my fruits. Alleluia, Alleluia.\end{Resp}
\begin{Resp}\Vsign In me is the favour of the way and the truth in me is the hope of life and strength.\end{Resp}
\begin{Resp}\Rsign Come unto me, all ye that be desirous of me, and fill yourselves with my fruits. Alleluia, Alleluia.\end{Resp}
\begin{Resp}\Vsign Glory be to the Father, and to the Son, \Star{} and to the Holy Ghost.\end{Resp}
\begin{Resp}\Rsign Come unto me, all ye that be desirous of me, and fill yourselves with my fruits. Alleluia, Alleluia.\end{Resp}
\switchcolumn*

\LA
\LessonHead{Lectio III}
\switchcolumn
\EN
\LessonHead{Lesson III}
\switchcolumn*

\LA
\Cite{Apo 16:1-6}
\switchcolumn
\EN
\Cite{Rev 16:1-6}
\switchcolumn*

\LA
\noindent Et audívi vocem magnam de templo, dicéntem septem ángelis: Ite, et effúndite septem phíalas iræ Dei in terram. Et ábiit primus, et effúdit phíalam suam in terram, et factum est vulnus sævum et péssimum in hómines, qui habébant caractérem béstiæ, et in eos qui adoravérunt imáginem ejus. Et secúndus ángelus effúdit phíalam suam in mare, et factus est sanguis tamquam mórtui: et omnis ánima vivens mórtua est in mari. Et tértius effúdit phíalam suam super flúmina, et super fontes aquárum, et factus est sanguis. Et audívi ángelum aquárum dicéntem: Justus es, Dómine, qui es, et qui eras sanctus, qui hæc judicásti: quia sánguinem sanctórum et prophetárum effudérunt, et sánguinem eis dedísti bíbere: digni enim sunt.\par
\switchcolumn
\EN
\noindent And I heard a great voice out of the temple, saying to the seven angels: Go, and pour out the seven vials of the wrath of God upon the earth. And the first went, and poured out his vial upon the earth, and there fell a sore and grievous wound upon men, who had the character of the beast; and upon them that adored the image thereof. And the second angel poured out his vial upon the sea, and there came blood as it were of a dead man; and every living soul died in the sea. And the third poured out his vial upon the rivers and the fountains of waters; and there was made blood. And I heard the angel of the waters saying: Thou art just, O Lord, who art, and who wast, the Holy One, because thou hast judged these things: For they have shed the blood of saints and prophets, and thou hast given them blood to drink; for they are worthy.\par
\switchcolumn*

\LA
\TeDeum{Te Deum}
\switchcolumn
\EN
\TeDeum{Te Deum}
\switchcolumn*

\end{paracol}

\Office{Feria Sexta infra Hebdomadam III post Octavam Paschæ}{Friday of the Third Week after the Octave of Easter}{Feria}
\addcontentsline{toc}{subsection}{Feria Sexta infra Hebdomadam III post Octavam Paschæ\enspace\textit{Friday of the Third Week after the Octave of Easter}}
\begin{paracol}{2}
\LA
\LessonHead{Lectio I}
\switchcolumn
\EN
\LessonHead{Lesson I}
\switchcolumn*

\LA
\LessonTitle{De libro Apocalýpsis beáti Joánnis Apóstoli}
\switchcolumn
\EN
\LessonTitle{Lesson from the book of Revelation}
\switchcolumn*

\LA
\Cite{Apo 19:1-5}
\switchcolumn
\EN
\Cite{Rev 19:1-5}
\switchcolumn*

\LA
\noindent Post hæc audívi quasi vocem turbárum multárum in cælo dicéntium: Allelúja: salus, et glória, et virtus Deo nostro est: quia vera et justa judícia sunt ejus, qui judicávit de meretríce magna, quæ corrúpit terram in prostitutióne sua, et vindicávit sánguinem servórum suórum de mánibus ejus. Et íterum dixérunt: Allelúja. Et fumus ejus ascéndit in sǽcula sæculórum. Et cecidérunt senióres vigínti quátuor, et quátuor animália, et adoravérunt Deum sedéntem super thronum, dicéntes: Amen: allelúja. Et vox de throno exívit, dicens: Laudem dícite Deo nostro omnes servi ejus: et qui timétis eum pusílli et magni.\par
\switchcolumn
\EN
\noindent After these things I heard as it were the voice of much people in heaven, saying: Alleluia. Salvation, and glory, and power is to our God. For true and just are his judgments, who hath judged the great harlot which corrupted the earth with her fornication, and hath revenged the blood of his servants, at her hands. And again they said: Alleluia. And her smoke ascendeth for ever and ever. And the four and twenty ancients, and the four living creatures fell down and adored God that sitteth upon the throne, saying: Amen; Alleluia. And a voice came out from the throne, saying: Give praise to our God, all ye his servants; and you that fear him, little and great.\par
\switchcolumn*

\LA
\begin{Resp}\Rsign Locútus est ad me unus ex septem Angelis, dicens: Veni, osténdam tibi novam nuptam, sponsam Agni: \Star{} Et vidi Jerúsalem descendéntem de cælo, ornátam monílibus suis, allelúja, allelúja, allelúja.\end{Resp}
\begin{Resp}\Vsign Et sústulit me in spíritu in montem magnum et altum.\end{Resp}
\begin{Resp}\Rsign Et vidi Jerúsalem descendéntem de cælo, ornátam monílibus suis, allelúja, allelúja, allelúja.\end{Resp}
\switchcolumn
\EN
\begin{Resp}\Rsign One of the seven Angels talked with me, saying: Come hither, I will show thee the bride, the Lamb's wife. \Star{} And I saw Jerusalem descending out of heaven, adorned with her jewels. Alleluia, Alleluia, Alleluia.\end{Resp}
\begin{Resp}\Vsign And he carried me away in the Spirit to a great and high mountain\end{Resp}
\begin{Resp}\Rsign And I saw Jerusalem descending out of heaven, adorned with her jewels. Alleluia, Alleluia, Alleluia.\end{Resp}
\switchcolumn*

\LA
\LessonHead{Lectio II}
\switchcolumn
\EN
\LessonHead{Lesson II}
\switchcolumn*

\LA
\Cite{Apo 19:6-10}
\switchcolumn
\EN
\Cite{Rev 19:6-10}
\switchcolumn*

\LA
\noindent Et audívi quasi vocem turbæ magnæ, et sicut vocem aquárum multárum, et sicut vocem tonitruórum magnórum, dicéntium: Allelúja: quóniam regnávit Dóminus Deus noster omnípotens. Gaudeámus, et exsultémus: et demus glóriam ei: quia venérunt núptiæ Agni, et uxor ejus præparávit se. Et datum est illi ut coopériat se býssino splendénti et cándido. Býssinum enim justificatiónes sunt sanctórum. Et dixit mihi: Scribe: Beáti qui ad cœnam nuptiárum Agni vocáti sunt; et dixit mihi: Hæc verba Dei vera sunt. Et cécidi ante pedes ejus, ut adorárem eum. Et dicit mihi: Vide ne féceris: consérvus tuus sum, et fratrum tuórum habéntium testimónium Jesu. Deum adóra. Testimónium enim Jesu est spíritus prophetíæ.\par
\switchcolumn
\EN
\noindent And I heard as it were the voice of a great multitude, and as the voice of many waters, and as the voice of great thunders, saying, Alleluia: for the Lord our God the Almighty hath reigned. Let us be glad and rejoice, and give glory to him; for the marriage of the Lamb is come, and his wife hath prepared herself. And it is granted to her that she should clothe herself with fine linen, glittering and white. For the fine linen are the justifications of saints. And he said to me: Write: Blessed are they that are called to the marriage supper of the Lamb. And he saith to me: These words of God are true. And I fell down before his feet, to adore him. And he saith to me: See thou do it not: I am thy fellow servant, and of thy brethren, who have the testimony of Jesus. Adore God. For the testimony of Jesus is the spirit of prophecy.\par
\switchcolumn*

\LA
\begin{Resp}\Rsign Audívi vocem in cælo Angelórum multórum dicéntium: \Star{} Timéte Dóminum, et date claritátem illi, et adoráte eum, qui fecit cælum et terram, mare et fontes aquárum, allelúja, allelúja.\end{Resp}
\begin{Resp}\Vsign Vidi Angelum Dei fortem, volántem per médium cæli, voce magna clamántem et dicéntem.\end{Resp}
\begin{Resp}\Rsign Timéte Dóminum, et date claritátem illi, et adoráte eum, qui fecit cælum et terram, mare et fontes aquárum, allelúja, allelúja.\end{Resp}
\begin{Resp}\Vsign Glória Patri, et Fílio, \Star{} et Spirítui Sancto.\end{Resp}
\begin{Resp}\Rsign Timéte Dóminum, et date claritátem illi, et adoráte eum, qui fecit cælum et terram, mare et fontes aquárum, allelúja, allelúja.\end{Resp}
\switchcolumn
\EN
\begin{Resp}\Rsign I heard in heaven the voice of many Angels, saying \Star{} Fear the Lord, and give glory to Him, and worship Him That made heaven and earth, the sea, and the fountains of waters. Alleluia, Alleluia.\end{Resp}
\begin{Resp}\Vsign I saw a strong Angel of God fly through the midst of heaven, crying with a loud voice and saying\end{Resp}
\begin{Resp}\Rsign Fear the Lord, and give glory to Him, and worship Him That made heaven and earth, the sea, and the fountains of waters. Alleluia, Alleluia.\end{Resp}
\begin{Resp}\Vsign Glory be to the Father, and to the Son, \Star{} and to the Holy Ghost.\end{Resp}
\begin{Resp}\Rsign Fear the Lord, and give glory to Him, and worship Him That made heaven and earth, the sea, and the fountains of waters. Alleluia, Alleluia.\end{Resp}
\switchcolumn*

\LA
\LessonHead{Lectio III}
\switchcolumn
\EN
\LessonHead{Lesson III}
\switchcolumn*

\LA
\Cite{Apo 19:11-16}
\switchcolumn
\EN
\Cite{Rev 19:11-16}
\switchcolumn*

\LA
\noindent Et vidi cælum apértum, et ecce equus albus, et qui sedébat super eum, vocabátur Fidélis, et Verax, et cum justítia júdicat et pugnat. Oculi autem ejus sicut flamma ignis, et in cápite ejus diadémata multa, habens nomen scriptum, quod nemo novit nisi ipse. Et vestítus erat veste aspérsa sánguine: et vocátur nomen ejus: Verbum Dei. Et exércitus qui sunt in cælo, sequebántur eum in equis albis, vestíti býssino albo et mundo. Et de ore ejus procédit gládius ex utráque parte acútus, ut in ipso percútiat gentes. Et ipse reget eas in virga férrea: et ipse calcat tórcular vini furóris iræ Dei omnipoténtis. Et habet in vestiménto et in fémore suo scriptum: Rex regum et Dóminus dominántium.\par
\switchcolumn
\EN
\noindent And I saw heaven opened, and behold a white horse; and he that sat upon him was called faithful and true, and with justice doth he judge and fight. And his eyes were as a flame of fire, and on his head were many diadems, and he had a name written, which no man knoweth but himself. And he was clothed with a garment sprinkled with blood; and his name is called, the Word of God. And the armies that are in heaven followed him on white horses, clothed in fine linen, white and clean. And out of his mouth proceedeth a sharp two edged sword; that with it he may strike the nations. And he shall rule them with a rod of iron; and he treadeth the winepress of the fierceness of the wrath of God the Almighty. And he hath on his garment, and on his thigh written: King of Kings, and Lord of Lords.\par
\switchcolumn*

\LA
\TeDeum{Te Deum}
\switchcolumn
\EN
\TeDeum{Te Deum}
\switchcolumn*

\end{paracol}

\Office{Sabbato infra Hebdomadam III post Octavam Paschæ}{Saturday of the Third Week after the Octave of Easter}{Feria}
\addcontentsline{toc}{subsection}{Sabbato infra Hebdomadam III post Octavam Paschæ\enspace\textit{Saturday of the Third Week after the Octave of Easter}}
\begin{paracol}{2}
\LA
\LessonHead{Lectio I}
\switchcolumn
\EN
\LessonHead{Lesson I}
\switchcolumn*

\LA
\LessonTitle{De libro Apocalýpsis beáti Joánnis Apóstoli}
\switchcolumn
\EN
\LessonTitle{Lesson from the book of Revelation}
\switchcolumn*

\LA
\Cite{Apo 22:1-7}
\switchcolumn
\EN
\Cite{Rev 22:1-7}
\switchcolumn*

\LA
\noindent Et osténdit mihi flúvium aquæ vitæ, spléndidum tamquam crystállum, procedéntem de sede Dei et Agni. In médio platéæ ejus, et ex utráque parte flúminis, lignum vitæ, áfferens fructus duódecim per menses síngulos, reddens fructum suum et fólia ligni ad sanitátem géntium. Et omne maledíctum non erit ámplius: sed sedes Dei et Agni in illa erunt, et servi ejus sérvient illi. Et vidébunt fáciem ejus: et nomen ejus in fróntibus eórum. Et nox ultra non erit: et non egébunt lúmine lucérnæ, neque lúmine solis, quóniam Dóminus Deus illuminábit illos, et regnábunt in sǽcula sæculórum. Et dixit mihi: Hæc verba fidelíssima sunt, et vera. Et Dóminus Deus spirítuum prophetárum misit ángelum suum osténdere servis suis quæ opórtet fíeri cito. Et ecce vénio velóciter. Beátus, qui custódit verba prophetíæ libri hujus.\par
\switchcolumn
\EN
\noindent And he showed me a river of water of life, clear as crystal, proceeding from the throne of God and of the Lamb. In the midst of the street thereof, and on both sides of the river, was the tree of life, bearing twelve fruits, yielding its fruits every month, and the leaves of the tree were for the healing of the nations. And there shall be no curse any more; but the throne of God and of the Lamb shall be in it, and his servants shall serve him. And they shall see his face: and his name shall be on their foreheads. And night shall be no more: and they shall not need the light of the lamp, nor the light of the sun, because the Lord God shall enlighten them, and they shall reign for ever and ever. And he said to me: These words are most faithful and true. And the Lord God of the spirits of the prophets sent his angel to show his servants the things which must be done shortly. And, Behold I come quickly. Blessed is he that keepeth the words of the prophecy of this book.\par
\switchcolumn*

\LA
\begin{Resp}\Rsign Decantábat pópulus Israël, allelúja: et univérsa multitúdo Jacob canébat legítime: \Star{} Et David cum cantóribus cítharam percutiébat in domo Dómini, et laudes Deo canébat, allelúja, allelúja.\end{Resp}
\begin{Resp}\Vsign Sanctificáti sunt ergo sacerdótes et levítæ: et univérsus Israël deducébat arcam fœ́deris Dómini in júbilo.\end{Resp}
\begin{Resp}\Rsign Et David cum cantóribus cítharam percutiébat in domo Dómini, et laudes Deo canébat, allelúja, allelúja.\end{Resp}
\switchcolumn
\EN
\begin{Resp}\Rsign The people of Israel sung: Alleluia, and all the multitude of Jacob sung in measure. \Star{} And David was with the singers, (and) played upon an harp in the house of the Lord, and sung praises unto God. Alleluia, Alleluia.\end{Resp}
\begin{Resp}\Vsign So the Priests and the Levites were sanctified, and all Israel brought up the ark of the covenant of the Lord with shouting.\end{Resp}
\begin{Resp}\Rsign And David was with the singers, (and) played upon an harp in the house of the Lord, and sung praises unto God. Alleluia, Alleluia.\end{Resp}
\switchcolumn*

\LA
\LessonHead{Lectio II}
\switchcolumn
\EN
\LessonHead{Lesson II}
\switchcolumn*

\LA
\Cite{Apo 22:8-12}
\switchcolumn
\EN
\Cite{Rev 22:8-12}
\switchcolumn*

\LA
\noindent Et ego Joánnes, qui audívi, et vidi hæc. Et postquam audíssem, et vidíssem, cécidi ut adorárem ante pedes ángeli, qui mihi hæc ostendébat: et dixit mihi: Vide ne féceris: consérvus enim tuus sum, et fratrum tuórum prophetárum, et eórum qui servant verba prophetíæ libri hujus: Deum adóra. Et dicit mihi: Ne signáveris verba prophetíæ libri hujus: tempus enim prope est. Qui nocet, nóceat adhuc: et qui in sórdibus est, sordéscat adhuc: et qui justus est, justificétur adhuc: et sanctus, sanctificétur adhuc. Ecce vénio cito, et merces mea mecum est, réddere unicuíque secúndum ópera sua.\par
\switchcolumn
\EN
\noindent And I, John, who have heard and seen these things. And after I had heard and seen, I fell down to adore before the feet of the angel, who showed me these things. And he said to me: See thou do it not: for I am thy fellow servant, and of thy brethren the prophets, and of them that keep the words of the prophecy of this book. Adore God. And he saith to me: Seal not the words of the prophecy of this book: for the time is at hand. He that hurteth, let him hurt still: and he that is filthy, let him be filthy still: and he that is just, let him be justified still: and he that is holy, let him be sanctified still. Behold, I come quickly; and my reward is with me, to render to every man according to his works.\par
\switchcolumn*

\LA
\begin{Resp}\Rsign Osténdit mihi Angelus fontem aquæ vivæ, et dixit ad me, allelúja: \Star{} Hic Deum adóra, allelúja, allelúja, allelúja.\end{Resp}
\begin{Resp}\Vsign Postquam audíssem et vidíssem, cécidi ut adorárem ante pedes Angeli, qui mihi hæc ostendébat, et dixit mihi.\end{Resp}
\begin{Resp}\Rsign Hic Deum adóra, allelúja, allelúja, allelúja.\end{Resp}
\switchcolumn
\EN
\begin{Resp}\Rsign The Angel showed me the fountain of the water of life and he said unto me, Alleluia. \Star{} Here worship God. Alleluia, Alleluia, Alleluia.\end{Resp}
\begin{Resp}\Vsign When I had heard and seen, I fell down to worship before the feet of the Angel, which showed me these things, and he said unto me\end{Resp}
\begin{Resp}\Rsign Here worship God. Alleluia, Alleluia, Alleluia.\end{Resp}
\switchcolumn*

\LA
\LessonHead{Lectio III}
\switchcolumn
\EN
\LessonHead{Lesson III}
\switchcolumn*

\LA
\Cite{Apo 22:13-21}
\switchcolumn
\EN
\Cite{Rev 22:13-21}
\switchcolumn*

\LA
\noindent Ego sum alpha et ómega, primus et novíssimus, princípium et finis. Beáti, qui lavant stolas suas in sánguine Agni: ut sit potéstas eórum in ligno vitæ, et per portas intrent in civitátem. Foris canes, et venéfici, et impudíci, et homicídæ, et idólis serviéntes, et omnis qui amat et facit mendácium. Ego Jesus misi ángelum meum testificári vobis hæc in ecclésiis. Ego sum radix, et genus David, stella spléndida et matutína. Et spíritus, et sponsa dicunt: Veni. Et qui audit, dicat: Veni. Et qui sitit, véniat: et qui vult, accípiat aquam vitæ, gratis. Contéstor enim omni audiénti verba prophetíæ libri hujus: si quis apposúerit ad hæc, appónet Deus super illum plagas scriptas in libro isto. Et si quis diminúerit de verbis libri prophetíæ hujus, áuferet Deus partem ejus de libro vitæ, et de civitáte sancta, et de his quæ scripta sunt in libro isto: dicit qui testimónium pérhibet istórum. Etiam vénio cito: amen. Veni, Dómine Jesu. Grátia Dómini nostri Jesu Christi cum ómnibus vobis. Amen.\par
\switchcolumn
\EN
\noindent I am Alpha and Omega, the first and the last, the beginning and the end. Blessed are they that wash their robes in the blood of the Lamb: that they may have a right to the tree of life, and may enter in by the gates into the city. Without are dogs, and sorcerers, and unchaste, and murderers, and servers of idols, and every one that loveth and maketh a lie. I Jesus have sent my angel, to testify to you these things in the churches. I am the root and stock of David, the bright and morning star. And the spirit and the bride say: Come. And he that heareth, let him say: Come. And he that thirsteth, let him come: and he that will, let him take the water of life, freely. For I testify to every one that heareth the words of the prophecy of this book: If any man shall add to these things, God shall add unto him the plagues written in this book. And if any man shall take away from the words of the book of this prophecy, God shall take away his part out of the book of life, and out of the holy city, and from these things that are written in this book. He that giveth testimony of these things, saith, Surely I come quickly: Amen. Come, Lord Jesus. The grace of our Lord Jesus Christ be with you all. Amen.\par
\switchcolumn*

\LA
\begin{Resp}\Rsign Vidi Jerúsalem descendéntem de cælo, ornátam auro mundo, et lapídibus pretiósis intéxtam: \Star{} Allelúja, allelúja.\end{Resp}
\begin{Resp}\Vsign Et erat structúra muri ejus ex lápide jáspide; ipsa vero aurum mundum, símile vitro mundo.\end{Resp}
\begin{Resp}\Rsign Allelúja, allelúja.\end{Resp}
\begin{Resp}\Vsign Glória Patri, et Fílio, \Star{} et Spirítui Sancto.\end{Resp}
\begin{Resp}\Rsign Allelúja, allelúja.\end{Resp}
\switchcolumn
\EN
\begin{Resp}\Rsign I saw Jerusalem coming down out of heaven, adorned with pure gold, and garnished with precious stones. \Star{} Alleluia, Alleluia.\end{Resp}
\begin{Resp}\Vsign The building of the wall of it was of jasper and the city was pure gold, like unto clear glass.\end{Resp}
\begin{Resp}\Rsign Alleluia, Alleluia.\end{Resp}
\begin{Resp}\Vsign Glory be to the Father, and to the Son, \Star{} and to the Holy Ghost.\end{Resp}
\begin{Resp}\Rsign Alleluia, Alleluia.\end{Resp}
\switchcolumn*

\end{paracol}

\Office{Dominica IV Post Pascha}{Fourth Sunday after Easter}{Semiduplex}
\addcontentsline{toc}{subsection}{Dominica IV Post Pascha\enspace\textit{Fourth Sunday after Easter}}
\begin{paracol}{2}
\LA
\LessonHead{Lectio I}
\switchcolumn
\EN
\LessonHead{Lesson I}
\switchcolumn*

\LA
\LessonTitle{Incipit Epístola cathólica beáti Jacóbi Apóstoli}
\switchcolumn
\EN
\LessonTitle{Lesson from the letter of St. James the Apostle}
\switchcolumn*

\LA
\Cite{Jas 1:1-6}
\switchcolumn
\EN
\Cite{Jas 1:1-6}
\switchcolumn*

\LA
\noindent Jacóbus, Dei et Dómini nostri Jesu Christi servus, duódecim tríbubus, quæ sunt in dispersióne, salútem. Omne gáudium existimáte fratres mei, cum in tentatiónes várias incidéritis: sciéntes quod probátio fídei vestræ patiéntiam operátur. Patiéntia autem opus perféctum habet: ut sitis perfécti et íntegri in nullo deficiéntes. Si quis autem vestrum índiget sapiéntia, póstulet a Deo, qui dat ómnibus affluénter, et non impróperat: et dábitur ei. Póstulet autem in fide nihil hǽsitans:\par
\switchcolumn
\EN
\noindent James the servant of God, and of our Lord Jesus Christ, to the twelve tribes which are scattered abroad, greeting. My brethren, count it all joy, when you shall fall into diverse temptations; Knowing that the trying of your faith worketh patience. And patience hath a perfect work; that you may be perfect and entire, failing in nothing. But if any of you want wisdom, let him ask of God, who giveth to all men abundantly, and upbraideth not; and it shall be given him. But let him ask in faith, nothing wavering.\par
\switchcolumn*

\LA
\begin{Resp}\Rsign Si oblítus fúero tui, allelúja, obliviscátur mei déxtera mea: \Star{} Adhǽreat lingua mea fáucibus meis, si non memínero tui, allelúja, allelúja.\end{Resp}
\begin{Resp}\Vsign Super flúmina Babylónis illic sédimus et flévimus, dum recordarémur tui, Sion.\end{Resp}
\begin{Resp}\Rsign Adhǽreat lingua mea fáucibus meis, si non memínero tui, allelúja, allelúja.\end{Resp}
\switchcolumn
\EN
\begin{Resp}\Rsign If I forget thee, Alleluia, let my right hand forget me. \Star{} If I do not remember thee, let my tongue cleave to the roof of my mouth, alleluia, alleluia.\end{Resp}
\begin{Resp}\Vsign By the rivers of Babylon there we sat down and wept, when we remembered thee, O Zion.\end{Resp}
\begin{Resp}\Rsign If I do not remember thee, let my tongue cleave to the roof of my mouth, alleluia, alleluia.\end{Resp}
\switchcolumn*

\LA
\LessonHead{Lectio II}
\switchcolumn
\EN
\LessonHead{Lesson II}
\switchcolumn*

\LA
\Cite{Jas 1:6-11}
\switchcolumn
\EN
\Cite{Jas 1:6-11}
\switchcolumn*

\LA
\noindent Qui enim hǽsitat, símilis est flúctui maris, qui a vento movétur et circumfértur: non ergo ǽstimet homo ille quod accípiat áliquid a Dómino. Vir duplex ánimo incónstans est in ómnibus viis suis. Gloriétur autem frater húmilis in exaltatióne sua: dives autem in humilitáte sua, quóniam sicut flos fœni transíbit; exórtus est enim sol cum ardóre, et arefécit fœnum, et flos ejus décidit, et decor vultus ejus depériit: ita et dives in itinéribus suis marcéscet.\par
\switchcolumn
\EN
\noindent For he that wavereth is like a wave of the sea, which is moved and carried about by the wind. Therefore let not that man think that he shall receive any thing of the Lord. A double minded man is inconstant in all his ways. But let the brother of low condition glory in his exaltation: And the rich, in his being low; because as the flower of the grass shall he pass away. For the sun rose with a burning heat, and parched the grass, and the flower thereof fell off, and the beauty of the shape thereof perished: so also shall the rich man fade away in his ways.\par
\switchcolumn*

\LA
\begin{Resp}\Rsign Vidérunt te aquæ, Deus, vidérunt te aquæ, et timuérunt: \Star{} Multitúdo sónitus aquárum vocem dedérunt nubes, allelúja, allelúja, allelúja.\end{Resp}
\begin{Resp}\Vsign Illuxérunt coruscatiónes tuæ orbi terræ: vidit et commóta est terra.\end{Resp}
\begin{Resp}\Rsign Multitúdo sónitus aquárum vocem dedérunt nubes, allelúja, allelúja, allelúja.\end{Resp}
\switchcolumn
\EN
\begin{Resp}\Rsign The waters saw thee, O God, the waters saw thee and they were afraid. \Star{} There was a noise as of many waters the clouds sent out a sound, alleluia, alleluia, alleluia.\end{Resp}
\begin{Resp}\Vsign thy lightnings lightened the world the earth saw it and shook.\end{Resp}
\begin{Resp}\Rsign There was a noise as of many waters the clouds sent out a sound, alleluia, alleluia, alleluia.\end{Resp}
\switchcolumn*

\LA
\LessonHead{Lectio III}
\switchcolumn
\EN
\LessonHead{Lesson III}
\switchcolumn*

\LA
\Cite{Jas 1:12-16}
\switchcolumn
\EN
\Cite{Jas 1:12-16}
\switchcolumn*

\LA
\noindent Beátus vir qui suffert tentatiónem: quóniam cum probátus fúerit, accípiet corónam vitæ, quam repromísit Deus diligéntibus se. Nemo cum tentátur, dicat quóniam a Deo tentátur: Deus enim intentátor malórum est: ipse autem néminem tentat. Unusquísque vero tentátur a concupiscéntia sua abstráctus, et illéctus. Deínde concupiscéntia cum concéperit, parit peccátum: peccátum vero cum consummátum fúerit, génerat mortem. Nolíte ítaque erráre, fratres mei dilectíssimi.\par
\switchcolumn
\EN
\noindent Blessed is the man that endureth temptation; for when he hath been proved, he shall receive a crown of life, which God hath promised to them that love him. Let no man, when he is tempted, say that he is tempted by God. For God is not a tempter of evils, and he tempteth no man. But every man is tempted by his own concupiscence, being drawn away and allured. Then when concupiscence hath conceived, it bringeth forth sin. But sin, when it is completed, begetteth death. Do not err, therefore, my dearest brethren.\par
\switchcolumn*

\LA
\begin{Resp}\Rsign Narrábo nomen tuum frátribus meis, allelúja: \Star{} In médio Ecclésiæ laudábo te, allelúja, allelúja.\end{Resp}
\begin{Resp}\Vsign Confitébor tibi in pópulis, Dómine, et psalmum dicam tibi in géntibus.\end{Resp}
\begin{Resp}\Rsign In médio Ecclésiæ laudábo te, allelúja, allelúja.\end{Resp}
\begin{Resp}\Vsign Glória Patri, et Fílio, \Star{} et Spirítui Sancto.\end{Resp}
\begin{Resp}\Rsign In médio Ecclésiæ laudábo te, allelúja, allelúja.\end{Resp}
\switchcolumn
\EN
\begin{Resp}\Rsign I will declare thy Name unto my brethren, alleluia. \Star{} In the midst of the congregation will I praise thee, alleluia, alleluia.\end{Resp}
\begin{Resp}\Vsign I will praise thee, O Lord, among the people, and sing unto thee among the nations.\end{Resp}
\begin{Resp}\Rsign In the midst of the congregation will I praise thee, alleluia, alleluia.\end{Resp}
\begin{Resp}\Vsign Glory be to the Father, and to the Son, \Star{} and to the Holy Ghost.\end{Resp}
\begin{Resp}\Rsign In the midst of the congregation will I praise thee, alleluia, alleluia.\end{Resp}
\switchcolumn*

\LA
\LessonHead{Lectio IV}
\switchcolumn
\EN
\LessonHead{Lesson IV}
\switchcolumn*

\LA
\LessonTitle{Ex Tractátu sancti Cypriáni Epíscopi et Mártyris de bono patiéntiæ}
\switchcolumn
\EN
\LessonTitle{From the exposition of St. Cyprian, Bishop and Martyr, on the good of patience.}
\switchcolumn*

\LA
\Source{Sermone 3 initio.}
\switchcolumn
\EN
\Source{Sermone 3. initio.}
\switchcolumn*

\LA
\noindent De patiéntia locutúrus, fratres dilectíssimi, et utilitátes ejus et cómmoda prædicatúrus, unde pótius incípiam, quam quod nunc quoque ad audiéntiam, vestram patiéntiam vídeo esse necessáriam: ut nec hoc ipsum, quod audítis et díscitis, sine patiéntia fácere possítis? Tunc enim demum sermo et rátio salutáris efficáciter díscitur, si patiénter, quod dícitur, audiátur. Nec invénio, fratres dilectíssimi, inter céteras cæléstis disciplínæ vias, quibus ac consequénda divínitus prǽmia spei ac fídei nostræ secta dirígitur, quid magis sit vel utílius ad vitam, vel majus ad glóriam, quam ut, qui præcéptis Domínicis obséquio timóris ac devotiónis innítimur, patiéntiam máxime tota observatióne tueámur. Hanc se sectári philósophi quoque profiténtur: sed tam illic patiéntia falsa est, quam et falsa sapiéntia est. Unde enim vel sápiens esse, vel pátiens possit, qui nec sapiéntiam nec patiéntiam Dei novit?\par
\switchcolumn
\EN
\noindent In speaking of patience, beloved brethren, and in preaching on its benefits and advantages, how can I better begin than by pointing out the fact that now, just for you to listen to me, I see that patience is necessary, as you could not even do this, namely, listen and learn, without patience. For only then is the word of God and way of salvation effectively learned, if one listens with patience to what is being said. Dearly beloved brethren, there are diverse paths of heavenly wisdom, wherein we are invited to walk, if we would reach in the end the reward which God hath prepared to crown hope and faith but I find no path more useful toward life, nor more sure toward glory than this, that while we humbly strive, in all fear, and in all godliness, to obey the commandments of the Lord, we should set our chiefest guard in an unceasing watch over our patience. The philosophers also say that they take this path, but their patience is as much a sham as their wisdom is a cheat, for who can be wise or patient who knoweth nothing of God's wisdom or God's patience are the lives of servers and worshippers of God. Let it be ours, then, to show forth by spiritual watchfulness that patience which is a part of the teaching which we have learnt from heaven. Patience is one of His Own virtues whereof God hath made us partakers with Him our Great Head is the Captain of the patient, and it is through patience that He hath crowned Himself with glory and honour.\par
\switchcolumn*

\LA
\begin{Resp}\Rsign In ecclésiis benedícite Deo, allelúja: \Star{} Dómino de fóntibus Israël, allelúja, allelúja.\end{Resp}
\begin{Resp}\Vsign Psalmum dícite nómini ejus, date glóriam laudi ejus.\end{Resp}
\begin{Resp}\Rsign Dómino de fóntibus Israël, allelúja, allelúja.\end{Resp}
\switchcolumn
\EN
\begin{Resp}\Rsign Bless ye God in the congregations, alleluia. \Star{} Even the Lord, ye that are of the fountains of Israel, alleluia, alleluia.\end{Resp}
\begin{Resp}\Vsign Sing forth the honour of His Name, make His praise glorious.\end{Resp}
\begin{Resp}\Rsign Even the Lord, ye that are of the fountains of Israel, alleluia, alleluia.\end{Resp}
\switchcolumn*

\LA
\LessonHead{Lectio V}
\switchcolumn
\EN
\LessonHead{Lesson V}
\switchcolumn*

\LA
\noindent Nos autem, fratres dilectíssimi, qui philósophi non verbis, sed factis sumus; nec vestítu sapiéntiam, sed veritáte præférimus: qui virtútum consciéntiam magis quam jactántiam nóvimus: qui non lóquimur magna, sed vívimus quasi servi et cultóres Dei: patiéntiam, quam magistériis cæléstibus díscimus, obséquiis spiritálibus præbeámus. Est enim nobis cum Deo virtus ista commúnis: inde patiéntia íncipit, inde cláritas ejus et dígnitas caput sumit. Orígo et magnitúdo patiéntiæ Deo auctóre procédit. Diligénda res hómini, quæ Deo cara est. Bonum quod amat, majéstas divína comméndat. Si Dóminus nobis et Pater Deus est, sectémur patiéntiam dómini páriter et patris; quia et servos opórtet esse obsequéntes, et fílios non decet esse degéneres.\par
\switchcolumn
\EN
\noindent But as for us, dearly beloved brethren, we are the real philosophers, whose wisdom lieth not in words but in deeds, and is manifested not in dresses but in the truth. We are they whose knowledge hath the inward consciousness, not the idle boasting, of strength. We are not speakers of high-sounding words, but our lives. Yea, God is Himself the Source, the Fountain, and the Greatness of patience, and it behoveth man to love what is beloved of God. That good thing which he loveth is commended unto him of God's Majesty. If God be our Lord and Father, let us follow after the example of our Lord and Father's patience, since it is the duty of servants to be obedient, and of sons to be home-minded.\par
\switchcolumn*

\LA
\begin{Resp}\Rsign In toto corde meo, allelúja, exquisívi te, allelúja: \Star{} Ne repéllas me a mandátis tuis, allelúja, allelúja.\end{Resp}
\begin{Resp}\Vsign Benedíctus es tu, Dómine, doce me justificatiónes tuas.\end{Resp}
\begin{Resp}\Rsign Ne repéllas me a mandátis tuis, allelúja, allelúja.\end{Resp}
\switchcolumn
\EN
\begin{Resp}\Rsign With my whole heart, alleluia, have I sought thee, alleluia. \Star{} O let me not wander from thy commandments, alleluia, alleluia.\end{Resp}
\begin{Resp}\Vsign Blessed art Thou, O Lord; teach me thy statutes.\end{Resp}
\begin{Resp}\Rsign O let me not wander from thy commandments, alleluia, alleluia.\end{Resp}
\switchcolumn*

\LA
\LessonHead{Lectio VI}
\switchcolumn
\EN
\LessonHead{Lesson VI}
\switchcolumn*

\LA
\noindent Patiéntia est, quæ nos Deo et comméndat et servat: ipsa est, quæ iram témperat, quæ linguam frenat, quæ mentem gubérnat, pacem custódit, disciplínam regit, libídinis ímpetum frangit, tumóris violéntiam cómprimit, incéndium simultátis exstínguit, coércet poténtiam dívitum, inópiam páuperum réfovet, tuétur in virgínibus beátam integritátem, in víduis laboriósam castitátem, in conjúnctis et maritátis indivíduam caritátem: facit húmiles in prósperis, in advérsis fortes, contra injúrias et contumélias mites: docet delinquéntibus cito ignóscere: si ipse delínquas, diu et multum rogáre: tentatiónes expúgnat, persecutiónes tólerat, passiónes et martýria consúmmat. Ipsa est, quæ fídei nostræ fundaménta fírmiter munit.\par
\switchcolumn
\EN
\noindent By our patience God draweth us toward Himself, and keepeth us His Own. Patience doth soothe anger, bridle the tongue, govern the mind, keep peace, set rules of self-control, break the onset of lust, still the swelling of temper, put out the fire begotten of hatred, make the rich meek, and relieve the need of the poor patience doth guard in virgins their blessed wholeness in widows, their careful purity in such as be married, their single-hearted love one toward the other. Patience doth teach such as be successful to be lowly-minded such as be unfortunate, to be brave and all to be gentle when they are wronged and insulted. Patience maketh a man soon to forgive them that trespass against him, and if he have trespassed against any, long and humbly to ask his pardon. Patience doth fight down temptations, bear persecution, and endure unto the end in suffering, and in uplifting of our testimony. Patience is the moat that guardeth the stout foundations of the castle of our faith.\par
\switchcolumn*

\LA
\begin{Resp}\Rsign Hymnum cantáte nobis, allelúja: \Star{} Quómodo cantábimus cánticum Dómini in terra aliéna? allelúja, allelúja.\end{Resp}
\begin{Resp}\Vsign Illic interrogavérunt nos, qui captívos duxérunt nos, verba cantiónum.\end{Resp}
\begin{Resp}\Rsign Quómodo cantábimus cánticum Dómini in terra aliéna? allelúja, allelúja.\end{Resp}
\begin{Resp}\Vsign Glória Patri, et Fílio, \Star{} et Spirítui Sancto.\end{Resp}
\begin{Resp}\Rsign Quómodo cantábimus cánticum Dómini in terra aliéna? allelúja, allelúja.\end{Resp}
\switchcolumn
\EN
\begin{Resp}\Rsign Sing us a song, alleluia. \Star{} How shall we sing the Lord's song in a strange land, alleluia, alleluia.\end{Resp}
\begin{Resp}\Vsign There they that carried us away captive required of us a song.\end{Resp}
\begin{Resp}\Rsign How shall we sing the Lord's song in a strange land, alleluia, alleluia.\end{Resp}
\begin{Resp}\Vsign Glory be to the Father, and to the Son, \Star{} and to the Holy Ghost.\end{Resp}
\begin{Resp}\Rsign How shall we sing the Lord's song in a strange land, alleluia, alleluia.\end{Resp}
\switchcolumn*

\LA
\LessonHead{Lectio VII}
\switchcolumn
\EN
\LessonHead{Lesson VII}
\switchcolumn*

\LA
\LessonTitle{Léctio sancti Evangélii secúndum Joánnem}
\switchcolumn
\EN
\LessonTitle{From the Holy Gospel according to John}
\switchcolumn*

\LA
\Cite{Joannes 16:5-14}
\switchcolumn
\EN
\Cite{John 16:5-14}
\switchcolumn*

\LA
\noindent In illo témpore: Dixit Jesus discípulis suis: Vado ad eum qui misit me; et nemo ex vobis intérrogat me, Quo vadis? Et réliqua.\par
\switchcolumn
\EN
\noindent At that time, Jesus said unto His disciples: I go My way to Him That sent Me, and none of you asketh Me: Whither goest Thou? And so on.\par
\switchcolumn*

\LA
\LessonTitle{Homilía sancti Augustíni Epíscopi}
\switchcolumn
\EN
\LessonTitle{Homily by St. Augustine, Bishop of Hippo}
\switchcolumn*

\LA
\Source{Tractatus 94 in Joannem, initio}
\switchcolumn
\EN
\Source{94th Tract on John}
\switchcolumn*

\LA
\noindent Cum Dóminus Jesus prædixísset discípulis suis persecutiónes, quas passúri erant post ejus abscéssum, subjúnxit atque ait: Hæc autem vobis ab inítio non dixi, quia vobíscum eram: nunc autem vado ad eum, qui me misit. Ubi primum vidéndum est, utrum eis futúras non prædíxerit ante passiónes. Sed álii tres Evangelístæ satis eum prædixísse ista demónstrant, ántequam ventum esset ad cenam: qua perácta, secúndum Joánnem ista locútus est, ubi ait: Hæc autem vobis ab inítio non dixi, quia vobíscum eram.\par
\switchcolumn
\EN
\noindent The Lord Jesus told His disciples what things they should suffer after that He was gone away from them, and then He said: “These things I said not unto you at the beginning, because I was with you; but now I go My way to Him That sent Me.” Let us first see whether it had been that He had not told them before this what they were to suffer in time coming. That He had done so amply before the night of the last Supper, is testified by the three first Evangelists, but it was when that Supper was ended that, according to John, He said: “These things I said not unto you at the beginning, because I was with you.”\par
\switchcolumn*

\LA
\begin{Resp}\Rsign Deus, cánticum novum cantábo tibi, allelúja: \Star{} In psaltério decem chordárum psallam tibi, allelúja, allelúja.\end{Resp}
\begin{Resp}\Vsign Deus meus es tu, et confitébor tibi: Deus meus es tu, et exaltábo te.\end{Resp}
\begin{Resp}\Rsign In psaltério decem chordárum psallam tibi, allelúja, allelúja.\end{Resp}
\switchcolumn
\EN
\begin{Resp}\Rsign I will sing a new song unto thee, O God, alleluia. \Star{} Upon a psaltery of ten strings will I sing praises unto thee, alleluia, alleluia.\end{Resp}
\begin{Resp}\Vsign Thou art my God, and I will praise thee Thou art my God, and I will exalt thee.\end{Resp}
\begin{Resp}\Rsign Upon a psaltery of ten strings will I sing praises unto thee, alleluia, alleluia.\end{Resp}
\switchcolumn*

\LA
\LessonHead{Lectio VIII}
\switchcolumn
\EN
\LessonHead{Lesson VIII}
\switchcolumn*

\LA
\noindent An forte hinc ista sólvitur quǽstio, quia et illi eum narrant passióni próximum fuísse cum hæc díceret? Non ergo ab inítio, quando cum illis erat: quia jam discessúrus, jamque ad Patrem perrectúrus hæc dixit. Et ídeo étiam secúndum illos Evangelístas verum est, quod hic dictum est: Hæc autem vobis ab inítio non dixi. Sed quid ágimus de fide Evangélii secúndum Matthǽum, qui hæc eis a Dómino non solum cum jam Pascha esset cum discípulis cœnatúrus, imminénte passióne, verum et ab inítio denuntiáta esse commémorat; ubi primum nominátim duódecim exprimúntur Apóstoli, et ad ópera divína mittúntur?\par
\switchcolumn
\EN
\noindent Are we then to try and loose the knot of this difficulty by asserting that, according to these three Evangelists, it was on the eve of the Passion, albeit before the Supper, that He had said these things unto them, and therefore not at the beginning, when He was with them, but when He was about to leave them, and go His way to the Father And in this way we might reconcile the truthfulness of what this Evangelist saith here “These things I said not unto you at the beginning” with the truthfulness of the! other three. But this explanation is rendered impossible by the Gospel according to Matthew, who telleth us how that the Lord spake to His Apostles concerning their sufferings to come, not only when He was on the point of eating the Passover with them, but at the very beginning, when the names of the twelve are first given, and they were sent forth to do the work of God. (Matth. x. 17-42.)\par
\switchcolumn*

\LA
\begin{Resp}\Rsign Bonum est confitéri Dómino, allelúja: \Star{} Et psállere, allelúja.\end{Resp}
\begin{Resp}\Vsign In decachórdo psaltério, cum cántico et cíthara.\end{Resp}
\begin{Resp}\Rsign Et psállere, allelúja.\end{Resp}
\begin{Resp}\Vsign Glória Patri, et Fílio, \Star{} et Spirítui Sancto.\end{Resp}
\begin{Resp}\Rsign Et psállere, allelúja.\end{Resp}
\switchcolumn
\EN
\begin{Resp}\Rsign It is a good thing to give thanks unto the Lord, alleluia. \Star{} And to sing praises, alleluia.\end{Resp}
\begin{Resp}\Vsign Upon an instrument of ten strings, upon the harp with a solemn sound.\end{Resp}
\begin{Resp}\Rsign And to sing praises, alleluia.\end{Resp}
\begin{Resp}\Vsign Glory be to the Father, and to the Son, \Star{} and to the Holy Ghost.\end{Resp}
\begin{Resp}\Rsign And to sing praises, alleluia.\end{Resp}
\switchcolumn*

\LA
\LessonHead{Lectio IX}
\switchcolumn
\EN
\LessonHead{Lesson IX}
\switchcolumn*

\LA
\noindent Quid sibi ergo vult, quod hic ait: Hæc autem vobis ab inítio non dixi, quia vobíscum eram: nisi quia ea, quæ hic dicit de Spíritu Sancto, quod sit ventúrus ad eos, et testimónium perhibitúrus, quando mala illa passúri sunt, hæc ab inítio eis non dixit, quia cum ipsis erat? Consolátor ergo ille, vel advocátus (utrúmque enim interpretátur, quod erat Græce Paráclitus), Christo abscedénte fúerat necessárius: et ídeo de illo non díxerat ab inítio, quando cum illis erat, quia ejus præséntia consolabántur.\par
\switchcolumn
\EN
\noindent It would seem then that when He said: “These things I said not unto you at the beginning, because I was with you,” He meant by “these things,” not the sufferings which they were to bear for His sake, but His promise of the Comforter Who should come to them, and testify while they suffered, (xv. 26, 27.) This Comforter then, or Advocate, (for the Greek word “Parakletos” will bear either interpretation,) would be needful to them when they saw Christ no more, and therefore it was that Christ spoke not of Him “at the beginning” (of the Gospel Dispensation) while He Himself “was with” His disciples, because His visible Presence was then their sufficient comfort.\par
\switchcolumn*

\LA
\TeDeum{Te Deum}
\switchcolumn
\EN
\TeDeum{Te Deum}
\switchcolumn*

\end{paracol}

\Office{Feria Secunda infra Hebdomadam IV post Octavam Paschæ}{Monday of the Fourth Week after the Octave of Easter}{Feria}
\addcontentsline{toc}{subsection}{Feria Secunda infra Hebdomadam IV post Octavam Paschæ\enspace\textit{Monday of the Fourth Week after the Octave of Easter}}
\begin{paracol}{2}
\LA
\LessonHead{Lectio I}
\switchcolumn
\EN
\LessonHead{Lesson I}
\switchcolumn*

\LA
\LessonTitle{De Epístola beáti Jacóbi Apóstoli}
\switchcolumn
\EN
\LessonTitle{Lesson from the letter of St. James the Apostle}
\switchcolumn*

\LA
\Cite{Jas 1:17-20}
\switchcolumn
\EN
\Cite{Jas 1:17-20}
\switchcolumn*

\LA
\noindent Omne datum óptimum, et omne donum perféctum desúrsum est, descéndens a Patre lúminum, apud quem non est transmutátio, nec vicissitúdinis obumbrátio. Voluntárie enim génuit nos verbo veritátis, ut simus inítium áliquod creatúræ ejus. Scitis, fratres mei dilectíssimi. Sit autem omnis homo velox ad audiéndum: tardus autem ad loquéndum, et tardus ad iram. Ira enim viri justítiam Dei non operátur.\par
\switchcolumn
\EN
\noindent Every best gift, and every perfect gift, is from above, coming down from the Father of lights, with whom there is no change, nor shadow of alteration. For of his own will hath he begotten us by the word of truth, that we might be some beginning of his creatures. You know, my dearest brethren. And let every man be swift to hear, but slow to speak, and slow to anger. For the anger of man worketh not the justice of God.\par
\switchcolumn*

\LA
\begin{Resp}\Rsign Dicant nunc, qui redémpti sunt, allelúja, \Star{} A Dómino, allelúja, allelúja.\end{Resp}
\begin{Resp}\Vsign Quos redémit de manu inimíci, et de regiónibus congregávit eos.\end{Resp}
\begin{Resp}\Rsign A Dómino, allelúja, allelúja.\end{Resp}
\switchcolumn
\EN
\begin{Resp}\Rsign Let now the redeemed of the Lord, alleluia. \Star{} Say, alleluia, alleluia, alleluia.\end{Resp}
\begin{Resp}\Vsign Let them whom He hath redeemed from the hand of the enemy, and gathered them out of the lands.\end{Resp}
\begin{Resp}\Rsign Say, alleluia, alleluia, alleluia.\end{Resp}
\switchcolumn*

\LA
\LessonHead{Lectio II}
\switchcolumn
\EN
\LessonHead{Lesson II}
\switchcolumn*

\LA
\Cite{Jas 1:21-24}
\switchcolumn
\EN
\Cite{Jas 1:21-24}
\switchcolumn*

\LA
\noindent Propter quod abiciéntes omnem immundítiam, et abundántiam malítiæ, in mansuetúdine suscípite ínsitum verbum, quod potest salváre ánimas vestras. Estóte autem factóres verbi, et non auditóres tantum: falléntes vosmetípsos. Quia si quis audítor est verbi, et non factor, hic comparábitur viro consideránti vultum nativitátis suæ in spéculo: considerávit enim se, et ábiit, et statim oblítus est qualis fúerit.\par
\switchcolumn
\EN
\noindent Wherefore casting away all uncleanness, and abundance of naughtiness, with meekness receive the ingrafted word, which is able to save your souls. But be ye doers of the word, and not hearers only, deceiving your own selves. For if a man be a hearer of the word, and not a doer, he shall be compared to a man beholding his own countenance in a glass. For he beheld himself, and went his way, and presently forgot what manner of man he was.\par
\switchcolumn*

\LA
\begin{Resp}\Rsign Cantáte Dómino, allelúja: \Star{} Psalmum dícite ei, allelúja.\end{Resp}
\begin{Resp}\Vsign Afférte Dómino glóriam et honórem, afférte Dómino glóriam nómini ejus.\end{Resp}
\begin{Resp}\Rsign Psalmum dícite ei, allelúja.\end{Resp}
\begin{Resp}\Vsign Glória Patri, et Fílio, \Star{} et Spirítui Sancto.\end{Resp}
\begin{Resp}\Rsign Psalmum dícite ei, allelúja.\end{Resp}
\switchcolumn
\EN
\begin{Resp}\Rsign O sing unto the Lord, alleluia. \Star{} Sing unto Him, alleluia.\end{Resp}
\begin{Resp}\Vsign Give unto the Lord glory and honour, give unto the Lord the glory due unto His Name.\end{Resp}
\begin{Resp}\Rsign Sing unto Him, alleluia.\end{Resp}
\begin{Resp}\Vsign Glory be to the Father, and to the Son, \Star{} and to the Holy Ghost.\end{Resp}
\begin{Resp}\Rsign Sing unto Him, alleluia.\end{Resp}
\switchcolumn*

\LA
\LessonHead{Lectio III}
\switchcolumn
\EN
\LessonHead{Lesson III}
\switchcolumn*

\LA
\Cite{Jas 1:25-27}
\switchcolumn
\EN
\Cite{Jas 1:25-27}
\switchcolumn*

\LA
\noindent Qui autem perspéxerit in legem perféctam libertátis, et permánserit in ea, non audítor obliviósus factus, sed factor óperis: hic beátus in facto suo erit. Si quis autem putat se religiósum esse, non refrénans linguam suam, sed sedúcens cor suum, hujus vana est relígio. Relígio munda et immaculáta apud Deum et Patrem, hæc est: visitáre pupíllos et víduas in tribulatióne eórum, et immaculátum se custodíre ab hoc sǽculo.\par
\switchcolumn
\EN
\noindent But he that hath looked into the perfect law of liberty, and hath continued therein, not becoming a forgetful hearer, but a doer of the work; this man shall be blessed in his deed. And if any man think himself to be religious, not bridling his tongue, but deceiving his own heart, this man's religion is vain. Religion clean and undefiled before God and the Father, is this: to visit the fatherless and widows in their tribulation: and to keep one's self unspotted from this world.\par
\switchcolumn*

\LA
\TeDeum{Te Deum}
\switchcolumn
\EN
\TeDeum{Te Deum}
\switchcolumn*

\end{paracol}

\Office{Feria Tertia infra Hebdomadam IV post Octavam Paschæ}{Tuesday of the Fourth Week after the Octave of Easter}{Feria}
\addcontentsline{toc}{subsection}{Feria Tertia infra Hebdomadam IV post Octavam Paschæ\enspace\textit{Tuesday of the Fourth Week after the Octave of Easter}}
\begin{paracol}{2}
\LA
\LessonHead{Lectio I}
\switchcolumn
\EN
\LessonHead{Lesson I}
\switchcolumn*

\LA
\LessonTitle{De Epístola beáti Jacóbi Apóstoli}
\switchcolumn
\EN
\LessonTitle{Lesson from the letter of St. James the Apostle}
\switchcolumn*

\LA
\Cite{Jas 2:1-4}
\switchcolumn
\EN
\Cite{Jas 2:1-4}
\switchcolumn*

\LA
\noindent Fratres mei, nolíte in personárum acceptióne habére fidem Dómini nostri Jesu Christi glóriæ. Etenim si introíerit in convéntum vestrum vir áureum ánnulum habens in veste cándida, introíerit autem et pauper in sórdido hábitu, et intendátis in eum qui indútus est veste præclára, et dixéritis ei: Tu sede hic bene: páuperi autem dicátis: Tu sta illic; aut sede sub scabéllo pedum meórum: nonne judicátis apud vosmetípsos, et facti estis júdices cogitatiónum iniquárum?\par
\switchcolumn
\EN
\noindent My brethren, have not the faith of our Lord Jesus Christ of glory with respect of persons. For if there shall come into your assembly a man having a golden ring, in fine apparel, and there shall come in also a poor man in mean attire, And you have respect to him that is clothed with the fine apparel, and shall say to him: Sit thou here well; but say to the poor man: Stand thou there, or sit under my footstool: Do you not judge within yourselves, and are become judges of unjust thoughts?\par
\switchcolumn*

\LA
\begin{Resp}\Rsign In ecclésiis benedícite Deo, allelúja: \Star{} Dómino de fóntibus Israël, allelúja, allelúja.\end{Resp}
\begin{Resp}\Vsign Psalmum dícite nómini ejus, date glóriam laudi ejus.\end{Resp}
\begin{Resp}\Rsign Dómino de fóntibus Israël, allelúja, allelúja.\end{Resp}
\switchcolumn
\EN
\begin{Resp}\Rsign Bless ye God in the congregations, alleluia. \Star{} Even the Lord, ye that are of the fountains of Israel, alleluia, alleluia.\end{Resp}
\begin{Resp}\Vsign Sing forth the honour of His Name, make His praise glorious.\end{Resp}
\begin{Resp}\Rsign Even the Lord, ye that are of the fountains of Israel, alleluia, alleluia.\end{Resp}
\switchcolumn*

\LA
\LessonHead{Lectio II}
\switchcolumn
\EN
\LessonHead{Lesson II}
\switchcolumn*

\LA
\Cite{Jas 2:5-9}
\switchcolumn
\EN
\Cite{Jas 2:5-9}
\switchcolumn*

\LA
\noindent Audíte, fratres mei dilectíssimi: nonne Deus elégit páuperes in hoc mundo, dívites in fide, et hærédes regni, quod repromísit Deus diligéntibus se? vos autem exhonorástis páuperem. Nonne dívites per poténtiam ópprimunt vos, et ipsi trahunt vos ad judícia? nonne ipsi blasphémant bonum nomen, quod invocátum est super vos? Si tamen legem perfícitis regálem secúndum Scriptúras: Díliges próximum tuum sicut teípsum: bene fácitis: si autem persónas accípitis, peccátum operámini, redargúti a lege quasi transgressóres.\par
\switchcolumn
\EN
\noindent Hearken, my dearest brethren: hath not God chosen the poor in this world, rich in faith, and heirs of the kingdom which God hath promised to them that love him? But you have dishonoured the poor man. Do not the rich oppress you by might? and do not they draw you before the judgment seats? Do not they blaspheme the good name that is invoked upon you? If then you fulfill the royal law, according to the scriptures, Thou shalt love thy neighbour as thyself; you do well. But if you have respect to persons, you commit sin, being reproved by the law as transgressors.\par
\switchcolumn*

\LA
\begin{Resp}\Rsign In toto corde meo, allelúja, exquisívi te, allelúja: \Star{} Ne repéllas me a mandátis tuis, allelúja, allelúja.\end{Resp}
\begin{Resp}\Vsign Benedíctus es tu, Dómine, doce me justificatiónes tuas.\end{Resp}
\begin{Resp}\Rsign Ne repéllas me a mandátis tuis, allelúja, allelúja.\end{Resp}
\begin{Resp}\Vsign Glória Patri, et Fílio, \Star{} et Spirítui Sancto.\end{Resp}
\begin{Resp}\Rsign Ne repéllas me a mandátis tuis, allelúja, allelúja.\end{Resp}
\switchcolumn
\EN
\begin{Resp}\Rsign With my whole heart, alleluia, have I sought thee, alleluia. \Star{} O let me not wander from thy commandments, alleluia, alleluia.\end{Resp}
\begin{Resp}\Vsign Blessed art Thou, O Lord, teach me thy statutes.\end{Resp}
\begin{Resp}\Rsign O let me not wander from thy commandments, alleluia, alleluia.\end{Resp}
\begin{Resp}\Vsign Glory be to the Father, and to the Son, \Star{} and to the Holy Ghost.\end{Resp}
\begin{Resp}\Rsign O let me not wander from thy commandments, alleluia, alleluia.\end{Resp}
\switchcolumn*

\LA
\LessonHead{Lectio III}
\switchcolumn
\EN
\LessonHead{Lesson III}
\switchcolumn*

\LA
\Cite{Jas 2:10-13}
\switchcolumn
\EN
\Cite{Jas 2:10-13}
\switchcolumn*

\LA
\noindent Quicúmque autem totam legem serváverit, offéndat autem in uno, factus est ómnium reus. Qui enim dixit: Non mœcháberis, dixit et: Non occídes. Quod si non mœcháberis, occídes autem, factus es transgréssor legis. Sic loquímini, et sic fácite sicut per legem libertátis incipiéntes judicári. Judícium enim sine misericórdia illi qui non fecit misericórdiam: superexáltat autem misericórdia judícium.\par
\switchcolumn
\EN
\noindent And whosoever shall keep the whole law, but offend in one point, is become guilty of all. For he that said, Thou shalt not commit adultery, said also, Thou shalt not kill. Now if thou do not commit adultery, but shalt kill, thou art become a transgressor of the law. So speak ye, and so do, as being to be judged by the law of liberty. For judgment without mercy to him that hath not done mercy. And mercy exalteth itself above judgment\par
\switchcolumn*

\LA
\TeDeum{Te Deum}
\switchcolumn
\EN
\TeDeum{Te Deum}
\switchcolumn*

\end{paracol}

\Office{Feria Quarta infra Hebdomadam IV post Octavam Paschæ}{Wednesday of the Fourth Week after the Octave of Easter}{Feria}
\addcontentsline{toc}{subsection}{Feria Quarta infra Hebdomadam IV post Octavam Paschæ\enspace\textit{Wednesday of the Fourth Week after the Octave of Easter}}
\begin{paracol}{2}
\LA
\LessonHead{Lectio I}
\switchcolumn
\EN
\LessonHead{Lesson I}
\switchcolumn*

\LA
\LessonTitle{De Epístola beáti Jacóbi Apóstoli}
\switchcolumn
\EN
\LessonTitle{Lesson from the letter of St. James the Apostle}
\switchcolumn*

\LA
\Cite{Jas 2:14-17}
\switchcolumn
\EN
\Cite{Jas 2:14-17}
\switchcolumn*

\LA
\noindent Quid próderit, fratres mei, si fidem quis dicat se habére, ópera autem non hábeat? numquid póterit fides salváre eum? Si autem frater et soror nudi sint, et indígeant victu cotidiáno, dicat autem áliquis ex vobis illis: Ite in pace, calefacímini et saturámini: non dedéritis autem eis quæ necessária sunt córpori, quid próderit? Sic et fides, si non hábeat ópera, mórtua est in semetípsa.\par
\switchcolumn
\EN
\noindent What shall it profit, my brethren, if a man say he hath faith, but hath not works? Shall faith be able to save him? And if a brother or sister be naked, and want daily food: And one of you say to them: Go in peace, be ye warmed and filled; yet give them not those things that are necessary for the body, what shall it profit? So faith also, if it have not works, is dead in itself.\par
\switchcolumn*

\LA
\begin{Resp}\Rsign Deus, cánticum novum cantábo tibi, allelúja: \Star{} In psaltério decem chordárum psallam tibi, allelúja, allelúja.\end{Resp}
\begin{Resp}\Vsign Deus meus es tu, et confitébor tibi: Deus meus es tu, et exaltábo te.\end{Resp}
\begin{Resp}\Rsign In psaltério decem chordárum psallam tibi, allelúja, allelúja.\end{Resp}
\switchcolumn
\EN
\begin{Resp}\Rsign I will sing a new song unto thee, O God, alleluia. \Star{} Upon a psaltery of ten strings will I sing praises unto thee, alleluia, alleluia.\end{Resp}
\begin{Resp}\Vsign Thou art my God, and I will praise thee Thou art my God, and I will exalt thee.\end{Resp}
\begin{Resp}\Rsign Upon a psaltery of ten strings will I sing praises unto thee, alleluia, alleluia.\end{Resp}
\switchcolumn*

\LA
\LessonHead{Lectio II}
\switchcolumn
\EN
\LessonHead{Lesson II}
\switchcolumn*

\LA
\Cite{Jas 2:18-22}
\switchcolumn
\EN
\Cite{Jas 2:18-22}
\switchcolumn*

\LA
\noindent Sed dicet quis: Tu fidem habes, et ego ópera hábeo: osténde mihi fidem tuam sine opéribus: et ego osténdam tibi ex opéribus fidem meam. Tu credis quóniam unus est Deus: bene facis: et dǽmones credunt, et contremíscunt. Vis autem scire, o homo inánis, quóniam fides sine opéribus mórtua est? Abraham pater noster nonne ex opéribus justificátus est, ófferens Isaac fílium suum super altáre? Vides quóniam fides cooperabátur opéribus illíus: et ex opéribus fides consummáta est?\par
\switchcolumn
\EN
\noindent But some man will say: Thou hast faith, and I have works: show me thy faith without works; and I will show thee, by works, my faith. Thou believest that there is one God. Thou dost well: the devils also believe and tremble. But wilt thou know, O vain man, that faith without works is dead? Was not Abraham our father justified by works, offering up Isaac his son upon the altar? Seest thou, that faith did co-operate with his works; and by works faith was made perfect?\par
\switchcolumn*

\LA
\begin{Resp}\Rsign Bonum est confitéri Dómino, allelúja: \Star{} Et psállere, allelúja.\end{Resp}
\begin{Resp}\Vsign In decachórdo psaltério, cum cántico et cíthara.\end{Resp}
\begin{Resp}\Rsign Et psállere, allelúja.\end{Resp}
\begin{Resp}\Vsign Glória Patri, et Fílio, \Star{} et Spirítui Sancto.\end{Resp}
\begin{Resp}\Rsign Et psállere, allelúja.\end{Resp}
\switchcolumn
\EN
\begin{Resp}\Rsign It is a good thing to give thanks unto the Lord, alleluia. \Star{} And to sing praises, alleluia.\end{Resp}
\begin{Resp}\Vsign Upon an instrument of ten strings, upon the harp with a solemn sound.\end{Resp}
\begin{Resp}\Rsign And to sing praises, alleluia.\end{Resp}
\begin{Resp}\Vsign Glory be to the Father, and to the Son, \Star{} and to the Holy Ghost.\end{Resp}
\begin{Resp}\Rsign And to sing praises, alleluia.\end{Resp}
\switchcolumn*

\LA
\LessonHead{Lectio III}
\switchcolumn
\EN
\LessonHead{Lesson III}
\switchcolumn*

\LA
\Cite{Jas 2:23-26}
\switchcolumn
\EN
\Cite{Jas 2:23-26}
\switchcolumn*

\LA
\noindent Et suppléta est Scriptúra, dicens: Crédidit Abraham Deo, et reputátum est illi ad justítiam, et amícus Dei appellátus est. Vidétis quóniam ex opéribus justificátur homo, et non ex fide tantum? Simíliter et Rahab méretrix, nonne ex opéribus justificáta est, suscípiens núntios, et ália via eíciens? Sicut enim corpus sine spíritu mórtuum est, ita et fides sine opéribus mórtua est.\par
\switchcolumn
\EN
\noindent And the scripture was fulfilled, saying: Abraham believed God, and it was reputed to him to justice, and he was called the friend of God. Do you see that by works a man is justified; and not by faith only? And in like manner also Rahab the harlot, was not she justified by works, receiving the messengers, and sending them out another way? For even as the body without the spirit is dead; so also faith without works is dead.\par
\switchcolumn*

\LA
\TeDeum{Te Deum}
\switchcolumn
\EN
\TeDeum{Te Deum}
\switchcolumn*

\end{paracol}

\Office{Feria Quinta infra Hebdomadam IV post Octavam Paschæ}{Thursday of the Fourth Week after the Octave of Easter}{Feria}
\addcontentsline{toc}{subsection}{Feria Quinta infra Hebdomadam IV post Octavam Paschæ\enspace\textit{Thursday of the Fourth Week after the Octave of Easter}}
\begin{paracol}{2}
\LA
\LessonHead{Lectio I}
\switchcolumn
\EN
\LessonHead{Lesson I}
\switchcolumn*

\LA
\LessonTitle{De Epístola beáti Jacóbi Apóstoli}
\switchcolumn
\EN
\LessonTitle{Lesson from the letter of St. James the Apostle}
\switchcolumn*

\LA
\Cite{Jas 3:1-3}
\switchcolumn
\EN
\Cite{Jas 3:1-3}
\switchcolumn*

\LA
\noindent Nolíte plures magístri fíeri fratres mei, sciéntes quóniam majus judícium súmitis. In multis enim offéndimus omnes. Si quis in verbo non offéndit, hic perféctus est vir: potest étiam freno circumdúcere totum corpus. Si autem equis frena in ora míttimus ad consentiéndum nobis, et omne corpus illórum circumférimus.\par
\switchcolumn
\EN
\noindent Be ye not many masters, my brethren, knowing that you receive the greater judgment. For in many things we all offend. If any man offend not in word, the same is a perfect man. He is able also with a bridle to lead about the whole body. For if we put bits into the mouths of horses, that they may obey us, and we turn about their whole body.\par
\switchcolumn*

\LA
\begin{Resp}\Rsign Si oblítus fúero tui, allelúja, obliviscátur mei déxtera mea: \Star{} Adhǽreat lingua mea fáucibus meis, si non memínero tui, allelúja, allelúja.\end{Resp}
\begin{Resp}\Vsign Super flúmina Babylónis illic sédimus et flévimus, dum recordarémur tui, Sion.\end{Resp}
\begin{Resp}\Rsign Adhǽreat lingua mea fáucibus meis, si non memínero tui, allelúja, allelúja.\end{Resp}
\switchcolumn
\EN
\begin{Resp}\Rsign If I forget thee, Alleluia, let my right hand forget me. \Star{} If I do not remember thee, let my tongue cleave to the roof of my mouth, alleluia, alleluia.\end{Resp}
\begin{Resp}\Vsign By the rivers of Babylon there we sat down and wept, when we remembered thee, O Zion.\end{Resp}
\begin{Resp}\Rsign If I do not remember thee, let my tongue cleave to the roof of my mouth, alleluia, alleluia.\end{Resp}
\switchcolumn*

\LA
\LessonHead{Lectio II}
\switchcolumn
\EN
\LessonHead{Lesson II}
\switchcolumn*

\LA
\Cite{Jas 3:4-6}
\switchcolumn
\EN
\Cite{Jas 3:4-6}
\switchcolumn*

\LA
\noindent Ecce et naves, cum magnæ sint, et a ventis válidis minéntur, circumferúntur a módico gubernáculo ubi ímpetus dirigéntis volúerit. Ita et lingua módicum quidem membrum est, et magna exáltat. Ecce quantus ignis quam magnam silvam incéndit! Et lingua ignis est, univérsitas iniquitátis.\par
\switchcolumn
\EN
\noindent Behold also ships, whereas they are great, and are driven by strong winds, yet are they turned about with a small helm, whithersoever the force of the governor willeth. Even so the tongue is indeed a little member, and boasteth great things. Behold how small a fire kindleth a great wood. And the tongue is a fire, a world of iniquity.\par
\switchcolumn*

\LA
\begin{Resp}\Rsign Vidérunt te aquæ, Deus, vidérunt te aquæ, et timuérunt: \Star{} Multitúdo sónitus aquárum vocem dedérunt nubes, allelúja, allelúja, allelúja.\end{Resp}
\begin{Resp}\Vsign Illuxérunt coruscatiónes tuæ orbi terræ: vidit et commóta est terra.\end{Resp}
\begin{Resp}\Rsign Multitúdo sónitus aquárum vocem dedérunt nubes, allelúja, allelúja, allelúja.\end{Resp}
\begin{Resp}\Vsign Glória Patri, et Fílio, \Star{} et Spirítui Sancto.\end{Resp}
\begin{Resp}\Rsign Multitúdo sónitus aquárum vocem dedérunt nubes, allelúja, allelúja, allelúja.\end{Resp}
\switchcolumn
\EN
\begin{Resp}\Rsign The waters saw thee, O God, the waters saw thee and they were afraid. \Star{} There was a noise as of many waters the clouds sent out a sound. Alleluia, Alleluia, Alleluia.\end{Resp}
\begin{Resp}\Vsign Thy lightnings lightened the world; the earth saw it and shook.\end{Resp}
\begin{Resp}\Rsign There was a noise as of many waters; the clouds sent out a sound. Alleluia, Alleluia, Alleluia.\end{Resp}
\begin{Resp}\Vsign Glory be to the Father, and to the Son, \Star{} and to the Holy Ghost.\end{Resp}
\begin{Resp}\Rsign There was a noise as of many waters; the clouds sent out a sound. Alleluia, Alleluia, Alleluia.\end{Resp}
\switchcolumn*

\LA
\LessonHead{Lectio III}
\switchcolumn
\EN
\LessonHead{Lesson III}
\switchcolumn*

\LA
\Cite{Jas 3:6-10}
\switchcolumn
\EN
\Cite{Jas 3:6-10}
\switchcolumn*

\LA
\noindent Lingua constitúitur in membris nostris, quæ máculat totum corpus, et inflámmat rotam nativitátis nostræ inflammáta a gehénna. Omnis enim natúra bestiárum, et vólucrum, et serpéntium, et ceterórum domántur, et dómita sunt a natúra humána: linguam autem nullus hóminum domáre potest: inquiétum malum, plena venéno mortífero. In ipsa benedícimus Deum et Patrem: et in ipsa maledícimus hómines, qui ad similitúdinem Dei facti sunt. Ex ipso ore procédit benedíctio et maledíctio.\par
\switchcolumn
\EN
\noindent The tongue is placed among our members, which defileth the whole body, and inflameth the wheel of our nativity, being set on fire by hell. For every nature of beasts, and of birds, and of serpents, and of the rest, is tamed, and hath been tamed, by the nature of man: But the tongue no man can tame, an unquiet evil, full of deadly poison. By it we bless God and the Father: and by it we curse men, who are made after the likeness of God. Out of the same mouth proceedeth blessing and cursing. My brethren, these things ought not so to be.\par
\switchcolumn*

\LA
\TeDeum{Te Deum}
\switchcolumn
\EN
\TeDeum{Te Deum}
\switchcolumn*

\end{paracol}

\Office{Feria Sexta infra Hebdomadam IV post Octavam Paschæ}{Friday of the Fourth Week after the Octave of Easter}{Feria}
\addcontentsline{toc}{subsection}{Feria Sexta infra Hebdomadam IV post Octavam Paschæ\enspace\textit{Friday of the Fourth Week after the Octave of Easter}}
\begin{paracol}{2}
\LA
\LessonHead{Lectio I}
\switchcolumn
\EN
\LessonHead{Lesson I}
\switchcolumn*

\LA
\LessonTitle{De Epístola beáti Jacóbi Apóstoli}
\switchcolumn
\EN
\LessonTitle{Lesson from the letter of St. James the Apostle}
\switchcolumn*

\LA
\Cite{Jas 4:1-4}
\switchcolumn
\EN
\Cite{Jas 4:1-4}
\switchcolumn*

\LA
\noindent Unde bella et lites in vobis? nonne hinc: ex concupiscéntiis vestris, quæ mílitant in membris vestris? concupíscitis, et non habétis: occíditis, et zelátis: et non potéstis adipísci: litigátis, et belligerátis, et non habétis, propter quod non postulátis. Pétitis, et non accípitis: eo quod male petátis: ut in concupiscéntiis vestris insumátis. Adúlteri, nescítis quia amicítia hujus mundi inimíca est Dei? quicúmque ergo volúerit amícus esse sǽculi hujus, inimícus Dei constitúitur.\par
\switchcolumn
\EN
\noindent From whence are wars and contentions among you? Are they not hence, from your concupiscences, which war in your members? You covet, and have not: you kill, and envy, and can not obtain. You contend and war, and you have not, because you ask not. You ask, and receive not; because you ask amiss: that you may consume it on your concupiscences. Adulterers, know you not that the friendship of this world is the enemy of God? Whosoever therefore will be a friend of this world, becometh an enemy of God.\par
\switchcolumn*

\LA
\begin{Resp}\Rsign In ecclésiis benedícite Deo, allelúja: \Star{} Dómino de fóntibus Israël, allelúja, allelúja.\end{Resp}
\begin{Resp}\Vsign Psalmum dícite nómini ejus, date glóriam laudi ejus.\end{Resp}
\begin{Resp}\Rsign Dómino de fóntibus Israël, allelúja, allelúja.\end{Resp}
\switchcolumn
\EN
\begin{Resp}\Rsign Bless ye God in the congregations, alleluia. \Star{} Even the Lord, ye that are of the fountains of Israel, alleluia, alleluia.\end{Resp}
\begin{Resp}\Vsign Sing forth the honour of His Name, make His praise glorious.\end{Resp}
\begin{Resp}\Rsign Even the Lord, ye that are of the fountains of Israel, alleluia, alleluia.\end{Resp}
\switchcolumn*

\LA
\LessonHead{Lectio II}
\switchcolumn
\EN
\LessonHead{Lesson II}
\switchcolumn*

\LA
\Cite{Jas 4:5-10}
\switchcolumn
\EN
\Cite{Jas 4:5-10}
\switchcolumn*

\LA
\noindent An putátis quia inániter Scriptúra dicat: Ad invídiam concupíscit spíritus qui hábitat in vobis? majórem autem dat grátiam. Propter quod dicit: Deus supérbis resístit, humílibus autem dat grátiam. Súbditi ergo estóte Deo, resístite autem diábolo, et fúgiet a vobis. Appropinquáte Deo, et appropinquábit vobis. Emundáte manus, peccatóres: et purificáte corda, dúplices ánimo. Míseri estóte, et lugéte, et ploráte: risus vester in luctum convertátur, et gáudium in mœrórem. Humiliámini in conspéctu Dómini, et exaltábit vos.\par
\switchcolumn
\EN
\noindent Or do you think that the scripture saith in vain: To envy doth the spirit covet which dwelleth in you? But he giveth greater grace. Wherefore he saith: God resisteth the proud, and giveth grace to the humble. Be subject therefore to God, but resist the devil, and he will fly from you. Draw nigh to God, and he will draw nigh to you. Cleanse your hands, ye sinners: and purify your hearts, ye double minded. Be afflicted, and mourn, and weep: let your laughter be turned into mourning, and your joy into sorrow. Be humbled in the sight of the Lord, and he will exalt you.\par
\switchcolumn*

\LA
\begin{Resp}\Rsign In toto corde meo, allelúja, exquisívi te, allelúja: \Star{} Ne repéllas me a mandátis tuis, allelúja, allelúja.\end{Resp}
\begin{Resp}\Vsign Benedíctus es tu, Dómine, doce me justificatiónes tuas.\end{Resp}
\begin{Resp}\Rsign Ne repéllas me a mandátis tuis, allelúja, allelúja.\end{Resp}
\begin{Resp}\Vsign Glória Patri, et Fílio, \Star{} et Spirítui Sancto.\end{Resp}
\begin{Resp}\Rsign Ne repéllas me a mandátis tuis, allelúja, allelúja.\end{Resp}
\switchcolumn
\EN
\begin{Resp}\Rsign With my whole heart, alleluia, have I sought thee, alleluia. \Star{} O let me not wander from thy commandments. Alleluia, Alleluia.\end{Resp}
\begin{Resp}\Vsign Blessed art Thou, O Lord, teach me thy statutes.\end{Resp}
\begin{Resp}\Rsign O let me not wander from thy commandments. Alleluia, Alleluia.\end{Resp}
\begin{Resp}\Vsign Glory be to the Father, and to the Son, \Star{} and to the Holy Ghost.\end{Resp}
\begin{Resp}\Rsign O let me not wander from thy commandments. Alleluia, Alleluia.\end{Resp}
\switchcolumn*

\LA
\LessonHead{Lectio III}
\switchcolumn
\EN
\LessonHead{Lesson III}
\switchcolumn*

\LA
\Cite{Jas 4:11-15}
\switchcolumn
\EN
\Cite{Jas 4:11-15}
\switchcolumn*

\LA
\noindent Nolíte detráhere altérutrum fratres. Qui détrahit fratri, aut qui júdicat fratrem suum, détrahit legi, et júdicat legem. Si autem júdicas legem, non es factor legis, sed judex. Unus est legislátor et judex, qui potest pérdere et liberáre. Tu autem quis es, qui júdicas próximum? Ecce nunc qui dícitis: Hódie, aut crástino íbimus in illam civitátem, et faciémus ibi quidem annum, et mercábimur, et lucrum faciémus: qui ignorátis quid erit in crástino. Quæ est enim vita vestra? vapor est ad módicum parens, et deínceps exterminábitur; pro eo ut dicátis: Si Dóminus volúerit. Et: Si vixérimus, faciémus hoc, aut illud.\par
\switchcolumn
\EN
\noindent Detract not one another, my brethren. He that detracteth his brother, or he that judgeth his brother, detracteth the law, and judgeth the law. But if thou judge the law, thou art not a doer of the law, but a judge. There is one lawgiver, and judge, that is able to destroy and to deliver. But who art thou that judgest thy neighbour? Behold, now you that say: Today or tomorrow we will go into such a city, and there we will spend a year, and will traffic, and make our gain. Whereas you know not what shall be on the morrow. For what is your life? It is a vapour which appeareth for a little while, and afterwards shall vanish away. For that you should say: If the Lord will, and if we shall live, we will do this or that.\par
\switchcolumn*

\LA
\TeDeum{Te Deum}
\switchcolumn
\EN
\TeDeum{Te Deum}
\switchcolumn*

\end{paracol}

\Office{Sabbato infra Hebdomadam IV post Octavam Paschæ}{Saturday of the Fourth Week after the Octave of Easter}{Feria}
\addcontentsline{toc}{subsection}{Sabbato infra Hebdomadam IV post Octavam Paschæ\enspace\textit{Saturday of the Fourth Week after the Octave of Easter}}
\begin{paracol}{2}
\LA
\LessonHead{Lectio I}
\switchcolumn
\EN
\LessonHead{Lesson I}
\switchcolumn*

\LA
\LessonTitle{De Epístola beáti Jacóbi Apóstoli}
\switchcolumn
\EN
\LessonTitle{Lesson from the letter of St. James the Apostle}
\switchcolumn*

\LA
\Cite{Jas 5:1-6}
\switchcolumn
\EN
\Cite{Jas 5:1-6}
\switchcolumn*

\LA
\noindent Agite nunc dívites, ploráte ululántes in misériis vestris, quæ advénient vobis. Divítiæ vestræ putrefáctæ sunt, et vestiménta vestra a tíneis comésta sunt. Aurum et argéntum vestrum æruginávit: et ærúgo eórum in testimónium vobis erit, et manducábit carnes vestras sicut ignis. Thesaurizástis vobis iram in novíssimis diébus. Ecce merces operariórum, qui messuérunt regiónes vestras, quæ fraudáta est a vobis, clamat: et clamor eórum in aures Dómini sábbaoth introívit. Epuláti estis super terram, et in luxúriis enutrístis corda vestra in die occisiónis. Addixístis, et occidístis justum, et non resístit vobis.\par
\switchcolumn
\EN
\noindent Go to now, ye rich men, weep and howl in your miseries, which shall come upon you. Your riches are corrupted: and your garments are moth-eaten. Your gold and silver is cankered: and the rust of them shall be for a testimony against you, and shall eat your flesh like fire. You have stored up to yourselves wrath against the last days. Behold the hire of the labourers, who have reaped down your fields, which by fraud has been kept back by you, crieth: and the cry of them hath entered into the ears of the Lord of Sabaoth. You have feasted upon earth: and in riotousness you have nourished your hearts, in the day of slaughter. You have condemned and put to death the Just One, and he resisted you not.\par
\switchcolumn*

\LA
\begin{Resp}\Rsign Deus, cánticum novum cantábo tibi, allelúja: \Star{} In psaltério decem chordárum psallam tibi, allelúja, allelúja.\end{Resp}
\begin{Resp}\Vsign Deus meus es tu, et confitébor tibi: Deus meus es tu, et exaltábo te.\end{Resp}
\begin{Resp}\Rsign In psaltério decem chordárum psallam tibi, allelúja, allelúja.\end{Resp}
\switchcolumn
\EN
\begin{Resp}\Rsign I will sing a new song unto thee, O God, alleluia. \Star{} Upon a psaltery of ten strings will I sing praises unto thee, alleluia, alleluia.\end{Resp}
\begin{Resp}\Vsign Thou art my God, and I will praise thee Thou art my God, and I will exalt thee.\end{Resp}
\begin{Resp}\Rsign Upon a psaltery of ten strings will I sing praises unto thee, alleluia, alleluia.\end{Resp}
\switchcolumn*

\LA
\LessonHead{Lectio II}
\switchcolumn
\EN
\LessonHead{Lesson II}
\switchcolumn*

\LA
\Cite{Jas 5:7-11}
\switchcolumn
\EN
\Cite{Jas 5:7-11}
\switchcolumn*

\LA
\noindent Patiéntes ígitur estóte, fratres, usque ad advéntum Dómini. Ecce agrícola exspéctat pretiósum fructum terræ, patiénter ferens donec accípiat temporáneum et serótinum. Patiéntes ígitur estóte et vos, et confirmáte corda vestra: quóniam advéntus Dómini appropinquávit. Nolíte ingemíscere, fratres, in altérutrum, ut non judicémini. Ecce judex ante jánuam assístit. Exémplum accípite, fratres, éxitus mali, labóris, et patiéntiæ, prophétas qui locúti sunt in nómine Dómini. Ecce beatificámus eos qui sustinuérunt. Sufferéntiam Job audístis, et finem Dómini vidístis, quóniam miséricors Dóminus est, et miserátor.\par
\switchcolumn
\EN
\noindent Be patient therefore, brethren, until the coming of the Lord. Behold, the husbandman waiteth for the precious fruit of the earth: patiently bearing till he receive the early and latter rain. Be you therefore also patient, and strengthen your hearts: for the coming of the Lord is at hand. Grudge not, brethren, one against another, that you may not be judged. Behold the judge standeth before the door. Take, my brethren, for an example of suffering evil, of labour and patience, the prophets, who spoke in the name of the Lord. Behold, we account them blessed who have endured. You have heard of the patience of Job, and you have seen the end of the Lord, that the Lord is merciful and compassionate\par
\switchcolumn*

\LA
\begin{Resp}\Rsign Bonum est confitéri Dómino, allelúja: \Star{} Et psállere, allelúja.\end{Resp}
\begin{Resp}\Vsign In decachórdo psaltério, cum cántico et cíthara.\end{Resp}
\begin{Resp}\Rsign Et psállere, allelúja.\end{Resp}
\switchcolumn
\EN
\begin{Resp}\Rsign It is a good thing to give thanks unto the Lord. Alleluia. \Star{} And to sing praises. Alleluia.\end{Resp}
\begin{Resp}\Vsign Upon an instrument of ten strings, upon the harp with a solemn sound.\end{Resp}
\begin{Resp}\Rsign And to sing praises. Alleluia.\end{Resp}
\switchcolumn*

\LA
\LessonHead{Lectio III}
\switchcolumn
\EN
\LessonHead{Lesson III}
\switchcolumn*

\LA
\Cite{Jas 5:12-16}
\switchcolumn
\EN
\Cite{Jas 5:12-16}
\switchcolumn*

\LA
\noindent Ante ómnia autem, fratres mei, nolíte juráre, neque per cælum, neque per terram, neque áliud quodcúmque juraméntum. Sit autem sermo vester: Est, est: Non, non: ut non sub judício decidátis. Tristátur áliquis vestrum? oret. Æquo ánimo est? psallat. Infirmátur quis in vobis? indúcat presbýteros ecclésiæ, et orent super eum, ungéntes eum óleo in nómine Dómini: et orátio fídei salvábit infírmum, et alleviábit eum Dóminus: et si in peccátis sit, remitténtur ei. Confitémini ergo altérutrum peccáta vestra, et oráte pro ínvicem ut salvémini: multum enim valet deprecátio justi assídua.\par
\switchcolumn
\EN
\noindent But above all things, my brethren, swear not, neither by heaven, nor by the earth, nor by any other oath. But let your speech be, yea, yea: no, no: that you fall not under judgment. Is any of you sad? Let him pray. Is he cheerful in mind? Let him sing. Is any man sick among you? Let him bring in the priests of the church, and let them pray over him, anointing him with oil in the name of the Lord. And the prayer of faith shall save the sick man: and the Lord shall raise him up: and if he be in sins, they shall be forgiven him. Confess therefore your sins one to another: and pray one for another, that you may be saved. For the continual prayer of a just man availeth much.\par
\switchcolumn*

\LA
\begin{Resp}\Rsign Cantáte Dómino, allelúja: \Star{} Psalmum dícite ei, allelúja.\end{Resp}
\begin{Resp}\Vsign Afférte Dómino glóriam et honórem, afférte Dómino glóriam nómini ejus.\end{Resp}
\begin{Resp}\Rsign Psalmum dícite ei, allelúja.\end{Resp}
\begin{Resp}\Vsign Glória Patri, et Fílio, \Star{} et Spirítui Sancto.\end{Resp}
\begin{Resp}\Rsign Psalmum dícite ei, allelúja.\end{Resp}
\switchcolumn
\EN
\begin{Resp}\Rsign O sing unto the Lord, alleluia. \Star{} Sing unto Him, alleluia.\end{Resp}
\begin{Resp}\Vsign Give unto the Lord glory and honour, give unto the Lord the glory due unto His Name.\end{Resp}
\begin{Resp}\Rsign Sing unto Him, alleluia.\end{Resp}
\begin{Resp}\Vsign Glory be to the Father, and to the Son, \Star{} and to the Holy Ghost.\end{Resp}
\begin{Resp}\Rsign Sing unto Him, alleluia.\end{Resp}
\switchcolumn*

\end{paracol}

\Office{Dominica V Post Pascha}{Fifth Sunday after Easter}{Semiduplex}
\addcontentsline{toc}{subsection}{Dominica V Post Pascha\enspace\textit{Fifth Sunday after Easter}}
\begin{paracol}{2}
\LA
\LessonHead{Lectio I}
\switchcolumn
\EN
\LessonHead{Lesson I}
\switchcolumn*

\LA
\LessonTitle{Incipit Epístola prima beáti Petri Apóstoli}
\switchcolumn
\EN
\LessonTitle{Lesson from the first letter of St. Peter the Apostle}
\switchcolumn*

\LA
\Cite{1 Pet 1:1-5}
\switchcolumn
\EN
\Cite{1 Pet 1:1-5}
\switchcolumn*

\LA
\noindent Petrus Apóstolus Jesu Christi, eléctis ádvenis dispersiónis Ponti, Galátiæ, Cappadóciæ, Asiæ, et Bithýniæ, secúndum præsciéntiam Dei Patris, in sanctificatiónem Spíritus, in obediéntiam, et aspersiónem sánguinis Jesu Christi. Grátia vobis, et pax multiplicétur. Benedíctus Deus et Pater Dómini nostri Jesu Christi, qui secúndum misericórdiam suam magnam regenerávit nos in spem vivam, per resurrectiónem Jesu Christi ex mórtuis, in hæreditátem incorruptíbilem, et incontaminátam, et immarcescíbilem, conservátam in cælis in vobis, qui in virtúte Dei custodímini per fidem in salútem, parátam revelári in témpore novíssimo.\par
\switchcolumn
\EN
\noindent Peter, an apostle of Jesus Christ, to the strangers dispersed through Pontus, Galatia, Cappadocia, Asia, and Bithynia, elect, According to the foreknowledge of God the Father, unto the sanctification of the Spirit, unto obedience and sprinkling of the blood of Jesus Christ: Grace unto you and peace be multiplied. Blessed be the God and Father of our Lord Jesus Christ, who according to his great mercy hath regenerated us unto a lively hope, by the resurrection of Jesus Christ from the dead, Unto an inheritance incorruptible, and undefiled, and that can not fade, reserved in heaven for you, Who, by the power of God, are kept by faith unto salvation, ready to be revealed in the last time.\par
\switchcolumn*

\LA
\begin{Resp}\Rsign Si oblítus fúero tui, allelúja, obliviscátur mei déxtera mea: \Star{} Adhǽreat lingua mea fáucibus meis, si non memínero tui, allelúja, allelúja.\end{Resp}
\begin{Resp}\Vsign Super flúmina Babylónis illic sédimus et flévimus, dum recordarémur tui, Sion.\end{Resp}
\begin{Resp}\Rsign Adhǽreat lingua mea fáucibus meis, si non memínero tui, allelúja, allelúja.\end{Resp}
\switchcolumn
\EN
\begin{Resp}\Rsign If I forget thee, Alleluia, let my right hand forget me. \Star{} If I do not remember thee, let my tongue cleave to the roof of my mouth, alleluia, alleluia.\end{Resp}
\begin{Resp}\Vsign By the rivers of Babylon there we sat down and wept, when we remembered thee, O Zion.\end{Resp}
\begin{Resp}\Rsign If I do not remember thee, let my tongue cleave to the roof of my mouth, alleluia, alleluia.\end{Resp}
\switchcolumn*

\LA
\LessonHead{Lectio II}
\switchcolumn
\EN
\LessonHead{Lesson II}
\switchcolumn*

\LA
\Cite{1 Pet 1:6-12}
\switchcolumn
\EN
\Cite{1 Pet 1:6-12}
\switchcolumn*

\LA
\noindent In quo exsultábitis, módicum nunc si opórtet contristári in váriis tentatiónibus: ut probátio vestræ fídei multo pretiósior auro (quod per ignem probátur) inveniátur in laudem, et glóriam, et honórem in revelatióne Jesu Christi: quem cum non vidéritis, dilígitis: in quem nunc quoque non vidéntes créditis: credéntes autem exsultábitis lætítia inenarrábili, et glorificáta: reportántes finem fídei vestræ, salútem animárum. De qua salúte exquisiérunt, atque scrutáti sunt prophétæ, qui de futúra in vobis grátia prophetavérunt: scrutántes in quod vel quale tempus significáret in eis Spíritus Christi: prænúntians eas quæ in Christo sunt passiónes, et posterióres glórias: quibus revelátum est quia non sibimetípsis, vobis autem ministrábant ea quæ nunc nuntiáta sunt vobis per eos qui evangelizavérunt vobis, Spíritu Sancto misso de cælo, in quem desíderant ángeli prospícere.\par
\switchcolumn
\EN
\noindent Wherein you shall greatly rejoice, if now you must be for a little time made sorrowful in diverse temptations: That the trial of your faith (much more precious than gold which is tried by the fire) may be found unto praise and glory and honour at the appearing of Jesus Christ: Whom having not seen, you love: in whom also now, though you see him not, you believe: and believing shall rejoice with joy unspeakable and glorified; Receiving the end of your faith, even the salvation of your souls. Of which salvation the prophets have inquired and diligently searched, who prophesied of the grace to come in you. Searching what or what manner of time the Spirit of Christ in them did signify: when it foretold those sufferings that are in Christ, and the glories that should follow: To whom it was revealed, that not to themselves, but to you they ministered those things which are now declared to you by them that have preached the gospel to you, the Holy Ghost being sent down from heaven, on whom the angels desire to look.\par
\switchcolumn*

\LA
\begin{Resp}\Rsign Vidérunt te aquæ, Deus, vidérunt te aquæ, et timuérunt: \Star{} Multitúdo sónitus aquárum vocem dedérunt nubes, allelúja, allelúja, allelúja.\end{Resp}
\begin{Resp}\Vsign Illuxérunt coruscatiónes tuæ orbi terræ: vidit et commóta est terra.\end{Resp}
\begin{Resp}\Rsign Multitúdo sónitus aquárum vocem dedérunt nubes, allelúja, allelúja, allelúja.\end{Resp}
\switchcolumn
\EN
\begin{Resp}\Rsign The waters saw thee, O God, the waters saw thee and they were afraid. \Star{} There was a noise as of many waters the clouds sent out a sound, alleluia, alleluia, alleluia.\end{Resp}
\begin{Resp}\Vsign thy lightnings lightened the world the earth saw it and shook.\end{Resp}
\begin{Resp}\Rsign There was a noise as of many waters the clouds sent out a sound, alleluia, alleluia, alleluia.\end{Resp}
\switchcolumn*

\LA
\LessonHead{Lectio III}
\switchcolumn
\EN
\LessonHead{Lesson III}
\switchcolumn*

\LA
\Cite{1 Pet 1:13-21}
\switchcolumn
\EN
\Cite{1 Pet 1:13-21}
\switchcolumn*

\LA
\noindent Propter quod succíncti lumbos mentis vestræ, sóbrii, perfécte speráte in eam, quæ offértur vobis, grátiam, in revelatiónem Jesu Christi: quasi fílii obediéntiæ, non configuráti prióribus ignorántiæ vestræ desidériis: sed secúndum eum qui vocávit vos, Sanctum: et ipsi in omni conversatióne sancti sitis: quóniam scriptum est: Sancti éritis, quóniam ego sanctus sum. Et si patrem invocátis eum, qui sine acceptióne personárum júdicat secúndum uniuscujúsque opus, in timóre incolátus vestri témpore conversámini. Sciéntes quod non corruptibílibus, auro vel argénto, redémpti estis de vana vestra conversatióne patérnæ traditiónis: sed pretióso sánguine quasi agni immaculáti Christi, et incontamináti: præcógniti quidem ante mundi constitutiónem, manifestáti autem novíssimis tempóribus propter vos, qui per ipsum fidéles estis in Deo, qui suscitávit eum a mórtuis, et dedit ei glóriam, ut fides vestra et spes esset in Deo:\par
\switchcolumn
\EN
\noindent Wherefore having the loins of your mind girt up, being sober, trust perfectly in the grace which is offered you in the revelation of Jesus Christ, As children of obedience, not fashioned according to the former desires of your ignorance: But according to him that hath called you, who is holy, be you also in all manner of conversation holy: Because it is written: You shall be holy, for I am holy. And if you invoke as Father him who, without respect of persons, judgeth according to every one's work: converse in fear during the time of your sojourning here. Knowing that you were not redeemed with corruptible things as gold or silver, from your vain conversation of the tradition of your fathers: But with the precious blood of Christ, as of a lamb unspotted and undefiled, Foreknown indeed before the foundation of the world, but manifested in the last times for you, Who through him are faithful in God, who raised him up from the dead, and hath given him glory, that your faith and hope might be in God.\par
\switchcolumn*

\LA
\begin{Resp}\Rsign Narrábo nomen tuum frátribus meis, allelúja: \Star{} In médio Ecclésiæ laudábo te, allelúja, allelúja.\end{Resp}
\begin{Resp}\Vsign Confitébor tibi in pópulis, Dómine, et psalmum dicam tibi in géntibus.\end{Resp}
\begin{Resp}\Rsign In médio Ecclésiæ laudábo te, allelúja, allelúja.\end{Resp}
\begin{Resp}\Vsign Glória Patri, et Fílio, \Star{} et Spirítui Sancto.\end{Resp}
\begin{Resp}\Rsign In médio Ecclésiæ laudábo te, allelúja, allelúja.\end{Resp}
\switchcolumn
\EN
\begin{Resp}\Rsign I will declare thy Name unto my brethren, alleluia. \Star{} In the midst of the congregation will I praise thee, alleluia, alleluia.\end{Resp}
\begin{Resp}\Vsign I will praise thee, O Lord, among the people, and sing unto thee among the nations.\end{Resp}
\begin{Resp}\Rsign In the midst of the congregation will I praise thee, alleluia, alleluia.\end{Resp}
\begin{Resp}\Vsign Glory be to the Father, and to the Son, \Star{} and to the Holy Ghost.\end{Resp}
\begin{Resp}\Rsign In the midst of the congregation will I praise thee, alleluia, alleluia.\end{Resp}
\switchcolumn*

\LA
\LessonHead{Lectio IV}
\switchcolumn
\EN
\LessonHead{Lesson IV}
\switchcolumn*

\LA
\LessonTitle{Ex libro sancti Ambrósii Epíscopi de fide resurrectiónis}
\switchcolumn
\EN
\LessonTitle{From the Book written by St. Ambrose, Bishop of Milan, on belief in the Resurrection.}
\switchcolumn*

\LA
\Source{Post medium}
\switchcolumn
\EN
\Source{After middle}
\switchcolumn*

\LA
\noindent Quóniam Dei mori non póterat Sapiéntia, resúrgere autem non póterat quod mórtuum non erat; assúmitur caro, quæ mori posset: ut dum móritur quod solet, quod mórtuum fúerat, hoc resúrgeret. Neque enim póterat esse, nisi per hóminem, resurréctio: quóniam sicut per hóminem mors, ita et per hóminem resurréctio mortuórum. Ergo resurréxit homo, quóniam homo mórtuus est: resuscitátus homo, sed resúscitans Deus. Tunc secúndum carnem homo, nunc per ómnia Deus. Nunc enim secúndum carnem jam nóvimus Christum, sed carnis grátiam tenémus, ut ipsum primítias quiescéntium, ipsum primogénitum ex mórtuis novérimus.\par
\switchcolumn
\EN
\noindent Since it was impossible that the Wisdom of God could die, and that which could not die could not rise from the dead, He took to Himself Flesh Which could die, that That Whose nature it was to die might die, and rise again. Neither was it possible that the resurrection of the dead should come otherwise than by man, “for since by man came death, by Man came also the resurrection of the dead.” (1 Cor. xv. 21.) Man He rose since Man He died, the Manhood quickened but the Godhead Quickener. Man then, as touching the Flesh God now, over all things. For now we know Christ no longer after the Flesh, but we owe it to the Flesh that we know Him as “become the First-fruits of them that slept” (1 Cor. xv. 23) “and the First-begotten of the dead” (Apoc. i. 5.)\par
\switchcolumn*

\LA
\begin{Resp}\Rsign In ecclésiis benedícite Deo, allelúja: \Star{} Dómino de fóntibus Israël, allelúja, allelúja.\end{Resp}
\begin{Resp}\Vsign Psalmum dícite nómini ejus, date glóriam laudi ejus.\end{Resp}
\begin{Resp}\Rsign Dómino de fóntibus Israël, allelúja, allelúja.\end{Resp}
\switchcolumn
\EN
\begin{Resp}\Rsign Bless ye God in the congregations, alleluia. \Star{} Even the Lord, ye that are of the fountains of Israel, alleluia, alleluia.\end{Resp}
\begin{Resp}\Vsign Sing forth the honour of His Name, make His praise glorious.\end{Resp}
\begin{Resp}\Rsign Even the Lord, ye that are of the fountains of Israel, alleluia, alleluia.\end{Resp}
\switchcolumn*

\LA
\LessonHead{Lectio V}
\switchcolumn
\EN
\LessonHead{Lesson V}
\switchcolumn*

\LA
\noindent Primítiæ útique ejúsdem sunt géneris atque natúræ, cujus et réliqui fructus: quorum pro lætióre provéntu primitíva Deo múnera deferúntur; sacrum munus pro ómnibus, et quasi reparátæ quædam liba natúræ. Primítiæ ergo quiescéntium Christus. Sed utrum suórum quiescéntium, qui quasi mortis exsórtes, dulci quodam sopóre tenéntur, an ómnium mortuórum? Sed sicut in Adam omnes moriúntur, ita et in Christo omnes vivificabúntur. Itaque sicut primítiæ mortis in Adam, ita étiam primítiæ resurrectiónis in Christo omnes resúrgent. Sed nemo despéret, neque justus dóleat commúne consórtium resurgéndi, cum præcípuum fructum virtútis exspéctet. Omnes quidem resúrgent, sed unusquísque, ut ait Apóstolus, in suo órdine. Commúnis est divínæ fructus cleméntiæ, sed distínctus ordo meritórum.\par
\switchcolumn
\EN
\noindent The first-fruits are of the same kind and nature as the other fruits, and they are brought as an offering to God, to win His blessing on the in-gathering, an holy offering made on behalf of all, and as it were the homage of restored nature. Christ then is the First-fruits of them that sleep. But is He the First-fruits of only His own loved ones that fall asleep in Him, and lie as it were untouched by death, wrapt in a sweet slumber Or is He the First-fruits of all the dead? But “as in Adam all die, even so in Christ shall all be made alive.” (1 Cor. xv. 22.) So that, as in Adam were the first-fruits of the death wherein all die, even so in Christ were the first-fruits of the resurrection, wherein all rise again. But let no man be hopeless, neither let it be a grief to the righteous to remember that to rise again will be common to all men, when he looketh for that day wherein the harvest of his life will nobly realise itself. All shall rise again, “but,” as saith the Apostle (23,) “every man in his own order.” The harvest of God's mercy will be for all, but in reward one man shall differ from another.\par
\switchcolumn*

\LA
\begin{Resp}\Rsign In toto corde meo, allelúja, exquisívi te, allelúja: \Star{} Ne repéllas me a mandátis tuis, allelúja, allelúja.\end{Resp}
\begin{Resp}\Vsign Benedíctus es tu, Dómine, doce me justificatiónes tuas.\end{Resp}
\begin{Resp}\Rsign Ne repéllas me a mandátis tuis, allelúja, allelúja.\end{Resp}
\switchcolumn
\EN
\begin{Resp}\Rsign With my whole heart, alleluia, have I sought thee, alleluia. \Star{} O let me not wander from thy commandments, alleluia, alleluia.\end{Resp}
\begin{Resp}\Vsign Blessed art Thou, O Lord; teach me thy statutes.\end{Resp}
\begin{Resp}\Rsign O let me not wander from thy commandments, alleluia, alleluia.\end{Resp}
\switchcolumn*

\LA
\LessonHead{Lectio VI}
\switchcolumn
\EN
\LessonHead{Lesson VI}
\switchcolumn*

\LA
\noindent Advértimus, quam grave sit sacrilégium, resurrectiónem non crédere. Si enim non resurgémus, ergo Christus gratis mórtuus est, ergo Christus non resurréxit. Si enim nobis non resurréxit, útique non resurréxit, qui sibi cur resúrgeret, non habébat. Resurréxit in eo mundus, resurréxit in eo cælum, resurréxit in eo terra. Erit enim cælum novum, et terra nova. Sibi autem non erat necessária resurréctio, quem mortis víncula non tenébant. Nam etsi secúndum hóminem mórtuus, in ipsis tamen erat liber inférnis. Vis scire quam liber? Factus sum sicut homo sine adjutório, inter mórtuos liber. Et bene liber, qui se póterat suscitáre, juxta quod scriptum est: Sólvite hoc templum, et in tríduo resuscitábo illud. Et bene liber, qui álios descénderat redemptúrus.\par
\switchcolumn
\EN
\noindent I tell you how grievous an outrage against God it is not to believe in the resurrection. If we shall Is this not rise again, then did Christ die in vain, then is Christ not risen For if [if He rose at all,] He rose for us, and if He had not us to rise for, then He is plainly not risen. In Him the world, in Him the heavens, in Him the earth rose again. For there shall be “a new heaven, and a new earth” (Apoc. xxi. 1.) For Himself He needed not to rise Whom the bands of death held not. For although He died as Man, yet was He free in the netherworld itself. Wouldest thou hear how free “I am as a man that hath no strength, free among the dead” (Ps. Ixxxvii. 6.) O how free Who was able to take up his life again at will (John x. 18), even as it is written that He said “Destroy this Temple, and in three days I will raise it up” (John ii. 19.) O how free Who descended into hell only to redeem others therefrom.\par
\switchcolumn*

\LA
\begin{Resp}\Rsign Hymnum cantáte nobis, allelúja: \Star{} Quómodo cantábimus cánticum Dómini in terra aliéna? allelúja, allelúja.\end{Resp}
\begin{Resp}\Vsign Illic interrogavérunt nos, qui captívos duxérunt nos, verba cantiónum.\end{Resp}
\begin{Resp}\Rsign Quómodo cantábimus cánticum Dómini in terra aliéna? allelúja, allelúja.\end{Resp}
\begin{Resp}\Vsign Glória Patri, et Fílio, \Star{} et Spirítui Sancto.\end{Resp}
\begin{Resp}\Rsign Quómodo cantábimus cánticum Dómini in terra aliéna? allelúja, allelúja.\end{Resp}
\switchcolumn
\EN
\begin{Resp}\Rsign Sing us a song, alleluia. \Star{} How shall we sing the Lord's song in a strange land, alleluia, alleluia.\end{Resp}
\begin{Resp}\Vsign There they that carried us away captive required of us a song.\end{Resp}
\begin{Resp}\Rsign How shall we sing the Lord's song in a strange land, alleluia, alleluia.\end{Resp}
\begin{Resp}\Vsign Glory be to the Father, and to the Son, \Star{} and to the Holy Ghost.\end{Resp}
\begin{Resp}\Rsign How shall we sing the Lord's song in a strange land, alleluia, alleluia.\end{Resp}
\switchcolumn*

\LA
\LessonHead{Lectio VII}
\switchcolumn
\EN
\LessonHead{Lesson VII}
\switchcolumn*

\LA
\LessonTitle{Léctio sancti Evangélii secúndum Joánnem}
\switchcolumn
\EN
\LessonTitle{From the Holy Gospel according to John}
\switchcolumn*

\LA
\Cite{Joannes 16:23-30}
\switchcolumn
\EN
\Cite{John 16:23-30}
\switchcolumn*

\LA
\noindent In illo témpore: Dixit Jesus discípulis suis: Amen, amen dico vobis: si quid petiéritis Patrem in nómine meo, dabit vobis. Et réliqua.\par
\switchcolumn
\EN
\noindent At that time, Jesus said unto His disciples: Amen, Amen, I say unto you: Whatsoever ye shall ask the Father in My Name, He will give it you. And so on.\par
\switchcolumn*

\LA
\LessonTitle{Homilía sancti Augustíni Epíscopi}
\switchcolumn
\EN
\LessonTitle{Homily by St. Augustine, Bishop of Hippo.}
\switchcolumn*

\LA
\Source{Tractatus 102 in Joannem}
\switchcolumn
\EN
\Source{102nd Tract on John.}
\switchcolumn*

\LA
\noindent Dómini verba nunc ista tractánda sunt: Amen, amen, dico vobis: Si quid petiéritis Patrem in nómine meo, dabit vobis. Jam dictum est in superióribus hujus Domínici sermónis pártibus, propter eos, qui nonnúlla petunt a Patre in Christi nómine, nec accípiunt: non peti in nómine Salvatóris, quidquid pétitur contra ratiónem salútis. Non enim sonum litterárum ac syllabárum, sed quod sonus ipse signíficat, et quod eo sono recte ac veráciter intellégitur, hoc accipiéndus est dícere, cum dicit: In nómine meo.\par
\switchcolumn
\EN
\noindent We have now to consider these words of the Lord “Amen, Amen, I say unto you, Whatsoever ye shall ask the Father in My Name, He will give it you.” It hath already been said in the earlier part of this discourse of the Lord, for the sake of some who ask the Father in Christ's Name and receive not, that whatsoever is asked, which tendeth not to salvation, is not asked in the Name of the Saviour. By the words “In My Name” we must not understand the vocalization of letters and syllables, but the meaning of what is said, the honest and true meaning.\par
\switchcolumn*

\LA
\begin{Resp}\Rsign Deus, cánticum novum cantábo tibi, allelúja: \Star{} In psaltério decem chordárum psallam tibi, allelúja, allelúja.\end{Resp}
\begin{Resp}\Vsign Deus meus es tu, et confitébor tibi: Deus meus es tu, et exaltábo te.\end{Resp}
\begin{Resp}\Rsign In psaltério decem chordárum psallam tibi, allelúja, allelúja.\end{Resp}
\switchcolumn
\EN
\begin{Resp}\Rsign I will sing a new song unto thee, O God, alleluia. \Star{} Upon a psaltery of ten strings will I sing praises unto thee, alleluia, alleluia.\end{Resp}
\begin{Resp}\Vsign Thou art my God, and I will praise thee Thou art my God, and I will exalt thee.\end{Resp}
\begin{Resp}\Rsign Upon a psaltery of ten strings will I sing praises unto thee, alleluia, alleluia.\end{Resp}
\switchcolumn*

\LA
\LessonHead{Lectio VIII}
\switchcolumn
\EN
\LessonHead{Lesson VIII}
\switchcolumn*

\LA
\noindent Unde qui hoc sentit de Christo, quod non est de único Dei Fílio sentiéndum, non petit in ejus nómine, etiámsi non táceat lítteris ac sýllabis Christum: quóniam in ejus nómine petit, quem cógitat cum petit. Qui vero quod est de illo sentiéndum sentit, ipse in ejus nómine petit: et áccipit quod petit, si non contra suam salútem sempitérnam petit. Accipit autem quando debet accípere. Quædam enim non negántur: sed ut cóngruo dentur témpore differúntur. Ita sane intelligéndum est quod ait: Dabit vobis: ut ea benefícia significáta sciántur his verbis, quæ ad eos, qui petunt, próprie pértinent. Exaudiúntur quippe omnes Sancti pro seípsis, non autem pro ómnibus exaudiúntur vel amícis, vel inimícis suis, vel quibúslibet áliis: quia non utcúmque dictum est, Dabit; sed, Dabit vobis.\par
\switchcolumn
\EN
\noindent Therefore, whosoever thinketh of Christ as he ought not to think of the Only Son of God, such an one doth not ask anything in Christ's Name, although he do actually utter letters and syllables to that effect, because by these sounds he meaneth not the Real Christ, but a fancied being who hath no existence except in the speaker's imagination. But on the other hand, whosoever thinketh of Christ as he ought to think, the same asketh in Christ's Name, and receiveth, provided only it be nothing against his own everlasting salvation but if it is good for him to receive, he receiveth. Some things are not given at once, but kept over till a more fitting season. Such is the true interpretation of the words “He will give it you” namely, that those things will be given which are good for them to ask. All the Saints also are heard when they ask for themselves, but not necessarily when they ask for their friends, or their enemies, or others, even as it is written, not simply “He will give it” but “He will give it you.”\par
\switchcolumn*

\LA
\begin{Resp}\Rsign Bonum est confitéri Dómino, allelúja: \Star{} Et psállere, allelúja.\end{Resp}
\begin{Resp}\Vsign In decachórdo psaltério, cum cántico et cíthara.\end{Resp}
\begin{Resp}\Rsign Et psállere, allelúja.\end{Resp}
\begin{Resp}\Vsign Glória Patri, et Fílio, \Star{} et Spirítui Sancto.\end{Resp}
\begin{Resp}\Rsign Et psállere, allelúja.\end{Resp}
\switchcolumn
\EN
\begin{Resp}\Rsign It is a good thing to give thanks unto the Lord, alleluia. \Star{} And to sing praises, alleluia.\end{Resp}
\begin{Resp}\Vsign Upon an instrument of ten strings, upon the harp with a solemn sound.\end{Resp}
\begin{Resp}\Rsign And to sing praises, alleluia.\end{Resp}
\begin{Resp}\Vsign Glory be to the Father, and to the Son, \Star{} and to the Holy Ghost.\end{Resp}
\begin{Resp}\Rsign And to sing praises, alleluia.\end{Resp}
\switchcolumn*

\LA
\LessonHead{Lectio IX}
\switchcolumn
\EN
\LessonHead{Lesson IX}
\switchcolumn*

\LA
\noindent Usque modo, inquit, non petístis quidquam in nómine meo. Pétite, et accipiétis, ut gáudium vestrum sit plenum. Hoc quod dicit, gáudium plenum, profécto non carnále, sed spiritále gáudium est: et quando tantum erit, ut áliquid ei jam non sit addéndum, proculdúbio tunc erit plenum. Quidquid ergo pétitur, quod pertíneat ad hoc gáudium consequéndum, hoc est in nómine Christi peténdum, si divínam intellígimus grátiam, si vere beátam póscimus vitam. Quidquid autem áliud pétitur, nihil pétitur: non quia nulla omníno res est, sed quia in tantæ rei comparatióne quidquid áliud concupíscitur, nihil est.\par
\switchcolumn
\EN
\noindent Hitherto saith the Lord, have ye asked nothing in My Name? “Ask, and ye shall receive, that your joy may be full.” This their joy, whereof He saith that it shall be full, is to be understood not of fleshly but of spiritual joy and when that joy is so great that it can be increased no more, then shall it without doubt be full. Whatsoever therefore we ask for the fulfilling of this joy, (that is, if we thereby mean grace, if we ask for that life which is the really blessed one,) that is a thing which it is meet to ask in Christ's Name. If we ask anything else than this, we ask nothing, although we do actually ask something, because all things are nothing in comparison with this.\par
\switchcolumn*

\LA
\TeDeum{Te Deum}
\switchcolumn
\EN
\TeDeum{Te Deum}
\switchcolumn*

\end{paracol}

\Office{Feria Secunda in Rogationibus}{Monday in Rogationtide}{Feria major}
\addcontentsline{toc}{subsection}{Feria Secunda in Rogationibus\enspace\textit{Monday in Rogationtide}}
\begin{paracol}{2}
\LA
\LessonHead{Lectio I}
\switchcolumn
\EN
\LessonHead{Lesson I}
\switchcolumn*

\LA
\LessonTitle{Léctio sancti Evangélii secúndum Lucam}
\switchcolumn
\EN
\LessonTitle{Continuation of the Holy Gospel according to Luke}
\switchcolumn*

\LA
\Cite{Luc 11:5-13}
\switchcolumn
\EN
\Cite{Luke 11:5-13}
\switchcolumn*

\LA
\noindent In illo témpore: Dixit Jesus discípulis suis: Quis vestrum habébit amícum, et ibit ad illum media nocte, et dicet illi: Amice, cómmoda mihi tres panes. Et réliqua.\par
\switchcolumn
\EN
\noindent At that time Jesus said unto His disciples Which of you shall have a friend, and shall go unto him at midnight, and say unto him Friend, lend me three loaves. And so on.\par
\switchcolumn*

\LA
\LessonTitle{Homilía sancti Ambrósii Epíscopi}
\switchcolumn
\EN
\LessonTitle{Homily by St. Ambrose, Bishop of Milan.}
\switchcolumn*

\LA
\Source{Liber 7 in Lucæ cap. 11}
\switchcolumn
\EN
\Source{Book vii. on Luke xi.}
\switchcolumn*

\LA
\noindent Alius præcépti locus est, ut ómnibus moméntis, non solum diébus, sed étiam noctibus, orátio deferátur. Vides enim, quod iste qui media nocte perréxit, tres panes ab amico suo póstulans, et in ipsa peténdi intentióne persístens, non defraudétur oratis. Qui sunt isti tres panes, nisi mysterii cæléstis aliméntum? Quod si diligas Dóminum Deum tuum, non solum tibi, sed étiam aliis póteris emereri. Quis autem amicior nobis, quam qui pro nobis corpus suum trádidit?\par
\switchcolumn
\EN
\noindent We gather from this commandment, among other things, that we ought to pray, not only by day, but also by night. Thou seest how that he which arose at midnight to ask three loaves of his friend, and endured in supplication, was not disappointed of that which he sought. Of what are these three loaves a figure, but of that our Mysterious Bread Which cometh down from heaven Thou seest that if thou lovest the Lord thy God, thou mayest win His bounty, not only for thyself, but for others likewise. And who can deserve more to be called our “Friend” than He Which gave His Own Body for us\par
\switchcolumn*

\LA
\begin{Resp}\Rsign Dicant nunc, qui redémpti sunt, allelúja, \Star{} A Dómino, allelúja, allelúja.\end{Resp}
\begin{Resp}\Vsign Quos redémit de manu inimíci, et de regiónibus congregávit eos.\end{Resp}
\begin{Resp}\Rsign A Dómino, allelúja, allelúja.\end{Resp}
\switchcolumn
\EN
\begin{Resp}\Rsign Let now the redeemed of the Lord, alleluia. \Star{} Say, alleluia, alleluia, alleluia.\end{Resp}
\begin{Resp}\Vsign Let them whom He hath redeemed from the hand of the enemy, and gathered them out of the lands.\end{Resp}
\begin{Resp}\Rsign Say, alleluia, alleluia, alleluia.\end{Resp}
\switchcolumn*

\LA
\LessonHead{Lectio II}
\switchcolumn
\EN
\LessonHead{Lesson II}
\switchcolumn*

\LA
\noindent Ab hoc media nocte panes David pétiit, et accépit. Petiit enim, quando dicebat: Media nocte surgébam ad confiténdum tibi. Ideo meruit hos panes, quos appósuit nobis edendos. Petiit cum dicit: Lavábo per síngulas noctes lectum meum. Neque enim tímuit, ne excitaret dormiéntem, quem scit semper vigilántem. Et ideo scriptórum mémores, noctibus ac diébus oratióni instántes, peccátis nostris véniam postulemus.\par
\switchcolumn
\EN
\noindent From this Friend it was that David asked bread at midnight, and received it, as he saith “At midnight I rise to give thanks unto thee.” (Ps. cxviii. 62.) Even thus did he obtain those loaves of spiritual nourishment which he still setteth before us for our refreshment. How he asked it, we know from that he saith “Every night wash I my bed.” (Ps. vi. 7.) He knew that there was no fear of waking Him Who sleepeth not. (Ps. cxx. 3.) Therefore let us keep in mind the things which are written for our learning, and be instant in prayer both by day and by night, to ask pardon of our sins.\par
\switchcolumn*

\LA
\begin{Resp}\Rsign Cantáte Dómino, allelúja: \Star{} Psalmum dícite ei, allelúja.\end{Resp}
\begin{Resp}\Vsign Afférte Dómino glóriam et honórem, afférte Dómino glóriam nómini ejus.\end{Resp}
\begin{Resp}\Rsign Psalmum dícite ei, allelúja.\end{Resp}
\switchcolumn
\EN
\begin{Resp}\Rsign O sing unto the Lord, alleluia. \Star{} Sing unto Him, alleluia.\end{Resp}
\begin{Resp}\Vsign Give unto the Lord glory and honour, give unto the Lord the glory due unto His Name.\end{Resp}
\begin{Resp}\Rsign Sing unto Him, alleluia.\end{Resp}
\switchcolumn*

\LA
\LessonHead{Lectio III}
\switchcolumn
\EN
\LessonHead{Lesson III}
\switchcolumn*

\LA
\noindent Nam si ille tam sanctus, et qui regni erat necessitátibus occupatus, septies in die laudem Dómino dicebat, matutínis, et vespertinis sacrifíciis semper intentus; quid nos fácere oportet, qui eo ámplius rogare debemus, quo frequéntius carnis ac mentis fragilitáte delinquimus, ut de via lassis, et istíus ævi cursu, ac vitæ hujus anfractu gráviter fatigátis, panis refectiónis deésse non possit, qui hóminis corda confírmet? Nec solum media nocte Dóminus, sed ómnibus prope docet vigilándum esse moméntis. Venit enim et vespertína, et secunda, et tertia vigília: et pulsáre consuévit. Beáti ítaque servi illi, quos cum venerit Dóminus, invénerit vigilántes.\par
\switchcolumn
\EN
\noindent If David, who was such a Saint, and whose time was so taken up by the cares of a kingdom, praised the Lord seven times a day, (Ps. cxviii. 164,) and was always present with godly zeal at the morning and evening sacrifice, what ought we to do, (who have so much the more need to pray, as the weakness of our body and mind doth so much oftener make us to fall, ) that we, wearied with this pilgrimage, and worn out by the gradual waning of our earthly day, and the changes of life, that we, I say, may not be | starved of that life-giving Bread Which strengtheneth man's heart. The Lord teacheth us to be watchful, all of us, and that, not at midnight only, but always. And if He shall come in the second watch, or come in the third watch, and find them so blessed are those servants whom the Lord, when He cometh, shall find watching. (Luke xii. 37.)\par
\switchcolumn*

\LA
\begin{Resp}\Rsign Narrábo nomen tuum frátribus meis, allelúja: \Star{} In médio Ecclésiæ laudábo te, allelúja, allelúja.\end{Resp}
\begin{Resp}\Vsign Confitébor tibi in pópulis, Dómine, et psalmum dicam tibi in géntibus.\end{Resp}
\begin{Resp}\Rsign In médio Ecclésiæ laudábo te, allelúja, allelúja.\end{Resp}
\begin{Resp}\Vsign Glória Patri, et Fílio, \Star{} et Spirítui Sancto.\end{Resp}
\begin{Resp}\Rsign In médio Ecclésiæ laudábo te, allelúja, allelúja.\end{Resp}
\switchcolumn
\EN
\begin{Resp}\Rsign I will declare thy Name unto my brethren, alleluia. \Star{} In the midst of the congregation will I praise thee, alleluia, alleluia.\end{Resp}
\begin{Resp}\Vsign I will praise thee, O Lord, among the people, and sing unto thee among the nations.\end{Resp}
\begin{Resp}\Rsign In the midst of the congregation will I praise thee, alleluia, alleluia.\end{Resp}
\begin{Resp}\Vsign Glory be to the Father, and to the Son, \Star{} and to the Holy Ghost.\end{Resp}
\begin{Resp}\Rsign In the midst of the congregation will I praise thee, alleluia, alleluia.\end{Resp}
\switchcolumn*

\end{paracol}

\Office{Feria Tertia in Rogationibus}{Tuesday in Rogationtide}{Feria}
\addcontentsline{toc}{subsection}{Feria Tertia in Rogationibus\enspace\textit{Tuesday in Rogationtide}}
\begin{paracol}{2}
\LA
\LessonHead{Lectio I}
\switchcolumn
\EN
\LessonHead{Lesson I}
\switchcolumn*

\LA
\LessonTitle{De Epístola prima beáti Petri Apóstoli}
\switchcolumn
\EN
\LessonTitle{Lesson from the first letter of St. Peter the Apostle}
\switchcolumn*

\LA
\Cite{1 Pet 4:1-7}
\switchcolumn
\EN
\Cite{1 Pet 4:1-7}
\switchcolumn*

\LA
\noindent Christo ígitur passo in carne, et vos eádem cogitatióne armámini: quia qui passus est in carne, désiit a peccátis: ut jam non desidériis hóminum, sed voluntáti Dei, quod réliquum est in carne vivat témporis. Súfficit enim prætéritum tempus ad voluntátem géntium consummándam his qui ambulavérunt in luxúriis, desidériis, vinoléntiis, comessatiónibus, potatiónibus, et illícitis idolórum cúltibus. In quo admirántur non concurréntibus vobis in eándem luxúriæ confusiónem, blasphemántes. Qui reddent ratiónem ei qui parátus est judicáre vivos et mórtuos. Propter hoc enim et mórtuis evangelizátum est: ut judicéntur quidem secúndum hómines in carne, vivant autem secúndum Deum in spíritu. Omnium autem finis appropinquávit.\par
\switchcolumn
\EN
\noindent Christ therefore having suffered in the flesh, be you also armed with the same thought: for he that hath suffered in the flesh, hath ceased from sins: That now he may live the rest of his time in the flesh, not after the desires of men, but according to the will of God. For the time past is sufficient to have fulfilled the will of the Gentiles, for them who have walked in riotousness, lusts, excess of wine, revellings, banquetings, and unlawful worshipping of idols. Wherein they think it strange, that you run not with them into the same confusion of riotousness, speaking evil of you. Who shall render account to him, who is ready to judge the living and the dead. For, for this cause was the gospel preached also to the dead: that they might be judged indeed according to men, in the flesh; but may live according to God, in the Spirit. But the end of all is at hand.\par
\switchcolumn*

\LA
\begin{Resp}\Rsign In ecclésiis benedícite Deo, allelúja: \Star{} Dómino de fóntibus Israël, allelúja, allelúja.\end{Resp}
\begin{Resp}\Vsign Psalmum dícite nómini ejus, date glóriam laudi ejus.\end{Resp}
\begin{Resp}\Rsign Dómino de fóntibus Israël, allelúja, allelúja.\end{Resp}
\switchcolumn
\EN
\begin{Resp}\Rsign Bless ye God in the congregations, alleluia. \Star{} Even the Lord, ye that are of the fountains of Israel, alleluia, alleluia.\end{Resp}
\begin{Resp}\Vsign Sing forth the honour of His Name, make His praise glorious.\end{Resp}
\begin{Resp}\Rsign Even the Lord, ye that are of the fountains of Israel, alleluia, alleluia.\end{Resp}
\switchcolumn*

\LA
\LessonHead{Lectio II}
\switchcolumn
\EN
\LessonHead{Lesson II}
\switchcolumn*

\LA
\Cite{1 Pet 4:7-11}
\switchcolumn
\EN
\Cite{1 Pet 4:7-11}
\switchcolumn*

\LA
\noindent Estóte ítaque prudéntes, et vigiláte in oratiónibus. Ante ómnia autem, mútuam in vobismetípsis caritátem contínuam habéntes: quia cáritas óperit multitúdinem peccatórum. Hospitáles ínvicem sine murmuratióne. Unusquísque, sicut accépit grátiam, in altérutrum illam administrántes, sicut boni dispensatóres multifórmis grátiæ Dei. Si quis lóquitur, quasi sermónes Dei: si quis minístrat, tamquam ex virtúte, quam adminístrat Deus: ut in ómnibus honorificétur Deus per Jesum Christum: cui est glória et impérium in sǽcula sæculórum. Amen.\par
\switchcolumn
\EN
\noindent Be prudent therefore, and watch in prayers. But before all things have a constant mutual charity among yourselves: for charity covereth a multitude of sins. Using hospitality one towards another, without murmuring, As every man hath received grace, ministering the same one to another: as good stewards of the manifold grace of God. If any man speak, let him speak, as the words of God. If any man minister, let him do it, as of the power, which God administereth: that in all things God may be honoured through Jesus Christ: to whom is glory and empire for ever and ever. Amen.\par
\switchcolumn*

\LA
\begin{Resp}\Rsign In toto corde meo, allelúja, exquisívi te, allelúja: \Star{} Ne repéllas me a mandátis tuis, allelúja, allelúja.\end{Resp}
\begin{Resp}\Vsign Benedíctus es tu, Dómine, doce me justificatiónes tuas.\end{Resp}
\begin{Resp}\Rsign Ne repéllas me a mandátis tuis, allelúja, allelúja.\end{Resp}
\begin{Resp}\Vsign Glória Patri, et Fílio, \Star{} et Spirítui Sancto.\end{Resp}
\begin{Resp}\Rsign Ne repéllas me a mandátis tuis, allelúja, allelúja.\end{Resp}
\switchcolumn
\EN
\begin{Resp}\Rsign With my whole heart, alleluia, have I sought thee, alleluia. \Star{} O let me not wander from thy commandments. Alleluia, Alleluia.\end{Resp}
\begin{Resp}\Vsign Blessed art Thou, O Lord teach me thy statutes.\end{Resp}
\begin{Resp}\Rsign O let me not wander from thy commandments. Alleluia, Alleluia.\end{Resp}
\begin{Resp}\Vsign Glory be to the Father, and to the Son, \Star{} and to the Holy Ghost.\end{Resp}
\begin{Resp}\Rsign O let me not wander from thy commandments. Alleluia, Alleluia.\end{Resp}
\switchcolumn*

\LA
\LessonHead{Lectio III}
\switchcolumn
\EN
\LessonHead{Lesson III}
\switchcolumn*

\LA
\Cite{1 Pet 4:12-17}
\switchcolumn
\EN
\Cite{1 Pet 4:12-17}
\switchcolumn*

\LA
\noindent Caríssimi, nolíte peregrinári in fervóre, qui ad tentatiónem vobis fit, quasi novi áliquid vobis contíngat: sed communicántes Christi passiónibus gaudéte, ut et in revelatióne glóriæ ejus gaudeátis exsultántes. Si exprobrámini in nómine Christi, beáti éritis: quóniam quod est honóris, glóriæ, et virtútis Dei, et qui est ejus Spíritus, super vos requiéscit. Nemo autem vestrum patiátur ut homicída, aut fur, aut malédicus, aut alienórum appetítor. Si autem ut Christiánus, non erubéscat: gloríficet autem Deum in isto nómine: quóniam tempus est ut incípiat judícium a domo Dei.\par
\switchcolumn
\EN
\noindent Dearly beloved, think not strange the burning heat which is to try you, as if some new thing happened to you; But if you partake of the sufferings of Christ, rejoice that when his glory shall be revealed, you may also be glad with exceeding joy. If you be reproached for the name of Christ, you shall be blessed: for that which is of the honour, glory, and power of God, and that which is his Spirit, resteth upon you. But let none of you suffer as a murderer, or a thief, or a railer, or a coveter of other men's things. But if as a Christian, let him not be ashamed, but let him glorify God in that name. For the time is, that judgment should begin at the house of God.\par
\switchcolumn*

\LA
\TeDeum{Te Deum}
\switchcolumn
\EN
\TeDeum{Te Deum}
\switchcolumn*

\end{paracol}

\Office{Feria Quarta in Rogationibus in Vigilia Ascensionis}{Wednesday in Rogationtide, Vigil of the Ascension}{Vigilia}
\addcontentsline{toc}{subsection}{Feria Quarta in Rogationibus in Vigilia Ascensionis\enspace\textit{Wednesday in Rogationtide, Vigil of the Ascension}}
\begin{paracol}{2}
\LA
\LessonHead{Lectio I}
\switchcolumn
\EN
\LessonHead{Lesson I}
\switchcolumn*

\LA
\LessonTitle{Léctio sancti Evangélii secúndum Joánnem}
\switchcolumn
\EN
\LessonTitle{Continuation of the Holy Gospel according to John}
\switchcolumn*

\LA
\Cite{Joannes 17:1-11}
\switchcolumn
\EN
\Cite{John 17:1-11}
\switchcolumn*

\LA
\noindent In illo témpore: Sublevátis Jesus óculis in cælum, dixit: Pater, venit hora, clarífica Fílium tuum. Et réliqua.\par
\switchcolumn
\EN
\noindent At that time, Jesus lifted up His Eyes to heaven, and spoke these words: Father, the hour is come; glorify thy Son. And so on.\par
\switchcolumn*

\LA
\LessonTitle{Homilía sancti Augustíni Epíscopi}
\switchcolumn
\EN
\LessonTitle{Homily by St. Augustine, Bishop of Hippo}
\switchcolumn*

\LA
\Source{Tractatus 104 in Joannem, sub med.}
\switchcolumn
\EN
\Source{104th Tract on John}
\switchcolumn*

\LA
\noindent Póterat Dóminus noster, unigénitus et coætérnus Patri, in forma servi, et ex forma servi, si hoc opus esset, oráre siléntio: sed ita se Patri exhibére vóluit precatórem, ut meminísset nostrum se esse doctórem. Proínde eam, quam fecit oratiónem pro nobis, notam fecit et nobis: quóniam tanti magístri non solum ad ipsos sermocinátio, sed étiam pro ipsis ad Patrem orátio, discipulórum est ædificátio: et si illórum, qui hæc dicta áderant auditúri, profécto et nostra, qui fuerámus conscrípta lectúri.\par
\switchcolumn
\EN
\noindent Our Lord, the Only-begotten and co-eternal Son of the Father, was able, if need were, in and from the form of a servant, to pray in silence, but He thus manifested Himself in prayer, remembering that He is our Teacher. Thus He made known unto us the prayer which He made for us since He was so great a Master that, not only His discourse to them, but His prayer to the Father for them, is an up-building to His disciples. And if it was so for them who were there to hear, truly it is so for us also, for whose instruction it hath been written down.\par
\switchcolumn*

\LA
\begin{Resp}\Rsign Deus, cánticum novum cantábo tibi, allelúja: \Star{} In psaltério decem chordárum psallam tibi, allelúja, allelúja.\end{Resp}
\begin{Resp}\Vsign Deus meus es tu, et confitébor tibi: Deus meus es tu, et exaltábo te.\end{Resp}
\begin{Resp}\Rsign In psaltério decem chordárum psallam tibi, allelúja, allelúja.\end{Resp}
\switchcolumn
\EN
\begin{Resp}\Rsign I will sing a new song unto thee, O God, alleluia. \Star{} Upon a psaltery of ten strings will I sing praises unto thee, alleluia, alleluia.\end{Resp}
\begin{Resp}\Vsign Thou art my God, and I will praise thee Thou art my God, and I will exalt thee.\end{Resp}
\begin{Resp}\Rsign Upon a psaltery of ten strings will I sing praises unto thee, alleluia, alleluia.\end{Resp}
\switchcolumn*

\LA
\LessonHead{Lectio II}
\switchcolumn
\EN
\LessonHead{Lesson II}
\switchcolumn*

\LA
\noindent Quaprópter hoc quod ait: Pater, venit hora, clarífica Fílium tuum: osténdit, omne tempus, et quid, quando fáceret vel fíeri síneret, ab illo esse dispósitum, qui témpori súbditus non est: quóniam quæ futúra erant per síngula témpora, in Dei sapiéntia causas efficiéntes habent, in qua nulla sunt témpora. Non ergo credátur hæc hora fato urgénte venísse, sed Deo pótius ordinánte. Nec sidérea necéssitas Christi connéxuit passiónem: absit enim ut sídera mori cógerent síderum Conditórem.\par
\switchcolumn
\EN
\noindent Wherefore, by these words: “Father, the hour is come, glorify thy Son,” He showeth that all time, and all whatsoever He doth, or alloweth to be done, and the season wherein He will do or allow it, is alike ordained of Him Who is Himself not subject to time. Yea, all things which were then to come, or are yet to come now, have the reason why they should be, in the Wisdom of God, Which is Itself independent of all time. “The hour is come.” We must not believe that that hour was brought on by the march of destiny, but was by ordination of God. No stars decreed irresistibly that the time was come for Christ to suffer God forbid that the revolutions of His planets should force death on Him Who made them.\par
\switchcolumn*

\LA
\begin{Resp}\Rsign Bonum est confitéri Dómino, allelúja: \Star{} Et psállere, allelúja.\end{Resp}
\begin{Resp}\Vsign In decachórdo psaltério, cum cántico et cíthara.\end{Resp}
\begin{Resp}\Rsign Et psállere, allelúja.\end{Resp}
\begin{Resp}\Vsign Glória Patri, et Fílio, \Star{} et Spirítui Sancto.\end{Resp}
\begin{Resp}\Rsign Et psállere, allelúja.\end{Resp}
\switchcolumn
\EN
\begin{Resp}\Rsign It is a good thing to give thanks unto the Lord, alleluia. \Star{} And to sing praises, alleluia.\end{Resp}
\begin{Resp}\Vsign Upon an instrument of ten strings, upon the harp with a solemn sound.\end{Resp}
\begin{Resp}\Rsign And to sing praises, alleluia.\end{Resp}
\begin{Resp}\Vsign Glory be to the Father, and to the Son, \Star{} and to the Holy Ghost.\end{Resp}
\begin{Resp}\Rsign And to sing praises, alleluia.\end{Resp}
\switchcolumn*

\LA
\LessonHead{Lectio III}
\switchcolumn
\EN
\LessonHead{Lesson III}
\switchcolumn*

\LA
\noindent Clarificátum a Patre Fílium nonnúlli accípiunt in hoc, quod ei non pepércit, sed pro nobis ómnibus trádidit eum. Sed si passióne clarificátus dícitur, quanto magis resurrectióne? Nam in passióne magis ejus humílitas quam cláritas commendátur, Apóstolo teste, qui dicit: Humiliávit semetípsum, factus obédiens usque ad mortem, mortem autem crucis. Deínde séquitur, et de ejus clarificatióne jam dicit: Propter quod et Deus illum exaltávit, et donávit ei nomen, quod est super omne nomen: ut in nómine Jesu omne genu flectátur, cæléstium, terréstrium, et infernórum. Et omnis lingua confiteátur, quia Dóminus Jesus Christus in glória est Dei Patris. Hæc est clarificátio Dómini nostri Jesu Christi, quæ ab ejus resurrectióne sumpsit exórdium.\par
\switchcolumn
\EN
\noindent Come think that the glorification of the Son by the Father was that “He spared Him not, but delivered Him up for us all.” (Rom. viii. 32.) But if we say that He was glorified by suffering, how much more shall we say that He was glorified by rising again While He suffered, His humbleness was more manifested than His glory, as witnesseth the Apostle, where he saith “He humbled Himself, and became obedient unto death, even the death of the cross” then he addeth touching His glorification “Wherefore God also hath highly exalted Him, and given Him a Name which is above every name, that at the Name of Jesus every knee should bow, of things in heaven, and things in earth, and things under the earth and that every tongue should confess that our Lord Jesus Christ is in the glory of God the Father.” (Phil. ii. 8-11.) This is the glorification of our Lord Jesus Christ, that glorification whose first rays dawned on the Resurrection morning.\par
\switchcolumn*

\LA
\TeDeum{Te Deum}
\switchcolumn
\EN
\TeDeum{Te Deum}
\switchcolumn*

\end{paracol}

\Office{In Ascensione Domini}{The Ascension of Our Lord}{Duplex I classis cum Octava privilegiata III classis}
\addcontentsline{toc}{subsection}{In Ascensione Domini\enspace\textit{The Ascension of Our Lord}}
\begin{paracol}{2}
\LA
\LessonHead{Lectio I}
\switchcolumn
\EN
\LessonHead{Lesson I}
\switchcolumn*

\LA
\LessonTitle{Incipit liber Actuum Apostolórum}
\switchcolumn
\EN
\LessonTitle{Lesson from the Acts of the Apostles}
\switchcolumn*

\LA
\Cite{Act. 1:1-5}
\switchcolumn
\EN
\Cite{Acts 1:1-5}
\switchcolumn*

\LA
\noindent Primum quidem sermónem feci de ómnibus, o Theóphile, quæ cœpit Jesus fácere, et docére, usque in diem, qua præcípiens Apóstolis per Spíritum Sanctum, quos elégit, assúmptus est: quibus et prǽbuit seípsum vivum post passiónem suam in multis arguméntis, per dies quadragínta appárens eis, et loquens de regno Dei. Et convéscens, præcépit eis ab Jerosólymis ne discéderent, sed exspectárent promissiónem Patris, quam audístis (inquit) per os meum: quia Joánnes quidem baptizávit aqua, vos autem baptizabímini Spíritu Sancto non post multos hos dies.\par
\switchcolumn
\EN
\noindent The former treatise I made, O Theophilus, of all things which Jesus began to do and to teach, Until the day on which, giving commandments by the Holy Ghost to the apostles whom he had chosen, he was taken up. To whom also he showed himself alive after his passion, by many proofs, for forty days appearing to them, and speaking of the kingdom of God. And eating together with them, he commanded them, that they should not depart from Jerusalem, but should wait for the promise of the Father, which you have heard (saith he) by my mouth. For John indeed baptized with water, but you shall be baptized with the Holy Ghost, not many days hence.\par
\switchcolumn*

\LA
\begin{Resp}\Rsign Post passiónem suam per dies quadragínta appárens eis, et loquens de regno Dei, allelúja: \Star{} Et, vidéntibus illis, elevátus est, allelúja: et nubes suscépit eum ab óculis eórum, allelúja.\end{Resp}
\begin{Resp}\Vsign Et convéscens, præcépit eis, ab Jerosólymis ne discéderent, sed exspectárent promissiónem Patris.\end{Resp}
\begin{Resp}\Rsign Et, vidéntibus illis, elevátus est, allelúja: et nubes suscépit eum ab óculis eórum, allelúja.\end{Resp}
\switchcolumn
\EN
\begin{Resp}\Rsign Being seen of them forty days after that He had suffered, and speaking of the kingdom of God. Alleluia. \Star{} And while they beheld, He was taken up, and a cloud received Him out of their sight. Alleluia.\end{Resp}
\begin{Resp}\Vsign And, eating together with them, He commanded them that they should not depart from Jerusalem, but wait for the Promise of the Father.\end{Resp}
\begin{Resp}\Rsign And while they beheld, He was taken up, and a cloud received Him out of their sight. Alleluia.\end{Resp}
\switchcolumn*

\LA
\LessonHead{Lectio II}
\switchcolumn
\EN
\LessonHead{Lesson II}
\switchcolumn*

\LA
\Cite{Act. 1:6-9}
\switchcolumn
\EN
\Cite{Acts 1:6-9}
\switchcolumn*

\LA
\noindent Igitur qui convénerant, interrogábant eum, dicéntes: Dómine, si in témpore hoc restítues regnum Israël? Dixit autem eis: Non est vestrum nosse témpora vel moménta, quæ Pater pósuit in sua potestáte: sed accipiétis virtútem superveniéntis Spíritus Sancti in vos, et éritis mihi testes in Jerúsalem, et in omni Judǽa, et Samaría, et usque ad últimum terræ. Et cum hæc dixísset, vidéntibus illis, elevátus est: et nubes suscépit eum ab óculis eórum.\par
\switchcolumn
\EN
\noindent They therefore who were come together, asked him, saying: Lord, wilt thou at this time restore again the kingdom to Israel? But he said to them: It is not for you to know the times or moments, which the Father hath put in his own power: But you shall receive the power of the Holy Ghost coming upon you, and you shall be witnesses unto me in Jerusalem, and in all Judea, and Samaria, and even to the uttermost part of the earth. And when he had said these things, while they looked on, he was raised up: and a cloud received him out of their sight.\par
\switchcolumn*

\LA
\begin{Resp}\Rsign Omnis pulchritúdo Dómini exaltáta est super sídera: \Star{} Spécies ejus in núbibus cæli, et nomen ejus in ætérnum pérmanet, allelúja.\end{Resp}
\begin{Resp}\Vsign A summo cælo egréssio ejus, et occúrsus ejus usque ad summum ejus.\end{Resp}
\begin{Resp}\Rsign Spécies ejus in núbibus cæli, et nomen ejus in ætérnum pérmanet, allelúja.\end{Resp}
\switchcolumn
\EN
\begin{Resp}\Rsign The Lord hath set His beauty above the stars; \Star{} His loveliness is in the clouds of heaven, and His Name endureth for ever. Alleluia.\end{Resp}
\begin{Resp}\Vsign His going forth is from the end of the heaven, and His circuit unto the ends of it.\end{Resp}
\begin{Resp}\Rsign His loveliness is in the clouds of heaven, and His Name endureth for ever. Alleluia.\end{Resp}
\switchcolumn*

\LA
\LessonHead{Lectio III}
\switchcolumn
\EN
\LessonHead{Lesson III}
\switchcolumn*

\LA
\Cite{Act. 1:10-14}
\switchcolumn
\EN
\Cite{Acts 1:10-14}
\switchcolumn*

\LA
\noindent Cumque intueréntur in cælum eúntem illum, ecce duo viri astitérunt juxta illos in véstibus albis, qui et dixérunt: Viri Galilǽi, quid statis aspiciéntes in cælum? Hic Jesus qui assúmptus est a vobis in cælum, sic véniet, quemádmodum vidístis eum eúntem in cælum. Tunc revérsi sunt Jerosólymam a monte, qui vocátur Olivéti, qui est juxta Jerúsalem, sábbati habens iter. Et cum introíssent in cœnáculum, ascendérunt ubi manébant Petrus et Joánnes, Jacóbus et Andréas, Philíppus et Thomas, Bartholomǽus et Matthǽus, Jacóbus Alphǽi et Simon Zelótes, et Judas Jacóbi. Hi omnes erant perseverántes unanímiter in oratióne cum muliéribus, et María Matre Jesu, et frátribus ejus.\par
\switchcolumn
\EN
\noindent And while they were beholding him going up to heaven, behold two men stood by them in white garments. Who also said: Ye men of Galilee, why stand you looking up to heaven? This Jesus who is taken up from you into heaven, shall so come, as you have seen him going into heaven. Then they returned to Jerusalem from the mount that is called Olivet, which is nigh Jerusalem, within a sabbath day's journey. And when they were come in, they went up into an upper room, where abode Peter and John, James and Andrew, Philip and Thomas, Bartholomew and Matthew, James of Alpheus, and Simon Zelotes, and Jude the brother of James. All these were persevering with one mind in prayer with the women, and Mary the mother of Jesus, and with his brethren.\par
\switchcolumn*

\LA
\begin{Resp}\Rsign Exaltáre, Dómine, allelúja, \Star{} In virtúte tua, allelúja.\end{Resp}
\begin{Resp}\Vsign Eleváta est, magnificéntia tua super cælos, Deus.\end{Resp}
\begin{Resp}\Rsign In virtúte tua, allelúja.\end{Resp}
\begin{Resp}\Vsign Glória Patri, et Fílio, \Star{} et Spirítui Sancto.\end{Resp}
\begin{Resp}\Rsign In virtúte tua, allelúja.\end{Resp}
\switchcolumn
\EN
\begin{Resp}\Rsign Be Thou exalted, O Lord. Alleluia. \Star{} In thine Own strength. Alleluia.\end{Resp}
\begin{Resp}\Vsign O God, Thou hast set thy glory above the heavens.\end{Resp}
\begin{Resp}\Rsign In thine Own strength. Alleluia.\end{Resp}
\begin{Resp}\Vsign Glory be to the Father, and to the Son, \Star{} and to the Holy Ghost.\end{Resp}
\begin{Resp}\Rsign In thine Own strength. Alleluia.\end{Resp}
\switchcolumn*

\LA
\LessonHead{Lectio IV}
\switchcolumn
\EN
\LessonHead{Lesson IV}
\switchcolumn*

\LA
\LessonTitle{Sermo sancti Leónis Papæ}
\switchcolumn
\EN
\LessonTitle{From the Sermons of Pope St. Leo the Great}
\switchcolumn*

\LA
\Source{Sermone 1 de Ascensione Domini}
\switchcolumn
\EN
\Source{1st on the Lords Ascension}
\switchcolumn*

\LA
\noindent Post beátam et gloriósam resurrectiónem Dómini nostri Jesu Christi, qua verum Dei templum, Judáica impietáte resolútum, divína in tríduo poténtia suscitávit, quadragenárius hódie, dilectíssimi, sanctórum diérum explétus est númerus, sacratíssima ordinatióne dispósitus, et ad utilitátem nostræ eruditiónis impénsus: ut, dum a Dómino in hoc spátio mora præséntiæ corporális exténditur, fides resurrectiónis documéntis necessáriis munirétur. Mors enim Christi multum discipulórum corda turbáverat: et de supplício crucis, de emissióne spíritus, de exanimáti córporis sepultúra gravátis mœstitúdine méntibus, quidam diffidéntiæ torpor obrépserat.\par
\switchcolumn
\EN
\noindent After the blessed and glorious Resurrection of our Lord Jesus Christ, wherein the Divine Power raised up in three days the true Temple of God Which the iniquity of the Jews had destroyed (John ii. 19,) God was pleased to ordain, by His Most Sacred Will, and in His Providence for our instruction and the profit of our souls, a season of forty days which season, dearly beloved brethren, doth end on this day. During that season the bodily Presence of the Lord still lingered on earth, that the reality of the fact of His having risen again from the dead might be armed with all needful proofs. The death of Christ had troubled the hearts of many of His disciples their thoughts were sad when they remembered His agony upon the Cross, His giving up of the Ghost, and the laying in the grave of His lifeless Body, and a sort of hesitation had begun to weigh on them.\par
\switchcolumn*

\LA
\begin{Resp}\Rsign Tempus est, ut revértar ad eum, qui me misit, dicit Dóminus: nolíte contristári, nec turbétur cor vestrum: \Star{} Rogo pro vobis Patrem, ut ipse vos custódiat, allelúja, allelúja.\end{Resp}
\begin{Resp}\Vsign Nisi ego abíero, Paráclitus non véniet: cum assúmptus fúero, mittam vobis eum.\end{Resp}
\begin{Resp}\Rsign Rogo pro vobis Patrem, ut ipse vos custódiat, allelúja, allelúja.\end{Resp}
\switchcolumn
\EN
\begin{Resp}\Rsign My time is come that I should return unto Him That sent Me, saith the Lord. Be not sorrowful, neither let your heart be troubled. \Star{} I pray the Father for you, that He may keep you. Alleluia, Alleluia.\end{Resp}
\begin{Resp}\Vsign If I go not away, the Comforter will not come unto you; where I am ascended, I will send Him unto you.\end{Resp}
\begin{Resp}\Rsign I pray the Father for you, that He may keep you. Alleluia, Alleluia.\end{Resp}
\switchcolumn*

\LA
\LessonHead{Lectio V}
\switchcolumn
\EN
\LessonHead{Lesson V}
\switchcolumn*

\LA
\noindent Unde beatíssimi Apóstoli, omnésque discípuli, qui et de éxitu crucis fúerant trépidi, et de fide resurrectiónis ambígui, ita sunt veritáte perspícua roboráti, ut, Dómino in cælórum eúnte sublímia, non solum nulla afficeréntur tristítia, sed étiam magno gáudio repleréntur. Et revéra magna et ineffábilis erat causa gaudéndi, cum in conspéctu sanctæ multitúdinis super ómnium creaturárum cæléstium dignitátem humáni géneris natúra conscénderet, supergressúra Angélicos órdines, et ultra Archangelórum altitúdines elevánda: nec ullis sublimitátibus modum suæ provectiónis habitúra, nisi ætérni Patris recépta conséssu, illíus glóriæ sociarétur in throno, cujus natúræ copulabátur in Fílio.\par
\switchcolumn
\EN
\noindent Hence the most blessed Apostles and all the disciples, who had been fearful at the finishing on the Cross, and doubtful of the trustworthiness of the rising again, were so strengthened by the clear demonstration of the fact, that, when they saw the Lord going up into the height of heaven, they sorrowed not, nay they were even filled with great joy And, in all verity, it was a great an unspeakable cause for joy to see the Manhood, in the presence of that the multitude of believers, exalted above all creatures even heavenly, rising above the ranks of the angelic armies and speeding Its glorious way where the most noble of the Archangels lie far behind, to rest no lower than that place where high above all principality and power, It taketh Its seat at the right hand of the Eternal Father, Sharer of His throne, and Partaker of His glory, and still of the very man's nature which the Son hath taken upon Him.\par
\switchcolumn*

\LA
\begin{Resp}\Rsign Non turbétur cor vestrum: ego vado ad Patrem; et cum assúmptus fúero a vobis, mittam vobis, allelúja, \Star{} Spíritum veritátis, et gaudébit cor vestrum, allelúja.\end{Resp}
\begin{Resp}\Vsign Ego rogábo Patrem, et álium Paráclitum dabit vobis.\end{Resp}
\begin{Resp}\Rsign Spíritum veritátis, et gaudébit cor vestrum, allelúja.\end{Resp}
\switchcolumn
\EN
\begin{Resp}\Rsign Let not your heart be troubled; I go unto the Father, and when I am taken from you, I will send unto you, alleluia, \Star{} The Spirit of truth, and your heart shall rejoice. Alleluia.\end{Resp}
\begin{Resp}\Vsign I will pray the Father, and He shall give you another Comforter.\end{Resp}
\begin{Resp}\Rsign The Spirit of truth, and your heart shall rejoice. Alleluia.\end{Resp}
\switchcolumn*

\LA
\LessonHead{Lectio VI}
\switchcolumn
\EN
\LessonHead{Lesson VI}
\switchcolumn*

\LA
\noindent Quia ígitur Christi ascénsio, nostra provéctio est; et quo præcéssit glória cápitis, eo spes vocátur et córporis: dignis, dilectíssimi, exsultémus gáudiis, et pia gratiárum actióne lætémur. Hódie enim non solum paradísi possessóres firmáti sumus, sed étiam cælórum in Christo supérna penetrávimus: amplióra adépti per ineffábilem Christi grátiam, quam per diáboli amiserámus invídiam. Nam quos viruléntus inimícus primi habitáculi felicitáte dejécit, eos sibi concorporátos Dei Fílius ad déxteram Patris collocávit: cum quo vivit et regnat in unitáte Spíritus Sancti Deus, per ómnia sǽcula sæculórum. Amen.\par
\switchcolumn
\EN
\noindent Therefore, dearly beloved brethren, let us also rejoice with worthy joy, for the Ascension of Christ is exaltation for us, and whither the glory of the Head of the Church is passed in, thither is the hope of the body of the Church called on to follow. Let us rejoice with exceeding great joy, and give God glad thanks. This day is not only the possession of Paradise made sure unto us, but in the Person of our Head we are actually begun to enter into the heavenly mansions above. Through the unspeakable goodness of Christ we have gained more than ever we lost by the envy of the devil. We, whom our venomous enemy thrust from our first happy home, we, being made of one body with the Son of God, have by Him been given a place at the right hand of the Father with Whom He liveth and reigneth, in the unity of the Holy Ghost, God, world without end. Amen.\par
\switchcolumn*

\LA
\begin{Resp}\Rsign Ascéndens Christus in altum, captívam duxit captivitátem, \Star{} Dedit dona homínibus, allelúja, allelúja, allelúja.\end{Resp}
\begin{Resp}\Vsign Ascéndit Deus in jubilatióne, et Dóminus in voce tubæ.\end{Resp}
\begin{Resp}\Rsign Dedit dona homínibus, allelúja, allelúja, allelúja.\end{Resp}
\begin{Resp}\Vsign Glória Patri, et Fílio, \Star{} et Spirítui Sancto.\end{Resp}
\begin{Resp}\Rsign Dedit dona homínibus, allelúja, allelúja, allelúja.\end{Resp}
\switchcolumn
\EN
\begin{Resp}\Rsign When Christ ascended up on high, He led captivity captive, \Star{} He gave gifts unto men. Alleluia, Alleluia, Alleluia.\end{Resp}
\begin{Resp}\Vsign God is gone up with a shout, and the Lord with the sound of a trumpet.\end{Resp}
\begin{Resp}\Rsign He gave gifts unto men. Alleluia, Alleluia, Alleluia.\end{Resp}
\begin{Resp}\Vsign Glory be to the Father, and to the Son, \Star{} and to the Holy Ghost.\end{Resp}
\begin{Resp}\Rsign He gave gifts unto men. Alleluia, Alleluia, Alleluia.\end{Resp}
\switchcolumn*

\LA
\LessonHead{Lectio VII}
\switchcolumn
\EN
\LessonHead{Lesson VII}
\switchcolumn*

\LA
\LessonTitle{Léctio sancti Evangélii secúndum Marcum}
\switchcolumn
\EN
\LessonTitle{From the Holy Gospel according to Mark}
\switchcolumn*

\LA
\Cite{Marc 16:14-20}
\switchcolumn
\EN
\Cite{Mark 16:14-20}
\switchcolumn*

\LA
\noindent In illo témpore: Recumbéntibus úndecim discípulis, appáruit illis Jesus: et exprobrávit incredulitátem eórum et durítiam cordis, quia iis, qui víderant eum resurrexísse, non credidérunt. Et réliqua.\par
\switchcolumn
\EN
\noindent At that time, Jesus appeared unto the eleven disciples as they sat at meat, and upbraided them with their unbelief and hardness of heart, because they believed not them which had seen Him after He was risen. And so on.\par
\switchcolumn*

\LA
\LessonTitle{Homilía sancti Gregórii Papæ}
\switchcolumn
\EN
\LessonTitle{Homily by Pope St. Gregory the Great}
\switchcolumn*

\LA
\Source{Homilia 29 in Evangelia}
\switchcolumn
\EN
\Source{29th on the Gospels}
\switchcolumn*

\LA
\noindent Quod resurrectiónem Domínicam discípuli tarde credidérunt, non tam illórum infírmitas, quam nostra, ut ita dicam, futúra fírmitas fuit. Ipsa namque resurréctio illis dubitántibus per multa arguménta monstráta est: quæ dum nos legéntes agnóscimus, quid áliud quam de illórum dubitatióne solidámur? Minus enim mihi María Magdaléne prǽstitit, quæ cítius crédidit, quam Thomas, qui diu dubitávit. Ille étenim dubitándo, vúlnerum cicatríces tétigit, et de nostro péctore dubitatiónis vulnus amputávit.\par
\switchcolumn
\EN
\noindent I may be allowed to say that the disciples' slowness to believe that the Lord had indeed risen from the dead, was not so much their weakness as our strength. In consequence of their doubts, the fact of the Resurrection was demonstrated by many infallible proofs. These proofs we read and acknowledge. What then assureth our faith, if not their doubt? For my part, I put my trust in Thomas, who doubted long, much more than in Mary Magdalene, who believed at once. Through his doubting, he came actually to handle the holes of the Wounds, and thereby closed up any wound of doubt in our hearts.\par
\switchcolumn*

\LA
\begin{Resp}\Rsign Ego rogábo Patrem, et álium Paráclitum dabit vobis, \Star{} Ut máneat vobíscum in ætérnum, Spíritum veritátis, allelúja.\end{Resp}
\begin{Resp}\Vsign Si enim non abíero, Paráclitus non véniet ad vos: si autem abíero, mittam eum ad vos.\end{Resp}
\begin{Resp}\Rsign Ut máneat vobíscum in ætérnum, Spíritum veritátis, allelúja.\end{Resp}
\switchcolumn
\EN
\begin{Resp}\Rsign I will pray the Father, and He shall give you another Comforter, \Star{} That He may abide with you for ever, even the Spirit of truth. Alleluia.\end{Resp}
\begin{Resp}\Vsign For if I go not away, the Comforter will not come unto you; but if I depart, I will send Him unto you.\end{Resp}
\begin{Resp}\Rsign That He may abide with you for ever, even the Spirit of truth. Alleluia.\end{Resp}
\switchcolumn*

\LA
\LessonHead{Lectio VIII}
\switchcolumn
\EN
\LessonHead{Lesson VIII}
\switchcolumn*

\LA
\noindent Ad insinuándam quoque veritátem Domínicæ resurrectiónis, notándum nobis est quid Lucas réferat, dicens: Convéscens, præcépit eis ab Jerosólymis ne discéderent. Et post pauca: Vidéntibus illis, elevátus est, et nubes suscépit eum ab óculis eórum. Notáte verba, signáte mystéria. Convéscens elevátus est. Comédit, et ascéndit: ut vidélicet per efféctum comestiónis, véritas patésceret carnis. Marcus vero, priúsquam cælum Dóminus ascéndat, eum de cordis atque infidelitátis durítia increpásse discípulos mémorat. Qua in re quid considerándum est, nisi quod idcírco Dóminus tunc discípulos increpávit, cum eos corporáliter relíquit, ut verba, quæ recédens díceret, in corde audiéntium árctius impréssa remanérent?\par
\switchcolumn
\EN
\noindent Now confirm to our minds the trustworthiness of the fact that our Lord did indeed rise again from the dead, it is well for us to remark one of the statements of Luke (Acts i. 4.) “Eating together with them, He commanded them that they should not (1 John xiv. 16, 17 xvi. 7.) depart from Jerusalem and a little afterward: “While they beheld, He was taken up, and a cloud received Him out of their sight.” Consider these words, note well these mysteries. After “eating together with them He was taken up.” He ate and ascended: that the fact of His eating might show the reality of the Body in Which He went up. But Mark telleth us that before the Lord ascended into heaven, He upbraided His disciples; with their unbelief and hardness of heart. From this I know not why we should gather, but that the Lord then upbraided His disciples, for whom He was about to be parted in the body, to the end that the words which He spoke unto them as He left them might be the deeper imprinted on their hearts.\par
\switchcolumn*

\LA
\begin{Resp}\Rsign Ponis nubem ascénsum tuum, Dómine: \Star{} Qui ámbulas super pennas ventórum, allelúja.\end{Resp}
\begin{Resp}\Vsign Confessiónem et decórem induísti, amíctus lumen sicut vestiméntum.\end{Resp}
\begin{Resp}\Rsign Qui ámbulas super pennas ventórum, allelúja.\end{Resp}
\begin{Resp}\Vsign Glória Patri, et Fílio, \Star{} et Spirítui Sancto.\end{Resp}
\begin{Resp}\Rsign Qui ámbulas super pennas ventórum, allelúja.\end{Resp}
\switchcolumn
\EN
\begin{Resp}\Rsign Thou makest the clouds the chariot, O Lord, \Star{} Thou walkest upon the wings of the wind. Alleluia.\end{Resp}
\begin{Resp}\Vsign Thou art clothed with honour and majesty, covering thyself with light as with a garment;\end{Resp}
\begin{Resp}\Rsign Thou walkest upon the wings of the wind. Alleluia.\end{Resp}
\begin{Resp}\Vsign Glory be to the Father, and to the Son, \Star{} and to the Holy Ghost.\end{Resp}
\begin{Resp}\Rsign Thou walkest upon the wings of the wind. Alleluia.\end{Resp}
\switchcolumn*

\LA
\LessonHead{Lectio IX}
\switchcolumn
\EN
\LessonHead{Lesson IX}
\switchcolumn*

\LA
\noindent Increpáta ígitur eórum durítia, quid admonéndo dicat, audiámus: Eúntes in mundum univérsum, prædicáte Evangélium omni creatúræ. Numquid, fratres mei, sanctum Evangélium vel insensátis rebus, vel brutis animálibus fúerat prædicándum, ut de eo discípulis dicátur: Prædicáte omni creatúræ? Sed omnis creatúræ nómine signátur homo. Omnis autem creatúræ áliquid habet homo. Habet namque commúne esse cum lapídibus, vívere cum arbóribus, sentíre cum animálibus, intellégere cum Angelis. Si ergo commúne habet áliquid cum omni creatúra homo, juxta áliquid omnis creatúra est homo. Omni ergo creatúræ prædicátur Evangélium, cum soli hómini prædicátur.\par
\switchcolumn
\EN
\noindent When then, He had rebuked the hardness of their heart, what command did He give them? Let us hear. “Go ye into all the world and preach the Gospel to every creature.” Was the Holy Gospel, then my brethren, to be preached to thing insensate, or to brute beasts, that the Lord said to His disciples: “Preach the Gospel to every creature”? Nay, but by the words “every creature” we must understand man, in whom are combined qualities of all creatures. Being he hath in common with stones, life in common with trees, feeling in common with beasts, understanding in common with angels. If, then, man hath something in common with every creature, man is to a certain extent every creature. The Gospel, then, if it be preached to man only, is preached to every creature.\par
\switchcolumn*

\LA
\TeDeum{Te Deum}
\switchcolumn
\EN
\TeDeum{Te Deum}
\switchcolumn*

\end{paracol}

\Office{Feria VI infra Octavam Ascensionis}{Friday within the Octave of the Ascension}{Semiduplex}
\addcontentsline{toc}{subsection}{Feria VI infra Octavam Ascensionis\enspace\textit{Friday within the Octave of the Ascension}}
\begin{paracol}{2}
\LA
\LessonHead{Lectio I}
\switchcolumn
\EN
\LessonHead{Lesson I}
\switchcolumn*

\LA
\LessonTitle{Incipit Epístola secúnda beáti Petri Apóstoli}
\switchcolumn
\EN
\LessonTitle{Lesson from the second letter of St. Peter the Apostle}
\switchcolumn*

\LA
\Cite{2 Pet 1:1-4}
\switchcolumn
\EN
\Cite{2 Pet 1:1-4}
\switchcolumn*

\LA
\noindent Simon Petrus, servus et apóstolus Jesu Christi, iis qui coæquálem nobíscum sortíti sunt fidem in justítia Dei nostri, et Salvatóris Jesu Christi. Grátia vobis, et pax adimpleátur in cognitióne Dei, et Christi Jesu Dómini nostri: Quómodo ómnia nobis divínæ virtútis suæ, quæ ad vitam et pietátem donáta sunt, per cognitiónem ejus, qui vocávit nos própria glória, et virtúte, per quem máxima, et pretiósa nobis promíssa donávit: ut per hæc efficiámini divínæ consórtes natúræ: fugiéntes ejus, quæ in mundo est, concupiscéntiæ corruptiónem.\par
\switchcolumn
\EN
\noindent Simon Peter, servant and apostle of Jesus Christ, to them that have obtained equal faith with us in the justice of our God and Saviour Jesus Christ. Grace to you and peace be accomplished in the knowledge of God and of Christ Jesus our Lord: As all things of his divine power which appertain to life and godliness, are given us, through the knowledge of him who hath called us by his own proper glory and virtue. By whom he hath given us most great and precious promises: that by these you may be made partakers of the divine nature: fleeing the corruption of that concupiscence which is in the world.\par
\switchcolumn*

\LA
\begin{Resp}\Rsign Post passiónem suam per dies quadragínta appárens eis, et loquens de regno Dei, allelúja: \Star{} Et, vidéntibus illis, elevátus est, allelúja: et nubes suscépit eum ab óculis eórum, allelúja.\end{Resp}
\begin{Resp}\Vsign Et convéscens, præcépit eis, ab Jerosólymis ne discéderent, sed exspectárent promissiónem Patris.\end{Resp}
\begin{Resp}\Rsign Et, vidéntibus illis, elevátus est, allelúja: et nubes suscépit eum ab óculis eórum, allelúja.\end{Resp}
\switchcolumn
\EN
\begin{Resp}\Rsign Being seen of them forty days after that He had suffered, and speaking of the kingdom of God. Alleluia. \Star{} And while they beheld, He was taken up, and a cloud received Him out of their sight. Alleluia.\end{Resp}
\begin{Resp}\Vsign And, eating together with them, He commanded them that they should not depart from Jerusalem, but wait for the Promise of the Father.\end{Resp}
\begin{Resp}\Rsign And while they beheld, He was taken up, and a cloud received Him out of their sight. Alleluia.\end{Resp}
\switchcolumn*

\LA
\LessonHead{Lectio II}
\switchcolumn
\EN
\LessonHead{Lesson II}
\switchcolumn*

\LA
\Cite{2 Pet 1:5-9}
\switchcolumn
\EN
\Cite{2 Pet 1:5-9}
\switchcolumn*

\LA
\noindent Vos autem curam omnem subinferéntes, ministráte in fide vestra virtútem, in virtúte autem sciéntiam, in sciéntia autem abstinéntiam, in abstinéntia autem patiéntiam, in patiéntia autem pietátem, in pietáte autem amórem fraternitátis, in amóre autem fraternitátis caritátem. Hæc enim si vobíscum adsint, et súperent, non vácuos nec sine fructu vos constítuent in Dómini nostri Jesu Christi cognitióne. Cui enim non præsto sunt hæc, cæcus est, et manu tentans, obliviónem accípiens purgatiónis véterum suórum delictórum.\par
\switchcolumn
\EN
\noindent And you, employing all care, minister in your faith, virtue; and in virtue, knowledge; And in knowledge, abstinence; and in abstinence, patience; and in patience, godliness; And in godliness, love of brotherhood; and in love of brotherhood, charity. For if these things be with you and abound, they will make you to be neither empty nor unfruitful in the knowledge of our Lord Jesus Christ. For he that hath not these things with him, is blind, and groping, having forgotten that he was purged from his old sins.\par
\switchcolumn*

\LA
\begin{Resp}\Rsign Omnis pulchritúdo Dómini exaltáta est super sídera: \Star{} Spécies ejus in núbibus cæli, et nomen ejus in ætérnum pérmanet, allelúja.\end{Resp}
\begin{Resp}\Vsign A summo cælo egréssio ejus, et occúrsus ejus usque ad summum ejus.\end{Resp}
\begin{Resp}\Rsign Spécies ejus in núbibus cæli, et nomen ejus in ætérnum pérmanet, allelúja.\end{Resp}
\switchcolumn
\EN
\begin{Resp}\Rsign The Lord hath set His beauty above the stars; \Star{} His loveliness is in the clouds of heaven, and His Name endureth for ever. Alleluia.\end{Resp}
\begin{Resp}\Vsign His going forth is from the end of the heaven, and His circuit unto the ends of it.\end{Resp}
\begin{Resp}\Rsign His loveliness is in the clouds of heaven, and His Name endureth for ever. Alleluia.\end{Resp}
\switchcolumn*

\LA
\LessonHead{Lectio III}
\switchcolumn
\EN
\LessonHead{Lesson III}
\switchcolumn*

\LA
\Cite{2 Pet 1:10-15}
\switchcolumn
\EN
\Cite{2 Pet 1:10-15}
\switchcolumn*

\LA
\noindent Quaprópter fratres, magis satágite ut per bona ópera certam vestram vocatiónem, et electiónem faciátis: hæc enim faciéntes, non peccábitis aliquándo. Sic enim abundánter ministrábitur vobis intróitus in ætérnum regnum Dómini nostri et Salvatóris Jesu Christi. Propter quod incípiam vos semper commonére de his: et quidem sciéntes et confirmátos vos in præsénti veritáte. Justum autem árbitror quámdiu sum in hoc tabernáculo, suscitáre vos in commonitióne: certus quod velox est deposítio tabernáculi mei secúndum quod et Dóminus noster Jesus Christus significávit mihi. Dabo autem óperam et frequénter habére vos post óbitum meum, ut horum memóriam faciátis.\par
\switchcolumn
\EN
\noindent Wherefore, brethren, labour the more, that by good works you may make sure your calling and election. For doing these things, you shall not sin at any time. For so an entrance shall be ministered to you abundantly into the everlasting kingdom of our Lord and Saviour Jesus Christ. For which cause I will begin to put you always in remembrance of these things: though indeed you know them, and are confirmed in the present truth. But I think it meet as long as I am in this tabernacle, to stir you up by putting you in remembrance. Being assured that the laying away of this my tabernacle is at hand, according as our Lord Jesus Christ also hath signified to me. And I will endeavour, that you frequently have after my decease, whereby you may keep a memory of these things.\par
\switchcolumn*

\LA
\begin{Resp}\Rsign Exaltáre, Dómine, allelúja, \Star{} In virtúte tua, allelúja.\end{Resp}
\begin{Resp}\Vsign Eleváta est, magnificéntia tua super cælos, Deus.\end{Resp}
\begin{Resp}\Rsign In virtúte tua, allelúja.\end{Resp}
\begin{Resp}\Vsign Glória Patri, et Fílio, \Star{} et Spirítui Sancto.\end{Resp}
\begin{Resp}\Rsign In virtúte tua, allelúja.\end{Resp}
\switchcolumn
\EN
\begin{Resp}\Rsign Be Thou exalted, O Lord. Alleluia. \Star{} In thine Own strength. Alleluia.\end{Resp}
\begin{Resp}\Vsign O God, Thou hast set thy glory above the heavens.\end{Resp}
\begin{Resp}\Rsign In thine Own strength. Alleluia.\end{Resp}
\begin{Resp}\Vsign Glory be to the Father, and to the Son, \Star{} and to the Holy Ghost.\end{Resp}
\begin{Resp}\Rsign In thine Own strength. Alleluia.\end{Resp}
\switchcolumn*

\LA
\LessonHead{Lectio IV}
\switchcolumn
\EN
\LessonHead{Lesson IV}
\switchcolumn*

\LA
\LessonTitle{Sermo sancti Leónis Papæ}
\switchcolumn
\EN
\LessonTitle{From the Sermons of Pope St. Leo the Great.}
\switchcolumn*

\LA
\Source{Sermo 2 de Ascensione Domini}
\switchcolumn
\EN
\Source{2nd for the Lord's Ascension.}
\switchcolumn*

\LA
\noindent Sacraméntum, dilectíssimi, salútis nostræ, quam prétio sánguinis sui universitátis Cónditor æstimávit, a die corporális ortus usque ad éxitum passiónis, per dispensatiónem humilitátis implétum est. Et licet multa, étiam in forma servi, divinitátis signa radiáverint; próprie tamen illíus témporis áctio ad demonstrándam suscépti hóminis pertínuit veritátem. Post passiónem vero ruptis mortis vínculis, quæ vim suam in eum, qui peccáti erat néscius, incedéndo perdíderat; infírmitas in virtútem, mortálitas in immortalitátem, contumélia transívit in glóriam: quam Dóminus Jesus Christus in multis manifestísque documéntis, multórum declarávit aspéctibus, donec triúmphum victóriæ, quem reportárat a mórtuis, inférret et cælis.\par
\switchcolumn
\EN
\noindent Dearly beloved brethren, that mysterious thing, our salvation, which the Maker of the universe thought worth purchasing with His Own Precious Blood, was aimed at by Him, in the dispensation of His humility, from the hour wherein He was born as touching the flesh, till the moment when, at the end of the Passion, He cried on the Cross: “It is finished.” Although from under the form of a servant many marks of His Godhead shone forth, yet, as a whole, the work of those three-and-thirty years was to manifest the verity of the Manhood Which the Son of God had taken into Himself. But when the suffering was all over, and the bands of death were broken, (that death which had lost all his power by seeking to bind Him Who knew no sin,) then was weakness changed into strength, mortality into immortality, insult into that glory which the Lord Jesus Christ, on so many occasions, made manifest by so many and infallible proofs, until the day came when that triumphant procession of victory, which He had led from the realms of shattered death, followed Him with unimaginable pomp into the heavens.\par
\switchcolumn*

\LA
\begin{Resp}\Rsign Tempus est, ut revértar ad eum, qui me misit, dicit Dóminus: nolíte contristári, nec turbétur cor vestrum: \Star{} Rogo pro vobis Patrem, ut ipse vos custódiat, allelúja, allelúja.\end{Resp}
\begin{Resp}\Vsign Nisi ego abíero, Paráclitus non véniet: cum assúmptus fúero, mittam vobis eum.\end{Resp}
\begin{Resp}\Rsign Rogo pro vobis Patrem, ut ipse vos custódiat, allelúja, allelúja.\end{Resp}
\switchcolumn
\EN
\begin{Resp}\Rsign My time is come that I should return unto Him That sent Me, saith the Lord. Be not sorrowful, neither let your heart be troubled. \Star{} I pray the Father for you, that He may keep you. Alleluia, Alleluia.\end{Resp}
\begin{Resp}\Vsign If I go not away, the Comforter will not come unto you; where I am ascended, I will send Him unto you.\end{Resp}
\begin{Resp}\Rsign I pray the Father for you, that He may keep you. Alleluia, Alleluia.\end{Resp}
\switchcolumn*

\LA
\LessonHead{Lectio V}
\switchcolumn
\EN
\LessonHead{Lesson V}
\switchcolumn*

\LA
\noindent Sicut ergo in solemnitáte pascháli resurréctio Dómini fuit nobis causa lætándi: ita ascénsio ejus in cælos præséntium nobis est matéria gaudiórum, recoléntibus illum diem, et rite venerántibus, quo natúra nostræ humilitátis in Christo super omnem cæli milítiam, super omnes órdines Angelórum, et ultra ómnium altitúdinem potestátum, ad Dei Patris est provécta conséssum. Quo órdine óperum divinórum nos fundáti, nos ædificáti sumus: ut mirabílior fíeret grátia Dei, cum remótis a conspéctu hóminum, quæ mérito reveréntiam sui sentiebántur indícere, fides non defíceret, spes non fluctuáret, cáritas non tepéret.\par
\switchcolumn
\EN
\noindent On the solemn Feast of the Passover the cause of our joy was that Christ was risen again. This day we rejoice because that He is ascended up into heaven. We call to mind and justly celebrate that day whereon our lowly nature was, in the Person of Christ, borne up high above all the heavenly armies, above all the circles of Angels, beyond the heights of all the Powers, even to where Christ is sitting on the right hand of the Father. Our foundations are laid, and our house is built upon this succession of the works of God and His grace is made more wonderful by this, that, though the visible Object of worship is removed from among men, the faith of the Church doth not grow weak, nor her hope wavering, nor her love cold.\par
\switchcolumn*

\LA
\begin{Resp}\Rsign Non turbétur cor vestrum: ego vado ad Patrem; et cum assúmptus fúero a vobis, mittam vobis, allelúja, \Star{} Spíritum veritátis, et gaudébit cor vestrum, allelúja.\end{Resp}
\begin{Resp}\Vsign Ego rogábo Patrem, et álium Paráclitum dabit vobis.\end{Resp}
\begin{Resp}\Rsign Spíritum veritátis, et gaudébit cor vestrum, allelúja.\end{Resp}
\switchcolumn
\EN
\begin{Resp}\Rsign Let not your heart be troubled; I go unto the Father, and when I am taken from you, I will send unto you, alleluia, \Star{} The Spirit of truth, and your heart shall rejoice. Alleluia.\end{Resp}
\begin{Resp}\Vsign I will pray the Father, and He shall give you another Comforter.\end{Resp}
\begin{Resp}\Rsign The Spirit of truth, and your heart shall rejoice. Alleluia.\end{Resp}
\switchcolumn*

\LA
\LessonHead{Lectio VI}
\switchcolumn
\EN
\LessonHead{Lesson VI}
\switchcolumn*

\LA
\noindent Magnárum hic vigor est méntium, et valde fidélium lumen est animárum, incunctánter crédere, quæ corpóreo non vidéntur intúitu, et ibi fígere desidérium, quo néqueas inférre conspéctum. Hæc autem píetas unde in nostris córdibus nascerétur, aut quómodo quisquam justificarétur per fidem, si in iis tantum salus nostra consísteret, quæ obtútibus subjacérent? Unde et illi viro, qui de resurrectióne Christi videbátur ambígere, nisi in ipsíus carne vestígia passiónis et visu explorásset et tactu: Quia vidísti me, inquit Dóminus credidísti: beáti qui non vidérunt, et credidérunt.\par
\switchcolumn
\EN
\noindent It is the back-bone of a strong mind and the eye of a trusty soul, to believe unhesitatingly that which is not seen with the bodily eyes, and to centre all love where there can be no experimental knowledge. This it is which is the only thing we can have of godliness for how could a man be justified through faith, if the saving objects were objects of sight There was a man who would not believe in the Resurrection of Christ until he had examined by sight, and touched the marks of the Passion in the Divine Body, and the Lord said to him Because thou hast seen Me, thou hast believed blessed are they that have not seen, and yet have believed.” (John xx. 29.)\par
\switchcolumn*

\LA
\begin{Resp}\Rsign Ascéndens Christus in altum, captívam duxit captivitátem, \Star{} Dedit dona homínibus, allelúja, allelúja, allelúja.\end{Resp}
\begin{Resp}\Vsign Ascéndit Deus in jubilatióne, et Dóminus in voce tubæ.\end{Resp}
\begin{Resp}\Rsign Dedit dona homínibus, allelúja, allelúja, allelúja.\end{Resp}
\begin{Resp}\Vsign Glória Patri, et Fílio, \Star{} et Spirítui Sancto.\end{Resp}
\begin{Resp}\Rsign Dedit dona homínibus, allelúja, allelúja, allelúja.\end{Resp}
\switchcolumn
\EN
\begin{Resp}\Rsign When Christ ascended up on high, He led captivity captive, \Star{} He gave gifts unto men. Alleluia, Alleluia, Alleluia.\end{Resp}
\begin{Resp}\Vsign God is gone up with a shout, and the Lord with the sound of a trumpet.\end{Resp}
\begin{Resp}\Rsign He gave gifts unto men. Alleluia, Alleluia, Alleluia.\end{Resp}
\begin{Resp}\Vsign Glory be to the Father, and to the Son, \Star{} and to the Holy Ghost.\end{Resp}
\begin{Resp}\Rsign He gave gifts unto men. Alleluia, Alleluia, Alleluia.\end{Resp}
\switchcolumn*

\LA
\LessonHead{Lectio VII}
\switchcolumn
\EN
\LessonHead{Lesson VII}
\switchcolumn*

\LA
\LessonTitle{Léctio sancti Evangélii secúndum Marcum}
\switchcolumn
\EN
\LessonTitle{From the Holy Gospel according to Mark}
\switchcolumn*

\LA
\Cite{Marc 16:14-20}
\switchcolumn
\EN
\Cite{Mark 16:14-20}
\switchcolumn*

\LA
\noindent In illo témpore: Recumbéntibus úndecim discípulis, appáruit illis Jesus: et exprobrávit incredulitátem eórum et durítiam cordis, quia iis, qui víderant eum resurrexísse, non credidérunt. Et réliqua.\par
\switchcolumn
\EN
\noindent At that time: Jesus appeared unto the eleven disciples as they sat at meat, and upbraided them with their unbelief and hardness of heart because they believed not them which had seen Him after He was risen. And so on.\par
\switchcolumn*

\LA
\LessonTitle{Homilía sancti Gregórii Papæ}
\switchcolumn
\EN
\LessonTitle{Homily by Pope St. Gregory the Great.}
\switchcolumn*

\LA
\Source{Ex eadem Homilía 29}
\switchcolumn
\EN
\Source{Same as before.}
\switchcolumn*

\LA
\noindent Qui credíderit, et baptizátus fúerit, salvus erit: qui vero non credíderit, condemnábitur. Fortásse unusquísque apud semetípsum dicat: Ego jam crédidi, salvus ero. Verum dicit, si fidem opéribus tenet. Vera étenim fides est, quæ in hoc, quod verbis dicit, móribus non contradícit. Hinc est enim quod de quibúsdam falsis fidélibus Paulus dicit: Qui confiténtur se nosse Deum, factis autem negant. Hinc Joánnes ait: Qui dicit se nosse Deum, et mandáta ejus non custódit, mendax est.\par
\switchcolumn
\EN
\noindent “He that believeth, and is baptized, shall be saved but he that believeth not shall be damned.” Perchance some man will say within himself: “I have already believed, and therefore I shall be saved.” Thou hast well said, if thou showest thy faith by thy works. He only hath a true faith whose life doth not give the lie to his confession. Hence it is that Paul saith, touching some who were falsely faithful: “They profess that they know God but in works they deny Him.” (Tit. i. 16.) And John likewise saith: “He that saith, I know Him and keepeth not His commandments, is a liar.” (I. ii. 4.)\par
\switchcolumn*

\LA
\begin{Resp}\Rsign Ego rogábo Patrem, et álium Paráclitum dabit vobis, \Star{} Ut máneat vobíscum in ætérnum, Spíritum veritátis, allelúja.\end{Resp}
\begin{Resp}\Vsign Si enim non abíero, Paráclitus non véniet ad vos: si autem abíero, mittam eum ad vos.\end{Resp}
\begin{Resp}\Rsign Ut máneat vobíscum in ætérnum, Spíritum veritátis, allelúja.\end{Resp}
\switchcolumn
\EN
\begin{Resp}\Rsign I will pray the Father, and He shall give you another Comforter, \Star{} That He may abide with you for ever, even the Spirit of truth. Alleluia.\end{Resp}
\begin{Resp}\Vsign For if I go not away, the Comforter will not come unto you; but if I depart, I will send Him unto you.\end{Resp}
\begin{Resp}\Rsign That He may abide with you for ever, even the Spirit of truth. Alleluia.\end{Resp}
\switchcolumn*

\LA
\LessonHead{Lectio VIII}
\switchcolumn
\EN
\LessonHead{Lesson VIII}
\switchcolumn*

\LA
\noindent Quod cum ita sit, fídei nostræ veritátem in vitæ nostræ consideratióne debémus agnóscere. Tunc enim veráciter fidéles sumus, si quod verbis promíttimus, opéribus complémus. In die quippe baptísmatis, ómnibus nos antíqui hostis opéribus, atque ómnibus pompis abrenuntiáre promísimus. Itaque unusquísque vestrum ad consideratiónem suam mentis óculos redúcat; et si servat post baptísmum, quod ante baptísmum spopóndit, certus jam, quia fidélis est, gáudeat.\par
\switchcolumn
\EN
\noindent Once, then, it so standeth, it is to our lives we must look for proof of the reality of our faith. Then only are we truly Christ's faithful people when our works are the fulfilment of our profession. The day whereon we were baptized we bound ourselves to renounce all the works of the old enemy, and all his pomps. Therefore let every one of you now turn his inward eye upon his own behaviour, and if, since his baptism, he hath kept that promise which he made before it, let him know that he is in very truth one of Christ's faithful ones and let him rejoice.\par
\switchcolumn*

\LA
\begin{Resp}\Rsign Ponis nubem ascénsum tuum, Dómine: \Star{} Qui ámbulas super pennas ventórum, allelúja.\end{Resp}
\begin{Resp}\Vsign Confessiónem et decórem induísti, amíctus lumen sicut vestiméntum.\end{Resp}
\begin{Resp}\Rsign Qui ámbulas super pennas ventórum, allelúja.\end{Resp}
\begin{Resp}\Vsign Glória Patri, et Fílio, \Star{} et Spirítui Sancto.\end{Resp}
\begin{Resp}\Rsign Qui ámbulas super pennas ventórum, allelúja.\end{Resp}
\switchcolumn
\EN
\begin{Resp}\Rsign Thou makest the clouds the chariot, O Lord, \Star{} Thou walkest upon the wings of the wind. Alleluia.\end{Resp}
\begin{Resp}\Vsign Thou art clothed with honour and majesty, covering thyself with light as with a garment;\end{Resp}
\begin{Resp}\Rsign Thou walkest upon the wings of the wind. Alleluia.\end{Resp}
\begin{Resp}\Vsign Glory be to the Father, and to the Son, \Star{} and to the Holy Ghost.\end{Resp}
\begin{Resp}\Rsign Thou walkest upon the wings of the wind. Alleluia.\end{Resp}
\switchcolumn*

\LA
\LessonHead{Lectio IX}
\switchcolumn
\EN
\LessonHead{Lesson IX}
\switchcolumn*

\LA
\noindent Sed ecce, si quod promísit, mínime servávit, si ad exercénda prava ópera, ad concupiscéndas mundi pompas dilápsus est: videámus, si jam scit plángere quod errávit. Apud misericórdem namque júdicem nec ille fallax habétur, qui ad veritátem revértitur, étiam postquam mentítur: quia omnípotens Deus, dum libénter nostram pœniténtiam súscipit, ipse suo judício hoc, quod errávimus, abscóndit.\par
\switchcolumn
\EN
\noindent But if he hath utterly broken his promise, if he hath fallen away to work iniquity, and to lust after the pomps of the world, let us see if he now knoweth how to weep over his backsliding. By the merciful Judge that man is not punished as a perjurer who in the end telleth the truth, even though he hath first lied. Because Almighty God doth, in His tender kindness, so receive our contrition, that, in His judgment, He declareth us not guilty of that which we have done amiss.\par
\switchcolumn*

\LA
\TeDeum{Te Deum}
\switchcolumn
\EN
\TeDeum{Te Deum}
\switchcolumn*

\end{paracol}

\Office{Sabbato infra Octavam Ascensionis}{Saturday within the Octave of the Ascension}{Semiduplex}
\addcontentsline{toc}{subsection}{Sabbato infra Octavam Ascensionis\enspace\textit{Saturday within the Octave of the Ascension}}
\begin{paracol}{2}
\LA
\LessonHead{Lectio I}
\switchcolumn
\EN
\LessonHead{Lesson I}
\switchcolumn*

\LA
\LessonTitle{De Epístola secúnda beáti Petri Apóstoli}
\switchcolumn
\EN
\LessonTitle{Lesson from the second letter of St. Peter the Apostle}
\switchcolumn*

\LA
\Cite{2 Pet 3:1-7}
\switchcolumn
\EN
\Cite{2 Pet 3:1-7}
\switchcolumn*

\LA
\noindent Hanc ecce vobis, caríssimi, secúndam scribo epístolam, in quibus vestram éxcito in commonitióne sincéram mentem: ut mémores sitis eórum, quæ prædíxi, verbórum, a sanctis prophétis et apostolórum vestrórum, præceptórum Dómini et Salvatóris. Hoc primum sciéntes, quod vénient in novíssimis diébus in deceptióne illusóres, juxta próprias concupiscéntias ambulántes, dicéntes: Ubi est promíssio, aut advéntus ejus? ex quo enim patres dormiérunt, ómnia sic perséverant ab inítio creatúræ. Latet enim eos hoc voléntes, quod cæli erant prius, et terra de aqua, et per aquam consístens Dei verbo: per quæ, ille tunc mundus aqua inundátus, périit. Cæli autem, qui nunc sunt, et terra eódem verbo repósiti sunt, igni reserváti in diem judícii, et perditiónis impiórum hóminum.\par
\switchcolumn
\EN
\noindent Behold this second epistle I write to you, my dearly beloved, in which I stir up by way of admonition your sincere mind: That you may be mindful of those words which I told you before from the holy prophets, and of your apostles, of the precepts of the Lord and Saviour. Knowing this first, that in the last days there shall come deceitful scoffers, walking after their own lusts, Saying: Where is his promise or his coming? for since the time that the fathers slept, all things continue as they were from the beginning of the creation. For this they are wilfully ignorant of, that the heavens were before, and the earth out of water, and through water, consisting by the word of God. Whereby the world that then was, being overflowed with water, perished. But the heavens and the earth which are now, by the same word are kept in store, reserved unto fire against the day of judgment and perdition of the ungodly men.\par
\switchcolumn*

\LA
\begin{Resp}\Rsign Post passiónem suam per dies quadragínta appárens eis, et loquens de regno Dei, allelúja: \Star{} Et, vidéntibus illis, elevátus est, allelúja: et nubes suscépit eum ab óculis eórum, allelúja.\end{Resp}
\begin{Resp}\Vsign Et convéscens, præcépit eis, ab Jerosólymis ne discéderent, sed exspectárent promissiónem Patris.\end{Resp}
\begin{Resp}\Rsign Et, vidéntibus illis, elevátus est, allelúja: et nubes suscépit eum ab óculis eórum, allelúja.\end{Resp}
\switchcolumn
\EN
\begin{Resp}\Rsign Being seen of them forty days after that He had suffered, and speaking of the kingdom of God. Alleluia. \Star{} And while they beheld, He was taken up, and a cloud received Him out of their sight. Alleluia.\end{Resp}
\begin{Resp}\Vsign And, eating together with them, He commanded them that they should not depart from Jerusalem, but wait for the Promise of the Father.\end{Resp}
\begin{Resp}\Rsign And while they beheld, He was taken up, and a cloud received Him out of their sight. Alleluia.\end{Resp}
\switchcolumn*

\LA
\LessonHead{Lectio II}
\switchcolumn
\EN
\LessonHead{Lesson II}
\switchcolumn*

\LA
\Cite{2 Pet 3:8-13}
\switchcolumn
\EN
\Cite{2 Pet 3:8-13}
\switchcolumn*

\LA
\noindent Unum vero hoc non láteat vos, caríssimi, quia unus dies apud Dóminum sicut mille anni, et mille anni sicut dies unus. Non tardat Dóminus promissiónem suam, sicut quidam exístimant: sed patiénter agit propter vos, nolens áliquos períre, sed omnes ad pœniténtiam revérti. Advéniet autem dies Dómini ut fur: in quo cæli magno ímpetu tránsient, eleménta vero calóre solvéntur, terra autem et quæ in ipsa sunt ópera, exuréntur. Cum ígitur hæc ómnia dissolvénda sunt, quales opórtet vos esse in sanctis conversatiónibus, et pietátibus, exspectántes, et properántes in advéntum diéi Dómini, per quem cæli ardéntes solvéntur, et eleménta ignis ardóre tabéscent? Novos vero cælos, et novam terram secúndum promíssa ipsíus exspectámus, in quibus justítia hábitat.\par
\switchcolumn
\EN
\noindent But of this one thing be not ignorant, my beloved, that one day with the Lord is as a thousand years, and a thousand years as one day. The Lord delayeth not his promise, as some imagine, but dealeth patiently for your sake, not willing that any should perish, but that all should return to penance. But the day of the Lord shall come as a thief, in which the heavens shall pass away with great violence, and the elements shall be melted with heat, and the earth and the works which are in it, shall be burnt up. Seeing then that all these things are to be dissolved, what manner of people ought you to be in holy conversation and godliness? Looking for and hasting unto the coming of the day of the Lord, by which the heavens being on fire shall be dissolved, and the elements shall melt with the burning heat? But we look for new heavens and a new earth according to his promises, in which justice dwelleth.\par
\switchcolumn*

\LA
\begin{Resp}\Rsign Omnis pulchritúdo Dómini exaltáta est super sídera: \Star{} Spécies ejus in núbibus cæli, et nomen ejus in ætérnum pérmanet, allelúja.\end{Resp}
\begin{Resp}\Vsign A summo cælo egréssio ejus, et occúrsus ejus usque ad summum ejus.\end{Resp}
\begin{Resp}\Rsign Spécies ejus in núbibus cæli, et nomen ejus in ætérnum pérmanet, allelúja.\end{Resp}
\switchcolumn
\EN
\begin{Resp}\Rsign The Lord hath set His beauty above the stars; \Star{} His loveliness is in the clouds of heaven, and His Name endureth for ever. Alleluia.\end{Resp}
\begin{Resp}\Vsign His going forth is from the end of the heaven, and His circuit unto the ends of it.\end{Resp}
\begin{Resp}\Rsign His loveliness is in the clouds of heaven, and His Name endureth for ever. Alleluia.\end{Resp}
\switchcolumn*

\LA
\LessonHead{Lectio III}
\switchcolumn
\EN
\LessonHead{Lesson III}
\switchcolumn*

\LA
\Cite{2 Pet 3:14-18}
\switchcolumn
\EN
\Cite{2 Pet 3:14-18}
\switchcolumn*

\LA
\noindent Propter quod, caríssimi, hæc exspectántes, satágite immaculáti, et invioláti ei inveníri in pace: et Dómini nostri longanimitátem, salútem arbitrémini: sicut et caríssimus frater noster Paulus secúndum datam sibi sapiéntiam scripsit vobis, sicut et in ómnibus epístolis, loquens in eis de his in quibus sunt quædam difficília intelléctu, quæ indócti et instábiles deprávant, sicut et céteras Scriptúras, ad suam ipsórum perditiónem. Vos ígitur fratres, præsciéntes custodíte, ne insipiéntium erróre tradúcti excidátis a própria firmitáte: créscite vero in grátia, et in cognitióne Dómini nostri, et Salvatóris Jesu Christi. Ipsi glória et nunc, et in diem æternitátis. Amen.\par
\switchcolumn
\EN
\noindent Wherefore, dearly beloved, waiting for these things, be diligent that you may be found before him unspotted and blameless in peace. And account the longsuffering of our Lord, salvation; as also our most dear brother Paul, according to the wisdom given him, hath written to you: As also in all his epistles, speaking in them of these things; in which are certain things hard to be understood, which the unlearned and unstable wrest, as they do also the other scriptures, to their own destruction. You therefore, brethren, knowing these things before, take heed, lest being led aside by the error of the unwise, you fall from your own steadfastness. But grow in grace, and in the knowledge of our Lord and Saviour Jesus Christ. To him be glory both now and unto the day of eternity. Amen.\par
\switchcolumn*

\LA
\begin{Resp}\Rsign Exaltáre, Dómine, allelúja, \Star{} In virtúte tua, allelúja.\end{Resp}
\begin{Resp}\Vsign Eleváta est, magnificéntia tua super cælos, Deus.\end{Resp}
\begin{Resp}\Rsign In virtúte tua, allelúja.\end{Resp}
\begin{Resp}\Vsign Glória Patri, et Fílio, \Star{} et Spirítui Sancto.\end{Resp}
\begin{Resp}\Rsign In virtúte tua, allelúja.\end{Resp}
\switchcolumn
\EN
\begin{Resp}\Rsign Be Thou exalted, O Lord. Alleluia. \Star{} In thine Own strength. Alleluia.\end{Resp}
\begin{Resp}\Vsign O God, Thou hast set thy glory above the heavens.\end{Resp}
\begin{Resp}\Rsign In thine Own strength. Alleluia.\end{Resp}
\begin{Resp}\Vsign Glory be to the Father, and to the Son, \Star{} and to the Holy Ghost.\end{Resp}
\begin{Resp}\Rsign In thine Own strength. Alleluia.\end{Resp}
\switchcolumn*

\LA
\LessonHead{Lectio IV}
\switchcolumn
\EN
\LessonHead{Lesson IV}
\switchcolumn*

\LA
\LessonTitle{Sermo sancti Leónis Papæ}
\switchcolumn
\EN
\LessonTitle{From the Sermons of Pope St. Leo the Great.}
\switchcolumn*

\LA
\Source{Sermo 2 de Ascensione Domini}
\switchcolumn
\EN
\Source{2nd on the Ascension.}
\switchcolumn*

\LA
\noindent Quod ítaque Redemptóris nostri conspícuum fuit, in sacraménta transívit: et ut fides excelléntior esset ac fírmior, visióni doctrína succéssit, cujus auctoritátem supérnis illumináta rádiis credéntium corda sequeréntur. Hanc fidem ascensióne Dómini auctam, et Spíritus Sancti múnere roborátam, non víncula, non cárceres, non exsília, non fames, non ignis, non laniátus ferárum, nec exquisíta persequéntium crudelitátibus supplícia terruérunt. Pro hac fide per univérsum mundum non solum viri, sed étiam féminæ: nec tantum impúbes púeri, sed étiam téneræ vírgines usque ad effusiónem sui sánguinis decertárunt. Hæc fides dæmónia ejécit, ægritúdines dépulit, mórtuos suscitávit.\par
\switchcolumn
\EN
\noindent And so the seen Presence of our Redeemer in the Body was changed for an unseen Presence in the Sacraments, and hearing was given to the Church in place of seeing, that her faith, rightly so called, might be the more victorious and stedfast and that teaching, which the hearts of all her children are called on to hear, is a teaching enlightened by rays from heaven. This faith, strengthened by the Ascension of the Lord, and established by the gift of the Holy Ghost, neither bonds, nor imprisonment, nor exile, nor famine, nor fire, nor savage beasts, nor those forms of death, finewrought in cruelty, wherein they that persecute us are well skilled, have been able to scare. For this faith there have striven throughout the whole world, even unto the out-pouring of their blood, not men only, but women also, not little lads only, but tender maidens. This is the faith which hath cast out devils, healed diseases, raised the dead.\par
\switchcolumn*

\LA
\begin{Resp}\Rsign Tempus est, ut revértar ad eum, qui me misit, dicit Dóminus: nolíte contristári, nec turbétur cor vestrum: \Star{} Rogo pro vobis Patrem, ut ipse vos custódiat, allelúja, allelúja.\end{Resp}
\begin{Resp}\Vsign Nisi ego abíero, Paráclitus non véniet: cum assúmptus fúero, mittam vobis eum.\end{Resp}
\begin{Resp}\Rsign Rogo pro vobis Patrem, ut ipse vos custódiat, allelúja, allelúja.\end{Resp}
\switchcolumn
\EN
\begin{Resp}\Rsign My time is come that I should return unto Him That sent Me, saith the Lord. Be not sorrowful, neither let your heart be troubled. \Star{} I pray the Father for you, that He may keep you. Alleluia, Alleluia.\end{Resp}
\begin{Resp}\Vsign If I go not away, the Comforter will not come unto you; where I am ascended, I will send Him unto you.\end{Resp}
\begin{Resp}\Rsign I pray the Father for you, that He may keep you. Alleluia, Alleluia.\end{Resp}
\switchcolumn*

\LA
\LessonHead{Lectio V}
\switchcolumn
\EN
\LessonHead{Lesson V}
\switchcolumn*

\LA
\noindent Unde et ipsi beáti Apóstoli, qui tot miráculis confirmáti, tot sermónibus erudíti, atrocitátem tamen Domínicæ passiónis expáverant, et veritátem resurrectiónis ejus non sine hæsitatióne suscéperant; tantum de ascensióne Dómini profecérunt, ut quidquid illis prius intúlerat metum, verterétur in gáudium. Totam enim contemplatiónem ánimi in divinitátem ad Patris déxteram considéntis eréxerant: nec jam corpóreæ visiónis tardabántur objéctu, quo minus in id áciem mentis inténderent, quod nec a Patre descendéndo abfúerat, nec a discípulis ascendéndo discésserat. Tunc ígitur, dilectíssimi, Fílius hóminis, Dei Fílius excelléntius sacratiúsque innótuit, cum in patérnæ majestátis glóriam se recépit: et ineffábili modo cœpit esse divinitáte præséntior, qui factus est humanitáte longínquior.\par
\switchcolumn
\EN
\noindent Hence even the blessed Apostles themselves, who had been comforted by so many miracles and taught by so many discourses, were sickened by the horrors of their Lord's Passion, and received but doubtfully the assurance of His Resurrection, till after the Lord's Ascension and then fared on so bravely, that all that had been fearful to them before became joyful then. The reason was that they had lifted up all their mind to think of the Godhead of Him Who sitteth at the right hand of the Father. They asked no longer for a seen Presence, when their spiritual eye had caught the fact that, even as, when He had come down to earth, He had not left His Father, so now that He was gone up into heaven, He had not left His disciples. So then it was, dearly beloved brethren, that the Son of man more excellently and more sacredly revealed Himself as the Son of God, when He had withdrawn Himself again into that glory which He had with the Father before the world was. In some unspeakable way He began to be more present, as touching His Godhead, when He removed Himself farther from us, as touching His Manhood.\par
\switchcolumn*

\LA
\begin{Resp}\Rsign Non turbétur cor vestrum: ego vado ad Patrem; et cum assúmptus fúero a vobis, mittam vobis, allelúja, \Star{} Spíritum veritátis, et gaudébit cor vestrum, allelúja.\end{Resp}
\begin{Resp}\Vsign Ego rogábo Patrem, et álium Paráclitum dabit vobis.\end{Resp}
\begin{Resp}\Rsign Spíritum veritátis, et gaudébit cor vestrum, allelúja.\end{Resp}
\switchcolumn
\EN
\begin{Resp}\Rsign Let not your heart be troubled; I go unto the Father, and when I am taken from you, I will send unto you, alleluia, \Star{} The Spirit of truth, and your heart shall rejoice. Alleluia.\end{Resp}
\begin{Resp}\Vsign I will pray the Father, and He shall give you another Comforter.\end{Resp}
\begin{Resp}\Rsign The Spirit of truth, and your heart shall rejoice. Alleluia.\end{Resp}
\switchcolumn*

\LA
\LessonHead{Lectio VI}
\switchcolumn
\EN
\LessonHead{Lesson VI}
\switchcolumn*

\LA
\noindent Tunc ad æquálem Patri Fílium erudítior fides gressu mentis cœpit accédere, et contrectatióne in Christo corpóreæ substántiæ, qua Patre minor est, non egére: quóniam glorificáti córporis manénte natúra, eo fides credéntium vocabátur, ubi non carnáli manu, sed spiritáli intelléctu par Genitóri Unigénitus tangerétur. Hinc illud est, quod post resurrectiónem suam Dóminus Maríæ Magdalénæ persónam Ecclésiæ gerénti, cum ad contáctum ipsíus properáret accédere, dicit: Noli me tángere, nondum enim ascéndi ad Patrem meum: hoc est, Nolo ut ad me corporáliter vénias, nec ut me sensu carnis agnóscas: ad sublimióra te díffero, majóra tibi prǽparo: cum ad Patrem ascéndero, tunc me perféctius veriúsque palpábis, apprehensúra quod non tangis, et creditúra quod non cernis.\par
\switchcolumn
\EN
\noindent Then it was that a better instructed faith began intellectually to approach the idea of a Son equal to the Father, and no longer to need to handle in Christ the bodily Matter, Which is of a nature as touching which He is inferior to the Father since, Its nature still remaining in the glorified Body, the faith of believers was summoned to that place where the Only-Begotten Son, Who is equal to the Father, is felt, not by the application of a bodily hand, but by the effort of a spiritualminded intellect. Hence it was that after His Resurrection, when Mary Magdalene, (in whom was there represented the Person of the whole Church,) wished to handle the Lord, He said: “Touch Me not, for I am not yet ascended to My Father”, that is, “I will no more that thy nearness to Me should be a nearness of body to Body, nor that thine experience of Me should henceforward be one proceeding from fleshly experiment for that, I appoint thee an higher world, I make ready for thee a nobler form of it than this after that I have ascended to My Father, a time will come when thou shalt indeed touch Me, but after a manner more perfect, more real than this, even a time when thou shalt lay hold on that which thou touchest not now, and believe that which thou seest not now.”\par
\switchcolumn*

\LA
\begin{Resp}\Rsign Ascéndens Christus in altum, captívam duxit captivitátem, \Star{} Dedit dona homínibus, allelúja, allelúja, allelúja.\end{Resp}
\begin{Resp}\Vsign Ascéndit Deus in jubilatióne, et Dóminus in voce tubæ.\end{Resp}
\begin{Resp}\Rsign Dedit dona homínibus, allelúja, allelúja, allelúja.\end{Resp}
\begin{Resp}\Vsign Glória Patri, et Fílio, \Star{} et Spirítui Sancto.\end{Resp}
\begin{Resp}\Rsign Dedit dona homínibus, allelúja, allelúja, allelúja.\end{Resp}
\switchcolumn
\EN
\begin{Resp}\Rsign When Christ ascended up on high, He led captivity captive, \Star{} He gave gifts unto men. Alleluia, Alleluia, Alleluia.\end{Resp}
\begin{Resp}\Vsign God is gone up with a shout, and the Lord with the sound of a trumpet.\end{Resp}
\begin{Resp}\Rsign He gave gifts unto men. Alleluia, Alleluia, Alleluia.\end{Resp}
\begin{Resp}\Vsign Glory be to the Father, and to the Son, \Star{} and to the Holy Ghost.\end{Resp}
\begin{Resp}\Rsign He gave gifts unto men. Alleluia, Alleluia, Alleluia.\end{Resp}
\switchcolumn*

\LA
\LessonHead{Lectio VII}
\switchcolumn
\EN
\LessonHead{Lesson VII}
\switchcolumn*

\LA
\LessonTitle{Léctio sancti Evangélii secúndum Marcum}
\switchcolumn
\EN
\LessonTitle{From the Holy Gospel according to Mark}
\switchcolumn*

\LA
\Cite{Marc 16:14-20}
\switchcolumn
\EN
\Cite{Mark 16:14-20}
\switchcolumn*

\LA
\noindent In illo témpore: Recumbéntibus úndecim discípulis, appáruit illis Jesus: et exprobrávit incredulitátem eórum et durítiam cordis, quia iis, qui víderant eum resurrexísse, non credidérunt. Et réliqua.\par
\switchcolumn
\EN
\noindent At that time Jesus appeared unto the eleven disciples as they sat at meat, and upbraided them with their unbelief and hardness of heart because they believed not them which had seen Him after He was risen. And so on.\par
\switchcolumn*

\LA
\LessonTitle{Homilía sancti Gregórii Papæ}
\switchcolumn
\EN
\LessonTitle{Homily by Pope St. Gregory the Great.}
\switchcolumn*

\LA
\Source{Ex eadem Homilía 29}
\switchcolumn
\EN
\Source{Same as before.}
\switchcolumn*

\LA
\noindent Signa autem eos qui creditúri sunt, hæc sequéntur: In nómine meo dæmónia eícient, linguis loquéntur novis, serpéntes tollent: et si mortíferum quid bíberint, non eis nocébit: super ægros manus impónent, et bene habébunt. Numquídnam, fratres mei, quia ista signa non fácitis, mínime créditis? Sed hæc necessária in exórdio Ecclésiæ fuérunt. Ut enim ad fidem crésceret multitúdo credéntium, miráculis fúerat nutriénda: quia et nos, cum arbústa plantámus, támdiu eis aquam infúndimus, quoúsque ea in terra jam coaluísse videámus: et si semel radícem fíxerint, irrigátio cessábit. Hinc est enim quod Paulus dicit: Linguæ in signum sunt non fidélibus, sed infidélibus.\par
\switchcolumn
\EN
\noindent “And these signs shall follow them that believe In My Name they shall cast out devils they shall speak with new tongues they shall take up serpents and if they drink any deadly thing, it shall not hurt them they shall lay hands on the sick, and they shall recover.” My brethren, these signs do not follow us. Do we, then, not believe Nay. The truth is, these things were needful when the Church was young. That she might grow by the increase of the faithful, she needed to be nourished with miracles. Even so we, when we plant a young tree, continually water and tend it till we see that it hath taken firm root in the earth but when once it hath taken firm root, it can grow of itself. Hence, Paul saith of tongues: “Tongues are for a sign, not to them that believe, but to them that believe not.” (1 Cor. xiv. 22.)\par
\switchcolumn*

\LA
\begin{Resp}\Rsign Ego rogábo Patrem, et álium Paráclitum dabit vobis, \Star{} Ut máneat vobíscum in ætérnum, Spíritum veritátis, allelúja.\end{Resp}
\begin{Resp}\Vsign Si enim non abíero, Paráclitus non véniet ad vos: si autem abíero, mittam eum ad vos.\end{Resp}
\begin{Resp}\Rsign Ut máneat vobíscum in ætérnum, Spíritum veritátis, allelúja.\end{Resp}
\switchcolumn
\EN
\begin{Resp}\Rsign I will pray the Father, and He shall give you another Comforter, \Star{} That He may abide with you for ever, even the Spirit of truth. Alleluia.\end{Resp}
\begin{Resp}\Vsign For if I go not away, the Comforter will not come unto you; but if I depart, I will send Him unto you.\end{Resp}
\begin{Resp}\Rsign That He may abide with you for ever, even the Spirit of truth. Alleluia.\end{Resp}
\switchcolumn*

\LA
\LessonHead{Lectio VIII}
\switchcolumn
\EN
\LessonHead{Lesson VIII}
\switchcolumn*

\LA
\noindent Habémus de his signis atque virtútibus, quæ adhuc subtílius consideráre debeámus. Sancta quippe Ecclésia quotídie spiritáliter facit, quod tunc per Apóstolos corporáliter faciébat. Nam sacerdótes ejus cum per exorcísmi grátiam manum credéntibus impónunt, et habitáre malígnos spíritus in eórum mente contradícunt, quid áliud fáciunt, nisi dæmónia ejíciunt? Et fidéles quique, qui jam vitæ véteris sæculária verba derelínquunt, sancta autem mystéria ínsonant, Conditóris sui laudes et poténtiam, quantum prǽvalent, narrant: quid áliud fáciunt, nisi novis linguis loquúntur? Qui dum bonis suis exhortatiónibus malítiam de aliénis córdibus áuferunt, serpéntes tollunt.\par
\switchcolumn
\EN
\noindent We have a deeper matter of thought touching these signs and mighty works. It is the work of the holy Church to do every day spiritually that which the Apostles then did carnally. When her Priests, armed with the power of exorcism, lay their hands upon believers, and command evil spirits to dwell no longer in their souls, what is it they do but cast out devils When Christ's faithful people themselves give up the language of their old life, and speak the wonderful works of God, the glory and power of their Maker, telling of them with all their strength, what is it they do then but speak with new tongues When either the one or the other doth by his exhortation charm the wickedness out of his neighbour's heart, what is it he doth but take up serpents?\par
\switchcolumn*

\LA
\begin{Resp}\Rsign Ponis nubem ascénsum tuum, Dómine: \Star{} Qui ámbulas super pennas ventórum, allelúja.\end{Resp}
\begin{Resp}\Vsign Confessiónem et decórem induísti, amíctus lumen sicut vestiméntum.\end{Resp}
\begin{Resp}\Rsign Qui ámbulas super pennas ventórum, allelúja.\end{Resp}
\begin{Resp}\Vsign Glória Patri, et Fílio, \Star{} et Spirítui Sancto.\end{Resp}
\begin{Resp}\Rsign Qui ámbulas super pennas ventórum, allelúja.\end{Resp}
\switchcolumn
\EN
\begin{Resp}\Rsign Thou makest the clouds the chariot, O Lord, \Star{} Thou walkest upon the wings of the wind. Alleluia.\end{Resp}
\begin{Resp}\Vsign Thou art clothed with honour and majesty, covering thyself with light as with a garment;\end{Resp}
\begin{Resp}\Rsign Thou walkest upon the wings of the wind. Alleluia.\end{Resp}
\begin{Resp}\Vsign Glory be to the Father, and to the Son, \Star{} and to the Holy Ghost.\end{Resp}
\begin{Resp}\Rsign Thou walkest upon the wings of the wind. Alleluia.\end{Resp}
\switchcolumn*

\LA
\LessonHead{Lectio IX}
\switchcolumn
\EN
\LessonHead{Lesson IX}
\switchcolumn*

\LA
\noindent Et dum pestíferas suasiónes áudiunt, sed tamen ad operatiónem pravam mínime pertrahúntur, mortíferum quidem est quod bibunt, sed non eis nocébit. Qui quóties próximos suos in bono ópere infirmári conspíciunt, dum eis tota virtúte concúrrunt, et exémplo suæ operatiónis illórum vitam róborant, qui in própria actióne títubant: quid áliud fáciunt, nisi super ægros manus impónunt, ut bene hábeant? Quæ nimírum mirácula tanto majóra sunt, quanto spiritália: tanto majóra sunt, quanto per hæc non córpora, sed ánimæ suscitántur.\par
\switchcolumn
\EN
\noindent When they hear the voice of temptation inviting to deadly sin, but are not drawn thereby to work iniquity, do they not then drink a deadly thing, and it doth not hurt them? As often as they see their neighbour fainting in well-doing, and run to help him with all their might, so that their example braceth the feeble life of the waverer, what do they but lay hands on the sick and they recover And indeed, such miracles as these are the greatest miracles, which are spiritual the greatest, for they bring health, not to the dying body, but to the immortal soul.\par
\switchcolumn*

\LA
\TeDeum{Te Deum}
\switchcolumn
\EN
\TeDeum{Te Deum}
\switchcolumn*

\end{paracol}

\Office{Dominica infra Octavam Ascensionis}{Sunday within the Octave of the Ascension}{Semiduplex}
\addcontentsline{toc}{subsection}{Dominica infra Octavam Ascensionis\enspace\textit{Sunday within the Octave of the Ascension}}
\begin{paracol}{2}
\LA
\LessonHead{Lectio I}
\switchcolumn
\EN
\LessonHead{Lesson I}
\switchcolumn*

\LA
\LessonTitle{Incipit Epístola prima beáti Joánnis Apóstoli}
\switchcolumn
\EN
\LessonTitle{Lesson from the first letter of St. John the Apostle}
\switchcolumn*

\LA
\Cite{1 Joannes 1:1-5}
\switchcolumn
\EN
\Cite{1 John 1:1-5}
\switchcolumn*

\LA
\noindent Quod fuit ab inítio, quod audívimus, quod vídimus óculis nostris, quod perspéximus, et manus nostræ contrectavérunt de verbo vitæ: et vita manifestáta est, et vídimus, et testámur, et annuntiámus vobis vitam ætérnam, quæ erat apud Patrem, et appáruit nobis: quod vídimus et audívimus, annuntiámus vobis, ut et vos societátem habeátis nobíscum, et socíetas nostra sit cum Patre, et cum Fílio ejus Jesu Christo. Et hæc scríbimus vobis ut gaudeátis, et gáudium vestrum sit plenum. Et hæc est annuntiátio, quam audívimus ab eo, et annuntiámus vobis: quóniam Deus lux est, et ténebræ in eo non sunt ullæ.\par
\switchcolumn
\EN
\noindent That which was from the beginning, which we have heard, which we have seen with our eyes, which we have looked upon, and our hands have handled, of the word of life: For the life was manifested; and we have seen and do bear witness, and declare unto you the life eternal, which was with the Father, and hath appeared to us: That which we have seen and have heard, we declare unto you, that you also may have fellowship with us, and our fellowship may be with the Father, and with his Son Jesus Christ. And these things we write to you, that you may rejoice, and your joy may be full. And this is the declaration which we have heard from him, and declare unto you: That God is light, and in him there is no darkness.\par
\switchcolumn*

\LA
\begin{Resp}\Rsign Post passiónem suam per dies quadragínta appárens eis, et loquens de regno Dei, allelúja: \Star{} Et, vidéntibus illis, elevátus est, allelúja: et nubes suscépit eum ab óculis eórum, allelúja.\end{Resp}
\begin{Resp}\Vsign Et convéscens, præcépit eis, ab Jerosólymis ne discéderent, sed exspectárent promissiónem Patris.\end{Resp}
\begin{Resp}\Rsign Et, vidéntibus illis, elevátus est, allelúja: et nubes suscépit eum ab óculis eórum, allelúja.\end{Resp}
\switchcolumn
\EN
\begin{Resp}\Rsign Being seen of them forty days after that He had suffered, and speaking of the kingdom of God. Alleluia. \Star{} And while they beheld, He was taken up, and a cloud received Him out of their sight. Alleluia.\end{Resp}
\begin{Resp}\Vsign And, eating together with them, He commanded them that they should not depart from Jerusalem, but wait for the Promise of the Father.\end{Resp}
\begin{Resp}\Rsign And while they beheld, He was taken up, and a cloud received Him out of their sight. Alleluia.\end{Resp}
\switchcolumn*

\LA
\LessonHead{Lectio II}
\switchcolumn
\EN
\LessonHead{Lesson II}
\switchcolumn*

\LA
\Cite{1 Joannes 1:6-10}
\switchcolumn
\EN
\Cite{John 1:6-10}
\switchcolumn*

\LA
\noindent Si dixérimus quóniam societátem habémus cum eo, et in ténebris ambulámus, mentímur, et veritátem non fácimus. Si autem in luce ambulámus sicut et ipse est in luce, societátem habémus ad ínvicem, et sanguis Jesu Christi, Fílii ejus, emúndat nos ab omni peccáto. Si dixérimus quóniam peccátum non habémus, ipsi nos sedúcimus, et véritas in nobis non est. Si confiteámur peccáta nostra: fidélis est, et justus, ut remíttat nobis peccáta nostra, et emúndet nos ab omni iniquitáte. Si dixérimus quóniam non peccávimus, mendácem fácimus eum, et verbum ejus non est in nobis.\par
\switchcolumn
\EN
\noindent If we say that we have fellowship with him, and walk in darkness, we lie, and do not the truth. But if we walk in the light, as he also is in the light, we have fellowship one with another, and the blood of Jesus Christ his Son cleanseth us from all sin. If we say that we have no sin, we deceive ourselves, and the truth is not in us. If we confess our sins, he is faithful and just, to forgive us our sins, and to cleanse us from all iniquity. If we say that we have not sinned, we make him a liar, and his word is not in us.\par
\switchcolumn*

\LA
\begin{Resp}\Rsign Omnis pulchritúdo Dómini exaltáta est super sídera: \Star{} Spécies ejus in núbibus cæli, et nomen ejus in ætérnum pérmanet, allelúja.\end{Resp}
\begin{Resp}\Vsign A summo cælo egréssio ejus, et occúrsus ejus usque ad summum ejus.\end{Resp}
\begin{Resp}\Rsign Spécies ejus in núbibus cæli, et nomen ejus in ætérnum pérmanet, allelúja.\end{Resp}
\switchcolumn
\EN
\begin{Resp}\Rsign The Lord hath set His beauty above the stars; \Star{} His loveliness is in the clouds of heaven, and His Name endureth for ever. Alleluia.\end{Resp}
\begin{Resp}\Vsign His going forth is from the end of the heaven, and His circuit unto the ends of it.\end{Resp}
\begin{Resp}\Rsign His loveliness is in the clouds of heaven, and His Name endureth for ever. Alleluia.\end{Resp}
\switchcolumn*

\LA
\LessonHead{Lectio III}
\switchcolumn
\EN
\LessonHead{Lesson III}
\switchcolumn*

\LA
\Cite{1 Joannes 2:1-6}
\switchcolumn
\EN
\Cite{1 John 2:1-6}
\switchcolumn*

\LA
\noindent Filíoli mei, hæc scribo vobis, ut non peccétis. Sed et si quis peccáverit, advocátum habémus apud Patrem, Jesum Christum justum: et ipse est propitiátio pro peccátis nostris: non pro nostris autem tantum, sed étiam pro totíus mundi. Et in hoc scimus quóniam cognóvimus eum, si mandáta ejus observémus. Qui dicit se nosse eum, et mandáta ejus non custódit, mendax est, et in hoc véritas non est. Qui autem servat verbum ejus, vere in hoc cáritas Dei perfécta est: et in hoc scimus quóniam in ipso sumus. Qui dicit se in ipso manére, debet, sicut ille ambulávit, et ipse ambuláre.\par
\switchcolumn
\EN
\noindent My little children, these things I write to you, that you may not sin. But if any man sin, we have an advocate with the Father, Jesus Christ the just: And he is the propitiation for our sins: and not for ours only, but also for those of the whole world. And by this we know that we have known him, if we keep his commandments. He who saith that he knoweth him, and keepeth not his commandments, is a liar, and the truth is not in him. But he that keepeth his word, in him in very deed the charity of God is perfected; and by this we know that we are in him. He that saith he abideth in him, ought himself also to walk, even as he walked.\par
\switchcolumn*

\LA
\begin{Resp}\Rsign Exaltáre, Dómine, allelúja, \Star{} In virtúte tua, allelúja.\end{Resp}
\begin{Resp}\Vsign Eleváta est, magnificéntia tua super cælos, Deus.\end{Resp}
\begin{Resp}\Rsign In virtúte tua, allelúja.\end{Resp}
\begin{Resp}\Vsign Glória Patri, et Fílio, \Star{} et Spirítui Sancto.\end{Resp}
\begin{Resp}\Rsign In virtúte tua, allelúja.\end{Resp}
\switchcolumn
\EN
\begin{Resp}\Rsign Be Thou exalted, O Lord. Alleluia. \Star{} In thine Own strength. Alleluia.\end{Resp}
\begin{Resp}\Vsign O God, Thou hast set thy glory above the heavens.\end{Resp}
\begin{Resp}\Rsign In thine Own strength. Alleluia.\end{Resp}
\begin{Resp}\Vsign Glory be to the Father, and to the Son, \Star{} and to the Holy Ghost.\end{Resp}
\begin{Resp}\Rsign In thine Own strength. Alleluia.\end{Resp}
\switchcolumn*

\LA
\LessonHead{Lectio IV}
\switchcolumn
\EN
\LessonHead{Lesson IV}
\switchcolumn*

\LA
\LessonTitle{Sermo sancti Augustíni Epíscopi}
\switchcolumn
\EN
\LessonTitle{From the Sermons of St. Augustine, Bishop of Hippo.}
\switchcolumn*

\LA
\Source{Sermo 2 de Ascensióne Dómini, qui est 175 de Tempore}
\switchcolumn
\EN
\Source{2nd on the Ascension}
\switchcolumn*

\LA
\noindent Salvátor noster, dilectíssimi fratres, ascéndit in cælum: non ergo turbémur in terra. Ibi sit mens, et hic erit réquies. Ascendámus cum Christo ínterim corde: cum dies ejus promíssus advénerit, sequémur et córpore. Scire tamen debémus, fratres, quia cum Christo non ascéndit supérbia, non avarítia, non luxúria: nullum vítium nostrum ascéndit cum médico nostro. Et ídeo si post médicum desiderámus ascéndere, debémus vítia et peccáta depónere. Omnes enim quasi quibúsdam compédibus nos premunt, et peccatórum nos rétibus ligáre conténdunt: et ídeo cum Dei adjutório, secúndum quod ait Psalmísta: Dirumpámus víncula eórum: ut secúri possímus dícere Dómino: Dirupísti víncula mea, tibi sacrificábo hóstiam laudis.\par
\switchcolumn
\EN
\noindent Dearly beloved brethren, our Saviour is gone up from us into heaven, but let us not be troubled on earth. Let only our heart be there with Him, and we shall have peace here. Let us in heart thither ascend with Christ in the mean while, and when that glad day which He hath promised cometh, our body will follow. But we must know, my brethren, that there are some things that cannot ascend with Christ pride cannot, nor covetousness, nor brutishness no one of our diseases can ascend thither where our Healer is. And, therefore, if we would follow our Healer, we must needs leave our diseases and sins behind us. All such things tie us down, as it were, with bands, and hamper us in the meshes of a net of sins but, with God's help, we will say with the Psalmist: “Let us break their bands asunder,” (ii. 3,) that we may be able honestly to say to the Lord: “Thou hast loosed my bonds, I will offer to thee the sacrifice of thanksgiving.” (cxv. 16, 17.)\par
\switchcolumn*

\LA
\begin{Resp}\Rsign Tempus est, ut revértar ad eum, qui me misit, dicit Dóminus: nolíte contristári, nec turbétur cor vestrum: \Star{} Rogo pro vobis Patrem, ut ipse vos custódiat, allelúja, allelúja.\end{Resp}
\begin{Resp}\Vsign Nisi ego abíero, Paráclitus non véniet: cum assúmptus fúero, mittam vobis eum.\end{Resp}
\begin{Resp}\Rsign Rogo pro vobis Patrem, ut ipse vos custódiat, allelúja, allelúja.\end{Resp}
\switchcolumn
\EN
\begin{Resp}\Rsign My time is come that I should return unto Him That sent Me, saith the Lord. Be not sorrowful, neither let your heart be troubled. \Star{} I pray the Father for you, that He may keep you. Alleluia, Alleluia.\end{Resp}
\begin{Resp}\Vsign If I go not away, the Comforter will not come unto you; where I am ascended, I will send Him unto you.\end{Resp}
\begin{Resp}\Rsign I pray the Father for you, that He may keep you. Alleluia, Alleluia.\end{Resp}
\switchcolumn*

\LA
\LessonHead{Lectio V}
\switchcolumn
\EN
\LessonHead{Lesson V}
\switchcolumn*

\LA
\noindent Resurréctio Dómini spes nostra est: ascénsio Dómini glorificátio nostra est. Ascensiónis hódie solémnia celebrámus. Si ergo recte, si fidéliter, si devóte, si sancte, si pie ascensiónem Dómini celebrámus, ascendámus cum illo, et sursum corda habeámus. Ascendéntes autem non extollámur, nec de nostris, quasi de própriis méritis præsumámus. Sursum autem corda habére debémus ad Dóminum. Sursum enim cor non ad Dóminum, supérbia vocátur: sursum autem cor ad Dóminum, refúgium vocátur. Vidéte, fratres, magnum miráculum. Altus est Deus: érigis te, et fugit a te: humílias te, et descéndit ad te. Quare hoc? Quia excélsus est, et humília réspicit, et alta de longe cognóscit. Humília de próximo réspicit, ut attóllat: alta, id est, supérba, de longe cognóscit, ut déprimat.\par
\switchcolumn
\EN
\noindent The Resurrection of the Lord is our hope the Ascension of the Lord is our glorification. To-day we keep the solemn holiday of the Ascension. If, therefore, our keeping of this holiday is to be a right, faithful, earnest, holy, godly keeping, we must in mind likewise ascend, and lift up our hearts unto the Lord. When we ascend we must not be high-minded, nor flatter ourselves with our good works, as though they were our own. We must lift up our hearts unto the Lord. When man's heart is lifted up, but not unto the Lord, such lifting-up is pride to lift up the heart unto the Lord, is to make the Most High our Refuge. Behold, my brethren, a great wonder. God is high, but if thou art lifted up He fleeth from thee, whereas, if thou humblest thyself, He cometh down to thee. Wherefore? “The Lord is high, yet hath He respect unto the lowly but the proud He knoweth from afar.” (Ps. cxxxvii. 6.) To the lowly He hath respect, that He may raise them up the proud He knoweth from afar, that He may thrust them down.\par
\switchcolumn*

\LA
\begin{Resp}\Rsign Non turbétur cor vestrum: ego vado ad Patrem; et cum assúmptus fúero a vobis, mittam vobis, allelúja, \Star{} Spíritum veritátis, et gaudébit cor vestrum, allelúja.\end{Resp}
\begin{Resp}\Vsign Ego rogábo Patrem, et álium Paráclitum dabit vobis.\end{Resp}
\begin{Resp}\Rsign Spíritum veritátis, et gaudébit cor vestrum, allelúja.\end{Resp}
\switchcolumn
\EN
\begin{Resp}\Rsign Let not your heart be troubled; I go unto the Father, and when I am taken from you, I will send unto you, alleluia, \Star{} The Spirit of truth, and your heart shall rejoice. Alleluia.\end{Resp}
\begin{Resp}\Vsign I will pray the Father, and He shall give you another Comforter.\end{Resp}
\begin{Resp}\Rsign The Spirit of truth, and your heart shall rejoice. Alleluia.\end{Resp}
\switchcolumn*

\LA
\LessonHead{Lectio VI}
\switchcolumn
\EN
\LessonHead{Lesson VI}
\switchcolumn*

\LA
\noindent Resurréxit enim Christus, ut spem nobis daret, quia surgit homo, qui móritur: ne moriéndo desperarémus, et vitam nostram in morte finítam putarémus, secúros nos fecit. Sollíciti enim erámus de ipsa ánima: et ille nobis resurgéndo, de carnis resurrectióne fidúciam dedit. Crede ergo, ut mundéris. Prius te opórtet crédere, ut póstea per fidem Deum mereáris aspícere. Deum enim vidére vis? Audi ipsum: Beáti mundo corde: quóniam ipsi Deum vidébunt. Prius ergo cógita de corde mundándo: quidquid ibi vides, quod dísplicet Deo, tolle.\par
\switchcolumn
\EN
\noindent Christ arose again, to give us hope that this mortal will yet put on immortality He hath assured against an hopeless death, and against the thought that death endeth life. We were troubled, eyen as touching the soul but Christ, arising from the grave, hath assured to us the resurrection of the body also. Believe therefore, that thou mayest be made pure. First it behoveth thee to believe, if by faith thou wouldest in the end worthily see God. And wouldest thou see God Give ear to His own words “Blessed are the pure in heart, for they shall see God.” (Matth. v. 8.) Think first, then, how to purify thine heart; take from it whatsoever thou seest in it which displeaseth God.\par
\switchcolumn*

\LA
\begin{Resp}\Rsign Ascéndens Christus in altum, captívam duxit captivitátem, \Star{} Dedit dona homínibus, allelúja, allelúja, allelúja.\end{Resp}
\begin{Resp}\Vsign Ascéndit Deus in jubilatióne, et Dóminus in voce tubæ.\end{Resp}
\begin{Resp}\Rsign Dedit dona homínibus, allelúja, allelúja, allelúja.\end{Resp}
\begin{Resp}\Vsign Glória Patri, et Fílio, \Star{} et Spirítui Sancto.\end{Resp}
\begin{Resp}\Rsign Dedit dona homínibus, allelúja, allelúja, allelúja.\end{Resp}
\switchcolumn
\EN
\begin{Resp}\Rsign When Christ ascended up on high, He led captivity captive, \Star{} He gave gifts unto men. Alleluia, Alleluia, Alleluia.\end{Resp}
\begin{Resp}\Vsign God is gone up with a shout, and the Lord with the sound of a trumpet.\end{Resp}
\begin{Resp}\Rsign He gave gifts unto men. Alleluia, Alleluia, Alleluia.\end{Resp}
\begin{Resp}\Vsign Glory be to the Father, and to the Son, \Star{} and to the Holy Ghost.\end{Resp}
\begin{Resp}\Rsign He gave gifts unto men. Alleluia, Alleluia, Alleluia.\end{Resp}
\switchcolumn*

\LA
\LessonHead{Lectio VII}
\switchcolumn
\EN
\LessonHead{Lesson VII}
\switchcolumn*

\LA
\LessonTitle{Léctio sancti Evangélii secúndum Joánnem}
\switchcolumn
\EN
\LessonTitle{From the Holy Gospel according to John}
\switchcolumn*

\LA
\Cite{Joannes 15:26-27;16:1-4}
\switchcolumn
\EN
\Cite{John 15:26-27; 16:1-4}
\switchcolumn*

\LA
\noindent In illo témpore: Dixit Jesus discípulis suis: Cum vénerit Paráclitus, quem ego mittam vobis a Patre, Spíritum veritátis, qui a Patre procédit, ille testimónium perhibébit de me. Et réliqua.\par
\switchcolumn
\EN
\noindent At that time, Jesus said unto His disciples: When the Comforter is come, Whom I will send unto you from the Father, even the Spirit of truth, Which proceedeth from the Father, He shall testify of Me. And so on.\par
\switchcolumn*

\LA
\LessonTitle{Homilía sancti Augustíni Epíscopi}
\switchcolumn
\EN
\LessonTitle{Homily by St. Augustine, Bishop of Hippo}
\switchcolumn*

\LA
\Source{Tractatus 92 in Joannem}
\switchcolumn
\EN
\Source{Tract 92 on John}
\switchcolumn*

\LA
\noindent Dóminus Jesus in sermóne, quem locútus est discípulis suis post cenam, próximus passióni, tamquam itúrus, et relictúrus eos præséntia corporáli, cum ómnibus autem suis usque in consummatiónem sǽculi futúrus præséntia spiritáli, exhortátus est eos ad perferéndas persecutiónes impiórum, quos mundi nómine nuncupávit. Ex quo tamen mundo étiam ipsos discípulos se elegísse dixit: ut scirent se Dei grátia esse, quod sunt; suis autem vítiis fuísse, quod fuérunt.\par
\switchcolumn
\EN
\noindent The Lord Jesus, in that discourse which He addressed to His disciples after the Last Supper, when He was on the very eve of the Passion, when He was, as it were, about to go away and leave them as touching His bodily Presence, albeit as touching His spiritual Presence, He is with us alway even unto the end of the world, (Matth. xxviii. 20,) in that discourse He exhorted them to bear patiently the persecution of wicked men, of whom He speaketh as “the world”, out of which world, nevertheless, He saith that He hath chosen even His disciples themselves, (xv. 19,) that they might know that it was by the grace of God that they were what they were, (1 Cor. xv. 10,) whereas it was by their own sins that they had been what they had been.\par
\switchcolumn*

\LA
\begin{Resp}\Rsign Ego rogábo Patrem, et álium Paráclitum dabit vobis, \Star{} Ut máneat vobíscum in ætérnum, Spíritum veritátis, allelúja.\end{Resp}
\begin{Resp}\Vsign Si enim non abíero, Paráclitus non véniet ad vos: si autem abíero, mittam eum ad vos.\end{Resp}
\begin{Resp}\Rsign Ut máneat vobíscum in ætérnum, Spíritum veritátis, allelúja.\end{Resp}
\switchcolumn
\EN
\begin{Resp}\Rsign I will pray the Father, and He shall give you another Comforter, \Star{} That He may abide with you for ever, even the Spirit of truth. Alleluia.\end{Resp}
\begin{Resp}\Vsign For if I go not away, the Comforter will not come unto you; but if I depart, I will send Him unto you.\end{Resp}
\begin{Resp}\Rsign That He may abide with you for ever, even the Spirit of truth. Alleluia.\end{Resp}
\switchcolumn*

\LA
\LessonHead{Lectio VIII}
\switchcolumn
\EN
\LessonHead{Lesson VIII}
\switchcolumn*

\LA
\noindent Deínde persecutóres et suos et ipsórum, Judǽos evidénter expréssit: ut omníno apparéret étiam ipsos mundi damnábilis appellatióne conclúsos, qui persequúntur sanctos. Cumque de illis díceret, quod ignorárent eum a quo missus est; et tamen odíssent et Fílium, et Patrem, hoc est, et eum qui missus est, et eum a quo missus est: (de quibus ómnibus in áliis semónibus jam disserúimus) ad hoc pervénit, ubi ait: Ut adimpleátur sermo, qui in lege eórum scriptus est: Quia ódio habuérunt me gratis.\par
\switchcolumn
\EN
\noindent If they have persecuted Me, they will also persecute you.” Here He clearly pointeth to the Jews, the persecutors both of Himself and of His disciples, so that we see that they which persecute His holy ones are as much citizens of the world of damnation as they which persecuted Himself. He saith: “They know not Him That sent Me,” (21,) and yet again, (24,) “They have hated both Me and My Father,” (xv. 24,) that is to say, both the Sender and the Sent, the meaning of which words we have already treated in other discourses and with that He cometh to the words “That the word might be fulfilled that is written in their law They hated Me without a cause.”\par
\switchcolumn*

\LA
\begin{Resp}\Rsign Si enim non abíero, Paráclitus non véniet ad vos: si autem abíero, mittam eum ad vos. \Star{} Cum autem vénerit ille, docébit vos omnem veritátem, allelúja.\end{Resp}
\begin{Resp}\Vsign Non enim loquétur a semetípso: sed quæcúmque áudiet, loquétur: et quæ ventúra sunt, annuntiábit vobis.\end{Resp}
\begin{Resp}\Rsign Cum autem vénerit ille, docébit vos omnem veritátem, allelúja.\end{Resp}
\begin{Resp}\Vsign Glória Patri, et Fílio, \Star{} et Spirítui Sancto.\end{Resp}
\begin{Resp}\Rsign Cum autem vénerit ille, docébit vos omnem veritátem, allelúja.\end{Resp}
\switchcolumn
\EN
\begin{Resp}\Rsign For if I go not away, the Comforter will not come unto you; but if I depart, I will send Him unto you; \Star{} And when He is come, He will guide you into all truth. Alleluia.\end{Resp}
\begin{Resp}\Vsign For He shall not speak of Himself; but whatsoever He shall hear, that shall He speak and He will show you things to come.\end{Resp}
\begin{Resp}\Rsign And when He is come, He will guide you into all truth. Alleluia.\end{Resp}
\begin{Resp}\Vsign Glory be to the Father, and to the Son, \Star{} and to the Holy Ghost.\end{Resp}
\begin{Resp}\Rsign And when He is come, He will guide you into all truth. Alleluia.\end{Resp}
\switchcolumn*

\LA
\LessonHead{Lectio IX}
\switchcolumn
\EN
\LessonHead{Lesson IX}
\switchcolumn*

\LA
\noindent Deínde tamquam consequénter adjúnxit, unde modo disputáre suscépimus: Cum autem vénerit Paráclitus, quem ego mittam vobis a Patre, Spíritum veritátis, qui a Patre procédit, ille testimónium perhibébit de me: et vos testimónium perhibébitis, quia ab inítio mecum estis. Quid hoc pértinet ad illud quod díxerat: Nunc autem et vidérunt, et odérunt et me, et Patrem meum: sed ut impleátur sermo, qui in lege eórum scriptus est: Quia ódio habuérunt me gratis? An quia Paráclitus quando venit Spíritus veritátis, eos, qui vidérunt, et odérunt, testimónio manifestióre convícit? Immo vero étiam áliquos ex illis qui vidérunt, et adhuc óderant, ad fidem, quæ per dilectiónem operátur, sui manifestatióne convértit.\par
\switchcolumn
\EN
\noindent Then saith the Lord, as though in continuation “But when the Comforter is come, Whom I will send unto you from the Father, even the Spirit of truth, Which proceedeth from the Father, He shall testify of Me. And ye also shall bear witness, because ye have been with Me from the beginning.” What connection hath this with the words: “Now have they both seen and hated both Me and My Father but that the word might be fulfilled that is written in their law They hated Me without a cause.” Is it that when the Comforter is come, even the Spirit of truth, He will confound by irrefragable testimony them who have both seen and hated both God the Son and God the Father? Yea, indeed, some there were who had seen and still hated, whom the testimony of the Comforter converted to the faith which worketh by love.\par
\switchcolumn*

\LA
\TeDeum{Te Deum}
\switchcolumn
\EN
\TeDeum{Te Deum}
\switchcolumn*

\end{paracol}

\Office{Feria II infra Octavam Ascensionis}{Monday within the Octave of the Ascension}{Semiduplex}
\addcontentsline{toc}{subsection}{Feria II infra Octavam Ascensionis\enspace\textit{Monday within the Octave of the Ascension}}
\begin{paracol}{2}
\LA
\LessonHead{Lectio I}
\switchcolumn
\EN
\LessonHead{Lesson I}
\switchcolumn*

\LA
\LessonTitle{De Epístola prima beáti Joánnis Apóstoli}
\switchcolumn
\EN
\LessonTitle{Lesson from the first letter of St. John the Apostle}
\switchcolumn*

\LA
\Cite{1 Joannes 3:1-6}
\switchcolumn
\EN
\Cite{1 John 3:1-6}
\switchcolumn*

\LA
\noindent Vidéte qualem caritátem dedit nobis Pater, ut fílii Dei nominémur et simus. Propter hoc mundus non novit nos: quia non novit eum. Caríssimi, nunc fílii Dei sumus: et nondum appáruit quid érimus. Scimus quóniam cum apparúerit, símiles ei érimus: quóniam vidébimus eum sícuti est. Et omnis qui habet hanc spem in eo, sanctíficat se, sicut et ille sanctus est. Omnis qui facit peccátum, et iniquitátem facit: et peccátum est iníquitas. Et scitis quia ille appáruit ut peccáta nostra tólleret: et peccátum in eo non est. Omnis qui in eo manet, non peccat: et omnis qui peccat, non vidit eum, nec cognóvit eum.\par
\switchcolumn
\EN
\noindent Behold what manner of charity the Father hath bestowed upon us, that we should be called, and should be the sons of God. Therefore the world knoweth not us, because it knew not him. Dearly beloved, we are now the sons of God; and it hath not yet appeared what we shall be. We know, that, when he shall appear, we shall be like to him: because we shall see him as he is. And every one that hath this hope in him, sanctifieth himself, as he also is holy. Whosoever committeth sin committeth also iniquity; and sin is iniquity. And you know that he appeared to take away our sins, and in him there is no sin. Whosoever abideth in him, sinneth not; and whosoever sinneth, hath not seen him, nor known him.\par
\switchcolumn*

\LA
\begin{Resp}\Rsign Post passiónem suam per dies quadragínta appárens eis, et loquens de regno Dei, allelúja: \Star{} Et, vidéntibus illis, elevátus est, allelúja: et nubes suscépit eum ab óculis eórum, allelúja.\end{Resp}
\begin{Resp}\Vsign Et convéscens, præcépit eis, ab Jerosólymis ne discéderent, sed exspectárent promissiónem Patris.\end{Resp}
\begin{Resp}\Rsign Et, vidéntibus illis, elevátus est, allelúja: et nubes suscépit eum ab óculis eórum, allelúja.\end{Resp}
\switchcolumn
\EN
\begin{Resp}\Rsign Being seen of them forty days after that He had suffered, and speaking of the kingdom of God. Alleluia. \Star{} And while they beheld, He was taken up, and a cloud received Him out of their sight. Alleluia.\end{Resp}
\begin{Resp}\Vsign And, eating together with them, He commanded them that they should not depart from Jerusalem, but wait for the Promise of the Father.\end{Resp}
\begin{Resp}\Rsign And while they beheld, He was taken up, and a cloud received Him out of their sight. Alleluia.\end{Resp}
\switchcolumn*

\LA
\LessonHead{Lectio II}
\switchcolumn
\EN
\LessonHead{Lesson II}
\switchcolumn*

\LA
\Cite{1 Joannes 3:7-12}
\switchcolumn
\EN
\Cite{John 3:7-12}
\switchcolumn*

\LA
\noindent Filíoli, nemo vos sedúcat. Qui facit justítiam, justus est, sicut et ille justus est. Qui facit peccátum, ex diábolo est: quóniam ab inítio diábolus peccat. In hoc appáruit Fílius Dei, ut dissólvat ópera diáboli. Omnis qui natus est ex Deo, peccátum non facit: quóniam semen ipsíus in eo manet, et non potest peccáre, quóniam ex Deo natus est. In hoc manifésti sunt fílii Dei, et fílii diáboli. Omnis qui non est justus, non est ex Deo, et qui non díligit fratrem suum: quóniam hæc est annuntiátio, quam audístis ab inítio, ut diligátis altérutrum. Non sicut Cain, qui ex malígno erat, et occídit fratrem suum. Et propter quid occídit eum? Quóniam ópera ejus malígna erant: fratris autem ejus, justa.\par
\switchcolumn
\EN
\noindent Little children, let no man deceive you. He that doth justice is just, even as he is just. He that committeth sin is of the devil: for the devil sinneth from the beginning. For this purpose, the Son of God appeared, that he might destroy the works of the devil. Whosoever is born of God, committeth not sin: for his seed abideth in him, and he can not sin, because he is born of God. In this the children of God are manifest, and the children of the devil. Whosoever is not just, is not of God, nor he that loveth not his brother. For this is the declaration, which you have heard from the beginning, that you should love one another. Not as Cain, who was of the wicked one, and killed his brother. And wherefore did he kill him? Because his own works were wicked: and his brother's just\par
\switchcolumn*

\LA
\begin{Resp}\Rsign Omnis pulchritúdo Dómini exaltáta est super sídera: \Star{} Spécies ejus in núbibus cæli, et nomen ejus in ætérnum pérmanet, allelúja.\end{Resp}
\begin{Resp}\Vsign A summo cælo egréssio ejus, et occúrsus ejus usque ad summum ejus.\end{Resp}
\begin{Resp}\Rsign Spécies ejus in núbibus cæli, et nomen ejus in ætérnum pérmanet, allelúja.\end{Resp}
\switchcolumn
\EN
\begin{Resp}\Rsign The Lord hath set His beauty above the stars; \Star{} His loveliness is in the clouds of heaven, and His Name endureth for ever. Alleluia.\end{Resp}
\begin{Resp}\Vsign His going forth is from the end of the heaven, and His circuit unto the ends of it.\end{Resp}
\begin{Resp}\Rsign His loveliness is in the clouds of heaven, and His Name endureth for ever. Alleluia.\end{Resp}
\switchcolumn*

\LA
\LessonHead{Lectio III}
\switchcolumn
\EN
\LessonHead{Lesson III}
\switchcolumn*

\LA
\Cite{1 Joannes 3:13-18}
\switchcolumn
\EN
\Cite{1 John 3:13-18}
\switchcolumn*

\LA
\noindent Nolíte mirári, fratres, si odit vos mundus. Nos scimus quóniam transláti sumus de morte ad vitam, quóniam dilígimus fratres. Qui non díligit, manet in morte: omnis qui odit fratrem suum, homicída est. Et scitis quóniam omnis homicída non habet vitam ætérnam in semetípso manéntem. In hoc cognóvimus caritátem Dei, quóniam ille ánimam suam pro nobis pósuit: et nos debémus pro frátribus ánimas pónere. Qui habúerit substántiam hujus mundi, et víderit fratrem suum necessitátem habére, et cláuserit víscera sua ab eo: quómodo cáritas Dei manet in eo? Filíoli mei, non diligámus verbo neque lingua, sed ópere et veritáte:\par
\switchcolumn
\EN
\noindent Wonder not, brethren, if the world hate you. We know that we have passed from death to life, because we love the brethren. He that loveth not, abideth in death. Whosoever hateth his brother is a murderer. And you know that no murderer hath eternal life abiding in himself. In this we have known the charity of God, because he hath laid down his life for us: and we ought to lay down our lives for the brethren. He that hath the substance of this world, and shall see his brother in need, and shall shut up his bowels from him: how doth the charity of God abide in him? My little children, let us not love in word, nor in tongue, but in deed, and in truth.\par
\switchcolumn*

\LA
\begin{Resp}\Rsign Exaltáre, Dómine, allelúja, \Star{} In virtúte tua, allelúja.\end{Resp}
\begin{Resp}\Vsign Eleváta est, magnificéntia tua super cælos, Deus.\end{Resp}
\begin{Resp}\Rsign In virtúte tua, allelúja.\end{Resp}
\begin{Resp}\Vsign Glória Patri, et Fílio, \Star{} et Spirítui Sancto.\end{Resp}
\begin{Resp}\Rsign In virtúte tua, allelúja.\end{Resp}
\switchcolumn
\EN
\begin{Resp}\Rsign Be Thou exalted, O Lord. Alleluia. \Star{} In thine Own strength. Alleluia.\end{Resp}
\begin{Resp}\Vsign O God, Thou hast set thy glory above the heavens.\end{Resp}
\begin{Resp}\Rsign In thine Own strength. Alleluia.\end{Resp}
\begin{Resp}\Vsign Glory be to the Father, and to the Son, \Star{} and to the Holy Ghost.\end{Resp}
\begin{Resp}\Rsign In thine Own strength. Alleluia.\end{Resp}
\switchcolumn*

\LA
\LessonHead{Lectio IV}
\switchcolumn
\EN
\LessonHead{Lesson IV}
\switchcolumn*

\LA
\LessonTitle{Sermo sancti Joánnis Chrysóstomi}
\switchcolumn
\EN
\LessonTitle{From the Sermons of St. John Chrysostom, Patriarch of Constantinople.}
\switchcolumn*

\LA
\Source{Sermo de Ascensione Domini, tom. 3}
\switchcolumn
\EN
\Source{On the Ascension, tom 3}
\switchcolumn*

\LA
\noindent Christus ascéndens in cælum, nostræ natúræ primítias óbtulit Patri, et oblátum donum mirátus est Pater, quod et tanta dígnitas offerébat; et quod offerebátur, nulla mácula fœdabátur. Nam et suis mánibus suscépit oblátum, et suæ sedis fecit esse partícipem, et quod plus est, ad partem suæ déxteræ collocávit. Cognoscámus, quis est ille qui audívit, Sede ad déxteram meam: quæ natúra est, cui Deus dixit, Esto meæ párticeps sedis. Illa natúra est, quæ audívit: Terra es, et in terram ibis.\par
\switchcolumn
\EN
\noindent Then Christ went up into heaven, He offered unto the Father the First-fruits of our nature, and the Father marvelled at the offering, seeing the Majesty of the Priest and the Spotlessness of the oblation. He received the Sacrifice into His Own hands, He made It to sit upon His Throne, nay, more, He gave It a place at His Own Right Hand. Let us ask what nature was His Who heard the words: “Sit Thou at My right hand,” (Ps. cix. 1,) what nature was His to Whom God said “Be Thou Partaker of My Throne?” It was the same nature as was his who heard the sentence “Dust thou art, and unto dust shalt thou return.” (Gen. iii. 19.) Archangels beheld our nature upon the Throne of the Lord, refulgent with eternal glory.\par
\switchcolumn*

\LA
\begin{Resp}\Rsign Tempus est, ut revértar ad eum, qui me misit, dicit Dóminus: nolíte contristári, nec turbétur cor vestrum: \Star{} Rogo pro vobis Patrem, ut ipse vos custódiat, allelúja, allelúja.\end{Resp}
\begin{Resp}\Vsign Nisi ego abíero, Paráclitus non véniet: cum assúmptus fúero, mittam vobis eum.\end{Resp}
\begin{Resp}\Rsign Rogo pro vobis Patrem, ut ipse vos custódiat, allelúja, allelúja.\end{Resp}
\switchcolumn
\EN
\begin{Resp}\Rsign My time is come that I should return unto Him That sent Me, saith the Lord. Be not sorrowful, neither let your heart be troubled. \Star{} I pray the Father for you, that He may keep you. Alleluia, Alleluia.\end{Resp}
\begin{Resp}\Vsign If I go not away, the Comforter will not come unto you; where I am ascended, I will send Him unto you.\end{Resp}
\begin{Resp}\Rsign I pray the Father for you, that He may keep you. Alleluia, Alleluia.\end{Resp}
\switchcolumn*

\LA
\LessonHead{Lectio V}
\switchcolumn
\EN
\LessonHead{Lesson V}
\switchcolumn*

\LA
\noindent Non enim ad omnem glóriam cælos transísse suffécerat, non cum Angelis stare: sed cælos transívit, supra Chérubim ascéndit, ultra Séraphim elevátur, nec ante stetit, nisi sedem Domínicam meruísset. Vide quo spátio cælum separátur a terra, immo terra quanto ab ínferis abest, et ipsum cælum quanto ab altióre cælo separátur, et de altióre cælo ad Angelos quantum spátii est, ad superióres étiam Potestátes, ad ipsam quoque Domínicam sedem. Super hæc ómnia natúra nostra eleváta est, ut homo, qui loco tam húmili tenebátur, ut descéndere non posset ultérius, ad tam excélsam sedem elevarétur, ut áltius non posset ascéndere.\par
\switchcolumn
\EN
\noindent It was not enough of glory for Him to be exalted above the heavens, nor to be ranked with angels but He was exalted above the heavens, He went up above the Cherubim, He ascended beyond the Seraphim, neither found He His rank beneath the Throne of the Lord of lords. Behold how high the heaven is above the earth, and the earth above hell, how high above the heaven is the heaven of heavens, how high above the heaven of heavens the Angels, above the Angels the Higher Powers, and above the Higher Powers the Throne of the Lord. Above all these hath One of our nature been exalted, so that man, which had fallen so low that there was no farther fall for him, is now in place so high, that there is thence no ascending.\par
\switchcolumn*

\LA
\begin{Resp}\Rsign Non turbétur cor vestrum: ego vado ad Patrem; et cum assúmptus fúero a vobis, mittam vobis, allelúja, \Star{} Spíritum veritátis, et gaudébit cor vestrum, allelúja.\end{Resp}
\begin{Resp}\Vsign Ego rogábo Patrem, et álium Paráclitum dabit vobis.\end{Resp}
\begin{Resp}\Rsign Spíritum veritátis, et gaudébit cor vestrum, allelúja.\end{Resp}
\switchcolumn
\EN
\begin{Resp}\Rsign Let not your heart be troubled; I go unto the Father, and when I am taken from you, I will send unto you, alleluia, \Star{} The Spirit of truth, and your heart shall rejoice. Alleluia.\end{Resp}
\begin{Resp}\Vsign I will pray the Father, and He shall give you another Comforter.\end{Resp}
\begin{Resp}\Rsign The Spirit of truth, and your heart shall rejoice. Alleluia.\end{Resp}
\switchcolumn*

\LA
\LessonHead{Lectio VI}
\switchcolumn
\EN
\LessonHead{Lesson VI}
\switchcolumn*

\LA
\noindent Et hæc osténdens Paulus dicébat: Qui descéndit, ipse est qui ascéndit. Et íterum: Descéndit ad inferióra terræ, et ascéndit super omnes cælos. Díscite ígitur quisnam ascéndit, et quæ natúra eleváta est. In hoc enim cúpio remorári sermóne, ut humáni géneris commemoratióne, divínam cleméntiam cum omni admiratióne discámus, quæ summum honórem, magnámque glóriam nostræ natúræ largíta est, quæ ómnibus hodiérna die méruit excélsior reperíri. Hódie Angeli atque Archángeli natúram nostram in sede Domínica immortáli glória fulgéntem vidérunt.\par
\switchcolumn
\EN
\noindent Paul also, dwelling on this, saith: “He That descended is the Same also That ascended up far above all heavens,” even as he had said: “Now, that He ascended, what is it but that He also descended first into the lower parts of the earth.” (Eph. iv. 9, 10.) Learn hence Who it was That ascended, and with what nature He was exalted. And with this thought I wish to bring my sermon to an end. From the thought of that glorified Manhood let us learn with amazement what the goodness of God is that goodness which hath crowned with an honour, higher than which is none, and a glory, greater than which is none, a Person Sharer of our nature, even That Person Which this day hath taken the place which is His of right, above all things other than Himself.\par
\switchcolumn*

\LA
\begin{Resp}\Rsign Ascéndens Christus in altum, captívam duxit captivitátem, \Star{} Dedit dona homínibus, allelúja, allelúja, allelúja.\end{Resp}
\begin{Resp}\Vsign Ascéndit Deus in jubilatióne, et Dóminus in voce tubæ.\end{Resp}
\begin{Resp}\Rsign Dedit dona homínibus, allelúja, allelúja, allelúja.\end{Resp}
\begin{Resp}\Vsign Glória Patri, et Fílio, \Star{} et Spirítui Sancto.\end{Resp}
\begin{Resp}\Rsign Dedit dona homínibus, allelúja, allelúja, allelúja.\end{Resp}
\switchcolumn
\EN
\begin{Resp}\Rsign When Christ ascended up on high, He led captivity captive, \Star{} He gave gifts unto men. Alleluia, Alleluia, Alleluia.\end{Resp}
\begin{Resp}\Vsign God is gone up with a shout, and the Lord with the sound of a trumpet.\end{Resp}
\begin{Resp}\Rsign He gave gifts unto men. Alleluia, Alleluia, Alleluia.\end{Resp}
\begin{Resp}\Vsign Glory be to the Father, and to the Son, \Star{} and to the Holy Ghost.\end{Resp}
\begin{Resp}\Rsign He gave gifts unto men. Alleluia, Alleluia, Alleluia.\end{Resp}
\switchcolumn*

\LA
\LessonHead{Lectio VII}
\switchcolumn
\EN
\LessonHead{Lesson VII}
\switchcolumn*

\LA
\LessonTitle{Léctio sancti Evangélii secúndum Marcum}
\switchcolumn
\EN
\LessonTitle{From the Holy Gospel according to Mark}
\switchcolumn*

\LA
\Cite{Marc 16:14-20}
\switchcolumn
\EN
\Cite{Mark 16:14-20}
\switchcolumn*

\LA
\noindent In illo témpore: Recumbéntibus úndecim discípulis, appáruit illis Jesus: et exprobrávit incredulitátem eórum et durítiam cordis, quia iis, qui víderant eum resurrexísse, non credidérunt. Et réliqua.\par
\switchcolumn
\EN
\noindent At that time, Jesus appeared unto the eleven disciples as they sat at meat, and upbraided them with their unbelief and hardness of heart because they believed not them which had seen Him after He was risen. And so on.\par
\switchcolumn*

\LA
\LessonTitle{Homilía sancti Gregórii Papæ}
\switchcolumn
\EN
\LessonTitle{Homily by Pope St. Gregory the Great.}
\switchcolumn*

\LA
\Source{Ex eadem Homilía 29}
\switchcolumn
\EN
\Source{Same as before.}
\switchcolumn*

\LA
\noindent Et Dóminus quidem Jesus, postquam locútus est eis, assúmptus est in cælum, et sedet a dextris Dei. In véteri Testaménto cognóvimus, quod Elias sit raptus in cælum. Sed áliud est cælum ǽreum, áliud æthéreum. Cælum quippe ǽreum terræ est próximum: unde et aves cæli dícimus, quia eas volitáre in aëre vidémus. In cælum ítaque ǽreum Elias sublevátus est, ut in secrétam quamdam terræ regiónem repénte ducerétur, ubi in magna jam carnis et spíritus quiéte víveret, quoúsque ad finem mundi rédeat, et mortis débitum solvat. Ille étenim mortem dístulit, non evásit: Redémptor autem noster, quia non dístulit, superávit, eámque resurgéndo consúmpsit, et resurrectiónis suæ glóriam ascendéndo declarávit.\par
\switchcolumn
\EN
\noindent “So then, after the Lord Jesus had spoken unto them, He was received up into heaven, and sat on the right hand of God.” We learn in the Old Testament, (4 Kings ii.,) that Elijah was taken up into heaven. But this word “heaven” may mean either the terrestrial atmosphere, or the space external to the sphere of this planet. Of these the atmosphere closely surrounds the earth, and we call the birds “the fowls of the heaven,” because we see them fly therein. It was only up into this that Elijah was taken, that he might be carried off suddenly into some part of the earth, to us unknown, and there live in profound peace of body and soul, until the end of the world, when he will return and pay the debt of nature. For him, therefore, death waiteth, but is not escaped. But our Redeemer made it not to wait for Him, but conquered it, and by rising again shattered it, and by His Ascension showed forth the glory of His Again-rising.\par
\switchcolumn*

\LA
\begin{Resp}\Rsign Ego rogábo Patrem, et álium Paráclitum dabit vobis, \Star{} Ut máneat vobíscum in ætérnum, Spíritum veritátis, allelúja.\end{Resp}
\begin{Resp}\Vsign Si enim non abíero, Paráclitus non véniet ad vos: si autem abíero, mittam eum ad vos.\end{Resp}
\begin{Resp}\Rsign Ut máneat vobíscum in ætérnum, Spíritum veritátis, allelúja.\end{Resp}
\switchcolumn
\EN
\begin{Resp}\Rsign I will pray the Father, and He shall give you another Comforter, \Star{} That He may abide with you for ever, even the Spirit of truth. Alleluia.\end{Resp}
\begin{Resp}\Vsign For if I go not away, the Comforter will not come unto you; but if I depart, I will send Him unto you.\end{Resp}
\begin{Resp}\Rsign That He may abide with you for ever, even the Spirit of truth. Alleluia.\end{Resp}
\switchcolumn*

\LA
\LessonHead{Lectio VIII}
\switchcolumn
\EN
\LessonHead{Lesson VIII}
\switchcolumn*

\LA
\noindent Notándum quoque est, quod Elias in curru légitur ascendísse: ut vidélicet apérte demonstrarétur, quia homo purus adjutório indigébat aliéno. Per Angelos quippe facta illa et osténsa sunt adjuménta: quia nec in cælum quidem aéreum per se ascéndere póterat, quem natúræ suæ infírmitas gravábat. Redémptor autem noster non curru, non Angelis sublevátus légitur: quia is, qui fécerat ómnia, nimírum super ómnia sua virtúte ferebátur. Illo étenim revertebátur, ubi erat: et inde redíbat, ubi remanébat: quia cum per humanitátem ascénderet in cælum, per divinitátem suam et terram páriter continébat et cælum.\par
\switchcolumn
\EN
\noindent We must mark also, how that Elijah was taken up in a chariot, as though to show plainly that for a mere man some outward help was needful. This help was given to him by Angels, as plainly appeareth, since it was impossible for one whom a weak nature yet weighed down earthward, to fly up even into the atmosphere. But of our Redeemer we read not that He was borne up in a chariot, or by Angels, since He by Whom all things were made, clearly rose above all things by His Own Power. He returned unto Him with Whom He was, and whither He returned, there He abode, for albeit as touching His Manhood He ascended up into heaven, yet, as touching His Godhead, He still comprehended both heaven and earth.\par
\switchcolumn*

\LA
\begin{Resp}\Rsign Ponis nubem ascénsum tuum, Dómine: \Star{} Qui ámbulas super pennas ventórum, allelúja.\end{Resp}
\begin{Resp}\Vsign Confessiónem et decórem induísti, amíctus lumen sicut vestiméntum.\end{Resp}
\begin{Resp}\Rsign Qui ámbulas super pennas ventórum, allelúja.\end{Resp}
\begin{Resp}\Vsign Glória Patri, et Fílio, \Star{} et Spirítui Sancto.\end{Resp}
\begin{Resp}\Rsign Qui ámbulas super pennas ventórum, allelúja.\end{Resp}
\switchcolumn
\EN
\begin{Resp}\Rsign Thou makest the clouds the chariot, O Lord, \Star{} Thou walkest upon the wings of the wind. Alleluia.\end{Resp}
\begin{Resp}\Vsign Thou art clothed with honour and majesty, covering thyself with light as with a garment;\end{Resp}
\begin{Resp}\Rsign Thou walkest upon the wings of the wind. Alleluia.\end{Resp}
\begin{Resp}\Vsign Glory be to the Father, and to the Son, \Star{} and to the Holy Ghost.\end{Resp}
\begin{Resp}\Rsign Thou walkest upon the wings of the wind. Alleluia.\end{Resp}
\switchcolumn*

\LA
\LessonHead{Lectio IX}
\switchcolumn
\EN
\LessonHead{Lesson IX}
\switchcolumn*

\LA
\noindent Sicut autem Joseph a frátribus vénditus venditiónem Redemptóris nostri figurávit: sic Henoch translátus, atque ad cælum aéreum Elias sublevátus, ascensiónem Domínicam utérque designávit. Ascensiónis ergo suæ Dóminus prænúntios et testes hábuit, unum ante legem, álium sub lege: ut quandóque veníret ipse, qui veráciter cælos penetráre potuísset. Unde et ipse ordo in eórum quoque utrorúmque sublevatióne per quædam increménta distínguitur. Nam Henoch translátus, Elías vero ad cælum subvéctus esse memorátur: ut veníret póstmodum, qui nec translátus nec subvéctus, cælum æthéreum sua virtúte penetráret.\par
\switchcolumn
\EN
\noindent But as the sale of Joseph by his brethren was a type of the sale of Christ, so were the translations of Enoch and Elijah types of His Ascension. The Lord therefore had had forerunners and witnesses of His Ascension, the one before the Law, the other under the Law, that Himself might one day come, Who was able indeed to pass into the heavens. Hence also there is some difference to be observed in the manner wherein each was translated. Enoch was seen no more, (Gen. v. 24,) for God took him Elijah was carried up by a whirlwind into heaven He That came after them was not taken up, nor carried up, but went up through space by His Own Power.\par
\switchcolumn*

\LA
\TeDeum{Te Deum}
\switchcolumn
\EN
\TeDeum{Te Deum}
\switchcolumn*

\end{paracol}

\Office{Feria III infra Octavam Ascensionis}{Tuesday within the Octave of the Ascension}{Semiduplex}
\addcontentsline{toc}{subsection}{Feria III infra Octavam Ascensionis\enspace\textit{Tuesday within the Octave of the Ascension}}
\begin{paracol}{2}
\LA
\LessonHead{Lectio I}
\switchcolumn
\EN
\LessonHead{Lesson I}
\switchcolumn*

\LA
\LessonTitle{De Epístola prima beáti Joánnis Apóstoli}
\switchcolumn
\EN
\LessonTitle{Lesson from the first letter of St. John the Apostle}
\switchcolumn*

\LA
\Cite{1 Joannes 4:1-6}
\switchcolumn
\EN
\Cite{1 John 4:1-6}
\switchcolumn*

\LA
\noindent Caríssimi, nolíte omni spirítui crédere, sed probáte spíritus si ex Deo sint: quóniam multi pseudoprophétæ exiérunt in mundum. In hoc cognóscitur Spíritus Dei: omnis spíritus qui confitétur Jesum Christum in carne venísse, ex Deo est: et omnis spíritus qui solvit Jesum, ex Deo non est, et hic est Antichrístus, de quo audístis quóniam venit, et nunc jam in mundo est. Vos ex Deo estis filíoli, et vicístis eum, quóniam major est qui in vobis est, quam qui in mundo. Ipsi de mundo sunt: ídeo de mundo loquúntur, et mundus eos audit. Nos ex Deo sumus. Qui novit Deum, audit nos; qui non est ex Deo, non audit nos: in hoc cognóscimus Spíritum veritátis, et spíritum erróris.\par
\switchcolumn
\EN
\noindent Dearly beloved, believe not every spirit, but try the spirits if they be of God: because many false prophets are gone out into the world. By this is the spirit of God known. Every spirit which confesseth that Jesus Christ is come in the flesh, is of God: And every spirit that dissolveth Jesus, is not of God: and this is Antichrist, of whom you have heard that he cometh, and he is now already in the world. You are of God, little children, and have overcome him. Because greater is he that is in you, than he that is in the world. They are of the world: therefore of the world they speak, and the world heareth them. We are of God. He that knoweth God, heareth us. He that is not of God, heareth us not. By this we know the spirit of truth, and the spirit of error.\par
\switchcolumn*

\LA
\begin{Resp}\Rsign Post passiónem suam per dies quadragínta appárens eis, et loquens de regno Dei, allelúja: \Star{} Et, vidéntibus illis, elevátus est, allelúja: et nubes suscépit eum ab óculis eórum, allelúja.\end{Resp}
\begin{Resp}\Vsign Et convéscens, præcépit eis, ab Jerosólymis ne discéderent, sed exspectárent promissiónem Patris.\end{Resp}
\begin{Resp}\Rsign Et, vidéntibus illis, elevátus est, allelúja: et nubes suscépit eum ab óculis eórum, allelúja.\end{Resp}
\switchcolumn
\EN
\begin{Resp}\Rsign Being seen of them forty days after that He had suffered, and speaking of the kingdom of God. Alleluia. \Star{} And while they beheld, He was taken up, and a cloud received Him out of their sight. Alleluia.\end{Resp}
\begin{Resp}\Vsign And, eating together with them, He commanded them that they should not depart from Jerusalem, but wait for the Promise of the Father.\end{Resp}
\begin{Resp}\Rsign And while they beheld, He was taken up, and a cloud received Him out of their sight. Alleluia.\end{Resp}
\switchcolumn*

\LA
\LessonHead{Lectio II}
\switchcolumn
\EN
\LessonHead{Lesson II}
\switchcolumn*

\LA
\Cite{1 Joannes 4:7-14}
\switchcolumn
\EN
\Cite{1 John 4:7-14}
\switchcolumn*

\LA
\noindent Caríssimi, diligámus nos ínvicem: quia cáritas ex Deo est. Et omnis qui díligit, ex Deo natus est, et cognóscit Deum. Qui non díligit, non novit Deum: quóniam Deus cáritas est. In hoc appáruit cáritas Dei in nobis, quóniam Fílium suum unigénitum misit Deus in mundum, ut vivámus per eum. In hoc est cáritas: non quasi nos dilexérimus Deum, sed quóniam ipse prior diléxit nos, et misit Fílium suum propitiatiónem pro peccátis nostris. Caríssimi, si sic Deus diléxit nos: et nos debémus altérutrum dilígere. Deum nemo vidit umquam. Si diligámus ínvicem, Deus in nobis manet, et cáritas ejus in nobis perfécta est. In hoc cognóscimus quóniam in eo manémus, et ipse in nobis: quóniam de Spíritu suo dedit nobis. Et nos vídimus, et testificámur quóniam Pater misit Fílium suum Salvatórem mundi.\par
\switchcolumn
\EN
\noindent Dearly beloved, let us love one another, for charity is of God. And every one that loveth, is born of God, and knoweth God. He that loveth not, knoweth not God: for God is charity. By this hath the charity of God appeared towards us, because God hath sent his only begotten Son into the world, that we may live by him. In this is charity: not as though we had loved God, but because he hath first loved us, and sent his Son to be a propitiation for our sins. My dearest, if God hath so loved us; we also ought to love one another. No man hath seen God at any time. If we love one another, God abideth in us, and his charity is perfected in us. In this we know that we abide in him, and he in us: because he hath given us of his spirit. And we have seen, and do testify, that the Father hath sent his Son to be the Saviour of the world.\par
\switchcolumn*

\LA
\begin{Resp}\Rsign Omnis pulchritúdo Dómini exaltáta est super sídera: \Star{} Spécies ejus in núbibus cæli, et nomen ejus in ætérnum pérmanet, allelúja.\end{Resp}
\begin{Resp}\Vsign A summo cælo egréssio ejus, et occúrsus ejus usque ad summum ejus.\end{Resp}
\begin{Resp}\Rsign Spécies ejus in núbibus cæli, et nomen ejus in ætérnum pérmanet, allelúja.\end{Resp}
\switchcolumn
\EN
\begin{Resp}\Rsign The Lord hath set His beauty above the stars; \Star{} His loveliness is in the clouds of heaven, and His Name endureth for ever. Alleluia.\end{Resp}
\begin{Resp}\Vsign His going forth is from the end of the heaven, and His circuit unto the ends of it.\end{Resp}
\begin{Resp}\Rsign His loveliness is in the clouds of heaven, and His Name endureth for ever. Alleluia.\end{Resp}
\switchcolumn*

\LA
\LessonHead{Lectio III}
\switchcolumn
\EN
\LessonHead{Lesson III}
\switchcolumn*

\LA
\Cite{1 Joannes 4:15-21}
\switchcolumn
\EN
\Cite{1 John 4:15-21}
\switchcolumn*

\LA
\noindent Quisquis conféssus fúerit quóniam Jesus est Fílius Dei, Deus in eo manet, et ipse in Deo. Et nos cognóvimus, et credídimus caritáti, quam habet Deus in nobis. Deus cáritas est: et qui manet in caritáte, in Deo manet, et Deus in eo. In hoc perfécta est cáritas Dei nobíscum, ut fidúciam habeámus in die judícii: quia sicut ille est, et nos sumus in hoc mundo. Timor non est in caritáte: sed perfécta cáritas foras mittit timórem, quóniam timor pœnam habet: qui autem timet, non est perféctus in caritáte. Nos ergo diligámus Deum, quóniam Deus prior diléxit nos. Si quis díxerit: Quóniam díligo Deum, et fratrem suum óderit, mendax est. Qui enim non díligit fratrem suum quem vidit, Deum, quem non vidit, quómodo potest dilígere? Et hoc mandátum habémus a Deo: ut qui díligit Deum, díligat et fratrem suum.\par
\switchcolumn
\EN
\noindent Whosoever shall confess that Jesus is the Son of God, God abideth in him, and he in God. And we have known, and have believed the charity, which God hath to us. God is charity: and he that abideth in charity, abideth in God, and God in him. In this is the charity of God perfected with us, that we may have confidence in the day of judgment: because as he is, we also are in this world. Fear is not in charity: but perfect charity casteth out fear, because fear hath pain. And he that feareth, is not perfected in charity. Let us therefore love God, because God first hath loved us. If any man say, I love God, and hateth his brother; he is a liar. For he that loveth not his brother, whom he seeth, how can he love God, whom he seeth not? And this commandment we have from God, that he, who loveth God, love also his brother.\par
\switchcolumn*

\LA
\begin{Resp}\Rsign Exaltáre, Dómine, allelúja, \Star{} In virtúte tua, allelúja.\end{Resp}
\begin{Resp}\Vsign Eleváta est, magnificéntia tua super cælos, Deus.\end{Resp}
\begin{Resp}\Rsign In virtúte tua, allelúja.\end{Resp}
\begin{Resp}\Vsign Glória Patri, et Fílio, \Star{} et Spirítui Sancto.\end{Resp}
\begin{Resp}\Rsign In virtúte tua, allelúja.\end{Resp}
\switchcolumn
\EN
\begin{Resp}\Rsign Be Thou exalted, O Lord. Alleluia. \Star{} In thine Own strength. Alleluia.\end{Resp}
\begin{Resp}\Vsign O God, Thou hast set thy glory above the heavens.\end{Resp}
\begin{Resp}\Rsign In thine Own strength. Alleluia.\end{Resp}
\begin{Resp}\Vsign Glory be to the Father, and to the Son, \Star{} and to the Holy Ghost.\end{Resp}
\begin{Resp}\Rsign In thine Own strength. Alleluia.\end{Resp}
\switchcolumn*

\LA
\LessonHead{Lectio IV}
\switchcolumn
\EN
\LessonHead{Lesson IV}
\switchcolumn*

\LA
\LessonTitle{Sermo sancti Máximi Epíscopi}
\switchcolumn
\EN
\LessonTitle{From the Sermons of St. Maximus, Bishop of Turin.}
\switchcolumn*

\LA
\Source{Homilia 43, quæ est 2 de Pentecoste, ante medium}
\switchcolumn
\EN
\Source{43rd, 2nd on Pentecost.}
\switchcolumn*

\LA
\noindent Méminit sánctitas vestra, quod áquilæ illi de Psaltério, cujus innovátam juventútem légimus, comparáverim Salvatórem. Est enim similitúdo non parva. Sicut enim áquila humília déserit, alta petit, cælórum vicína conscéndit: ita et Salvátor humília inférni deséruit, paradísi altióra pétiit, cælórum fastígia penetrávit. Et sicut áquila relíctis terrénis sórdibus, sublíme volans, purióris áëris salubritáte perfrúitur: ita et Dóminus, terrenórum fæcem déserens peccatórum, in Sanctis suis vólitans, purióris vitæ simplicitáte lætátur.\par
\switchcolumn
\EN
\noindent My holy brethren, ye remember that I have likened the Saviour to that eagle, touching which it is written in the Book of Psalms, cii.\par
5, “thy youth is renewed like the eagle's.” There are many points of likeness. The eagle riseth above ground, wingeth his way aloft, and mounteth skyward even so did the Saviour rise from the depth of the grave, mount up unto the exalted mansions of Paradise, and enter the heights of heaven. The eagle leaveth below him the foul mists of earth, flieth above, and drinketh in health from a purer air even so did the Lord leave below Him the filthy slough of sinners on earth, and rejoice Himself with the honesty of a purer life, when He soared again into His Own holy home.\par
\switchcolumn*

\LA
\begin{Resp}\Rsign Tempus est, ut revértar ad eum, qui me misit, dicit Dóminus: nolíte contristári, nec turbétur cor vestrum: \Star{} Rogo pro vobis Patrem, ut ipse vos custódiat, allelúja, allelúja.\end{Resp}
\begin{Resp}\Vsign Nisi ego abíero, Paráclitus non véniet: cum assúmptus fúero, mittam vobis eum.\end{Resp}
\begin{Resp}\Rsign Rogo pro vobis Patrem, ut ipse vos custódiat, allelúja, allelúja.\end{Resp}
\switchcolumn
\EN
\begin{Resp}\Rsign My time is come that I should return unto Him That sent Me, saith the Lord. Be not sorrowful, neither let your heart be troubled. \Star{} I pray the Father for you, that He may keep you. Alleluia, Alleluia.\end{Resp}
\begin{Resp}\Vsign If I go not away, the Comforter will not come unto you; where I am ascended, I will send Him unto you.\end{Resp}
\begin{Resp}\Rsign I pray the Father for you, that He may keep you. Alleluia, Alleluia.\end{Resp}
\switchcolumn*

\LA
\LessonHead{Lectio V}
\switchcolumn
\EN
\LessonHead{Lesson V}
\switchcolumn*

\LA
\noindent Per ómnia ígitur áquilæ comparátio cónvenit Salvatóri. Sed quid fácimus, quod áquila prædam frequénter díripit, tollit frequénter aliénum? Nec in hoc tamen dissímilis est Salvátor. Prædam enim quodámmodo sústulit, cum hóminem, quem suscépit, inférni raptum fáucibus portávit ad cælum, et aliénæ dominatiónis, id est, diabólicæ potestátis servum de captivitáte érutum, duxit ad altióra captívum, sicut scriptum est in prophéta: Ascéndens in altum, captívam duxit captivitátem, dedit dona homínibus.\par
\switchcolumn
\EN
\noindent In all ways, therefore, is the Saviour aptly likened to an eagle. But what can we make of this, that the eagle is a bird of prey, oft-times a plunderer. Even in this he is like to the Saviour. He bore off His prey, when He carried off from the jaws of hell to heaven the Manhood Which He had swooped to take to Himself, yea, when He led captive to an higher home him whom He had delivered from the mastership of another lord, namely the devil, even as it is written in the Prophet, (Ps. lxvii. 19,) “Thou hast ascended on high, Thou hast led captivity captive, Thou hast received gifts among men.”\par
\switchcolumn*

\LA
\begin{Resp}\Rsign Non turbétur cor vestrum: ego vado ad Patrem; et cum assúmptus fúero a vobis, mittam vobis, allelúja, \Star{} Spíritum veritátis, et gaudébit cor vestrum, allelúja.\end{Resp}
\begin{Resp}\Vsign Ego rogábo Patrem, et álium Paráclitum dabit vobis.\end{Resp}
\begin{Resp}\Rsign Spíritum veritátis, et gaudébit cor vestrum, allelúja.\end{Resp}
\switchcolumn
\EN
\begin{Resp}\Rsign Let not your heart be troubled; I go unto the Father, and when I am taken from you, I will send unto you, alleluia, \Star{} The Spirit of truth, and your heart shall rejoice. Alleluia.\end{Resp}
\begin{Resp}\Vsign I will pray the Father, and He shall give you another Comforter.\end{Resp}
\begin{Resp}\Rsign The Spirit of truth, and your heart shall rejoice. Alleluia.\end{Resp}
\switchcolumn*

\LA
\LessonHead{Lectio VI}
\switchcolumn
\EN
\LessonHead{Lesson VI}
\switchcolumn*

\LA

\switchcolumn
\EN
\LessonTitle{“Thou hast ascended on high, Thou hast led captivity captive.” O how nobly doth the Prophet paint the Triumph of the Lord We hear how that of old time, when kings marched in triumph, the procession of prisoners walked before the chariot of their conqueror. Lo, the Lord entereth the heavens, not after, but amid a most glorious band of captives. That band are not led before His chariot, but themselves bear up their Saviour. In some mystic sense, when the Son of God bore to heaven the Son of man, captivity both led and was led.}
\switchcolumn*

\LA
\noindent Ascéndit, inquit, in altum, captívam duxit captivitátem. Quam bene triúmphum Dómini prophéta descríbit! Solébat, sicut dicunt, regum triumphántium currus captivórum pompa præcédere. Ecce Dóminum eúntem ad cælos non præcédit, sed comitátur gloriósa captívitas; non ante vehículum dúcitur, sed ipsa évehit Salvatórem. Quodam enim mystério, dum Fílius Dei Fílium hóminis sústulit ad cælum, ipsa captívitas portátur, et portat.\par
\switchcolumn
\EN

\switchcolumn*

\LA
\begin{Resp}\Rsign Ascéndens Christus in altum, captívam duxit captivitátem, \Star{} Dedit dona homínibus, allelúja, allelúja, allelúja.\end{Resp}
\begin{Resp}\Vsign Ascéndit Deus in jubilatióne, et Dóminus in voce tubæ.\end{Resp}
\begin{Resp}\Rsign Dedit dona homínibus, allelúja, allelúja, allelúja.\end{Resp}
\begin{Resp}\Vsign Glória Patri, et Fílio, \Star{} et Spirítui Sancto.\end{Resp}
\begin{Resp}\Rsign Dedit dona homínibus, allelúja, allelúja, allelúja.\end{Resp}
\switchcolumn
\EN
\begin{Resp}\Rsign When Christ ascended up on high, He led captivity captive, \Star{} He gave gifts unto men. Alleluia, Alleluia, Alleluia.\end{Resp}
\begin{Resp}\Vsign God is gone up with a shout, and the Lord with the sound of a trumpet.\end{Resp}
\begin{Resp}\Rsign He gave gifts unto men. Alleluia, Alleluia, Alleluia.\end{Resp}
\begin{Resp}\Vsign Glory be to the Father, and to the Son, \Star{} and to the Holy Ghost.\end{Resp}
\begin{Resp}\Rsign He gave gifts unto men. Alleluia, Alleluia, Alleluia.\end{Resp}
\switchcolumn*

\LA
\LessonHead{Lectio VII}
\switchcolumn
\EN
\LessonHead{Lesson VII}
\switchcolumn*

\LA
\LessonTitle{Léctio sancti Evangélii secúndum Marcum}
\switchcolumn
\EN
\LessonTitle{From the Holy Gospel according to Mark}
\switchcolumn*

\LA
\Cite{Marc 16:14-20}
\switchcolumn
\EN
\Cite{Mark 16:14-20}
\switchcolumn*

\LA
\noindent In illo témpore: Recumbéntibus úndecim discípulis, appáruit illis Jesus: et exprobrávit incredulitátem eórum et durítiam cordis, quia iis, qui víderant eum resurrexísse, non credidérunt. Et réliqua.\par
\switchcolumn
\EN
\noindent At that time, Jesus appeared unto the eleven disciples as they sat at meat, and upbraided them with their unbelief and hardness of heart, because they believed not them which had seen Him after He was risen. And so on.\par
\switchcolumn*

\LA
\LessonTitle{Homilía sancti Gregórii Papæ}
\switchcolumn
\EN
\LessonTitle{Homily by Pope St. Gregory the Great.}
\switchcolumn*

\LA
\Source{Ex eadem Homilía 29}
\switchcolumn
\EN
\Source{Same as before.}
\switchcolumn*

\LA
\noindent Considerándum nobis est, quid est quod Marcus ait: Sedet a dextris Dei: et Stéphanus dicit: Vídeo cælos apértos, et Fílium hóminis stantem a dextris Dei. Quid est quod hunc Marcus sedéntem, Stéphanus vero stantem se vidére testátur? Sed scitis, fratres, quia sedére judicántis est, stare vero pugnántis, vel adjuvántis.\par
\switchcolumn
\EN
\noindent We must ponder the meaning of these words of Mark, “He sat on the right hand of God,” and how that Stephen said, (Acts vii. 56,) “Behold, I see the heavens opened, and the Son of man standing on the right hand of God.” Wherefore doth Mark say that He sat, whereas Stephen testifieth that he saw Him standing But ye know, my brethren, that to sit is for him that judgeth, to stand, for him that fighteth, or helpeth.\par
\switchcolumn*

\LA
\begin{Resp}\Rsign Ego rogábo Patrem, et álium Paráclitum dabit vobis, \Star{} Ut máneat vobíscum in ætérnum, Spíritum veritátis, allelúja.\end{Resp}
\begin{Resp}\Vsign Si enim non abíero, Paráclitus non véniet ad vos: si autem abíero, mittam eum ad vos.\end{Resp}
\begin{Resp}\Rsign Ut máneat vobíscum in ætérnum, Spíritum veritátis, allelúja.\end{Resp}
\switchcolumn
\EN
\begin{Resp}\Rsign I will pray the Father, and He shall give you another Comforter, \Star{} That He may abide with you for ever, even the Spirit of truth. Alleluia.\end{Resp}
\begin{Resp}\Vsign For if I go not away, the Comforter will not come unto you; but if I depart, I will send Him unto you.\end{Resp}
\begin{Resp}\Rsign That He may abide with you for ever, even the Spirit of truth. Alleluia.\end{Resp}
\switchcolumn*

\LA
\LessonHead{Lectio VIII}
\switchcolumn
\EN
\LessonHead{Lesson VIII}
\switchcolumn*

\LA
\noindent Quia ergo Redémptor noster assúmptus in cælum et nunc ómnia júdicat, et ad extrémum judex ómnium véniet, hunc post assumptiónem Marcus sedére descríbit, quia post ascensiónis suæ glóriam judex in fine vidébitur. Stéphanus vero hunc in labóre certáminis pósitus stantem vidit, quem adjutórem hábuit: quia ut iste in terra persecutórum infidelitátem vínceret, pro illo de cælo illíus grátia pugnávit.\par
\switchcolumn
\EN
\noindent Since therefore, our Redeemer is ascended up into heaven, and even now is Judge of all, beside that at the end of the world He will so come, therefore doth Mark say that He sitteth where He hath gone up, because we look for Him, after that His glorious Ascension, that He will come again at the end to be our Judge. But Stephen, while yet he was in the throes of the battle, saw Him That was helping him standing. Stephen on earth was overcoming the unbelief of his persecutors, but it was the grace of Him That is in heaven that fought in him all the while.\par
\switchcolumn*

\LA
\begin{Resp}\Rsign Ponis nubem ascénsum tuum, Dómine: \Star{} Qui ámbulas super pennas ventórum, allelúja.\end{Resp}
\begin{Resp}\Vsign Confessiónem et decórem induísti, amíctus lumen sicut vestiméntum.\end{Resp}
\begin{Resp}\Rsign Qui ámbulas super pennas ventórum, allelúja.\end{Resp}
\begin{Resp}\Vsign Glória Patri, et Fílio, \Star{} et Spirítui Sancto.\end{Resp}
\begin{Resp}\Rsign Qui ámbulas super pennas ventórum, allelúja.\end{Resp}
\switchcolumn
\EN
\begin{Resp}\Rsign Thou makest the clouds the chariot, O Lord, \Star{} Thou walkest upon the wings of the wind. Alleluia.\end{Resp}
\begin{Resp}\Vsign Thou art clothed with honour and majesty, covering thyself with light as with a garment;\end{Resp}
\begin{Resp}\Rsign Thou walkest upon the wings of the wind. Alleluia.\end{Resp}
\begin{Resp}\Vsign Glory be to the Father, and to the Son, \Star{} and to the Holy Ghost.\end{Resp}
\begin{Resp}\Rsign Thou walkest upon the wings of the wind. Alleluia.\end{Resp}
\switchcolumn*

\LA
\LessonHead{Lectio IX}
\switchcolumn
\EN
\LessonHead{Lesson IX}
\switchcolumn*

\LA

\switchcolumn
\EN
\LessonTitle{“And they went forth and preached everywhere, the Lord working with them, and confirming the word with signs following.” What are we to see in this, what are we to remember, but that obedience followed commandment, and signs obedience But now, since, by the will of God, we have lightly run over our reading from the Gospel, it remaineth that we should say somewhat by way of reflection on this great Festival.}
\switchcolumn*

\LA
\noindent Séquitur: Illi autem profécti prædicavérunt ubíque, Dómino cooperánte et sermónem confirmánte sequéntibus signis. Quid in his considerándum est, quid memóriæ commendándum: nisi quod præcéptum obediéntia, obediéntiam vero signa secúta sunt? Sed quia auctóre Deo bréviter lectiónem evangélicam exponéndo transcúrrimus: restat, ut áliquid de ipsa tantæ consideratióne solemnitátis dicámus.\par
\switchcolumn
\EN

\switchcolumn*

\LA
\TeDeum{Te Deum}
\switchcolumn
\EN
\TeDeum{Te Deum}
\switchcolumn*

\end{paracol}

\Office{Feria IV infra Octavam Ascensionis}{Wednesday within the Octave of the Ascension}{Semiduplex}
\addcontentsline{toc}{subsection}{Feria IV infra Octavam Ascensionis\enspace\textit{Wednesday within the Octave of the Ascension}}
\begin{paracol}{2}
\LA
\LessonHead{Lectio I}
\switchcolumn
\EN
\LessonHead{Lesson I}
\switchcolumn*

\LA
\LessonTitle{Incipit Epístola secúnda beáti Joánnis Apóstoli}
\switchcolumn
\EN
\LessonTitle{Lesson from the second letter of St. John the Apostle}
\switchcolumn*

\LA
\Cite{2 Joannes 1:1-5}
\switchcolumn
\EN
\Cite{2 John 1:1-5}
\switchcolumn*

\LA
\noindent Sénior Eléctæ dóminæ, et natis ejus, quos ego díligo in veritáte, et non ego solus, sed et omnes qui cognovérunt veritátem, propter veritátem, quæ pérmanet in nobis, et nobíscum erit in ætérnum. Sit vobíscum grátia, misericórdia, pax a Deo Patre, et a Christo Jesu Fílio Patris in veritáte, et caritáte. Gavísus sum valde, quóniam invéni de fíliis tuis ambulántes in veritáte, sicut mandátum accépimus a Patre. Et nunc rogo te dómina, non tamquam mandátum novum scribens tibi, sed quod habúimus ab inítio, ut diligámus altérutrum.\par
\switchcolumn
\EN
\noindent The ancient to the lady Elect, and her children, whom I love in the truth, and not I only, but also all they that have known the truth, For the sake of the truth which dwelleth in us, and shall be with us for ever. Grace be with you, mercy, and peace from God the Father, and from Christ Jesus the Son of the Father; in truth and charity. I was exceeding glad, that I found of thy children walking in truth, as we have received a commandment from the Father. And now I beseech thee, lady, not as writing a new commandment to thee, but that which we have had from the beginning, that we love one another.\par
\switchcolumn*

\LA
\begin{Resp}\Rsign Post passiónem suam per dies quadragínta appárens eis, et loquens de regno Dei, allelúja: \Star{} Et, vidéntibus illis, elevátus est, allelúja: et nubes suscépit eum ab óculis eórum, allelúja.\end{Resp}
\begin{Resp}\Vsign Et convéscens, præcépit eis, ab Jerosólymis ne discéderent, sed exspectárent promissiónem Patris.\end{Resp}
\begin{Resp}\Rsign Et, vidéntibus illis, elevátus est, allelúja: et nubes suscépit eum ab óculis eórum, allelúja.\end{Resp}
\switchcolumn
\EN
\begin{Resp}\Rsign Being seen of them forty days after that He had suffered, and speaking of the kingdom of God. Alleluia. \Star{} And while they beheld, He was taken up, and a cloud received Him out of their sight. Alleluia.\end{Resp}
\begin{Resp}\Vsign And, eating together with them, He commanded them that they should not depart from Jerusalem, but wait for the Promise of the Father.\end{Resp}
\begin{Resp}\Rsign And while they beheld, He was taken up, and a cloud received Him out of their sight. Alleluia.\end{Resp}
\switchcolumn*

\LA
\LessonHead{Lectio II}
\switchcolumn
\EN
\LessonHead{Lesson II}
\switchcolumn*

\LA
\Cite{2 Joannes 1:6-9}
\switchcolumn
\EN
\Cite{2 John 1:6-9}
\switchcolumn*

\LA
\noindent Et hæc est cáritas, ut ambulémus secúndum mandáta ejus. Hoc est enim mandátum, ut quemádmodum audístis ab inítio, in eo ambulétis. Quóniam multi seductóres exiérunt in mundum, qui non confiténtur Jesum Christum venísse in carnem: hic est sedúctor, et Antichrístus. Vidéte vosmetípsos, ne perdátis quæ operáti estis: sed ut mercédem plenam accipiátis. Omnis qui recédit, et non pérmanet in doctrína Christi, Deum non habet: qui pérmanet in doctrína, hic et Patrem et Fílium habet.\par
\switchcolumn
\EN
\noindent And this is charity, that we walk according to his commandments. For this is the commandment, that, as you have heard from the beginning, you should walk in the same: For many seducers are gone out into the world, who confess not that Jesus Christ is come in the flesh: this is a seducer and an antichrist. Look to yourselves, that you lose not the things which you have wrought: but that you may receive a full reward. Whosoever revolteth, and continueth not in the doctrine of Christ, hath not God. He that continueth in the doctrine, the same hath both the Father and the Son.\par
\switchcolumn*

\LA
\begin{Resp}\Rsign Omnis pulchritúdo Dómini exaltáta est super sídera: \Star{} Spécies ejus in núbibus cæli, et nomen ejus in ætérnum pérmanet, allelúja.\end{Resp}
\begin{Resp}\Vsign A summo cælo egréssio ejus, et occúrsus ejus usque ad summum ejus.\end{Resp}
\begin{Resp}\Rsign Spécies ejus in núbibus cæli, et nomen ejus in ætérnum pérmanet, allelúja.\end{Resp}
\switchcolumn
\EN
\begin{Resp}\Rsign The Lord hath set His beauty above the stars; \Star{} His loveliness is in the clouds of heaven, and His Name endureth for ever. Alleluia.\end{Resp}
\begin{Resp}\Vsign His going forth is from the end of the heaven, and His circuit unto the ends of it.\end{Resp}
\begin{Resp}\Rsign His loveliness is in the clouds of heaven, and His Name endureth for ever. Alleluia.\end{Resp}
\switchcolumn*

\LA
\LessonHead{Lectio III}
\switchcolumn
\EN
\LessonHead{Lesson III}
\switchcolumn*

\LA
\Cite{2 Joannes 1:10-13}
\switchcolumn
\EN
\Cite{2 John 1:10-13}
\switchcolumn*

\LA
\noindent Si quis venit ad vos, et hanc doctrínam non affert, nolíte recípere eum in domum, nec Ave ei dixéritis. Qui enim dicit illi Ave, commúnicat opéribus ejus malígnis. Plura habens vobis scríbere, nólui per cartam et atraméntum: spero enim me futúrum apud vos, et os ad os loqui: ut gáudium vestrum plenum sit. Salútant te fílii soróris tuæ Eléctæ.\par
\switchcolumn
\EN
\noindent If any man come to you, and bring not this doctrine, receive him not into the house nor say to him, God speed you. For he that saith unto him, God speed you, communicateth with his wicked works. Having more things to write unto you, I would not by paper and ink: for I hope that I shall be with you, and speak face to face: that your joy may be full. The children of thy sister Elect salute thee.\par
\switchcolumn*

\LA
\begin{Resp}\Rsign Exaltáre, Dómine, allelúja, \Star{} In virtúte tua, allelúja.\end{Resp}
\begin{Resp}\Vsign Eleváta est, magnificéntia tua super cælos, Deus.\end{Resp}
\begin{Resp}\Rsign In virtúte tua, allelúja.\end{Resp}
\begin{Resp}\Vsign Glória Patri, et Fílio, \Star{} et Spirítui Sancto.\end{Resp}
\begin{Resp}\Rsign In virtúte tua, allelúja.\end{Resp}
\switchcolumn
\EN
\begin{Resp}\Rsign Be Thou exalted, O Lord. Alleluia. \Star{} In thine Own strength. Alleluia.\end{Resp}
\begin{Resp}\Vsign O God, Thou hast set thy glory above the heavens.\end{Resp}
\begin{Resp}\Rsign In thine Own strength. Alleluia.\end{Resp}
\begin{Resp}\Vsign Glory be to the Father, and to the Son, \Star{} and to the Holy Ghost.\end{Resp}
\begin{Resp}\Rsign In thine Own strength. Alleluia.\end{Resp}
\switchcolumn*

\LA
\LessonHead{Lectio IV}
\switchcolumn
\EN
\LessonHead{Lesson IV}
\switchcolumn*

\LA
\LessonTitle{Sermo sancti Gregórii Nysséni}
\switchcolumn
\EN
\LessonTitle{From the Sermons of St. Gregory, Bishop of Nyssa.}
\switchcolumn*

\LA
\Source{Oratio de Ascensione Domini}
\switchcolumn
\EN
\Source{Discourse on the Lords Ascension.}
\switchcolumn*

\LA
\noindent Hodiérnam celebritátem satis per se magnam, prophéta David majórem éfficit, dum illi gáudium e Psalmis adjúngit. Hic enim excélsus Prophéta supra seípsum egrédiens, tamquam córporis ónere nihil premátur, infert se cæléstibus potestátibus, et voces eárum nobis expónit, cum in cælum redeúntem Dóminum ipsæ comitántes, Angelis, qui versántur in terris, quibúsque in humánam vitam ingréssus commíssus est, ímperant ad hunc modum: Tóllite portas, príncipes, vestras; et elevámini, portæ æternáles, et introíbit Rex glóriæ.\par
\switchcolumn
\EN
\noindent The very thought of this day's Festival is great enough in itself, but the Prophet David hath much inflamed our joyful enthusiasm by the Psalms. This noble Prophet hath, as it were, gone out of himself, as though the body were a weight duller than his spirit could bear he joineth company with the Powers of heaven, and telleth what they said when they went with the Lord heavenward, and cried in tones of command to those Angels who work on earth, and by whose heralding the Birth of the Incarnate One had been proclaimed “Lift up your gates, O ye princes, and be ye lift up, ye everlasting doors, and the King of glory shall come in.” (Ps. xxiii. 7, 9.)\par
\switchcolumn*

\LA
\begin{Resp}\Rsign Tempus est, ut revértar ad eum, qui me misit, dicit Dóminus: nolíte contristári, nec turbétur cor vestrum: \Star{} Rogo pro vobis Patrem, ut ipse vos custódiat, allelúja, allelúja.\end{Resp}
\begin{Resp}\Vsign Nisi ego abíero, Paráclitus non véniet: cum assúmptus fúero, mittam vobis eum.\end{Resp}
\begin{Resp}\Rsign Rogo pro vobis Patrem, ut ipse vos custódiat, allelúja, allelúja.\end{Resp}
\switchcolumn
\EN
\begin{Resp}\Rsign My time is come that I should return unto Him That sent Me, saith the Lord. Be not sorrowful, neither let your heart be troubled. \Star{} I pray the Father for you, that He may keep you. Alleluia, Alleluia.\end{Resp}
\begin{Resp}\Vsign If I go not away, the Comforter will not come unto you; where I am ascended, I will send Him unto you.\end{Resp}
\begin{Resp}\Rsign I pray the Father for you, that He may keep you. Alleluia, Alleluia.\end{Resp}
\switchcolumn*

\LA
\LessonHead{Lectio V}
\switchcolumn
\EN
\LessonHead{Lesson V}
\switchcolumn*

\LA
\noindent Et quóniam ubicúmque fúerit ille, qui in seípso ómnia cóntinet, pro suscipiéntium captu seípsum dimetítur (neque enim solum inter hómines homo fit, verum étiam dum inter Angelos versátur, ad illórum vocem sese demíttit) idcírco janitóres intérrogant: Quis est iste Rex glóriæ? Respóndent ipsis, demonstrántque fortem et poténtem in prǽlio, qui pugnatúrus erat contra illum, qui natúram humánam in servitúte captívam detinébat, et eversúrus eum, qui mortis habébat impérium: ut gravíssimo hoste superáto, genus hóminum in libertátem et pacem vindicáret.\par
\switchcolumn
\EN
\noindent He, Who containeth all things, is everywhere, but for the sake of them which receive Him, He is pleased to make Himself a local Presence which hath bounds. Not only did He become a Man among men, but when conversing among Angels, He alloweth that title also to be given Him. The gatekeepers therefore ask “Who is this King of glory?” and it is answered them that He is “The Lord, strong and mighty, the Lord, mighty in battle,” the Lord, Whose work it had been to fight him who held mankind in bondage, and to “destroy him that had the power of death, that is, the devil” (Heb. ii. 14) that now that dark enemy was trampled down, and man had had won for him freedom and peace.\par
\switchcolumn*

\LA
\begin{Resp}\Rsign Non turbétur cor vestrum: ego vado ad Patrem; et cum assúmptus fúero a vobis, mittam vobis, allelúja, \Star{} Spíritum veritátis, et gaudébit cor vestrum, allelúja.\end{Resp}
\begin{Resp}\Vsign Ego rogábo Patrem, et álium Paráclitum dabit vobis.\end{Resp}
\begin{Resp}\Rsign Spíritum veritátis, et gaudébit cor vestrum, allelúja.\end{Resp}
\switchcolumn
\EN
\begin{Resp}\Rsign Let not your heart be troubled; I go unto the Father, and when I am taken from you, I will send unto you, alleluia, \Star{} The Spirit of truth, and your heart shall rejoice. Alleluia.\end{Resp}
\begin{Resp}\Vsign I will pray the Father, and He shall give you another Comforter.\end{Resp}
\begin{Resp}\Rsign The Spirit of truth, and your heart shall rejoice. Alleluia.\end{Resp}
\switchcolumn*

\LA
\LessonHead{Lectio VI}
\switchcolumn
\EN
\LessonHead{Lesson VI}
\switchcolumn*

\LA
\noindent Occúrrunt ei custódes, et portas jubent reclúdi, ut in ipsis rursum glóriam assequátur. Verum non agnóscunt eum, qui sórdidam vitæ nostræ stolam indútus est: cujus rubra sunt vestiménta ex humanórum malórum torculári. Itaque rursus cómites ejus vócibus illis interrogántur: Quis est iste Rex glóriæ? Respondétur autem non ámplius, Fortis et potens in prǽlio; sed, Dóminus virtútum, qui mundi principátum obtínuit, qui summátim ómnia in se collégit, qui prístinum in statum cuncta restítuit: ipse est Rex glóriæ.\par
\switchcolumn
\EN
\noindent The keepers run to the gates, and bid the doors unfold, that the Lord may enter in, to take again the glory which He had there among them before. But when they see Him, clad in the likeness of sinful flesh, (Rom.viii. 3,) they know Him not, even Him Who is red in His apparel, because that He hath trodden Alone the winepress of human pain, and the blood is sprinkled upon His garments. Therefore they cry again to their fellows that bear Him company: “Who is this King of glory?” And they answer them no more: “The Lord, strong and mighty, the Lord mighty in battle” but “The Lord of hosts the Lord, Whose Own are become the kingdoms of the world the Lord, Who hath made Himself the Head of all things the Lord, Who hath made all things new (Apoc. xxi. 5.)” He is the King of glory!\par
\switchcolumn*

\LA
\begin{Resp}\Rsign Ascéndens Christus in altum, captívam duxit captivitátem, \Star{} Dedit dona homínibus, allelúja, allelúja, allelúja.\end{Resp}
\begin{Resp}\Vsign Ascéndit Deus in jubilatióne, et Dóminus in voce tubæ.\end{Resp}
\begin{Resp}\Rsign Dedit dona homínibus, allelúja, allelúja, allelúja.\end{Resp}
\begin{Resp}\Vsign Glória Patri, et Fílio, \Star{} et Spirítui Sancto.\end{Resp}
\begin{Resp}\Rsign Dedit dona homínibus, allelúja, allelúja, allelúja.\end{Resp}
\switchcolumn
\EN
\begin{Resp}\Rsign When Christ ascended up on high, He led captivity captive, \Star{} He gave gifts unto men. Alleluia, Alleluia, Alleluia.\end{Resp}
\begin{Resp}\Vsign God is gone up with a shout, and the Lord with the sound of a trumpet.\end{Resp}
\begin{Resp}\Rsign He gave gifts unto men. Alleluia, Alleluia, Alleluia.\end{Resp}
\begin{Resp}\Vsign Glory be to the Father, and to the Son, \Star{} and to the Holy Ghost.\end{Resp}
\begin{Resp}\Rsign He gave gifts unto men. Alleluia, Alleluia, Alleluia.\end{Resp}
\switchcolumn*

\LA
\LessonHead{Lectio VII}
\switchcolumn
\EN
\LessonHead{Lesson VII}
\switchcolumn*

\LA
\LessonTitle{Léctio sancti Evangélii secúndum Marcum}
\switchcolumn
\EN
\LessonTitle{From the Holy Gospel according to Mark}
\switchcolumn*

\LA
\Cite{Marc 16:14-20}
\switchcolumn
\EN
\Cite{Mark 16:14-20}
\switchcolumn*

\LA
\noindent In illo témpore: Recumbéntibus úndecim discípulis, appáruit illis Jesus: et exprobrávit incredulitátem eórum et durítiam cordis, quia iis, qui víderant eum resurrexísse, non credidérunt. Et réliqua.\par
\switchcolumn
\EN
\noindent At that time, Jesus appeared unto the eleven disciples as they sat at meat, and upbraided them with their unbelief and hardness of heart, because they believed not them which had seen Him after He was risen. And so on.\par
\switchcolumn*

\LA
\LessonTitle{Homilía sancti Gregórii Papæ}
\switchcolumn
\EN
\LessonTitle{Homily by Pope St. Gregory the Great.}
\switchcolumn*

\LA
\Source{Ex eadem Homilía 29}
\switchcolumn
\EN
\Source{Same as before.}
\switchcolumn*

\LA
\noindent Hoc autem nobis primum quæréndum est, quidnam sit, quod, nato Dómino, apparuérunt Angeli, et tamen non legúntur in albis véstibus apparuísse: ascendénte autem Dómino, missi Angeli in albis legúntur véstibus apparuísse. Sic étenim scriptum est: Vidéntibus illis elevátus est, et nubes suscépit eum ab óculis eórum. Cumque intueréntur in cælum eúntem illum, ecce duo viri stetérunt juxta illos in véstibus albis. In albis autem véstibus gáudium et solémnitas mentis osténditur. Quid est ergo, quod, nato Dómino, non in albis véstibus; ascendénte autem Dómino, in albis véstibus Angeli appárent: nisi quod tunc magna solémnitas Angelis facta est, cum cælum Deus homo penetrávit? Quia nascénte Dómino, videbátur divínitas humiliáta; ascendénte vero Dómino, est humánitas exaltáta. Albæ étenim vestes exaltatióni magis cóngruunt, quam humiliatióni.\par
\switchcolumn
\EN
\noindent The first question we have to ask is why we read that Angels appeared at the time of the Birth of the Lord, but we read not that they appeared in white apparel whereas, when the Lord ascended into heaven, it is written that the angels which appeared were clad in white. “While they beheld, He was taken up, and a cloud received Him out of their sight. And while they looked steadfastly toward heaven, as He went up, behold, two men stood by them in white apparel,” Acts i. 9,\par
10. White raiment is an outward sign of solemn inward joy. That the occasion of God-made-Man entering into heaven was a great Festival for Angels, is the reason which we see why angels are specially named as robed in white at His Ascension, and not at His Birth. At the Birth of the Lord the Godhead was manifested veiled under the form of a servant, but at His Ascension the Manhood was seen exalted and white vestments are more apt to exaltation than humiliation.\par
\switchcolumn*

\LA
\begin{Resp}\Rsign Ego rogábo Patrem, et álium Paráclitum dabit vobis, \Star{} Ut máneat vobíscum in ætérnum, Spíritum veritátis, allelúja.\end{Resp}
\begin{Resp}\Vsign Si enim non abíero, Paráclitus non véniet ad vos: si autem abíero, mittam eum ad vos.\end{Resp}
\begin{Resp}\Rsign Ut máneat vobíscum in ætérnum, Spíritum veritátis, allelúja.\end{Resp}
\switchcolumn
\EN
\begin{Resp}\Rsign I will pray the Father, and He shall give you another Comforter, \Star{} That He may abide with you for ever, even the Spirit of truth. Alleluia.\end{Resp}
\begin{Resp}\Vsign For if I go not away, the Comforter will not come unto you; but if I depart, I will send Him unto you.\end{Resp}
\begin{Resp}\Rsign That He may abide with you for ever, even the Spirit of truth. Alleluia.\end{Resp}
\switchcolumn*

\LA
\LessonHead{Lectio VIII}
\switchcolumn
\EN
\LessonHead{Lesson VIII}
\switchcolumn*

\LA
\noindent In ascensióne ergo ejus Angeli in albis véstibus vidéri debuérunt: quia qui in nativitáte sua appáruit Deus húmilis, in ascensióne sua osténsus est homo sublímis. Sed hoc nobis magnópere, fratres caríssimi, in hac solemnitáte pensándum est: quia delétum est hodiérna die chirógraphum damnatiónis nostræ, mutáta est senténtia corruptiónis nostræ. Illa enim natúra, cui dictum est: Terra es, et in terram ibis; hódie in cælum ivit. Pro hac ipsa namque carnis nostræ sublevatióne, per figúram beátus Job Dóminum avem vocat. Quia enim ascensiónis ejus mystérium Judǽam non intellígere conspéxit, de infidelitáte ejus per figúram beátus Job senténtiam prótulit, dicens: Sémitam ignorávit avis.\par
\switchcolumn
\EN
\noindent Therefore were the angels bound to appear in white apparel at the Ascension at His Birth He Who thought it not robbery to be equal with God, was seen in the form in which He had humbled Himself; at His Ascension the Manhood Which He had taken into God was seen glorified. Again, dearly beloved brethren, we must remember today, how that Christ hath “blotted out the hand-writing that was against us,” and reversed the sentence which doomed us to corruption. That same nature to which it was said, “Dust thou art, and unto dust shalt thou return,” that same nature is His Who hath this day ascended up into heaven. It is because of this up-lifting of our flesh that blessed Job, by a figure, calleth the Lord a bird. The Jews could not understand the Mystery of the Ascension, and in view of this their unbelief, blessed Job said mystically “He knew not the path of the bird.”\par
\switchcolumn*

\LA
\begin{Resp}\Rsign Ponis nubem ascénsum tuum, Dómine: \Star{} Qui ámbulas super pennas ventórum, allelúja.\end{Resp}
\begin{Resp}\Vsign Confessiónem et decórem induísti, amíctus lumen sicut vestiméntum.\end{Resp}
\begin{Resp}\Rsign Qui ámbulas super pennas ventórum, allelúja.\end{Resp}
\begin{Resp}\Vsign Glória Patri, et Fílio, \Star{} et Spirítui Sancto.\end{Resp}
\begin{Resp}\Rsign Qui ámbulas super pennas ventórum, allelúja.\end{Resp}
\switchcolumn
\EN
\begin{Resp}\Rsign Thou makest the clouds the chariot, O Lord, \Star{} Thou walkest upon the wings of the wind. Alleluia.\end{Resp}
\begin{Resp}\Vsign Thou art clothed with honour and majesty, covering thyself with light as with a garment;\end{Resp}
\begin{Resp}\Rsign Thou walkest upon the wings of the wind. Alleluia.\end{Resp}
\begin{Resp}\Vsign Glory be to the Father, and to the Son, \Star{} and to the Holy Ghost.\end{Resp}
\begin{Resp}\Rsign Thou walkest upon the wings of the wind. Alleluia.\end{Resp}
\switchcolumn*

\LA
\LessonHead{Lectio IX}
\switchcolumn
\EN
\LessonHead{Lesson IX}
\switchcolumn*

\LA
\noindent Avis enim recte appellátus est Dóminus, quia corpus cárneum ad ǽthera librávit. Cujus avis sémitam ignorávit, quisquis eum ad cælum ascendísse non crédidit. De hac solemnitáte per Psalmístam dícitur: Eleváta est magnificéntia tua super cælos. De hac rursus ait: Ascéndit Deus in jubilatióne, et Dóminus in voce tubæ. De hac íterum dicit: Ascéndens in altum captívam duxit captivitátem, dedit dona homínibus. Ascéndens quippe in altum, captívam duxit captivitátem: quia corruptiónem nostram virtúte suæ incorruptiónis absórbuit. Dedit vero dona homínibus: quia misso désuper Spíritu, álii sermónem sapiéntiæ, álii sermónem sciéntiæ, álii grátiam virtútum, álii grátiam curatiónum, álii génera linguárum, álii interpretatiónem tríbuit sermónum.\par
\switchcolumn
\EN
\noindent The name of a bird is well given to the Lord, Who bodily soared up into heaven. And the path of that Bird knoweth no man, who believeth not in the Ascension into heaven. It is of this glorious occasion that the Psalmist saith: “Who hast set thy glory above the heavens,” (viii. 2,) and again “God is gone up with a shout, and the Lord with the sound of a trumpet,” (xlvi. 6.) And yet again he saith “Thou hast ascended on high, Thou hast led captivity captive,” (lxvii. 19.) “When Christ ascended up on high, He led captivity captive,” (Eph. iv. 8,) because by His Own incorruptibility He swallowed up our corruptibility. “He gave gifts unto men,” because by sending the Spirit from above, He gave “to one, the word of wisdom to another, the word of knowledge to another, the working of miracles to another, the gifts of healing; to another, diverse kinds of tongues to another, the interpretation of tongues,” (1 Cor. xii. 8-10.)\par
\switchcolumn*

\LA
\TeDeum{Te Deum}
\switchcolumn
\EN
\TeDeum{Te Deum}
\switchcolumn*

\end{paracol}

\Office{In Octava Ascensionis}{Thursday, Octave Day of the Ascension}{Duplex majus}
\addcontentsline{toc}{subsection}{In Octava Ascensionis\enspace\textit{Thursday, Octave Day of the Ascension}}
\begin{paracol}{2}
\LA
\LessonHead{Lectio I}
\switchcolumn
\EN
\LessonHead{Lesson I}
\switchcolumn*

\LA
\LessonTitle{De Epístola beáti Pauli Apóstoli ad Ephésios}
\switchcolumn
\EN
\LessonTitle{Lesson from the letter of St. Paul the Apostle to the Ephesians}
\switchcolumn*

\LA
\Cite{Eph 4:1-8}
\switchcolumn
\EN
\Cite{Eph 4:1-8}
\switchcolumn*

\LA
\noindent Obsecro ítaque vos ego vinctus in Dómino, ut digne ambulétis vocatióne, qua vocáti estis, cum omni humilitáte, et mansuetúdine, cum patiéntia, supportántes ínvicem in caritáte, sollíciti serváre unitátem Spíritus in vínculo pacis. Unum corpus, et unus Spíritus, sicut vocáti estis in una spe vocatiónis vestræ. Unus Dóminus, una fides, unum baptísma. Unus Deus et Pater ómnium, qui est super omnes, et per ómnia, et in ómnibus nobis. Unicuíque autem nostrum data est grátia secúndum mensúram donatiónis Christi. Propter quod dicit: Ascéndens in altum, captívam duxit captivitátem: dedit dona homínibus.\par
\switchcolumn
\EN
\noindent I therefore, a prisoner in the Lord, beseech you that you walk worthy of the vocation in which you are called, With all humility and mildness, with patience, supporting one another in charity. Careful to keep the unity of the Spirit in the bond of peace. One body and one Spirit; as you are called in one hope of your calling. One Lord, one faith, one baptism. One God and Father of all, who is above all, and through all, and in us all. But to every one of us is given grace, according to the measure of the giving of Christ. Wherefore he saith: Ascending on high, he led captivity captive; he gave gifts to men.\par
\switchcolumn*

\LA
\begin{Resp}\Rsign Post passiónem suam per dies quadragínta appárens eis, et loquens de regno Dei, allelúja: \Star{} Et, vidéntibus illis, elevátus est, allelúja: et nubes suscépit eum ab óculis eórum, allelúja.\end{Resp}
\begin{Resp}\Vsign Et convéscens, præcépit eis, ab Jerosólymis ne discéderent, sed exspectárent promissiónem Patris.\end{Resp}
\begin{Resp}\Rsign Et, vidéntibus illis, elevátus est, allelúja: et nubes suscépit eum ab óculis eórum, allelúja.\end{Resp}
\switchcolumn
\EN
\begin{Resp}\Rsign Being seen of them forty days after that He had suffered, and speaking of the kingdom of God. Alleluia. \Star{} And while they beheld, He was taken up, and a cloud received Him out of their sight. Alleluia.\end{Resp}
\begin{Resp}\Vsign And, eating together with them, He commanded them that they should not depart from Jerusalem, but wait for the Promise of the Father.\end{Resp}
\begin{Resp}\Rsign And while they beheld, He was taken up, and a cloud received Him out of their sight. Alleluia.\end{Resp}
\switchcolumn*

\LA
\LessonHead{Lectio II}
\switchcolumn
\EN
\LessonHead{Lesson II}
\switchcolumn*

\LA
\Cite{Eph 4:9-14}
\switchcolumn
\EN
\Cite{Eph 4:9-14}
\switchcolumn*

\LA
\noindent Quod autem ascéndit, quid est, nisi quia et descéndit primum in inferióres partes terræ? Qui descéndit, ipse est et qui ascéndit super omnes cælos, ut impléret ómnia. Et ipse dedit quosdam quidem apóstolos, quosdam autem prophétas, álios vero evangelístas, álios autem pastóres et doctóres, ad consummatiónem sanctórum in opus ministérii, in ædificatiónem córporis Christi: donec occurrámus omnes in unitátem fídei, et agnitiónis Fílii Dei, in virum perféctum, in mensúram ætátis plenitúdinis Christi: ut jam non simus párvuli fluctuántes, et circumferámur omni vento doctrínæ in nequítia hóminum, in astútia ad circumventiónem erróris.\par
\switchcolumn
\EN
\noindent Now that he ascended, what is it, but because he also descended first into the lower parts of the earth? He that descended is the same also that ascended above all the heavens, that he might fill all things. And he gave some apostles, and some prophets, and other some evangelists, and other some pastors and doctors, For the perfecting of the saints, for the work of the ministry, for the edifying of the body of Christ: Until we all meet into the unity of faith, and of the knowledge of the Son of God, unto a perfect man, unto the measure of the age of the fulness of Christ; That henceforth we be no more children tossed to and fro, and carried about with every wind of doctrine by the wickedness of men, by cunning craftiness, by which they lie in wait to deceive.\par
\switchcolumn*

\LA
\begin{Resp}\Rsign Omnis pulchritúdo Dómini exaltáta est super sídera: \Star{} Spécies ejus in núbibus cæli, et nomen ejus in ætérnum pérmanet, allelúja.\end{Resp}
\begin{Resp}\Vsign A summo cælo egréssio ejus, et occúrsus ejus usque ad summum ejus.\end{Resp}
\begin{Resp}\Rsign Spécies ejus in núbibus cæli, et nomen ejus in ætérnum pérmanet, allelúja.\end{Resp}
\switchcolumn
\EN
\begin{Resp}\Rsign The Lord hath set His beauty above the stars; \Star{} His loveliness is in the clouds of heaven, and His Name endureth for ever. Alleluia.\end{Resp}
\begin{Resp}\Vsign His going forth is from the end of the heaven, and His circuit unto the ends of it.\end{Resp}
\begin{Resp}\Rsign His loveliness is in the clouds of heaven, and His Name endureth for ever. Alleluia.\end{Resp}
\switchcolumn*

\LA
\LessonHead{Lectio III}
\switchcolumn
\EN
\LessonHead{Lesson III}
\switchcolumn*

\LA
\Cite{Eph 4:15-21}
\switchcolumn
\EN
\Cite{Eph 4:15-21}
\switchcolumn*

\LA
\noindent Veritátem autem faciéntes in caritáte, crescámus in illo per ómnia, qui est caput Christus: ex quo totum corpus compáctum et connéxum per omnem junctúram subministratiónis, secúndum operatiónem in mensúram uniuscujúsque membri, augméntum córporis facit in ædificatiónem sui in caritáte. Hoc ígitur dico, et testíficor in Dómino, ut jam non ambulétis, sicut et gentes ámbulant in vanitáte sensus sui, ténebris obscurátum habéntes intelléctum, alienáti a vita Dei per ignorántiam, quæ est in illis, propter cæcitátem cordis ipsórum, qui desperántes, semetípsos tradidérunt impudicítiæ, in operatiónem immundítiæ omnis in avarítiam. Vos autem non ita didicístis Christum, si tamen illum audístis, et in ipso edócti estis.\par
\switchcolumn
\EN
\noindent But doing the truth in charity, we may in all things grow up in him who is the head, even Christ: From whom the whole body, being compacted and fitly joined together, by what every joint supplieth, according to the operation in the measure of every part, maketh increase of the body, unto the edifying of itself in charity. This then I say and testify in the Lord: That henceforward you walk not as also the Gentiles walk in the vanity of their mind, Having their understanding darkened, being alienated from the life of God through the ignorance that is in them, because of the blindness of their hearts. Who despairing, have given themselves up to lasciviousness, unto the working of all uncleanness, unto the working of all uncleanness, unto covetousness. But you have not so learned Christ; If so be that you have heard him, and have been taught in him.\par
\switchcolumn*

\LA
\begin{Resp}\Rsign Exaltáre, Dómine, allelúja, \Star{} In virtúte tua, allelúja.\end{Resp}
\begin{Resp}\Vsign Eleváta est, magnificéntia tua super cælos, Deus.\end{Resp}
\begin{Resp}\Rsign In virtúte tua, allelúja.\end{Resp}
\begin{Resp}\Vsign Glória Patri, et Fílio, \Star{} et Spirítui Sancto.\end{Resp}
\begin{Resp}\Rsign In virtúte tua, allelúja.\end{Resp}
\switchcolumn
\EN
\begin{Resp}\Rsign Be Thou exalted, O Lord. Alleluia. \Star{} In thine Own strength. Alleluia.\end{Resp}
\begin{Resp}\Vsign O God, Thou hast set thy glory above the heavens.\end{Resp}
\begin{Resp}\Rsign In thine Own strength. Alleluia.\end{Resp}
\begin{Resp}\Vsign Glory be to the Father, and to the Son, \Star{} and to the Holy Ghost.\end{Resp}
\begin{Resp}\Rsign In thine Own strength. Alleluia.\end{Resp}
\switchcolumn*

\LA
\LessonHead{Lectio IV}
\switchcolumn
\EN
\LessonHead{Lesson IV}
\switchcolumn*

\LA
\LessonTitle{Sermo sancti Augustíni Epíscopi}
\switchcolumn
\EN
\LessonTitle{From the Sermons of St. Augustine, Bishop of Hippo.}
\switchcolumn*

\LA
\Source{Sermo 3 de Ascensione Domini, qui est 176 de Tempore}
\switchcolumn
\EN
\Source{3rd on the Ascension, 176th on the Season.}
\switchcolumn*

\LA
\noindent Omnia, caríssimi, quæ Dóminus Jesus Christus in hoc mundo sub fragilitáte nostra mirácula édidit, nobis profíciunt: qui dum humánam conditiónem sidéribus importávit, credéntibus cælum patére posse monstrávit: et dum victórem mortis in cæléstia elevávit, victóribus quo sequántur osténdit. Ascénsio ergo Dómini cathólicæ fídei confirmátio fuit: ut secúri in pósterum crederémus miráculi illíus donum, cujus jam in præsénti percepissémus efféctum: et fidélis quisque cum jam tanta percéperit, per ea, quæ cognóscit prǽstita, discat speráre promíssa, ac Dei sui prætéritam præsentémque bonitátem, quasi futurórum téneat cautiónem.\par
\switchcolumn
\EN
\noindent Dearly beloved brethren, all the wonderful works which our Lord Jesus Christ did in this world, under the weakness of our nature, are profitable for us when He exalted His Manhood above the stars, He showed that heaven may open for a believer and while He, the Conqueror of death, went up into the heavenly mansions, He showed to him that overcometh, whither he also may follow. Therefore, the ascension of the Lord is the seal of the Catholic Faith, which assureth in us the hope of the gift which is yet to come to us, from a miracle whereof we already feel the fruits. Thus let every one that is faithful, having already received so much, learn to hope for that which is promised, on the ground of that which he knoweth to have been given, and hold the goodness of God in times which have been, and times which now are, as a sure pledge of the same in times to come.\par
\switchcolumn*

\LA
\begin{Resp}\Rsign Tempus est, ut revértar ad eum, qui me misit, dicit Dóminus: nolíte contristári, nec turbétur cor vestrum: \Star{} Rogo pro vobis Patrem, ut ipse vos custódiat, allelúja, allelúja.\end{Resp}
\begin{Resp}\Vsign Nisi ego abíero, Paráclitus non véniet: cum assúmptus fúero, mittam vobis eum.\end{Resp}
\begin{Resp}\Rsign Rogo pro vobis Patrem, ut ipse vos custódiat, allelúja, allelúja.\end{Resp}
\switchcolumn
\EN
\begin{Resp}\Rsign My time is come that I should return unto Him That sent Me, saith the Lord. Be not sorrowful, neither let your heart be troubled. \Star{} I pray the Father for you, that He may keep you. Alleluia, Alleluia.\end{Resp}
\begin{Resp}\Vsign If I go not away, the Comforter will not come unto you; where I am ascended, I will send Him unto you.\end{Resp}
\begin{Resp}\Rsign I pray the Father for you, that He may keep you. Alleluia, Alleluia.\end{Resp}
\switchcolumn*

\LA
\LessonHead{Lectio V}
\switchcolumn
\EN
\LessonHead{Lesson V}
\switchcolumn*

\LA
\noindent Super excélsa ergo cæli terrénum corpus impónitur: ossa, intra sepúlcri angústias paulo ante conclúsa, Angelórum cœ́tibus interúntur: in grémium immortalitátis mortális natúra transfúnditur: et ídeo sacra apostólicæ lectiónis testátur história: Cum hæc dixísset, inquit, vidéntibus illis, elevátus est. Dum audis elevátum, agnósce milítiæ cæléstis obséquium: unde hodiérna festívitas hóminis nobis et Dei sacraménta manifestávit. Sub una eadémque persóna, in eo qui élevat, divínam poténtiam; in eo autem, qui elevátur, humánam cognósce substántiam.\par
\switchcolumn
\EN
\noindent An earthly Body, then, is now lifted up above the heights of heaven the Bones, Which but a little while before had lain within the narrow walls of the grave, have made their entry among the angelic hosts human nature hath been given a place in the lap of immortality and therefore the Apostle whose account we have heard read, saith “When He had spoken these things, while they beheld, He was taken up.” (Acts i. 9.) While thou hearest these words, “taken up,” thou must understand thereby the ministry of the angelic army whereby this Festival revealeth to us the Mystery of Him who is both God and Man. United in One Person, we see in Him who lifted up, Divine Power, and in Him Who was lifted up, very Man.\par
\switchcolumn*

\LA
\begin{Resp}\Rsign Non turbétur cor vestrum: ego vado ad Patrem; et cum assúmptus fúero a vobis, mittam vobis, allelúja, \Star{} Spíritum veritátis, et gaudébit cor vestrum, allelúja.\end{Resp}
\begin{Resp}\Vsign Ego rogábo Patrem, et álium Paráclitum dabit vobis.\end{Resp}
\begin{Resp}\Rsign Spíritum veritátis, et gaudébit cor vestrum, allelúja.\end{Resp}
\switchcolumn
\EN
\begin{Resp}\Rsign Let not your heart be troubled; I go unto the Father, and when I am taken from you, I will send unto you, alleluia, \Star{} The Spirit of truth, and your heart shall rejoice. Alleluia.\end{Resp}
\begin{Resp}\Vsign I will pray the Father, and He shall give you another Comforter.\end{Resp}
\begin{Resp}\Rsign The Spirit of truth, and your heart shall rejoice. Alleluia.\end{Resp}
\switchcolumn*

\LA
\LessonHead{Lectio VI}
\switchcolumn
\EN
\LessonHead{Lesson VI}
\switchcolumn*

\LA
\noindent Ideóque omnímodis detestánda sunt venéna Orientális erróris, qui ímpia novitáte præsúmit assérere Fílium Dei ac Fílium hóminis uníus esse natúræ. In altérutra enim parte, vel qui solum hóminem fuísse díxerit, negábit Conditóris glóriam: vel qui solum Deum, negábit misericórdiam Redemptóris. Quo génere non fácile Ariánus evangélicam póterit habére veritátem, ubi Fílium Dei nunc æquálem légimus, nunc minórem. Qui enim uníus natúræ Salvatórem nostrum mortífera persuasióne credíderit, solum hóminem, aut solum Deum cogétur dícere crucifíxum. Sed non ita est. Mortem enim nec solus Deus sentíre, nec solus homo superáre potuísset.\par
\switchcolumn
\EN
\noindent Therefore are utterly to be loathed those pestiferous teachings of Eastern falsehood, those brand-new inventions of ungodliness which dare to assert that He Who in One Person is both Son of God and Son of Man, hath but one nature. On the one hand, if a man say that Christ is not Partaker of the Divine nature, he hath denied the glory of his Maker on the other, he who saith that the Manhood is not of the nature of man, hath denied the mercy of his Saviour. As touching these points, it is well-nigh impossible for an Arian to believe that the Gospel writers are any better than liars, since they distinctly assert in some places that the Son of God is equal, and, in others, that He is inferior, to the Father. Further, if a man be given over to this soul-slaying delusion of believing that our Saviour hath only one nature, he must of necessity admit either that it was only God, or that it was only man who was crucified. But it was not so. If He had been of no nature but the Divine, He could not have suffered, and if He had been of no nature but the human, He could not have conquered death.\par
\switchcolumn*

\LA
\begin{Resp}\Rsign Ascéndens Christus in altum, captívam duxit captivitátem, \Star{} Dedit dona homínibus, allelúja, allelúja, allelúja.\end{Resp}
\begin{Resp}\Vsign Ascéndit Deus in jubilatióne, et Dóminus in voce tubæ.\end{Resp}
\begin{Resp}\Rsign Dedit dona homínibus, allelúja, allelúja, allelúja.\end{Resp}
\begin{Resp}\Vsign Glória Patri, et Fílio, \Star{} et Spirítui Sancto.\end{Resp}
\begin{Resp}\Rsign Dedit dona homínibus, allelúja, allelúja, allelúja.\end{Resp}
\switchcolumn
\EN
\begin{Resp}\Rsign When Christ ascended up on high, He led captivity captive, \Star{} He gave gifts unto men. Alleluia, Alleluia, Alleluia.\end{Resp}
\begin{Resp}\Vsign God is gone up with a shout, and the Lord with the sound of a trumpet.\end{Resp}
\begin{Resp}\Rsign He gave gifts unto men. Alleluia, Alleluia, Alleluia.\end{Resp}
\begin{Resp}\Vsign Glory be to the Father, and to the Son, \Star{} and to the Holy Ghost.\end{Resp}
\begin{Resp}\Rsign He gave gifts unto men. Alleluia, Alleluia, Alleluia.\end{Resp}
\switchcolumn*

\LA
\LessonHead{Lectio VII}
\switchcolumn
\EN
\LessonHead{Lesson VII}
\switchcolumn*

\LA
\LessonTitle{Léctio sancti Evangélii secúndum Marcum}
\switchcolumn
\EN
\LessonTitle{From the Holy Gospel according to Mark}
\switchcolumn*

\LA
\Cite{Marc 16:14-20}
\switchcolumn
\EN
\Cite{Mark 16:14-20}
\switchcolumn*

\LA
\noindent In illo témpore: Recumbéntibus úndecim discípulis, appáruit illis Jesus: et exprobrávit incredulitátem eórum et durítiam cordis, quia iis, qui víderant eum resurrexísse, non credidérunt. Et réliqua.\par
\switchcolumn
\EN
\noindent At that time, Jesus appeared unto the eleven disciples as they sat at meat, and upbraided them with their unbelief and hardness of heart, because they believed not them which had seen Him after He was risen. And so on.\par
\switchcolumn*

\LA
\LessonTitle{Homilía sancti Gregórii Papæ}
\switchcolumn
\EN
\LessonTitle{Homily by Pope St. Gregory the Great.}
\switchcolumn*

\LA
\Source{Ex eadem Homilía 29.}
\switchcolumn
\EN
\Source{Same as before.}
\switchcolumn*

\LA
\noindent De hac ascensiónis ejus glória étiam Hábacuc ait: Elevátus est sol, et luna stetit in órdine suo. Quis enim solis nómine, nisi Dóminus; et quæ lunæ nómine, nisi Ecclésia designátur? Quoúsque enim Dóminus ascéndit ad cælos, sancta ejus Ecclésia advérsa mundi omnímodo formidávit: at postquam ejus ascensióne roboráta est, apérte prædicávit quod occúlte crédidit. Elevátus est ergo sol, et luna stetit in órdine suo: quia cum Dóminus cælum pétiit, sancta ejus Ecclésia in auctoritáte prædicatiónis excrévit.\par
\switchcolumn
\EN
\noindent The Prophet Habakkuk also hath spoken of the glory of Christ's Ascension in the words “The sun was lifted up on high, and the moon stood still in her habitation,” Who is here signified by the Sun, if not the Saviour or by the Moon, if not the Church? Until the Lord was withdrawn from her sight, (that is, by His Ascension,) His Holy Church was pale before the hostile glare of the world, but after He was ascended, she waxed stronger, and distinctly shed forth the beams of that faith which had hitherto dwelt hiddenly in her. “The sun was lifted up, and the moon stood still in her habitation” when the Lord was gone away into heaven, His holy Church waxed stronger in her enlightening power.\par
\switchcolumn*

\LA
\begin{Resp}\Rsign Ego rogábo Patrem, et álium Paráclitum dabit vobis, \Star{} Ut máneat vobíscum in ætérnum, Spíritum veritátis, allelúja.\end{Resp}
\begin{Resp}\Vsign Si enim non abíero, Paráclitus non véniet ad vos: si autem abíero, mittam eum ad vos.\end{Resp}
\begin{Resp}\Rsign Ut máneat vobíscum in ætérnum, Spíritum veritátis, allelúja.\end{Resp}
\switchcolumn
\EN
\begin{Resp}\Rsign I will pray the Father, and He shall give you another Comforter, \Star{} That He may abide with you for ever, even the Spirit of truth. Alleluia.\end{Resp}
\begin{Resp}\Vsign For if I go not away, the Comforter will not come unto you; but if I depart, I will send Him unto you.\end{Resp}
\begin{Resp}\Rsign That He may abide with you for ever, even the Spirit of truth. Alleluia.\end{Resp}
\switchcolumn*

\LA
\LessonHead{Lectio VIII}
\switchcolumn
\EN
\LessonHead{Lesson VIII}
\switchcolumn*

\LA
\noindent Hinc ejúsdem Ecclésiæ voce per Salomónem dícitur: Ecce iste venit sáliens in móntibus, et transíliens colles. Considerávit namque tantórum óperum cúlmina, et ait: Ecce iste venit sáliens in móntibus. Veniéndo quippe ad redemptiónem nostram, quosdam, ut ita dicam, saltus dedit. Vultis, fratres caríssimi, ipsos ejus saltus agnóscere? De cælo venit in úterum, de útero venit in præsépe, de præsépe venit in crucem, de cruce venit in sepúlcrum, de sepúlcro rédiit in cælum. Ecce, ut nos post se cúrrere fáceret, quosdam pro nobis saltus manifestáta per carnem Véritas dedit: quia exsultávit ut gigas ad curréndam viam suam, ut nos ei dicerémus ex corde: Trahe nos: post te currémus in odórem unguentórum tuórum.\par
\switchcolumn
\EN
\noindent Hence it is that Solomon hath put into the mouth of the, (same) Church the words: “Behold,; He cometh! leaping upon the mountains, skipping upon the hills These hills are his lofty and; noble achievements. “Behold, He cometh leaping upon the mountains.” When He came to redeem us, He came, if I may so say, in leaps. My dearly beloved brethren, would you know what His leaps were? From heaven he leapt into the womb of the Virgin, from the womb into the manger, from the manger on to the Cross, from the Cross into the grave, and from the grave up to heaven. Lo, how the Truth made manifest in the Flesh did leap for our sakes, that He might draw us to run after Him for this end did He rejoice, as a strong man to run a race,\par
\switchcolumn*

\LA
\begin{Resp}\Rsign Ponis nubem ascénsum tuum, Dómine: \Star{} Qui ámbulas super pennas ventórum, allelúja.\end{Resp}
\begin{Resp}\Vsign Confessiónem et decórem induísti, amíctus lumen sicut vestiméntum.\end{Resp}
\begin{Resp}\Rsign Qui ámbulas super pennas ventórum, allelúja.\end{Resp}
\begin{Resp}\Vsign Glória Patri, et Fílio, \Star{} et Spirítui Sancto.\end{Resp}
\begin{Resp}\Rsign Qui ámbulas super pennas ventórum, allelúja.\end{Resp}
\switchcolumn
\EN
\begin{Resp}\Rsign Thou makest the clouds the chariot, O Lord, \Star{} Thou walkest upon the wings of the wind. Alleluia.\end{Resp}
\begin{Resp}\Vsign Thou art clothed with honour and majesty, covering thyself with light as with a garment;\end{Resp}
\begin{Resp}\Rsign Thou walkest upon the wings of the wind. Alleluia.\end{Resp}
\begin{Resp}\Vsign Glory be to the Father, and to the Son, \Star{} and to the Holy Ghost.\end{Resp}
\begin{Resp}\Rsign Thou walkest upon the wings of the wind. Alleluia.\end{Resp}
\switchcolumn*

\LA
\LessonHead{Lectio IX}
\switchcolumn
\EN
\LessonHead{Lesson IX}
\switchcolumn*

\LA
\noindent Unde, fratres caríssimi, opórtet, ut illuc sequámur corde, ubi eum córpore ascendísse crédimus. Desidéria terréna fugiámus: nihil nos jam deléctet in ínfimis, qui Patrem habémus in cælis. Et hoc nobis est magnópere perpendéndum: quia is qui plácidus ascéndit, terríbilis redíbit; et quidquid nobis cum mansuetúdine præcépit, hoc a nobis cum distinctióne éxiget. Nemo ergo indúlta pœniténtiæ témpora parvipéndat, nemo curam sui, dum valet, ágere négligat: quia Redémptor noster tanto tunc in judícium distríctior véniet, quanto nobis ante judícium magnam patiéntiam prærogávit.\par
\switchcolumn
\EN
\noindent Therefore, dearly beloved brethren, it behoveth us in heart and mind thither to ascend, where we believe Him to have already ascended bodily. Let us fly earthly lusts for us, who have a Father in heaven, let nothing be sweet below And very much must we keep in our minds this thought, that He Which ascended up in peace, will return in dreadful Majesty and will require from us with justice an account of our keeping of those commandments which He gave us in mercy. Let no man therefore reckon lightly this season which is given unto us that we may repent ourselves, nor be reckless touching the state of his soul; our Redeemer will be all the sterner, when He cometh to judgment, as He hath been wondrously long-suffering before.\par
\switchcolumn*

\LA
\TeDeum{Te Deum}
\switchcolumn
\EN
\TeDeum{Te Deum}
\switchcolumn*

\end{paracol}

\Office{Feria V infra Hebdomadam post Ascensionem}{Thursday after the Ascension}{Feria}
\addcontentsline{toc}{subsection}{Feria V infra Hebdomadam post Ascensionem\enspace\textit{Thursday after the Ascension}}
\begin{paracol}{2}
\LA
\LessonHead{Lectio I}
\switchcolumn
\EN
\LessonHead{Lesson I}
\switchcolumn*

\LA
\LessonTitle{Incipit Epístola tértia beáti Joánnis Apóstoli}
\switchcolumn
\EN
\LessonTitle{Lesson from the third letter of St. John the Apostle}
\switchcolumn*

\LA
\Cite{3 Joannes 1:1-4}
\switchcolumn
\EN
\Cite{3 John 1:1-4}
\switchcolumn*

\LA
\noindent Sénior Gaio caríssimo, quem ego díligo in veritáte. Caríssime, de ómnibus oratiónem fácio próspere te íngredi, et valére sicut próspere agit ánima tua. Gavísus sum valde veniéntibus frátribus, et testimónium perhibéntibus veritáti tuæ, sicut tu in veritáte ámbulas. Majórem horum non hábeo grátiam, quam ut áudiam fílios meos in veritáte ambuláre.\par
\switchcolumn
\EN
\noindent The ancient to the dearly beloved Gaius, whom I love in truth. Dearly beloved, concerning all things I make it my prayer that thou mayest proceed prosperously, and fare well as thy soul doth prosperously. I was exceedingly glad when the brethren came and gave testimony to the truth in thee, even as thou walkest in the truth. I have no greater grace than this, to hear that my children walk in truth.\par
\switchcolumn*

\LA
\begin{Resp}\Rsign Post passiónem suam per dies quadragínta appárens eis, et loquens de regno Dei, allelúja: \Star{} Et, vidéntibus illis, elevátus est, allelúja: et nubes suscépit eum ab óculis eórum, allelúja.\end{Resp}
\begin{Resp}\Vsign Et convéscens, præcépit eis, ab Jerosólymis ne discéderent, sed exspectárent promissiónem Patris.\end{Resp}
\begin{Resp}\Rsign Et, vidéntibus illis, elevátus est, allelúja: et nubes suscépit eum ab óculis eórum, allelúja.\end{Resp}
\switchcolumn
\EN
\begin{Resp}\Rsign Being seen of them forty days after that He had suffered, and speaking of the kingdom of God. Alleluia. \Star{} And while they beheld, He was taken up, and a cloud received Him out of their sight. Alleluia.\end{Resp}
\begin{Resp}\Vsign And, eating together with them, He commanded them that they should not depart from Jerusalem, but wait for the Promise of the Father.\end{Resp}
\begin{Resp}\Rsign And while they beheld, He was taken up, and a cloud received Him out of their sight. Alleluia.\end{Resp}
\switchcolumn*

\end{paracol}

\Office{Feria VI post Octavam Ascensionis}{Friday after the Octave of the Ascension}{Semiduplex}
\addcontentsline{toc}{subsection}{Feria VI post Octavam Ascensionis\enspace\textit{Friday after the Octave of the Ascension}}
\begin{paracol}{2}
\LA
\LessonHead{Lectio I}
\switchcolumn
\EN
\LessonHead{Lesson I}
\switchcolumn*

\LA
\LessonTitle{Incipit Epístola tértia beáti Joánnis Apóstoli}
\switchcolumn
\EN
\LessonTitle{Lesson from the third letter of St. John the Apostle}
\switchcolumn*

\LA
\Cite{3 Joannes 1:1-4}
\switchcolumn
\EN
\Cite{3 John 1:1-4}
\switchcolumn*

\LA
\noindent Sénior Gaio caríssimo, quem ego díligo in veritáte. Caríssime, de ómnibus oratiónem fácio próspere te íngredi, et valére sicut próspere agit ánima tua. Gavísus sum valde veniéntibus frátribus, et testimónium perhibéntibus veritáti tuæ, sicut tu in veritáte ámbulas. Majórem horum non hábeo grátiam, quam ut áudiam fílios meos in veritáte ambuláre.\par
\switchcolumn
\EN
\noindent The ancient to the dearly beloved Gaius, whom I love in truth. Dearly beloved, concerning all things I make it my prayer that thou mayest proceed prosperously, and fare well as thy soul doth prosperously. I was exceedingly glad when the brethren came and gave testimony to the truth in thee, even as thou walkest in the truth. I have no greater grace than this, to hear that my children walk in truth.\par
\switchcolumn*

\LA
\begin{Resp}\Rsign Post passiónem suam per dies quadragínta appárens eis, et loquens de regno Dei, allelúja: \Star{} Et, vidéntibus illis, elevátus est, allelúja: et nubes suscépit eum ab óculis eórum, allelúja.\end{Resp}
\begin{Resp}\Vsign Et convéscens, præcépit eis, ab Jerosólymis ne discéderent, sed exspectárent promissiónem Patris.\end{Resp}
\begin{Resp}\Rsign Et, vidéntibus illis, elevátus est, allelúja: et nubes suscépit eum ab óculis eórum, allelúja.\end{Resp}
\switchcolumn
\EN
\begin{Resp}\Rsign Being seen of them forty days after that He had suffered, and speaking of the kingdom of God. Alleluia. \Star{} And while they beheld, He was taken up, and a cloud received Him out of their sight. Alleluia.\end{Resp}
\begin{Resp}\Vsign And, eating together with them, He commanded them that they should not depart from Jerusalem, but wait for the Promise of the Father.\end{Resp}
\begin{Resp}\Rsign And while they beheld, He was taken up, and a cloud received Him out of their sight. Alleluia.\end{Resp}
\switchcolumn*

\LA
\LessonHead{Lectio II}
\switchcolumn
\EN
\LessonHead{Lesson II}
\switchcolumn*

\LA
\Cite{3 Joannes 1:5-10}
\switchcolumn
\EN
\Cite{3 John 1:5-10}
\switchcolumn*

\LA
\noindent Caríssime, fidéliter facis quidquid operáris in fratres, et hoc in peregrínos, qui testimónium reddidérunt caritáti tuæ in conspéctu ecclésiæ: quos, benefáciens, dedúces digne Deo. Pro nómine enim ejus profécti sunt, nihil accipiéntes a géntibus. Nos ergo debémus suscípere hujúsmodi, ut cooperatóres simus veritátis. Scripsíssem fórsitan ecclésiæ: sed is qui amat primátum gérere in eis, Diótrephes, non récipit nos: propter hoc si vénero, commonébo ejus ópera, quæ facit, verbis malígnis gárriens in nos: et quasi non ei ista suffíciant, neque ipse súscipit fratres: et eos qui suscípiunt, próhibet, et de ecclésia éicit.\par
\switchcolumn
\EN
\noindent Dearly beloved, thou dost faithfully whatever thou dost for the brethren, and that for strangers, Who have given testimony to thy charity in the sight of the church: whom thou shalt do well to bring forward on their way in a manner worthy of God. Because, for his name they went out, taking nothing of the Gentiles. We therefore ought to receive such, that we may be fellow helpers of the truth. I had written perhaps to the church: but Diotrephes, who loveth to have the pre-eminence among them, doth not receive us. For this cause, if I come, I will advertise his works which he doth, with malicious words prating against us. And as if these things were not enough for him, neither doth he himself receive the brethren, and them that do receive them he forbiddeth, and casteth out of the church.\par
\switchcolumn*

\LA
\begin{Resp}\Rsign Omnis pulchritúdo Dómini exaltáta est super sídera: \Star{} Spécies ejus in núbibus cæli, et nomen ejus in ætérnum pérmanet, allelúja.\end{Resp}
\begin{Resp}\Vsign A summo cælo egréssio ejus, et occúrsus ejus usque ad summum ejus.\end{Resp}
\begin{Resp}\Rsign Spécies ejus in núbibus cæli, et nomen ejus in ætérnum pérmanet, allelúja.\end{Resp}
\switchcolumn
\EN
\begin{Resp}\Rsign The Lord hath set His beauty above the stars; \Star{} His loveliness is in the clouds of heaven, and His Name endureth for ever. Alleluia.\end{Resp}
\begin{Resp}\Vsign His going forth is from the end of the heaven, and His circuit unto the ends of it.\end{Resp}
\begin{Resp}\Rsign His loveliness is in the clouds of heaven, and His Name endureth for ever. Alleluia.\end{Resp}
\switchcolumn*

\LA
\LessonHead{Lectio III}
\switchcolumn
\EN
\LessonHead{Lesson III}
\switchcolumn*

\LA
\Cite{3 Joannes 1:11-14}
\switchcolumn
\EN
\Cite{3 John 1:11-14}
\switchcolumn*

\LA
\noindent Caríssime, noli imitári malum, sed quod bonum est. Qui benéfacit, ex Deo est: qui maléfacit, non vidit Deum. Demétrio testimónium rédditur ab ómnibus, et ab ipsa veritáte, sed et nos testimónium perhibémus: et nosti quóniam testimónium nostrum verum est. Multa hábui tibi scríbere: sed nólui per atraméntum et cálamum scríbere tibi. Spero autem prótinus te vidére, et os ad os loquémur. Pax tibi. Salútant te amíci. Salúta amícos nominátim.\par
\switchcolumn
\EN
\noindent Dearly beloved, follow not that which is evil, but that which is good. He that doth good, is of God: he that doth evil, hath not seen God. To Demetrius testimony is given by all, and by the truth itself, yea and we also give testimony: and thou knowest that our testimony is true. I had many things to write unto thee: but I would not by ink and pen write to thee. But I hope speedily to see thee, and we will speak mouth to mouth. Peace be to thee. Our friends salute thee. Salute the friends by name.\par
\switchcolumn*

\LA
\begin{Resp}\Rsign Exaltáre, Dómine, allelúja, \Star{} In virtúte tua, allelúja.\end{Resp}
\begin{Resp}\Vsign Eleváta est, magnificéntia tua super cælos, Deus.\end{Resp}
\begin{Resp}\Rsign In virtúte tua, allelúja.\end{Resp}
\begin{Resp}\Vsign Glória Patri, et Fílio, \Star{} et Spirítui Sancto.\end{Resp}
\begin{Resp}\Rsign In virtúte tua, allelúja.\end{Resp}
\switchcolumn
\EN
\begin{Resp}\Rsign Be Thou exalted, O Lord. Alleluia. \Star{} In thine Own strength. Alleluia.\end{Resp}
\begin{Resp}\Vsign O God, Thou hast set thy glory above the heavens.\end{Resp}
\begin{Resp}\Rsign In thine Own strength. Alleluia.\end{Resp}
\begin{Resp}\Vsign Glory be to the Father, and to the Son, \Star{} and to the Holy Ghost.\end{Resp}
\begin{Resp}\Rsign In thine Own strength. Alleluia.\end{Resp}
\switchcolumn*

\LA
\LessonHead{Lectio IV}
\switchcolumn
\EN
\LessonHead{Lesson IV}
\switchcolumn*

\LA
\LessonTitle{De sermóne sancti Augustíni Epíscopi}
\switchcolumn
\EN
\LessonTitle{From the Sermons of St. Augustine, Bishop of Hippo.}
\switchcolumn*

\LA
\Source{Idem Sermone 176 in fine}
\switchcolumn
\EN
\Source{Same as before.}
\switchcolumn*

\LA
\noindent Salvátor noster, caríssimi, si non in nostra carne diábolum triumphávit, se exércuit, non nobis vicit. Si non in nostro córpore resurréxit, conditióni nostræ resurgéndo nihil cóntulit. Hæc qui dicit, suscéptæ et assúmptæ carnis non intélligit ratiónem, confúndit órdinem, evácuat utilitátem. Si non in nostra carne perégit medicínam, solam ergo ex hómine nascéndi elégit injúriam. Recédat a sénsibus nostris tam periculósa persuásio. De nostro est quod appéndit, de suo est quod donávit. Meum testor esse quod cécidit, ut meum sit quod resurréxit. Meum testor esse quod jácuit intra túmulum, ut meum sit quod ascéndit in cælum.\par
\switchcolumn
\EN
\noindent Dearly beloved brethren, if the Flesh wherein our Saviour trampled down the devil had not been of our nature, He would indeed have exercised Himself, but He would not have conquered for us. If the Body wherein He rose from the grave had not been of our nature, His Resurrection would not have affected our state. Whoso asserteth this, that Christ hath but one nature, he doth not understand why Christ took Flesh upon Him, he confoundeth the order, and maketh void the benefit of the Incarnation. If the Flesh wherein our Healer came was not sharer in human | nature, then all that by His Birth He took from man would have been degradation. O may such dangerous! dreams be far from our thoughts What He took is ours, what He gave is His. I testify that the first Adam, who fell, and the second Adam, Who rose from the dead, are both of the same human nature that I am of. I testify that What lay in the grave, and What ascended into heaven, is of the same human nature that I am of.\par
\switchcolumn*

\LA
\begin{Resp}\Rsign Tempus est, ut revértar ad eum, qui me misit, dicit Dóminus: nolíte contristári, nec turbétur cor vestrum: \Star{} Rogo pro vobis Patrem, ut ipse vos custódiat, allelúja, allelúja.\end{Resp}
\begin{Resp}\Vsign Nisi ego abíero, Paráclitus non véniet: cum assúmptus fúero, mittam vobis eum.\end{Resp}
\begin{Resp}\Rsign Rogo pro vobis Patrem, ut ipse vos custódiat, allelúja, allelúja.\end{Resp}
\switchcolumn
\EN
\begin{Resp}\Rsign My time is come that I should return unto Him That sent Me, saith the Lord. Be not sorrowful, neither let your heart be troubled. \Star{} I pray the Father for you, that He may keep you. Alleluia, Alleluia.\end{Resp}
\begin{Resp}\Vsign If I go not away, the Comforter will not come unto you; where I am ascended, I will send Him unto you.\end{Resp}
\begin{Resp}\Rsign I pray the Father for you, that He may keep you. Alleluia, Alleluia.\end{Resp}
\switchcolumn*

\LA
\LessonHead{Lectio V}
\switchcolumn
\EN
\LessonHead{Lesson V}
\switchcolumn*

\LA
\noindent In illo ítaque nostri géneris córpore nos Christi mors vivificávit, nos resurréctio eréxit, nos ascénsio consecrávit. In illo nostræ oríginis córpore cæléstibus regnis arrham nostræ conditiónis impósuit. In illo ítaque nostri géneris córpore nos Christi mors vivificávit, nos resurréctio eréxit, nos ascénsio consecrávit. In illo nostræ oríginis córpore cæléstibus regnis arrham nostræ conditiónis impósuit.\par
Elaborémus ergo, caríssimi, ut quemádmodum Dóminus in hac die nostro cum córpore ad supérna conscéndit, ita nos post illum quómodo póssumus, spe ascendámus, et corde sequámur. Ipso afféctu páriter et proféctu ascendámus post illum: étiam per vítia ac passiónes nostras. Si útique unusquísque nostrum súbdere eas sibi stúdeat, ac super eas stare consuéscat, ex ipsis sibi gradum cónstruit, quo possit ad superióra conscéndere. Elevábunt nos, si fúerint infra nos.\par
\switchcolumn
\EN
\noindent It was therefore just because His Body was of our nature, that Christ's Death hath quickened us, His Resurrection raised us up, His Ascension sanctified us. It was just because His Body was of our nature, that His Presence in the heavenly kingdoms is a pledge that we also shall one day be there. Let us therefore strive, dearly beloved brethren, since the Lord hath on this day gone up on high in a Body of our nature, ourselves, as far as we can, to ascend thither in hope, to follow Him with our heart. Let us ascend to Him in love, and speed keeping pace with love, even by our very sins and passions. If every one of us would strive to get above them, and accustom himself to tread on them, he might make of even them a stepping-stone to mount to higher things. Such things lift us up if they are underneath us.\par
\switchcolumn*

\LA
\begin{Resp}\Rsign Non turbétur cor vestrum: ego vado ad Patrem; et cum assúmptus fúero a vobis, mittam vobis, allelúja, \Star{} Spíritum veritátis, et gaudébit cor vestrum, allelúja.\end{Resp}
\begin{Resp}\Vsign Ego rogábo Patrem, et álium Paráclitum dabit vobis.\end{Resp}
\begin{Resp}\Rsign Spíritum veritátis, et gaudébit cor vestrum, allelúja.\end{Resp}
\switchcolumn
\EN
\begin{Resp}\Rsign Let not your heart be troubled; I go unto the Father, and when I am taken from you, I will send unto you, alleluia, \Star{} The Spirit of truth, and your heart shall rejoice. Alleluia.\end{Resp}
\begin{Resp}\Vsign I will pray the Father, and He shall give you another Comforter.\end{Resp}
\begin{Resp}\Rsign The Spirit of truth, and your heart shall rejoice. Alleluia.\end{Resp}
\switchcolumn*

\LA
\LessonHead{Lectio VI}
\switchcolumn
\EN
\LessonHead{Lesson VI}
\switchcolumn*

\LA
\noindent De vítiis nostris scalam nobis fácimus, si vítia ipsa calcámus. Nam cum bonitátis auctóre non ascéndit malítia, nec cum Fílio Vírginis libído atque luxúria. Non, inquam, ascéndunt vítia post virtútum paréntem, peccáta post justum, nec infirmitátes ac morbi possunt ire post médicum. Igitur, si intráre ipsíus médici regnum vólumus, prius vúlnera nostra curémus. Ordinémus et custodiámus in nobis statum utriúsque substántiæ; ne ánimam, nobiliórem útique hóminis portiónem, tártaro pars devólvat inférior: sed secum pótius cælo sanctificátum corpus acquírat natúra gloriósior: ipso adjuvánte, qui vivit et regnat in sǽcula sæculórum. Amen.\par
\switchcolumn
\EN
\noindent We make our vices a ladder, if we tread them down. With the Author of goodness there ascended no spite with the Son of the Virgin, no lust or sensuality. I say vices do not follow to heaven the Father of perfection, sin the Holy One of God, neither weakness nor disease the Divine Healer. If therefore we would enter into the kingdom of that Healer, we must first take heed to our sores. We must so order and guard in us the mutual relations of our soul and body, that the soul, the nobler part of man, may not be dragged down to hell by her grovelling companion, but may rather, being herself of a nature more glorious, bear with her to heaven at the last a sanctified body, by the help of Him Who liveth and reigneth for ever and ever. Amen.\par
\switchcolumn*

\LA
\begin{Resp}\Rsign Ascéndens Christus in altum, captívam duxit captivitátem, \Star{} Dedit dona homínibus, allelúja, allelúja, allelúja.\end{Resp}
\begin{Resp}\Vsign Ascéndit Deus in jubilatióne, et Dóminus in voce tubæ.\end{Resp}
\begin{Resp}\Rsign Dedit dona homínibus, allelúja, allelúja, allelúja.\end{Resp}
\begin{Resp}\Vsign Glória Patri, et Fílio, \Star{} et Spirítui Sancto.\end{Resp}
\begin{Resp}\Rsign Dedit dona homínibus, allelúja, allelúja, allelúja.\end{Resp}
\switchcolumn
\EN
\begin{Resp}\Rsign When Christ ascended up on high, He led captivity captive, \Star{} He gave gifts unto men. Alleluia, Alleluia, Alleluia.\end{Resp}
\begin{Resp}\Vsign God is gone up with a shout, and the Lord with the sound of a trumpet.\end{Resp}
\begin{Resp}\Rsign He gave gifts unto men. Alleluia, Alleluia, Alleluia.\end{Resp}
\begin{Resp}\Vsign Glory be to the Father, and to the Son, \Star{} and to the Holy Ghost.\end{Resp}
\begin{Resp}\Rsign He gave gifts unto men. Alleluia, Alleluia, Alleluia.\end{Resp}
\switchcolumn*

\LA
\LessonHead{Lectio VII}
\switchcolumn
\EN
\LessonHead{Lesson VII}
\switchcolumn*

\LA
\LessonTitle{Léctio sancti Evangélii secúndum Joánnem}
\switchcolumn
\EN
\LessonTitle{From the Holy Gospel according to John}
\switchcolumn*

\LA
\Cite{Joannes 15:26-27; 16:1-4}
\switchcolumn
\EN
\Cite{John 15:26-27; 16:1-4}
\switchcolumn*

\LA
\noindent In illo témpore: Dixit Jesus discípulis suis: Cum vénerit Paráclitus, quem ego mittam vobis a Patre, Spíritum veritátis, qui a Patre procédit, ille testimónium perhibébit de me. Et réliqua.\par
\switchcolumn
\EN
\noindent At that time, Jesus said unto His disciples: When the Comforter is come, whom I will send unto you from the Father, even the Spirit of truth, Which proceedeth from the Father, He shall testify of Me. And so on.\par
\switchcolumn*

\LA
\LessonTitle{Homilía sancti Augustíni Epíscopi}
\switchcolumn
\EN
\LessonTitle{Homily by St. Augustine, Bishop of Hippo.}
\switchcolumn*

\LA
\Source{Tract. 92 in Joann., circa medium}
\switchcolumn
\EN
\Source{92nd Tract on John.}
\switchcolumn*

\LA
\noindent Venit die Pentecóstes Spíritus Sanctus in centum vigínti hómines congregátos, in quibus et Apóstoli omnes erant: qui illo adimpléti, cum linguis ómnium géntium loqueréntur, plures ex his, qui áderant, tanto miráculo stupefácti (quandóquidem vidérunt, loquénte Petro, tam magnum atque divínum testimónium perhibéri de Christo, ut ille, qui occísus ab eis inter mórtuos deputabátur resurrexísse et vívere probarétur) compúncti corde convérsi sunt, et tanti sánguinis, tam ímpie atque immániter fusi, indulgéntiam percepérunt, ipso redémpti sánguine quem fudérunt.\par
\switchcolumn
\EN
\noindent Upon the day of Pentecost the Holy Ghost came down upon a congregation of an hundred and twenty men, among whom were all the Apostles. These men, after they had been filled with the Spirit, began to speak with the tongues of all nations, and many of the bystanders, amazed at the marvel, when they saw in the discourse of Peter, how great and how Divine a witness was borne to the fact that the Christ, Whom they had murdered, and Whom they reckoned among the dead, had risen again and was alive, many of these bystanders were pricked in their heart (Acts ii. 37) and were converted. They received pardon from that noble Blood, Which they had so sacrilegiously and so brutally shed, seeing that that Blood had redeemed even Its Own out-pourers.\par
\switchcolumn*

\LA
\begin{Resp}\Rsign Ego rogábo Patrem, et álium Paráclitum dabit vobis, \Star{} Ut máneat vobíscum in ætérnum, Spíritum veritátis, allelúja.\end{Resp}
\begin{Resp}\Vsign Si enim non abíero, Paráclitus non véniet ad vos: si autem abíero, mittam eum ad vos.\end{Resp}
\begin{Resp}\Rsign Ut máneat vobíscum in ætérnum, Spíritum veritátis, allelúja.\end{Resp}
\switchcolumn
\EN
\begin{Resp}\Rsign I will pray the Father, and He shall give you another Comforter, \Star{} That He may abide with you for ever, even the Spirit of truth. Alleluia.\end{Resp}
\begin{Resp}\Vsign For if I go not away, the Comforter will not come unto you; but if I depart, I will send Him unto you.\end{Resp}
\begin{Resp}\Rsign That He may abide with you for ever, even the Spirit of truth. Alleluia.\end{Resp}
\switchcolumn*

\LA
\LessonHead{Lectio VIII}
\switchcolumn
\EN
\LessonHead{Lesson VIII}
\switchcolumn*

\LA
\noindent Christi enim sanguis sic in remissiónem peccatórum ómnium fusus est, ut ipsum étiam peccátum posset delére, quo fusus est. Hoc ergo íntuens Dóminus dicébat: Odio habuérunt me gratis: cum autem vénerit Paráclitus, ille testimónium perhibébit de me. Tamquam díceret: Odio me habuérunt, et occidérunt vidéntes: sed tale de me Paráclitus testimónium perhibébit, ut eos fáciat in me crédere non vidéntes. Et vos, inquit, testimónium perhibébitis, quia ab inítio mecum estis. Perhibébit Spíritus Sanctus, perhibébitis et vos. Quia enim ab inítio mecum estis, potéstis prædicáre quod nostis: quod ut modo non faciátis, illíus Spíritus plenitúdo nondum adest vobis.\par
\switchcolumn
\EN
\noindent The Blood of Christ “Which is shed for many for the remission of sins” was so effectually shed, that It could remit even the very sin that shed It. Toward this looked the Lord when He said “They hated Me without a cause but when the Comforter is come, Whom I will send unto you from the Father, He shall testify of Me.” This was as though He had said They have hated Me and slain Me while they see Me, but when they shall see Me no more, the Comforter shall bear such testimony of Me, as will compel them to believe in Me. “And ye also,” saith He, “shall bear witness, because ye have been with Me from the beginning,” the Holy Ghost shall bear witness, and ye also shall bear witness. “Because ye have been with Me from the beginning,” ye are able to speak that ye do know, (John iii. 11,) which ye do not now, while as yet the fulness of the Spirit is not come upon you.\par
\switchcolumn*

\LA
\begin{Resp}\Rsign Si enim non abíero, Paráclitus non véniet ad vos: si autem abíero, mittam eum ad vos. \Star{} Cum autem vénerit ille, docébit vos omnem veritátem, allelúja.\end{Resp}
\begin{Resp}\Vsign Non enim loquétur a semetípso: sed quæcúmque áudiet, loquétur: et quæ ventúra sunt, annuntiábit vobis.\end{Resp}
\begin{Resp}\Rsign Cum autem vénerit ille, docébit vos omnem veritátem, allelúja.\end{Resp}
\begin{Resp}\Vsign Glória Patri, et Fílio, \Star{} et Spirítui Sancto.\end{Resp}
\begin{Resp}\Rsign Cum autem vénerit ille, docébit vos omnem veritátem, allelúja.\end{Resp}
\switchcolumn
\EN
\begin{Resp}\Rsign For if I go not away, the Comforter will not come unto you; but if I depart, I will send Him unto you; \Star{} And when He is come, He will guide you into all truth. Alleluia.\end{Resp}
\begin{Resp}\Vsign For He shall not speak of Himself; but whatsoever He shall hear, that shall He speak and He will show you things to come.\end{Resp}
\begin{Resp}\Rsign And when He is come, He will guide you into all truth. Alleluia.\end{Resp}
\begin{Resp}\Vsign Glory be to the Father, and to the Son, \Star{} and to the Holy Ghost.\end{Resp}
\begin{Resp}\Rsign And when He is come, He will guide you into all truth. Alleluia.\end{Resp}
\switchcolumn*

\LA
\LessonHead{Lectio IX}
\switchcolumn
\EN
\LessonHead{Lesson IX}
\switchcolumn*

\LA
\noindent Ille ergo testimónium perhibébit de me, et vos testimónium perhibébitis: dabit enim vobis fidúciam testimónium perhibéndi cáritas Dei diffúsa in córdibus vestris per Spíritum Sanctum, qui dábitur vobis. Quæ útique Petro adhuc défuit, quando mulíeris ancíllæ interrogatióne pertérritus, non pótuit verum testimónium perhibére, sed contra suam pollicitatiónem timóre magno compúlsus est ter negáre. Timor autem iste non est in caritáte: sed perfécta cáritas foras mittit timórem. Dénique ante passiónem Dómini, servílis timor ejus interrogátus est a fémina servitútis; post resurrectiónem vero Dómini, liberális ejus amor ab ipso príncipe libertátis: et ídeo ibi turbabátur, hic tranquillabátur; ibi quem diléxerat, negábat; hic quem negáverat, diligébat. Sed adhuc étiam tunc amor ipse infírmus fúerat et angústus, donec eum roboráret et dilatáret Spíritus Sanctus.\par
\switchcolumn
\EN
\noindent He shall testify of Me and ye also shall bear witness when the love of God is shed abroad in your hearts by the Holy Ghost Which shall be given unto you, and maketh you not ashamed to lift up your testimony. This love had not been so shed abroad in Peter's heart when he was frightened by the questioning of the maid-servant, and could not bear witness to the truth, but brake his promise, and was driven by strong fear to deny Christ thrice. “There is no such fear in love; but perfect love casteth out fear.” Before the Passion of the Lord, Peter's slavish fear was questioned by a bondwoman but after the Resurrection of the Lord, his free love was asked by the very Prince of freedom, and therefore the first questioning shook him, but under the second he was at peace at the first he denied Him Whom he had loved at the second he loved Him Whom he had denied. But, even so, his love was weak and narrow, until the Holy Ghost had strengthened and widened it.\par
\switchcolumn*

\LA
\TeDeum{Te Deum}
\switchcolumn
\EN
\TeDeum{Te Deum}
\switchcolumn*

\end{paracol}

\Office{Sabbato in Vigilia Pentecostes}{Vigil of Pentecost}{Semiduplex I classis}
\addcontentsline{toc}{subsection}{Sabbato in Vigilia Pentecostes\enspace\textit{Vigil of Pentecost}}
\begin{paracol}{2}
\LA
\LessonHead{Lectio I}
\switchcolumn
\EN
\LessonHead{Lesson I}
\switchcolumn*

\LA
\LessonTitle{Incipit Epístola cathólica beáti Judæ Apóstoli}
\switchcolumn
\EN
\LessonTitle{Beginning of the letter of the blessed Apostle Jude}
\switchcolumn*

\LA
\Cite{Judas 1:1-4}
\switchcolumn
\EN
\Cite{Jude 1:1-4}
\switchcolumn*

\LA
\noindent Judas, Jesu Christi servus, frater autem Jacóbi, his qui sunt in Deo Patre diléctis, et Christo Jesu conservátis, et vocátis. Misericórdia vobis, et pax, et cáritas adimpleátur. Caríssimi, omnem sollicitúdinem fáciens scribéndi vobis de commúni vestra salúte, necésse hábui scríbere vobis: déprecans supercertári semel tráditæ sanctis fídei. Subintroiérunt enim quidam hómines (qui olim præscrípti sunt in hoc judícium) ímpii, Dei nostri grátiam transferéntes in luxúriam, et solum Dominatórem, et Dóminum nostrum Jesum Christum negántes.\par
\switchcolumn
\EN
\noindent Jude, the servant of Jesus Christ, and brother of James: to them that are beloved in God the Father, and preserved in Jesus Christ, and called. Mercy unto you, and peace, and charity be fulfilled. Dearly beloved, taking all care to write unto you concerning your common salvation, I was under a necessity to write unto you: to beseech you to contend earnestly for the faith once delivered to the saints. For certain men are secretly entered in, (who were written of long ago unto this judgment,) ungodly men, turning the grace of our Lord God into riotousness, and denying the only sovereign Ruler, and our Lord Jesus Christ.\par
\switchcolumn*

\LA
\begin{Resp}\Rsign Post passiónem suam per dies quadragínta appárens eis, et loquens de regno Dei, allelúja: \Star{} Et, vidéntibus illis, elevátus est, allelúja: et nubes suscépit eum ab óculis eórum, allelúja.\end{Resp}
\begin{Resp}\Vsign Et convéscens, præcépit eis, ab Jerosólymis ne discéderent, sed exspectárent promissiónem Patris.\end{Resp}
\begin{Resp}\Rsign Et, vidéntibus illis, elevátus est, allelúja: et nubes suscépit eum ab óculis eórum, allelúja.\end{Resp}
\switchcolumn
\EN
\begin{Resp}\Rsign Being seen of them forty days after that He had suffered, and speaking of the kingdom of God. Alleluia. \Star{} And while they beheld, He was taken up, and a cloud received Him out of their sight. Alleluia.\end{Resp}
\begin{Resp}\Vsign And, eating together with them, He commanded them that they should not depart from Jerusalem, but wait for the Promise of the Father.\end{Resp}
\begin{Resp}\Rsign And while they beheld, He was taken up, and a cloud received Him out of their sight. Alleluia.\end{Resp}
\switchcolumn*

\LA
\LessonHead{Lectio II}
\switchcolumn
\EN
\LessonHead{Lesson II}
\switchcolumn*

\LA
\Cite{Judas 1:5-8}
\switchcolumn
\EN
\Cite{Jude 1:5-8}
\switchcolumn*

\LA
\noindent Commonére autem vos volo, sciéntes semel ómnia, quóniam Jesus pópulum de terra Ægýpti salvans, secúndo eos, qui non credidérunt, pérdidit: ángelos vero, qui non servavérunt suum principátum, sed dereliquérunt suum domicílium, in judícium magni diéi, vínculis ætérnis sub calígine reservávit. Sicut Sódoma, et Gomórrha, et finítimæ civitátes símili modo exfornicátæ, et abeúntes post carnem álteram, factæ sunt exémplum, ignis ætérni pœnam sustinéntes. Simíliter et hi carnem quidem máculant, dominatiónem autem spernunt, majestátem autem blasphémant.\par
\switchcolumn
\EN
\noindent I will therefore admonish you, though ye once knew all things, that Jesus, having saved the people out of the land of Egypt, did afterwards destroy them that believed not: And the angels who kept not their principality, but forsook their own habitation, he hath reserved under darkness in everlasting chains, unto the judgment of the great day. As Sodom and Gomorrha, and the neighbouring cities, in like manner, having given themselves to fornication, and going after other flesh, were made an example, suffering the punishment of eternal fire. In like manner these men also defile the flesh, and despise dominion, and blaspheme majesty.\par
\switchcolumn*

\LA
\begin{Resp}\Rsign Omnis pulchritúdo Dómini exaltáta est super sídera: \Star{} Spécies ejus in núbibus cæli, et nomen ejus in ætérnum pérmanet, allelúja.\end{Resp}
\begin{Resp}\Vsign A summo cælo egréssio ejus, et occúrsus ejus usque ad summum ejus.\end{Resp}
\begin{Resp}\Rsign Spécies ejus in núbibus cæli, et nomen ejus in ætérnum pérmanet, allelúja.\end{Resp}
\switchcolumn
\EN
\begin{Resp}\Rsign The Lord hath set His beauty above the stars; \Star{} His loveliness is in the clouds of heaven, and His Name endureth for ever. Alleluia.\end{Resp}
\begin{Resp}\Vsign His going forth is from the end of the heaven, and His circuit unto the ends of it.\end{Resp}
\begin{Resp}\Rsign His loveliness is in the clouds of heaven, and His Name endureth for ever. Alleluia.\end{Resp}
\switchcolumn*

\LA
\LessonHead{Lectio III}
\switchcolumn
\EN
\LessonHead{Lesson III}
\switchcolumn*

\LA
\Cite{Judas 1:9-13}
\switchcolumn
\EN
\Cite{Jude 1:9-13}
\switchcolumn*

\LA
\noindent Cum Míchael Archángelus cum diábolo dísputans altercarétur de Móysi córpore, non est ausus judícium inférre blasphémiæ: sed dixit: Imperet tibi Dóminus. Hi autem quæcúmque quidem ignórant, blasphémant: quæcúmque autem naturáliter, tamquam muta animália, norunt, in his corrumpúntur. Væ illis, quia in via Cain abiérunt, et erróre Bálaam mercéde effúsi sunt, et in contradictióne Core periérunt! Hi sunt in épulis suis máculæ, convivántes sine timóre, semetípsos pascéntes, nubes sine aqua, quæ a ventis circumferéntur, árbores autumnáles, infructuósæ, bis mórtuæ, eradicátæ, fluctus feri maris, despumántes suas confusiónes, sídera errántia: quibus procélla tenebrárum serváta est in ætérnum.\par
\switchcolumn
\EN
\noindent When Michael the archangel, disputing with the devil, contended about the body of Moses, he durst not bring against him the judgment of railing speech, but said: The Lord command thee. But these men blaspheme whatever things they know not: and what things soever they naturally know, like dumb beasts, in these they are corrupted. Woe unto them, for they have gone in the way of Cain: and after the error of Balaam they have for reward poured out themselves, and have perished in the contradiction of Core. These are spots in their banquets, feasting together without fear, feeding themselves, clouds without water, which are carried about by winds, trees of the autumn, unfruitful, twice dead, plucked up by the roots, Raging waves of the sea, foaming out their own confusion; wandering stars, to whom the storm of darkness is reserved for ever.\par
\switchcolumn*

\LA
\begin{Resp}\Rsign Exaltáre, Dómine, allelúja, \Star{} In virtúte tua, allelúja.\end{Resp}
\begin{Resp}\Vsign Eleváta est, magnificéntia tua super cælos, Deus.\end{Resp}
\begin{Resp}\Rsign In virtúte tua, allelúja.\end{Resp}
\begin{Resp}\Vsign Glória Patri, et Fílio, \Star{} et Spirítui Sancto.\end{Resp}
\begin{Resp}\Rsign In virtúte tua, allelúja.\end{Resp}
\switchcolumn
\EN
\begin{Resp}\Rsign Be Thou exalted, O Lord. Alleluia. \Star{} In thine Own strength. Alleluia.\end{Resp}
\begin{Resp}\Vsign O God, Thou hast set thy glory above the heavens.\end{Resp}
\begin{Resp}\Rsign In thine Own strength. Alleluia.\end{Resp}
\begin{Resp}\Vsign Glory be to the Father, and to the Son, \Star{} and to the Holy Ghost.\end{Resp}
\begin{Resp}\Rsign In thine Own strength. Alleluia.\end{Resp}
\switchcolumn*

\LA
\LessonHead{Lectio IV}
\switchcolumn
\EN
\LessonHead{Lesson IV}
\switchcolumn*

\LA
\LessonTitle{Ex Tractátu sancti Augustíni Epíscopi de Symbolo ad Catechúmenos}
\switchcolumn
\EN
\LessonTitle{From the Treatise upon the Creed, addressed to Catechumens by St. Augustine, Bishop of Hippo}
\switchcolumn*

\LA
\Source{Lib. 4. cap. 1. tom. 9.}
\switchcolumn
\EN
\Source{Bk. iv. ch. I. tom. 9}
\switchcolumn*

\LA
\noindent Dum per sacratíssimum crucis signum vos suscépit in útero sancta mater Ecclésia, quæ sicut et fratres vestros cum summa lætítia spiritáliter páriet, nova proles futúra tantæ matris, quoúsque per lavácrum sanctum regenerátos veræ luci restítuat, cóngruis aliméntis eos, quos portat, pascat in útero, et ad diem partus sui lætos læta perdúcat: quóniam non tenétur hæc senténtia Hevæ, quæ in tristítia et gémitu parit fílios; nec ipsos gaudéntes, sed pótius flentes. Hæc enim solvit, quod illa ligáverat: ut prolem, quam per inobediéntiam sui, morti donávit, hæc per obediéntiam restítuat vitæ. Omnia sacraménta, quæ acta sunt et agúntur in vobis per ministérium servórum Dei, exorcísmis, oratiónibus, cánticis spirituálibus, insufflatiónibus, cilício, inclinatióne cervícum, humilitáte pedum, pavor ipse omni securitáte appeténdus: hæc ómnia, ut dixi, escæ sunt, quæ vos refíciunt in útero, ut renátos ex baptísmo hílares vos mater exhíbeat Christo.\par
\switchcolumn
\EN
\noindent We are yet the unborn offspring of a great Mother. Our Holy Mother the Church hath by the most sacred sign of the Cross received you into her womb, and from thence she is now just about to bring you forth, as she hath already brought forth your brethren, with thrills of spiritual joy. But until, through the washing of regeneration, she bringeth you forth into true light, she feedeth you in her womb with such food as becometh your condition, and in gladness matureth her children for the glad moment of her delivery. This Mother is not stricken by the doom of Eve, to bring forth children in sorrow, and they themselves oftenertimes weeping than laughing. Rather doth your spiritual Mother annul the sentence of your earthly Eve, by disobedience, endowed her offspring with death the Church, by obedience, giveth them newness of life. All the mystic prayers and ceremonies which have been and are still being performed over you by the ministry of the servants of God, exorcisms, prayers, spiritual songs, onbreathings, haircloth, prostrations, baring of the feet, the dread which ye feel, albeit so safe, all these things, I say unto you, are the nourishment which ye are ever. drawing from your Mother while yet ye are in her womb, that at the baptismal birth she may be able to present you strong and laughing babes unto Christ.\par
\switchcolumn*

\LA
\begin{Resp}\Rsign Tempus est, ut revértar ad eum, qui me misit, dicit Dóminus: nolíte contristári, nec turbétur cor vestrum: \Star{} Rogo pro vobis Patrem, ut ipse vos custódiat, allelúja, allelúja.\end{Resp}
\begin{Resp}\Vsign Nisi ego abíero, Paráclitus non véniet: cum assúmptus fúero, mittam vobis eum.\end{Resp}
\begin{Resp}\Rsign Rogo pro vobis Patrem, ut ipse vos custódiat, allelúja, allelúja.\end{Resp}
\switchcolumn
\EN
\begin{Resp}\Rsign My time is come that I should return unto Him That sent Me, saith the Lord. Be not sorrowful, neither let your heart be troubled. \Star{} I pray the Father for you, that He may keep you. Alleluia, Alleluia.\end{Resp}
\begin{Resp}\Vsign If I go not away, the Comforter will not come unto you; where I am ascended, I will send Him unto you.\end{Resp}
\begin{Resp}\Rsign I pray the Father for you, that He may keep you. Alleluia, Alleluia.\end{Resp}
\switchcolumn*

\LA
\LessonHead{Lectio V}
\switchcolumn
\EN
\LessonHead{Lesson V}
\switchcolumn*

\LA
\noindent Accepístis et sýmbolum, protectiónem parturiéntis contra venéna serpéntis. In Apocalýpsi Joánnis Apóstoli scriptum est hoc, quod staret draco in conspéctu mulíeris, quæ paritúra erat, ut cum peperísset, natum ejus coméderet. Dracónem diábolum esse, nullus vestrum ignórat: mulíerem illam Vírginem Maríam significásse, quæ caput nostrum íntegra íntegrum péperit; quæ étiam ipsa figúram in se sanctæ Ecclésiæ demonstrávit: ut quómodo Fílium páriens, Virgo permánsit, ita et hæc omni témpore membra ejus páriat, et virginitátem non amíttat. Ipsas senténtias sacratíssimi sýmboli adjuvánte Dómino exponéndas suscépimus, ut, quid síngulæ contíneant, vestris sénsibus intimémus. Paráta sunt corda vestra, quia exclúsus est inimícus de córdibus vestris.\par
\switchcolumn
\EN
\noindent We have also received the Creed, which is the shield of the travailing Mother against the venom of the dragon. In the Apocalypse of the Apostle John (xii. 4) it is written “And the dragon stood before the woman which was ready to be delivered, for to devour her child as soon as it was born.” That this dragon is the devil ye all know. Ye know likewise that by the woman is signified the Virgin Mary, who, herself a Virgin, bore our Virgin Head, and who is revealed unto us as a type of the Holy Church, in that, even as Mary, though she bore a Son, remained a Virgin, so the Church doth in all times give birth to all her members, and yet is ever presented a chaste virgin to Christ. I have undertaken, with the help of the Lord, to expound every clause of the Creed, that I may bring home to your understandings what each containeth. Your hearts are ready, for the enemy hath been shut out of your hearts.\par
\switchcolumn*

\LA
\begin{Resp}\Rsign Non turbétur cor vestrum: ego vado ad Patrem; et cum assúmptus fúero a vobis, mittam vobis, allelúja, \Star{} Spíritum veritátis, et gaudébit cor vestrum, allelúja.\end{Resp}
\begin{Resp}\Vsign Ego rogábo Patrem, et álium Paráclitum dabit vobis.\end{Resp}
\begin{Resp}\Rsign Spíritum veritátis, et gaudébit cor vestrum, allelúja.\end{Resp}
\switchcolumn
\EN
\begin{Resp}\Rsign Let not your heart be troubled; I go unto the Father, and when I am taken from you, I will send unto you, alleluia, \Star{} The Spirit of truth, and your heart shall rejoice. Alleluia.\end{Resp}
\begin{Resp}\Vsign I will pray the Father, and He shall give you another Comforter.\end{Resp}
\begin{Resp}\Rsign The Spirit of truth, and your heart shall rejoice. Alleluia.\end{Resp}
\switchcolumn*

\LA
\LessonHead{Lectio VI}
\switchcolumn
\EN
\LessonHead{Lesson VI}
\switchcolumn*

\LA
\noindent Huic vos renuntiáre proféssi estis: in qua professióne, non homínibus, sed Deo et Angelis ejus conscribéntibus dixístis, Renúntio. Renuntiáte non solum vócibus, sed étiam móribus: non tantum sono linguæ, sed et actu vitæ: nec tantum lábiis sonántibus, sed opéribus pronuntiántibus. Scitóte vos cum cállido, antíquo, et veternóso inimíco suscepísse certámen: non in vobis post renuntiatiónem invéniat ópera sua, non jure vos áttrahat in servitútem suam. Deprehénderis enim, et detégeris Christiáne, quando áliud agis, et áliud profitéris: fidélis in nómine, áliud demónstrans in ópere, non tenens promissiónis tuæ fidem: modo ingrédiens ecclésiam oratiónes fúndere, post módicum in spectáculis cum histriónibus impudíce clamáre. Quid tibi cum pompis diáboli, quibus renuntiásti?\par
\switchcolumn
\EN
\noindent We have made profession of renouncing the enemy. At the moment of that profession it was not before men only, but in the presence of God and His Angels that ye said: “I do renounce him.” Renounce him, not only in your words, but in your ways not only with your voices, but with your lives not only with your lips, but in your works. Know ye well that the wrestling which ye have undertaken is a strife with an enemy who is subtle, and old, and patient now that ye have once renounced him, let him never again find in you his works never again give him the right to bring you into bondage. O Christian thou wilt be caught and exposed, if thou dost one thing and professest another, if thou art faithful in name, and makest it to be evident by thy works that thou hast broken the faith pledged by this promise if some while thou goest into a church to pray, and anon to the shows to join in applauding obscene representations. What hast thou to do any more with the pomps of the devil, which thou hast renounced?\par
\switchcolumn*

\LA
\begin{Resp}\Rsign Ascéndens Christus in altum, captívam duxit captivitátem, \Star{} Dedit dona homínibus, allelúja, allelúja, allelúja.\end{Resp}
\begin{Resp}\Vsign Ascéndit Deus in jubilatióne, et Dóminus in voce tubæ.\end{Resp}
\begin{Resp}\Rsign Dedit dona homínibus, allelúja, allelúja, allelúja.\end{Resp}
\begin{Resp}\Vsign Glória Patri, et Fílio, \Star{} et Spirítui Sancto.\end{Resp}
\begin{Resp}\Rsign Dedit dona homínibus, allelúja, allelúja, allelúja.\end{Resp}
\switchcolumn
\EN
\begin{Resp}\Rsign When Christ ascended up on high, He led captivity captive, \Star{} He gave gifts unto men. Alleluia, Alleluia, Alleluia.\end{Resp}
\begin{Resp}\Vsign God is gone up with a shout, and the Lord with the sound of a trumpet.\end{Resp}
\begin{Resp}\Rsign He gave gifts unto men. Alleluia, Alleluia, Alleluia.\end{Resp}
\begin{Resp}\Vsign Glory be to the Father, and to the Son, \Star{} and to the Holy Ghost.\end{Resp}
\begin{Resp}\Rsign He gave gifts unto men. Alleluia, Alleluia, Alleluia.\end{Resp}
\switchcolumn*

\LA
\LessonHead{Lectio VII}
\switchcolumn
\EN
\LessonHead{Lesson VII}
\switchcolumn*

\LA
\LessonTitle{Léctio sancti Evangélii secúndum Joánnem}
\switchcolumn
\EN
\LessonTitle{From the Holy Gospel according to John}
\switchcolumn*

\LA
\Cite{Joannes 14:15-21}
\switchcolumn
\EN
\Cite{John 14:15-21}
\switchcolumn*

\LA
\noindent In illo témpore: Dixit Jesus discípulis suis: Si dilígitis me, mandáta mea serváte. Et ego rogábo Patrem, et álium Paráclitum dabit vobis. Et réliqua.\par
\switchcolumn
\EN
\noindent At that time, Jesus said unto His disciples: If ye love Me, keep My commandments. And I will pray the Father, and He shall give you another Comforter. And so on.\par
\switchcolumn*

\LA
\LessonTitle{Homilía sancti Augustíni Epíscopi}
\switchcolumn
\EN
\LessonTitle{Homily by St. Augustine, Bishop of Hippo}
\switchcolumn*

\LA
\Source{Tractatus 74 in Joannem, sub finem, et 75}
\switchcolumn
\EN
\Source{74th and 75th Tracts on John}
\switchcolumn*

\LA
\noindent Quod ait, Rogábo Patrem, et álium Paráclitum dabit vobis: osténdit et seípsum esse Paráclitum. Paráclitus enim Latíne dícitur advocátus: et dictum est de Christo: Advocátum habémus ad Patrem, Jesum Christum justum. Sic autem mundum dixit non posse accípere Spíritum Sanctum, sicut étiam dictum est: Prudéntia carnis inimíca est Deo: legi enim Dei non est subjécta, nec enim potest: velut si dicámus: Injustítia justítia esse non potest. Mundum quippe ait hoc loco, mundi signíficans dilectóres: quæ diléctio non est a Patre. Et ídeo dilectióni hujus mundi, de qua satágimus ut minuátur et consumátur in nobis, contrária est diléctio Dei, quæ diffúnditur in córdibus nostris per Spíritum Sanctum, qui datus est nobis.\par
\switchcolumn
\EN
\noindent By these words of the Lord: “I will pray the Father, and He shall give you another Comforter”, He doth imply that Himself is a Comforter. The Greek word used, namely “Parakletos,” signifieth also an Advocate, and is used in that sense where it is written “We have an Advocate Parakleton with the Father, Jesus Christ the Righteous.” “Even the Spirit of truth, Whom the world cannot receive,” because as we read elsewhere, “the carnal mind is enmity against God; for it is not subject to the law of God, neither indeed can it be” as we may say plainly nothing can make unrighteousness righteous. By “the world,” in this place, we must understand the lovers of the world, a love which cometh not of the Father. And therefore it is that this love of the world, which we strive to lessen and to destroy in ourselves, is contrary to “the love of God, which is shed abroad in our hearts by the Holy Ghost which is given unto us.”\par
\switchcolumn*

\LA
\begin{Resp}\Rsign Ego rogábo Patrem, et álium Paráclitum dabit vobis, \Star{} Ut máneat vobíscum in ætérnum, Spíritum veritátis, allelúja.\end{Resp}
\begin{Resp}\Vsign Si enim non abíero, Paráclitus non véniet ad vos: si autem abíero, mittam eum ad vos.\end{Resp}
\begin{Resp}\Rsign Ut máneat vobíscum in ætérnum, Spíritum veritátis, allelúja.\end{Resp}
\switchcolumn
\EN
\begin{Resp}\Rsign I will pray the Father, and He shall give you another Comforter, \Star{} That He may abide with you for ever, even the Spirit of truth. Alleluia.\end{Resp}
\begin{Resp}\Vsign For if I go not away, the Comforter will not come unto you; but if I depart, I will send Him unto you.\end{Resp}
\begin{Resp}\Rsign That He may abide with you for ever, even the Spirit of truth. Alleluia.\end{Resp}
\switchcolumn*

\LA
\LessonHead{Lectio VIII}
\switchcolumn
\EN
\LessonHead{Lesson VIII}
\switchcolumn*

\LA
\noindent Mundus ergo eum accípere non potest, quia non videt eum, neque scit eum. Non enim habet invisíbiles óculos mundána diléctio, per quos vidéri Spíritus Sanctus potest, qui vidéri nisi invisibíliter non potest. Vos autem, inquit, cognoscétis eum: quia apud vos manébit, et in vobis erit. Erit in eis, ut máneat; non manébit, ut sit: prius est enim esse alícubi, quam manére. Sed ne putárent quod dictum est, Apud vos manébit; ita dictum, quemádmodum apud hóminem hospes visibíliter manére consuévit, expósuit quid díxerit: Apud vos manébit, cum adjúnxit et dixit, In vobis erit.\par
\switchcolumn
\EN
\noindent The Spirit of truth Whom the world cannot receive, because it seeth Him not, neither knoweth Him “for to love the world is to lack those spiritual eyes, which are able to see Him Who is invisible, the Holy Ghost. “But ye know Him,” saith the Lord to His disciples, “for He shall dwell with you, and shall be in you.” He will be in them to dwell in them, not dwell in them to be in them for one must first be in a place before one dwell there. But lest the Apostles should think that the words, “He shall dwell with you,” signified that He should visibly abide with them for a while, as do guests in the houses of men, the Lord saith in explanation “He shall be in you.”\par
\switchcolumn*

\LA
\begin{Resp}\Rsign Si enim non abíero, Paráclitus non véniet ad vos: si autem abíero, mittam eum ad vos. \Star{} Cum autem vénerit ille, docébit vos omnem veritátem, allelúja.\end{Resp}
\begin{Resp}\Vsign Non enim loquétur a semetípso: sed quæcúmque áudiet, loquétur: et quæ ventúra sunt, annuntiábit vobis.\end{Resp}
\begin{Resp}\Rsign Cum autem vénerit ille, docébit vos omnem veritátem, allelúja.\end{Resp}
\begin{Resp}\Vsign Glória Patri, et Fílio, \Star{} et Spirítui Sancto.\end{Resp}
\begin{Resp}\Rsign Cum autem vénerit ille, docébit vos omnem veritátem, allelúja.\end{Resp}
\switchcolumn
\EN
\begin{Resp}\Rsign For if I go not away, the Comforter will not come unto you; but if I depart, I will send Him unto you; \Star{} And when He is come, He will guide you into all truth. Alleluia.\end{Resp}
\begin{Resp}\Vsign For He shall not speak of Himself; but whatsoever He shall hear, that shall He speak and He will show you things to come.\end{Resp}
\begin{Resp}\Rsign And when He is come, He will guide you into all truth. Alleluia.\end{Resp}
\begin{Resp}\Vsign Glory be to the Father, and to the Son, \Star{} and to the Holy Ghost.\end{Resp}
\begin{Resp}\Rsign And when He is come, He will guide you into all truth. Alleluia.\end{Resp}
\switchcolumn*

\LA
\LessonHead{Lectio IX}
\switchcolumn
\EN
\LessonHead{Lesson IX}
\switchcolumn*

\LA
\noindent Ergo invisibíliter vidétur. Nec, si non sit in nobis, potest esse in nobis ejus sciéntia: sic enim a nobis vidétur in nobis et nostra consciéntia. Nam fáciem vidémus altérius, nostram vidére non póssumus: consciéntiam vero nostram vidémus, altérius non vidémus. Sed consciéntia nunquam est nisi in nobis: Spíritus autem Sanctus potest esse étiam sine nobis. Datur quippe ut sit et in nobis: sed vidéri et sciri, quemádmodum vidéndus et sciéndus est, non potest a nobis, si non sit in nobis. Post promissiónem Spíritus Sancti, ne quisquam putáret, quod ita eum Dóminus datúrus fúerit velut pro seípso, ut non et ipse cum eis esset futúrus, adjécit atque ait: Non relínquam vos órphanos, véniam ad vos. Quamvis ergo nos Fílius Dei suo Patri adoptáverit fílios, et eúndem Patrem nos volúerit habére per grátiam, qui ejus Pater est per natúram: tamen étiam ipse circa nos patérnum afféctum quodámmodo demónstrat, cum dicit: Non relínquam vos órphanos.\par
\switchcolumn
\EN
\noindent Therefore is He seen That is invisible. If He were not in us we could have in us no knowledge of Him but He is seen in us, as we see our conscience. We see the faces of other men, but we cannot see our Own but of consciences we see none save that within ourselves. But our conscience is never elsewhere but within us whereas the Holy Ghost may be without us, as well as within us. He is given to be within us, and, unless He be within us, we can neither see nor know Him, either within or without us. Then, after that He had promised the Holy Ghost, the Lord, lest they should deem that He was to give them that other Comforter instead of Himself, and that He Himself was to be no longer with them, said also “I will not leave you orphans I will come to you.” Therefore, although the Son of God hath made us by adoption sons of His Own Father, and hath willed that the Same Who is His Father by nature should be our Father by grace, nevertheless, He showeth that Himself hath toward us a love as of a Father, where He saith “I will not leave you orphans.”\par
\switchcolumn*

\LA
\TeDeum{Te Deum}
\switchcolumn
\EN
\TeDeum{Te Deum}
\switchcolumn*

\end{paracol}

\SeasonTitle{Octava Pentecostes}{The Octave of Pentecost}

\addcontentsline{toc}{section}{Octava Pentecostes\enspace\textit{The Octave of Pentecost}}

\Office{Dominica Pentecostes}{Pentecost Sunday}{Duplex I classis cum Octava privilegiata I ordinis}
\addcontentsline{toc}{subsection}{Dominica Pentecostes\enspace\textit{Pentecost Sunday}}
\begin{paracol}{2}
\LA
\LessonHead{Lectio I}
\switchcolumn
\EN
\LessonHead{Lesson I}
\switchcolumn*

\LA
\LessonTitle{Léctio sancti Evangélii secúndum Joánnem}
\switchcolumn
\EN
\LessonTitle{Continuation of the Holy Gospel according to John}
\switchcolumn*

\LA
\Cite{Joannes 14:23-31}
\switchcolumn
\EN
\Cite{John 14:13-31}
\switchcolumn*

\LA
\noindent In illo témpore: Dixit Jesus discípulis suis: Si quis díligit me, sermónem meum servábit, et Pater meus díliget eum, et ad eum veniémus, et mansiónem apud eum faciémus. Et réliqua.\par
\switchcolumn
\EN
\noindent At that time, Jesus said unto His disciples: If a man love Me, He will keep My word, and My Father will love him, and We will come unto him, and make Our abode with him. And so on.\par
\switchcolumn*

\LA
\LessonTitle{Homilía sancti Gregórii Papæ}
\switchcolumn
\EN
\LessonTitle{Homily by Pope St. Gregory the Great}
\switchcolumn*

\LA
\Source{Homilia 30 in Evangelia}
\switchcolumn
\EN
\Source{30th on the Gospels}
\switchcolumn*

\LA
\noindent Libet, fratres caríssimi, evangélicæ verba lectiónis sub brevitáte transcúrrere, ut post diútius líceat in contemplatióne tantæ solemnitátis immorári. Hódie namque Spíritus Sanctus repentíno sónitu super discípulos venit, mentésque carnálium in sui amórem permutávit, et foris apparéntibus linguis ígneis, intus facta sunt corda flammántia; quia dum Deum in ignis visióne suscepérunt, per amórem suáviter arsérunt. Ipse namque Spíritus Sanctus amor est: unde et Joánnes dicit Deus cáritas est. Qui ergo mente íntegra Deum desíderat, profécto jam habet, quem amat. Neque enim quisquam posset Deum dilígere, si eum quem díligit, non habéret.\par
\switchcolumn
\EN
\noindent Dearly beloved brethren, our best way will be to run briefly through the words which have been read from the Holy Gospel, and thereafter rest for a while quietly gazing upon the solemn subject of this great Festival. This is the day whereon “suddenly there came a sound from heaven,” and the Holy Ghost descended upon the Apostles, and, for fleshly minds, gave them minds wherein the love of God was shed abroad and, while without “there appeared unto them cloven tongues, like as of fire, and it sat upon each of them,” within, their hearts were enkindled. While they received the visible presence of God in the form of fire, the flames of His love enwrapped them. The Holy Ghost Himself is love whence it is that John saith “God is love.” Whosoever therefore loveth God with all his soul, already hath obtained Him Whom he loveth, for no man is able to love God, if He have not gained Him Whom he loveth.\par
\switchcolumn*

\LA
\begin{Resp}\Rsign Cum compleréntur dies Pentecóstes, erant omnes páriter in eódem loco, allelúja: et súbito factus est sonus de cælo, allelúja, \Star{} Tamquam spíritus veheméntis, et replévit totam domum, allelúja, allelúja.\end{Resp}
\begin{Resp}\Vsign Dum ergo essent in unum discípuli congregáti propter metum Judæórum, sonus repénte de cælo, venit super eos.\end{Resp}
\begin{Resp}\Rsign Tamquam spíritus veheméntis, et replévit totam domum, allelúja, allelúja.\end{Resp}
\switchcolumn
\EN
\begin{Resp}\Rsign When the day of Pentecost was fully come, they were all with one accord in one place, alleluia, and suddenly there came a sound from heaven, alleluia. \Star{} as of a mighty rushing wind, and it filled all the house, alleluia, alleluia.\end{Resp}
\begin{Resp}\Vsign Where the disciples were assembled for fear of the Jews, there suddenly came upon them a sound from heaven.\end{Resp}
\begin{Resp}\Rsign As of a rushing mighty wind, and it filled all the house, alleluia, alleluia.\end{Resp}
\switchcolumn*

\LA
\LessonHead{Lectio II}
\switchcolumn
\EN
\LessonHead{Lesson II}
\switchcolumn*

\LA
\noindent Sed ecce, si unusquísque vestrum requirátur an díligat Deum: tota fidúcia et secúra mente respóndet, Díligo. In ipso autem lectiónis exórdio audístis quid Véritas dicit: Si quis díligit me, sermónem meum servábit. Probátio ergo dilectiónis, exhibítio est óperis. Hinc in epístola sua idem Joánnes dicit: Qui dicit: Díligo Deum, et mandáta ejus non custódit, mendax est. Vere étenim Deum dilígimus et mandáta ejus custodímus, si nos a nostris voluptátibus coarctámus. Nam qui adhuc per illícita desidéria díffluit, profécto Deum non amat: quia ei in sua voluntáte contradícit.\par
\switchcolumn
\EN
\noindent But, behold, now, if I shall ask any one of you whether he loveth God, he will answer with all boldness and quietness of spirit: “I do love him.” But at the very beginning of this day's Lesson from the Gospel, ye have heard what the Truth saith: “If a man love Me, he will keep My word.” The test, then, of love, is whether it is showed by works. Hence the same John hath said in his Epistle (I. iv. 20, v. 3): “If a man say, I love God, and keepeth not His commandments, He is a liar.” Then do we indeed love God, and keep His commandments, if we deny ourselves the gratification of our appetites. Whosoever still wandereth after unlawful desires, such an one plainly loveth not God, for he saith, Nay, to that which God willeth.\par
\switchcolumn*

\LA
\begin{Resp}\Rsign Repléti sunt omnes Spíritu Sancto: et cœpérunt loqui, prout Spíritus Sanctus dabat éloqui illis: \Star{} Et convénit multitúdo dicéntium, allelúja.\end{Resp}
\begin{Resp}\Vsign Loquebántur váriis linguis Apóstoli magnália Dei.\end{Resp}
\begin{Resp}\Rsign Et convénit multitúdo dicéntium, allelúja.\end{Resp}
\begin{Resp}\Vsign Glória Patri, et Fílio, \Star{} et Spirítui Sancto.\end{Resp}
\begin{Resp}\Rsign Et convénit multitúdo dicéntium, allelúja.\end{Resp}
\switchcolumn
\EN
\begin{Resp}\Rsign They were all filled with the Holy Ghost, and began to speak, as the Holy Ghost gave them utterance; \Star{} And the multitude came together, saying: Alleluia.\end{Resp}
\begin{Resp}\Vsign The Apostles spoke in diverse tongues the wonderful works of God.\end{Resp}
\begin{Resp}\Rsign And the multitude came together, saying: Alleluia.\end{Resp}
\begin{Resp}\Vsign Glory be to the Father, and to the Son, \Star{} and to the Holy Ghost.\end{Resp}
\begin{Resp}\Rsign And the multitude came together, saying: Alleluia.\end{Resp}
\switchcolumn*

\LA
\LessonHead{Lectio III}
\switchcolumn
\EN
\LessonHead{Lesson III}
\switchcolumn*

\LA

\switchcolumn
\EN
\LessonTitle{“And My Father will love him, and We will come unto him, and make Our abode with him.” O my dearly beloved brethren, think what a dignity is that, to have God abiding as a guest in our heart Surely if some rich man or some powerful friend were to come into our house, we would hasten to have our whole house cleaned, lest, perchance, when he came in, he should see aught to displease his eye. So let him that would make his mind an abode for God, cleanse it from all the filth of works of iniquity. Lo, again, what saith the Truth: “We will come unto him, and make Our abode with him.” There are some hearts whereunto God cometh, but maketh not His abode therein with a certain pricking they feel His Presence, but in time of temptation they forget that which hath pricked them and so they turn again to work unrighteousness, even as though they had never repented.}
\switchcolumn*

\LA
\noindent Et Pater meus díliget eum, et ad eum veniémus, et mansiónem apud eum faciémus. Pensáte, fratres caríssimi, quanta sit ista dígnitas, habére in cordis hospítio advéntum Dei. Certe, si domum nostram quisquam dives aut prǽpotens amícus intráret, omni festinántia domus tota mundarétur, ne quid fortásse esset, quod óculos amíci intrántis offénderet. Tergat ergo sordes pravi óperis, qui Deo prǽparat domum mentis. Sed vidéte quid Véritas dicat: Veniémus, et mansiónem apud eum faciémus. In quorúmdam étenim corda venit, et mansiónem non facit: quia, per compunctiónem quidem, Dei respéctum percípiunt, sed tentatiónis témpore hoc ipsum quo compúncti fúerant, obliviscúntur; sicque ad perpetránda peccáta rédeunt, ac si hæc mínime planxíssent.\par
\switchcolumn
\EN

\switchcolumn*

\LA
\TeDeum{Te Deum}
\switchcolumn
\EN
\TeDeum{Te Deum}
\switchcolumn*

\end{paracol}

\Office{Die II infra octavam Pentecostes}{Monday within the Octave of Pentecost}{Duplex I classis}
\addcontentsline{toc}{subsection}{Die II infra octavam Pentecostes\enspace\textit{Monday within the Octave of Pentecost}}
\begin{paracol}{2}
\LA
\LessonHead{Lectio I}
\switchcolumn
\EN
\LessonHead{Lesson I}
\switchcolumn*

\LA
\LessonTitle{Léctio sancti Evangélii secúndum Joánnem}
\switchcolumn
\EN
\LessonTitle{Continuation of the Holy Gospel according to John}
\switchcolumn*

\LA
\Cite{Joannes 3:16-21}
\switchcolumn
\EN
\Cite{John 3:16-21}
\switchcolumn*

\LA
\noindent In illo témpore: Dixit Jesus Nicodémo: Sic Deus diléxit mundum, ut Fílium suum unigénitum daret: ut omnis, qui credit in eum, non péreat, sed hábeat vitam ætérnam. Et réliqua.\par
\switchcolumn
\EN
\noindent At that time, Jesus said unto Nicodemus: God so loved the world that He gave His Only-begotten Son, that whosoever believeth in Him should not perish, but have everlasting life. And so on.\par
\switchcolumn*

\LA
\LessonTitle{Homilía sancti Augustíni Epíscopi}
\switchcolumn
\EN
\LessonTitle{Homily by St. Augustine, Bishop of Hippo}
\switchcolumn*

\LA
\Source{Tract. 12 in Joann., sub fin.}
\switchcolumn
\EN
\Source{12th Tract on John}
\switchcolumn*

\LA
\noindent Quantum in médico est, sanáre venit ægrótum. Ipse se intérimit, qui præcépta médici observáre non vult. Venit Salvátor in mundum. Quare Salvátor dictus est mundi, nisi ut salvet mundum, non ut júdicet mundum? Salvári non vis ab ipso: ex te judicáberis. Et quid dicam, Judicáberis? Vide quid ait: Qui credit in eum, non judicátur. Qui autem non credit: quid dictúrum sperábas, nisi: Judicátur? Quod addit: Jam, inquit, judicátus est: nondum appáruit judícium, et jam factum est judícium.\par
\switchcolumn
\EN
\noindent The Physician cometh that, as often as in him lieth, he may heal the sick man. He is his own destroyer who will not keep the commandments of the Physician. Into the world came the Saviour. Why is He called the Saviour of the world, but because He came “into the world not to condemn the world, but that the world through Him might be saved”? If thou willest not be saved through Him, thou wilt be condemned of thyself. And why say I that thou wilt be condemned? Because it is written: “He that believeth in Him is not condemned.” What then canst thou hope that He will say of “him that believeth not,” but that He will be condemned And indeed He doth say farther: “He that believeth not is condemned already.” He is condemned already, though the condemnation be not yet openly pronounced.\par
\switchcolumn*

\LA
\begin{Resp}\Rsign Jam non dicam vos servos, sed amícos meos; quia ómnia cognovístis, quæ operátus sum in médio vestri, allelúja: \Star{} Accípite Spíritum Sanctum in vobis Paráclitum: ille est, quem Pater mittet vobis, allelúja.\end{Resp}
\begin{Resp}\Vsign Vos amíci mei estis, si fecéritis quæ ego præcípio vobis.\end{Resp}
\begin{Resp}\Rsign Accípite Spíritum Sanctum in vobis Paráclitum: ille est, quem Pater mittet vobis, allelúja.\end{Resp}
\switchcolumn
\EN
\begin{Resp}\Rsign Henceforth I call you not servants, but I have called you My friends, because ye have known all things, whatsoever I have done among you. Alleluia. \Star{} Receive ye the Holy Ghost, Who is your Comforter within you; the Same is He Whom the Father will send unto you. Alleluia.\end{Resp}
\begin{Resp}\Vsign Ye are My friends, if ye do whatsoever I command you.\end{Resp}
\begin{Resp}\Rsign Receive ye the Holy Ghost Who is your Comforter within you; the Same is He Whom the Father will send unto you. Alleluia.\end{Resp}
\switchcolumn*

\LA
\LessonHead{Lectio II}
\switchcolumn
\EN
\LessonHead{Lesson II}
\switchcolumn*

\LA
\noindent Novit enim Dóminus qui sunt ejus: novit qui permáneant ad corónam, qui permáneant ad flammam. Novit in área sua tríticum, novit et páleam: novit ségetem, novit et zizánia. Jam judicátus est, qui non credit. Quare judicátus? Quia non crédidit in nómine unigéniti Fílii Dei. Hoc est autem judícium: quia lux venit in mundum, et dilexérunt hómines magis ténebras, quam lucem: erant enim mala ópera eórum. Fratres mei, quorum ópera bona invénit Dóminus? Nullórum. Omnia ópera mala invénit. Quómodo ergo quidam fecérunt veritátem, et venérunt ad lucem? Et hoc enim séquitur: Qui autem facit veritátem, venit ad lucem.\par
\switchcolumn
\EN
\noindent He is condemned already, for “the Lord knoweth them that are His.” He knoweth them for whom is laid up the crown, and likewise them that are reserved unto the fire. His eye seeth in the field of the world the distinction of the wheat and of the straw, of the good corn and of the tares. “He that believeth not is condemned already.” And why? “Because he hath not believed in the Name of the Only-begotten Son of God. And this is the condemnation that light is come into the world, and men loved darkness rather than light, because their deeds were evil.” “Because their deeds were evil”, but, my brethren, is there one man of whom God findeth that his works are good? No, not one. God findeth all works to be (in themselves) bad. How then do we hear that some there be who do truth, and come to the light? For these words come anon: “But he that doeth truth, cometh to the light.”\par
\switchcolumn*

\LA
\begin{Resp}\Rsign Spíritus Sanctus, procédens a throno, Apostolórum péctora invisibíliter penetrávit novo sanctificatiónis signo: \Star{} Ut in ore eórum ómnium génera nasceréntur linguárum, allelúja.\end{Resp}
\begin{Resp}\Vsign Advénit ignis divínus, non combúrens, sed illúminans, et tríbuit eis charísmatum dona.\end{Resp}
\begin{Resp}\Rsign Ut in ore eórum ómnium génera nasceréntur linguárum, allelúja.\end{Resp}
\begin{Resp}\Vsign Glória Patri, et Fílio, \Star{} et Spirítui Sancto.\end{Resp}
\begin{Resp}\Rsign Ut in ore eórum ómnium génera nasceréntur linguárum, allelúja.\end{Resp}
\switchcolumn
\EN
\begin{Resp}\Rsign The Holy Ghost, Which proceedeth from the Throne, entered unseen into the hearts of the Apostles, with a new token of sanctification, even \Star{} That all manner of tongues should spring to their lips. Alleluia.\end{Resp}
\begin{Resp}\Vsign The fire of God fell, not to burn them, but to enlighten them, and gave them gifts of grace.\end{Resp}
\begin{Resp}\Rsign That all manner of tongues should spring to their lips. Alleluia.\end{Resp}
\begin{Resp}\Vsign Glory be to the Father, and to the Son, \Star{} and to the Holy Ghost.\end{Resp}
\begin{Resp}\Rsign That all manner of tongues should spring to their lips. Alleluia.\end{Resp}
\switchcolumn*

\LA
\LessonHead{Lectio III}
\switchcolumn
\EN
\LessonHead{Lesson III}
\switchcolumn*

\LA
\noindent Sed dilexérunt, inquit, ténebras magis quam lucem. Ibi pósuit vim. Multi enim dilexérunt peccáta sua, multi conféssi sunt peccáta sua: quia qui confitétur peccáta sua, et accúsat peccáta sua, jam cum Deo facit. Accúsat Deus peccáta tua: si et tu accúsas, conjúngeris Deo. Quasi duæ res sunt, homo et peccátor. Quod audis homo, Deus fecit: quod audis peccátor, ipse homo fecit. Dele, quod fecísti, ut Deus salvet, quod fecit. Opórtet ut óderis in te opus tuum, et ames in te opus Dei. Cum autem cœ́perit tibi displicére quod fecísti, inde incípiunt bona ópera tua, quia accúsas mala ópera tua. Inítium óperum bonórum, conféssio est óperum malórum.\par
\switchcolumn
\EN
\noindent But the Lord saith [of such as these, who are condemned already, because they believe not in Him]: “They loved darkness rather than light.” And here He maketh the great point [of difference between such, and them that do the truth.] There are many who have loved their sins there are many who have confessed their sins and he that confessed and denounceth his sin, is working already with God. God denounceth thy sins, and if thou denounce them likewise, then dost thou join thyself with God in His act. The man and the sinner are two different things. God made the man, and the man made the sinner. Put away thy work, and God will save His. Thou art behoven to hate in thyself thine own work, and to love God's work. When thine own works begin to displease thee, then is it that thou beginnest to do well, because thou denouncest thine own evil works. The first thing to do, if thou wouldest do good works, is to acknowledge thine evil ones.\par
\switchcolumn*

\LA
\TeDeum{Te Deum}
\switchcolumn
\EN
\TeDeum{Te Deum}
\switchcolumn*

\end{paracol}

\Office{Die III infra octavam Pentecostes}{Tuesday within the Octave of Pentecost}{Duplex I classis}
\addcontentsline{toc}{subsection}{Die III infra octavam Pentecostes\enspace\textit{Tuesday within the Octave of Pentecost}}
\begin{paracol}{2}
\LA
\LessonHead{Lectio I}
\switchcolumn
\EN
\LessonHead{Lesson I}
\switchcolumn*

\LA
\LessonTitle{Léctio sancti Evangélii secúndum Joánnem}
\switchcolumn
\EN
\LessonTitle{Continuation of the Holy Gospel according to John}
\switchcolumn*

\LA
\Cite{Joannes 10:1-10}
\switchcolumn
\EN
\Cite{John 10:1-10}
\switchcolumn*

\LA
\noindent In illo témpore: Dixit Jesus pharisǽis: Amen, amen dico vobis: qui non intrat per óstium in ovíle óvium, sed ascéndit aliúnde, ille fur est, et latro. Qui autem intrat per óstium, pastor est óvium. Et réliqua.\par
\switchcolumn
\EN
\noindent At that time: Jesus said unto the Pharisees: Amen, amen, I say unto you, he that entereth not by the door into the sheep-fold, but climbeth up some other way, the same is a thief and a robber but he that entereth by the door is the shepherd of the sheep. And so on.\par
\switchcolumn*

\LA
\LessonTitle{Homilía sancti Augustíni Epíscopi}
\switchcolumn
\EN
\LessonTitle{Homily by St. Augustine, Bishop of Hippo}
\switchcolumn*

\LA
\Source{Tract. 45 in Joann. post initium}
\switchcolumn
\EN
\Source{45th Tract on John}
\switchcolumn*

\LA
\noindent Dóminus de grege suo, et de óstio quo intrátur ad ovíle, similitúdinem propósuit in hodiérna lectióne. Dicant ergo pagáni: Bene vívimus! Si per óstium non intrant, quid prodest eis unde gloriántur? Ad hoc enim debet unicuíque prodésse bene vívere, ut detur illi semper vívere: nam cui non datur semper vívere, quid prodest bene vívere? Quia nec bene vívere dicéndi sunt, qui finem bene vivéndi vel cæcitáte nésciunt, vel inflatióne contémnunt. Non est autem cuíquam spes vera et certa semper vivéndi, nisi agnóscat vitam, quod est Christus, et per jánuam intret in ovíle.\par
\switchcolumn
\EN
\noindent In the words of the Gospel which are this day read, the Lord has spoken unto us in similitudes, touching His flock, and the Door whereby entry is made into their fold. The Pagans therefore may say, “We have good lives,” but if they enter not in the Door, what doth that profit them whereof they make their boast? Good life is profitable to a man if it lead unto life everlasting, but if he does not have life everlasting, what shall his good life profit him? Neither indeed can it be truly said that they live good lives, who are either so blinded as not to know, or so puffed up as to despise, the end of a good life. And no man can have a true and certain hope of life everlasting, unless he know the true Life, Which is Christ, and enter in by that Door into the sheepfold.\par
\switchcolumn*

\LA
\begin{Resp}\Rsign Apparuérunt Apóstolis dispertítæ linguæ tamquam ignis, allelúja: \Star{} Sedítque supra síngulos eórum Spíritus Sanctus, allelúja, allelúja.\end{Resp}
\begin{Resp}\Vsign Et cœpérunt loqui váriis linguis, prout Spíritus Sanctus dabat éloqui illis.\end{Resp}
\begin{Resp}\Rsign Sedítque supra síngulos eórum Spíritus Sanctus, allelúja, allelúja.\end{Resp}
\switchcolumn
\EN
\begin{Resp}\Rsign There appeared unto the Apostles cloven tongues, like as of fire. Alleluia. \Star{} And the Holy Ghost sat upon each of them. Alleluia, Alleluia.\end{Resp}
\begin{Resp}\Vsign And they began to speak with other tongues, as the Holy Ghost gave them utterance.\end{Resp}
\begin{Resp}\Rsign And the Holy Ghost sat upon each of them. Alleluia, Alleluia.\end{Resp}
\switchcolumn*

\LA
\LessonHead{Lectio II}
\switchcolumn
\EN
\LessonHead{Lesson II}
\switchcolumn*

\LA
\noindent Quærunt ergo plerúmque tales hómines étiam persuadére homínibus, ut bene vivant, et Christiáni non sint. Per áliam partem volunt ascéndere, rápere et occídere; non, ut bonus pastor, conserváre atque salváre. Fuérunt ergo quidam philósophi de virtútibus et vítiis subtília multa tractántes, dividéntes, definiéntes, ratiocinatiónes acutíssimas concludéntes, libros impléntes, suam sapiéntiam buccis crepántibus ventilántes, qui étiam dícere audérent homínibus: Nos sequímini, sectam nostram tenéte, si vultis beáte vívere. Sed non intrábant per óstium: pérdere volébant, mactáre et occídere.\par
\switchcolumn
\EN
\noindent There are many such, who try to persuade men to live good lives but not to be Christians. These are they who would fain “climb up some other way,” “for to kill and to destroy,” and are not as the Good Shepherd, Who is come to keep and to save. There have been philosophers who have treated many subtle questions of right and wrong, who have been the authors of many distinctions and definitions, who have completed many exceedingly clever arguments, who have filled many books, and have proclaimed their own wisdom with braying trumpets. These dared to say to men: “Follow us embrace our school of thought, and you will find therein the secret of an happy life.” But these were not of them who enter in by the Door; they came not but for to steal, and to kill, and to destroy.\par
\switchcolumn*

\LA
\begin{Resp}\Rsign Loquebántur váriis linguis Apóstoli magnália Dei, \Star{} Prout Spíritus Sanctus dabat éloqui illis, allelúja.\end{Resp}
\begin{Resp}\Vsign Repléti sunt omnes Spíritu Sancto, et cœpérunt loqui.\end{Resp}
\begin{Resp}\Rsign Prout Spíritus Sanctus dabat éloqui illis, allelúja.\end{Resp}
\begin{Resp}\Vsign Glória Patri, et Fílio, \Star{} et Spirítui Sancto.\end{Resp}
\begin{Resp}\Rsign Prout Spíritus Sanctus dabat éloqui illis, allelúja.\end{Resp}
\switchcolumn
\EN
\begin{Resp}\Rsign The Apostles spake in diverse tongues the wonderful works of God, \Star{} As the Holy Ghost gave them utterance. Alleluia.\end{Resp}
\begin{Resp}\Vsign They were all filled with the Holy Ghost, and began to speak\end{Resp}
\begin{Resp}\Rsign As the Holy Ghost gave them utterance. Alleluia.\end{Resp}
\begin{Resp}\Vsign Glory be to the Father, and to the Son, \Star{} and to the Holy Ghost.\end{Resp}
\begin{Resp}\Rsign As the Holy Ghost gave them utterance. Alleluia.\end{Resp}
\switchcolumn*

\LA
\LessonHead{Lectio III}
\switchcolumn
\EN
\LessonHead{Lesson III}
\switchcolumn*

\LA
\noindent Quid de istis dicam? Ecce ipsi pharisǽi legébant, et in eo quod legébant, Christum sonábant, ventúrum sperábant, et præséntem non agnoscébant. Jactábant se étiam ipsi inter Vidéntes, hoc est, inter sapiéntes, et negábant Christum, et non intrábant per óstium. Ergo et ipsi, si quos forte sedúcerent, mactándos et occidéndos, non liberándos sedúcerent. Et hos dimittámus. Videámus illos, si forte ipsi intrant per óstium, qui ipsíus Christi nómine gloriántur. Innumerábiles enim sunt, qui se Vidéntes non solum jactant, sed a Christo illuminátos vidéri volunt: sunt autem hærétici.\par
\switchcolumn
\EN
\noindent Touching these, what shall I say? Behold, the Pharisees themselves read of Christ, and therefore talked of Christ they looked for His coming, and when He came, they knew Him not. They boasted that they themselves were among the Seers, that is, of the wise ones, and they denied Christ, and entered not in by the Door. Therefore they, if they led away any, led them away only to kill and to destroy, not to free them. So much for them. Now let us see if all they who boast the name of Christian enter in by the Door. Some there are, and their number cannot be reckoned, who not only boast that they themselves are among the Seers, but would fain appear as though their hearts were enlightened by Christ but they are heretics.\par
\switchcolumn*

\LA
\TeDeum{Te Deum}
\switchcolumn
\EN
\TeDeum{Te Deum}
\switchcolumn*

\end{paracol}

\Office{Feria Quarta Quattuor Temporum Pentecostes}{Ember Wednesday of Pentecost}{Semiduplex I classis}
\addcontentsline{toc}{subsection}{Feria Quarta Quattuor Temporum Pentecostes\enspace\textit{Ember Wednesday of Pentecost}}
\begin{paracol}{2}
\LA
\LessonHead{Lectio I}
\switchcolumn
\EN
\LessonHead{Lesson I}
\switchcolumn*

\LA
\LessonTitle{Léctio sancti Evangélii secúndum Joánnem}
\switchcolumn
\EN
\LessonTitle{Continuation of the Holy Gospel according to John}
\switchcolumn*

\LA
\Cite{Joannes 6:44-52}
\switchcolumn
\EN
\Cite{John 6:44-52}
\switchcolumn*

\LA
\noindent In illo témpore: Dixit Jesus turbis Judæórum: Nemo potest veníre ad me, nisi Pater, qui misit me, tráxerit eum. Et réliqua.\par
\switchcolumn
\EN
\noindent At that time, Jesus said unto the multitudes of the Jews: No man can come to Me, except the Father, Which hath sent me, draw him. And so on.\par
\switchcolumn*

\LA
\LessonTitle{Homilía sancti Augustíni Epíscopi}
\switchcolumn
\EN
\LessonTitle{Homily by St. Augustine, Bishop of Hippo}
\switchcolumn*

\LA
\Source{Tract. 26 in Joann., post init.}
\switchcolumn
\EN
\Source{26th Tract on John}
\switchcolumn*

\LA
\noindent Noli cogitáre te invítum trahi: tráhitur ánimus et amóre. Nec timére debémus, ne ab homínibus, qui verba perpéndunt, et a rebus máxime divínis intellegéndis longe remóti sunt, in hoc Scripturárum sanctárum evangélico verbo fórsitan reprehendámur, et dicátur nobis: Quómodo voluntáte credo, si trahor? Ego dico: Parum est voluntáte: étiam voluptáte tráheris. Quid est, trahi voluptáte? Delectáre in Dómino: et dabit tibi petitiónes cordis tui. Est quædam volúptas cordis, cui panis dulcis est ille cæléstis. Porro si poétæ dícere lícuit: Trahit sua quemque volúptas: non necéssitas, sed volúptas; non obligátio, sed delectátio: quanto fórtius nos dícere debémus, trahi hóminem ad Christum, qui delectátur veritáte, delectátur beatitúdine, delectátur justítia, delectátur sempitérna vita, quod totum Christus est? An vero habent córporis sensus voluptátes suas, et ánimus deséritur a voluptátibus suis? Si ánimus non habet voluptátes suas, unde dícitur: Fílii autem hóminum sub tégmine alárum tuárum sperábunt: inebriabúntur ab ubertáte domus tuæ, et torrénte voluptátis tuæ potábis eos. Quóniam apud te est fons vitæ: et in lúmine tuo vidébimus lumen.\par
\switchcolumn
\EN
\noindent Think not that thou art drawn against thy will; the soul is drawn, not willingly only, but lovingly. Neither must we be afraid lest men who are great weighers of words, and very far from understanding the things of God, should catch us up upon this Gospel doctrine of the Holy Scriptures, and should say to us: How can my faith be willing if am drawn? I answer: Thou art not drawn as touching thy will, but by pleasure. And, now, what is being drawn by pleasure? Delight thyself in the Lord, and He shall give thee the desires of thine heart. (Ps. xxxvi. 4.) There is pleasure in that heart to which the Bread That came down from heaven is sweet. The poet is allowed to say His special pleasure draweth each, but pleasure, which so draweth, is not a necessity, not a bond, but a delight how much more strongly, may we say that men are drawn to Christ, who delight in truth, who delight in blessedness, who delight in righteousness, who delight in life everlasting, since truth and blessedness, and righteousness and everlasting life are all to be found in Christ? Or have the bodily senses pleasure, and the spiritual senses none? If the spiritual sense have no pleasures, wherefore is it written: And the children of men shall put their trust under the shadow of thy wings. They shall be abundantly satisfied with the fatness of thy house, and Thou shalt make them drink of the river of thy pleasures. For with thee is the fountain of life, and in thy light shall we see light. (Ps. xxxv. 8.)\par
\switchcolumn*

\LA
\begin{Resp}\Rsign Disciplínam et sapiéntiam dócuit eos Dóminus, allelúja: firmávit in illis grátiam Spíritus sui, \Star{} Et intelléctu implévit corda eórum, allelúja.\end{Resp}
\begin{Resp}\Vsign Repentíno namque sónitu Spíritus Sanctus super eos venit.\end{Resp}
\begin{Resp}\Rsign Et intelléctu implévit corda eórum, allelúja.\end{Resp}
\switchcolumn
\EN
\begin{Resp}\Rsign The Lord taught them good judgment and knowledge, Alleluia. He established in them the grace of His Spirit, \Star{} And filled their hearts with understanding. Alleluia.\end{Resp}
\begin{Resp}\Vsign For with a sudden sound the Holy Ghost came upon them.\end{Resp}
\begin{Resp}\Rsign And filled their hearts with understanding. Alleluia.\end{Resp}
\switchcolumn*

\LA
\LessonHead{Lectio II}
\switchcolumn
\EN
\LessonHead{Lesson II}
\switchcolumn*

\LA
\noindent Da amántem, et sentit quod dico: da desiderántem, da esuriéntem, da in ista solitúdine peregrinántem, atque sitiéntem, et fontem ætérnæ pátriæ suspirántem: da talem, et scit quid dicam. Si autem frígido loquor, nescit, quid loquor. Tales erant isti, qui ínvicem murmurábant. Pater, inquit, quem tráxerit, venit ad me. Quid est autem, Pater quem tráxerit, cum ipse Christus trahat? Quare vóluit dícere, Pater quem tráxerit? Si trahéndi sumus, ab illo trahámur, cui dicit quædam, quæ díligit: Post odórem unguentórum tuórum currémus. Sed quid intéllegi vóluit, advertámus fratres, et, quantum póssumus, capiámus. Trahit Pater ad Fílium eos, qui proptérea credunt in Fílium, quia eum cógitant Patrem habére Deum. Deus enim Pater æquálem sibi génuit Fílium; et qui cógitat, atque in fide sua sentit et rúminat, æquálem esse Patri eum, in quem credit, ipsum trahit Pater ad Fílium.\par
\switchcolumn
\EN
\noindent Give me a lover, and he will catch my meaning; give me one who longs, give me an hungerer, give me a wanderer in this desert, a thirst and gasping for the fountains of the eternal Fatherland; give me such an one, and he will catch my meaning. If I talk to some cold creature, he will not. Such cold creatures were they of whom it is written: The Jews then murmured at Him because He said, I am the Bread Which came down from heaven. And they said: Is not this Jesus the son of Joseph, whose father and Mother we know? How is it then that He saith: I came down from heaven? Jesus therefore answered and said unto them: Murmur not among yourselves. No man can come to Me, except the Father, Which hath sent Me, draw him. (41-44.) But wherefore speaketh Christ of them whom the Father draweth, since He Himself draweth. Why was it His will to say: No man can come to Me except the Father draw him? If we are to be drawn, let us be drawn by Him to Whom one that loved much said: Draw me, we will run after the savour of thy good ointments. (Cant. i. 4.) But let us consider, my brethren, what He meant, and understand it as well as we can. The Father draweth to the Son them who believe in the Son, because they are persuaded that He hath God to His Father. God the Father begetteth to Himself a coequal Son; and whosoever is persuaded, and realiseth unto himself by faith, and thinketh, that He in Whom he believeth is equal to the Father, him the Father is drawing unto the Son.\par
\switchcolumn*

\LA
\begin{Resp}\Rsign Ite in univérsum orbem, et prædicáte Evangélium, allelúja: \Star{} Qui credíderit et baptizátus fúerit, salvus erit, allelúja, allelúja, allelúja.\end{Resp}
\begin{Resp}\Vsign In nómine meo dæmónia ejícient, linguis loquéntur novis, serpéntes tollent.\end{Resp}
\begin{Resp}\Rsign Qui credíderit et baptizátus fúerit, salvus erit, allelúja, allelúja, allelúja.\end{Resp}
\begin{Resp}\Vsign Glória Patri, et Fílio, \Star{} et Spirítui Sancto.\end{Resp}
\begin{Resp}\Rsign Qui credíderit et baptizátus fúerit, salvus erit, allelúja, allelúja, allelúja.\end{Resp}
\switchcolumn
\EN
\begin{Resp}\Rsign Go ye unto all the world and preach the Gospel. Alleluia. \Star{} He that believeth, and is baptized, shall be saved. Alleluia, Alleluia, Alleluia.\end{Resp}
\begin{Resp}\Vsign In My Name shall they cast out devils; they shall speak with new tongues; they shall take up serpents.\end{Resp}
\begin{Resp}\Rsign He that believeth, and is baptized, shall be saved. Alleluia, Alleluia, Alleluia.\end{Resp}
\begin{Resp}\Vsign Glory be to the Father, and to the Son, \Star{} and to the Holy Ghost.\end{Resp}
\begin{Resp}\Rsign He that believeth, and is baptized, shall be saved. Alleluia, Alleluia, Alleluia.\end{Resp}
\switchcolumn*

\LA
\LessonHead{Lectio III}
\switchcolumn
\EN
\LessonHead{Lesson III}
\switchcolumn*

\LA
\noindent Aríus crédidit creatúram, non eum traxit Pater; quia non consíderat Patrem, qui Fílium non credit æquálem. Quid dicis, o Ari? quid dicis, hærétice? quid loquéris? Quid est Christus? Non, inquit, Deus verus, sed quem fecit Deus verus. Non te traxit Pater; non enim intellexísti Patrem, cujus Fílium negas. Aliud cógitas, non est ipse Fílius: nec a Patre tráheris, nec ad Fílium tráheris. Aliud est enim Fílius, áliud quod tu dicis. Photínus dicit: Homo solum est Christus, non est et Deus. Qui sic credit, non Pater eum traxit. Quem traxit Pater? Illum qui dicit: Tu es Christus Fílius Dei vivi. Ramum víridem osténdis ovi, et trahis illam. Nuces púero demonstrántur, et tráhitur: et quod currit, tráhitur, amándo tráhitur, sine læsióne córporis tráhitur, cordis vínculo tráhitur. Si ergo ista, quæ inter delícias et voluptátes terrénas revelántur amántibus, trahunt, quóniam verum est, Trahit sua quemque volúptas; non trahit revelátus Christus a Patre? Quid enim fórtius desíderat ánima, quam veritátem?\par
\switchcolumn
\EN
\noindent Arius, who believed that the Son was made, was not one of them whom the Father draweth since whosoever believeth not that the Father is a Father by the begetting of a coequal Son, such an one knoweth not the Father. What sayest thou, O Arius? What sayest thou, O thou heretic? What is thy profession? What is Christ? He is not, saith Arius, Himself Very God. Then, O Arius, the Father hath not drawn thee; thou hast not understood His dignity as a Father, to Whom thou deniest His Son. Thou dost deny the existence of the Son of God, the Father draweth thee not, and thou art not drawn to the Son, since the Son of whom thou speakest is another son, [existing only in thine imagination,] and not the really existent Son. Photinus said: Christ is a mere man, and not God at all. He who uttered those words was not one of them whom the Father draweth. But whom hath the Father drawn? The Father drew him who said: Thou art the Christ, the Son of the living God. (Matth. xvi. 16, 17.) Show a sheep a green bough, and thou drawest him. Let a boy see some nuts, and he is drawn by them. As they run, they are drawn, drawn by taste, drawn without bodily hurt, drawn by a line bound to their heart. If, then, among earthly things, such as be sweet and pleasant draw such as love them, as soon as they see them, so that it is truth to say, His special pleasure draweth each, doth not that Christ, Whom the Father hath revealed, draw? What stronger object of love can a soul have than the Truth?\par
\switchcolumn*

\LA
\TeDeum{Te Deum}
\switchcolumn
\EN
\TeDeum{Te Deum}
\switchcolumn*

\end{paracol}

\Office{Feria Quinta infra octavam Pentecostes}{Thursday within the Octave of Pentecost}{Semiduplex I classis}
\addcontentsline{toc}{subsection}{Feria Quinta infra octavam Pentecostes\enspace\textit{Thursday within the Octave of Pentecost}}
\begin{paracol}{2}
\LA
\LessonHead{Lectio I}
\switchcolumn
\EN
\LessonHead{Lesson I}
\switchcolumn*

\LA
\LessonTitle{Léctio sancti Evangélii secúndum Lucam}
\switchcolumn
\EN
\LessonTitle{Continuation of the Holy Gospel according to Luke}
\switchcolumn*

\LA
\Cite{Luc 9:1-6}
\switchcolumn
\EN
\Cite{Luke 9:1-6}
\switchcolumn*

\LA
\noindent In illo témpore: Convocátis Jesus duódecim Apóstolis, dedit illis virtútem et potestátem super ómnia dæmónia, et ut languóres curárent. Et réliqua.\par
\switchcolumn
\EN
\noindent At that time: Jesus called His twelve disciples together, and gave them power and authority over all devils, and to cure diseases. And so on.\par
\switchcolumn*

\LA
\LessonTitle{Homilía sancti Ambrósii Epíscopi}
\switchcolumn
\EN
\LessonTitle{Homily by St. Ambrose, Bishop of Milan}
\switchcolumn*

\LA
\Source{Liber 6 in cap. 9 Lucæ}
\switchcolumn
\EN
\Source{Bk. 6 on Luke ix}
\switchcolumn*

\LA
\noindent Qualis débeat esse, qui evangelízat regnum Dei, præcéptis evangélicis designátur: ut sine virga, sine pera, sine calceaménto, sine pane, sine pecúnia, hoc est, subsídii sæculáris adminícula non requírens, fidéque tutus, putet sibi quo minus ea requírat, magis posse suppétere. Quæ possunt, qui volunt, ad eum deriváre tractátum, ut spiritálem tantúmmodo locus iste formáre videátur afféctum: qui velut induméntum quoddam videátur córporis exuísse, non solum potestáte rejécta contemptísque divítiis, sed étiam carnis ipsíus illécebris abdicátis. Quibus primo ómnium datur pacis atque constántiæ generále mandátum, ut pacem ferant, constántiam servent, hospitális necessitúdinis jura custódiant: aliénum a prædicatóre regni cæléstis ástruens cursitáre per domos, et inviolábilis hospítii jura mutáre.\par
\switchcolumn
\EN
\noindent We learn from the commandments of the Gospel what manner of men they ought to be who preach the glad tidings of the kingdom of God “Take nothing for your journey neither staves nor scrip, neither bread neither money.” Thus let the Apostle destitute of earthly help, and panoplies in faith, deem himself able to do all the more, as he needeth all the less Such as please may also put upon these words a spiritual interpretation in that a man may be said to lay as the encumbrances of the body, not only by abdicating power, and casting away riches, but also by denying the very body itself its pleasures. The first general commandment given to the Apostles touching their manners was to be bringers of peace, (Matth. x. 13,) and to be no gadders about, but keepers of the laws of guests. To wander from house to house, and to abuse the rights of hospitality, are things alien to a preacher of the kingdom of heaven.\par
\switchcolumn*

\LA
\begin{Resp}\Rsign Advénit ignis divínus, non combúrens sed illúminans, non consúmens sed lucens: et invénit corda discipulórum receptácula munda: \Star{} Et tríbuit eis charísmatum dona, allelúja, allelúja.\end{Resp}
\begin{Resp}\Vsign Invénit eos concórdes caritáte, et collustrávit eos inúndans grátia Deitátis.\end{Resp}
\begin{Resp}\Rsign Et tríbuit eis charísmatum dona, allelúja, allelúja.\end{Resp}
\switchcolumn
\EN
\begin{Resp}\Rsign The fire of God fell, not to burn them, but to enlighten them not to devour them, but to illuminate them and found the hearts of the disciples clean vessels. \Star{} And He gave them gifts of His grace. Alleluia, Alleluia.\end{Resp}
\begin{Resp}\Vsign He found them one in love, and the out-poured grace of the Godhead shone through them.\end{Resp}
\begin{Resp}\Rsign And He gave them gifts of His grace. Alleluia, Alleluia.\end{Resp}
\switchcolumn*

\LA
\LessonHead{Lectio II}
\switchcolumn
\EN
\LessonHead{Lesson II}
\switchcolumn*

\LA
\noindent Sed, ut hospítii grátia deferénda censétur: ita étiam, si non recipiántur, excutiéndum púlverem, et egrediéndum de civitáte mandátur. Quo non medíocris boni remunerátio docétur hospítii: ut non solum pacem tribuámus hospítibus, verum étiam, si qua eos terrénæ obúmbrant delícta levitátis, recéptis apostólicæ prædicatiónis vestígiis auferántur. Nec otióse secúndum Matthǽum, domus, quam ingrediántur Apóstoli, eligénda decérnitur: ut mutándi hospítii, necessitudinísque violándæ causa non súppetat. Non tamen éadem cáutio receptóri mandátur hospítii: ne, dum hospes elígitur, hospitálitas ipsa minuátur.\par
\switchcolumn
\EN
\noindent But as the kindness of hospitality is to be met with courtesy, so also is it said “Whosoever will not receive you, when ye go out of that city, shake off the very dust from your feet, for a testimony against them.” Hereby is it taught that hospitality doth meet with a good reward, since not only do we bring peace to such as receive us, but also, if they be shadowed by some earthly vanities, these defects are taken away, where enter the feet of them that bear the glad tidings of Apostolic preachment. It is well written in Matthew (x. 11): “Into whatsoever city or town ye shall enter, inquire who in it is worthy and there abide till ye go thence” thus avoiding any possible need of going from house to house. But no such selection is commanded to him that giveth hospitality, lest his hospitality itself should be lessened, while he picketh his guests.\par
\switchcolumn*

\LA
\begin{Resp}\Rsign Spíritus Sanctus replévit totam domum, ubi erant Apóstoli: et apparuérunt illis dispertítæ linguæ, tamquam ignis, sedítque supra síngulos eórum: \Star{} Et repléti sunt omnes Spíritu Sancto, et cœpérunt loqui váriis linguis, prout Spíritus Sanctus dabat éloqui illis, allelúja, allelúja, allelúja.\end{Resp}
\begin{Resp}\Vsign Dum ergo essent in unum discípuli congregáti propter metum Judæórum, sonus repénte de cælo venit super eos.\end{Resp}
\begin{Resp}\Rsign Et repléti sunt omnes Spíritu Sancto, et cœpérunt loqui váriis linguis, prout Spíritus Sanctus dabat éloqui illis, allelúja, allelúja, allelúja.\end{Resp}
\begin{Resp}\Vsign Glória Patri, et Fílio, \Star{} et Spirítui Sancto.\end{Resp}
\begin{Resp}\Rsign Et repléti sunt omnes Spíritu Sancto, et cœpérunt loqui váriis linguis, prout Spíritus Sanctus dabat éloqui illis, allelúja, allelúja, allelúja.\end{Resp}
\switchcolumn
\EN
\begin{Resp}\Rsign The Holy Ghost filled all the house where the Apostles were and there appeared unto them cloven tongues like as of fire, and it sat upon each of them. \Star{} And they were all filled with the Holy Ghost, and began to speak in diverse tongues as the Holy Spirit gave them utterance. Alleluia, Alleluia, Alleluia.\end{Resp}
\begin{Resp}\Vsign When the disciples were all with one accord in one place, for fear of the Jews, suddenly there came a sound from heaven upon them.\end{Resp}
\begin{Resp}\Rsign And they were all filled with the Holy Ghost, and began to speak in diverse tongues as the Holy Spirit gave them utterance. Alleluia, Alleluia, Alleluia.\end{Resp}
\begin{Resp}\Vsign Glory be to the Father, and to the Son, \Star{} and to the Holy Ghost.\end{Resp}
\begin{Resp}\Rsign And they were all filled with the Holy Ghost, and began to speak in diverse tongues, as the Holy Spirit gave them utterance. Alleluia, Alleluia, Alleluia.\end{Resp}
\switchcolumn*

\LA
\LessonHead{Lectio III}
\switchcolumn
\EN
\LessonHead{Lesson III}
\switchcolumn*

\LA
\noindent Sed hæc, ut secúndum lítteram de hospítii religióne venerábilis est forma præcépti: ita étiam de mystério senténtia cæléstis arrídet. Etenim cum domus elígitur, dignus hospes inquíritur. Videámus ígitur, ne forte Ecclésia præferénda designétur, et Christus. Quæ enim dígnior domus apostólicæ prædicatiónis ingréssu, quam sancta Ecclésia? Aut quis præferéndus magis ómnibus vidétur esse quam Christus, qui pedes suis laváre consuévit hospítibus: et quoscúmque sua recéperit domo, pollútis non patiátur habitáre vestígiis; sed maculósos licet vitæ prióris, in réliquum tamen dignétur mundáre procéssus? Hic est ígitur solus, quem nemo debet desérere, nemo mutáre. Cui bene dícitur: Dómine, ad quem íbimus? verba vitæ ætérnæ habes, et nos crédimus.\par
\switchcolumn
\EN
\noindent This passage, taken according to the plain meaning, is a sacred commandment touching the religious duty of hospitality, but its heavenly words likewise hint at a mystery. When the house is chosen, it is asked if the master thereof be worthy. Let us see if this be not perchance a figure of the Church, and her Master, Christ. What worthier house can the Apostolic preacher enter, than the Holy Church? Or what host is more to be preferred before all others, than Christ, Whose use it is to wash the feet of His guests yea, Who suffereth not that any whom He receiveth into His house should dwell there with foul feet, but, defiled as they are by their former wanderings, doth vouchsafe to change them into new and clean livers. He Alone is He, from Whose house no man ought ever to go forth, nor change His roof for any other shelter, for unto Him it is well said “Lord, to whom shall we go? Thou hast the words of eternal life, and we believe.” (John vi. 68, 69.)\par
\switchcolumn*

\LA
\TeDeum{Te Deum}
\switchcolumn
\EN
\TeDeum{Te Deum}
\switchcolumn*

\end{paracol}

\Office{Feria Sexta Quattuor Temporum Pentecostes}{Ember Friday of Pentecost}{Semiduplex I classis}
\addcontentsline{toc}{subsection}{Feria Sexta Quattuor Temporum Pentecostes\enspace\textit{Ember Friday of Pentecost}}
\begin{paracol}{2}
\LA
\LessonHead{Lectio I}
\switchcolumn
\EN
\LessonHead{Lesson I}
\switchcolumn*

\LA
\LessonTitle{Léctio sancti Evangélii secúndum Lucam}
\switchcolumn
\EN
\LessonTitle{Continuation of the Holy Gospel according to Luke}
\switchcolumn*

\LA
\Cite{Luc 5:17-26}
\switchcolumn
\EN
\Cite{Luke 5:17-26}
\switchcolumn*

\LA
\noindent In illo témpore: Factum est in una diérum, et Jesus sedébat docens. Et erant pharisǽi sedéntes, et legis doctóres, qui vénerant ex omni castéllo Galilǽæ et Judǽæ, et Jerúsalem: et virtus Dómini erat ad sanándum eos. Et réliqua.\par
\switchcolumn
\EN
\noindent At that time, it came to pass on a certain day, as Jesus sat and taught, that there were Pharisees, and Doctors of the law sitting by, which were come out of every town in Galilee, and Judaea and Jerusalem and the power of the Lord was present to heal them. And so on.\par
\switchcolumn*

\LA
\LessonTitle{Homilía sancti Ambrósii Epíscopi}
\switchcolumn
\EN
\LessonTitle{Homily by St. Ambrose, Bishop of Milan}
\switchcolumn*

\LA
\Source{Liber 5 in cap. 5 Lucæ, post initium}
\switchcolumn
\EN
\Source{Book 5 on Luke v}
\switchcolumn*

\LA
\noindent Non otiósa hujus paralýtici, nec angústa medicína est, quando Dóminus et orásse præmíttitur; non útique propter suffrágium, sed propter exémplum. Imitándi enim spécimen dedit, non precándi ámbitum requisívit. Et conveniéntibus ex omni Galilǽa, et Judǽa, et Jerúsalem, legis doctóribus, inter ceterórum remédia debílium, paralýtici istíus medicína descríbitur. Primum ómnium, quod ante díximus, unusquísque æger peténdæ precatóres salútis debet adhibére, per quos nostræ vitæ compágo resolúta, actuúmque nostrórum clauda vestígia, verbi cæléstis remédio reforméntur.\par
\switchcolumn
\EN
\noindent “And, behold, men brought in a bed a man which was taken with a palsy.” The healing of this paralytic was not idle, nor its fruits limited to himself. The Lord healed him, or ever he could ask, not because of the entreaties of others, but for example's sake. He gave a pattern to be followed, and sought not the intercession of prayer. In the presence of the Pharisees and doctors of the law, which were come out of every town of Galilee, and Judaea, and Jerusalem, many sick folk were healed, but among them is specially described the healing of this paralytic. First of all, as we have before said, every sick man ought to engage his friends to offer up prayers for his recovery, that so the tottering framework of this our life, and the distorted feet of our works, may be righted by the healing power of the word from heaven.\par
\switchcolumn*

\LA
\begin{Resp}\Rsign Non vos me elegístis, sed ego elégi vos, et pósui vos: \Star{} Ut eátis, et fructum afferátis, et fructus vester máneat, allelúja, allelúja.\end{Resp}
\begin{Resp}\Vsign Sicut misit me Pater, et ego mitto vos.\end{Resp}
\begin{Resp}\Rsign Ut eátis, et fructum afferátis, et fructus vester máneat, allelúja, allelúja.\end{Resp}
\switchcolumn
\EN
\begin{Resp}\Rsign Ye have not chosen Me, but I have chosen you, and ordained you \Star{} That ye should go and bring forth fruit, and that your fruit should remain. Alleluia, Alleluia.\end{Resp}
\begin{Resp}\Vsign As My Father hath sent Me, even so send I you.\end{Resp}
\begin{Resp}\Rsign That ye should go and bring forth fruit, and that your fruit should remain. Alleluia, Alleluia.\end{Resp}
\switchcolumn*

\LA
\LessonHead{Lectio II}
\switchcolumn
\EN
\LessonHead{Lesson II}
\switchcolumn*

\LA
\noindent Sint ígitur áliqui monitóres mentis, qui ánimum hóminis, quamvis exterióris córporis debilitáte torpéntem, ad superióra érigant. Quorum rursus adminículis et attóllere et humiliáre se fácilis ante Jesum locétur, Domínico vidéri dignus aspéctu. Humilitátem enim réspicit Dóminus: quia respéxit humilitátem ancíllæ suæ. Quorum fidem ut vidit, dixit: Homo, remittúntur tibi peccáta tua. Magnus Dóminus, qui aliórum mérito ignóscit áliis: et dum álios probat, áliis reláxat erráta. Cur apud te, homo, colléga non váleat, cum apud Deum servus et interveniéndi méritum, et jus hábeat impetrándi?\par
\switchcolumn
\EN
\noindent Here ought therefore to be advisers, who should rouse up the minds of the sick to higher things, since when the body becometh languid with sickness, the mind is apt to follow its example. With the help of such friends he can be brought and laid on the ground before the Feet of Jesus, and seem worthy of a glance from the Lord for the Lord looketh upon such as lie lowly before Him, “for He hath regarded the lowliness of His handmaiden” (Luke ii. 48.) “And when He saw their faith, He said unto him Man, thy sins are forgiven thee.” Great is the Lord, Who, for the sake of some, forgiveth the sins of others Who trieth some, and pardoneth the wanderings of others. Why should thine equal, O man, avail not with thee, if a slave have won power to intercede, and right to obtain, with God?\par
\switchcolumn*

\LA
\begin{Resp}\Rsign Spíritus Dómini replévit orbem terrárum: \Star{} Et hoc quod cóntinet ómnia, sciéntiam habet vocis, allelúja, allelúja.\end{Resp}
\begin{Resp}\Vsign Omnium est enim ártifex, omnem habens virtútem, ómnia prospíciens.\end{Resp}
\begin{Resp}\Rsign Et hoc quod cóntinet ómnia, sciéntiam habet vocis, allelúja, allelúja.\end{Resp}
\begin{Resp}\Vsign Glória Patri, et Fílio, \Star{} et Spirítui Sancto.\end{Resp}
\begin{Resp}\Rsign Et hoc quod cóntinet ómnia, sciéntiam habet vocis, allelúja, allelúja.\end{Resp}
\switchcolumn
\EN
\begin{Resp}\Rsign The Spirit of the Lord filleth the world, \Star{} And That Which containeth all things hath knowledge of the voice. Alleluia, Alleluia.\end{Resp}
\begin{Resp}\Vsign For [Wisdom] is the worker of all things, having all power, overseeing all things.\end{Resp}
\begin{Resp}\Rsign And That Which containeth all things hath knowledge of the voice. Alleluia, Alleluia.\end{Resp}
\begin{Resp}\Vsign Glory be to the Father, and to the Son, \Star{} and to the Holy Ghost.\end{Resp}
\begin{Resp}\Rsign And That Which containeth all things hath knowledge of the voice. Alleluia, Alleluia.\end{Resp}
\switchcolumn*

\LA
\LessonHead{Lectio III}
\switchcolumn
\EN
\LessonHead{Lesson III}
\switchcolumn*

\LA
\noindent Disce, qui júdicas, ignóscere: disce, qui æger es, impetráre. Si grávium peccatórum diffídis véniam, ádhibe precatóres, ádhibe Ecclésiam, quæ pro te precétur, cujus contemplatióne, quod tibi Dóminus negáre posset, ignóscat. Et quamvis históriæ fidem non debeámus omíttere, ut vere paralýtici istíus corpus credámus esse sanátum: cognósce tamen interióris hóminis sanitátem, cui peccáta donántur. Cum Judǽi ásserunt peccáta a solo Deo posse concédi, Deum útique eum confiténtur; suóque judício perfídiam suam produnt, qui opus ástruunt, ut persónam negent.\par
\switchcolumn
\EN
\noindent O Thou that judgest, learn to forgive thou that art sick, to pray. If thou doubt of the pardon of thy sins, because of their grievousness, get thee to the Church, that she may pray for thee, and that the Lord, accepting her countenance, may grant to her petitions what He refuses to thine. And although we are bound to accept this history as one of fact, and to believe that the body of the paralytic was healed yet remember thou also his inward cure, unto whom his sins were forgiven. The Jews said: “Who can forgive sins but God alone?” And in these words they confessed the Godhead of Him Who forgave the sins of the paralytic, and themselves condemned their own unbelief in Him Whose work they acknowledged, but Whose Person they denied.\par
\switchcolumn*

\LA
\TeDeum{Te Deum}
\switchcolumn
\EN
\TeDeum{Te Deum}
\switchcolumn*

\end{paracol}

\Office{Sabbato Quattuor Temporum Pentecostes}{Ember Saturday of Pentecost}{Semiduplex I classis}
\addcontentsline{toc}{subsection}{Sabbato Quattuor Temporum Pentecostes\enspace\textit{Ember Saturday of Pentecost}}
\begin{paracol}{2}
\LA
\LessonHead{Lectio I}
\switchcolumn
\EN
\LessonHead{Lesson I}
\switchcolumn*

\LA
\LessonTitle{Léctio sancti Evangélii secúndum Lucam}
\switchcolumn
\EN
\LessonTitle{Continuation of the Holy Gospel according to Luke}
\switchcolumn*

\LA
\Cite{Luc 4:38-44}
\switchcolumn
\EN
\Cite{Luke 4:38}
\switchcolumn*

\LA
\noindent In illo témpore: Surgens Jesus de synagóga, introívit in domum Simónis. Socrus autem Simónis tenebátur magnis fébribus. Et réliqua.\par
\switchcolumn
\EN
\noindent At that time Jesus arose out of the synagogue, and entered into Simon's house. And Simon's wife's mother was taken with a great fever. And so on.\par
\switchcolumn*

\LA
\LessonTitle{Homilía sancti Ambrósii Epíscopi}
\switchcolumn
\EN
\LessonTitle{Homily by St. Ambrose, Bishop of Milan}
\switchcolumn*

\LA
\Source{Liber 4 in cap. 4 Lucæ, circa finem}
\switchcolumn
\EN
\Source{Bk. 4 on Luke iv}
\switchcolumn*

\LA
\noindent Vide cleméntiam Dómini Salvatóris: nec indignatióne commótus nec scélere offénsus, nec injúria violátus Judǽam déserit: quin étiam ímmemor injúriæ, memor cleméntiæ, nunc docéndo, nunc liberándo, nunc sanándo, infídæ plebis corda demúlcet. Et bene sanctus Lucas virum ab spíritu nequítiæ liberátum ante præmísit, et subdit féminæ sanitátem. Utrúmque enim sexum Dóminus curatúrus advénerat: sed prior sanári débuit, qui prior creátus est; nec prætermítti illa, quæ mobilitáte magis ánimi, quam pravitáte peccáverat.\par
\switchcolumn
\EN
\noindent Behold here how long-suffering is the Lord our Redeemer! Neither moved to anger against them, nor sickened at their guilt, nor outraged by their attacks, did He leave the Jews' country. Nay, forgetting their iniquity, and mindful only of His mercy, He strove to soften their hard and unbelieving hearts, sometimes by His teaching, and sometimes by freeing some of them, and sometimes by healing them. St. Luke doth well to tell us first of the man who was delivered from an unclean spirit, and then of the healing of a woman. The Lord indeed came to heal both sexes, but that must be healed first which was created first, and then must not she be passed by whose first sin arose rather from fickleness of heart than from depraved will.\par
\switchcolumn*

\LA
\begin{Resp}\Rsign Repléti sunt omnes Spíritu Sancto: et cœpérunt loqui, prout Spíritus Sanctus dabat éloqui illis: \Star{} Et convénit multitúdo dicéntium, allelúja.\end{Resp}
\begin{Resp}\Vsign Loquebántur váriis linguis Apóstoli magnália Dei.\end{Resp}
\begin{Resp}\Rsign Et convénit multitúdo dicéntium, allelúja.\end{Resp}
\switchcolumn
\EN
\begin{Resp}\Rsign They were all filled with the Holy Ghost, and began to speak, as the Holy Ghost gave them utterance; \Star{} And the multitude came together, saying: Alleluia.\end{Resp}
\begin{Resp}\Vsign The Apostles spoke in diverse tongues the wonderful works of God.\end{Resp}
\begin{Resp}\Rsign And the multitude came together, saying: Alleluia.\end{Resp}
\switchcolumn*

\LA
\LessonHead{Lectio II}
\switchcolumn
\EN
\LessonHead{Lesson II}
\switchcolumn*

\LA
\noindent Sábbato medicínæ Domínicæ ópera cœpta signíficat, ut inde nova creatúra cœ́perit, ubi vetus creatúra ante desívit: nec sub lege esse Dei Fílium, sed supra legem in ipso princípio designáret: nec solvi legem, sed impléri. Neque enim per legem, sed verbo factus est mundus, sicut légimus: Verbo Dómini cæli firmáti sunt. Non sólvitur ergo lex, sed implétur: ut fiat renovátio hóminis jam labéntis. Unde et Apóstolus ait: Exspoliántes vos véterem hóminem, indúite novum, qui secúndum Deum creátus est.\par
\switchcolumn
\EN
\noindent That the Lord began to heal on the Sabbath-day showeth in a figure how that the new creation beginneth where the old creation ended. It showeth, moreover, that the Son of God, Who is come not to destroy the law but to fulfill the law, (Matth. v. 17,) is not under the law, but above the law. Neither was it by the law, but by the Word, that the world was created, as it is written “By the Word of the Lord were the heavens made.” (Ps. xxxii. 6.) The law, then, is not destroyed, but fulfilled, in the Redemption of fallen man. Whence also the Apostle saith: “Put off, concerning the former conversation, the old man, which is corrupt according to the deceitful lusts and be renewed in the spirit of your mind and put on the new man, which after God is created in righteousness and true holiness.” (Eph. iv. 22.)\par
\switchcolumn*

\LA
\begin{Resp}\Rsign Jam non dicam vos servos, sed amícos meos; quia ómnia cognovístis, quæ operátus sum in médio vestri, allelúja: \Star{} Accípite Spíritum Sanctum in vobis Paráclitum: ille est, quem Pater mittet vobis, allelúja.\end{Resp}
\begin{Resp}\Vsign Vos amíci mei estis, si fecéritis quæ ego præcípio vobis.\end{Resp}
\begin{Resp}\Rsign Accípite Spíritum Sanctum in vobis Paráclitum: ille est, quem Pater mittet vobis, allelúja.\end{Resp}
\begin{Resp}\Vsign Glória Patri, et Fílio, \Star{} et Spirítui Sancto.\end{Resp}
\begin{Resp}\Rsign Accípite Spíritum Sanctum in vobis Paráclitum: ille est, quem Pater mittet vobis, allelúja.\end{Resp}
\switchcolumn
\EN
\begin{Resp}\Rsign Henceforth I call you not servants, but I have called you My friends, because ye have known all things, whatsoever I have done among you. Alleluia. \Star{} Receive ye the Holy Ghost, Who is your Comforter within you; the Same is He Whom the Father will send unto you. Alleluia.\end{Resp}
\begin{Resp}\Vsign Ye are My friends, if ye do whatsoever I command you.\end{Resp}
\begin{Resp}\Rsign Receive ye the Holy Ghost Who is your Comforter within you; the Same is He Whom the Father will send unto you. Alleluia.\end{Resp}
\begin{Resp}\Vsign Glory be to the Father, and to the Son, \Star{} and to the Holy Ghost.\end{Resp}
\begin{Resp}\Rsign Receive ye the Holy Ghost Who is your Comforter within you; the Same is He Whom the Father will send unto you. Alleluia.\end{Resp}
\switchcolumn*

\LA
\LessonHead{Lectio III}
\switchcolumn
\EN
\LessonHead{Lesson III}
\switchcolumn*

\LA
\noindent Et bene sábbato cœpit, ut ipsum se osténderet Creatórem, qui ópera opéribus intéxeret, et prosequerétur opus, quod ipse jam cœ́perat: ut si domum faber renováre dispónat, non a fundaméntis, sed a culmínibus íncipit sólvere vetustátem. Itaque ibi prius manum ádmovet, ubi ante desíerat: deínde a minóribus íncipit, ut ad majóra pervéniat. Liberáre a dǽmone et hómines, sed in verbo Dei possunt: resurrectiónem mórtuis imperáre, divínæ solíus est potestátis. Fortássis étiam in typo mulíeris illíus socrus Simónis et Andréæ, váriis críminum fébribus caro nostra languébat, et diversárum cupiditátum immódicis æstuábat illécebris. Nec minórem febrem amóris esse díxerim, quam calóris. Itaque illa ánimum, hæc corpus inflámmat. Febris enim nostra, avarítia est: febris nostra, libído est: febris nostra, luxúria est: febris nostra, ambítio est: febris nostra, iracúndia est.\par
\switchcolumn
\EN
\noindent It was well that He began to heal on the Sabbath, that He might show Himself to be the Creator, weaving in one with another of His works, and continuing that which He had already begun, even as a workman, being to repair an house, beginneth not to take down that which is old from the foundations, but from the roof. Thus doth the Lord begin to lay to His hand again, in that place whence last He hath lifted it; then He beginneth which after God is created in righteousness and true holiness: (Eph. iv. 22.) with things lesser, that He may go on to things greater. Even men are able to deliver other men from evil spirits, albeit with the word of God: to command the dead to rise again is for God's power alone. Perchance, also, this woman, the mother-in-law of Simon and Andrew, was a type of our nature, stricken down with the great fever of sin, and burning with unlawful lusts after diverse objects. Nor would I say that the passion which rageth in the mind is a lesser fire than that fever which burneth the body. Covetousness, and lust, and uncleanness, and vain desires, and strivings, and anger, these be our fevers.\par
\switchcolumn*

\LA
\TeDeum{Te Deum}
\switchcolumn
\EN
\TeDeum{Te Deum}
\switchcolumn*

\end{paracol}

\PartTitle{Proprium de Sanctis}{Proper of Saints}

\addcontentsline{toc}{chapter}{Proprium de Sanctis\enspace\textit{Proper of Saints}}

\SeasonTitle{Mensis Februarius}{February}

\addcontentsline{toc}{section}{Mensis Februarius\enspace\textit{February}}

\par\needspace{10\baselineskip}\addvspace{1.2em}{\centering\small\textsc{Die 7 Februarii} · \textsc{7 February}\par}
\Office{S. Romualdi Abbatis}{St. Romuald, Abbot}{Duplex}
\addcontentsline{toc}{subsection}{Die 7 Februarii · S. Romualdi Abbatis\enspace\textit{St. Romuald, Abbot}}
\begin{paracol}{2}
\LA
\RefLine{Lectiones I–III: de Scriptura occurrente.}
\switchcolumn
\EN
\RefLine{Lessons I–III: from the Scripture of the day.}
\switchcolumn*

\LA
\LessonHead{Lectio IV}
\switchcolumn
\EN
\LessonHead{Lesson IV}
\switchcolumn*

\LA
\noindent Romuáldus, Ravénnæ, Sérgio patre, nóbili génere natus, adoléscens in propínquum monastérium Classénse pœniténtiæ causa secéssit: ubi religiósi hóminis sermóne ad pietátis stúdium veheméntius incénsus, viso étiam semel et íterum per noctem in ecclésia beáto Apollinári, quod Dei servus illi futúrum promíserat, mónachus effícitur. Mox ad Marínum, vitæ sanctitáte ac severióre disciplína in fínibus Venetórum eo témpore célebrem, se cóntulit, ut ad arctam et sublímem perfectiónis viam eo magístro ac duce uterétur.\par
\switchcolumn
\EN
\noindent The holy Abbot Romuald was the son of one Sergius, of a noble family of Ravenna. While he was still very young, he went to a neighbouring monastery at Classis to do penance. While he was there he heard a discourse by a monk, which stirred him up strongly to aim at godliness of living; and he had afterwards in the Church by night two visions in which the blessed servant of God Apollinaris foretold to him that he should become a monk himself. He accordingly did so; and soon afterwards betook himself to one Marinus, whose holy life and strict discipline were then much noised about in all the coasts of the Venetians, that he might by his teaching and guidance attain towards the hard and lofty point of perfection.\par
\switchcolumn*

\LA
\begin{Resp}\Rsign Honéstum fecit illum Dóminus, et custodívit eum ab inimícis, et a seductóribus tutávit illum: \Star{} Et dedit illi claritátem ætérnam.\end{Resp}
\begin{Resp}\Vsign Justum dedúxit Dóminus per vias rectas, et osténdit illi regnum Dei.\end{Resp}
\begin{Resp}\Rsign Et dedit illi claritátem ætérnam.\end{Resp}
\switchcolumn
\EN
\begin{Resp}\Rsign The Lord made him honourable, and defended him from his enemies, and kept him safe from those that lay in wait for him; \Star{} And gave him perpetual glory.\end{Resp}
\begin{Resp}\Vsign He went down with him into the pit, and left him not in bonds.\end{Resp}
\begin{Resp}\Rsign And gave him perpetual glory.\end{Resp}
\switchcolumn*

\LA
\LessonHead{Lectio V}
\switchcolumn
\EN
\LessonHead{Lesson V}
\switchcolumn*

\LA
\noindent Multis sátanæ insídiis et hóminum invídia oppugnátus, tanto humílior se assídue jejúniis et oratiónibus exercébat, et rerum cæléstium meditatióne, vim lacrimárum profúndens, fruebátur: vultu tamen ádeo læto semper erat, ut intuéntes exhilaráret. Magno apud príncipes et reges in honóre fuit; multíque ejus consílio, mundi illécebris abjéctis, solitúdinem petiérunt. Martýrii quoque cupiditáte flagrávit, cujus causa dum in Pannóniam proficíscitur, morbo, quo afflictabátur cum progrederétur, levabátur cum recéderet, revérti cógitur.\par
\switchcolumn
\EN
\noindent The more he was assailed by the wiles of Satan and the unkindness of men, the more did he exercise himself in lowliness, with continual fasting and prayer, and rejoice in thinking of heavenly things, with abundance of tears. And all the while he bore so bright a face as gladdened all who looked on him. He was held in great honour by princes and kings, and his counsel moved many to leave the blandishments of the world and withdraw to the desert. He had such a burning desire to obtain the crown of martyrdom that he set out for Pannonia on purpose to seek it, but, falling into sickness whenever he went forward though growing strong again whenever he drew back, he behoved to return home.\par
\switchcolumn*

\LA
\begin{Resp}\Rsign Amávit eum Dóminus, et ornávit eum: stolam glóriæ índuit eum, \Star{} Et ad portas paradísi coronávit eum.\end{Resp}
\begin{Resp}\Vsign Índuit eum Dóminus lorícam fídei, et ornávit eum.\end{Resp}
\begin{Resp}\Rsign Et ad portas paradísi coronávit eum.\end{Resp}
\switchcolumn
\EN
\begin{Resp}\Rsign The Lord loved him and beautified him; He clothed him with a robe of glory, \Star{} And crowned him at the gates of Paradise.\end{Resp}
\begin{Resp}\Vsign The Lord hath put on him the breast-plate of faith, and hath adorned him.\end{Resp}
\begin{Resp}\Rsign And crowned him at the gates of Paradise.\end{Resp}
\switchcolumn*

\LA
\LessonHead{Lectio VI}
\switchcolumn
\EN
\LessonHead{Lesson VI}
\switchcolumn*

\LA
\noindent In vita et post mortem miráculis clarus, spíritu étiam prophetíæ non cáruit. Scalam a terra cælum pertingéntem, in similitúdinem Jacob Patriárchæ, per quam hómines in veste cándida ascendébant et descendébant, per visum conspéxit; eóque Camaldulénses mónachos, quorum institúti auctor fuit, designári mirabíliter agnóvit. Dénique cum annos centum et vigínti ágeret, et centum ipsos in summa vitæ asperitáte Deo servísset, ad eum migrávit, anno salútis millésimo vigésimo séptimo. Ejus corpus quinquénnio postquam sepúltum fúerat, íntegrum repértum, Fabriáni in ecclésia sui órdinis honorífice cónditum est.\par
\switchcolumn
\EN
\noindent God worked miracles by him both during his life and after his death, and likewise gave him the gift of prophecy. Like the Patriarch Jacob, he saw a ladder reaching from earth to heaven, and men in white garments ascending and descending upon it, in whom he marvellously knew were represented the monks of the Camaldolese Institute, of which he was the founder. At the age of 120 years, of which he had spent 100 in serving God in great hardness, he passed into His Presence, in the year of Salvation 1027. Five years after his death his body was found incorrupt, and laid in a magnificent grave in the Church of his order at Fabriano.\par
\switchcolumn*

\LA
\begin{Resp}\Rsign Iste homo perfécit ómnia quæ locútus est ei Deus, et dixit ad eum: Ingrédere in réquiem meam: \Star{} Quia te vidi justum coram me ex ómnibus géntibus.\end{Resp}
\begin{Resp}\Vsign Iste est, qui contémpsit vitam mundi, et pervénit ad cæléstia regna.\end{Resp}
\begin{Resp}\Rsign Quia te vidi justum coram me ex ómnibus géntibus.\end{Resp}
\begin{Resp}\Vsign Glória Patri, et Fílio, \Star{} et Spirítui Sancto.\end{Resp}
\begin{Resp}\Rsign Quia te vidi justum coram me ex ómnibus géntibus.\end{Resp}
\switchcolumn
\EN
\begin{Resp}\Rsign This is he which did according unto all that God commanded him; and God said unto him: Enter thou into My rest, \Star{} For thee have I seen righteous before Me among all people.\end{Resp}
\begin{Resp}\Vsign This is he which loved not his life in this world, and is come unto an everlasting kingdom.\end{Resp}
\begin{Resp}\Rsign For thee have I seen righteous before Me among all people.\end{Resp}
\begin{Resp}\Vsign Glory be to the Father, and to the Son, \Star{} and to the Holy Ghost.\end{Resp}
\begin{Resp}\Rsign For thee have I seen righteous before Me among all people.\end{Resp}
\switchcolumn*

\LA
\LessonHead{Lectio VII}
\switchcolumn
\EN
\LessonHead{Lesson VII}
\switchcolumn*

\LA
\LessonTitle{Léctio sancti Evangélii secúndum Matthǽum}
\switchcolumn
\EN
\LessonTitle{From the Holy Gospel according to Matthew}
\switchcolumn*

\LA
\Cite{Matt 19:27-29}
\switchcolumn
\EN
\Cite{Matt 19:27-29}
\switchcolumn*

\LA
\noindent In illo témpore: Dixit Petrus ad Jesum: Ecce nos relíquimus ómnia, et secúti sumus te: quid ergo erit nobis? Et réliqua.\par
\switchcolumn
\EN
\noindent In that time, Peter said to Jesus said to him: Behold we have left all things, and have followed thee: what therefore shall we have? And so on.\par
\switchcolumn*

\LA
\LessonTitle{Homilía sancti Bedæ Venerábilis Presbýteri}
\switchcolumn
\EN
\LessonTitle{Sermon of the Venerable Bede, Priest (at Jarrow)}
\switchcolumn*

\LA
\Source{Homilia in Natali S. Benedicti Ep.}
\switchcolumn
\EN
\Source{Homily for St. Benedict's Birth-day}
\switchcolumn*

\LA
\noindent Perféctus ille est, qui ábiens vendit ómnia quæ habet, et dat paupéribus, ac véniens séquitur Christum; habébit enim thesáurum non deficiéntem in cælis. Unde bene, interrogánte Petro, dixit tálibus Jesus: Amen dico vobis, quod vos qui secúti estis me, in regeneratióne, cum séderit Fílius hóminis in sede majestátis suæ, sedébitis et vos super sedes duódecim, judicántes duódecim tribus Israël. In hac quippe vita pro ejus nómine laborántes, in ália prǽmium speráre dócuit, id est, in regeneratióne; cum vidélicet in vitam immortálem fuérimus resurgéndo regeneráti, qui in vitam cadúcam mortáliter erámus géniti.\par
\switchcolumn
\EN
\noindent If thou wilt be perfect, saith Christ, go and sell that thou hast, and give to the poor, and come and follow Me and thou shalt have treasure in heaven. (Matth. xix. 21.) Yea, treasure that passeth not away! Unto such saith Jesus, at the questioning of Peter Amen I say unto you, that ye which have followed Me, in the regeneration, when the Son of Man shall sit in the throne of His glory, ye also shall sit upon twelve thrones, judging the twelve tribes of Israel. He taught them, which work for His Name's sake in this life, to look for their reward in another life that is, in the regeneration. In the regeneration! when we who have been born dying creatures into a dying life, shall in the resurrection be born again into an undying life.\par
\switchcolumn*

\LA
\begin{Resp}\Rsign Iste est qui ante Deum magnas virtútes operátus est, et de omni corde suo laudávit Dóminum: \Star{} Ipse intercédat pro peccátis ómnium populórum.\end{Resp}
\begin{Resp}\Vsign Ecce homo sine queréla, verus Dei cultor, ábstinens se ab omni ópere malo, et pérmanens in innocéntia sua.\end{Resp}
\begin{Resp}\Rsign Ipse intercédat pro peccátis ómnium populórum.\end{Resp}
\switchcolumn
\EN
\begin{Resp}\Rsign This is he which wrought great wonders before God, and praised the Lord with all his heart. \Star{} May he pray for all people, that their sins may be forgiven unto them.\end{Resp}
\begin{Resp}\Vsign Behold a man without blame, a worshipper of God in truth, keeping himself clean from every evil work, and abiding still in his innocency.\end{Resp}
\begin{Resp}\Rsign May he pray for all people, that their sins may be forgiven unto them.\end{Resp}
\switchcolumn*

\LA
\LessonHead{Lectio VIII}
\switchcolumn
\EN
\LessonHead{Lesson VIII}
\switchcolumn*

\LA
\noindent Et justa prorsus retribútio, ut, qui hic pro Christo humánæ glóriam celsitúdinis neglexérunt, illic a Christo júdices glorificáti singuláriter cum eo assídeant qui a sequéndis ejus vestígiis nulla ratióne póterant avélli. Nemo autem putet duódecim tantum Apóstolos, quia pro Juda prævaricánte Matthías eléctus est, tunc esse judicatúros; sicut nec duódecim solæ sunt tribus Israël judicándæ: alióquin tribus Levi, quæ tertiadécima est, injudicáta recédet.\par
\switchcolumn
\EN
\noindent And soothly, it is a just retribution, that they, who, while they were yet here, have for Christ's sake set no store by being great among men, should there by Christ be singularly glorified to be the assessors of His judgment-seat, even they whom nothing here could turn aside from being the followers of His footsteps. Nevertheless, let there be no man that believeth that the twelve Apostles only, among whom Matthias holdeth that place from which Judas by transgression fell, (Acts i. 25,) that they only shall judge, even as the twelve tribes of Israel shall not alone be judged for then were the tribe of Levi, which is the thirteenth, unjudged.\par
\switchcolumn*

\LA
\begin{Resp}\Rsign Sint lumbi vestri præcíncti, et lucérnæ ardéntes in mánibus vestris: \Star{} Et vos símiles homínibus exspectántibus dóminum suum, quando revertátur a núptiis.\end{Resp}
\begin{Resp}\Vsign Vigiláte ergo, quia nescítis qua hora Dóminus vester ventúrus sit.\end{Resp}
\begin{Resp}\Rsign Et vos símiles homínibus exspectántibus dóminum suum, quando revertátur a núptiis.\end{Resp}
\begin{Resp}\Vsign Glória Patri, et Fílio, \Star{} et Spirítui Sancto.\end{Resp}
\begin{Resp}\Rsign Et vos símiles homínibus exspectántibus dóminum suum, quando revertátur a núptiis.\end{Resp}
\switchcolumn
\EN
\begin{Resp}\Rsign Let your loins be girded about, and your lights burning; \Star{} And ye yourselves like unto men that wait for their lord, when he will return from the wedding.\end{Resp}
\begin{Resp}\Vsign Watch therefore, for ye know not what hour your Lord doth come.\end{Resp}
\begin{Resp}\Rsign And ye yourselves like unto men that wait for their lord, when he will return from the wedding.\end{Resp}
\begin{Resp}\Vsign Glory be to the Father, and to the Son, \Star{} and to the Holy Ghost.\end{Resp}
\begin{Resp}\Rsign And ye yourselves like unto men that wait for their lord, when he will return from the wedding.\end{Resp}
\switchcolumn*

\LA
\LessonHead{Lectio IX}
\switchcolumn
\EN
\LessonHead{Lesson IX}
\switchcolumn*

\LA
\noindent Et Paulus, qui tertiusdécimus est Apóstolus, judicándi sorte privábitur? cum ipse dicat: Nescítis quóniam ángelos judicábimus, quanto magis sæculária? Sciéndum namque est, omnes, qui ad exémplum Apostolórum, sua reliquérunt ómnia et secúti sunt Christum, júdices cum eo ventúros, sicut étiam omne mortálium genus esse judicándum. Quia enim duodenário sæpe número solet in Scriptúris univérsitas designári, per duódecim sedes Apostolórum ómnium numerósitas judicántium, et, per duódecim tribus Israël, univérsitas eórum qui judicándi sunt, osténditur.\par
\switchcolumn
\EN
\noindent Moreover, then, were Paul, who is the thirteenth Apostle, deprived of all part in the judgment; whereas he saith of himself: Know ye not that we shall judge angels? How much more things that pertain to this life? But it behoveth us to know that every one who hath forsaken all and followed Christ, as did the Apostles, shall also come with Him to judgment, even as every man shall stand at His judgment seat. And the Scriptures use often to signify all by this number twelve; by the twelve thrones of the Apostles are signified the thrones of all them that shall judge; and by the twelve tribes of Israel, the whole number of them that shall be judged.\par
\switchcolumn*

\LA
\TeDeum{Te Deum}
\switchcolumn
\EN
\TeDeum{Te Deum}
\switchcolumn*

\end{paracol}

\par\needspace{10\baselineskip}\addvspace{1.2em}{\centering\small\textsc{Die 8 Februarii} · \textsc{8 February}\par}
\Office{S. Joannis de Matha Confessoris}{St. John of Matha, Confessor}{Duplex}
\addcontentsline{toc}{subsection}{Die 8 Februarii · S. Joannis de Matha Confessoris\enspace\textit{St. John of Matha, Confessor}}
\begin{paracol}{2}
\LA
\RefLine{Lectiones I–III: de Scriptura occurrente.}
\switchcolumn
\EN
\RefLine{Lessons I–III: from the Scripture of the day.}
\switchcolumn*

\LA
\LessonHead{Lectio IV}
\switchcolumn
\EN
\LessonHead{Lesson IV}
\switchcolumn*

\LA
\noindent Joánnes de Matha, ordinis sanctíssimæ Trinitátis redemptiónis captivórum institutor, Falcone in Provincia natus est paréntibus pietáte et nobilitate conspicuis. Studiórum causa Aquas Sextias, mox Parisios profectus, confectóque theologíæ curriculo, magisterii lauream adeptus, doctrinæ et virtútum splendore enituit. Quibus motus Parisiénsis antístes, ad sacrum presbyterátus ordinem præ humilitate reluctántem promovit, eo consílio, ut in ea civitáte cómmorans, sapiéntia et móribus studiosæ juventúti prælucéret. Cum autem in sacello ejusdem epíscopi, ipso cum aliis astante, primum Deo Sacrum offerret, cælésti favore meruit recreari. Nam Angelus cándida et fulgénti veste indútus, cui in péctore crux rúbei et cærúlei coloris assuta erat, brácchiis cancellátis et super duos captivos ad látera positos, Christianum unum, álterum Maurum, exténsis, appáruit. Qua visióne in éxtasim raptus, intelléxit prótinus vir Dei, se ad redimendos ab infidelibus captivos destinari.\par
\switchcolumn
\EN
\noindent John de la Mata, the founder of the Order of the Most Holy Trinity for the Ransom of Prisoners, was born at Faucon, in Provence, (upon Midsummer's Day, in the year 1169,) and was the child of parents equally distinguished for their rank and their godly life. He went for his education first to Aix and then to Paris. At the University of Paris, where he went through the course of Divinity and took the degree of Doctor, he became eminent for learning and virtue. For this reason the Bishop of Paris ordained him Priest, an honour from which his lowliness caused him to shrink, in the hope that he should induce him to remain at Paris, and be a bright example of wisdom and manners to the students who resorted thither. He offered up the Holy Sacrifice to God for the first time in the private Chapel of the Bishop, and in the presence of that Prelate and diverse other persons. In the midst of the ceremony, a vision from God appeared to John. There appeared to him an angel, clad in raiment white and glistering; having sewn on his breast a cross of red and blue. His arms were crossed before him, and his hands were upon the heads of two slaves, one a Christian and the other a Moor. And immediately the man of God was in the spirit, and knew that he was called to the work of ransoming bondsmen from the power of the unbelievers.\par
\switchcolumn*

\LA
\begin{Resp}\Rsign Honéstum fecit illum Dóminus, et custodívit eum ab inimícis, et a seductóribus tutávit illum: \Star{} Et dedit illi claritátem ætérnam.\end{Resp}
\begin{Resp}\Vsign Justum dedúxit Dóminus per vias rectas, et osténdit illi regnum Dei.\end{Resp}
\begin{Resp}\Rsign Et dedit illi claritátem ætérnam.\end{Resp}
\switchcolumn
\EN
\begin{Resp}\Rsign The Lord made him honourable, and defended him from his enemies, and kept him safe from those that lay in wait for him; \Star{} And gave him perpetual glory.\end{Resp}
\begin{Resp}\Vsign He went down with him into the pit, and left him not in bonds.\end{Resp}
\begin{Resp}\Rsign And gave him perpetual glory.\end{Resp}
\switchcolumn*

\LA
\LessonHead{Lectio V}
\switchcolumn
\EN
\LessonHead{Lesson V}
\switchcolumn*

\LA
\noindent Quo vero matúrius in re tanti moménti procéderet, in solitúdinem secessit, ibique divino nutu factum est, ut Felicem Valesium in ipsa erémo jam multis annis degentem repérerit: cum quo ínita societáte, se per triennium in oratióne, et contemplatióne, omniúmque virtútum studio exercuit. Cóntigit autem, ut, dum secum de rebus divinis prope fontem colloqueréntur, cervus ad eos accesserit, crucem inter córnua gerens rúbei et cærúlei coloris. Cumque Felix ob rei novitátem mirarétur, narrávit ei Joánnes visiónem in prima Missa hábitam; et exínde ferventius oratióni incumbéntes, ter in somnis admoniti, Romam proficísci decrevérunt, ut a summo Pontifice novi ordinis pro rediméndis captivis institutiónem impetrarent. Electus fuerat eo témpore Innocentius tertius; qui, illis benígne acceptis, dum secum de re proposita deliberat, in festo sanctæ Agnetis secundo, Laterani intra Missárum solemnia ad sacræ Hostiæ elevatiónem, Angelus ei cándida veste, cruce bicolori, spécie redimentis captivos appáruit. Quo viso, Pontifex institutum approbávit, et novum ordinem sanctíssimæ Trinitátis redemptiónis captivórum vocari jussit, ejusque professóribus albas vestes cum cruce rúbei et cærúlei coloris præbuit.\par
\switchcolumn
\EN
\noindent That he might set himself with due forethought to the carrying out of his work, he withdrew into a certain desert, and there, by the will of God, he found Felix de Valois, who had already spent many years in that place. With him he joined company, and they passed three years together in continual prayer, meditation, and all spiritual exercises. It came to pass, one day, when they were sitting on the bank of a spring, that there came to them a stag having between his horns a cross of red and blue. Felix cried out in wonder at that sight, and John then told him of the vision that had appeared to him when he was saying his first Mass. Thenceforth they gave themselves with redoubled fervour to prayer, and, being three times warned in sleep, they determined to go to Rome, and pray the Pope to institute an Order for the ransom of prisoners. They arrived at the time of the election of Innocent III., who received them courteously, and entertained in his mind their petition. While he was in consideration, he went to the Lateran Cathedral, on the second Feast of St. Agnes, and there, while Mass was being solemnly sung, at the moment of the elevation of the Sacred Host, there appeared to him an angel, clad in raiment white and glistering, having sewn on his breast a cross of red and blue, and making as though he would free prisoners. Thereupon the Pope founded the Order, commanding that it should be called the Order of the Most Holy Trinity for the Ransom of prisoners, and that they who professed in it should be clad in white raiment, having sewn on their breasts a cross of red and blue.\par
\switchcolumn*

\LA
\begin{Resp}\Rsign Amávit eum Dóminus, et ornávit eum: stolam glóriæ índuit eum, \Star{} Et ad portas paradísi coronávit eum.\end{Resp}
\begin{Resp}\Vsign Índuit eum Dóminus lorícam fídei, et ornávit eum.\end{Resp}
\begin{Resp}\Rsign Et ad portas paradísi coronávit eum.\end{Resp}
\switchcolumn
\EN
\begin{Resp}\Rsign The Lord loved him and beautified him; He clothed him with a robe of glory, \Star{} And crowned him at the gates of Paradise.\end{Resp}
\begin{Resp}\Vsign The Lord hath put on him the breast-plate of faith, and hath adorned him.\end{Resp}
\begin{Resp}\Rsign And crowned him at the gates of Paradise.\end{Resp}
\switchcolumn*

\LA
\LessonHead{Lectio VI}
\switchcolumn
\EN
\LessonHead{Lesson VI}
\switchcolumn*

\LA
\noindent Sic stabilito ordine, sancti fundatores in Gálliam rediérunt, primoque cœnobio Cervi Frígidi in diœcesi Meldénsi constructo, ad ejus regimen Felix remánsit; et Joánnes Romam cum aliquot sociis reversus est, ubi Innocentius domum, ecclésiam et hospitale sancti Thomæ de Formis in monte Cælio eis donávit cum multis redítibus et possessiónibus. Datis quoque litteris ad Miramolínum regem Marochii, opus redemptiónis felici auspicio inchoátum fuit. Tum ad Hispánias, sub jugo Saracenórum magna ex parte oppressas, Joánnes profectus est, regumque, príncipum atque aliórum fidelium animos ad captivórum et páuperum commiseratiónem commovit. Monasteria ædificávit, hospitalia eréxit, magnoque lucro animárum plures captivos redémit. Romam tandem reversus, sanctisque opéribus incumbens, assiduis labóribus attritus et morbo confectus, ardentíssimo Dei et próximi amóre exæstuans, ad extremum devenit. Quare frátribus convocatis, eisque ad opus redemptiónis cælitus præmonstrátum efficaciter cohortátis, obdormívit in Dómino sextodecimo Kalendas Januarii, anno salútis millesimo ducentésimo décimo tertio; ejusque corpus in ipsa ecclésia sancti Thomæ de Formis condigno honóre tumulátum fuit.\par
\switchcolumn
\EN
\noindent The Order being thus established, the holy Founders returned into France, and built their first Convent at Cerfroid, in the diocese of Meaux. Felix remained in charge of this house, and John went back to Rome with several companions. To them Innocent gave the house, Church, and hospital of St. Thomas de Formis on the Coelian Mount, with great endowments and property. Moreover he gave them a letter of introduction to Miramolin, King of Morocco, and they began with bright hopes the work of ransoming prisoners. John next betook himself to Spain, great part of which was then in the hands of the Saracens, and stirred up the hearts of the kings, princes, and all the faithful to have pity on slaves and the poor. He built Convents, founded Hospitals, and ransomed many bondsmen, to the great gain of souls. At last he returned to Rome, still busied in good works, but worn out by unceasing toil, and weakened by sickness. As he drew near the end of his earthly pilgrimage, his burning love for God and for his neighbour suffered no diminution. He called together his brethren, and earnestly exhorted them to go on with that work of ransom which had been pointed out to them from heaven, and then fell asleep in the Lord, on the 21st day of December, 1213. His body was buried with due honour in the Church of St. Thomas de Formis.\par
\switchcolumn*

\LA
\begin{Resp}\Rsign Iste homo perfécit ómnia quæ locútus est ei Deus, et dixit ad eum: Ingrédere in réquiem meam: \Star{} Quia te vidi justum coram me ex ómnibus géntibus.\end{Resp}
\begin{Resp}\Vsign Iste est, qui contémpsit vitam mundi, et pervénit ad cæléstia regna.\end{Resp}
\begin{Resp}\Rsign Quia te vidi justum coram me ex ómnibus géntibus.\end{Resp}
\begin{Resp}\Vsign Glória Patri, et Fílio, \Star{} et Spirítui Sancto.\end{Resp}
\begin{Resp}\Rsign Quia te vidi justum coram me ex ómnibus géntibus.\end{Resp}
\switchcolumn
\EN
\begin{Resp}\Rsign This is he which did according unto all that God commanded him; and God said unto him: Enter thou into My rest, \Star{} For thee have I seen righteous before Me among all people.\end{Resp}
\begin{Resp}\Vsign This is he which loved not his life in this world, and is come unto an everlasting kingdom.\end{Resp}
\begin{Resp}\Rsign For thee have I seen righteous before Me among all people.\end{Resp}
\begin{Resp}\Vsign Glory be to the Father, and to the Son, \Star{} and to the Holy Ghost.\end{Resp}
\begin{Resp}\Rsign For thee have I seen righteous before Me among all people.\end{Resp}
\switchcolumn*

\LA
\LessonHead{Lectio VII}
\switchcolumn
\EN
\LessonHead{Lesson VII}
\switchcolumn*

\LA
\LessonTitle{Léctio sancti Evangélii secúndum Lucam}
\switchcolumn
\EN
\LessonTitle{From the Holy Gospel according to Luke}
\switchcolumn*

\LA
\Cite{Luc 12:35-40}
\switchcolumn
\EN
\Cite{Luke 12:35-40}
\switchcolumn*

\LA
\noindent In illo témpore: Dixit Jesus discípulis suis: Sint lumbi vestri præcincti, et lucernæ ardentes in mánibus vestris. Et réliqua.\par
\switchcolumn
\EN
\noindent At that time, Jesus said unto His disciples: Let your loins be girded about, and your lights burning. And so on.\par
\switchcolumn*

\LA
\LessonTitle{Homilía sancti Gregórii Papæ}
\switchcolumn
\EN
\LessonTitle{Homily by Pope St. Gregory (the Great.)}
\switchcolumn*

\LA
\Source{Homilia 13 in Evang.}
\switchcolumn
\EN
\Source{13th on the Gospels.}
\switchcolumn*

\LA
\noindent Sancti Evangélii, fratres caríssimi, apérta vobis est léctio recitata. Sed ne alíquibus ipsa ejus planíties alta fortásse videátur, eam sub brevitáte transcúrrimus, quátenus ejus exposítio ita nesciéntibus fiat cógnita, ut tamen sciéntibus non sit onerósa. Dóminus dicit: Sint lumbi vestri præcíncti. Lumbos enim præcíngimus, cum carnis luxúriam per continéntiam coarctámus. Sed quia minus est, mala non ágere, nisi étiam quisque stúdeat, et bonis opéribus insudáre, prótinus ádditur: Et lucérnæ ardéntes in mánibus vestris. Lucérnas quippe ardéntes in mánibus tenémus, cum per bona ópera próximis nostris lucis exémpla monstrámus. De quibus profécto opéribus Dóminus dicit: Lúceat lux vestra coram homínibus, ut vídeant ópera vestra bona, et gloríficent Patrem vestrum qui in cælis est.\par
\switchcolumn
\EN
\noindent Dearly beloved brethren, the words of the Holy Gospel, which have just been read, lie open before you, and, lest their very plainness should make them seem to some to be hard, we will go through them with such shortness as that neither may they which understand not remain unenlightened, nor they which understand be wearied. The Lord saith Let your loins be girded about. Now, we gird our loins about, when by continency we master the lustful inclination of the flesh. But, forasmuch as it sufficeth not for a man to abstain from evil deeds, if he strive not to join thereto the earnest doing of good works, it is immediately added And your lights burning. Our lights burn when, by good works, we give bright example to our neighbour; concerning which works the Lord saith Let your light so shine before men, that they may see your good works, and glorify your Father Which is in heaven. (Matth. v. 16.)\par
\switchcolumn*

\LA
\begin{Resp}\Rsign Iste est qui ante Deum magnas virtútes operátus est, et de omni corde suo laudávit Dóminum: \Star{} Ipse intercédat pro peccátis ómnium populórum.\end{Resp}
\begin{Resp}\Vsign Ecce homo sine queréla, verus Dei cultor, ábstinens se ab omni ópere malo, et pérmanens in innocéntia sua.\end{Resp}
\begin{Resp}\Rsign Ipse intercédat pro peccátis ómnium populórum.\end{Resp}
\switchcolumn
\EN
\begin{Resp}\Rsign This is he which wrought great wonders before God, and praised the Lord with all his heart. \Star{} May he pray for all people, that their sins may be forgiven unto them.\end{Resp}
\begin{Resp}\Vsign Behold a man without blame, a worshipper of God in truth, keeping himself clean from every evil work, and abiding still in his innocency.\end{Resp}
\begin{Resp}\Rsign May he pray for all people, that their sins may be forgiven unto them.\end{Resp}
\switchcolumn*

\LA
\LessonHead{Lectio VIII}
\switchcolumn
\EN
\LessonHead{Lesson VIII}
\switchcolumn*

\LA
\noindent Duo autem sunt, quæ jubéntur, et lumbos restríngere, et lucérnas tenére: ut et mundítia sit castitátis in córpore, et lumen veritátis in operatióne. Redemptóri étenim nostro unum sine áltero placére nequáquam potest: si aut is qui bona agit, adhuc luxúriæ inquinaménta non déserit: aut is qui castitáte præéminet, necdum se per bona ópera exércet. Nec cástitas ergo magna est sine bono ópere, nec opus bonum est áliquod sine castitáte. Sed et si utrúmque ágitur, restat, ut quisquis ille est, spe ad supérnam pátriam tendat, et nequáquam se a vítiis pro mundi hujus honestáte contíneat.\par
\switchcolumn
\EN
\noindent Here, then, are two commandments, to gird our loins about, and to keep our lights burning the cleanness of purity in our body, and the light of the truth in our works. Whoso hath the one and not the other, pleaseth not thereby our Redeemer; that is, he pleaseth Him not which doth good works, but bridleth not himself from the pollutions of lust, neither he which is eminent in chastity, but exerciseth not himself in good works. Neither is chastity a great thing without good works, nor good works anything without chastity. And if any man do both, it remaineth that he must look by hope toward our Fatherland above, and not have for his reason wherethrough he turneth himself away from vice, the love of honour in this present world.\par
\switchcolumn*

\LA
\begin{Resp}\Rsign Sint lumbi vestri præcíncti, et lucérnæ ardéntes in mánibus vestris: \Star{} Et vos símiles homínibus exspectántibus dóminum suum, quando revertátur a núptiis.\end{Resp}
\begin{Resp}\Vsign Vigiláte ergo, quia nescítis qua hora Dóminus vester ventúrus sit.\end{Resp}
\begin{Resp}\Rsign Et vos símiles homínibus exspectántibus dóminum suum, quando revertátur a núptiis.\end{Resp}
\begin{Resp}\Vsign Glória Patri, et Fílio, \Star{} et Spirítui Sancto.\end{Resp}
\begin{Resp}\Rsign Et vos símiles homínibus exspectántibus dóminum suum, quando revertátur a núptiis.\end{Resp}
\switchcolumn
\EN
\begin{Resp}\Rsign Let your loins be girded about, and your lights burning; \Star{} And ye yourselves like unto men that wait for their lord, when he will return from the wedding.\end{Resp}
\begin{Resp}\Vsign Watch therefore, for ye know not what hour your Lord doth come.\end{Resp}
\begin{Resp}\Rsign And ye yourselves like unto men that wait for their lord, when he will return from the wedding.\end{Resp}
\begin{Resp}\Vsign Glory be to the Father, and to the Son, \Star{} and to the Holy Ghost.\end{Resp}
\begin{Resp}\Rsign And ye yourselves like unto men that wait for their lord, when he will return from the wedding.\end{Resp}
\switchcolumn*

\LA
\LessonHead{Lectio IX}
\switchcolumn
\EN
\LessonHead{Lesson IX}
\switchcolumn*

\LA
\noindent Et vos símiles homínibus exspectántibus dóminum suum, quando revertátur a núptiis: ut cum vénerit et pulsáverit, conféstim apériant ei. Venit quippe Dóminus, cum ad judícium próperat: pulsat vero, cum jam per ægritúdinis moléstias esse mortem vicínam desígnat. Cui conféstim aperímus, si hunc cum amóre suscípimus. Aperíre enim júdici pulsánti non vult, qui exíre de córpore trépidat: et vidére eum, quem contempsísse se méminit, júdicem formídat. Qui autem de sua spe et operatióne secúrus est, pulsanti conféstim áperit, quia lætus júdicem sústinet; et, cum tempus propínquæ mortis advénerit, de glória retributiónis hilaréscit.\par
\switchcolumn
\EN
\noindent And ye yourselves like unto men that wait for their lord, when he will return from the wedding that, when he cometh and knocketh, they may open unto him immediately. The Lord cometh at the hour of judgment He knocketh when, by the pains of sickness, He biddeth us know that death is nigh. To Him open we immediately, if we receive Him in love. Whoso feareth to leave this body, will not open to the Judge when He knocketh, for he dreadeth to see that Judge, Whom he knoweth that he hath despised. But whosoever knoweth that his hope and works are built upon a good foundation, when he heareth the Judge knock, openeth to Him immediately, for to such an one that coming is blessed, yea, when the hour of death is at hand, such an one haileth with gladness a glorious reward.\par
\switchcolumn*

\LA
\TeDeum{Te Deum}
\switchcolumn
\EN
\TeDeum{Te Deum}
\switchcolumn*

\end{paracol}

\par\needspace{10\baselineskip}\addvspace{1.2em}{\centering\small\textsc{Die 9 Februarii} · \textsc{9 February}\par}
\Office{S. Cyrilli Episc. Alexandrini Confessoris et Ecclesiæ Doctoris}{St. Cyril of Alexandria, Bishop, Confessor, and Doctor of the Church}{Duplex}
\addcontentsline{toc}{subsection}{Die 9 Februarii · S. Cyrilli Episc. Alexandrini Confessoris et Ecclesiæ Doctoris\enspace\textit{St. Cyril of Alexandria, Bishop, Confessor, and Doctor of the Church}}
\begin{paracol}{2}
\LA
\RefLine{Lectiones I–III: de Scriptura occurrente.}
\switchcolumn
\EN
\RefLine{Lessons I–III: from the Scripture of the day.}
\switchcolumn*

\LA
\RefLine{Lectiones VII–VIII: ex Commúni Doctoris Pontificis (C4a).}
\switchcolumn
\EN
\RefLine{Lessons VII–VIII: from the Common of a Doctor Bishop (C4a).}
\switchcolumn*

\LA
\RefLine{Lectio IX: de S. Apolloniæ Virginis et Martyris (02-09o).}
\switchcolumn
\EN
\RefLine{Lesson IX: of S. Apollonia, Virgin and Martyr (02-09o).}
\switchcolumn*

\LA
\LessonHead{Lectio IV}
\switchcolumn
\EN
\LessonHead{Lesson IV}
\switchcolumn*

\LA
\noindent Cyríllus Alexandrínus, cujus præcónia non uníus tantum vel altérius sunt comprobáta testimónio, sed étiam œcumenicórum conciliórum Ephesíni et Chalcedonénsis actis celebráta, claris ortus paréntibus, ac Theóphili epíscopi Alexandríni nepos, adhuc adoléscens præcelléntis ingénii clara specímina dedit. Lítteris ac sciéntiis egrégie imbútus, ad Joánnem epíscopum, Jerosolymitánum se cóntulit, ut in Christiána fide perficerétur. Alexandríam deínde cum rediísset, Theóphilo vita functo, ad illíus sedem evéctus est: quo in múnere ita óptimi pastóris formam ab Apóstolo definítam constánter præ se tulit, ut sanctíssimi præsulis glóriam mérito sit adéptus.\par
\switchcolumn
\EN
\noindent The praises of Cyril of Alexandria have been celebrated not only by one writer or another, but have even been registered in the acts of the Ecumenical Councils of Ephesus and Chalcedon. He was born of distinguished parents, and was the nephew of Theophilus, Pope of Alexandria. While he was still young he displayed marks of his excellent understanding. After giving a deep study to letters and science he betook himself to John, Bishop of Jerusalem, to be perfected in the Christian faith. After his return to Alexandria, and the death of Theophilus, he was raised to that see. In this office he kept ever before his eyes the type of the Shepherd of souls as it had been laid down by the Apostle; and by ever adhering thereto deservedly earned the glory of an holy Bishop.\par
\switchcolumn*

\LA
\begin{Resp}\Rsign Invéni David servum meum, óleo sancto meo unxi eum: \Star{} Manus enim mea auxiliábitur ei.\end{Resp}
\begin{Resp}\Vsign Nihil profíciet inimícus in eo, et fílius iniquitátis non nocébit ei.\end{Resp}
\begin{Resp}\Rsign Manus enim mea auxiliábitur ei.\end{Resp}
\switchcolumn
\EN
\begin{Resp}\Rsign I have found David My servant, with My holy oil have I anointed him \Star{} For My hand shall help him.\end{Resp}
\begin{Resp}\Vsign The enemy shall prevail nothing against him, nor the son of wickedness afflict him.\end{Resp}
\begin{Resp}\Rsign For My hand shall help him.\end{Resp}
\switchcolumn*

\LA
\LessonHead{Lectio V}
\switchcolumn
\EN
\LessonHead{Lesson V}
\switchcolumn*

\LA
\noindent Salútis animárum zelo incénsus curas omnes inténdit, ut sibi commíssum gregem in fídei et morum integritáte serváret, atque a venenátis infidélium et hæreticórum páscuis defénderet. Hinc tum Nováti ásseclas e civitáte expélli, tum Judæos, qui furóre acti in cædem Christianórum conspiráverant, juxta leges puníri satégit. Singuláre vero Cyrílli pro cathólicæ fídei incolumitáte enítuit stúdium contra Nestórium Constantinopolitánum epíscopum, asseréntem Jesum Christum ex María Vírgine hóminem tantum et non Deum natum, eíque divinitátem pro méritis esse collátam; cujus emendatiónem cum frustra tentásset, eum sancto Cælestíno Pontífici máximo denuntiávit.\par
\switchcolumn
\EN
\noindent Zeal for the salvation of souls was kindled in him, and he undertook all care to keep in the faith and in soundness of life the flock unto him committed, and to preserve them from the poisonous pastures of infidelity and heresy; hence, in accordance with the laws, he caused the followers of Novatus to be expelled from the city, and those Jews to be punished who had been induced by rage to plan a massacre of the Christians. His eminent care for the preservation of the Catholic faith pure and undenled shone forth especially in his controversy against Nestorius, Patriarch of Constantinople, who asserted that Jesus Christ had been born of the Virgin Mary as man only and not as God, and that the Godhead had been bestowed upon Him because of His merits. Cyril first attempted to convert Nestorius, but when he found this hopeless he denounced him to the Supreme Pontiff the holy Celestine.\par
\switchcolumn*

\LA
\begin{Resp}\Rsign Pósui adjutórium super poténtem, et exaltávi eléctum de plebe mea: \Star{} Manus enim mea auxiliábitur ei.\end{Resp}
\begin{Resp}\Vsign Invéni David servum meum, óleo sancto meo unxi eum.\end{Resp}
\begin{Resp}\Rsign Manus enim mea auxiliábitur ei.\end{Resp}
\switchcolumn
\EN
\begin{Resp}\Rsign I have laid help upon one that is mighty, and have exalted one chosen out of My people \Star{} For My hand shall help him.\end{Resp}
\begin{Resp}\Vsign I have found David My servant, with My holy oil have I anointed him.\end{Resp}
\begin{Resp}\Rsign For My hand shall help him.\end{Resp}
\switchcolumn*

\LA
\LessonHead{Lectio VI}
\switchcolumn
\EN
\LessonHead{Lesson VI}
\switchcolumn*

\LA
\noindent Cælestíni delegáta auctoritáte, concílio Ephesíno præfuit, in quo hæresis Nestoriána pénitus proscrípta est, damnátus Nestórius et a sua Sede dejéctus, ac dogma cathólicum de una in Christo eáque divína persóna, et divína gloriósæ Vírginis Maríæ maternitáte assértum; plaudénte pópulo univérso, qui incredíbili gáudio géstiens, collucéntibus fácibus domum dedúxit epíscopus. Sed hac de causa Cyríllus calúmniis, injúriis et persecutiónibus, plúrimis a Nestório ejúsque fautóribus impetítus fuit; quas ipse patientíssime tulit, ita ut, de sola fide sollícitus, quidquid advérsus eum effutiébant ac moliebántur hærétici, pro níhilo habéret. Tandem pro Ecclésia Dei máximis perfúnctus labóribus, plurimísque scriptis éditis tum ad éthnicos et hæréticos confutándos, tum ad sacras Scriptúras et cathólica explanánda dógmata, sancto fine quiévit anno quadringentésimo quadragésimo quarto, episcopátus trigésimo secúndo. Leo décimus tértius Póntifex máximus Offícium et Missam præclaríssimi hujus fídei cathólicæ propugnatóris et Orientális ecclésiæ lúminis, ad Ecclésiam univérsam exténdit.\par
\switchcolumn
\EN
\noindent As delegate of Pope Celestine, Cyril presided at the Council of Ephesus where the Nestorian heresy was condemned; Nestorius deprived of his see; and the Catholic doctrine as to the unity of Person in Christ and the divine Motherhood of the glorious Virgin Mary was laid down amid the rejoicings of all the people, who escorted the bishops to their lodgings with a torch-light procession. For this reason Nestorius and his followers made Cyril the object of slanders, insults, and persecutions which he bore with profound patience, having all his care for the purity of the faith, and taking no heed to what the heretics might say or try against him. At length he died a holy death, in the year of salvation 444 and of his own papacy the 32nd. After vast work for the Church of God, and leaving behind him diverse writings directed either against heathens and heretics or to the exposition of the holy Scriptures and of Catholic doctrine, the Supreme Pontiff Leo XIII. extended to the Universal Church the Office and Mass of this most eminent champion of the Catholic faith, and light of the Eastern Church.\par
\switchcolumn*

\LA
\begin{Resp}\Rsign Iste est, qui ante Deum magnas virtútes operátus est, et omnis terra doctrína ejus repléta est: \Star{} Ipse intercédat pro peccátis ómnium populórum.\end{Resp}
\begin{Resp}\Vsign Iste est, qui contémpsit vitam mundi, et pervénit ad cæléstia regna.\end{Resp}
\begin{Resp}\Rsign Ipse intercédat pro peccátis ómnium populórum.\end{Resp}
\begin{Resp}\Vsign Glória Patri, et Fílio, \Star{} et Spirítui Sancto.\end{Resp}
\begin{Resp}\Rsign Ipse intercédat pro peccátis ómnium populórum.\end{Resp}
\switchcolumn
\EN
\begin{Resp}\Rsign This is he which wrought great wonders before God, and the whole earth is full of his teaching \Star{} May he pray for all people, that their sins may be forgiven unto them\end{Resp}
\begin{Resp}\Vsign This is he which loved not his life in this world, and hath attained unto the kingdom of heaven.\end{Resp}
\begin{Resp}\Rsign May he pray for all people, that their sins may be forgiven unto them\end{Resp}
\begin{Resp}\Vsign Glory be to the Father, and to the Son, \Star{} and to the Holy Ghost.\end{Resp}
\begin{Resp}\Rsign May he pray for all people, that their sins may be forgiven unto them\end{Resp}
\switchcolumn*

\end{paracol}

\par\needspace{10\baselineskip}\addvspace{1.2em}{\centering\small\textsc{Die 9 Februarii} · \textsc{9 February}\par}
\Office{S. Apolloniæ Virginis et Martyris}{S. Apollonia, Virgin and Martyr}{Simplex}
\addcontentsline{toc}{subsection}{Die 9 Februarii · S. Apolloniæ Virginis et Martyris\enspace\textit{S. Apollonia, Virgin and Martyr}}
\begin{paracol}{2}
\LA
\LessonHead{Lectio IX (pro commemoratione)}
\switchcolumn
\EN
\LessonHead{Lesson IX (when commemorated)}
\switchcolumn*

\LA
\LessonTitle{Commemoratio: S. Apolloniæ Virginis et Martyris}
\switchcolumn
\EN
\LessonTitle{Commemoration of: S. Apollonia, Virgin and Martyr}
\switchcolumn*

\LA
\noindent Apollónia virgo Alexandrína, sub Décio imperatóre, cum ingravescénte jam ætáte ad idóla sisterétur, ut eis veneratiónem adhibéret; illis contémptis, Jesum Christum verum Deum coléndum esse prædicábat. Quam ob rem omnes ei contúsi sunt et evúlsi dentes; ac, nisi Christum detestáta deos cóleret, accénso rogo combustúros vivam mináti sunt ímpii carnífices. Quibus illa, se quamvis mortem pro Jesu Christi fide subitúram, respóndit. Itaque comprehénsa ut comburerétur, cum paulísper, quasi delíberans quid agéndum esset, stetísset, ex illórum mánibus elápsa, álacris in ignem sibi parátum, majóri Spíritus Sancti flamma intus accénsa, se injécit. Unde brevi, consúmpto córpore, puríssimus spíritus in cælum ad sempitérnam martyrii corónam evolávit.\par
\switchcolumn
\EN
\noindent Apollonia was an aged virgin of Alexandria, who, (in the year of salvation 249,) in the reign of the Emperor Decius, was brought before the idols to worship them, but refused, declaring that Christ Jesus is True God, and that to Him worship is due. The cruel executioners beat and pulled out all her teeth, and threatened to burn her alive if she would not deny Christ. To whom she answered, that for Christ Jesus' sake she was ready to die. Being taken to the place of execution she stood for a few moments as if in doubt, and then, the fire of the Holy Ghost burning up in her heart, she broke from those that held her, and leapt of her own accord into the flames. Her body was quickly consumed, and her soul departed pure to obtain the eternal crown of martyrdom.\par
\switchcolumn*

\LA
\TeDeum{Te Deum}
\switchcolumn
\EN
\TeDeum{Te Deum}
\switchcolumn*

\end{paracol}

\par\needspace{10\baselineskip}\addvspace{1.2em}{\centering\small\textsc{Die 10 Februarii} · \textsc{10 February}\par}
\Office{S. Scholasticæ Virginis}{St. Scholastica, Virgin}{Duplex}
\addcontentsline{toc}{subsection}{Die 10 Februarii · S. Scholasticæ Virginis\enspace\textit{St. Scholastica, Virgin}}
\begin{paracol}{2}
\LA
\RefLine{Lectiones I–III: de Scriptura occurrente.}
\switchcolumn
\EN
\RefLine{Lessons I–III: from the Scripture of the day.}
\switchcolumn*

\LA
\RefLine{Lectiones VII–IX: ex Commúni Virginum tantum (C6a).}
\switchcolumn
\EN
\RefLine{Lessons VII–IX: from the Common of a Virgin not a Martyr (C6a).}
\switchcolumn*

\LA
\LessonHead{Lectio IV}
\switchcolumn
\EN
\LessonHead{Lesson IV}
\switchcolumn*

\LA
\LessonTitle{Ex libro Dialogórum sancti Gregórii Papæ}
\switchcolumn
\EN
\LessonTitle{From the Second Book of the Dialogues of Pope St. Gregory (the Great.)}
\switchcolumn*

\LA
\Source{Lib. 2, Cap. 33.}
\switchcolumn
\EN
\Source{Ch. 33}
\switchcolumn*

\LA
\noindent Scholástica, venerábilis patris Benedícti soror, omnipoténti Dómino ab ipso infántiæ témpore dedicáta, ad eum semel per annum veníre consuéverat; ad quam vir Dei non longe extra jánuam in possessióne monastérii descendébat. Quadam vero die venit ex more, atque ad eam, cum discípulis venerábilis ejus descéndit frater; qui totum diem in Dei láudibus sacrísque collóquiis ducéntes, incumbéntibus jam noctis ténebris, simul accepérunt cibum. Cumque adhuc ad mensam sedérent, et inter sacra collóquia tárdior se hora protráheret, éadem sanctimoniális fémina soror ejus eum rogávit, dicens: Quæso te, ut ista nocte me non déseras, ut usque mane de cæléstis vitæ gáudiis loquámur. Cui ille respóndit: Quid est quod lóqueris, soror? manére extra cellam nullátenus possum. Tanta vero erat cæli serénitas, ut nulla in áëre nubes apparéret. Sanctimoniális autem fémina, cum verba fratris negántis pósuit, et caput in mánibus omnipoténtem Dóminum rogatúra declinávit. Cumque leváret de mensa caput, tanta coruscatiónis et tonítrui virtus, tántaque inundátio plúviæ erúpit, ut neque venerábilis Benedíctus, neque fratres, qui cum eo áderant, extra loci limen, quo conséderant, pedem movére potúerint.\par
\switchcolumn
\EN
\noindent The worshipful Scholastica, the sister of our Father Benedict, was hallowed unto the Lord Almighty from a child. Her custom was to come to see her brother once every year. And when she came, the man of God went down unto her, not far from the gate, but, as it were, within the borders of his monastery. And there was a day when she came, as her custom was, and her worshipful brother went down to her, and his disciples with him. Then they passed the whole day together, praising God, and speaking one to the other of spiritual things. And when the night came, they brake bread together. And while they were yet at table, and conversed together on spiritual things, the hour was late. Then the holy woman his sister besought him, saying Leave me not, I pray thee, this night, but let us speak even until morning of the gladness of the eternal life. He answered her: What is it that thou sayest, my sister? I can by no means remain out of my cell. Now the firmament was so clear that there were no clouds in the sky. Then the holy nun, when she had heard the words of her brother, that he would not abide with her, clasped her hands on the table, and laid her face on her hands, and besought the Lord Almighty. And it came to pass that when she lifted up her head from the table, there were great thunderings and lightnings, and a flood of rain, insomuch that neither the worshipful Benedict nor the brethren that were with him could move as much as a foot over the threshold of the place where they sat.\par
\switchcolumn*

\LA
\begin{Resp}\Rsign Propter veritátem, et mansuetúdinem, et justítiam: \Star{} Et dedúcet te mirabíliter déxtera tua.\end{Resp}
\begin{Resp}\Vsign Spécie tua et pulchritúdine tua inténde, próspere procéde, et regna.\end{Resp}
\begin{Resp}\Rsign Et dedúcet te mirabíliter déxtera tua.\end{Resp}
\switchcolumn
\EN
\begin{Resp}\Rsign Because of truth, and meekness, and righteousness; \Star{} And thy right hand shall lead thee wonderfully.\end{Resp}
\begin{Resp}\Vsign In thy comeliness, and thy beauty, go forward, fare prosperously, and reign.\end{Resp}
\begin{Resp}\Rsign And thy right hand shall lead thee wonderfully.\end{Resp}
\switchcolumn*

\LA
\LessonHead{Lectio V}
\switchcolumn
\EN
\LessonHead{Lesson V}
\switchcolumn*

\LA
\noindent Sanctimoniális quippe fémina caput in mánibus declínans, lacrimárum flúvium in mensam fúderat, per quas serenitátem áëris ad plúviam traxit. Nec paulo tárdius post oratiónem inundátio illa secúta est; sed tanta fuit conveniéntia oratiónis et inundatiónis, ut de mensa caput jam cum tonítruo leváret; quátenus unum idémque esset moméntum, et leváre caput, et plúviam depónere. Tunc vir Dei inter corúscos et tonítruos atque ingéntis plúviæ inundatiónem, videns se ad monastérium non posse remeáre, cœpit cónqueri contristátus, dicens: Parcat tibi omnípotens Deus, soror: quid est quod fecísti? Cui illa respóndit: Ecce rogávi te et audíre me noluísti; rogávi Deum meum, et audívit me: modo ergo, si potes, egrédere et, me dimíssa, ad monastérium recéde. Ipse autem exíre extra tectum non valens, qui remanére sponte nóluit in loco, mansit invítus. Sicque factum est, ut totam noctem pervígilem dúcerent, atque per sacra spiritális vitæ collóquia, sese vicária relatióne satiárent.\par
\switchcolumn
\EN
\noindent Now when the holy woman laid her head in her hands upon the table, she wept bitterly, and as she wept, the clearness of the sky was turned to a tempest. As she prayed, immediately the flood followed. And the time was so, that she lifted up her head when it thundered, and when she had lifted up her head, the rain came. When the man of God saw that he could not return to his monastery, because of the lightnings, and thunderings, and the great rain, he was sorrowful and grieved, saying Almighty God forgive thee, my sister; what is this that thou hast done? She answered him Behold, I besought thee, and thou wouldest not hear; I besought my God, and He hath heard me; if, therefore, thou wilt, go forth, leave me alone, and go thy way to thy monastery. But he could not, and so he tarried in the same place, not willingly, but of necessity. And so it came to pass that they slept not all that night, but fed one another with discourse on spiritual things.\par
\switchcolumn*

\LA
\begin{Resp}\Rsign Dilexísti justítiam, et odísti iniquitátem: \Star{} Proptérea unxit te Deus, Deus tuus, óleo lætítiæ.\end{Resp}
\begin{Resp}\Vsign Propter veritátem, et mansuetúdinem, et justítiam.\end{Resp}
\begin{Resp}\Rsign Proptérea unxit te Deus, Deus tuus, óleo lætítiæ.\end{Resp}
\switchcolumn
\EN
\begin{Resp}\Rsign Thou hast loved righteousness, and hated iniquity; \Star{} Therefore God, thy God, hath anointed thee with the oil of gladness.\end{Resp}
\begin{Resp}\Vsign Because of truth, and meekness, and righteousness.\end{Resp}
\begin{Resp}\Rsign Therefore God, thy God, hath anointed thee with the oil of gladness.\end{Resp}
\switchcolumn*

\LA
\LessonHead{Lectio VI}
\switchcolumn
\EN
\LessonHead{Lesson VI}
\switchcolumn*

\LA
\noindent Cumque die áltero éadem venerábilis fémina ad cellam própriam recessísset, vir Dei ad monastérium rédiit. Cum ecce post tríduum, in cella consístens, elevátis in áëra óculis vidit ejúsdem soróris suæ ánimam de córpore egréssam in colúmbæ spécie cæli secréta penetráre. Qui, tantæ ejus glóriæ congáudens, omnipoténti Deo in hymnis et láudibus grátias réddidit, ejúsque óbitum frátribus denuntiávit. Quos étiam prótinus misit, ut ejus corpus ad monastérium deférrent, atque in sepúlcro, quod sibi ipsi paráverat, pónerent. Quo facto cóntigit, ut, quorum mens una semper in Deo fúerat, eórum quoque córpora nec sepultúra separáret.\par
\switchcolumn
\EN
\noindent And when the morning was come, the worshipful woman arose, and went unto her own cell, and the man of God went back to his monastery. And, behold, after three days he was sitting in his cell, and he lifted up his eyes to heaven, and saw the soul of his sister, delivered from the body, fly to heaven in a bodily shape like a dove. Wherefore he rejoiced because of the glory that was revealed in her, and gave thanks to Almighty God in hymns and praises, and made known to the brethren that she was dead. He commanded them also to go and take up her body, and bring it to his monastery, and lay it in the grave which he had made ready for himself. Whereby it came to pass that they twain who had ever been of one mind in the Lord, even in death were not divided.\par
\switchcolumn*

\LA
\begin{Resp}\Rsign Afferéntur Regi vírgines post eam, próximæ ejus; \Star{} Afferéntur tibi in lætítia et exsultatióne.\end{Resp}
\begin{Resp}\Vsign Spécie tua et pulchritúdine tua; inténde, próspere procéde, et regna.\end{Resp}
\begin{Resp}\Rsign Afferéntur tibi in lætítia et exsultatióne.\end{Resp}
\begin{Resp}\Vsign Glória Patri, et Fílio, \Star{} et Spirítui Sancto.\end{Resp}
\begin{Resp}\Rsign Afferéntur tibi in lætítia et exsultatióne.\end{Resp}
\switchcolumn
\EN
\begin{Resp}\Rsign After her shall virgins be brought unto the King, her fellows \Star{} They shall be brought unto thee with gladness and rejoicing.\end{Resp}
\begin{Resp}\Vsign In thy comeliness and thy beauty, go forward, fare prosperously, and reign.\end{Resp}
\begin{Resp}\Rsign They shall be brought unto thee with gladness and rejoicing.\end{Resp}
\begin{Resp}\Vsign Glory be to the Father, and to the Son, \Star{} and to the Holy Ghost.\end{Resp}
\begin{Resp}\Rsign They shall be brought unto thee with gladness and rejoicing.\end{Resp}
\switchcolumn*

\end{paracol}

\par\needspace{10\baselineskip}\addvspace{1.2em}{\centering\small\textsc{Die 11 Februarii} · \textsc{11 February}\par}
\Office{In Apparitione Beatæ Mariæ Virginis Immaculatæ}{In Apparitione Beatæ Mariæ Virginis Immaculatæ}{Duplex majus}
\addcontentsline{toc}{subsection}{Die 11 Februarii · In Apparitione Beatæ Mariæ Virginis Immaculatæ\enspace\textit{In Apparitione Beatæ Mariæ Virginis Immaculatæ}}
\begin{paracol}{2}
\LA
\LessonHead{Lectio I}
\switchcolumn
\EN
\LessonHead{Lesson I}
\switchcolumn*

\LA
\LessonTitle{De Parábolis Salomónis}
\switchcolumn
\EN
\LessonTitle{Lesson from the book of Proverbs}
\switchcolumn*

\LA
\Cite{Prov 8:12-17}
\switchcolumn
\EN
\Cite{Prov 8:12-17}
\switchcolumn*

\LA
\noindent Ego sapiéntia hábito in consílio et erudítis intérsum cogitatiónibus. Timor Dómini odit malum: arrogántiam, et supérbiam, et viam pravam, et os bilíngue detéstor. Meum est consílium et ǽquitas, mea est prudéntia, mea est fortitúdo. Per me reges regnant, et legum conditóres justa decérnunt; per me príncipes ímperant, et poténtes decérnunt justítiam. Ego diligéntes me díligo; et qui mane vígilant ad me, invénient me.\par
\switchcolumn
\EN
\noindent I wisdom dwell in counsel, and am present in learned thoughts. The fear of the Lord hateth evil: I hate arrogance, and pride, and every wicked way, and a mouth with a double tongue. Counsel and equity is mine, prudence is mine, strength is mine. By me kings reign, and lawgivers decree just things, By me princes rule, and the mighty decree justice. I love them that love me: and they that in the morning early watch for me, shall find me.\par
\switchcolumn*

\LA
\begin{Resp}\Rsign Sapiéntia quæ attíngit a fine usque ad finem fórtiter, et dispónit ómnia suáviter, ædificávit sibi domum: \Star{} Ecce tabernáculum Dei cum homínibus.\end{Resp}
\begin{Resp}\Vsign Vidi sanctam civitátem Jerúsalem novam, parátam sicut sponsam ornátam viro suo.\end{Resp}
\begin{Resp}\Rsign Ecce tabernáculum Dei cum homínibus.\end{Resp}
\switchcolumn
\EN
\begin{Resp}\Rsign Wisdom which reacheth from end to end mightily, and ordereth all things sweetly, hath built herself a house: \Star{} Behold the tabernacle of God with men.\end{Resp}
\begin{Resp}\Vsign I saw the holy city, the new Jerusalem, prepared as a bride adorned for her husband.\end{Resp}
\begin{Resp}\Rsign Behold the tabernacle of God with men.\end{Resp}
\switchcolumn*

\LA
\LessonHead{Lectio II}
\switchcolumn
\EN
\LessonHead{Lesson II}
\switchcolumn*

\LA
\Cite{Prov 8:18-25}
\switchcolumn
\EN
\Cite{Prov 8:18-25}
\switchcolumn*

\LA
\noindent Mecum sunt divítiæ et glória, opes supérbæ et justítia. Mélior est enim fructus meus auro et lápide pretióso, et genímina mea argénto elécto. In viis justítiæ ámbulo, in médio semitárum judícii, ut ditem diligéntes me et thesáuros eórum répleam. Dóminus possidébit me in inítio viárum suárum, ántequam quidquam fáceret a princípio. Ab ætérno ordináta sum et ex antíquis, ántequam terra fíeret. Nondum erant abýssi, et ego jam concépta eram; necdum fontes aquárum erúperant, necdum montes gravi mole constíterant; ante colles ego parturiébar.\par
\switchcolumn
\EN
\noindent With me are riches and glory, glorious riches and justice. For my fruit is better than gold and the precious stone, and my blossoms than choice silver. I walk in the way of justice, in the midst of the paths of judgment, That I may enrich them that love me, and may fill their treasures. The Lord possessed me in the beginning of his ways, before he made any thing from the beginning. I was set up from eternity, and of old before the earth was made. The depths were not as yet, and I was already conceived. neither had the fountains of waters as yet sprung out: The mountains with their huge bulk had not as yet been established: before the hills I was brought forth:\par
\switchcolumn*

\LA
\begin{Resp}\Rsign Quasi arcus refúlgens inter nébulas, et quasi flos rosárum in diébus vernis, et quasi lília in tránsitu aquæ, \Star{} Sic fulget Virgo immaculáta.\end{Resp}
\begin{Resp}\Vsign Arcum meum ponam in núbibus, et erit signum fœ́deris mei vobíscum.\end{Resp}
\begin{Resp}\Rsign Sic fulget Virgo immaculáta.\end{Resp}
\switchcolumn
\EN
\begin{Resp}\Rsign As the rainbow giving light in the clouds, and as the flower of roses in the days of the spring, and as the lilies on the brink of the water, \Star{} Thus shineth the Virgin immaculate.\end{Resp}
\begin{Resp}\Vsign I will set my bow in the clouds, and it shall be a sign of my covenant with you.\end{Resp}
\begin{Resp}\Rsign Thus shineth the Virgin immaculate.\end{Resp}
\switchcolumn*

\LA
\LessonHead{Lectio III}
\switchcolumn
\EN
\LessonHead{Lesson III}
\switchcolumn*

\LA
\Cite{Prov 8:34-36; 9:1-5}
\switchcolumn
\EN
\Cite{Prov 8:34-36; 9:1-5}
\switchcolumn*

\LA
\noindent Beátus homo qui audit me, et qui vígilat ad fores meas cotídie, et obsérvat ad postes óstii mei. Qui me invénerit, invéniet vitam, et háuriet salútem a Dómino; qui autem in me peccáverit, lædet ánimam suam. Omnes, qui me odérunt, díligunt mortem. Sapiéntia ædificávit sibi domum, excídit colúmnas septem. Immolávit víctimas suas, míscuit vinum et propósuit mensam suam. Misit ancíllas suas ut vocárent ad arcem et ad mœ́nia civitátis: Si quis est párvulus, véniat ad me. Et insipiéntibus locúta est: Veníte, comédite panem meum, et bíbite vinum quod míscui vobis.\par
\switchcolumn
\EN
\noindent Blessed is the man that heareth me, and that watcheth daily at my gates, and waiteth at the posts of my doors. He that shall find me, shall find life, and shall have salvation from the Lord: But he that shall sin against me, shall hurt his own soul. All that hate me love death. Wisdom hath built herself a house, she hath hewn her out seven pillars. She hath slain her victims, mingled her wine, and set forth her table. She hath sent her maids to invite to the tower, and to the walls of the city: Whosoever is a little one, let him come to me. And to the unwise she said: Come, eat my bread, and drink the wine which I have mingled for you.\par
\switchcolumn*

\LA
\begin{Resp}\Rsign Surge, amíca mea, speciósa mea, et veni, colúmba mea: \Star{} Osténde mihi fáciem tuam, sonet vox tua in áuribus meis.\end{Resp}
\begin{Resp}\Vsign Vox túrturis audíta est in terra nostra.\end{Resp}
\begin{Resp}\Rsign Osténde mihi fáciem tuam, sonet vox tua in áuribus meis.\end{Resp}
\begin{Resp}\Vsign Glória Patri, et Fílio, \Star{} et Spirítui Sancto.\end{Resp}
\begin{Resp}\Rsign Osténde mihi fáciem tuam, sonet vox tua in áuribus meis.\end{Resp}
\switchcolumn
\EN
\begin{Resp}\Rsign Arise, my love, my beautiful one, and come, my dove: \Star{} Shew me thy face, let thy voice sound in my ears.\end{Resp}
\begin{Resp}\Vsign The voice of the turtle is heard in our land.\end{Resp}
\begin{Resp}\Rsign Shew me thy face, let thy voice sound in my ears\end{Resp}
\begin{Resp}\Vsign Glory be to the Father, and to the Son, \Star{} and to the Holy Ghost.\end{Resp}
\begin{Resp}\Rsign Shew me thy face, let thy voice sound in my ears\end{Resp}
\switchcolumn*

\LA
\LessonHead{Lectio IV}
\switchcolumn
\EN
\LessonHead{Lesson IV}
\switchcolumn*

\LA
\noindent Anno quarto a dogmática definitióne de immaculáto beátæ Vírginis Concéptu, ad Gavi flúminis oram prope óppidum Lourdes diœcésis Tarbiénsis in Gállia, ipsa Virgo in rupis sinu super specum Massabiélle puéllæ cuidam, vernácula lingua Bernadétte nuncupátæ, paupérrimæ quidem sed ingénuæ ac piæ, plúries se conspiciéndam óbtulit. Immaculáta Virgo juveníli ac benígno videbátur aspéctu, nívea veste niveóque pállio contécta, ac zona cærúlea succíncta; nudos pedes áurea rosa ornábat. Primo apparitiónis die, qui fuit undécimus Februárii anno millésimo octingentésimo quinquagésimo octávo, puéllam signum crucis rite piéque faciéndum edócuit, atque ad sacri rosárii recitatiónem, exémplo suo, corónam, quæ prius ex brácchio demíssa pendébat, manu advólvens, excitávit: quod in céteris étiam apparitiónibus prǽstitit. Altero autem apparitiónis die, puélla in simplicitáte cordis sui, diabólicam fraudem timens, lustrálem aquam in Vírginem effúdit; sed beáta Virgo, léniter arrídens, benigniórem illi vultum osténdit. Cum vero tértio apparuísset, puéllam ad specum per quíndecim dies invitávit. Exínde eam sǽpius est allocúta, ac pro peccatóribus oráre, terram deosculári, pœnitentiámque ágere est hortáta; deínde imperávit, ut sacerdótibus edíceret, ædificándum ibi esse sacéllum, solemnísque supplicatiónibus more illo accedéndum. Mandávit ínsuper ut e fonte, qui sub aréna adhuc latébat sed mox erat eruptúrus, aquam bíberet, eáque se abstérgeret. Dénique die festo Annuntiatiónis, percontánti eníxe puéllæ illíus nomen, cujus aspéctu tóties dignáta fúerat, Virgo, admótis péctori mánibus elatísque in cælum óculis, respóndit: Immaculáta Concéptio ego sum.\par
\switchcolumn
\EN
\noindent In the fourth year from the dogmatic definition of the Immaculate Conception of the Blessed Virgin, at the bank of the river Gave near the town of Lourdes of the diocese of Tarbes in France, the Virgin herself in a bend of the rock above the grotto of Massabiele, often shewed herself to a certain girl, called in the vernacular tongue Bernadette, indeed most poor but noble and pious, to be seen. The Immaculate Virgin appeared with a young and kind appearance, clothed with a white garment and white veil, and girt with a blue girdle; she adorned her bare feet with a golden rose. On the first day of the apparition, which was the eleventh of February in the one thousand eight hundred fifty-eighth year, she taught the girl the sign of the cross to be duly and piously made, and she incited her, by her example, to the recitation of the sacred rosary, turning over with her hand the chaplet, which before was hanging down from her arm: which she supplied also in the other apparitions. And on the second day of the apparition, the girl in the simplicity of her heart, fearing diabolic fraud, flung holy water on the Virgin; but the blessed Virgin, smiling gently, shewed her face more kindly to her. And when she appeared a third time, she invited the girl to the grotto for fifteen days. Thence she often addressed her; then commanded, that it might be declared to the priests, a chapel to be built there, and to be approached there for supplications in a manner of solemnity. Moreover she commanded that from the spring, which thus far was hidden under the sand but now was about to erupt, she might drink the water, and by it cleanse herself. Finally on the feast day of the Annunciation, to the girl earnestly inquiring her name, the Virgin, for whose appearance she was so often deemed worthy, the Virgin, her hands having been moved to her breast and her eyes raised to heaven, answered: I am the Immaculate Conception.\par
\switchcolumn*

\LA
\begin{Resp}\Rsign Quæ est ista, quæ progréditur quasi auróra consúrgens, \Star{} Pulchra ut luna, elécta ut sol?\end{Resp}
\begin{Resp}\Vsign Ipsa est colúmba mea, perfécta mea, immaculáta mea.\end{Resp}
\begin{Resp}\Rsign Pulchra ut luna, elécta ut sol.\end{Resp}
\switchcolumn
\EN
\begin{Resp}\Rsign Who is she, that cometh forth as the morning rising, \Star{} Fair as the moon, bright as the sun?\end{Resp}
\begin{Resp}\Vsign She is my dove, my perfect one, my undefiled.\end{Resp}
\begin{Resp}\Rsign Fair as the moon, bright as the sun?\end{Resp}
\switchcolumn*

\LA
\LessonHead{Lectio V}
\switchcolumn
\EN
\LessonHead{Lesson V}
\switchcolumn*

\LA
\noindent Percrebrescénte fama beneficiórum, quæ in sacro specu recepísse fidéles dicebántur, augebátur in dies hóminum concúrsus, quos loci relígio ad specum advocábat. Itaque prodigiórum fama puellǽque candóre mortus Tarbiénsis epíscopus, quarto ab enarrátis anno, post jurídiam factórum inquisitiónem, supernaturáles esse apparitiónis notas sua senténtia probávit, cultúmque Vírginis immaculátæ in eódem specu permísit. Mox ædificátum sacéllum: ex illa die pene innúmeræ fidélium turbæ, voti ac supplicatiónis causa, ex Gállia, Bélgio, Itália, Hispánia ceterísque Európæ provínciis necnon ex longínquis Américæ regiónibus quovis anno illuc advéniunt, noménque Immaculátæ de Lourdes ubíque terrárum inclaréscit. Fontis aqua, in cunctas orbis partes deláta, ægris sanitátem restítuit. Orbis vero cathólicus tantórum memor benefactórum, ædes sacras mirábili ópere ibi exstrúxit. Vexílla innúmera, acceptórum beneficiórum véluti monuménta, illuc a civitátibus ac géntibus missa, ædem Vírginis miro ornátu décorant. In hac sua véluti sede Immaculáta Virgo júgiter cólitur; intérdiu quidem précibus, religióso cantu solemnibúsque áliis cæremóniis; noctu vero sacris illis supplicatiónibus, quibus infinítæ propémodum peregrinántium turbæ céreis facibúsque accénsis procédunt et laudes beátæ Vírginis cóncinunt.\par
\switchcolumn
\EN
\noindent The fame of favours becoming widespread, which the faithful were said to have received in the sacred grotto, in time the assembly of men increased, whom the sanctity of the place called to the grotto. Therefore the bishop of Tarbes moved by the fame of the wonders and the purity of the girl, after a juridical inquiry of the happenings, approved by his decree the signs of the apparition to be supernatural, and permitted the worship of the immaculate Virgin in the same grotto. Soon a chapel is built: from that day the almost innumerable crowds of the faithful, for the purpose of prayer and supplication, from France, Belgium, Italy, Spain, and the other provinces of Europe and also from the remote regions of America arrived there in every year, and the name of the Immaculate of Lourdes becometh famous of countries everywhere. The water of the spring, having been carried into all parts of the world, restored health to the sick. And the Catholic world mindful of such great favours, built holy shrines there with wonderful work. Innumerable banners, as if monuments of received benefits, sent there by cities and nations, decorate the shrine of the Virgin with wonderful adornment. In this as if her own seat the immaculate Virgin is continually honoured: indeed by day with prayers, religious chant and other solemn ceremonies; and by night with those solemn supplications, by which the almost infinite crowds of pilgrims proceed with candles and lit torches and sing the praises of the blessed Virgin.\par
\switchcolumn*

\LA
\begin{Resp}\Rsign Erit in novíssimis diébus mons præparátus Vírgini Maríæ in vértice móntium, et elevábitur super cælos, et ibunt pópuli multi et dicent: \Star{} Veníte et ascendámus ad montem.\end{Resp}
\begin{Resp}\Vsign Sicut lætántium ómnium habitátio est in te.\end{Resp}
\begin{Resp}\Rsign Veníte et ascendámus ad montem.\end{Resp}
\switchcolumn
\EN
\begin{Resp}\Rsign In the last days the mountain a house for the Virgin Mary shall be prepared on the top of mountains, and it shall be exalted above the heavens, and many people shall go and say: \Star{} Come and let us go up to the mountain.\end{Resp}
\begin{Resp}\Vsign The dwelling in thee is as it were of all rejoicing.\end{Resp}
\begin{Resp}\Rsign Come and let us go up to the mountain.\end{Resp}
\switchcolumn*

\LA
\LessonHead{Lectio VI}
\switchcolumn
\EN
\LessonHead{Lesson VI}
\switchcolumn*

\LA
\noindent Peregrinatiónes hujúsmodi fidem, frigescénte sǽculo, excitásse, ánimum ad christiánam legem profiténdam addidísse, cultúmque Vírginis immaculátæ mirum in modum auxísse, ómnibus compértum est. In qua mirábili fídei professióne christiánus pópulus sacerdótes véluti duces habet, qui illuc suas plebes addúcunt. Ipse étiam Sacrórum antístites sanctum locum frequénter ádeunt, peregrinatiónibus præsunt, solemnioribúsque festis intérsunt. Nec ádeo rarum est ipsos Románæ Ecclésiæ purpurátos patres húmili peregrinórum more accedéntes conspícere. Ipsi quoque Románi Pontífices, pro sua erga Immaculátam de Lourdes pietáte, sacram ædem donis nobilíssimis cumulárunt. Pius nonus, sacris indulgéntiis, archiconfraternitátis privilégio ac minóris Basílicæ título ipsam insignívit; ac Deíparæ imáginem íbidem cultam, solémni ritu per legátum suum apostólicum in Gállia, diadémate distínctam vóluit. Leo vero décimus tértius innúmera étiam cóntulit benefícia, indulgéntias ad modum jubilǽi vigésimo quinto Apparitiónis anno verténte concéssit, peregrinatiónes sua auctoritáte verbóque provéxit, ac solémnem Ecclésiæ sub título Rosárii dedicatiónem suo nómine péragi curávit. Quorum beneficiórum amplitúdinem cumulávit, cum, plúrium episcopórum rogátu, solémne festum sub título Apparitiónis beátæ Maríæ Vírginis immaculátæ, próprio Offício et própria Missa celebrándum benígne concéssit. Tandem Pius décimus Póntifex máximus pro sua erga Deíparam pietáte, ac plurimórum votis ánnuens Sacrórum antístitum, idem festum ad Ecclésiam univérsam exténdit.\par
\switchcolumn
\EN
\noindent It is known by all for certain the pilgrimages of this sort to have stirred up faith, in a world becoming cold, the soul to have increased to profess the Christian law, and to have spread the worship of the immaculate Virgin in a wonderful manner. In which wonderful profession of faith the Christian people have the priests as leaders, who lead their people thence. The holy Bishops themselves also frequently visit the holy place, lead pilgrimages, and take part in more solemn feasts. Nor indeed is it rare to observe these empurpled fathers of the Roman Church approaching humble in the manner of pilgrims. Even these Roman Pontiffs, for their piety towards the Immaculate of Lourdes, enhanced the sacred shrine with most noble gifts. Pius the ninth, with sacred indulgences, distinguished it with the privilege of an archconfraternity and title of a minor Basilica; and willed the image honoured there of the Godbearer, with solemn rite through his apostolic legate in France, to be adorned with a crown. And Leo the thirteenth also conferred innumerable benefits, granted indulgences in the manner of a jubilee in the twenty-fifth year of the Apparition occurring, promoted pilgrimages by his authority and word, and arranged to complete by his name the solemn dedication of the Church under the title of the Rosary. He accumulated the extent of which benefits, when, with the asking of many bishops, he kindly granted a solemn feast to be celebrated under the title of the Apparition of the blessed Virgin Mary immaculate with a proper Office and proper Mass. Finally supreme Pontiff Pius the tenth for his piety towards the Godbearer, and agreeing to the wishes of many holy Bishops, extended the same feast to the universal Church.\par
\switchcolumn*

\LA
\begin{Resp}\Rsign Prævenísti eam, Dómine, in benedictiónibus dulcédinis, posuísti in cápite ejus \Star{} Corónam de lápide pretióso.\end{Resp}
\begin{Resp}\Vsign Magna est glória ejus in salutári tuo, glóriam et magnum decórem impónes super eam.\end{Resp}
\begin{Resp}\Rsign Corónam de lápide pretióso.\end{Resp}
\begin{Resp}\Vsign Glória Patri, et Fílio, \Star{} et Spirítui Sancto.\end{Resp}
\begin{Resp}\Rsign Corónam de lápide pretióso.\end{Resp}
\switchcolumn
\EN
\begin{Resp}\Rsign Thou hast prevented her, O Lord, with blessings of sweetness, thou hast set on her head \Star{} A crown of precious stones.\end{Resp}
\begin{Resp}\Vsign Her glory is great in thy salvation, glory and great beauty shalt thou lay upon her.\end{Resp}
\begin{Resp}\Rsign A crown of precious stones.\end{Resp}
\begin{Resp}\Vsign Glory be to the Father, and to the Son, \Star{} and to the Holy Ghost.\end{Resp}
\begin{Resp}\Rsign A crown of precious stones.\end{Resp}
\switchcolumn*

\LA
\LessonHead{Lectio VII}
\switchcolumn
\EN
\LessonHead{Lesson VII}
\switchcolumn*

\LA
\LessonTitle{Léctio sancti Evangélii secúndum Lucam}
\switchcolumn
\EN
\LessonTitle{From the Holy Gospel according to Luke}
\switchcolumn*

\LA
\Cite{Luc 1:26-28}
\switchcolumn
\EN
\Cite{Luke 1:26-28}
\switchcolumn*

\LA
\noindent In illo témpore: Missus est Angelus Gábriel a Deo in civitátem Galilææ, cui nomen Názareth, ad Vírginem desponsátam viro, cui nomen erat Joseph, de domo David, et nomen Vírginis María. Et réliqua.\par
\switchcolumn
\EN
\noindent In that time the angel Gabriel was sent from God into a city of Galilee, called Nazareth, To a virgin espoused to a man whose name was Joseph, of the house of David; and the virgin's name was Mary. And so on.\par
\switchcolumn*

\LA
\LessonTitle{Homilía sancti Bernárdi Abbátis}
\switchcolumn
\EN
\LessonTitle{Homily by St. Bernard, Abbot (of Clairvaux.)}
\switchcolumn*

\LA
\Source{Ex Homilía 2 super Missus}
\switchcolumn
\EN
\Source{2nd on this text.}
\switchcolumn*

\LA
\noindent Lætáre, pater Adam, sed magis tu, o Heva mater, exsúlta, qui sicut ómnium paréntes, ita ómnium fuístis peremptóres; et quod infelícius est, prius peremptóres quam paréntes. Ambo, inquam, consolámini super fília, et tali fília, sed illa ámplius de qua malum ortum est prius, cujus oppróbrium in omnes pertransívit mulíeres. Instat namque tempus, quo jam tollátur oppróbrium, nec hábeat vir quid causétur advérsus féminam: qui útique, dum se imprudénter excusáre conarétur, crudéliter illam accusáre non cunctátus est, dicens: Múlier quam dedísti mihi, dedit mihi de ligno, et comédi. Proptérea curre Heva ad Maríam; curre mater ad fíliam; fília pro matre respóndeat: ipsa matris oppróbrium áuferat; ipsa patri pro matre satisfáciat: quia ecce si vir cécidit per féminam, jam non erígitur nisi per féminam.\par
\switchcolumn
\EN
\noindent Rejoice, father Adam, and yet more thou mother Eve, ye that are the source of all, and the ruin of all, and the unhappy cause of their ruin before ye gave them birth. Be comforted both in your daughter, and such a daughter; but chiefly thou, O woman, of whom the first evil came, and who hast cast thy slur upon all women. The time is come for the slur to be taken away, and for the man to have nothing to say against the woman. At the first, when he unwisely began to make excuse, he scrupled not to throw the blame upon her, saying, The woman whom Thou gavest to be with me, she gave me of the tree, and I did eat. Wherefore, O Eve, betake thyself to Mary Mother, betake thyself to thy daughter let the daughter answer for the mother let her take away her mother's reproach; let her make up to her father for her mother's fault for if man be fallen by means of woman, it is by means of woman that he is raised up again.\par
\switchcolumn*

\LA
\begin{Resp}\Rsign Tu ergo ínvoca Dóminum, lóquere Regi pro nobis, \Star{} Et líbera nos de morte.\end{Resp}
\begin{Resp}\Vsign Omnes sitiéntes, veníte ad aquas, et hauriétis salútem a Dómino.\end{Resp}
\begin{Resp}\Rsign Et líbera nos de morte.\end{Resp}
\switchcolumn
\EN
\begin{Resp}\Rsign Therefore do thou call upon the Lord, speak to the King for us, \Star{} And deliver us from death.\end{Resp}
\begin{Resp}\Vsign All you that thirst, come to the waters, and draw salvation from the Lord.\end{Resp}
\begin{Resp}\Rsign And deliver us from death.\end{Resp}
\switchcolumn*

\LA
\LessonHead{Lectio VIII}
\switchcolumn
\EN
\LessonHead{Lesson VIII}
\switchcolumn*

\LA
\noindent Quid dicébas, o Adam? Múlier quam dedísti mihi, dedit mihi de ligno, et comédi. Verba malítiæ sunt hæc quibus magis áugeas quam déleas culpam. Verúmtamen Sapiéntia vicit malítiam, cum occasiónem véniæ, quam a te Deus interrogándo elícere tentávit, sed non pótuit, in thesáuro indeficiéntis suæ pietátis invénit. Rédditur nempe fémina pro fémina, prudens pro fátua, húmilis pro supérba, quæ pro ligno mortis gustum tibi pórrigat vitæ, et pro venenóso cibo illo amaritúdinis, dulcédinem páriat fructus ætérni. Muta ergo iníquæ excusatiónis verbum in vocem gratiárum actiónis, et dic: Dómine, múlier, quam dedísti mihi, dedit mihi de ligno vitæ, et comédi; et dulce factum est super mel ori meo, quia in ipso vivificásti me. Ecce enim ad hoc missus est Angelus ad Vírginem. O admirándam et omni honóre digníssimam Vírginem; o féminam singuláriter venerándam, super omnes féminas admirábilem, paréntum reparatrícem, posterórum vivificatrícem.\par
\switchcolumn
\EN
\noindent What didst thou say, O Adam? The woman whom Thou gavest to be with me, she gave me of the tree, and I did eat. These are wrathful words, by the which thou dost rather magnify than diminish thine offence. Nevertheless, Wisdom hath defeated thy malice. God asked thee that He might find in thee an occasion of pardon, but, in that He found it not, He hath sought and found it in the Treasure of His Own mercy. One woman answereth for another; the wise for the foolish; the lowly for the proud; for her that gave thee of the tree of death, another that giveth thee to taste of the tree of life; for her that brought thee the bitter food of sin, another that giveth thee of the sweet fruits of righteousness. Wherefore accuse the woman no more, but speak in thanksgiving, and say, Lord, the woman whom Thou hast given me, she hath given me of the tree of life, and I have eaten; and it is in my mouth sweeter than honey, for thereby hast Thou quickened me. (Ps. cxviii. 103, 93.) Behold, it was for this that the angel Gabriel was sent to the Virgin, to the most worshipful of women, a woman more wonderful than all women, the restorer of them that went before, and the quickener of them that come after her.\par
\switchcolumn*

\LA
\begin{Resp}\Rsign Plantávit Dóminus Deus paradísum voluptátis, produxítque lignum vitæ in médio ejus: \Star{} Et flúvius egrediebátur de loco voluptátis.\end{Resp}
\begin{Resp}\Vsign Emissiónes tuæ paradisus, Virgo María.\end{Resp}
\begin{Resp}\Rsign Et flúvius egrediebátur de loco voluptátis.\end{Resp}
\begin{Resp}\Vsign Glória Patri, et Fílio, \Star{} et Spirítui Sancto.\end{Resp}
\begin{Resp}\Rsign Et flúvius egrediebátur de loco voluptátis.\end{Resp}
\switchcolumn
\EN
\begin{Resp}\Rsign The Lord God planted a paradise of pleasure, and brought forth the tree of life in its midst: \Star{} And a river went out of the place of pleasure.\end{Resp}
\begin{Resp}\Vsign Thy plants are a paradise, O Virgin Mary.\end{Resp}
\begin{Resp}\Rsign And a river went out of the place of pleasure.\end{Resp}
\begin{Resp}\Vsign Glory be to the Father, and to the Son, \Star{} and to the Holy Ghost.\end{Resp}
\begin{Resp}\Rsign And a river went out of the place of pleasure.\end{Resp}
\switchcolumn*

\LA
\LessonHead{Lectio IX}
\switchcolumn
\EN
\LessonHead{Lesson IX}
\switchcolumn*

\LA
\noindent Quam tibi áliam prædixísse Deus vidétur, quando ad serpéntem ait: Inimicítias ponam inter te et mulíerem? Et si adhuc dúbitas, quod de María díxerit, audi quod séquitur: Ipsa cónteret caput tuum. Cui hæc serváta victória est, nisi Maríæ? Ipsa procul dúbio caput contrívit venenátum, quæ omnimódam malígni suggestiónem tam de carnis illécebra, quam de mentis supérbia dedúxit ad níhilum. Quam vero áliam Sálomon requirébat, cum dicébat: Mulíerem fortem quis invéniet? Nóverat quippe vir sápiens hujus sexus infirmitátem, frágile corpus, lúbricam mentem. Quia tamen et Deum légerat promisísse, et ita vidébat congrúere, ut qui vícerat per féminam, vincerétur per ipsam, veheménter admírans ajébat: Mulíerem fortem quis invéniet? Quod est dícere: si ita de manu féminæ pendet et nostra ómnium salus, et innocéntiæ restitútio, et de hoste victória: fortis omníno necésse est ut provideátur, quæ ad tantum opus possit esse idónea.\par
\switchcolumn
\EN
\noindent Has it not of this thy daughter, O Adam, that God spake when He said unto the serpent, I will put enmity between thee and the woman And if thou wilt still doubt that He speaketh of Mary, hear what followeth She shall bruise thy head. Who won this conquest but Mary? She brought to nought the whole wiles of Satan, whether for the pollution of her body or the injury of her soul. Was it not of her that Solomon spake, where he saith, Who shall find a virtuous woman? (Prov. xxxi. 10.) The wise man knew the weaknesses of women, how frail they are in body, and how changeable in mind. But he had read that God had promised that the enemy, who had prevailed by means of a woman, was by a woman to be overthrown, and he believed. But he wondered greatly, and said, Who shall find a virtuous woman? that is to say If our salvation, and the bringing back of that which is lost, and the final triumph over the enemy, is in the hand of a woman, it must needs be that a virtuous woman be found, meet to work in that matter\par
\switchcolumn*

\LA
\TeDeum{Te Deum}
\switchcolumn
\EN
\TeDeum{Te Deum}
\switchcolumn*

\end{paracol}

\par\needspace{10\baselineskip}\addvspace{1.2em}{\centering\small\textsc{Die 12 Februarii} · \textsc{12 February}\par}
\Office{Ss. Septem Fundatorum Ordinis Servorum B. M. V.}{Seven Holy Founders of the Order of the Servants of the Blessed Virgin Mary}{Duplex}
\addcontentsline{toc}{subsection}{Die 12 Februarii · Ss. Septem Fundatorum Ordinis Servorum B. M. V.\enspace\textit{Seven Holy Founders of the Order of the Servants of the Blessed Virgin Mary}}
\begin{paracol}{2}
\LA
\RefLine{Lectiones I–III: de Scriptura occurrente.}
\switchcolumn
\EN
\RefLine{Lessons I–III: from the Scripture of the day.}
\switchcolumn*

\LA
\LessonHead{Lectio IV}
\switchcolumn
\EN
\LessonHead{Lesson IV}
\switchcolumn*

\LA
\noindent Sǽculo tértio décimo cum Frideríci secúndi diro schísmate cruentísque factiónibus cultióres Itáliæ pópuli scinderéntur, próvidens Dei misericórdia, præter álios sanctitáte illústres, septem e Florentína nobilitáte viros suscitávit, qui in caritáte conjúncti, præclárum fratérnæ dilectiónis præbérent exémplum. Hi, nimírum Bonfílius Monáldius, Bonajúncta Manéttus, Manéttus Antellénsis, Amidéus de Amidéis, Ugúccio Ugucciónum, Sostenéus de Sostenéis et Aléxius Falconérius, cum, anno trigésimo tértio ejus sǽculi, die sacra Vírgini cælo recéptæ, in quodam piórum hóminum convéntu, Laudántium nuncupáto, fervéntius orárent; ab eádem Deípara síngulis apparénte sunt admóniti, ut sánctius perfectiúsque vitæ genus amplecteréntur. Re ítaque prius cum Florentíno prǽsule colláta, hi septem viri, géneris nobilitáte divitiísque posthábitis, sub vilíssimis detritísque véstibus cilício indúti, octáva die Septémbris in rurálem quamdam ædículam secessére, ut ea die primórdia vitæ sanctióris auspicaréntur, qua ipsa Dei Génetrix mortálibus orta sanctíssimam vitam incéperat.\par
\switchcolumn
\EN
\noindent In the thirteenth century, when the more cultured parts of Italy were rent by the dread dissension of the Emperor Frederick the Second and by bloody civil wars, the mercy of God set forth diverse men eminent for holiness, and among others raised up seven nobles of Florence, who were bound one to another in charity and gave an illustrious example of brotherly love. Their names were Bonfiglio Monaldi, Bonajuncta Manetti, Manetto Antalli, Amadeo de' Amidei, Uguccio de' Uguccioni, Sosteneo de' Sostenei, and Alexis de' Falconieri. Upon the holiday of the Assumption of the Virgin into heaven in the year\par
1233 They were praying in the oratory of a guild called the Guild of Praise, when the same Mother of God appeared to each one of them, and bade them embrace a life of greater holiness and perfection. These seven men discussed the matter with the Bishop of Florence, and then, considering neither the nobility of their birth nor their wealth, and clad in haircloth under vile and worn - out garments, withdrew into a little house in the country upon the th day of September, that they might begin their holier life upon the same day whereon the Mother of God herself had by her birth begun her life of holiness upon earth.\par
\switchcolumn*

\LA
\begin{Resp}\Rsign Honéstum fecit illum Dóminus, et custodívit eum ab inimícis, et a seductóribus tutávit illum: \Star{} Et dedit illi claritátem ætérnam.\end{Resp}
\begin{Resp}\Vsign Justum dedúxit Dóminus per vias rectas, et osténdit illi regnum Dei.\end{Resp}
\begin{Resp}\Rsign Et dedit illi claritátem ætérnam.\end{Resp}
\switchcolumn
\EN
\begin{Resp}\Rsign The Lord made him honourable, and defended him from his enemies, and kept him safe from those that lay in wait for him; \Star{} And gave him perpetual glory.\end{Resp}
\begin{Resp}\Vsign He went down with him into the pit, and left him not in bonds.\end{Resp}
\begin{Resp}\Rsign And gave him perpetual glory.\end{Resp}
\switchcolumn*

\LA
\LessonHead{Lectio V}
\switchcolumn
\EN
\LessonHead{Lesson V}
\switchcolumn*

\LA
\noindent Hoc vitæ institútum quam sibi foret accéptum Deus miráculo osténdit. Nam cum paulo deínceps hi septem viri per Florentínam urbem ostiátim eleemósynam emendicárent, áccidit, ut repénte infántium voce, quos inter fuit sanctus Philíppus Benítius quintum ætátis mensem vix ingréssus, beátæ Maríæ Servi acclamaréntur: quo deínde nómine semper appelláti sunt. Quare, vitándi pópuli occúrsus ac solitúdinis amóre ducti, in Senárii montis recéssu omnes convenére, ibíque cæléste quoddam vitæ genus aggréssi sunt. Victitábant enim in spelúncis, sola aqua herbísque conténti; vigíliis aliísque asperitátibus corpus atterébant, Christi passiónem ac mæstíssimæ ejúsdem Genitrícis dolóres assídue meditántes. Quod cum olim sacra Parascéves die impénsius exsequeréntur, ipsa beáta Virgo illis iteráto appárens, lúgubrem vestem, quam indúerent, osténdit; sibíque acceptíssimum fore significávit, ut novum in Ecclésia regulárem órdinem excitárent, qui jugem recóleret ac promovéret memóriam dolórum, quos ipsa pértulit sub cruce Dómini. Hæc sanctus Petrus, ínclytus órdinis Prædicatórum Martyr, ex familiári cum sanctis illis viris consuetúdine ac peculiári étiam Deíparæ visióne cum didicísset; iis auctor fuit, ut órdinem regulárem sub appellatióne Servórum beátæ Vírginis institúerent: qui póstea ab Innocéntio quarto Pontífice máximo approbátus fuit.\par
\switchcolumn
\EN
\noindent She showed by a miracle how acceptable in His sight should be their manner of life, for a short while after, when these seven men were begging alms from door to door through the city of Florence, it came to pass that some children, among whom was holy Philip Benizi, who had then scarcely entered the fifth month of his age, called them blessed Mary's servants, by the which name they were called ever after. To avoid meeting people, and in the desire to be alone, they all withdrew together to the solitude of Monte Senario, and there began a kind of heavenly life. They lived in caves and upon herbs and water only, while they wore out their bodies with watching and other hardships, while they contemplated unweariedly the sufferings of Christ and the woes of His most sorrowful Mother. One Good Friday, when their thoughts were fixed thereon more than ever, the Blessed Virgin appeared to them twice, and showed them her garments of mourning as those wherein they should clothe themselves. She bade them know that she would take it right well that they should raise up in the Church a new order to recall the memory of the sorrows which she bore beneath the Cross of the Lord. Holy Peter, the illustrious martyr of the Order of Friars Preachers, learnt this not only from his familiar converse with these holy men, but also from a special vision of the Mother of God, and it was on his incitement that they founded the regular Order called that of the Servites, or servants of the Blessed Virgin, the which Order was afterward approved by the Supreme Pontiff Innocent IV.\par
\switchcolumn*

\LA
\begin{Resp}\Rsign Amávit eum Dóminus, et ornávit eum: stolam glóriæ índuit eum, \Star{} Et ad portas paradísi coronávit eum.\end{Resp}
\begin{Resp}\Vsign Índuit eum Dóminus lorícam fídei, et ornávit eum.\end{Resp}
\begin{Resp}\Rsign Et ad portas paradísi coronávit eum.\end{Resp}
\switchcolumn
\EN
\begin{Resp}\Rsign The Lord loved him and beautified him; He clothed him with a robe of glory, \Star{} And crowned him at the gates of Paradise.\end{Resp}
\begin{Resp}\Vsign The Lord hath put on him the breast-plate of faith, and hath adorned him.\end{Resp}
\begin{Resp}\Rsign And crowned him at the gates of Paradise.\end{Resp}
\switchcolumn*

\LA
\LessonHead{Lectio VI}
\switchcolumn
\EN
\LessonHead{Lesson VI}
\switchcolumn*

\LA
\noindent Porro sancti illi viri, cum plures sibi sócios adjunxíssent, Itáliæ civitátes atque óppida, præsértim Etrúriæ excúrrere cœpérunt, prædicántes ubíque Christum crucifíxum, civíles discórdias compescéntes et innúmeros fere dévios ad virtútis sémitam revocántes. Neque Itáliam modo, sed et Gálliam, Germániam ac Polóniam suis evangélicis labóribus excoluérunt. Dénique cum bonum Christi odórem longe latéque diffudíssent, portentórum quoque glória illústres, migrárunt ad Dóminum. Sed quos unus veræ fraternitátis ac religiónis amor in vita sociáverat, unum páriter demórtuos contéxit sepúlcrum, únaque pópuli venerátio prosecúta est. Quaprópter Clemens undécimus et Benedíctus décimus tértius Pontífices máximi delátum iísdem a plúribus sǽculis indivíduum cultum confírmarunt; ac Leo décimus tértius, approbátis ántea miráculis, post indúltam veneratiónem ad collectívam eorúmdem invocatiónem a Deo patrátis, eósdem anno quinquagésimo sacerdótii sui Sanctórum honóribus cumulávit, eorúmque memóriam Offício et Missa in univérsa Ecclésia quotánnis recoléndam instítuit.\par
\switchcolumn
\EN
\noindent These holy men, when they had gathered to themselves some companions, began to go through the cities and towns of Italy, and especially of Tuscany, everywhere preaching Christ crucified, stilling contests among the citizens, and calling back almost countless backsliders into the path of grace. Neither did they make Italy only the field of their Gospel labours, but also France, Germany, and Poland. They passed away to be ever with the Lord when they had spread far and wide a sweet savour of Christ, and were famous also for the glory of signs and wonders. As one love of brotherhood and of the monastic life had joined them together upon earth, so one grave held their dead bodies, and one honour was paid them by the people. For this reason the Supreme Pontiffs Clement XI. and Benedict XIII. confirmed the honour whicfl had for centuries been paid to them individually, and Leo XIII., after proof of their miracles which had been wrought by God on the common invocation of these saints, after their veneration had been sanctioned in the jubilee year of his priesthood, decreed to them the honours paid to Saints, and ordered that their memory should every year be kept throughout the universal Church with an office and Mass.\par
\switchcolumn*

\LA
\begin{Resp}\Rsign Iste homo perfécit ómnia quæ locútus est ei Deus, et dixit ad eum: Ingrédere in réquiem meam: \Star{} Quia te vidi justum coram me ex ómnibus géntibus.\end{Resp}
\begin{Resp}\Vsign Iste est, qui contémpsit vitam mundi, et pervénit ad cæléstia regna.\end{Resp}
\begin{Resp}\Rsign Quia te vidi justum coram me ex ómnibus géntibus.\end{Resp}
\begin{Resp}\Vsign Glória Patri, et Fílio, \Star{} et Spirítui Sancto.\end{Resp}
\begin{Resp}\Rsign Quia te vidi justum coram me ex ómnibus géntibus.\end{Resp}
\switchcolumn
\EN
\begin{Resp}\Rsign This is he which did according unto all that God commanded him; and God said unto him: Enter thou into My rest, \Star{} For thee have I seen righteous before Me among all people.\end{Resp}
\begin{Resp}\Vsign This is he which loved not his life in this world, and is come unto an everlasting kingdom.\end{Resp}
\begin{Resp}\Rsign For thee have I seen righteous before Me among all people.\end{Resp}
\begin{Resp}\Vsign Glory be to the Father, and to the Son, \Star{} and to the Holy Ghost.\end{Resp}
\begin{Resp}\Rsign For thee have I seen righteous before Me among all people.\end{Resp}
\switchcolumn*

\LA
\LessonHead{Lectio VII}
\switchcolumn
\EN
\LessonHead{Lesson VII}
\switchcolumn*

\LA
\LessonTitle{Léctio sancti Evangélii secúndum Matthǽum.}
\switchcolumn
\EN
\LessonTitle{From the Holy Gospel according to Matthew}
\switchcolumn*

\LA
\Cite{Matt 19:27-29}
\switchcolumn
\EN
\Cite{Matt 19:27-29}
\switchcolumn*

\LA
\noindent In illo témpore: Dixit Petrus ad Jesum: Ecce nos relíquimus ómnia, et secúti sumus te: quid ergo erit nobis? Et réliqua.\par
\switchcolumn
\EN
\noindent At that time, Peter said unto Jesus: Behold, we have forsaken all, and followed thee what shall we have therefore? And so on.\par
\switchcolumn*

\LA
\LessonTitle{Homilía sancti Hierónymi Presbýteri.}
\switchcolumn
\EN
\LessonTitle{Homily by St. Jerome, Priest (at Bethlehem.)}
\switchcolumn*

\LA
\Source{Lib. 3. in Matth. cap. 19.}
\switchcolumn
\EN
\Source{Bk. iii. on Matth xix.}
\switchcolumn*

\LA
\noindent Grandis fidúcia! Petrus piscátor erat, dives non fúerat, cibos manu et arte quærébat: et tamen lóquitur confidénter: Relíquimus ómnia. Et quia non súfficit tantum relínquere, jungit quod perféctum est: Et secúti sumus te. Fécimus quod jussísti: quid ígitur nobis dabis prǽmii? Jesus autem dixit illis: Amen dico vobis, quod vos qui secúti estis me, in regeneratióne, cum séderit Fílius hóminis in sede majestátis suæ, sedébitis et vos super sedes duódecim, judicántes duódecim tribus Ísraël. Non dixit: Qui reliquístis ómnia; hoc enim et Crates fecit philósophus, et multi álii divítias contempsérunt: sed, Qui secúti estis me: quod próprie Apostolórum est, atque credéntium.\par
\switchcolumn
\EN
\noindent Peter was a fisherman, he was not rich, he earned his bread by his hand and skill, and nevertheless he is thus bold, and saith confidently: We have forsaken all. And because it sufficeth not to forsake only, he addeth that which to do is to be perfect: and followed thee. We have done that which Thou hast commanded us, what reward therefore wilt Thou give us? And Jesus said unto them Amen I say unto you, that ye which have followed Me, in the regeneration, when the Son of Man shall sit in the throne of His glory, ye also shall sit upon twelve thrones, judging the twelve tribes of Israel. He said not, Ye which have forsaken all, for this did even Crates the philosopher, and they which have set nothing by riches are many, but, Ye which have followed Me. This did the Apostles, and this do believers do.\par
\switchcolumn*

\LA
\begin{Resp}\Rsign Iste est qui ante Deum magnas virtútes operátus est, et de omni corde suo laudávit Dóminum: \Star{} Ipse intercédat pro peccátis ómnium populórum.\end{Resp}
\begin{Resp}\Vsign Ecce homo sine queréla, verus Dei cultor, ábstinens se ab omni ópere malo, et pérmanens in innocéntia sua.\end{Resp}
\begin{Resp}\Rsign Ipse intercédat pro peccátis ómnium populórum.\end{Resp}
\switchcolumn
\EN
\begin{Resp}\Rsign This is he which wrought great wonders before God, and praised the Lord with all his heart. \Star{} May he pray for all people, that their sins may be forgiven unto them.\end{Resp}
\begin{Resp}\Vsign Behold a man without blame, a worshipper of God in truth, keeping himself clean from every evil work, and abiding still in his innocency.\end{Resp}
\begin{Resp}\Rsign May he pray for all people, that their sins may be forgiven unto them.\end{Resp}
\switchcolumn*

\LA
\LessonHead{Lectio VIII}
\switchcolumn
\EN
\LessonHead{Lesson VIII}
\switchcolumn*

\LA
\noindent In regeneratióne, cum séderit Fílius hóminis in sede majestátis suæ (quando et mórtui de corruptióne resúrgent incorrúpti) sedébitis et vos in sóliis judicántium, condemnántes duódecim tribus Ísraël: quia vobis credéntibus, illi crédere noluérunt. Et omnis qui relíquerit domum, vel fratres, aut soróres, aut patrem, aut matrem, aut uxórem, aut fílios, aut agros propter nomen meum, céntuplum accípiet, et vitam ætérnam possidébit. Locus iste cum illa senténtia cóngruit, in qua Salvátor lóquitur: Non veni pacem míttere, sed gládium. Veni enim separáre hóminem a patre suo, et matrem a fília, et nurum a socru: et inimíci hóminis doméstici ejus. Qui ergo propter fidem Christi, et prædicatiónem Evangélii, omnes afféctus contémpserint, atque divítias et sǽculi voluptátes: isti céntuplum recípient, et vitam ætérnam possidébunt.\par
\switchcolumn
\EN
\noindent In the regeneration, when the Son of Man shall sit in the throne of His glory, and when the dead shall rise again from corruption incorruptible, (i Cor. xv. 53,) ye also shall sit upon twelve thrones of judgment, condemning the twelve tribes of Israel, because, when ye believed in Me, they would not. (John iii. 18.) And every one that hath forsaken houses, or brethren, or sisters, or father, or mother, or wife, or children, or lands, for My Name's sake, shall receive an hundredfold, and shall inherit everlasting life. This place agreeth well with that other where the Saviour saith I came not to send peace, but a sword. For I am come to set a man at variance against his father, and the daughter against her mother, and the daughter-in-law against her mother-in-law; and a man's foes shall be they of his own household. (Matth. x. 34.) Every one, therefore, that hath set no store by affection, and riches, and the pleasures of the world, for Christ's faith's sake, and the preaching of the Gospel, shall receive an hundred-fold, and shall inherit everlasting life.\par
\switchcolumn*

\LA
\begin{Resp}\Rsign Sint lumbi vestri præcíncti, et lucérnæ ardéntes in mánibus vestris: \Star{} Et vos símiles homínibus exspectántibus dóminum suum, quando revertátur a núptiis.\end{Resp}
\begin{Resp}\Vsign Vigiláte ergo, quia nescítis qua hora Dóminus vester ventúrus sit.\end{Resp}
\begin{Resp}\Rsign Et vos símiles homínibus exspectántibus dóminum suum, quando revertátur a núptiis.\end{Resp}
\begin{Resp}\Vsign Glória Patri, et Fílio, \Star{} et Spirítui Sancto.\end{Resp}
\begin{Resp}\Rsign Et vos símiles homínibus exspectántibus dóminum suum, quando revertátur a núptiis.\end{Resp}
\switchcolumn
\EN
\begin{Resp}\Rsign Let your loins be girded about, and your lights burning; \Star{} And ye yourselves like unto men that wait for their lord, when he will return from the wedding.\end{Resp}
\begin{Resp}\Vsign Watch therefore, for ye know not what hour your Lord doth come.\end{Resp}
\begin{Resp}\Rsign And ye yourselves like unto men that wait for their lord, when he will return from the wedding.\end{Resp}
\begin{Resp}\Vsign Glory be to the Father, and to the Son, \Star{} and to the Holy Ghost.\end{Resp}
\begin{Resp}\Rsign And ye yourselves like unto men that wait for their lord, when he will return from the wedding.\end{Resp}
\switchcolumn*

\LA
\LessonHead{Lectio IX}
\switchcolumn
\EN
\LessonHead{Lesson IX}
\switchcolumn*

\LA
\noindent Ex occasióne hujus senténtiæ quidam introdúcunt mille annos post resurrectiónem, dicéntes, tunc nobis céntuplum ómnium rerum quas dimísimus, et vitam ætérnam esse reddéndam: non intelligéntes, quod, si in céteris digna sit repromíssio, in uxóribus appáreat turpitúdo: ut qui unam pro Dómino dimíserit centum recípiat in futúro. Sensus ergo iste est: Qui carnália pro Salvatóre dimíserit, spirituália recípiet: quæ comparatióne et mérito sui ita erunt, quasi si parvo número centenárius númerus comparétur.\par
\switchcolumn
\EN
\noindent By reason of these words, an hundredfold, some will have that there shall be a thousand years after the resurrection, wherein they that have forsaken all things shall receive an hundredfold of those things which they have forsaken, and shall inherit everlasting life. Such men consider not that though in other things this were worthy, as touching wives it is unseemly for it becometh us not to think that he that hath forsaken one wife in this world, shall receive an hundred wives in that which is to come. But the meaning is this, that every one that for the Saviour's sake hath forsaken earthly things, shall receive spiritual things which things, being rightly weighed against earthly things, are as though an hundredfold were weighed against one.\par
\switchcolumn*

\LA
\TeDeum{Te Deum}
\switchcolumn
\EN
\TeDeum{Te Deum}
\switchcolumn*

\end{paracol}

\par\needspace{10\baselineskip}\addvspace{1.2em}{\centering\small\textsc{Die 14 Februarii} · \textsc{14 February}\par}
\Office{S. Valentini Presbyteri et Martyris}{St. Valentine, Priest and Martyr}{Simplex}
\addcontentsline{toc}{subsection}{Die 14 Februarii · S. Valentini Presbyteri et Martyris\enspace\textit{St. Valentine, Priest and Martyr}}
\begin{paracol}{2}
\LA
\RefLine{Lectiones I–II: de Scriptura occurrente.}
\switchcolumn
\EN
\RefLine{Lessons I–II: from the Scripture of the day.}
\switchcolumn*

\LA
\LessonHead{Lectio III (pro commemoratione)}
\switchcolumn
\EN
\LessonHead{Lesson III (when commemorated)}
\switchcolumn*

\LA
\LessonTitle{Sermo sancti Augustíni Epíscopi}
\switchcolumn
\EN
\LessonTitle{From the Sermons of St Augustine, Bishop (of Hippo.)}
\switchcolumn*

\LA
\Source{Sermo 44 de Sanctis}
\switchcolumn
\EN
\Source{44 on the Saints.}
\switchcolumn*

\LA
\noindent Triumphális beáti Mártyris Valentini dies hódie nobis anniversária celebritáte recúrrit: cujus glorificatióni sicut congáudet Ecclésia, sic ejus propónit sequénda vestígia. Si enim compátimur, et conglorificábimur. In cujus glorióso agóne duo nobis præcípue consideránda sunt: induráta vidélicet tortóris sævítia, et Mártyris invícta patiéntia. Sævítia tortóris, ut eam detestémur: patiéntia Mártyris, ut eam imitémur. Audi Psalmístam advérsus malítiam increpántem: Noli æmulári in malignántibus quóniam tamquam fœnum velóciter aréscent. Quod autem advérsus malignántes patiéntia exhibénda sit, audi Apóstolum suadéntem: Patiéntia vobis necessária est, ut reportétis promissiónes.\par
\switchcolumn
\EN
\noindent The illustrious day whereon the blessed Martyr Valentine conquered, doth this day come round to us again and as the Church doth rejoice with him in his glory, so doth she set before us his footsteps to be followed. For if we suffer, we shall also reign with him. In his glorious battle we have two things chiefly to consider the hardened cruelty of the tormentor, and the unconquered patience of the Martyr the cruelty of the tormentor, that we may abhor it; the patience of the Martyr, that we may imitate it. Hear what the Psalmist saith, complaining against sin: “Fret not thyself because of the evil-doers, for they shall soon dry up like the grass.” (xxxvi. 1.) But touching the patience which is to be shown against the evil-doers, hear the word wherewith the Apostle moveth us Ye have need of patience, that ye may receive the promise. (Heb. x. 36.)\par
\switchcolumn*

\LA
\TeDeum{Te Deum}
\switchcolumn
\EN
\TeDeum{Te Deum}
\switchcolumn*

\end{paracol}

\par\needspace{10\baselineskip}\addvspace{1.2em}{\centering\small\textsc{Die 15 Februarii} · \textsc{15 February}\par}
\Office{SS. Faustini et Jovitæ}{Sts. Faustinus and Jovita, Martyrs}{Simplex}
\addcontentsline{toc}{subsection}{Die 15 Februarii · SS. Faustini et Jovitæ\enspace\textit{Sts. Faustinus and Jovita, Martyrs}}
\begin{paracol}{2}
\LA
\RefLine{Lectiones I–II: de Scriptura occurrente.}
\switchcolumn
\EN
\RefLine{Lessons I–II: from the Scripture of the day.}
\switchcolumn*

\LA
\LessonHead{Lectio III (pro commemoratione)}
\switchcolumn
\EN
\LessonHead{Lesson III (when commemorated)}
\switchcolumn*

\LA
\noindent Faustinus et Jovita fratres, nobiles Brixiani, in multis Italiæ urbibus, quo vincti sæviente Trajani persecutione ducebantur, acerbissima supplicia perpessi, fortes in christianæ fidei confessione perstiterunt. Nam Brixiæ diu vinculis constricti, feris étiam objecti in ignemque conjecti, et a bestiis et a flamma íntegri et incolumes servati sunt. Inde vero eisdem catenis colligati Mediolanum venerunt; ubi eorum fides tentata exquisitissimis tormentis, tamquam igne aurum, in cruciatibus magis enituit. Postea Romam missi, ab Evaristo Pontifice confirmati, ibi quoque crudelissime torquentur. Denique perducti Neápolim, in ea étiam urbe varie cruciati, vinctis manibus pedibusque in mare demerguntur; unde per Angelos mirabiliter erepti sunt. Quare multos et constantia in tormentis et miraculorum virtute ad Christi fidem converterunt. Postrémo reducti Brixiam initio suscepti ab Hadriano imperii, securi percussi, illustrem martyrii coronam acceperunt.\par
\switchcolumn
\EN
\noindent Faustinus and Jovita were brothers, born of a noble family at Brescia. While Trajan's persecution was raging, they were taken about in chains from one city of Italy to another, and exhibited in torture in each. This cruelty utterly failed to silence their confession of Christ, Whom they preached by their sufferings in every place where they were shown. They were afterwards kept for a long time at Brescia, where they were exhibited with wild beasts, and tormented with fire. Being both still alive, they were brought to Milan, without their chains having ever been taken off. At Milan they were tortured again with every invention of cruelty that could be devised. Nevertheless the great power of their faith made them more than conquerors, shining even as gold tried in the furnace. From Milan they were brought to Rome, where they were confirmed by Pope Evaristus, and where they were put to the torture again with extreme barbarity. They were afterwards shown in public at Naples, where the tormentors displayed their skill in diverse ways upon them. Here they were thrown chained into the sea, but the angels delivered them. Their stations of suffering, by their God-like patience, and the wonderful Power displayed in them, had now turned many souls to Jesus. In the end they were carried back to Brescia, and, when Hadrian took the empire, they were put to death by the axe at that place. The crown of martyrdom which they won is glorious.\par
\switchcolumn*

\LA
\TeDeum{Te Deum}
\switchcolumn
\EN
\TeDeum{Te Deum}
\switchcolumn*

\end{paracol}

\par\needspace{10\baselineskip}\addvspace{1.2em}{\centering\small\textsc{Die 18 Februarii} · \textsc{18 February}\par}
\Office{S. Simeonis Episcopi et Martyris}{St. Simeon, Bishop and Martyr}{Simplex}
\addcontentsline{toc}{subsection}{Die 18 Februarii · S. Simeonis Episcopi et Martyris\enspace\textit{St. Simeon, Bishop and Martyr}}
\begin{paracol}{2}
\LA
\RefLine{Lectiones I–II: de Scriptura occurrente.}
\switchcolumn
\EN
\RefLine{Lessons I–II: from the Scripture of the day.}
\switchcolumn*

\LA
\LessonHead{Lectio III (pro commemoratione)}
\switchcolumn
\EN
\LessonHead{Lesson III (when commemorated)}
\switchcolumn*

\LA
\noindent Simeon, fílius Cléophæ, post Jacobum próximus Jerosolymis ordinátus epíscopus, Trajano imperatóre apud Atticum consularem est accusatus, quod christianus esset et Christi propinquus. Comprehendebántur enim omnes eo témpore, quicúmque ex genere David orti essent. Quare multis cruciátus tormentis eodem passiónis genere, quod Salvátor noster subíerat, affícitur: mirántibus ómnibus, quod homo ætate confectus (erat enim centum et viginti annórum) acerbíssimos crucis dolóres fortiter constanterque paterétur.\par
\switchcolumn
\EN
\noindent Simeon, the son of Cleophas, (Matth. xiii. 55,) was chosen the second Bishop of Jerusalem, (in the year 62,) being the first after James. Under the Emperor Trajan he was accused before the Pro-Consul Atticus, as being both a Christian and a relation of Christ, this being the time when all were arrested that were of the lineage of David. He underwent with great suffering the same things that were inflicted on our Saviour, and all men marvelled to see with how great boldness and firmness he endured the grievous torment of the cross, at his great age, for he was an hundred and twenty years old. (a.d. 107 or 116.)\par
\switchcolumn*

\LA
\TeDeum{Te Deum}
\switchcolumn
\EN
\TeDeum{Te Deum}
\switchcolumn*

\end{paracol}

\par\needspace{10\baselineskip}\addvspace{1.2em}{\centering\small\textsc{Die 22 Februarii} · \textsc{22 February}\par}
\Office{In Cathedra S. Petri Apostoli Antiochiæ}{Chair of St. Peter at Antioch}{Duplex majus}
\addcontentsline{toc}{subsection}{Die 22 Februarii · In Cathedra S. Petri Apostoli Antiochiæ\enspace\textit{Chair of St. Peter at Antioch}}
\begin{paracol}{2}
\LA
\LessonHead{Lectio I}
\switchcolumn
\EN
\LessonHead{Lesson I}
\switchcolumn*

\LA
\LessonTitle{Incipit Epístola prima beáti Petri Apóstoli}
\switchcolumn
\EN
\LessonTitle{Lesson from the first letter of St. Peter the Apostle}
\switchcolumn*

\LA
\Cite{1 Pet 1:1-5}
\switchcolumn
\EN
\Cite{1 Pet 1:1-5}
\switchcolumn*

\LA
\noindent Petrus Apóstolus Jesu Christi, eléctis ádvenis dispersiónis Ponti, Galátiæ, Cappadóciæ, Asiæ, et Bithýniæ, secúndum præsciéntiam Dei Patris, in sanctificatiónem Spíritus, in obediéntiam, et aspersiónem sánguinis Jesu Christi: Grátia vobis, et pax multiplicétur. Benedíctus Deus et Pater Dómini nostri Jesu Christi, qui secúndum misericórdiam suam magnam regenerávit nos in spem vivam, per resurrectiónem Jesu Christi ex mórtuis, in hereditátem incorruptíbilem, et incontaminátam, et immarcescíbilem, conservátam in cælis in vobis, qui in virtúte Dei custodímini per fidem in salútem, parátam revelári in témpore novíssimo.\par
\switchcolumn
\EN
\noindent Peter, an apostle of Jesus Christ, to the strangers dispersed through Pontus, Galatia, Cappadocia, Asia, and Bithynia, elect, According to the foreknowledge of God the Father, unto the sanctification of the Spirit, unto obedience and sprinkling of the blood of Jesus Christ: Grace unto you and peace be multiplied. Blessed be the God and Father of our Lord Jesus Christ, who according to his great mercy hath regenerated us unto a lively hope, by the resurrection of Jesus Christ from the dead, Unto an inheritance incorruptible, and undefiled, and that can not fade, reserved in heaven for you, Who, by the power of God, are kept by faith unto salvation, ready to be revealed in the last time.\par
\switchcolumn*

\LA
\begin{Resp}\Rsign Simon Petre, ántequam de navi vocárem te, novi te et super plebem meam príncipem te constítui: \Star{} Et claves regni cælórum trádidi tibi.\end{Resp}
\begin{Resp}\Vsign Quodcúmque ligáveris super terram erit ligátum et in cælis; et quodcúmque sólveris super terram, erit solútum et in cælis.\end{Resp}
\begin{Resp}\Rsign Et claves regni cælórum trádidi tibi.\end{Resp}
\switchcolumn
\EN
\begin{Resp}\Rsign Simon Peter, before I called thee out of the ship, I knew thee, and appointed thee for a ruler over My people. \Star{} And I have given unto thee the keys of the kingdom of heaven.\end{Resp}
\begin{Resp}\Vsign Whatsoever, thou shalt bind on earth, shall be bound in heaven; and whatsoever thou shalt loose on earth, shall be loosed in heaven.\end{Resp}
\begin{Resp}\Rsign And I have given unto thee the keys of the kingdom of heaven.\end{Resp}
\switchcolumn*

\LA
\LessonHead{Lectio II}
\switchcolumn
\EN
\LessonHead{Lesson II}
\switchcolumn*

\LA
\Cite{1 Pet 1:6-9}
\switchcolumn
\EN
\Cite{1 Pet 1:6-9}
\switchcolumn*

\LA
\noindent In quo exsultábitis, módicum nunc si opórtet contristári in váriis tentatiónibus: ut probátio vestræ fídei multo pretiósior auro (quod per ignem probátur) inveniátur in laudem, et glóriam, et honórem, in revelatióne Jesu Christi: quem cum non vidéritis, dilígitis: in quem nunc quoque non vidéntes créditis: credéntes autem exsultábitis lætítia inenarrábili, et glorificáta: reportántes finem fídei vestræ, salútem animárum.\par
\switchcolumn
\EN
\noindent Wherein you shall greatly rejoice, if now you must be for a little time made sorrowful in diverse temptations: That the trial of your faith much more precious than gold which is tried by the fire may be found unto praise and glory and honour at the appearing of Jesus Christ: Whom having not seen, you love: in whom also now, though you see him not, you believe: and believing shall rejoice with joy unspeakable and glorified; Receiving the end of your faith, even the salvation of your souls.\par
\switchcolumn*

\LA
\begin{Resp}\Rsign Si díligis me, Simon Petre, pasce oves meas. Dómine, tu nosti quia amo te, \Star{} Et ánimam meam pono pro te.\end{Resp}
\begin{Resp}\Vsign Si oportúerit me mori tecum, non te negábo.\end{Resp}
\begin{Resp}\Rsign Et ánimam meam pono pro te.\end{Resp}
\switchcolumn
\EN
\begin{Resp}\Rsign Simon Peter, if thou lovest Me, feed My sheep. Lord, Thou knowest that I love thee; \Star{} I will lay down my life for thy sake.\end{Resp}
\begin{Resp}\Vsign If I should die with thee, I will not deny thee.\end{Resp}
\begin{Resp}\Rsign I will lay down my life for thy sake.\end{Resp}
\switchcolumn*

\LA
\LessonHead{Lectio III}
\switchcolumn
\EN
\LessonHead{Lesson III}
\switchcolumn*

\LA
\Cite{1 Pet 1:10-12}
\switchcolumn
\EN
\Cite{1 Pet 1:10-12}
\switchcolumn*

\LA
\noindent De qua salúte exquisiérunt atque scrutáti sunt prophétæ, qui de futúra in vobis grátia prophetavérunt; scrutántes in quod vel quale tempus significáret in eis Spíritus Christi: prænúntians eas quæ in Christo sunt passiónes et posterióres glórias: quibus revelátum est, quia non sibimetípsis, vobis autem ministrábant ea, quæ nunc nuntiáta sunt vobis per eos, qui evangelizavérunt vobis, Spíritu Sancto misso de cælo, in quem desíderant Angeli prospícere.\par
\switchcolumn
\EN
\noindent Of which salvation the prophets have inquired and diligently searched, who prophesied of the grace to come in you. Searching what or what manner of time the Spirit of Christ in them did signify: when it foretold those sufferings that are in Christ, and the glories that should follow: To whom it was revealed, that not to themselves, but to you they ministered those things which are now declared to you by them that have preached the gospel to you, the Holy Ghost being sent down from heaven, on whom the angels desire to look.\par
\switchcolumn*

\LA
\begin{Resp}\Rsign Tu es Petrus, et super hanc petram ædificábo Ecclésiam meam, et portæ ínferi non prævalébunt advérsus eam: \Star{} Et tibi dabo claves regni cælórum.\end{Resp}
\begin{Resp}\Vsign Quodcúmque ligáveris super terram, erit ligátum et in cælis; et quodcúmque sólveris super terram, erit solútum et in cælis.\end{Resp}
\begin{Resp}\Rsign Et tibi dabo claves regni cælórum.\end{Resp}
\begin{Resp}\Vsign Glória Patri, et Fílio, \Star{} et Spirítui Sancto.\end{Resp}
\begin{Resp}\Rsign Et tibi dabo claves regni cælórum.\end{Resp}
\switchcolumn
\EN
\begin{Resp}\Rsign Thou art Peter, and upon this rock I will build My Church, and the gates of hell shall not prevail against it. \Star{} And I will give unto thee the keys of the kingdom of heaven.\end{Resp}
\begin{Resp}\Vsign Whatsoever thou shalt bind on earth, shall be bound in heaven; and whatsoever thou shalt loose on earth, shall be loosed in heaven.\end{Resp}
\begin{Resp}\Rsign And I will give unto thee the keys of the kingdom of heaven.\end{Resp}
\begin{Resp}\Vsign Glory be to the Father, and to the Son, \Star{} and to the Holy Ghost.\end{Resp}
\begin{Resp}\Rsign And I will give unto thee the keys of the kingdom of heaven.\end{Resp}
\switchcolumn*

\LA
\LessonHead{Lectio IV}
\switchcolumn
\EN
\LessonHead{Lesson IV}
\switchcolumn*

\LA
\LessonTitle{Sermo sancti Augustíni Epíscopi}
\switchcolumn
\EN
\LessonTitle{From the Sermons of St. Augustine, Bishop (of Hippo.)}
\switchcolumn*

\LA
\Source{Sermo 15 de Sanctis}
\switchcolumn
\EN
\Source{16th on the Saints}
\switchcolumn*

\LA
\noindent Institútio solemnitátis hodiérnæ a senióribus nostris Cáthedræ nomen accépit, ídeo quod primus Apostolórum Petrus hódie episcopátus cáthedram suscepísse referátur. Recte ergo ecclésiæ natálem sedis illíus colunt, quam Apóstolus pro ecclesiárum salúte suscépit, dicénte Dómino: Tu es Petrus, et super hanc petram ædificábo Ecclésiam meam.\par
\switchcolumn
\EN
\noindent The solemn Feast of today received from our forefathers the name of that of St Peter's Chair at Antioch, because there is a tradition that it was on this day that Peter, first of the Apostles, was enthroned in a Bishop's Chair. Rightly, therefore, do the Churches observe the first day of that Chair, the right to which the Apostle received for the salvation of the Churches from the Lord of the Churches Himself, with the words: Thou art Peter, and upon this rock I will build My Church.\par
\switchcolumn*

\LA
\begin{Resp}\Rsign Tu es pastor óvium, princeps Apostolórum: tibi trádidit Deus ómnia regna mundi: \Star{} Et ídeo tibi tráditæ sunt claves regni cælórum.\end{Resp}
\begin{Resp}\Vsign Quodcúmque ligáveris super terram, erit ligátum et in cælis; et quodcúmque sólveris super terram, erit solútum et in cælis.\end{Resp}
\begin{Resp}\Rsign Et ídeo tibi tráditæ sunt claves regni cælórum.\end{Resp}
\switchcolumn
\EN
\begin{Resp}\Rsign Thou art the Shepherd of the sheep, and the Prince of the Apostles, and unto thee hath God given all the kingdoms of the world. \Star{} Therefore unto thee hath He given the keys of the kingdom of heaven.\end{Resp}
\begin{Resp}\Vsign Whatsoever thou shalt bind on earth shall be bound in heaven; and whatsoever thou shalt loose on earth shall be loosed in heaven.\end{Resp}
\begin{Resp}\Rsign Therefore unto thee hath He given the keys of the kingdom of heaven.\end{Resp}
\switchcolumn*

\LA
\LessonHead{Lectio V}
\switchcolumn
\EN
\LessonHead{Lesson V}
\switchcolumn*

\LA
\noindent Petrum ítaque fundaméntum Ecclésiæ Dóminus nominávit: et ídeo digne fundaméntum hoc Ecclésia colit, supra quod ecclesiástici ædifícii altitúdo consúrgit. Unde conveniénter psalmus, qui lectus est, dicit: Exáltent eum in ecclésia plebis, et in cáthedra seniórum laudent eum. Benedíctus Deus, qui beátum Petrum Apóstolum in Ecclésia exaltári præcépit: quia dignum est ut fundaméntum hoc in Ecclésia honorétur, per quod ad cælum conscénditur.\par
\switchcolumn
\EN
\noindent It was the Lord Himself Who called Peter the foundation of the Church, and therefore it is right that the Church should reverence this foundation whereon her mighty structure riseth. Justly is it written in the Psalm which we have just heard: Let them exalt him in the congregation of the people, and praise him in the assembly of the elders. Blessed be God, Who hath commanded that the Blessed Apostle Peter should be exalted in the congregation! Worthy to be honoured by the Church is that foundation from which her goodly towers rise, pointing to heaven!\par
\switchcolumn*

\LA
\begin{Resp}\Rsign Ego pro te rogávi, Petre, ut non defíciat fides tua: \Star{} Et tu aliquándo convérsus confírma fratres tuos.\end{Resp}
\begin{Resp}\Vsign Caro et sanguis non revelávit tibi, sed Pater meus, qui est in cælis.\end{Resp}
\begin{Resp}\Rsign Et tu aliquándo convérsus confírma fratres tuos.\end{Resp}
\switchcolumn
\EN
\begin{Resp}\Rsign Peter, I have prayed for thee, that thy faith fail not \Star{} And when thou art converted, strengthen thy brethren.\end{Resp}
\begin{Resp}\Vsign Flesh and blood hath not revealed it unto thee, but My Father Which is in heaven.\end{Resp}
\begin{Resp}\Rsign And when thou art converted, strengthen thy brethren.\end{Resp}
\switchcolumn*

\LA
\LessonHead{Lectio VI}
\switchcolumn
\EN
\LessonHead{Lesson VI}
\switchcolumn*

\LA
\noindent Quod natális ergo Cáthedræ hódie cólitur, sacerdotále honorátur offícium. Sibi hoc Ecclésiæ ínvicem præstant, quia tanto necésse plus habet Ecclésia dignitatis, quanto sacerdotále offícium plus honóris. Cum solemnitátem hanc Ecclésiis mérito religiósa observátio introdúxerit, miror, cur apud quosdam infidéles hódie tam perniciósus error incréverit, ut super túmulos defunctórum cibos et vina conférant, quasi egréssæ de corpóribus ánimæ carnáles cibos requírant.\par
\switchcolumn
\EN
\noindent In the honour which is this day paid to the inauguration of the first Bishop's throne, an honour is paid to the office of all Bishops. The Churches testify one to another, that, the greater the Church's dignity, the greater the reverence due to her priests. While I confess how rightly godly custom hath exalted this Feast in the estimation of all the Churches, the more do I wonder at the growth of that unhealthy error which at this day causeth some unbelievers to lay food and wine upon the graves of the dead, as if souls once rid of the body had any longer any need of bodily refreshment.\par
\switchcolumn*

\LA
\begin{Resp}\Rsign Petre, amas me? Tu scis, Dómine, quia amo te. \Star{} Pasce oves meas.\end{Resp}
\begin{Resp}\Vsign Simon Joánnis, díligis me plus his? Tu scis, Dómine, quia amo te.\end{Resp}
\begin{Resp}\Rsign Pasce oves meas.\end{Resp}
\begin{Resp}\Vsign Glória Patri, et Fílio, \Star{} et Spirítui Sancto.\end{Resp}
\begin{Resp}\Rsign Pasce oves meas.\end{Resp}
\switchcolumn
\EN
\begin{Resp}\Rsign Peter, lovest thou Me? Lord, Thou knowest that I love thee. \Star{} Feed My sheep.\end{Resp}
\begin{Resp}\Vsign Simon, son of Jonas, lovest thou Me more than these? Lord, Thou knowest that I love thee.\end{Resp}
\begin{Resp}\Rsign Feed My sheep.\end{Resp}
\begin{Resp}\Vsign Glory be to the Father, and to the Son, \Star{} and to the Holy Ghost.\end{Resp}
\begin{Resp}\Rsign Feed My sheep.\end{Resp}
\switchcolumn*

\LA
\LessonHead{Lectio VII}
\switchcolumn
\EN
\LessonHead{Lesson VII}
\switchcolumn*

\LA
\LessonTitle{Léctio sancti Evangélii secúndum Matthǽum}
\switchcolumn
\EN
\LessonTitle{From the Holy Gospel according to Matthew}
\switchcolumn*

\LA
\Cite{Matt 16:13-19}
\switchcolumn
\EN
\Cite{Matt 16:13-19}
\switchcolumn*

\LA
\noindent In illo témpore: Venit Jesus in partes Cæsaréæ Philíppi, et interrogábat discípulos suos, dicens: Quem dicunt hómines esse Fílium hóminis? Et réliqua.\par
\switchcolumn
\EN
\noindent At that time, Jesus came into the coasts of Caesarea Philippi, and He asked His disciples, saying: Who do men say that I, the Son of Man, am? And so on.\par
\switchcolumn*

\LA
\LessonTitle{Homilía sancti Leónis Papæ}
\switchcolumn
\EN
\LessonTitle{Homily by Pope St. Leo (the Great.)}
\switchcolumn*

\LA
\Source{Sermo 3 in annivers. assumpt. suæ, post initium}
\switchcolumn
\EN
\Source{3rd on the Anniversary of his own election}
\switchcolumn*

\LA
\noindent Apóstolos Dóminus, quid de se hómines opinéntur, intérrogat: et tamdiu sermo respondéntium commúnis est, quámdiu humánæ intellegéntiæ ambigúitas explicátur. At ubi quid hábeat sensus discipulórum exígitur, primus est in Dómini confessióne, qui primus est in apostólica dignitáte. Qui cum dixísset: Tu es Christus Fílius Dei vivi: respóndit ei Jesus: Beátus es Simon Bar-Jona, quia caro et sanguis non revelávit tibi, sed Pater meus, qui in cælis est. Id est, Ideo beátus es, quia Pater meus te dócuit: nec terréna opínio te feféllit, sed inspirátio cæléstis instrúxit: et non caro et sanguis, sed ille me tibi, cujus sum unigénitus Fílius, indicávit.\par
\switchcolumn
\EN
\noindent The Lord asked His disciples Who men said that He was, and their answers were human as long as they were the answers of human reason, unilluminated by Divine light. At last, when the glimmerings of earthly conjecture were spoken, he whose Apostleship is the first in dignity, was the first to confess his Lord. And Simon Peter answered and said: Thou art the Christ, the Son of the living God. And Jesus answered and said unto him: Blessed art thou, Simon Bar-Jona, for flesh and blood hath not revealed it unto thee, but My Father Which is in heaven. That is to say, For this cause art thou blessed, because My Father Himself hath taught thee; the opinions of men have not beguiled thee, the voices of angels have not taught thee, not flesh and blood, but He, Whose Only begotten Son I am, hath revealed Me unto thee.\par
\switchcolumn*

\LA
\begin{Resp}\Rsign Quem dicunt hómines esse Filium hóminis? dixit Jesus discípulis suis. Respóndens Petrus dixit: Tu es Christus Fílius Dei vivi. \Star{} Et ego dico tibi, quia tu es Petrus, et super hanc petram ædificábo Ecclésiam meam.\end{Resp}
\begin{Resp}\Vsign Beátus es, Simon Bar-Jona, quia caro et sanguis non revelávit tibi, sed Pater meus qui est in cælis.\end{Resp}
\begin{Resp}\Rsign Et ego dico tibi, quia tu es Petrus, et super hanc petram ædificábo Ecclésiam meam.\end{Resp}
\switchcolumn
\EN
\begin{Resp}\Rsign Jesus asked His disciples, saying: Who do men say that I, the Son of Man, am? Peter answered, and said: Thou art the Christ, the Son of the living God. \Star{} And I say unto thee, that thou art Peter, and upon this rock I will build My Church.\end{Resp}
\begin{Resp}\Vsign Blessed art thou, Simon Bar-Jona, for flesh and blood hath not revealed it unto thee, but My Father Which is in heaven.\end{Resp}
\begin{Resp}\Rsign And I say unto thee, that thou art Peter, and upon this rock I will build My Church.\end{Resp}
\switchcolumn*

\LA
\LessonHead{Lectio VIII}
\switchcolumn
\EN
\LessonHead{Lesson VIII}
\switchcolumn*

\LA
\noindent Et ego, inquit, dico tibi: hoc est, sicut Pater meus tibi manifestávit divinitátem meam, ita et ego tibi notam fácio excelléntiam tuam, quia tu es Petrus: id est, cum ego sim inviolábilis petra, ego lapis anguláris, qui fácio útraque unum, ego fundaméntum præter quod nemo potest áliud pónere: tamen tu quoque petra es, quia mea virtúte solidáris, ut quæ mihi potestáte sunt própria, sint tibi mecum participatióne commúnia. Et super hanc petram ædificábo Ecclésiam meam, et portæ ínferi non prævalébunt advérsus eam. Super hanc, inquit, fortitúdinem ætérnum éxstruam templum: et Ecclésiæ meæ cælo inserénda sublímitas, in hujus fídei firmitáte consúrget.\par
\switchcolumn
\EN
\noindent Thus saith the Lord unto Simon Peter: And I say also unto thee, That thou art Peter. That is to say, Even as My Father hath revealed unto thee concerning Me that I am God, even so now will I also reveal unto thee that thou art Peter; I am the sure Rock of defence, the Corner Stone, Who make both one, (Eph. ii. 20, 15,) I am the Foundation, beside Which other can no man lay, (1 Cor. iii. 11,) and thou also art a rock, in My Strength made hard, and those things whereof I by right am Lord, into thy hand do I give them, that thou mayst bear rule over them, for Me, and with Me. And upon this rock I will build My Church, and the gates of hell shall not prevail against it. Upon this strength of thine, whereof I am the Strength, I will build My eternal temple, and upon the truth of thy confession of Me I will make to rise My glorious Church whose spires shall pierce to heaven.\par
\switchcolumn*

\LA
\begin{Resp}\Rsign Elégit te Dóminus sacerdótem sibi, ad sacrificándum ei \Star{} Hóstiam laudis.\end{Resp}
\begin{Resp}\Vsign Immola Deo sacrifícium laudis, et redde Altíssimo vota tua.\end{Resp}
\begin{Resp}\Rsign Hóstiam laudis.\end{Resp}
\begin{Resp}\Vsign Glória Patri, et Fílio, \Star{} et Spirítui Sancto.\end{Resp}
\begin{Resp}\Rsign Hóstiam laudis.\end{Resp}
\switchcolumn
\EN
\begin{Resp}\Rsign The Lord hath chosen thee for a priest unto Himself, to offer up unto Him \Star{} The sacrifice of praise.\end{Resp}
\begin{Resp}\Vsign Offer unto God thanksgiving, and pay thy vows unto the Most High.\end{Resp}
\begin{Resp}\Rsign The sacrifice of praise.\end{Resp}
\begin{Resp}\Vsign Glory be to the Father, and to the Son, \Star{} and to the Holy Ghost.\end{Resp}
\begin{Resp}\Rsign The sacrifice of praise.\end{Resp}
\switchcolumn*

\LA
\LessonHead{Lectio IX}
\switchcolumn
\EN
\LessonHead{Lesson IX}
\switchcolumn*

\LA
\noindent Hanc confessiónem portæ ínferi non tenébunt, mortis víncula non ligábunt: vox enim ista vox vitæ est. Et sicut confessóres suos in cæléstia próvehit, ita negatóres ad inférna demérgit. Propter quod dicit beatíssimo Petro: Tibi dabo claves regni cælórum: Et quæcúmque ligáveris super terram, erunt ligáta et in cælis: et quæcúmque sólveris super terram, erunt solúta et in cælis. Transívit quidem étiam in álios Apóstolos vis potestátis istíus, et ad omnes Ecclésiæ príncipes decréti hujus constitútio commeávit: sed non frustra uni commendátur, quod ómnibus intimátur. Petro enim ídeo hoc singuláriter créditur, quia cunctis Ecclésiæ rectóribus Petri forma præpónitur. Manet ergo Petri privilégium, ubicúmque ex ipsíus fertur æquitáte judícium. Nec nímia est vel sevéritas vel remíssio, ubi nihil erit ligátum, nihil solútum, nisi quod beátus Petrus aut sólverit, aut ligáverit.\par
\switchcolumn
\EN
\noindent Against this confession the gates of hell shall never prevail, neither shall the bonds of death take hold upon it. Thus saith He That is faithful and true. And as this confession hath power to lift up to heaven them that make it, so is it able to thrust down to hell them that gainsay it. Wherefore it is said unto the most blessed Peter: And I will give unto thee the keys of the kingdom of heaven and whatsoever thou shalt bind on earth, shall be bound in heaven; and whatsoever thou shalt loose on earth, shall be loosed in heaven. This power passed indeed to the other Apostles also; this the Lord's will had effect in them; but it is not in vain that it is written that that was given to one which passed from him to all. To Peter alone were the keys given, and Peter is set as the pattern for all them that bear rule in the Church to follow. There remaineth therefore the right of Peter, wheresoever his judgment decreeth justice. Neither is there anything too hard, or too lax, where there is nothing bound and nothing loosed, save when Peter bindeth or looseth.\par
\switchcolumn*

\LA
\TeDeum{Te Deum}
\switchcolumn
\EN
\TeDeum{Te Deum}
\switchcolumn*

\end{paracol}

\par\needspace{10\baselineskip}\addvspace{1.2em}{\centering\small\textsc{Die 23 Februarii} · \textsc{23 February}\par}
\Office{In Vigilia S. Matthiæ Ap.}{Vigil of St. Matthias, Apostle}{Vigilia}
\addcontentsline{toc}{subsection}{Die 23 Februarii · In Vigilia S. Matthiæ Ap.\enspace\textit{Vigil of St. Matthias, Apostle}}
\begin{paracol}{2}
\LA
\RefLine{Lectiones I–III: de Scriptura occurrente.}
\switchcolumn
\EN
\RefLine{Lessons I–III: from the Scripture of the day.}
\switchcolumn*

\end{paracol}

\par\needspace{10\baselineskip}\addvspace{1.2em}{\centering\small\textsc{Die 23 Februarii} · \textsc{23 February}\par}
\Office{S. Petri Damiani Episcopi Confessoris et Ecclesiæ Doctoris}{St. Peter Damian, Confessor and Doctor of the Church}{Duplex}
\addcontentsline{toc}{subsection}{Die 23 Februarii · S. Petri Damiani Episcopi Confessoris et Ecclesiæ Doctoris\enspace\textit{St. Peter Damian, Confessor and Doctor of the Church}}
\begin{paracol}{2}
\LA
\RefLine{Lectiones I–III: ex Commúni Doctoris Pontificis (C4a).}
\switchcolumn
\EN
\RefLine{Lessons I–III: from the Common of a Doctor Bishop (C4a).}
\switchcolumn*

\LA
\RefLine{Lectiones VII–VIII: ex Commúni Doctoris Pontificis (C4a).}
\switchcolumn
\EN
\RefLine{Lessons VII–VIII: from the Common of a Doctor Bishop (C4a).}
\switchcolumn*

\LA
\RefLine{Lectio IX: de Scriptura occurrente.}
\switchcolumn
\EN
\RefLine{Lesson IX: from the Scripture of the day.}
\switchcolumn*

\LA
\LessonHead{Lectio IV}
\switchcolumn
\EN
\LessonHead{Lesson IV}
\switchcolumn*

\LA
\noindent Petrus, Ravennæ honestis paréntibus natus, adhuc lactens a matre, numerosæ prolis pertæsa, abjícitur; sed domesticæ mulíeris ópera semivivus exceptus ac recreatus, genitrici ad humanitátis sensum revocátæ redditur. Utroque orbátus parente, tamquam vile mancipium sub aspera fratris tutela duram servitútem exercuit. Religiónis in Deum ac pietátis erga patrem, egregium tunc specimen dedit; invéntum síquidem forte nummum, non propriæ inediæ sublevandæ, sed sacerdoti, qui divinum sacrifícium ad illíus expiatiónem offerret, erogávit. A Damiáno fratre, a quo, uti fertur, cognoméntum accepit, benígne recéptus, ejus cura litteris erudítur in quibus brevi tantum profecit, ut magistris admiratióni esset. Cum autem liberálibus sciéntiis floréret et nómine, eas cum laude docuit. Interim ut corpus ratióni súbderet, sub mollibus vestibus cilícium adhibuit; jejuniis, vigiliis et oratiónibus solerter insistens. Calénte juvénta dum carnis stimulis acriter urgerétur, insultántium libidinum faces rigéntibus fluvii mersus aquis noctu exstinguebat; tum venerabília quæque loca obire, totumque Psaltérium recitare consueverat. Ope assidua páuperes levabat, quibus frequenter pastis convivio, propriis ipse mánibus ministrabat.\par
\switchcolumn
\EN
\noindent The holy Doctor Peter Damian was born of respectable parents at Ravenna, (about the year of our Lord 988.) While he was still a suckling, his mother, overcome with the care of many children, cast him out to perish, but one of the women servants saved him when he was nigh to death, and fed him until natural affection appeared again in his mother, to whom she then gave him back. After the death of both his parents he lived with a brother who treated him like the lowest slave, and in whose house he underwent a hard bondage. Even while he was in this condition he gave a wonderful proof of his faith toward God, and his dutiful love toward his father. It chanced that one day he found a considerable sum of money, but instead of using it to relieve his own poverty, he gave it all to a priest to offer God's sacrifice for the forgiveness of his father's sins. He had happily another brother called Damian, the same from whom he seemeth afterwards to have taken his surname. By him he was affectionately adopted, and put in the way of being educated. He made such progress in learning as astonished his teachers, and when he had won an eminent name in letters, he began to teach on his own accord with general applause. Meanwhile, lest his body should get the better of his mind, he constantly wore a hair-shirt under his softer clothes, and exercised himself in fasting, watching, and prayer. In the spring-time of his age he was grievously tormented by the stings of the flesh; and sometimes, when the rebellions of lust seemed about to get the mastery over him at night, he threw himself into a freezing stream to check them. After this he would go about visiting consecrated places, and repeat the whole book of Psalms. He was most careful in relieving the poor, on whom he would wait with his own hands.\par
\switchcolumn*

\LA
\begin{Resp}\Rsign Invéni David servum meum, óleo sancto meo unxi eum: \Star{} Manus enim mea auxiliábitur ei.\end{Resp}
\begin{Resp}\Vsign Nihil profíciet inimícus in eo, et fílius iniquitátis non nocébit ei.\end{Resp}
\begin{Resp}\Rsign Manus enim mea auxiliábitur ei.\end{Resp}
\switchcolumn
\EN
\begin{Resp}\Rsign I have found David My servant, with My holy oil have I anointed him \Star{} For My hand shall help him.\end{Resp}
\begin{Resp}\Vsign The enemy shall prevail nothing against him, nor the son of wickedness afflict him.\end{Resp}
\begin{Resp}\Rsign For My hand shall help him.\end{Resp}
\switchcolumn*

\LA
\LessonHead{Lectio V}
\switchcolumn
\EN
\LessonHead{Lesson V}
\switchcolumn*

\LA
\noindent Perficiendæ magis vitæ causa, in Avellanénsi Eugubinæ diœcesis cœnobio, ordini monachórum sanctæ Crucis Fontis Avellanæ a beato Ludulpho, sancti Romualdi discipulo, fundato, nomen dedit. Non ita multo post in monastérium Pomposianum, mox in cœnobium sancti Vincentii Petræ Pertusæ ab abbate suo missus, utrumque ascetérium verbo sacro, præcláris institutiónibus et móribus excoluit. Ad suos revocatus, post præsidis obitum, Avellanitárum famíliæ præfícitur, quam, novis variis in locis exstructis domiciliis et sanctíssimis institútis ita auxit, ut alter ejus ordinis parens ac præcipuum ornaméntum jure sit habitus. Salutarem Petri sollicitúdinem alia quoque diversi institúti cœnobia, canonicórum convéntus, et pópuli sunt experti. Urbináti diœcesi non uno nómine profuit: Theuzóni episcopo in causa gravíssima assedit, ipsumque in recte administrando episcopatu consílio et ópera juvit. Divinórum contemplatióne, corporis maceratiónibus, ceterisque spectátæ sanctimoniæ exemplis excelluit. His motus Stephanus nonus Pontifex maximus eum, licet invitum et reluctántem, sanctæ Romanæ Ecclésiæ cardinalem creávit et Ostiensem episcopum. Quas Petrus dignitates splendidíssimis virtútibus et consentáneis episcopali ministerio opéribus gessit.\par
\switchcolumn
\EN
\noindent Desiring to attain to perfection of life he betook himself to the convent of Font-Avellano, in the diocese of Gubbio, in Umbria, a house founded by the blessed Ludolph, the disciple of St. Romuald, for the monks of the Holy Cross. He dwelt there not long before he was sent by his Abbat, first to the Abbey of Pomposia, and, secondly, to that of St. Vincent at Pietra Pertusa, both which brotherhoods he greatly profited by his godly exhortations, discreet rules, and grave manners. After his return home, and the death of his Superior, he was chosen to rule the brethren of Avellano. Here he founded diverse new hermitages, and made the community so to flourish under his saintly direction, that he is esteemed the second Father and chief ornament of that Order. This healthful care of Peter was made a blessing to convents of other Rules than his own, to houses of Canons, and to the people. He was many ways profitable to the diocese of Urbino. He sat with Theuzo the Bishop of that See to judge of a most weighty matter, and led him by his counsel and assistance rightly to administer his Bishopric. He was foremost in contemplation of the things of God, in severity toward his own body, and in other things whereby to set a bright example of godliness. In consideration of these things the Supreme Pontiff Stephen IX., (in the year 1057,) created him, in spite of his own unwillingness and objections, a Cardinal of the Holy Roman Church, and appointed him Bishop of Ostia. This dignity Peter bore with the highest reputation for piety, and adorned with works meet for a Bishop.\par
\switchcolumn*

\LA
\begin{Resp}\Rsign Pósui adjutórium super poténtem, et exaltávi eléctum de plebe mea: \Star{} Manus enim mea auxiliábitur ei.\end{Resp}
\begin{Resp}\Vsign Invéni David servum meum, óleo sancto meo unxi eum.\end{Resp}
\begin{Resp}\Rsign Manus enim mea auxiliábitur ei.\end{Resp}
\switchcolumn
\EN
\begin{Resp}\Rsign I have laid help upon one that is mighty, and have exalted one chosen out of My people \Star{} For My hand shall help him.\end{Resp}
\begin{Resp}\Vsign I have found David My servant, with My holy oil have I anointed him.\end{Resp}
\begin{Resp}\Rsign For My hand shall help him.\end{Resp}
\switchcolumn*

\LA
\LessonHead{Lectio VI}
\switchcolumn
\EN
\LessonHead{Lesson VI}
\switchcolumn*

\LA
\noindent Dificillimo témpore Romanæ Ecclésiæ, summisque Pontificibus doctrina, legatiónibus, aliisque susceptis labóribus mirifice adfuit. Adversus Nicolaitárum et simoníacam hæreses ad mortem usque strenue decertávit. Hujusmodi depulsis malis, Mediolanensem Ecclésiæ Romanæ conciliávit. Benedicto et Cadalóo falsis pontificibus fortiter restitit: Henricum quartum Germaniæ regem ab iníquo uxoris divortio deterruit: Ravennates ad débita Romano Pontifici obsequia revocatos sacris restituit: canonicos Veliternos ad sanctioris vitæ leges compósuit. In provincia præsertim Urbinate vix ulla fuit episcopalis ecclésia, de qua Petrus non sit bene meritus: Eugubínam, quam aliquándo creditam hábuit, multis levávit incommodis; alias álibi, quando oportuit, perinde curávit, ac suæ essent tutelæ commissæ. Cardinalatu et episcopali dignitate depositis, nihil de prístina juvándi próximos sedulitate remisit. Jejunium sextæ fériæ in honórem sanctæ Crucis Jesu Christi, horarias beátæ Dei Genitrícis preces, ejusque die Sabbato cultum propagávit. Inferendæ quoque sibi verberatiónis morem ad patratórum scelerum expiatiónem provéxit. Demum sanctitáte, doctrina, miraculis et præcláre actis illustris, dum e Ravennate legatióne rediret, Favéntiæ octavo Kalendas Martii migrávit ad Christum. Ejus corpus ibidem apud Cistercienses, multis miraculis clarum, frequénti populórum veneratióne cólitur. Ipsum Faventini, non semel in præsénti discrimine propitium apud Deum delegérunt. Leo vero duodecimus Pontifex maximus Offícium Missamque in ejus honórem tamquam Confessóris Pontificis, quæ aliquibus in diœcesibus atque in ordine Camalulensium jam celebrabántur, ex sacrórum Rituum Congregatiónis consulto, áddita Doctoris qualitate, ad universam exténdit Ecclésiam.\par
\switchcolumn
\EN
\noindent At the most anxious times he greatly sustained the Church of Rome and the Supreme Pontiffs by his teaching, by missions which he discharged, and by diverse other labours which he undertook on their behalf. He strove manfully even unto death against the heresies of the Nicola'itans and the Simoniacs, by putting down which evils he reconciled the Church of Milan to that of Rome. He was one of the stoutest opponents of the false Popes Benedict and Cadalous. He deterred Henry IV., King of Germany, from his wicked scheme for putting away his wife. He recalled the people of Ravenna to their bounden duty to the Bishop of Rome, and restored them to the communion of the Church. He reformed the Canons of Velletri, and brought them to lead more godly lives. There were hardly any Cathedral Churches, especially in the province of Urbino, of which he did not deserve well. In Gubbio, of which he had at one time the management, he abolished many things unseemly. He brought about improvements in many and diverse places, as if each were his special charge. (In 1062) he gave up his dignities of Cardinal and Bishop, but he allowed his love toward his neighbours to know no diminution. He was particularly zealous in spreading abroad four devout practices 1st, To fast every Friday in honour of the Holy Cross of Jesus Christ; 2nd, To recite the Hours of the Blessed Mother of God, called also her Little Office; 3rd, To sanctify Saturday in her honour; and 4th, and especially, to scourge oneself in punishment for sin committed. At length he departed to be with Christ, at Faenza, on his way back from his mission to Ravenna, on the 22nd of February, (in the year 1072,) at the height of his reputation for holiness, learning, miracles, and good works. His body is buried in the house of the Cistercians at Faenza, where the people resort often to his grave with great reverence. The citizens of Faenza, to whom he hath been found good at need even to this day, have chosen him for their Patron in the presence of God. The supreme Pontiff Leo XII., finding that an Office and Mass in memory of him, as a Confessor and Bishop, was in use in some dioceses, and in the Camaldolese Order, by advice of the Sacred Congregation of Rites, added the title of Doctor, and extended the use of the said Office and Mass to the whole Church.\par
\switchcolumn*

\LA
\begin{Resp}\Rsign Iste est, qui ante Deum magnas virtútes operátus est, et omnis terra doctrína ejus repléta est: \Star{} Ipse intercédat pro peccátis ómnium populórum.\end{Resp}
\begin{Resp}\Vsign Iste est, qui contémpsit vitam mundi, et pervénit ad cæléstia regna.\end{Resp}
\begin{Resp}\Rsign Ipse intercédat pro peccátis ómnium populórum.\end{Resp}
\begin{Resp}\Vsign Glória Patri, et Fílio, \Star{} et Spirítui Sancto.\end{Resp}
\begin{Resp}\Rsign Ipse intercédat pro peccátis ómnium populórum.\end{Resp}
\switchcolumn
\EN
\begin{Resp}\Rsign This is he which wrought great wonders before God, and the whole earth is full of his teaching \Star{} May he pray for all people, that their sins may be forgiven unto them\end{Resp}
\begin{Resp}\Vsign This is he which loved not his life in this world, and hath attained unto the kingdom of heaven.\end{Resp}
\begin{Resp}\Rsign May he pray for all people, that their sins may be forgiven unto them\end{Resp}
\begin{Resp}\Vsign Glory be to the Father, and to the Son, \Star{} and to the Holy Ghost.\end{Resp}
\begin{Resp}\Rsign May he pray for all people, that their sins may be forgiven unto them\end{Resp}
\switchcolumn*

\end{paracol}

\par\needspace{10\baselineskip}\addvspace{1.2em}{\centering\small\textsc{Die 24 Februarii} · \textsc{24 February}\par}
\Office{S. Matthiæ Apostoli}{St. Matthias the Apostle}{Duplex II classis}
\addcontentsline{toc}{subsection}{Die 24 Februarii · S. Matthiæ Apostoli\enspace\textit{St. Matthias the Apostle}}
\begin{paracol}{2}
\LA
\RefLine{Lectio IX: de Scriptura occurrente.}
\switchcolumn
\EN
\RefLine{Lesson IX: from the Scripture of the day.}
\switchcolumn*

\LA
\LessonHead{Lectio I}
\switchcolumn
\EN
\LessonHead{Lesson I}
\switchcolumn*

\LA
\LessonTitle{De Actibus Apostolórum}
\switchcolumn
\EN
\LessonTitle{From the Acts of the Apostles}
\switchcolumn*

\LA
\Cite{Acts 1:15-18}
\switchcolumn
\EN
\Cite{Acts 1:15-18}
\switchcolumn*

\LA
\noindent In diébus illis, exsúrgens Petrus in médio fratrum, dixit (erat autem turba hóminum simul, fere centum vigínti): Viri fratres, opórtet impléri Scriptúram quam prædíxit Spíritus Sanctus per os David de Juda, qui fuit dux eórum qui comprehendérunt Jesum: qui connumerátus erat in nobis, et sortítus est sortem ministérii hujus. Et hic quidem possédit agrum de mercéde iniquitátis, et suspénsus crépuit médius: et diffúsa sunt ómnia víscera ejus.\par
\switchcolumn
\EN
\noindent In those days Peter rising up in the midst of the brethren, said: now the number of persons together was about an hundred and twenty: Men, brethren, the scripture must needs be fulfilled, which the Holy Ghost spoke before by the mouth of David concerning Judas, who was the leader of them that apprehended Jesus: Who was numbered with us, and had obtained part of this ministry. And he indeed hath possessed a field of the reward of iniquity, and being hanged, burst asunder in the midst: and all his bowels gushed out.\par
\switchcolumn*

\LA
\begin{Resp}\Rsign Ecce ego mitto vos sicut oves in médio lupórum, dicit Dóminus: \Star{} Estóte ergo prudéntes sicut serpéntes, et símplices sicut colúmbæ.\end{Resp}
\begin{Resp}\Vsign Dum lucem habétis, crédite in lucem, ut fílii lucis sitis.\end{Resp}
\begin{Resp}\Rsign Estóte ergo prudéntes sicut serpéntes, et símplices sicut colúmbæ.\end{Resp}
\switchcolumn
\EN
\begin{Resp}\Rsign Behold, I send you as sheep in the midst of wolves, said the Lord: \Star{} Be ye, therefore, wise as serpents and simple as doves.\end{Resp}
\begin{Resp}\Vsign Whilst you have the light, believe in the light, that you may be the children of light.\end{Resp}
\begin{Resp}\Rsign Be ye therefore wise as serpents and simple as doves.\end{Resp}
\switchcolumn*

\LA
\LessonHead{Lectio II}
\switchcolumn
\EN
\LessonHead{Lesson II}
\switchcolumn*

\LA
\Cite{Acts 1:19-22}
\switchcolumn
\EN
\Cite{Acts 1:19-22}
\switchcolumn*

\LA
\noindent Et notum factum est ómnibus habitántibus Jerúsalem, ita ut appellarétur ager ille, lingua eórum, Hacéldama, hoc est, ager sánguinis. Scriptum est enim in libro Psalmórum: Fiat commorátio eórum desérta, et non sit qui inhábitet in ea: et episcopátum ejus accípiat alter. Opórtet ergo ex his viris qui nobíscum sunt congregáti in omni témpore quo intrávit et exívit inter nos Dóminus Jesus, incípiens a baptísmate Joánnis usque in diem qua assúmptus est a nobis, testem resurrectiónis ejus nobíscum fíeri unum ex istis.\par
\switchcolumn
\EN
\noindent And it became known to all the inhabitants of Jerusalem: so that the same field was called in their tongue, Haceldama, that is to say, The field of blood. For it is written in the book of Psalms: Let their habitation become desolate, and let there be none to dwell therein. And his bishopric let another take. Wherefore of these men who have companied with us all the time that the Lord Jesus came in and went out among us, Beginning from the baptism of John, until the day wherein he was taken up from us, one of these must be made a witness with us of his resurrection.\par
\switchcolumn*

\LA
\begin{Resp}\Rsign Tóllite jugum meum super vos, dicit Dóminus, et díscite a me, quia mitis sum et húmilis corde: \Star{} Jugum enim meum suáve est, et onus meum leve.\end{Resp}
\begin{Resp}\Vsign Et inveniétis réquiem animábus vestris.\end{Resp}
\begin{Resp}\Rsign Jugum enim meum suáve est, et onus meum leve.\end{Resp}
\switchcolumn
\EN
\begin{Resp}\Rsign Take up my yoke upon you, said the Lord, and learn of me, because I am meek, and humble of heart. \Star{} For my yoke is sweet and my burden light.\end{Resp}
\begin{Resp}\Vsign And you shall find rest to your souls.\end{Resp}
\begin{Resp}\Rsign For my yoke is sweet and my burden light.\end{Resp}
\switchcolumn*

\LA
\LessonHead{Lectio III}
\switchcolumn
\EN
\LessonHead{Lesson III}
\switchcolumn*

\LA
\Cite{Acts 1:23-26}
\switchcolumn
\EN
\Cite{Acts 1:23-26}
\switchcolumn*

\LA
\noindent Et statuérunt duos, Joseph, qui vocabátur Bársabas, qui cognominátus est Justus, et Matthíam. Et orántes dixérunt: Tu Dómine, qui corda nosti ómnium, osténde quem elégeris ex his duóbus unum, accípere locum ministérii hujus et apostolátus, de quo prævaricátus est Judas ut abíret in locum suum. Et dedérunt sortes eis, et cécidit sors super Matthíam: et annumerátus est cum úndecim Apóstolis.\par
\switchcolumn
\EN
\noindent And they appointed two, Joseph, called Barsabas, who was surnamed Justus, and Matthias. And praying, they said: Thou, Lord, who knowest the hearts of all men, shew whether of these two thou hast chosen, To take the place of this ministry and apostleship, from which Judas hath by transgression fallen, that he might go to his own place. And they gave them lots, and the lot fell upon Matthias, and he was numbered with the eleven apostles.\par
\switchcolumn*

\LA
\begin{Resp}\Rsign Dum stetéritis ante reges et prǽsides, nolíte cogitáre quómodo aut quid loquámini: \Star{} Dábitur enim vobis in illa hora, quid loquámini.\end{Resp}
\begin{Resp}\Vsign Non enim vos estis qui loquímini: sed Spíritus Patris vestri, qui lóquitur in vobis.\end{Resp}
\begin{Resp}\Rsign Dábitur enim vobis in illa hora, quid loquámini.\end{Resp}
\begin{Resp}\Vsign Glória Patri, et Fílio, \Star{} et Spirítui Sancto.\end{Resp}
\begin{Resp}\Rsign Dábitur enim vobis in illa hora, quid loquámini.\end{Resp}
\switchcolumn
\EN
\begin{Resp}\Rsign But when they shall deliver you up to the judges, take no thought how or what to speak: \Star{} for it shall be given you in that hour what to speak:\end{Resp}
\begin{Resp}\Vsign For it is not you that speak, but the spirit of your Father that speaketh in you.\end{Resp}
\begin{Resp}\Rsign For it shall be given you in that hour what to speak.\end{Resp}
\begin{Resp}\Vsign Glory be to the Father, and to the Son, \Star{} and to the Holy Ghost.\end{Resp}
\begin{Resp}\Rsign For it shall be given you in that hour what to speak.\end{Resp}
\switchcolumn*

\LA
\LessonHead{Lectio IV}
\switchcolumn
\EN
\LessonHead{Lesson IV}
\switchcolumn*

\LA
\LessonTitle{De Expositióne sancti Augustíni Epíscopi super psalmum octogésimum sextum}
\switchcolumn
\EN
\LessonTitle{From the Exposition of the Eighty-sixth Psalm by St. Augustine, Bishop (of Hippo.)}
\switchcolumn*

\LA
\Source{Ante medium}
\switchcolumn
\EN
\Source{Ante medium}
\switchcolumn*

\LA
\noindent Fundaménta ejus in móntibus sanctis: díligit Dóminus portas Sion. Quare sunt fundaménta Apóstoli et Prophétæ? Quia eórum auctóritas portat infirmitátem nostram. Quare sunt portæ? Quia per ipsos intrámus ad regnum Dei. Prǽdicant enim nobis: et, cum per ipsos intrámus, per Christum intrámus; ipse est enim jánua. Et cum dicúntur duódecim portæ Jerúsalem, et una porta Christus et duódecim portæ Christus, quia in duódecim portis Christus; et ídeo duodenárius númerus Apostolórum. Sacraméntum magnum hujus duodenárii significátio est númeri. Sedébitis, inquit, super duódecim sedes, judicántes duódecim tribus Israël.\par
\switchcolumn
\EN
\noindent Her foundation is in the holy mountains; the Lord loveth the gates of Zion. Wherefore hath the city twelve foundations, and in them the names of the Prophets and of the Apostles of the Lamb? Because their authority is the foundation whereon our weakness resteth. Wherefore are they the gates? Because through them we enter in unto the kingdom of God, since they have preached the same unto us, and when we enter in through their preaching, we enter in by Christ, Who is Himself The Door. (John x. 7.) And, whereas it is written that the city hath twelve gates, and, again, that Christ is the one Door, Christ is all the twelve, for He is in all the twelve and therefore were twelve Apostles chosen. There lieth a great mystery in the signification of this number twelve: Ye shall sit, said the Lord, upon twelve thrones, judging the twelve tribes of Israel.\par
\switchcolumn*

\LA
\begin{Resp}\Rsign Vidi conjúnctos viros, habéntes spléndidas vestes, et Ángelus Dómini locútus est ad me, dicens: \Star{} Isti sunt viri sancti facti amíci Dei.\end{Resp}
\begin{Resp}\Vsign Vidi Ángelum Dei fortem, volántem per médium cælum, voce magna clamántem et dicéntem.\end{Resp}
\begin{Resp}\Rsign Isti sunt viri sancti facti amíci Dei.\end{Resp}
\switchcolumn
\EN
\begin{Resp}\Rsign I saw men standing together, clad in shining raiment, and the Angel of the Lord spoke unto me, saying: \Star{} These men are holy, for they are the friends of God.\end{Resp}
\begin{Resp}\Vsign I saw a strong Angel of God fly into the midst of heaven, saying with a loud voice:\end{Resp}
\begin{Resp}\Rsign These men are holy, for they are the friends of God.\end{Resp}
\switchcolumn*

\LA
\LessonHead{Lectio V}
\switchcolumn
\EN
\LessonHead{Lesson V}
\switchcolumn*

\LA
\noindent Si duódecim sellæ ibi sunt, non est ubi sédeat tértius décimus Paulus Apóstolus, et non erit quómodo júdicet; et ipse se judicatúrum dixit, non hómines tantum, sed et ángelos. Quos ángelos, nisi apóstatas ángelos? Nescítis, inquit, quia ángelos judicábimus? Respondéret ergo turba: Quid te jactas judicatúrum? Ubi sedébis? Duódecim sedes dixit Dóminus duódecim Apóstolis, unus cécidit Judas, in locum ipsíus sanctus Matthías ordinátus est; implétus est duodenárius númerus sédium. Primo locum invéni, ubi sédeas; et sic te mináre judicatúrum. Duódecim ergo sedes quid sibi velint, videámus. Sacraméntum est cujúsdam universitátis; quia per totum orbem terrárum futúra erat Ecclésia, unde vocátur hoc ædifícium ad Christi compágem.\par
\switchcolumn
\EN
\noindent If then there be set there twelve thrones of judgment, (Ps. cxxi. 5,) Paul, in that he is the thirteenth Apostle, hath not where to sit, nor wherein to judge. Nevertheless, he hath said of himself that he will judge not men only, but angels. Know ye not, saith he, that we shall judge angels? (i Cor. vi. 3,) that is, the fallen angels. Then might they have answered him: Wherefore boastest thou thyself to be a judge? For where is thy seat? The Lord hath said that for the twelve Apostles there shall be twelve thrones; one of the twelve, even Judas, is indeed fallen, but holy Matthias is chosen into his place; for the twelve thrones there are still twelve to sit thereon first find whereon thou shalt sit, and afterward give thyself out for a judge. Let us see, then, what is the meaning of these twelve thrones. By them is signified in a mystery the whole world, since the Church shall be through all the earth, whence this building is called to be built up together in Christ.\par
\switchcolumn*

\LA
\begin{Resp}\Rsign Beáti estis, cum maledíxerint vobis hómines, et persecúti vos fúerint, et díxerint omne malum advérsum vos, mentiéntes, propter me: \Star{} Gaudéte et exsultáte, quóniam merces vestra copiósa est in cælis.\end{Resp}
\begin{Resp}\Vsign Cum vos óderint hómines, et cum separáverint vos, et exprobráverint, et ejécerint nomen vestrum tamquam malum propter Fílium hóminis.\end{Resp}
\begin{Resp}\Rsign Gaudéte et exsultáte, quóniam merces vestra copiósa est in cælis.\end{Resp}
\switchcolumn
\EN
\begin{Resp}\Rsign Blessed are ye when men shall revile you, and persecute you, and shall say all manner of evil against you falsely, for My sake; \Star{} Rejoice, and be exceeding glad, for great is your reward in heaven.\end{Resp}
\begin{Resp}\Vsign When men shall hate you, and when they shall separate you from their company, and shall reproach you, and cast out your name as evil, for the Son of Man's sake.\end{Resp}
\begin{Resp}\Rsign Rejoice, and be exceeding glad, for great is your reward in heaven.\end{Resp}
\switchcolumn*

\LA
\LessonHead{Lectio VI}
\switchcolumn
\EN
\LessonHead{Lesson VI}
\switchcolumn*

\LA
\noindent Et ídeo, quia úndique venítur ad judicándum, duódecim sedes sunt; sicut, quia úndique intrátur in illam civitátem, duódecim portæ sunt. Non solum ergo illi duódecim et Apóstolus Paulus, sed quotquot judicatúri sunt, propter significatiónem universitátis, ad sedes duódecim pértinent; quemádmodum quotquot intrábunt, ad duódecim portas pértinent. Partes enim mundi quátuor sunt, Oriens, Occidens, Aquilo et Merídies. Istæ quátuor partes assídue nominántur in Scriptúris. Ab istis quátuor ventis, sicut dicit Dóminus in Evangélio, a quátuor ventis se collectúrum eléctos suos; ab ómnibus ergo istis quátuor ventis vocátur Ecclésia. Quómodo vocátur? Undique in Trinitáte vocátur. Non vocátur nisi per baptísmum in nómine Patris, et Fílii, et Spíritus Sancti. Quátuor ergo ter ducta duódecim inveniúntur.\par
\switchcolumn
\EN
\noindent Therefore is it said that there shall be twelve thrones, because from all quarters shall there come men to be judged; even as it is said that the city hath twelve gates, because from all quarters shall the nations of them which are saved, enter into it. So, not the twelve only, and the Apostle Paul, but all, as many as shall judge, have part in these twelve thrones, this signifying, that they shall judge all men; even as all that enter into the city, have part in her twelve gates. For there are four quarters of the world, the East, and the West, and the North, and the South of which four quarters is mention often made in the Scriptures. From the four winds shall the elect be gathered together, as saith the Lord in the Gospel: And He shall send His Angels with a great sound of a trumpet; and they shall gather together His elect from the four winds, from one end of heaven to the other. (Matth. xxiv. 31.) From the four winds, therefore, is the Church called together; and how are they called? Everywhere are they called in the Trinity; for they are called no otherwise than by baptizing them in the Name of the Father, and of the Son, and of the Holy Ghost. (Matth. xxvii. 19.) Now four being multiplied by three is twelve.\par
\switchcolumn*

\LA
\begin{Resp}\Rsign Isti sunt triumphatóres et amíci Dei, qui contemnéntes jussa príncipum, meruérunt prǽmia ætérna: \Star{} Modo coronántur, et accípiunt palmam.\end{Resp}
\begin{Resp}\Vsign Isti sunt qui venérunt ex magna tribulatióne, et lavérunt stolas suas in sánguine Agni.\end{Resp}
\begin{Resp}\Rsign Modo coronántur, et accípiunt palmam.\end{Resp}
\begin{Resp}\Vsign Glória Patri, et Fílio, \Star{} et Spirítui Sancto.\end{Resp}
\begin{Resp}\Rsign Modo coronántur, et accípiunt palmam.\end{Resp}
\switchcolumn
\EN
\begin{Resp}\Rsign These are they which have conquered, and are become the friends of God, who recked not of the commandments of princes, and earned the everlasting reward. \Star{} And now have they crowns on their heads, and palms in their hands.\end{Resp}
\begin{Resp}\Vsign These are they which came out of great tribulation, and have washed their robes in the blood of the Lamb.\end{Resp}
\begin{Resp}\Rsign And now have they crowns on their heads, and palms in their hands.\end{Resp}
\begin{Resp}\Vsign Glory be to the Father, and to the Son, \Star{} and to the Holy Ghost.\end{Resp}
\begin{Resp}\Rsign And now have they crowns on their heads, and palms in their hands.\end{Resp}
\switchcolumn*

\LA
\LessonHead{Lectio VII}
\switchcolumn
\EN
\LessonHead{Lesson VII}
\switchcolumn*

\LA
\LessonTitle{Léctio sancti Evangélii secúndum Matthǽum}
\switchcolumn
\EN
\LessonTitle{From the holy Gospel according to Matthew}
\switchcolumn*

\LA
\Cite{Matt 11:25-30}
\switchcolumn
\EN
\Cite{Matt 11:25-30}
\switchcolumn*

\LA
\noindent In illo témpore: Respóndens Jesus dixit: Confíteor tibi, Pater, Dómine cæli et terræ, quia abscondísti hæc a sapiéntibus et prudéntibus, et revelásti ea párvulis. Et réliqua.\par
\switchcolumn
\EN
\noindent At that time, Jesus answered and said: I confess to thee, O Father, Lord of heaven and earth, because thou hast hid these things from the wise and prudent, and hast revealed them to the little ones. And so on.\par
\switchcolumn*

\LA
\LessonTitle{Homilía sancti Augustíni Epíscopi}
\switchcolumn
\EN
\LessonTitle{Homily by St. Augustine, Bishop (of Hippo.)}
\switchcolumn*

\LA
\Source{Sermo 10 de verbis Domini}
\switchcolumn
\EN
\Source{10th Sermon on the Words of the Lord}
\switchcolumn*

\LA
\noindent Veníte ad me, omnes qui laborátis. Quare enim omnes laborámus, nisi quia sumus hómines mortáles, frágiles, infírmi, lútea vasa portántes, quæ fáciunt ínvicem angústias? Sed, si angustiántur vasa carnis, dilaténtur spátia caritátis. Quid ergo dicit, Veníte ad me, omnes qui laborátis, nisi ut non laborétis? Dénique promíssio ejus in promptu est; quóniam laborántes vocávit, quærent forte, qua mercéde vocáti sunt. Et ego vos, inquit, refíciam. Tóllite jugum meum super vos, et díscite a me, non mundum fabricáre, non cuncta visibília et invisibília creáre, non in ipso mundo mirabília fácere et mórtuos suscitáre; sed, Quóniam mitis sum et húmilis corde.\par
\switchcolumn
\EN
\noindent Come unto Me, all ye that labour! And wherefore labour we all, but because we are frail, sickly, dying creatures, burdened with earthen vessels which distress us? But if these fleshly vessels be distressful, let the open expanse of love be free and wide. Come unto Me, all ye that labour! and why? That we may labour no more. His promise is an instant promise, for He calleth such as are labouring. Perchance they will ask Him what shall be their reward? And I, saith He, will give you rest. Take My yoke upon you, and learn of me not how to make the world, not how to create all things visible and invisible, not to work wonders in the earth, nor to raise the dead but for I am meek and lowly in heart.\par
\switchcolumn*

\LA
\begin{Resp}\Rsign Isti sunt qui vivéntes in carne, plantavérunt Ecclésiam sánguine suo: \Star{} Cálicem Dómini bibérunt, et amíci Dei facti sunt.\end{Resp}
\begin{Resp}\Vsign In omnem terram exívit sonus eórum, et in fines orbis terræ verba eórum.\end{Resp}
\begin{Resp}\Rsign Cálicem Dómini bibérunt, et amíci Dei facti sunt.\end{Resp}
\switchcolumn
\EN
\begin{Resp}\Rsign These are they who while yet they lived in the flesh, planted the Church in their own blood; \Star{} They drank of the Lord's cup, and became the friends of God.\end{Resp}
\begin{Resp}\Vsign Their sound is gone out through all the earth, and their words to the ends of the world.\end{Resp}
\begin{Resp}\Rsign They drank of the Lord's cup, and became the friends of God.\end{Resp}
\switchcolumn*

\LA
\LessonHead{Lectio VIII}
\switchcolumn
\EN
\LessonHead{Lesson VIII}
\switchcolumn*

\LA
\noindent Magnus esse vis? a mínimo íncipe. Cógitas magnam fábricam constrúere celsitúdinis? de fundaménto prius cógita humilitátis. Et quantam quisque vult et dispónit superimpónere molem ædifícii, quanto erit majus ædifícium, tanto áltius fodit fundaméntum. Et fábrica quidem cum constrúitur, in supérna consúrgit; qui autem fodit fundaméntum, ad ima deprímitur. Ergo et fábrica ante celsitúdinem humiliátur, et fastígium post humiliatiónem erígitur.\par
\switchcolumn
\EN
\noindent Wilt thou be great? Begin by being little. Dost thou think to raise up a lofty building? Then lay the foundations thereof in lowliness. The greater soever, and the more massy, be that which any man thinketh to build, so much the deeper doth he dig his foundation. And when the house is built, it towereth heavenward; but he which layeth the foundation goeth down into the earth. The building, therefore, is low before it is high, and, after it is low, it riseth high to the roof.\par
\switchcolumn*

\LA
\begin{Resp}\Rsign Isti sunt viri sancti, quos elégit Dóminus in caritáte non ficta, et dedit illis glóriam sempitérnam: \Star{} Quorum doctrína fulget Ecclésia, ut sole luna.\end{Resp}
\begin{Resp}\Vsign Sancti per fidem vicérunt regna: operáti sunt justítiam.\end{Resp}
\begin{Resp}\Rsign Quorum doctrína fulget Ecclésia, ut sole luna.\end{Resp}
\begin{Resp}\Vsign Glória Patri, et Fílio, \Star{} et Spirítui Sancto.\end{Resp}
\begin{Resp}\Rsign Quorum doctrína fulget Ecclésia, ut sole luna.\end{Resp}
\switchcolumn
\EN
\begin{Resp}\Rsign These men are saints, whom the Lord hath chosen in love unfeigned, and hath given them glory everlasting. These are they \Star{} By the light of whose teaching the Church is glorified, even as the moon is glorified by the light of the sun.\end{Resp}
\begin{Resp}\Vsign The saints through faith subdued kingdoms, wrought righteousness.\end{Resp}
\begin{Resp}\Rsign By the light of whose teaching the Church is glorified, even as the moon is glorified by the light of the sun.\end{Resp}
\begin{Resp}\Vsign Glory be to the Father, and to the Son, \Star{} and to the Holy Ghost.\end{Resp}
\begin{Resp}\Rsign By the light of whose teaching the Church is glorified, even as the moon is glorified by the light of the sun.\end{Resp}
\switchcolumn*

\LA
\LessonHead{Lectio IX}
\switchcolumn
\EN
\LessonHead{Lesson IX}
\switchcolumn*

\LA
\noindent Quod est fastígium construéndæ fábricæ, quam molímur? quo perventúrum est cacúmen ædifícii? Cito dico, usque ad conspéctum Dei. Vidétis quam excélsum est, quanta res est conspícere Deum. Qui desíderat, et quod dico et quod audit intéllegit. Promíttitur nobis conspéctus Dei, veri Dei, summi Dei. Hoc enim bonum est, vidéntem vidére. Nam, qui colunt falsos deos, fácile illos vident; sed eos vident, qui óculos habent et non vident. Nobis autem promíttitur vísio Dei vivéntis et vidéntis.\par
\switchcolumn
\EN
\noindent What is the roof of the house on which we labour? Whither do its spires rise? I answer you at once; to the presence of God. You see how high it is, yea, what it is to see God. He that will, understandeth what I say, and he heareth. What is promised you is to see God, God, the True, God, the Supreme. Blessed is he who seeth Him by Whom he is seen. Such as worship false gods see them easily, but they see them who have eyes and see not. But unto us it is promised that we shall see that God Who liveth and seeth. (Gen. xvi. 14.)\par
\switchcolumn*

\LA
\TeDeum{Te Deum}
\switchcolumn
\EN
\TeDeum{Te Deum}
\switchcolumn*

\end{paracol}

\par\needspace{10\baselineskip}\addvspace{1.2em}{\centering\small\textsc{Die 27 Februarii} · \textsc{27 February}\par}
\Office{S. Gabrielis a Virgine Perdolente Confessoris}{St. Gabriel of the Sorrowful Virgin, Confessor}{Duplex}
\addcontentsline{toc}{subsection}{Die 27 Februarii · S. Gabrielis a Virgine Perdolente Confessoris\enspace\textit{St. Gabriel of the Sorrowful Virgin, Confessor}}
\begin{paracol}{2}
\LA
\RefLine{Lectiones I–III: ex Commúni Confessoris non pontificis (C5), in 2 loco.}
\switchcolumn
\EN
\RefLine{Lessons I–III: from the Common of a Confessor not a Bishop (C5), set 2.}
\switchcolumn*

\LA
\RefLine{Lectio IX: de Scriptura occurrente.}
\switchcolumn
\EN
\RefLine{Lesson IX: from the Scripture of the day.}
\switchcolumn*

\LA
\LessonHead{Lectio IV}
\switchcolumn
\EN
\LessonHead{Lesson IV}
\switchcolumn*

\LA
\noindent Gabriel, Assisii in Umbria, honesto genere natus, et Franciscus ob seráphici civis memóriam vocatus, egregiam animi índolem a púero osténdit. Adoléscens, cum Spoleti litteris operam daret, inani sæculi spécie et pompa aliquantulum állici visus est. Sed miserentis Dei munere, qui eum ad perfectiónem christianæ vitæ jamdúdum invitabat, cum in morbum incidísset, sæculi vanitátem fastidire cœpit, atque immortalia dumtáxat bona appétere. Quo autem citius Deo vocanti obtemperaret, factum est, ut insignem illam beatíssimæ Vírginis Icónem, solemni pompa extra Spoletinæ ecclésiæ septa delatam intúitus, divini amoris flammam conciperet, simulque Institutum Clericórum a Passióne Jesu amplecti statúeret. Itaque non exíguas difficultates eluctatus, in recessu Morrovallénsi, lúgubrem vestem lætus induit, et Gabriel a Vírgine perdolénte maluit appellari; ad ejusdem gaudiórum et dolórum memóriam perpetuo recoléndam.\par
\switchcolumn
\EN
\noindent Gabriel, born at Assisi in Umbria of a respectable family, and called Francis in memory of his seraphic fellow-townsman, showed from boyhood an excellent disposition of soul. As a youth, when studying letters at Spoleto, he seemed for a time to be allured by the empty beauty and pomp of the world. But by the gift of the merciful God, who had already called him to the perfection of a Christian life when he had fallen sick, he began to tire of the vanity of the world, and to desire immortal treasures alone. But to quicken his obedience to the call of God, it happened that as he saw the celebrated Image of the Blessed Virgin being carried with solemn pomp outside the precincts of the church of Spoleto, he experienced the flame of divine love, and at the same time decided to enter the Institute of the Clerks of the Passion of Jesus. Therefore, after overcoming no slight difficulties, he joyfully donned the somber habit in the secluded place of Morrovalle, and chose to be called Gabriel of Our Lady of Sorrows, to recall forever the memory of her joys and griefs.\par
\switchcolumn*

\LA
\begin{Resp}\Rsign Honéstum fecit illum Dóminus, et custodívit eum ab inimícis, et a seductóribus tutávit illum: \Star{} Et dedit illi claritátem ætérnam.\end{Resp}
\begin{Resp}\Vsign Justum dedúxit Dóminus per vias rectas, et osténdit illi regnum Dei.\end{Resp}
\begin{Resp}\Rsign Et dedit illi claritátem ætérnam.\end{Resp}
\switchcolumn
\EN
\begin{Resp}\Rsign The Lord made him honourable, and defended him from his enemies, and kept him safe from those that lay in wait for him; \Star{} And gave him perpetual glory.\end{Resp}
\begin{Resp}\Vsign He went down with him into the pit, and left him not in bonds.\end{Resp}
\begin{Resp}\Rsign And gave him perpetual glory.\end{Resp}
\switchcolumn*

\LA
\LessonHead{Lectio V}
\switchcolumn
\EN
\LessonHead{Lesson V}
\switchcolumn*

\LA
\noindent In tirocinio, cum regulari observántia et ómnium exercitatióne virtútum cotidie magis eminéret, brevi eo pérvenit, ut absolutæ sanctimoniæ exemplar haberétur non modo a sodálibus, vel provectis, sed étiam ultra cœnobii septa, factus bonus odor Christi in omni loco. Dominicæ passiónis cultor assiduus, in ea meditanda dies noctesque insumebat. In divinam Eucharistiam, quæ ejusdem Passiónis memóriam prodit, incredibili quodam studio ferebátur; qua cum se refíceret, seraphico ardore flagrábat. Nihil autem insígnius quam ejus erga magnam Dei Parentem pietas fuit. Ipsam omni obsequii genere percólere consuevit; sed præsertim confectam afflictamque cruciátibus Jesu tam dolenter contemplari, ut vim lacrimárum profúnderet. Pérdolens Virgo quasi tota ei vitæ rátio fuit, adeptæque ab eo sanctitátis magistra; ita ut inter æquales una fúerit senténtia, ideo excitátum Dei famulum divinitus fuísse, ut cultus Maríæ perdolentis magnum exemplo ejus cáperet increméntum.\par
\switchcolumn
\EN
\noindent In the novitiate, day by day he became conspicuous for regular observance and for the exercise of all the virtues, and in a short time he came to be considered a pattern of perfect holiness, not only by his companions and his seniors, but also beyond the confines of the monastery; he became a sweet odour in Christ in every place. An assiduous devotee of the Lord's Passion, he spent days and nights meditating upon it. He was drawn by unbelievable zeal towards the Holy Eucharist, a memorial of that Passion; and when he nourished himself with it, he burned with seraphic ardour. There was nothing more noticeable than his filial piety towards the great Mother of God. He was accustomed to pay her honour for every type of devotion, but especially to contemplate her stricken and afflicted by the sufferings of Jesus, with such sorrow that he shed floods of tears. The sorrowful Virgin was, as it were, the whole reason of his being, and the teacher of the holiness that he had acquired. As a result all his associates shared the one opinion that this servant of God had been inspired from on high so that the cult of St. Mary of Sorrows through his example might receive a great increase.\par
\switchcolumn*

\LA
\begin{Resp}\Rsign Amávit eum Dóminus, et ornávit eum: stolam glóriæ índuit eum, \Star{} Et ad portas paradísi coronávit eum.\end{Resp}
\begin{Resp}\Vsign Índuit eum Dóminus lorícam fídei, et ornávit eum.\end{Resp}
\begin{Resp}\Rsign Et ad portas paradísi coronávit eum.\end{Resp}
\switchcolumn
\EN
\begin{Resp}\Rsign The Lord loved him and beautified him; He clothed him with a robe of glory, \Star{} And crowned him at the gates of Paradise.\end{Resp}
\begin{Resp}\Vsign The Lord hath put on him the breast-plate of faith, and hath adorned him.\end{Resp}
\begin{Resp}\Rsign And crowned him at the gates of Paradise.\end{Resp}
\switchcolumn*

\LA
\LessonHead{Lectio VI}
\switchcolumn
\EN
\LessonHead{Lesson VI}
\switchcolumn*

\LA
\noindent Inter ceteras virtútes christianam humilitátem et obœdiéntiam maxime diléxit: nam inter omnes se minimum exístimans, abjectíssima quæque ministeria domus cupide affectabat, et antistitum suórum non modo jussa, sed et optáta diligentíssime perficiebat. Idem, refrenátis sensibus et vitæ asperitáte usus, illibátum retinuit florem virginitátis ac plane mundo crucifixus unice Deo vixit, intima Dómini sui fruítus consuetúdine. Ita brevem vitæ cursum tot virtútibus nobilitátum conficiens, Insulæ in Aprutio, caritátis incendio verius quam vi morbi consumptus, divinæque Matris ope recreatus, placidíssimo éxitu ad Superos evolávit, anno millesimo octingentésimo sexagesimo secundo, ætátis suæ vigesimo quarto. Eum deinceps, a Deo miraculis illustrátum, Pius Papa décimus Cælitum beatórum número accensuit. Benedíctus vero décimus quintus, Pontifex maximus, anno millesimo nongentésimo vigesimo, post cónditum Institutum a Passióne ducentésimo, in solemnitate Ascensiónis Dómini, beato júveni Sanctórum honores decrevit; et Pius undecimus ejus Offícium et Missam ad universam Ecclésiam extendit.\par
\switchcolumn
\EN
\noindent Among other virtues, he especially loved Christian humility and obedience; for he considered himself the least of all. He therefore strove eagerly to do all the most menial work of the house, and he most diligently performed, not only the direct commands, but even the unexpressed wishes of his superiors. Curbing his senses, and accustoming himself to a life of austerity, he retained unfaded the flower of his virginity, and completely crucified to the world, he lived to God alone, enjoying an intimate familiarity with his Lord. And so, at Isola in the Abruzzi, filling the short span of his life with so many noble virtues, consumed by the fire of charity rather than by disease, and refreshed by the aid of the Mother of God, his soul flew to heaven in a most peaceful journey in the year 1862, at the age of twenty-four. Then, as he had been made illustrious by God through miracles, Pope Pius X added him to the number of the Blessed in heaven. Likewise, the Supreme Pontiff Benedict XV in 1920, two hundred years after the foundation of the Institute of the Passion, on the feast of the Ascension of the Lord, decreed the honours of the Saints to the blessed youth; and Pius XI extended his Office and Mass to the Universal Church.\par
\switchcolumn*

\LA
\begin{Resp}\Rsign Iste homo perfécit ómnia quæ locútus est ei Deus, et dixit ad eum: Ingrédere in réquiem meam: \Star{} Quia te vidi justum coram me ex ómnibus géntibus.\end{Resp}
\begin{Resp}\Vsign Iste est, qui contémpsit vitam mundi, et pervénit ad cæléstia regna.\end{Resp}
\begin{Resp}\Rsign Quia te vidi justum coram me ex ómnibus géntibus.\end{Resp}
\begin{Resp}\Vsign Glória Patri, et Fílio, \Star{} et Spirítui Sancto.\end{Resp}
\begin{Resp}\Rsign Quia te vidi justum coram me ex ómnibus géntibus.\end{Resp}
\switchcolumn
\EN
\begin{Resp}\Rsign This is he which did according unto all that God commanded him; and God said unto him: Enter thou into My rest, \Star{} For thee have I seen righteous before Me among all people.\end{Resp}
\begin{Resp}\Vsign This is he which loved not his life in this world, and is come unto an everlasting kingdom.\end{Resp}
\begin{Resp}\Rsign For thee have I seen righteous before Me among all people.\end{Resp}
\begin{Resp}\Vsign Glory be to the Father, and to the Son, \Star{} and to the Holy Ghost.\end{Resp}
\begin{Resp}\Rsign For thee have I seen righteous before Me among all people.\end{Resp}
\switchcolumn*

\LA
\LessonHead{Lectio VII}
\switchcolumn
\EN
\LessonHead{Lesson VII}
\switchcolumn*

\LA
\LessonTitle{Léctio sancti Evangélii secúndum Marcum}
\switchcolumn
\EN
\LessonTitle{From the Holy Gospel according to Mark}
\switchcolumn*

\LA
\Cite{Marc 10:13-21}
\switchcolumn
\EN
\Cite{Mark 10:13-21}
\switchcolumn*

\LA
\noindent In illo témpore: Offerébant Jesu párvulos, ut tángeret illos: discípuli autem comminabántur offeréntibus. Et réliqua.\par
\switchcolumn
\EN
\noindent At that time: They brought young children to Jesus, that he should touch them: and his disciples rebuked those that brought them. And so on, and that which followeth.\par
\switchcolumn*

\LA
\LessonTitle{Homilía sancti Bedæ Venerábilis Presbýteri}
\switchcolumn
\EN
\LessonTitle{A Homily by St. Venerable Bede the Priest}
\switchcolumn*

\LA
\Source{Commentárium in Marcum 10:13-21}
\switchcolumn
\EN
\Source{Commentarium in Marcum 10:13-21}
\switchcolumn*

\LA
\noindent Ait discípulis Jesus: Sínite párvulos veníre ad me, et ne prohibuéritis eos; talium enim est regnum Dei. Significanter dixit: Talium est; non: Istórum; ut osténderet, non ætátem regnare, sed mores; et his, qui similem haberent innocéntiam et simplicitátem, præmium repromitti: Apóstolo quoque in eandem senténtiam congruénte: Fratres, nolíte fíeri púeri sensibus; sed malítia párvuli estote, sensu autem ut perfecti sitis. Amen, dico vobis: quisquis non recéperit regnum Dei velut párvulus, non intrábit in illud. Sicut puer non persevérat in iracúndia, non læsus méminit, non videns pulchram mulierem delectátur, non aliud cogitat, aliud lóquitur; sic et vos, nisi talem habuéritis innocéntiam et animi puritátem, regnum cælórum non poteritis intrare. Aliter, regnum Dei, id est doctrinam Evangélii, sicut párvuli accípere jubemur; quia quómodo párvulus in discéndo non contradicit doctóribus neque ratiónes et verba componit advérsum eos resistens, sed fideliter suscipit quod docétur et cum metu obtemperat et quiescit; ita et nos, in obœdiéndo simpliciter et sine ulla rectractatióne verbis Dómini, fácere debemus. Et complexans eos, et impónens manus super illos, benedicebat eos. Complexus benedicit párvulos, ut húmiles spíritu sua benedictióne, grátia et dilectióne dignos esse significet.\par
\switchcolumn
\EN
\noindent Jesus saith unto his disciples: Suffer the little children to come unto me, and forbid them not: for of such is the kingdom of God. It is noteworthy that he saith: Of such; not: Of these. That is, he is concerned with conduct, not with age; for the reward is promised to all such as are like unto children, not in age, but in innocence and simplicity; as saith the Apostle: Be not children in understanding; howbeit, in malice be ye children, but in understanding be men. Verily I say unto you, Whosoever shall not receive the kingdom of God as a little child, he shall not enter therein. A child doth not for long remain angry or remember an injury; he taketh no lustful pleasure in looking at a beautiful woman; he doth not think one thing and say another. In like manner, ye also, unless ye have this innocence and purity of soul, cannot enter the kingdom of heaven. We are commanded to receive the kingdom of God (that is, the teaching of the Gospel) like as children, because they, when learning, do not contradict their teacher, nor adduce reasons and arguments against them, but trustingly accept what they are taught in respectful silence and obedience. And the Lord took them up in his arms, put his hands upon them, and blessed them. He blesseth the children by taking them up in his arms, to shew that such as are poor in spirit do merit his blessing, grace, and love.\par
\switchcolumn*

\LA
\begin{Resp}\Rsign Iste est qui ante Deum magnas virtútes operátus est, et de omni corde suo laudávit Dóminum: \Star{} Ipse intercédat pro peccátis ómnium populórum.\end{Resp}
\begin{Resp}\Vsign Ecce homo sine queréla, verus Dei cultor, ábstinens se ab omni ópere malo, et pérmanens in innocéntia sua.\end{Resp}
\begin{Resp}\Rsign Ipse intercédat pro peccátis ómnium populórum.\end{Resp}
\switchcolumn
\EN
\begin{Resp}\Rsign This is he which wrought great wonders before God, and praised the Lord with all his heart. \Star{} May he pray for all people, that their sins may be forgiven unto them.\end{Resp}
\begin{Resp}\Vsign Behold a man without blame, a worshipper of God in truth, keeping himself clean from every evil work, and abiding still in his innocency.\end{Resp}
\begin{Resp}\Rsign May he pray for all people, that their sins may be forgiven unto them.\end{Resp}
\switchcolumn*

\LA
\LessonHead{Lectio VIII}
\switchcolumn
\EN
\LessonHead{Lesson VIII}
\switchcolumn*

\LA
\noindent Et, cum egréssus esset in viam, procurrens quidam, genu flexo ante eum, rogábat eum: Magister bone, quid fáciam, ut vitam ætérnam percipiam? Audierat, credo, iste quæsítor vitæ ætérnæ a Dómino, tantum eos, qui parvulórum velint esse similes, dignos esse introítu regni cæléstis: atque ideo curam gerens tractátus certioris, poscit sibi non per parábolas, sed aperte, quibus óperum meritis vitam ætérnam cónsequi possit, exponi. Jesus autem dixit ei: Præcepta nosti. Hæc est puerílis innocéntiæ castitas, quæ nobis imitanda propónitur, si regnum Dei volumus intrare. At ille respondens, ait ille: Magister, hæc ómnia observávi a juventúte mea. Non est putandus homo iste vel voto tentántis (ut quidam putavére) Dóminum interrogasse, vel de sua esse vita mentitus, cum se legis mandáta custodisse dicebat; sed, simpliciter, ut víxerit esse conféssus. Quia si mendacii aut simulatiónis noxa reus tenerétur, nequáquam intuitus arcana cordis ejus, eum diligere dicerétur Jesus.\par
\switchcolumn
\EN
\noindent And as he was going forth on his journey, a certain man came running up, and fell on his knees before him, and asked him: Good Master, what shall I do that I may inherit eternal life? Doubtless this seeker after eternal life was one of those who had heard the Lord say that whoso shall not receive the kingdom of God as little child shall not enter therein; and therefore, taking the safer course, he asketh for himself not in parables, but simply and openly, for teaching concerning the meritorious works by which eternal life is to be gained. Whereupon Christ said: Thou knowest the commandments. This is the chastity of childlike innocence, which is proposed for our imitation, if we wish to enter the kingdom of God. But he answering, said unto him: Master, all these have I observed from my youth. This man should not be considered as having asked the Lord with a view to tempting him (as some have thought), or to have lied about his life when he said he had kept the commandments of the law; but that rather he was trying to explain how he had hitherto lived. Because, if he be considered guilty of lies and pretence, it would not have been said of Jesus, who looketh on the secrets of the heart, that he, beholding him, loved him.\par
\switchcolumn*

\LA
\begin{Resp}\Rsign Sint lumbi vestri præcíncti, et lucérnæ ardéntes in mánibus vestris: \Star{} Et vos símiles homínibus exspectántibus dóminum suum, quando revertátur a núptiis.\end{Resp}
\begin{Resp}\Vsign Vigiláte ergo, quia nescítis qua hora Dóminus vester ventúrus sit.\end{Resp}
\begin{Resp}\Rsign Et vos símiles homínibus exspectántibus dóminum suum, quando revertátur a núptiis.\end{Resp}
\begin{Resp}\Vsign Glória Patri, et Fílio, \Star{} et Spirítui Sancto.\end{Resp}
\begin{Resp}\Rsign Et vos símiles homínibus exspectántibus dóminum suum, quando revertátur a núptiis.\end{Resp}
\switchcolumn
\EN
\begin{Resp}\Rsign Let your loins be girded about, and your lights burning; \Star{} And ye yourselves like unto men that wait for their lord, when he will return from the wedding.\end{Resp}
\begin{Resp}\Vsign Watch therefore, for ye know not what hour your Lord doth come.\end{Resp}
\begin{Resp}\Rsign And ye yourselves like unto men that wait for their lord, when he will return from the wedding.\end{Resp}
\begin{Resp}\Vsign Glory be to the Father, and to the Son, \Star{} and to the Holy Ghost.\end{Resp}
\begin{Resp}\Rsign And ye yourselves like unto men that wait for their lord, when he will return from the wedding.\end{Resp}
\switchcolumn*

\LA
\LessonHead{Lectio IX}
\switchcolumn
\EN
\LessonHead{Lesson IX}
\switchcolumn*

\LA
\noindent Díligit enim Dóminus eos, qui mandáta legis, quamvis minora, custódiunt: sed nihilóminus, quod in lege minus fuerat, iis qui perfecti esse desiderant, osténdit, quia non venit solvere legem aut Prophétas, sed adimplere. Ad quam profecto adimpletiónem pertinet, quod hic consequenter adjungitur: Vade, quæcúmque habes vende et da paupéribus, et habebis thesáurum in cælo, et veni, séquere me. Quicúmque perfectus esse volúerit, debet vendere quæ habet; et non ex parte véndere, sicut Ananías et Saphira, sed totum véndere: et cum vendiderit, dare omne paupéribus, et sic sibi præparáre thesáurum in regno cælórum. Nec hoc ad perfectiónem sufficit, nisi, post contemptas divítias, Salvatórem sequátur; id est relictis malis, fáciat bona. Facílius enim sæculum contemnitur quam volúntas. Multi divítias relinquéntes, Dóminum non sequúntur. Sequitur autem Dóminum, qui imitator ejus est et per vestígia illíus graditur. Qui enim dicit se in Christo credere, debet, quómodo ille ambulávit, et ipse ambuláre.\par
\switchcolumn
\EN
\noindent For the love of the Lord goeth out to all those who not only keep the commandments of the law, but have regard even for things which are not of strict precept. Which latter things he who came not to destroy the Law and the Prophets, but to fulfill them, showeth forth unto such as would be perfect, saying: Sell whatsoever thou hast, and give to the poor, and thou shalt have treasure in heaven; and come, follow me. Whosoever would be perfect must sell all (not merely in part, as did Ananias and Sapphira), and thereafter give all to the poor, and thus prepare for himself a treasure in the kingdom of heaven. But contempt for riches is not enough to achieve perfection. He must also follow the Saviour, leaving evil and doing good. For it is easier to spurn the world than the will; and many spurn their riches, but fail thereafter to follow the Lord. Whosoever would follow him must be his imitator, and walk in his footsteps. For whoso saith he believeth in Christ ought also to walk as Christ walked.\par
\switchcolumn*

\LA
\TeDeum{Te Deum}
\switchcolumn
\EN
\TeDeum{Te Deum}
\switchcolumn*

\end{paracol}

\SeasonTitle{Mensis Martius}{March}

\addcontentsline{toc}{section}{Mensis Martius\enspace\textit{March}}

\par\needspace{10\baselineskip}\addvspace{1.2em}{\centering\small\textsc{Die 4 Martii} · \textsc{4 March}\par}
\Office{S. Casimiri Confessoris}{St. Casimir, Confessor}{Semiduplex}
\addcontentsline{toc}{subsection}{Die 4 Martii · S. Casimiri Confessoris\enspace\textit{St. Casimir, Confessor}}
\begin{paracol}{2}
\LA
\RefLine{Lectiones I–III: ex Commúni Confessoris non pontificis (C5), in 2 loco.}
\switchcolumn
\EN
\RefLine{Lessons I–III: from the Common of a Confessor not a Bishop (C5), set 2.}
\switchcolumn*

\LA
\RefLine{Lectiones VII–VIII: ex Commúni Confessoris non pontificis (C5).}
\switchcolumn
\EN
\RefLine{Lessons VII–VIII: from the Common of a Confessor not a Bishop (C5).}
\switchcolumn*

\LA
\RefLine{Lectio IX: de Scriptura occurrente.}
\switchcolumn
\EN
\RefLine{Lesson IX: from the Scripture of the day.}
\switchcolumn*

\LA
\LessonHead{Lectio IV}
\switchcolumn
\EN
\LessonHead{Lesson IV}
\switchcolumn*

\LA
\noindent Casimírus, patre Casimiro, matre Elisabetha Austriaca, Poloniæ regibus ortus, a pueritia sub optimis magistris, pietate et bonis artibus instructus, juveniles artus aspero domabat cilicio, et assiduis extenuabat jejuniis. Regii spreta lecti mollitie, dura cubabat humo, et clam intempesta nocte, præ foribus templorum pronus in terra divinam exorabat clementiam. In Christi contemplanda passione assiduus, Missarum solemniis adeo erecta in Deum mente solebat adesse, ut extra se rapi videretur.\par
\switchcolumn
\EN
\noindent This Casimir was the son of Casimir III, King of Poland, by Elizabeth of Austria, his wife, (and was born upon the 3rd day of October, in the year 1458.) From his childhood he was taught by the best masters, and was trained in all godliness and good learning. While he was still a boy he wore rough haircloth, and chastened himself with much fasting. He forsook the softness of his princely bed, and lay upon the hard ground, and on stormy nights he would go out secretly and prostrate himself before the doors of the churches, crying to God for mercy. He was unwearied in contemplating the Passion of Christ, and when he was present at Mass, so profound was his recollection, that he seemed to be altogether beside himself.\par
\switchcolumn*

\LA
\begin{Resp}\Rsign Honéstum fecit illum Dóminus, et custodívit eum ab inimícis, et a seductóribus tutávit illum: \Star{} Et dedit illi claritátem ætérnam.\end{Resp}
\begin{Resp}\Vsign Justum dedúxit Dóminus per vias rectas, et osténdit illi regnum Dei.\end{Resp}
\begin{Resp}\Rsign Et dedit illi claritátem ætérnam.\end{Resp}
\switchcolumn
\EN
\begin{Resp}\Rsign The Lord made him honourable, and defended him from his enemies, and kept him safe from those that lay in wait for him; \Star{} And gave him perpetual glory.\end{Resp}
\begin{Resp}\Vsign He went down with him into the pit, and left him not in bonds.\end{Resp}
\begin{Resp}\Rsign And gave him perpetual glory.\end{Resp}
\switchcolumn*

\LA
\LessonHead{Lectio V}
\switchcolumn
\EN
\LessonHead{Lesson V}
\switchcolumn*

\LA
\noindent Catholicam promovere fidem summopere studuit, et Ruthenorum schisma abolere: quapropter Casimírum patrem induxit, ut legem ferret, ne schismatici nova templa construerent, nec vetera collabentia restaurarent. Erga pauperes et calamitatibus oppressos beneficus et misericors, patris et defensoris egenorum nomen obtinuit. Virginitatem, quam ab incunabulis servavit illæsam, sub extremo vitæ termino fortiter asseruit, dum gravi pressus infirmitate, mori potius quam castitatis jacturam ex medicorum consilio subire, constanter decrevit.\par
\switchcolumn
\EN
\noindent He made the propagation of the Catholic faith one of the chief works of his life, and strove hard against the schism in Ruthenia. He persuaded his father to forbid by law that the schismatics should build any new churches, or repair the existing ones when they fell into decay. So great was his liberality and tenderness toward the needy and the afflicted, that he came to be called the father and guardian of the poor. From his infancy he never soiled his purity, and in his last illness, when his physicians advised him to seek for relief from his grievous sufferings by the sacrifice of his chastity, he cheerfully determined rather to die.\par
\switchcolumn*

\LA
\begin{Resp}\Rsign Amávit eum Dóminus, et ornávit eum: stolam glóriæ índuit eum, \Star{} Et ad portas paradísi coronávit eum.\end{Resp}
\begin{Resp}\Vsign Índuit eum Dóminus lorícam fídei, et ornávit eum.\end{Resp}
\begin{Resp}\Rsign Et ad portas paradísi coronávit eum.\end{Resp}
\switchcolumn
\EN
\begin{Resp}\Rsign The Lord loved him and beautified him; He clothed him with a robe of glory, \Star{} And crowned him at the gates of Paradise.\end{Resp}
\begin{Resp}\Vsign The Lord hath put on him the breast-plate of faith, and hath adorned him.\end{Resp}
\begin{Resp}\Rsign And crowned him at the gates of Paradise.\end{Resp}
\switchcolumn*

\LA
\LessonHead{Lectio VI}
\switchcolumn
\EN
\LessonHead{Lesson VI}
\switchcolumn*

\LA
\noindent Consummatus in brevi, virtutibus et meritis plenus, prænuntiato mortis die, inter sacerdotum et religiosorum choros' spiritum Deo reddidit, anno ætatis vigesimo quinto. Corpus Vilnam delatum multis claret miraculis. Etenim præterquam quod puella defuncta vitam, cæci visum, claudi gressum, et varii infirmi sanitatem ad ejus sepulchrum recuperarunt; Lithuanis exiguo numero ad potentissimi hostis insperatam irruptionem trepidantibus in aëre apparens, insignem tribuit victoriam. Quibus permotus Leo decimus eumdem Sanctorum catalogo adscripsit.\par
\switchcolumn
\EN
\noindent Being made perfect in a short space, and full of piety and good works, he foretold the day of his own death, and, gathering round him a choir of priests and monks, he rendered his soul into the hands of God Whom they were praising, \textasciitilde{}(upon the 4th day of March, in the year of our Lord 1482, and) the 25th of his own age. His body was carried to Wilna, where many miracles are reputed to have been wrought around it. At his grave a dead girl is said to have received her life again, blind men their sight, cripples the power of walking, and many sick folk health. Moreover, on an occasion when the Lithuanians in scanty numbers were exposed to the shock of a powerful enemy, they believed that he appeared in the air, and gave them the signal victory which they won. On the assurance of these things, Leo X. was moved to add his name to those of the Saints.\par
\switchcolumn*

\LA
\begin{Resp}\Rsign Iste homo perfécit ómnia quæ locútus est ei Deus, et dixit ad eum: Ingrédere in réquiem meam: \Star{} Quia te vidi justum coram me ex ómnibus géntibus.\end{Resp}
\begin{Resp}\Vsign Iste est, qui contémpsit vitam mundi, et pervénit ad cæléstia regna.\end{Resp}
\begin{Resp}\Rsign Quia te vidi justum coram me ex ómnibus géntibus.\end{Resp}
\begin{Resp}\Vsign Glória Patri, et Fílio, \Star{} et Spirítui Sancto.\end{Resp}
\begin{Resp}\Rsign Quia te vidi justum coram me ex ómnibus géntibus.\end{Resp}
\switchcolumn
\EN
\begin{Resp}\Rsign This is he which did according unto all that God commanded him; and God said unto him: Enter thou into My rest, \Star{} For thee have I seen righteous before Me among all people.\end{Resp}
\begin{Resp}\Vsign This is he which loved not his life in this world, and is come unto an everlasting kingdom.\end{Resp}
\begin{Resp}\Rsign For thee have I seen righteous before Me among all people.\end{Resp}
\begin{Resp}\Vsign Glory be to the Father, and to the Son, \Star{} and to the Holy Ghost.\end{Resp}
\begin{Resp}\Rsign For thee have I seen righteous before Me among all people.\end{Resp}
\switchcolumn*

\end{paracol}

\par\needspace{10\baselineskip}\addvspace{1.2em}{\centering\small\textsc{Die 6 Martii} · \textsc{6 March}\par}
\Office{Ss. Perpetuæ et Felicitatis Martyrum}{Sts. Perpetua and Felicity, Martyrs}{Duplex}
\addcontentsline{toc}{subsection}{Die 6 Martii · Ss. Perpetuæ et Felicitatis Martyrum\enspace\textit{Sts. Perpetua and Felicity, Martyrs}}
\begin{paracol}{2}
\LA
\RefLine{Lectiones I–III: ex Commúni plurimarum non Virginum Martyrum (C7b).}
\switchcolumn
\EN
\RefLine{Lessons I–III: from the Common of multiple widowed Martyrs (C7b).}
\switchcolumn*

\LA
\RefLine{Lectiones VII–VIII: ex Commúni plurimarum non Virginum Martyrum (C7b).}
\switchcolumn
\EN
\RefLine{Lessons VII–VIII: from the Common of multiple widowed Martyrs (C7b).}
\switchcolumn*

\LA
\RefLine{Lectio IX: de Scriptura occurrente.}
\switchcolumn
\EN
\RefLine{Lesson IX: from the Scripture of the day.}
\switchcolumn*

\LA
\LessonHead{Lectio IV}
\switchcolumn
\EN
\LessonHead{Lesson IV}
\switchcolumn*

\LA
\noindent Perpetua et Felicitas, in persecutione Severi imperatoris, in Africa, una cum Revocato, Saturnino et Secundulo comprehensæ sunt et in tenebricosum carcerem detrusæ; quibus ultra adjunctus est Satyrus. Erant adhuc catechumenæ, sed paulo post baptizatæ sunt. Paucis diebus interjectis, e carcere ad forum deductæ cum sociis, post gloriosam confessionem, ab Hilarione procuratore damnantur ad bestias. Inde hilares descendunt ad carcerem, ubi variis visionibus recreantur et ad martyrii palmam accenduntur. Perpetuam, nec patris senio confecti iteratæ preces et lacrimæ, nec erga filium infantem pendentem ad ubera maternus amor, nec supplicii atrocitas, a Christi fide dimovere umquam potuerunt\par
\switchcolumn
\EN
\noindent Perpetua and Felicity were arrested during the persecution of the emperor Severus in Africa, together with Revocatus, Saturninus and Secundulus, and were cast into a dark dungeon, where Satyrus was added to their company. They were as yet catechumens but shortly afterwards they were baptized. After a few days they and their companions were led forth from prison to the court, and, after a glorious profession of faith, were condemned by the procurator Hilarion to the beasts. Thereupon they went down to the prison rejoicing, and there they were refreshed with many visions, and fired with longing for the martyr's palm. Neither the repeated prayers and tears of Perpetua's father, a man of extreme old age, nor her motherly love for her baby son still at the breast, nor the atrocity of the torture, could shake her faith in Christ.\par
\switchcolumn*

\LA
\begin{Resp}\Rsign Propter veritátem, et mansuetúdinem, et justítiam: \Star{} Et dedúcet te mirabíliter déxtera tua.\end{Resp}
\begin{Resp}\Vsign Spécie tua et pulchritúdine tua inténde, próspere procéde, et regna.\end{Resp}
\begin{Resp}\Rsign Et dedúcet te mirabíliter déxtera tua.\end{Resp}
\switchcolumn
\EN
\begin{Resp}\Rsign Because of truth, and meekness, and righteousness; \Star{} And thy right hand shall lead thee wonderfully.\end{Resp}
\begin{Resp}\Vsign In thy comeliness, and thy beauty, go forward, fare prosperously, and reign.\end{Resp}
\begin{Resp}\Rsign And thy right hand shall lead thee wonderfully.\end{Resp}
\switchcolumn*

\LA
\LessonHead{Lectio V}
\switchcolumn
\EN
\LessonHead{Lesson V}
\switchcolumn*

\LA
\noindent Felicitas vero instante spectaculo die, cum octo jam menses prægnans esset, in magno erat luctu, ne differretur; leges quippe vetabant prægnantes supplicio affici. At precibus commartyrum accelerato partu, enixa est filiam. Cumque in partu laborans doleret, ait illi quidam de custodibus: Quæ sic modo doles, quid facies objecta bestiis? Cui illa: Modo ego patior; illic autem alius erit in me, qui patietur pro me, quia et ego pro illo passura sum.\par
\switchcolumn
\EN
\noindent As the day of the spectacle came close, Felicity was afflicted with great sorrow lest it should be put off, since she was eight months with child; and the law forbade expectant mothers to be put to the torture. But at the prayers of her fellow-martyrs her delivery was hastened, and she gave birth to a daughter. While she was groaning in the pains of childbirth, one of the jailers said to her: What wilt thou do when thrown to the beasts, if thou groanest thus now? She replied: Now it is I who suffer; but then Another will be within me, who will suffer on my behalf, seeing that it is for him that I am to suffer.\par
\switchcolumn*

\LA
\begin{Resp}\Rsign Dilexísti justítiam, et odísti iniquitátem: \Star{} Proptérea unxit te Deus, Deus tuus, óleo lætítiæ.\end{Resp}
\begin{Resp}\Vsign Propter veritátem, et mansuetúdinem, et justítiam.\end{Resp}
\begin{Resp}\Rsign Proptérea unxit te Deus, Deus tuus, óleo lætítiæ.\end{Resp}
\switchcolumn
\EN
\begin{Resp}\Rsign Thou hast loved righteousness, and hated iniquity; \Star{} Therefore God, thy God, hath anointed thee with the oil of gladness.\end{Resp}
\begin{Resp}\Vsign Because of truth, and meekness, and righteousness.\end{Resp}
\begin{Resp}\Rsign Therefore God, thy God, hath anointed thee with the oil of gladness.\end{Resp}
\switchcolumn*

\LA
\LessonHead{Lectio VI}
\switchcolumn
\EN
\LessonHead{Lesson VI}
\switchcolumn*

\LA
\noindent In amphitheatrum, toto inspectante populo, producuntur tandem generosæ mulieres, Nonis Martii, ac primum flagellis cæduntur. Tunc a ferocissima vacca aliquamdiu jactatæ, plagis concisæ et in terram elisæ sunt. Demum cum sociis, qui a variis bestiis vexati fuerant, gladiorum ictibus conficiuntur. Harum sanctarum Mártyrum festum Pius decimus Pontifex maximus ad ritum duplicem pro universa Ecclesia evexit ac diei sextæ Martii adsignari mandavit.\par
\switchcolumn
\EN
\noindent At length the noble-hearted women were brought into the amphitheatre, in the sight of all the people, on the 7th day of March, and were first beaten with scourges. Then they were tossed for some time by a ferocious cow, beaten with lashes, and dashed on the ground. Lastly, together with their companions, who had been attacked by various beasts, they were slain with blows of the sword. Pope Pius X raised the feast of these holy Martyrs to the rite of a double for the Universal Church, and ordered it to be kept on March 6th.\par
\switchcolumn*

\LA
\begin{Resp}\Rsign Fallax grátia, et vana est pulchritúdo: \Star{} Múlier timens Dóminum, ipsa laudábitur.\end{Resp}
\begin{Resp}\Vsign Date ei de fructu mánuum suárum, et laudent eam in portis ópera ejus.\end{Resp}
\begin{Resp}\Rsign Múlier timens Dóminum, ipsa laudábitur.\end{Resp}
\begin{Resp}\Vsign Glória Patri, et Fílio, \Star{} et Spirítui Sancto.\end{Resp}
\begin{Resp}\Rsign Múlier timens Dóminum, ipsa laudábitur.\end{Resp}
\switchcolumn
\EN
\begin{Resp}\Rsign Favour is deceitful, and beauty is vain \Star{} A woman that feareth God she shall be praised.\end{Resp}
\begin{Resp}\Vsign Give her of the fruit of her hands, and let her own works praise her in the gates.\end{Resp}
\begin{Resp}\Rsign A woman that feareth God she shall be praised.\end{Resp}
\begin{Resp}\Vsign Glory be to the Father, and to the Son, \Star{} and to the Holy Ghost.\end{Resp}
\begin{Resp}\Rsign A woman that feareth God she shall be praised.\end{Resp}
\switchcolumn*

\end{paracol}

\par\needspace{10\baselineskip}\addvspace{1.2em}{\centering\small\textsc{Die 7 Martii} · \textsc{7 March}\par}
\Office{S. Thomæ de Aquino Confessoris et Ecclesiæ Doctoris}{St. Thomas Aquinas, Confessor and Doctor of the Church}{Duplex}
\addcontentsline{toc}{subsection}{Die 7 Martii · S. Thomæ de Aquino Confessoris et Ecclesiæ Doctoris\enspace\textit{St. Thomas Aquinas, Confessor and Doctor of the Church}}
\begin{paracol}{2}
\LA
\RefLine{Lectiones I–III: ex Commúni Doctoris non Pontificis (C5a).}
\switchcolumn
\EN
\RefLine{Lessons I–III: from the Common of a Doctor not a Bishop (C5a).}
\switchcolumn*

\LA
\RefLine{Lectiones VII–VIII: ex Commúni Doctoris non Pontificis (C5a).}
\switchcolumn
\EN
\RefLine{Lessons VII–VIII: from the Common of a Doctor not a Bishop (C5a).}
\switchcolumn*

\LA
\RefLine{Lectio IX: de Scriptura occurrente.}
\switchcolumn
\EN
\RefLine{Lesson IX: from the Scripture of the day.}
\switchcolumn*

\LA
\LessonHead{Lectio IV}
\switchcolumn
\EN
\LessonHead{Lesson IV}
\switchcolumn*

\LA
\noindent Præclárum christiáni orbis decus et Ecclésiæ lumen, beatíssimus vir Thomas, Landulpho comite Aquinate et Theodora Neapolitana, nobílibus paréntibus natus, futuræ in Deiparam devotiónis afféctum adhuc infantulus osténdit. Nam chártulam ab eo invéntam, in qua salutátio Angelica scripta erat, frustra adniténte nutrice, compressa manu valide retinuit, et a matre per vim abreptam, ploratu et gestu repétiit, ac mox redditam deglutívit. Quintum annum agens, monachis sancti Benedicti Cassinátibus custodiendus traditur. Inde Neapolim studiórum causa missus, jam adoléscens Fratrum Prædicatórum Ordinem suscépit. Sed matre ac frátribus id indigne feréntibus, Lutétiam Parisiórum mittitur. Quem fratres in itinere per vim raptum, in arcem castri sancti Joánnis perducunt: ubi varie exagitatus, ut sanctum proponsitum mutaret, mulierem etiam, quæ ad labefactándam ejus constantiam introducta fuerat, titióne fugávit. Mox beátus juvenis, flexis genibus ante signum crucis orans, ibique somno correptus, per quietem sentire visus est sibi ab Angelis constringi lumbos; quo ex témpore omni póstea libidinis sensu cáruit. Soróribus, quæ ut eum a pio consílio removerent, in castrum venerant, persuasit, ut, contemptis curis sæcularibus, ad exercitatiónem cæléstis vitæ se conferrent.\par
\switchcolumn
\EN
\noindent That splendid adornment of the Christian world and light of the Church, blessed Thomas of Aquino, was the son of Landulph, Earl of Aquino, and Theodora of Naples, his wife, being nobly descended on both sides. (He was born in the year of salvation 1226,) and even as an infant gave token of the love which he afterwards bore to the Mother of God. He found a little bit of paper upon which was written the Angelic Salutation, and held it firm in his hand in spite of the efforts of his wet-nurse; his mother took it away by force, but he cried and stretched out for it, and when she gave it back to him, he swallowed it. When he was only four years old, he was given into the keeping of the Benedictine monks of Monte Cassino. He was thence sent to Naples to study, and there, while very young, entered the Order of Friars Preachers. This displeased his mother and brothers, and he left Naples for Paris. When he was on his journey his brothers met him, and carried him off by force to the castle of Monte San Giovanni, where they imprisoned him in the keep. Here they used every means to break him of his intention, and at last brought a woman into his room to try to overcome his purity. The lad drove her out with a fire-brand. When he was alone he knelt down before the figure of the Cross, and there he fell asleep. As he slept, it seemed to him that angels came and girded his loins and from this time he never felt the least sexual inclination. His sisters came to the castle to beseech him to give up his purpose of leaving the world, but he so worked on them by his godly exhortations, that both of them ever after set no value on earthly things, and busied themselves rather with heavenly.\par
\switchcolumn*

\LA
\begin{Resp}\Rsign Honéstum fecit illum Dóminus, et custodívit eum ab inimícis, et a seductóribus tutávit illum: \Star{} Et dedit illi claritátem ætérnam.\end{Resp}
\begin{Resp}\Vsign Justum dedúxit Dóminus per vias rectas, et osténdit illi regnum Dei.\end{Resp}
\begin{Resp}\Rsign Et dedit illi claritátem ætérnam.\end{Resp}
\switchcolumn
\EN
\begin{Resp}\Rsign The Lord made him honourable, and defended him from his enemies, and kept him safe from those that lay in wait for him; \Star{} And gave him perpetual glory.\end{Resp}
\begin{Resp}\Vsign He went down with him into the pit, and left him not in bonds.\end{Resp}
\begin{Resp}\Rsign And gave him perpetual glory.\end{Resp}
\switchcolumn*

\LA
\LessonHead{Lectio V}
\switchcolumn
\EN
\LessonHead{Lesson V}
\switchcolumn*

\LA
\noindent Emissus e castro per fenestram, Neapolim redúcitur; unde Romam, póstea Parisium a fratre Joánne Theutónico, ordinis Prædicatórum generali magistro, ductus, Alberto Magno doctore, philosophíæ ac theologíæ operam dedit. Viginti quinque annos natus, magister est appellatus, publiceque philósophos ac Theólogos summa cum laude est interpretatus. Numquam se lectióni aut scriptióni dedit, nisi post oratiónem. In difficultátibus locórum sacræ Scripturæ, ad oratiónem jejunium adhibebat. Quin étiam sodali suo fratri Regínaldo dicere solebat, quidquid sciret non tam studio aut labóre suo peperisse, quam divínitus tráditum accepisse. Neapoli cum ad imáginem Crucifixi vehementius oraret, hanc vocem audívit: Bene scripsísti de me, Thoma; quam ergo mercedem accípies? Cui ille: Non aliam, Dómine, nisi teípsum. Collatiónes Patrum assidue pervolutábat; et nullum fuit scriptórum genus, in quo non esset diligentíssime versatus. Scripta ejus et multitúdine, et varietáte, et facilitate explicándi res difficiles adeo excellunt, ut uberrima atque incorrupta illíus doctrina, cum revelátis veritátibus mire consentiens, aptíssima sit ad ómnium témporum erróres pervincendos.\par
\switchcolumn
\EN
\noindent Being let down from a window, Thomas escaped out of the castle of Monte San Giovanni, and returned to Naples. Thence he went first to Rome, and then to Paris, in company of Brother John the German, then Master-General of the Friars Preachers. At Paris he studied Philosophy and Theology under Albert the Great Doctor. At the age of twenty-five years he took the degree of Master, and gave public disquisitions on the Philosophers and Theologians with great distinction. He never set himself to read or write till he had first prayed, and when he was about to take in hand a hard passage of the Holy Scriptures, he fasted also. Hence he was wont to say to Brother Reginald his comrade, that whatever he knew, he had learnt, not so much from his own labour and study, as from the inspiration of God. At Naples he was once kneeling in very earnest prayer before an image of Christ Crucified, when he heard a voice which said Thomas, thou hast written well of Me what reward wilt thou that I give thee? He answered: Lord, thyself. He studied most carefully the works of the Fathers, and there was no kind of author in which he was not well read. His own writings are so wonderful, both because of their number, their variety, and the clearness of his explanations of hard things, that his rich and pure teaching, marvellously consonant with revealed truth, is an admirable antidote for the errors of all times.\par
\switchcolumn*

\LA
\begin{Resp}\Rsign Amávit eum Dóminus, et ornávit eum: stolam glóriæ índuit eum, \Star{} Et ad portas paradísi coronávit eum.\end{Resp}
\begin{Resp}\Vsign Índuit eum Dóminus lorícam fídei, et ornávit eum.\end{Resp}
\begin{Resp}\Rsign Et ad portas paradísi coronávit eum.\end{Resp}
\switchcolumn
\EN
\begin{Resp}\Rsign The Lord loved him and beautified him; He clothed him with a robe of glory, \Star{} And crowned him at the gates of Paradise.\end{Resp}
\begin{Resp}\Vsign The Lord hath put on him the breast-plate of faith, and hath adorned him.\end{Resp}
\begin{Resp}\Rsign And crowned him at the gates of Paradise.\end{Resp}
\switchcolumn*

\LA
\LessonHead{Lectio VI}
\switchcolumn
\EN
\LessonHead{Lesson VI}
\switchcolumn*

\LA
\noindent A summo Pontifice Urbano quarto Romam vocatus, ejus jussu ecclesiásticum lucubrávit Offícium in Corporis Christi solemnitate celebrándum: oblatos vero honores et Neapolitanum archiepiscopátum, étiam deferénte Cleménte quarto, recusávit. A prædicatióne divini verbi non desistebat; quod cum fáceret per octavam Paschæ in basilica sancti Petri, mulierem, quæ ejus fimbriam tetígerat, a fluxu sánguinis liberávit. Missus a beato Gregório décimo ad concílium Lugdunense, in monasterio Fossæ Novæ in morbum íncidit, ubi ægrotus Cantica canticórum explanávit. Ibidem obiit quinquagenarius, anno salútis millesimo ducentésimo septuagesimo quarto, Nonis Martii. Miraculis étiam mórtuus cláruit; quibus probátis, a Joánne vigesimo secundo in Sanctórum númerum relátus est, anno millesimo tercentésimo vigesimo tertio, translato póstea ejus corpore Tolosam, ex mandato beáti Urbani quinti. Cum sanctis angelicis spirítibus non minus innocéntia quam ingenio comparatus, Doctoris Angelici nomen jure est adeptus, eidem auctoritate sancti Pii quinti confirmátum. Leo autem décimus tertius, libentíssime excípiens postulatiónes et vota ómnium pene Sacrórum antistitum orbis catholici, ad tot præcipue philosophicórum systématum a veritáte aberrántium luem propulsandam, ad increménta scientiárum et communem humani generis utilitátem, eum, ex sacrórum Rituum Congregatiónis consulto, per apostolicas litteras cælestem patronum scholárum ómnium catholicárum declarávit et instituit.\par
\switchcolumn
\EN
\noindent The Supreme Pontiff Urban IV. sent for him to Rome, and at his command he composed the Church Office for the feast of Corpus Christi. The Pope could not persuade him to accept any dignity. Pope Clement IV. also offered him the Archbishopric of Naples, but he refused it. He did not neglect the preaching of the Word of God. Once while he was giving a course of sermons in the Basilica of St Peter, during the octave of Easter, a woman who had an issue of blood was healed by touching the hem of his garment. He was sent by blessed Gregory X. to the Council of Lyons, but fell sick on his way to the Abbey of Fossa Nuovo, and there during his illness he made an exposition of the Song of Songs. There he died on the 7th day of March, in the year of salvation 1274, aged fifty years. He was distinguished for miracles even after his death, and on proof of these Pope John XXII. added his name to those of the Saints in the year 1323. His body was afterwards carried to Toulouse by command of blessed Urban V. He has been compared to an angel, both on account of his innocency and of his intellectual power, and has hence been deservedly termed the Angelic Doctor. The use of which title as applied to him was approved by the authority of holy Pius V. Leo XIII. cheerfully agreeing to the prayers and wishes of nearly all the bishops of the Catholic world, and in conformity with a vote of the Congregation of Sacred Rites, by his Apostolic letters declared and recognized Thomas of Aquino as the patron in heaven of all Catholic schools, as an antidote to the plague of so many false systems, especially of philosophy, for the increase of scientific knowledge, and for the common good of all mankind.\par
\switchcolumn*

\LA
\begin{Resp}\Rsign Iste homo perfécit ómnia quæ locútus est ei Deus, et dixit ad eum: Ingrédere in réquiem meam: \Star{} Quia te vidi justum coram me ex ómnibus géntibus.\end{Resp}
\begin{Resp}\Vsign Iste est, qui contémpsit vitam mundi, et pervénit ad cæléstia regna.\end{Resp}
\begin{Resp}\Rsign Quia te vidi justum coram me ex ómnibus géntibus.\end{Resp}
\begin{Resp}\Vsign Glória Patri, et Fílio, \Star{} et Spirítui Sancto.\end{Resp}
\begin{Resp}\Rsign Quia te vidi justum coram me ex ómnibus géntibus.\end{Resp}
\switchcolumn
\EN
\begin{Resp}\Rsign This is he which did according unto all that God commanded him; and God said unto him: Enter thou into My rest, \Star{} For thee have I seen righteous before Me among all people.\end{Resp}
\begin{Resp}\Vsign This is he which loved not his life in this world, and is come unto an everlasting kingdom.\end{Resp}
\begin{Resp}\Rsign For thee have I seen righteous before Me among all people.\end{Resp}
\begin{Resp}\Vsign Glory be to the Father, and to the Son, \Star{} and to the Holy Ghost.\end{Resp}
\begin{Resp}\Rsign For thee have I seen righteous before Me among all people.\end{Resp}
\switchcolumn*

\end{paracol}

\par\needspace{10\baselineskip}\addvspace{1.2em}{\centering\small\textsc{Die 8 Martii} · \textsc{8 March}\par}
\Office{S. Joannis de Deo Confessoris}{St. John of God, Confessor}{Duplex}
\addcontentsline{toc}{subsection}{Die 8 Martii · S. Joannis de Deo Confessoris\enspace\textit{St. John of God, Confessor}}
\begin{paracol}{2}
\LA
\RefLine{Lectiones I–III: ex Commúni Confessoris non pontificis (C5).}
\switchcolumn
\EN
\RefLine{Lessons I–III: from the Common of a Confessor not a Bishop (C5).}
\switchcolumn*

\LA
\RefLine{Lectio IX: de Scriptura occurrente.}
\switchcolumn
\EN
\RefLine{Lesson IX: from the Scripture of the day.}
\switchcolumn*

\LA
\LessonHead{Lectio IV}
\switchcolumn
\EN
\LessonHead{Lesson IV}
\switchcolumn*

\LA
\noindent Joánnes de Deo, ex catholicis piisque paréntibus in oppido Montis Majoris Junióris, regni Lusitaniæ, natus, quam sublimiter in sortem Dómini fúerit eléctus, insuetus splendor super ejus domo refulgens, sonitúsque æris campani sua sponte emissus, ab ipso nativitátis témpore non obscure prænuntiárunt. A laxiori vivéndi ratióne, divina operante virtúte, revocatus, magnæ sanctitátis exhibere specimen cœpit, et ob audítam prædicatiónem verbi Dei sic ad melióra se excitátum sensit, ut jam ab ipso sanctioris vitæ rudiménto consummátum aliquid perfectumque visus sit attigisse. Bonis ómnibus in páuperes carcéribus inclusos erogatis, admirábilis pœniténtiæ suique ipsíus contemptus cuncto pópulo spectaculum factus, a plerisque, ceu demens, gráviter afflíctus, in carcerem améntibus destinátum conjícitur. At Joánnes cælésti caritate magis incénsus, gemino atque amplo valetudinario ex piórum eleemosynis in civitáte Granaténsi exstructo, jactoque novi ordinis fundaménto, Ecclésiam nova prole fœcundávit fratrum Hospitalitátis, infirmis præcláro animárum corporúmque profectu inserviéntium, et longe lateque per orbem diffusórum.\par
\switchcolumn
\EN
\noindent John of God was born of Catholic and godly parents in the town of Montemor in Portugal, (in the year 1495.) The lot to which God had elected him was foreshown at his birth by a light shining over the house, and by the ringing of a bell untouched by human hands. He fell at one time into a loose habit of life, but was recalled by the grace of God, and began to show tokens of true reformation. By hearing the Word of God, he so felt himself stirred up to strive after nobler things, that he considered not that to which he had already attained, and yearned to be perfect, as our Father in heaven is perfect. He gave away all his property to the poor and prisoners, and became a gazing-stock to all that knew him, by the strength of his repentance, and the depth of his selfcontempt. On this account he was commonly supposed to be mad, and was once shut up in a lunatic asylum. He was only the more filled with schemes of charity, and collected, by begging, funds sufficient to build a large double Hospital in the town ot Granada. Here he founded the new Order of Hospital Brethren with which he enriched the Church. These Brethren are now spread throughout all parts of the world, and engaged in ministering to the souls and bodies of the sick.\par
\switchcolumn*

\LA
\begin{Resp}\Rsign Honéstum fecit illum Dóminus, et custodívit eum ab inimícis, et a seductóribus tutávit illum: \Star{} Et dedit illi claritátem ætérnam.\end{Resp}
\begin{Resp}\Vsign Justum dedúxit Dóminus per vias rectas, et osténdit illi regnum Dei.\end{Resp}
\begin{Resp}\Rsign Et dedit illi claritátem ætérnam.\end{Resp}
\switchcolumn
\EN
\begin{Resp}\Rsign The Lord made him honourable, and defended him from his enemies, and kept him safe from those that lay in wait for him; \Star{} And gave him perpetual glory.\end{Resp}
\begin{Resp}\Vsign He went down with him into the pit, and left him not in bonds.\end{Resp}
\begin{Resp}\Rsign And gave him perpetual glory.\end{Resp}
\switchcolumn*

\LA
\LessonHead{Lectio V}
\switchcolumn
\EN
\LessonHead{Lesson V}
\switchcolumn*

\LA
\noindent Paupéribus ægrotis, quos propriis quandoque humeris domum deferebat, nulla re ad ánimæ corporisque salútem profícua déerit. Effusa quoque extra nosocomíum caritate indigéntibus muliéribus viduis, et præcipue virgínibus periclitántibus clam aliménta subministrabat, curamque indefessam adhibebat, ut carnis concupiscéntiam a próximis hujusmodi vítio inquinátis exterminaret. Cum autem maximum in regio Granaténsi valetudinario excitátum fuísset incendium, Joánnes impávidus prosiliit in ignem, huc illuc discurrens, quoúsque tum infirmos humeris exportatos, tum lectulos e fenestris projectos ab igne vindicávit, ac per dimidiam horam inter flammas, jam in immensum succrescéntes, versatus, exínde divínitus incólumis, univérsis civibus admirántibus exívit; in schola caritátis édocens, segniórem in eum fuísse ignem qui foris ússerat, quam qui intus accenderat.\par
\switchcolumn
\EN
\noindent He strove to get for the sick poor, whom he sometimes brought to the Hospital on his own shoulders, whatever was needful for their souls or bodies. His charity was extended to the poor outside of his institution, and he used to supply food privately to necessitous widows, and more so to young women whose virtue was tempted on account of their poverty. He was most careful in encouraging the virtue of purity in all whom he knew. On one occasion when there was a great fire in the hospital at Granada, John bravely entered the burning house, ran from one part of it to another, carried out the sick on his shoulders, and threw the beds out of the windows, and finally, after passing half-an-hour in the midst of the flames, which were now raging with great violence, by the mercy of God left the building uninjured, to the great wonder of all the citizens; thereby to teach all them that love God that the fire which burnt in his heart gave him strength to risk the fire which threatened him from without.\par
\switchcolumn*

\LA
\begin{Resp}\Rsign Amávit eum Dóminus, et ornávit eum: stolam glóriæ índuit eum, \Star{} Et ad portas paradísi coronávit eum.\end{Resp}
\begin{Resp}\Vsign Índuit eum Dóminus lorícam fídei, et ornávit eum.\end{Resp}
\begin{Resp}\Rsign Et ad portas paradísi coronávit eum.\end{Resp}
\switchcolumn
\EN
\begin{Resp}\Rsign The Lord loved him and beautified him; He clothed him with a robe of glory, \Star{} And crowned him at the gates of Paradise.\end{Resp}
\begin{Resp}\Vsign The Lord hath put on him the breast-plate of faith, and hath adorned him.\end{Resp}
\begin{Resp}\Rsign And crowned him at the gates of Paradise.\end{Resp}
\switchcolumn*

\LA
\LessonHead{Lectio VI}
\switchcolumn
\EN
\LessonHead{Lesson VI}
\switchcolumn*

\LA
\noindent Multiplici asperitátum genere, demissíssima obediéntia, extrema paupertáte, orándi studio, rerum divinárum contemplatióne, ac in beátam Vírginem pietáte mirifice excelluit, et lacrimárum dono enituit. Denique gravi morbo correptus, ómnibus Ecclésiæ sacramentis rite sancteque refectus, viribus licet destitútus, propriis indútus vestibus, e léctulo surgens ac provolutus in genua, manu et corde Christum Dóminum e cruce pendentem perstringens, octavo Idus Martii anno millesimo quingentésimo quinquagesimo obiit in osculo Dómini; quem étiam mórtuus ténuit, nec dimísit, et in eádem corporis constitutióne sex circiter horas, quoúsque inde dimotus fuísset, tota civitáte inspectante, mirabíliter permansit, odórem mire fragrántem diffúndens. Quem ante et post óbitum plurimis miraculis clarum, Alexander octavus Pontifex maximus in Sanctórum númerum rétulit; et Leo décimus tertius, ex Sacrórum catholici orbis antístitum voto ac Rituum Congregatiónis consulto, cælestem ómnium hospitalium et infirmórum ubíque degéntium patronum declarávit, ipsiúsque nomen in agonizántium litaníis invocari præcepit.\par
\switchcolumn
\EN
\noindent He was a marked example of every kind of austerity, of the most lowly obedience, of the deepest voluntary poverty, of the most constant prayer, of ghostly contemplation, and of love towards the Blessed Virgin. He was distinguished for the gift of tears. Being at last seized by deadly sickness, he duly received, with saintly affection, all the Sacraments of the Church. After all strength seemed to have left him, he got out of his bed, put on his own clothes, and knelt down before an image of the Lord Christ hanging on the Cross. Round it he threw his arms and pressed it against his heart, and in this position, as it were in the kiss of the Lord, he died, on the 8th day of March 1550. After his death his body did not leave its grip of the crucifix until it was forcibly taken away, six hours after. During these six hours all the inhabitants of the city came to see it, and noticed a savour of strange sweetness proceeding from it. His name was illustrious as a worker of miracles both before and after his death, and the Supreme Pontiff Alexander VIII added it to those of the Saints, and Leo XIII, at the desire of the Bishops of the Catholic world, and in accordance with a vote of the Congregation of Rites, declared him the patron in heaven of all the sick and those who nurse them, wheresoever dwelling, and ordered that his name should be called upon in the Litany for the dying.\par
\switchcolumn*

\LA
\begin{Resp}\Rsign Iste homo perfécit ómnia quæ locútus est ei Deus, et dixit ad eum: Ingrédere in réquiem meam: \Star{} Quia te vidi justum coram me ex ómnibus géntibus.\end{Resp}
\begin{Resp}\Vsign Iste est, qui contémpsit vitam mundi, et pervénit ad cæléstia regna.\end{Resp}
\begin{Resp}\Rsign Quia te vidi justum coram me ex ómnibus géntibus.\end{Resp}
\begin{Resp}\Vsign Glória Patri, et Fílio, \Star{} et Spirítui Sancto.\end{Resp}
\begin{Resp}\Rsign Quia te vidi justum coram me ex ómnibus géntibus.\end{Resp}
\switchcolumn
\EN
\begin{Resp}\Rsign This is he which did according unto all that God commanded him; and God said unto him: Enter thou into My rest, \Star{} For thee have I seen righteous before Me among all people.\end{Resp}
\begin{Resp}\Vsign This is he which loved not his life in this world, and is come unto an everlasting kingdom.\end{Resp}
\begin{Resp}\Rsign For thee have I seen righteous before Me among all people.\end{Resp}
\begin{Resp}\Vsign Glory be to the Father, and to the Son, \Star{} and to the Holy Ghost.\end{Resp}
\begin{Resp}\Rsign For thee have I seen righteous before Me among all people.\end{Resp}
\switchcolumn*

\LA
\LessonHead{Lectio VII}
\switchcolumn
\EN
\LessonHead{Lesson VII}
\switchcolumn*

\LA
\LessonTitle{Léctio sancti Evangélii secúndum Matthǽum}
\switchcolumn
\EN
\LessonTitle{From the Holy Gospel according to Matthew}
\switchcolumn*

\LA
\Cite{Matt 22:34-46}
\switchcolumn
\EN
\Cite{Matt 22:34-46}
\switchcolumn*

\LA
\noindent In illo témpore: Accessérunt ad Jesum pharisæi, et interrogávit eum unus ex eis legis doctor tentans eum: Magister, quod est mandátum magnum in lege? Et réliqua.\par
\switchcolumn
\EN
\noindent At that time The Pharisees came together: And one of them, a doctor of the law, asking him, tempting him: Master, which is the greatest commandment in the law? And so on.\par
\switchcolumn*

\LA
\LessonTitle{Homilía sancti Joánnis Chrysóstomi}
\switchcolumn
\EN
\LessonTitle{Homily by St. John Chrysostom, Patriarch (of Constantinople.)}
\switchcolumn*

\LA
\Source{Homilia 72 in Matthæum}
\switchcolumn
\EN
\Source{72nd on Matthew.}
\switchcolumn*

\LA
\noindent Sadducæis confúsis, pharisæi rursus aggrediúntur; cumque quiéscere oporteret, decertare voluérunt: et legis perítiam profitentem præmittunt, non discere, sed tentare cupiéntes; ac ita intérrogant: Quodnam primum mandátum in lege sit. Nam cum primum illud sit, Diliges Dóminum Deum tuum: putántes causas sibi allatúrum ad mandátum hoc corrigéndum, aliquid addendo, quóniam Deum se faciebat, hoc modo intérrogant. Quid ígitur Christus? Ut ostendat idcirco ad hæc eos devenisse, quia nulla in eis esset caritas, sed invidiæ livore tabéscerent: Diliges, inquit, Dóminum Deum tuum: hoc primum et magnum mandátum est. Secúndum autem simile huic: Diliges próximum tuum sicut teípsum.\par
\switchcolumn
\EN
\noindent When the Pharisees had heard that Christ had put the Sadducees to silence, they gathered themselves together for a fresh attack; just when it behoved them to be quiet, they willed to contend; and so they put forward one of themselves, who professed skill in the law, not wishing to learn, but to lay a snare. This person therefore proposed the question Which is the great commandment in the law? The first and great commandment is Thou shalt love the Lord thy God, but they expected that He would make some exception or addition to this in His Own case, since He made Himself God. (John x. 33.) With this expectation they asked Him the question: But what said Christ? To show that they had adopted this course, because they were loveless, and sick with envy, He answered Thou shalt love the Lord thy God with all thy heart, and with all thy soul, and with all thy mind. This is the first and great commandment. And the second is like unto it Thou shalt love thy neighbour as thyself.\par
\switchcolumn*

\LA
\begin{Resp}\Rsign Iste est qui ante Deum magnas virtútes operátus est, et de omni corde suo laudávit Dóminum: \Star{} Ipse intercédat pro peccátis ómnium populórum.\end{Resp}
\begin{Resp}\Vsign Ecce homo sine queréla, verus Dei cultor, ábstinens se ab omni ópere malo, et pérmanens in innocéntia sua.\end{Resp}
\begin{Resp}\Rsign Ipse intercédat pro peccátis ómnium populórum.\end{Resp}
\switchcolumn
\EN
\begin{Resp}\Rsign This is he which wrought great wonders before God, and praised the Lord with all his heart. \Star{} May he pray for all people, that their sins may be forgiven unto them.\end{Resp}
\begin{Resp}\Vsign Behold a man without blame, a worshipper of God in truth, keeping himself clean from every evil work, and abiding still in his innocency.\end{Resp}
\begin{Resp}\Rsign May he pray for all people, that their sins may be forgiven unto them.\end{Resp}
\switchcolumn*

\LA
\LessonHead{Lectio VIII}
\switchcolumn
\EN
\LessonHead{Lesson VIII}
\switchcolumn*

\LA
\noindent Quam ob rem simile est huic? Quóniam hoc illud inducit, et ab illo rursus munítur. Quicúmque enim male agit, ódio habet lucem, et non venit ad lucem. Et rursus: Dixit insípiens in corde suo, Non est Deus. Deínde sequitur: Corrúpti sunt, et abominábiles facti sunt in stúdiis suis. Et íterum: Radix ómnium malórum avarítia est; quam quidam appeténtes, erravérunt a fide. Et, Qui díligit me, mandáta mea servábit: quorum caput et radix est: Diliges Dóminum Deum tuum, et próximum tuum sicut teípsum.\par
\switchcolumn
\EN
\noindent Mighty is this second commandment like unto the first? Because the first is the second's source and sanction. For every one that doeth evil hateth the light, neither cometh to the light. (John. iii. 20.) And again The fool hath said in his heart There is no God and there folio weth They are corrupt, and become abominable in their works. (Ps. xiii. I.) And yet again: The love of money is the root of all evil; which while some coveted after, they have erred from the faith. (i Tim. vi. 10.) And yet once more: If ye love Me, keep My commandments. (John xiv. 15,) of which commandments the head and root is Thou shalt love the Lord thy God; and thy neighbour as thyself.\par
\switchcolumn*

\LA
\begin{Resp}\Rsign Sint lumbi vestri præcíncti, et lucérnæ ardéntes in mánibus vestris: \Star{} Et vos símiles homínibus exspectántibus dóminum suum, quando revertátur a núptiis.\end{Resp}
\begin{Resp}\Vsign Vigiláte ergo, quia nescítis qua hora Dóminus vester ventúrus sit.\end{Resp}
\begin{Resp}\Rsign Et vos símiles homínibus exspectántibus dóminum suum, quando revertátur a núptiis.\end{Resp}
\begin{Resp}\Vsign Glória Patri, et Fílio, \Star{} et Spirítui Sancto.\end{Resp}
\begin{Resp}\Rsign Et vos símiles homínibus exspectántibus dóminum suum, quando revertátur a núptiis.\end{Resp}
\switchcolumn
\EN
\begin{Resp}\Rsign Let your loins be girded about, and your lights burning; \Star{} And ye yourselves like unto men that wait for their lord, when he will return from the wedding.\end{Resp}
\begin{Resp}\Vsign Watch therefore, for ye know not what hour your Lord doth come.\end{Resp}
\begin{Resp}\Rsign And ye yourselves like unto men that wait for their lord, when he will return from the wedding.\end{Resp}
\begin{Resp}\Vsign Glory be to the Father, and to the Son, \Star{} and to the Holy Ghost.\end{Resp}
\begin{Resp}\Rsign And ye yourselves like unto men that wait for their lord, when he will return from the wedding.\end{Resp}
\switchcolumn*

\LA
\LessonHead{Lectio IX}
\switchcolumn
\EN
\LessonHead{Lesson IX}
\switchcolumn*

\LA
\noindent Si ergo diligere Deum, diligere proximum est (nam si diligis me, o Petre, inquit, pasce oves meas): si étiam dilectio proximi facit ut mandata custodias: merito ait in his totam legem et prophetas pendere. Et quemadmodum in superioribus, cum de resurrectione interrogaretur, plus docuit quam tentantes petebant: sic in hoc loco de primo interrogatus mandato, secundum étiam non valde quam primum inferius, sponte attulit: secundum enim est primo simile. Ita occulte insinuavit, odio illos ad quærendum incitari. Caritas enim, inquit, non gemulatur.\par
\switchcolumn
\EN
\noindent If therefore, to love God is to love our neighbour also, as it appeareth where it is written Simon, son of Jonas, lovest thou Me? And he said unto Him Lord, Thou knowest all things; Thou knowest that I love thee. Jesus saith unto him Feed My sheep, (John xxi. 17,) and if love is the fulfilling of the law, \textasciitilde{}(Rom. xiii. 10,) justly doth the Lord say that on these two commandments hang all the law and the Prophets. And even as when, before this, (23-32,) being interrogated about the Resurrection, He answered them more than they asked, so, now, being interrogated concerning the first and great commandment, He answereth them, of His own accord, touching that second one also, which is little lower than the first, for the second is like unto it. Herein He would have them understand that it was hatred stirred them up to question Him. For Charity, saith the Apostle, envieth not. (1 Cor. xiii. 4.)\par
\switchcolumn*

\LA
\TeDeum{Te Deum}
\switchcolumn
\EN
\TeDeum{Te Deum}
\switchcolumn*

\end{paracol}

\par\needspace{10\baselineskip}\addvspace{1.2em}{\centering\small\textsc{Die 9 Martii} · \textsc{9 March}\par}
\Office{S. Franciscæ Romanæ Viduæ}{St. Frances of Rome, Widow}{Duplex}
\addcontentsline{toc}{subsection}{Die 9 Martii · S. Franciscæ Romanæ Viduæ\enspace\textit{St. Frances of Rome, Widow}}
\begin{paracol}{2}
\LA
\RefLine{Lectiones I–III: ex Commúni non Virginum non Martyrum (C7a).}
\switchcolumn
\EN
\RefLine{Lessons I–III: from the Common of Widows (C7a).}
\switchcolumn*

\LA
\RefLine{Lectiones VII–VIII: ex Commúni non Virginum non Martyrum (C7a).}
\switchcolumn
\EN
\RefLine{Lessons VII–VIII: from the Common of Widows (C7a).}
\switchcolumn*

\LA
\RefLine{Lectio IX: de Scriptura occurrente.}
\switchcolumn
\EN
\RefLine{Lesson IX: from the Scripture of the day.}
\switchcolumn*

\LA
\LessonHead{Lectio IV}
\switchcolumn
\EN
\LessonHead{Lesson IV}
\switchcolumn*

\LA
\noindent Francisca nobilis matrona Romana, ab ineunte ætate illustria dedit virtutum exempla: etenim pueriles ludos et illecebras mundi respuens, solitudine et oratione magnopere delectabatur. Undecim annos nata, virginitatem suam Deo consecrare, et monasterium ingredi proposuit: parentum tamen voluntati humiliter obtemperans, Laurentio de Pontianis, juveni æque diviti ac nobili nupsit. In matrimonio arctioris vitæ propositum, quantum licuit, semper retinuit: a spectaculis, conviviis, aliisque hujusmodi oblectamentis abhorrens, lanea ac vulgari veste utens, et quidquid a domesticis curis supererat temporis, orationi aut proximorum utilitati tribuens: in id vero maxima sollicitudine incumbens, ut matronas Romanas a pompis sæculi et ornatus vanitate revocaret. Quapropter domum Oblatarum, sub regula sancti Benedicti, congregationis Montis Oliveti, adhuc viro alligata, in Urbe instituit. Viri exsilium, bonorum jacturam, ac universæ domus mœrorem non modo constantissime toleravit: sed gratias agens cum beato Job, illud frequenter usurpabat: Dominus dedit, Dominus abstulit: sit nomen Domini benedictum.\par
\switchcolumn
\EN
\noindent The noble Roman matron Frances A (was born in the year 1384, and) was a pattern of godliness from her earliest years. As a child she shrank from games, and set no store by the amusements of the world, but delighted to be continually alone and engaged in prayer. At the age of eleven years she desired to consecrate her virginity to God, and to enter a convent, but humbly yielded obedience to the wishes of her parents, and was married to Lawrence de' Pontiani, a young man whose rank was equal to his wealth. As a wife she persevered, as far as she lawfully could, in her determination to lead an austere life; she abstained as much as possible from going to shows, feasts, and such like amusements, dressed plainly in woollen stuffs, and spent in prayer or the service of her neighbour whatever time she did not occupy with her duties as mistress of her husband's house. She strove earnestly to wean the married women of Rome from the vanities of the world and the frivolities of dress. To this end she founded during her husband's lifetime the Sisterhood of the Oblates, under the rule of the Benedictine congregation called of the Mount of Olives. When it pleased God, (in the year 1413,) that her husband should be banished, all her goods taken away, and her home ruined, she meekly bowed down before His holy will, often repeating the words of the blessed Job The Lord gave, and the Lord hath taken away; blessed be the name of the Lord. (i. 21.)\par
\switchcolumn*

\LA
\begin{Resp}\Rsign Propter veritátem, et mansuetúdinem, et justítiam: \Star{} Et dedúcet te mirabíliter déxtera tua.\end{Resp}
\begin{Resp}\Vsign Spécie tua et pulchritúdine tua inténde, próspere procéde, et regna.\end{Resp}
\begin{Resp}\Rsign Et dedúcet te mirabíliter déxtera tua.\end{Resp}
\switchcolumn
\EN
\begin{Resp}\Rsign Because of truth, and meekness, and righteousness; \Star{} And thy right hand shall lead thee wonderfully.\end{Resp}
\begin{Resp}\Vsign In thy comeliness, and thy beauty, go forward, fare prosperously, and reign.\end{Resp}
\begin{Resp}\Rsign And thy right hand shall lead thee wonderfully.\end{Resp}
\switchcolumn*

\LA
\LessonHead{Lectio V}
\switchcolumn
\EN
\LessonHead{Lesson V}
\switchcolumn*

\LA
\noindent Viro defuncto, ad prædictam Oblatarum domum convolans, nudis pedibus, fune ad collum alligato, humi prostrata, multis cum lacrimis earum numero adscribi suppliciter postulavit. Voti compos facta, licet esset omnium mater, non alio tamen quam ancillæ, vilissimæque feminæ, et immunditiæ vasculi titulo gloriabatur. Quam vilem sui existimationem et verbo declaravit et exemplo: sæpe enim e suburbana vinea revertens, et lignorum fascem proprio capiti impositum deferens, vel eisdem onustum agens per Urbem asellum, pauperibus subveniebat, in quos étiam largas eleemosynas erogabat: ægrotantesque in xenodochiis visitans, non corporali tantum cibo, sed salutaribus monitis recreabat. Corpus suum vigiliis, jejuniis, cilicio, ferreo cingulo, crebrisque flagellis in servitutem redigere jugiter satagebat. Cibum illi semel in die herbæ et legumina, aqua potum præbuit. Hos tamen corporis cruciatus aliquando confessarii mandato, a cujus ore nutuque pendebat, modice temperavit.\par
\switchcolumn
\EN
\noindent On her husband's death she (in 1437) betook herself immediately to the house of the Oblates, and, with her feet bare and a rope round her neck, threw herself down on the threshold, entreating the sisters with tears to receive her into their number. When she obtained her wish, although she was the mother of them all, she would be among them only as one that served, glorying rather to be called the most degraded of women and a vessel of uncleanness. Her lowly esteem of herself was shown both by her word and example. She passed often through the city from a vineyard in the country carrying a bundle of sticks on her head, or driving an ass laden with faggots; she succoured the needy, for whom she collected large alms, and visited the sick in the hospitals, ministering to them both food for the body and exhortations healthful for their souls. She strove continually to bring her body into subjection by watchings, fastings, haircloth, the wearing of an iron girdle, and the often use of a scourge. She never ate but once a day, and then only vegetables, and she took no drink but water. These severities she however sometimes relaxed, in obedience to her confessor, on whose word and wishes she framed her customs.\par
\switchcolumn*

\LA
\begin{Resp}\Rsign Dilexísti justítiam, et odísti iniquitátem: \Star{} Proptérea unxit te Deus, Deus tuus, óleo lætítiæ.\end{Resp}
\begin{Resp}\Vsign Propter veritátem, et mansuetúdinem, et justítiam.\end{Resp}
\begin{Resp}\Rsign Proptérea unxit te Deus, Deus tuus, óleo lætítiæ.\end{Resp}
\switchcolumn
\EN
\begin{Resp}\Rsign Thou hast loved righteousness, and hated iniquity; \Star{} Therefore God, thy God, hath anointed thee with the oil of gladness.\end{Resp}
\begin{Resp}\Vsign Because of truth, and meekness, and righteousness.\end{Resp}
\begin{Resp}\Rsign Therefore God, thy God, hath anointed thee with the oil of gladness.\end{Resp}
\switchcolumn*

\LA
\LessonHead{Lectio VI}
\switchcolumn
\EN
\LessonHead{Lesson VI}
\switchcolumn*

\LA
\noindent Divina mysteria, præsertim vero Christi Domini passionem, tanto mentis ardore, tantaque lacrimarum vi contemplabatur, ut præ doloris magnitudine pene confici videretur. Sæpe étiam cum oraret, maxime sumpto sanctissimæ Eucharistiæ sacramento, spiritu in Deum elevata, ac cælestium contemplatione rapta, immobilis permanebat. Quapropter humani generis hostis variis eam contumeliis ac verberibus a proposito dimovere conabatur: quem tamen illa imperterrita semper elusit; Angeli præsertim præsidio, cujus familiari consuetudine gloriosum de eo triumphum reportavit. Gratia curationum et prophetiæ dono enituit, quo et futura prædixit, et cordium secreta penetravit. Non semel aquæ, vel per rivum decurrentes, vel e cælo labentes, intactam prorsus, dum Deo vacaret, reliquerunt. Modica panis fragmenta, quæ vix tribus sororibus reficiendis fuíssent satis, sic ejus precibus Dominus multiplicavit, ut quindecim inde exsaturatis, tantum superfuerit, ut canistrum impleverit: et aliquando earumdem sororum, extra Urbem mense Januario ligna parantium, sitim recentis uvæ racemis ex vitæ in arbore pendentibus mirabiliter obtentis, abunde expleverit. Denique meritis et miraculis clara, migravit ad Dominum, anno ætatis suæ quinquagesimo sexto. Quam Paulus quintus Pontifex Maximus in Sanctorum numerum retulit.\par
\switchcolumn
\EN
\noindent So great was her mental realisation of the things of God, and chiefly of the sufferings of the Lord Christ, and so abundant her tears in contemplating them, that she seemed sometimes about to sink under her grief. Often when she was engaged in prayer, and principally after she had received the Most Holy Sacrament of the Eucharist, her spirit became altogether lifted up to God, and she remained motionless, carried away by the thought of heavenly things. The enemy of man assailed her with diverse reproaches and buffetings to break her off her intent, but she feared him not, and with the help of an Angel whom God gave her to be her familiar friend, she won a noble victory over the tempter. God glorified her with the gifts of healing and of prophecy, whereby she foretold things to come, and saw the secrets of the hearts of men. More than once while her thoughts were busy in God she remained unwet by streams or rain. When there was left only bread enough for three sisters, the Lord at her prayers was pleased so to multiply it, that fifteen had enough, and the basket was filled again with the fragments. In the month of January also, when the sisters were gathering sticks in the country, and were thirsty, she satisfied them abundantly with bunches of fresh grapes from a tree. She departed to be with the Lord, famous for good works and miracles, in the fifty-sixth year of her age, (upon the 9th day of March, in the year of our Lord 1440.) The Supreme Pontiff Paul V. caused her to be numbered among the saints.\par
\switchcolumn*

\LA
\begin{Resp}\Rsign Fallax grátia, et vana est pulchritúdo: \Star{} Múlier timens Dóminum, ipsa laudábitur.\end{Resp}
\begin{Resp}\Vsign Date ei de fructu mánuum suárum, et laudent eam in portis ópera ejus.\end{Resp}
\begin{Resp}\Rsign Múlier timens Dóminum, ipsa laudábitur.\end{Resp}
\begin{Resp}\Vsign Glória Patri, et Fílio, \Star{} et Spirítui Sancto.\end{Resp}
\begin{Resp}\Rsign Múlier timens Dóminum, ipsa laudábitur.\end{Resp}
\switchcolumn
\EN
\begin{Resp}\Rsign Favour is deceitful, and beauty is vain \Star{} A woman that feareth God she shall be praised.\end{Resp}
\begin{Resp}\Vsign Give her of the fruit of her hands, and let her own works praise her in the gates.\end{Resp}
\begin{Resp}\Rsign A woman that feareth God she shall be praised.\end{Resp}
\begin{Resp}\Vsign Glory be to the Father, and to the Son, \Star{} and to the Holy Ghost.\end{Resp}
\begin{Resp}\Rsign A woman that feareth God she shall be praised.\end{Resp}
\switchcolumn*

\end{paracol}

\par\needspace{10\baselineskip}\addvspace{1.2em}{\centering\small\textsc{Die 10 Martii} · \textsc{10 March}\par}
\Office{Ss. Quadraginta Martyrum}{Forty Holy Martyrs}{Semiduplex}
\addcontentsline{toc}{subsection}{Die 10 Martii · Ss. Quadraginta Martyrum\enspace\textit{Forty Holy Martyrs}}
\begin{paracol}{2}
\LA
\RefLine{Lectiones I–III: ex Commúni Plurimorum Martyrum Pontificum (C3).}
\switchcolumn
\EN
\RefLine{Lessons I–III: from the Commune Plurimorum Martyrum Pontificum (C3).}
\switchcolumn*

\LA
\RefLine{Lectio IX: de Scriptura occurrente.}
\switchcolumn
\EN
\RefLine{Lesson IX: from the Scripture of the day.}
\switchcolumn*

\LA
\LessonHead{Lectio IV}
\switchcolumn
\EN
\LessonHead{Lesson IV}
\switchcolumn*

\LA
\noindent Licinio imperatore, et Agricolao præside, ad Sebasten Armeniæ urbem, quadraginta militum fides in Jesum Christum, et fortitudo in cruciatibus perferendis enituit. Qui sæpius in horribilem carcerem detrusi, vinculisque constricti, cum ora ipsorum lapidibus contusa fuíssent, hiemis tempore frigidissimo nudi sub aperto aëre supra stagnum rigens pernoctare jussi sunt, ut frigore congelati necarentur. Una autem erat omnium oratio: Quadraginta in stadium ingressi sumus, quadraginta item Domine corona donemur, ne una quidem huic numero desit. Est in honore hic numerus, quem tu quadraginta dierum jejunio decorasti, per quem divina lex ingressa est in orbem terrarum. Elias quadraginta dierum jejunio Deum quærens, ejus visionem consecutus est. Et hæc quidem illorum erat oratio.\par
\switchcolumn
\EN
\noindent While Licinius was Emperor and Agricolaus President, (in the year of our Lord 320,) forty soldiers at Sebaste, a city of Armenia, gave a singular instance of faith in Jesus Christ, and bravery under suffering. After being often remanded to an horrid prison-house, bound in fetters, and their mouths bruised with stones, they were ordered out in the depth of winter, stripped naked, and put upon a frozen pool, to die of cold during the night. The prayer of them all was the same O Lord, forty of us have begun to run in the race, grant that all forty may receive the crown, let not one be wanting at the last. Behold, is it not an honourable number in thy sight, Who didst bless the fast of forty days, and at the end thy Divine Law came forth to the earth? When also Elias sought thee, Thou, O God, didst reveal thyself unto him when he had fasted for forty days. Even so was their petition.\par
\switchcolumn*

\LA
\begin{Resp}\Rsign Sancti tui, Dómine, mirábile consecúti sunt iter, serviéntes præcéptis tuis, ut inveniréntur illǽsi in aquis válidis: \Star{} Terra appáruit árida: et in Mari Rubro via sine impediménto.\end{Resp}
\begin{Resp}\Vsign Quóniam percússit petram, et fluxérunt aquæ, et torréntes inundavérunt.\end{Resp}
\begin{Resp}\Rsign Terra appáruit árida: et in Mari Rubro via sine impediménto.\end{Resp}
\switchcolumn
\EN
\begin{Resp}\Rsign thy Saints, O Lord, have passed a wonderful way, serving thy commandments, that they might be found without hurt in the midst of the mighty waters. \Star{} Dry land appeared, and, out of the Red Sea, a way without impediment.\end{Resp}
\begin{Resp}\Vsign He smote the rock, and the waters gushed out, and the streams overflowed.\end{Resp}
\begin{Resp}\Rsign Dry land appeared, and, out of the Red Sea, a way without impediment.\end{Resp}
\switchcolumn*

\LA
\LessonHead{Lectio V}
\switchcolumn
\EN
\LessonHead{Lesson V}
\switchcolumn*

\LA
\noindent Ceteris autem custodibus somno deditis, solus vigilabat janitor, qui et illos orantes, et luce circumfusos, et quosdam e cælo descendentes Angelos tamquam a Rege missos, qui coronas triginta novem militibus distribuerent intuens ita secum loquebatur: Quadraginta hi sunt: quadragesimi corona ubi est? Quæ dum cogitaret, unus ex illo numero, cui animus ad frigus ferendum defecerat in proximum tepefactum balneum desiliens, sanctos illos summo dolore affecit. Verum Deus illorum preces irritas esse non est passus: nam rei eventum admiratus janitor, mox custodibus e somno excitatis, detractisque sibi vestibus, ac se Christianum esse clara voce professus, Martyribus se adjunxit. Cum vero præsidis satellites janitorem quoque Christianum esse cognovissent bacillis comminuta omnium eorum crura fregerunt.\par
\switchcolumn
\EN
\noindent When the keepers were all asleep and the watchman only was awake, he heard them praying and saw a light shining round about them, and Angels coming down from heaven, as the messengers of the King, bearing nine-and-thirty crowns, and distributing them to the soldiers. Then he said within himself: Are not forty here? Where is the crown of the fortieth? And as he looked he saw one of them whose courage could not bear the cold, come and leap into a warm bath that stood by; and the Saints were grievously afflicted. Nevertheless God suffered not that their prayer should return unto them void; for the watchman wondered, and called the keepers, and stripped himself of his clothes; and, when with a loud voice he had confessed himself a Christian, he joined the Martyrs. When the servants of the President knew that the watchman also was a Christian, they brake the legs of them all with staves.\par
\switchcolumn*

\LA
\begin{Resp}\Rsign Vérbera carníficum non timuérunt Sancti Dei, moriéntes pro Christi nómine: \Star{} Ut herédes fíerent in domo Dómini.\end{Resp}
\begin{Resp}\Vsign Tradidérunt córpora sua propter Deum ad supplícia.\end{Resp}
\begin{Resp}\Rsign Ut herédes fíerent in domo Dómini.\end{Resp}
\switchcolumn
\EN
\begin{Resp}\Rsign The Saints of God shrank not from the stripes of the executioners, but died for Christ's Name's sake \Star{} That they might be made joint-heirs in the house of the Lord.\end{Resp}
\begin{Resp}\Vsign They gave their bodies for God's sake to death.\end{Resp}
\begin{Resp}\Rsign That they might be made joint-heirs in the house of the Lord.\end{Resp}
\switchcolumn*

\LA
\LessonHead{Lectio VI}
\switchcolumn
\EN
\LessonHead{Lesson VI}
\switchcolumn*

\LA
\noindent In eo supplicio mortui sunt omnes, præter Melithonem natu minimum. Quem cum præsens mater ejus, fractis cruribus, adhuc viventem vidisset, sic cohortata est: Fili, paulisper sustine; ecce Christus ad januam stat adjuvans te. Cum vero reliquorum corpora plaustris imponi cerneret, ut in rogum inférrentur, ac filium suum relinqui, quod speraret impia turba puerum, si vixisset, ad idolorum cultum revocari posse; ipso in humeros sublato, sancta mater vehicula Mártyrum corporibus onusta strenue prosequebatur: in cujus amplexu Melithon spiritum Deo reddidit, ejusque corpus in eumdem illum ceterorum Mártyrum rogum pia mater injecit: ut qui fide et virtute conjunctissimi fuerant, funeris étiam societate copulati, una in cælum pervenirent. Combustis illis, eorum reliquiæ projectæ in profluentem, cum mirabiliter in unum confluxissent locum, salvæ et integræ repertæ, honorifico sepulchro conditæ sunt.\par
\switchcolumn
\EN
\noindent Under this torment died they all, saving Melithon, who was the youngest. Now, his mother stood by, and when she saw that his legs were broken, but that he was yet alive, she cried, and said My son, have patience but a little longer. Behold how Christ standeth at the door to help thee. When she saw the bodies of all the others put upon carts and taken away to be burned, and that her son was left behind, because the multitude wickedly hoped that being but a lad, if he lived, he might yet be drawn to commit idolatry, the holy mother took him on her own shoulders and bravely followed behind the carts laden with the bodies of the Martyrs. In her arms Melithon gave up his soul to God, and the mother who loved him so well laid his body with her own hands upon the pile, with those of the other Martyrs, that, as they had all been one in faith and strength, in death they might not be divided, and might enter heaven together. After the burning, what remained of them was thrown into a running stream, but the ashes were all washed together into one place, and being found and rescued, they were laid in an honourable sepulchre.\par
\switchcolumn*

\LA
\begin{Resp}\Rsign Tamquam aurum in fornáce probávit eléctos Dóminus, et quasi holocáusti hóstiam accépit illos: et in témpore erit respéctus illórum: \Star{} Quóniam donum et pax est eléctis Dei.\end{Resp}
\begin{Resp}\Vsign Qui confídunt in illum, intélligent veritátem, et fidéles in dilectióne acquiéscent illi.\end{Resp}
\begin{Resp}\Rsign Quóniam donum et pax est eléctis Dei.\end{Resp}
\begin{Resp}\Vsign Glória Patri, et Fílio, \Star{} et Spirítui Sancto.\end{Resp}
\begin{Resp}\Rsign Quóniam donum et pax est eléctis Dei.\end{Resp}
\switchcolumn
\EN
\begin{Resp}\Rsign As gold in the furnace hath the Lord tried His chosen ones, and received them as a burnt-offering, and yet a while, and they shall be regarded; \Star{} For the grace of God, and His peace, are with His chosen.\end{Resp}
\begin{Resp}\Vsign They that put their trust in Him shall understand the truth and such as be faithful in love shall abide with Him.\end{Resp}
\begin{Resp}\Rsign For the grace of God, and His peace, are with His chosen.\end{Resp}
\begin{Resp}\Vsign Glory be to the Father, and to the Son, \Star{} and to the Holy Ghost.\end{Resp}
\begin{Resp}\Rsign For the grace of God, and His peace, are with His chosen.\end{Resp}
\switchcolumn*

\LA
\LessonHead{Lectio VII}
\switchcolumn
\EN
\LessonHead{Lesson VII}
\switchcolumn*

\LA
\LessonTitle{Léctio sancti Evangélii secúndum Lucam}
\switchcolumn
\EN
\LessonTitle{From the Holy Gospel according to Luke}
\switchcolumn*

\LA
\Cite{Luc 6:17-23}
\switchcolumn
\EN
\Cite{Luke 6:17-23}
\switchcolumn*

\LA
\noindent In illo témpore: Descéndens Jesus de monte, stetit in loco campéstri, et turba discipulórum ejus, et multitúdo copiósa plebis ab omni Judǽa, et Jerúsalem, et marítima, et Tyri, et Sidónis. Et réliqua.\par
\switchcolumn
\EN
\noindent At that time, Jesus came down from the mountain, and stood in the plain, and the company of His disciples, and a great multitude of people out of all Judea, and Jerusalem, and from the sea coast of Tyre and Sidon. And so on.\par
\switchcolumn*

\LA
\LessonTitle{Homilía sancti Ambrósii Epíscopi}
\switchcolumn
\EN
\LessonTitle{Homily by St Ambrose, Bishop (of Milan.)}
\switchcolumn*

\LA
\Source{Lib. 5. in Lucam. Cap. 6., post init.}
\switchcolumn
\EN
\Source{Bk. v. on Luke vi.}
\switchcolumn*

\LA
\noindent Advérte ómnia diligénter, quómodo et cum Apóstolis ascéndat et descéndat ad turbas. Quómodo enim turba nisi in húmili Christum vidéret? Non séquitur ad excélsa, non ascéndit ad sublímia. Dénique ubi descéndit, invénit infírmos: in excélsis enim infírmi esse non possunt. Hinc étiam Matthǽus docet in inferióribus débiles esse sanátos. Prius enim unusquísque sanándus est, ut paulátim virtútibus procedéntibus ascéndere possit ad montem; et ídeo quemque in inferióribus sanat, hoc est, a libídine révocat, injúriam cæcitátis avértit. Ad vúlnera nostra descéndit: ut usu quodam et cópia suæ natúræ, compartícipes nos fáciat esse regni cæléstis.\par
\switchcolumn
\EN
\noindent Mark well how Jesus goeth upward with His disciples, and downward to the multitude. How should the multitude behold Christ, save in a lower place? Such go not up to the things which are above; such attain not to the things which are high. And when Jesus cometh down, He findeth such as are diseased for such like go not up to the heights. Hence also Matthew saith that there were there all sick people, (iv. 23.) Of these every man had need of healing, that, when he had received strength, by and by, he might go up into the mountain. And therefore, being Himself come down, He healeth them in the plain, that is to say, He calleth them away from their lust, and freeth them of their blindness. He cometh down to our wounds, to the end that by a certain use of His nature, and by the abundance thereof, He might make us joint-heirs of the kingdom of heaven.\par
\switchcolumn*

\LA
\begin{Resp}\Rsign Propter testaméntum Dómini, et leges patérnas, Sancti Dei perstitérunt in amóre fraternitátis: \Star{} Quia unus fuit semper spíritus in eis, et una fides.\end{Resp}
\begin{Resp}\Vsign Ecce quam bonum et quam jucúndum habitáre fratres in unum.\end{Resp}
\begin{Resp}\Rsign Quia unus fuit semper spíritus in eis, et una fides.\end{Resp}
\switchcolumn
\EN
\begin{Resp}\Rsign Because of the covenant of the Lord, and the laws of their fathers, the Saints of God abode in brotherly love, \Star{} For one spirit and one faith was ever in them.\end{Resp}
\begin{Resp}\Vsign Behold how good and how pleasant it is for brethren to dwell together in unity.\end{Resp}
\begin{Resp}\Rsign For one spirit and one faith was ever in them.\end{Resp}
\switchcolumn*

\LA
\LessonHead{Lectio VIII}
\switchcolumn
\EN
\LessonHead{Lesson VIII}
\switchcolumn*

\LA
\noindent Beáti páuperes, quia vestrum est regnum Dei. Quátuor tantum beatitúdines sanctus Lucas Domínicas posuit, octo vero sanctus Matthǽus: sed in illis octo istæ quátuor sunt, et in quátuor istis illæ octo. Hic enim quátuor velut virtútes ampléxus est cardináles: ille in illis octo mýsticum númerum reserávit. Pro octáva enim multi inscribúntur Psalmi: et mandátum áccipis octo illis partem dare fortásse benedictiónibus. Sicut enim spei nostræ octáva perféctio est, ita octáva summa virtútum est.\par
\switchcolumn
\EN
\noindent Blessed be ye poor, for your's is the kingdom of God. Saint Luke giveth us but four of the Lord's Beatitudes, and Saint Matthew eight but in those eight are contained these four, and in these four those eight. For in these four are embraced the cardinal virtues and in those eight they are set forth in a number full of mystery. It is written at the head of more than one of the Psalms that they are for the octave, and thou hast received the commandment Give a portion to seven, and also to eight to seven or eight what? Perchance degrees of blessedness. For as this eighth Beatitude doth name the most glorious realization of our hope the kingdom of Heaven so doth it also name the most royal exertion of our strength blessed are they which are persecuted.\par
\switchcolumn*

\LA
\begin{Resp}\Rsign Sancti mei, qui in carne pósiti, certámen habuístis: \Star{} Mercédem labóris ego reddam vobis.\end{Resp}
\begin{Resp}\Vsign Veníte benedícti Patris mei, percípite regnum.\end{Resp}
\begin{Resp}\Rsign Mercédem labóris ego reddam vobis.\end{Resp}
\begin{Resp}\Vsign Glória Patri, et Fílio, \Star{} et Spirítui Sancto.\end{Resp}
\begin{Resp}\Rsign Mercédem labóris ego reddam vobis.\end{Resp}
\switchcolumn
\EN
\begin{Resp}\Rsign O ye My Saints, who, being in the flesh, didst have striving \Star{} I will render unto you a reward of your labours.\end{Resp}
\begin{Resp}\Vsign Come, ye blessed of My Father, inherit the kingdom\end{Resp}
\begin{Resp}\Rsign I will render unto you a reward of your labours.\end{Resp}
\begin{Resp}\Vsign Glory be to the Father, and to the Son, \Star{} and to the Holy Ghost.\end{Resp}
\begin{Resp}\Rsign I will render unto you a reward of your labours.\end{Resp}
\switchcolumn*

\end{paracol}

\par\needspace{10\baselineskip}\addvspace{1.2em}{\centering\small\textsc{Die 12 Martii} · \textsc{12 March}\par}
\Office{S. Gregorii Papæ Confessoris et Ecclesiæ Doctoris}{S. Gregory the Great, Pope, Confessor and Doctor of the Church}{Duplex}
\addcontentsline{toc}{subsection}{Die 12 Martii · S. Gregorii Papæ Confessoris et Ecclesiæ Doctoris\enspace\textit{S. Gregory the Great, Pope, Confessor and Doctor of the Church}}
\begin{paracol}{2}
\LA
\RefLine{Lectiones I–III: ex Commúni Doctoris Pontificis (C4a).}
\switchcolumn
\EN
\RefLine{Lessons I–III: from the Common of a Doctor Bishop (C4a).}
\switchcolumn*

\LA
\RefLine{Lectio IX: de Scriptura occurrente.}
\switchcolumn
\EN
\RefLine{Lesson IX: from the Scripture of the day.}
\switchcolumn*

\LA
\LessonHead{Lectio IV}
\switchcolumn
\EN
\LessonHead{Lesson IV}
\switchcolumn*

\LA
\noindent Gregórius Magnus, Románus, Gordiáni senatóris fílius, adoléscens philosophíæ óperam dedit; et prætório offício functus, patre mórtuo, sex monastéria in Sicília ædificávit, Romæ séptimum sancti Andréæ nómine in suis ǽdibus, prope basílicam sanctórum Joánnis et Pauli ad Clivum Scauri, ubi Hilarióne ac Maximiáno magístris, monáchi vitam proféssus, póstea abbas fuit. Mox diáconus cardinális creátus, Constantinópolim a Pelágio Pontífice ad Tibérium Constantínum imperatórem legátus míttitur: apud quem memorábile étiam illud effécit, quod Eutýchium patriárcham, qui scrípserat contra veram ac tractábilem córporum resurrectiónem, ita convícit, ut ejus librum imperátor in ignem injíceret. Quare Eutýchius paulo post cum in morbum incidísset, instánte morte, pellem manus suæ tenébat, multis præséntibus, dicens: Confíteor quia omnes in hac carne resurgémus.\par
\switchcolumn
\EN
\noindent Gregory the Great was a Roman, the son of Gordian the Senator, (and was born about the year of our Lord 540.) As a young man he studied philosophy, and afterwards discharged the office of Praetor. After his father's death he built six monasteries in Sicily, and a seventh in honour of St. Andrew, in his own house at Rome, hard by the Church of Saints John and Paul at the ascent of the hill Scaurus. In this monastery of St. Andrew, he and his masters, Hilarion and Maximian, professed themselves monks, and Gregory was afterwards Abbot. Later on, he was created a Cardinal Deacon, and sent to Constantinople as legate from Pope Pelagius to the Emperor Tiberius Constantine. Before the Emperor he so successfully disputed against the Patriarch Eutychius, who had denied that our bodies shall verily and indeed rise again, that the Prince threw the book of the said Patriarch into the fire. Eutychius himself also, soon after fell sick, and when he felt death coming on him, he took hold of the skin of his own hand and said in the hearing of many that stood by: I acknowledge that we shall all rise again in this flesh.\par
\switchcolumn*

\LA
\begin{Resp}\Rsign Invéni David servum meum, óleo sancto meo unxi eum: \Star{} Manus enim mea auxiliábitur ei.\end{Resp}
\begin{Resp}\Vsign Nihil profíciet inimícus in eo, et fílius iniquitátis non nocébit ei.\end{Resp}
\begin{Resp}\Rsign Manus enim mea auxiliábitur ei.\end{Resp}
\switchcolumn
\EN
\begin{Resp}\Rsign I have found David My servant, with My holy oil have I anointed him \Star{} For My hand shall help him.\end{Resp}
\begin{Resp}\Vsign The enemy shall prevail nothing against him, nor the son of wickedness afflict him.\end{Resp}
\begin{Resp}\Rsign For My hand shall help him.\end{Resp}
\switchcolumn*

\LA
\LessonHead{Lectio V}
\switchcolumn
\EN
\LessonHead{Lesson V}
\switchcolumn*

\LA
\noindent Romam rédiens, Pelágio pestiléntia subláto, summo ómnium consénsu Póntifex elígitur: quem honórem ne accíperet, quámdiu pótuit, recusávit; nam aliéno vestítu in spelúnca delítuit: ubi deprehénsus indício ígnæa colúmnæ, ad sanctum Petrum consecrátur. In pontificátu multa successóribus doctrínæ ac sanctitátis exémpla relíquit. Peregrínos quotídie ad mensam adhibébat: in quibus et Ángelum, et Dóminum Angelórum peregríni fácie accépit. Páuperes et urbános et extérnos, quorum númerum descríptum habébat, benígne sustentábat. Cathólicam fidem multis locis labefactátam restítuit. Nam Donatístas in Africa, Ariános in Hispánia représsit: Agnoítas Alexandría ejécit. Pállium Syágrio Augustodunénsi epíscopo dare nóluit, nisi neóphytos hæréticos expélleret ex Gállia. Gothos hǽresim Ariánam relínquere coégit. Missis in Británniam doctis et sanctis viris Augustíno et áliis mónachis, ínsulam ad Jesu Christi fidem convértit, vere a Beda presbýtero Angliæ vocátus Apóstolus. Joánnis patriárchæ Constantinopolitáni audáciam fregit, qui sibi universális Ecclésiæ epíscopi nomen arrogábat. Maurítium imperatórem, eos qui mílites fuíssent, mónachos fíeri prohibéntem, a senténtia detérruit.\par
\switchcolumn
\EN
\noindent Gregory returned to Rome, and, Pelagius being dead of a plague, he was unanimously chosen Pope. This honour he refused as long as he could. He disguised himself and took refuge in a cave, but was betrayed by a fiery pillar. Being discovered and overruled, he was consecrated at the grave of St. Peter, \textasciitilde{}(upon the 3rd day of September, in the year 590.) He left behind him many examples of doctrine and holiness to them that have followed him in the Papacy. Every day he brought pilgrims to his table, and among them he entertained not an Angel only, but the very Lord of Angels in the guise of a pilgrim. He tenderly cared for the poor, of whom he kept a list, as well without as within the city. He restored the Catholic faith in many places where it had been overthrown. He fought successfully against the Donatists in Africa and the Arians in Spain. He cleansed Alexandria of the Agnoites. He refused to give the Pall to Syagrius, Bishop of Autun, unless he would expel the Neophyte heretics from Gaul. He caused the Goths to abandon the Arian heresy. He sent into Britain Augustine and diverse other learned and holy monks, who brought the inhabitants of that island to believe in Jesus Christ. Hence Gregory is justly called by Bede, the Priest of Jarrow, the Apostle of England. He rebuked the presumption of John, Patriarch of Constantinople, who had taken to himself the title of Bishop of the Universal Church, and he dissuaded the Emperor Maurice from forbidding soldiers to become monks.\par
\switchcolumn*

\LA
\begin{Resp}\Rsign Pósui adjutórium super poténtem, et exaltávi eléctum de plebe mea: \Star{} Manus enim mea auxiliábitur ei.\end{Resp}
\begin{Resp}\Vsign Invéni David servum meum, óleo sancto meo unxi eum.\end{Resp}
\begin{Resp}\Rsign Manus enim mea auxiliábitur ei.\end{Resp}
\switchcolumn
\EN
\begin{Resp}\Rsign I have laid help upon one that is mighty, and have exalted one chosen out of My people \Star{} For My hand shall help him.\end{Resp}
\begin{Resp}\Vsign I have found David My servant, with My holy oil have I anointed him.\end{Resp}
\begin{Resp}\Rsign For My hand shall help him.\end{Resp}
\switchcolumn*

\LA
\LessonHead{Lectio VI}
\switchcolumn
\EN
\LessonHead{Lesson VI}
\switchcolumn*

\LA
\noindent Ecclésiam ornávit sanctíssimis institútis et légibus. Apud sanctum Petrum coácta sýnodo, multa constítuit: in iis, Ut in Missa Kýrie eléison nóvies repeterétur: ut extra id tempus, quod continétur Septuagésima et Pascha, Alle­lúja dicerétur: ut adderétur in Cánone, Diésque nostros in tua pace dispónas. Litanías, Statiónes, et ecclesiásticum offícium auxit. Quátuor concíliis, Nicǽno, Constantinopolitáno, Ephesíno, Chalcedonénsi, tamquam quátuor Evangéliis honórem habéri vóluit. Epíscopis Sicíliæ, qui ex antíqua ecclesiárum consuetúdine Romam síngulis triénniis conveniébant, quinto quoque anno semel veníre indúlsit. Multos libros confécit: quos cum dictáret, testátus est Petrus diáconus, se Spíritum Sanctum colúmbæ spécie in ejus cápite sæpe vidísse. Admirabília sunt quæ dixit, fecit, scripsit, decrévit, præsértim infírma semper et ægra valetúdine. Qui dénique multis éditis miráculis, pontificátus anno décimo tértio, mense sexto, die décimo, quarto Idus Mártii, qui dies festus a Græcis étiam propter insígnem hujus Pontíficis sapiéntiam ac sanctitátem præcípuo honóre celebrátur, ad cæléstem beatitúdinem evocátus est. Cujus corpus sepúltum est in basílica sancti Petri, prope Secretárium.\par
\switchcolumn
\EN
\noindent Gregory adorned the Church with holy customs and laws. He called together a Synod in the Church of St. Peter, and therein ordained many things; among others, the ninefold repetition of the words Kyrie eleison in the Mass, the saying of the word Alle­luja in the Church service except between Septuagesima inclusive and Easter exclusive, and the addition to the Canon of the Mass of the words Do Thou order all our days in thy peace. He increased the Litanies, the number of the Churches where is held the observance called a Station; and the length of the Church Service. He would that the four Councils of Nice, Constantinople, Ephesus, and Chalcedon should be honoured like four Gospels. He released the Sicilian Bishops from visiting Rome every three years, willing them to come instead once every five years. He was the author of many books, and Peter the Deacon declareth that he often saw the Holy Ghost on his head in the form of a dove when he was dictating them. It is a marvel how much he spoke, did, wrote, and legislated, suffering all the while from a weak and sickly body. He worked many miracles. At last God called him away to be blessed for ever in heaven, in the thirteenth year, sixth month, and tenth day of his Pontificate, being the 12 th day of March, (in the year of salvation 604.) This day is observed by the Greeks, as well as by us, as a festival, on account of the eminent wisdom and holiness of this Pope. His body was buried in the Church of St. Peter, hard by the Private Chapel.\par
\switchcolumn*

\LA
\begin{Resp}\Rsign Iste est, qui ante Deum magnas virtútes operátus est, et omnis terra doctrína ejus repléta est: \Star{} Ipse intercédat pro peccátis ómnium populórum.\end{Resp}
\begin{Resp}\Vsign Iste est, qui contémpsit vitam mundi, et pervénit ad cæléstia regna.\end{Resp}
\begin{Resp}\Rsign Ipse intercédat pro peccátis ómnium populórum.\end{Resp}
\begin{Resp}\Vsign Glória Patri, et Fílio, \Star{} et Spirítui Sancto.\end{Resp}
\begin{Resp}\Rsign Ipse intercédat pro peccátis ómnium populórum.\end{Resp}
\switchcolumn
\EN
\begin{Resp}\Rsign This is he which wrought great wonders before God, and the whole earth is full of his teaching \Star{} May he pray for all people, that their sins may be forgiven unto them\end{Resp}
\begin{Resp}\Vsign This is he which loved not his life in this world, and hath attained unto the kingdom of heaven.\end{Resp}
\begin{Resp}\Rsign May he pray for all people, that their sins may be forgiven unto them\end{Resp}
\begin{Resp}\Vsign Glory be to the Father, and to the Son, \Star{} and to the Holy Ghost.\end{Resp}
\begin{Resp}\Rsign May he pray for all people, that their sins may be forgiven unto them\end{Resp}
\switchcolumn*

\LA
\LessonHead{Lectio VII}
\switchcolumn
\EN
\LessonHead{Lesson VII}
\switchcolumn*

\LA
\LessonTitle{Léctio sancti Evangélii secúndum Matthǽum}
\switchcolumn
\EN
\LessonTitle{From the Holy Gospel according to Matthew}
\switchcolumn*

\LA
\Cite{Matt 16:13-19}
\switchcolumn
\EN
\Cite{Matt 16:13-19}
\switchcolumn*

\LA
\noindent In illo témpore: Venit Jesus in partes Cæsaréæ Philíppi, et interrogábat discípulos suos, dicens: Quem dicunt hómines esse Fílium hóminis? Et réliqua.\par
\switchcolumn
\EN
\noindent At that time: When Jesus came into the coasts of Caesarea Philippi, he asked his disciples, saying, Whom do men say that I the Son of man am? And so on.\par
\switchcolumn*

\LA
\LessonTitle{Homilía sancti Leónis Papæ}
\switchcolumn
\EN
\LessonTitle{A Homily by St. Leo the Pope}
\switchcolumn*

\LA
\Source{Sermo 2 in anniversario assumpt. suæ, ante medium}
\switchcolumn
\EN
\Source{Sermo 2 in anniversario assumpt. suæ ante medium}
\switchcolumn*

\LA
\noindent Cum, sicut evangélica lectióne reserátum est, interrogásset Dóminus discípulos, quem ipsum (multis divérsa opinántibus) créderent; respondissétque beátus Petrus, dicens: Tu es Christus Fílius Dei vivi; Dóminus ait: Beátus es, Simon Bar-Jona, quia caro et sanguis non revelávit tibi, sed Pater meus, qui in cælis est: et ego dico tibi, quia tu es Petrus, et super hanc petram ædificábo Ecclésiam meam, et portæ ínferi non prævalébunt advérsus eam. Et tibi dabo claves regni cælórum: et quodcúmque ligáveris super terram, erit ligátum et in cælis: et quodcúmque sólveris super terram, erit solútum et in cælis. Manet ergo disposítio veritátis, et beátus Petrus, in accépta fortitúdine petræ persevérans, suscépta Ecclésiæ gubernácula non relíquit.\par
\switchcolumn
\EN
\noindent When the Lord, as we read in the Gospel, asked his disciples who did men, amid their diverse speculations, believe him the Son of Man to be, blessed Peter answered and said: Thou art the Christ, the Son of the living God. And the Lord answered and said unto him: Blessed art thou, Simon Bar-Jona: for flesh and blood hath not revealed it unto thee, but my Father, which is in heaven: and I say also unto thee: That thou art Peter, and upon this rock I will build my Church, and the gates of hell shall not prevail against it; and I will give unto thee the keys of the kingdom of heaven; and whatsoever thou shalt bind on earth shall be bound in heaven; and whatsoever thou shalt loose on earth shall be loosed in heaven. But the dispensation of truth perdures, and blessed Peter, persevering in the strength of the rock which he hath received, hath not relinquished the position he assumed at the helm of the Church.\par
\switchcolumn*

\LA
\begin{Resp}\Rsign Amávit eum Dóminus, et ornávit eum: stolam glóriæ índuit eum, \Star{} Et ad portas paradísi coronávit eum.\end{Resp}
\begin{Resp}\Vsign Induit eum Dóminus lorícam fídei, et ornávit eum.\end{Resp}
\begin{Resp}\Rsign Et ad portas paradísi coronávit eum.\end{Resp}
\switchcolumn
\EN
\begin{Resp}\Rsign The Lord loved him and beautified him He clothed him with a robe of glory \Star{} And crowned him at the gates of Paradise.\end{Resp}
\begin{Resp}\Vsign The Lord hath put on him the breast-plate of faith, and hath adorned him.\end{Resp}
\begin{Resp}\Rsign And crowned him at the gates of Paradise.\end{Resp}
\switchcolumn*

\LA
\LessonHead{Lectio VIII}
\switchcolumn
\EN
\LessonHead{Lesson VIII}
\switchcolumn*

\LA
\noindent In univérsa namque Ecclésia, Tu es Christus Fílius Dei vivi, quotídie Petrus dicit; et omnis lingua, quæ confitétur Dóminum, magistério hujus vocis imbúitur. Hæc fides diábolum vincit et captivórum ejus víncula dissólvit. Hæc érutos mundo, ínserit cælo, et portæ ínferi advérsus eam prævalére non possunt. Tanta enim divínitus soliditáte muníta est, ut eam neque hærética umquam corrúmpere právitas, nec pagána potúerit superáre perfídia. His ítaque modis, dilectíssimi, rationábile obséquio celebrétur hodiérna festívitas: ut in persóna humilitátis meæ ille intelligátur, ille honorétur, in quo et ómnium pastórum sollicitúdo, cum commendatárum sibi óvium custódia persevérat, et cujus étiam dígnitas in indígno heréde non déficit.\par
\switchcolumn
\EN
\noindent In the universal Church it is as if Peter were still saying every day: Thou art the Christ, the Son of the living God. For every tongue which confesseth the Lord is taught that confession by the teaching of Peter. This is the Faith that overcometh the devil and looseth the bonds of his prisoners. This is the Faith which maketh men free of the world and bringeth them to heaven, and the gates of hell are impotent to prevail against it. This is the rock which God hath fortified with such ramparts of salvation, that the contagion of heresy will never be able to infect it, nor idolatry and unbelief to overcome it. And therefore, dearly beloved, we celebrate today’s festival with reasonable obedience, that in my humble person he may be acknowledged and honoured who doth continue to care for all the shepherds as well as sheep entrusted unto him, and who doth lose none of his dignity even in an unworthy successor.\par
\switchcolumn*

\LA
\begin{Resp}\Rsign In médio Ecclésiæ apéruit os ejus, \Star{} Et implévit eum Dóminus spíritu sapiéntiæ et intelléctus.\end{Resp}
\begin{Resp}\Vsign Jucunditátem et exsultatiónem thesaurizávit super eum.\end{Resp}
\begin{Resp}\Rsign Et implévit eum Dóminus spíritu sapiéntiæ et intelléctus.\end{Resp}
\begin{Resp}\Vsign Glória Patri, et Fílio, \Star{} et Spirítui Sancto.\end{Resp}
\begin{Resp}\Rsign Et implévit eum Dóminus spíritu sapiéntiæ et intelléctus.\end{Resp}
\switchcolumn
\EN
\begin{Resp}\Rsign In the midst of the congregation did the Lord open his mouth. \Star{} And filled him with the spirit of wisdom and understanding.\end{Resp}
\begin{Resp}\Vsign He made him rich with joy and gladness.\end{Resp}
\begin{Resp}\Rsign And filled him with the spirit of wisdom and understanding.\end{Resp}
\begin{Resp}\Vsign Glory be to the Father, and to the Son, \Star{} and to the Holy Ghost.\end{Resp}
\begin{Resp}\Rsign And filled him with the spirit of wisdom and understanding.\end{Resp}
\switchcolumn*

\end{paracol}

\par\needspace{10\baselineskip}\addvspace{1.2em}{\centering\small\textsc{Die 17 Martii} · \textsc{17 March}\par}
\Office{S. Patricii Episcopi et Confessoris}{St. Patrick, Bishop and Confessor}{Duplex}
\addcontentsline{toc}{subsection}{Die 17 Martii · S. Patricii Episcopi et Confessoris\enspace\textit{St. Patrick, Bishop and Confessor}}
\begin{paracol}{2}
\LA
\RefLine{Lectiones I–III: ex Commúni unius Confessoris Pontificis (C4).}
\switchcolumn
\EN
\RefLine{Lessons I–III: from the Common of a Confessor Bishop (C4).}
\switchcolumn*

\LA
\RefLine{Lectiones VII–VIII: ex Commúni unius Confessoris Pontificis (C4).}
\switchcolumn
\EN
\RefLine{Lessons VII–VIII: from the Common of a Confessor Bishop (C4).}
\switchcolumn*

\LA
\RefLine{Lectio IX: de Scriptura occurrente.}
\switchcolumn
\EN
\RefLine{Lesson IX: from the Scripture of the day.}
\switchcolumn*

\LA
\LessonHead{Lectio IV}
\switchcolumn
\EN
\LessonHead{Lesson IV}
\switchcolumn*

\LA
\noindent Patricius, Hiberniæ dictus Apostolus, Calphurnio patre, matre Conchessa sancti Martini Turonensis episcopi, ut perhibent, consanguinea, majori in Britannia natus, puer in barbarorum sæpius incidit captivitatem. Eo in statu, pascendis gregibus præpositus jam tum futuræ sanctitatis specimen præbuit: fidei namque, divinique timoris et amoris spiritu repletus, antelucano tempore per nives, gelu ac pluvias ad preces Deo fundendas, impiger consurgebat; solitus centies interdiu, centiesque noctu Deum orare. A servitute tertio exemptus, et inter clericos relatus, in divinis lectionibus longo se tempore exercuit. Galliis, Italia, insulisque Tyrrheni maris labore summo peragratis, divino tandem monitu ad Hibernorum salutem advocatur; et facta a beato Cælestino Papa Evangelii nuntiandi potestate, consecratusque episcopus, in Hiberniam perrexit.\par
\switchcolumn
\EN
\noindent Patrick, called the Apostle of Ireland, was born in Great Britain. The name of his father was Calphurnius, and that of his mother Conchessa. She is said to have been a relation of St. Martin, Bishop of Tours. When Patrick was a lad, he was several times taken prisoner by savages, and while being in their hands he was employed as a shepherd, he already showed marks of his saintliness to come. His spirit was filled with faith, and love, and fear of God, so that he would rise before the light, in snow, and frost, and rain, to make his prayers to God, being accustomed to address God in prayer an hundred times every day, and an hundred times every night. After being rescued from his third captivity, he was placed among the clergy, and for a long time exercised himself in sacred learning. To this end he travelled with much labour, through Gaul, Italy, and the islands of the Tyrrhenian Sea, but at last being called of God to work for the salvation of the Irish, and, having received from the Blessed Pope Celestine a commission to preach the gospel, and likewise being consecrated a Bishop, he betook himself to Ireland.\par
\switchcolumn*

\LA
\begin{Resp}\Rsign Invéni David servum meum, óleo sancto meo unxi eum: \Star{} Manus enim mea auxiliábitur ei.\end{Resp}
\begin{Resp}\Vsign Nihil profíciet inimícus in eo, et fílius iniquitátis non nocébit ei.\end{Resp}
\begin{Resp}\Rsign Manus enim mea auxiliábitur ei.\end{Resp}
\switchcolumn
\EN
\begin{Resp}\Rsign I have found David My servant, with My holy oil have I anointed him \Star{} For My hand shall help him.\end{Resp}
\begin{Resp}\Vsign The enemy shall prevail nothing against him, nor the son of wickedness afflict him.\end{Resp}
\begin{Resp}\Rsign For My hand shall help him.\end{Resp}
\switchcolumn*

\LA
\LessonHead{Lectio V}
\switchcolumn
\EN
\LessonHead{Lesson V}
\switchcolumn*

\LA
\noindent Eo in munere mirum quot vir apostolicus mala, quot ærumnas et labores, quot pertulerit adversarios. Verum Dei afflante benignitate, terra illa, idolorum antea cultrix, eum mox prædicante Patricio fructum dedit, ut sanctorum insula deínde fuerit appellata. Frequentissimi ab eo populi sacro sunt regenerati lavacro: episcopi, clericique plurimi ordinati;\par
virgines ac viduæ ad continentiæ leges institutæ. Armachanam sedem, Romani Pontificis auctoritate, totius insulæ principem metropolim constituit, sanctorumque reliquiis ab Urbe relatis decoravit. Supernis visionibus, prophetiæ dono, ingentibusque signis et prodigiis a Deo exornatus adeo refulsit, ut longe lateque celebrior Patricii se fama diffuderit.\par
\switchcolumn
\EN
\noindent In the discharge of his calling it is a marvel with how many evils, with how many sufferings and labours, and with how many adversaries the Apostolic Patrick had to bear. Nevertheless, by the goodness of God, that island, which had up to that time been given over to the serving of idols, was, through the preaching of Patrick, so wrought on that she soon brought forth the fruit which won her the name of the Island of Saints. Patrick caused many of her people to be born again by the washing of regeneration; he ordained many Bishops and clerks; he decreed rules for virgins and widows living in continency. By the authority of the Bishop of Rome he established the See of Armagh as the Primatial See of all Ireland, and enriched the Church with relics of the Saints brought from Rome. Patrick, moreover, was so eminently adorned with heavenly visions, with the gift of prophecy, and with great signs and wonders from God, that the fame of him spread itself abroad more and more, day by day.\par
\switchcolumn*

\LA
\begin{Resp}\Rsign Pósui adjutórium super poténtem, et exaltávi eléctum de plebe mea: \Star{} Manus enim mea auxiliábitur ei.\end{Resp}
\begin{Resp}\Vsign Invéni David servum meum, óleo sancto meo unxi eum.\end{Resp}
\begin{Resp}\Rsign Manus enim mea auxiliábitur ei.\end{Resp}
\switchcolumn
\EN
\begin{Resp}\Rsign I have laid help upon one that is mighty, and have exalted one chosen out of My people \Star{} For My hand shall help him.\end{Resp}
\begin{Resp}\Vsign I have found David My servant, with My holy oil have I anointed him.\end{Resp}
\begin{Resp}\Rsign For My hand shall help him.\end{Resp}
\switchcolumn*

\LA
\LessonHead{Lectio VI}
\switchcolumn
\EN
\LessonHead{Lesson VI}
\switchcolumn*

\LA
\noindent Præter quotidianam ecclesiarum sollicitudinem, invictum ab oratione spiritum numquam relaxabat. Ajunt enim, integrum quotidie Psalterium, una cum Canticis et Hymnis, ducentisque orationibus, consuevisse recitare, tercenties per dies singulos flexis genibus Deum adorare, ac in qualibet hora diei canonica centies se crucis signo munire. Noctem tria in spatia distribuens, primum in centum Psalmis percurrendis et bis centies genuflectendo alterum in reliquis quinquaginta Psalmis, algidis aquis immersus, ac corde, oculis manibusque ad cælum erectus, absolvendis insumebat: tertium vero super nudum lapidem stratus tenui dabat quieti. Humilitatis eximius cultor, apostolico more a manuum suarum labore non abstinuit. Assiduis tandem curis pro Ecclesia consumptus, verbo et opere clarus, in extrema senectute divinis mysteriis refectus, obdormivit in Domino, sepultusque est apud Dunum in Ultonia, a christiana salute sæculo quinto.\par
\switchcolumn
\EN
\noindent Besides that which came upon him daily, the care of all the Churches of Ireland, he never suffered his spirit to weary in constant prayer. They say that it was his custom to repeat every day the whole Book of Psalms, together with Songs and Hymns, and two hundred Prayers; that he bent his knees to God in worship three hundred times every day, and that he made on himself the sign of the Cross an hundred times at each of the Seven Hours of the Church Service. He divided the night into three portions; during the first he repeated the first hundred Psalms, and bent his knees two hundred times; during the second he remained plunged in cold water, with heart, eyes, and hands lifted up to heaven, and in that state repeated the remaining fifty Psalms; during the third he took his short rest, lying upon a bare stone. He was a great practicer of lowliness, and, after the pattern of the Apostle, he always continued to work with his own hands. At last he fell asleep in the Lord in extreme old age, refreshed with the Divine Mysteries, worn out with unceasing care for the Churches, and glorious both in word and work. His body is buried in Down in Ulster. He passed away in the fifth century after the giving of salvation by Christ.\par
\switchcolumn*

\LA
\begin{Resp}\Rsign Iste est, qui ante Deum magnas virtútes operátus est, et omnis terra doctrína ejus repléta est: \Star{} Ipse intercédat pro peccátis ómnium populórum.\end{Resp}
\begin{Resp}\Vsign Iste est, qui contémpsit vitam mundi, et pervénit ad cæléstia regna.\end{Resp}
\begin{Resp}\Rsign Ipse intercédat pro peccátis ómnium populórum.\end{Resp}
\begin{Resp}\Vsign Glória Patri, et Fílio, \Star{} et Spirítui Sancto.\end{Resp}
\begin{Resp}\Rsign Ipse intercédat pro peccátis ómnium populórum.\end{Resp}
\switchcolumn
\EN
\begin{Resp}\Rsign This is he which wrought great wonders before God, and the whole earth is full of his teaching \Star{} May he pray for all people, that their sins may be forgiven unto them\end{Resp}
\begin{Resp}\Vsign This is he which loved not his life in this world, and hath attained unto the kingdom of heaven.\end{Resp}
\begin{Resp}\Rsign May he pray for all people, that their sins may be forgiven unto them\end{Resp}
\begin{Resp}\Vsign Glory be to the Father, and to the Son, \Star{} and to the Holy Ghost.\end{Resp}
\begin{Resp}\Rsign May he pray for all people, that their sins may be forgiven unto them\end{Resp}
\switchcolumn*

\end{paracol}

\par\needspace{10\baselineskip}\addvspace{1.2em}{\centering\small\textsc{Die 18 Martii} · \textsc{18 March}\par}
\Office{S. Cyrilli Episcopi Hierosolymitani Confessoris et Ecclesiæ Doctoris}{St. Cyril of Jerusalem, Confessor and Doctor of the Church}{Duplex}
\addcontentsline{toc}{subsection}{Die 18 Martii · S. Cyrilli Episcopi Hierosolymitani Confessoris et Ecclesiæ Doctoris\enspace\textit{St. Cyril of Jerusalem, Confessor and Doctor of the Church}}
\begin{paracol}{2}
\LA
\RefLine{Lectiones I–III: ex Commúni Doctoris Pontificis (C4a).}
\switchcolumn
\EN
\RefLine{Lessons I–III: from the Common of a Doctor Bishop (C4a).}
\switchcolumn*

\LA
\RefLine{Lectio IX: de Scriptura occurrente.}
\switchcolumn
\EN
\RefLine{Lesson IX: from the Scripture of the day.}
\switchcolumn*

\LA
\LessonHead{Lectio IV}
\switchcolumn
\EN
\LessonHead{Lesson IV}
\switchcolumn*

\LA
\noindent Cyrillus Hierosolymitanus a teneris annis divinarum scripturarum studio summopere deditus, adeo in earum scientia profecit, ut orthodoxæ fidei strenuus assertor evaserit. Monasticis institutis imbutus, perpetuæ continentiæ, omnique severiori vivendi rationi se addictum voluit. Postquam a sancto Maximo Hierosolymæ episcopo presbyter ordinatus fuit, munus verbi divini fidelibus prædicandi et catechumenos edocendi summa cum laude implevit, atque illas vere mirandas conscripsit catecheses, quibus totam ecclesiasticam doctrinam dilucide et copiose complexus, singula religionis dogmata contra fidei hostes solide propugnavit. Ita vero in his enucleate et distincte disseruit, ut non solum jam exortas hæreses, sed futuras étiam quasi præsagiens everterit, quemadmodum præstitit asserendo Corporis et Sanguinis Christi realem præsentiam in mirabili Eucharistiæ sacramento. Vita autem functo sancto Maximo, a provinciæ episcopis in illius locum suffectus est.\par
\switchcolumn
\EN
\noindent Cyril of Jerusalem was given to the study of the Holy Scriptures from a child, and so learnt therein that he became an eminent champion of the orthodox faith. He embraced the monastic institute in perpetual continency, and all hardship of living. He was ordained Priest by holy Maximus, Patriarch of Jerusalem, and undertook with eminent success the task of preaching the word of God to the faithful and of instructing the catechumens. Thus did he compose those truly wonderful Catecheses, wherein he has embraced, clearly and fully, all the teaching of the Church, and stoutly defended every one of her doctrines against the enemies of the faith. His treatment of these subjects was such that he has overthrown therein, not only the heresies which had then come into being, but, by a kind of foreknowledge, even those which were to arise in later times. Of this an instance is his contention for the real Presence of the Body and Blood of Christ in the wondrous Sacrament of the Eucharist. After the death of holy Maximus, the bishops of the province chose Cyril in his place.\par
\switchcolumn*

\LA
\begin{Resp}\Rsign Invéni David servum meum, óleo sancto meo unxi eum: \Star{} Manus enim mea auxiliábitur ei.\end{Resp}
\begin{Resp}\Vsign Nihil profíciet inimícus in eo, et fílius iniquitátis non nocébit ei.\end{Resp}
\begin{Resp}\Rsign Manus enim mea auxiliábitur ei.\end{Resp}
\switchcolumn
\EN
\begin{Resp}\Rsign I have found David My servant, with My holy oil have I anointed him \Star{} For My hand shall help him.\end{Resp}
\begin{Resp}\Vsign The enemy shall prevail nothing against him, nor the son of wickedness afflict him.\end{Resp}
\begin{Resp}\Rsign For My hand shall help him.\end{Resp}
\switchcolumn*

\LA
\LessonHead{Lectio V}
\switchcolumn
\EN
\LessonHead{Lesson V}
\switchcolumn*

\LA
\noindent In episcopatu injurias multas et calamitates, non secus ac beatus Athanasius, cui cœvus erat, ab Arianorum factionibus fidei causa perpessus fuit. Hi enim ægre ferentes Cyrillum vehementer hæresibus obsistere, ipsum calumniis aggrediuntur, et in conciliabulo depositum e sua sede deturbant. Quorum furori ut se subtraheret, Tarsum Ciliciæ aufugit, et quoad vixit Constantius, exsilii rigorem pertulit. Post illius mortem, Juliano Apostata ad imperium evecto, Hierosolymam redire potuit, ubi ardenti zelo gregi suo ab erroribus et a vitiis revocando operam navavit. Sed iterum, Valente imperatore, exsulare coactus est, donec, reddita Ecclesiæ pace per Theodosium Magnum, et Arianorum crudelitate audaciaque repressa, ab eodem imperatore tamquam fortissimus Christi athleta honoribus susceptus suæ sedi restitutus fuit. Quam strenue et sancte sublimis officii sui munia impleverit, luculenter apparet ex florenti tunc temporis Hierosolymitanæ ecclesiæ statu, quem sanctus Basilius loca sancta veneraturus, ibi aliquamdiu commoratus, describit\par
\switchcolumn
\EN
\noindent In his office of Bishop he had for the faith's sake, like his blessed contemporary Athanasius, to endure many wrongs and sufferings at the hands of the Arian sect. The Arians could not bear that Cyril should steadfastly withstand their heresy. They assailed him with calumnies, deposed him in a pretended council, and drove him out of his see. To escape their rage he fled to Tarsus in Cilicia, and as long as Constantius lived he bore the hardships of exile. After his death and the accession to the imperial throne of the Apostate Julian, Cyril was able to return to Jerusalem, where he set himself with burning zeal to deliver his flock from false doctrine and from sin. He was driven into exile a second time under the Emperor Valens. But when peace was restored to the Church by Theodosius the Great, and the cruelty and insolence of the Arians were restrained, Cyril was received with honour by the Emperor as one of Christ's most eminent soldiers, and was restored to his see. With what earnestness and holiness he fulfilled the duties of his exalted office was made manifest by the flourishing state of the church of Jerusalem at that time, of which a picture hath been left for us by holy Basil, who dwelt there for a while when he went to worship at the holy places.\par
\switchcolumn*

\LA
\begin{Resp}\Rsign Pósui adjutórium super poténtem, et exaltávi eléctum de plebe mea: \Star{} Manus enim mea auxiliábitur ei.\end{Resp}
\begin{Resp}\Vsign Invéni David servum meum, óleo sancto meo unxi eum.\end{Resp}
\begin{Resp}\Rsign Manus enim mea auxiliábitur ei.\end{Resp}
\switchcolumn
\EN
\begin{Resp}\Rsign I have laid help upon one that is mighty, and have exalted one chosen out of My people \Star{} For My hand shall help him.\end{Resp}
\begin{Resp}\Vsign I have found David My servant, with My holy oil have I anointed him.\end{Resp}
\begin{Resp}\Rsign For My hand shall help him.\end{Resp}
\switchcolumn*

\LA
\LessonHead{Lectio VI}
\switchcolumn
\EN
\LessonHead{Lesson VI}
\switchcolumn*

\LA
\noindent Venerandi hujus præsulis sanctitatem cælestibus signis a Deo fuísse illustratam memoriæ traditum accepimus. Inter hæc recensetur præclara crucis, solis radiis fulgentioris, apparitio, quæ episcopatus ejus initia decoravit. Hujusmodi prodigii ethnici et Christiani testes oculares fuerunt cum ipso Cyrillo qui gratiis primum in Ecclesia Deo redditis, illud per epistolam Constantio imperatori narravit. Nec minus admiratione dignum, quod Judæis templum a Tito eversum restaurare ex impio imperatoris Juliani jussu conantibus, evenit. Vehementi enim terræmotu oborto, et ingentibus flammarum globis e terra erumpentibus, omnia opera ignis consumpsit, ita ut Judæi et Julianus deterriti, ab incepto destiterint; prout scilicet indubitanter futurum Cyrillus prædixerat. Qui demum paulo ante obitum Concilio œcumenico Constantinopolitano interfuit, in quo Macedonii hæresis, et iterum Ariana condemnata est. Ac Jerúsalem inde reversus, fere septuagenarius, trigesimo quinto sui episcopatus anno sancto fine quievit. Ejus Officium ac Missam Leo decimus tertius Pontifex Maximus ab universa Ecclesia celebrari mandavit.\par
\switchcolumn
\EN
\noindent Tradition hath handed down that God Himself crowned with signs from heaven the holiness of this venerable Patriarch. Among these signs is numbered an apparition of a cross, more resplendent than the beams of the sun, which appeared at the beginning of his Patriarchate. Not only Cyril himself, but heathens and Christians alike were eye-witnesses of this marvel, and Cyril first gave thanks to God therefore in the church, and then sent news thereof by letter to the Emperor Constantius. A thing no less wonderful came to pass when the Jews were commanded by the profane Emperor Julian to attempt the restoration of the temple which had been destroyed by Titus. A great earthquake arose, and great masses of fire broke forth from the earth and consumed all the works, so that the Jews and Julian were dismayed and stayed their hand, all the which it can be proved that Cyril had foretold. A little while before his death he was present at the second Council of Constantinople; herein was condemned the heresy of Macedonius, and once more the Arian heresy. After his return to Jerusalem he died a holy death in the 69th year of his age and the 35 th of his episcopate. The Supreme Pontiff Leo XIII. commanded that his office and Mass should be celebrated throughout the universal Church.\par
\switchcolumn*

\LA
\begin{Resp}\Rsign Iste est, qui ante Deum magnas virtútes operátus est, et omnis terra doctrína ejus repléta est: \Star{} Ipse intercédat pro peccátis ómnium populórum.\end{Resp}
\begin{Resp}\Vsign Iste est, qui contémpsit vitam mundi, et pervénit ad cæléstia regna.\end{Resp}
\begin{Resp}\Rsign Ipse intercédat pro peccátis ómnium populórum.\end{Resp}
\begin{Resp}\Vsign Glória Patri, et Fílio, \Star{} et Spirítui Sancto.\end{Resp}
\begin{Resp}\Rsign Ipse intercédat pro peccátis ómnium populórum.\end{Resp}
\switchcolumn
\EN
\begin{Resp}\Rsign This is he which wrought great wonders before God, and the whole earth is full of his teaching \Star{} May he pray for all people, that their sins may be forgiven unto them\end{Resp}
\begin{Resp}\Vsign This is he which loved not his life in this world, and hath attained unto the kingdom of heaven.\end{Resp}
\begin{Resp}\Rsign May he pray for all people, that their sins may be forgiven unto them\end{Resp}
\begin{Resp}\Vsign Glory be to the Father, and to the Son, \Star{} and to the Holy Ghost.\end{Resp}
\begin{Resp}\Rsign May he pray for all people, that their sins may be forgiven unto them\end{Resp}
\switchcolumn*

\LA
\LessonHead{Lectio VII}
\switchcolumn
\EN
\LessonHead{Lesson VII}
\switchcolumn*

\LA
\LessonTitle{Léctio sancti Evangélii secúndum Matthǽum.}
\switchcolumn
\EN
\LessonTitle{From the Holy Gospel according to Matthew}
\switchcolumn*

\LA
\Cite{Matt 10:23-28}
\switchcolumn
\EN
\Cite{Matt 10:23-28}
\switchcolumn*

\LA
\noindent In illo témpore: Dixit Jesus discipulis suis: Cum persequentur vos in civitate ista, fugite in aliam. Et réliqua.\par
\switchcolumn
\EN
\noindent At that time, Jesus said to his disciples: And when they shall persecute you in this city, flee into another. And so on.\par
\switchcolumn*

\LA
\LessonTitle{Homilia sancti Athanasii Epíscopi.}
\switchcolumn
\EN
\LessonTitle{Homily by St. Athanasius, Pope (of Alexandria.)}
\switchcolumn*

\LA
\Source{In Apologia de fuga sua, ante med.}
\switchcolumn
\EN
\Source{Defence of his own flight.}
\switchcolumn*

\LA
\noindent In lege præceptum erat ut constituerentur civitates refugiorum, ut qui quomodocumque ad necem quærerentur, servari possent. In consummatione porro sæculorum cum advenisset illud ipsum Verbum Patris, quod Moysi antea locutum fuerat, rursus hoc præceptum dedit: Cum vos, inquiens, persecuti fuerint in civitate ista, fugite in aliam. Pauloque post subjicit: Cum videritis illam abominationem desolationis quæ dicta est per Danielem prophetam, consistentem in loco sancto (qui legit, intelligat), tunc qui in Judǽa sunt fugiant ad montes et qui in tecto est, ne descendat tollere aliquid de domo sua: et qui in agro est, non revertatur tollere tunicam suam.\par
\switchcolumn
\EN
\noindent It is written in the Law, (Num. xxxv. 11,) Ye shall appoint you cities to be cities of refuge for you, that in these cities they which were pursued to put them to death might enter and be safe. And in the latter days when He was come, even that very Word of the Father, Which had spoken aforetime unto Moses, He gave again the same commandment When they persecute you in this city, flee ye into another. And, a while afterward, He said: “When ye shall see the abomination of desolation, spoken of by Daniel the Prophet, stand in the Holy Place, whoso readeth, let him understand, then let them which be in Judaea flee unto the mountains; let him which is on the house-top not come down to take anything out of his house; neither let him which is in the field return back to take his clothes.” (Matth. xxiv. 15-18.)\par
\switchcolumn*

\LA
\begin{Resp}\Rsign Amávit eum Dóminus, et ornávit eum: stolam glóriæ índuit eum, \Star{} Et ad portas paradísi coronávit eum.\end{Resp}
\begin{Resp}\Vsign Induit eum Dóminus lorícam fídei, et ornávit eum.\end{Resp}
\begin{Resp}\Rsign Et ad portas paradísi coronávit eum.\end{Resp}
\switchcolumn
\EN
\begin{Resp}\Rsign The Lord loved him and beautified him He clothed him with a robe of glory \Star{} And crowned him at the gates of Paradise.\end{Resp}
\begin{Resp}\Vsign The Lord hath put on him the breast-plate of faith, and hath adorned him.\end{Resp}
\begin{Resp}\Rsign And crowned him at the gates of Paradise.\end{Resp}
\switchcolumn*

\LA
\LessonHead{Lectio VIII}
\switchcolumn
\EN
\LessonHead{Lesson VIII}
\switchcolumn*

\LA
\noindent Hæc cum scirent Sancti, ejusmodi tenuérunt suæ conversatiónis institutum. Quæ enim nunc præcepit Dóminus, éadem quoque ante suum in carne advéntum locútus est in Sanctis; et hoc institutum hómines ad perfectiónem ducit. Nam quod Deus jusserit, id omnino faciéndum est. Ideoque et ipsum Verbum propter nos homo factum, non indignum putávit, cum quærerétur, quemádmodum et nos, abscóndere se; et cum persecutiónem paterétur, fugere, et insidias declináre: cum autem a se definitum tempus ipse adduxísset, in quo corporáliter pro ómnibus pati volebat, ultro seípsum trádidit insidiántibus.\par
\switchcolumn
\EN
\noindent The Saints, therefore, knowing these words of the Lord, have obeyed them in their lives. What the Lord hath now commanded by His Own Mouth He commanded through His Saints before that He Himself was come in the flesh, and to obey this commandment worketh in a man perfection, since whatever God commandeth is a thing which it behoveth man to do. For this cause, that very Word of God Which was made flesh for our sake thought it meet when they sought Him, even as at this present time they are seeking us, to hide Himself, (John viii. 59,) and, when they persecuted Him, to fly and escape from their laying in wait for Him although when that time came which He had Himself decreed, and wherein He willed, as touching the Body, to suffer for us all, He willingly gave Himself up to His enemies.\par
\switchcolumn*

\LA
\begin{Resp}\Rsign In médio Ecclésiæ apéruit os ejus, \Star{} Et implévit eum Dóminus spíritu sapiéntiæ et intelléctus.\end{Resp}
\begin{Resp}\Vsign Jucunditátem et exsultatiónem thesaurizávit super eum.\end{Resp}
\begin{Resp}\Rsign Et implévit eum Dóminus spíritu sapiéntiæ et intelléctus.\end{Resp}
\begin{Resp}\Vsign Glória Patri, et Fílio, \Star{} et Spirítui Sancto.\end{Resp}
\begin{Resp}\Rsign Et implévit eum Dóminus spíritu sapiéntiæ et intelléctus.\end{Resp}
\switchcolumn
\EN
\begin{Resp}\Rsign In the midst of the congregation did the Lord open his mouth. \Star{} And filled him with the spirit of wisdom and understanding.\end{Resp}
\begin{Resp}\Vsign He made him rich with joy and gladness.\end{Resp}
\begin{Resp}\Rsign And filled him with the spirit of wisdom and understanding.\end{Resp}
\begin{Resp}\Vsign Glory be to the Father, and to the Son, \Star{} and to the Holy Ghost.\end{Resp}
\begin{Resp}\Rsign And filled him with the spirit of wisdom and understanding.\end{Resp}
\switchcolumn*

\end{paracol}

\par\needspace{10\baselineskip}\addvspace{1.2em}{\centering\small\textsc{Die 19 Martii} · \textsc{19 March}\par}
\Office{S. Joseph Sponsi B.M.V. Confessoris}{St. Joseph, Spouse of the Blessed Virgin Mary, Confessor}{Duplex I classis}
\addcontentsline{toc}{subsection}{Die 19 Martii · S. Joseph Sponsi B.M.V. Confessoris\enspace\textit{St. Joseph, Spouse of the Blessed Virgin Mary, Confessor}}
\begin{paracol}{2}
\LA
\RefLine{Lectio IX: de Scriptura occurrente.}
\switchcolumn
\EN
\RefLine{Lesson IX: from the Scripture of the day.}
\switchcolumn*

\LA
\LessonHead{Lectio I}
\switchcolumn
\EN
\LessonHead{Lesson I}
\switchcolumn*

\LA
\LessonTitle{De libro Génesis}
\switchcolumn
\EN
\LessonTitle{Lesson from the book of Genesis}
\switchcolumn*

\LA
\Cite{Gen 39:1-5}
\switchcolumn
\EN
\Cite{Gen 39:1-5}
\switchcolumn*

\LA
\noindent Joseph ígitur ductus est in Ægýptum, emítque eum Putíphar eunúchus Pharaónis, princeps exércitus, vir Ægýptius, de manu Ismaelitárum, a quibus perdúctus erat. Fuítque Dóminus cum eo, et erat vir in cunctis próspere agens: habitavítque in domo dómini sui, qui óptime nóverat Dóminum esse cum eo, et ómnia, quæ gerébat, ab eo dírigi in manu illíus. Invenítque Joseph grátiam coram dómino suo, et ministrábat ei: a quo præpósitus ómnibus gubernábat créditam sibi domum, et univérsa quæ ei trádita fúerant: benedixítque Dóminus dómui Ægýptii propter Joseph.\par
\switchcolumn
\EN
\noindent And Joseph was brought into Egypt, and Putiphar an eunuch of Pharao, chief captain of the army, an Egyptian, bought him of the Ismaelites, by whom he was brought. And the Lord was with him, and he was a prosperous man in all things: and he dwelt in his master's house, Who knew very well that the Lord was with him, and made all that he did to prosper in his hand. And Joseph found favour in the sight of his master, and ministered to him: and being set over all by him, he governed the house committed to him, and all things that were delivered to him: And the Lord blessed the house of the Egyptian for Joseph's sake.\par
\switchcolumn*

\LA
\begin{Resp}\Rsign Fuit Dóminus cum Joseph, et dedit ei grátiam in conspéctu príncipis cárceris: \Star{} Qui trádidit in manu illíus univérsos vinctos.\end{Resp}
\begin{Resp}\Vsign Quidquid fiébat, sub ipso erat: Dóminus enim erat cum illo, et ómnia ópera ejus dirigébat.\end{Resp}
\begin{Resp}\Rsign Qui trádidit in manu illíus univérsos vinctos.\end{Resp}
\switchcolumn
\EN
\begin{Resp}\Rsign The Lord was with Joseph, and gave him favour in the sight of the keeper of the prison. \Star{} And the keeper of the prison committed to Joseph's hand all the prisoners that were in the prison.\end{Resp}
\begin{Resp}\Vsign And whatsoever they did there, he was the doer of it because the Lord was with him, and that which he did, the Lord made it to prosper.\end{Resp}
\begin{Resp}\Rsign And the keeper of the prison committed to Joseph's hand all the prisoners that were in the prison.\end{Resp}
\switchcolumn*

\LA
\LessonHead{Lectio II}
\switchcolumn
\EN
\LessonHead{Lesson II}
\switchcolumn*

\LA
\Cite{Gen 41:37-40}
\switchcolumn
\EN
\Cite{Gen 41:37-40}
\switchcolumn*

\LA
\noindent Plácuit Pharaóni consílium et cunctis minístris ejus: locutúsque est ad eos: Num inveníre potérimus talem virum, qui spíritu Dei plenus sit? Dixit ergo ad Joseph: Quia osténdit tibi Deus ómnia quæ locútus es, numquid sapientiórem et consímilem tui inveníre pótero? Tu eris super domum meam, et ad tui oris impérium cunctus pópulus obédiet: uno tantum regni sólio te præcédam.\par
\switchcolumn
\EN
\noindent The counsel pleased Pharao and all his servants. And he said to them: Can we find such another man, that is full of the spirit of God? He said therefore to Joseph: Seeing God hath shewn thee all that thou hast said, can I find one wiser and one like unto thee? Thou shalt be over my house, and at the commandment of thy mouth all the people shall obey: only in the kingly throne will I be above thee.\par
\switchcolumn*

\LA
\begin{Resp}\Rsign Esuriénte terra Ægýpti, clamávit pópulus ad regem, aliménta petens. Quibus ille respóndit: \Star{} Ite ad Joseph, et quidquid vobis díxerit, fácite.\end{Resp}
\begin{Resp}\Vsign Crescébat cotídie fames in omni terra, aperuítque Joseph univérsa hórrea, et vendébat Ægýptiis.\end{Resp}
\begin{Resp}\Rsign Ite ad Joseph, et quidquid vobis díxerit, fácite.\end{Resp}
\switchcolumn
\EN
\begin{Resp}\Rsign When all the land of Egypt was famished, the people cried to the king for bread. And the king said unto all the Egyptians: \Star{} Go unto Joseph; and what he saith to you, do.\end{Resp}
\begin{Resp}\Vsign The famine was more grievous every day over all the face of the earth, and Joseph opened all the storehouses, and sold unto the Egyptians.\end{Resp}
\begin{Resp}\Rsign Go unto Joseph; and what he saith to you, do.\end{Resp}
\switchcolumn*

\LA
\LessonHead{Lectio III}
\switchcolumn
\EN
\LessonHead{Lesson III}
\switchcolumn*

\LA
\Cite{Gen 41:41-44}
\switchcolumn
\EN
\Cite{Gen 41:41-44}
\switchcolumn*

\LA
\noindent Dixítque rursus Phárao ad Joseph: Ecce, constítui te super univérsam terram Ægýpti. Tulítque ánnulum de manu sua, et dedit eum in manu ejus: vestivítque eum stola býssina, et collo torquem áuream circumpósuit. Fecítque eum ascéndere super currum suum secúndum, clamánte præcóne, ut omnes coram eo genu flécterent, et præpósitum esse scirent univérsæ terræ Ægýpti. Dixit quoque rex ad Joseph: Ego sum Phárao: absque tuo império non movébit quisquam manum aut pedem in omni terra Ægýpti.\par
\switchcolumn
\EN
\noindent And again Pharao said to Joseph: Behold, I have appointed thee over the whole land of Egypt. And he took his ring from his own hand, and gave it into his hand: and he put upon him a robe of silk, and put a chain of gold about his neck. And he made him go up into his second chariot, the crier proclaiming that all should bow their knee before him, and that they should know he was made govenor over the whole land of Egypt. And the king said to Joseph: I am Pharao; without thy commandment no man shall move hand or foot in all the land of Egypt.\par
\switchcolumn*

\LA
\begin{Resp}\Rsign Fecit me Dóminus quasi patrem regis, et dóminum univérsæ domus ejus: nolíte pavére; \Star{} Pro salúte enim vestra misit me Deus ante vos in Ægýptum.\end{Resp}
\begin{Resp}\Vsign Veníte ad me, et ego dabo vobis ómnia bona Ægýpti, et comedétis medúllam terræ.\end{Resp}
\begin{Resp}\Rsign Pro salúte enim vestra misit me Deus ante vos in Ægýptum.\end{Resp}
\begin{Resp}\Vsign Glória Patri, et Fílio, \Star{} et Spirítui Sancto.\end{Resp}
\begin{Resp}\Rsign Pro salúte enim vestra misit me Deus ante vos in Ægýptum.\end{Resp}
\switchcolumn
\EN
\begin{Resp}\Rsign The Lord hath made me as a father to Pharaoh, and lord of all his house fear not \Star{} For God sent me before you into Egypt, to save your lives.\end{Resp}
\begin{Resp}\Vsign Come unto me, and I will give you all the good things of Egypt, and ye shall eat the fat of the land.\end{Resp}
\begin{Resp}\Rsign For God sent me before you into Egypt, to save your lives.\end{Resp}
\begin{Resp}\Vsign Glory be to the Father, and to the Son, \Star{} and to the Holy Ghost.\end{Resp}
\begin{Resp}\Rsign For God sent me before you into Egypt, to save your lives.\end{Resp}
\switchcolumn*

\LA
\LessonHead{Lectio IV}
\switchcolumn
\EN
\LessonHead{Lesson IV}
\switchcolumn*

\LA
\LessonTitle{Sermo sancti Bernárdi Abbátis}
\switchcolumn
\EN
\LessonTitle{From the Sermons of St. Bernard, Abbot of Clairvaux.}
\switchcolumn*

\LA
\Source{Homil. 2 super Missus est, prope finem}
\switchcolumn
\EN
\Source{2nd on Luke i. 26}
\switchcolumn*

\LA
\noindent Quis et qualis homo fúerit beátus Joseph, cónice ex appellatióne, qua, licet dispensatória, méruit honorári ádeo, ut pater Dei et dictus et créditus sit; cónice et ex próprio vocábulo, quod augméntum non dúbitas interpretári. Simul et meménto magni illíus quondam Patriárchæ vénditi in Ægýpto; et scito ipsíus istum non solum vocábulum fuísse sortítum, sed et castimóniam adéptum, innocéntiam assecútum et grátiam.\par
\switchcolumn
\EN
\noindent What and what manner of man the blessed Joseph was, we may gather from that title wherewith, albeit only as a deputy, God deemed him fit to be honoured he was both called, and supposed to be the Father of God. We may gather it from his very name, which, being interpreted, signifieth Increase. Remember likewise that great Patriarch who was sold into Egypt, and know that the Husband of Mary not only received his name, but inherited his purity, and was likened to him in innocence and in grace.\par
\switchcolumn*

\LA
\begin{Resp}\Rsign Ascéndit Joseph a Galilǽa de civitáte Názareth in Judǽam, in civitátem David, quæ vocátur Béthlehem: \Star{} Eo quod esset de domo et família David.\end{Resp}
\begin{Resp}\Vsign Ut profiterétur cum María desponsáta sibi uxóre.\end{Resp}
\begin{Resp}\Rsign Eo quod esset de domo et família David.\end{Resp}
\switchcolumn
\EN
\begin{Resp}\Rsign Joseph went up from Galilee, out of the city of Nazareth, into Judea, unto the city of David, which is called Bethlehem \Star{} Because he was of the house and lineage of David.\end{Resp}
\begin{Resp}\Vsign To be enrolled with Mary his espoused wife.\end{Resp}
\begin{Resp}\Rsign Because he was of the house and lineage of David.\end{Resp}
\switchcolumn*

\LA
\LessonHead{Lectio V}
\switchcolumn
\EN
\LessonHead{Lesson V}
\switchcolumn*

\LA
\noindent Síquidem ille Joseph, fratérna ex invídia vénditus et ductus in Ægýptum, Christi venditiónem præfigurávit: iste Joseph, Herodiánam invídiam fúgiens, Christum in Ægýptum portávit. Ille dómino suo fidem servans, dóminæ nóluit commiscéri: iste Dóminam suam, Dómini sui matrem, vírginem agnóscens et ipse cóntinens, fidéliter custodívit. Illi data est intellegéntia in mystériis somniórum: isti datum est cónscium fíeri atque partícipem cæléstium sacramentórum.\par
\switchcolumn
\EN
\noindent If then, that Joseph that was sold by his brethren through envy, and was brought down to Egypt, was a type of Christ sold by a disciple, and handed over to the Gentiles, the other Joseph flying from the envy of Herod carried Christ into Egypt. That first Joseph kept loyal to his master, and would not carnally know his master's wife; that second Joseph knew that the Lady, the Mother of his Lord, was a virgin, and he himself remained faithfully virgin toward her. To that first Joseph it was given to know dark things in interpreting of dreams; to the second Joseph it was given in sleep to know the mysteries of the kingdom of heaven.\par
\switchcolumn*

\LA
\begin{Resp}\Rsign Surge, et áccipe Púerum et Matrem ejus, et fuge in Ægýptum; \Star{} Et esto ibi, usque dum dicam tibi.\end{Resp}
\begin{Resp}\Vsign Ut adimplerétur quod dictum est a Dómino per prophétam dicéntem: Ex Ægýpto vocávi Fílium meum.\end{Resp}
\begin{Resp}\Rsign Et esto ibi, usque dum dicam tibi.\end{Resp}
\switchcolumn
\EN
\begin{Resp}\Rsign Arise, and take the young Child and His mother, and flee into Egypt; \Star{} And be thou there until I bring thee word.\end{Resp}
\begin{Resp}\Vsign That it might be fulfilled which was spoken of the Lord by the Prophet, saying Out of Egypt have I called My Son.\end{Resp}
\begin{Resp}\Rsign And be thou there until I bring thee word.\end{Resp}
\switchcolumn*

\LA
\LessonHead{Lectio VI}
\switchcolumn
\EN
\LessonHead{Lesson VI}
\switchcolumn*

\LA
\noindent Ille fruménta servávit non sibi, sed omni pópulo: iste Panem vivum e cælo servándum accépit tam sibi quam toti mundo. Non est dúbium, quin bonus et fidélis homo fúerit iste Joseph, cui Mater desponsáta est Salvatóris. Fidélis, inquam, servus, et prudens, quem constítuit Dóminus suæ Matris solátium, suæ carnis nutrítium, solum dénique in terris magni consílii coadjutórem fidelíssimum.\par
\switchcolumn
\EN
\noindent The first Joseph laid by bread, not for himself, but for all people; the second Joseph received into his keeping that Living Bread Which came down from heaven, not for him only, but for the whole world. We cannot doubt but that that Joseph was good and faithful to whom was espoused the Mother of the Saviour. Yea, I say, he was a faithful and wise servant, whom the Lord appointed to be the comfort of His own Mother, the keeper of His own Body, and the only and trusty helper in the Eternal Counsels.\par
\switchcolumn*

\LA
\begin{Resp}\Rsign Cum indúcerent púerum Jesum paréntes ejus, ut fácerent secúndum consuetúdinem legis pro eo, \Star{} Accépit eum Símeon in ulnas suas, et benedíxit Deum.\end{Resp}
\begin{Resp}\Vsign Et erat Pater ejus et Mater mirántes super his, quæ dicebántur de illo.\end{Resp}
\begin{Resp}\Rsign Accépit eum Símeon in ulnas suas, et benedíxit Deum.\end{Resp}
\begin{Resp}\Vsign Glória Patri, et Fílio, \Star{} et Spirítui Sancto.\end{Resp}
\begin{Resp}\Rsign Accépit eum Símeon in ulnas suas, et benedíxit Deum.\end{Resp}
\switchcolumn
\EN
\begin{Resp}\Rsign When His parents brought the Child Jesus into the temple, to do for Him after the custom of the law, \Star{} Simeon took Him up in his arms, and blessed God.\end{Resp}
\begin{Resp}\Vsign And His father and mother marvelled at those things which were spoken of Him.\end{Resp}
\begin{Resp}\Rsign Simeon took Him up in his arms, and blessed God.\end{Resp}
\begin{Resp}\Vsign Glory be to the Father, and to the Son, \Star{} and to the Holy Ghost.\end{Resp}
\begin{Resp}\Rsign Simeon took Him up in his arms, and blessed God.\end{Resp}
\switchcolumn*

\LA
\LessonHead{Lectio VII}
\switchcolumn
\EN
\LessonHead{Lesson VII}
\switchcolumn*

\LA
\LessonTitle{Léctio sancti Evangélii secúndum Matthǽum}
\switchcolumn
\EN
\LessonTitle{Continuation of the Holy Gospel according to Matthew}
\switchcolumn*

\LA
\Cite{Matt 1:18-21}
\switchcolumn
\EN
\Cite{Matt 1:18-21}
\switchcolumn*

\LA
\noindent Cum esset desponsáta Mater Jesu María Joseph, ántequam convenírent, invénta est in útero habens de Spíritu Sancto. Et réliqua.\par
\switchcolumn
\EN
\noindent When as his mother Mary was espoused to Joseph, before they came together, she was found with child, of the Holy Ghost. And so on.\par
\switchcolumn*

\LA
\LessonTitle{Homilía sancti Hierónymi Presbýteri}
\switchcolumn
\EN
\LessonTitle{Homily by St. Jerome, Priest (at Bethlehem.)}
\switchcolumn*

\LA
\Source{Lib. 1 Comment. in c. 1 Matth.}
\switchcolumn
\EN
\Source{Book of Commentaries, on Matth. i.}
\switchcolumn*

\LA
\noindent Quare non de símplici vírgine, sed de desponsáta concípitur? Primum, ut per generatiónem Joseph, orígo Maríæ monstrarétur: secúndo, ne lapidarétur a Judǽis ut adúltera: tértio, ut in Ægýptum fúgiens habéret solátium. Martyr Ignátius étiam quartam áddidit causam, cur a desponsáta concéptus sit: Ut partus, ínquiens, ejus celarétur diábolo, dum eum putat non de vírgine, sed de uxóre generátum.\par
\switchcolumn
\EN
\noindent Why was the Lord conceived of an espoused virgin rather than of a free? First, for the sake of the genealogy of Mary, which we have obtained by that of Joseph. Secondly, because she was thus saved from being stoned by the Jews as an adulteress. Thirdly, that Himself and His mother might have a guardian on their journey into Egypt. To these, Ignatius, the martyr of Antioch, has added a fourth reason namely, that the birth might take place unknown to the devil, who would naturally suppose that Mary had conceived by Joseph.\par
\switchcolumn*

\LA
\begin{Resp}\Rsign Dicit Mater Jesu ad illum: Fili, quid fecísti nobis sic? \Star{} Ecce pater tuus et ego doléntes quærebámus te.\end{Resp}
\begin{Resp}\Vsign Et ait ad illos: Quid est, quod me quærebátis? nesciebátis, quia in his quæ Patris mei sunt, opórtet me esse?\end{Resp}
\begin{Resp}\Rsign Ecce pater tuus et ego doléntes quærebámus te.\end{Resp}
\switchcolumn
\EN
\begin{Resp}\Rsign The Mother of Jesus said unto Him Son, why hast Thou thus dealt with us? \Star{} Behold, thy father and I have sought thee sorrowing.\end{Resp}
\begin{Resp}\Vsign And He said unto them: How is it that ye sought Me? Wist ye not that I must be about My Father's business?\end{Resp}
\begin{Resp}\Rsign Behold, thy father and I have sought thee sorrowing.\end{Resp}
\switchcolumn*

\LA
\LessonHead{Lectio VIII}
\switchcolumn
\EN
\LessonHead{Lesson VIII}
\switchcolumn*

\LA
\noindent Antequam convenírent, invénta est in útero habens de Spíritu Sancto. Non ab álio invénta est, nisi a Joseph, qui pene licéntia maritáli futúræ uxóris ómnia nóverat. Quod autem dícitur, Antequam convenírent, non séquitur ut póstea convénerint: sed Scriptúra quod factum non sit, osténdit.\par
\switchcolumn
\EN
\noindent Before they came together, she was found with child of the Holy Ghost. She was found, that is, by Joseph, but by no one else. He had already almost an husband's privilege to know all that concerned her. Before they came together. This doth not imply that they ever did come together the Scripture merely showeth the absolute fact that up to this time they had not done so.\par
\switchcolumn*

\LA
\begin{Resp}\Rsign Descéndit Jesus cum eis, et venit Názareth: \Star{} Et erat súbditus illis.\end{Resp}
\begin{Resp}\Vsign Proficiébat sapiéntia, et ætáte, et grátia apud Deum et hómines.\end{Resp}
\begin{Resp}\Rsign Et erat súbditus illis.\end{Resp}
\begin{Resp}\Vsign Glória Patri, et Fílio, \Star{} et Spirítui Sancto.\end{Resp}
\begin{Resp}\Rsign Et erat súbditus illis.\end{Resp}
\switchcolumn
\EN
\begin{Resp}\Rsign Jesus went down with them, and came to Nazareth \Star{} And was subject unto them.\end{Resp}
\begin{Resp}\Vsign He increased in wisdom and stature, and in favour with God and man.\end{Resp}
\begin{Resp}\Rsign And was subject unto them.\end{Resp}
\begin{Resp}\Vsign Glory be to the Father, and to the Son, \Star{} and to the Holy Ghost.\end{Resp}
\begin{Resp}\Rsign And was subject unto them.\end{Resp}
\switchcolumn*

\LA
\LessonHead{Lectio IX}
\switchcolumn
\EN
\LessonHead{Lesson IX}
\switchcolumn*

\LA
\noindent Joseph autem vir ejus, cum esset justus, et nollet eam tradúcere, vóluit occúlte dimíttere eam. Si quis fornicáriæ conjúngitur, unum corpus effícitur, et in lege præcéptum est, non solum reos, sed et cónscios críminum obnóxios esse peccáti: quómodo Joseph, cum crimen celáret uxóris, justus scríbitur? Sed hoc testimónium Maríæ est, quod Joseph sciens illíus castitátem, et admírans quod evénerat, celat siléntio, cujus mystérium nesciébat.\par
\switchcolumn
\EN
\noindent Then Joseph her husband, being a just man, and not willing to make her a public example, was minded to put her away privily. If any man be joined to a fornicatress they become one body; and according to the law they that are privy to a crime are thereby guilty. How then can it be that Joseph is described as a just man, at the very time he was compounding the criminality of his espoused? It must have been that he knew her to be pure, and yet understood not the mystery of her pregnancy, but, while he wondered at that which had happened, was willing to hold his peace.\par
\switchcolumn*

\LA
\TeDeum{Te Deum}
\switchcolumn
\EN
\TeDeum{Te Deum}
\switchcolumn*

\end{paracol}

\par\needspace{10\baselineskip}\addvspace{1.2em}{\centering\small\textsc{Die 21 Martii} · \textsc{21 March}\par}
\Office{S. Benedicti Abbatis}{St. Benedict, Abbot}{Duplex majus}
\addcontentsline{toc}{subsection}{Die 21 Martii · S. Benedicti Abbatis\enspace\textit{St. Benedict, Abbot}}
\begin{paracol}{2}
\LA
\RefLine{Lectio IX: de Scriptura occurrente.}
\switchcolumn
\EN
\RefLine{Lesson IX: from the Scripture of the day.}
\switchcolumn*

\LA
\LessonHead{Lectio I}
\switchcolumn
\EN
\LessonHead{Lesson I}
\switchcolumn*

\LA
\LessonTitle{De libro Ecclesiastici}
\switchcolumn
\EN
\LessonTitle{Lesson from the book of Ecclesiasticus}
\switchcolumn*

\LA
\Cite{Sir 44:1-5}
\switchcolumn
\EN
\Cite{Sir 44:1-5}
\switchcolumn*

\LA
\noindent Laudémus viros gloriósos, et paréntes nostros in generatióne sua. Multam glóriam fecit Dóminus: magnificéntia sua a sǽculo. Dominántes in potestátibus suis, hómines magni virtúte et prudéntia sua prǽditi, nuntiántes in prophétis dignitátem prophetárum: et imperántes in præsénti pópulo, et virtúte prudéntiæ pópulis sanctíssima verba: in perítia sua requiréntes modos músicos, et narrántes cármina Scripturárum:\par
\switchcolumn
\EN
\noindent Let us now praise men of renown, and our fathers in their generation. The Lord hath wrought great glory through his magnificence from the beginning. Such as have borne rule in their dominions, men of great power, and endued with their wisdom, shewing forth in the prophets the dignity of prophets, And ruling over the present people, and by the strength of wisdom instructing the people in most holy words. Such as by their skill sought out musical tunes, and published canticles of the scriptures.\par
\switchcolumn*

\LA
\begin{Resp}\Rsign Euge, serve bone et fidélis, quia in pauca fuísti fidélis, supra multa te constítuam: \Star{} Intra in gáudium Dómini tui.\end{Resp}
\begin{Resp}\Vsign Dómine, quinque talénta tradidísti mihi, ecce ália quinque superlucrátus sum.\end{Resp}
\begin{Resp}\Rsign Intra in gáudium Dómini tui.\end{Resp}
\switchcolumn
\EN
\begin{Resp}\Rsign Well done, thou good and faithful servant, thou hast been faithful over a few things. I will make thee ruler over many things: \Star{} Enter thou into the joy of thy Lord.\end{Resp}
\begin{Resp}\Vsign Lord, thou deliveredst unto me five talents; behold, I have gained beside them five talents more.\end{Resp}
\begin{Resp}\Rsign Enter thou into the joy of thy Lord.\end{Resp}
\switchcolumn*

\LA
\LessonHead{Lectio II}
\switchcolumn
\EN
\LessonHead{Lesson II}
\switchcolumn*

\LA
\Cite{Sir 44:6-9}
\switchcolumn
\EN
\Cite{Sir 44:6-9}
\switchcolumn*

\LA
\noindent Hómines dívites in virtúte, pulchritúdinis stúdium habéntes, pacificántes in dómibus suis. Omnes isti in generatiónibus gentis suæ glóriam adépti sunt, et in diébus suis habéntur in láudibus. Qui de illis nati sunt reliquérunt nomen narrándi laudes eórum. Et sunt quorum non est memória: periérunt quasi qui non fúerint: et nati sunt quasi non nati, et fílii ipsórum cum ipsis.\par
\switchcolumn
\EN
\noindent Rich men in virtue, studying beautifulness: living at peace in their houses. All these have gained glory in their generations, and were praised in their days. They that were born of them have left a name behind them, that their praises might be related: And there are some, of whom there is no memorial: who are perished, as if they had never been: and are become as if they had never been born, and their children with them.\par
\switchcolumn*

\LA
\begin{Resp}\Rsign Justus germinábit sicut lílium: \Star{} Et florébit in ætérnum ante Dóminum.\end{Resp}
\begin{Resp}\Vsign Plantátus in domo Dómini, in átriis domus Dei nostri.\end{Resp}
\begin{Resp}\Rsign Et florébit in ætérnum ante Dóminum.\end{Resp}
\switchcolumn
\EN
\begin{Resp}\Rsign The righteous shall grow as the lily; \Star{} Yea, he shall flourish in the presence of the Lord for ever.\end{Resp}
\begin{Resp}\Vsign Those that be planted in the house of the Lord, shall flourish in the courts of the house of our God.\end{Resp}
\begin{Resp}\Rsign Yea, he shall flourish in the presence of the Lord for ever.\end{Resp}
\switchcolumn*

\LA
\LessonHead{Lectio III}
\switchcolumn
\EN
\LessonHead{Lesson III}
\switchcolumn*

\LA
\Cite{Sir 44:10-15}
\switchcolumn
\EN
\Cite{Sir 44:10-15}
\switchcolumn*

\LA
\noindent Sed illi viri misericórdiæ sunt, quorum pietátes non defuérunt. Cum sémine eórum pérmanent bona: hæréditas sancta nepótes eórum, et in testaméntis stetit semen eórum: et fílii eórum propter illos usque in ætérnum manent: semen eórum et glória eórum non derelinquétur. Córpora ipsórum in pace sepúlta sunt, et nomen eórum vivit in generatiónem et generatiónem. Sapiéntiam ipsórum narrent pópuli, et laudem eórum núntiet ecclésia.\par
\switchcolumn
\EN
\noindent But these were men of mercy, whose godly deeds have not failed: Good things continue with their seed, Their posterity are a holy inheritance, and their seed hath stood in the covenants. And their children for their sakes remain for ever: their seed and their glory shall not be forsaken. Their bodies are buried in peace, and their name liveth unto generation and generation. Let the people shew forth their wisdom, and the church declare their praise.\par
\switchcolumn*

\LA
\begin{Resp}\Rsign Iste cognóvit justítiam, et vidit mirabília magna, et exorávit Altíssimum: \Star{} Et invéntus est in número Sanctórum.\end{Resp}
\begin{Resp}\Vsign Iste est qui contémpsit vitam mundi, et pervénit ad cæléstia regna.\end{Resp}
\begin{Resp}\Rsign Et invéntus est in número Sanctórum.\end{Resp}
\begin{Resp}\Vsign Glória Patri, et Fílio, \Star{} et Spirítui Sancto.\end{Resp}
\begin{Resp}\Rsign Et invéntus est in número Sanctórum.\end{Resp}
\switchcolumn
\EN
\begin{Resp}\Rsign This is he which knew righteousness, and saw great wonders, and made his prayer unto the Most High; \Star{} And he is numbered among the Saints.\end{Resp}
\begin{Resp}\Vsign This is he which loved not his life in this world, and is come unto an everlasting kingdom.\end{Resp}
\begin{Resp}\Rsign And he is numbered among the Saints.\end{Resp}
\begin{Resp}\Vsign Glory be to the Father, and to the Son, \Star{} and to the Holy Ghost.\end{Resp}
\begin{Resp}\Rsign And he is numbered among the Saints.\end{Resp}
\switchcolumn*

\LA
\LessonHead{Lectio IV}
\switchcolumn
\EN
\LessonHead{Lesson IV}
\switchcolumn*

\LA
\noindent Benedíctus, Núrsiæ nóbili génere ortus, Romæ liberálibus disciplínis erudítus, ut totum se Jesu Christo daret, ad eum locum, qui Sublácus dícitur, in altíssimam spelúncam penetrávit; in qua sic per triénnium delítuit, ut unus id sciret Románus mónachus, quo ad vitæ necessitátem minístro utebátur. Dum ígitur ei quodam die ardéntes ad libídinem faces a diábolo subjiceréntur, se in vépribus támdiu volutávit, dum, laceráto córpore, voluptátis sensus dolóre opprimerétur. Sed jam erumpénte ex illis látebris fama ejus sanctitátis, quidam mónachi se illi instituéndos tradidérunt: quorum vivéndi licéntia cum ejus objurgatiónes ferre non posset, venénum in potióne ei dare constítuunt. Verum, póculum ei præbéntibus, crucis signo vas confrégit, ac relícto monastério in solitúdinem se recépit.\par
\switchcolumn
\EN
\noindent Benedict was born of a noble family at Norcia, (about the year of our Lord 480,) and studied letters at Rome. Desiring to give himself altogether to Christ Jesus, he betook himself to a very deep cave at the place now called Subiaco. In this place he lay hid for three years, unknown to all except the monk Romanus, by means of whom he received the necessaries of life. While he was in the cave at Subiaco, the devil one day assailed him with an extraordinary storm of impure temptation, and to get it under, he rolled himself in brambles till his whole body was lacerated, and the sting of pain drove out the sallies of lust. At last the fame of his holiness spread itself abroad from the desert, and some monks came to him for guidance, but the looseness of their lives was such that they could not bear his exhortations, and they plotted together to poison him in his drink. When they gave him the cup, he made the sign of the Cross over it, whereupon it immediately broke, and Benedict left that monastery, and retired to a desert place alone.\par
\switchcolumn*

\LA
\begin{Resp}\Rsign Honéstum fecit illum Dóminus, et custodívit eum ab inimícis, et a seductóribus tutávit illum: \Star{} Et dedit illi claritátem ætérnam.\end{Resp}
\begin{Resp}\Vsign Justum dedúxit Dóminus per vias rectas, et osténdit illi regnum Dei.\end{Resp}
\begin{Resp}\Rsign Et dedit illi claritátem ætérnam.\end{Resp}
\switchcolumn
\EN
\begin{Resp}\Rsign The Lord made him honourable, and defended him from his enemies, and kept him safe from those that lay in wait for him; \Star{} And gave him perpetual glory.\end{Resp}
\begin{Resp}\Vsign He went down with him into the pit, and left him not in bonds.\end{Resp}
\begin{Resp}\Rsign And gave him perpetual glory.\end{Resp}
\switchcolumn*

\LA
\LessonHead{Lectio V}
\switchcolumn
\EN
\LessonHead{Lesson V}
\switchcolumn*

\LA
\noindent Sed cum multi ad eum quotídie discípuli convenírent, duódecim monastéria ædificávit, éaque sanctíssimis légibus communívit. Póstea Cassínum migrávit, ubi simulácrum Apóllinis, qui adhuc ibi colebátur, commínuit, aram evértit et lucos succéndit; ibíque sancti Martíni sacéllum et sancti Joánnis ædículam exstrúxit, oppidános autem et íncolas christiánis præcéptis ímbuit. Quare augebátur in dies magis divína grátia Benedíctus, ut étiam prophético spíritu ventúra prædíceret. Quod ubi accépit Tótila Gothórum rex, exploratúrus, an res ita esset, spathárium suum régio ornátu et comitátu præmíttit, qui se regem simuláret. Quem ut ille vidit, Depóne, inquit, fili, depóne quod geris; nam tuum non est. Tótilæ vero prædíxit advéntum ejus in Urbem, maris transmissiónem, et post novem annos mortem.\par
\switchcolumn
\EN
\noindent Neverthless his disciples followed him daily, and for them he built twelve monasteries, and set holy laws to govern them. Afterwards he went to Cassino, and brake the image of Apollo which was still worshipped there, overturned the altar, and burnt the groves. There, (in the year 529,) he built the Church of St Martin and the little chapel of St. John; and instilled Christianity into the townspeople and inhabitants. He grew in the grace of God day by day, so that being endowed with the spirit of prophecy he foretold things to come. When Totila, King of the Goths, heard of it, and would see whether it really were so, he sent his Spatharius before him, with the kingly ensigns and attendance, and feigning himself to be Totila. But as soon as Benedict saw him he said: My son, put off that which thou wearest, for it is not thine. To Totila himself he foretold that he would go to Rome, would cross the sea, and would die after nine years.\par
\switchcolumn*

\LA
\begin{Resp}\Rsign Amávit eum Dóminus, et ornávit eum: stolam glóriæ índuit eum, \Star{} Et ad portas paradísi coronávit eum.\end{Resp}
\begin{Resp}\Vsign Índuit eum Dóminus lorícam fídei, et ornávit eum.\end{Resp}
\begin{Resp}\Rsign Et ad portas paradísi coronávit eum.\end{Resp}
\switchcolumn
\EN
\begin{Resp}\Rsign The Lord loved him and beautified him; He clothed him with a robe of glory, \Star{} And crowned him at the gates of Paradise.\end{Resp}
\begin{Resp}\Vsign The Lord hath put on him the breast-plate of faith, and hath adorned him.\end{Resp}
\begin{Resp}\Rsign And crowned him at the gates of Paradise.\end{Resp}
\switchcolumn*

\LA
\LessonHead{Lectio VI}
\switchcolumn
\EN
\LessonHead{Lesson VI}
\switchcolumn*

\LA
\noindent Qui áliquot ménsibus ántequam e vita migráret, præmónuit discípulos, quo die esset moritúrus; ac sepúlcrum, in quo suum corpus condi vellet, sex diébus ántequam eo inferrétur, aperíri jussit: sextóque die deférri vóluit in ecclésiam; ubi, sumpta Eucharístia, sublátis in cælum óculis orans, inter manus discipulórum efflávit ánimam: quam duo mónachi eúntem in cælum vidérunt, pállio ornátum pretiosíssimo, circum eam fulgéntibus lampádibus, et claríssima et gravíssima spécie virum, stantem supra caput ipsíus, dicéntem audiérunt: Hæc est via, qua diléctus Dómini Benedíctus in cælum ascéndit.\par
\switchcolumn
\EN
\noindent Some months before he departed this life, Benedict forewarned his disciples on what day he was to die; and he ordered his grave to be opened six days before he was carried to it. On the sixth day, (being the 21st of March, in the year 543,) he would be carried into the Church, where he received the Eucharist, and then, in the arms of his disciples, with his eyes lifted up to heaven, and wrapt in prayer, he gave up the ghost. Two monks saw his soul rising to heaven, clothed in a most precious garment, and surrounded with lights, and One of a most glorious and awful aspect standing above, Whom they heard saying This is the way whereby Benedict, the beloved of the Lord, goeth up to heaven.\par
\switchcolumn*

\LA
\begin{Resp}\Rsign Iste homo perfécit ómnia quæ locútus est ei Deus, et dixit ad eum: Ingrédere in réquiem meam: \Star{} Quia te vidi justum coram me ex ómnibus géntibus.\end{Resp}
\begin{Resp}\Vsign Iste est, qui contémpsit vitam mundi, et pervénit ad cæléstia regna.\end{Resp}
\begin{Resp}\Rsign Quia te vidi justum coram me ex ómnibus géntibus.\end{Resp}
\begin{Resp}\Vsign Glória Patri, et Fílio, \Star{} et Spirítui Sancto.\end{Resp}
\begin{Resp}\Rsign Quia te vidi justum coram me ex ómnibus géntibus.\end{Resp}
\switchcolumn
\EN
\begin{Resp}\Rsign This is he which did according unto all that God commanded him; and God said unto him: Enter thou into My rest, \Star{} For thee have I seen righteous before Me among all people.\end{Resp}
\begin{Resp}\Vsign This is he which loved not his life in this world, and is come unto an everlasting kingdom.\end{Resp}
\begin{Resp}\Rsign For thee have I seen righteous before Me among all people.\end{Resp}
\begin{Resp}\Vsign Glory be to the Father, and to the Son, \Star{} and to the Holy Ghost.\end{Resp}
\begin{Resp}\Rsign For thee have I seen righteous before Me among all people.\end{Resp}
\switchcolumn*

\LA
\LessonHead{Lectio VII}
\switchcolumn
\EN
\LessonHead{Lesson VII}
\switchcolumn*

\LA
\LessonTitle{Léctio sancti Evangélii secúndum Matthǽum.}
\switchcolumn
\EN
\LessonTitle{From the Holy Gospel according to Matthew}
\switchcolumn*

\LA
\Cite{Matt 19:27-29}
\switchcolumn
\EN
\Cite{Matt 19:27-29}
\switchcolumn*

\LA
\noindent In illo témpore: Dixit Petrus ad Jesum: Ecce nos relíquimus ómnia, et secúti sumus te: quid ergo erit nobis? Et réliqua.\par
\switchcolumn
\EN
\noindent At that time, Peter said unto Jesus: Behold, we have forsaken all, and followed thee what shall we have therefore? And so on.\par
\switchcolumn*

\LA
\LessonTitle{Homilía sancti Hierónymi Presbýteri.}
\switchcolumn
\EN
\LessonTitle{Homily by St. Jerome, Priest (at Bethlehem.)}
\switchcolumn*

\LA
\Source{Lib. 3. in Matth. cap. 19.}
\switchcolumn
\EN
\Source{Bk. iii. on Matth xix.}
\switchcolumn*

\LA
\noindent Grandis fidúcia! Petrus piscátor erat, dives non fúerat, cibos manu et arte quærébat: et tamen lóquitur confidénter: Relíquimus ómnia. Et quia non súfficit tantum relínquere, jungit quod perféctum est: Et secúti sumus te. Fécimus quod jussísti: quid ígitur nobis dabis prǽmii? Jesus autem dixit illis: Amen dico vobis, quod vos qui secúti estis me, in regeneratióne, cum séderit Fílius hóminis in sede majestátis suæ, sedébitis et vos super sedes duódecim, judicántes duódecim tribus Ísraël. Non dixit: Qui reliquístis ómnia; hoc enim et Crates fecit philósophus, et multi álii divítias contempsérunt: sed, Qui secúti estis me: quod próprie Apostolórum est, atque credéntium.\par
\switchcolumn
\EN
\noindent Peter was a fisherman, he was not rich, he earned his bread by his hand and skill, and nevertheless he is thus bold, and saith confidently: We have forsaken all. And because it sufficeth not to forsake only, he addeth that which to do is to be perfect: and followed thee. We have done that which Thou hast commanded us, what reward therefore wilt Thou give us? And Jesus said unto them Amen I say unto you, that ye which have followed Me, in the regeneration, when the Son of Man shall sit in the throne of His glory, ye also shall sit upon twelve thrones, judging the twelve tribes of Israel. He said not, Ye which have forsaken all, for this did even Crates the philosopher, and they which have set nothing by riches are many, but, Ye which have followed Me. This did the Apostles, and this do believers do.\par
\switchcolumn*

\LA
\begin{Resp}\Rsign Isti sunt qui vivéntes in carne, plantavérunt Ecclésiam sánguine suo: \Star{} Cálicem Dómini bibérunt, et amíci Dei facti sunt.\end{Resp}
\begin{Resp}\Vsign In omnem terram exívit sonus eórum, et in fines orbis terræ verba eórum.\end{Resp}
\begin{Resp}\Rsign Cálicem Dómini bibérunt, et amíci Dei facti sunt.\end{Resp}
\switchcolumn
\EN
\begin{Resp}\Rsign These are they who while yet they lived in the flesh, planted the Church in their own blood; \Star{} They drank of the Lord's cup, and became the friends of God.\end{Resp}
\begin{Resp}\Vsign Their sound is gone out through all the earth, and their words to the ends of the world.\end{Resp}
\begin{Resp}\Rsign They drank of the Lord's cup, and became the friends of God.\end{Resp}
\switchcolumn*

\LA
\LessonHead{Lectio VIII}
\switchcolumn
\EN
\LessonHead{Lesson VIII}
\switchcolumn*

\LA
\noindent In regeneratióne, cum séderit Fílius hóminis in sede majestátis suæ (quando et mórtui de corruptióne resúrgent incorrúpti) sedébitis et vos in sóliis judicántium, condemnántes duódecim tribus Ísraël: quia vobis credéntibus, illi crédere noluérunt. Et omnis qui relíquerit domum, vel fratres, aut soróres, aut patrem, aut matrem, aut uxórem, aut fílios, aut agros propter nomen meum, céntuplum accípiet, et vitam ætérnam possidébit. Locus iste cum illa senténtia cóngruit, in qua Salvátor lóquitur: Non veni pacem míttere, sed gládium. Veni enim separáre hóminem a patre suo, et matrem a fília, et nurum a socru: et inimíci hóminis doméstici ejus. Qui ergo propter fidem Christi, et prædicatiónem Evangélii, omnes afféctus contémpserint, atque divítias et sǽculi voluptátes: isti céntuplum recípient, et vitam ætérnam possidébunt.\par
\switchcolumn
\EN
\noindent In the regeneration, when the Son of Man shall sit in the throne of His glory, and when the dead shall rise again from corruption incorruptible, (i Cor. xv. 53,) ye also shall sit upon twelve thrones of judgment, condemning the twelve tribes of Israel, because, when ye believed in Me, they would not. (John iii. 18.) And every one that hath forsaken houses, or brethren, or sisters, or father, or mother, or wife, or children, or lands, for My Name's sake, shall receive an hundredfold, and shall inherit everlasting life. This place agreeth well with that other where the Saviour saith I came not to send peace, but a sword. For I am come to set a man at variance against his father, and the daughter against her mother, and the daughter-in-law against her mother-in-law; and a man's foes shall be they of his own household. (Matth. x. 34.) Every one, therefore, that hath set no store by affection, and riches, and the pleasures of the world, for Christ's faith's sake, and the preaching of the Gospel, shall receive an hundred-fold, and shall inherit everlasting life.\par
\switchcolumn*

\LA
\begin{Resp}\Rsign Isti sunt viri sancti, quos elégit Dóminus in caritáte non ficta, et dedit illis glóriam sempitérnam: \Star{} Quorum doctrína fulget Ecclésia, ut sole luna.\end{Resp}
\begin{Resp}\Vsign Sancti per fidem vicérunt regna: operáti sunt justítiam.\end{Resp}
\begin{Resp}\Rsign Quorum doctrína fulget Ecclésia, ut sole luna.\end{Resp}
\begin{Resp}\Vsign Glória Patri, et Fílio, \Star{} et Spirítui Sancto.\end{Resp}
\begin{Resp}\Rsign Quorum doctrína fulget Ecclésia, ut sole luna.\end{Resp}
\switchcolumn
\EN
\begin{Resp}\Rsign These men are saints, whom the Lord hath chosen in love unfeigned, and hath given them glory everlasting. These are they \Star{} By the light of whose teaching the Church is glorified, even as the moon is glorified by the light of the sun.\end{Resp}
\begin{Resp}\Vsign The saints through faith subdued kingdoms, wrought righteousness.\end{Resp}
\begin{Resp}\Rsign By the light of whose teaching the Church is glorified, even as the moon is glorified by the light of the sun.\end{Resp}
\begin{Resp}\Vsign Glory be to the Father, and to the Son, \Star{} and to the Holy Ghost.\end{Resp}
\begin{Resp}\Rsign By the light of whose teaching the Church is glorified, even as the moon is glorified by the light of the sun.\end{Resp}
\switchcolumn*

\end{paracol}

\par\needspace{10\baselineskip}\addvspace{1.2em}{\centering\small\textsc{Die 24 Martii} · \textsc{24 March}\par}
\Office{S. Gabrielis Archangeli}{St. Gabriel the Archangel}{Duplex majus}
\addcontentsline{toc}{subsection}{Die 24 Martii · S. Gabrielis Archangeli\enspace\textit{St. Gabriel the Archangel}}
\begin{paracol}{2}
\LA
\RefLine{Lectio IX: de Scriptura occurrente.}
\switchcolumn
\EN
\RefLine{Lesson IX: from the Scripture of the day.}
\switchcolumn*

\LA
\LessonHead{Lectio I}
\switchcolumn
\EN
\LessonHead{Lesson I}
\switchcolumn*

\LA
\LessonTitle{De Daniéle Prophéta}
\switchcolumn
\EN
\LessonTitle{Lesson from the book of Daniel}
\switchcolumn*

\LA
\Cite{Dan 9:20-23}
\switchcolumn
\EN
\Cite{Dan 9:20-23}
\switchcolumn*

\LA
\noindent Ego Dániel, cum adhuc lóquerer, et orárem, et confitérer peccáta mea, et peccáta pópuli mei Israël, et prostérnerem preces meas in conspéctu Dei mei, pro monte sancto Dei mei: Adhuc me loquénte in oratióne, ecce vir Gábriel, quem víderam in visióne a princípio, cito volans tétigit me in témpore sacrifícii vespertíni, et dócuit me, et locútus est mihi, dixítque: Dániel, nunc egréssus sum ut docérem te, et intellígeres: ab exórdio precum tuárum egréssus est sermo; ego autem veni ut indicárem tibi, quia vir desideriórum es; tu ergo animadvérte sermónem, et intéllige visiónem.\par
\switchcolumn
\EN
\noindent Now while I was yet speaking, and praying, and confessing my sins, and the sins of my people of Israel, and presenting my supplications in the sight of my God, for the holy mountain of my God: As I was yet speaking in prayer, behold the man Gabriel, whom I had seen in the vision at the beginning, flying swiftly touched me at the time of the evening sacrifice. And he instructed me, and spoke to me, and said: O Daniel, I am now come forth to teach thee, and that thou mightest understand. From the beginning of thy prayers the word came forth: and I am come to show it to thee, because thou art a man of desires: therefore do thou mark the word, and understand the vision.\par
\switchcolumn*

\LA
\begin{Resp}\Rsign Cum oráret Dániel et confiterétur peccáta sua, et peccáta pópuli sui: \Star{} Ecce Archángelus Gábriel cito volans tétigit eum in témpore sacrifícii vespertíni.\end{Resp}
\begin{Resp}\Vsign Cumque prostérneret preces suas in conspéctu Dei sui,\end{Resp}
\begin{Resp}\Rsign Ecce Archángelus Gábriel cito volans tétigit eum in témpore sacrifícii vespertíni.\end{Resp}
\switchcolumn
\EN
\begin{Resp}\Rsign Whiles Daniel was praying, and confessing his sins, and the sins of his people \Star{} Behold, the Archangel Gabriel, being caused to fly swiftly, touched him about the time of the evening oblation.\end{Resp}
\begin{Resp}\Vsign Whiles he was presenting his supplication before his God\end{Resp}
\begin{Resp}\Rsign Behold, the Archangel Gabriel, being caused to fly swiftly, touched him about the time of the evening oblation.\end{Resp}
\switchcolumn*

\LA
\LessonHead{Lectio II}
\switchcolumn
\EN
\LessonHead{Lesson II}
\switchcolumn*

\LA
\Cite{Dan 9:24-25}
\switchcolumn
\EN
\Cite{Dan 9:24-25}
\switchcolumn*

\LA
\noindent Septuagínta hebdómades abbreviátæ sunt super pópulum tuum et super urbem sanctam tuam, ut consummétur prævaricátio, et finem accípiat peccátum, et deleátur iníquitas, et adducátur justítia sempitérna, et impleátur vísio et prophetía, et ungátur Sanctus sanctórum. Scito ergo, et animadvérte: Ab éxitu sermónis, ut íterum ædificétur Jerúsalem, usque ad Christum ducem, hebdómades septem et hebdómades sexagínta duæ erunt: et rursum ædificábitur platéa, et muri in angústia témporum.\par
\switchcolumn
\EN
\noindent Seventy weeks are shortened upon thy people, and upon thy holy city, that transgression may be finished, and sin may have an end, and iniquity may be abolished; and everlasting justice may be brought; and vision and prophecy may be fulfilled; and the saint of saints may be anointed. Know thou therefore, and take notice: that from the going forth of the word, to build up Jerusalem again, unto Christ the prince, there shall be seven weeks, and sixty-two weeks: and the street shall be built again, and the walls in straitness of times.\par
\switchcolumn*

\LA
\begin{Resp}\Rsign Locútus est Gábriel Daniéli et dixit: Ab exórdio precum tuárum egréssus est sermo. \Star{} Ego autem veni ut indicárem tibi, quia vir desideriórum es.\end{Resp}
\begin{Resp}\Vsign Tu autem animadvérte sermónem, et intéllige visiónem.\end{Resp}
\begin{Resp}\Rsign Ego autem veni ut indicárem tibi, quia vir desideriórum es.\end{Resp}
\switchcolumn
\EN
\begin{Resp}\Rsign And Gabriel talked with Daniel, and said: At the beginning of thy supplications the commandment came forth \Star{} And I am come to show thee, for thou art greatly beloved.\end{Resp}
\begin{Resp}\Vsign Therefore consider the matter and understand the vision.\end{Resp}
\begin{Resp}\Rsign And I am come to show thee, for thou art greatly beloved.\end{Resp}
\switchcolumn*

\LA
\LessonHead{Lectio III}
\switchcolumn
\EN
\LessonHead{Lesson III}
\switchcolumn*

\LA
\Cite{Dan 9:26-27}
\switchcolumn
\EN
\Cite{Dan 9:26-27}
\switchcolumn*

\LA
\noindent Et post hebdómades sexagínta duas occidétur Christus, et non erit ejus pópulus qui eum negatúrus est; et civitátem et sanctuárium dissipábit pópulus cum duce ventúro: et finis ejus vástitas, et post finem belli statúta desolátio. Confirmábit autem pactum multis hebdómada una: et in dimídio hebdómadis defíciet hóstia et sacrifícium, et erit in templo abominátio desolatiónis: et usque ad consummatiónem et finem perseverábit desolátio.\par
\switchcolumn
\EN
\noindent And after sixty-two weeks Christ shall be slain: and the people that shall deny him shall not be his. And a people with their leader that shall come, shall destroy the city and the sanctuary: and the end thereof shall be waste, and after the end of the war the appointed desolation. And he shall confirm the covenant with many, in one week: and in the half of the week the victim and the sacrifice shall fall: and there shall be in the temple the abomination of desolation: and the desolation shall continue even to the consummation, and to the end.\par
\switchcolumn*

\LA
\begin{Resp}\Rsign Ecce vir Gábriel quem víderam in visióne a princípio cito volans tétigit me in témpore sacrifícii vespertíni et dócuit me et locútus est mihi dixítque \Star{} Dániel nunc egréssus sum ut docérem te et intellígeres\end{Resp}
\begin{Resp}\Vsign Gábriel fac me intellígere istam visiónem: et venit et stetit juxta ubi ego stabam, et ait ad me:\end{Resp}
\begin{Resp}\Rsign Dániel nunc egréssus sum ut docérem te et intellígeres\end{Resp}
\begin{Resp}\Vsign Glória Patri, et Fílio, \Star{} et Spirítui Sancto.\end{Resp}
\begin{Resp}\Rsign Dániel nunc egréssus sum ut docérem te et intellígeres\end{Resp}
\switchcolumn
\EN
\begin{Resp}\Rsign Behold, the man Gabriel, whom I had seen in the vision at the beginning, being caused to fly swiftly, touched me about the time of the evening oblation; and he informed me, and talked with me, and said: \Star{} O Daniel, I am now come forth to give thee skill and understanding.\end{Resp}
\begin{Resp}\Vsign O Gabriel, make me to understand the vision. So he came near where I stood. And he said unto me\end{Resp}
\begin{Resp}\Rsign O Daniel, I am now come forth to give thee skill and understanding.\end{Resp}
\begin{Resp}\Vsign Glory be to the Father, and to the Son, \Star{} and to the Holy Ghost.\end{Resp}
\begin{Resp}\Rsign O Daniel, I am now come forth to give thee skill and understanding.\end{Resp}
\switchcolumn*

\LA
\LessonHead{Lectio IV}
\switchcolumn
\EN
\LessonHead{Lesson IV}
\switchcolumn*

\LA
\LessonTitle{Sermo sancti Bedæ venerábilis Presbýteri}
\switchcolumn
\EN
\LessonTitle{A Sermon by St. Venerable Bede, the Priest (at Jarrow.)}
\switchcolumn*

\LA
\Source{Expositio in Luc. 1, 11-20.}
\switchcolumn
\EN
\Source{Exposition on Luke 1:11-20}
\switchcolumn*

\LA
\noindent Appáruit Zacharíæ Angelus, stans a dextris altáris incénsi. Bene Angelus et in templo, et juxta altáre, et dextris appáret; quia vidélicet et veri sacerdótis advéntum, et mystérium sacrifícii universális, et cæléstis doni gáudium prǽdicat. Nam, sicut per sinístram præséntia, sic per déxteram sæpe bona prænuntiántur ætérna. Juxta quod in sapiéntiæ laude cánitur: Longitúdo diérum in déxtera ejus; in sinístra illíus divítiæ et glória. Treméntem Zacharíam confórtet Angelus, quia, sicut humánæ fragilitátis est spiritális creatúræ visióne turbári, ita et angélicas benignitátis est pavéntes de aspéctu suo mortáles mox blandiéndo solári. At contra dæmónicæ est ferocitátis, quos sui præséntia térritos sénserit, amplióri semper horróre concútere, qua nulla mélius ratióne quam fide superátur intrépida.\par
\switchcolumn
\EN
\noindent The Angel appeared to Zacharias in the sanctuary of the temple at the right side of the altar of incense. Which same was fitting, thus: he appeared in the sanctuary because he came to proclaim sacrifice; and at the right side thereof, to indicate how joyous was the honour about to be bestowed on mankind by the heavenly gift. The right side is the side of honour, and therefore words indicating a position at the right hand are often used to signify an eternal good, and by the same token, to be at the left side doth sometimes signify only present good. As for example where the Book of Proverbs singeth thus in praise of wisdom: Length of days in in her right hand, and in her left hand riches and honour. First of all the Angel comforteth the trembling Zacharias. Fear not, saith he. For just as it is natural for human frailty to fear spiritual manifestations, so it is natural for Angels to comfort with good words the mortals that be in this wise fearful. Contrariwise, when the devil perceiveth that his audacious manifestations do frighten, he proceedeth to frighten as much as he can, and that with an increasing fearsomeness. There is no better way to overcome his workings than by a courageous faith.\par
\switchcolumn*

\LA
\begin{Resp}\Rsign Factum est, cum sacerdótio fungerétur Zacharías in órdine vicis sufflánte Deum, \Star{} Appáruit ei Angelus Gábriel, stans a dextris altáris incénsi.\end{Resp}
\begin{Resp}\Vsign Cumque, ingréssus in templum Dómini, incénsum póneret secúndum consuetúdinem sacerdótii sui.\end{Resp}
\begin{Resp}\Rsign Appáruit ei Angelus Gábriel, stans a dextris altáris incénsi.\end{Resp}
\switchcolumn
\EN
\begin{Resp}\Rsign It came to pass that while Zacharias executed the Priest's office before God, in the order of his course \Star{} There appeared unto him the Angel Gabriel, standing on the right side of the Altar of incense.\end{Resp}
\begin{Resp}\Vsign When he went into the temple of the Lord to burn incense, according to the custom of the Priest's office.\end{Resp}
\begin{Resp}\Rsign There appeared unto him the Angel Gabriel, standing on the right side of the Altar of incense.\end{Resp}
\switchcolumn*

\LA
\LessonHead{Lectio V}
\switchcolumn
\EN
\LessonHead{Lesson V}
\switchcolumn*

\LA
\noindent Angelus, deprecatiónem dicens exaudítam, partum contínuo promíttit uxóris. Non quod ille, qui pro pópulo oblatúrus intráverat, relíctis públicis votis, pro accipiéndis fíliis oráre potúerit, præsértim quia nemo orat quod se acceptúrum despérat. Adeo autem ille, jam suæ ætátis et infœcúndæ cónjugis memor, nasci sibi fílios desperábat, ut nec Angelo hoc promitténti créderet. Sed quod dicit: Exaudíta est deprecátio tua, pro pópuli redemptióne signíficat; Et uxor tua páriet tibi fílium, ejúsdem redemptiónis órdinem pandit, quod vidélicet natus Zacharíæ fílius Redemptóri illíus pópuli præconándo sit iter factúrus. Quia enim Zacharíam, pro plebe supplicántem, díxerat exaudítum, docet quo órdine plebs eádem salvári et pérfici débeat, ad prædicatiónem vidélicet Joánnis pæniténdo et credéndo in Christum.\par
\switchcolumn
\EN
\noindent Next, the Angel saith that the prayer of Zacharias was heard, and then straightway promiseth that the wife of Zacharias should bear a child. We are not to understand that he had been praying for the birth of a son whilst he was offering the sacrifice according to the liturgy of that time, for we are told that he had given up hope of a son, and no one prayeth for that which he hath no hope of obtaining. Yea, so hopeless was he of ever having children of his own, because Elizabeth was barren, and they were both now well stricken in years, that he did not even believe the Angel's promise. Therefore the words of the Angel: Thy prayer is heard: refer to the redemption of the people, for which Zacharias had prayed in the pleading of the sacrifice. And the words: Thy wife shall bear a son: do shew the manner of that redemption, for he addeth that the son of Zacharias shall go before the Redeemer as a herald, to make ready his way amongst the people. Thus, in this saying that the prayer of supplication offered by Zacharias was heard of God, the Angel sheweth in what manner the people can be brought to salvation and perfection; namely, by repentance at the preaching of John, whereby they are to be led to faith in Christ.\par
\switchcolumn*

\LA
\begin{Resp}\Rsign Descéndit Gábriel Angelus ad Zacharíam dicens: \Star{} Ne tímeas quóniam exaudíta est deprecátio tua, et uxor tua Elísabeth paret tibi fílium, et vocábis nomen ejus Joánnem.\end{Resp}
\begin{Resp}\Vsign Et Zacharías turbátus est videns, et timor írruit super eum: ait autem ad illum Angelus:\end{Resp}
\begin{Resp}\Rsign Ne tímeas quóniam exaudíta est deprecátio tua, et uxor tua Elísabeth paret tibi fílium, et vocábis nomen ejus Joánnem.\end{Resp}
\switchcolumn
\EN
\begin{Resp}\Rsign The Angel Gabriel came down unto Zacharias, and said unto him: \Star{} Fear not; for thy prayer is heard; and thy wife Elizabeth shall bear thee a son, and thou shalt call his name John.\end{Resp}
\begin{Resp}\Vsign And when Zacharias saw him, he was troubled, and fear fell upon him; but the Angel said unto him\end{Resp}
\begin{Resp}\Rsign Fear not; for thy prayer is heard; and thy wife Elizabeth shall bear thee a son, and thou shalt call his name John.\end{Resp}
\switchcolumn*

\LA
\LessonHead{Lectio VI}
\switchcolumn
\EN
\LessonHead{Lesson VI}
\switchcolumn*

\LA
\noindent Zacharías ob altitúdinem promissórum hǽsitans, signum quo crédere váleat, inquírit cui sola Angeli vísio vel allocútio pro signo suffícere debúerat. Unde méritam diffidéntiæ tacéndo pœnam luit, cui tacitúrnitas eádem et signum fídei quod quæsívit, et infidelitátis esset pœna quem méruit. Vult intélligi quod si homo tália promítteret, impúne signum flagitáre licéret, at cum ángelus promíttat, jam dubitári non déceat. Datque signum quod rogátur, ut qui discredéndo locútus est, jam tacéndo crédere discat. Ubi notándo quod Angelus se et ante Deum adstáre, et ad evangelizándum Zacharíæ missum esse testátur; quia et cum ad nos véniunt Angeli, sic extérius implent ministérium, ut tamen nunquam desint intérius per contemplatiónem. Et mittúntur ígitur et assístunt, quia, et si circumscríptus est angélicus spíritus, summus tamen spíritus ipse qui Deus est, circumscríptus non est. Angeli ítaque et missi ante ipsum sunt, quia, quólibet missi véniunt, intra ipsum currunt. -- Festum sancti Gabriélis Archángeli Benedíctus Papa décimus quintus ad univérsam ecclésiam exténdit.\par
\switchcolumn
\EN
\noindent But Zacharias hesitateth because of the sublime things which have been promised. Wherefore he asketh for a sign, that he may believe, albeit the coming of the Angel and his words of promise ought to have been a sufficient sign. Hence he was stricken dumb as a just penalty for his slowness of belief: to be dumb was both a sign to stir him up to the faith which he sought, and the penance which he deserved for his unbelief. We may thus understand that if a man of earth had promised such things, it would be lawful to seek for a sign, but that when an Angel is sent from heaven to give God's promise, there should have been no occasion for doubt. And yet the Angel giveth the desired sign, so that he who spake from disbelief may learn from silence to believe. Note that the Angel saith: I am Gabriel, that stand in the presence of God, and am sent to speak unto thee these glad tidings. Doubtless when Angels come to us they fulfill this active and outward ministry in such a way that they yet do always remain in God's presence by contemplation. Wherefore, they stand in his presence even though they be sent from him on a mission. An Angel is a created spirit, and therefore hath many limitations. But God hath no limitations, and is everywhere. Thus when he sendeth his Angels from his presence, they yet do stand therein, for whithersoever they go on a mission, they go in him. - This Feast of the Archangel Gabriel was extended to the universal Church by Pope Benedict XV.\par
\switchcolumn*

\LA
\begin{Resp}\Rsign Ego sum Gábriel Angelus, qui asto ante Dóminum; et missus sum loqui te, et hæc evangelizáre: \Star{} Ecce eris tacens, et non póteris loqui usque in diem quo hæc fiant.\end{Resp}
\begin{Resp}\Vsign Pro eo quod non credidísti verbis meis, quæ implebúntur témpore suo,\end{Resp}
\begin{Resp}\Rsign Ecce eris tacens, et non póteris loqui usque in diem quo hæc fiant.\end{Resp}
\begin{Resp}\Vsign Glória Patri, et Fílio, \Star{} et Spirítui Sancto.\end{Resp}
\begin{Resp}\Rsign Ecce eris tacens, et non póteris loqui usque in diem quo hæc fiant.\end{Resp}
\switchcolumn
\EN
\begin{Resp}\Rsign I am Gabriel, that stand in the presence of God; and am sent to speak unto thee, and to show thee these glad tidings. \Star{} Behold, thou shalt be dumb, and not able to speak, until the day that these things shall be performed.\end{Resp}
\begin{Resp}\Vsign Because thou believest not my words, which shall be fulfilled in their season.\end{Resp}
\begin{Resp}\Rsign Behold, thou shalt be dumb, and not able to speak, until the day that these things shall be performed.\end{Resp}
\begin{Resp}\Vsign Glory be to the Father, and to the Son, \Star{} and to the Holy Ghost.\end{Resp}
\begin{Resp}\Rsign Behold, thou shalt be dumb, and not able to speak, until the day that these things shall be performed.\end{Resp}
\switchcolumn*

\LA
\LessonHead{Lectio VII}
\switchcolumn
\EN
\LessonHead{Lesson VII}
\switchcolumn*

\LA
\LessonTitle{Léctio sancti Evangélii secúndum Lucam}
\switchcolumn
\EN
\LessonTitle{From the Holy Gospel according to Luke}
\switchcolumn*

\LA
\Cite{Luc 1:26-38}
\switchcolumn
\EN
\Cite{Luke 1:26-38}
\switchcolumn*

\LA
\noindent In illo témpore: Missus est Angelus Gábriel a Deo in civitátem Galilǽæ, cui nomen Názareth, ad Vírginem desponsátam viro, cui nomen erat Joseph de domo David, et nomen Vírginis María. Et réliqua.\par
\switchcolumn
\EN
\noindent At that time: The Angel Gabriel was sent from God unto a city of Galilee named Nazareth, to a virgin espoused to a man whose name was Joseph, of the house of David: and the Virgin's name was Mary. And so on.\par
\switchcolumn*

\LA
\LessonTitle{Homilía Sancti Bernárdi Abbátis}
\switchcolumn
\EN
\LessonTitle{A Homily by St. Bernard the Abbot}
\switchcolumn*

\LA
\Source{Homilia 1 supra Missus est n. 2}
\switchcolumn
\EN
\Source{Homily 1 on Missus est, no. 2}
\switchcolumn*

\LA
\noindent Non árbitror hunc ángelum de minóribus esse, qui quálibet ex causa crebra sóleant ad terras fungi legatióne. Quod ex ejus nómine palam intélligi dícitur, quod interpretátio Fortitúdo Dei dícitur; et quia non ab álio áliquo forte excellentióre se (ut ássolet) spíritu, sed ab ipso Deo mitti perhibétur. Propter hoc ergo pósitum est, A Deo. Vel ídeo dictum est, A Deo, ne cui, vel beatórum spirítuum suum Deus, ántequam Vírgini, revelásse putétur consílium, excépto dumtáxat Archángelo Gabriéle; qui útique tantæ inter suos inveníri potúerit excelléntiæ, ut tali et nómine dignus haberétur et núntio.\par
\switchcolumn
\EN
\noindent Consider that this Angel was not one of lesser rank, even though such are, on one account or another, often sent on embassies to this earth. That he was an Angel of greater rank is indicated by his name which signifieth: Strength of God: and by the fact that he was sent, not by some Angel perhaps more excellent than he (as is usual), but from God himself. Therefore for this reason it is said: From God. Or, it might be for another reason, namely, lest it should be thought that God had discourteously revealed his counsel to any of the blessed spirits, except only the Archangel Gabriel, before he did so to the Virgin. For Gabriel alone was found so eminent among his compeers as to be held worthy both of such a name and message.\par
\switchcolumn*

\LA
\begin{Resp}\Rsign Missus est Gábriel Angelus ad Maríam Vírginem desponsátam Joseph, núntians ei verbum: et expavéscit Virgo de lúmine. \Star{} Ne tímeas, María, invenísti enim grátiam apud Dóminum: ecce concípies, et páries, et vocábitur Altíssimi Fílius\end{Resp}
\begin{Resp}\Vsign Quæ cum audísset turbáta est in sermóne ejus, et cogitábat qualis esset ista salutátio; et ait Angelus ei.\end{Resp}
\begin{Resp}\Rsign Ne tímeas, María, invenísti enim grátiam apud Dóminum: ecce concípies, et páries, et vocábitur Altíssimi Fílius.\end{Resp}
\switchcolumn
\EN
\begin{Resp}\Rsign The Angel Gabriel was sent to Mary, the virgin espoused to Joseph, and the Virgin was afraid of the light. \Star{} Fear not, Mary, for thou hast found grace with the Lord. Behold, thou shalt conceive and bring forth a Son, and He shall be called the Son of the Highest.\end{Resp}
\begin{Resp}\Vsign And when she heard it, she was troubled at his saying, and cast in her mind what manner of salutation this should be. And the Angel said unto her:\end{Resp}
\begin{Resp}\Rsign Fear not, Mary, for thou hast found grace with the Lord. Behold, thou shalt conceive, and bring forth a Son, and He shall be called the Son of the Highest.\end{Resp}
\switchcolumn*

\LA
\LessonHead{Lectio VIII}
\switchcolumn
\EN
\LessonHead{Lesson VIII}
\switchcolumn*

\LA
\noindent Nec discórdat nomen a núntio. Dei quippe virtútem Christum quem mélius nuntiáre decébat quam hunc, quem símile nomen honórat? Nam quid est áliud fortitúdo, quam virtus? Non autem dedécens aut incóngruum videátur, Dóminum et núntium commúni censéri vocábulo; cum símilis in utróque appellatiónis, non sit tamen utriúsque símilis causa. Aliter quippe Christus fortitúdo vel virtus Dei dícitur, áliter Angelus; Angelus enim tantum nuncupatíve, Christus autem étiam substantíve.\par
\switchcolumn
\EN
\noindent Neither do his name and his message disagree. For whom did it more behove to announce Christ, who is the Power of God, than him who is honoured by a like name? For what else is power than strength? Neither doth it appear to be unbecoming or unseemly that the Lord and his messenger should be known by a like title. Christ is called the power or strength of God in a very different sense from that in which this appellation is given to the Angel. In the Angel it is but a name. In the case of Christ, it is also an eternal attribute.\par
\switchcolumn*

\LA
\begin{Resp}\Rsign Gaude, María Virgo cunctas hǽreses sola interemísti: \Star{} Quæ Gabriéli Archángeli dictis credidísti, dum Virgo Deum a hóminem genuísti, et post partum, Virgo, invioláta permansísti.\end{Resp}
\begin{Resp}\Vsign Benedícta tu in muliéribus, et benedíctus fructus ventris tui.\end{Resp}
\begin{Resp}\Rsign Quæ Gabriéli Archángeli dictis credidísti, dum Virgo Deum a hóminem genuísti, et post partum, Virgo, invioláta permansísti.\end{Resp}
\begin{Resp}\Vsign Glória Patri, et Fílio, \Star{} et Spirítui Sancto.\end{Resp}
\begin{Resp}\Rsign Quæ Gabriéli Archángeli dictis credidísti, dum Virgo Deum a hóminem genuísti, et post partum, Virgo, invioláta permansísti.\end{Resp}
\switchcolumn
\EN
\begin{Resp}\Rsign Rejoice, O Mary, by whose mighty hand the Church hath victory o'er her foes achieved \Star{} Since thou to Gabriel's word of quickening power in lowliness hast listened and believed thou, still a virgin, in thy blessed womb hast God Incarnate of thy flesh conceived, and, still of heaven, of that virginity remainest after childbirth unbereaved.\end{Resp}
\begin{Resp}\Vsign Blessed art thou among women, and blessed is the fruit of thy womb.\end{Resp}
\begin{Resp}\Rsign Since thou to Gabriel's word of quickening power in lowliness hast listened and believed thou, still a virgin, in thy blessed womb hast God Incarnate of thy flesh conceived, and, still of heaven, of that virginity remainest after childbirth unbereaved.\end{Resp}
\begin{Resp}\Vsign Glory be to the Father, and to the Son, \Star{} and to the Holy Ghost.\end{Resp}
\begin{Resp}\Rsign Since thou to Gabriel's word of quickening power in lowliness hast listened and believed thou, still a virgin, in thy blessed womb hast God Incarnate of thy flesh conceived, and, still of heaven, of that virginity remainest after childbirth unbereaved.\end{Resp}
\switchcolumn*

\end{paracol}

\par\needspace{10\baselineskip}\addvspace{1.2em}{\centering\small\textsc{Die 25 Martii} · \textsc{25 March}\par}
\Office{In Annuntiatione Beatæ Mariæ Virginis}{In Annuntiatione Beatæ Mariæ Virginis}{Duplex I classis}
\addcontentsline{toc}{subsection}{Die 25 Martii · In Annuntiatione Beatæ Mariæ Virginis\enspace\textit{In Annuntiatione Beatæ Mariæ Virginis}}
\begin{paracol}{2}
\LA
\RefLine{Lectio IX: de Scriptura occurrente.}
\switchcolumn
\EN
\RefLine{Lesson IX: from the Scripture of the day.}
\switchcolumn*

\LA
\LessonHead{Lectio I}
\switchcolumn
\EN
\LessonHead{Lesson I}
\switchcolumn*

\LA
\LessonTitle{De Isaía Prophéta}
\switchcolumn
\EN
\LessonTitle{Lesson from the book of Isaias}
\switchcolumn*

\LA
\Cite{Isa 7:10-15}
\switchcolumn
\EN
\Cite{Isa 7:10-15}
\switchcolumn*

\LA
\noindent Et adjécit Dóminus loqui ad Achaz, dicens: Pete tibi signum a Dómino Deo tuo in profúndum inférni, sive in excélsum supra. Et dixit Achaz: Non petam, et non tentábo Dóminum. Et dixit: Audíte ergo, domus David: Numquid parum vobis est, moléstos esse homínibus, quia molésti estis et Deo meo? Propter hoc dabit Dóminus ipse vobis signum. Ecce virgo concípiet, et páriet fílium, et vocábitur nomen ejus Emmánuel. Butýrum et mel cómedet, ut sciat reprobáre malum, et elígere bonum.\par
\switchcolumn
\EN
\noindent And the Lord spoke again to Achaz, saying: Ask thee a sign of the Lord thy God either unto the depth of hell, or unto the height above. And Achaz said: I will not ask, and I will not tempt the Lord. And he said: Hear ye therefore, O house of David: Is it a small thing for you to be grievous to men, that you are grievous to my God also? Therefore the Lord himself shall give you a sign. Behold a virgin shall conceive, and bear a son, and his name shall be called Emmanuel. He shall eat butter and honey, that he may know to refuse the evil, and to choose the good.\par
\switchcolumn*

\LA
\begin{Resp}\Rsign Missus est Gábriel Angelus ad Maríam Vírginem desponsátam Joseph, núntians ei verbum; et expavéscit Virgo de lúmine: ne tímeas, María, invenísti grátiam apud Dóminum: \Star{} Ecce concípies et páries, et vocábitur Altíssimi Fílius. (Allelúja.)\end{Resp}
\begin{Resp}\Vsign Dabit ei Dóminus Deus sedem David, patris ejus, et regnábit in domo Jacob in ætérnum.\end{Resp}
\begin{Resp}\Rsign Ecce concípies et páries, et vocábitur Altíssimi Fílius. (Allelúja.)\end{Resp}
\switchcolumn
\EN
\begin{Resp}\Rsign The angel Gabriel was sent to Mary, a Virgin espoused to Joseph, to bring unto her the word of the Lord and when the Virgin saw the light she was afraid. Fear not, Mary, for thou hast found grace from the Lord. \Star{} Behold, thou shalt conceive and bring forth a son, and He shall be called the Son of the Highest.\end{Resp}
\begin{Resp}\Vsign The Lord God shall give unto Him the throne of His father David, and He shall reign over the house of Jacob for ever.\end{Resp}
\begin{Resp}\Rsign Behold, thou shalt conceive and bring forth a son, and He shall be called the Son of the Highest.\end{Resp}
\switchcolumn*

\LA
\LessonHead{Lectio II}
\switchcolumn
\EN
\LessonHead{Lesson II}
\switchcolumn*

\LA
\Cite{Isa 11:1-5}
\switchcolumn
\EN
\Cite{Isa 11:1-5}
\switchcolumn*

\LA
\noindent Et egrediétur virga de radíce Jesse, et flos de radíce ejus ascéndet. Et requiéscet super eum Spíritus Dómini: spíritus sapiéntiæ et intelléctus, spíritus consílii et fortitúdinis, spíritus sciéntiæ et pietátis; Et replébit eum spíritus timóris Dómini: non secúndum visiónem oculórum judicábit, neque secúndum audítum áurium árguet: Sed judicábit in justítia páuperes, et árguet in æquitáte pro mansuétis terræ: et percútiet terram virga oris sui, et spíritu labiórum suórum interfíciet ímpium. Et erit justítia cíngulum lumbórum ejus: et fides cinctórium renum ejus.\par
\switchcolumn
\EN
\noindent And there shall come forth a rod out of the root of Jesse, and a flower shall rise up out of his root. And the spirit of the Lord shall rest upon him: the spirit of wisdom, and of understanding, the spirit of counsel, and of fortitude, the spirit of knowledge, and of godliness. And he shall be filled with the spirit of the fear of the Lord. He shall not judge according to the sight of the eyes, nor reprove according to the hearing of the ears. But he shall judge the poor with justice, and shall reprove with equity for the meek of the earth: and he shall strike the earth with the rod of his mouth, and with the breath of his lips he shall slay the wicked. And justice shall be the girdle of his loins: and faith the girdle of his reins.\par
\switchcolumn*

\LA
\begin{Resp}\Rsign Ave, María, grátia plena; Dóminus tecum: \Star{} Spíritus Sanctus supervéniet in te, et virtus Altíssimi obumbrábit tibi: quod enim ex te nascétur Sanctum, vocábitur Fílius Dei. (Allelúja.)\end{Resp}
\begin{Resp}\Vsign Quómodo fiet istud, quóniam virum non cognósco? Et respóndens Angelus, dixit ei.\end{Resp}
\begin{Resp}\Rsign Spíritus Sanctus supervéniet in te, et virtus Altíssimi obumbrábit tibi: quod enim ex te nascétur Sanctum, vocábitur Fílius Dei. (Allelúja.)\end{Resp}
\switchcolumn
\EN
\begin{Resp}\Rsign Hail, Mary, full of grace; the Lord is with thee; \Star{} The Holy Ghost shall come upon thee, and the power of the Highest shall overshadow thee therefore also that Holy Thing Which shall be born of thee shall be called the Son of God.\end{Resp}
\begin{Resp}\Vsign How shall this be, seeing I know not a man? And the Angel answered and said unto her,\end{Resp}
\begin{Resp}\Rsign The Holy Ghost shall come upon thee, and the power of the Highest shall overshadow thee; therefore also that Holy Thing Which shall be born of thee shall be called the Son of God.\end{Resp}
\switchcolumn*

\LA
\LessonHead{Lectio III}
\switchcolumn
\EN
\LessonHead{Lesson III}
\switchcolumn*

\LA
\Cite{Isa 35:1-7}
\switchcolumn
\EN
\Cite{Isa 35:1-7}
\switchcolumn*

\LA
\noindent Lætábitur desérta et ínvia, et exsultábit solitúdo, et florébit quasi lílium. Gérminans germinábit, et exsultábit lætabúnda et laudans: glória Líbani data est ei: decor Carméli et Saron, ipsi vidébunt glóriam Dómini, et decórem Dei nostri. Confortáte manus dissolútas, et génua debília roboráte. Dícite pusillánimis: Confortámini, et nolíte timére: ecce Deus vester ultiónem addúcet retributiónis: Deus ipse véniet, et salvábit vos. Tunc aperiéntur óculi cæcórum, et aures surdórum patébunt. Tunc sáliet sicut cervus claudus, et apérta erit lingua mutórum: quia scissæ sunt in desérto aquæ, et torréntes in solitúdine. Et quæ erat árida, erit in stagnum, et sítiens in fontes aquárum.\par
\switchcolumn
\EN
\noindent The land that was desolate and impassable shall be glad, and the wilderness shall rejoice, and shall flourish like the lily. It shall bud forth and blossom, and shall rejoice with joy and praise: the glory of Libanus is given to it: the beauty of Carmel, and Saron, they shall see the glory of the Lord, and the beauty of our God. Strengthen ye the feeble hands, and confirm the weak knees. Say to the fainthearted: Take courage, and fear not: behold your God will bring the revenge of recompense: God himself will come and will save you. Then shall the eyes of the blind be opened, and the ears of the deaf shall be unstopped. Then shall the lame man leap as a hart, and the tongue of the dumb shall be free: for waters are broken out in the desert, and streams in the wilderness. And that which was dry land, shall become a pool, and the thirsty land springs of water.\par
\switchcolumn*

\LA
\begin{Resp}\Rsign Súscipe verbum, Virgo María, quod tibi a Dómino per Angelum transmíssum est: concípies et páries Deum páriter et hóminem, \Star{} Ut benedícta dicáris inter omnes mulíeres. (Allelúja.)\end{Resp}
\begin{Resp}\Vsign Páries quidem fílium, et virginitátis non patiéris detriméntum: efficiéris grávida, et eris mater semper intácta.\end{Resp}
\begin{Resp}\Rsign Ut benedícta dicáris inter omnes mulíeres. (Allelúja.)\end{Resp}
\begin{Resp}\Vsign Glória Patri, et Fílio, \Star{} et Spirítui Sancto.\end{Resp}
\begin{Resp}\Rsign Ut benedícta dicáris inter omnes mulíeres. (Allelúja.)\end{Resp}
\switchcolumn
\EN
\begin{Resp}\Rsign Receive, O Virgin Mary, receive the word of the Lord, which is sent thee by His Angel; thou shalt conceive, and shalt bring forth God and Man together. \Star{} And thou shalt be called blessed among all women.\end{Resp}
\begin{Resp}\Vsign Thou shalt bring forth a son, and remain a maiden undefiled; thou shalt conceive and be a Mother, still Virgin unspotted.\end{Resp}
\begin{Resp}\Rsign And thou shalt be called blessed among all women.\end{Resp}
\begin{Resp}\Vsign Glory be to the Father, and to the Son, \Star{} and to the Holy Ghost.\end{Resp}
\begin{Resp}\Rsign And thou shalt be called blessed among all women.\end{Resp}
\switchcolumn*

\LA
\LessonHead{Lectio IV}
\switchcolumn
\EN
\LessonHead{Lesson IV}
\switchcolumn*

\LA
\LessonTitle{Sermo sancti Leónis Papæ}
\switchcolumn
\EN
\LessonTitle{From the Sermons of Pope St. Leo (the Great.)}
\switchcolumn*

\LA
\Source{Sermo 2 de Nativ. Domini}
\switchcolumn
\EN
\Source{2nd for Christmas}
\switchcolumn*

\LA
\noindent Deus omnípotens et clemens, cujus natúra bónitas, cujus volúntas poténtia, cujus opus misericórdia est, statim ut nos diabólica malígnitas venéno suæ mortificávit invídiæ, prædestináta renovándis mortálibus suæ pietátis remédia, inter ipsa mundi primórdia præsignávit; denúntians serpénti futúrum semen mulíeris, quod nóxii cápitis elatiónem sua virtúte contéreret, Christum scílicet in carne ventúrum; Deum hominémque signans, qui natus ex Vírgine, violatórem humánæ propáginis incorrúpta nativitáte damnáret.\par
\switchcolumn
\EN
\noindent The Almighty and merciful God, Whose nature is goodness, Whose will is power, and Whose work is mercy, did, at the very beginning of the world, as soon as the devil's hatred had mortally poisoned us with the venom of his envy, foretell those remedies which His mercy had foreordained for our healing. He bade the serpent know that there was to be a Seed of the woman Who should yet bruise the swelling of his pestilential head; this Seed was none other than the Christ to come in the flesh, that God and Man in one Person, Who, being born of a Virgin, should, by His undefiled birth, damn the seducer of man.\par
\switchcolumn*

\LA
\begin{Resp}\Rsign Ecce virgo concípiet, et páriet fílium, dicit Dóminus: \Star{} Et vocábitur nomen ejus Admirábilis, Deus, Fortis. (Allelúja.)\end{Resp}
\begin{Resp}\Vsign Super sólium David, et super regnum ejus sedébit in ætérnum.\end{Resp}
\begin{Resp}\Rsign Et vocábitur nomen ejus Admirábilis, Deus, Fortis. (Allelúja.)\end{Resp}
\switchcolumn
\EN
\begin{Resp}\Rsign Behold, the Virgin shall conceive, and bear a son, saith the Lord; \Star{} And His name shall be called Wonderful, the Mighty God.\end{Resp}
\begin{Resp}\Vsign He shall sit upon the throne of David, and upon his kingdom for ever.\end{Resp}
\begin{Resp}\Rsign And His name shall be called Wonderful, the Mighty God.\end{Resp}
\switchcolumn*

\LA
\LessonHead{Lectio V}
\switchcolumn
\EN
\LessonHead{Lesson V}
\switchcolumn*

\LA
\noindent Nam quia gloriabátur diábolus hóminem sua fraude decéptum divínis caruísse munéribus, et immortalitátis dote nudátum, duram mortis subísse senténtiam, seque in malis suis quoddam de prævaricatóris consórtio invenísse solátium; Deum quoque, justæ severitátis exigénte ratióne, erga hóminem, quem in tanto honóre condíderat, antíquam mutásse senténtiam: opus fuit, dilectíssimi, secréti dispensatióne consílii ut incommutábilis Deus, cujus volúntas non potest sua benignitáte privári, primam pietátis suæ dispositiónem sacraménto occultióre compléret; et homo diabólicæ iniquitátis versútia actus in culpam, contra Dei propósitum non períret.\par
\switchcolumn
\EN
\noindent The devil rejoiced that by his fraud he had so deceived man as to make him lose the gifts of God, forfeit his privilege of eternal life, bring himself under the hard sentence of death, and find in his misery a certain comfort in the accomplice of his guilt; he rejoiced also that God, in His just anger, was changed towards man, whom He had made in such honour. But, dearly beloved brethren, that Unchangeable God, Whose Will cannot be divorced from His goodness, by His own secret counsel carried out in a mysterious way His original purpose of goodness, and man, who had been led into sin by the wicked craft of the devil, perished not to disappoint that gracious purpose of God.\par
\switchcolumn*

\LA
\begin{Resp}\Rsign Egrediétur virga de radíce Jesse, et flos de radíce ejus ascéndet: \Star{} Et erit justítia cíngulum lumbórum ejus, et fides cinctórium renum ejus. (Allelúja.)\end{Resp}
\begin{Resp}\Vsign Et requiéscet super eum Spíritus Dómini: spíritus sapiéntiæ et intelléctus, spíritus consílii et fortitúdinis.\end{Resp}
\begin{Resp}\Rsign Et erit justítia cíngulum lumbórum ejus, et fides cinctórium renum ejus. (Allelúja.)\end{Resp}
\switchcolumn
\EN
\begin{Resp}\Rsign There shall come forth a rod out of the stem of Jesse, and a Flower shall grow out of his roots. \Star{} And righteousness shall be the girdle of his loins, and faithfulness the girdle of his reins.\end{Resp}
\begin{Resp}\Vsign And the Spirit of the Lord shall rest upon him, the Spirit of wisdom and understanding, the Spirit of counsel and might.\end{Resp}
\begin{Resp}\Rsign And righteousness shall be the girdle of his loins, and faithfulness the girdle of his reins.\end{Resp}
\switchcolumn*

\LA
\LessonHead{Lectio VI}
\switchcolumn
\EN
\LessonHead{Lesson VI}
\switchcolumn*

\LA
\noindent Adveniéntibus ergo tempóribus, dilectíssimi, quæ redemptióni hóminum fúerant præstitúta, ingréditur hæc ínfima Jesus Christus, Dóminus noster, de cælésti sede descéndens, et a patérna glória non recédens, novo órdine, nova nativitáte generátus: novo órdine, quia, invisíbilis in suis, visíbilis factus est in nostris; incomprehensíbilis, vóluit comprehéndi; ante témpora manens, esse cœpit ex témpore; universitátis Dóminus servílem formam, obumbráta majestátis suæ dignitáte, suscépit; impassíbilis Deus, non dedignátus est homo esse passíbilis; et immortális, mortis légibus subjacére.\par
\switchcolumn
\EN
\noindent Then therefore, dearly beloved brethren, the fulness of that time came, which God had appointed for our Redemption, our Lord Jesus Christ entered this lower world, came down from His heavenly throne, and, while He left not that glory which He hath with the Father before the world was, was incarnate by a new order and a new birth new, in that He Who is Invisible among His own, was made visible among us; He Who is Incomprehensible, willed to be comprehended; He Who is before the ages, began to be in time; the Lord of all shadowed the glory of His Majesty, and took upon Him the form of a servant; the Impassible God vouchsafed to become a man subject to suffering; and the Immortal laid Himself under the laws of death.\par
\switchcolumn*

\LA
\begin{Resp}\Rsign Sancta et immaculáta virgínitas, quibus te láudibus éfferam, néscio: \Star{} Quia quem cæli cápere non póterant, tuo grémio contulísti. (Allelúja.)\end{Resp}
\begin{Resp}\Vsign Benedícta tu in muliéribus, et benedíctus fructus ventris tui.\end{Resp}
\begin{Resp}\Rsign Quia quem cæli cápere non póterant, tuo grémio contulísti. (Allelúja.)\end{Resp}
\begin{Resp}\Vsign Glória Patri, et Fílio, \Star{} et Spirítui Sancto.\end{Resp}
\begin{Resp}\Rsign Quia quem cæli cápere non póterant, tuo grémio contulísti. (Allelúja.)\end{Resp}
\switchcolumn
\EN
\begin{Resp}\Rsign O Mary, how holy and how spotless is thy virginity! I am too dull to praise thee: \Star{} For thou hast borne in thy breast Him Whom the heavens cannot contain.\end{Resp}
\begin{Resp}\Vsign Blessed art thou among women, and blessed is the fruit of thy womb.\end{Resp}
\begin{Resp}\Rsign For thou hast borne in thy breast Him Whom the heavens cannot contain.\end{Resp}
\begin{Resp}\Vsign Glory be to the Father, and to the Son, \Star{} and to the Holy Ghost.\end{Resp}
\begin{Resp}\Rsign For thou hast borne in thy breast Him Whom the heavens cannot contain.\end{Resp}
\switchcolumn*

\LA
\LessonHead{Lectio VII}
\switchcolumn
\EN
\LessonHead{Lesson VII}
\switchcolumn*

\LA
\LessonTitle{Léctio sancti Evangélii secúndum Lucam}
\switchcolumn
\EN
\LessonTitle{From the Holy Gospel according to Luke}
\switchcolumn*

\LA
\Cite{Luc 1:26-38}
\switchcolumn
\EN
\Cite{Luke 1:26-38}
\switchcolumn*

\LA
\noindent In illo témpore: Missus est Angelus Gábriel a Deo in civitátem Galilǽæ, cui nomen Názareth, ad Vírginem desponsátam viro, cui nomen erat Joseph, de domo David, et nomen Vírginis María. Et réliqua.\par
\switchcolumn
\EN
\noindent At that time the angel Gabriel was sent from God into a city of Galilee, called Nazareth, To a virgin espoused to a man whose name was Joseph, of the house of David; and the virgin's name was Mary. And so on.\par
\switchcolumn*

\LA
\LessonTitle{Homilía sancti Ambrósii Epíscopi}
\switchcolumn
\EN
\LessonTitle{Homily by St. Ambrose, Bishop (of Milan.)}
\switchcolumn*

\LA
\Source{Liber 2 in Lucam}
\switchcolumn
\EN
\Source{Bk. ii on Luke}
\switchcolumn*

\LA
\noindent Latent quidem divína mystéria, nec fácile, juxta prophéticum dictum, quisquam hóminum potest scire consílium Dei. Sed tamen ex céteris factis, atque præcéptis Dómini Salvatóris póssumus intellégere, et hoc propensióris fuísse consílii, quod ea potíssimum elécta est, ut Dóminum páreret, quæ erat desponsáta viro. Cur autem non ántequam desponsarétur, impléta est? Fortásse ne dicerétur quod concéperat ex adultério.\par
\switchcolumn
\EN
\noindent The mysteries of God are unsearchable, and it is especially declared by a Prophet, that a man can hardly know His counsels. (Wisd. ix. 13.) Nevertheless, some things have been revealed to us, and we may gather from some of the words and works of the Lord our Saviour, that there was a special purpose of God, in the fact that she who was chosen to be the mother of the Lord was espoused to a man. Why did not the power of the Highest overshadow her before she was so espoused? Perhaps it was lest any might blasphemously say that she had conceived in fornication the Holy One.\par
\switchcolumn*

\LA
\begin{Resp}\Rsign Congratulámini mihi, omnes qui dilígitis Dóminum: quia cum essem párvula, plácui Altíssimo, \Star{} Et de meis viscéribus génui Deum et hóminem. (Allelúja.)\end{Resp}
\begin{Resp}\Vsign Beátam me dicent omnes generatiónes, quia ancíllam húmilem respéxit Deus.\end{Resp}
\begin{Resp}\Rsign Et de meis viscéribus génui Deum et hóminem. (Allelúja.)\end{Resp}
\switchcolumn
\EN
\begin{Resp}\Rsign Rejoice with me, all ye that love the Lord, for while I was yet a little one, I pleased the Most High. \Star{} And I have brought forth from my bowels God and man.\end{Resp}
\begin{Resp}\Vsign All generations shall call me blessed, since the Lord hath regarded the lowliness of His handmaiden.\end{Resp}
\begin{Resp}\Rsign And I have brought forth from my bowels God and man.\end{Resp}
\switchcolumn*

\LA
\LessonHead{Lectio VIII}
\switchcolumn
\EN
\LessonHead{Lesson VIII}
\switchcolumn*

\LA
\noindent Et ingréssus ad eam Angelus. Disce vírginem móribus, disce vírginem verecúndia, disce oráculo, disce mystério. Trepidáre vírginum est, et ad omnes viri ingréssus pavére, omnes viri affátus veréri. Discant mulíeres propósitum pudóris imitári. Sola in penetrálibus, quam nemo virórum víderit, solus Angelus repérerit: sola sine cómite, sola sine teste, ne quo degénere depravarétur affátu, ab Angelo salutátur.\par
\switchcolumn
\EN
\noindent And the angel came in unto her. Let us learn from this Virgin how to bear ourselves, let us learn her modesty, let us learn by her devout utterance, above all let us learn by the holy mystery enacted. It is the part of a maiden to be timid, to avoid the advances of men, and to shrink from men's addresses. Would that our women would learn from the example of modesty here set before us. She upon whom the stare of men had never been fixed was alone in her chamber, and was found only by an angel. There was neither companion nor witness there, that what passed might not be debased in gossip and the angel saluted her.\par
\switchcolumn*

\LA
\begin{Resp}\Rsign Gaude, María Virgo, cunctas hǽreses sola interemísti, quæ Gabriélis Archángeli dictis credidísti: \Star{} Dum Virgo Deum et hóminem genuísti, et post partum, Virgo, invioláta permansísti. (Allelúja.)\end{Resp}
\begin{Resp}\Vsign Beáta es quæ credidísti: quia perfécta sunt ea, quæ dicta sunt tibi a Dómino.\end{Resp}
\begin{Resp}\Rsign Dum Virgo Deum et hóminem genuísti, et post partum, Virgo, invioláta permansísti. (Allelúja.)\end{Resp}
\begin{Resp}\Vsign Glória Patri, et Fílio, \Star{} et Spirítui Sancto.\end{Resp}
\begin{Resp}\Rsign Dum Virgo Deum et hóminem genuísti, et post partum, Virgo, invioláta permansísti. (Allelúja.)\end{Resp}
\switchcolumn
\EN
\begin{Resp}\Rsign Rejoice, O Mary, by whose mighty hand the Church hath victory o'er her foes achieved, since thou to Gabriel's word of quickening power in lowliness hast listened, and believed; \Star{} Thou, still a virgin, in thy blessed womb hast God Incarnate of thy flesh conceived, and, still of heaven, of that virginity remainest after childbirth unbereaved.\end{Resp}
\begin{Resp}\Vsign Blessed art thou that hast believed, for there is a performance of those things which were told thee from the Lord.\end{Resp}
\begin{Resp}\Rsign Thou, still a virgin, in thy blessed womb hast God Incarnate of thy flesh conceived, and, still of heaven, of that virginity remainest after childbirth unbereaved.\end{Resp}
\begin{Resp}\Vsign Glory be to the Father, and to the Son, \Star{} and to the Holy Ghost.\end{Resp}
\begin{Resp}\Rsign Thou, still a virgin, in thy blessed womb hast God Incarnate of thy flesh conceived, and, still of heaven, of that virginity remainest after childbirth unbereaved.\end{Resp}
\switchcolumn*

\LA
\LessonHead{Lectio IX}
\switchcolumn
\EN
\LessonHead{Lesson IX}
\switchcolumn*

\LA
\noindent Tanti namque mandáti mystérium non hóminis fuit, sed Angeli ore proméndum. Hódie primum audítur: Spíritus Sanctus supervéniet in te. Et audítur, et créditur. Dénique, Ecce, inquit, ancílla Dómini: contíngat mihi secúndum verbum tuum. Vide humilitátem, vide devotiónem. Ancíllam se dicit Dómini, quæ mater elígitur: nec repentíno exaltáta promísso est.\par
\switchcolumn
\EN
\noindent The message of God to the Virgin was a mystery, which it was not lawful for the mouth of men, but only of angels, to utter. For the first time on earth the words are spoken The Holy Ghost shall come upon thee. The holy maiden heareth, and believeth. At length she saith Behold the handmaid of the Lord be it unto me according to thy word. Here is an example of lowliness, here is a pattern of true devotion. At the very moment that she is told she is chosen to be the mother of the Lord she at once declareth herself His handmaid. The knowledge that she was Mother of God caused in the heart of Mary only an act of humility.\par
\switchcolumn*

\LA
\TeDeum{Te Deum}
\switchcolumn
\EN
\TeDeum{Te Deum}
\switchcolumn*

\end{paracol}

\par\needspace{10\baselineskip}\addvspace{1.2em}{\centering\small\textsc{Die 27 Martii} · \textsc{27 March}\par}
\Office{S. Joannis Damasceni Confessoris et Ecclesiæ Doctoris}{St. John Damascene, Confessor}{Duplex}
\addcontentsline{toc}{subsection}{Die 27 Martii · S. Joannis Damasceni Confessoris et Ecclesiæ Doctoris\enspace\textit{St. John Damascene, Confessor}}
\begin{paracol}{2}
\LA
\RefLine{Lectio IX: de Scriptura occurrente.}
\switchcolumn
\EN
\RefLine{Lesson IX: from the Scripture of the day.}
\switchcolumn*

\LA
\LessonHead{Lectio I}
\switchcolumn
\EN
\LessonHead{Lesson I}
\switchcolumn*

\LA
\LessonTitle{De libro Ecclesiástici}
\switchcolumn
\EN
\LessonTitle{Lesson from the book of Ecclesiasticus}
\switchcolumn*

\LA
\Cite{Sir 39:1-5}
\switchcolumn
\EN
\Cite{Sir 39:1-5}
\switchcolumn*

\LA
\noindent Sapiéntiam ómnium antiquórum exquíret sápiens, et in prophétis vacábit. Narratiónem virórum nominatórum conservábit, et in versútias parabolárum simul introíbit. Occúlta proverbiórum exquíret, et in abscónditis parabolárum conversábitur. In médio magnatórum ministrábit, et in conspéctu prǽsidis apparébit. In terram alienigenárum géntium pertránsiet; bona enim et mala in homínibus tentábit.\par
\switchcolumn
\EN
\noindent The wise men will seek out the wisdom of all the ancients, and will be occupied in the prophets. He will keep the sayings of renowned men, and will enter withal into the subtilties of parables. He will search out the hidden meanings of proverbs, and will be conversant in the secrets of parables. He shall serve among great men, and: appear before the governor. He shall pass into strange countries: for he shall try good and evil among men.\par
\switchcolumn*

\LA
\begin{Resp}\Rsign Euge, serve bone et fidélis, quia in pauca fuísti fidélis, supra multa te constítuam: \Star{} Intra in gáudium Dómini tui.\end{Resp}
\begin{Resp}\Vsign Dómine, quinque talénta tradidísti mihi, ecce ália quinque superlucrátus sum.\end{Resp}
\begin{Resp}\Rsign Intra in gáudium Dómini tui.\end{Resp}
\switchcolumn
\EN
\begin{Resp}\Rsign Well done, thou good and faithful servant, thou hast been faithful over a few things. I will make thee ruler over many things: \Star{} Enter thou into the joy of thy Lord.\end{Resp}
\begin{Resp}\Vsign Lord, thou deliveredst unto me five talents; behold, I have gained beside them five talents more.\end{Resp}
\begin{Resp}\Rsign Enter thou into the joy of thy Lord.\end{Resp}
\switchcolumn*

\LA
\LessonHead{Lectio II}
\switchcolumn
\EN
\LessonHead{Lesson II}
\switchcolumn*

\LA
\Cite{Sir 39:6-10}
\switchcolumn
\EN
\Cite{Sir 39:6-10}
\switchcolumn*

\LA
\noindent Cor suum tradet ad vigilándum dilúculo ad Dóminum, qui fecit illum, et in conspéctu Altíssimi deprecábitur. Apériet os suum in oratióne, et pro delíctis suis deprecábitur. Si enim Dóminus magnus volúerit, spíritu intellegéntiæ replébit illum: et ipse tamquam imbres mittet elóquia sapiéntiæ suæ, et in oratióne confitébitur Dómino: et ipse díriget consílium ejus, et disciplínam, et in abscónditis suis consiliábitur.\par
\switchcolumn
\EN
\noindent He will give his heart to resort early to the Lord that made him, and he will pray in the sight of the most High. He will open his mouth in prayer, and will make supplication for his sins. For if it shall please the great Lord, he will fill him with the spirit of understanding: And he will pour forth the words of his wisdom as showers, and in his prayer he will confess to the Lord. And he shall direct his counsel, and his knowledge, and in his secrets shall he meditate.\par
\switchcolumn*

\LA
\begin{Resp}\Rsign Justus germinábit sicut lílium: \Star{} Et florébit in ætérnum ante Dóminum.\end{Resp}
\begin{Resp}\Vsign Plantátus in domo Dómini, in átriis domus Dei nostri.\end{Resp}
\begin{Resp}\Rsign Et florébit in ætérnum ante Dóminum.\end{Resp}
\switchcolumn
\EN
\begin{Resp}\Rsign The righteous shall grow as the lily; \Star{} Yea, he shall flourish in the presence of the Lord for ever.\end{Resp}
\begin{Resp}\Vsign Those that be planted in the house of the Lord, shall flourish in the courts of the house of our God.\end{Resp}
\begin{Resp}\Rsign Yea, he shall flourish in the presence of the Lord for ever.\end{Resp}
\switchcolumn*

\LA
\LessonHead{Lectio III}
\switchcolumn
\EN
\LessonHead{Lesson III}
\switchcolumn*

\LA
\Cite{Sir 39:11-14}
\switchcolumn
\EN
\Cite{Sir 39:11-14}
\switchcolumn*

\LA
\noindent Ipse palam fáciet disciplínam doctrínæ suæ, et in lege testaménti Dómini gloriábitur. Collaudábunt multi sapiéntiam ejus, et usque in sǽculum non delébitur. Non recédet memória ejus, et nomen ejus requirétur a generatióne in generatiónem. Sapiéntiam ejus enarrábunt gentes, et laudem ejus enuntiábit Ecclésia.\par
\switchcolumn
\EN
\noindent He shall show forth the discipline he hath learned, and shall glory in the law of the covenant of the Lord. Many shall praise his wisdom, and it shall never be forgotten. The memory of him shall not depart away, and his name shall be in request from generation to generation. Nations shall declare his wisdom, and the church shall show forth his praise.\par
\switchcolumn*

\LA
\begin{Resp}\Rsign Iste cognóvit justítiam, et vidit mirabília magna, et exorávit Altíssimum: \Star{} Et invéntus est in número Sanctórum.\end{Resp}
\begin{Resp}\Vsign Iste est qui contémpsit vitam mundi, et pervénit ad cæléstia regna.\end{Resp}
\begin{Resp}\Rsign Et invéntus est in número Sanctórum.\end{Resp}
\begin{Resp}\Vsign Glória Patri, et Fílio, \Star{} et Spirítui Sancto.\end{Resp}
\begin{Resp}\Rsign Et invéntus est in número Sanctórum.\end{Resp}
\switchcolumn
\EN
\begin{Resp}\Rsign This is he which knew righteousness, and saw great wonders, and made his prayer unto the Most High; \Star{} And he is numbered among the Saints.\end{Resp}
\begin{Resp}\Vsign This is he which loved not his life in this world, and is come unto an everlasting kingdom.\end{Resp}
\begin{Resp}\Rsign And he is numbered among the Saints.\end{Resp}
\begin{Resp}\Vsign Glory be to the Father, and to the Son, \Star{} and to the Holy Ghost.\end{Resp}
\begin{Resp}\Rsign And he is numbered among the Saints.\end{Resp}
\switchcolumn*

\LA
\LessonHead{Lectio IV}
\switchcolumn
\EN
\LessonHead{Lesson IV}
\switchcolumn*

\LA
\noindent Joánnes, a pátrio loco Damascenus dictus, nobili genere natus, humanis divinisque litteris a Cosma monacho Constantinopoli fuit excultus; cumque ea tempestáte imperátor Leo Isáuricus nefario bello sacrárum imaginum cultum insectarétur, Joánnes, hortatu Gregórii tertii Romani Pontificis, et sermóne et scriptis sanctitátem illíus cultus sédulo propugnávit. Quo facto tantam Leónis advérsum se invidiam concitávit, ut hic confictis litteris ipsum tamquam proditórem accusarit apud Damasci calípham, qui Joánne consiliario et administro utebátur. Credulus fraudi princeps Joánni nequidquam calumniam ejuranti præcídi déxteram jussit. Verum innocéntiæ vindex ádfuit cliénti suo sanctíssima Virgo, cujus opem precibus enixe imploraverat, ejusque beneficio trunca manus restituta ita brácchio coáluit, ac si divisa numquam fuísset. Quo maxime miraculo permotus Joánnes, quod pridem animo conceperat, éxsequi státuit. Itaque ægre a calípha impetrato secessu, suas omnes facultates in egénos distríbuit, et servos libertáte donávit; tum sacra Palæstinæ loca peregrínus lustrávit, ac demum una cum Cosma institutore suo in lauram sancti Sabbæ prope Hierosolymam concessit, ibique présbyter initiátus est.\par
\switchcolumn
\EN
\noindent This John is called John of Damascus, from his native place. He was of noble birth, and studied sacred and profane letters at Constantinople, under the monk Cosmas. At what time the Emperor Leo the Isaurian was making a wicked attack upon the honouring of holy images, John, at the desire of the Roman Pontiff, Gregory III., earnestly defended both by his words and his writings, the holiness of this honour. By this he roused against him so great a hatred on the part of Leo, that that Prince, by forged letters, accused John as a traitor to the Caliph of Damascus, whom he was serving as a councillor and minister. John denied the charge, but the Caliph was deceived by it, and caused his right hand to be cut off. He called earnestly for the help of the. most holy Virgin, and she manifested the innocency of her servant by reuniting his hand to his arm, as though it had never been cut off. This miracle moved John to carry out a design which he had long had in mind. He obtained from the Caliph, albeit with difficulty, leave to go away, distributed all his goods to feed the poor, and freed all his slaves, then visited as a pilgrim the holy places in Palestine, and at length withdrew, along with his teacher Cosmas, to the monastery of St. Saba, between Jerusalem and the Dead Sea. There he was ordained priest.\par
\switchcolumn*

\LA
\begin{Resp}\Rsign Honéstum fecit illum Dóminus, et custodívit eum ab inimícis, et a seductóribus tutávit illum: \Star{} Et dedit illi claritátem ætérnam.\end{Resp}
\begin{Resp}\Vsign Justum dedúxit Dóminus per vias rectas, et osténdit illi regnum Dei.\end{Resp}
\begin{Resp}\Rsign Et dedit illi claritátem ætérnam.\end{Resp}
\switchcolumn
\EN
\begin{Resp}\Rsign The Lord made him honourable, and defended him from his enemies, and kept him safe from those that lay in wait for him; \Star{} And gave him perpetual glory.\end{Resp}
\begin{Resp}\Vsign He went down with him into the pit, and left him not in bonds.\end{Resp}
\begin{Resp}\Rsign And gave him perpetual glory.\end{Resp}
\switchcolumn*

\LA
\LessonHead{Lectio V}
\switchcolumn
\EN
\LessonHead{Lesson V}
\switchcolumn*

\LA
\noindent In religióne vitæ palæstra præclarióra virtútum exempla monachis præbuit, demissiónis potíssimum et obediéntiæ. Abjectíssima quæque cœnobii munia véluti sibi propria deposcebat, ac sédulo obíbat. Contextas a se spórtulas venditare Damasci jussus, in ea nimirum civitáte ubi olim summis honóribus perfunctus fuerat, irrisiónes ac ludíbria vulgi ávide captábat. Obediéntiam ádeo cóluit, ut non modo ad quémlibet præsidum nutum præsto esset; sed ne causam quidem eórum quæ præcipiebántur, quamvis ardua essent et insolita, quæréndam sibi umquam putarit. Inter has virtútum exercitatiónes, catholicum dogma de sanctárum imaginum cultu impense tueri numquam déstitit. Quare ut ante Leónis Isáurici, ita postmodum Constantini Coprónymi advérsum se odia vexatiónesque provocávit; eo vel magis quod líbere arrogantiam imperatórum retúnderet, qui fidei negotia pertractare, deque his senténtiam arbitratu suo ferre audebant.\par
\switchcolumn
\EN
\noindent As a monk John set a bright example to all the others, especially as regarded lowliness and obedience. He sought for the lowest offices in the community, as though they were in a peculiar sense his own, and fulfilled them with the greatest care. When he was sent to Damascus to sell baskets made by himself, he welcomed the mockery and jests of the lowest classes in that city where he had before time been charged with the most honourable offices. He was so devoted to obedience that he not only started up to obey every nod of his superiors, but also never thought it right to ask the reason of any duty laid upon him, however difficult or however strange it might be. While thus living he never ceased earnestly to defend the Catholic doctrine as to the honouring of holy images. For this reason he drew upon himself the hatred and persecution of the Emperor Constantine Copronymus, as he had first done that of the Emperor Leo the Isaurian, and this all the more because he freely rebuked the arrogance of these Emperors, who must needs take in hand matters concerning the faith, and pronounce sentence upon them according to their own judgment.\par
\switchcolumn*

\LA
\begin{Resp}\Rsign Amávit eum Dóminus, et ornávit eum: stolam glóriæ índuit eum, \Star{} Et ad portas paradísi coronávit eum.\end{Resp}
\begin{Resp}\Vsign Índuit eum Dóminus lorícam fídei, et ornávit eum.\end{Resp}
\begin{Resp}\Rsign Et ad portas paradísi coronávit eum.\end{Resp}
\switchcolumn
\EN
\begin{Resp}\Rsign The Lord loved him and beautified him; He clothed him with a robe of glory, \Star{} And crowned him at the gates of Paradise.\end{Resp}
\begin{Resp}\Vsign The Lord hath put on him the breast-plate of faith, and hath adorned him.\end{Resp}
\begin{Resp}\Rsign And crowned him at the gates of Paradise.\end{Resp}
\switchcolumn*

\LA
\LessonHead{Lectio VI}
\switchcolumn
\EN
\LessonHead{Lesson VI}
\switchcolumn*

\LA
\noindent Mirum sane est quam multa tum ad fidem tutandam, tum ad pietátem fovéndam, et soluta et adstricta numeris oratióne, Joánnes elucubráverit; dignus sane qui ab áltera Nicæna synodo amplíssimis laudibus celebrarétur, et ob aureum oratiónis flumen Chrysórrhoas appellarétur. Neque solum contra Iconómachos orthodoxam fidem deféndit; sed omnes ferme hæreticos, præsertim Acéphalos, Monothelitas, Theopaschitas strenue impugnávit. Ecclésiæ jura potestátemque egregie vindicávit. Primátum Principis Apostolórum disertíssimis verbis asseruit; ipsumque ecclesiárum cólumen, infractam petram, orbis terrárum magistrum et moderatórem sæpius nóminat. Universa autem ejus scripta non modo eruditióne et doctrina præstant, sed étiam quemdam ingenuæ pietátis sensum præférunt, præcipue cum Genitrícis Dei laudes prædicat, quam singulari cultu et amóre prosequebátur. Illud vero maxime in laudem Joánnis cedit, quod primus universam theologíam recto ordine comprehénderit et sancti Thomæ viam complanáverit ad sacram doctrinam tam præclára méthodo tractandam. Tandem vir sanctíssimus meritis plenus devexáque jam ætate, in pace Christi quiévit anno circiter septingentésimo quinquagesimo quarto. Ejus Offícium et Missam Leo décimus tertius Pontifex maximus, áddito Doctoris título, univérsæ Ecclésiæ concessit.\par
\switchcolumn
\EN
\noindent It is a marvel how many things John devised both for the protection of the faith, and for the encouragement of godliness, and expressed in his writings both in prose and verse. He was worthy of the high praise which was given him by the Second Council of Nicea. On account of the golden streams of his eloquence, he was surnamed Chrysorrhoas, or John of the golden streams. It was not against the enemies of holy images alone that he defended the orthodox faith. He fought stoutly against the Acephali, the Monothelites, and the Theopaschites. He maintained the laws and the power of the Church. He taught with great learning the Primacy of the Prince of the Apostles, and many times calleth him the Pillar of the Churches, the unbroken rock, and the Teacher and Ruler of the world. The whole of his writings are not only steeped in learning and teaching, but have a certain savour of simple piety, especially when he is praising the Mother of God, toward whom he was filled with a special reverence and love. But the greatest praise of John is that he was the first who arranged in order a complete course of theology, and prepared the way in which holy Thomas of Aquino has so clearly dealt with the whole body of sacred doctrine. This truly holy man, full of days and good works, fell asleep in the peace of Christ about the year of salvation 754. The supreme Pontiff, Leo XIII., established his office and Mass throughout the universal church, whereof he also gave him the title of doctor.\par
\switchcolumn*

\LA
\begin{Resp}\Rsign Iste homo perfécit ómnia quæ locútus est ei Deus, et dixit ad eum: Ingrédere in réquiem meam: \Star{} Quia te vidi justum coram me ex ómnibus géntibus.\end{Resp}
\begin{Resp}\Vsign Iste est, qui contémpsit vitam mundi, et pervénit ad cæléstia regna.\end{Resp}
\begin{Resp}\Rsign Quia te vidi justum coram me ex ómnibus géntibus.\end{Resp}
\begin{Resp}\Vsign Glória Patri, et Fílio, \Star{} et Spirítui Sancto.\end{Resp}
\begin{Resp}\Rsign Quia te vidi justum coram me ex ómnibus géntibus.\end{Resp}
\switchcolumn
\EN
\begin{Resp}\Rsign This is he which did according unto all that God commanded him; and God said unto him: Enter thou into My rest, \Star{} For thee have I seen righteous before Me among all people.\end{Resp}
\begin{Resp}\Vsign This is he which loved not his life in this world, and is come unto an everlasting kingdom.\end{Resp}
\begin{Resp}\Rsign For thee have I seen righteous before Me among all people.\end{Resp}
\begin{Resp}\Vsign Glory be to the Father, and to the Son, \Star{} and to the Holy Ghost.\end{Resp}
\begin{Resp}\Rsign For thee have I seen righteous before Me among all people.\end{Resp}
\switchcolumn*

\LA
\LessonHead{Lectio VII}
\switchcolumn
\EN
\LessonHead{Lesson VII}
\switchcolumn*

\LA
\LessonTitle{Léctio sancti Evangélii secúndum Lucam}
\switchcolumn
\EN
\LessonTitle{From the Holy Gospel according to Luke}
\switchcolumn*

\LA
\Cite{Luc 6:6-11}
\switchcolumn
\EN
\Cite{Luke 6:6-11}
\switchcolumn*

\LA
\noindent In illo témpore: Factum est et in alio sabbato, ut intraret Jesus in synagógam, et docéret: et erat ibi homo, et manus ejus déxtera erat árida. Et réliqua.\par
\switchcolumn
\EN
\noindent At that time: It came to pass also on another Sabbath, that Jesus entered into the synagogue, and taught; and there was there a man whose right hand was withered. And so on.\par
\switchcolumn*

\LA
\LessonTitle{Homilía sancti Petri Chrysologi}
\switchcolumn
\EN
\LessonTitle{Homily by St. Peter Chrysologus.}
\switchcolumn*

\LA
\Source{Sermo 32}
\switchcolumn
\EN
\Source{Sermon 32}
\switchcolumn*

\LA
\noindent In hoc hómine ómnium hóminum imago figurátur, in hoc geritur cura cunctórum, in hoc universórum sánitas diu exspectáta reparátur. Arúerat enim manus hóminis magis stupore fidei quam siccitate nervórum, et plus culpa consciéntiæ quam debilitate carnali. Antíqua ista nimis erat, et quæ in ipso mundi principio contígerat ægritúdo; nec arte hóminis aut beneficio póterat hæc curari, quæ Dei fuerat indignatióne contracta. Tetigerat vétita, inconcessa præsumpserat, cum se ad árborem sciéndi bonum malumque porrexerat: Auctore indigebat, non qui malagma impóneret, sed qui posset illatam relaxare senténtiam, et ignoscéndo resólvere quod religáverat indignando.\par
\switchcolumn
\EN
\noindent This man is a figure of all men. His healing is a type of their healing, and his soundness is a pledge of that soundness for which all have looked so long. The hand of man hath withered through the deadness of faith rather than through the drying up of the sinews, and by the fault of the conscience rather than by the weakness of the flesh. The withering up of man's hand hath been of old time, and a sickness which smote him at the very beginning of the world, and no art or benefit of man could heal that which had been blasted by the wrath of God. That hand had touched the forbidden thing, it had sought that which was unlawful when it had been stretched out to the tree of the knowledge of good and evil. It had need of Him who had made it, not to lay a plaster upon it, but to cancel the sentence which He had uttered, and to loosen by pardon that which He had bound by judgment.\par
\switchcolumn*

\LA
\begin{Resp}\Rsign Iste est qui ante Deum magnas virtútes operátus est, et de omni corde suo laudávit Dóminum: \Star{} Ipse intercédat pro peccátis ómnium populórum.\end{Resp}
\begin{Resp}\Vsign Ecce homo sine queréla, verus Dei cultor, ábstinens se ab omni ópere malo, et pérmanens in innocéntia sua.\end{Resp}
\begin{Resp}\Rsign Ipse intercédat pro peccátis ómnium populórum.\end{Resp}
\switchcolumn
\EN
\begin{Resp}\Rsign This is he which wrought great wonders before God, and praised the Lord with all his heart. \Star{} May he pray for all people, that their sins may be forgiven unto them.\end{Resp}
\begin{Resp}\Vsign Behold a man without blame, a worshipper of God in truth, keeping himself clean from every evil work, and abiding still in his innocency.\end{Resp}
\begin{Resp}\Rsign May he pray for all people, that their sins may be forgiven unto them.\end{Resp}
\switchcolumn*

\LA
\LessonHead{Lectio VIII}
\switchcolumn
\EN
\LessonHead{Lesson VIII}
\switchcolumn*

\LA
\noindent In hoc hómine nostræ tantum géritur umbra sanitátis; perfécta autem salus nobis reservátur in Christo: quia tunc aríditas nostræ manus miseranda dissolvitur, cum cruore perfunditur Dominicæ passiónis, cum in illo vitali ligno crucis exténditur, cum carpit fructuosam de dolóre virtútem, cum totam árborem salútis amplectitur, cum clavis Dómini corpus affígitur, quo numquam ad árborem concupiscéntiæ et áridæ redeat voluptátis. Et ait hómini habénti manum áridam: Surge in médium, professor debilitátis própriæ, supernæ pietátis exactor, testis divinæ virtútis, Judáicæ incredulitátis assertor: surge in médium; ut quos non compungit virtus tanta signórum, quos non ópera tantæ salútis inclínant, vel debilitátis tantæ miserátio constringat et mítiget.\par
\switchcolumn
\EN
\noindent This man's healing is a type of the healing of all men, our perfect health is to be found in Christ, then shall our miserable hand be withered no more when there droppeth thereon the Blood of the Suffering Lord, when it is stretched forth to the Tree of Life, which is the Cross. When it gathereth the mighty fruit of His suffering, when it layeth hold upon the Tree of Salvation, when the body is so nailed thereto with the nails of the Lord that it can never return again to the tree of lust and barren enjoyment. And He said to the man which had the withered hand, Rise up, and stand forth in the midst. Rise up and stand forth in the midst, O Thou that dost confess thine own weakness, thou that dost call for pity from on high, thou that canst witness to the power of God; rise up and stand forth in the midst, thou that tellest of the unbelief of the Jews; the power of so many signs hath not pierced them, so many works of healing hath not beset them; let the pity shown to such misery constrain them and soften them.\par
\switchcolumn*

\LA
\begin{Resp}\Rsign In médio Ecclésiæ apéruit os ejus, \Star{} Et implévit eum Dóminus spíritu sapiéntiæ et intelléctus.\end{Resp}
\begin{Resp}\Vsign Jucunditátem et exsultatiónem thesaurizávit super eum.\end{Resp}
\begin{Resp}\Rsign Et implévit eum Dóminus spíritu sapiéntiæ et intelléctus.\end{Resp}
\begin{Resp}\Vsign Glória Patri, et Fílio, \Star{} et Spirítui Sancto.\end{Resp}
\begin{Resp}\Rsign Et implévit eum Dóminus spíritu sapiéntiæ et intelléctus.\end{Resp}
\switchcolumn
\EN
\begin{Resp}\Rsign In the midst of the congregation did the Lord open his mouth. \Star{} And filled him with the spirit of wisdom and understanding.\end{Resp}
\begin{Resp}\Vsign He made him rich with joy and gladness.\end{Resp}
\begin{Resp}\Rsign And filled him with the spirit of wisdom and understanding.\end{Resp}
\begin{Resp}\Vsign Glory be to the Father, and to the Son, \Star{} and to the Holy Ghost.\end{Resp}
\begin{Resp}\Rsign And filled him with the spirit of wisdom and understanding.\end{Resp}
\switchcolumn*

\end{paracol}

\par\needspace{10\baselineskip}\addvspace{1.2em}{\centering\small\textsc{Die 28 Martii} · \textsc{28 March}\par}
\Office{S. Joannis a Capistrano Confessoris}{St. John of Capistrano, Confessor}{Semiduplex}
\addcontentsline{toc}{subsection}{Die 28 Martii · S. Joannis a Capistrano Confessoris\enspace\textit{St. John of Capistrano, Confessor}}
\begin{paracol}{2}
\LA
\RefLine{Lectiones I–III: ex Commúni Confessoris non pontificis (C5).}
\switchcolumn
\EN
\RefLine{Lessons I–III: from the Common of a Confessor not a Bishop (C5).}
\switchcolumn*

\LA
\RefLine{Lectio IX: de Scriptura occurrente.}
\switchcolumn
\EN
\RefLine{Lesson IX: from the Scripture of the day.}
\switchcolumn*

\LA
\LessonHead{Lectio IV}
\switchcolumn
\EN
\LessonHead{Lesson IV}
\switchcolumn*

\LA
\noindent Joannes, Capistrani in Pelignis ortus et Perusium studiorum causa missus, in christianis et liberalibus disciplinis adeo profecit, ut ob egregiam juris scientiam aliquot civitatibus a Neapolis rege Ladislao præfectus fuerit. Dum autem earum rempublicam sanctissime gerens perturbatis rebus tranquillitatem revocare studet capitur ipse et in vincula conjicitur; quibus mirabiliter ereptus, Francisci Assisiensis regulam inter fratres Minores profitetur. Ad divinarum litterarum studium progressus præceptorem nactus est sanctum Bernardinum Senensem, cujus et virtutis exempla, in cultu potissimum sanctissimi nominis Jesu ac Deiparæ propagando, egregie est imitatus. Aquilanum episcopatum recusavit, et severiore disciplina atque scriptis quæ plurima edidit ad mores reformandos, maxime enituit.\par
\switchcolumn
\EN
\noindent This John was born at Capistrano, in the Abruzzi. He was educated at Perugia, and became so expert in letters, both sacred and profane, that on account of his eminent knowledge of law, Ladislaus, King of Naples, set him over several cities. He was seeking in righteousness to bring the affairs of these places out of trouble into peace, when he himself was kidnapped and put in chains. From this captivity he marvellously escaped, and then professed himself a Friar Minor under the rule of Francis of Assisi. Here he went forward in the study of divinity, and had as a teacher the holy Bernardine of Sienna, of whom he was one of the most marked followers, especially in spreading abroad the honour paid to the Most Holy Name of Jesus, and to the Mother of God. The bishopric of Aquila was offered to him, but he refused it. He was chiefly known by the hardship of his self-denial, and by the writings which he published in large numbers for the reform of manners.\par
\switchcolumn*

\LA
\begin{Resp}\Rsign Honéstum fecit illum Dóminus, et custodívit eum ab inimícis, et a seductóribus tutávit illum: \Star{} Et dedit illi claritátem ætérnam.\end{Resp}
\begin{Resp}\Vsign Justum dedúxit Dóminus per vias rectas, et osténdit illi regnum Dei.\end{Resp}
\begin{Resp}\Rsign Et dedit illi claritátem ætérnam.\end{Resp}
\switchcolumn
\EN
\begin{Resp}\Rsign The Lord made him honourable, and defended him from his enemies, and kept him safe from those that lay in wait for him; \Star{} And gave him perpetual glory.\end{Resp}
\begin{Resp}\Vsign He went down with him into the pit, and left him not in bonds.\end{Resp}
\begin{Resp}\Rsign And gave him perpetual glory.\end{Resp}
\switchcolumn*

\LA
\LessonHead{Lectio V}
\switchcolumn
\EN
\LessonHead{Lesson V}
\switchcolumn*

\LA
\noindent Prædicationi verbi Dei sedulo incumbens, Italiam fere universam lustravit; quo in munere et virtute sermonis et miraculorum frequentia innumeras prope animas in viam salutis reduxit. Eum Martinus quintus ad exstinguendam Fraticellorum sectam inquisitorem instituit. A Nicolao quinto contra Judæos et Saracenos generalis inquisitor in Italia constitutus, plurimos ad Christi fidem convertit. In Oriente multa optime constituit, et in concilio Florentino, ubi veluti sol quidam fulsit, Armenos Ecclesiæ catholicæ restituit. Idem Pontifex, postulante Friderico tertio imperatore, illum apostolicæ Sedis nuntium in Germaniam legavit, ut hæreticos ad catholicam fidem et principum animos ad concordiam revocaret. In Germania aliisque provinciis Dei gloriam sexennali ministerio mirifice auxit, Hussitis, Adamitis, Thaboritis, Hebræisque innumeris doctrinæ veritate ac miraculorum luce ad Ecclesiæ sinum traductis.\par
\switchcolumn
\EN
\noindent He devoted himself without ceasing to the preaching of the Word of God, in the which work he travelled throughout nearly all Italy, and by the power of eloquence and of miracles not a few, he recalled souls almost countless into the path of salvation. Martin V. appointed him Inquisitor to stamp out the sect of the Fraticelli. Nicolas V. appointed him InquisitorGeneral in Italy against Judaism and Mohammadanism, and he brought many such misbelievers to believe in Christ. He did much good work in the affairs of the Eastern Church, and at the Council of Florence, wherein he shone like a sun, he brought back the Armenians to the Catholic church. The same Pope Nicolas V., at the request of the Emperor Frederick III., sent him into Germany as Nuncio of the Apostolic See, in order that he might bring back the heretics to the Catholic faith and the minds of the princes to peace and agreement. He did a wonderful work for God's glory during the six years that he laboured in Germany and other countries, and by his teaching of the truth and the striking evidence of his miracles brought back to the bosom of the Church almost countless numbers of Hussites, Adamites, Taborites, and Jews.\par
\switchcolumn*

\LA
\begin{Resp}\Rsign Amávit eum Dóminus, et ornávit eum: stolam glóriæ índuit eum, \Star{} Et ad portas paradísi coronávit eum.\end{Resp}
\begin{Resp}\Vsign Índuit eum Dóminus lorícam fídei, et ornávit eum.\end{Resp}
\begin{Resp}\Rsign Et ad portas paradísi coronávit eum.\end{Resp}
\switchcolumn
\EN
\begin{Resp}\Rsign The Lord loved him and beautified him; He clothed him with a robe of glory, \Star{} And crowned him at the gates of Paradise.\end{Resp}
\begin{Resp}\Vsign The Lord hath put on him the breast-plate of faith, and hath adorned him.\end{Resp}
\begin{Resp}\Rsign And crowned him at the gates of Paradise.\end{Resp}
\switchcolumn*

\LA
\LessonHead{Lectio VI}
\switchcolumn
\EN
\LessonHead{Lesson VI}
\switchcolumn*

\LA
\noindent Cum Callistus tertius, ipso potissimum deprecante, cruce signatos mittere decrevisset, Joannes per Pannoniam aliasque provincias volitavit, qua verbo, qua litteris principum animos ita ad bellum accendit, ut brevi millia Christianorum septuaginta conscripta sint. Ejus consilio et virtute potissimum Taurunensis victoria relata est, centum ac viginti Turcarum millibus partim cæsis, partim fugatis. Cujus victoriæ cum Romam nuntius venisset octavo Idus Augusti, idem Callistus ejus diei memoriæ solemnia Transfigurationis Christi Domini perpetuo consecravit. Lethali morbo ægrotum et Villacum delatum viri principes plures visitarunt; quos ipse ad tuendam religionem hortatus, animam Deo sancte reddidit anno salutis millesimo quadringentesimo quinquagesimo sexto. Ejus gloriam post mortem Deus multis miraculis confirmavit: quibus rite probatis, Alexander octavus anno millesimo sexcentesimo nonagesimo Joánnem in Sanctorum numerum retulit; ejusque Officium ac Missam Leo decimus tertius, altero ab ejus canonizatione sæculo, ad universam extendit Ecclesiam.\par
\switchcolumn
\EN
\noindent It was mainly at the entreaty of John that Calistus III. proclaimed a Crusade, and John hastened about through Pannonia and other provinces, where by his words and his letters he so roused the minds of princes to that holy war, that in a short while seventy thousand Christian soldiers were enrolled. It was mainly through his advice and by his power that victory was gained at Belgrade, when one hundred and twenty thousand Turks were either slain or put to flight. The news of this victory reached Rome upon the sixth day of August, and Pope Calistus thereupon consecrated that day for ever to the solemn commemoration of the transfiguration of the Lord Christ. As John lay sick unto death at Illak, many princes came to see him, and he exhorted them to protect religion. He gave up his soul in holiness to God, (upon the 23rd day of October,) in the year of salvation 1456. God confirmed his glory by many miracles after his death, and when these had been duly proved Alexander VIII. enrolled his name with those of the saints in the year 1690, and two hundred years after his canonization Leo XIII. extended his Office and Mass to the whole Church.\par
\switchcolumn*

\LA
\begin{Resp}\Rsign Iste homo perfécit ómnia quæ locútus est ei Deus, et dixit ad eum: Ingrédere in réquiem meam: \Star{} Quia te vidi justum coram me ex ómnibus géntibus.\end{Resp}
\begin{Resp}\Vsign Iste est, qui contémpsit vitam mundi, et pervénit ad cæléstia regna.\end{Resp}
\begin{Resp}\Rsign Quia te vidi justum coram me ex ómnibus géntibus.\end{Resp}
\begin{Resp}\Vsign Glória Patri, et Fílio, \Star{} et Spirítui Sancto.\end{Resp}
\begin{Resp}\Rsign Quia te vidi justum coram me ex ómnibus géntibus.\end{Resp}
\switchcolumn
\EN
\begin{Resp}\Rsign This is he which did according unto all that God commanded him; and God said unto him: Enter thou into My rest, \Star{} For thee have I seen righteous before Me among all people.\end{Resp}
\begin{Resp}\Vsign This is he which loved not his life in this world, and is come unto an everlasting kingdom.\end{Resp}
\begin{Resp}\Rsign For thee have I seen righteous before Me among all people.\end{Resp}
\begin{Resp}\Vsign Glory be to the Father, and to the Son, \Star{} and to the Holy Ghost.\end{Resp}
\begin{Resp}\Rsign For thee have I seen righteous before Me among all people.\end{Resp}
\switchcolumn*

\LA
\LessonHead{Lectio VII}
\switchcolumn
\EN
\LessonHead{Lesson VII}
\switchcolumn*

\LA
\LessonTitle{Léctio sancti Evangélii secúndum Lucam}
\switchcolumn
\EN
\LessonTitle{From the Holy Gospel according to Luke}
\switchcolumn*

\LA
\Cite{Luc 9:1-6}
\switchcolumn
\EN
\Cite{Luke 9:1-6}
\switchcolumn*

\LA
\noindent In illo témpore: Convocatis Jesus duodecim Apostolis, dedit illis virtutem et potestatem super omnia dæmonia, et ut languores curarent. Et réliqua.\par
\switchcolumn
\EN
\noindent At that time: Jesus called the twelve Apostles together and gave power and authority over all devils and to cure diseases. And so on.\par
\switchcolumn*

\LA
\LessonTitle{Homilia sancti Bonaventuræ Epíscopi}
\switchcolumn
\EN
\LessonTitle{Homily by St. Bonaventure, (Cardinal) Bishop (of Albano.)}
\switchcolumn*

\LA
\Source{Expositio in cap. 9 Lucæ}
\switchcolumn
\EN
\Source{On Luke ix.}
\switchcolumn*

\LA
\noindent Apostoli ideo nominati sunt, ut eorum commendaretur auctoritas. Apostolus enim missus interpretatur: missi autem fuerunt ad prædicandum, secundum illud: Non misit me Christus baptizare, sed evangelizare. Fuerunt ad prædicandum missi, non rem parvam, sed magnam scilicet regnum Dei, per quod potest intelligi doctrina veritatis, juxta illud: Auferetur a vobis regnum Dei, et dabitur genti facienti fructus ejus. Potest étiam dici gratia Spiritus Sancti, secundum illud: Non est regnum Dei esca et potus, sed justitia et pax, et gaudium in Spiritu Sancto; et infra: Ecce regnum Dei intra vos est. Potest étiam dici gloria æterna, juxta illud: Amen dico vobis, nisi quis renatus fuerit ex aqua et Spiritu Sancto, non potest introire in regnum Dei.\par
\switchcolumn
\EN
\noindent Apostles are so called as a mark of their authority, for this word Apostle signifieth sent, and they were sent out to preach, as it is written: (1 Cor. i. 17,) “Christ sent me not to baptize but to preach the gospel.” They were sent to preach not any small thing but a very great thing, even the Kingdom of God, whereby we may understand the teaching of the truth, as it is said, (Matth. xxi. 43,) The Kingdom of God should be taken from you and given to a nation bringing forth the fruits thereof. The Kingdom of God may also be understood to signify the grace of the Holy Ghost, as it is written: (Rom. xiv. 17,) “The Kingdom of God is not meat and drink, but righteousness, and peace, and joy in the Holy Ghost”; as also it was said: (Luke xvii. 21,) “The Kingdom of God is within you.” The Kingdom of God may also be understood to signify eternal glory, as it is said: (John iii. 5,) “Amen, I say unto thee, except a man be born again of water and of the Holy Ghost he cannot enter into the Kingdom of God.”\par
\switchcolumn*

\LA
\begin{Resp}\Rsign Iste est qui ante Deum magnas virtútes operátus est, et de omni corde suo laudávit Dóminum: \Star{} Ipse intercédat pro peccátis ómnium populórum.\end{Resp}
\begin{Resp}\Vsign Ecce homo sine queréla, verus Dei cultor, ábstinens se ab omni ópere malo, et pérmanens in innocéntia sua.\end{Resp}
\begin{Resp}\Rsign Ipse intercédat pro peccátis ómnium populórum.\end{Resp}
\switchcolumn
\EN
\begin{Resp}\Rsign This is he which wrought great wonders before God, and praised the Lord with all his heart. \Star{} May he pray for all people, that their sins may be forgiven unto them.\end{Resp}
\begin{Resp}\Vsign Behold a man without blame, a worshipper of God in truth, keeping himself clean from every evil work, and abiding still in his innocency.\end{Resp}
\begin{Resp}\Rsign May he pray for all people, that their sins may be forgiven unto them.\end{Resp}
\switchcolumn*

\LA
\LessonHead{Lectio VIII}
\switchcolumn
\EN
\LessonHead{Lesson VIII}
\switchcolumn*

\LA
\noindent Omnibus his modis Apostoli sunt missi prædicare regnum Dei, scilicet veram doctrinam, divinam gratiam, et ætérnam gloriam. Et quia propter auctoritatem prædicationis concesserat potestatem curationis, ideo subdit: Et sanare infirmos; scilicet misit ad confirmationem veritatis prædicatæ secundum illud: Illi autem profecti, prædicaverunt ubique, Domino cooperante et sermonem confirmante sequentibus signis. Unde signum missionis spiritualis ad prædicandum est sanatio audientium a morbis vitiorum.\par
\switchcolumn
\EN
\noindent The Apostles were sent to preach the Kingdom of God in all these three senses, that is to say, as the true teaching, as the grace of God, and as eternal glory. In order to invest their teaching with authority He gave them the power to cure diseases, whence where it is written And He sent them to preach the Kingdom of God it is also said And to heal the sick. This power He gave in order to confirm the truth of their preaching, as it is written: (Mark xvi. 20,) “And they went forth, and preached everywhere, the Lord working with them, and confirming the Word with signs following.” The sign that a preacher is indeed sent forth by the Spirit of God is that they that hear him should be cured of the disease of sin.\par
\switchcolumn*

\LA
\begin{Resp}\Rsign Sint lumbi vestri præcíncti, et lucérnæ ardéntes in mánibus vestris: \Star{} Et vos símiles homínibus exspectántibus dóminum suum, quando revertátur a núptiis.\end{Resp}
\begin{Resp}\Vsign Vigiláte ergo, quia nescítis qua hora Dóminus vester ventúrus sit.\end{Resp}
\begin{Resp}\Rsign Et vos símiles homínibus exspectántibus dóminum suum, quando revertátur a núptiis.\end{Resp}
\begin{Resp}\Vsign Glória Patri, et Fílio, \Star{} et Spirítui Sancto.\end{Resp}
\begin{Resp}\Rsign Et vos símiles homínibus exspectántibus dóminum suum, quando revertátur a núptiis.\end{Resp}
\switchcolumn
\EN
\begin{Resp}\Rsign Let your loins be girded about, and your lights burning; \Star{} And ye yourselves like unto men that wait for their lord, when he will return from the wedding.\end{Resp}
\begin{Resp}\Vsign Watch therefore, for ye know not what hour your Lord doth come.\end{Resp}
\begin{Resp}\Rsign And ye yourselves like unto men that wait for their lord, when he will return from the wedding.\end{Resp}
\begin{Resp}\Vsign Glory be to the Father, and to the Son, \Star{} and to the Holy Ghost.\end{Resp}
\begin{Resp}\Rsign And ye yourselves like unto men that wait for their lord, when he will return from the wedding.\end{Resp}
\switchcolumn*

\end{paracol}

\SeasonTitle{Mensis Aprilis}{April}

\addcontentsline{toc}{section}{Mensis Aprilis\enspace\textit{April}}

\par\needspace{10\baselineskip}\addvspace{1.2em}{\centering\small\textsc{Die 2 Aprilis} · \textsc{2 April}\par}
\Office{S. Francisci de Paula Confessoris}{St. Francis of Paula, Confessor}{Duplex}
\addcontentsline{toc}{subsection}{Die 2 Aprilis · S. Francisci de Paula Confessoris\enspace\textit{St. Francis of Paula, Confessor}}
\begin{paracol}{2}
\LA
\RefLine{Lectiones I–III: ex Commúni Confessoris non pontificis (C5).}
\switchcolumn
\EN
\RefLine{Lessons I–III: from the Common of a Confessor not a Bishop (C5).}
\switchcolumn*

\LA
\RefLine{Lectio IX: de Scriptura occurrente.}
\switchcolumn
\EN
\RefLine{Lesson IX: from the Scripture of the day.}
\switchcolumn*

\LA
\LessonHead{Lectio IV}
\switchcolumn
\EN
\LessonHead{Lesson IV}
\switchcolumn*

\LA
\noindent Francíscus Paulæ, quod est Calábriæ óppidum, loco húmili natus est; quem paréntes, cum diu prole caruíssent, voto facto, beáti Francísci précibus suscepérunt. Is adoléscens divíno ardóre succénsus in erémum secéssit, ubi annis sex, victu ásperam, sed meditatiónibus cæléstibus suávem vitam duxit. Sed cum virtútum ejus fama lóngius manáret, multíque ad eum pietátis stúdio concúrrerent, fratérnæ caritátis causa e solitúdine egréssus, ecclésiam prope Paulam ædificávit, ibique prima sui órdinis fundaménta jecit.\par
\switchcolumn
\EN
\noindent This Francis was born of humble parents at Paola, a town in Calabria, (about the year of our Lord 1416.) His parents, who had long been childless, obtained him, after making a vow, by the prayers of blessed Francis. While he was yet a lad, the love of God moved him to withdraw into a desert place, where he lived for six years, hardly as to the body, but sumptuously in meditation on things heavenly. Nevertheless, when the fame of his holy life was noised abroad, and many betook themselves to him, that they might learn godliness, he was drawn out of the desert by love to his neighbour, and built a church near Paola, wherein he laid the first foundations of his Order.\par
\switchcolumn*

\LA
\begin{Resp}\Rsign Honéstum fecit illum Dóminus, et custodívit eum ab inimícis, et a seductóribus tutávit illum: \Star{} Et dedit illi claritátem ætérnam.\end{Resp}
\begin{Resp}\Vsign Justum dedúxit Dóminus per vias rectas, et osténdit illi regnum Dei.\end{Resp}
\begin{Resp}\Rsign Et dedit illi claritátem ætérnam.\end{Resp}
\switchcolumn
\EN
\begin{Resp}\Rsign The Lord made him honourable, and defended him from his enemies, and kept him safe from those that lay in wait for him; \Star{} And gave him perpetual glory.\end{Resp}
\begin{Resp}\Vsign He went down with him into the pit, and left him not in bonds.\end{Resp}
\begin{Resp}\Rsign And gave him perpetual glory.\end{Resp}
\switchcolumn*

\LA
\LessonHead{Lectio V}
\switchcolumn
\EN
\LessonHead{Lesson V}
\switchcolumn*

\LA
\noindent Erat in eo mirífica loquéndi grátia: perpétuam virginitátem servávit: humilitátem sic cóluit, ut se ómnium mínimum díceret, suósque alúmnos Mínimos appellári volúerit. Rudi amíctu, nudis pédibus incédens, humi cubábat. Cibi abstinéntia fuit admirábili: semel in die post solis occásum reficiebátur, et ad panem, et aquæ potum, vix áliquid ejusmodi obsónii adhibébat, quo vesci in Quadragésima licet: quam consuetúdinem, ut fratres sui toto anni témpore retinérent, quarto eos voto adstrínxit.\par
\switchcolumn
\EN
\noindent In his words there was a wonderful charm he kept his virginity always inviolate he was so great a lover of lowliness that he used to call himself the last of all, and would that his disciples should be called the Minimi, which is, being interpreted, the Least of the brethren. His raiment was coarse; he went always bare-footed; and he slept on the ground. The extreme smallness of the amount of food which he took was extraordinary. He ate only once a day, and that after sunset. Then he took only bread and water, with scarcely any of such condiment as is allowed in Lent. He bound his disciples by a fourth vow, added to those of Poverty, Chastity, and Obedience, to observe the same rule of eating as himself.\par
\switchcolumn*

\LA
\begin{Resp}\Rsign Amávit eum Dóminus, et ornávit eum: stolam glóriæ índuit eum, \Star{} Et ad portas paradísi coronávit eum.\end{Resp}
\begin{Resp}\Vsign Índuit eum Dóminus lorícam fídei, et ornávit eum.\end{Resp}
\begin{Resp}\Rsign Et ad portas paradísi coronávit eum.\end{Resp}
\switchcolumn
\EN
\begin{Resp}\Rsign The Lord loved him and beautified him; He clothed him with a robe of glory, \Star{} And crowned him at the gates of Paradise.\end{Resp}
\begin{Resp}\Vsign The Lord hath put on him the breast-plate of faith, and hath adorned him.\end{Resp}
\begin{Resp}\Rsign And crowned him at the gates of Paradise.\end{Resp}
\switchcolumn*

\LA
\LessonHead{Lectio VI}
\switchcolumn
\EN
\LessonHead{Lesson VI}
\switchcolumn*

\LA
\noindent Multis miráculis servi sui sanctitátem Deus testári vóluit, quorum illud in primis célebre, quod a nautis rejéctus, Sicíliæ fretum strato super flúctibus pállio cum sócio transmísit. Multa étiam futúra prophético spíritu prædíxit. A Ludovíco undécimo Francórum rege expetítus, magnóque in honóre est hábitus. Dénique annum primum et nonagésimum agens, Turónis migrávit ad Dominum, anno salutis millésimo quingentésimo séptimo: cujus corpus dies úndecim insepúltum, ita incorrúptum permánsit, ut suávem étiam odórem effláret. Eum Leo Papa décimus in Sanctórum númerum rétulit.\par
\switchcolumn
\EN
\noindent It was the will of God to make the holiness of His servant manifest by many miracles. The most notorious of these is that on one occasion when some seamen refused to take him over the Straits of Messina, he spread his cloak upon the sea, and crossed over on it with his companion. In the spirit of prophecy he foretold many things to come. Louis XI, King of France, held him in great worship, and bade him to his court. At last, at Tours, in the ninety-first year of his age, and the 1507th of our salvation, he departed hence to be ever buried for eleven days after his death, but it not only showed no signs of corruption but even gave forth a sweet savour. Pope Leo X caused him to be numbered among the Saints.\par
\switchcolumn*

\LA
\begin{Resp}\Rsign Iste homo perfécit ómnia quæ locútus est ei Deus, et dixit ad eum: Ingrédere in réquiem meam: \Star{} Quia te vidi justum coram me ex ómnibus géntibus.\end{Resp}
\begin{Resp}\Vsign Iste est, qui contémpsit vitam mundi, et pervénit ad cæléstia regna.\end{Resp}
\begin{Resp}\Rsign Quia te vidi justum coram me ex ómnibus géntibus.\end{Resp}
\begin{Resp}\Vsign Glória Patri, et Fílio, \Star{} et Spirítui Sancto.\end{Resp}
\begin{Resp}\Rsign Quia te vidi justum coram me ex ómnibus géntibus.\end{Resp}
\switchcolumn
\EN
\begin{Resp}\Rsign This is he which did according unto all that God commanded him; and God said unto him: Enter thou into My rest, \Star{} For thee have I seen righteous before Me among all people.\end{Resp}
\begin{Resp}\Vsign This is he which loved not his life in this world, and is come unto an everlasting kingdom.\end{Resp}
\begin{Resp}\Rsign For thee have I seen righteous before Me among all people.\end{Resp}
\begin{Resp}\Vsign Glory be to the Father, and to the Son, \Star{} and to the Holy Ghost.\end{Resp}
\begin{Resp}\Rsign For thee have I seen righteous before Me among all people.\end{Resp}
\switchcolumn*

\LA
\LessonHead{Lectio VII}
\switchcolumn
\EN
\LessonHead{Lesson VII}
\switchcolumn*

\LA
\LessonTitle{Léctio sancti Evangélii secúndum Lucam}
\switchcolumn
\EN
\LessonTitle{From the Holy Gospel according to Luke}
\switchcolumn*

\LA
\Cite{Luc 12:32-34}
\switchcolumn
\EN
\Cite{Luke 12:32-35}
\switchcolumn*

\LA
\noindent In illo témpore: Dixit Jesus discípulis suis: Nolíte timére pusíllus grex, quia complácuit Patri vestro dare vobis regnum. Et réliqua.\par
\switchcolumn
\EN
\noindent At that time, Jesus said unto His disciples: Fear not, little flock, for it is your Father's good pleasure to give you the kingdom. And so on.\par
\switchcolumn*

\LA
\LessonTitle{Homilía sancti Bedæ Venerábilis Presbýteri}
\switchcolumn
\EN
\LessonTitle{Homily by the Venerable Bede, Priest at Jarrow, and Doctor of the Church.}
\switchcolumn*

\LA
\Source{Lib. 4, Cap. 54, in Luc. 12}
\switchcolumn
\EN
\Source{Bk. iv. Ch. 54 on Luke xii.}
\switchcolumn*

\LA
\noindent Pusíllum gregem electórum, vel ob comparatiónem majóris númeri reprobórum, vel pótius ob humilitátis devotiónem nóminat: quia vidélicet Ecclésiam suam quantálibet numerositáte jam dilatátam, tamen usque ad finem mundi humilitáte vult créscere, et ad promíssum regnum humilitáte perveníre. Ideóque ejus labóres blande consolátus, quam regnum Dei tantum quærére prǽcipit, eídem regnum a Patre dandum complácita benignitáte promíttit.\par
\switchcolumn
\EN
\noindent The elect are called a little flock, perchance because the reprobate are far more in number than they, but, more probably, because they love to be lowly, since it is God's will that however much His Church should grow in numbers, she should grow with lowliness even unto the end of the world, and should enter lowly into that kingdom which is hers by His promise. That kingdom He promiseth to her here, when He biddeth her to seek only the kingdom of God, and, to comfort her in her travail, He doth so sweetly and so graciously say that her Father will give it to her.\par
\switchcolumn*

\LA
\begin{Resp}\Rsign Iste est qui ante Deum magnas virtútes operátus est, et de omni corde suo laudávit Dóminum: \Star{} Ipse intercédat pro peccátis ómnium populórum.\end{Resp}
\begin{Resp}\Vsign Ecce homo sine queréla, verus Dei cultor, ábstinens se ab omni ópere malo, et pérmanens in innocéntia sua.\end{Resp}
\begin{Resp}\Rsign Ipse intercédat pro peccátis ómnium populórum.\end{Resp}
\switchcolumn
\EN
\begin{Resp}\Rsign This is he which wrought great wonders before God, and praised the Lord with all his heart. \Star{} May he pray for all people, that their sins may be forgiven unto them.\end{Resp}
\begin{Resp}\Vsign Behold a man without blame, a worshipper of God in truth, keeping himself clean from every evil work, and abiding still in his innocency.\end{Resp}
\begin{Resp}\Rsign May he pray for all people, that their sins may be forgiven unto them.\end{Resp}
\switchcolumn*

\LA
\LessonHead{Lectio VIII}
\switchcolumn
\EN
\LessonHead{Lesson VIII}
\switchcolumn*

\LA
\noindent Véndite quæ possidétis, et date eleemósynam. Nolíte, inquit, timére, ne propter regnum Dei militántibus, hujus vitæ necessária desint: quin étiam posséssa propter eleemósynam véndite. Quod tunc digne fit, quando quis semel pro Dómino suis ómnibus spretis, nihilóminus post hæc labóre mánuum, unde et victum transígere, et eleemósynam dare queat, operátur. Unde gloriátur Apóstolus, dicens: Argéntum, et aurum, aut vestem nullíus concupívi: ipsi scitis, quóniam ad ea quæ mihi opus erant, et his qui mecum sunt, ministravérunt manus istæ. Ómnia osténdi vobis, quóniam sic laborántes opórtet suscípere infírmos.\par
\switchcolumn
\EN
\noindent Sell what ye have and give alms. Fear not, He saith, lest, while ye fight for the kingdom of God, ye should lack such things as are needful for this life, nay rather, sell even that which ye have, and give alms. This doth, whosoever for the Lord's sake leaveth all that he hath, and then worketh with his hands, that so he may have to eat, and withal to give alms. In this doth the Apostle boast himself, saying I have coveted no man's silver, or gold, or apparel, as ye yourselves know for these hands have ministered unto my necessities, and to them that were with me. I have showed you all things, how that so labouring ye ought to support the weak. (Acts xx. 33, 34, 35.)\par
\switchcolumn*

\LA
\begin{Resp}\Rsign Sint lumbi vestri præcíncti, et lucérnæ ardéntes in mánibus vestris: \Star{} Et vos símiles homínibus exspectántibus dóminum suum, quando revertátur a núptiis.\end{Resp}
\begin{Resp}\Vsign Vigiláte ergo, quia nescítis qua hora Dóminus vester ventúrus sit.\end{Resp}
\begin{Resp}\Rsign Et vos símiles homínibus exspectántibus dóminum suum, quando revertátur a núptiis.\end{Resp}
\begin{Resp}\Vsign Glória Patri, et Fílio, \Star{} et Spirítui Sancto.\end{Resp}
\begin{Resp}\Rsign Et vos símiles homínibus exspectántibus dóminum suum, quando revertátur a núptiis.\end{Resp}
\switchcolumn
\EN
\begin{Resp}\Rsign Let your loins be girded about, and your lights burning; \Star{} And ye yourselves like unto men that wait for their lord, when he will return from the wedding.\end{Resp}
\begin{Resp}\Vsign Watch therefore, for ye know not what hour your Lord doth come.\end{Resp}
\begin{Resp}\Rsign And ye yourselves like unto men that wait for their lord, when he will return from the wedding.\end{Resp}
\begin{Resp}\Vsign Glory be to the Father, and to the Son, \Star{} and to the Holy Ghost.\end{Resp}
\begin{Resp}\Rsign And ye yourselves like unto men that wait for their lord, when he will return from the wedding.\end{Resp}
\switchcolumn*

\LA
\LessonHead{Lectio IX}
\switchcolumn
\EN
\LessonHead{Lesson IX}
\switchcolumn*

\LA
\noindent Fácite vobis sácculos qui non veteráscunt: eleemósynas vidélicet operándo quarum merces in ætérnum máneat. Ubi non hoc præcéptum esse putándum est, ut nil pecúniæ reservétur a sanctis, vel suis scílicet, vel páuperum úsibus suggeréndæ: cum et ipse Dóminus, cui ministrábant Ángeli, tamen ad informándam Ecclésiam suam lóculos habuísse legátur, et a fidélibus obláta consérvans, et suórum necessitátibus aliísque indigéntibus tríbuens; sed ne Deo propter ista serviátur, et ob inópiæ timórem justítia deserátur.\par
\switchcolumn
\EN
\noindent Provide yourselves bags which wax not old that is to say, by almsgiving, the reward thereof remaineth for ever. Nevertheless, we must not think here that this commandment forbiddeth the Saints to keep money for their own use, and for helping of the poor. The Lord Himself, to Whom Angels ministered, had a bag, and kept therein that which the faithful people gave unto Him (John xii. 6,) to relieve therewith the need of His disciples, and other poor folk. But we are commanded not to serve God for gain, nor to work unrighteousness for fear of poverty.\par
\switchcolumn*

\LA
\TeDeum{Te Deum}
\switchcolumn
\EN
\TeDeum{Te Deum}
\switchcolumn*

\LA
\LessonHead{Lectio IX (pro commemoratione)}
\switchcolumn
\EN
\LessonHead{Lesson IX (when commemorated)}
\switchcolumn*

\LA
\LessonTitle{Commemoratio: S. Francisci de Paula Confessoris}
\switchcolumn
\EN
\LessonTitle{Commemoration of: St. Francis of Paula, Confessor}
\switchcolumn*

\LA
\noindent Francíscus Paulæ in Calábria natus est. Aduléscens, divíno ardóre succénsus, in erémum secéssit, ubi annis sex victu ásperam sed meditatiónibus cæléstibus suávem vitam duxit. Cum autem virtútum ejus fama longe manáret, multíque ad eum pietátis stúdio concúrrerent, caritátis causa e solitúdine egréssus, ecclésiam prope Paulam ædificávit, ibíque prima sui órdinis fundaménta jecit. Perpétuam virginitátem servávit: humilitátem sic cóluit, ut se ómnium mínimum díceret, suósque alúmnos Mínimos appellári volúerit. Rudi amíctu, nudis pédibus incédens, humi cubábat. Caritáti ita addíctus fuit, ut sui Ordinis tésseram esse jússerit. Multis miráculis servi sui sanctitátem Deus comprobávit, quorum illud in primis célebre, quod a nautis rejéctus, Sicíliæ fretum, strato super flúctibus pállio, cum sócio transmísit. Multa étiam prophético spíritu prædíxit. Turónis migrávit ad Dóminum, anno salútis millésimo quingentésimo séptimo, ætátis anno nonagésimo primo.\par
\switchcolumn
\EN
\noindent Francis of Paula, born in Calabria, even as a young man burned with love of God and withdrew to an hermitage where, for six years, he lived a life hard in its austerity but sweetened by meditating on heavenly things. When the fame of his virtues spread abroad, many came to him to be trained in the love of God; and, emerging from his solitude for the sake of charity, he built a church near Paola, and there laid the first foundations of his Order. Her preserved his virginity all his life; and so cultivated humility that he called himself the least of all men and wished his followers to be named Minims. Clad in poor garments, he went about barefoot and slept on the ground. He cherished charity so greatly that he made it the motto of his Order. God witnessed to the holiness of his servant by many miracles, one of which is particularly famous: when he was refused passage on a ship by the sailors, he spread out his cloak on the waves and so crossed the straits of Sicily with his companion. He also predicted many things in the spirit of prophecy. At Tours, he went to the Lord, in the year of salvation 1507, in the ninety-first year of his age.\par
\switchcolumn*

\end{paracol}

\par\needspace{10\baselineskip}\addvspace{1.2em}{\centering\small\textsc{Die 4 Aprilis} · \textsc{4 April}\par}
\Office{S. Isidori Episcopi Confessoris et Ecclesiæ Doctoris}{S. Isidore of Seville, Bishop, Confessor and Doctor of the Church}{Duplex}
\addcontentsline{toc}{subsection}{Die 4 Aprilis · S. Isidori Episcopi Confessoris et Ecclesiæ Doctoris\enspace\textit{S. Isidore of Seville, Bishop, Confessor and Doctor of the Church}}
\begin{paracol}{2}
\LA
\RefLine{Lectiones I–III: ex Commúni Doctoris Pontificis (C4a).}
\switchcolumn
\EN
\RefLine{Lessons I–III: from the Common of a Doctor Bishop (C4a).}
\switchcolumn*

\LA
\RefLine{Lectio IX: de Scriptura occurrente.}
\switchcolumn
\EN
\RefLine{Lesson IX: from the Scripture of the day.}
\switchcolumn*

\LA
\LessonHead{Lectio IV}
\switchcolumn
\EN
\LessonHead{Lesson IV}
\switchcolumn*

\LA
\noindent Isidorus, natione Hispanus, Doctor egregius, ex nova Carthagine Severiano patre provinciæ duce natus, a sanctis episcopis Leandro Hispalensi, et Fulgentio Carthaginiensi fratribus suis pie et liberaliter educatus, Latinis, Græcis, et Hebraicis litteris, divinisque et humanis legibus instructus, omni scientiarum, atque christianarum virtutum genere præstantissimus evasit. Adhuc adolescens hæresim Arianam, quæ gentem Gotharum Hispaniæ latissime dominantem jampridem invaserat, tanta constantia palam oppugnavit, ut parum abfuerit, quin ab hæreticis necaretur. Leandro vita functo ad Hispalensem cathedram invitus quidem, sed urgente in primis Reccaredo rege, magnoque étiam cleri populique consensu assumitur, ejtisque electionem sanctus Gregorius Magnus, nedum auctoritate apostolica confirmasse, sed et electum transmisso de more pallio decorasse, quin étiam suum et apostolicæ Sedis in universa Hispania Vicarium constituisse perhibetur.\par
\switchcolumn
\EN
\noindent Isidore, the admirable teacher, was a Spaniard by birth, being the son of Severian, governor of the Province of Carthagena. He was trained up in all godliness and learning by his holy brethren Leander, Archbishop of Seville, and Fulgentius, Bishop of Carthagena. He was well instructed in the Latin, Greek, and Hebrew letters, and he came from his masters a most eminent scholar in all human knowledge, and a pattern of all Christian graces. While yet he was very young, he attacked with such firmness the Arian heresy, which had of former times polluted the Gothic nation, who then were the chief rulers of Spain, that he was near being murdered by the heretics. After that Leander was departed this life, Isidore was chosen to the See of Seville, against his own will, but at the vehement instance of King Reccared, and with the strong assent of the clergy and people. Holy Gregory the Great not only confirmed his election by his own Apostolic authority, and caused him to be adorned, as is the custom, with a Pallium sent from the body of Blessed Peter, but is also stated to have appointed him Vicar of the Apostolic See for all Spain.\par
\switchcolumn*

\LA
\begin{Resp}\Rsign Invéni David servum meum, óleo sancto meo unxi eum: \Star{} Manus enim mea auxiliábitur ei.\end{Resp}
\begin{Resp}\Vsign Nihil profíciet inimícus in eo, et fílius iniquitátis non nocébit ei.\end{Resp}
\begin{Resp}\Rsign Manus enim mea auxiliábitur ei.\end{Resp}
\switchcolumn
\EN
\begin{Resp}\Rsign I have found David My servant, with My holy oil have I anointed him \Star{} For My hand shall help him.\end{Resp}
\begin{Resp}\Vsign The enemy shall prevail nothing against him, nor the son of wickedness afflict him.\end{Resp}
\begin{Resp}\Rsign For My hand shall help him.\end{Resp}
\switchcolumn*

\LA
\LessonHead{Lectio V}
\switchcolumn
\EN
\LessonHead{Lesson V}
\switchcolumn*

\LA
\noindent In episcopatu quantum fuerit constans, humilis, patiens, misericors, in christiana et ecclesiastica disciplina instauranda sollicitus, edque verbo et scriptis stabilienda indefessus, atque omni demum virtutum ornamento insignitus, nullius lingua enarrare sufficeret. Monastici quoque instituti per Hispaniam promotor et amplificator eximius plura construxit monasteria, collegia itidem ædificavit, ubi studiis sacris et lectionibus vacans, plurimos discipulos, qui ad eum confluebant, erudivit, quos inter sancti Ildefonsus Toletanus, et Brauho Cæsaraugustanus episcopi emicuerunt. Coacto Hispali concilio, Acephalorum hæresim Hispaniæ jam minitantem, acri et eloquenti disputatidne fregit atque contrivit. Tantam apud omnes sanctitatis et doctrinæ famam adeptus est, ut elapso vix ab ejus obitu sextodecimo anno, universa Toletana synodo duorum supra quinquaginta episcoporum plaudente, ipsoque étiam sancto Ildefonso suffragante, Doctor egregius, catholicæ Ecclesiæ novissimum decus, in sæculorum fine doctissimus, et cum reverentia nominandus appellari meruerit; eumque sanctus Braulio non modo Gregorio Magno comparaverit, sed et erudiendæ Hispaniæ loco Jacobi Apostoli cælitus datum esse censuerit.\par
\switchcolumn
\EN
\noindent When he was Archbishop no tongue can tell how leal he was, how lowly, and meek, and merciful, how careful to restore the laws of Christianity and the Church, and how unwearied in establishing the same by his word and writings, yea, how brightly he shone in all graces. He was a leading promoter and spreader of monastic institutions throughout Spain. He built many monasteries. He founded colleges in which, when his duty allowed him spare time for sacred study and reading, he taught the many disciples who betook themselves to him from all quarters. Among these, two of the most distinguished were the holy Bishops Ildephonsus of Toledo, and Braulio of Saragossa. He called the Council of Seville, wherein, in a most incisive and eloquent discourse, he shattered and crushed the heresy of the Acephali, by which Spain was then threatened. So great was his fame among all men for the holiness of his life and doctrine, that scarcely sixteen years after his death the whole Council of Toledo, by the acclamation of more than fifty Bishops, among whom was the holy Ildephonsus himself, declared him to be worthy to be called the excellent Teacher, the newest ornament of the Catholic Church, one whose learning would endure to the end of the world, and of worshipful memory. It was the opinion of the holy Braulio that he was not only fit to be compared to Gregory the Great, but also that he was a gift from God to Spain instead of the Apostle James.\par
\switchcolumn*

\LA
\begin{Resp}\Rsign Pósui adjutórium super poténtem, et exaltávi eléctum de plebe mea: \Star{} Manus enim mea auxiliábitur ei.\end{Resp}
\begin{Resp}\Vsign Invéni David servum meum, óleo sancto meo unxi eum.\end{Resp}
\begin{Resp}\Rsign Manus enim mea auxiliábitur ei.\end{Resp}
\switchcolumn
\EN
\begin{Resp}\Rsign I have laid help upon one that is mighty, and have exalted one chosen out of My people \Star{} For My hand shall help him.\end{Resp}
\begin{Resp}\Vsign I have found David My servant, with My holy oil have I anointed him.\end{Resp}
\begin{Resp}\Rsign For My hand shall help him.\end{Resp}
\switchcolumn*

\LA
\LessonHead{Lectio VI}
\switchcolumn
\EN
\LessonHead{Lesson VI}
\switchcolumn*

\LA
\noindent Scripsit Isidorus libros etymologiarum, et de ecclesiasticis officiis, aliosque quam plurimos christianæ, et ecclesiasticæ disciplinæ adeo utiles, ut sanctus Leo Papa quartus ad episcopos Britanniæ scribere non dubitaverit, sicut Hierónymi et Augustíni, ita Isidori dicta retinenda esse, ubi contigerit inusitatum negotium, quod per Canones minime definiri possit. Plures étiam ex ejusdem scriptis sententiæ inter canonicas Ecclesiæ leges relatæ conspiciuntur. Præfuit concilio Toletano quarto, omnium Hispaniæ celeberrimo. Denique cum ab Hispania Arianam hæresim eliminasset, morte sua, et regni vastatione a Saracenorum armis publice prænuntiata, postquam quadraginta circiter annos suam rexisset ecclesiam, Hispali migravit in cælum anno sexcentesimo trigesimo sexto. Ejus corpus inter Leandrum fratrem, et Florentinam sororem, ut ipse mandaverat, primo conditum, Ferdinandus primus Castellæ et Legionis rex, ab Eneto Saraceno Hispali dominante, magno pretio redemptum Legionem transtulit, et in ejus honorem templum ædificatum est, ubi miraculis clarus magna populi devotione colitur.\par
\switchcolumn
\EN
\noindent Isidore wrote Books of Etymologies and on Church Offices, and likewise many others, so useful in the administration of Christian and Church Law, that the holy Pope Leo IV felt no scruple in writing to the Bishops of Britain, that the sayings of Isidore were worthy to be kept like those of Jerome and Augustine, whenever there is to be done some strange work, wherein the rules of the Canon Law are not enough defined. Many sentences from his writings may also be discovered embedded in the Canon Law of the Church itself. He presided over the Fourth Council of Toledo, the most celebrated that ever met in Spain. Before his death he had purged Spain of the Arian heresy, and publicly foretold his own dissolution and the wasting of the kingdom by the Saracens which was to come. He passed away to heaven, at Seville, where he had ruled his Church for forty years, \textasciitilde{}(upon the 4th day of April,) in the year of our Lord 636. In accordance with his own commands, his body was first buried between his brother Leander and his sister Florentina, but Ferdinand I, King of Castille and Leon, bought it for a great price from Enet, the Saracen, who then ruled at Seville, carried it to Leon, and there built a Church in honour of him, wherein his said body lieth, illustrious through miracles, and reverenced with great worship by the people.\par
\switchcolumn*

\LA
\begin{Resp}\Rsign Iste est, qui ante Deum magnas virtútes operátus est, et omnis terra doctrína ejus repléta est: \Star{} Ipse intercédat pro peccátis ómnium populórum.\end{Resp}
\begin{Resp}\Vsign Iste est, qui contémpsit vitam mundi, et pervénit ad cæléstia regna.\end{Resp}
\begin{Resp}\Rsign Ipse intercédat pro peccátis ómnium populórum.\end{Resp}
\begin{Resp}\Vsign Glória Patri, et Fílio, \Star{} et Spirítui Sancto.\end{Resp}
\begin{Resp}\Rsign Ipse intercédat pro peccátis ómnium populórum.\end{Resp}
\switchcolumn
\EN
\begin{Resp}\Rsign This is he which wrought great wonders before God, and the whole earth is full of his teaching \Star{} May he pray for all people, that their sins may be forgiven unto them\end{Resp}
\begin{Resp}\Vsign This is he which loved not his life in this world, and hath attained unto the kingdom of heaven.\end{Resp}
\begin{Resp}\Rsign May he pray for all people, that their sins may be forgiven unto them\end{Resp}
\begin{Resp}\Vsign Glory be to the Father, and to the Son, \Star{} and to the Holy Ghost.\end{Resp}
\begin{Resp}\Rsign May he pray for all people, that their sins may be forgiven unto them\end{Resp}
\switchcolumn*

\LA
\LessonHead{Lectio VII}
\switchcolumn
\EN
\LessonHead{Lesson VII}
\switchcolumn*

\LA
\LessonTitle{Léctio sancti Evangélii secúndum Matthǽum.}
\switchcolumn
\EN
\LessonTitle{From the Holy Gospel according to Matthew}
\switchcolumn*

\LA
\Cite{Matt 5:13-19}
\switchcolumn
\EN
\Cite{Matt 5:13-19}
\switchcolumn*

\LA
\noindent In illo témpore: Dixit Jesus discipulis suis: Vos estis sal terræ. Quod si sal evamierit, in quo salietur? Ad nihilum valet ultra, nisi ut mittatur foras, et conculcetur ab hominibus. Et réliqua.\par
\switchcolumn
\EN
\noindent At that time, Jesus said unto His disciples: “Ye are the salt of the earth; but if the salt have lost his savour, wherewith shall it be salted? It is thenceforth good for nothing, but to be cast out, and to be trodden under foot of men.” And so on.\par
\switchcolumn*

\LA
\LessonTitle{Homilia sancti Isidori Epíscopi.}
\switchcolumn
\EN
\LessonTitle{Homily by St. Isidore, Archbishop (of Seville.)}
\switchcolumn*

\LA
\Source{Ex libro 2. Officiorum ad S. Fulgentium cap. 5.}
\switchcolumn
\EN
\Source{Bk. ii. to St. Fulgentius in Offices, c. 5}
\switchcolumn*

\LA
\noindent Qui in erudiendis, atque instituendis ad virtutem populis præerit, necesse est, ut in omnibus sanctus sit, et in nullo reprehensibihs habeatur. Qui enim alium de peccatis arguit, ipse a peccato debet esse alienus. Nam qua fronte subjectos arguere poterit, cum illi statim possit correctus ingerere: Ante doce te quæ recta sunt? Primitus quippe semetipsum corrigere debet, qui alios ad bene vivendum admonere studet; ita ut in omnibus semetipsum formam vivendi præbeat, cunctosque ad bonum opus, et doctrina et opere provocet. Cui étiam scientia Scripturarum necessaria est: quia si episcopi tantum sancta sit vita, sibi soli prodest, sic vivens. Porro si et doctrina et sermone fuerit eruditus, potest ceteros quoque instruiere, et docere suos, et adversarios repercutere, qui nisi refutati fuerint atque convicti, facile possunt simplicium corda pervertere.\par
\switchcolumn
\EN
\noindent Whosoever is set over the people to teach them and to catechize them in good works, him it behooveth in all things to be holy, and in nothing to be held blameworthy. For he which rebuketh another for sin, should have no dealings with sin himself. Since with what face can he rebuke them which are under him, if he which is rebuked of him be able to answer him straightway, saying Begin by teaching thyself to do well? Verily, whosoever setteth himself to teach others to live well, him it behooveth first of all to correct his own life, so that in all things he may be able to give the same his own life for an example, and may provoke all to good living by his works as well as by his words. Likewise also he must needs be learned in the Scriptures, since if the life of a Bishop be holy only, then is he profitable to himself only. But if he be learned also in his teaching and discourse, he is able to edify his neighbors, both teaching such as are his own, and confounding the gainsayers, who, unless they be confounded and unmasked, are easily able to lead astray the hearts of the simple.\par
\switchcolumn*

\LA
\begin{Resp}\Rsign Amávit eum Dóminus, et ornávit eum: stolam glóriæ índuit eum, \Star{} Et ad portas paradísi coronávit eum.\end{Resp}
\begin{Resp}\Vsign Induit eum Dóminus lorícam fídei, et ornávit eum.\end{Resp}
\begin{Resp}\Rsign Et ad portas paradísi coronávit eum.\end{Resp}
\switchcolumn
\EN
\begin{Resp}\Rsign The Lord loved him and beautified him He clothed him with a robe of glory \Star{} And crowned him at the gates of Paradise.\end{Resp}
\begin{Resp}\Vsign The Lord hath put on him the breast-plate of faith, and hath adorned him.\end{Resp}
\begin{Resp}\Rsign And crowned him at the gates of Paradise.\end{Resp}
\switchcolumn*

\LA
\LessonHead{Lectio VIII}
\switchcolumn
\EN
\LessonHead{Lesson VIII}
\switchcolumn*

\LA
\noindent Hujus sermo debet esse purus, simplex, apertus, plenus gravitatis et honestatis, plenus suavitatis et gratiæ, tractans de mysterio legis, de doctrina fidei, de virtute continentiæ, de disciplina justitiæ: unumquemque admonens diversa exhortatione, juxta professionem, morumque qualitatem: scilicet ut prænoscat quid, cui, quando, vel quomodo proferat. Cujus præ ceteris speciale officium est, Scripturas legere, percurrere Canones, exempla sanctorum imitari, vigiliis, jejuniis, orationibus incumbere, cum fratribus pacem habere, nec quemquam de membris suis discerpere; nullum damnare, nisi comprobatum, nullum excommunicare nisi discussum. Quique ita humilitate pariter, et auctoritate præesse debet, ut neque per nimiam humilitatem suam subditorum vitia convalescere faciat, neque per immoderantiam severitatis potestatem exerceat, sed tanto cautius erga commissos sibi, quanto durius a Christo indagari formidat.\par
\switchcolumn
\EN
\noindent Such an one it behooveth, that his discourse should be pure, plain, open, very weighty, and seemly, full of sweetness and comeliness, touching often the mystery of Law, the teaching of faith, the manliness of self-control, and the training of righteousness. Such an one it behooveth to exhort all men with varying exhortation, according to the profession and way of life of each, that is to say, such an one must know what, to whom, when, and how to speak. His duty is, before all others, to read the Scriptures, to know the Canons, to copy the examples of the Saints, to be instant in watching, fasting, and prayer, to keep peace with his brethren, to separate himself from none of the members of Christ, to condemn no man untried, and to excommunicate no man unheard. Such an one it behoveth, as he is the first in authority, so also to be the first in lowliness, yet ever so, that, by misplaced lowliness, he suffer not nor encourage the sins of those that are under him, nor use his authority hardly and with violence, but as one that is the more careful of the flock committed unto him, as being mindful of that stricter account which he will have to give at the fearful judgment seat of Christ.\par
\switchcolumn*

\LA
\begin{Resp}\Rsign In médio Ecclésiæ apéruit os ejus, \Star{} Et implévit eum Dóminus spíritu sapiéntiæ et intelléctus.\end{Resp}
\begin{Resp}\Vsign Jucunditátem et exsultatiónem thesaurizávit super eum.\end{Resp}
\begin{Resp}\Rsign Et implévit eum Dóminus spíritu sapiéntiæ et intelléctus.\end{Resp}
\begin{Resp}\Vsign Glória Patri, et Fílio, \Star{} et Spirítui Sancto.\end{Resp}
\begin{Resp}\Rsign Et implévit eum Dóminus spíritu sapiéntiæ et intelléctus.\end{Resp}
\switchcolumn
\EN
\begin{Resp}\Rsign In the midst of the congregation did the Lord open his mouth. \Star{} And filled him with the spirit of wisdom and understanding.\end{Resp}
\begin{Resp}\Vsign He made him rich with joy and gladness.\end{Resp}
\begin{Resp}\Rsign And filled him with the spirit of wisdom and understanding.\end{Resp}
\begin{Resp}\Vsign Glory be to the Father, and to the Son, \Star{} and to the Holy Ghost.\end{Resp}
\begin{Resp}\Rsign And filled him with the spirit of wisdom and understanding.\end{Resp}
\switchcolumn*

\LA
\LessonHead{Lectio IX}
\switchcolumn
\EN
\LessonHead{Lesson IX}
\switchcolumn*

\LA
\noindent Tenebit quoque illam supereminentem donis omnibus caritatem, sine qua omnis virtus nihil est. Custos enim castitatis, caritas; locus autem hujus custodis humilitas. Habebit étiam inter hæc omnia castitatis eminentiam: ita ut mens Christo dedita, ab omni inquinamento carnis sit munda et libera. Inter hæc oportebit cum sollicita dispensatione curam pauperum gerere, esurientes pascere, vestire nudos, suscipere peregrinos, captivos redimere, viduas ac pupillos tueri, pervigilem in cunctis exhibere curam, providentiam habere distributione discreta. In quo étiam hospitalitas ita erit præcipua, ut omnes cum benignitate et caritate suscipiat. Si enim omnes fideles illud Evangelium audire desiderant: Hospes fui, et suscepistis me; quanto magis episcopus, cujus diversorium cunctorum debet esse receptaculum?\par
\switchcolumn
\EN
\noindent Such an one must have firm hold on charity, that gift which surpassed all others, and without which all others are nothing worth. Charity is the keeper of chastity, and that keeper's home is lowly-mindedness. With all other gifts he must needs be eminent for purity, yea, his must be a mind belonging utterly to Christ, and clean and free from any fleshly defilement. But these are not all his needful gifts. Besides these, it behoveth him to undertake the care of the poor, and to do the same with zeal and likewise with prudence, to feed the hungry, to clothe the naked, to entertain strangers, to ransom prisoners, to be the guardian of the widow and the orphan, to watch over all without ceasing, and to be heedful that his alms be neither foolish nor wasteful. In him hospitality must shine, entertaining all men with courtesy and brotherly love; for if it be the duty of all the faithful to listen to that Gospel which saith: “I was a stranger, and ye took Me in”, (Matth. xxv. 35,) how much more is it the duty of Bishops, whose house it behooveth to be an home for all men?\par
\switchcolumn*

\LA
\TeDeum{Te Deum}
\switchcolumn
\EN
\TeDeum{Te Deum}
\switchcolumn*

\LA
\LessonHead{Lectio IX (pro commemoratione)}
\switchcolumn
\EN
\LessonHead{Lesson IX (when commemorated)}
\switchcolumn*

\LA
\LessonTitle{Commemoratio: S. Isidori Episcopi Confessoris et Ecclesiæ Doctoris}
\switchcolumn
\EN
\LessonTitle{Commemoration of: S. Isidore of Seville, Bishop, Confessor and Doctor of the Church}
\switchcolumn*

\LA
\noindent Isidórus, natióne Hispánus, ex nova Carthágine Severiáno patre, provínciæ duce, natus, a sanctis epíscopis Leándro Hispalénsi et Fulgéntio Carthaginiénsi, frátribus suis, pie et liberáliter educátus, omni scientiárum et christianárum virtútum génere præstantíssimus evásit. Leándro vita functo, ad Hispalénsem Cáthedram assúmptus, Sedis apostólicæ in tota Hispánia vicárius constitútus est. In episcopátu se ómnium bonórum óperum præbuit exémplar, et instaurándæ ecclesiásticæ discíplinæ sollícitus máxime fuit. Coácto Híspali concílio, Acephalórum hǽresim, Hispániæ jam minitántem, acri et eloquénti disputatióne fregit atque contrívit. Tantam apud omnes sanctitátis et doctrínæ famam adéptus est, ut, elápso vix ab ejus óbitu anno sextodécimo, Doctor egrégius merúerit appellári. Libros scripsit perútiles et eruditióne plenos. Præfuit concílio Toletáno quarto, ómnium Hispániæ celebérrimo. Dénique, postquam quadragínta círciter annos suam rexísset ecclésiam, Híspali migrávit in cælum, anno sexcentésimo trigésimo sexto.\par
\switchcolumn
\EN
\noindent Isidore of Spain was born at Carthagena, his father, Severianus, being governor of the province. His brothers, Leander of Seville and Fulgentius of Carthagena, both holy bishops, educated him in the love of God and in the liberal arts, and he became outstanding in all forms of knowledge and Christian virtue. When Leander died, Isidore was raised to the bishopric of Seville and made vicar apostolic for the whole of Spain. In his episcopal ministry he gave an example of all good works, and was especially concerned with the restoration of ecclesiastical discipline. When a Council was convoked at Seville, he broke up and stamped out the heresy of the Acephali, then threatening Spain, by the force and eloquence of his arguments. The fame of his holiness and teaching became so widespread that, hardly sixteen years after his death, he was called the illustrious Doctor. He wrote many useful books filled with learning, and presided over the fourth Council of Toledo, the most celebrated of those held in Spain. Finally, after he had ruled his church for about forty years, he died at Seville in the year 636.\par
\switchcolumn*

\LA
\TeDeum{Te Deum}
\switchcolumn
\EN
\TeDeum{Te Deum}
\switchcolumn*

\end{paracol}

\par\needspace{10\baselineskip}\addvspace{1.2em}{\centering\small\textsc{Die 5 Aprilis} · \textsc{5 April}\par}
\Office{S. Vincentii Ferrerii Confessoris}{St. Vincent Ferrer, Confessor}{Duplex}
\addcontentsline{toc}{subsection}{Die 5 Aprilis · S. Vincentii Ferrerii Confessoris\enspace\textit{St. Vincent Ferrer, Confessor}}
\begin{paracol}{2}
\LA
\RefLine{Lectiones I–III: ex Commúni Confessoris non pontificis (C5).}
\switchcolumn
\EN
\RefLine{Lessons I–III: from the Common of a Confessor not a Bishop (C5).}
\switchcolumn*

\LA
\RefLine{Lectiones VII–VIII: ex Commúni Confessoris non pontificis (C5).}
\switchcolumn
\EN
\RefLine{Lessons VII–VIII: from the Common of a Confessor not a Bishop (C5).}
\switchcolumn*

\LA
\RefLine{Lectio IX: de Scriptura occurrente.}
\switchcolumn
\EN
\RefLine{Lesson IX: from the Scripture of the day.}
\switchcolumn*

\LA
\LessonHead{Lectio IV}
\switchcolumn
\EN
\LessonHead{Lesson IV}
\switchcolumn*

\LA
\noindent Vincentius, honesta stirpe Valentiæ in Hispania natus, ab ineunte ætate cor gessit senile. Qui dum caliginosi hujus sæculi labilem cursum pro ingenii sui modulo consideraret, religionis habitum in ordine Prædicatorum decimo octavo ætatis suæ anno suscepit; et emissa solemni professione, sacris litteris sedulo incumbens, theologiæ lauream summa cum laude consecutus est. Mox obtenta a superioribus licentia, verbum Dei prædicare, Judæorum perfidiam arguere, Saracenorum errores confutare tanta virtute et efficacia cœpit, ut ingentem ipsorum infiddlium multitudinem ad Christi fidem perduxerit, et multa Christianorum millia a peccatis ad pœnitentiam, a vitus ad virtutem revocarit. Electus enim a Deo, ut monita salutis in omnes gentes, tribus et linguas difunderet, et extremi tremendique judicii diem appropinquare ostenderet, omnium auditorum animos terrore concussos, atque a terrenis affectibus avulsos, ad Dei amórem excitabat.\par
\switchcolumn
\EN
\noindent This Vincent was born of respectable parents, at Valencia in Spain, (upon the 23rd day of January, in the year of our Lord 1357.) Even as a child he had an heart like the heart of an old man. Considering, to the utmost of his young understanding, how fleeting is the course of this dark world, he, in the eighteenth year of his age, took the habit of a Friar in the Order of Preachers. After he had made his solemn profession, he devoted himself to sacred learning, and took the degree of Master in Divinity with much distinction. He soon after received permission from his superiors to preach the word of God, on which duty he entered with such power and success, striving against the unbelief of the Jews, and overthrowing the errors of the Saracens, that he brought an exceeding great multitude of unbelievers to believe in Christ, and turned many thousands of Christians from sin to sorrow, and from vice to virtue. He was a chosen vessel unto God to proclaim the tidings of salvation among all nations, and tribes, and tongues, crying out that the last day, that awful day of judgment, is at hand, smiting consternation into the minds of all, as many as heard him, weaning their love from a perishing world, and turning it to God.\par
\switchcolumn*

\LA
\begin{Resp}\Rsign Honéstum fecit illum Dóminus, et custodívit eum ab inimícis, et a seductóribus tutávit illum: \Star{} Et dedit illi claritátem ætérnam.\end{Resp}
\begin{Resp}\Vsign Justum dedúxit Dóminus per vias rectas, et osténdit illi regnum Dei.\end{Resp}
\begin{Resp}\Rsign Et dedit illi claritátem ætérnam.\end{Resp}
\switchcolumn
\EN
\begin{Resp}\Rsign The Lord made him honourable, and defended him from his enemies, and kept him safe from those that lay in wait for him; \Star{} And gave him perpetual glory.\end{Resp}
\begin{Resp}\Vsign He went down with him into the pit, and left him not in bonds.\end{Resp}
\begin{Resp}\Rsign And gave him perpetual glory.\end{Resp}
\switchcolumn*

\LA
\LessonHead{Lectio V}
\switchcolumn
\EN
\LessonHead{Lesson V}
\switchcolumn*

\LA
\noindent In hoc autem apostolico munere hic vitæ ejus tenor perpetuus fuit quotidie Missam summo mane cum cantu celebravit. Quotidie ad populum concionem habuit; inviolabile semper jejunium, nisi urgens adesset necessitas, servavit; sancta et recta consilia nulli denegavit, carnes numquam comedit, nec vestem lineam induit: populorum jurgia sedavit, dissidentia regna pace composuit: et cum vestis inconsutilis Ecclesiæ diro schismate scinderetur, ut uniretur, et unita servaretur, plurimum laboravit. Virtutibus omnibus claruit, suosque detractores et persecutores, in simplicitate et humilitate ambulans, cum mansuetudine recepit, et amplexus est.\par
\switchcolumn
\EN
\noindent While Vincent wrought the Apostolic work of preaching committed to him, he lived ever as follows j Every morning he sang a solemn Mass, and every day he preached in public. He fasted every day, unless prevented by some absolute necessity. He refused to no one his holy and just advice. He never ate meat, nor wore linen. He quieted public disturbances, and negotiated the peace of kingdoms. When the seamless garment of the Church was rent by an horrid schism, he worked his every nerve to unite it again, and keep it one. He was a burning and a shining light of all virtues, walking always in lowliness and simpleness, so that he meekly welcomed and embraced them which spake evil against him and persecuted him.\par
\switchcolumn*

\LA
\begin{Resp}\Rsign Amávit eum Dóminus, et ornávit eum: stolam glóriæ índuit eum, \Star{} Et ad portas paradísi coronávit eum.\end{Resp}
\begin{Resp}\Vsign Índuit eum Dóminus lorícam fídei, et ornávit eum.\end{Resp}
\begin{Resp}\Rsign Et ad portas paradísi coronávit eum.\end{Resp}
\switchcolumn
\EN
\begin{Resp}\Rsign The Lord loved him and beautified him; He clothed him with a robe of glory, \Star{} And crowned him at the gates of Paradise.\end{Resp}
\begin{Resp}\Vsign The Lord hath put on him the breast-plate of faith, and hath adorned him.\end{Resp}
\begin{Resp}\Rsign And crowned him at the gates of Paradise.\end{Resp}
\switchcolumn*

\LA
\LessonHead{Lectio VI}
\switchcolumn
\EN
\LessonHead{Lesson VI}
\switchcolumn*

\LA
\noindent Per ipsum divina virtus, in confirmationem vitæ et prædicationis ejus, multa signa et miracula fecit. Nam frequentissime super ægros manus imposuit, et sanitatem adepti sunt: spiritus immundos e corporibus expulit: surdis auditum, mutis loquelam, cæcis visum restituit, leprosos mundavit, mortuos suscitavit. Senio tandem et morbo confectus infatigabilis Evangelii præco, plurimis Europæ provinciis cum ingenti animarum fructu peragratis, Venetiæ in Britannia minori prædicationis et vitæ cursum feliciter consummavit, anno salutis millesimo quadringentesimo decimo nono: quem Callistus tertius sanctorum numero ascripsit.\par
\switchcolumn
\EN
\noindent The Power of God confirmed his life and doctrine with many great signs and wonders. He often laid his hands upon the sick and they recovered. He cast out unclean spirits, and made the deaf to hear, the dumb to speak, and the blind to see. He cleansed the lepers, and raised the dead. After passing through many countries of Europe with exceeding profit to souls, worn out with age and disease, but still ever the same unwearied herald of the Gospel, he brought his life and his preaching together to an happy end, at Vannes in Brittany, (upon the 5th day of April,) in the year of salvation 1419. Pope Callistus III. numbered him with the Saints.\par
\switchcolumn*

\LA
\begin{Resp}\Rsign Iste homo perfécit ómnia quæ locútus est ei Deus, et dixit ad eum: Ingrédere in réquiem meam: \Star{} Quia te vidi justum coram me ex ómnibus géntibus.\end{Resp}
\begin{Resp}\Vsign Iste est, qui contémpsit vitam mundi, et pervénit ad cæléstia regna.\end{Resp}
\begin{Resp}\Rsign Quia te vidi justum coram me ex ómnibus géntibus.\end{Resp}
\begin{Resp}\Vsign Glória Patri, et Fílio, \Star{} et Spirítui Sancto.\end{Resp}
\begin{Resp}\Rsign Quia te vidi justum coram me ex ómnibus géntibus.\end{Resp}
\switchcolumn
\EN
\begin{Resp}\Rsign This is he which did according unto all that God commanded him; and God said unto him: Enter thou into My rest, \Star{} For thee have I seen righteous before Me among all people.\end{Resp}
\begin{Resp}\Vsign This is he which loved not his life in this world, and is come unto an everlasting kingdom.\end{Resp}
\begin{Resp}\Rsign For thee have I seen righteous before Me among all people.\end{Resp}
\begin{Resp}\Vsign Glory be to the Father, and to the Son, \Star{} and to the Holy Ghost.\end{Resp}
\begin{Resp}\Rsign For thee have I seen righteous before Me among all people.\end{Resp}
\switchcolumn*

\LA
\LessonHead{Lectio IX (pro commemoratione)}
\switchcolumn
\EN
\LessonHead{Lesson IX (when commemorated)}
\switchcolumn*

\LA
\LessonTitle{Commemoratio: S. Vincentii Ferrerii Confessoris}
\switchcolumn
\EN
\LessonTitle{Commemoration of: St. Vincent Ferrer, Confessor}
\switchcolumn*

\LA
\noindent Vincéntius, honésta stirpe Valéntiæ in Hispánia natus, décimo octávo ætátis anno religiónis hábitum in órdine Prædicatórum índuit; in quo, sacris lítteris sédulo incúmbens, theologíæ láuream summa cum laude obtínuit. Mox verbum Dei tanta virtúte et efficácia prædicáre cœpit, ut ingéntem infidélium multitúdinem ad Christi fidem perdúxerit, et multa Christianórum míllia a peccátis ad pœniténtiam revocáverit. Quotídie Missam cum cantu celebrávit, quotídie ad pópulum conciónem hábuit, carnes numquam comédit, populórum júrgia sedávit; et, cum vestis inconsútilis Ecclésiæ diro schísmate discinderétur, ut unirétur et uníta servarétur, plúrimum laborávit. Sénio tandem et morbo conféctus, ac miráculis appríme clarus, Venétiæ in Británnia minóri sanctíssime óbiit. Quem Callístus Papa tértius inter Sanctos rétulit.\par
\switchcolumn
\EN
\noindent Born of a good family at Valencia in Spain, Vincent was clothed in the habit of the Friars Preachers at the age of eighteen. He applied himself zealously to sacred studies, and received the degree of Doctor of Theology with great distinction. Soon he began to preach the word of God with such force and effectiveness that he brought a vast number of unbelievers to the faith of Christ and called back many thousands of Christians from sin to repentance. He sang Mass every day and daily gave a sermon to the people, never ate meat, calmed quarrels and dissensions and, when the seamless garment of the Church was rent by harmful schism, laboured unceasingly to unite it and to preserve its unity. At length, worn out by old age and illness, famous for his miracles, he died a most holy death at Vannes in Britanny, and Pope Callistus III numbered him among the Saints.\par
\switchcolumn*

\end{paracol}

\par\needspace{10\baselineskip}\addvspace{1.2em}{\centering\small\textsc{Die 11 Aprilis} · \textsc{11 April}\par}
\Office{S. Leonis I Papæ Confessoris et Ecclesiæ Doctoris}{St. Leo the Great, Pope, Confessor and Doctor of the Church}{Duplex}
\addcontentsline{toc}{subsection}{Die 11 Aprilis · S. Leonis I Papæ Confessoris et Ecclesiæ Doctoris\enspace\textit{St. Leo the Great, Pope, Confessor and Doctor of the Church}}
\begin{paracol}{2}
\LA
\RefLine{Lectiones I–III: de Scriptura occurrente.}
\switchcolumn
\EN
\RefLine{Lessons I–III: from the Scripture of the day.}
\switchcolumn*

\LA
\RefLine{Lectiones VII–IX: ex Commúni Summorum Pontificum Confessorum (C4b).}
\switchcolumn
\EN
\RefLine{Lessons VII–IX: from the Common of Popes (C4b).}
\switchcolumn*

\LA
\LessonHead{Lectio IV}
\switchcolumn
\EN
\LessonHead{Lesson IV}
\switchcolumn*

\LA
\noindent Leo primus, Etruscus, eo tempore præfuit Ecclesiæ, cum rex Hunnorum Attila, cognomento Flagellum Dei, in Italiam invadens, Aquilejam triennii obsidione captam diripuit et incendit. Unde cum Romam ardenti furore raperetur, jamque copias, ubi Mincius in Padum influit, trajicere pararet, occurrit ei Leo, malorum Italiæ impendentium misericordia permotus: cujus divina eloquentia persuasum est Attilæ, ut regrederetur. Qui interrogatus a suis, quid esset quod præter consuetudinem tam humiliter Romani Pontificis imperata faceret; respondit, se astantem quemdam alium, illo loquente, sacerdotali habitu veritum esse, sibi stricto gladio minitantem mortem, nisi Leoni obtemperaret. Quare in Pannoniam reversus est.\par
\switchcolumn
\EN
\noindent Leo was an Etruscan who ruled the Church at the time when Attila, king of the Huns, whose surname is the Scourge of God, invaded Italy, and after a siege of three years, took, sacked, and burnt Aquileia. Thence he was hurrying to Rome, on fire with anger, and his troops were already preparing to cross the Po, at the place where that river is joined by the Mincio, when he was met by Leo, moved with compassion at the thought of the ruin which hung over Italy. By his God-given eloquence, Attila was persuaded to turn back, and when he was afterwards asked by his servants why, contrary to his custom, he had so meekly yielded to the entreaties of the Bishop of Rome, he answered that he had been alarmed by a figure dressed like a Priest, which had appeared at the side of Leo while he was speaking, holding a drawn sword, and had made as though to kill the king unless he consented. And so he returned into Pannonia.\par
\switchcolumn*

\LA
\begin{Resp}\Rsign Invéni David servum meum, óleo sancto meo unxi eum: \Star{} Manus enim mea auxiliábitur ei, allelúja.\end{Resp}
\begin{Resp}\Vsign Nihil profíciet inimícus in eo, et fílius iniquitátis non nocébit ei.\end{Resp}
\begin{Resp}\Rsign Manus enim mea auxiliábitur ei, allelúja.\end{Resp}
\switchcolumn
\EN
\begin{Resp}\Rsign I have found David My servant, with My holy oil have I anointed him \Star{} For My hand shall help him, alleluia.\end{Resp}
\begin{Resp}\Vsign The enemy shall prevail nothing against him, nor the son of wickedness afflict him.\end{Resp}
\begin{Resp}\Rsign For My hand shall help him, alleluia.\end{Resp}
\switchcolumn*

\LA
\LessonHead{Lectio V}
\switchcolumn
\EN
\LessonHead{Lesson V}
\switchcolumn*

\LA
\noindent Leo autem Romæ singulari omnium lætitia exceptus, paulo post invadenti Urbem Genserico, eadem eloquentiæ vi et sanctitatis opinione persuasit, ut ab incendio, ignominiis, ac cædibus abstineret. Sed cum Ecclesiam a multis hæresibus oppugnari, maximeque a Nestorianis et Eutychianis exagitari videret; ad eam purgandam, et in fide catholica confirmandam, concilium Chalcedonense indixit, ubi sexcentis triginta coactis episcopis, Eutyches et Dioscorus, et iterum Nestorius condemnati sunt: ejusdemque concilii decreta sua auctoritate confirmavit.\par
\switchcolumn
\EN
\noindent While Leo went back to Rome, where he was received with rejoicing by all men. A while later, Genseric entered the city, but Leo, by the power of his eloquence and the authority of his holy life, persuaded him to abstain from fire, insult, and slaughter. When Leo beheld how the Church was assailed by many heresies, and in dire trouble through the Nestorians and Eutychians, to purify the same and establish her in the Catholic Faith, he called the Council of Chalcedon, where, in an assembly of six hundred and thirty Bishops Nestorius was again condemned, along with Eutyches and Dioscorus; the decrees of which Council were confirmed by the authority of Leo.\par
\switchcolumn*

\LA
\begin{Resp}\Rsign Pósui adjutórium super poténtem, et exaltávi eléctum de plebe mea: \Star{} Manus enim mea auxiliábitur ei, allelúja.\end{Resp}
\begin{Resp}\Vsign Invéni David servum meum, óleo sancto meo unxi eum.\end{Resp}
\begin{Resp}\Rsign Manus enim mea auxiliábitur ei, allelúja.\end{Resp}
\switchcolumn
\EN
\begin{Resp}\Rsign I have laid help upon one that is mighty, and have exalted one chosen out of My people \Star{} For My hand shall help him, alleluia.\end{Resp}
\begin{Resp}\Vsign I have found David My servant, with My holy oil have I anointed him.\end{Resp}
\begin{Resp}\Rsign For My hand shall help him, alleluia.\end{Resp}
\switchcolumn*

\LA
\LessonHead{Lectio VI}
\switchcolumn
\EN
\LessonHead{Lesson VI}
\switchcolumn*

\LA
\noindent His actis, sanctus Pontifex se ad reficiendas et ædificandas ecclesias convertit. Cujus suasu Demetria, pia femina, sancti Stephani ecclesiam construxit in suo fundo via Latina, tertio ab Urbe milliario: ipse via Appia sub nomine sancti Cornelii alteram condidit. Multas præterea et sacras ædes, et sacra earum vasa restituit. In tribus basilicis, Petri, Pauli, et Constantiniana, cameras exstruxit: ædificavit monasterium vicinum basilicæ sancti Petri: sepulcris Apostolorum custodes adhibuit, quos cubicularios appellavit. Statuit, ut in actione mysterii diceretur: Sanctum sacrificium, immaculatam hostiam. Sanxit, ne monacha benedictum capitis velamen reciperet, nisi quadraginta annorum virginitatem probasset. His et aliis præclare gestis, cum multa sancte et luculenter scripsisset, quarto Idus Novembris obdormivit in Domino. Sedit in pontificatu annos viginti unum, mensem unum, dies tredecim.\par
\switchcolumn
\EN
\noindent After these matters, this holy Pope set himself to the restoration and building of Churches. By his advice that godly woman Demetria built the Church of St. Stephen upon her farm on the Latin Road, at the third milestone from the city. He himself built another Church upon the Appian Way, which Church is called that of St. Cornelius. He restored likewise many other Churches, and the holy vessels used therein. He built Clergy -houses at the three Basilicas of Peter, Paul, and Constantine. He built a monastery hard by the Basilica of St. Peter. He appointed for the graves of the Apostles certain keepers, whom he called the Chamberlains of the said Apostles. He ordained that in the action of the Mystery should be uttered the words An holy sacrifice, an offering without spot. He ordered that no nun should have the covering of her head blessed 4 until she had made trial of her virginity for forty years. After doing all these and other illustrious works, and after he had written much that is both godly and easy to be understood, he fell asleep in the Lord on the eleventh day of April, (in the year 461.) He held the Papal See for twenty years, one month, and thirteen days.\par
\switchcolumn*

\LA
\begin{Resp}\Rsign Iste est, qui ante Deum magnas virtútes operátus est, et omnis terra doctrína ejus repléta est: \Star{} Ipse intercédat pro peccátis ómnium populórum, allelúja.\end{Resp}
\begin{Resp}\Vsign Iste est, qui contémpsit vitam mundi, et pervénit ad cæléstia regna.\end{Resp}
\begin{Resp}\Rsign Ipse intercédat pro peccátis ómnium populórum, allelúja.\end{Resp}
\begin{Resp}\Vsign Glória Patri, et Fílio, \Star{} et Spirítui Sancto.\end{Resp}
\begin{Resp}\Rsign Ipse intercédat pro peccátis ómnium populórum, allelúja.\end{Resp}
\switchcolumn
\EN
\begin{Resp}\Rsign This is he which wrought great wonders before God, and the whole earth is full of his teaching \Star{} May he pray for all people, that their sins may be forgiven unto them, alleluia.\end{Resp}
\begin{Resp}\Vsign This is he which loved not his life in this world, and hath attained unto the kingdom of heaven.\end{Resp}
\begin{Resp}\Rsign May he pray for all people, that their sins may be forgiven unto them, alleluia.\end{Resp}
\begin{Resp}\Vsign Glory be to the Father, and to the Son, \Star{} and to the Holy Ghost.\end{Resp}
\begin{Resp}\Rsign May he pray for all people, that their sins may be forgiven unto them, alleluia.\end{Resp}
\switchcolumn*

\LA
\LessonHead{Lectio IX (pro commemoratione)}
\switchcolumn
\EN
\LessonHead{Lesson IX (when commemorated)}
\switchcolumn*

\LA
\LessonTitle{Commemoratio: S. Leonis I Papæ Confessoris et Ecclesiæ Doctoris}
\switchcolumn
\EN
\LessonTitle{Commemoration of: St. Leo the Great, Pope, Confessor and Doctor of the Church}
\switchcolumn*

\LA
\noindent Leo primus, Etrúscus, eo témpore prǽfuit Ecclésiæ, cum rex Hunnórum Attila, cognoménto Flagéllum Dei, in Itáliam invádens, post captam et incénsam Aquiléjam, Romam petitúrus cópias paráverat. Cui occúrrens Leo, divína persuásit eloquéntia, ut regrederétur. Tum Romæ singulári ómnium lætítia excéptus, paulo post, invadénti Urbem Genseríco eádem eloquéntiæ vi persuásit, ut ab incéndiis, ignomíniis et cǽdibus abstinéret. Cum autem Ecclésiam a multis hærésibus maximéque a Nestoriánis et Eutychiánis exagitári vidéret, Concílium Chalcedonénse indíxit; ubi, sexcéntis trigínta coáctis epíscopis, Eutyches et Dióscorus, et íterum Nestórius condemnáti sunt; cujus Concílii decréta auctoritáte sua firmávit. Multas sacras ædes exstrúxit, ac monastérium vicínum basílicæ sancti Petri ædificávit. His et áliis præcláre gestis, cum multa sancte et luculénter scripsísset, quarto Idus Novémbris obdormívit in Dómino, anno pontificátus sui vigésimo primo.\par
\switchcolumn
\EN
\noindent Leo I, an Etruscan, ruled over the Church at the time when Attila, King of the Huns and called the Scourge of God, was invading Italy; he had taken and burned Aquileia and was preparing his forces to attack Rome. Leo went out to meet him and, by God-given eloquence, persuaded him to withdraw; then Leo was welcomed back to Rome with great rejoicing. A little later, when Genseric was invading the city, Leo persuaded him, with the same forceful eloquence, to abstain from burning, outrages and slaughter. When Leo saw the Church harassed by many heresies, and especially by the Nestorians and the Eutychians, he called the Council of Chalcedon at which, with six hundred and thirty bishops assembled, Eutyches and Dioscorus were condemned and the condemnation of Nestorius repeated. The decrees of this Council were then confirmed by Leo's authority. He constructed many churches and built a monastery near the Basilica of St. Peter. After a life filled with these and other admirable works, including a great number of holy and eloquent writings, he fell asleep in the Lord on the tenth day of November, in the twenty-first year of his pontificate.\par
\switchcolumn*

\end{paracol}

\par\needspace{10\baselineskip}\addvspace{1.2em}{\centering\small\textsc{Die 13 Aprilis} · \textsc{13 April}\par}
\Office{S. Hermenegildi Martyris}{St. Hermenegild, Martyr}{Semiduplex}
\addcontentsline{toc}{subsection}{Die 13 Aprilis · S. Hermenegildi Martyris\enspace\textit{St. Hermenegild, Martyr}}
\begin{paracol}{2}
\LA
\RefLine{Lectiones I–III: de Scriptura occurrente.}
\switchcolumn
\EN
\RefLine{Lessons I–III: from the Scripture of the day.}
\switchcolumn*

\LA
\LessonHead{Lectio I}
\switchcolumn
\EN
\LessonHead{Lesson I}
\switchcolumn*

\LA
\LessonTitle{De Epístola beáti Pauli Apóstoli ad Romános}
\switchcolumn
\EN
\LessonTitle{Lesson from the letter of St Paul the Apostle to the Romans}
\switchcolumn*

\LA
\Cite{Rom 8:12-19}
\switchcolumn
\EN
\Cite{Rom 8:12-19}
\switchcolumn*

\LA
\noindent Fratres: Debitóres sumus non carni, ut secúndum carnem vivámus. Si enim secúndum carnem vixéritis, moriémini: si autem spíritu facta carnis mortificavéritis, vivétis. Quicúmque enim spíritu Dei agúntur, ii sunt fílii Dei. Non enim accepístis spíritum servitútis íterum in timóre, sed accepístis spíritum adoptiónis filiórum, in quo clamámus: Abba (Pater). Ipse enim Spíritus testimónium reddit spirítui nostro quod sumus fílii Dei. Si autem fílii, et herédes: herédes, quidem Dei, coherédes autem Christi: si tamen compátimur ut et conglorificémur. Exístimo enim quod non sunt condígnæ passiónes hujus témporis ad futúram glóriam, quæ revelábitur in nobis. Nam exspectátio creatúræ revelatiónem filiórum Dei exspéctat.\par
\switchcolumn
\EN
\noindent Therefore, brethren, we are debtors, not to the flesh, to live according to the flesh. For if you live according to the flesh, you shall die: but if by the Spirit you mortify the deeds of the flesh, you shall live. For whosoever are led by the Spirit of God, they are the sons of God. For you have not received the spirit of bondage again in fear; but you have received the spirit of adoption of sons, whereby we cry: Abba, Father. For the Spirit himself giveth testimony to our spirit, that we are the sons of God. And if sons, heirs also; heirs indeed of God, and joint heirs with Christ: yet so, if we suffer with him, that we may be also glorified with him. For I reckon that the sufferings of this time are not worthy to be compared with the glory to come, that shall be revealed in us. For the expectation of the creature waiteth for the revelation of the sons of God\par
\switchcolumn*

\LA
\begin{Resp}\Rsign Iste Sanctus pro lege Dei sui certávit usque ad mortem, et a verbis impiórum non tímuit: \Star{} Fundátus enim erat supra firmam petram.\end{Resp}
\begin{Resp}\Vsign Iste est qui contémpsit vitam mundi, et pervénit ad cæléstia regna.\end{Resp}
\begin{Resp}\Rsign Fundátus enim erat supra firmam petram.\end{Resp}
\switchcolumn
\EN
\begin{Resp}\Rsign This man is holy, for he hath striven for the law of his God even unto death, and hath not feared for the words of the ungodly; \Star{} For he had his foundation upon a strong rock.\end{Resp}
\begin{Resp}\Vsign This is he which loved not his life in this world, and is come unto an everlasting kingdom.\end{Resp}
\begin{Resp}\Rsign For he had his foundation upon a strong rock.\end{Resp}
\switchcolumn*

\LA
\LessonHead{Lectio II}
\switchcolumn
\EN
\LessonHead{Lesson II}
\switchcolumn*

\LA
\Cite{Rom 8:28-34}
\switchcolumn
\EN
\Cite{Rom 8:28-34}
\switchcolumn*

\LA
\noindent Scimus autem quóniam diligéntibus Deum ómnia cooperántur in bonum iis qui secúndum propósitum vocáti sunt sancti. Nam quos præscívit et prædestinávit confórmes fíeri imáginis Fílii sui ut sit ipse primogénitus in multis frátribus. Quos autem prædestinávit, hos et vocávit: et quos vocávit, hos et justificávit: quos autem justificávit, illos et glorificávit. Quid ergo dicémus ad hæc? Si Deus pro nobis, quis contra nos? Qui étiam próprio Fílio suo non pepércit, sed pro nobis ómnibus trádidit illum, quómodo non étiam cum illo ómnia nobis donávit? Quis accusábit advérsus eléctos Dei? Deus qui justíficat, quis est qui condémnet? Christus Jesus, qui mórtuus est, immo qui et resurréxit, qui est ad déxteram Dei, qui étiam interpéllat pro nobis.\par
\switchcolumn
\EN
\noindent And we know that to them that love God, all things work together unto good, to such as, according to his purpose, are called to be saints. For whom he foreknew, he also predestinated to be made conformable to the image of his Son; that he might be the firstborn amongst many brethren. And whom he predestinated, them he also called. And whom he called, them he also justified. And whom he justified, them he also glorified. What shall we then say to these things? If God be for us, who is against us? He that spared not even his own Son, but delivered him up for us all, how hath he not also, with him, given us all things? Who shall accuse against the elect of God? God that justifieth. Who is he that shall condemn? Christ Jesus that died, yea that is risen also again; who is at the right hand of God, who also maketh intercession for us.\par
\switchcolumn*

\LA
\begin{Resp}\Rsign Justus germinábit sicut lílium: \Star{} Et florébit in ætérnum ante Dóminum.\end{Resp}
\begin{Resp}\Vsign Plantátus in domo Dómini, in átriis domus Dei nostri.\end{Resp}
\begin{Resp}\Rsign Et florébit in ætérnum ante Dóminum.\end{Resp}
\switchcolumn
\EN
\begin{Resp}\Rsign The righteous shall grow as the lily \Star{} Yea, he shall flourish in the presence of the Lord for ever.\end{Resp}
\begin{Resp}\Vsign Those that be planted in the house of the Lord, shall flourish in the courts of the house of our God.\end{Resp}
\begin{Resp}\Rsign Yea, he shall flourish in the presence of the Lord for ever.\end{Resp}
\switchcolumn*

\LA
\LessonHead{Lectio III}
\switchcolumn
\EN
\LessonHead{Lesson III}
\switchcolumn*

\LA
\Cite{Rom 8:35-39}
\switchcolumn
\EN
\Cite{Rom 8:35-39}
\switchcolumn*

\LA
\noindent Quis ergo nos separábit a caritáte Christi? tribulátio? an angústia? an fames? an núditas? an perículum? an persecútio? an gládius? (sicut scriptum est: Quia propter te mortificámur tota die: æstimáti sumus sicut oves occisiónis.) Sed in his ómnibus superámus propter eum qui diléxit nos. Certus sum enim quia neque mors, neque vita, neque Ángeli, neque Principátus, neque Virtútes, neque instántia, neque futúra, neque fortitúdo, neque altitúdo, neque profúndum, neque creatúra ália póterit nos separáre a caritáte Dei, quæ est in Christo Jesu Dómino nostro.\par
\switchcolumn
\EN
\noindent Who then shall separate us from the love of Christ? Shall tribulation? or distress? or famine? or nakedness? or danger? or persecution? or the sword? As it is written: For thy sake we are put to death all the day long. We are accounted as sheep for the slaughter. But in all these things we overcome, because of him that hath loved us. For I am sure that neither death, nor life, nor angels, nor principalities, nor powers, nor things present, nor things to come, nor might, Nor height, nor depth, nor any other creature, shall be able to separate us from the love of God, which is in Christ Jesus our Lord.\par
\switchcolumn*

\LA
\begin{Resp}\Rsign Iste cognóvit justítiam, et vidit mirabília magna, et exorávit Altíssimum \Star{} Et invéntus est in número Sanctórum.\end{Resp}
\begin{Resp}\Vsign Iste est qui contémpsit vitam mundi, et pervénit ad cæléstia regna.\end{Resp}
\begin{Resp}\Rsign Et invéntus est in número Sanctórum.\end{Resp}
\begin{Resp}\Vsign Glória Patri, et Fílio, \Star{} et Spirítui Sancto.\end{Resp}
\begin{Resp}\Rsign Et invéntus est in número Sanctórum.\end{Resp}
\switchcolumn
\EN
\begin{Resp}\Rsign This is he which knew righteousness, and saw great wonders, and made his prayer unto the Most High \Star{} And he is numbered among the Saints.\end{Resp}
\begin{Resp}\Vsign This is he which loved not his life in this world, and is come unto an everlasting kingdom.\end{Resp}
\begin{Resp}\Rsign And he is numbered among the Saints.\end{Resp}
\begin{Resp}\Vsign Glory be to the Father, and to the Son, \Star{} and to the Holy Ghost.\end{Resp}
\begin{Resp}\Rsign And he is numbered among the Saints.\end{Resp}
\switchcolumn*

\LA
\LessonHead{Lectio IV}
\switchcolumn
\EN
\LessonHead{Lesson IV}
\switchcolumn*

\LA
\LessonTitle{Ex libro Dialogorum sancti Gregorii Papæ}
\switchcolumn
\EN
\LessonTitle{From the Book of Dialogues written by Pope St. Gregory}
\switchcolumn*

\LA
\Source{Lib. 3 Cap. 31}
\switchcolumn
\EN
\Source{Lib. 3. Cap. 31}
\switchcolumn*

\LA
\noindent Hermenegildus rex, Leovigildi regis Visigothorum filius, ab Ariana hæresi ad fidem catholicam viro reverendissimo Leandro Hispalensi episcopo, dudum mihi in amicitiis familiariter juncto, prædicante, conversus est. Quem pater Arianus, ut ad eamdem hæresim rediret, et præmiis suadere, et minis terrere conatus est. Cumque ille constantissime responderet, numquam se veram fidem posse relinquere, quam semel agnovisset: iratus pater eum privavit regno, rebusque exspoliavit omnibus. Cumque nec sic virtutem mentis illius emollire valuisset, in arcta illum custodia concludens, collum manusque illius ferro ligavit. Cœpit ítaque Hermenegildus rex juvenis terrenum regnum despicere, et forti desiderio cæleste quærens, in ciliciis vinculatus jacens, omnipotenti Deo ad confortandum se, preces effundere; tantoque sublimius gloriam transeuntis mundi despicere, quanto et religatus agnoverat nihil fuísse, quod potuerit auferri.\par
\switchcolumn
\EN
\noindent Hermenegild, the son of Leovigild, King of the Visigoths, was turned from the Arian heresy to the Catholic Faith by the preaching of that most worshipful man Leander, Bishop of Seville, the same who was for a long season mine own familiar friend. Then his father, being himself an Arian, strove to bring him back to that heresy, first by offering him gifts, and then seeking to awe him by threatening. And when he answered alway that, having once had knowledge of the true faith, he never could forsake it, his father was wroth, and took away his kingdom from him, and plundered him of all his goods. And when not even so could he sap the manliness of his soul, he cast him into a most strait prison, having his neck and his hands in fetters of iron. And so that young King Hermenegild began to hold in little esteem an earthly kingdom, and to long exceedingly for an heavenly. Yea, he clothed himself in sackcloths in the prison, and as he lay bound therein, he poured forth supplications to Almighty God to give him strength. There he lay bound, having suffered the loss of all things, but his suffering made him but to esteem more worthless the glory of this world, which passeth away so easily.\par
\switchcolumn*

\LA
\begin{Resp}\Rsign Lux perpétua lucébit Sanctis tuis, Dómine, \Star{} Et ætérnitas témporum, allelúja, allelúja.\end{Resp}
\begin{Resp}\Vsign Lætítia sempitérna erit super cápita eórum: gáudium et exsultatiónem obtinébunt.\end{Resp}
\begin{Resp}\Rsign Et ætérnitas témporum, allelúja, allelúja.\end{Resp}
\switchcolumn
\EN
\begin{Resp}\Rsign Eternal light will shine over your saints, O Lord, \Star{} And the eternity of the times, alleluia, alleluia.\end{Resp}
\begin{Resp}\Vsign Everlasting joy shall be upon their heads, they shall obtain joy and gladness\end{Resp}
\begin{Resp}\Rsign And the eternity of the times, alleluia, alleluia.\end{Resp}
\switchcolumn*

\LA
\LessonHead{Lectio V}
\switchcolumn
\EN
\LessonHead{Lesson V}
\switchcolumn*

\LA
\noindent Superveniente autem Paschalis festivitatis die, intempestæ noctis silentio, ad eum perfidus pater Arianum episcopum misit, ut ex ejus manu sacrilegæ consecrationis communionem perciperet, atque per hoc ad patris gratiam redire mereretur. Sed vir Deo deditus, Ariano episcopo venienti exprobravit, ut debuit, ejusque a se perfidiam dignis increpationibus repulit: quia etsi exterius jacebat ligatus, apud se tamen in magno mentis culmine stabat securus. Ad se ítaque reverso episcopo, Arianus pater infremuit, statimque suos apparitores misit, qui constantissimum Confessorem Dei illic ubi jacebat, occiderent; quod et factum est. Nam mox ut ingressi sunt, securim cerebro ejus infigentes, vitam corporis abstulerunt: hocque in eo valuerunt perimere, quod ipsum quoque, qui peremptus est, in se constiterat despexisse. Sed pro ostendenda vera ejus gloria, superna quoque non defuere miracula. Nam cœpit in nocturno silentio psalmodiæ cantus ad corpus ejusdem Regis et Martyris audiri: atque ideo veraciter Regis, quia et Martyris.\par
\switchcolumn
\EN
\noindent But when the day of the glad Passover came, at dead of night, the unbelieving father sent to his son the Arian Bishop, to offer him, as the price of his favour, to receive at the hands of the said Bishop the Communion which was the result of a sacrilegious consecration. But when the Arian Bishop came into the prison, the servant of God, remembering that he was not his own but God's man, rebuked the unbeliever as he deserved, and drave him from his presence with just reproaches for though he was weak and bound as touching this outer body, yet was he strong in the mighty castle of his soul. The Bishop, therefore, went away again to that Arian father. And when he came to Leovigild, he waxed exceeding wroth, and sent his servants to kill God's faithful witness where he lay. Which thing was done for as soon as they came to him into the prison, they clave his head with an axe, and freed him from the dying life of this house of our tabernacle. And so they did to him all that which they that kill the body are able to do, and it was a thing which now of a long season he feared not, seeing that when they have done that, they have no more that they can do, but fearing rather Him Who, when He hath killed, hath power to cast both body and soul into hell. But God, to make manifest the glory of His servant, was pleased to work signs from heaven, for of a sudden the solemn swell of singing of Psalms was heard at that dead hour of night from round about the place where lay the body of the kingly martyr, kingly now in an higher and truer sense than the sense of earthly kingship, since he had witnessed a good confession for the truth, sealing it with his blood.\par
\switchcolumn*

\LA
\begin{Resp}\Rsign In servis suis, allelúja, \Star{} Consolábitur Deus, allelúja.\end{Resp}
\begin{Resp}\Vsign Judicábit Dóminus pópulum suum, et in servis suis.\end{Resp}
\begin{Resp}\Rsign Consolábitur Deus, allelúja.\end{Resp}
\switchcolumn
\EN
\begin{Resp}\Rsign His servants, alleluia \Star{} God will comfort, alleluia, alleluia.\end{Resp}
\begin{Resp}\Vsign The Lord will judge His people and, His servants,\end{Resp}
\begin{Resp}\Rsign God will comfort, alleluia, alleluia.\end{Resp}
\switchcolumn*

\LA
\LessonHead{Lectio VI}
\switchcolumn
\EN
\LessonHead{Lesson VI}
\switchcolumn*

\LA
\noindent Quidam étiam ferunt, quod illic nocturno tempore accensæ lampades apparebant; unde et factum est, quatenus corpus illius, ut videlicet Martyris, jure a cunctis fidelibus venerari debuisset. Pater vero perfidus et parricida, commotus pœnitentia, hoc fecisse se doluit, nec tamen usque ad obtinendam salutem pœnituit. Nam quia vera esset catholica fides agnovit, sed gentis suæ timore perterritus, ad hanc pervenire non meruit. Qui, oborta ægritudine, ad extrema perductus est, et Leandro episcopo, quem prius vehementer afflixerat, Reccaredum regem filium suum, quem in sua hæresi relinquebat, commendare curavit, ut in ipso quoque talia faceret, qualia et in fratre suis exhortationibus fecisset. Qua commendatione expleta, defunctus est. Post cujus mortem Reccaredus rex non patrem perfidum, sed fratrem Martyrem sequens, ab Arianæ hæreseos pravitate conversus est, totamque Visigothorum gentem ita ad veram perduxit fidem, ut nullum in suo regno militare permitteret, qui regni Dei hostis exsistere per hæreticam pravitatem non timeret Nec mirum quod veræ fidei prædicator factus est, qui frater est Martyris: cujus hunc quoque merita adjuvant, ut ad omnipotentis Dei gremium tam multos reducat.\par
\switchcolumn
\EN
\noindent Come say, too, that lights were seen there that night. Wherefore it came to pass that the body of the martyr became the rightful object of reverence to all God's faithful people. The unbelieving father, murderer of his own child, was seized with remorse, and repented him of what he had done, but he sorrowed not unto salvation. For though he knew that the Catholic faith was true, he stood in fear of his people, and deserved not to attain unto it. He fell sick, and, when he was at the point of death, he made it his duty to recommend King Reccared his surviving son to the care of the Bishop Leander, whom aforetime he had grievously persecuted, that though Reccared was now left in heresy, the Bishop might work in him by his exhortations the same change that he had worked in his brother, which when Leovigild had said, he died. After his death, King Reccared took for his example not his unbelieving father but his martyred brother. He forsook the Arian heresy, and brought the whole nation of the Visigoths to believe in the true faith, so that he allowed no man in his kingdom to be an officer, who dared any longer range himself through heresy as an enemy of the Kingdom of God. Neither need we marvel that Reccared was a preacher of the faith, since he had had to his brother a martyr, for whose sake Almighty God hath helped him to bring back so many to the bosom of their Father Who is in heaven.\par
\switchcolumn*

\LA
\begin{Resp}\Rsign Fíliæ Jerúsalem, veníte et vidéte Mártyres cum corónis, quibus coronávit eos Dóminus \Star{} In die solemnitátis et lætítiæ, allelúja.\end{Resp}
\begin{Resp}\Vsign Quóniam confortávit seras portárum tuárum, benedíxit fílios tuos in te.\end{Resp}
\begin{Resp}\Rsign In die solemnitátis et lætítiæ, allelúja.\end{Resp}
\begin{Resp}\Vsign Glória Patri, et Fílio, \Star{} et Spirítui Sancto.\end{Resp}
\begin{Resp}\Rsign In die solemnitátis et lætítiæ, allelúja.\end{Resp}
\switchcolumn
\EN
\begin{Resp}\Rsign Come forth, O ye daughters of Jerusalem, and behold the Martyrs with the crowns wherewith the Lord crowned them; \Star{} In the day of His feasting, and of His gladness. Alleluia.\end{Resp}
\begin{Resp}\Vsign For He hath strengthened the bars of thy gates He hath blessed thy children within thee.\end{Resp}
\begin{Resp}\Rsign In the day of His feasting, and of His gladness. Alleluia.\end{Resp}
\begin{Resp}\Vsign Glory be to the Father, and to the Son, \Star{} and to the Holy Ghost.\end{Resp}
\begin{Resp}\Rsign In the day of His feasting, and of His gladness. Alleluia.\end{Resp}
\switchcolumn*

\LA
\LessonHead{Lectio VII}
\switchcolumn
\EN
\LessonHead{Lesson VII}
\switchcolumn*

\LA
\LessonTitle{Léctio sancti Evangélii secúndum Lucam}
\switchcolumn
\EN
\LessonTitle{From the Holy Gospel according to Luke}
\switchcolumn*

\LA
\Cite{Luc 14:26-33}
\switchcolumn
\EN
\Cite{Luke 14:26-33}
\switchcolumn*

\LA
\noindent In illo témpore: Dixit Jesus turbis: Si quis venit ad me, et non odit patrem suum, et matrem, et uxórem, et fílios, et fratres, et soróres, adhuc autem et ánimam suam, non potest meus esse discípulus. Et réliqua.\par
\switchcolumn
\EN
\noindent At that time, Jesus said unto the multitudes: If any man come to Me, and hate not his father, and mother, and wife, and children, and brethren, and sisters, yea, and his own life also, he cannot be My disciple. And so on.\par
\switchcolumn*

\LA
\LessonTitle{Homilía sancti Gregórii Papæ}
\switchcolumn
\EN
\LessonTitle{Homily by Pope St Gregory the Great}
\switchcolumn*

\LA
\Source{Homil. 37 in Evangelia}
\switchcolumn
\EN
\Source{Homilia 37th on the Gospels.}
\switchcolumn*

\LA
\noindent Si considerémus fratres caríssimi, quæ et quanta sunt, quæ nobis promittúntur in cælis, viléscunt ánimo ómnia quæ habéntur in terris. Terréna namque substántia supérnæ felicitáti comparáta, pondus est, non subsídium. Temporális vita ætérnæ vitæ comparáta, mors est pótius dicénda quam vita. Ipse enim quotidiánus deféctus corruptiónis quid est áliud, quam quædam prolíxitas mortis? Quæ autem lingua dícere, vel quis intelléctus cápere súfficit, illa supérnæ civitátis quanta sint gáudia; Angelórum choris interésse, cum beatíssimis spirítibus glóriæ Conditóris assístere, præséntem Dei vultum cérnere, incircumscríptum lumen vidére, nullo mortis metu áffici, incorruptiónis perpétuæ múnere lætári?\par
\switchcolumn
\EN
\noindent Dearly beloved brethren, if we consider what and how great things are promised unto us in heaven, all things which are upon earth grow poor to our mind. For when this world's goods are reckoned against the gladness above, they are found to be a clog rather than an help. This present life being compared to life eternal, ought rather to be called death than life. For what is the daily failing of our corruption but, as it were, a creeping death But what tongue is there that can tell, or what understanding that can comprehend how great is the rejoicing in the city above, where they have part with the choirs of Angels, where they stand with the most blessed spirits before the glory of the Creator, where they see the face of God present, where they behold the Incomprehensible Light, where they have no fear of death, and where they rejoice eternally incorruptible\par
\switchcolumn*

\LA
\begin{Resp}\Rsign Ego sum vitis vera, et vos pálmites: \Star{} Qui manet in me, et ego in eo, hic fert fructum multum, allelúja, allelúja.\end{Resp}
\begin{Resp}\Vsign Sicut diléxit me Pater, et ego diléxi vos.\end{Resp}
\begin{Resp}\Rsign Qui manet in me, et ego in eo, hic fert fructum multum, allelúja, allelúja.\end{Resp}
\switchcolumn
\EN
\begin{Resp}\Rsign I am the vine: you the branches. \Star{} He that abideth in me, and I in him, the same beareth much fruit, alleluia, alleluia.\end{Resp}
\begin{Resp}\Vsign As the Father hath loved me, I also have loved you.\end{Resp}
\begin{Resp}\Rsign He that abideth in me, and I in him, the same beareth much fruit, alleluia, alleluia.\end{Resp}
\switchcolumn*

\LA
\LessonHead{Lectio VIII}
\switchcolumn
\EN
\LessonHead{Lesson VIII}
\switchcolumn*

\LA
\noindent Sed ad hæc audíta inardéscit ánimus, jamque illic cupit assístere, ubi se sperat sine fine gaudére. Sed ad magna prǽmia perveníri non potest, nisi per magnos labóres. Unde et Paulus egrégius prædicátor dicit: Non coronábitur, nisi qui legítime certáverit. Deléctet ergo mentem magnitúdo præmiórum, sed non detérreat certámen labórum. Unde ad se veniéntibus Véritas dicit: Si quis venit ad me, et non odit patrem suum, et matrem, et uxórem, et fílios, et fratres, et soróres, adhuc autem et ánimam suam, non potest meus esse discípulus.\par
\switchcolumn
\EN
\noindent Then we hear these things our hearts burn within us and we long to be already there, where we hope to rejoice for ever. But we cannot attain unto great rewards, save through great labour. Therefore saith the excellent preacher Paul: “He is not crowned, except he strive lawfully.” (2 Tim. ii. 5.) The greatness of the reward doth delight our mind let not the throes of the struggle dishearten us. Therefore the Truth saith unto every one that cometh unto Him “If any man come to Me, and hate not his father and mother, and wife, and children, and brethren, and sisters, yea, and his own life also, he cannot be My disciple.”\par
\switchcolumn*

\LA
\begin{Resp}\Rsign Cándidi facti sunt Nazarǽi ejus, allelúja: splendórem Deo dedérunt, allelúja: \Star{} Et sicut lac coaguláti sunt, allelúja, allelúja.\end{Resp}
\begin{Resp}\Vsign Candidióres nive, nitidióres lacte, rubicundióres ébore antíquo, sapphíro pulchrióres.\end{Resp}
\begin{Resp}\Rsign Et sicut lac coaguláti sunt, allelúja, allelúja.\end{Resp}
\begin{Resp}\Vsign Glória Patri, et Fílio, \Star{} et Spirítui Sancto.\end{Resp}
\begin{Resp}\Rsign Et sicut lac coaguláti sunt, allelúja, allelúja.\end{Resp}
\switchcolumn
\EN
\begin{Resp}\Rsign Her Nazarites are become pure, Alleluia: they reflect the glory of God, Alleluia. \Star{} They are whiter than milk. Alleluia, Alleluia.\end{Resp}
\begin{Resp}\Vsign They are purer than snow, they are whiter than milk, they are more ruddy in body than coral, their polishing is of sapphire.\end{Resp}
\begin{Resp}\Rsign They are whiter than milk. Alleluia, Alleluia.\end{Resp}
\begin{Resp}\Vsign Glory be to the Father, and to the Son, \Star{} and to the Holy Ghost.\end{Resp}
\begin{Resp}\Rsign They are whiter than milk. Alleluia, Alleluia.\end{Resp}
\switchcolumn*

\LA
\LessonHead{Lectio IX}
\switchcolumn
\EN
\LessonHead{Lesson IX}
\switchcolumn*

\LA
\noindent Sed percontári libet, quómodo paréntes et carnáliter propínquos præcípimur odísse, qui jubémur et inimícos dilígere? Et certe Véritas de uxóre dicit: Quod Deus conjúnxit, homo non séparet. Et Paulus ait: Viri, dilígite uxóres vestras, sicut et Christus Ecclésiam. Ecce, discípulus uxórem diligéndam prǽdicat, cum magíster dicat: Qui uxórem non odit, non potest meus esse discípulus. Numquid áliud judex núntiat, áliud præco clamat? An simul et odísse póssumus, et dilígere? Sed si vim præcépti perpéndimus, utrúmque ágere per discretiónem valémus: ut uxórem, et eos, qui nobis carnis cognatióne conjúncti sunt, et quos próximos nóvimus, diligámus; et quos adversários in via Dei pátimur, odiéndo et fugiéndo nesciámus.\par
\switchcolumn
\EN
\noindent Or it may be asked how we are commanded in one place to hate our parents, and them that are near us in the flesh, and in another place to love even our enemies. And, verily, the Truth hath said, as touching a wife “What God hath joined together, let not man put asunder.” (Matth. xix. 6.) And Paul saith “ Husbands, love your wives, even as Christ also loved the Church.” (Eph. v. 25.) Behold, the disciple commandeth a man to love his wife, and the Master saith “If any man hate not his wife, he cannot be My disciple.” Doth the judge, then, order one proclamation, and the crier make another? or can the man both love and hate? If we consider well the force of the commandment, we shall be able in wisdom to do both. Let us love wife, and kindred, and neighbour, as touching their nearness in the flesh but as touching the way of God, if they withstand us therein, let us not know them, but hate them and flee from them.\par
\switchcolumn*

\LA
\TeDeum{Te Deum}
\switchcolumn
\EN
\TeDeum{Te Deum}
\switchcolumn*

\LA
\LessonHead{Lectio IX (pro commemoratione)}
\switchcolumn
\EN
\LessonHead{Lesson IX (when commemorated)}
\switchcolumn*

\LA
\LessonTitle{Commemoratio: S. Hermenegildi Martyris}
\switchcolumn
\EN
\LessonTitle{Commemoration of: St. Hermenegild, Martyr}
\switchcolumn*

\LA
\Source{Liber 3, cap. 31}
\switchcolumn
\EN
\Source{Book 3, chap. 31}
\switchcolumn*

\LA
\noindent Ex libro Dialogórum sancti Gregórii Papæ\par
\switchcolumn
\EN
\noindent From the Book of Dialogues written by Pope St. Gregory\par
\switchcolumn*

\LA
Hermenegíldus rex, Leovigíldi regis Visigothórum fílius, ab Ariána hǽresi ad fidem cathólicam viro reverendíssimo Leándro Hispalénsi epíscopo, dudum mihi in amicítiis familiáriter juncto, prædicánte, convérsus est. Quem pater Ariánus, ut ad eándem hǽresim redíret, et prǽmiis suadére, et minis terrére conátus est. Cumque ille constantíssime respondéret, nunquam se veram fidem posse relínquere, quam semel agnovísset: irátus pater eum privávit regno, rebúsque exspoliávit ómnibus; et in arcta illum custódia conclúdens, collum manúsque illíus ferro ligávit. Cœpit ítaque Hermenegíldus rex júvenis terrénum regnum despícere, et forti desidério cæléste quærens, in cilíciis vinculátus jacens, omnipoténti Deo ad confortándum se preces effúndere. Superveniénte autem Paschális festivitátis die, intempéstæ noctis siléntio ad eum pérfidus pater Ariánum epíscopum misit, ut ex ejus manu sacrílegæ consecratiónis communiónem percíperet, atque per hoc ad patris grátiam redíre mererétur. Sed vir Deo déditus, Ariáno epíscopo veniénti exprobrávit, ut débuit, ejúsque a se perfídiam dignis increpatiónibus répulit. Ad se ítaque revérso epíscopo, Ariánus pater infrémuit, statímque suos apparitóres misit, qui constantíssimum Confessórem Dei illic, ubi jacébat, occidérunt.\par
\switchcolumn
\EN
King Hermenegild, son of Leovigild, King of the Visigoths, was converted from the Arian heresy to the Catholic faith by the preaching of that most worthy man, Leander, Bishop of Seville, and one of my greatest friends. Hermenegild's father was an Arian, and tried first to persuade him by promises and then to terrify him by threats to return to that heresy. But when Hermenegild continued to answer that he could never abandon the true faith now that he had come to know it, his angry father deprived him of his kingdom, took away all his possessions, and shut him up under strict guard, with his neck and hands chained. And so the young King Hermenegild began to despise the kingdom of earth and eagerly to seek the kingdom of heaven. Lying bound and in sackcloth, he poured forth prayers to Almighty God to give him strength. When Easter came, his treacherous father sent an Arian bishop to him in the middle of the night, to have Hermenegild receive sacrilegiously consecrated Communion from this bishop's hands and so be restored to his father's favour. But, as a man devoted to God, he gave the Arian bishop the rebuke he deserved and refused his treacherous offer with fitting indignation. When the bishop returned, Hermenegild's Arian father was enraged and at once sent his servants, who killed this staunch Confessor of God where he lay imprisoned.\par
\switchcolumn*

\end{paracol}

\par\needspace{10\baselineskip}\addvspace{1.2em}{\centering\small\textsc{Die 14 Aprilis} · \textsc{14 April}\par}
\Office{S. Justini Martyris}{St. Justini Martyr}{Duplex}
\addcontentsline{toc}{subsection}{Die 14 Aprilis · S. Justini Martyris\enspace\textit{St. Justini Martyr}}
\begin{paracol}{2}
\LA
\RefLine{Lectiones I–III: de Scriptura occurrente.}
\switchcolumn
\EN
\RefLine{Lessons I–III: from the Scripture of the day.}
\switchcolumn*

\LA
\RefLine{Lectio IX: de Ss. Tiburtii, Valeriani, et Maximi Martyrum (04-14t).}
\switchcolumn
\EN
\RefLine{Lesson IX: of Sts. Tiburtius, Valerian et Maximus, Martyrs (04-14t).}
\switchcolumn*

\LA
\LessonHead{Lectio I}
\switchcolumn
\EN
\LessonHead{Lesson I}
\switchcolumn*

\LA
\LessonTitle{De Epístola beáti Pauli Apóstoli ad Romános}
\switchcolumn
\EN
\LessonTitle{Lesson from the letter of St Paul the Apostle to the Romans}
\switchcolumn*

\LA
\Cite{Rom 8:12-19}
\switchcolumn
\EN
\Cite{Rom 8:12-19}
\switchcolumn*

\LA
\noindent Fratres: Debitóres sumus non carni, ut secúndum carnem vivámus. Si enim secúndum carnem vixéritis, moriémini: si autem spíritu facta carnis mortificavéritis, vivétis. Quicúmque enim spíritu Dei agúntur, ii sunt fílii Dei. Non enim accepístis spíritum servitútis íterum in timóre, sed accepístis spíritum adoptiónis filiórum, in quo clamámus: Abba (Pater). Ipse enim Spíritus testimónium reddit spirítui nostro quod sumus fílii Dei. Si autem fílii, et herédes: herédes, quidem Dei, coherédes autem Christi: si tamen compátimur ut et conglorificémur. Exístimo enim quod non sunt condígnæ passiónes hujus témporis ad futúram glóriam, quæ revelábitur in nobis. Nam exspectátio creatúræ revelatiónem filiórum Dei exspéctat.\par
\switchcolumn
\EN
\noindent Therefore, brethren, we are debtors, not to the flesh, to live according to the flesh. For if you live according to the flesh, you shall die: but if by the Spirit you mortify the deeds of the flesh, you shall live. For whosoever are led by the Spirit of God, they are the sons of God. For you have not received the spirit of bondage again in fear; but you have received the spirit of adoption of sons, whereby we cry: Abba, Father. For the Spirit himself giveth testimony to our spirit, that we are the sons of God. And if sons, heirs also; heirs indeed of God, and joint heirs with Christ: yet so, if we suffer with him, that we may be also glorified with him. For I reckon that the sufferings of this time are not worthy to be compared with the glory to come, that shall be revealed in us. For the expectation of the creature waiteth for the revelation of the sons of God\par
\switchcolumn*

\LA
\begin{Resp}\Rsign Iste Sanctus pro lege Dei sui certávit usque ad mortem, et a verbis impiórum non tímuit: \Star{} Fundátus enim erat supra firmam petram.\end{Resp}
\begin{Resp}\Vsign Iste est qui contémpsit vitam mundi, et pervénit ad cæléstia regna.\end{Resp}
\begin{Resp}\Rsign Fundátus enim erat supra firmam petram.\end{Resp}
\switchcolumn
\EN
\begin{Resp}\Rsign This man is holy, for he hath striven for the law of his God even unto death, and hath not feared for the words of the ungodly; \Star{} For he had his foundation upon a strong rock.\end{Resp}
\begin{Resp}\Vsign This is he which loved not his life in this world, and is come unto an everlasting kingdom.\end{Resp}
\begin{Resp}\Rsign For he had his foundation upon a strong rock.\end{Resp}
\switchcolumn*

\LA
\LessonHead{Lectio II}
\switchcolumn
\EN
\LessonHead{Lesson II}
\switchcolumn*

\LA
\Cite{Rom 8:28-34}
\switchcolumn
\EN
\Cite{Rom 8:28-34}
\switchcolumn*

\LA
\noindent Scimus autem quóniam diligéntibus Deum ómnia cooperántur in bonum iis qui secúndum propósitum vocáti sunt sancti. Nam quos præscívit et prædestinávit confórmes fíeri imáginis Fílii sui ut sit ipse primogénitus in multis frátribus. Quos autem prædestinávit, hos et vocávit: et quos vocávit, hos et justificávit: quos autem justificávit, illos et glorificávit. Quid ergo dicémus ad hæc? Si Deus pro nobis, quis contra nos? Qui étiam próprio Fílio suo non pepércit, sed pro nobis ómnibus trádidit illum, quómodo non étiam cum illo ómnia nobis donávit? Quis accusábit advérsus eléctos Dei? Deus qui justíficat, quis est qui condémnet? Christus Jesus, qui mórtuus est, immo qui et resurréxit, qui est ad déxteram Dei, qui étiam interpéllat pro nobis.\par
\switchcolumn
\EN
\noindent And we know that to them that love God, all things work together unto good, to such as, according to his purpose, are called to be saints. For whom he foreknew, he also predestinated to be made conformable to the image of his Son; that he might be the firstborn amongst many brethren. And whom he predestinated, them he also called. And whom he called, them he also justified. And whom he justified, them he also glorified. What shall we then say to these things? If God be for us, who is against us? He that spared not even his own Son, but delivered him up for us all, how hath he not also, with him, given us all things? Who shall accuse against the elect of God? God that justifieth. Who is he that shall condemn? Christ Jesus that died, yea that is risen also again; who is at the right hand of God, who also maketh intercession for us.\par
\switchcolumn*

\LA
\begin{Resp}\Rsign Justus germinábit sicut lílium: \Star{} Et florébit in ætérnum ante Dóminum.\end{Resp}
\begin{Resp}\Vsign Plantátus in domo Dómini, in átriis domus Dei nostri.\end{Resp}
\begin{Resp}\Rsign Et florébit in ætérnum ante Dóminum.\end{Resp}
\switchcolumn
\EN
\begin{Resp}\Rsign The righteous shall grow as the lily \Star{} Yea, he shall flourish in the presence of the Lord for ever.\end{Resp}
\begin{Resp}\Vsign Those that be planted in the house of the Lord, shall flourish in the courts of the house of our God.\end{Resp}
\begin{Resp}\Rsign Yea, he shall flourish in the presence of the Lord for ever.\end{Resp}
\switchcolumn*

\LA
\LessonHead{Lectio III}
\switchcolumn
\EN
\LessonHead{Lesson III}
\switchcolumn*

\LA
\Cite{Rom 8:35-39}
\switchcolumn
\EN
\Cite{Rom 8:35-39}
\switchcolumn*

\LA
\noindent Quis ergo nos separábit a caritáte Christi? tribulátio? an angústia? an fames? an núditas? an perículum? an persecútio? an gládius? (sicut scriptum est: Quia propter te mortificámur tota die: æstimáti sumus sicut oves occisiónis.) Sed in his ómnibus superámus propter eum qui diléxit nos. Certus sum enim quia neque mors, neque vita, neque Ángeli, neque Principátus, neque Virtútes, neque instántia, neque futúra, neque fortitúdo, neque altitúdo, neque profúndum, neque creatúra ália póterit nos separáre a caritáte Dei, quæ est in Christo Jesu Dómino nostro.\par
\switchcolumn
\EN
\noindent Who then shall separate us from the love of Christ? Shall tribulation? or distress? or famine? or nakedness? or danger? or persecution? or the sword? As it is written: For thy sake we are put to death all the day long. We are accounted as sheep for the slaughter. But in all these things we overcome, because of him that hath loved us. For I am sure that neither death, nor life, nor angels, nor principalities, nor powers, nor things present, nor things to come, nor might, Nor height, nor depth, nor any other creature, shall be able to separate us from the love of God, which is in Christ Jesus our Lord.\par
\switchcolumn*

\LA
\begin{Resp}\Rsign Iste cognóvit justítiam, et vidit mirabília magna, et exorávit Altíssimum \Star{} Et invéntus est in número Sanctórum.\end{Resp}
\begin{Resp}\Vsign Iste est qui contémpsit vitam mundi, et pervénit ad cæléstia regna.\end{Resp}
\begin{Resp}\Rsign Et invéntus est in número Sanctórum.\end{Resp}
\begin{Resp}\Vsign Glória Patri, et Fílio, \Star{} et Spirítui Sancto.\end{Resp}
\begin{Resp}\Rsign Et invéntus est in número Sanctórum.\end{Resp}
\switchcolumn
\EN
\begin{Resp}\Rsign This is he which knew righteousness, and saw great wonders, and made his prayer unto the Most High \Star{} And he is numbered among the Saints.\end{Resp}
\begin{Resp}\Vsign This is he which loved not his life in this world, and is come unto an everlasting kingdom.\end{Resp}
\begin{Resp}\Rsign And he is numbered among the Saints.\end{Resp}
\begin{Resp}\Vsign Glory be to the Father, and to the Son, \Star{} and to the Holy Ghost.\end{Resp}
\begin{Resp}\Rsign And he is numbered among the Saints.\end{Resp}
\switchcolumn*

\LA
\LessonHead{Lectio IV}
\switchcolumn
\EN
\LessonHead{Lesson IV}
\switchcolumn*

\LA
\noindent Justínus, Prisci fílius, ex Græco genere Flaviæ Neápolis in Syria Palæstina natus, adolescéntiam in litterárum ómnium stúdiis transegit. Vir factus adeo philosophíæ amóre correptus est, ut ad veritátem assequéndam, quotquot aderant, philosophórum sectis nomen déderit, eorúmque præcépta scrutátus sit. Cum in his fallacem tantum sapiéntiam errorémque reperísset, superna illustratióne per senem quemdam ignotum aspectuque venerábilem edoctus, veræ christianæ fídei philosophíam amplexus est. Hinc sacræ Scripturæ libros diu noctuque præ mánibus habens, ita ex eórum meditatióne divinus ignis in ánima ejus exársit, ut ea qua pollébat eruditiónis vi, eminentem Jesu Christi sciéntiam adeptus, plurima conscripserit volúmina ad christiánam fidem exponéndam magisque propagandam.\par
\switchcolumn
\EN
\noindent Justin, the son of Priscus, was a Greek by race, but was born at Nablus in Palestine. He passed his youth in the study of letters. When he became a man he was so taken with the love of philosophy and the desire of truth that he became a student in the schools of all the philosophers and examined the teaching of them all. In them he found only deceitful wisdom and error. The light of heaven was given him, through an old man of worshipful aspect whom he knew not, and he embraced the philosophy of the true Christian faith. Henceforth he had the books of the Holy Scriptures in his hands by day and by night, and by meditating thereon the fire of God was so kindled in his soul that, himself possessing the excellency of the knowledge of Christ Jesus our Lord, he wrote many books, with all the learning which he possessed, to set forth the Christian faith and to spread it abroad.\par
\switchcolumn*

\LA
\begin{Resp}\Rsign Lux perpétua lucébit Sanctis tuis, Dómine, \Star{} Et ætérnitas témporum, allelúja, allelúja.\end{Resp}
\begin{Resp}\Vsign Lætítia sempitérna erit super cápita eórum: gáudium et exsultatiónem obtinébunt.\end{Resp}
\begin{Resp}\Rsign Et ætérnitas témporum, allelúja, allelúja.\end{Resp}
\switchcolumn
\EN
\begin{Resp}\Rsign Eternal light will shine over your saints, O Lord, \Star{} And the eternity of the times, alleluia, alleluia.\end{Resp}
\begin{Resp}\Vsign Everlasting joy shall be upon their heads, they shall obtain joy and gladness\end{Resp}
\begin{Resp}\Rsign And the eternity of the times, alleluia, alleluia.\end{Resp}
\switchcolumn*

\LA
\LessonHead{Lectio V}
\switchcolumn
\EN
\LessonHead{Lesson V}
\switchcolumn*

\LA
\noindent Inter præclaríssima Justini ópera binæ eminent fídei christianæ apologíæ, quas cum coram senatu, imperatóribus Antonino Pio ejúsque fíliis nec non Marco Antonino Vero et Lucio Aurelio Cómmodo Christi ásseclas sævíssime divexántibus, porrexísset, eamdemque fidem disputándo strenue propugnasset, obtinuit, ut a Christianórum cæde publico príncipum edícto temperátum fúerit. Verum Justino haud parcitum est. Nam Crescéntis cynici, cujus vitam et mores nefarios redargúerat, insídiis accusatus, a satellítibus comprehénsus est. Adductus autem ad Romæ præsidem nómine Rústicum, cum hic ab eo quæsivísset, quænam essent Christianórum præcépta, hanc bonam confessiónem coram multis testibus conféssus est: Rectum dogma, quod nos christiáni hómines cum pietáte servamus, hoc est: ut Deum unum existimémus Factórem atque Creatórem ómnium quæ vidéntur, quæque corpóreis óculis non cernúntur; et Dóminum Jesum Christum Dei Fílium confiteámur, olim a prophétis prænuntiátum, qui et humani generis Judex ventúrus est.\par
\switchcolumn
\EN
\noindent Amongthe most famous of the works of Justin are his two Apologies or Defences of the Christian faith. These he brought before the Senate when the Emperors Antoninus Pius, and his sons, as also Marcus Antoninus Verus and Lucius Aurelius Commodus, were savagely persecuting the followers of Christ, and by their means, and his vigorous disputations in favour of the same faith, he obtained a public edict from the government to stay the slaughter of the Christians. But Justin himself did not escape he had rebuked the life and infamous manners of the Cynic Crescens, and was accused and arrested through that person's plottings. He was brought before Rusticus, the President of Rome, who asked him what were the doctrines of the Christians, whereto he answered, in the presence of many witnesses, with this good confession: The right doctrine which we Christian men do keep with godliness is this, that we should believe that there is one God, Who is the Maker and Creator of all things, both those things which are seen and those things which bodily eyes do not see, and that we should confess the Lord Jesus Christ, the Son of God, Who was foretold of old time by the prophets, and Who will come to be the Judge of all mankind.\par
\switchcolumn*

\LA
\begin{Resp}\Rsign In servis suis, allelúja, \Star{} Consolábitur Deus, allelúja.\end{Resp}
\begin{Resp}\Vsign Judicábit Dóminus pópulum suum, et in servis suis.\end{Resp}
\begin{Resp}\Rsign Consolábitur Deus, allelúja.\end{Resp}
\switchcolumn
\EN
\begin{Resp}\Rsign His servants, alleluia \Star{} God will comfort, alleluia, alleluia.\end{Resp}
\begin{Resp}\Vsign The Lord will judge His people and, His servants,\end{Resp}
\begin{Resp}\Rsign God will comfort, alleluia, alleluia.\end{Resp}
\switchcolumn*

\LA
\LessonHead{Lectio VI}
\switchcolumn
\EN
\LessonHead{Lesson VI}
\switchcolumn*

\LA
\noindent Quóniam Justinus in prima sua apología exposúerat quómodo Christiáni convenírent ad Sacra celebranda, et quænam fúerint sacri hujus convéntus mysteria, ad repelléndas ethnicórum calumnias; exquisívit ab eo præses, in quonam loco conveníret ipse et ceteri hujus Urbis Christifidéles. Justinus autem réticens convéntuum loca, ne sancta et fratres próderet cánibus, domicílium tantum suum indicávit, ubi manére et discípulos excólere solébat penes célebrem títulum Pastoris in ædibus Pudéntis. Demum præses optiónem ei dedit vel ut diis sacrificaret, vel per totum corpus flagellis cædi perferret. Cum invictus fídei vindex assereret se in votis semper habuísse cruciátus pérpeti propter Dóminum Jesum Christum, a quo magnam in cælis mercédem cónsequi expectabat, præses in eum capitálem senténtiam pronuntiávit. Itaque mirábilis philósophus Deum collaudans, post vérbera, fuso pro Christo sánguine, glorióso martyrio coronátus est. Quidam vero fidéles clam illíus sustulérunt corpus, et in loco idoneo condidérunt. Leo décimus tertius Póntifex maximus ejúsdem Offícium et Missam ab univérsa Ecclésia celebrári præcepit.\par
\switchcolumn
\EN
\noindent In order to rebut the slanders of the heathen, Justin had in his first Apology given an open account of the gathering of the Christians for divine worship, and what were the holy Mysteries celebrated in these assemblies. The President therefore asked him what was the place where he and Christ's other faithful ones in the city were accustomed to meet. Justin, lest he should betray that which was holy unto God and his brethren, told only where was his own lodging, where he was used to abide and to teach his disciples, hard by the famous Church of the Shepherd, in the house of Pudens. The President then gave him the choice whether to sacrifice to the gods or to be hided with scourges over his whole body. The unconquered champion of the faith answered that he had always desired to suffer in the Name of the Lord Jesus Christ, from Whom he looked to receive a mighty reward in heaven. The President thereupon sentenced him to death, and then this excellent philosopher, giving praise to God, was first beaten and afterwards shed his blood for Christ's sake, and so received the crown of a glorious martyrdom. Some of the faithful secretly stole away his body, and buried it in a fitting place. The Supreme Pontiff Leo XIII. commanded that his Office and Mass should be used throughout the whole Church.\par
\switchcolumn*

\LA
\begin{Resp}\Rsign Fíliæ Jerúsalem, veníte et vidéte Mártyres cum corónis, quibus coronávit eos Dóminus \Star{} In die solemnitátis et lætítiæ, allelúja.\end{Resp}
\begin{Resp}\Vsign Quóniam confortávit seras portárum tuárum, benedíxit fílios tuos in te.\end{Resp}
\begin{Resp}\Rsign In die solemnitátis et lætítiæ, allelúja.\end{Resp}
\begin{Resp}\Vsign Glória Patri, et Fílio, \Star{} et Spirítui Sancto.\end{Resp}
\begin{Resp}\Rsign In die solemnitátis et lætítiæ, allelúja.\end{Resp}
\switchcolumn
\EN
\begin{Resp}\Rsign Come forth, O ye daughters of Jerusalem, and behold the Martyrs with the crowns wherewith the Lord crowned them; \Star{} In the day of His feasting, and of His gladness. Alleluia.\end{Resp}
\begin{Resp}\Vsign For He hath strengthened the bars of thy gates He hath blessed thy children within thee.\end{Resp}
\begin{Resp}\Rsign In the day of His feasting, and of His gladness. Alleluia.\end{Resp}
\begin{Resp}\Vsign Glory be to the Father, and to the Son, \Star{} and to the Holy Ghost.\end{Resp}
\begin{Resp}\Rsign In the day of His feasting, and of His gladness. Alleluia.\end{Resp}
\switchcolumn*

\LA
\LessonHead{Lectio VII}
\switchcolumn
\EN
\LessonHead{Lesson VII}
\switchcolumn*

\LA
\LessonTitle{Léctio sancti Evangélii secúndum Lucam.}
\switchcolumn
\EN
\LessonTitle{From the Holy Gospel according to Luke}
\switchcolumn*

\LA
\Cite{Luc 12:2-8}
\switchcolumn
\EN
\Cite{Luke 12:2-8}
\switchcolumn*

\LA
\noindent In illo témpore: Dixit Jesus discípulis suis: Nihil opertum est, quod non revelétur; neque abscónditum, quod non sciátur. Et réliqua.\par
\switchcolumn
\EN
\noindent In the meantime, when there were gathered together an innumerable multitude of people, in so much that they trod one upon another, Jesus began to say unto His disciples first of all, Beware ye of the leaven of the Pharisees, which is hypocrisy, for there is nothing covered, that shall not be revealed; neither hid, that shall not be known. And so on.\par
\switchcolumn*

\LA
\LessonTitle{Homilía sancti Joánnis Chrysostomi.}
\switchcolumn
\EN
\LessonTitle{Homily by St. John Chrysostom, Patriarch of Constantinople)}
\switchcolumn*

\LA
\Source{Homilia in cap. x Matth. v. 26. et seq.}
\switchcolumn
\EN
\Source{On Matthew x. 26.}
\switchcolumn*

\LA
\noindent Nihil est opertum quod non revelábitur, nec occultum quod non sciétur. Quod autem dicit, hujúsmodi est: Sufficit quidem vobis ad consolatiónem, si ego Magister et Dóminus consors sim conviciórum. Si vero adhuc doletis hæc audiéntes, illud quoque animo reputate, vos non multum póstea ab hac suspicióne liberátum iri. Cur enim id ægre fertis? quia præstigiatores et deceptores vos vocant? At paululum exspectáte, et servatores benefactoresque orbis vos prædicábunt omnes. At enim tempus illa ómnia, quæ subobscura erant, revelabit, et illórum calumniam déteget, virtutémque vestram conspicuam reddet. Cum enim ex rebus ipsis comprobabímini salvatores esse et benefici, et omni virtúte conspicui, illórum dictis hómines non attendent, sed rei veritáti; ac illi quidam sycophantæ, mendáces, malédici, vos vero ipso sole splendidióres deprehendémini. Multum quippe témporis spatium vos notos reddet, prædicabit, et tuba clariórem emíttet vocem, vestræque virtútis testes univérsos hómines exhibétur. Ne ítaque ea, quæ nunc dicúntur, vos dejíciant, sed spes futurórum bonórum érigat. Non possunt enim ea, quæ ad vos spectant, occultári.\par
\switchcolumn
\EN
\noindent The disciple is not above his master, nor the servant above his lord. It is enough for the disciple that he be as his master, and the servant as his lord. If they have called the master of the house Beelzebub, how much more shall they call them of his household? Fear them not, therefore: for] there is nothing covered, that shall not be revealed; and hid, that shall not be known. It is as though He would say: It is comfort enough for you, if I, your Master and Lord, am a partaker in your reproach. But if it grieve you unto this present to hear these things, bethink you likewise that it is but a little while, and ye shall be free from that reproach. For what is it that grieveth you? is it that they call you tricksters and deceivers? Wait but a little while and all men shall call you the preservers and benefactors of the world. In a little while all the things which are dark now shall be made clear, and the falsehood of them that reproach you and your own goodness shall be shown in the light. For when that which cometh to pass shall itself show that ye are preservers and benefactors, and filled with all goodness, men will regard not the words of your gainsayers but the truth. They that now speak evil of you will be found out in the slanderers, liars, and calumniators, and ye shall be seen to be brighter than the sun; time shall make you known and shall preach you with a voice louder than the voice of a trumpet, and shall bring forward all men as the witnesses of your goodness. Let not, therefore, those things which are now spoken cast you down, but rather let the hope of the good things which are to come lift you up. For the things which regard you cannot be hidden.\par
\switchcolumn*

\LA
\begin{Resp}\Rsign Ego sum vitis vera, et vos pálmites: \Star{} Qui manet in me, et ego in eo, hic fert fructum multum, allelúja, allelúja.\end{Resp}
\begin{Resp}\Vsign Sicut diléxit me Pater, et ego diléxi vos.\end{Resp}
\begin{Resp}\Rsign Qui manet in me, et ego in eo, hic fert fructum multum, allelúja, allelúja.\end{Resp}
\switchcolumn
\EN
\begin{Resp}\Rsign I am the vine: you the branches. \Star{} He that abideth in me, and I in him, the same beareth much fruit, alleluia, alleluia.\end{Resp}
\begin{Resp}\Vsign As the Father hath loved me, I also have loved you.\end{Resp}
\begin{Resp}\Rsign He that abideth in me, and I in him, the same beareth much fruit, alleluia, alleluia.\end{Resp}
\switchcolumn*

\LA
\LessonHead{Lectio VIII}
\switchcolumn
\EN
\LessonHead{Lesson VIII}
\switchcolumn*

\LA
\noindent Deínde, postquam illos omni angore, timóre et sollicitúdine liberávit, et probris ómnibus superióres réddidit, demum illos opportune de libertáte prædicándi allóquitur; nam dicit: Quod dico vobis in ténebris, dícite in lúmine; et quod in aure audítis, prædicate super tecta. Quamquam non erant ténebræ cum hæc diceret, neque ad aurem loquebátur: sed hæc hyperbólice dicta sunt. Quia enim solos alloquebátur, et in parvo Palæstinæ angulo, ídeo dicit, In ténebris et In aure; hunc loquéndi modum cómparans cum loquéndi fiducia, qua illos póstea instructurus erat. Ne in una, duabus tribusque civitátibus, sed per totum orbem prædicate, terram máreque peragrántes, habitátam, non habitátam; ac tyrannis, pópulus, philosophis, rhetóribus cum magna fiducia ómnia dícite. Ideo dixit super tecta et in lúmine; sine ullo subterfugio, et cum omni libertáte.\par
\switchcolumn
\EN
\noindent And when He had freed them from pain, fear, and care, and set them above the reproaches of men, He spake unto them in due season concerning the freedom of preaching, What I tell you in darkness, that speak ye in light and what ye hear in the ear, that preach ye upon the house-tops. It was not darkness when He uttered these words, neither was He speaking into their ear. These words were a figure He was speaking to them alone and in a little corner of Palestine, and therefore He saith in darkness and in the ear, as comparing this manner of speech with that boldness of speaking wherewith He was afterwards to inspire them. Preach, He saith, not in one nor two nor three cities, but throughout the whole world go over the earth and the sea, the land that is dwelt in and the land that is not dwelt in speak all things with great boldness to kings and to peoples, to philosophers and to rhetoricians therefore without any subtlety, but with all freedom, He saith, What I tell you in darkness, that speak ye in light and what ye hear in the ear, that preach ye upon the house-tops.\par
\switchcolumn*

\LA
\begin{Resp}\Rsign Cándidi facti sunt Nazarǽi ejus, allelúja: splendórem Deo dedérunt, allelúja: \Star{} Et sicut lac coaguláti sunt, allelúja, allelúja.\end{Resp}
\begin{Resp}\Vsign Candidióres nive, nitidióres lacte, rubicundióres ébore antíquo, sapphíro pulchrióres.\end{Resp}
\begin{Resp}\Rsign Et sicut lac coaguláti sunt, allelúja, allelúja.\end{Resp}
\begin{Resp}\Vsign Glória Patri, et Fílio, \Star{} et Spirítui Sancto.\end{Resp}
\begin{Resp}\Rsign Et sicut lac coaguláti sunt, allelúja, allelúja.\end{Resp}
\switchcolumn
\EN
\begin{Resp}\Rsign Her Nazarites are become pure, Alleluia: they reflect the glory of God, Alleluia. \Star{} They are whiter than milk. Alleluia, Alleluia.\end{Resp}
\begin{Resp}\Vsign They are purer than snow, they are whiter than milk, they are more ruddy in body than coral, their polishing is of sapphire.\end{Resp}
\begin{Resp}\Rsign They are whiter than milk. Alleluia, Alleluia.\end{Resp}
\begin{Resp}\Vsign Glory be to the Father, and to the Son, \Star{} and to the Holy Ghost.\end{Resp}
\begin{Resp}\Rsign They are whiter than milk. Alleluia, Alleluia.\end{Resp}
\switchcolumn*

\LA
\LessonHead{Lectio IX (pro commemoratione)}
\switchcolumn
\EN
\LessonHead{Lesson IX (when commemorated)}
\switchcolumn*

\LA
\LessonTitle{Commemoratio: S. Justini Martyris}
\switchcolumn
\EN
\LessonTitle{Commemoration of: St. Justini Martyr}
\switchcolumn*

\LA
\noindent Justínus, Prisci fílius, ex Græco génere Fláviæ Neápolis in Syria Palæstína natus, ádeo philosophíæ amóre corréptus est, ut ad veritátem assequéndam, quotquot áderant, philosophórum sectis nomen déderit. In quibus tamen cum fallácem tantum sapiéntiam reperísset, supérna illustratióne edóctus, christiánæ fídei philosophíam ampléxus est. Hinc sacræ Scriptúræ libros diu noctúque versans ibíque eminéntem Jesu Christi sciéntiam adéptus, plura conscrípsit ad christiánam fidem exponéndam magísque propagándam; quæ inter binæ præstant pro fide christiána apologíæ. Quas cum imperatóribus Antoníno Pio ejúsque fíliis porrexísset, et fidem disputándo strénue propugnásset, obtínuit, ut a Christianórum cæde público príncipem edícto temperarétur. Ipse tamen Crescéntis Cýnici, cujus et ímpios mores redargúerat, insídiis accusátus, a satellítibus captus est; et ad Rústicum præféctum addúctus, cum in confessióne fídei strénue permanéret, cápitis damnátus, glorióso martýrio coronátus occúbuit.\par
\switchcolumn
\EN
\noindent Justin, son of Priscus and of Greek extraction, was born at Nablus in Palestine. He was so possessed by love for philosophy that, in his search for truth, he enrolled in every philosophical school he could find. But in these he found nothing but fallacious wisdom, and, taught by enlightenment sent him from heaven, he embraced the philosophy of the Christian faith. Thereafter he studied the books of holy Scripture day and night and, when he had thus become learned in the supreme knowledge of Jesus Christ, he wrote many books to explain and to spread the Christian faith, among which his two Apologies are particularly well known. When he had presented these to the emperors Antoninus Pius and his sons, and fought strenuously for the faith by disputations as well, he obtained a public edict from the princes to restrain the slaughter of Christians. But he himself was accused through the plotting of Crescens the Cynic, whose wicked ways he had criticized. He was taken prisoner by Crescens' followers and brought before Rusticus the prefect. When he firmly held to his confession of the faith, he was condemned to death, and went to his rest crowned with the glory of martyrdom.\par
\switchcolumn*

\LA
\TeDeum{Te Deum}
\switchcolumn
\EN
\TeDeum{Te Deum}
\switchcolumn*

\end{paracol}

\par\needspace{10\baselineskip}\addvspace{1.2em}{\centering\small\textsc{Die 14 Aprilis} · \textsc{14 April}\par}
\Office{Ss. Tiburtii, Valeriani, et Maximi Martyrum}{Sts. Tiburtius, Valerian et Maximus, Martyrs}{Simplex}
\addcontentsline{toc}{subsection}{Die 14 Aprilis · Ss. Tiburtii, Valeriani, et Maximi Martyrum\enspace\textit{Sts. Tiburtius, Valerian et Maximus, Martyrs}}
\begin{paracol}{2}
\LA
\LessonHead{Lectio IX (pro commemoratione)}
\switchcolumn
\EN
\LessonHead{Lesson IX (when commemorated)}
\switchcolumn*

\LA
\LessonTitle{Commemoratio: Ss. Tiburtii, Valeriani, et Maximi Martyrum}
\switchcolumn
\EN
\LessonTitle{Commemoration of: Sts. Tiburtius, Valerian et Maximus, Martyrs}
\switchcolumn*

\LA
\noindent Valerianus Romanus nobili genere ortus, Alexandro Sevéro imperatóre, hortatu beátæ Cæcíliæ Vírginis, quam sibi pari nobilitate uxórem desponderat, una cum Tiburtio fratre a sancto Urbano Papa baptizátur. Quos ubi præfectus Urbis Almachius Christiános esse cognóvit, et patrimonio paupéribus distributo, Christianórum córpora sepelire; accersítos gráviter reprehéndit: atque ubi Christum Deum constanter confiténtes, deos autem dæmoniórum inánia simulácra prædicántes videt, virgis cædi jubet. Sed cum verbéribus cogi non possent, ut Jovis simulacrum veneraréntur, immo fortes in fídei veritáte permanérent, ad quartum ab Urbe lápidem secúri feriúntur. Quorum virtútem admirátus Maximus præfecti cubicularius, qui eos ad supplícium perduxerat, Christiánum se esse proféssus est, cum multis præterea præfecti minístris: qui paulo post plumbátis contusi, omnes ex diaboli minístris, Christi Dómini Mártyres evasérunt.\par
\switchcolumn
\EN
\noindent Valerian was a Roman, of a family as noble as that of the blessed maiden Cecily, to whom he was contracted in marriage, in the reign of the Emperor Alexander Severus. At her persuasion he and his brother Tiburtius were baptized by the holy Pope Urban. When it came to the knowledge of Almachius, the Prefect of the city, that they were become Christians, had given their substance to the poor, and were burying the bodies of the faithful, he sent for them and strongly rebuked them but as they constantly confessed that Christ is God, and that the gods of the heathen are but vain images of devils, he commanded them to be beaten with rods. But, forasmuch as no blows could force them to worship the image of Jupiter, but they seemed rather to wax strong in witnessing to the truth of the faith that was in them, they were beheaded at the fourth mile-stone from the city. One of the clerks of the Prefect, named Maximus, who had led them out to die, was so moved at the sight of their courage that he himself, with many other servants of the Prefect, owned to being a Christian they were sentenced to be scourged to death with whips loaded with lead, under the which torment, in a little while, all these, who had once been the devil's ministers, passed away as martyrs of Christ the Lord.\par
\switchcolumn*

\LA
\TeDeum{Te Deum}
\switchcolumn
\EN
\TeDeum{Te Deum}
\switchcolumn*

\end{paracol}

\par\needspace{10\baselineskip}\addvspace{1.2em}{\centering\small\textsc{Die 17 Aprilis} · \textsc{17 April}\par}
\Office{S. Aniceti Papæ et Martyris}{St. Anicetus, Pope and Martyr}{Simplex}
\addcontentsline{toc}{subsection}{Die 17 Aprilis · S. Aniceti Papæ et Martyris\enspace\textit{St. Anicetus, Pope and Martyr}}
\begin{paracol}{2}
\LA
\RefLine{Lectiones I–II: de Scriptura occurrente.}
\switchcolumn
\EN
\RefLine{Lessons I–II: from the Scripture of the day.}
\switchcolumn*

\LA
\LessonHead{Lectio III (pro commemoratione)}
\switchcolumn
\EN
\LessonHead{Lesson III (when commemorated)}
\switchcolumn*

\LA
\noindent Anicétus Syrus, imperatóre Marco Aurélio Antoníno, præfuit Ecclésiæ. Decrévit ne clérici comam nutrírent. Quínquies mense Decémbri ordinávit presbyteros decem et septem, diáconos quátuor, epíscopos per divérsa loca novem. Vixit in pontificátu annos octo, menses octo, dies vigínti quátuor. Propter Christi fidem martyrio coronátus, sepúltus est via Appia in cœmetério, quod póstea Callísti appellátum est, décimo quinto Kaléndas Majas.\par
\switchcolumn
\EN
\noindent Anicetus Syrus, imperatore Marco Aurelio Antonino, praefuit Ecclesiae. Decrevit ne clerici comam nutrirent. Quinquies mense Decembri ordinavit presbyteros decem et septem, diaconos quatuor, episcopos per diversa loca novem. Vixit in pontificatu annos octo, menses octo, dies viginti quatuor. Propter Christi fidem martyrio coronatus, sepultus est via Appia in ccemeterio, quod postea Callisti appellatum est, decimo quinto Kalendas Majas.\par
\switchcolumn*

\LA
\TeDeum{Te Deum}
\switchcolumn
\EN
\TeDeum{Te Deum}
\switchcolumn*

\end{paracol}

\par\needspace{10\baselineskip}\addvspace{1.2em}{\centering\small\textsc{Die 21 Aprilis} · \textsc{21 April}\par}
\Office{S. Anselmi Episcopi Confessoris et Ecclesiæ Doctoris}{St. Anselm, Bishop, Confessor and Doctor of the Church}{Duplex}
\addcontentsline{toc}{subsection}{Die 21 Aprilis · S. Anselmi Episcopi Confessoris et Ecclesiæ Doctoris\enspace\textit{St. Anselm, Bishop, Confessor and Doctor of the Church}}
\begin{paracol}{2}
\LA
\RefLine{Lectiones I–III: de Scriptura occurrente.}
\switchcolumn
\EN
\RefLine{Lessons I–III: from the Scripture of the day.}
\switchcolumn*

\LA
\LessonHead{Lectio IV}
\switchcolumn
\EN
\LessonHead{Lesson IV}
\switchcolumn*

\LA
\noindent Ansélmus, Augustæ Prætoriæ in finibus Itáliæ, Gundulpho et Ermemberga, nobílibus et catholicis parentibus natus, a teneris annis assiduo litterarum studio, atque perfectióris vitæ desiderio, non obscurum futuræ sanctitatis et doctrinæ specimen dedit. Et licet juvenili ardore aliquando ad sæculi illecebras traheretur, brevi tamen in pristinam viam revocatus, patria et bonis omnibus derelictis, ad monasterium Beccense ordinis sancti Benedicti se contulit: ubi emissa regulari professione sub Herluino abbate observantissimo, et Lanfranco viro doctissimo, tanto animi fervore, et jugi studio in litteris et virtutibus assequendis profecit, ut mirum in modum tamquam sanctitatis et doctrinæ exemplar ab omnibus haberetur.\par
\switchcolumn
\EN
\noindent Anselm was born of noble and Catholic parents, named Gundulph and Hermenberga, at Aosta, in Piedmont, (about the year of our Lord 1033.) From his tenderest years his diligence in study, and his aspirations to a more perfect state of life, gave no indistinct foreshadowing of the holiness and learning to which he afterwards attained. The heat of youth drew him for a while into the snares of the world, but he soon returned to his first courses, and, forsaking his country and his goods, betook himself (in 1060.) to the monastery of Bee, under the rule of St. Benedict. There he made his profession as a monk, and under the rigid discipline of Herluin, the Abbot, and the learned instruction of the profound Lanfranc, with great zeal of spirit and eager obedience to the Rule, he made such progress in learning and godliness, that he shone before all others as an example of holiness of life, and power of doctrine.\par
\switchcolumn*

\LA
\begin{Resp}\Rsign Invéni David servum meum, óleo sancto meo unxi eum: \Star{} Manus enim mea auxiliábitur ei, allelúja.\end{Resp}
\begin{Resp}\Vsign Nihil profíciet inimícus in eo, et fílius iniquitátis non nocébit ei.\end{Resp}
\begin{Resp}\Rsign Manus enim mea auxiliábitur ei, allelúja.\end{Resp}
\switchcolumn
\EN
\begin{Resp}\Rsign I have found David My servant, with My holy oil have I anointed him \Star{} For My hand shall help him, alleluia.\end{Resp}
\begin{Resp}\Vsign The enemy shall prevail nothing against him, nor the son of wickedness afflict him.\end{Resp}
\begin{Resp}\Rsign For My hand shall help him, alleluia.\end{Resp}
\switchcolumn*

\LA
\LessonHead{Lectio V}
\switchcolumn
\EN
\LessonHead{Lesson V}
\switchcolumn*

\LA
\noindent Abstinentiæ et continentiæ tantæ fuit, ut assiduitate jejunii omnis pene ciborum sensus in eo videretur exstinctus. Diurno enim tempore in exercitiis monasticis docendo, et respondendo variis de religione quæsitis, emenso, quod reliquum erat noctis, somno subtrahebat, ut divinis meditationibus, quas perenni lacrimarum imbre fovebat, mentem recrearet Electus in priorem monasterii invidos fratres ita caritate, humilitate, et prudentia lenivit, ut quos æmulos acceperat, sibi et Deo amicos, maximo cum regularis observantiæ emolumento redderet Mortuo abbate, et in ejus locum (licet invitus) suffectus, tanta doctrinæ et sanctitatis fama ubique refulsit, ut non modo regibus et episcopis venerationi esset, sed sancto Gregorio septimo étiam acceptus, qui tunc magnis persecutionibus agitatus, litteras amoris plenas ad eum dedit, quibus se et Ecclesiam catholicam ejus orationibus commendabat.\par
\switchcolumn
\EN
\noindent Certification and purity were his marked characteristics, and by constant fasting all taste for food seemed to have died in him. He spent the day in the monastic work, in teaching, and in answering hard questions upon religion, and he took away from sleep during what remained to him of the night, that he might refresh his soul by thoughts of God, wherein he was alway comforted by an unceasing flow of tears. When he was chosen Prior of the monastery, he so won over, by his charity, lowliness, and wisdom, some brethren who looked ill upon him, that from enviers, as he had found them, he turned them into lovers of God and of himself likewise, with exceeding gain to the strictness of observance in that Abbey. After the death of the Abbot, (in 1078,) Anselm, though against his own will, was chosen to succeed him. In this high place the light of his learning and holiness so shone all round about, that he was reverenced not only by Kings and Bishops, but was taken up by the holy Pope Gregory VII., who, amid the great persecutions which were then trying him, wrote with words of great love to Anselm to recommend himself and the Catholic Church to his prayers.\par
\switchcolumn*

\LA
\begin{Resp}\Rsign Pósui adjutórium super poténtem, et exaltávi eléctum de plebe mea: \Star{} Manus enim mea auxiliábitur ei, allelúja.\end{Resp}
\begin{Resp}\Vsign Invéni David servum meum, óleo sancto meo unxi eum.\end{Resp}
\begin{Resp}\Rsign Manus enim mea auxiliábitur ei, allelúja.\end{Resp}
\switchcolumn
\EN
\begin{Resp}\Rsign I have laid help upon one that is mighty, and have exalted one chosen out of My people \Star{} For My hand shall help him, alleluia.\end{Resp}
\begin{Resp}\Vsign I have found David My servant, with My holy oil have I anointed him.\end{Resp}
\begin{Resp}\Rsign For My hand shall help him, alleluia.\end{Resp}
\switchcolumn*

\LA
\LessonHead{Lectio VI}
\switchcolumn
\EN
\LessonHead{Lesson VI}
\switchcolumn*

\LA
\noindent Defuncto Lanfranco archiepiscopo Cantuariensi, ejus olim præceptore, Anselmus, urgente Willelmo Angliæ rege, et instantibus clero ac populo, ipso tamen repugnante, ad ejusdem ecclesiæ regimen vocatus, statim (ut corruptos populi mores reformaret) verbo et exemplo prius, dein scriptis, et conciliis celebratis, pristinam pietatem, et ecclesiasticam disciplinam reduxit. Sed cum mox idem Willelmus rex vi et minis Ecclesiæ jura usurpare tentasset, ipse sacerdotali constantia restitit; bonorumque direptionem, et exsilium passus, Romam ad Urbanum secundum se contulit: a quo honorifice exceptus, et summis laudibus ornatus est, cum in Barensi concilio Spiritum Sanctum étiam a Filio procedentem contra Græcorum errorem innumeris Scripturarum, et sanctorum Patrum testimoniis propugnasset. E vivis Willelmo sublato, ab Henrico rege ejus fratre in Angliam revocatus, obdormivit in Domino, famam non solum miraculorum, et sanctitatis (præcipue ob insignem devotionem erga Domini nostri passionem, et beatam Virginem ejus Matrem) assecutus, sed étiam doctrinæ, quam ad defensionem christianæ religionis, animarum profectum, et omnium theologorum, qui sacras litteras scholastica methodo tradiderunt, normam cælitus hausisse ex ejus libris omnibus apparet.\par
\switchcolumn
\EN
\noindent After the death of Lanfranc, Archbishop of Canterbury, (in 1089,) Anselm, whose teacher Lanfranc had formerly been, was driven by William II., King of England, supported by the entreaties of the clergy and people, though sorely against his own wishes to take upon him the government of that Church. Raised to that See \textasciitilde{}(upon the 4th day of December, in the year 1093,) he straightway set himself to reform the corrupt manners of the people, and, first by his word and example, and then by his writings and the Councils which he held, succeeded in restoring the ancient godliness and discipline of the Church. But when the aforesaid King William tried by force and threats to seize on the rights of the Church, Anselm withstood him as beseemed a Priest, and after that he had suffered the plundering of all his goods and been sent into banishment, he betook himself to Rome to Urban II. There he was received with great worship, and won high praise for that in the Council of Bari, (in 1098,) he maintained by countless proofs from Scripture and the holy Fathers, against the error of the Greeks, that the Holy Ghost proceedeth from the Son also. When William lived no more, his brother Henry I., King of England, (in the year 1100,) called back Anselm thither, and there he fell asleep in the Lord, (upon the 21 st day of April,1109.) His is a name illustrious not for miracles only, nor for holiness, and indeed he had a wondrous love for his Lord Who had suffered for him, and for the blessed Maiden Mother of the Same our Lord, but also for the deep learning which he used for the defence of the Christian Religion and the good of souls. That wonderful knowledge of theology which he had, and which is shown in all the books which he wrote, seemeth to have been given him from heaven for the teaching of all writers on the same subject, who have used what is called the Scholastic method.\par
\switchcolumn*

\LA
\begin{Resp}\Rsign Iste est, qui ante Deum magnas virtútes operátus est, et omnis terra doctrína ejus repléta est: \Star{} Ipse intercédat pro peccátis ómnium populórum, allelúja.\end{Resp}
\begin{Resp}\Vsign Iste est, qui contémpsit vitam mundi, et pervénit ad cæléstia regna.\end{Resp}
\begin{Resp}\Rsign Ipse intercédat pro peccátis ómnium populórum, allelúja.\end{Resp}
\begin{Resp}\Vsign Glória Patri, et Fílio, \Star{} et Spirítui Sancto.\end{Resp}
\begin{Resp}\Rsign Ipse intercédat pro peccátis ómnium populórum, allelúja.\end{Resp}
\switchcolumn
\EN
\begin{Resp}\Rsign This is he which wrought great wonders before God, and the whole earth is full of his teaching \Star{} May he pray for all people, that their sins may be forgiven unto them, alleluia.\end{Resp}
\begin{Resp}\Vsign This is he which loved not his life in this world, and hath attained unto the kingdom of heaven.\end{Resp}
\begin{Resp}\Rsign May he pray for all people, that their sins may be forgiven unto them, alleluia.\end{Resp}
\begin{Resp}\Vsign Glory be to the Father, and to the Son, \Star{} and to the Holy Ghost.\end{Resp}
\begin{Resp}\Rsign May he pray for all people, that their sins may be forgiven unto them, alleluia.\end{Resp}
\switchcolumn*

\LA
\LessonHead{Lectio VII}
\switchcolumn
\EN
\LessonHead{Lesson VII}
\switchcolumn*

\LA
\LessonTitle{Léctio sancti Evangélii secúndum Matthǽum}
\switchcolumn
\EN
\LessonTitle{From the Holy Gospel according to Matthew}
\switchcolumn*

\LA
\Cite{Matt 5:13-19}
\switchcolumn
\EN
\Cite{Matt 5:13-19}
\switchcolumn*

\LA
\noindent In illo témpore: Dixit Jesus discípulis suis: Vos estis sal terræ. Quod si sal evanúerit, in quo saliétur? Et réliqua.\par
\switchcolumn
\EN
\noindent At that time, Jesus said unto His disciples: Ye are the salt of the earth. But if the salt have lost his savour, wherewith shall it be salted. And so on.\par
\switchcolumn*

\LA
\LessonTitle{Homilía sancti Hilárii Epíscopi}
\switchcolumn
\EN
\LessonTitle{Homily by St. Hilary, Bishop of Poitiers.}
\switchcolumn*

\LA
\Source{Comm. in Matthæum, cap. 4, n. 10}
\switchcolumn
\EN
\Source{Comment, on Matth. V}
\switchcolumn*

\LA
\noindent Vos estis sal terræ: Quod si sal infatuátum fúerit, ad níhilum valet id quod saliétur. Sal, ut árbitror, terræ nullum est. Quómodo ergo Apóstolos sal terræ nuncupávit? Sed propríetas est quærénda dictórum, quam et Apostolórum offícium, et ipsíus salis natúra monstrábit. Sal est in se uno cóntinens aquæ et ignis eleméntum, et hoc ex duóbus est unum.\par
\switchcolumn
\EN
\noindent “Ye are the salt of the earth. But if the salt have lost his savour, wherewith shall it be salted It is thenceforth good for nothing, but to be cast out, and to be trodden under foot of men.” There is, I take it, no such thing as salt of the earth. How, then, can the Apostles be called the salt of the earth? But the true meaning of these words will be made plain, when we consider the duty of Apostles, and the nature of salt itself. Now, salt is a compound of the elements of water and fire, out of the which two things in salt there is made one.\par
\switchcolumn*

\LA
\begin{Resp}\Rsign Amávit eum Dóminus, et ornávit eum: stolam glóriæ índuit eum, \Star{} Et ad portas paradísi coronávit eum, allelúja.\end{Resp}
\begin{Resp}\Vsign Induit eum Dóminus lorícam fídei, et ornávit eum.\end{Resp}
\begin{Resp}\Rsign Et ad portas paradísi coronávit eum, allelúja.\end{Resp}
\switchcolumn
\EN
\begin{Resp}\Rsign The Lord loved him and beautified him He clothed him with a robe of glory \Star{} And crowned him at the gates of Paradise, alleluia.\end{Resp}
\begin{Resp}\Vsign The Lord hath put on him the breast-plate of faith, and hath adorned him.\end{Resp}
\begin{Resp}\Rsign And crowned him at the gates of Paradise, alleluia.\end{Resp}
\switchcolumn*

\LA
\LessonHead{Lectio VIII}
\switchcolumn
\EN
\LessonHead{Lesson VIII}
\switchcolumn*

\LA
\noindent Hic ígitur in omnem usum humáni géneris efféctus, incorruptiónem corpóribus, quibus fúerit aspérsus, impértit, et ad omnem sensum cónditi sapóris aptíssimus est. Apóstoli autem sunt rerum cæléstium prædicatóres, et æternitátis velut satóres, immortalitátem ómnibus corpóribus, quibus eórum sermo aspérsus fúerit, conferéntes. Mérito ígitur sal terræ nuncupáti sunt, per doctrínæ virtútem saliéndi modo æternitáti córpora reservántes.\par
\switchcolumn
\EN
\noindent This thing, therefore, thus made to serve in diverse ways the use of men, doth keep from corruption bodies whereon it is sprinkled, and doth readily yield to all the senses the perception of its inborn savour. And thus are the Apostles, seeing that they are the preachers of the kingdom of heaven, and in a certain sense the sowers of the seed of life everlasting, since that Word of God which they scatter hath power to make this mortal put on immortality. Meetly then are they called salt, the savour of whose teaching doth keep sweet the receiver thereof even unto life everlasting.\par
\switchcolumn*

\LA
\begin{Resp}\Rsign In médio Ecclésiæ apéruit os ejus, \Star{} Et implévit eum Dóminus spíritu sapiéntiæ et intelléctus, allelúja.\end{Resp}
\begin{Resp}\Vsign Jucunditátem et exsultatiónem thesaurizávit super eum.\end{Resp}
\begin{Resp}\Rsign Et implévit eum Dóminus spíritu sapiéntiæ et intelléctus, allelúja.\end{Resp}
\begin{Resp}\Vsign Glória Patri, et Fílio, \Star{} et Spirítui Sancto.\end{Resp}
\begin{Resp}\Rsign Et implévit eum Dóminus spíritu sapiéntiæ et intelléctus, allelúja.\end{Resp}
\switchcolumn
\EN
\begin{Resp}\Rsign In the midst of the congregation did the Lord open his mouth. \Star{} And filled him with the spirit of wisdom and understanding, alleluia.\end{Resp}
\begin{Resp}\Vsign He made him rich with joy and gladness.\end{Resp}
\begin{Resp}\Rsign And filled him with the spirit of wisdom and understanding, alleluia.\end{Resp}
\begin{Resp}\Vsign Glory be to the Father, and to the Son, \Star{} and to the Holy Ghost.\end{Resp}
\begin{Resp}\Rsign And filled him with the spirit of wisdom and understanding, alleluia.\end{Resp}
\switchcolumn*

\LA
\LessonHead{Lectio IX}
\switchcolumn
\EN
\LessonHead{Lesson IX}
\switchcolumn*

\LA
\noindent Sed natúra salis semper éadem est, nec immutári umquam potest. Verum quia conversióni homo súbjacet, et solus beátus, qui usque ad finem in ómnibus Dei opéribus permánserit: ídeo eos sal terræ nuncupátos monet in tráditæ sibi potestátis virtúte persístere, ne infatuáti nihil sáliant, et ipsi, sensu accépti sapóris amísso, vivificáre corrúpta non possint, et projécti de Ecclésiæ promptuáriis, cum his quos salíerint, pédibus incedéntium proterántur.\par
\switchcolumn
\EN
\noindent But the nature of salt is to be ever the same, and unchanging, and, on the other hand, the nature of man hath this weakness, to be changeable. He only is blessed who hath continued even unto the end in all the works which God hath commanded. Therefore dot! the Lord warn them whom He calleth the salt of the earth, that they are behoven to remain strong in that strength which He hath given unto them, lest, becoming themselves savourless, they should be impotent to season others losing the freshness of their own saltness, be unable to stop the corruption round about them and so the Church cast them out of her buttery, and they and those that they should have salted, be together trodden under foot of such as enter in.\par
\switchcolumn*

\LA
\TeDeum{Te Deum}
\switchcolumn
\EN
\TeDeum{Te Deum}
\switchcolumn*

\LA
\LessonHead{Lectio IX (pro commemoratione)}
\switchcolumn
\EN
\LessonHead{Lesson IX (when commemorated)}
\switchcolumn*

\LA
\LessonTitle{Commemoratio: S. Anselmi Episcopi Confessoris et Ecclesiæ Doctoris}
\switchcolumn
\EN
\LessonTitle{Commemoration of: St. Anselm, Bishop, Confessor and Doctor of the Church}
\switchcolumn*

\LA
\noindent Ansélmus, Augústæ Prætóriæ in fínibus Itáliæ, nobílibus et cathólicis paréntibus natus, adoléscens, pátria et bonis ómnibus derelíctis, in monastério Beccénsi órdinis sancti Benedícti emíssa regulári professióne, in lítteris et virtútibus assequéndis mirum in modum profécit. Régibus, epíscopis veneratióni fuit, et sancto Gregório séptimo étiam accéptus, qui tunc, persecutiónibus agitátus, lítteras amóris plenas ad eum dedit, se et Ecclésiam ejus oratiónibus comméndans. Defúncto Lanfránco archiepíscopo Cantuariénsi, ejus olim præceptóre, ad ejúsdem ecclésiæ régimen vocátus, verbo et exémplo, scriptis et concíliis celebrátis, prístinam pietátem et ecclesiásticam disciplínam redúxit. Sed cum mox Willélmus rex vi et minis jura Ecclésiæ usurpáre tentásset, ípseque invícte restitísset, bonórum direptiónem et exsílium passus, Romam ad Urbánum secúndum se cóntulit. A quo honorífice excéptus et summis láudibus ornátus, in Barénsi concílio Spíritum Sanctum étiam a Fílio procedéntem, contra Græcórum errórem, innúmeris Scripturárum et sanctórum Patrum testimóniis propugnávit. Post mortem Willélmi, ab Henríco rege ejus fratre in Angliam revocátus obdormívit in Dómino.\par
\switchcolumn
\EN
\noindent Anselm, born of noble Catholic parents at Aosta on the borders of Italy, as a young man abandoned his homeland and all his possessions and was professed at the Benedictine monastery of Bec, where he advanced in a most wonderful way in the attainment of learning and virtue. He was held in honour by kings and bishops, and was a friend of St. Gregory VII, at the time much troubled by persecutions, who wrote him letters filled with affection, commending himself and the Church to Anselm's prayers. After the death of Lanfranc, Archbishop of Canterbury and Anselm's former teacher, he was called to rule over that Church, and, by word and example, by writings and by holding councils, he restored it to its pristine state of piety and ecclesiastical discipline. But, soon after, when King William tried by force and threats to usurp the rights of the Church and Anselm steadfastly resisted, his possessions were confiscated and he himself exiled. He went to Urban II in Rome, who welcomed him with honour and the highest praise. At the Council of Bari, he defended against the errors of the Greeks the doctrine that the Holy Ghost proceedeth also from the Son, by countless proofs taken from the Scriptures and the holy Fathers. After King William's death, his brother Henry recalled Anselm to England, and there he fell asleep in the Lord.\par
\switchcolumn*

\LA
\TeDeum{Te Deum}
\switchcolumn
\EN
\TeDeum{Te Deum}
\switchcolumn*

\end{paracol}

\par\needspace{10\baselineskip}\addvspace{1.2em}{\centering\small\textsc{Die 22 Aprilis} · \textsc{22 April}\par}
\Office{SS. Soteris et Caji Summorum Pontificum et Martyrum}{Sts. Soter and Caius, Popes and Martyrs}{Semiduplex}
\addcontentsline{toc}{subsection}{Die 22 Aprilis · SS. Soteris et Caji Summorum Pontificum et Martyrum\enspace\textit{Sts. Soter and Caius, Popes and Martyrs}}
\begin{paracol}{2}
\LA
\RefLine{Lectiones I–III: de Scriptura occurrente.}
\switchcolumn
\EN
\RefLine{Lessons I–III: from the Scripture of the day.}
\switchcolumn*

\LA
\RefLine{Lectiones VII–IX: ex Commúni plurium Summorum Pontificum Martyrum Tempore Paschali (C3bp).}
\switchcolumn
\EN
\RefLine{Lessons VII–IX: from the Commune plurium Summorum Pontificum Martyrum Tempore Paschali (C3bp).}
\switchcolumn*

\LA
\LessonHead{Lectio IV}
\switchcolumn
\EN
\LessonHead{Lesson IV}
\switchcolumn*

\LA
\noindent Soter, Fundis in Campánia natus, sancívit ne sacræ vírgines vasa sacra et pallas attingerent, neve thuris ministerio in ecclésia uteréntur. Idem státuit ut Christi corpus in Cœna Dómini sumerétur ab ómnibus, iis exceptis, qui propter grave peccátum id fácere prohiberéntur. Sedit in pontificatu annos tres, menses undecim, dies decem et octo. Martyrio coronátur sub Marco Aurelio imperatóre, et in cœmeterio, quod póstea Callísti dictum est, sepelitur; more majórum mense Decembri creátis presbyteris decem et octo, diáconis novem, epíscopis per diversa loca undecim.\par
\switchcolumn
\EN
\noindent Soter, a countryman of Fondi in Campania, succeeded the holy martyr Anicete, (in the year 173.) It was he who ordained that nuns should not touch the sacred vessels and linen of the Altar, nor serve with the incense in the Church. He ordained likewise, that on the anniversary of the Lord's Supper, every one should receive the Body of Christ, except those who were forbidden to do so on account of grievous sin. He sat as Pope three years, eleven months, and twenty-eight days. He ordained in the month of December eighteen Priests, nine Deacons, and eleven Bishops for diverse places. He was crowned with martyrdom under the Emperor Marcus Aurelius, (in 177,) and was buried after the manner of them that had gone before him, in the Cemetery, which was afterwards called that of St. Calixtus.\par
\switchcolumn*

\LA
\begin{Resp}\Rsign Lux perpétua lucébit Sanctis tuis, Dómine, \Star{} Et ætérnitas témporum, allelúja, allelúja.\end{Resp}
\begin{Resp}\Vsign Lætítia sempitérna erit super cápita eórum: gáudium et exsultatiónem obtinébunt.\end{Resp}
\begin{Resp}\Rsign Et ætérnitas témporum, allelúja, allelúja.\end{Resp}
\switchcolumn
\EN
\begin{Resp}\Rsign Eternal light will shine over your saints, O Lord, \Star{} And the eternity of the times, alleluia, alleluia.\end{Resp}
\begin{Resp}\Vsign Everlasting joy shall be upon their heads, they shall obtain joy and gladness\end{Resp}
\begin{Resp}\Rsign And the eternity of the times, alleluia, alleluia.\end{Resp}
\switchcolumn*

\LA
\LessonHead{Lectio V}
\switchcolumn
\EN
\LessonHead{Lesson V}
\switchcolumn*

\LA
\noindent Cajus Dálmata ex genere Diocletiáni imperatóris, constítuit ut his ordinum et honórum gradibus in Ecclésia ad episcopátum ascenderétur: Ostiarii, Lectóris, Exorcistæ, Acólythi, Subdiáconi, Diáconi, Presbýteri. Hic Diocletiáni crudelitátem in Christiános fúgiens, aliquámdiu in spelúnca delituit; verum octo post annis una cum Gabino fratre martyrii corónam consecutus est, cum sedísset annos duodecim, menses quatuor, dies quinque; creátis mense Decembri presbyteris vigintiquinque, diáconis octo, epíscopis quinque. Sepultus est in cœmeterio Callísti, décimo Kalendas Maji. Ejus memóriam Urbanus octavus in Urbe renovávit, dírutam ecclésiam restituit; titulo, statióne et ipsíus relíquiis decorávit.\par
\switchcolumn
\EN
\noindent Gaius was a Dalmatian and a kinsman of the Emperor Diocletian and succeeded holy Eutychian in the year 283. It was he who ordained that the following should be the order of degrees in the Church through which all should pass before they be made Bishop First, Doorkeeper second, Reader third, Exorcist fourth, Acolyte fifth, Subdeacon sixth, Deacon seventh, Priest. Caius fled from the cruelties practiced by Diocletian against the Christians, and lay hid for a while in a cave, but after eight years he and his brother Gabinus won the crown of martyrdom, (upon the 21st day of April, in the year 296.) At that time he had sat in the chair of Peter twelve years, four months, and five days, and had ordained in the month of December twenty - five Priests, eight Deacons, and five Bishops. He was buried in the Cemetery of Calixtus upon the twenty-second day of April. It was Urban VIII. who renewed the memorial of him in the city, rebuilt his Church, which had been in ruins, and distinguished it by making it one of those whence the Cardinals take their titles, and of those which are called “Stations,” and enriching it with the reliques of the Saint.\par
\switchcolumn*

\LA
\begin{Resp}\Rsign In servis suis, allelúja, \Star{} Consolábitur Deus, allelúja.\end{Resp}
\begin{Resp}\Vsign Judicábit Dóminus pópulum suum, et in servis suis.\end{Resp}
\begin{Resp}\Rsign Consolábitur Deus, allelúja.\end{Resp}
\switchcolumn
\EN
\begin{Resp}\Rsign His servants, alleluia \Star{} God will comfort, alleluia, alleluia.\end{Resp}
\begin{Resp}\Vsign The Lord will judge His people and, His servants,\end{Resp}
\begin{Resp}\Rsign God will comfort, alleluia, alleluia.\end{Resp}
\switchcolumn*

\LA
\LessonHead{Lectio VI}
\switchcolumn
\EN
\LessonHead{Lesson VI}
\switchcolumn*

\LA
\LessonTitle{Sermo sancti Ambrósii Epíscopi}
\switchcolumn
\EN
\LessonTitle{From the sermon of Saint Ambrose Bishop (of Milan)}
\switchcolumn*

\LA
\Source{Sermo 22.}
\switchcolumn
\EN
\Source{Sermon 22.}
\switchcolumn*

\LA
\noindent Dignum et congruum est, fratres, ut post lætítiam Paschæ, quam in Ecclésia celebrávimus, gaudia nostra cum sanctis Martyribus conferámus; et iis annuntiémus Dominicæ resurrectiónis glóriam, qui consortes sunt Dominicæ passiónis. Qui enim socii sunt contumeliæ, debent et participes esse lætítiæ. Ita enim dicit beátus Apóstolus: Sicut socii passiónum estis, et resurrectiónis eritis; si tolerábimus, inquit, et conregnábimus. Qui ergo toleravérunt mala propter Christum, debent et glóriam habere cum Christo.\par
\switchcolumn
\EN
\noindent Dearly beloved brethren, it is very meet and right that after the gladness of Easter, which we have celebrated in the Church, we should mingle our own joy with the joy of the holy Martyrs; yea, that we should tell of the glory of the Lord's rising again, to them that have been made partakers of the Lord's sufferings. It truly must needs be that they which have been partakers of His sufferings, should be also of His joy. For thus saith the blessed Apostle: “As ye are partakers of the sufferings, so shall ye be also of the consolation.” (2 Cor. i. 7.) And again: “If we suffer, we shall also reign with Him.” (2 Tim. ii. 12.) He, therefore, that endureth sorrow for Christ, must needs also have glory with Christ.\par
\switchcolumn*

\LA
\begin{Resp}\Rsign Fíliæ Jerúsalem, veníte et vidéte Mártyres cum corónis, quibus coronávit eos Dóminus \Star{} In die solemnitátis et lætítiæ, allelúja.\end{Resp}
\begin{Resp}\Vsign Quóniam confortávit seras portárum tuárum, benedíxit fílios tuos in te.\end{Resp}
\begin{Resp}\Rsign In die solemnitátis et lætítiæ, allelúja.\end{Resp}
\begin{Resp}\Vsign Glória Patri, et Fílio, \Star{} et Spirítui Sancto.\end{Resp}
\begin{Resp}\Rsign In die solemnitátis et lætítiæ, allelúja.\end{Resp}
\switchcolumn
\EN
\begin{Resp}\Rsign Come forth, O ye daughters of Jerusalem, and behold the Martyrs with the crowns wherewith the Lord crowned them; \Star{} In the day of His feasting, and of His gladness. Alleluia.\end{Resp}
\begin{Resp}\Vsign For He hath strengthened the bars of thy gates He hath blessed thy children within thee.\end{Resp}
\begin{Resp}\Rsign In the day of His feasting, and of His gladness. Alleluia.\end{Resp}
\begin{Resp}\Vsign Glory be to the Father, and to the Son, \Star{} and to the Holy Ghost.\end{Resp}
\begin{Resp}\Rsign In the day of His feasting, and of His gladness. Alleluia.\end{Resp}
\switchcolumn*

\LA
\LessonHead{Lectio IX (pro commemoratione)}
\switchcolumn
\EN
\LessonHead{Lesson IX (when commemorated)}
\switchcolumn*

\LA
\LessonTitle{Commemoratio: SS. Soteris et Caji Summorum Pontificum et Martyrum}
\switchcolumn
\EN
\LessonTitle{Commemoration of: Sts. Soter and Caius, Popes and Martyrs}
\switchcolumn*

\LA
\noindent Soter, Fundis in Campánia natus, sancívit ne sacræ vírgines vasa sacra et pallas attíngerent, neve thuris ministério in ecclésia uteréntur. Idem státuit, ut Christi corpus in Cena Dómini sumerétur ab ómnibus, iis excéptis, qui propter grave peccátum id fácere prohiberéntur. Martýrio coronátur sub Marco Aurélio imperatóre, sepúltus est in cœmetério, quod póstea Callísti dictum est. Cajus Dálmata ex génere Diocletiáni imperatóris, constítuit ut his órdinum et honórum grádibus in Ecclésia ad episcopátum ascenderétur: Ostiárii nempe, Lectóris, Exorcístæ, Acólythi, Subdiáconi, Diáconi, Presbýteri. Hic Diocletiáni crudelitátem in Christiános fúgiens, aliquámdiu in spelúnca delítuit; verum octo post annis una cum Gabíno fratre martýrii corónam consecútus, in cœmetério Callísti páriter sepúltus est.\par
\switchcolumn
\EN
\noindent Soter, a countryman of Fondi in Campania, succeeded the holy martyr Anicetus. It was he who ordained that nuns should not touch the sacred vessels and linen of the Altar, nor serve with the incense in the Church. He ordained likewise, that on the anniversary of the Lord's Supper, every one should receive the Body of Christ, except those who were forbidden to do so on account of grievous sin. He was crowned with martyrdom under the Emperor Marcus Aurélius, and was buried in the Cemetery, which was afterwards called that of St. Calixtus. Caius was a Dalmatian and a kinsman of the Emperor Diocletian. It was he who ordained that the following should be the order of degrees in the Church through which all should pass before they be made bishop: first, Porter; second, Lector; third, Exorcist; fourth, Acolyte; fifth, Subdeacon; sixth, Deacon; seventh, Priest. Caius fled from the cruelties practiced by Diocletian against the Christians, and lay hid for a while in a cave, but after eight years he and his brother Gabinus won the crown of martyrdom, and was likewise buried in the Cemetery of Callistus.\par
\switchcolumn*

\LA
\TeDeum{Te Deum}
\switchcolumn
\EN
\TeDeum{Te Deum}
\switchcolumn*

\end{paracol}

\par\needspace{10\baselineskip}\addvspace{1.2em}{\centering\small\textsc{Die 23 Aprilis} · \textsc{23 April}\par}
\Office{S. Georgii Martyris}{St. George, Martyr}{Semiduplex}
\addcontentsline{toc}{subsection}{Die 23 Aprilis · S. Georgii Martyris\enspace\textit{St. George, Martyr}}
\begin{paracol}{2}
\LA
\RefLine{Lectiones I–III: de Scriptura occurrente.}
\switchcolumn
\EN
\RefLine{Lessons I–III: from the Scripture of the day.}
\switchcolumn*

\LA
\RefLine{Lectiones IV–VI: ex Commúni unius Martyris Tempore Paschali (C2ap), in 2 loco.}
\switchcolumn
\EN
\RefLine{Lessons IV–VI: from the Common of a Single Martyr in Paschaltide (C2ap), set 2.}
\switchcolumn*

\LA
\RefLine{Lectiones VII–IX: ex Commúni unius Martyris Tempore Paschali (C2ap).}
\switchcolumn
\EN
\RefLine{Lessons VII–IX: from the Common of a Single Martyr in Paschaltide (C2ap).}
\switchcolumn*

\LA
\LessonHead{Lectio IX (pro commemoratione)}
\switchcolumn
\EN
\LessonHead{Lesson IX (when commemorated)}
\switchcolumn*

\LA
\LessonTitle{Commemoratio: S. Georgii Martyris}
\switchcolumn
\EN
\LessonTitle{Commemoration of: St. George, Martyr}
\switchcolumn*

\LA
\noindent The martyr George beareth among the Easterns the title of the holy and glorious Archmartyr. He suffered a glorious death, for Christ's sake, in the persecution under Diocletian. When peace was given to the Church soon after, under Constantine, the memory of the martyr began to be celebrated, and churches were built under his invocation at Lydda in Palestine and at Constantinople. For thenceforth an extraordinary enthusiasm with regard to him grew up among the faithful, first in all parts of the East, and afterwards in the West. Of old time, when Christian armies had been about to fight, they have been used to call as patrons upon holy George, Maurice, and Sebastian. There had been already special honour paid in England to the holy martyr George, and the supreme Pontiff Benedict XIV declared him the protector of the whole kingdom.\par
\switchcolumn
\EN
\noindent The martyr George beareth among the Easterns the title of the holy and glorious Archmartyr, He suffered a glorious death, for Christ's sake, in the persecution under Diocletian. When peace was given to the Church soon after, under Constantine, the memory of the martyr began to be celebrated, and churches were built under his invocation at Lydda in Palestine and at Constantinople. For thenceforth an extraordinary enthusiasm with regard to him grew up among the faithful, first in all parts of the East, and afterwards in the West. Of old time, when Christian armies had been about to fight, they have been used to call as patrons upon holy George, Maurice, and Sebastian. There had been already special honour paid in England to the holy martyr George, and the supreme Pontiff Benedict XIV. declared him the protector of the whole kingdom.\par
\switchcolumn*

\LA
\TeDeum{Te Deum}
\switchcolumn
\EN
\TeDeum{Te Deum}
\switchcolumn*

\end{paracol}

\par\needspace{10\baselineskip}\addvspace{1.2em}{\centering\small\textsc{Die 24 Aprilis} · \textsc{24 April}\par}
\Office{S. Fidelis de Sigmaringa Martyris}{St. Fidelis of Sigmaringen, Martyr}{Duplex}
\addcontentsline{toc}{subsection}{Die 24 Aprilis · S. Fidelis de Sigmaringa Martyris\enspace\textit{St. Fidelis of Sigmaringen, Martyr}}
\begin{paracol}{2}
\LA
\RefLine{Lectiones I–III: de Scriptura occurrente.}
\switchcolumn
\EN
\RefLine{Lessons I–III: from the Scripture of the day.}
\switchcolumn*

\LA
\RefLine{Lectiones VII–IX: ex Commúni unius Martyris Tempore Paschali (C2ap).}
\switchcolumn
\EN
\RefLine{Lessons VII–IX: from the Common of a Single Martyr in Paschaltide (C2ap).}
\switchcolumn*

\LA
\LessonHead{Lectio IV}
\switchcolumn
\EN
\LessonHead{Lesson IV}
\switchcolumn*

\LA
\noindent Fidélis, in oppido Sueviæ Sigmaringa ex honesta Reyórum família natus, ab ineunte ætate singularibus natúræ et gratiæ donis ornatus præfulsit. Egregiam quippe sortitus índolem, morúmque optima imbutus disciplína, dum Friburgi philosophíæ et juris utriusque lauream emeruit, in schola étiam Christi ad perfectiónis ápicem sedulo virtútum exercítio conténdere cœpit. Nobílium exínde virórum varias Europæ provincias lustrántium comes ascitus, eos ad christianam pietátem sectandam, tam verbis quam opéribus, excitare non déstitit. Quinimmo in eodem itinere crebris austeritátibus desidéria carnis mortificare, ac ita seípsum regere studuit, ut in tanta rerum vicissitúdine nullo umquam visus fúerit iræ motu perturbari. Juris præterea et justítiæ strenuus propugnator, post réditum in Germaniam celebre sibi nomen acquisívit in advocáti munere; in quo tamen cum fori pericula esset expertus, tutiórem ætérnæ salútis viam íngredi deliberávit, et superna vocatióne illustratus, paulo post ordini Seraphico inter fratres Minores Capuccinos adscribi pétiit.\par
\switchcolumn
\EN
\noindent Fidelis was born of the respectable family of Rey in the town of Sigmaringen in Swabia, (in the year of our Lord 1577.) From his childhood he was adorned with many bright gifts of nature and grace. Intellectually distinguished, and assisted by all the advantages of education, he took at Fribourg the degrees of Philosophy and of Civil and Canon Law, and it was while engaged in these studies, that he began to strive after the height of perfection in the school of Christ, to which end he earnestly trained himself in all the exercises of godliness. He ceased not to exhort to Christian godliness, both by his words and works, the noblemen who made him their companion, and who were drawn from the chief families of diverse parts of Europe. While on his travels, he was careful to mortify the lusts of the flesh by frequent austerities, and so to get the command of himself, that he was never seen under any circumstances to be moved to anger. He was a zealous champion of law and justice, and when he returned into Germany, he won a most distinguished name in his profession as an advocate. After a while, however, in view of the dangers which beset him at the Bar, he determined to enter on a path safer as regarded his eternal salvation, and, in obedience to an inward call from above, he sought admission into the Seraphic Order, among the Capuchin Friars Minor, (in the year 1612.)\par
\switchcolumn*

\LA
\begin{Resp}\Rsign Lux perpétua lucébit Sanctis tuis, Dómine, \Star{} Et ætérnitas témporum, allelúja, allelúja.\end{Resp}
\begin{Resp}\Vsign Lætítia sempitérna erit super cápita eórum: gáudium et exsultatiónem obtinébunt.\end{Resp}
\begin{Resp}\Rsign Et ætérnitas témporum, allelúja, allelúja.\end{Resp}
\switchcolumn
\EN
\begin{Resp}\Rsign Eternal light will shine over your saints, O Lord, \Star{} And the eternity of the times, alleluia, alleluia.\end{Resp}
\begin{Resp}\Vsign Everlasting joy shall be upon their heads, they shall obtain joy and gladness\end{Resp}
\begin{Resp}\Rsign And the eternity of the times, alleluia, alleluia.\end{Resp}
\switchcolumn*

\LA
\LessonHead{Lectio V}
\switchcolumn
\EN
\LessonHead{Lesson V}
\switchcolumn*

\LA
\noindent Piæ petitiónes compos redditus, mundi suique contemptor insígnis, in ipso statim tirocinio, magisque cum solemnis professióne vota in gáudio spíritus Dómino nuncupasset, in regulari observántia ómnibus admiratióni fuit et exemplo. Oratióni maxime et sacris litteris vacans, in verbi quoque ministerio singulari grátia excellens, nedum Catholicos ad meliórem frugem, verum étiam heterodoxos ad veritátis cognitiónem attraxit. Plúribus in locis cœnobii præfectus constitútus, prudéntia, justítia, mansuetúdine, discretióne et humilitátis laude, munus sibi demandátum exercuit. Arctíssimæ paupertátis zelator egregius, quidquid vel minus necessárium viderétur, e cœnobio pénitus eliminávit. Inter austera jejunia, vigílias et flagélla salutári seípsum prósequens ódio, in alios amórem, quasi mater in fílios, osténdit. Cum pestífera febris Austríacas militares copias dire affligeret, ipse in extremis infirmórum indigentiis ad assidua caritátis offícia toto spíritu incúbuit. In componéndis étiam animórum dissidiis, aliisque próximi necessitátibus sublevandis consílio et opere adeo præcláre se gessit, ut pater pátriæ merúerit appellari.\par
\switchcolumn
\EN
\noindent After he had obtained his holy wish, he showed himself even in his noviceship a singular despiser of the world and of himself, and still more so when with great spiritual joy he had made his solemn profession to the Lord. By his observance of the Rule, he became the wonder and the example of all. He gave himself chiefly to prayer and sacred learning, but he excelled, by a remarkable grace, in the ministry of the Word, and thereby not only stirred up the Catholics to bring forth more fruit, but also drew misbelievers to the knowledge of the truth. He was set at the head of communities of Friars in diverse places, and discharged the duty so laid upon him with great praise for prudence, justice, meekness, wisdom, and lowliness. He was animated by a vehement love of the strictest poverty, and cleansed the convent of whatever was not altogether needful. While he pursued himself with an healthy hatred, and most stern fastings, watchings, and scourgings, he showed to all others a love like the love of a mother for her sons. When a contagious fever made horrid ravages among the Austrian soldiers, he gave himself up with his whole soul to unwearied offices of tenderness toward the helpless sick. In allaying quarrels and relieving the temporal distress of his neighbour, he bore himself with such wisdom and zeal as to earn the name of Father of his country.\par
\switchcolumn*

\LA
\begin{Resp}\Rsign In servis suis, allelúja, \Star{} Consolábitur Deus, allelúja.\end{Resp}
\begin{Resp}\Vsign Judicábit Dóminus pópulum suum, et in servis suis.\end{Resp}
\begin{Resp}\Rsign Consolábitur Deus, allelúja.\end{Resp}
\switchcolumn
\EN
\begin{Resp}\Rsign His servants, alleluia \Star{} God will comfort, alleluia, alleluia.\end{Resp}
\begin{Resp}\Vsign The Lord will judge His people and, His servants,\end{Resp}
\begin{Resp}\Rsign God will comfort, alleluia, alleluia.\end{Resp}
\switchcolumn*

\LA
\LessonHead{Lectio VI}
\switchcolumn
\EN
\LessonHead{Lesson VI}
\switchcolumn*

\LA
\noindent Deíparæ Vírginis et rosarii cultor exímius, illíus præcipue aliorúmque Sanctórum patrocinio a Deo postulávit, ut in catholicæ fidei obsequium vitam sibi et sánguinem fúndere licéret. Cumque ardens hoc desidérium in quotidiana Sacri devota celebratióne magis accenderétur, mira Dei providéntia factum est, ut fortis Christi athleta præses eligerétur illárum missiónum, quas Congregátio de Propaganda Fide pro Rhætia tunc témporis institúerat. Quod arduum munus prompto hilaríque animo suscípiens, tanto fervore exsecutus est, ut, plúribus hæreticis ad orthodoxam fidem conversis, spes non modica effulserit totius illíus gentis Ecclésiæ et Christo reconciliandæ. Prophetíæ dono præditus, futuras Rhætiæ calamitátes, suique necem ab hæreticis inferéndam sæpius prædixit. Postquam vero insidiárum probe cónscius impendénti agoni se præparasset, die vigesima quarta Aprilis anno millesimo sexcentésimo vigesimo secundo, ad ecclésiam loci, Sevisium nuncupati, se cóntulit; ubi ab hæreticis, qui pridie conversiónem simulántes, eum dolóse ad prædicándum invitáverant, concióne tumultuárie interrupta, per verbera ac vulnera eídem crudeliter inflicta gloriósam mortem magno et álacri corde perpéssus, primítias Mártyrum memorátæ Congregatiónis proprio sánguine consecrávit; plúribus signis et miraculis exínde clarus, præsertim Curiæ et Feldkírchii, ubi summa pópuli veneratióne illíus relíquiæ asservántur.\par
\switchcolumn
\EN
\noindent Tenderly and warmly loved the maiden Mother of God and her Rosary, and he besought God under the patronage of many of His holy servants, but especially under that of the same blessed Mother, to vouchsafe to let him offer his life and his blood together for the sake of the Catholic faith. This burning desire came upon him more and more, day by day, as he celebrated with great ardour of spirit the Holy Liturgy and by the unexpected Providence of God it came to pass that this brave He was a travelling tutor. soldier of Christ was chosen President of the Missions which the Congregation for the Propagation of the Faith had at that time just founded for the Grisons. He accepted this hard task with a willing and joyful heart, and discharged it with such zeal, that many heretics were turned to the orthodox faith, and great hope was engendered that the whole of that people would return to the peace of Christ and His Church. Fidelis, who was gifted with the spirit of Prophecy, often foretold the great woes which afterwards came upon the Grisons, and that he himself would be murdered by the heretics. At last, on a certain 23rd of April, some of the heretics, who pretended to be converted, entreated him to come and preach the following day at the Church of a place which is called Sevis. He complied with the treacherous invitation, but, as he knew that plots were being laid against him, he had made himself ready beforehand for the last conflict. On the 24th day of April, in the year 1622, he went to Sevis, and began to preach, but his discourse was interrupted by a riot, and on his way back, he was [met by a party of Calvinists, and) brutally murdered. By this glorious death, which he suffered with a willing and cheerful heart, he offered to God in his own blood the first-fruits of martyrdom from the abovementioned Congregation. God hath since glorified him by many signs and wonders, especially at Chur and Feldkirch, where his relics are kept with much popular veneration.\par
\switchcolumn*

\LA
\begin{Resp}\Rsign Fíliæ Jerúsalem, veníte et vidéte Mártyres cum corónis, quibus coronávit eos Dóminus \Star{} In die solemnitátis et lætítiæ, allelúja.\end{Resp}
\begin{Resp}\Vsign Quóniam confortávit seras portárum tuárum, benedíxit fílios tuos in te.\end{Resp}
\begin{Resp}\Rsign In die solemnitátis et lætítiæ, allelúja.\end{Resp}
\begin{Resp}\Vsign Glória Patri, et Fílio, \Star{} et Spirítui Sancto.\end{Resp}
\begin{Resp}\Rsign In die solemnitátis et lætítiæ, allelúja.\end{Resp}
\switchcolumn
\EN
\begin{Resp}\Rsign Come forth, O ye daughters of Jerusalem, and behold the Martyrs with the crowns wherewith the Lord crowned them; \Star{} In the day of His feasting, and of His gladness. Alleluia.\end{Resp}
\begin{Resp}\Vsign For He hath strengthened the bars of thy gates He hath blessed thy children within thee.\end{Resp}
\begin{Resp}\Rsign In the day of His feasting, and of His gladness. Alleluia.\end{Resp}
\begin{Resp}\Vsign Glory be to the Father, and to the Son, \Star{} and to the Holy Ghost.\end{Resp}
\begin{Resp}\Rsign In the day of His feasting, and of His gladness. Alleluia.\end{Resp}
\switchcolumn*

\LA
\LessonHead{Lectio IX (pro commemoratione)}
\switchcolumn
\EN
\LessonHead{Lesson IX (when commemorated)}
\switchcolumn*

\LA
\LessonTitle{Commemoratio: S. Fidelis de Sigmaringa Martyris}
\switchcolumn
\EN
\LessonTitle{Commemoration of: St. Fidelis of Sigmaringen, Martyr}
\switchcolumn*

\LA
\noindent Fidélis, in óppido Suéviæ Sigmarínga ex honésta Reyórum família natus, célebre sibi nomen acquisívit in advocáti múnere; quo tamen, cum fori perícula esset expértus, déstitit, et supérna vocatióne illustrátus, inter fratres Minóres Capuccínos adscríbi pétiit. Voti compos factus, regulári observántia ómnibus admiratióni fuit et exémplo. Deíparæ Vírginis et rosárii cultor exímius, a Deo postulávit, ut pro cathólica fide martyr occúmberet, quod et consecútus est. Eléctus enim præses Missiónum, quas Congregátio de Propagánda Fide pro Rhǽtia tunc témporis institúerat, cum nulli labóri parcens plures hæréticos ad Christi fidem convertísset, malórum invídiam súbiit. Itaque die vigésima quarta Aprílis anno millésimo sexcentésimo vigésimo secúndo, ad ecclésiam loci, Sevísium nuncupáti, verbéribus ac vulnéribus cæsus ab hæréticis, qui conversiónem simulántes dolóse eum invitáverant, primítias Mártyrum memorátæ Congregatiónis próprio sánguine consecrávit.\par
\switchcolumn
\EN
\noindent Fidelis was born of a good family called Rey at Sigmaringen in Swabia, and gained a famous name for himself in legal work. But when experience shewed him its dangers, he abandoned it and, enlightened by a call from heaven, asked to be admitted among the Capuchins. When he had taken his vows, he was an inspiration and example to all by his observance of the rule and zealously promoted the cult of the Virgin Mother of God and of the Rosary. He asked God that he might die as a martyr for the Catholic faith, and this was granted him. For he was appointed leader of the missions which the Congregation for the Propagation of the Faith had established at this time for the Grisons and, since he spared himself no toil and converted many heretics to Christ, he incurred the hatred of evil men. So, on the 24th of April in the year 1622, at a church in a place called Sevis, he was struck down with blows and wounds by some heretics who had deceitfully invited him to meet them, pretending to be converts. And thus he consecrated with his own blood the first-fruits of the Martyrs of the Congregation for the Propagation of the Faith.\par
\switchcolumn*

\LA
\TeDeum{Te Deum}
\switchcolumn
\EN
\TeDeum{Te Deum}
\switchcolumn*

\end{paracol}

\par\needspace{10\baselineskip}\addvspace{1.2em}{\centering\small\textsc{Die 25 Aprilis} · \textsc{25 April}\par}
\Office{S. Marci Evangelistæ}{St. Mark the Evangelist}{Duplex II classis}
\addcontentsline{toc}{subsection}{Die 25 Aprilis · S. Marci Evangelistæ\enspace\textit{St. Mark the Evangelist}}
\begin{paracol}{2}
\LA
\RefLine{Lectiones I–III: ex Commúni Evangelistarum tempore Paschali (C1ap).}
\switchcolumn
\EN
\RefLine{Lessons I–III: from the Commune Evangelistarum tempore Paschali (C1ap).}
\switchcolumn*

\LA
\RefLine{Lectiones VII–IX: ex Commúni Evangelistarum tempore Paschali (C1ap).}
\switchcolumn
\EN
\RefLine{Lessons VII–IX: from the Commune Evangelistarum tempore Paschali (C1ap).}
\switchcolumn*

\LA
\LessonHead{Lectio IV}
\switchcolumn
\EN
\LessonHead{Lesson IV}
\switchcolumn*

\LA
\LessonTitle{Ex libro sancti Hierónymi Presbýteri de Scriptóribus ecclesiásticis}
\switchcolumn
\EN
\LessonTitle{From the Book upon Church Writers, composed by St. Jerome, Priest at Bethlehem.}
\switchcolumn*

\LA
\Source{De viris illustribus, cap 8}
\switchcolumn
\EN
\Source{On Illustrious Men, chapter 8}
\switchcolumn*

\LA
\noindent Marcus, discípulus et intérpres Petri, juxta quod Petrum referéntem audíerat, rogátus Romæ a frátribus, breve scripsit Evangélium. Quod cum Petrus audísset, et probávit, et Ecclésiæ legéndum sua auctoritáte dedit. Assúmpto ítaque Evangélio quod ipse confécerat, perréxit in Ægýptum, et primus Alexandríæ Christum annúntians, constítuit ecclésiam tanta doctrína et vitæ continéntia, ut omnes sectatóres Christi ad exémplum sui cógeret.\par
\switchcolumn
\EN
\noindent Mark was the disciple and interpreter of Peter, and it was from what he had heard Peter tell, that, at the request of the brethren at Rome, he wrote the shortest of the Gospels. When Peter had heard it, he approved it, and gave it to the Church to be read, by his authority. Mark betook himself to Egypt, with the Gospel which he had compiled, and was the first man who preached Christ at Alexandria. There he founded a Church with such teaching and austerity of life, that all who followed Christ were constrained to imitate him.\par
\switchcolumn*

\LA
\begin{Resp}\Rsign Lux perpétua lucébit Sanctis tuis, Dómine, \Star{} Et ætérnitas témporum, allelúja, allelúja.\end{Resp}
\begin{Resp}\Vsign Lætítia sempitérna erit super cápita eórum: gáudium et exsultatiónem obtinébunt.\end{Resp}
\begin{Resp}\Rsign Et ætérnitas témporum, allelúja, allelúja.\end{Resp}
\switchcolumn
\EN
\begin{Resp}\Rsign Eternal light will shine over your saints, O Lord, \Star{} And the eternity of the times, alleluia, alleluia.\end{Resp}
\begin{Resp}\Vsign Everlasting joy shall be upon their heads, they shall obtain joy and gladness\end{Resp}
\begin{Resp}\Rsign And the eternity of the times, alleluia, alleluia.\end{Resp}
\switchcolumn*

\LA
\LessonHead{Lectio V}
\switchcolumn
\EN
\LessonHead{Lesson V}
\switchcolumn*

\LA
\noindent Dénique Philo, disertíssimus Judæórum, videns Alexandríæ primam ecclésiam adhuc judaizántem, quasi in laudem gentis suæ, librum super eórum conversatióne scripsit. Et quómodo Lucas narrat Jerosólymæ credéntes ómnia habuísse commúnia: sic et ille, quod Alexandríæ sub Marco fíeri doctóre cernébat, memóriæ trádidit. Mórtuus est autem octávo Nerónis anno, et sepúltus Alexandríæ, succedénte sibi Aniáno.\par
\switchcolumn
\EN
\noindent Last of all, Philo, that most learned Jew, observing that the first Church of Alexandria still kept the law of Moses, wrote a book concerning their manners, as if in praise of his own nation, wherein he saith that under the teaching of Mark, the Christians of Alexandria had all things in common, just as Luke telleth us was the case with all them that believed at Jerusalem. Mark died in the eighth year of Nero, and was buried at Alexandria. Anianus succeeded him.\par
\switchcolumn*

\LA
\begin{Resp}\Rsign Virtúte magna reddébant Apóstoli \Star{} Testimónium resurrectiónis Jesu Christi Dómini nostri, allelúja, allelúja.\end{Resp}
\begin{Resp}\Vsign Repléti quidem Spíritu Sancto, loquebántur cum fidúcia verbum Dei.\end{Resp}
\begin{Resp}\Rsign Testimónium resurrectiónis Jesu Christi Dómini nostri, allelúja, allelúja.\end{Resp}
\switchcolumn
\EN
\begin{Resp}\Rsign With great power did the Apostles \Star{} Give testimony of the resurrection of Jesus Christ our Lord, alleluia, alleluia.\end{Resp}
\begin{Resp}\Vsign And they were all filled with the Holy Ghost: and they spoke the word of God with confidence.\end{Resp}
\begin{Resp}\Rsign Give testimony of the resurrection of Jesus Christ our Lord, alleluia, alleluia.\end{Resp}
\switchcolumn*

\LA
\LessonHead{Lectio VI}
\switchcolumn
\EN
\LessonHead{Lesson VI}
\switchcolumn*

\LA
\LessonTitle{De Expositióne sancti Gregórii Papæ super Ezechiélem Prophétam}
\switchcolumn
\EN
\LessonTitle{From the Exposition of the Book of the Prophet Ezekiel by Pope St. Gregory (the Great.)}
\switchcolumn*

\LA
\Source{Homilia 3, Lib. 1}
\switchcolumn
\EN
\Source{Hom. 3, Bk. I}
\switchcolumn*

\LA
\noindent Sancta quátuor animália, quæ prophetíæ spíritu futúra prævidéntur, subtíli narratióne describúntur, cum dícitur: Quátuor fácies uni, et quátuor pennæ uni. Quid per fáciem, nisi notítia; et quid per pennas, nisi volátus exprímitur? Per fáciem quippe unusquísque cognóscitur: per pennas vero in altum ávium córpora sublevántur. Fácies ítaque ad fidem pértinet, penna ad contemplatiónem. Per fidem namque ab omnipoténti Deo cognóscimur, sicut ipse de suis óvibus dicit: Ego sum pastor bonus, et cognósco oves meas, et cognóscunt me meæ. Qui rursus ait: Ego scio quos elégerim. Per contemplatiónem vero, qua super nosmetípsos tóllimur, quasi in áëra levámur.\par
\switchcolumn
\EN
\noindent The Prophet writeth very minutely touching the four holy living creatures, whom he saw in the spirit as being to come. He saith: “Every one had four faces, and every one had four wings.” What signifieth the face save likeness whereby we are known? or wings, save the power to fly since it is by the face that man is known from man, and by their wings that the birds' bodies are carried up into the air. So the face pertaineth to certitude, and the wings to contemplation. With certitude we are known of God Almighty, Who saith: “I am the Good Shepherd, and know My sheep, and am known of Mine.” (John x. 14.) And again: “I know whom I have chosen.” (xiii. 18.) And by contemplation, whereby we rise above ourselves, we as it were fly heavenwards.\par
\switchcolumn*

\LA
\begin{Resp}\Rsign Isti sunt agni novélli, qui annuntiavérunt, allelúja, modo venérunt ad fontes, \Star{} Repléti sunt claritáte, allelúja, allelúja.\end{Resp}
\begin{Resp}\Vsign In conspéctu Agni amícti sunt stolis albis, et palmæ in mánibus eórum.\end{Resp}
\begin{Resp}\Rsign Repléti sunt claritáte, allelúja, allelúja.\end{Resp}
\begin{Resp}\Vsign Glória Patri, et Fílio, \Star{} et Spirítui Sancto.\end{Resp}
\begin{Resp}\Rsign Repléti sunt claritáte, allelúja, allelúja.\end{Resp}
\switchcolumn
\EN
\begin{Resp}\Rsign These are the young lambs, who were promised, they are just coming to the fountains. \Star{} They are filled with glory, alleluia, alleluia.\end{Resp}
\begin{Resp}\Vsign In sight of the Lamb, clothed with white robes, and palms in their hands.\end{Resp}
\begin{Resp}\Rsign They are filled with glory, alleluia, alleluia.\end{Resp}
\begin{Resp}\Vsign Glory be to the Father, and to the Son, \Star{} and to the Holy Ghost.\end{Resp}
\begin{Resp}\Rsign They are filled with glory, alleluia, alleluia.\end{Resp}
\switchcolumn*

\end{paracol}

\par\needspace{10\baselineskip}\addvspace{1.2em}{\centering\small\textsc{Die 26 Aprilis} · \textsc{26 April}\par}
\Office{SS. Cleti et Marcellini Summorum Pontificum et Martyrum}{Sts. Cletus and Marcellinus, Popes and Martyrs}{Semiduplex}
\addcontentsline{toc}{subsection}{Die 26 Aprilis · SS. Cleti et Marcellini Summorum Pontificum et Martyrum\enspace\textit{Sts. Cletus and Marcellinus, Popes and Martyrs}}
\begin{paracol}{2}
\LA
\RefLine{Lectiones I–III: de Scriptura occurrente.}
\switchcolumn
\EN
\RefLine{Lessons I–III: from the Scripture of the day.}
\switchcolumn*

\LA
\RefLine{Lectiones VII–IX: ex Commúni plurium Summorum Pontificum Martyrum Tempore Paschali (C3bp).}
\switchcolumn
\EN
\RefLine{Lessons VII–IX: from the Commune plurium Summorum Pontificum Martyrum Tempore Paschali (C3bp).}
\switchcolumn*

\LA
\LessonHead{Lectio IV}
\switchcolumn
\EN
\LessonHead{Lesson IV}
\switchcolumn*

\LA
\noindent Cletus Romanus, patre Æmiliano, de regione quinta, e vico Patricio, imperatoribus Vespasiano et Tito, Ecclesiam gubernavit. Is ex præcepto Principis Apostolorum in Urbe viginti quinque presbyteros ordinavit. Primus in litteris verbis illis usus est: Salutem et apostolicam benedictionem. Qui Ecclesia optime constituta, cum ei præfuísset annos duodecim, menses septem, dies duos, Domitiano imperatore, secunda post Neronem persecutione, martyrio coronatus est, et in Vaticano juxta corpus beati Petri sepultus.\par
\switchcolumn
\EN
\noindent Cletus was a Roman, the son of Emilian, of the Fifth Region of the city, and the street called Noble. He ruled the Church in the time of the Emperors Vespasian and Titus. In accordance with the precept of the Prince of the Apostles He ordained twenty-five Priests for the city. He was the first Pope who made use in his letters of the phrase “Health and Apostolic Benediction.” When he had ruled the Church for twelve years, seven months, and two days, and brought it into an excellent state of order, in the reign of the Emperor Domitian, and the second persecution since the time of Nero, he was crowned with martyrdom, and buried on the Vatican mount, hard by the body of blessed Peter.\par
\switchcolumn*

\LA
\begin{Resp}\Rsign Lux perpétua lucébit Sanctis tuis, Dómine, \Star{} Et ætérnitas témporum, allelúja, allelúja.\end{Resp}
\begin{Resp}\Vsign Lætítia sempitérna erit super cápita eórum: gáudium et exsultatiónem obtinébunt.\end{Resp}
\begin{Resp}\Rsign Et ætérnitas témporum, allelúja, allelúja.\end{Resp}
\switchcolumn
\EN
\begin{Resp}\Rsign Eternal light will shine over your saints, O Lord, \Star{} And the eternity of the times, alleluia, alleluia.\end{Resp}
\begin{Resp}\Vsign Everlasting joy shall be upon their heads, they shall obtain joy and gladness\end{Resp}
\begin{Resp}\Rsign And the eternity of the times, alleluia, alleluia.\end{Resp}
\switchcolumn*

\LA
\LessonHead{Lectio V}
\switchcolumn
\EN
\LessonHead{Lesson V}
\switchcolumn*

\LA
\noindent Marcellinus Romanus, ab anno ducentesimo nonagesimo sexto ad annum trecentesimum quartum in immani imperatoris Diocletiani persecutione Ecclesiæ præfuit. Multas pertulit angustias ob improbam eorum severitatem, qui eum redarguebant de nimia indulgentia erga lapsos in idololatriam, quæque causa fuit, ut per calumniam infamatus fuerit, quasi thus idolis adhibuisset. Verum hic beatus Pontifex in confessione fidei, una cum tribus aliis Christianis, Claudio, Cyrino et Antonino, capite plexus est. Quorum projecta corpora, cum triginta sex dies jussu imperatoris sepultura caruissent, beatus Marcellus a sancto Petro in somnis admonitus, cum presbyteris et diaconis, hymnis et luminibus adhibitis, honorifice sepelienda curavit in cœmeterio Priscillæ via Salaria. Rexit Ecclesiam annos septem, menses undecim, dies viginti tres: quo tempore fecit ordinationes duas mense Decembri, quibus creavit presbyteros quatuor, episcopos per diversa loca quinque.\par
\switchcolumn
\EN
\noindent Marcellinus was a Roman; he ruled the Church from the year 296 to the year 304, during the savage persecution which was ordered by the Emperor Diocletian. He suffered through the false severity of those who blamed him as being too indulgent toward them who had fallen into idolatry, and for this reason also hath been slandered to the effect that he himself burnt incense to idols but this blessed Pope, on account of his confession of the faith, was put to death along with three other Christians, whose names are Claudius, Cyrinus, and Antoninus. At the command of the Emperor their bodies were cast out unburied, and lay so for thirty-six days. At the end of that time St. Peter appeared in a dream to Blessed Marcellus, and in obedience to his command the said Marcellus went with certain Priests and Deacons, singing hymns, and carrying lights, and buried these four bodies honourably in the Cemetery of Priscilla upon the Salarian Way. Marcellinus ruled the Church for seven years, eleven months, and twenty-three days. During this time he held two Advent ordinations, and ordained at them four Priests, and five Bishops for diverse Sees.\par
\switchcolumn*

\LA
\begin{Resp}\Rsign In servis suis, allelúja, \Star{} Consolábitur Deus, allelúja.\end{Resp}
\begin{Resp}\Vsign Judicábit Dóminus pópulum suum, et in servis suis.\end{Resp}
\begin{Resp}\Rsign Consolábitur Deus, allelúja.\end{Resp}
\switchcolumn
\EN
\begin{Resp}\Rsign His servants, alleluia \Star{} God will comfort, alleluia, alleluia.\end{Resp}
\begin{Resp}\Vsign The Lord will judge His people and, His servants,\end{Resp}
\begin{Resp}\Rsign God will comfort, alleluia, alleluia.\end{Resp}
\switchcolumn*

\LA
\LessonHead{Lectio VI}
\switchcolumn
\EN
\LessonHead{Lesson VI}
\switchcolumn*

\LA
\LessonTitle{Sermo sancti Ambrósii Epíscopi.}
\switchcolumn
\EN
\LessonTitle{From the Sermons of St. Ambrose, Bishop of Milan}
\switchcolumn*

\LA
\Source{Serm. 22.}
\switchcolumn
\EN
\Source{xxii.}
\switchcolumn*

\LA
\noindent Dignum et congruum est, fratres, ut post lætitiam Paschæ, quam in Ecclesia celebravimus, gaudia nostra cum sanctis Martyribus conferamus: et iis annuntiemus Dominicæ resurrectionis gloriam, qui consortes sunt Dominicæ passionis. Qui enim socii sunt contumeliæ, debent et participes esse lætitiæ. Ita enim dicit beatus Apostolus: Sicut socii passionum estis, et resurrectionis eritis; si tolerabimus, inquit, et conregnabimus. Qui ergo toleraverunt mala propter Christum, debent et gloriam habere cum Christo.\par
\switchcolumn
\EN
\noindent Dearly beloved brethren, it is very meet and right that after the gladness of Easter, which we have celebrated in the Church, we should mingle our own joy with the joy of the holy Martyrs; yea, that we should tell of the glory of the Lord's rising again, to them that have been made partakers of the Lord's sufferings. It truly must needs be that they which have been partakers of His sufferings, should be also of His joy. For thus saith the blessed Apostle: “As ye are partakers of the sufferings, so shall ye be also of the consolation.” (2 Cor. i. 7.) And again: “If we suffer, we shall also reign with Him.” (2 Tim. ii. 12.) He, therefore, that endureth sorrow for Christ, must needs also have glory with Christ.\par
\switchcolumn*

\LA
\begin{Resp}\Rsign Fíliæ Jerúsalem, veníte et vidéte Mártyres cum corónis, quibus coronávit eos Dóminus \Star{} In die solemnitátis et lætítiæ, allelúja.\end{Resp}
\begin{Resp}\Vsign Quóniam confortávit seras portárum tuárum, benedíxit fílios tuos in te.\end{Resp}
\begin{Resp}\Rsign In die solemnitátis et lætítiæ, allelúja.\end{Resp}
\begin{Resp}\Vsign Glória Patri, et Fílio, \Star{} et Spirítui Sancto.\end{Resp}
\begin{Resp}\Rsign In die solemnitátis et lætítiæ, allelúja.\end{Resp}
\switchcolumn
\EN
\begin{Resp}\Rsign Come forth, O ye daughters of Jerusalem, and behold the Martyrs with the crowns wherewith the Lord crowned them; \Star{} In the day of His feasting, and of His gladness. Alleluia.\end{Resp}
\begin{Resp}\Vsign For He hath strengthened the bars of thy gates He hath blessed thy children within thee.\end{Resp}
\begin{Resp}\Rsign In the day of His feasting, and of His gladness. Alleluia.\end{Resp}
\begin{Resp}\Vsign Glory be to the Father, and to the Son, \Star{} and to the Holy Ghost.\end{Resp}
\begin{Resp}\Rsign In the day of His feasting, and of His gladness. Alleluia.\end{Resp}
\switchcolumn*

\LA
\LessonHead{Lectio IX (pro commemoratione)}
\switchcolumn
\EN
\LessonHead{Lesson IX (when commemorated)}
\switchcolumn*

\LA
\LessonTitle{Commemoratio: SS. Cleti et Marcellini Summorum Pontificum et Martyrum}
\switchcolumn
\EN
\LessonTitle{Commemoration of: Sts. Cletus and Marcellinus, Popes and Martyrs}
\switchcolumn*

\LA
\noindent Cletus Románus, imperatóribus Vespasiáno et Tito Ecclésiam gubernávit. Ex præcépto Príncipis Apostolórum in Urbe vigínti quinque presbýteros ordinávit. Primus in lítteris verbis illis usus est: Salútem et apostólicam benedictiónem. Ecclésia óptime constitúta, Domitiáno imperatóre, secúnda post Nerónem persecutióne, martýrio coronátus est, et in Vaticáno juxta corpus beáti Petri sepúltus. Marcellínus Románus, in immáni imperatóris Diocletiáni persecutióne Ecclésiæ præfuit. Multas pértulit angústias ob ímprobam eórum severitátem, qui eum redarguébant de nímia indulgéntia erga lapsos in idololatríam, quæque causa fuit, ut per calúmniam infamátus fúerit, quasi thus idólis adhibuísset. Verum hic beátus Póntifex in confessióne fídei, una cum tribus áliis Christiánis, Cláudio, Cyríno et Antoníno, cápite plexus est.\par
\switchcolumn
\EN
\noindent Cletus was a Roman, and ruled the Church in the time of the Emperors Vespasian and Titus. In accordance with the precept of the Prince of the Apostles he ordained twenty-five priests for the City. He was the first Pope who made use in his letters of the phrase, Health and Apostolic Benediction. Having brought the Church into an excellent state of order, he was crowned with martyrdom in the reign of the Emperor Domitian, and the second persecution since the time of Nero, and was buried on the Vatican mount, hard by the body of blessed Peter. Marcellinus was a Roman; he ruled the Church during the savage persecution which was ordered by the Emperor Diocletian. He suffered through the false severity of those who blamed him as being too indulgent toward them who had fallen into idolatry, and for this reason also hath been slandered to the effect that he himself burnt incense to idols; but this blessed Pope, on account of his confession of the faith, was put to death along with three other Christians, whose names are Claudius, Cyrinus and Antoninus.\par
\switchcolumn*

\LA
\TeDeum{Te Deum}
\switchcolumn
\EN
\TeDeum{Te Deum}
\switchcolumn*

\end{paracol}

\par\needspace{10\baselineskip}\addvspace{1.2em}{\centering\small\textsc{Die 27 Aprilis} · \textsc{27 April}\par}
\Office{S. Petri Canisii Confessoris et Ecclesiæ Doctoris}{St. Peter Canisius, Confessor and Doctor of the Church}{Duplex}
\addcontentsline{toc}{subsection}{Die 27 Aprilis · S. Petri Canisii Confessoris et Ecclesiæ Doctoris\enspace\textit{St. Peter Canisius, Confessor and Doctor of the Church}}
\begin{paracol}{2}
\LA
\RefLine{Lectiones I–III: de Scriptura occurrente.}
\switchcolumn
\EN
\RefLine{Lessons I–III: from the Scripture of the day.}
\switchcolumn*

\LA
\LessonHead{Lectio IV}
\switchcolumn
\EN
\LessonHead{Lesson IV}
\switchcolumn*

\LA
\noindent Petrus Canisius, Noviomagi in Gelria eo ipso anno natus est, quo Lutherus in Germania aperta rebellione ab Ecclesia descivit, et Ignatius de Loyola, in Hispania, terrestri militia abdicata, ad prælianda prælia Domini se convertit; Deo nimirum portendente, quos ille posthac adversarios, quem sacræ militiæ ducem esset habiturus. Coloniæ Agrippinæ, quo studiorum causa concesserat, perpetuo castitatis voto se Deo obstrinxit, et paulo post Societati Jesu nomen dedit. Sacerdotio auctus, catholicam fidem contra novatorum insidias legationibus, sermonibus, scriptis libris statim defendendam suscepit. Ob præclaram sapientiam et exploratum rerum usum a Cardinali Augustano et a pontificiis Legatis magnopere expetitus, semel atque iterum concilio Tridentino interfuit; cujus étiam decreta ex auctoritate Pii quarti Pontificis maximi rite per Germaniam promulganda et in morem inducenda curavit. A Paulo quarto ad conventum Petricoviensem ire jussus, aliisque a Gregorio decimo tertio legationibus obeundis adhibitus, alacri semper et numquam fracto difficultatibus animo, gravissima, religionis negotia tractávit, ac vel inter præsentia vitæ discrimina ad felicem exitum perduxit.\par
\switchcolumn
\EN
\noindent Peter Canisius was born at Nijmegen in Gelderland, the Netherlands, in the very year in which Luther openly rebelled against the Church in Germany, and in which Ignatius Loyola in Spain gave up earthly warfare to fight the battles of the Lord; God thus shewed what adversaries he was to encounter, and under whose leadership he was to fight. He made his studies at Cologne, where he took a vow to God of perpetual chastity, and shortly afterwards entered the Society of Jesus. After his ordination as priest, he began at once to defend the Catholic faith against the wiles of the innovators by missions, sermons, and writing books. His eminent wisdom and experience caused the Cardinal of Augsburg and the papal legates to invite him to the Council of Trent, and he was present at its sittings more than once. Moreover, by the authority of the Supreme Pontiff, Pius IV, he was entrusted with the charge of making its decrees known in Germany and carrying them into effect. Paul IV sent him to the Diet of Petrikau, and Gregory XIII entrusted him with the carrying out of other missions, all of which he undertook with an eager spirit, never conquered by any difficulties, and carried the most important affairs of religion through all the crises of this present life to a successful end.\par
\switchcolumn*

\LA
\begin{Resp}\Rsign Honéstum fecit illum Dóminus, et custodívit eum ab inimícis, et a seductóribus tutávit illum: \Star{} Et dedit illi claritátem ætérnam, allelúja.\end{Resp}
\begin{Resp}\Vsign Justum dedúxit Dóminus per vias rectas, et osténdit illi regnum Dei.\end{Resp}
\begin{Resp}\Rsign Et dedit illi claritátem ætérnam, allelúja.\end{Resp}
\switchcolumn
\EN
\begin{Resp}\Rsign The Lord made him honourable, and defended him from his enemies, and kept him safe from those that lay in wait for him; \Star{} And gave him perpetual glory, alleluia.\end{Resp}
\begin{Resp}\Vsign He went down with him into the pit, and left him not in bonds.\end{Resp}
\begin{Resp}\Rsign And gave him perpetual glory, alleluia.\end{Resp}
\switchcolumn*

\LA
\LessonHead{Lectio V}
\switchcolumn
\EN
\LessonHead{Lesson V}
\switchcolumn*

\LA
\noindent Superno caritátis igne, quem in Basilica Vaticana e penetralibus Cordis Jesu olim copiose hauserat, inflammatus. et divinæ gloriæ amplificandæ unice intentus, dici vix potest, quot per annos amplius quadraginta labores susceperit. ærumnasque pertulerit, ut complures Germaniæ civitates ac provincias vel ab hæreseos contagione defenderet, vel hæresi infectas catholicæ fidei restitueret. In Ratisbonensi et Augustano conventu imperii proceres ad jura Ecclesiæ tuenda et mores populorum emendandos excitavit: in Vormatiensi insolescentes impietatis magistros ad silentium adegit. A sancto Ignatio Germaniæ superioris provinciæ præfectus, domos et collegia multis locis condidit. Collegium Germanicum, Romæ constitutum omni ope provehere atque amplificare studuit; in academiis sacrarum humanarumque disciplinarum studia, miserandum in modum collapsa instauravit; contra Centuriatores Magdeburgenses duo volumina egregie conscripsit: et summam doctrinæ christianæ theologorum judicio et publico trium sæculorum usu ubique probatissimam, alíaque complura ad populorum institutionem valde accommodata in vulgus edidit. Quam ob rem hæreticorum malleus et alter Germaniæ apostolus appellatus, plane dignus habitus est, qui ad tutandam in Germania religionem divinitus electus putaretur\par
\switchcolumn
\EN
\noindent Inflamed with the heavenly fire of charity, which he had once received in the Vatican basilica from the sanctuary of the Heart of Jesus, and intent only on increasing the glory of God, it is almost impossible to describe how, for more than forty years, he took upon himself laborious tasks, and endured hardship, that he might defend very many cities and provinces of Germany from the contagion of heresy, or restore to the Catholic faith those that were infected with heresy. At the Diets of Ratisbon and Augsburg, he exhorted the princes of the Empire to defend the rights of the Church and reform the lives of their subjects. At Worms he reduced the insolent teachers of impiety to silence. St. Ignatius made him prefect of the province of Upper Germany, where he founded houses and colleges in many places. He used every effort to advance and enlarge the German College founded at Rome; he restored the study of sacred and profane learning in academies, which had fallen into a wretched condition. He wrote two excellent volumes against the Centuriators of Magdeburg; and he edited a summary of Christian doctrine, which has been thoroughly approved by the judgment of theologians and by common use everywhere for three centuries, as well as very many other works useful for public instruction in the vulgar tongue. For all these reasons he was called the Hammer of the Heretics, and the Second Apostle of Germany, and is rightly thought to have been worthy of having been chosen by God to protect religion in Germany.\par
\switchcolumn*

\LA
\begin{Resp}\Rsign Amávit eum Dóminus, et ornávit eum: stolam glóriæ índuit eum, \Star{} Et ad portas paradísi coronávit eum, allelúja.\end{Resp}
\begin{Resp}\Vsign Índuit eum Dóminus lorícam fídei, et ornávit eum.\end{Resp}
\begin{Resp}\Rsign Et ad portas paradísi coronávit eum, allelúja.\end{Resp}
\switchcolumn
\EN
\begin{Resp}\Rsign The Lord loved him and beautified him; He clothed him with a robe of glory, \Star{} And crowned him at the gates of Paradise, alleluia.\end{Resp}
\begin{Resp}\Vsign The Lord hath put on him the breast-plate of faith, and hath adorned him.\end{Resp}
\begin{Resp}\Rsign And crowned him at the gates of Paradise, alleluia.\end{Resp}
\switchcolumn*

\LA
\LessonHead{Lectio VI}
\switchcolumn
\EN
\LessonHead{Lesson VI}
\switchcolumn*

\LA
\noindent Inter hæc, precatione crebra et assidua rerum supernarum commentatione, lacrimis sæpe perfusus et ánimo interdum a sensibus abducto. Deo se conjungere solitus erat. A viris principibus vel sanctitate clarissimis et a quatuor summis Pontificibus magno in honore habitus, adeo de se demisse sentiebat ut se omnium minimum et diceret et haberet. Vindobonensem episcopatum semel iterum ac tertio recusavit. Moderatoribus suis obsequentissimus paratus erat ad ipsorum nutum omnia relinquere aut aggredi étiam cum valetudinis et vitæ periculo. Voluntaria sui ipsius coërcitione castitatem perpetuo sepsit. Demum Friburgi Helvetiorum, ubi plurimum pro Dei gloria et salute animarum ultimis vitæ suæ annis desudaverat, migravit ad Deum die vigesima prima Decembris anno millesimo quingentesimo nonagesimo septimo, ætatis suæ septimo supra septuagesimum. Hunc vero strenuum catholicæ veritatis propugnatorem Pius Papa nonus cælitum Beatorum honoribus adauxit; novis autem fulgentem signis Pius undecimus Pontifex maximus, anno jubilæi Sanctorum numero accensuit, simulque Doctorem universalis Ecclesiæ declaravit.\par
\switchcolumn
\EN
\noindent In these activities he was accustomed to unite himself to God by frequent prayer and assiduous meditation on heavenly things, often bathed in tears and sometimes with his soul rapt in ecstasy. He was held in great honor by men of rank, or of most distinguished holiness, and by four of the Supreme Pontiffs, but he thought so humbly of himself, that he spoke of and held himself as the least of all. He refused the bishopric of Vienna no less than three times. He was most obedient to his superiors, and ready at their mere nod to stop or to undertake all labors, even at the risk of his health and life. He guarded his chastity with perpetual voluntary self-mortification. At length, at Fribourg in Switzerland, where during the last years of his life he had labored much for the glory of God and the salvation of souls, he passed to God on the 21st day of December, 1597, in the seventy-seventh year of his age. This zealous champion of Catholic truth was adorned with the heavenly honors of the blessed by Pope Pius IX; and, as fresh miracles added to his renown, the Supreme Pontiff Pius XI, in the year of the Jubilee, included him among the Saints, and at the same time declared him a Doctor of the Universal Church.\par
\switchcolumn*

\LA
\begin{Resp}\Rsign Iste homo perfécit ómnia quæ locútus est ei Deus, et dixit ad eum: Ingrédere in réquiem meam: \Star{} Quia te vidi justum coram me ex ómnibus géntibus, allelúja.\end{Resp}
\begin{Resp}\Vsign Iste est, qui contémpsit vitam mundi, et pervénit ad cæléstia regna.\end{Resp}
\begin{Resp}\Rsign Quia te vidi justum coram me ex ómnibus géntibus, allelúja.\end{Resp}
\begin{Resp}\Vsign Glória Patri, et Fílio, \Star{} et Spirítui Sancto.\end{Resp}
\begin{Resp}\Rsign Quia te vidi justum coram me ex ómnibus géntibus, allelúja.\end{Resp}
\switchcolumn
\EN
\begin{Resp}\Rsign This is he which did according unto all that God commanded him; and God said unto him: Enter thou into My rest, \Star{} For thee have I seen righteous before Me among all people, alleluia.\end{Resp}
\begin{Resp}\Vsign This is he which loved not his life in this world, and is come unto an everlasting kingdom.\end{Resp}
\begin{Resp}\Rsign For thee have I seen righteous before Me among all people, alleluia.\end{Resp}
\begin{Resp}\Vsign Glory be to the Father, and to the Son, \Star{} and to the Holy Ghost.\end{Resp}
\begin{Resp}\Rsign For thee have I seen righteous before Me among all people, alleluia.\end{Resp}
\switchcolumn*

\LA
\LessonHead{Lectio VII}
\switchcolumn
\EN
\LessonHead{Lesson VII}
\switchcolumn*

\LA
\LessonTitle{Léctio sancti Evangélii secúndum Matthǽum.}
\switchcolumn
\EN
\LessonTitle{From the Holy Gospel according to Matthew}
\switchcolumn*

\LA
\Cite{Matt 5:13-19}
\switchcolumn
\EN
\Cite{Matt 5:13-19}
\switchcolumn*

\LA
\noindent In illo témpore: Dixit Jesus discípulis suis: Vos estis sal terræ. Quod si sal evanúerit, in quo saliétur? Et réliqua.\par
\switchcolumn
\EN
\noindent At that time: Jesus said unto his disciples: Ye are the salt of the earth: But if the salt have lost his savor, wherewith shall it be salted? And so forth.\par
\switchcolumn*

\LA
\LessonTitle{Homilia sancti Petri Canisii Presbýteri}
\switchcolumn
\EN
\LessonTitle{Homily of St. Peter Canisius, Priest}
\switchcolumn*

\LA
\Source{Notæ In Evangelia In festo S. Martini Ep, post initium}
\switchcolumn
\EN
\Source{Note in the Gospel of the feast of St. Martin the Bishop, after the beginning}
\switchcolumn*

\LA
\noindent Amabo et colam missos a Christo Apostolos horumque successores in Evangelii semine spargendo sedulos et indefessos propagandi verbi cooperatores, qui jure testari possunt: Sic nos existimet homo ut ministros Christi et dispensatores mysteriorum Dei. Voluit enim Christus ut vigilantissimus ac fidelissimus paterfamilias per tales ministros ac legatos lucernam evangelicam igne cælitus demisso accendi et accensam non modio supponi, sed super candelabrum constitui, quæ suum splendorem longe lateque diffunderet omnesque tum Judæorum tum gentium vigentes tenebras et errores profligaret.\par
\switchcolumn
\EN
\noindent I shall always love and reverence the Apostles sent by Christ, and their successors as well, in their work of sowing the Gospel seed. For all such may justly say of themselves: Let a man so account of us, as of the ministers of Christ, and stewards of the mysteries of God. It was Christ himself, like a watchful and most faithful householder, who wished that the Gospel-lamp should be lighted by his ministers and stewards with fire sent down from heaven; and once lighted, that it should not be put under a bushel, but set upon a candlestick, so as to spread its brightness far and wide, and put to flight all darkness and error rife among both Jews and Gentiles.\par
\switchcolumn*

\LA
\begin{Resp}\Rsign Iste est qui ante Deum magnas virtútes operátus est, et de omni corde suo laudávit Dóminum: \Star{} Ipse intercédat pro peccátis ómnium populórum, allelúja.\end{Resp}
\begin{Resp}\Vsign Ecce homo sine queréla, verus Dei cultor, ábstinens se ab omni ópere malo, et pérmanens in innocéntia sua.\end{Resp}
\begin{Resp}\Rsign Ipse intercédat pro peccátis ómnium populórum, allelúja.\end{Resp}
\switchcolumn
\EN
\begin{Resp}\Rsign This is he which wrought great wonders before God, and praised the Lord with all his heart. \Star{} May he pray for all people, that their sins may be forgiven unto them, alleluia.\end{Resp}
\begin{Resp}\Vsign Behold a man without blame, a worshipper of God in truth, keeping himself clean from every evil work, and abiding still in his innocency.\end{Resp}
\begin{Resp}\Rsign May he pray for all people, that their sins may be forgiven unto them, alleluia.\end{Resp}
\switchcolumn*

\LA
\LessonHead{Lectio VIII}
\switchcolumn
\EN
\LessonHead{Lesson VIII}
\switchcolumn*

\LA
\noindent Etenim evangelico Doctori sat non est, verbo lucere populo, et vocem in deserto clamantem præstare multisque in pietate juvandis lingua operam dare, ne alioquin, si verbi ministerium prætermittat, canis mutus non valens latrare a propheta dicatur. Sed et ardere illum oportet, ut opere atque caritate instructus, munus suum ornet evangelicum. Paulumque ducem sequatur. Is quippe non contentus Ephesiorum episcopo demandare: Præcipe hoc et doce: labora sicut bonus miles Christi Jesu: constanter étiam apud amicos et inimicos evangelizávit, ac episcopis apud Ephesum collectis bona dixit conscientia: Vos scitis, quomodo nihil subtraxerim utilium, quominus annuntiarem vobis, et docerem vos publice et per domos, testificans Judæis atque Gentilibus in Deum pœnitentiam et fidem in Dominum nostrum Jesum Christum.\par
\switchcolumn
\EN
\noindent It is not enough for the Gospel-teacher to please the people with his speaking. He must also be the voice of one crying in the wilderness, and so by his eloquence call many to the good life. He must not be a dumb dog, not even able to bark, as spoken of by the Prophet Isaiah. Yea, he should also burn in such a way that, equipped with good works and love, he may adorn his evangelical office, and follow the leadership of Paul. He indeed was not satisfied with bidding the bishop of the Ephesians: This command and teach: conduct thyself in work as a good soldier of Christ Jesus; but he unflaggingly preached the Gospel to friend and foe alike, and said with a good conscience to the bishops gathered at Ephesus: Ye know how I kept back nothing that was profitable unto you, but have shewed you, and have taught you publicly, and from house to house, urging Jews and Gentiles to turn unto God in repentance and to believe in our Lord Jesus Christ.\par
\switchcolumn*

\LA
\begin{Resp}\Rsign In médio Ecclésiæ apéruit os ejus, \Star{} Et implévit eum Dóminus spíritu sapiéntiæ et intelléctus, allelúja.\end{Resp}
\begin{Resp}\Vsign Jucunditátem et exsultatiónem thesaurizávit super eum.\end{Resp}
\begin{Resp}\Rsign Et implévit eum Dóminus spíritu sapiéntiæ et intelléctus, allelúja.\end{Resp}
\begin{Resp}\Vsign Glória Patri, et Fílio, \Star{} et Spirítui Sancto.\end{Resp}
\begin{Resp}\Rsign Et implévit eum Dóminus spíritu sapiéntiæ et intelléctus, allelúja.\end{Resp}
\switchcolumn
\EN
\begin{Resp}\Rsign In the midst of the congregation did the Lord open his mouth. \Star{} And filled him with the spirit of wisdom and understanding, alleluia.\end{Resp}
\begin{Resp}\Vsign He made him rich with joy and gladness.\end{Resp}
\begin{Resp}\Rsign And filled him with the spirit of wisdom and understanding, alleluia.\end{Resp}
\begin{Resp}\Vsign Glory be to the Father, and to the Son, \Star{} and to the Holy Ghost.\end{Resp}
\begin{Resp}\Rsign And filled him with the spirit of wisdom and understanding, alleluia.\end{Resp}
\switchcolumn*

\LA
\LessonHead{Lectio IX}
\switchcolumn
\EN
\LessonHead{Lesson IX}
\switchcolumn*

\LA
\noindent Talem enim pastorem decet esse in Ecclesia, qui more Pauli, omnibus omnia fiat, ut in illo reperiat æger curationem, mæstus lætitiam, desperans fiduciam, imperitus doctrinam, dubius consilium, pœnitens veniam atque solatium, et quidquid tandem ad salutem est cuique necessarium. Quocirca pulchre Christus, cum primarios mundi Ecclesiæque, doctores constituere vellet, non sat habuit discipulis dicere: Vos estis lux mundi: sed étiam illud subjecit: Non potest civitas abscondi supra montem posita, neque accendunt lucernam et ponunt eam sub modio, sed super candelabrum, ut luceat omnibus qui in domo sunt. Errant enim opinantes ecclesiastæ, quod muneri suo doctrinæ splendore magis quam vitæ integritate et caritatis ardore possint satisfacere.\par
\switchcolumn
\EN
\noindent Such should be the shepherd in the Church who, like Paul, becometh all things to all men, so that the sick may find healing in him; the sad, joy; the desperate, hope; the ignorant, instruction; those in doubt, advice; the penitent, forgiveness and comfort, and finally, every one whatever is necessary unto salvation. And so Christ, when he wished to appoint the chief teachers of the world and of the Church, did not limit himself to saying to his disciples: Ye are the light of the world: but also added: A city that is set on an hill cannot be hid; neither do men light a candle, and put it under a bushel, but on a candlestick, and it giveth light unto all that are in the house. Those churchmen err who imagine that it is by brilliant preaching that they fulfill their office; rather, it is by holiness of life and an all-embracing love.\par
\switchcolumn*

\LA
\TeDeum{Te Deum}
\switchcolumn
\EN
\TeDeum{Te Deum}
\switchcolumn*

\LA
\LessonHead{Lectio IX (pro commemoratione)}
\switchcolumn
\EN
\LessonHead{Lesson IX (when commemorated)}
\switchcolumn*

\LA
\LessonTitle{Commemoratio: S. Petri Canisii Confessoris et Ecclesiæ Doctoris}
\switchcolumn
\EN
\LessonTitle{Commemoration of: St. Peter Canisius, Confessor and Doctor of the Church}
\switchcolumn*

\LA
\noindent Petrus Canísius, Noviómagi in Gélria natus, statim ac Societáti Jesu nomen déderat, cathólicam fidem contra novatórum insídias, legatiónibus, sermónibus, scriptis libris defendéndam suscépit. Divínæ glóriæ amplificándæ únice inténtus, dici vix potest, quot per annos ámplius quadragínta labóres ærumnásque pertúlerit. Semel atque íterum Concílio Tridentíno intérfuit, Germániæ regiónes bene multas peragrávit, omne hóminum genus salubérrimis institútis públice et privátim excóluit, complurésque urbes atque províncias vel ab hæréseos contagióne deféndit vel hǽresi inféctas cathólicæ fídei restítuit. A sancto Ignátio Germániæ Superióris provínciæ præféctus, domos et collégia multis locis cóndidit. Contra Centuriatóres Magdeburgénses duo egrégia volúmina conscrípsit, et summam doctrínæ christiánæ theologórum judício et longo populórum usu probatíssimam alíaque complúra in vulgus édidit; dignus proínde qui hæreticórum málleus et alter Germániæ apóstolus vocarétur. Dénique Fribúrgi Helvetiórum, séptimum supra septuagésimum annum agens, die vicésima prima decémbris anno millésimo quingentésimo nonagésimo séptimo quiévit in Dómino. Eum Pius Papa undécimus Sanctórum fastis adjúnxit et simul universális Ecclésiæ Doctórem declarávit.\par
\switchcolumn
\EN
\noindent Peter Canisius was born at Nijmegen in Gelderland, the Netherlands. After joining the Society of Jesus, he began at once to defend the Catholic faith against the wiles of the innovators by missions, sermons, and writing books. His one idea was to promote the greater glory of God and it would be impossible to relate how many labors and hardships he went through for more than forty years. He took part more than once in the Council of Trent, travelled on successful missions through many parts of Germany, instructed all classes of society with sound teaching, both public and private, and defended many cities and provinces against the contagion of heresy or brought back to the Catholic faith those tainted by heresy. He was made head of the province of Germany by St. Ignatius, and built houses and colleges in many places. Against the Centuriators of Magdeburg, he wrote two famous volumes, and brought out a summary of Christian doctrine highly approved both by the judgment of theologians and by its long continued use among the people, and many other works in the mother-tongue. From all this, he earned the title of Hammer of the Heretics and second Apostle of Germany. Finally, at Fribourg in Switzerland, at the age of seventy-seven, he rested in the Lord on the 21st day of December, 1597. Pius XI added him to the list of the Saints and at the same time proclaimed him a Doctor of the Universal Church.\par
\switchcolumn*

\LA
\TeDeum{Te Deum}
\switchcolumn
\EN
\TeDeum{Te Deum}
\switchcolumn*

\end{paracol}

\par\needspace{10\baselineskip}\addvspace{1.2em}{\centering\small\textsc{Die 28 Aprilis} · \textsc{28 April}\par}
\Office{S. Pauli a Cruce Confessoris}{St. Paul of the Cross, Confessor}{Duplex}
\addcontentsline{toc}{subsection}{Die 28 Aprilis · S. Pauli a Cruce Confessoris\enspace\textit{St. Paul of the Cross, Confessor}}
\begin{paracol}{2}
\LA
\RefLine{Lectiones I–III: de Scriptura occurrente.}
\switchcolumn
\EN
\RefLine{Lessons I–III: from the Scripture of the day.}
\switchcolumn*

\LA
\RefLine{Lectio IX: de S. Vitalis Martyris (04-28o).}
\switchcolumn
\EN
\RefLine{Lesson IX: of St. Vitalis, Martyr (04-28o).}
\switchcolumn*

\LA
\LessonHead{Lectio IV}
\switchcolumn
\EN
\LessonHead{Lesson IV}
\switchcolumn*

\LA
\noindent Paulus a Cruce, Uvádæ in Liguria natus, sed a Castellátio prope Alexandríam Statiellórum nobili genere oriundus, qua futurus esset sanctitáte clarus, innotuit miro splendore, qui noctu implévit pariéntis matris cubiculum, et insígni augustæ cæli Reginæ beneficio, quæ púerum in flumen delapsum a certo naufragio illæsum erípuit. A primo ratiónis usu Jesu Christi crucifixi amóre flagrans, ejus contemplatióni prolixius vacare cœpit, et carnem innocentíssimam vigiliis, flagellis, jejuniis, potu in sexta feria ex acéto cum felle mixto, ac dura quavis castigatióne contérere. Martyrii desidério incénsus, exercitui se adjunxit, qui Venetiis, ad bellum Turcis inferéndum, comparabátur; cógnita vero inter orándum Dei voluntáte, arma ultro réddidit, præstantiori milítiæ operam datúrus, quæ Ecclésiæ præsidio esse æternámque hóminum salútem procurare totis viribus niterétur. Reversus in pátriam, honestíssimis nuptiis sibique deláta patrui hereditáte recusatis, arctiórem inire crucis sémitam ac rudi túnica a suo episcopo indui vóluit. Tum ejus jussu, ob eminentem vitæ sanctimóniam et rerum divinárum sciéntiam, nondum clericus, Dominicum agrum, máximo cum animárum fructu, divini verbi prædicatióne excoluit.\par
\switchcolumn
\EN
\noindent Paul of the Cross was sprung of a noble family of the Danei, at Castellazzo, hard by Alessandria, in the Province of Acqui, in the territory of the then Republic of Genoa, but was born at Ovada, in the same province. The holiness with which he was afterwards to shine was foreshown by a strange light which filled his mother's chamber while she was in labour, and by the remarkable help which was bestowed upon him by the great Queen of Heaven, who delivered him unhurt from certain destruction when he was fallen into a river as a lad. From the first use of reason he burnt with love for Jesus crucified, and began to spend long times in contemplating Him. He chastised his innocent flesh with watching, scourging, fasting, and all severe hardships, and on Friday he drank vinegar mingled with gall. He was seized with a desire for martyrdom, and enlisted in the army which was being raised at Venice to fight against the Turks but in consequence of the Will of God, made known to him while he was in prayer, he left the army in order to serve in a more exalted regiment whose duty it should be to defend the Church and to toil for the eternal salvation of men. When he returned home he refused a very honourable marriage, and also the inheritance which was bequeathed to him by his father's brother, and would fain enter upon a straiter way of the cross and be clad by his own Bishop with a rough tunic. By command of the Bishop, on account of his eminent holiness of life and knowledge of the things of God, he began, even before he became a clerk, to toil in the Lord's field with great profit of souls by preaching the Word.\par
\switchcolumn*

\LA
\begin{Resp}\Rsign Honéstum fecit illum Dóminus, et custodívit eum ab inimícis, et a seductóribus tutávit illum: \Star{} Et dedit illi claritátem ætérnam, allelúja.\end{Resp}
\begin{Resp}\Vsign Justum dedúxit Dóminus per vias rectas, et osténdit illi regnum Dei.\end{Resp}
\begin{Resp}\Rsign Et dedit illi claritátem ætérnam, allelúja.\end{Resp}
\switchcolumn
\EN
\begin{Resp}\Rsign The Lord made him honourable, and defended him from his enemies, and kept him safe from those that lay in wait for him; \Star{} And gave him perpetual glory, alleluia.\end{Resp}
\begin{Resp}\Vsign He went down with him into the pit, and left him not in bonds.\end{Resp}
\begin{Resp}\Rsign And gave him perpetual glory, alleluia.\end{Resp}
\switchcolumn*

\LA
\LessonHead{Lectio V}
\switchcolumn
\EN
\LessonHead{Lesson V}
\switchcolumn*

\LA
\noindent Romam profectus, theologicis disciplinis rite imbutus, a summo Pontifice Benedicto décimo tertio ex obediéntia sacerdotio auctus est. Facta sibi ab eodem potestate aggregándi socios, in solitúdinem recéssit Argentarii montis, quo eum beáta Virgo jamprídem invitaverat, veste illi simul ostensa atri coloris, passiónis Fílii sui insígnibus, decorata, ibique fundaménta jecit novæ congregatiónis. Quæ brevi, plurimis ab eo tolerátis labóribus, præcláris aucta viris, cum Dei benedictióne valde succrevit; a Sede apostolica non semel confirmáta, una cum regulis, quas orando ipse a Deo acceperat, et quarto addito voto, pergratam Dominicæ passiónis memóriam promovendi. Sacras vírgines quoque instituit, quæ excéssum caritátis divini Sponsi sedulo meditaréntur. Hæc inter, animárum inexhausta aviditate ab Evangélii prædicatióne numquam deficiens, hómines pene innumeros, étiam perditíssimos aut in hæresim lapsos, in salútis trámitem adduxit. Præsertim Christi enárranda passióne, mirífica ejus oratiónis vis erat, qua una cum astántibus in fletum effúsus, quælibet obduráta corda ad pœniténtiam scindebat.\par
\switchcolumn
\EN
\noindent He betook himself to Rome, and when he had there studied a regular course of theology he was ordained Priest in obedience to the command of the Supreme Pontiff Benedict XIII., who also gave him permission to gather comrades around him. He withdrew to the solitude of Mount Argentaro, whither he had been already called by the Blessed Virgin, at which same time she also showed him in vision a black habit marked with the emblems of the sufferings of her Son. At Mount Argentaro, he laid the foundations of his new Congregation, which under the blessing of God grew quickly, through the labors of Paul, and attracted to it eminent men. It received the confirmation of the Apostolic See more than once, with the rules which Paul himself had received from God in prayer and the addition of a fourth vow, that, namely, to promote the blessed remembrance of the sufferings of the Lord. He founded a congregation of holy virgins also, who should dwell constantly upon the overflowing love of the Divine Bridegroom. Amid all these works his untiring love for souls caused him never to weary in the preaching of the Gospel, and he led into the path of salvation men almost countless, among whom were some of the most lost, or those who had fallen into heresy. The greatest and most wonderful power of his preaching was how he told of the sufferings of Christ, so that he himself and his hearers would alike burst into tears, and hardened hearts were cloven by repentance.\par
\switchcolumn*

\LA
\begin{Resp}\Rsign Amávit eum Dóminus, et ornávit eum: stolam glóriæ índuit eum, \Star{} Et ad portas paradísi coronávit eum, allelúja.\end{Resp}
\begin{Resp}\Vsign Índuit eum Dóminus lorícam fídei, et ornávit eum.\end{Resp}
\begin{Resp}\Rsign Et ad portas paradísi coronávit eum, allelúja.\end{Resp}
\switchcolumn
\EN
\begin{Resp}\Rsign The Lord loved him and beautified him; He clothed him with a robe of glory, \Star{} And crowned him at the gates of Paradise, alleluia.\end{Resp}
\begin{Resp}\Vsign The Lord hath put on him the breast-plate of faith, and hath adorned him.\end{Resp}
\begin{Resp}\Rsign And crowned him at the gates of Paradise, alleluia.\end{Resp}
\switchcolumn*

\LA
\LessonHead{Lectio VI}
\switchcolumn
\EN
\LessonHead{Lesson VI}
\switchcolumn*

\LA
\noindent Tanta in ejus péctore alebátur divinæ caritátis flamma, ut indusium quod erat cordi propius, sæpe véluti igne adustum, et binæ cóstulæ elátæ apparúerint. Sacrum præsertim fáciens non poterat a lácrimis temperare: frequénti quoque éxtasi cum mira intérdum corporis elevatióne frui, vultúque superna luce radiante conspiciebátur. Quandoque, dum concionarétur, cæléstis vox verba ei suggerentis audíta fuit, aut sermo ejus ad plura míllia passuum intónuit. Prophetíæ et linguárum dono, cordium scrutatióne, potestate in dæmones, in morbos, in eleménta enituit. Cumque ipsis summis Pontificibus carus ac venerandus esset, servum inutilem, peccatórem nequíssimum, a dæmóniis quoque conculcándum se judicábat. Tandem, asperrimi vitæ generis ad longam usque senectútem tenacíssimus, anno millesimo septingentésimo septuagesimo quinto, cum præclára mónita, véluti sui spíritus transmissa hereditáte, alumnis tradidísset, Ecclésiæ sacramentis ac cælésti visióne recreatus, Romæ, qua prædixerat die, migrávit in cælum. Eum Pius nonus Pontifex maximus in Beatórum, novisque deínde fulgentem signis, in Sanctórum númerum retulit.\par
\switchcolumn
\EN
\noindent The fire of the love of God burnt so in his heart that the part of his under-garment which was next thereto often presented the appearance of having been scorched, and two of his ribs seemed to be raised. He could not withhold his tears, more especially when he was saying Mass, and when he was in a state of trance, as oftentimes befell, his body was sometimes seen to be raised into the air, and his face to shine as with light from heaven. Sometimes when he was preaching a heavenly voice was heard prompting him, or his words became audible at the distance of several miles. He was eminent for the gifts of prophecy, of speaking with tongues, of reading the heart, and of power over evil spirits, over diseases, and over the inanimate elements of nature. The Supreme Pontiffs themselves regarded him as dear and venerable, but he held himself to be but an unprofitable servant, and a sinful wretch upon whom devils might well trample. He held to the bitter hardships of his life, even unto a great age, and passed to heaven from Rome, (upon the 18th day of October,) being the day which he had himself foretold, in the year 1775, after he had addressed to his disciples noble exhortations which are as the heritage of his spirit, and had been comforted by the sacraments of the Church, and by an heavenly vision. The Supreme Pontiff Pius IX. numbered his name among those of the blessed, and then, after renewed signs and wonders, among those of the Saints.\par
\switchcolumn*

\LA
\begin{Resp}\Rsign Iste homo perfécit ómnia quæ locútus est ei Deus, et dixit ad eum: Ingrédere in réquiem meam: \Star{} Quia te vidi justum coram me ex ómnibus géntibus, allelúja.\end{Resp}
\begin{Resp}\Vsign Iste est, qui contémpsit vitam mundi, et pervénit ad cæléstia regna.\end{Resp}
\begin{Resp}\Rsign Quia te vidi justum coram me ex ómnibus géntibus, allelúja.\end{Resp}
\begin{Resp}\Vsign Glória Patri, et Fílio, \Star{} et Spirítui Sancto.\end{Resp}
\begin{Resp}\Rsign Quia te vidi justum coram me ex ómnibus géntibus, allelúja.\end{Resp}
\switchcolumn
\EN
\begin{Resp}\Rsign This is he which did according unto all that God commanded him; and God said unto him: Enter thou into My rest, \Star{} For thee have I seen righteous before Me among all people, alleluia.\end{Resp}
\begin{Resp}\Vsign This is he which loved not his life in this world, and is come unto an everlasting kingdom.\end{Resp}
\begin{Resp}\Rsign For thee have I seen righteous before Me among all people, alleluia.\end{Resp}
\begin{Resp}\Vsign Glory be to the Father, and to the Son, \Star{} and to the Holy Ghost.\end{Resp}
\begin{Resp}\Rsign For thee have I seen righteous before Me among all people, alleluia.\end{Resp}
\switchcolumn*

\LA
\LessonHead{Lectio VII}
\switchcolumn
\EN
\LessonHead{Lesson VII}
\switchcolumn*

\LA
\LessonTitle{Léctio sancti Evangélii secúndum Lucam}
\switchcolumn
\EN
\LessonTitle{From the Holy Gospel according to Luke}
\switchcolumn*

\LA
\Cite{Luc 10:1-9}
\switchcolumn
\EN
\Cite{Luke 10:1-9}
\switchcolumn*

\LA
\noindent In illo témpore: Designávit Dóminus et álios septuagínta duos: et misit illos binos ante fáciem suam, in omnem civitátem et locum, quo erat ipse ventúrus. Et réliqua.\par
\switchcolumn
\EN
\noindent At that time, the Lord appointed other seventy-two also, and sent them two and two before His face into every city and place, whither He Himself would come. And so on.\par
\switchcolumn*

\LA
\LessonTitle{Homilía sancti Gregórii Papæ}
\switchcolumn
\EN
\LessonTitle{Homily by Pope St. Gregory (the Great.)}
\switchcolumn*

\LA
\Source{Homilia 17 in Evangelia}
\switchcolumn
\EN
\Source{17th Homily on the Gospels}
\switchcolumn*

\LA
\noindent Dóminus et Salvátor noster, fratres caríssimi, aliquándo nos sermónibus, aliquándo vero opéribus ádmonet. Ipsa étenim facta ejus præcépta sunt: quia dum áliquid tácitus facit, quid ágere debeámus innotéscit. Ecce enim binos in prædicatiónem discípulos mittit: quia duo sunt præcépta caritátis, Dei vidélicet amor, et próximi: et minus quam inter duos cáritas habéri non potest. Nemo enim próprie ad semetípsum habére caritátem dícitur: sed diléctio in álterum tendit, ut cáritas esse possit.\par
\switchcolumn
\EN
\noindent Dearly beloved brethren, our Lord and Saviour doth sometimes admonish us by words, and sometimes by works. Yea, His very works do themselves teach us for that which He doth silently His example still moveth us to copy. Behold how He sendeth forth His disciples to preach by two and two since there are two commandments to love, that is, a commandment to love God, and a commandment to love our neighbour and where there are not two, the one, being alone, hath not whereon to do the Lord's commandment. And no man can properly be said to love himself: for love tendeth outward toward our neighbour, if it be the love whereto the Gospel doth oblige us.\par
\switchcolumn*

\LA
\begin{Resp}\Rsign Iste est qui ante Deum magnas virtútes operátus est, et de omni corde suo laudávit Dóminum: \Star{} Ipse intercédat pro peccátis ómnium populórum, allelúja.\end{Resp}
\begin{Resp}\Vsign Ecce homo sine queréla, verus Dei cultor, ábstinens se ab omni ópere malo, et pérmanens in innocéntia sua.\end{Resp}
\begin{Resp}\Rsign Ipse intercédat pro peccátis ómnium populórum, allelúja.\end{Resp}
\switchcolumn
\EN
\begin{Resp}\Rsign This is he which wrought great wonders before God, and praised the Lord with all his heart. \Star{} May he pray for all people, that their sins may be forgiven unto them, alleluia.\end{Resp}
\begin{Resp}\Vsign Behold a man without blame, a worshipper of God in truth, keeping himself clean from every evil work, and abiding still in his innocency.\end{Resp}
\begin{Resp}\Rsign May he pray for all people, that their sins may be forgiven unto them, alleluia.\end{Resp}
\switchcolumn*

\LA
\LessonHead{Lectio VIII}
\switchcolumn
\EN
\LessonHead{Lesson VIII}
\switchcolumn*

\LA
\noindent Ecce enim binos ad prædicándum discípulos Dóminus mittit: quátenus hoc nobis tácitus ínnuat, quia qui caritátem erga álterum non habet, prædicatiónis offícium suscípere nullátenus debet. Bene autem dícitur, quia misit eos ante fáciem suam in omnem civitátem et locum, quo erat ipse ventúrus. Prædicatóres enim suos Dóminus séquitur: quia prædicátio prǽvenit, et tunc ad mentis nostræ habitáculum Dóminus venit, quando verba exhortatiónis præcúrrunt: atque per hoc véritas in mente suscípitur.\par
\switchcolumn
\EN
\noindent Behold, the Lord sendeth forth His disciples to preach by two and two and thus doing, He doth silently teach us that whosoever loveth not his neighbour, such a one it behoveth not to take upon him the office of a preacher. Well also is it said that He sent them before His face into every city and place whither He Himself would come. The Lord followeth His preachers first cometh preaching, and then the Lord Himself cometh to the house of our mind, whither the word of exhortation hath come before and so cometh the truth into our mind.\par
\switchcolumn*

\LA
\begin{Resp}\Rsign Sint lumbi vestri præcíncti, et lucérnæ ardéntes in mánibus vestris: \Star{} Et vos símiles homínibus exspectántibus dóminum suum, quando revertátur a núptiis, allelúja.\end{Resp}
\begin{Resp}\Vsign Vigiláte ergo, quia nescítis qua hora Dóminus vester ventúrus sit.\end{Resp}
\begin{Resp}\Rsign Et vos símiles homínibus exspectántibus dóminum suum, quando revertátur a núptiis, allelúja.\end{Resp}
\begin{Resp}\Vsign Glória Patri, et Fílio, \Star{} et Spirítui Sancto.\end{Resp}
\begin{Resp}\Rsign Et vos símiles homínibus exspectántibus dóminum suum, quando revertátur a núptiis, allelúja.\end{Resp}
\switchcolumn
\EN
\begin{Resp}\Rsign Let your loins be girded about, and your lights burning; \Star{} And ye yourselves like unto men that wait for their lord, when he will return from the wedding, alleluia.\end{Resp}
\begin{Resp}\Vsign Watch therefore, for ye know not what hour your Lord doth come.\end{Resp}
\begin{Resp}\Rsign And ye yourselves like unto men that wait for their lord, when he will return from the wedding, alleluia.\end{Resp}
\begin{Resp}\Vsign Glory be to the Father, and to the Son, \Star{} and to the Holy Ghost.\end{Resp}
\begin{Resp}\Rsign And ye yourselves like unto men that wait for their lord, when he will return from the wedding, alleluia.\end{Resp}
\switchcolumn*

\LA
\LessonHead{Lectio IX (pro commemoratione)}
\switchcolumn
\EN
\LessonHead{Lesson IX (when commemorated)}
\switchcolumn*

\LA
\LessonTitle{Commemoratio: S. Pauli a Cruce Confessoris}
\switchcolumn
\EN
\LessonTitle{Commemoration of: St. Paul of the Cross, Confessor}
\switchcolumn*

\LA
\noindent Paulus a Cruce, Uvádæ in Ligúria natus, a primo ratiónis usu, Jesu Christi crucifíxi amóre flagrávit. Martýrii desidério incénsus, exercítui se adjúnxit, qui Venétiis ad bellum Turcis inferéndum comparabátur. Cógnita vero Dei voluntáte, honestíssimis núptiis sibíque deláta pátrui hereditáte recusátis, rudi túnica a suo epíscopo indútus, nondum cléricus, Domínicum agrum divíni verbi prædicatióne excóluit. Romæ a Summo Pontífice Benedícto décimo tértio ex obediéntia sacerdótio auctus, in solitúdinem recéssit Argentárii montis, quo eum beáta Virgo jamprídem invitáverat, veste illi simul osténsa atri colóris passiónis Fílii sui insígnibus decoráta, ibíque fundaménta jecit novæ congregatiónis, cujus sodáles voto adstringeréntur Domínicæ passiónis memóriam promovéndi. Sacras Vírgines quoque instítuit, quæ Dómini passiónem sédulo meditaréntur. Prædicatióne, virtútibus et divínis charismátibus clarus, Romæ obdormívit in Dómino, anno millésimo septingentésimo septuagésimo quinto. Eum Pius nonus Póntifex máximus in Beatórum, dein in Sanctórum númerum rétulit.\par
\switchcolumn
\EN
\noindent Paul of the Cross was born at Ovada in Liguria and, as soon as he came to the use of reason, burned with love for Jesus Christ crucified. Fired with the desire for martyrdom, he joined the army which was being assembled in Venice to fight against the Turks. But when the will of God was made known to him, and he had refused a most honourable marriage and an inheritance left to him by his uncle, he received a coarse tunic as a habit from his bishop and, although not yet a cleric, cultivated the field of the Lord by preaching the word of God. In Rome, out of obedience to Pope Benedict XIII, he was raised to the priesthood. Then he retired into the solitude of Monte Argentario, where the Blessed Virgin had already invited him to go, at the same time shewing unto him a black habit adorned with the insignia of her Son's Passion. There he laid the foundations of a new congregation, whose members bind themselves by vow to promote the memory of the Lord's Passion, and he also established one for nuns to meditate continually upon this mystery. Renowned for his preaching, virtues, and divine charisms, he fell asleep in the Lord at Rome, in the year 1775. Pope Pius IX enrolled him among the Blessed and then among the Saints.\par
\switchcolumn*

\LA
\TeDeum{Te Deum}
\switchcolumn
\EN
\TeDeum{Te Deum}
\switchcolumn*

\end{paracol}

\par\needspace{10\baselineskip}\addvspace{1.2em}{\centering\small\textsc{Die 28 Aprilis} · \textsc{28 April}\par}
\Office{S. Vitalis Martyris}{St. Vitalis, Martyr}{Simplex}
\addcontentsline{toc}{subsection}{Die 28 Aprilis · S. Vitalis Martyris\enspace\textit{St. Vitalis, Martyr}}
\begin{paracol}{2}
\LA
\LessonHead{Lectio IX (pro commemoratione)}
\switchcolumn
\EN
\LessonHead{Lesson IX (when commemorated)}
\switchcolumn*

\LA
\LessonTitle{Commemoratio: S. Vitalis Martyris}
\switchcolumn
\EN
\LessonTitle{Commemoration of: St. Vitalis, Martyr}
\switchcolumn*

\LA
\noindent Vitalis miles, sanctorum Gervasii et Protasii pater, una cum Paulino judice Ravennam ingressus, cum vidisset Ursicinum medicum ob christianæ fidei confessionem ductum ad supplicium, paululum in tormentis titubare, exclamavit: Ursicine medice, qui alios curare solitus es, cave ne te mortis æternæ jaculo conficias. Qua voce confirmatus Ursicinus, martyrium fortiter subivit. Quare Paulinus incensus Vitalem comprehendi jubet, et equuleo tortum, atque in profundam foveam demersum, lapidibus obrui. Quo facto quidam Apollinis sacerdos, qui Paulinum in Vitalem incitarat, oppressus a dæmone, clamare cœpit: Tu me nimium Vitalis Christi Martyr incendis: et illo æstu jactatus, se præcipitavit in flumen.\par
\switchcolumn
\EN
\noindent Vitalis was a soldier, and the father of the holy Martyrs Gervase and Protase. He went to Ravenna with Paulinus the judge, and there saw the physician Ursicinus led out to die, because he owned to being a believer in Christ. As the torments went on, Ursicinus seemed to waver a little, and Vitalis cried out to him, Ursicinus! as a physician thou hast been used to heal other men's bodies, take heed lest thou let thine own soul die eternally. These words encouraged Ursicinus, and he endured bravely in his testimony even unto the end but Paulinus was filled with fury, and caused Vitalis to be seized, tormented on the rack, and finally thrown into a pit and buried under an heap of stones. When it was over, a certain priest of Apollo, who had urged on Paulinus against Vitalis, was seized by the devil, and began to cry out, Vitalis Vitalis thou art Christ's Martyr, but thou makest me to burn thou makest me to burn until in that phrenzy he threw himself into the river.\par
\switchcolumn*

\LA
\TeDeum{Te Deum}
\switchcolumn
\EN
\TeDeum{Te Deum}
\switchcolumn*

\end{paracol}

\par\needspace{10\baselineskip}\addvspace{1.2em}{\centering\small\textsc{Die 29 Aprilis} · \textsc{29 April}\par}
\Office{S. Petri Martyris}{St. Peter the Martyr}{Duplex}
\addcontentsline{toc}{subsection}{Die 29 Aprilis · S. Petri Martyris\enspace\textit{St. Peter the Martyr}}
\begin{paracol}{2}
\LA
\RefLine{Lectiones I–III: de Scriptura occurrente.}
\switchcolumn
\EN
\RefLine{Lessons I–III: from the Scripture of the day.}
\switchcolumn*

\LA
\RefLine{Lectiones VII–IX: ex Commúni unius Martyris Tempore Paschali (C2ap).}
\switchcolumn
\EN
\RefLine{Lessons VII–IX: from the Common of a Single Martyr in Paschaltide (C2ap).}
\switchcolumn*

\LA
\LessonHead{Lectio IV}
\switchcolumn
\EN
\LessonHead{Lesson IV}
\switchcolumn*

\LA
\noindent Petrus, Veronæ paréntibus Manichæórum hæresi infectis natus, ab ipsa pæne infántia contra hæreses pugnávit. Puer annórum septem, cum scholas frequentaret, aliquándo a patruo hæretico interrogátus quid tandem in iis didicísset, christianæ fidei symbolum se didicisse respóndit; neque ullis umquam patris patruíve blanditiis aut minis a fidei constanti dimoveri pótuit. Adoléscens Bonóniam studiórum causa venit; ubi, a Spíritu Sancto ad sublimioris vitæ formam vocatus, ordinis Prædicatórum institutum suscépit.\par
\switchcolumn
\EN
\noindent Peter was born at Verona, (in the year of our Lord 1205,) of parents polluted with the Manichaean heresy, but he himself began his lifelong strife against error when he was but a little child. When he was seven years old he went to school, and was asked by his heretic uncle what he learnt there he answered that he had learnt the Christian Creed and neither his father nor his uncle were ever able to shake his constancy in the faith, either by cajolements or threats. When he was a young lad he went to Bologna to study, and there he was called by the Holy Ghost to an higher state of life, and entered the Order of Friars Preachers, at fifteen years of age.\par
\switchcolumn*

\LA
\begin{Resp}\Rsign Lux perpétua lucébit Sanctis tuis, Dómine, \Star{} Et ætérnitas témporum, allelúja, allelúja.\end{Resp}
\begin{Resp}\Vsign Lætítia sempitérna erit super cápita eórum: gáudium et exsultatiónem obtinébunt.\end{Resp}
\begin{Resp}\Rsign Et ætérnitas témporum, allelúja, allelúja.\end{Resp}
\switchcolumn
\EN
\begin{Resp}\Rsign Eternal light will shine over your saints, O Lord, \Star{} And the eternity of the times, alleluia, alleluia.\end{Resp}
\begin{Resp}\Vsign Everlasting joy shall be upon their heads, they shall obtain joy and gladness\end{Resp}
\begin{Resp}\Rsign And the eternity of the times, alleluia, alleluia.\end{Resp}
\switchcolumn*

\LA
\LessonHead{Lectio V}
\switchcolumn
\EN
\LessonHead{Lesson V}
\switchcolumn*

\LA
\noindent Magno virtútum splendore in religióne eluxit: corpus et ánimum ab omni impuritate ita custodívit, ut nullius mortíferi peccáti labe se inquinátum umquam sénserit. Carnem jejuniis et vigiliis macerábat; mentem divinis contemplatiónibus exercebat. In salúte animárum procuranda assidue versabátur; peculiáris gratiæ dono hæreticos ácriter confutábat. Tantam in concionando vim hábuit, ut innumerábilis hóminum multitúdo ad eum audiéndum conflúeret, multique ad pœniténtiam converteréntur.\par
\switchcolumn
\EN
\noindent He was marked by great perfection as a Friar so watchful was he over the purity of his body and soul, that he never felt himself defiled by a mortal sin. He chastened his body by fasting and watching, and ennobled his soul by the contemplation of the things of God. He was constantly busied in works for furthering the salvation of souls and had a peculiar gift of grace for clearly convincing heretics. Such was his power as a preacher, that countless crowds were drawn together to hear him, and many were moved to repentance.\par
\switchcolumn*

\LA
\begin{Resp}\Rsign In servis suis, allelúja, \Star{} Consolábitur Deus, allelúja.\end{Resp}
\begin{Resp}\Vsign Judicábit Dóminus pópulum suum, et in servis suis.\end{Resp}
\begin{Resp}\Rsign Consolábitur Deus, allelúja.\end{Resp}
\switchcolumn
\EN
\begin{Resp}\Rsign His servants, alleluia \Star{} God will comfort, alleluia, alleluia.\end{Resp}
\begin{Resp}\Vsign The Lord will judge His people and, His servants,\end{Resp}
\begin{Resp}\Rsign God will comfort, alleluia, alleluia.\end{Resp}
\switchcolumn*

\LA
\LessonHead{Lectio VI}
\switchcolumn
\EN
\LessonHead{Lesson VI}
\switchcolumn*

\LA
\noindent Tanto fidei ardore incénsus erat, ut pro ea mortem subire optaret, eamque a Deo grátiam enixe precarétur. Itaque hæretici necem, quam is paulo ante concionando prædixerat, illi intulérunt. Nam cum sanctæ Inquisitiónis munus géreret, illum Como Mediolanum redeuntem ímpius sicarius semel atque íterum in cápite gládio vulnerávit; jamque pene mórtuus, symbolórum fidei, quam infans virili fortitúdine conféssus fuerat, in ipso supremo spíritu pronuntiávit; iterúmque látera mucrone transverberatus, ad martyrii palmam migrávit in cælum, anno salútis millesimo ducentésimo quinquagesimo secundo. Quem multis, illustrem miraculis Innocentius quartus anno sequénti sanctórum Mártyrum número adscripsit.\par
\switchcolumn
\EN
\noindent The faith which was in him burnt so hotly, that he longed to seal his confession with his blood, and oftentimes he earnestly besought from God the grace to do so. It was but a little while before the heretics murdered him, that he foretold, in preaching, his own approaching death. While he was intrusted with the duties of the Holy Inquisition, he was returning from Como to Milan, when an ungodly ruffian assailed him and wounded him once and again in the head with a sword. Peter, to whom these blows were nearly fatal, began with his last breath to recite that Profession of the Faith, to which as a little child he had clung with such manly courage, but the murderer thrust the weapon into his side, and he passed away to receive a Martyr's palm in heaven. It was (the 6th day of April, in) the year of salvation 1252. In the following year, Innocent IV., seeing by how many miracles God had been pleased to glorify him, added his name to the sacred roll of Martyrs.\par
\switchcolumn*

\LA
\begin{Resp}\Rsign Fíliæ Jerúsalem, veníte et vidéte Mártyres cum corónis, quibus coronávit eos Dóminus \Star{} In die solemnitátis et lætítiæ, allelúja.\end{Resp}
\begin{Resp}\Vsign Quóniam confortávit seras portárum tuárum, benedíxit fílios tuos in te.\end{Resp}
\begin{Resp}\Rsign In die solemnitátis et lætítiæ, allelúja.\end{Resp}
\begin{Resp}\Vsign Glória Patri, et Fílio, \Star{} et Spirítui Sancto.\end{Resp}
\begin{Resp}\Rsign In die solemnitátis et lætítiæ, allelúja.\end{Resp}
\switchcolumn
\EN
\begin{Resp}\Rsign Come forth, O ye daughters of Jerusalem, and behold the Martyrs with the crowns wherewith the Lord crowned them; \Star{} In the day of His feasting, and of His gladness. Alleluia.\end{Resp}
\begin{Resp}\Vsign For He hath strengthened the bars of thy gates He hath blessed thy children within thee.\end{Resp}
\begin{Resp}\Rsign In the day of His feasting, and of His gladness. Alleluia.\end{Resp}
\begin{Resp}\Vsign Glory be to the Father, and to the Son, \Star{} and to the Holy Ghost.\end{Resp}
\begin{Resp}\Rsign In the day of His feasting, and of His gladness. Alleluia.\end{Resp}
\switchcolumn*

\LA
\LessonHead{Lectio IX (pro commemoratione)}
\switchcolumn
\EN
\LessonHead{Lesson IX (when commemorated)}
\switchcolumn*

\LA
\LessonTitle{Commemoratio: S. Petri Martyris}
\switchcolumn
\EN
\LessonTitle{Commemoration of: St. Peter the Martyr}
\switchcolumn*

\LA
\noindent Petrus, Verónæ paréntibus Manichæórum hǽresi inféctis natus, ab ipsa pene infántia contra hǽreses pugnávit, nullis umquam patris patruíve blandítiis a fídei constánti dimótus. Adoléscens Bonóniam studiórum causa venit, ibíque órdinis Prædicatórum institútum suscépit; in quo et magno virtútum splendóre, præcípue córporis animǽque puritáte nullo umquam letháli peccáto fœdáta, et miro pœniténtiæ et contemplatiónis stúdio excélluit. In salútem animárum procurándam máximo cum fructu incúbuit, tanto fídei ardóre incénsus, ut pro ea mortem subeúndi grátiam a Deo precarétur; quam et obtínuit. Nam cum sanctæ Inquisitiónis munus gerens Como Mediolánum redíret, ab ímpio sicário in cápite gládio est vulnerátus; jamque pene mórtuus, sýmbolum fídei, quam ab infántia viríli fortitúdine defénderat, in ipso suprémo spíritu pronúntians, ad palmam martýrii evolávit, anno salútis millésimo ducentésimo quinquagésimo secúndo. Quem Innocéntius quartus anno sequénti, sanctórum Mártyrum número adscrípsit.\par
\switchcolumn
\EN
\noindent Peter was born in Verona of parents infected with Manichaeism, but almost from his infancy fought against heresies, and no persuasions of his father or his uncle could move him from holding to the faith. As a young man he went to Bologna to study and there joined the Order of Preachers. As a religious he gave a shining example of virtue, especially by his purity of body and soul, never sullied by mortal sin, and excelled in his wonderful zeal for penance and contemplation. He devoted himself most fruitfully to gaining the salvation of souls, and was kindled with such zeal for the faith that he asked God for the grace of dying for it, a favour which was granted him. For, while he was carrying out the work of the Holy Inquisition and on his way back from Como to Milan, he was wounded in the head by the sword of a wicked assassin. Nearly dead, he repeated with his last breath the symbol of the faith which he had defended with manly fortitude from his childhood, and gained the palm of martyrdom in the year of salvation 1252. A year later, Innocent IV enrolled him among the holy Martyrs.\par
\switchcolumn*

\LA
\TeDeum{Te Deum}
\switchcolumn
\EN
\TeDeum{Te Deum}
\switchcolumn*

\end{paracol}

\par\needspace{10\baselineskip}\addvspace{1.2em}{\centering\small\textsc{Die 30 Aprilis} · \textsc{30 April}\par}
\Office{S. Catharinæ Senensis Virginis}{St. Catherine of Siena, Virgin}{Duplex}
\addcontentsline{toc}{subsection}{Die 30 Aprilis · S. Catharinæ Senensis Virginis\enspace\textit{St. Catherine of Siena, Virgin}}
\begin{paracol}{2}
\LA
\RefLine{Lectiones I–III: de Scriptura occurrente.}
\switchcolumn
\EN
\RefLine{Lessons I–III: from the Scripture of the day.}
\switchcolumn*

\LA
\RefLine{Lectiones VII–IX: ex Commúni Virginum tantum (C6a).}
\switchcolumn
\EN
\RefLine{Lessons VII–IX: from the Common of a Virgin not a Martyr (C6a).}
\switchcolumn*

\LA
\LessonHead{Lectio IV}
\switchcolumn
\EN
\LessonHead{Lesson IV}
\switchcolumn*

\LA
\noindent Catharina, virgo Senensis piis orta parentibus, beati Dominici habitum, quem sorores de Pœnitentia gestant, impetravit. Summa ejus fuit abstinentia, et admirabilis vitæ austeritas. Inventa est aliquando a die Cinerum usque ad Ascensionem Domini jejunium perduxisse, sola Eucharistiæ communione contenta. Luctabatur quam frequentissime cum dæmonibus, multisque illorum molestiis vexabatur, æstuabat febribus, nec aliorum morborum cruciatu carebat. Magnum et sanctum erat Catharinæ nomen, et undique ad eam ægroti et malignis vexati spiritibus deducebantur. Languoribus et febribus in Christi nomine imperabat, et dæmones cogebat ab obsessis abire corporibus.\par
\switchcolumn
\EN
\noindent This Katharine was a maiden of Siena, and was born of godly parents, (in the year 1347.) She took the habit of the Third Order of St. Dominic. Her fasts were most severe, and the austerity of her life wonderful. It was discovered that on some occasions she took no food at all from Ash Wednesday till Ascension Day, receiving all needful strength by taking the Holy Communion. She was engaged oftentimes in a wrestling with devils, and was sorely tried by them with diverse assaults: she was consumed by fevers, and suffered likewise from other diseases. Great and holy was the name of Katharine, and sick folk, and such as were vexed with evil spirits, were brought to her from all quarters. Through the Name of Christ, she had command over sickness and fever, and forced the foul spirits to leave the bodies of the tormented.\par
\switchcolumn*

\LA
\begin{Resp}\Rsign Propter veritátem, et mansuetúdinem, et justítiam: \Star{} Et dedúcet te mirabíliter déxtera tua, allelúja.\end{Resp}
\begin{Resp}\Vsign Spécie tua et pulchritúdine tua inténde, próspere procéde, et regna.\end{Resp}
\begin{Resp}\Rsign Et dedúcet te mirabíliter déxtera tua, allelúja.\end{Resp}
\switchcolumn
\EN
\begin{Resp}\Rsign Because of truth, and meekness, and righteousness; \Star{} And thy right hand shall lead thee wonderfully, alleluia.\end{Resp}
\begin{Resp}\Vsign In thy comeliness, and thy beauty, go forward, fare prosperously, and reign.\end{Resp}
\begin{Resp}\Rsign And thy right hand shall lead thee wonderfully, alleluia.\end{Resp}
\switchcolumn*

\LA
\LessonHead{Lectio V}
\switchcolumn
\EN
\LessonHead{Lesson V}
\switchcolumn*

\LA
\noindent Cum Pisis immoraretur die Dominico, refecta cibo cælesti, et in extasim rapta, vidit Dominum crucifixum magno cum lumine advenientem, et ex ejus vulnerum cicatricibus quinque radios ad quinque loca sui corporis descendentes; ideoque mysterium advertens, Dominum precata, ne cicatrices apparerent; continuo radii colorem sanguineum mutaverunt in splendidum, et in formam puræ lucis pervenerunt ad manus, pedes, et cor ejus: ac tantus erat dolor, quem sensibiliter patiebatur, ut nisi Deus minuisset, brevi se crederet morituram. Hanc ítaque gratiam amantissimus Dominus nova gratia cumulavit, ut sentiret dolorem illapsa vi vulnerum et cruenta signa non apparerent. Quod ita contigisse cum Dei famula confessario suo Raymundo retulisset, ut oculis étiam repræsentaretur, radios in imaginibus beatæ Catharinæ ad dicta quinque loca pertingentes, pia fidelium cura pictis coloribus expressit.\par
\switchcolumn
\EN
\noindent While she dwelt at Pisa, on a certain Lord's Day, after she had received the Living Bread Which came down from heaven, she was in the spirit; and saw the Lord nailed to the Cross advancing towards her. There was a great light round about Him, and five rays of light streaming from the five marks of the Wounds in His Feet, and Hands, and Side, which smote her upon the five corresponding places in her body. When Katharine perceived this vision, she besought the Lord that no marks might become manifest upon her flesh, and straightway the five beams of light changed from the colour of blood into that of gold, and touched in the form of pure light her feet, and hands, and side. At this moment the agony which she felt was so piercing, that she believed that if God had not lessened it, she would have died. Thus the Lord in His great love for her, gave her this great grace, in a new and twofold manner, namely, that she felt all the pain of the wounds, but without there being any bloody marks to meet the gaze of men. This was the account given by the handmaiden of God to her Confessor, Raymund, and it is for this reason that when the godly wishes of the faithful lead them to make pictures of the blessed Katharine, they paint her with golden rays of light proceeding from those five places in her body which correspond to the five places wherein our Lord was wounded by the nails and spear.\par
\switchcolumn*

\LA
\begin{Resp}\Rsign Dilexísti justítiam, et odísti iniquitátem: \Star{} Proptérea unxit te Deus, Deus tuus, óleo lætítiæ, allelúja.\end{Resp}
\begin{Resp}\Vsign Propter veritátem, et mansuetúdinem, et justítiam.\end{Resp}
\begin{Resp}\Rsign Proptérea unxit te Deus, Deus tuus, óleo lætítiæ, allelúja.\end{Resp}
\switchcolumn
\EN
\begin{Resp}\Rsign Thou hast loved righteousness, and hated iniquity; \Star{} Therefore God, thy God, hath anointed thee with the oil of gladness, alleluia.\end{Resp}
\begin{Resp}\Vsign Because of truth, and meekness, and righteousness.\end{Resp}
\begin{Resp}\Rsign Therefore God, thy God, hath anointed thee with the oil of gladness, alleluia.\end{Resp}
\switchcolumn*

\LA
\LessonHead{Lectio VI}
\switchcolumn
\EN
\LessonHead{Lesson VI}
\switchcolumn*

\LA
\noindent Doctrina ejus infusa, non acquisita fuit; sacrarum litterarum professoribus difficillimas de divinitate quæstiones proponentibus respondit Nemo ad eam accessit, qui non melior abierit: multa exstinxit odia, et mortales sedavit inimicitias. Pro pace Florentinorum, qui cum Ecclesia dissidebant, et interdicto ecclesiastico suppositi erant, Avenionem ad Gregorium undecimum Pontificem maximum profecta est. Cui étiam votum ejus de petenda Urbe soli Deo notum, sese divinitus cognovisse monstravit: deliberavitque Pontifex, ea étiam suadente, ad sedem suam Romanam personaliter accedere: quod et fecit. Eidem Gregorio et Urbano sexto ejus successori acceptissima fuit, adeo ut legationibus eorum fungeretur. Denique post innumera virtutum insignia, dono prophetiæ, et pluribus clara miraculis, anno ætatis suæ tertio circiter et trigesimo, migravit ad Sponsum. Quam Pius secundus Pontifex maximus sanctarum Virginum numero adscripsit.\par
\switchcolumn
\EN
\noindent The learning which Katharine had was not acquired but inspired. She answered Professors of Divinity upon the very hardest questions concerning God. No one was ever in her company without going away better. She healed many hatreds, and quieted the most deadly feuds. To make peace for the Florentines, who had quarrelled with the Church, and were under an Ecclesiastical Interdict, she travelled to Avignon, (in 1376) to see the Chief Pontiff Gregory XI. To him she showed that she had had revealed to her from heaven his secret purpose of going back to Rome, which had been known only to God and himself. It was at her persuasion, as well as by his own judgment, that the Pope did in the end return to his own See. She was much respected by this Gregory, as well as by his successor Urban VI, who even employed her in their embassies. The Bridegroom took her home, (upon the 29th day of April, in the year of salvation 1380,) when she was about thirty-three years old, after she had given almost countless proofs of extraordinary Christian graces, and manifestly displayed the gifts of Prophecy and miracles. Pope Pius II enrolled her among the Virgin Saints.\par
\switchcolumn*

\LA
\begin{Resp}\Rsign Afferéntur Regi vírgines post eam, próximæ ejus; \Star{} Afferéntur tibi in lætítia et exsultatióne, allelúja.\end{Resp}
\begin{Resp}\Vsign Spécie tua et pulchritúdine tua; inténde, próspere procéde, et regna.\end{Resp}
\begin{Resp}\Rsign Afferéntur tibi in lætítia et exsultatióne, allelúja.\end{Resp}
\begin{Resp}\Vsign Glória Patri, et Fílio, \Star{} et Spirítui Sancto.\end{Resp}
\begin{Resp}\Rsign Afferéntur tibi in lætítia et exsultatióne, allelúja.\end{Resp}
\switchcolumn
\EN
\begin{Resp}\Rsign After her shall virgins be brought unto the King, her fellows \Star{} They shall be brought unto thee with gladness and rejoicing, alleluia.\end{Resp}
\begin{Resp}\Vsign In thy comeliness and thy beauty, go forward, fare prosperously, and reign.\end{Resp}
\begin{Resp}\Rsign They shall be brought unto thee with gladness and rejoicing, alleluia.\end{Resp}
\begin{Resp}\Vsign Glory be to the Father, and to the Son, \Star{} and to the Holy Ghost.\end{Resp}
\begin{Resp}\Rsign They shall be brought unto thee with gladness and rejoicing, alleluia.\end{Resp}
\switchcolumn*

\LA
\LessonHead{Lectio IX (pro commemoratione)}
\switchcolumn
\EN
\LessonHead{Lesson IX (when commemorated)}
\switchcolumn*

\LA
\LessonTitle{Commemoratio: S. Catharinæ Senensis Virginis}
\switchcolumn
\EN
\LessonTitle{Commemoration of: St. Catherine of Siena, Virgin}
\switchcolumn*

\LA
\noindent Catharína, virgo Senénsis, piis orta paréntibus, beáti Domínici hábitum, quem Soróres de Pœniténtia gestant, impetrávit. Summa ejus fuit abstinéntia et admirábilis vitæ austéritas. Cum Pisis morarétur, die Domínico, refécta cibo cælésti et in éxtasim rapta, vidit Dóminum crucifíxum magno cum lúmine adveniéntem, et ex ejus vúlnerum cicatrícibus quinque rádios ad quinque loca sui córporis descendéntes. Mystérium advértens, Dóminum precáta, ne cicatríces apparérent, contínuo rádii, colóre sanguíneo mutáto in spléndidum, in formam puræ lucis pervenérunt ad manus, pedes et cor ejus. Tantus vero erat dolor, quem sensibíliter patiebátur, licet vúlnerum cruénta signa non apparérent, ut, nisi Deus minuísset, brevi se créderet moritúram. Doctrína ejus infúsa, non acquisíta fuit. Aveniónem ad Gregórium Papam undécimum profécta, illi votum ejus de peténda Urbe, soli Deo notum, sese divínitus cognovísse monstrávit, et auctor fuit, ut Póntifex ad Sedem Románam personáliter accéderet. Anno ætátis suæ tértio círciter et trigésimo migrávit ad Sponsum. Quam Pius secúndus sanctárum Vírginum número adscrípsit.\par
\switchcolumn
\EN
\noindent Catharine, a virgin of Siena, born of devout parents, was granted the habit of St. Dominic worn by the Sisters of Penance. Her abstinence was most strict, and her whole life was one of marvellous austerity. When she was staying at Pisa on the Lord's Day, refreshed by the Bread of heaven and rapt in ecstasy, she saw the crucified Lord coming with a great light and, from the marks of his wounds, five rays coming down to the same places in her body. Aware of the mystery, she implored the Lord that the wounds would not be visible, and the colour of the rays immediately changed from that of blood to brightness, in the form of pure light touching her hands and feet and heart. But such was the pain she suffered, even though the signs of the bleeding wounds could not be seen, that she believed she would soon have died if God had not lessened it. Her learning was infused, not acquired. She went to Pope Gregory XI at Avignon and shewed him that she knew by divine means of the vow he had made to return to the City, a vow known to God alone, and she was the cause of the Pope's going to occupy in person his See at Rome. In about the thirty-third year of her age she went to her Bridegroom, and Pius II enrolled her among the holy Virgins.\par
\switchcolumn*

\LA
\TeDeum{Te Deum}
\switchcolumn
\EN
\TeDeum{Te Deum}
\switchcolumn*

\end{paracol}

\SeasonTitle{Mensis Majus}{May}

\addcontentsline{toc}{section}{Mensis Majus\enspace\textit{May}}

\par\needspace{10\baselineskip}\addvspace{1.2em}{\centering\small\textsc{Die 1 Maii} · \textsc{1 May}\par}
\Office{Ss. Philippi et Jacobi Apostolorum}{Sts. Philip and James, Apostles}{Duplex II classis}
\addcontentsline{toc}{subsection}{Die 1 Maii · Ss. Philippi et Jacobi Apostolorum\enspace\textit{Sts. Philip and James, Apostles}}
\begin{paracol}{2}
\LA
\LessonHead{Lectio I}
\switchcolumn
\EN
\LessonHead{Lesson I}
\switchcolumn*

\LA
\LessonTitle{Incipit Epístola cathólica beáti Jacóbi Apóstoli}
\switchcolumn
\EN
\LessonTitle{Lesson from the letter of St. James the Apostle}
\switchcolumn*

\LA
\Cite{Jas 1:1-6}
\switchcolumn
\EN
\Cite{Jas 1:1-6}
\switchcolumn*

\LA
\noindent Jacóbus, Dei et Dómini nostri Jesu Christi servus, duódecim tríbubus, quæ sunt in dispersióne, salútem. Omne gáudium existimáte fratres mei, cum in tentatiónes várias incidéritis: sciéntes quod probátio fídei vestræ patiéntiam operátur. Patiéntia autem opus perféctum habet: ut sitis perfécti et íntegri in nullo deficiéntes. Si quis autem vestrum índiget sapiéntia, póstulet a Deo, qui dat ómnibus affluénter, et non impróperat: et dábitur ei. Póstulet autem in fide nihil hǽsitans:\par
\switchcolumn
\EN
\noindent James the servant of God, and of our Lord Jesus Christ, to the twelve tribes which are scattered abroad, greeting. My brethren, count it all joy, when you shall fall into diverse temptations; Knowing that the trying of your faith worketh patience. And patience hath a perfect work; that you may be perfect and entire, failing in nothing. But if any of you want wisdom, let him ask of God, who giveth to all men abundantly, and upbraideth not; and it shall be given him. But let him ask in faith, nothing wavering.\par
\switchcolumn*

\LA
\begin{Resp}\Rsign Beátus vir qui métuit Dóminum, allelúja: \Star{} In mandátis ejus cupit nimis, allelúja, allelúja, allelúja.\end{Resp}
\begin{Resp}\Vsign Glória et divítiæ in domo ejus, et justítia ejus manet in sǽculum sǽculi.\end{Resp}
\begin{Resp}\Rsign In mandátis ejus cupit nimis, allelúja, allelúja, allelúja.\end{Resp}
\switchcolumn
\EN
\begin{Resp}\Rsign Blessed is the man that feareth the Lord, alleluia. \Star{} He shall delight exceedingly in his commandments, alleluia, alleluia, alleluia.\end{Resp}
\begin{Resp}\Vsign Glory and wealth shall be in his house: and his justice remaineth for ever and ever.\end{Resp}
\begin{Resp}\Rsign He shall delight exceedingly in his commandments, alleluia, alleluia, alleluia.\end{Resp}
\switchcolumn*

\LA
\LessonHead{Lectio II}
\switchcolumn
\EN
\LessonHead{Lesson II}
\switchcolumn*

\LA
\Cite{Jas 1:6-11}
\switchcolumn
\EN
\Cite{Jas 1:6-11}
\switchcolumn*

\LA
\noindent Qui enim hǽsitat, símilis est flúctui maris, qui a vento movétur et circumfértur: non ergo ǽstimet homo ille quod accípiat áliquid a Dómino. Vir duplex ánimo incónstans est in ómnibus viis suis. Gloriétur autem frater húmilis in exaltatióne sua: dives autem in humilitáte sua, quóniam sicut flos fœni transíbit; exórtus est enim sol cum ardóre, et arefécit fœnum, et flos ejus décidit, et decor vultus ejus depériit: ita et dives in itinéribus suis marcéscet.\par
\switchcolumn
\EN
\noindent For he that wavereth is like a wave of the sea, which is moved and carried about by the wind. Therefore let not that man think that he shall receive any thing of the Lord. A double minded man is inconstant in all his ways. But let the brother of low condition glory in his exaltation: And the rich, in his being low; because as the flower of the grass shall he pass away. For the sun rose with a burning heat, and parched the grass, and the flower thereof fell off, and the beauty of the shape thereof perished: so also shall the rich man fade away in his ways.\par
\switchcolumn*

\LA
\begin{Resp}\Rsign Tristítia vestra, allelúja, \Star{} Convertétur in gáudium, allelúja, allelúja.\end{Resp}
\begin{Resp}\Vsign Mundus autem gaudébit, vos vero contristabímini, sed tristítia vestra.\end{Resp}
\begin{Resp}\Rsign Convertétur in gáudium, allelúja, allelúja.\end{Resp}
\switchcolumn
\EN
\begin{Resp}\Rsign Your sorrow, alleluia. \Star{} Shall be turned into joy, alleluia.\end{Resp}
\begin{Resp}\Vsign The world shall rejoice: and you shall be made sorrowful, but your sorrow.\end{Resp}
\begin{Resp}\Rsign Shall be turned into joy, alleluia.\end{Resp}
\switchcolumn*

\LA
\LessonHead{Lectio III}
\switchcolumn
\EN
\LessonHead{Lesson III}
\switchcolumn*

\LA
\Cite{Jas 1:12-16}
\switchcolumn
\EN
\Cite{Jas 1:12-16}
\switchcolumn*

\LA
\noindent Beátus vir qui suffert tentatiónem: quóniam cum probátus fúerit, accípiet corónam vitæ, quam repromísit Deus diligéntibus se. Nemo cum tentátur, dicat quóniam a Deo tentátur: Deus enim intentátor malórum est: ipse autem néminem tentat. Unusquísque vero tentátur a concupiscéntia sua abstráctus, et illéctus. Deínde concupiscéntia cum concéperit, parit peccátum: peccátum vero cum consummátum fúerit, génerat mortem. Nolíte ítaque erráre, fratres mei dilectíssimi.\par
\switchcolumn
\EN
\noindent Blessed is the man that endureth temptation; for when he hath been proved, he shall receive a crown of life, which God hath promised to them that love him. Let no man, when he is tempted, say that he is tempted by God. For God is not a tempter of evils, and he tempteth no man. But every man is tempted by his own concupiscence, being drawn away and allured. Then when concupiscence hath conceived, it bringeth forth sin. But sin, when it is completed, begetteth death. Do not err, therefore, my dearest brethren.\par
\switchcolumn*

\LA
\begin{Resp}\Rsign Pretiósa in conspéctu Dómini, allelúja, \Star{} Mors Sanctórum ejus, allelúja.\end{Resp}
\begin{Resp}\Vsign Custódit Dóminus ómnia ossa eórum, unum ex his non conterétur.\end{Resp}
\begin{Resp}\Rsign Mors Sanctórum ejus, allelúja.\end{Resp}
\begin{Resp}\Vsign Glória Patri, et Fílio, \Star{} et Spirítui Sancto.\end{Resp}
\begin{Resp}\Rsign Mors Sanctórum ejus, allelúja.\end{Resp}
\switchcolumn
\EN
\begin{Resp}\Rsign Precious in the sight of the Lord, alleluia. \Star{} Is the death of his saints, alleluia.\end{Resp}
\begin{Resp}\Vsign The Lord is nigh unto them that are of a contrite heart: and he will save the humble of spirit.\end{Resp}
\begin{Resp}\Rsign Is the death of his saints, alleluia.\end{Resp}
\begin{Resp}\Vsign Glory be to the Father, and to the Son, \Star{} and to the Holy Ghost.\end{Resp}
\begin{Resp}\Rsign Is the death of his saints, alleluia.\end{Resp}
\switchcolumn*

\LA
\LessonHead{Lectio IV}
\switchcolumn
\EN
\LessonHead{Lesson IV}
\switchcolumn*

\LA
\noindent Philíppus Bethsáidæ natus, unus ex duódecim Apóstolis, qui primum a Christo Dómino vocáti sunt: a quo cum accepísset Nathánaël, venísse Messíam in lege promíssum, ad Dóminum dedúctus est. Quam vero Christus eum familiáriter adhibéret, illud fácile declárat, quod gentíles Salvatórem vidére cupiéntes, ad Philíppum accessérunt; et Dóminus, cum in solitúdine hóminum multitúdinem páscere vellet, sic Philíppum affátus est: Unde emémus panes, ut mandúcent hi? Is, accépto Spíritu Sancto, cum ei Scýthia ad prædicándum Evangélium obtigísset, omnem fere illam gentem ad christiánam fidem convértit. Postrémo, cum Hierápolim Phrýgiæ venísset, pro Christi nómine cruci affíxus lapidibúsque óbrutus est, Kaléndis Maji. Ejus corpus ibídem a Christiánis sepúltum, póstea Romam delátum, in basílica duódecim Apostolórum una cum córpore beáti Jacóbi Apóstoli cónditum est.\par
\switchcolumn
\EN
\noindent Philip was born in the town of Bethsaida, and was one of the first of the twelve Apostles who were called by the Lord Christ. Then Philip findeth Nathanael, and saith unto him: “We have found Him of Whom Moses in the Law, and the Prophets, did write.” And so he brought him to the Lord. How familiarly he was in the company of Christ, is manifest from that which is written: “There were certain Greeks among them that came up to worship at the Feast the same came therefore to Philip, and desired him, saying: Sir, we would see Jesus.” When the Lord was in the wilderness, and was about to feed a great multitude, He said unto Philip: “Whence shall we buy bread, that these may eat?” Philip, after that he had received the Holy Ghost, took Scythia, by lot, as the land wherein he was to preach the Gospel, and brought nearly all that people to believe in Christ. At the last he came to Hierapolis in Phrygia, and there, for Christ’s Name’s sake, he was fastened to a cross and stoned to death. The day was the first of May. The Christians of Hierapolis buried his body at that place, but it was afterwards brought to Rome and laid in the Basilica of the Twelve Apostles, beside that of the blessed Apostle James.\par
\switchcolumn*

\LA
\begin{Resp}\Rsign Lux perpétua lucébit Sanctis tuis, Dómine, \Star{} Et ætérnitas témporum, allelúja, allelúja.\end{Resp}
\begin{Resp}\Vsign Lætítia sempitérna erit super cápita eórum: gáudium et exsultatiónem obtinébunt.\end{Resp}
\begin{Resp}\Rsign Et ætérnitas témporum, allelúja, allelúja.\end{Resp}
\switchcolumn
\EN
\begin{Resp}\Rsign Eternal light will shine over your saints, O Lord, \Star{} And the eternity of the times, alleluia, alleluia.\end{Resp}
\begin{Resp}\Vsign Everlasting joy shall be upon their heads, they shall obtain joy and gladness\end{Resp}
\begin{Resp}\Rsign And the eternity of the times, alleluia, alleluia.\end{Resp}
\switchcolumn*

\LA
\LessonHead{Lectio V}
\switchcolumn
\EN
\LessonHead{Lesson V}
\switchcolumn*

\LA
\noindent Jacóbus frater Dómini, cognoménto Justus, ab ineúnte ætáte vinum et síceram non bibit, carne abstínuit, nunquam tonsus est, nec unguénto nec bálneo usus. Huic uni licébat íngredi in Sancta sanctórum. Idem líneis véstibus utebátur: cui étiam assidúitas orándi ita callum génibus obdúxerat, ut durítie caméli pellem imitarétur. Eum post Christi ascensiónem Apóstoli Jerosolymórum epíscopum creavérunt; ad quem étiam Princeps Apostolórum misit qui nuntiáret se e cárcere ab Angelo edúctum fuísse. Cum autem in concílio Jerosólymis controvérsia esset orta de lege et circumcisióne; Jacóbus, Petri senténtiam secútus, ad fratres hábuit conciónem, in qua vocatiónem géntium probávit, fratribúsque abséntibus scribéndum esse dixit, ne géntibus jugum Mosáicæ legis impónerent. De quo et lóquitur Apóstolus ad Gálatas: Alium autem Apostolórum vidi néminem, nisi Jacóbum fratrem Dómini.\par
\switchcolumn
\EN
\noindent James, surnamed the Just, the brother of our Lord Jesus Christ, was a Nazarite from the womb. During his whole life he never drank wine or strong drink, never ate meat, never shaved, and never took a bath. He was the only man who was allowed to go into the Holy of Holies. His raiment was always linen. So continually did he kneel in prayer, that the skin of his knees became horny, like a camel’s knees. After Christ was ascended, the Apostles made James Bishop of Jerusalem and even the Prince of the Apostles gave special intelligence to him after that he was delivered from prison by an angel. When in the Council of Jerusalem certain questions were mooted touching the law and circumcision, James, following the opinion of Peter, addressed a discourse to the brethren, wherein he proved the call of the Gentiles, and commanded letters to be sent to such brethren as were absent, that they might take heed not to lay upon the Gentiles the yoke of the Law of Moses. It is of him that the Apostle Paul saith, writing to the Galatians: “Other of the Apostles saw I none, save James the Lord’s brother.”\par
\switchcolumn*

\LA
\begin{Resp}\Rsign Virtúte magna reddébant Apóstoli \Star{} Testimónium resurrectiónis Jesu Christi Dómini nostri, allelúja, allelúja.\end{Resp}
\begin{Resp}\Vsign Repléti quidem Spíritu Sancto, loquebántur cum fidúcia verbum Dei.\end{Resp}
\begin{Resp}\Rsign Testimónium resurrectiónis Jesu Christi Dómini nostri, allelúja, allelúja.\end{Resp}
\switchcolumn
\EN
\begin{Resp}\Rsign With great power did the Apostles \Star{} Give testimony of the resurrection of Jesus Christ our Lord, alleluia, alleluia.\end{Resp}
\begin{Resp}\Vsign And they were all filled with the Holy Ghost: and they spoke the word of God with confidence.\end{Resp}
\begin{Resp}\Rsign Give testimony of the resurrection of Jesus Christ our Lord, alleluia, alleluia.\end{Resp}
\switchcolumn*

\LA
\LessonHead{Lectio VI}
\switchcolumn
\EN
\LessonHead{Lesson VI}
\switchcolumn*

\LA
\noindent Tanta autem erat Jacóbi vitæ sánctitas, ut fímbriam vestiménti ejus certátim hómines cúperent attíngere. Nam is nonagínta sex annos natus, cum trigínta annis illi Ecclésiæ sanctíssime præfuísset, Christum Dei Fílium constantíssime prǽdicans, lapídibus primum appétitur; mox in altíssimum Templi locum addúctus, inde præcipitátus est. Qui, confráctis crúribus, jacens semivívus, manus tendébat ad cælum, Deúmque pro illórum salúte deprecabátur his verbis: Ignósce eis, Dómine, quia nésciunt quid fáciunt. Qua in oratióne, gráviter ejus cápite fullónis fuste percússo, ánimam Deo réddidit, séptimo Nerónis anno, et juxta Templum, ubi præcipitátus fúerat, sepúltus est. Unam scripsit Epístolam, quæ de septem cathólicis est.\par
\switchcolumn
\EN
\noindent So great was James’ holiness of life that men strove one with another to touch the hem of his garment. When he was ninety-six years old, and had most holily governed the Church of Jerusalem for thirty years, ever most constantly preaching Christ the Son of God, he laid down his life for the faith. He was first stoned, and afterward taken up on to a pinnacle of the Temple and cast down from thence. His legs were broken by the fall, and he was well-nigh dead, but he lifted up his hands towards heaven, and prayed to God for the salvation of his murderers, saying: “Lord, forgive them, for they know not what they do.” As he said this, one that stood by smote him grievously upon the head with a fuller’s club, and he resigned his spirit to God. He testified in the seventh year of Nero, and was buried hard by the Temple, in the place where he had fallen. He wrote one of the Seven Epistles which are called Catholic.\par
\switchcolumn*

\LA
\begin{Resp}\Rsign Isti sunt agni novélli, qui annuntiavérunt, allelúja, modo venérunt ad fontes, \Star{} Repléti sunt claritáte, allelúja, allelúja.\end{Resp}
\begin{Resp}\Vsign In conspéctu Agni amícti sunt stolis albis, et palmæ in mánibus eórum.\end{Resp}
\begin{Resp}\Rsign Repléti sunt claritáte, allelúja, allelúja.\end{Resp}
\begin{Resp}\Vsign Glória Patri, et Fílio, \Star{} et Spirítui Sancto.\end{Resp}
\begin{Resp}\Rsign Repléti sunt claritáte, allelúja, allelúja.\end{Resp}
\switchcolumn
\EN
\begin{Resp}\Rsign These are the young lambs, who were promised, they are just coming to the fountains. \Star{} They are filled with glory, alleluia, alleluia.\end{Resp}
\begin{Resp}\Vsign In sight of the Lamb, clothed with white robes, and palms in their hands.\end{Resp}
\begin{Resp}\Rsign They are filled with glory, alleluia, alleluia.\end{Resp}
\begin{Resp}\Vsign Glory be to the Father, and to the Son, \Star{} and to the Holy Ghost.\end{Resp}
\begin{Resp}\Rsign They are filled with glory, alleluia, alleluia.\end{Resp}
\switchcolumn*

\LA
\LessonHead{Lectio VII}
\switchcolumn
\EN
\LessonHead{Lesson VII}
\switchcolumn*

\LA
\LessonTitle{Léctio sancti Evangélii secúndum Joánnem}
\switchcolumn
\EN
\LessonTitle{From the Holy Gospel according to John}
\switchcolumn*

\LA
\Cite{Joannes 14:1-13}
\switchcolumn
\EN
\Cite{John 14:1-13}
\switchcolumn*

\LA
\noindent In illo témpore: Dixit Jesus discípulis suis: Non turbétur cor vestrum. Créditis in Deum, et in me crédite. In domo Patris mei mansiónes multæ sunt. Et réliqua.\par
\switchcolumn
\EN
\noindent At that time, Jesus said unto His disciples: “Let not your heart be troubled. Ye believe in God, believe also in Me. In My Father’s house there are many mansions.” And so on.\par
\switchcolumn*

\LA
\LessonTitle{Homilía sancti Augustíni Epíscopi}
\switchcolumn
\EN
\LessonTitle{Homily by St. Augustine, Bishop of Hippo}
\switchcolumn*

\LA
\Source{Tract. 67 in Joannem}
\switchcolumn
\EN
\Source{67th Tract on John}
\switchcolumn*

\LA
\noindent Erigénda est nobis, fratres, ad Deum major inténtio, ut verba sancti Evangélii, quæ modo in nostris áuribus sonuérunt, étiam mente cápere utcúmque possímus. Ait enim Dóminus Jesus: Non turbétur cor vestrum. Créditis in Deum, et in me crédite. Ne mortem tamquam hómines timérent, et ídeo turbaréntur, consolátur eos, étiam se Deum esse contéstans. Créditis, inquit, in Deum, et in me crédite. Cónsequens est enim, ut, si in Deum créditis, et in me crédere debeátis: quod non esset cónsequens, si Christus non esset Deus.\par
\switchcolumn
\EN
\noindent It behoveth us, my brethren, to have our minds more given toward God, if we would that those words of the Holy Gospel which have just sounded in our ears, should become a living reality for our understandings. The Lord Jesus saith: “Let not your heart be troubled. Ye believe in God, believe also in Me.” Lest, being but men, their heart should be troubled by the fear of death, He strengtheneth them, even by the reminder that He is God. He saith: “Ye believe in God, believe also in Me”, for if ye believe in God, ye must needs believe in Me. And this were not so, if Christ were not God.\par
\switchcolumn*

\LA
\begin{Resp}\Rsign Ego sum vitis vera, et vos pálmites: \Star{} Qui manet in me, et ego in eo, hic fert fructum multum, allelúja, allelúja.\end{Resp}
\begin{Resp}\Vsign Sicut diléxit me Pater, et ego diléxi vos.\end{Resp}
\begin{Resp}\Rsign Qui manet in me, et ego in eo, hic fert fructum multum, allelúja, allelúja.\end{Resp}
\switchcolumn
\EN
\begin{Resp}\Rsign I am the vine: you the branches. \Star{} He that abideth in me, and I in him, the same beareth much fruit, alleluia, alleluia.\end{Resp}
\begin{Resp}\Vsign As the Father hath loved me, I also have loved you.\end{Resp}
\begin{Resp}\Rsign He that abideth in me, and I in him, the same beareth much fruit, alleluia, alleluia.\end{Resp}
\switchcolumn*

\LA
\LessonHead{Lectio VIII}
\switchcolumn
\EN
\LessonHead{Lesson VIII}
\switchcolumn*

\LA
\noindent Créditis in Deum, et in eum crédite, cui natúra est, non rapína, esse æquálem Deo; semetípsum enim exinanívit, non tamen formam Dei amíttens, sed formam servi accípiens. Mortem metúitis huic formæ servi: non turbétur cor vestrum; suscitábit illam forma Dei. Sed quid est, quod séquitur: In domo Patris mei mansiónes multæ sunt; nisi quia et sibi metuébant? Unde audíre debuérunt: Non turbétur cor vestrum. Quis enim eórum non metúeret, cum Petro dictum esset, fidentióri atque promptióri: Non cantábit gallus, donec ter me neges?\par
\switchcolumn
\EN
\noindent We believe in God, believe also in Him Who is by nature and not by robbery equal with God, for in that He emptied Himself, He did it not by laying aside the form of God, but by taking upon Him the form of a servant. Ye fear death for this form of a servant, but let not your heart be troubled, the form of God will raise it up again. But what signifieth that which followeth: “In My Father’s house there are many mansions.” Was it not that they had fear on their own account, and needed for themselves to hear Him say, “Let not your heart be troubled” Which of them trembled not when they had heard Him say to Peter, the lealest and boldest of them all: “The cock shall not crow this day, before that thou shalt thrice deny that thou knowest Me.”\par
\switchcolumn*

\LA
\begin{Resp}\Rsign Cándidi facti sunt Nazarǽi ejus, allelúja: splendórem Deo dedérunt, allelúja: \Star{} Et sicut lac coaguláti sunt, allelúja, allelúja.\end{Resp}
\begin{Resp}\Vsign Candidióres nive, nitidióres lacte, rubicundióres ébore antíquo, sapphíro pulchrióres.\end{Resp}
\begin{Resp}\Rsign Et sicut lac coaguláti sunt, allelúja, allelúja.\end{Resp}
\begin{Resp}\Vsign Glória Patri, et Fílio, \Star{} et Spirítui Sancto.\end{Resp}
\begin{Resp}\Rsign Et sicut lac coaguláti sunt, allelúja, allelúja.\end{Resp}
\switchcolumn
\EN
\begin{Resp}\Rsign Her Nazarites are become pure, Alleluia: they reflect the glory of God, Alleluia. \Star{} They are whiter than milk. Alleluia, Alleluia.\end{Resp}
\begin{Resp}\Vsign They are purer than snow, they are whiter than milk, they are more ruddy in body than coral, their polishing is of sapphire.\end{Resp}
\begin{Resp}\Rsign They are whiter than milk. Alleluia, Alleluia.\end{Resp}
\begin{Resp}\Vsign Glory be to the Father, and to the Son, \Star{} and to the Holy Ghost.\end{Resp}
\begin{Resp}\Rsign They are whiter than milk. Alleluia, Alleluia.\end{Resp}
\switchcolumn*

\LA
\LessonHead{Lectio IX}
\switchcolumn
\EN
\LessonHead{Lesson IX}
\switchcolumn*

\LA
\noindent Tamquam ergo essent ab illo peritúri, mérito turbabántur; sed cum áudiunt, In domo Patris mei mansiónes multæ sunt; si quo minus, dixíssem vobis: Quia vado paráre vobis locum; a perturbatióne recreántur, certi ac fidéntes, étiam post perícula tentatiónum, se apud Deum cum Christo esse mansúros. Quia etsi álius est álio fórtior, álius álio sapiéntior, álius álio jústior, álius álio sánctior, in domo Patris mei mansiónes multæ sunt. Nullus eórum alienábitur ab illa domo, ubi mansiónem pro suo quisque acceptúrus est mérito.\par
\switchcolumn
\EN
\noindent Meetly were they troubled, for that they were about to be scattered from Him, but when they heard Him say, “In My Father’s house are many mansions,” they had been comforted even if He had not also said, “I go to prepare a place for you,” for then they believed and knew, that, when all dangers and all trials were for ever over, they should be for ever with the Lord, with Christ and with God. Yea, though one man be stronger than another, though one be wiser than another, though one be holier than another, yet “in My Father’s house are many mansions.” That house is an house wherein none are strangers, but every man shall receive a mansion therein according as his work shall be.\par
\switchcolumn*

\LA
\TeDeum{Te Deum}
\switchcolumn
\EN
\TeDeum{Te Deum}
\switchcolumn*

\end{paracol}

\par\needspace{10\baselineskip}\addvspace{1.2em}{\centering\small\textsc{Die 2 Maii} · \textsc{2 May}\par}
\Office{S. Athanasii Episcopi Confessoris et Ecclesiæ Doctoris}{S. Athanasius, Confessor and Doctor of the Church}{Duplex}
\addcontentsline{toc}{subsection}{Die 2 Maii · S. Athanasii Episcopi Confessoris et Ecclesiæ Doctoris\enspace\textit{S. Athanasius, Confessor and Doctor of the Church}}
\begin{paracol}{2}
\LA
\RefLine{Lectiones I–III: de Scriptura occurrente.}
\switchcolumn
\EN
\RefLine{Lessons I–III: from the Scripture of the day.}
\switchcolumn*

\LA
\LessonHead{Lectio IV}
\switchcolumn
\EN
\LessonHead{Lesson IV}
\switchcolumn*

\LA
\noindent Athanasius Alexandrinus, catholicæ religiónis propugnator acerrimus, ab Alexandro episcopo Alexandrino diaconus factus est, in cujus locum successit. Quem étiam antea secutus fuerat ad Nicænum Concílium; ubi cum Aríi impietátem repressísset, tantum ódium Arianórum suscépit, ut ex eo témpore ei insidias moliri numquam destiterint. Nam coacto ad Tyrum concílio magna ex parte Arianórum episcopórum, subornarunt muliérculam, quæ accusaret Athansium, quod hospítio acceptus sibi stuprum per vim intulísset. Introductus ígitur est Athanasius, et una cum eo Timótheus presbyter; qui simulans se esse Athanasium, Ego ne, inquit, mulier, apud te sum diversatus? ego te violávi? Cui illa petulanter: Tu mihi vim attulísti; idque jurejurando affirmans, judicum fidem obtestabátur, ut tantum flagitium vindicárent. Qua cógnita fraude, rejecta est mulíeris impudéntia.\par
\switchcolumn
\EN
\noindent The great Athanasius, the greatest soldier that the Catholic Religion hath perhaps ever had, was an Alexandrian. He was ordained Deacon by Alexander, (in the year 326,) Bishop of that city, whom he afterwards succeeded. (In 325) he had followed Alexander to the Council of Nicea, where he wrestled triumphantly against the blasphemy of Arius. For this reason he was honoured with so much of their hatred by the Arians, that their vindictiveness never forsook him from that time forward. (In the year 335,) they called together a Council at Tyre, composed for the most part of Arian Bishops, where they suborned a wretched woman to charge Athanasius with having raped her when she had received him as a guest into her house. Athanasius therefore came into the assembly, and with him a certain priest whose name was Timothy. This Timothy arose as though he were Athanasius, and asked her, saying Woman, was it I that was thy guest was it I that raped thee She cried out indignantly Yea, thou it was that didst rape me, the which she attested with an oath, and called on the honour of the judges to punish such iniquity. Upon this discovery of her perjury, they drave the shameless woman from their presence.\par
\switchcolumn*

\LA
\begin{Resp}\Rsign Invéni David servum meum, óleo sancto meo unxi eum: \Star{} Manus enim mea auxiliábitur ei, allelúja.\end{Resp}
\begin{Resp}\Vsign Nihil profíciet inimícus in eo, et fílius iniquitátis non nocébit ei.\end{Resp}
\begin{Resp}\Rsign Manus enim mea auxiliábitur ei, allelúja.\end{Resp}
\switchcolumn
\EN
\begin{Resp}\Rsign I have found David My servant, with My holy oil have I anointed him \Star{} For My hand shall help him, alleluia.\end{Resp}
\begin{Resp}\Vsign The enemy shall prevail nothing against him, nor the son of wickedness afflict him.\end{Resp}
\begin{Resp}\Rsign For My hand shall help him, alleluia.\end{Resp}
\switchcolumn*

\LA
\LessonHead{Lectio V}
\switchcolumn
\EN
\LessonHead{Lesson V}
\switchcolumn*

\LA
\noindent Arsénium quoque episcopum ab Athanasio interfectum Ariáni pervulgarunt; quem dum occulte détinent, manum mórtui déferunt in judícium, ab Athanasio ad usum mágicæ artis Arsenio amputatam criminantes. At Arsenius noctu aufúgiens, cum se in conspectu totius concilii statuísset, Athanasii inimicórum impudentíssimum scelus apéruit. Quod illi nihilóminus mágicis artibus Athanasii tribuéntes, vitæ ejus insidiari non desistebant. Quam ob rem in exsílium actus, in Gállia apud Tréviros exsulávit. Gravibus deinceps ac diuturnis sub Constantio imperatóre, Arianórum fautore, tempestátibus jactátus et incredibiles calamitátes perpéssus, magnam orbis terræ partem peragrávit; ac sæpe e sua ecclésia ejéctus, sæpe étiam in eamdem et Julii, Romani Pontificis, auctoritate, et Constántis imperatóris, Constantii fratris, patrocinio, decretis quoque concilii Sardicénsis ac Jerosolymitani, restitútus est, Ariánis interea illi semper infestis; quorum pertinacem iram, et summum vitæ discrimen fúgiens, in sicca cisterna quinque annis se ábdidit, ejus rei tantum conscio quodam Athanasii amico, qui eum clam sustentabat.\par
\switchcolumn
\EN
\noindent The Arians also accused Athanasius of having murdered the (schismatic) Bishop Arsenius. This Arsenius they kept shut up, and brought into the court a dead man's hand, which they declared had been his, and had been cut off by Athanasius to use in sorcery. But Arsenius escaped in the night, and when he appeared before all the Council whole and sound, the brazen-faced crime of the enemies of Athanasius was exposed. This appearance nevertheless they attributed to Athanasius being a warlock, and persisted still in their attack on him. He was driven into exile, and banished to Treves in Gaul. Thenceforth, under authority of the Emperor Constantius, that abettor of Arians, he was hunted to and fro with unceasing persecutions. He suffered hardships which it is difficult to believe. He was sent wandering all about the Roman world. He was twice more thrust out of his See, and again restored through the authority of Pope Julius of Rome, and with the protection of the Emperor Constans, the brother of Constantius, by decrees of the Councils of Sardica and of Jerusalem. The vindictiveness of the Arians never let him alone. In his third exile so great was the danger of his life from the pursuit of their undying hatred, that he had to lie hid for five years in a dry cistern, unknown to all men, save one of his friends who brought him food.\par
\switchcolumn*

\LA
\begin{Resp}\Rsign Pósui adjutórium super poténtem, et exaltávi eléctum de plebe mea: \Star{} Manus enim mea auxiliábitur ei, allelúja.\end{Resp}
\begin{Resp}\Vsign Invéni David servum meum, óleo sancto meo unxi eum.\end{Resp}
\begin{Resp}\Rsign Manus enim mea auxiliábitur ei, allelúja.\end{Resp}
\switchcolumn
\EN
\begin{Resp}\Rsign I have laid help upon one that is mighty, and have exalted one chosen out of My people \Star{} For My hand shall help him, alleluia.\end{Resp}
\begin{Resp}\Vsign I have found David My servant, with My holy oil have I anointed him.\end{Resp}
\begin{Resp}\Rsign For My hand shall help him, alleluia.\end{Resp}
\switchcolumn*

\LA
\LessonHead{Lectio VI}
\switchcolumn
\EN
\LessonHead{Lesson VI}
\switchcolumn*

\LA
\noindent Constantio mortuo, cum Julianus Apóstata, qui ei in imperio successit, éxsules episcopos ad suas ecclésias redire permisísset, Athanasius Alexandríam reversus, summo honóre exceptus est. Sed non multo post, iísdem Ariánis impelléntibus, a Juliáno exagitatus, rursus discedere cogitur. Cumque ab ejus satellítibus ad necem conquirerétur, qua fugiebat navícula conversa in contrariam flúminis partem, iis qui se insequebántur, ex industria occurrit; et quæréntibus quantum inde abesset Athanasius, respóndit eum non longe abesse: ítaque illos contrárium tenéntes cursum effugit, atque Alexandríam rédiens, ibidem usque ad Juliáni óbitum occultus permansit. Qui paulo post, Alexandríæ alia exorta tempestáte, quatuor menses in paterno sepúlcro delituit. Ac denique ex tot tantisque periculis divinitus ereptus, Alexandríæ mórtuus est in suo léctulo, sub Valénte: cujus vita et mors magnis nobilitáta est miraculis. Multa pie et ad illustrándum catholicam fidem præcláre scripsit, sexque et quadragínta annos in summa témporum varietáte Alexandrinam ecclésiam sanctíssime gubernávit.\par
\switchcolumn
\EN
\noindent After the death of Constantius, Julian the Apostate, who succeeded him, allowed every sort of Bishop who had been banished to return to their own Churches. Athanasius therefore returned to Alexandria, and was received with profound reverence. But it was not long before the same Arians got Julian to hunt him down again, and again it behoved him to fly. A band of soldiers were sent in pursuit of him to kill him, and as he fled up the Nile, their boat pressed hard on his. Athanasius, before they were yet in sight, had his own boat turned round, and went down the stream to meet them. As the vessels passed one another the murderers called out to ask if they knew where Athanasius was, and the servant of God himself cried to them in answer, Ye are close to him! whereupon they redoubled their exertions to ascend the stream, and Athanasius went peacefully down to Alexandria, and found means of concealment till the death of Julian. Yet once again he had to fly from another persecution at Alexandria, and in this his fifth and last exile he hid himself for four months in his own father's sepulchre. From all these so many and so great dangers did God deliver him, and at last he died in his own bed at Alexandria, (upon the 2nd day of May, in the year of salvation 373,) in the reign of Valens. He wrote much that is both godly and luminous in explaining the Catholic Faith, and governed the Church of Alexandria in great holiness, amid all changes of weather, for six and forty years.\par
\switchcolumn*

\LA
\begin{Resp}\Rsign Iste est, qui ante Deum magnas virtútes operátus est, et omnis terra doctrína ejus repléta est: \Star{} Ipse intercédat pro peccátis ómnium populórum, allelúja.\end{Resp}
\begin{Resp}\Vsign Iste est, qui contémpsit vitam mundi, et pervénit ad cæléstia regna.\end{Resp}
\begin{Resp}\Rsign Ipse intercédat pro peccátis ómnium populórum, allelúja.\end{Resp}
\begin{Resp}\Vsign Glória Patri, et Fílio, \Star{} et Spirítui Sancto.\end{Resp}
\begin{Resp}\Rsign Ipse intercédat pro peccátis ómnium populórum, allelúja.\end{Resp}
\switchcolumn
\EN
\begin{Resp}\Rsign This is he which wrought great wonders before God, and the whole earth is full of his teaching \Star{} May he pray for all people, that their sins may be forgiven unto them, alleluia.\end{Resp}
\begin{Resp}\Vsign This is he which loved not his life in this world, and hath attained unto the kingdom of heaven.\end{Resp}
\begin{Resp}\Rsign May he pray for all people, that their sins may be forgiven unto them, alleluia.\end{Resp}
\begin{Resp}\Vsign Glory be to the Father, and to the Son, \Star{} and to the Holy Ghost.\end{Resp}
\begin{Resp}\Rsign May he pray for all people, that their sins may be forgiven unto them, alleluia.\end{Resp}
\switchcolumn*

\LA
\LessonHead{Lectio VII}
\switchcolumn
\EN
\LessonHead{Lesson VII}
\switchcolumn*

\LA
\LessonTitle{Léctio sancti Evangélii secúndum Matthǽum.}
\switchcolumn
\EN
\LessonTitle{From the Holy Gospel according to Matthew}
\switchcolumn*

\LA
\Cite{Matt 10:23-28}
\switchcolumn
\EN
\Cite{Matt 10:23-28}
\switchcolumn*

\LA
\noindent In illo témpore: Dixit Jesus discípulis suis: Cum persequétur vos in civitáte ista, fugite in aliam. Et réliqua.\par
\switchcolumn
\EN
\noindent At that time Jesus said to his disciples: And when they shall persecute you in this city, flee into another. And so on.\par
\switchcolumn*

\LA
\LessonTitle{Homilía sancti Athanasii Epíscopi.}
\switchcolumn
\EN
\LessonTitle{Homily by St. Athanasius, Pope (of Alexandria.)}
\switchcolumn*

\LA
\Source{In Apologia de fuga sua, ante medium}
\switchcolumn
\EN
\Source{Defence of his own flight.}
\switchcolumn*

\LA
\noindent In lege præcéptum erat ut constitueréntur civitátes refugiórum, ut, qui quomodocúmque ad necem quæreréntur, servari possent. In consummatióne porro sæculórum cum advenísset illud ipsum Verbum Patris, quod Móysi antea locútum fuerat, rursus hoc præcéptum dedit, Cum vos, inquiens, persecúti fúerint in civitáte ista, fugite in aliam. Pauloque post súbjicit: Cum vidéritis illam abominatiónem desolatiónis, quæ dicta est per Danielem prophétam, consistentem in loco sancto (qui legit, intélligat), tunc qui in Judǽa sunt, fúgiant ad montes; et qui in tecto est, ne descendat tollere aliquid de domo sua; et qui in agro est, non revertátur tollere túnicam suam.\par
\switchcolumn
\EN
\noindent It is written in the Law, (Num. xxxv. 11,) Ye shall appoint you cities to be cities of refuge for you, that in these cities they which were pursued to put them to death might enter and be safe. And in the latter days when He was come, even that very Word of the Father, Which had spoken aforetime unto Moses, He gave again the same commandment When they persecute you in this city, flee ye into another. And, a while afterward, He said When ye shall see the abomination of desolation, spoken of by Daniel the Prophet, stand in the Holy Place, whoso readeth, let him understand, then let them which be in Judaea flee unto the mountains; let him which is on the house-top not come down to take anything out of his house; neither let him which is in the field return back to take his clothes. (Matth. xxiv. 15-18.)\par
\switchcolumn*

\LA
\begin{Resp}\Rsign Amávit eum Dóminus, et ornávit eum: stolam glóriæ índuit eum, \Star{} Et ad portas paradísi coronávit eum, allelúja.\end{Resp}
\begin{Resp}\Vsign Induit eum Dóminus lorícam fídei, et ornávit eum.\end{Resp}
\begin{Resp}\Rsign Et ad portas paradísi coronávit eum, allelúja.\end{Resp}
\switchcolumn
\EN
\begin{Resp}\Rsign The Lord loved him and beautified him He clothed him with a robe of glory \Star{} And crowned him at the gates of Paradise, alleluia.\end{Resp}
\begin{Resp}\Vsign The Lord hath put on him the breast-plate of faith, and hath adorned him.\end{Resp}
\begin{Resp}\Rsign And crowned him at the gates of Paradise, alleluia.\end{Resp}
\switchcolumn*

\LA
\LessonHead{Lectio VIII}
\switchcolumn
\EN
\LessonHead{Lesson VIII}
\switchcolumn*

\LA
\noindent Hæc cum scirent Sancti, ejusmodi tenuérunt suæ conversatiónis institutum. Quæ enim nunc præcepit Dóminus, éadem quoque ante suum in carne advéntum locútus est in Sanctis; et hoc institutum hómines ad perfectiónem ducit. Nam quod Deus jusserit, id omnino faciéndum est. Ideoque et ipsum Verbum propter nos homo factum, non indignum putávit, cum quærerétur, quemádmodum et nos, abscóndere se; et cum persecutiónem paterétur, fugere, et insidias declináre: cum autem a se definitum tempus ipse adduxísset, in quo corporáliter pro ómnibus pati volebat, ultro seípsum trádidit insidiántibus.\par
\switchcolumn
\EN
\noindent The Saints, therefore, knowing these words of the Lord, have obeyed them in their lives. What the Lord hath now commanded by His Own Mouth He commanded through His Saints before that He Himself was come in the flesh, and to obey this commandment worketh in a man perfection, since whatever God commandeth is a thing which it behoveth man to do. For this cause, that very Word of God Which was made flesh for our sake thought it meet when they sought Him, even as at this present time they are seeking us, to hide Himself, (John viii. 59,) and, when they persecuted Him, to fly and escape from their laying in wait for Him although when that time came which He had Himself decreed, and wherein He willed, as touching the Body, to suffer for us all, He willingly gave Himself up to His enemies.\par
\switchcolumn*

\LA
\begin{Resp}\Rsign In médio Ecclésiæ apéruit os ejus, \Star{} Et implévit eum Dóminus spíritu sapiéntiæ et intelléctus, allelúja.\end{Resp}
\begin{Resp}\Vsign Jucunditátem et exsultatiónem thesaurizávit super eum.\end{Resp}
\begin{Resp}\Rsign Et implévit eum Dóminus spíritu sapiéntiæ et intelléctus, allelúja.\end{Resp}
\begin{Resp}\Vsign Glória Patri, et Fílio, \Star{} et Spirítui Sancto.\end{Resp}
\begin{Resp}\Rsign Et implévit eum Dóminus spíritu sapiéntiæ et intelléctus, allelúja.\end{Resp}
\switchcolumn
\EN
\begin{Resp}\Rsign In the midst of the congregation did the Lord open his mouth. \Star{} And filled him with the spirit of wisdom and understanding, alleluia.\end{Resp}
\begin{Resp}\Vsign He made him rich with joy and gladness.\end{Resp}
\begin{Resp}\Rsign And filled him with the spirit of wisdom and understanding, alleluia.\end{Resp}
\begin{Resp}\Vsign Glory be to the Father, and to the Son, \Star{} and to the Holy Ghost.\end{Resp}
\begin{Resp}\Rsign And filled him with the spirit of wisdom and understanding, alleluia.\end{Resp}
\switchcolumn*

\LA
\LessonHead{Lectio IX}
\switchcolumn
\EN
\LessonHead{Lesson IX}
\switchcolumn*

\LA
\noindent At vero sancti hómines cum hanc quoque formam a Salvatore didicissent, (ab ipso enim et antea et semper omnes docebántur) advérsus persecutores ut legitime certarent, fugiébant, et ab illis quæsíti se abscóndebant. Cum enim præstitúti sibi a divina providéntia témporis finem ignorarent, nolébant insidiántibus se témere trádere: sed contra, cum scirent quod scriptum est, in mánibus Dei esse hóminum sortes, et Dóminum mortificáre et vivificáre; potius in finem usque perseverábant, circumeuntes, ut ait Apóstolus, in melótis et pellibus caprínis, egéntes, angustiati, in solitudínibus errántes, et in speluncis et cavernis terræ laténtes, quoad vel definitum mortis tempus veníret, vel qui tempus ipsum definíerat, Deus cum eis loquerétur, et insidiántes cohibéret, aut certe persecutóribus eos traderet, utcúmque illi placuísset.\par
\switchcolumn
\EN
\noindent Holy men of God, therefore, have learnt to take example from their Saviour, (and the Same is and hath been the Teacher of all such, whether of old time, or in these latter days,) and know how that it is lawful to baffle their persecutors by flying from them, and by lying hid when they seek them. For since they know not the day nor the hour wherein an all-seeing God hath ordained their end, they do not daringly give themselves into the power of such as hate them, but rather, knowing it to be written, “My times are in Thy hand,” (Ps. xxx. 16) and that “the Lord killeth and maketh alive,” (1 Kings ii. 6,) they “endure unto the end,” (Matth. xxiv. 13.) “they wander about,” as saith the Apostle, “in sheepskins and goatskins, being destitute, afflicted, [tormented (of whom the world is not worthy,)] they wander in deserts, [and in mountains,] and” hide “in dens and caves of the earth,” (Heb. xi. 37,) until either their appointed time come, or until more plainly God, the real Appointer of times, speaketh unto them, and chaineth up the persecutors, or manifestly giveth them over into the hands of the same, as may be His Own good pleasure.\par
\switchcolumn*

\LA
\TeDeum{Te Deum}
\switchcolumn
\EN
\TeDeum{Te Deum}
\switchcolumn*

\LA
\LessonHead{Lectio IX (pro commemoratione)}
\switchcolumn
\EN
\LessonHead{Lesson IX (when commemorated)}
\switchcolumn*

\LA
\LessonTitle{Commemoratio: S. Athanasii Episcopi Confessoris et Ecclesiæ Doctoris}
\switchcolumn
\EN
\LessonTitle{Commemoration of: S. Athanasius, Confessor and Doctor of the Church}
\switchcolumn*

\LA
\noindent Athanásius, epíscopus Alexandrínus, cathólicæ religiónis propugnátor acérrimus, cum, adhuc diáconus, in Concílio Nicǽno Aríi impietátem repressísset, tantum ódium Arianórum suscépit, ut ex eo témpore ei insídias molíri nunquam destíterint. In exsílium actus, in Gállia apud Tréviros exsulávit. Incredíbiles dein calamitátes perpéssus, magnam orbis partem peragrávit; ac sæpe e sua ecclésia ejéctus, sæpe étiam in eándem, Júlii, Románi Pontificis, auctoritáte atque decrétis concílii Sardicénsis ac Jerosolymitáni, restitútus est, Ariánis intérea illi semper inféstis. Dénique ex tot tantísque perículis divínitus eréptus, Alexandríæ mórtuus est sub Valénte. Ejus vita et mors magnis nobilitáta est miráculis. Multa pie et ad illustrándam cathólicam fidem præcláre scripsit, sexque et quadragínta annos in summa témporum varietáte Alexandrínam ecclésiam sanctíssime gubernávit.\par
\switchcolumn
\EN
\noindent Athanasius was Bishop of Alexandria, and a most vigorous defender of the Catholic religion. When he was still a deacon, he refuted the impiety of Arius at the Council of Nicaea, and earned such hatred from the Arians that, from that time on, they never ceased to lay snares for him. Driven into exile, he went to Treves in Gaul. He endured unbelievable hardships and wandered over a great part of the world, being often driven from his Church, and often restored by the authority of Pope Julius and the decrees of the Councils of Sardica and Jerusalem. All this while, he was persecuted by the Arians. Finally, rescued, by the help of God, from so many great dangers, he died at Alexandria during the reign of Emperor Valens. His life and death are marked by great miracles. He wrote many works, both of devotion and of catechetics, and, with great holiness, he ruled the Church of Alexandria, in those most troubled times, for forty-six years.\par
\switchcolumn*

\LA
\TeDeum{Te Deum}
\switchcolumn
\EN
\TeDeum{Te Deum}
\switchcolumn*

\end{paracol}

\par\needspace{10\baselineskip}\addvspace{1.2em}{\centering\small\textsc{Die 3 Maii} · \textsc{3 May}\par}
\Office{Inventione Sanctæ Crucis}{Inventione Sanctæ Crucis}{Duplex II classis}
\addcontentsline{toc}{subsection}{Die 3 Maii · Inventione Sanctæ Crucis\enspace\textit{Inventione Sanctæ Crucis}}
\begin{paracol}{2}
\LA
\RefLine{Lectio IX: de Ss. Alexandri I Papæ, Eventii et Theoduli Mm. ac Juvenalis Ep. et Conf. (05-03r).}
\switchcolumn
\EN
\RefLine{Lesson IX: of Sts. Alexander and Companions, Martyrs (05-03r).}
\switchcolumn*

\LA
\LessonHead{Lectio I}
\switchcolumn
\EN
\LessonHead{Lesson I}
\switchcolumn*

\LA
\LessonTitle{De Epístola beáti Pauli Apóstoli ad Gálatas}
\switchcolumn
\EN
\LessonTitle{Lesson from the letter of St. Paul the Apostle to the Galatians}
\switchcolumn*

\LA
\Cite{Gal 3:10-14}
\switchcolumn
\EN
\Cite{Gal 3:10-14}
\switchcolumn*

\LA
\noindent Quicúmque ex opéribus legis sunt, sub maledícto sunt. Scriptum est enim: Maledíctus omnis qui non permánserit in ómnibus quæ scripta sunt in libro legis ut fáciat ea. Quóniam autem in lege nemo justificátur apud Deum, maniféstum est: quia justus ex fide vivit. Lex autem non est ex fide, sed, Qui fécerit ea, vivet in illis. Christus nos redémit de maledícto legis, factus pro nobis maledíctum: quia scriptum est: Maledíctus omnis qui pendet in ligno: ut in géntibus benedíctio Ábrahæ fíeret in Christo Jesu, ut pollicitatiónem Spíritus accipiámus per fidem.\par
\switchcolumn
\EN
\noindent For as many as are of the works of the law, are under a curse. For it is written: Cursed is every one, that abideth not in all things, which are written in the book of the law to do them. But that in the law no man is justified with God, it is manifest: because the just man liveth by faith. But the law is not of faith: but, He that doth those things, shall live in them. Christ hath redeemed us from the curse of the law, being made a curse for us: for it is written: Cursed is every one that hangeth on a tree: That the blessing of Abraham might come on the Gentiles through Christ Jesus: that we may receive the promise of the Spirit by faith.\par
\switchcolumn*

\LA
\begin{Resp}\Rsign Gloriósum diem sacra venerátur Ecclésia, dum triumphále reserátur lignum: \Star{} In quo Redémptor noster, mortis víncula rumpens, cállidum áspidem superávit. (Allelúja, allelúja, allelúja.)\end{Resp}
\begin{Resp}\Vsign In ligno pendens nostræ salútis sémitam Verbum Patris invénit.\end{Resp}
\begin{Resp}\Rsign In quo Redémptor noster, mortis víncula rumpens, cállidum áspidem superávit. (Allelúja, allelúja, allelúja.)\end{Resp}
\switchcolumn
\EN
\begin{Resp}\Rsign Lo the Church, with solemn gladness, hails the day for ever glorious, when the opening earth revealeth that dread tree of mystic triumph. \Star{} On whose boughs her dying Saviour shattered death and crushed the serpent. Alleluia, Alleluia, Alleluia.\end{Resp}
\begin{Resp}\Vsign He the Word of God eternal, on those stately branches hanging, hath for us a new way opened.\end{Resp}
\begin{Resp}\Rsign On whose boughs her dying Saviour shattered death and crushed the serpent. Alleluia, Alleluia, Alleluia.\end{Resp}
\switchcolumn*

\LA
\LessonHead{Lectio II}
\switchcolumn
\EN
\LessonHead{Lesson II}
\switchcolumn*

\LA
\LessonTitle{De Epístola ad Philippénses}
\switchcolumn
\EN
\LessonTitle{From the letter to the Philippines}
\switchcolumn*

\LA
\Cite{Phil 2:5-11}
\switchcolumn
\EN
\Cite{Phil 2:5-11}
\switchcolumn*

\LA
\noindent Hoc enim sentíte in vobis, quod et in Christo Jesu: qui, cum in forma Dei esset, non rapínam arbitrátus est esse se æquálem Deo: sed semetípsum exinanívit, formam servi accípiens, in similitúdinem hóminum factus, et hábitu invéntus ut homo. Humiliávit semetípsum factus obédiens usque ad mortem, mortem autem crucis. Propter quod et Deus exaltávit illum, et donávit illi nomen, quod est super omne nomen: ut in nómine Jesu omne genu flectátur cæléstium, terréstrium et infernórum, et omnis lingua confiteátur, quia Dóminus Jesus Christus in glória est Dei Patris.\par
\switchcolumn
\EN
\noindent For let this mind be in you, which was also in Christ Jesus: Who being in the form of God, thought it not robbery to be equal with God: But emptied himself, taking the form of a servant, being made in the likeness of men, and in habit found as a man. He humbled himself, becoming obedient unto death, even to the death of the cross. For which cause God also hath exalted him, and hath given him a name which is above all names: That in the name of Jesus every knee should bow, of those that are in heaven, on earth, and under the earth: And that every tongue should confess that the Lord Jesus Christ is in the glory of God the Father.\par
\switchcolumn*

\LA
\begin{Resp}\Rsign Crux fidélis, inter omnes arbor una nóbilis: nulla silva talem profert, fronde, flore, gérmine: \Star{} Dulce lignum, dulces clavos, dulce pondus sustínuit. (Allelúja.)\end{Resp}
\begin{Resp}\Vsign Super ómnia ligna cedrórum tu sola excélsior.\end{Resp}
\begin{Resp}\Rsign Dulce lignum, dulces clavos, dulce pondus sustínuit. (Allelúja.)\end{Resp}
\switchcolumn
\EN
\begin{Resp}\Rsign Faithful Cross, above all other, one and only noble tree None in foliage, none in blossom, none in fruit thy peers may be \Star{} Sweetest wood and sweetest iron, Sweetest weight is hung on thee. Alleluia.\end{Resp}
\begin{Resp}\Vsign Thou art higher than all cedars.\end{Resp}
\begin{Resp}\Rsign Sweetest wood, and sweetest iron, Sweetest weight is hung on thee. Alleluia.\end{Resp}
\switchcolumn*

\LA
\LessonHead{Lectio III}
\switchcolumn
\EN
\LessonHead{Lesson III}
\switchcolumn*

\LA
\LessonTitle{De Epístola ad Colossénses}
\switchcolumn
\EN
\LessonTitle{From the letter to the Colossians.}
\switchcolumn*

\LA
\Cite{Col 2:9-15}
\switchcolumn
\EN
\Cite{Col 2:9-15}
\switchcolumn*

\LA
\noindent In Christo inhábitat omnis plenitúdo divinitátis corporáliter: et estis in illo repléti, qui est caput omnis principátus et potestátis; in quo et circumcísi estis circumcisióne non manu facta in exspoliatióne córporis carnis, sed in circumcisióne Christi; consepúlti ei in baptísmo, in quo et resurrexístis per fidem operatiónis Dei, qui suscitávit illum a mórtuis. Et vos, cum mórtui essétis in delíctis, et præpútio carnis vestræ, convivificávit cum illo, donans vobis ómnia delícta: delens quod advérsus nos erat chirógraphum decréti, quod erat contrárium nobis, et ipsum tulit de médio, affígens illud cruci: et expólians principátus, et potestátes tradúxit confidénter, palam triúmphans illos in semetípso.\par
\switchcolumn
\EN
\noindent For in him dwelleth all the fulness of the Godhead corporeally; And you are filled in him, who is the head of all principality and power: In whom also you are circumcised with circumcision not made by hand, in despoiling of the body of the flesh, but in the circumcision of Christ: Buried with him in baptism, in whom also you are risen again by the faith of the operation of God, who hath raised him up from the dead. And you, when you were dead in your sins, and the uncircumcision of your flesh; he hath quickened together with him, forgiving you all offences: Blotting out the handwriting of the decree that was against us, which was contrary to us. And he hath taken the same out of the way, fastening it to the cross: And despoiling the principalities and powers, he hath exposed them confidently in open show, triumphing over them in himself.\par
\switchcolumn*

\LA
\begin{Resp}\Rsign Hæc est arbor digníssima, in paradísi médio situáta, \Star{} In qua salútis auctor própria morte mortem ómnium superávit. (Allelúja, allelúja.)\end{Resp}
\begin{Resp}\Vsign Crux præcellénti decóre fúlgida, quam Hélena Constantíni mater concupiscénti ánimo requisívit.\end{Resp}
\begin{Resp}\Rsign In qua salútis auctor própria morte mortem ómnium superávit. (Allelúja, allelúja.)\end{Resp}
\begin{Resp}\Vsign Glória Patri, et Fílio, \Star{} et Spirítui Sancto.\end{Resp}
\begin{Resp}\Rsign In qua salútis auctor própria morte mortem ómnium superávit. (Allelúja, allelúja.)\end{Resp}
\switchcolumn
\EN
\begin{Resp}\Rsign This is that noble tree, planted in the midst of the garden \Star{} Whereon the Author of our salvation did by His Own death openly triumph over the death of all men. Alleluia, Alleluia.\end{Resp}
\begin{Resp}\Vsign Even the Cross, whereof the glory is so excellent, and whereafter Helen, the mother of Constantine, did so diligently search until she found it.\end{Resp}
\begin{Resp}\Rsign Whereon the Author of our salvation did by His Own death openly triumph over the death of all men. Alleluia, Alleluia.\end{Resp}
\begin{Resp}\Vsign Glory be to the Father, and to the Son, \Star{} and to the Holy Ghost.\end{Resp}
\begin{Resp}\Rsign Whereon the Author of our salvation did by His Own death openly triumph over the death of all men. Alleluia, Alleluia.\end{Resp}
\switchcolumn*

\LA
\LessonHead{Lectio IV}
\switchcolumn
\EN
\LessonHead{Lesson IV}
\switchcolumn*

\LA
\noindent Post insígnem victóriam, quam Constantínus imperátor, divínitus accépto signo Domínicæ Crucis, ex Maxéntio reportávit, Hélena Constantíni mater, in somnis admónita, conquiréndæ Crucis stúdio Jerosólymam venit; ubi marmóream Véneris státuam, in Crucis loco a Géntibus collocátam ad tolléndam Christi Dómini passiónis memóriam, post centum círciter octogínta annos, everténdam curávit. Quod item fecit ad præsépe Salvatóris et in loco resurrectiónis, inde Adónidis hinc Jovis subláto simulácro.\par
\switchcolumn
\EN
\noindent After that famous victory which the Emperor Constantine gained over Maxentius, (in the year 312,) on the eve of which the banner of the Cross of the Lord had been given to him from heaven, Helen, the mother of Constantine, being warned in a dream, came to Jerusalem, (in 326,) to seek for the Cross. There it was her care to cause to be overthrown the marble statue of Venus, which had stood on Calvary for about one hundred and eighty years, and which had originally been put there to desecrate and destroy the memorial of the sufferings of the Lord Christ. The like work Helen did at Bethlehem, by cleansing from an image of Adonis the stable where the Saviour was born, and from an idol of Jupiter, the place where He had arisen from the dead.\par
\switchcolumn*

\LA
\begin{Resp}\Rsign Nos autem gloriári opórtet in Cruce Dómini nostri Jesu Christi, in quo est salus, vita, et resurréctio nostra: \Star{} Per quem salváti et liberáti sumus. (Allelúja.)\end{Resp}
\begin{Resp}\Vsign Tuam Crucem adorámus, Dómine, et recólimus tuam gloriósam passiónem.\end{Resp}
\begin{Resp}\Rsign Per quem salváti et liberáti sumus. (Allelúja.)\end{Resp}
\switchcolumn
\EN
\begin{Resp}\Rsign But us it behoveth to glory in the Cross of our Lord Jesus Christ, in Whom is our salvation, life, and resurrection \Star{} Who hath saved us and redeemed us. Alleluia.\end{Resp}
\begin{Resp}\Vsign O Lord, we worship thy Cross, and make memorial of thy glorious passion.\end{Resp}
\begin{Resp}\Rsign Who hath saved us and redeemed us. Alleluia.\end{Resp}
\switchcolumn*

\LA
\LessonHead{Lectio V}
\switchcolumn
\EN
\LessonHead{Lesson V}
\switchcolumn*

\LA
\noindent Itaque loco Crucis purgáto, alte defóssæ tres cruces érutæ sunt, repertúsque seórsum ab illis Crucis Domínicæ títulus: qui cum ex tribus cui affíxus fuísset, non apparéret, eam dubitatiónem sústulit miráculum. Nam Macárius Jerosolymórum epíscopus, factis Deo précibus, síngulas cruces cuídam féminæ, gravi morbo laboránti, admóvit; cui cum réliquæ nihil profuíssent, adhíbita tértia Crux statim eam sanávit.\par
\switchcolumn
\EN
\noindent Then she had thus cleansed the place where the Cross had stood, Helen caused deep excavations to be made, which resulted in the discovery of three crosses, and, apart from them, the writing which had been nailed on that of the Lord. But which of the crosses had been His was unknown, and was only manifested by a miracle. Macarius, Bishop of Jerusalem, after offering solemn prayers to God, touched with each of the three a woman who was afflicted with a grievous disease. The two first had no effect, but at the touch of the third she was immediately healed.\par
\switchcolumn*

\LA
\begin{Resp}\Rsign Dum sacrum pignus cǽlitus exaltátur, Christi fides roborátur: \Star{} Adsunt prodígia divína in virga Móysi prímitus figuráta. (Allelúja, allelúja.)\end{Resp}
\begin{Resp}\Vsign Ad Crucis contáctum resúrgunt mórtui, et Dei magnália reserántur.\end{Resp}
\begin{Resp}\Rsign Adsunt prodígia divína in virga Móysi prímitus figuráta. (Allelúja, allelúja.)\end{Resp}
\switchcolumn
\EN
\begin{Resp}\Rsign The Relique true from heaven revealed, hath now the Gospel's figure sealed \Star{} As by the serpent Moses reared, so by the Cross the sick are healed. Alleluia, alleluia.\end{Resp}
\begin{Resp}\Vsign When the dead touch the Cross they arise, and the wonderful works of God are made manifest.\end{Resp}
\begin{Resp}\Rsign As by the serpent Moses reared, so by the Cross the sick are healed. Alleluia, alleluia.\end{Resp}
\switchcolumn*

\LA
\LessonHead{Lectio VI}
\switchcolumn
\EN
\LessonHead{Lesson VI}
\switchcolumn*

\LA
\noindent Hélena, salutári Cruce invénta, magnificentíssimam ibi exstrúxit ecclésiam, in qua partem Crucis relíquit thecis argénteis inclúsam, partem Constantíno fílio détulit; quæ Romæ repósita fuit in ecclésia sanctæ Crucis in Jerúsalem, ædificáta in ædibus Sessoriánis. Clavos étiam áttulit fílio, quibus sanctíssimum Jesu Christi corpus fixum fúerat. Quo ex témpore Constantínus legem sancívit, ne crux ad supplícium cuíquam adhiberétur. Ita res, quæ ántea homínibus probro ac ludíbrio fúerat, veneratióni et glóriæ esse cœpit.\par
\switchcolumn
\EN
\noindent Helen, after she had found the life-giving Cross, built over the site of the Passion a Church of extraordinary splendour, wherein she deposited part of the Cross, shut up in a silver case. Another part which she gave to her son, Constantine, was laid up in the Church of the Holy Cross of Jerusalem, which he built at Rome on the site of the Sessorian Palace. She also gave to her son the nails with which the Most Holy Body of Jesus Christ had been pierced. Constantine established a law abolishing the punishment of crucifixion for all time coming and thenceforth what had hitherto been a hissing and a curse among men, began to be esteemed worshipful and glorious.\par
\switchcolumn*

\LA
\begin{Resp}\Rsign Hoc signum Crucis erit in cælo, cum Dóminus ad judicándum vénerit: \Star{} Tunc manifésta erunt abscóndita cordis nostri. (Allelúja, allelúja.)\end{Resp}
\begin{Resp}\Vsign Cum séderit Fílius hóminis in sede majestátis suæ, et cœ́perit judicáre sǽculum per ignem.\end{Resp}
\begin{Resp}\Rsign Tunc manifésta erunt abscóndita cordis nostri. (Allelúja, allelúja.)\end{Resp}
\begin{Resp}\Vsign Glória Patri, et Fílio, \Star{} et Spirítui Sancto.\end{Resp}
\begin{Resp}\Rsign Tunc manifésta erunt abscóndita cordis nostri. (Allelúja, allelúja.)\end{Resp}
\switchcolumn
\EN
\begin{Resp}\Rsign This Sign of the Cross shall be in heaven, when the Lord cometh to judgment. \Star{} Then shall the secrets of our hearts be made manifest. Alleluia, alleluia.\end{Resp}
\begin{Resp}\Vsign When the Son of Man shall sit in the throne of His glory, and shall begin to judge the world by fire.\end{Resp}
\begin{Resp}\Rsign Then shall the secrets of our hearts be made manifest. Alleluia, alleluia.\end{Resp}
\begin{Resp}\Vsign Glory be to the Father, and to the Son, \Star{} and to the Holy Ghost.\end{Resp}
\begin{Resp}\Rsign Then shall the secrets of our hearts be made manifest. Alleluia, alleluia.\end{Resp}
\switchcolumn*

\LA
\LessonHead{Lectio VII}
\switchcolumn
\EN
\LessonHead{Lesson VII}
\switchcolumn*

\LA
\LessonTitle{Léctio sancti Evangélii secúndum Joánnem}
\switchcolumn
\EN
\LessonTitle{From the Holy Gospel according to John}
\switchcolumn*

\LA
\Cite{Joannes 3:1-15}
\switchcolumn
\EN
\Cite{John 3:1-15}
\switchcolumn*

\LA
\noindent In illo témpore: Erat homo ex pharisǽis, Nicodémus nómine, princeps Judæórum. Hic venit ad Jesum nocte, et dixit ei: Rabbi, scimus quia a Deo venísti magíster. Et réliqua.\par
\switchcolumn
\EN
\noindent At that time There was a man of the Pharisees named Nicodemus, a ruler of the Jews. The same came to Jesus by night, and said unto Him Rabbi, we know that Thou art a Teacher come from God. And so on.\par
\switchcolumn*

\LA
\LessonTitle{Homilía sancti Augustíni Epíscopi}
\switchcolumn
\EN
\LessonTitle{Homily by St. Augustine, Bishop of Hippo}
\switchcolumn*

\LA
\Source{Tract. 11 in Joann., post init.}
\switchcolumn
\EN
\Source{11th Tract on John.}
\switchcolumn*

\LA
\noindent Nicodémus ex his erat, qui credíderant in nómine Jesu, vidéntes signa et prodígia quæ faciébat. Supérius enim hoc dixit: Cum autem esset Jerosólymis in Pascha in die festo, multi credidérunt in nómine ejus. Quare credidérunt in nómine ejus? Séquitur, et dicit: Vidéntes signa ejus, quæ faciébat. Et de Nicodémo quid dicit? Erat princeps Judæórum, nómine Nicodémus. Hic venit ad eum nocte, et ait illi: Rabbi, scimus quia a Deo venísti magíster. Et iste ergo credíderat in nómine ejus. Et ipse unde credíderat? Séquitur: Nemo enim potest hæc signa fácere, quæ tu facis, nisi fúerit Deus cum eo.\par
\switchcolumn
\EN
\noindent Nicodemus was one of them which believed in the Name of Jesus, when they saw the signs and wonders which He did. So hath John given us to understand a few words before our text: “Now when He was in Jerusalem at the Passover, in the feast-day, many believed in His Name” And wherefore did they believe in His Name John telleth us immediately: “When they saw the miracles which He did.” And now, what saith he touching Nicodemus? “There was a man of the Pharisees, named Nicodemus, a ruler of the Jews. The same came to Jesus by night, and said unto Him Rabbi, we know that Thou art a Teacher come from God.” Nicodemus therefore believed in His Name. And why did he believe He saith “For no man can do these miracles that Thou doest, except God be with him.”\par
\switchcolumn*

\LA
\begin{Resp}\Rsign Dulce lignum, dulces clavos, dulce pondus sustínuit: \Star{} Quæ sola digna fuit portáre prétium hujus sǽculi. (Allelúja.)\end{Resp}
\begin{Resp}\Vsign Hoc signum Crucis erit in cælo cum Dóminus ad judicándum vénerit.\end{Resp}
\begin{Resp}\Rsign Quæ sola digna fuit portáre prétium hujus sǽculi. (Allelúja.)\end{Resp}
\switchcolumn
\EN
\begin{Resp}\Rsign Sweetest wood and sweetest iron, Sweetest weight is hung on thee \Star{} Thou alone wast counted worthy this world's ransom to uphold. Alleluia.\end{Resp}
\begin{Resp}\Vsign The sign of the Cross shall be in heaven when the Lord cometh to judgment.\end{Resp}
\begin{Resp}\Rsign Thou alone wast counted worthy this world's ransom to uphold. Alleluia.\end{Resp}
\switchcolumn*

\LA
\LessonHead{Lectio VIII}
\switchcolumn
\EN
\LessonHead{Lesson VIII}
\switchcolumn*

\LA
\noindent Si ergo Nicodémus de illis multis erat, qui credíderant in nómine ejus, jam in isto Nicodémo attendámus, quare Jesus non se credébat eis. Respóndit Jesus et dixit ei: Amen, amen dico tibi, nisi quis renátus fúerit dénuo, non potest vidére regnum Dei. Ipsis ergo se credit Jesus, qui nati fúerint dénuo. Ecce illi credíderant in eum, et Jesus non se credébat eis. Tales sunt omnes catechúmeni: ipsi jam credunt in nómine Christi, sed Jesus non se credit eis. Inténdat et intélligat cáritas vestra. Si dixérimus catechúmeno: Credis in Christum? respóndet, Credo, et signat se Cruce Christi: portat in fronte, et non erubéscit de Cruce Dómini sui. Ecce credit in nómine ejus. Interrogémus eum: Mandúcas carnem Fílii hóminis, et bibis sánguinem Fílii hóminis? nescit quid dícimus, quia Jesus non se crédidit ei.\par
\switchcolumn
\EN
\noindent Lo, then, Nicodemus was one of the many which had believed in His Name, let us seek to find in Nicodemus why “Jesus did not commit Himself unto them” Jesus answered and said unto him “Amen, Amen, I say unto thee except a man be born again, he cannot see the kingdom of God.” Jesus therefore committeth Himself unto such as be born again. Behold, Nicodemus and they that were with him believed in Jesus, but Jesus did not commit Himself unto them. Just so are all Catechumens they believe in the Name of Christ, but Jesus hath not yet committed Himself unto them. Now I trust ye will be good enough to pay attention, and understand what I am going to say. If ye ask of a Catechumen “Dost thou believe in Christ” he saith “I believe,” and he signeth himself with the sign of the Cross. The Cross of his Lord is marked upon his forehead, and he is not ashamed of it. Behold, he believeth in the Name of Christ. But let us ask him “Dost thou eat the flesh of the Son of Man” and he knoweth not what we mean, for Jesus hath not yet committed Himself unto him.\par
\switchcolumn*

\LA
\begin{Resp}\Rsign Sicut Móyses exaltávit serpéntem in desérto, ita exaltári opórtet Fílium hóminis: \Star{} Ut omnis qui credit in ipsum, non péreat, sed hábeat vitam ætérnam. (Allelúja.)\end{Resp}
\begin{Resp}\Vsign Non misit Deus Fílium suum in mundum ut júdicet mundum, sed ut salvétur mundus per ipsum.\end{Resp}
\begin{Resp}\Rsign Ut omnis qui credit in ipsum, non péreat, sed hábeat vitam ætérnam. (Allelúja.)\end{Resp}
\begin{Resp}\Vsign Glória Patri, et Fílio, \Star{} et Spirítui Sancto.\end{Resp}
\begin{Resp}\Rsign Ut omnis qui credit in ipsum, non péreat, sed hábeat vitam ætérnam. (Allelúja.)\end{Resp}
\switchcolumn
\EN
\begin{Resp}\Rsign As Moses lifted up the serpent in the wilderness, even so must the Son of Man be lifted up \Star{} That whosoever believeth in Him should not perish, but have everlasting life. Alleluia.\end{Resp}
\begin{Resp}\Vsign God sent not His Son into the world to condemn the world, but that the world through Him might be saved.\end{Resp}
\begin{Resp}\Rsign That whosoever believeth in Him should not perish, but have everlasting life. Alleluia.\end{Resp}
\begin{Resp}\Vsign Glory be to the Father, and to the Son, \Star{} and to the Holy Ghost.\end{Resp}
\begin{Resp}\Rsign That whosoever believeth in Him should not perish, but have everlasting life. Alleluia.\end{Resp}
\switchcolumn*

\end{paracol}

\par\needspace{10\baselineskip}\addvspace{1.2em}{\centering\small\textsc{Die 3 Maii} · \textsc{3 May}\par}
\Office{Ss. Alexandri I Papæ, Eventii et Theoduli Mm. ac Juvenalis Ep. et Conf.}{Sts. Alexander and Companions, Martyrs}{Simplex}
\addcontentsline{toc}{subsection}{Die 3 Maii · Ss. Alexandri I Papæ, Eventii et Theoduli Mm. ac Juvenalis Ep. et Conf.\enspace\textit{Sts. Alexander and Companions, Martyrs}}
\begin{paracol}{2}
\LA
\LessonHead{Lectio IX (pro commemoratione)}
\switchcolumn
\EN
\LessonHead{Lesson IX (when commemorated)}
\switchcolumn*

\LA
\LessonTitle{Commemoratio: Ss. Alexandri I Papæ, Eventii et Theoduli Mm. ac Juvenalis Ep. et Conf.}
\switchcolumn
\EN
\LessonTitle{Commemoration of: Sts. Alexander and Companions, Martyrs}
\switchcolumn*

\LA
\noindent Alexander Romanus, Hadriáno imperatóre regens Ecclésiam, magnam partem Romanæ nobilitátis ad Christum convértit. Is constítuit ut tantummodo panis et vinum in mysterio offerrétur: vinum autem aqua miscéri jussit, propter sánguinem et aquam, quæ ex Jesu Christi látere profluxérunt; et in Cánone Missæ addidit, Qui pridie quam paterétur. Idem decrevit ut aqua benedicta, sale admixto, perpetuo in ecclésia asservarétur, et in cubiculis adhiberétur ad fugándos dæmones. Sedit annos decem, menses quinque et dies viginti, vitæ sanctitáte et salutáribus institútis illustris. Martyrio coronátus est una cum Eventio et Theodúlo presbyteris, sepultúsque est via Nomentana, tertio ab Urbe lápide, eodem in loco, ubi secúri percússus fuerat; creátis diverso témpore mense Decembri presbyteris sex, diáconis duobus et epíscopis per diversa loca quinque. Eórum corpora, póstea in Urbem translata, in ecclésia sanctæ Sabinæ condita sunt. In eumdem diem incidit beáta mors sancti Juvenalis, Narniénsis epíscopi; qui, cum plurimos in ea urbe sanctitáte et doctrina Christo peperísset, clarus miraculis in pace quiévit ibique honorifice sepúltus est.\par
\switchcolumn
\EN
\noindent Alexander was a Roman, who ruled the Church during the reign of the Emperor Hadrian. He turned to Christ a great number of the Roman nobility. He ordained that nothing but bread and wine should be offered at the mystery, but that some water should be mingled with the wine, in memory of the Blood and Water Which flowed from the Side of Jesus Christ. He added to the Canon of the Mass the words “Who, the night before He suffered.” He also ordained that blessed water mingled with salt, should be kept always in Churches, and should be used in private rooms to scare away devils. He sat in the throne of Peter ten years, five months, and twenty days. He hath great renown on account of the holiness of his life, and the usefulness of his institutions. He was crowned with martyrdom, \textasciitilde{}(in the year 119,) together with the Priests Eventius and Theodulus, and was buried beside the Nomentan Way, at the third mile-stone from the city, in the same place where he had been beheaded. During his Popedom he held diverse Advent ordinations, and at them ordained six Priests, two Deacons, and five Bishops for diverse places. The bodies of these three Martyrs, Alexander, Eventius, and Theodulus, were afterwards brought into the city, and buried in the Church of St Sabina. On this day likewise, (about the year 367,) occurred the blessed death of Juvenal, the holy Bishop of Narni, who by the holiness of his life and teaching, became the father in Christ of so many of the dwellers in that city. He fell asleep very peacefully, with great fame for miracles, and was there honourably buried.\par
\switchcolumn*

\LA
\TeDeum{Te Deum}
\switchcolumn
\EN
\TeDeum{Te Deum}
\switchcolumn*

\end{paracol}

\par\needspace{10\baselineskip}\addvspace{1.2em}{\centering\small\textsc{Die 4 Maii} · \textsc{4 May}\par}
\Office{S. Monicæ Viduæ}{St. Monica, Widow}{Duplex}
\addcontentsline{toc}{subsection}{Die 4 Maii · S. Monicæ Viduæ\enspace\textit{St. Monica, Widow}}
\begin{paracol}{2}
\LA
\RefLine{Lectiones I–III: de Scriptura occurrente.}
\switchcolumn
\EN
\RefLine{Lessons I–III: from the Scripture of the day.}
\switchcolumn*

\LA
\LessonHead{Lectio IV}
\switchcolumn
\EN
\LessonHead{Lesson IV}
\switchcolumn*

\LA
\noindent Monica, sancti Augustíni dupliciter mater, quia eum et mundo et cælo péperit, maríto mortuo, quem senectute confectum Jesu Christo conciliávit, castam et opéribus misericórdiæ exercitam viduitátem agebat; in assiduis vero ad Deum oratiónibus pro fílio, qui in Manichæórum sectam inciderat, lácrimas effundébat. Quem étiam Mediolanum secuta est; ubi ipsum frequenter hortabátur, ut ad episcopum Ambrosium se conferret. Quod cum ille fecísset, ejus et publicis conciónibus et privátis colloquiis catholicæ fidei veritátem edoctus, ab eodem baptizátus est.\par
\switchcolumn
\EN
\noindent Monica was twice over the mother of St. Augustine, for, under God, he owed to her both earth and heaven. When her husband was very old she made him a friend of Jesus Christ, and after his death she lived a widow in all purity and constantly occupied in works of mercy. Her son Augustine had fallen into the heresy of the Manichaeans, and for his conversion she earnestly pleaded with God for years, with strong crying and tears. She followed Augustine to Milan, and tenderly and constantly besought him to confer with Ambrose the Bishop. This he consented to do, and at last, through the public sermons and private conversations of Ambrose, his eyes were opened to see the truth of the Catholic Religion, and he received baptism at the Bishop's hands, (at the Passover of the year 387.)\par
\switchcolumn*

\LA
\begin{Resp}\Rsign Propter veritátem, et mansuetúdinem, et justítiam: \Star{} Et dedúcet te mirabíliter déxtera tua, allelúja.\end{Resp}
\begin{Resp}\Vsign Spécie tua et pulchritúdine tua inténde, próspere procéde, et regna.\end{Resp}
\begin{Resp}\Rsign Et dedúcet te mirabíliter déxtera tua, allelúja.\end{Resp}
\switchcolumn
\EN
\begin{Resp}\Rsign Because of truth, and meekness, and righteousness; \Star{} And thy right hand shall lead thee wonderfully, alleluia.\end{Resp}
\begin{Resp}\Vsign In thy comeliness, and thy beauty, go forward, fare prosperously, and reign.\end{Resp}
\begin{Resp}\Rsign And thy right hand shall lead thee wonderfully, alleluia.\end{Resp}
\switchcolumn*

\LA
\LessonHead{Lectio V}
\switchcolumn
\EN
\LessonHead{Lesson V}
\switchcolumn*

\LA
\noindent Mox in Africam redeuntes, cum ad Ostia Tiberina constitissent, incidit in febrem. Quo in morbo cum eam quodam die ánima defecísset, ut se collegit, Ubi, inquit, eram? Et astántes íntuens, Ponite hic matrem vestram: tantum vos rogo, ut ad altáre Dómini meminéritis mei. Nono autem die beáta mulier ánimam Deo réddidit. Ejus corpus ibi in ecclésia sanctæ Aureæ sepúltum est: quod póstea, Martino quinto summo Pontifice, Romam translátum, in ecclésia sancti Augustíni honorifice cónditum est.\par
\switchcolumn
\EN
\noindent The mother and son set out to return to their home in Africa, but after they had reached Ostia at the mouth of the Tiber, she was stricken down by a fever. One day as she lay sick, she came to herself after her mind had been long wandering, and said: “Where am I?” Then she saw who were standing by, and said “Let your mother lie here only, remember me at the altar of the Lord.” On the ninth day this blessed lady surrendered her spirit to God. Her body was buried there at Ostia in the Church of St. Aurea, but, long after, in the papacy of Martin V, it was carried to Rome and honourably buried again in the Church of St. Augustine.\par
\switchcolumn*

\LA
\begin{Resp}\Rsign Dilexísti justítiam, et odísti iniquitátem: \Star{} Proptérea unxit te Deus, Deus tuus, óleo lætítiæ, allelúja.\end{Resp}
\begin{Resp}\Vsign Propter veritátem, et mansuetúdinem, et justítiam.\end{Resp}
\begin{Resp}\Rsign Proptérea unxit te Deus, Deus tuus, óleo lætítiæ, allelúja.\end{Resp}
\switchcolumn
\EN
\begin{Resp}\Rsign Thou hast loved righteousness, and hated iniquity; \Star{} Therefore God, thy God, hath anointed thee with the oil of gladness, alleluia.\end{Resp}
\begin{Resp}\Vsign Because of truth, and meekness, and righteousness.\end{Resp}
\begin{Resp}\Rsign Therefore God, thy God, hath anointed thee with the oil of gladness, alleluia.\end{Resp}
\switchcolumn*

\LA
\LessonHead{Lectio VI}
\switchcolumn
\EN
\LessonHead{Lesson VI}
\switchcolumn*

\LA
\Source{Ex libro 9 Conf. cap. 12.}
\switchcolumn
\EN
\Source{Confessions of St. Augustine. Bk. ix. ch. 12}
\switchcolumn*

\LA
\noindent Subdit vero Augustinus de matris morte disserens: Neque enim decére arbitrabámur funus illud quéstubus lacrimosis gemitibúsque celebrare, quia illa nec mísere, nec omnino moriebátur: hoc et documéntis morum ejus, et fide non ficta rationibúsque certis tenebamus. Atque inde paulatim reducébam in prístinum sensum ancíllam tuam, conversationémque ejus piam in te et sanctam, in nos blandam atque morígeram, qua súbito destitútus sum; et libuit flere de illa et pro illa. Et si quis peccátum invenerit, flevisse me matrem meam exígua parte horæ, matrem óculis meis mortuam, quæ me multos annos fleverat, ut óculis suis viverem, non irrideat; sed potius, si est grandi caritate, pro peccátis meis fleat ipse ad te Patrem ómnium fratrum Christi tui.\par
\switchcolumn
\EN
\noindent Augustine added these words after describing his mother's death: “We did not think that hers was a death which it was seemly to mark with repining, or tears, or lamentations, seeing that she died not sorrowfully, nor at all as touching her best and noblest part. This we knew, because we knew what her life had been, her faith unfeigned, her sure and certain hope. And then, nevertheless, I remembered again what thine handmaid was used to be, her walk with thee, how godly and holy it was, and with us so gentle and long-suffering and that it was all, gone away from me now. And I wept, over her and for her. And if any man will make it blame to me that I wept for a little while, when I saw lying dead before my eyes my mother, who had wept over me so many years, that she might see me live, I say, if any man will make it blame to me, I pray him not to sneer at me, but rather \textasciitilde{}(if his charity be so great) himself to weep over my sins before thee, Who art a Father to all them to whom thy Christ is a Brother.”\par
\switchcolumn*

\LA
\begin{Resp}\Rsign Fallax grátia, et vana est pulchritúdo: \Star{} Múlier timens Dóminum, ipsa laudábitur, allelúja.\end{Resp}
\begin{Resp}\Vsign Date ei de fructu mánuum suárum, et laudent eam in portis ópera ejus.\end{Resp}
\begin{Resp}\Rsign Múlier timens Dóminum, ipsa laudábitur, allelúja.\end{Resp}
\begin{Resp}\Vsign Glória Patri, et Fílio, \Star{} et Spirítui Sancto.\end{Resp}
\begin{Resp}\Rsign Múlier timens Dóminum, ipsa laudábitur, allelúja.\end{Resp}
\switchcolumn
\EN
\begin{Resp}\Rsign Favour is deceitful, and beauty is vain \Star{} A woman that feareth God she shall be praised, alleluia.\end{Resp}
\begin{Resp}\Vsign Give her of the fruit of her hands, and let her own works praise her in the gates.\end{Resp}
\begin{Resp}\Rsign A woman that feareth God she shall be praised, alleluia.\end{Resp}
\begin{Resp}\Vsign Glory be to the Father, and to the Son, \Star{} and to the Holy Ghost.\end{Resp}
\begin{Resp}\Rsign A woman that feareth God she shall be praised, alleluia.\end{Resp}
\switchcolumn*

\LA
\LessonHead{Lectio VII}
\switchcolumn
\EN
\LessonHead{Lesson VII}
\switchcolumn*

\LA
\LessonTitle{Léctio sancti Evangélii secúndum Lucam}
\switchcolumn
\EN
\LessonTitle{From the Holy Gospel according to Luke}
\switchcolumn*

\LA
\Cite{Luc 7:11-16}
\switchcolumn
\EN
\Cite{Luke 7:11-16}
\switchcolumn*

\LA
\noindent In illo témpore: Ibat Jesus in civitátem, quæ vocátur Naim: et ibant cum eo discípuli ejus, et turba copiosa. Et réliqua.\par
\switchcolumn
\EN
\noindent At that time Jesus went into a city called Naim; and His disciples went with Him, and much people. And so on.\par
\switchcolumn*

\LA
\LessonTitle{Homilía sancti Augustíni Epíscopi}
\switchcolumn
\EN
\LessonTitle{Homily by St. Augustine, (Bishop of Hippo.)}
\switchcolumn*

\LA
\Source{Sermo 44 de Verbis Dómini, circa initium}
\switchcolumn
\EN
\Source{44th Discourse on the Words of the Lord.}
\switchcolumn*

\LA
\noindent De juvene illo resuscitato gavisa est mater vidua; de homínibus in spíritu quotídie suscitátis gaudet mater Ecclésia. Ille quidem mórtuus erat corpore; illi autem mente. Illius mors visibilis visibíliter plangebátur; illórum mors invisibilis nec quærebátur, nec videbátur. Quæsívit ille, qui nóverat mórtuos. Ille solus nóverat mórtuos, qui poterat fácere vivos. Nisi enim ad mórtuos suscitándos venísset, Apóstolus non diceret: Surge, qui dormis, et exsúrge a mórtuis, et illuminábit te Christus.\par
\switchcolumn
\EN
\noindent That her son was called, again to life was the joy of that widowed mother that souls of men are every day called to life is the joy of our Mother the Church. He was dead in body they have been dead in mind. His death was outward, and was outwardly bewailed their inward death hath been neither mourned for nor seen. But He hath sought for them, Who hath seen that they are dead, and He only hath seen that they are dead, Who hath been able to make them alive. If He had not come to raise the dead, the Apostle had not said: “Awake, thou that sleepest, and arise from the dead, and Christ shall give thee light.” (Eph. v. 14.)\par
\switchcolumn*

\LA
\begin{Resp}\Rsign Os suum apéruit sapiéntiæ, et lex cleméntiæ in lingua ejus: considerávit sémitas domus suæ, \Star{} Et panem otiósa non comédit, allelúja.\end{Resp}
\begin{Resp}\Vsign Gustávit et vidit quia bona est negotiátio ejus: non exstinguétur in nocte lucérna ejus.\end{Resp}
\begin{Resp}\Rsign Et panem otiósa non comédit, allelúja.\end{Resp}
\switchcolumn
\EN
\begin{Resp}\Rsign She openeth her mouth with wisdom, and in her tongue is the law of kindness. She looketh well to the ways of her household \Star{} And eateth not the bread of idleness, alleluia.\end{Resp}
\begin{Resp}\Vsign She tasteth and perceiveth that her merchandise is good. Her candle goeth not out by night.\end{Resp}
\begin{Resp}\Rsign And she eateth not the bread of idleness, alleluia.\end{Resp}
\switchcolumn*

\LA
\LessonHead{Lectio VIII}
\switchcolumn
\EN
\LessonHead{Lesson VIII}
\switchcolumn*

\LA
\noindent Tres autem mórtuos invenímus a Dómino resuscitatos visibíliter, míllia invisibíliter. Quot autem mórtuos visibíliter suscitaverit, quis novit? Non enim ómnia, quæ fecit, scripta sunt. Joánnes hoc dixit: Multa alia fecit Jesus, quæ si scripta essent, árbitror totum mundum non posse libros cápere. Multi ergo sunt alii sine dubio suscitati, sed non tres frustra commemorati. Dóminus enim noster Jesus Christus ea, quæ faciebat corporáliter, étiam spiritáliter volebat intelligi. Neque enim tantum miracula propter miracula faciebat; sed ut illa, quæ faciebat, mira essent vidéntibus, vera essent intelligéntibus.\par
\switchcolumn
\EN
\noindent It is written how the Lord raised from the dead three persons visibly, but thousands invisibly. But how many they may have been whom He raised visibly, who knoweth? For all the things which He did are not written. John saith thus: “There are also many other things which Jesus did, the which, if they should be written every one, I suppose that even the world itself could not contain the books that should be written.” (xxi. 25.) There were then, doubtless, many more raised to life, but it is not meaningless that three are recorded. For our Lord Jesus Christ hath willed that those things which He did carnally, we should understand also spiritually. He worked not miracles only for the sake of working wonders, but that His works might be at once wonderful to them that beheld, and true to them that understand them.\par
\switchcolumn*

\LA
\begin{Resp}\Rsign Regnum mundi et omnem ornátum sǽculi contémpsi, propter amórem Dómini mei Jesu Christi: \Star{} Quem vidi, quem amávi, in quem crédidi, quem diléxi, allelúja.\end{Resp}
\begin{Resp}\Vsign Eructávit cor meum verbum bonum: dico ego ópera mea Regi.\end{Resp}
\begin{Resp}\Rsign Quem vidi, quem amávi, in quem crédidi, quem diléxi, allelúja.\end{Resp}
\begin{Resp}\Vsign Glória Patri, et Fílio, \Star{} et Spirítui Sancto.\end{Resp}
\begin{Resp}\Rsign Quem vidi, quem amávi, in quem crédidi, quem diléxi, allelúja.\end{Resp}
\switchcolumn
\EN
\begin{Resp}\Rsign The kingdom of this world and all the beauty of life I have esteemed as nothing, for the excellency of the love of Jesus Christ my Lord \Star{} Whom, having seen, I loved Whom, having believed, I longed after, alleluia.\end{Resp}
\begin{Resp}\Vsign My heart is overflowing with a good matter I speak of my works unto the King.\end{Resp}
\begin{Resp}\Rsign Whom, having seen, I loved Whom, having believed, I longed after, alleluia.\end{Resp}
\begin{Resp}\Vsign Glory be to the Father, and to the Son, \Star{} and to the Holy Ghost.\end{Resp}
\begin{Resp}\Rsign Whom, having seen, I loved Whom, having believed, I longed after, alleluia.\end{Resp}
\switchcolumn*

\LA
\LessonHead{Lectio IX}
\switchcolumn
\EN
\LessonHead{Lesson IX}
\switchcolumn*

\LA
\noindent Quemádmodum qui videt litteras in códice óptime scripto, et non novit légere, laudat quidem antiquarii manum, admirans ápicum pulchritúdinem; sed quid sibi velint, quid índicent illi ápices, nescit, et est óculis laudator, mente non cógnitor. Alius autem et laudat artifícium, et capit intelléctum: ille útique, qui non solum vidére quod commune est ómnibus potest, sed étiam légere; quod qui non dídicit, non potest. Ita qui vidérunt Christi miracula, et non intellexérunt quid sibi vellent, et quid intelligéntibus quodámmodo innúerent, miráti sunt tantum quia facta sunt; alii vero et facta mirati, et intellecta assecuti. Tales nos in schola Christi esse debemus.\par
\switchcolumn
\EN
\noindent Even as one that looketh upon a scroll right fairly written, and knoweth not how to read therein, praiseth the hand of the old scribe when he seeth the beauty of the points, but what it saith, what those points mean, he knoweth not, and praiseth by the eye, without understanding by the mind, and as, on the other hand, he that can not only gaze on it, as can all men, but also can read it, praiseth the penmanship, and catcheth the sense likewise, which the unlearned cannot do even so, there were some that saw the miracles which Christ did, and understood not what they meant, nor what they, as it were, hinted to such as did understand them, and these only marvelled to see them wrought. And other some there were which saw the works, and marvelled, and understood them, and profited by them. And it is as these last that we ought to be in the school of Christ.\par
\switchcolumn*

\LA
\TeDeum{Te Deum}
\switchcolumn
\EN
\TeDeum{Te Deum}
\switchcolumn*

\LA
\LessonHead{Lectio IX (pro commemoratione)}
\switchcolumn
\EN
\LessonHead{Lesson IX (when commemorated)}
\switchcolumn*

\LA
\LessonTitle{Commemoratio: S. Monicæ Viduæ}
\switchcolumn
\EN
\LessonTitle{Commemoration of: St. Monica, Widow}
\switchcolumn*

\LA
\noindent Mónica, sancti Augustíni piíssima mater, uxórum et viduárum exémplar, fílium Manichǽum Mediolánum secúta, postquam assíduis oratiónibus, lácrimis et jejúniis, ope Ambrósii epíscopi, eum Christo lucrifécit, cum ipso in Africam rédiens, ad Ostia Tiberína in febrim íncidit, et nono die placidíssime óbiit. De cujus morte mærens fílius dísserens subdit: Neque enim decére arbitrabámur funus illud quéstubus lacrimósis gemitibúsque celebráre, quia illa nec mísere, nec omníno moriebátur: hoc et documéntis morum ejus, et fide non ficta rationibúsque certis tenebámus. Atque inde paulátim reducébam in prístinum sensum ancíllam tuam, conversationémque ejus piam in te et sanctam, in nos blandam atque morígeram, qua súbito destitútus sum; et líbuit flere de illa et pro illa. Et si quis peccátum invénerit, flevísse me matrem óculis meis mórtuam, quæ me multos annos fléverat, ut óculis suis víverem, non irrídeat; sed pótius, si est grandi caritáte, pro peccátis meis fleat ipse ad te Patrem ómnium fratrum Christi tui. Ejus corpus, in ecclésia sanctæ Aureæ prius sepúltum, póstea, Martíno quinto summo Pontífice, Romam translátum, in ecclésia sancti Augustíni honorífice cónditum est.\par
\switchcolumn
\EN
\noindent Monica, the devout mother of St. Augustine and a shining example to wives and mothers, followed her son to Milan when he become a Manichaean and, by her constant prayers, tears and fasting, gained him for Christ with the help of Ambrose the bishop. When she was returning with him to Africa, she fell ill of fever at Ostia on the Tiber, and, after nine days, died peacefully. After telling about her death, her sorrowing son adds: We did not think that hers was a death which it was seemly to mark with repining, or tears, or lamentations, seeing that she died not sorrowfully, nor at all as touching her best and noblest part. This we knew, because we knew what her life had been, her faith unfeigned, her sure and certain hope. And then, nevertheless, I remembered again what thine handmaid was used to be, her walk with thee, how godly and holy it was, and with us so gentle and long-suffering; and that it was all gone away from me now. And I wept, over her and for her. And if any man will make it blame to me that I wept for a little while, when I saw lying dead before mine eyes my mother, who had wept over me so many years, that she might see me live, I say, if any man will make it blame to me, I pray him not to sneer at me, but rather (if his charity be so great) himself to weep over my sins before thee, who art a Father to all them to whom thy Christ is a Brother. St. Monica's body was first buried in the Church of St. Aurea, but, long after, in the papacy of Martin V, it was carried to Rome and honourably buried again in the Church of St. Augustine.\par
\switchcolumn*

\LA
\TeDeum{Te Deum}
\switchcolumn
\EN
\TeDeum{Te Deum}
\switchcolumn*

\end{paracol}

\par\needspace{10\baselineskip}\addvspace{1.2em}{\centering\small\textsc{Die 5 Maii} · \textsc{5 May}\par}
\Office{S. Pii V Papæ et Confessoris}{St. Pius V, Pope and Confessor}{Duplex}
\addcontentsline{toc}{subsection}{Die 5 Maii · S. Pii V Papæ et Confessoris\enspace\textit{St. Pius V, Pope and Confessor}}
\begin{paracol}{2}
\LA
\RefLine{Lectiones I–III: de Scriptura occurrente.}
\switchcolumn
\EN
\RefLine{Lessons I–III: from the Scripture of the day.}
\switchcolumn*

\LA
\RefLine{Lectiones VII–IX: ex Commúni Summorum Pontificum Confessorum (C4b).}
\switchcolumn
\EN
\RefLine{Lessons VII–IX: from the Common of Popes (C4b).}
\switchcolumn*

\LA
\LessonHead{Lectio IV}
\switchcolumn
\EN
\LessonHead{Lesson IV}
\switchcolumn*

\LA
\noindent Pius in oppido Insubriæ, quod Boschum vocant, natus, sed e Bonónia oriundus ex nobili Ghisleriórum família, cum quatuordecim esset annórum, ordinem Prædicatórum ingréssus est. Erat in eo admirábilis patiéntia, profúnda humilitas, summa vitæ austeritas, continuum oratiónis studium, et regularis observantiæ ac divini honoris ardentíssimus zelus. Philosophíæ vero ac theologíæ incumbens, adeo in iis excelluit, ut illas docéndi munus magna cum laude per multos annos exercúerit. Sacras conciónes plúribus in locis cum ingénti auditórum fructu hábuit. Inquisitoris offícium inviolábili animi fortitúdine diu sustinuit; multasque civitátes, non sine vitæ discrimine, ab hæresi tunc grassante immunes servávit.\par
\switchcolumn
\EN
\noindent Michael Ghislieri afterwards proclaimed Pope under the name of Pius V, was born \textasciitilde{}(on the 27th of January, in the year 1504,) at the town of Bosco in the Milanese, but his family was a noble one of Bologna. At the age of fourteen years he entered the order of Friars Preachers. He was a man marked by a wonderful long-suffering, a deep lowliness, a great hardness of living, an unwavering earnestness in prayer, and a most strong zeal for the perfect observance of the Rule of his Order, and for the greater glory of God. He gave himself to the study of Philosophy and Theology, and was so learned in both, that he discharged for many years with great reputation the duties of a Professor of those sciences. He preached publicly in many places, to the great profit of his hearers. He long did the work of Inquisitor with unflinching spirit, and preserved many cities, not without risk to his own life, from the heresy which was then creeping in everywhere.\par
\switchcolumn*

\LA
\begin{Resp}\Rsign Invéni David servum meum, óleo sancto meo unxi eum: \Star{} Manus enim mea auxiliábitur ei, allelúja.\end{Resp}
\begin{Resp}\Vsign Nihil profíciet inimícus in eo, et fílius iniquitátis non nocébit ei.\end{Resp}
\begin{Resp}\Rsign Manus enim mea auxiliábitur ei, allelúja.\end{Resp}
\switchcolumn
\EN
\begin{Resp}\Rsign I have found David My servant, with My holy oil have I anointed him \Star{} For My hand shall help him, alleluia.\end{Resp}
\begin{Resp}\Vsign The enemy shall prevail nothing against him, nor the son of wickedness afflict him.\end{Resp}
\begin{Resp}\Rsign For My hand shall help him, alleluia.\end{Resp}
\switchcolumn*

\LA
\LessonHead{Lectio V}
\switchcolumn
\EN
\LessonHead{Lesson V}
\switchcolumn*

\LA
\noindent A Paulo quarto, cui ob exímias virtútes caríssimus erat, ad Nepesínum et Sutrinum episcopátum promotus, et post biennium inter Romanæ Ecclésiæ presbyteros cardinales adscriptus fuit. Tum ad ecclésiam Montis Regalis in Subalpinis a Pio quarto translatus, cum plures in eam abusus irrepsisse cognovísset, totam diœcesim lustrávit; rebusque compositis Romam reversus, gravíssimis expediéndis negotiis applicatus, quod justum erat, apostolica libertáte et constántia decernebat. Mortuo autem Pio, præter ómnium exspectatiónem electus Pontifex, nihil in vitæ ratióne, excepto exteriori habitu, immutávit. Fuit in eo religiónis propagandæ perpétuum studium, in ecclesiástica disciplína restituénda indeféssus labor, in exstirpandis erróribus assidua vigilantia, in sublevandis egéntium necessitátibus indeficiens beneficéntia, in Sedis apostolicæ júribus vindicándis robur invictum.\par
\switchcolumn
\EN
\noindent Paul IV, to whom his virtues had greatly endeared him, raised him (in 1556,) to the united Bishoprics of Nepi and Sutri, and after two years he was enrolled among the Cardinal Priests of the Roman Church. Pius IV. translated him to the Church of Mondovi in Piedmont, wherein, on his coming, he found that many corruptions had crept in. He reformed the whole of his diocese, and, after settling his affairs, returned to Rome, where his attention was called to matters of the gravest business, in determining which he used Apostolic boldness and firmness. After the death of Pius IV, the fifth Pius, to the astonishment of all men, was elected to succeed him, (on the 7th of January, 1566.) On becoming Pope he changed his way of life in no respect except as regarded his raiment. The Propagation of Religion was to him the object of unceasing care the restoration of the Discipline of the Church, of unwearied toil the uprooting of error, of sleepless watchfulness the relieving the needs of the poor, of unfailing charity the maintenance of the rights of the Apostolic See, of adamantine firmness.\par
\switchcolumn*

\LA
\begin{Resp}\Rsign Pósui adjutórium super poténtem, et exaltávi eléctum de plebe mea: \Star{} Manus enim mea auxiliábitur ei, allelúja.\end{Resp}
\begin{Resp}\Vsign Invéni David servum meum, óleo sancto meo unxi eum.\end{Resp}
\begin{Resp}\Rsign Manus enim mea auxiliábitur ei, allelúja.\end{Resp}
\switchcolumn
\EN
\begin{Resp}\Rsign I have laid help upon one that is mighty, and have exalted one chosen out of My people \Star{} For My hand shall help him, alleluia.\end{Resp}
\begin{Resp}\Vsign I have found David My servant, with My holy oil have I anointed him.\end{Resp}
\begin{Resp}\Rsign For My hand shall help him, alleluia.\end{Resp}
\switchcolumn*

\LA
\LessonHead{Lectio VI}
\switchcolumn
\EN
\LessonHead{Lesson VI}
\switchcolumn*

\LA
\noindent Selímum Turcárum tyrannum multis elátum victoriis, ingénti comparáta classe, ad Echínadas ínsulas non tam armis quam fusis ad Deum precibus devicit. Quam victoriam ea ipsa hora, qua obténta fuit, Deo revelante, cognóvit suisque familiaribus indicávit. Dum vero novam in ipsos Turcas expeditiónem molirétur, in gravem morbum incidit; et, acerbíssimis dolóribus patientíssime tolerátis, ad extrema deveniens, cum sacraménta de more suscepísset, ánimam Deo placidíssime réddidit, anno millesimo quingentésimo septuagesimo secundo, ætátis suæ sexagesimo octavo, cum sedísset annos sex, menses tres, dies viginti quatuor. Corpus ejus in basilica sanctæ Maríæ ad Præsépe summa fidelium veneratióne colitur, multis a Deo ejus intercessióne patrátis miraculis. Quibus rite probátis, a Cleménte undecimo, Pontifice máximo, Sanctórum número adscriptus est.\par
\switchcolumn
\EN
\noindent The Turkish Sultan Selim was bloated with many victories, and had got together an huge fleet in the Gulf of Lepanto, but Pius V. crushed him, (on the 7th of October, 1571,) not so much by force of arms as by dint of the prayers wherein he pleaded with God. At the hour that the victory was won, Pius knew it by the inward revelation of God, and stated the fact to his servants. He was busied with the preparations for a new expedition against the Turks, when he was laid down by grievous sickness. He bore most sharp sufferings with the gentlest patience, and when the end came, he received the Sacraments as is usual, and with great peace yielded his spirit to God, (on the 1st of May,) in the year of salvation 1572, and of his own age the 68th, having sat as Pope six years, three months, and twenty-four days. His body is buried in the Church of St. Mary, where the Manger from Bethlehem is, and is there held in great respect by the faithful, who have obtained from God by his prayers, many evident miracles. The said miracles having been proved by a judicial investigation Pope Clement XI enrolled his name among those of the Saints.\par
\switchcolumn*

\LA
\begin{Resp}\Rsign Iste est, qui ante Deum magnas virtútes operátus est, et omnis terra doctrína ejus repléta est: \Star{} Ipse intercédat pro peccátis ómnium populórum, allelúja.\end{Resp}
\begin{Resp}\Vsign Iste est, qui contémpsit vitam mundi, et pervénit ad cæléstia regna.\end{Resp}
\begin{Resp}\Rsign Ipse intercédat pro peccátis ómnium populórum, allelúja.\end{Resp}
\begin{Resp}\Vsign Glória Patri, et Fílio, \Star{} et Spirítui Sancto.\end{Resp}
\begin{Resp}\Rsign Ipse intercédat pro peccátis ómnium populórum, allelúja.\end{Resp}
\switchcolumn
\EN
\begin{Resp}\Rsign This is he which wrought great wonders before God, and the whole earth is full of his teaching \Star{} May he pray for all people, that their sins may be forgiven unto them, alleluia.\end{Resp}
\begin{Resp}\Vsign This is he which loved not his life in this world, and hath attained unto the kingdom of heaven.\end{Resp}
\begin{Resp}\Rsign May he pray for all people, that their sins may be forgiven unto them, alleluia.\end{Resp}
\begin{Resp}\Vsign Glory be to the Father, and to the Son, \Star{} and to the Holy Ghost.\end{Resp}
\begin{Resp}\Rsign May he pray for all people, that their sins may be forgiven unto them, alleluia.\end{Resp}
\switchcolumn*

\LA
\LessonHead{Lectio IX (pro commemoratione)}
\switchcolumn
\EN
\LessonHead{Lesson IX (when commemorated)}
\switchcolumn*

\LA
\LessonTitle{Commemoratio: S. Pii V Papæ et Confessoris}
\switchcolumn
\EN
\LessonTitle{Commemoration of: St. Pius V, Pope and Confessor}
\switchcolumn*

\LA
\noindent Pius, Boschi in Insúbria natus, cum quatuórdecim esset annórum, órdinem Prædicatórum ingréssus est. Sacerdótio auctus, sacras conciónes plúribus in locis cum ingénti animárum fructu hábuit. Inquisitóris offícium diu fórtiter ac laudabíliter sustínuit, et a Paulo quarto ad Nepesínum et Sutrínum episcopátum promótus, post biénnium inter patres cardináles adscríptus est. A Pio quarto translátus ad ecclésiam Montis Regális in Subalpínis, diœcésim vísitans, abúsus représsit; et Romam revérsus, in gravíssimis negótiis expediéndis fuit occupátus. Mórtuo autem Pio, præter ómnium exspectatiónem eléctus Póntifex, nihil in vitæ ratióne, excépto exterióri hábitu, immutávit. Selímum Turcárum tyránnum, ingénti comparáta classe, ad Echínadas ínsulas non tam armis quam fusis précibus devícit. Dum vero novam in ipsos Turcas expeditiónem paráret, piíssime óbiit in Dómino, anno millésimo quingentésimo septuagésimo secúndo, ætátis suæ sexagésimo octávo. Corpus ejus in basílica sanctæ Maríæ majóris summa fidélium veneratióne cólitur. Eum Clemens Papa undécimus catálogo Sanctórum adscrípsit.\par
\switchcolumn
\EN
\noindent Pius, born at Bosco in Lombardy, entered the Order of Preachers at the age of fourteen. When he had been ordained to the priesthood, he preached in many churches, and his sermons bore much spiritual fruit. For a long time he carried out the office of Inquisitor in a way that was both vigorous and praiseworthy. He was promoted by Paul IV to the bishopric of Nepi and Sutri, and two years later he was enrolled among the Cardinals of the Church. Pius IV transferred him to the See of Mondóvi in Piedmont, where he visited his whole diocese and remedied many abuses. Thereafter, he returned to Rome, where he was engaged in many important affairs. When Pius IV died, he was elected Pope, contrary to all expectation; and thereafter he changed nothing in his way of life except his outer garments. Not so much by arms as by the prayers he poured out, he overcame Selim, the Sultan of the Turks, who had gathered together a huge fleet at Lepanto. In the year 1572, while preparing a new expedition against the Turks, he died peacefully in the Lord, at the age of sixty-eight. His body is venerated by the faithful in the basilica of St. Mary Major. Clement XI enrolled him among the Saints.\par
\switchcolumn*

\LA
\TeDeum{Te Deum}
\switchcolumn
\EN
\TeDeum{Te Deum}
\switchcolumn*

\end{paracol}

\par\needspace{10\baselineskip}\addvspace{1.2em}{\centering\small\textsc{Die 6 Maii} · \textsc{6 May}\par}
\Office{S. Joannis Apostoli ante Portam Latinam}{St. John the Apostle Before the Latin Gate}{Duplex majus}
\addcontentsline{toc}{subsection}{Die 6 Maii · S. Joannis Apostoli ante Portam Latinam\enspace\textit{St. John the Apostle Before the Latin Gate}}
\begin{paracol}{2}
\LA
\LessonHead{Lectio I}
\switchcolumn
\EN
\LessonHead{Lesson I}
\switchcolumn*

\LA
\LessonTitle{Incipit Epístola prima beáti Joánnis Apóstoli}
\switchcolumn
\EN
\LessonTitle{Lesson from the first letter of St. John the Apostle}
\switchcolumn*

\LA
\Cite{1 Joannes 1:1-5}
\switchcolumn
\EN
\Cite{1 John 1:1-5}
\switchcolumn*

\LA
\noindent Quod fuit ab inítio, quod audívimus, quod vídimus óculis nostris, quod perspéximus, et manus nostræ contrectavérunt de verbo vitæ: et vita manifestáta est, et vídimus, et testámur, et annuntiámus vobis vitam ætérnam, quæ erat apud Patrem, et appáruit nobis: quod vídimus, et audívimus, annuntiámus vobis, ut et vos societátem habeátis nobíscum, et socíetas nostra sit cum Patre et cum Fílio ejus Jesu Christo. Et hæc scríbimus vobis ut gaudeátis, et gáudium vestrum sit plenum. Et hæc est annuntiátio, quam audívimus ab eo, et annuntiámus vobis: Quóniam Deus lux est, et ténebræ in eo non sunt ullæ.\par
\switchcolumn
\EN
\noindent That which was from the beginning, which we have heard, which we have seen with our eyes, which we have looked upon, and our hands have handled, of the word of life: For the life was manifested; and we have seen and do bear witness, and declare unto you the life eternal, which was with the Father, and hath appeared to us: That which we have seen and have heard, we declare unto you, that you also may have fellowship with us, and our fellowship may be with the Father, and with his Son Jesus Christ. And these things we write to you, that you may rejoice, and your joy may be full. And this is the declaration which we have heard from him, and declare unto you: That God is light, and in him there is no darkness.\par
\switchcolumn*

\LA
\begin{Resp}\Rsign Beátus vir qui métuit Dóminum, allelúja: \Star{} In mandátis ejus cupit nimis, allelúja, allelúja, allelúja.\end{Resp}
\begin{Resp}\Vsign Glória et divítiæ in domo ejus, et justítia ejus manet in sǽculum sǽculi.\end{Resp}
\begin{Resp}\Rsign In mandátis ejus cupit nimis, allelúja, allelúja, allelúja.\end{Resp}
\switchcolumn
\EN
\begin{Resp}\Rsign Blessed is the man that feareth the Lord, alleluia. \Star{} He shall delight exceedingly in his commandments, alleluia, alleluia, alleluia.\end{Resp}
\begin{Resp}\Vsign Glory and wealth shall be in his house: and his justice remaineth for ever and ever.\end{Resp}
\begin{Resp}\Rsign He shall delight exceedingly in his commandments, alleluia, alleluia, alleluia.\end{Resp}
\switchcolumn*

\LA
\LessonHead{Lectio II}
\switchcolumn
\EN
\LessonHead{Lesson II}
\switchcolumn*

\LA
\Cite{1 Joannes 1:6-10}
\switchcolumn
\EN
\Cite{1 John 1:6-10}
\switchcolumn*

\LA
\noindent Si dixérimus quóniam societátem habémus cum eo, et in ténebris ambulámus, mentímur, et veritátem non fácimus. Si autem in luce ambulámus sicut et ipse est in luce, societátem habémus ad ínvicem, et sanguis Jesu Christi, Fílii ejus, emúndat nos ab omni peccáto. Si dixérimus quóniam peccátum non habémus, ipsi nos sedúcimus, et véritas in nobis non est. Si confiteámur peccáta nostra, fidélis est et justus, ut remíttat nobis peccáta nostra, et emúndet nos ab omni iniquitáte. Si dixérimus quóniam non peccávimus, mendácem fácimus eum, et verbum ejus non est in nobis.\par
\switchcolumn
\EN
\noindent If we say that we have fellowship with him, and walk in darkness, we lie, and do not the truth. But if we walk in the light, as he also is in the light, we have fellowship one with another, and the blood of Jesus Christ his Son cleanseth us from all sin. If we say that we have no sin, we deceive ourselves, and the truth is not in us. If we confess our sins, he is faithful and just, to forgive us our sins, and to cleanse us from all iniquity. If we say that we have not sinned, we make him a liar, and his word is not in us.\par
\switchcolumn*

\LA
\begin{Resp}\Rsign Tristítia vestra, allelúja, \Star{} Convertétur in gáudium, allelúja, allelúja.\end{Resp}
\begin{Resp}\Vsign Mundus autem gaudébit, vos vero contristabímini, sed tristítia vestra.\end{Resp}
\begin{Resp}\Rsign Convertétur in gáudium, allelúja, allelúja.\end{Resp}
\switchcolumn
\EN
\begin{Resp}\Rsign Your sorrow, alleluia. \Star{} Shall be turned into joy, alleluia.\end{Resp}
\begin{Resp}\Vsign The world shall rejoice: and you shall be made sorrowful, but your sorrow.\end{Resp}
\begin{Resp}\Rsign Shall be turned into joy, alleluia.\end{Resp}
\switchcolumn*

\LA
\LessonHead{Lectio III}
\switchcolumn
\EN
\LessonHead{Lesson III}
\switchcolumn*

\LA
\Cite{1 Joannes 2:1-5}
\switchcolumn
\EN
\Cite{1 John 2:1-5}
\switchcolumn*

\LA
\noindent Filíoli mei, hæc scribo vobis, ut non peccétis. Sed et si quis peccáverit, advocátum habémus apud Patrem, Jesum Christum justum: et ipse est propitiátio pro peccátis nostris: non pro nostris autem tantum, sed étiam pro totíus mundi. Et in hoc scimus quóniam cognóvimus eum, si mandáta ejus observémus. Qui dicit se nosse eum, et mandáta ejus non custódit, mendax est, et in hoc véritas non est. Qui autem servat verbum ejus, vere in hoc cáritas Dei perfécta est.\par
\switchcolumn
\EN
\noindent My little children, these things I write to you, that you may not sin. But if any man sin, we have an advocate with the Father, Jesus Christ the just: And he is the propitiation for our sins: and not for ours only, but also for those of the whole world. And by this we know that we have known him, if we keep his commandments. He who saith that he knoweth him, and keepeth not his commandments, is a liar, and the truth is not in him. But he that keepeth his word, in him in very deed the charity of God is perfected.\par
\switchcolumn*

\LA
\begin{Resp}\Rsign Pretiósa in conspéctu Dómini, allelúja, \Star{} Mors Sanctórum ejus, allelúja.\end{Resp}
\begin{Resp}\Vsign Custódit Dóminus ómnia ossa eórum, unum ex his non conterétur.\end{Resp}
\begin{Resp}\Rsign Mors Sanctórum ejus, allelúja.\end{Resp}
\begin{Resp}\Vsign Glória Patri, et Fílio, \Star{} et Spirítui Sancto.\end{Resp}
\begin{Resp}\Rsign Mors Sanctórum ejus, allelúja.\end{Resp}
\switchcolumn
\EN
\begin{Resp}\Rsign Precious in the sight of the Lord, alleluia. \Star{} Is the death of his saints, alleluia.\end{Resp}
\begin{Resp}\Vsign The Lord is nigh unto them that are of a contrite heart: and he will save the humble of spirit.\end{Resp}
\begin{Resp}\Rsign Is the death of his saints, alleluia.\end{Resp}
\begin{Resp}\Vsign Glory be to the Father, and to the Son, \Star{} and to the Holy Ghost.\end{Resp}
\begin{Resp}\Rsign Is the death of his saints, alleluia.\end{Resp}
\switchcolumn*

\LA
\LessonHead{Lectio IV}
\switchcolumn
\EN
\LessonHead{Lesson IV}
\switchcolumn*

\LA
\LessonTitle{Ex libro sancti Hierónymi Presbýteri contra Jovinianum.}
\switchcolumn
\EN
\LessonTitle{From the Book against Jovinian written by St. Jerome, Priest at Bethlehem.}
\switchcolumn*

\LA
\Source{Liber 1, n. 26}
\switchcolumn
\EN
\Source{Bk. 1}
\switchcolumn*

\LA
\noindent Joánnes Apóstolus, unus ex discípulis Dómini, qui mínimus natu tráditur fuísse inter Apóstolos, et quem fides Christi vírginem repérerat, virgo permánsit; et ídeo plus amátur a Dómino, et recúmbit super pectus Jesu; et, quod Petrus, qui uxórem habúerat, interrogáre non audet, illum rogat, ut intérroget; et post resurrectiónem, nuntiánte María Magdaléne quod Dóminus resurrexísset, utérque cucúrrit ad sepúlcrum, sed ille prævénit; cumque essent in navi, et piscaréntur in lacu Genésareth, Jesus stabat in líttore, nec sciébant Apóstoli quem vidérent; solus virgo Vírginem agnóscit, et dicit Petro: Dóminus est.\par
\switchcolumn
\EN
\noindent The Apostle John was one of the disciples of the Lord. There is a tradition that he was the youngest of the Apostles. He was a virgin when the Faith of Christ found him, and he hath remained a virgin for ever. This is why he was “the disciple whom Jesus loved” more than any of the others, and why he “leaned on Jesus' Breast.” When Peter, who had been married, wished to ask the Lord who it was that was about to betray Him, he dared not ask for himself, but beckoned to John, that he should ask it. (John xiii. 23, 24.) After the Resurrection, when “Mary Magdalene came and told the disciples that the Lord was risen, Peter and John ran both together to the sepulchre, but John did outrun Peter.” (xx. 2-4.) Later on, when the Apostles were on the Sea of Galilee, in a ship, fishing, “Jesus stood on the shore, but the disciples knew not that it was Jesus,” till virgin knew Virgin, and “that disciple whom Jesus loved saith unto Peter It is the Lord.”\par
\switchcolumn*

\LA
\begin{Resp}\Rsign Lux perpétua lucébit Sanctis tuis, Dómine, \Star{} Et ætérnitas témporum, allelúja, allelúja.\end{Resp}
\begin{Resp}\Vsign Lætítia sempitérna erit super cápita eórum: gáudium et exsultatiónem obtinébunt.\end{Resp}
\begin{Resp}\Rsign Et ætérnitas témporum, allelúja, allelúja.\end{Resp}
\switchcolumn
\EN
\begin{Resp}\Rsign Eternal light will shine over your saints, O Lord, \Star{} And the eternity of the times, alleluia, alleluia.\end{Resp}
\begin{Resp}\Vsign Everlasting joy shall be upon their heads, they shall obtain joy and gladness\end{Resp}
\begin{Resp}\Rsign And the eternity of the times, alleluia, alleluia.\end{Resp}
\switchcolumn*

\LA
\LessonHead{Lectio V}
\switchcolumn
\EN
\LessonHead{Lesson V}
\switchcolumn*

\LA
\noindent Fuit autem Joánnes et Apóstolus, et Evangelísta, et Prophéta: Apóstolus, quia scripsit ad ecclésias ut magíster; Evangelísta, quia librum Evangélii cóndidit, quod, excépto Matthǽo, álii ex duódecim Apóstoli non fecérunt; Prophéta, vidit enim in Patmos ínsula, in qua fúerat a Domitiáno príncipe ob Dómini martýrium relegátus, Apocalýpsim, infiníta futurórum mystéria continéntem. Refert autem Tertulliánus, quod Romæ missus in fervéntis ólei dólium, púrior et vegétior exíverit, quam intráverit.\par
\switchcolumn
\EN
\noindent John was both an Apostle, and an Evangelist, and a Prophet. He was an Apostle, in that he wrote to the Churches, as their Teacher. He was an Evangelist, in that he wrote one of the Gospels, the like whereto was not done by any other of the twelve Apostles, save Matthew. He was a Prophet, in that when he was in the isle of Patmos, whither he had been banished by Domitian on account of his uplifting of his testimony for the Lord, he saw there that Apocalypse which containeth such unfathomable mysteries concerning “things which shall be hereafter.” (Apoc. i. 19.) Tertullian also saith that when he was at Rome, he was put into a vessel of boiling oil, but that he came out cleaner and healthier than he went in.\par
\switchcolumn*

\LA
\begin{Resp}\Rsign Virtúte magna reddébant Apóstoli \Star{} Testimónium resurrectiónis Jesu Christi Dómini nostri, allelúja, allelúja.\end{Resp}
\begin{Resp}\Vsign Repléti quidem Spíritu Sancto, loquebántur cum fidúcia verbum Dei.\end{Resp}
\begin{Resp}\Rsign Testimónium resurrectiónis Jesu Christi Dómini nostri, allelúja, allelúja.\end{Resp}
\switchcolumn
\EN
\begin{Resp}\Rsign With great power did the Apostles \Star{} Give testimony of the resurrection of Jesus Christ our Lord, alleluia, alleluia.\end{Resp}
\begin{Resp}\Vsign And they were all filled with the Holy Ghost: and they spoke the word of God with confidence.\end{Resp}
\begin{Resp}\Rsign Give testimony of the resurrection of Jesus Christ our Lord, alleluia, alleluia.\end{Resp}
\switchcolumn*

\LA
\LessonHead{Lectio VI}
\switchcolumn
\EN
\LessonHead{Lesson VI}
\switchcolumn*

\LA
\noindent Sed et ipsum ejus Evangélium multum distat a céteris. Matthǽus quasi de hómine íncipit scríbere: Liber generatiónis Jesu Christi, fílii David, fílii Abraham; Lucas a sacerdótio Zacharíæ; Marcus a prophetía Malachíæ et Isaíæ. Primus habet fáciem hóminis, propter genealogíam; secúndus fáciem vítuli, propter sacerdótium; tértius fáciem leónis, propter vocem clamántis in desérto, Paráte viam Dómini, rectas fácite sémitas ejus. Joánnes vero noster quasi áquila ad supérna volat, et ad ipsum Patrem pervénit, dicens: In princípio erat Verbum, et Verbum erat apud Deum, et Deus erat Verbum.\par
\switchcolumn
\EN
\noindent There is a great difference between his Gospel and the three others. Matthew beginneth to write as of a man “The Book of the Generation of Jesus Christ, the Son of David, the son of Abraham.” Luke's first words of history relate to the priesthood of Zacharias Mark commenceth with the prophecies of Malachi and Isaiah. The first hath the face of a man, with an human genealogy the second hath the face of a calf, being a victim offered by priests the third hath the face of a lion, even “the voice of one crying in the wilderness, Prepare ye the way of the Lord, make His paths straight” but the John of whom I write is like a flying eagle, (Apoc. iv. 7,) whose kingly flight beareth him up above earth-gathered clouds, an eagle that wingeth his way toward the Father Himself, and which crieth: “In the beginning was the Word, and the Word was with God, and the Word was God.”\par
\switchcolumn*

\LA
\begin{Resp}\Rsign Isti sunt agni novélli, qui annuntiavérunt, allelúja, modo venérunt ad fontes, \Star{} Repléti sunt claritáte, allelúja, allelúja.\end{Resp}
\begin{Resp}\Vsign In conspéctu Agni amícti sunt stolis albis, et palmæ in mánibus eórum.\end{Resp}
\begin{Resp}\Rsign Repléti sunt claritáte, allelúja, allelúja.\end{Resp}
\begin{Resp}\Vsign Glória Patri, et Fílio, \Star{} et Spirítui Sancto.\end{Resp}
\begin{Resp}\Rsign Repléti sunt claritáte, allelúja, allelúja.\end{Resp}
\switchcolumn
\EN
\begin{Resp}\Rsign These are the young lambs, who were promised, they are just coming to the fountains. \Star{} They are filled with glory, alleluia, alleluia.\end{Resp}
\begin{Resp}\Vsign In sight of the Lamb, clothed with white robes, and palms in their hands.\end{Resp}
\begin{Resp}\Rsign They are filled with glory, alleluia, alleluia.\end{Resp}
\begin{Resp}\Vsign Glory be to the Father, and to the Son, \Star{} and to the Holy Ghost.\end{Resp}
\begin{Resp}\Rsign They are filled with glory, alleluia, alleluia.\end{Resp}
\switchcolumn*

\LA
\LessonHead{Lectio VII}
\switchcolumn
\EN
\LessonHead{Lesson VII}
\switchcolumn*

\LA
\LessonTitle{Léctio sancti Evangélii secúndum Matthǽum.}
\switchcolumn
\EN
\LessonTitle{From the Holy Gospel according to Matthew}
\switchcolumn*

\LA
\Cite{Matt 20:20-23}
\switchcolumn
\EN
\Cite{Matt 20:20-23}
\switchcolumn*

\LA
\noindent In illo témpore: Accéssit ad Jesum mater filiórum Zebedǽi cum fíliis suis, adórans et petens áliquid ab eo. Et réliqua.\par
\switchcolumn
\EN
\noindent AT that time: Came to Jesus the mother of Zebedee's children, with her sons, worshipping Him, and desiring a certain thing of Him. And so on.\par
\switchcolumn*

\LA
\LessonTitle{Homilía sancti Hierónymi Presbýteri.}
\switchcolumn
\EN
\LessonTitle{Homily by St. Jerome, Priest at Bethlehem.}
\switchcolumn*

\LA
\Source{Liber 3 Comm. in Matth. c. 20.}
\switchcolumn
\EN
\Source{Comment, on Matt. xx. Bk. 3}
\switchcolumn*

\LA
\noindent Unde opiniónem regni habet mater filiórum Zebedǽi, ut cum Dóminus díxerit: Fílius hóminis tradétur princípibus sacerdótum et scribis, et condemnábunt eum morte, et tradent Géntibus ad illudéndum et flagellándum et crucifigéndum; et ignomíniam passiónis timéntibus discípulis nuntiáret, illa glóriam póstulet triumphántis? Hac, ut reor, ex causa, quia post ómnia díxerat Dóminus, Et tértia die resúrget; putávit eum múlier post resurrectiónem íllico regnatúrum, et hoc, quod in secúndo advéntu promíttitur, in primo esse compléndum; et, aviditáte femínea præséntia cupit, ímmemor futurórum.\par
\switchcolumn
\EN
\noindent Whence had the mother of Zebedee's children gotten her idea of the Lord's kingdom He had but just said “The Son of man shall be betrayed unto the chief priests and unto the scribes, and they shall condemn Him to death, and shall deliver Him to the Gentiles to mock, (and to scourge,) and to crucify Him.” He had told His trembling disciples of the outrages that awaited Him in His Passion and yet that mother came to Him to ask for her sons a share in the glory of His Triumph. I think it was because the Lord, after He had said all the rest, had said also “And the third day He shall rise again.” The woman supposed that after His resurrection His kingdom would immediately be established, and that that would be fulfilled at His first coming which is promised at His second. And so, with womanly haste, she forgetteth the future, and catcheth at the present.\par
\switchcolumn*

\LA
\begin{Resp}\Rsign Ego sum vitis vera, et vos pálmites: \Star{} Qui manet in me, et ego in eo, hic fert fructum multum, allelúja, allelúja.\end{Resp}
\begin{Resp}\Vsign Sicut diléxit me Pater, et ego diléxi vos.\end{Resp}
\begin{Resp}\Rsign Qui manet in me, et ego in eo, hic fert fructum multum, allelúja, allelúja.\end{Resp}
\switchcolumn
\EN
\begin{Resp}\Rsign I am the vine: you the branches. \Star{} He that abideth in me, and I in him, the same beareth much fruit, alleluia, alleluia.\end{Resp}
\begin{Resp}\Vsign As the Father hath loved me, I also have loved you.\end{Resp}
\begin{Resp}\Rsign He that abideth in me, and I in him, the same beareth much fruit, alleluia, alleluia.\end{Resp}
\switchcolumn*

\LA
\LessonHead{Lectio VIII}
\switchcolumn
\EN
\LessonHead{Lesson VIII}
\switchcolumn*

\LA
\noindent Mater póstulat, et Dóminus discípulis lóquitur, intélligens preces ejus ex filiórum descéndere voluntáte. Potéstis bíbere cálicem, quem ego bibitúrus sum? Cálicem in Scriptúris divínis passiónem intellígimus, juxta illud: Pater, si possíbile est, tránseat a me calix iste; et in Psalmo: Quid retríbuam Dómino pro ómnibus quæ retríbuit mihi? Cálicem salutáris accípiam, et nomen Dómini invocábo. Statímque infert quis iste sit calix: Pretiósa in conspéctu Dómini mors Sanctórum ejus.\par
\switchcolumn
\EN
\noindent It was the mother who asked, but the Lord addressed His answer to the disciples, understanding that she had made her prayer in obedience to their wishes. “Are ye able to drink of the cup that I shall drink of?” From God's written Word we gather that by this cup, He meant the Passion, touching the which we read that He said “O My Father, if it be possible, let this cup pass from Me!” Likewise is it written in the hundred-and-fifteenth Psalm “I will take the cup of salvation, and call upon the Name of the Lord,” and what that life-giving cup was, the words which soon follow tell us “Precious in the sight of the Lord is the death of His Saints.”\par
\switchcolumn*

\LA
\begin{Resp}\Rsign Cándidi facti sunt Nazarǽi ejus, allelúja: splendórem Deo dedérunt, allelúja: \Star{} Et sicut lac coaguláti sunt, allelúja, allelúja.\end{Resp}
\begin{Resp}\Vsign Candidióres nive, nitidióres lacte, rubicundióres ébore antíquo, sapphíro pulchrióres.\end{Resp}
\begin{Resp}\Rsign Et sicut lac coaguláti sunt, allelúja, allelúja.\end{Resp}
\begin{Resp}\Vsign Glória Patri, et Fílio, \Star{} et Spirítui Sancto.\end{Resp}
\begin{Resp}\Rsign Et sicut lac coaguláti sunt, allelúja, allelúja.\end{Resp}
\switchcolumn
\EN
\begin{Resp}\Rsign Her Nazarites are become pure, Alleluia: they reflect the glory of God, Alleluia. \Star{} They are whiter than milk. Alleluia, Alleluia.\end{Resp}
\begin{Resp}\Vsign They are purer than snow, they are whiter than milk, they are more ruddy in body than coral, their polishing is of sapphire.\end{Resp}
\begin{Resp}\Rsign They are whiter than milk. Alleluia, Alleluia.\end{Resp}
\begin{Resp}\Vsign Glory be to the Father, and to the Son, \Star{} and to the Holy Ghost.\end{Resp}
\begin{Resp}\Rsign They are whiter than milk. Alleluia, Alleluia.\end{Resp}
\switchcolumn*

\LA
\LessonHead{Lectio IX}
\switchcolumn
\EN
\LessonHead{Lesson IX}
\switchcolumn*

\LA
\noindent Quǽritur quómodo cálicem martýrii fílii Zebedǽi, Jacóbus vidélicet et Joánnes, bíberint; cum Scriptúra narret Jacóbum tantum Apóstolum ab Heróde cápite truncátum, Joánnes autem própria morte vitam finíerit. Sed, si legámus ecclesiásticas histórias, in quibus fertur, quod et ipse propter martýrium sit missus in fervéntis ólei dólium, et inde ad suscipiéndam corónam Christi athléta procésserit, statímque relegátus in Patmos ínsulam sit; vidébimus martýrio ánimum non defuísse, et bibísse Joánnem cálicem confessiónis, quem et tres púeri in camíno ignis bibérunt, licet persecútor non fúderit sánguinem.\par
\switchcolumn
\EN
\noindent The question ariseth, how the two sons of Zebedee, James namely and John, drank of the cup of contention even unto blood against sin, seeing that though we know by the Scriptures that “Herod the king killed James the brother of John with the sword,” yet John ended his earthly life by a natural death. But if we read the Records of the Church, we shall find there told how that John, on account of his testifying to the truth, was cast into a vessel of boiling oil, and although the holy champion came out unhurt and continued his pilgrimage here for a while longer, before he received his crown from Christ's hand, being straightway banished into the isle of Patmos, yet we see that he had the soul of a martyr, and drank the same cup of martyrdom that was drunk by the three children in the burning fiery furnace, albeit the persecutor did not actually shed his blood.\par
\switchcolumn*

\LA
\TeDeum{Te Deum}
\switchcolumn
\EN
\TeDeum{Te Deum}
\switchcolumn*

\LA
\LessonHead{Lectio IX (pro commemoratione)}
\switchcolumn
\EN
\LessonHead{Lesson IX (when commemorated)}
\switchcolumn*

\LA
\LessonTitle{Commemoratio: S. Joannis Apostoli ante Portam Latinam}
\switchcolumn
\EN
\LessonTitle{Commemoration of: St. John the Apostle Before the Latin Gate}
\switchcolumn*

\LA
\noindent Ex libro sancti Hierónymi Presbýteri contra Joviniánum\par
\switchcolumn
\EN
\noindent The Lesson is taken from the Book of St. Jerome the Priest against Jovinian\par
\switchcolumn*

\LA
Joánnes Apóstolus, unus ex discípulis Dómini, qui mínimus natu tráditur fuísse inter Apóstolos, et quem fides Christi vírginem repérerat, virgo permánsit; et ídeo plus amátur a Dómino, et recúmbit super pectus Jesu; et, quod Petrus, qui uxórem habúerat, interrogáre non audet, illum rogat, ut intérroget; et post resurrectiónem, nuntiánte María Magdaléne quod Dóminus resurrexísset, utérque cucúrrit ad sepúlcrum, sed ille prævénit; cumque essent in navi, et piscaréntur in lacu Genésareth, Jesus stabat in líttore, nec sciébant Apóstoli quem vidérent; solus virgo Vírginem agnóscit, et dicit Petro: Dóminus est. Fuit autem Joánnes et Apóstolus, et Evangelísta, et Prophéta. Refert autem Tertulliánus, quod Romæ missus in fervéntis ólei dólium, púrior et vegétior exíverit, quam intráverit.\par
\switchcolumn
\EN
The Apostle John was one of the first disciples of the Lord, and there is a tradition that he was the youngest of the Apostles. He was a virgin when the Faith of Christ found him, and he hath remained a virgin for ever. This is why he was the disciple whom Jesus loved more than any of the others, and why he leaned on Jesus' breast. When Peter, who had been married, wished to ask the Lord who it was that was about to betray him, he dared not ask for himself, but beckoned to John, that he should ask it. After the resurrection, when Mary Magdalene came and told the disciples that the Lord was risen, Peter and John ran both together to the sepulchre, but John did outrun Peter. Later on, when the Apostles were on the Sea of Galilee, in a ship, fishing, Jesus stood on the shore, but the disciples knew not that it was Jesus, till virgin knew Virgin, and that disciple whom Jesus loved saith unto Peter: It is the Lord. John was both an Apostle and an Evangelist, and also a Prophet. Also Tertullian saith that when he was at Rome, he was put into a vessel of boiling oil, but that he came out cleaner and healthier than he went in.\par
\switchcolumn*

\LA
\TeDeum{Te Deum}
\switchcolumn
\EN
\TeDeum{Te Deum}
\switchcolumn*

\end{paracol}

\par\needspace{10\baselineskip}\addvspace{1.2em}{\centering\small\textsc{Die 7 Maii} · \textsc{7 May}\par}
\Office{S. Stanislai Episcopi et Martyris}{St. Stanislaus, Bishop and Martyr}{Duplex}
\addcontentsline{toc}{subsection}{Die 7 Maii · S. Stanislai Episcopi et Martyris\enspace\textit{St. Stanislaus, Bishop and Martyr}}
\begin{paracol}{2}
\LA
\RefLine{Lectiones I–III: de Scriptura occurrente.}
\switchcolumn
\EN
\RefLine{Lessons I–III: from the Scripture of the day.}
\switchcolumn*

\LA
\RefLine{Lectiones VII–IX: ex Commúni unius Martyris Tempore Paschali (C2p).}
\switchcolumn
\EN
\RefLine{Lessons VII–IX: from the Common of a Single Martyr in Paschaltide (C2p).}
\switchcolumn*

\LA
\LessonHead{Lectio IV}
\switchcolumn
\EN
\LessonHead{Lesson IV}
\switchcolumn*

\LA
\noindent Stanislaus Polonus, apud Cracoviam nobili genere natus et piis paréntibus, qui antea per annos trigínta stériles, illum a Deo precibus impetrarunt, ab ineunte ætate futuræ sanctitátis specimen dedit. Adoléscens bonis artibus operam navávit, multumque in sacra cánonum et theologíæ doctrina profecit. Paréntibus mórtuis, amplum patrimónium paupéribus distríbuit, vitæ monasticæ desidério. Sed Dei providéntia canonicus Cracoviénsis et concionátor factus a Lampérto epíscopo, in ejus póstea locum, quamvis invitus, suffícitur. Quo in munere, ómnium pastorálium virtútum laude, et præcípue misericórdia in páuperes, enítuit.\par
\switchcolumn
\EN
\noindent This Stanislaw was a Pole. He was born of a noble family, (on the 26th day of July, in the year of our Lord 1030,) at (Szcrepanow, in the diocese of) Cracow. His godly parents, who had been childless for thirty years, obtained him from God by prayer, and from his earliest years he gave token of the holiness of life which afterwards marked him. When he was a young man he applied himself heartily to all useful learning, and was deeply read in the sacred teaching of the Canons and of Theological science. After the death of his parents he inherited great possessions, but he sold them, and distributed the price to the poor, purposing himself to become a monk. However, by the Providence of God, Lampert, Bishop of Cracow, named him Canon of the Cathedral Church of that diocese, and Preacher in the same and afterwards, (in 1072,) he was elected, against his own will, to succeed to Lampert's place. In this office he was a bright and shining light of all virtues that become a shepherd of souls, especially of tenderness toward the poor.\par
\switchcolumn*

\LA
\begin{Resp}\Rsign Lux perpétua lucébit Sanctis tuis, Dómine, \Star{} Et ætérnitas témporum, allelúja, allelúja.\end{Resp}
\begin{Resp}\Vsign Lætítia sempitérna erit super cápita eórum: gáudium et exsultatiónem obtinébunt.\end{Resp}
\begin{Resp}\Rsign Et ætérnitas témporum, allelúja, allelúja.\end{Resp}
\switchcolumn
\EN
\begin{Resp}\Rsign Eternal light will shine over your saints, O Lord, \Star{} And the eternity of the times, alleluia, alleluia.\end{Resp}
\begin{Resp}\Vsign Everlasting joy shall be upon their heads, they shall obtain joy and gladness\end{Resp}
\begin{Resp}\Rsign And the eternity of the times, alleluia, alleluia.\end{Resp}
\switchcolumn*

\LA
\LessonHead{Lectio V}
\switchcolumn
\EN
\LessonHead{Lesson V}
\switchcolumn*

\LA
\noindent Erat tum Poloniæ rex Bolesláus, quem gráviter offendit, quod illíus notam libidinem publice arguebat. Quare in solemni regni convéntu Stanisláum per calumniam in judícium coram se vocari curat, tamquam pagum occuparet, quem ecclésiæ suæ nómine coémerat. Quod cum neque tabulis probare posset, et testes veritátem cidere timerent, spondet epíscopus, se Petrum pagi venditórem, qui triennio ante obíerat, intra dies tres in judícium adducturum. Conditióne cum risu accepta, vir Dei toto triduo jejuniis et oratióni incumbit; ipso sponsiónis die, post oblátum Missæ sacrifícium, Petrum e sepúlcro surgere jubet; qui statim redivivus, episcopum ad regium tribunal euntem sequitur, ibique, rege et ceteris stupore attónitis, de agro a se véndito et pretio rite sibi ab episcopo persoluto testimónium dicit, atque íterum in Dómino obdormívit.\par
\switchcolumn
\EN
\noindent AT that time Boleslaw II. was King of Poland, and him Stanislaw grievously offended, because he openly rebuked him for his shameless lust. Wherefore, in a solemn Parliament of his kingdom, he made Stanislaw to be brought before him on a false accusation of having taken wrongfully a certain village, which he had bought in the name of his Church. The Bishop could not rebut this charge by documents, and the witnesses were in too great fear to speak the truth. Stanislaw therefore said that in three days he would produce before the judgment-seat one Peter, from whom he had bought the village, and who had been dead three years. His enemies laughed thereat, and closed with his proposal, and the man of God gave himself up to fasting and prayer for three days. On the day which he had promised, after he had offered up the Holy Sacrifice of the Mass, he commanded Peter to rise from the grave. Peter then immediately came to life, arose, and followed Stanislaw to the King's judgment-seat, where before the King and all others, who were struck dumb with amazement, he bore witness of the sale of the village, and the honest payment of the price by the Bishop, and then again fell asleep in the Lord.\par
\switchcolumn*

\LA
\begin{Resp}\Rsign In servis suis, allelúja, \Star{} Consolábitur Deus, allelúja.\end{Resp}
\begin{Resp}\Vsign Judicábit Dóminus pópulum suum, et in servis suis.\end{Resp}
\begin{Resp}\Rsign Consolábitur Deus, allelúja.\end{Resp}
\switchcolumn
\EN
\begin{Resp}\Rsign His servants, alleluia \Star{} God will comfort, alleluia, alleluia.\end{Resp}
\begin{Resp}\Vsign The Lord will judge His people and, His servants,\end{Resp}
\begin{Resp}\Rsign God will comfort, alleluia, alleluia.\end{Resp}
\switchcolumn*

\LA
\LessonHead{Lectio VI}
\switchcolumn
\EN
\LessonHead{Lesson VI}
\switchcolumn*

\LA
\noindent At Boleslaum, frustra sæpe admónitum, Stanisláus tandem a fidelium communióne rémovet. Ille, iracúndia furens, milites in ecclésiam immíttit, ut sanctum episcopum confodiant; qui, ter conati, occúlta vi, tertio divinitus sunt depulsi. Postremo ímpius rex sacerdotem Dei, Hostiam immaculátam ad altáre offeréntem, sua manu obtruncat. Corpus, membratim concisum et per agros projectum, áquilæ a feris mirabíliter defendunt. Mox canonici Cracovienses sparsa membra, nocturni de cælo splendoris indicio cólligunt, et suis locis apte disponunt; quæ súbito ita inter se copuláta sunt, ut nulla vúlnerum vestígia exstarent. Multis præterea miraculis servi sui sanctitátem Deus declarávit post ejus mortem: quibus permotus Innocentius quartus, summus Pontifex, illum in Sanctórum númerum rétulit. Clemens vero octavus, Pontifex maximus, sancti Stanislai festo die in Románum Breviárium reláto, gloriósi Mártyris memóriam dúplici Officio ubíque celebrári jussit.\par
\switchcolumn
\EN
\noindent Stanislaw often rebuked Boleslaw, but when he found it was in vain, he at last cut him off from the communion of Christ's faithful people. Thereupon Boleslaw became frenzied with rage, and (on the 8th of May, in the year 1079,) sent soldiers to the Church to murder the holy Bishop. This they thrice essayed to do, but God was pleased that they should be held back by some unseen power. In the end, the ungodly King with his own hand cut off the head of the Priest of God as he was standing at the Altar offering up the Sacrifice without spot. His body was hewn into pieces and strewn about the fields, but the eagles strangely kept the beasts of prey off it. The Canons of the Cathedral of Cracow soon gathered together the mutilated and scattered limbs, which they were enabled to see by a lightness which overspread the sky at night and they fitted them together, each into his place. The relics immediately so joined themselves one to the other, that no marks of wounds remained. Moreover, God was pleased to manifest the holiness of His servant by many wonders after his death, by the which being moved, Pope Innocent IV. added his name to those of the Saints, and the Supreme Pontiff Clement VIII. gave his Feast a place in the Service Book of the Church of Rome, commanding that the memory of so glorious a Martyr should be everywhere celebrated under the Double rite.\par
\switchcolumn*

\LA
\begin{Resp}\Rsign Fíliæ Jerúsalem, veníte et vidéte Mártyres cum corónis, quibus coronávit eos Dóminus \Star{} In die solemnitátis et lætítiæ, allelúja.\end{Resp}
\begin{Resp}\Vsign Quóniam confortávit seras portárum tuárum, benedíxit fílios tuos in te.\end{Resp}
\begin{Resp}\Rsign In die solemnitátis et lætítiæ, allelúja.\end{Resp}
\begin{Resp}\Vsign Glória Patri, et Fílio, \Star{} et Spirítui Sancto.\end{Resp}
\begin{Resp}\Rsign In die solemnitátis et lætítiæ, allelúja.\end{Resp}
\switchcolumn
\EN
\begin{Resp}\Rsign Come forth, O ye daughters of Jerusalem, and behold the Martyrs with the crowns wherewith the Lord crowned them; \Star{} In the day of His feasting, and of His gladness. Alleluia.\end{Resp}
\begin{Resp}\Vsign For He hath strengthened the bars of thy gates He hath blessed thy children within thee.\end{Resp}
\begin{Resp}\Rsign In the day of His feasting, and of His gladness. Alleluia.\end{Resp}
\begin{Resp}\Vsign Glory be to the Father, and to the Son, \Star{} and to the Holy Ghost.\end{Resp}
\begin{Resp}\Rsign In the day of His feasting, and of His gladness. Alleluia.\end{Resp}
\switchcolumn*

\LA
\LessonHead{Lectio IX (pro commemoratione)}
\switchcolumn
\EN
\LessonHead{Lesson IX (when commemorated)}
\switchcolumn*

\LA
\LessonTitle{Commemoratio: S. Stanislai Episcopi et Martyris}
\switchcolumn
\EN
\LessonTitle{Commemoration of: St. Stanislaus, Bishop and Martyr}
\switchcolumn*

\LA
\noindent Stanisláus, apud Cracóviam nóbili génere natus, quem pii paréntes, per trigínta annos stériles, a Deo précibus impetrárunt, ab ineúnte ætáte futúræ sanctitátis spécimen dedit. Adoléscens in sacra cánonum et theologíæ doctrína multum profécit. Paréntibus mórtuis, amplum patrimónium paupéribus distríbuit, vitæ monásticæ desidério. Sed, Deo áliter disponénte, canónicus Cracoviénsis et concionátor factus a Lampérto epíscopo, in ejus póstea locum, quamvis invítus, sufféctus est. Quo in múnere ómnium pastorálium virtútum laude, et præcípue misericórdia in páuperes enítuit. Bolesláum Polóniæ regem, sǽpius ob mores corrúptos frustra admónitum, a fidélium communióne remóvit. Qui idcírco iracúndia furens, mílites in ecclésiam immíttit, ut sanctum epíscopum confódiant; sed cum divínitus fuíssent depúlsi, ímpius rex sacerdótem Dei, Hóstiam immaculátam ad altáre offeréntem, sua manu obtrúncat. Multis miráculis Servi sui sanctitátem Deus declarávit post ejus mortem, quibus permótus Innocéntius quartus illum in Sanctórum númerum rétulit.\par
\switchcolumn
\EN
\noindent Stanislas was born of a noble family at Krakow. His birth was God's answer to the devout prayers of his parents, who in thirty years of marriage had remained childless. From his earliest days, he gave tokens of future holiness, and, as a young man, made great progress in the study of canon law and theology. When his parents died, he distributed his large inheritance to the poor, intending to embrace the monastic life. But God arranged otherwise, and he was made Canon and Preacher in Krakow by Bishop Lampert. Against his wishes, he later became Bishop Lampert's successor, and in this office he was a shining example of every pastoral virtue, especially of kindness to the poor. He excommunicated Boleslas, King of Poland, whom he had often rebuked in vain for his corrupt way of life. In a violent rage the King sent soldiers to the church to kill the holy Bishop. But when these were driven off by the power of God, the King with his own hand beheaded Bishop Stanislas as he was offering the immaculate Victim on the altar. God proclaimed the holiness of his servant by many miracles after his death; and on the strength of them Innocent IV enrolled him among the Saints.\par
\switchcolumn*

\LA
\TeDeum{Te Deum}
\switchcolumn
\EN
\TeDeum{Te Deum}
\switchcolumn*

\end{paracol}

\par\needspace{10\baselineskip}\addvspace{1.2em}{\centering\small\textsc{Die 8 Maii} · \textsc{8 May}\par}
\Office{In Apparitione S. Michaëlis Archangeli}{Apparition of St. Michael the Archangel}{Duplex majus}
\addcontentsline{toc}{subsection}{Die 8 Maii · In Apparitione S. Michaëlis Archangeli\enspace\textit{Apparition of St. Michael the Archangel}}
\begin{paracol}{2}
\LA
\LessonHead{Lectio I}
\switchcolumn
\EN
\LessonHead{Lesson I}
\switchcolumn*

\LA
\LessonTitle{De Daniéle Prophéta}
\switchcolumn
\EN
\LessonTitle{Lesson from the book of Daniel}
\switchcolumn*

\LA
\Cite{Dan 7:9-11}
\switchcolumn
\EN
\Cite{Dan 7:9-11}
\switchcolumn*

\LA
\noindent Aspiciébam donec throni pósiti sunt, et antíquus diérum sedit. Vestiméntum ejus cándidum quasi nix, et capílli cápitis ejus quasi lana munda, thronus ejus flammæ ignis, rotæ ejus ignis accénsus. Flúvius ígneus rapidúsque egrediebátur a fácie ejus; míllia míllium ministrábant ei, et décies míllies centéna míllia assistébant ei. Judícium sedit, et libri apérti sunt. Aspiciébam propter vocem sermónum grándium, quos cornu illud loquebátur; et vidi quóniam interfécta esset béstia, et perísset corpus ejus, et tráditum esset ad comburéndum igni.\par
\switchcolumn
\EN
\noindent I beheld till thrones were placed, and the Ancient of days sat: his garment was white as snow, and the hair of his head like clean wool: his throne like flames of fire: the wheels of it like a burning fire. A swift stream of fire issued forth from before him: thousands of thousands ministered to him, and ten thousand times a hundred thousand stood before him: the judgment sat, and the books were opened. I beheld because of the voice of the great words which that horn spoke: and I saw that the beast was slain, and the body thereof was destroyed, and given to the fire to be burnt:\par
\switchcolumn*

\LA
\begin{Resp}\Rsign Factum est siléntium in cælo, dum commítteret bellum draco cum Michaéle Archángelo: \Star{} Audíta est vox míllia míllium dicéntium: Salus, honor et virtus omnipoténti Deo. (Allelúja.)\end{Resp}
\begin{Resp}\Vsign Míllia míllium ministrábant ei, et décies centéna míllia assistébant ei.\end{Resp}
\begin{Resp}\Rsign Audíta est vox míllia míllium dicéntium: Salus, honor et virtus omnipoténti Deo. (Allelúja.)\end{Resp}
\switchcolumn
\EN
\begin{Resp}\Rsign There was silence in heaven while the dragon fought against Michael the Archangel. \Star{} I heard the voice of thousands of thousands, saying: Salvation, and honour, and power unto God the Almighty. Alleluia.\end{Resp}
\begin{Resp}\Vsign Thousands of thousands ministered unto Him, and ten thousand times hundreds of thousands stood before Him.\end{Resp}
\begin{Resp}\Rsign I heard the voice of thousands of thousands, saying: Salvation, and honour, and power unto God, the Almighty. Alleluia.\end{Resp}
\switchcolumn*

\LA
\LessonHead{Lectio II}
\switchcolumn
\EN
\LessonHead{Lesson II}
\switchcolumn*

\LA
\Cite{Dan 10:4-8}
\switchcolumn
\EN
\Cite{Dan 10:4-8}
\switchcolumn*

\LA
\noindent Die autem vigésima et quarta mensis primi, eram juxta flúvium magnum, qui est Tigris. Et levávi óculos meos, et vidi: et ecce vir unus vestítus líneis, et renes ejus accíncti auro obrýzo; et corpus ejus quasi chrysólithus, et fácies ejus velut spécies fúlguris, et óculi ejus ut lampas ardens, et brácchia ejus, et quæ deórsum sunt usque ad pedes, quasi spécies æris candéntis; et vox sermónum ejus ut vox multitúdinis. Vidi autem ego Dániel solus visiónem; porro viri qui erant mecum non vidérunt; sed terror nímius írruit super eos, et fugérunt in abscónditum. Ego autem, relíctus solus, vidi visiónem grandem hanc, et non remánsit in me fortitúdo, sed et spécies mea immutáta est in me, et emárcui nec hábui quidquam vírium.\par
\switchcolumn
\EN
\noindent And in the four and twentieth day of the first month I was by the great river which is the Tigris. And I lifted up my eyes, and I saw: and behold a man clothed in linen, and his loins were girded with the finest gold: And his body was like the chrysolite, and his face as the appearance of lightning, and his eyes as a burning lamp: and his arms, and all downward even to the feet, like in appearance to glittering brass: and the voice of his word like the voice of a multitude. And I Daniel alone saw the vision: for the men that were with me saw it not: but an exceeding great terror fell upon them, and they fled away, and hid themselves. And I being left alone saw this great vision: and there remained no strength in me, and the appearance of my countenance was changed in me, and I fainted away, and retained no strength.\par
\switchcolumn*

\LA
\begin{Resp}\Rsign Stetit Angelus juxta aram templi, habens thuríbulum áureum in manu sua, et data sunt ei incénsa multa: \Star{} Et ascéndit fumus arómatum de manu Angeli in conspéctu Dómini. (Allelúja.)\end{Resp}
\begin{Resp}\Vsign In conspéctu Angelórum psallam tibi: adorábo ad templum sanctum tuum, et confitébor nómini tuo, Dómine.\end{Resp}
\begin{Resp}\Rsign Et ascéndit fumus arómatum de manu Angeli in conspéctu Dómini. (Allelúja.)\end{Resp}
\switchcolumn
\EN
\begin{Resp}\Rsign An Angel stood at the Altar of the temple, having a golden censer in his hand and there was given unto him much incense; \Star{} And the smoke of the incense ascended up before the Lord, out of the Angel's hand. Alleluia.\end{Resp}
\begin{Resp}\Vsign Before the Angels will I sing praise unto thee; I will worship toward thy holy temple, and praise thy Name, O Lord.\end{Resp}
\begin{Resp}\Rsign And the smoke of the incense ascended up before the Lord, out of the Angel's hand. Alleluia.\end{Resp}
\switchcolumn*

\LA
\LessonHead{Lectio III}
\switchcolumn
\EN
\LessonHead{Lesson III}
\switchcolumn*

\LA
\Cite{Dan 10:9-14}
\switchcolumn
\EN
\Cite{Dan 10:9-14}
\switchcolumn*

\LA
\noindent Et audívi vocem sermónum ejus: et áudiens jacébam consternátus super fáciem meam, et vultus meus hærébat terræ. Et ecce manus tétigit me, et eréxit me super génua mea et super artículos mánuum meárum. Et dixit ad me: Dániel, vir desideriórum, intéllege verba quæ ego loquor ad te, et sta in gradu tuo; nunc enim sum missus ad te. Cumque dixísset mihi sermónem istum, steti tremens. Et ait ad me: Noli metúere, Dániel; quia, ex die primo quo posuísti cor tuum ad intellegéndum, ut te afflígeres in conspéctu Dei tui, exaudíta sunt verba tua, et ego veni propter sermónes tuos. Princeps autem regni Persárum réstitit mihi vigínti et uno diébus; et ecce Míchaël, unus de princípibus primis, venit in adjutórium meum, et ego remánsi ibi juxta regem Persárum. Veni autem ut docérem te quæ ventúra sunt pópulo tuo in novíssimis diébus, quóniam adhuc vísio in dies.\par
\switchcolumn
\EN
\noindent And I heard the voice of his words: and when I heard, I lay in a consternation, upon my face, and my face was close to the ground. And behold a hand touched me, and lifted me up upon my knees, and upon the joints of my hands. And he said to me: Daniel, thou man of desires, understand the words that I speak to thee, and stand upright: for I am sent now to thee. And when he had said this word to me, I stood trembling. And he said to me: Fear not, Daniel: for from the first day that thou didst set thy heart to understand, to afflict thyself in the sight of thy God, thy words have been heard: and I am come for thy words. But the prince of the kingdom of the Persians resisted me one and twenty days: and behold Michael, one of the chief princes, came to help me, and I remained there by the king of the Persians. But I am come to teach thee what things shall befall thy people in the latter days, for as yet the vision is for days.\par
\switchcolumn*

\LA
\begin{Resp}\Rsign In conspéctu Angelórum psallam tibi, et adorábo ad templum sanctum tuum: \Star{} Et confitébor nómini tuo, Dómine. (Allelúja.)\end{Resp}
\begin{Resp}\Vsign Super misericórdia tua et veritáte tua: quóniam magnificásti super nos nomen sanctum tuum.\end{Resp}
\begin{Resp}\Rsign Et confitébor nómini tuo, Dómine. (Allelúja.)\end{Resp}
\begin{Resp}\Vsign Glória Patri, et Fílio, \Star{} et Spirítui Sancto.\end{Resp}
\begin{Resp}\Rsign Et confitébor nómini tuo, Dómine. (Allelúja.)\end{Resp}
\switchcolumn
\EN
\begin{Resp}\Rsign Before the Angels will I sing praise unto thee, and will worship toward thy holy temple \Star{} And I will praise thy Name, O Lord. Alleluia.\end{Resp}
\begin{Resp}\Vsign For thy loving kindness, and for thy truth; for Thou hast glorified thine holy Name on us.\end{Resp}
\begin{Resp}\Rsign And I will praise thy Name, O Lord. Alleluia.\end{Resp}
\begin{Resp}\Vsign Glory be to the Father, and to the Son, \Star{} and to the Holy Ghost.\end{Resp}
\begin{Resp}\Rsign And I will praise thy Name, O Lord. Alleluia.\end{Resp}
\switchcolumn*

\LA
\LessonHead{Lectio IV}
\switchcolumn
\EN
\LessonHead{Lesson IV}
\switchcolumn*

\LA
\noindent Beátum Michaélem Archángelum sǽpius homínibus apparuísse, et sacrórum Librórum auctoritáte et véteri Sanctórum traditióne comprobátur. Quam ob rem multis in locis facti memória celebrátur. Eum, ut olim synagóga Judæórum, sic nunc custódem et patrónum Dei venerátur Ecclésia. Gelásio autem primo, Pontífice máximo, in Apúlia in vértice Gargáni montis, ad cujus radíces íncolunt Sipontíni, Archángeli Michaélis fuit illústris apparítio.\par
\switchcolumn
\EN
\noindent That the blessed Archangel Michael hath oftentimes been seen of men is attested on the authority of the Holy Bible, and also by the ancient traditions of the Saints. For this reason such visions are held in remembrance in many places. As of old time did the Synagogue of the Jews, so now doth the Church of God venerate Michael as her watcher and defender. But during the papacy of Gelasius I. the summit of Mount Gargano in Apulia, at whose foot lieth the town of Siponto, was the scene of an extraordinary appearance of this same Archangel Michael.\par
\switchcolumn*

\LA
\begin{Resp}\Rsign Hic est Míchaël Archángelus, princeps milítiæ Angelórum, \Star{} Cujus honor præstat benefícia populórum, et orátio perdúcit ad regna cælórum. (Allelúja.)\end{Resp}
\begin{Resp}\Vsign Archángelus Míchaël præpósitus paradísi, quem honoríficant Angelórum cives.\end{Resp}
\begin{Resp}\Rsign Cujus honor præstat benefícia populórum, et orátio perdúcit ad regna cælórum. (Allelúja.)\end{Resp}
\switchcolumn
\EN
\begin{Resp}\Rsign This is Michael, who to battle leads the armies of the skies \Star{} Whosoever on him calleth blessed within his wardship lies. His a prayer whose voice availing aids from earth toward heaven to rise. Alleluia.\end{Resp}
\begin{Resp}\Vsign The Archangel Michael is the Vice-Roy of Paradise, and the Angels that are the dwellers therein, do hold him in worship.\end{Resp}
\begin{Resp}\Rsign Whosoever on him calleth blessed within his wardship lies. His a prayer whose voice availing aids from earth toward heaven to rise. Alleluia.\end{Resp}
\switchcolumn*

\LA
\LessonHead{Lectio V}
\switchcolumn
\EN
\LessonHead{Lesson V}
\switchcolumn*

\LA
\noindent Factum est enim, ut ex grégibus armentórum Gargáni cujúsdam taurus longe discéderet; quem diu conquisítum, in áditu spelúncæ hæréntem invenérunt. Cum vero quidam ex illis, ut taurum confígeret, sagíttam emisísset retórta sagítta in ipsum récidit sagittárium. Quæ res cum præséntes ac deínceps céteros tanto timóre affecísset, ut ad eam spelúncam própius accédere nemo audéret, Sipontíni epíscopum cónsulunt; qui, indícto trium diérum jejúnio et oratióne, rem a Deo respóndit quæri oportére.\par
\switchcolumn
\EN
\noindent And it came to pass on this wise. A certain man had a bull grazing with the flock upon Mount Gargano, and it strayed. And when they had sought it for a long while they found it jammed in the mouth of a cavern. Then one that stood there shot an arrow at it to slay it, but the arrow turned round and came back against him that had shot it. They therefore that saw it, and all those that heard it, were sore afraid because of that which had come to pass, so that no man dared any more to draw near to the cavern. But when they had sought counsel of the Bishop of Siponto, he answered, that it behoved to seek the interpretation from God, and proclaimed three days of fasting and prayer.\par
\switchcolumn*

\LA
\begin{Resp}\Rsign Venit Míchaël Archángelus cum multitúdine Angelórum, cui trádidit Deus ánimas Sanctórum, \Star{} Ut perdúcat eas in paradísum exsultatiónis. (Allelúja.)\end{Resp}
\begin{Resp}\Vsign Emítte, Dómine, Spíritum sanctum tuum de cælis, spíritum sapiéntiæ et intelléctus.\end{Resp}
\begin{Resp}\Rsign Ut perdúcat eas in paradísum exsultatiónis. (Allelúja.)\end{Resp}
\switchcolumn
\EN
\begin{Resp}\Rsign Where Angels lead the spirits of the blessed dead the glad procession moves with Michael at its head \Star{} To lead them into the garden of Eden. Alleluia.\end{Resp}
\begin{Resp}\Vsign Lord, send thy Holy Spirit from heaven the Spirit of wisdom and understanding.\end{Resp}
\begin{Resp}\Rsign To lead them into the garden of Eden. Alleluia.\end{Resp}
\switchcolumn*

\LA
\LessonHead{Lectio VI}
\switchcolumn
\EN
\LessonHead{Lesson VI}
\switchcolumn*

\LA
\noindent Post tríduum Míchaël Archángelus epíscopum monet in sua tutéla esse eum locum, eóque indício demonstrásse, velle ibi cultum Deo in sui et Angelórum memóriam adhibéri. Quare epíscopus una cum cívibus ad eam spelúncam ire pergit. Quam cum in templi cujúsdam similitúdinem conformátam vidíssent, locum illum divínis offíciis celebráre cœpérunt:: qui multis póstea miráculis illustrátus est. Nec ita multo post Bonifátius Papa, Romæ in summo circo, sancti Michaélis ecclésiam dedicávit tértio Kaléndas Octóbris; quo die étiam ómnium Angelórum memóriam Ecclésia celébrat. Hodiérnus autem dies Archángeli Michaélis apparitióne consecrátus est.\par
\switchcolumn
\EN
\noindent After three days the Archangel Michael gave warning to the Bishop that that place was under his protection, and that he had thus pointed out by a sign that he wished that worship should be offered to God there, with remembrance of himself and of the Angels. Then the Bishop and the citizens made haste and came to the cavern and when they found that the form thereof was somewhat after the fashion of a Church, they began to perform the public worship of God therein which sanctuary hath been glorified with many miracles. It was not long after these things that Pope Boniface IV. hallowed the Church of St. Michael on Hadrian's Mole at Rome, on the 29th day of September, on the which day the Church also holdeth in remembrance All Angels. But this present day is hallowed in remembrance of the manifestation of the Archangel Michael.\par
\switchcolumn*

\LA
\begin{Resp}\Rsign In témpore illo consúrget Míchaël, qui stat pro fíliis vestris: \Star{} Et véniet tempus, quale non fuit, ex quo gentes esse cœpérunt, usque ad illud. (Allelúja.)\end{Resp}
\begin{Resp}\Vsign In témpore illo salvábitur pópulus tuus omnis, qui invéntus fúerit scriptus in libro vitæ.\end{Resp}
\begin{Resp}\Rsign Et véniet tempus, quale non fuit, ex quo gentes esse cœpérunt, usque ad illud. (Allelúja.)\end{Resp}
\begin{Resp}\Vsign Glória Patri, et Fílio, \Star{} et Spirítui Sancto.\end{Resp}
\begin{Resp}\Rsign Et véniet tempus, quale non fuit, ex quo gentes esse cœpérunt, usque ad illud. (Allelúja.)\end{Resp}
\switchcolumn
\EN
\begin{Resp}\Rsign At that time shall Michael stand up, which standeth for your children. \Star{} And there shall be a time, such as never was since there was a nation even to that same time. Alleluia.\end{Resp}
\begin{Resp}\Vsign At that time thy people shall be delivered, every one that shall be found written in the Book of Life.\end{Resp}
\begin{Resp}\Rsign And there shall be a time, such as never was since there was a nation even to that same time. Alleluia.\end{Resp}
\begin{Resp}\Vsign Glory be to the Father, and to the Son, \Star{} and to the Holy Ghost.\end{Resp}
\begin{Resp}\Rsign And there shall be a time, such as never was since there was a nation even to that same time. Alleluia.\end{Resp}
\switchcolumn*

\LA
\LessonHead{Lectio VII}
\switchcolumn
\EN
\LessonHead{Lesson VII}
\switchcolumn*

\LA
\LessonTitle{Léctio sancti Evangélii secúndum Matthǽum}
\switchcolumn
\EN
\LessonTitle{From the Holy Gospel according to Matthew}
\switchcolumn*

\LA
\Cite{Matt 18:1-10}
\switchcolumn
\EN
\Cite{Matt 18:1-10}
\switchcolumn*

\LA
\noindent In illo témpore: Accessérunt discípuli ad Jesum, dicéntes: Quis, putas, major est in regno cælórum? Et réliqua.\par
\switchcolumn
\EN
\noindent At that time came the disciples unto Jesus, saying: Who is the greatest in the kingdom of heaven? And so on.\par
\switchcolumn*

\LA
\LessonTitle{Homilía sancti Hilárii Epíscopi}
\switchcolumn
\EN
\LessonTitle{Homily by St. Hilary, Bishop of Poitiers.}
\switchcolumn*

\LA
\Source{Comment in Matth. can. 18}
\switchcolumn
\EN
\Source{Com. on Matth. xviii.}
\switchcolumn*

\LA
\noindent Nónnisi revérsos in natúram puerórum introíre regnum cælórum Dóminus docet: id est, per simplicitátem puerílem vítia córporum nostrórum animǽque revocánda. Púeros autem, credéntes omnes per audiéntiæ fidem nuncupávit. Hi enim patrem sequúntur, matrem amant, próximo velle malum nésciunt, curam opum négligunt; non insoléscunt, non odérunt, non mentiúntur, dictis credunt, et quod áudiunt, verum habent. Reverténdum ígitur est ad simplicitátem infántium; quia, in ea collocáti, spéciem humilitátis Domínicæ circumferémus.\par
\switchcolumn
\EN
\noindent Unless ye become as little children,” saith the Lord, “ye shall not enter into the kingdom of heaven,” that is, unless by the uprooting of bodily and mental depravity, we bring our souls to the innocency of childhood. But He giveth the name of children to all such as believe by the hearing of faith. Children follow their father, love their mother, know not how to wish evil to their neighbours, are not careful for earthly riches they insult not, they hate not, they lie not, they believe what they are told, and take for truth what they hear. Us then it behoveth to return to the simpleness of little children, for when we are well rooted therein, we shall so far bear about in ourselves an image of the sublime simpleness of the Lord Jesus.\par
\switchcolumn*

\LA
\begin{Resp}\Rsign In conspéctu géntium nolíte timére; vos enim in córdibus vestris adoráte et timéte Dóminum; \Star{} Angelus enim ejus vobíscum est. (Allelúja.)\end{Resp}
\begin{Resp}\Vsign Stetit Angelus juxta aram templi, habens thuríbulum áureum in manu sua.\end{Resp}
\begin{Resp}\Rsign Angelus enim ejus vobíscum est. (Allelúja.)\end{Resp}
\switchcolumn
\EN
\begin{Resp}\Rsign Be not ye afraid before the Gentiles, but in your hearts, worship ye the Lord, and fear Him \Star{} For His Angel is with you. Alleluia.\end{Resp}
\begin{Resp}\Vsign An Angel stood at the Altar of the Temple, having a golden censer in his hand.\end{Resp}
\begin{Resp}\Rsign For His Angel is with you. Alleluia.\end{Resp}
\switchcolumn*

\LA
\LessonHead{Lectio VIII}
\switchcolumn
\EN
\LessonHead{Lesson VIII}
\switchcolumn*

\LA
\noindent Væ huic mundo ab scándalis. Humílitas passiónis scándalum mundo est. In hoc enim máxime ignorántia detinétur humána, quod sub deformitáte crucis, ætérnæ glóriæ Dóminum nóluit accípere. Et quid mundo tam periculósum, quam non recepísse Christum? Ideo vero necésse esse ait veníre scándala; quia, ad sacraméntum reddéndæ nobis æternitátis, omnis in eo passiónis humílitas esset explénda.\par
\switchcolumn
\EN
\noindent Woe unto the world because of offences. The lowliness of the Passion is an offence unto the world. Such is the state of stupidity to which man's ignorance hath reduced itself, that it turneth away from the Lord of Eternal Glory, because of the unsightliness of the Cross. And what is so certain to bring woe unto the world as to turn away from Christ. And therefore He saith: “It must needs be that offences come,” because His fulfilling the lowliness of the Passion was the predestined mean whereby He was to give us eternal life.\par
\switchcolumn*

\LA
\begin{Resp}\Rsign Míchaël Archángelus venit in adjutórium pópulo Dei, \Star{} Stetit in auxílium pro animábus justis. (Allelúja.)\end{Resp}
\begin{Resp}\Vsign Stetit Angelus juxta aram templi, habens thuríbulum áureum in manu sua.\end{Resp}
\begin{Resp}\Rsign Stetit in auxílium pro animábus justis. (Allelúja.)\end{Resp}
\begin{Resp}\Vsign Glória Patri, et Fílio, \Star{} et Spirítui Sancto.\end{Resp}
\begin{Resp}\Rsign Stetit in auxílium pro animábus justis. (Allelúja.)\end{Resp}
\switchcolumn
\EN
\begin{Resp}\Rsign The Archangel Michael came to help God's people. \Star{} He arose to succour the spirits of the righteous. Alleluia.\end{Resp}
\begin{Resp}\Vsign An Angel stood at the Altar of the Temple, having a golden censer in his hand.\end{Resp}
\begin{Resp}\Rsign He arose to succour the spirits of the righteous. Alleluia.\end{Resp}
\begin{Resp}\Vsign Glory be to the Father, and to the Son, \Star{} and to the Holy Ghost.\end{Resp}
\begin{Resp}\Rsign He arose to succour the spirits of the righteous. Alleluia.\end{Resp}
\switchcolumn*

\LA
\LessonHead{Lectio IX}
\switchcolumn
\EN
\LessonHead{Lesson IX}
\switchcolumn*

\LA
\noindent Vidéte ne contemnátis unum de pusíllis istis, qui credunt in me. Aptíssimum vínculum mútui amóris impósuit, ad eos præcípue qui vere in Dómino credidíssent. Pusillórum enim Angeli quotídie Deum vident: quia Fílius hóminis venit salváre quæ pérdita sunt. Ergo et Fílius hóminis salvat, et Deum Angeli vident, et Angeli pusillórum præsunt fidélium oratiónibus. Præésse Angelos absolúta auctóritas est. Salvatórum ígitur per Christum oratiónes Angeli quotídie Deo ófferunt. Ergo periculóse ille contémnitur, cujus desidéria ac postulatiónes ad ætérnum et invisíbilem Deum, ambitióso Angelórum famulátu ac ministério, pervehúntur.\par
\switchcolumn
\EN
\noindent Make heed that ye despise not one of these little ones that believe in Me.« He hath laid on us a most meet tie to constrain us to love one another, especially such as indeed believe in the Lord. »For I say unto you that in heaven their Angels do always behold the face of My Father Which is in heaven. For the Son of Man is come to save that which was lost.« From these words we see, first, that the Son of Man saveth secondly, that the Angels do see God and thirdly, that the Angels of these little ones have the wardship over the prayers of the faithful. That the Angels have this wardship is taught us absolutely. The Angels therefore do every day offer to God the prayers which they which are saved do make to Him in the Name of Christ. Therefore it is dangerous for a man to despise them, seeing that these are they by whose watchful service and ministry, his wishes and requests are presented before the throne of the eternal and unseen God.\par
\switchcolumn*

\LA
\TeDeum{Te Deum}
\switchcolumn
\EN
\TeDeum{Te Deum}
\switchcolumn*

\LA
\LessonHead{Lectio IX (pro commemoratione)}
\switchcolumn
\EN
\LessonHead{Lesson IX (when commemorated)}
\switchcolumn*

\LA
\LessonTitle{Commemoratio: In Apparitione S. Michaëlis Archangeli}
\switchcolumn
\EN
\LessonTitle{Commemoration of: Apparition of St. Michael the Archangel}
\switchcolumn*

\LA
\noindent Beátum Michaélem Archángelum sǽpius homínibus apparuísse, et sacrórum Librórum auctoritáte et véteri Sanctórum traditióne comprobátur. Quam ob rem multis in locis facti memória celebrátur. Eum, ut olim synagóga Judæórum, sic nunc custódem et patrónum Dei venerátur Ecclésia. Gelásio autem primo, Pontífice máximo, in Apúlia in vértice Gargáni montis, ad cujus radíces íncolunt Sipontíni, Archángeli Michaélis fuit illústris apparítio.\par
Nec ita multo post Bonifátius Papa, Romæ in summo circo, sancti Michaélis ecclésiam dedicávit tértio Kaléndas Octóbris; quo die étiam ómnium Angelórum memóriam Ecclésia celébrat. Hodiérnus autem dies Archángeli Michaélis apparitióne consecrátus est.\par
\switchcolumn
\EN
\noindent That the blessed Archangel Michael hath oftentimes been seen of men is attested on the authority of the Holy Bible, and also by the ancient traditions of the Saints. For this reason such visions are held in remembrance in many places. As of old time did the Synagogue of the Jews, so now doth the Church of God venerate Michael as her watcher and defender. But during the Popedom of Gelasius I, the summit of Mount Gargano in Apulia, at whose foot lieth the town of Siponto, was the scene of an extraordinary appearance of this same Archangel Michael. It was not long after these things that Pope Boniface IV hallowed the Church of St. Michael on Hadrian's Mole at Rome, on the 29th day of September, on the which day the Church also holdeth in remembrance All Angels. But this present day is hallowed in remembrance of the manifestation of the Archangel Michael.\par
\switchcolumn*

\LA
\TeDeum{Te Deum}
\switchcolumn
\EN
\TeDeum{Te Deum}
\switchcolumn*

\end{paracol}

\par\needspace{10\baselineskip}\addvspace{1.2em}{\centering\small\textsc{Die 9 Maii} · \textsc{9 May}\par}
\Office{S. Gregorii Nazianzeni Episcopi Confessoris et Ecclesiæ Doctoris}{St. Gregory Nazianzen, Bishop, Confessor and Doctor of the Church}{Duplex}
\addcontentsline{toc}{subsection}{Die 9 Maii · S. Gregorii Nazianzeni Episcopi Confessoris et Ecclesiæ Doctoris\enspace\textit{St. Gregory Nazianzen, Bishop, Confessor and Doctor of the Church}}
\begin{paracol}{2}
\LA
\RefLine{Lectiones I–III: de Scriptura occurrente.}
\switchcolumn
\EN
\RefLine{Lessons I–III: from the Scripture of the day.}
\switchcolumn*

\LA
\RefLine{Lectiones VII–IX: ex Commúni Doctoris Pontificis (C4a).}
\switchcolumn
\EN
\RefLine{Lessons VII–IX: from the Common of a Doctor Bishop (C4a).}
\switchcolumn*

\LA
\LessonHead{Lectio IV}
\switchcolumn
\EN
\LessonHead{Lesson IV}
\switchcolumn*

\LA
\noindent Gregórius, nobilis Cappadox, ex singulari divinárum Litterárum sciéntia Theólogi cognomen consecutus, Nazianzi in Cappadocia natus, Athenis in omni disciplinárum genere una cum sancto Basilío eruditus, ad stúdia sacrárum Litterárum se convértit; in quibus se in cœnobio per aliquot annos exercuérunt, illárum senténtiam non ex proprio ingenio, sed ex majórum ratióne et auctoritate interpretantes. Qui cum doctrina et vitæ sanctitáte florérent, vocáti ad munus prædicandæ evangelicæ veritátis, plurimos Jesu Christo fílios peperérunt.\par
\switchcolumn
\EN
\noindent Gregory, to whom, is commonly given, on account of his extraordinary depth of sacred learning, the title of “the Divine”, was a noble Cappadocian, born at Nazianzus in that country, and educated at Athens along with St. Basil, with whom likewise, when they had acquired knowledge in diverse branches of earthly learning, he gave himself up to learn the things of God. This they did for some years in a Monastery, framing their opinions, not out of their own heads, but according to the interpretation arrived at by the wisdom and decision of the ancients. They were both distinguished by power of doctrine and holiness of life, they were both called to the duty of preaching the Gospel of truth and, through the Gospel. they both begat many sons unto Christ.\par
\switchcolumn*

\LA
\begin{Resp}\Rsign Invéni David servum meum, óleo sancto meo unxi eum: \Star{} Manus enim mea auxiliábitur ei, allelúja.\end{Resp}
\begin{Resp}\Vsign Nihil profíciet inimícus in eo, et fílius iniquitátis non nocébit ei.\end{Resp}
\begin{Resp}\Rsign Manus enim mea auxiliábitur ei, allelúja.\end{Resp}
\switchcolumn
\EN
\begin{Resp}\Rsign I have found David My servant, with My holy oil have I anointed him \Star{} For My hand shall help him, alleluia.\end{Resp}
\begin{Resp}\Vsign The enemy shall prevail nothing against him, nor the son of wickedness afflict him.\end{Resp}
\begin{Resp}\Rsign For My hand shall help him, alleluia.\end{Resp}
\switchcolumn*

\LA
\LessonHead{Lectio V}
\switchcolumn
\EN
\LessonHead{Lesson V}
\switchcolumn*

\LA
\noindent Gregórius ígitur, aliquándo domum reversus, primum Sasimórum epíscopus creátus est, deínde Nazianzenam ecclésiam administrávit. Tum Constantinopolim ad eam regéndam ecclésiam accersítus, cum civitátem hæresum purgatam erróribus ad catholicam fidem reduxísset, quod ei summum ómnium amórem conciliare debebat, multórum parávit invidiam. Itaque, cum inter episcopos magna proptérea esset facta seditio, sponte cedens episcopatu, illud prophetæ dictum usurpávit: Si propter me commóta est ista tempéstas, dejicite me in mare, ut vos jactari desinatis. Quare Nazianzum reversus, cum illi ecclésiæ Eulálium præficiéndum curasset, totum se ad contemplatiónem et scriptiónem divinárum rerum cóntulit.\par
\switchcolumn
\EN
\noindent Gregory after a while returned home. He was first made Bishop of Sasima, and afterwards administered the Church at Nazianzus. Then he was called to rule the Church of Constantinople. That city, which he found reeking with heresy, he purged, and brought again to the Catholic faith. But this, which deserved for him the warmest love of all men, raised up many enemies. Among the Bishops themselves there was a great party against him, and to still their contentions, he, of his own free will, gave up his see, saying with the Prophet Jonah: “Take me up, and cast me forth into the sea so shall the sea be calm unto you for I know that for my sake this great tempest is upon you,” (i. 12.) So he went his way back again to Nazianzus, and when he had seen that Eulalius was set over that Church, he gave himself up altogether to think and write concerning the things of God.\par
\switchcolumn*

\LA
\begin{Resp}\Rsign Pósui adjutórium super poténtem, et exaltávi eléctum de plebe mea: \Star{} Manus enim mea auxiliábitur ei, allelúja.\end{Resp}
\begin{Resp}\Vsign Invéni David servum meum, óleo sancto meo unxi eum.\end{Resp}
\begin{Resp}\Rsign Manus enim mea auxiliábitur ei, allelúja.\end{Resp}
\switchcolumn
\EN
\begin{Resp}\Rsign I have laid help upon one that is mighty, and have exalted one chosen out of My people \Star{} For My hand shall help him, alleluia.\end{Resp}
\begin{Resp}\Vsign I have found David My servant, with My holy oil have I anointed him.\end{Resp}
\begin{Resp}\Rsign For My hand shall help him, alleluia.\end{Resp}
\switchcolumn*

\LA
\LessonHead{Lectio VI}
\switchcolumn
\EN
\LessonHead{Lesson VI}
\switchcolumn*

\LA
\noindent Scripsit autem multa, et soluta oratióne et versibus, mirábili pietáte et eloquéntia; quibus doctórum hóminum sanctorúmque judício id assecutus est, ut nihil in illis, nisi ex veræ pietátis et catholicæ religiónis regula, reperiátur, nemo quidquam jure vocare possit in dubium. Consubstantialitátis Fílii fuit acerrimus propugnator. Ut autem vitæ laude nemo ei præpositus est; sic et oratiónis gravitáte omnes facile superávit. In iis scribéndi ac legéndi stúdiis ruri vitam monachi exercens, imperatóre Theodosio ad cælestem vitam senio confectus migrávit.\par
\switchcolumn
\EN
\noindent He wrote much, both in prose and verse, with wonderful godliness and eloquence. According to the judgment of learned and holy men, there is nothing in his writings which anywhere strays from the line of true godliness and Catholic truth, and not a single word which any one can justly call in doubt. He was one of the latest champions of the doctrine that the Son is of one substance with the Father. No one has ever won greater praise for goodness of life, neither was any man more earnest in prayer. During the reign of the Emperor Theodosius, he dwelt in the country after the manner of a monk, and unceasingly taken up with writing and reading, until, in a good old age, he laid down his earthly, to enter on an heavenly life.\par
\switchcolumn*

\LA
\begin{Resp}\Rsign Iste est, qui ante Deum magnas virtútes operátus est, et omnis terra doctrína ejus repléta est: \Star{} Ipse intercédat pro peccátis ómnium populórum, allelúja.\end{Resp}
\begin{Resp}\Vsign Iste est, qui contémpsit vitam mundi, et pervénit ad cæléstia regna.\end{Resp}
\begin{Resp}\Rsign Ipse intercédat pro peccátis ómnium populórum, allelúja.\end{Resp}
\begin{Resp}\Vsign Glória Patri, et Fílio, \Star{} et Spirítui Sancto.\end{Resp}
\begin{Resp}\Rsign Ipse intercédat pro peccátis ómnium populórum, allelúja.\end{Resp}
\switchcolumn
\EN
\begin{Resp}\Rsign This is he which wrought great wonders before God, and the whole earth is full of his teaching \Star{} May he pray for all people, that their sins may be forgiven unto them, alleluia.\end{Resp}
\begin{Resp}\Vsign This is he which loved not his life in this world, and hath attained unto the kingdom of heaven.\end{Resp}
\begin{Resp}\Rsign May he pray for all people, that their sins may be forgiven unto them, alleluia.\end{Resp}
\begin{Resp}\Vsign Glory be to the Father, and to the Son, \Star{} and to the Holy Ghost.\end{Resp}
\begin{Resp}\Rsign May he pray for all people, that their sins may be forgiven unto them, alleluia.\end{Resp}
\switchcolumn*

\LA
\LessonHead{Lectio IX (pro commemoratione)}
\switchcolumn
\EN
\LessonHead{Lesson IX (when commemorated)}
\switchcolumn*

\LA
\LessonTitle{Commemoratio: S. Gregorii Nazianzeni Episcopi Confessoris et Ecclesiæ Doctoris}
\switchcolumn
\EN
\LessonTitle{Commemoration of: St. Gregory Nazianzen, Bishop, Confessor and Doctor of the Church}
\switchcolumn*

\LA
\noindent Gregórius Nazianzénus, nóbilis Cáppadox, ob singulárem divinárum Litterárum sciéntiam, Theólogi cognómen consecútus, Athénis in omni disciplinárum génere una cum sancto Basilío erudítus, ad stúdia sacrárum Litterárum se convértit. Primum Sasimórum epíscopus creátus est, deínde Nazianzénam ecclésiam administrávit. Tum Constantinópolim ad eam regéndam ecclésiam accersítus, cum civitátem erróribus hǽresum purgátam, ad cathólicam fidem reduxísset, quod ei ómnium amórem conciliáre debébat, multórum parávit invídiam. Itaque, cum inter epíscopos magna proptérea esset facta sedítio, sponte cedens episcopátu, illud prophétæ dictum usurpávit: Si propter me commóta est ista tempéstas, deícite me in mare, ut vos jactári desinátis. Naziánzum revérsus, cum illi ecclésiæ Eulálium præficiéndum curásset, se totum ad oratiónem et stúdium rerum divinárum cóntulit. Egrégie multa scripsit solúta oratióne ac vérsibus, et consubstantialitátis Fílii fuit acérrimus propugnátor. Imperatóre Theodósio, ad cæléstem vitam sénio conféctus migrávit.\par
\switchcolumn
\EN
\noindent Gregory Nazianzus, a noble Cappadocian, earned the name of The Divine from his extraordinary knowledge of the sacred sciences. It was to these that he turned after being educated at Athens, together with St. Basil, in every branch of learning. He was first made Bishop of Sosima and then administered the Church of Nazianzus. Summoned to rule over the Church of Constantinople, he purged the city of heretical errors and brought it back to the Catholic faith. Although this deed should have won him the love of all, it earned him the hatred of many; so that, when a great quarrel had arisen among the bishops on his account, he resigned his See voluntarily, making his own the words of the prophet Jonah: If this storm hath arisen on my account, then throw me into the sea, that you may cease to be tossed about. He returned to Nazianzus, and having arranged that Eulalius should be its bishop, devoted himself wholly to prayer and the study of divine things. He wrote many famous works, both in prose and in verse, and was a most ardent defender of the doctrine of the consubstantiality of the Son with the Father. When Theodosius was emperor, Gregory, now grown old, departed to the life of heaven.\par
\switchcolumn*

\LA
\TeDeum{Te Deum}
\switchcolumn
\EN
\TeDeum{Te Deum}
\switchcolumn*

\end{paracol}

\par\needspace{10\baselineskip}\addvspace{1.2em}{\centering\small\textsc{Die 10 Maii} · \textsc{10 May}\par}
\Office{S. Antonini Episcopi et Confessoris}{St. Antoninus, Bishop and Confessor}{Duplex}
\addcontentsline{toc}{subsection}{Die 10 Maii · S. Antonini Episcopi et Confessoris\enspace\textit{St. Antoninus, Bishop and Confessor}}
\begin{paracol}{2}
\LA
\RefLine{Lectiones I–III: de Scriptura occurrente.}
\switchcolumn
\EN
\RefLine{Lessons I–III: from the Scripture of the day.}
\switchcolumn*

\LA
\RefLine{Lectiones VII–VIII: ex Commúni unius Confessoris Pontificis (C4).}
\switchcolumn
\EN
\RefLine{Lessons VII–VIII: from the Common of a Confessor Bishop (C4).}
\switchcolumn*

\LA
\RefLine{Lectio IX: de Ss. Gordiani et Epimachi Martyrum (05-10o).}
\switchcolumn
\EN
\RefLine{Lesson IX: of Sts. Gordian and Epimachus, Martyrs (05-10o).}
\switchcolumn*

\LA
\LessonHead{Lectio IV}
\switchcolumn
\EN
\LessonHead{Lesson IV}
\switchcolumn*

\LA
\noindent Antoninus, Florentiæ honestis parentibus natus, ab ipsa jam pueritia egregium futuræ sanctitatis specimen exhibuit. Annum agens sextum decimum, religionem Prædicatorum amplexus, cœpit exinde maximis clarere virtutibus. Otio perpetuum bellum indixit. Post nocturnum brevem somnum primus matutinis precibus aderat; quibus persolutis, reliquum tempus noctis orationibus, aut certe lectioni, et scriptioni librorum tribuebat; et si quando importunior fessis membris somnus obreperet, ad parietem paululum declinato capite, ac tantisper discusso somno mox sacras vigilias avidius repetebat.\par
\switchcolumn
\EN
\noindent Antonine was born of respectable parents at Florence, (in the year of grace 1389,) and the holiness of his after life was foreshadowed in him even as a little child. When he was sixteen years of age he entered the Order of Friars Preachers, and from that time forth he was a burning and a shining light to all the godly. He proclaimed a truceless war against idleness after a short night's rest, he was the first to come to the service of Mattins when they were over he spent the rest of the night in prayer, or at least in reading, or writing out books, or if sleep altogether overcame his weary body, he would rest against the wall with his head bowed down, and then shake off slumber again, and set himself anew with fresh eagerness to his sacred watch.\par
\switchcolumn*

\LA
\begin{Resp}\Rsign Invéni David servum meum, óleo sancto meo unxi eum: \Star{} Manus enim mea auxiliábitur ei, allelúja.\end{Resp}
\begin{Resp}\Vsign Nihil profíciet inimícus in eo, et fílius iniquitátis non nocébit ei.\end{Resp}
\begin{Resp}\Rsign Manus enim mea auxiliábitur ei, allelúja.\end{Resp}
\switchcolumn
\EN
\begin{Resp}\Rsign I have found David My servant, with My holy oil have I anointed him \Star{} For My hand shall help him, alleluia.\end{Resp}
\begin{Resp}\Vsign The enemy shall prevail nothing against him, nor the son of wickedness afflict him.\end{Resp}
\begin{Resp}\Rsign For My hand shall help him, alleluia.\end{Resp}
\switchcolumn*

\LA
\LessonHead{Lectio V}
\switchcolumn
\EN
\LessonHead{Lesson V}
\switchcolumn*

\LA
\noindent Disciplinæ regularis sui ipsius severissimus exactor, carnes, nisi in gravi ægritudine, numquam edit. Humi, aut in nudo tabulato cubabat. Cilicio semper usus, et interdum zona ferrea ad vivam cutem incinctus, virginitatem integerrime semper coluit. In explicandis consiliis tantæ dexteritatis fuit, ut communi elogio Antoninus consiliorum diceretur. Adeo autem in eo humilitas enituit, ut étiam cœnobiis, ac provinciis præfectus abjectissima monasterii officia demississime obiret. Ab Eugenio quarto Florentinus archiepiscopus renuntiatus, ægerrime tandem nec nisi apostolicis minis perterrefactus ut episcopatum acciperet, acquievit.\par
\switchcolumn
\EN
\noindent He required of himself the most unflinching observance of the Rule of his Order, and never ate meat unless he were grievously ill. He slept upon the ground or upon bare boards. He always wore haircloth, and sometimes an iron girdle which bit into his naked skin. His virginity he kept ever undimmed by the least breath or shadow. He was so skilful in giving advice that he gained the common nickname of “Counsel Antonine.” At the same time so beautifully brilliant was his lowliness, that even when he was at the head of houses and provinces of his Order, he most cheerfully undertook all the meanest services of the houses where he was. Eugenius IV. appointed him Archbishop of Florence, and he took it so ill, that it was only when awed by the threats of the Apostolic See that he obeyed, and accepted the dignity, (in the year 1446.)\par
\switchcolumn*

\LA
\begin{Resp}\Rsign Pósui adjutórium super poténtem, et exaltávi eléctum de plebe mea: \Star{} Manus enim mea auxiliábitur ei, allelúja.\end{Resp}
\begin{Resp}\Vsign Invéni David servum meum, óleo sancto meo unxi eum.\end{Resp}
\begin{Resp}\Rsign Manus enim mea auxiliábitur ei, allelúja.\end{Resp}
\switchcolumn
\EN
\begin{Resp}\Rsign I have laid help upon one that is mighty, and have exalted one chosen out of My people \Star{} For My hand shall help him, alleluia.\end{Resp}
\begin{Resp}\Vsign I have found David My servant, with My holy oil have I anointed him.\end{Resp}
\begin{Resp}\Rsign For My hand shall help him, alleluia.\end{Resp}
\switchcolumn*

\LA
\LessonHead{Lectio VI}
\switchcolumn
\EN
\LessonHead{Lesson VI}
\switchcolumn*

\LA
\noindent In eo munere vix dici potest, quantum prudentia, pietate, caritate, mansuetudine, et sacerdotali zelo excelluerit. Illud mirandum, tantum ingenio valuisse, ut omnes ferme scientias per se nullo adhibito præceptore, absolutissime didicerit. Tandem post multos labores, multis étiam editis insignis doctrinæ libris, sacra Eucharistia et Unctione percepta, complexus Crucifixi imaginem, mortem lætus aspexit sexto Nonas Maji, anno millesimo quadringentesimo quinquagesimo nono. Miraculis vivens et post mortem conspicuus, Sanctorum numero adscriptus est ab Hadriano sexto, anno Domini millesimo quingentesimo vigesimo tertio.\par
\switchcolumn
\EN
\noindent As Archbishop it can hardly be told how noble he was, in wisdom, in godliness, in love, in meekness, in Priestly zeal. It was wonderful to see how thoroughly he taught himself nearly all the sciences, without the help of a master. At last, after much work, and publishing many valuable books on Doctrine, he received the Holy Eucharist and was anointed, and then, clasping the image of his crucified Saviour to his heart, joyfully welcomed death, on the 2nd day of May, in the year 1459. He was remarkable for the working of miracles, both during his life and after his death, and Adrian VI. enrolled his name among those of the Saints, in the year 1523.\par
\switchcolumn*

\LA
\begin{Resp}\Rsign Iste est, qui ante Deum magnas virtútes operátus est, et omnis terra doctrína ejus repléta est: \Star{} Ipse intercédat pro peccátis ómnium populórum, allelúja.\end{Resp}
\begin{Resp}\Vsign Iste est, qui contémpsit vitam mundi, et pervénit ad cæléstia regna.\end{Resp}
\begin{Resp}\Rsign Ipse intercédat pro peccátis ómnium populórum, allelúja.\end{Resp}
\begin{Resp}\Vsign Glória Patri, et Fílio, \Star{} et Spirítui Sancto.\end{Resp}
\begin{Resp}\Rsign Ipse intercédat pro peccátis ómnium populórum, allelúja.\end{Resp}
\switchcolumn
\EN
\begin{Resp}\Rsign This is he which wrought great wonders before God, and the whole earth is full of his teaching \Star{} May he pray for all people, that their sins may be forgiven unto them, alleluia.\end{Resp}
\begin{Resp}\Vsign This is he which loved not his life in this world, and hath attained unto the kingdom of heaven.\end{Resp}
\begin{Resp}\Rsign May he pray for all people, that their sins may be forgiven unto them, alleluia.\end{Resp}
\begin{Resp}\Vsign Glory be to the Father, and to the Son, \Star{} and to the Holy Ghost.\end{Resp}
\begin{Resp}\Rsign May he pray for all people, that their sins may be forgiven unto them, alleluia.\end{Resp}
\switchcolumn*

\LA
\LessonHead{Lectio IX (pro commemoratione)}
\switchcolumn
\EN
\LessonHead{Lesson IX (when commemorated)}
\switchcolumn*

\LA
\LessonTitle{Commemoratio: S. Antonini Episcopi et Confessoris}
\switchcolumn
\EN
\LessonTitle{Commemoration of: St. Antoninus, Bishop and Confessor}
\switchcolumn*

\LA
\noindent Antonínus, Floréntiæ honéstis paréntibus natus, ab ipsa puerítia egrégium futúræ sanctitátis spécimen exhíbuit. Annum agens sextum décimum, religiónem Prædicatórum ampléxus, cœpit exínde máximis virtútibus clarére. Singulári fuit abstinéntia, et virginitátem integérrime semper cóluit. In explicándis consíliis tantæ fuit dexteritátis, ut Antonínus consiliórum dicerétur. Ab Eugénio quarto Florentínus archiepíscopus renuntiátus, ægérrime tamen, nec nisi apostólicis minis perterrefáctus, ut episcopátum accíperet, acquiévit. In eo múnere prudéntia, pietáte, caritáte, mansuetúdine et sacerdotáli zelo excélluit. Omnes fere sciéntias, nullo adhíbito præceptóre, absolutíssimas dídicit, et multos insígnis doctrínæ libros scripsit. Obiit in Dómino sexto nonas maji, anno millésimo quadringentésimo quinquagésimo nono, et ab Hadriáno sexto in album Sanctórum fuit relátus.\par
\switchcolumn
\EN
\noindent Antoninus, born in Florence of good parents, from his boyhood gave remarkable evidence of his future sanctity. At the age of sixteen, he entered the Order of Friars Preachers, and from then on, he shone with the greatest virtues. His living was singularly abstemious, and he preserved his virginity intact. He was so skillful in giving advice that he was called Antoninus the Counsellor. Named Archbishop of Florence by Eugenius IV, he reluctantly accepted the post only at the Pope's insistence. In this office he excelled in prudence, piety, charity, gentleness and priestly zeal. With no teacher to help him, he gained a thorough mastery of almost all the sciences and wrote many famous books expounding them. He died in the Lord on May 2, 1459, and was canonized by Adrian VI.\par
\switchcolumn*

\LA
\TeDeum{Te Deum}
\switchcolumn
\EN
\TeDeum{Te Deum}
\switchcolumn*

\end{paracol}

\par\needspace{10\baselineskip}\addvspace{1.2em}{\centering\small\textsc{Die 10 Maii} · \textsc{10 May}\par}
\Office{Ss. Gordiani et Epimachi Martyrum}{Sts. Gordian and Epimachus, Martyrs}{Simplex}
\addcontentsline{toc}{subsection}{Die 10 Maii · Ss. Gordiani et Epimachi Martyrum\enspace\textit{Sts. Gordian and Epimachus, Martyrs}}
\begin{paracol}{2}
\LA
\LessonHead{Lectio IX (pro commemoratione)}
\switchcolumn
\EN
\LessonHead{Lesson IX (when commemorated)}
\switchcolumn*

\LA
\LessonTitle{Commemoratio: Ss. Gordiani et Epimachi Martyrum}
\switchcolumn
\EN
\LessonTitle{Commemoration of: Sts. Gordian and Epimachus, Martyrs}
\switchcolumn*

\LA
\noindent Gordianus judex, cum ad eum Januarius presbyter, ut condemnaretur, sub Juliano Apostata ductus esset, ab eodem in christiana fide instructus, cum uxore et quinquaginta tribus aliis ex eadem famiha Romæ baptizatur. Quare præfectus, relegdto Januario, Gordianum a Clementiano vicario includi jubet in carcerem: qui postea eumdem Gordianum vinctum catenis ad se accersitum, cum a fidei proposito deterrere non posset, plumbatis diu cæsum, capite plecti imperat. Cujus corpus ante Apollinis templum canibus objectum, noctu a Christianis via Latina sepelitur in eadem crypta, in quam reliquiæ beati Epimachi Martyris translatæ fuerant ab Alexandria: ubi is diu propter Christi confessionem constrictus in carcere, postremo combustus, martyrio coronatus est.\par
\switchcolumn
\EN
\noindent Gordian was a judge before whom, in the reign of Julian the Apostate, Januarius the Priest was brought to be condemned. Januarius instructed Gordian in the Christian faith, and himself, with his wife, and fifty-three other persons of the same household, were all baptized at Rome. On this account the Praetor sent back Januarius, and ordered Clementian the Deputy to cast Gordian into prison. Afterward he caused the same Gordian to be brought before him in chains, and when he found he could not shake him in his will to cleave to the faith, he commanded that he should first be hided with whips loaded with lead, and thereafter beheaded. His body was thrown out before the temple of Apollo for dogs to eat, but the Christians buried it at night in the catacombs upon the Latin Way, in the same vault where were already lying the remains of the blessed Martyr Epimachus. These had been brought from Alexandria, in which city Epimachus had long been imprisoned for owning Christ, and had in the end grasped the crown of his testimony by being burnt alive.\par
\switchcolumn*

\LA
\TeDeum{Te Deum}
\switchcolumn
\EN
\TeDeum{Te Deum}
\switchcolumn*

\end{paracol}

\par\needspace{10\baselineskip}\addvspace{1.2em}{\centering\small\textsc{Die 12 Maii} · \textsc{12 May}\par}
\Office{Ss. Nerei, Achillei et Domitillæ Virg. atque Pancratii Martyrum}{Ss. Nereus, Achilleus and Domitilla the Virgin and Pancras, Martyrs}{Semiduplex}
\addcontentsline{toc}{subsection}{Die 12 Maii · Ss. Nerei, Achillei et Domitillæ Virg. atque Pancratii Martyrum\enspace\textit{Ss. Nereus, Achilleus and Domitilla the Virgin and Pancras, Martyrs}}
\begin{paracol}{2}
\LA
\RefLine{Lectiones I–III: de Scriptura occurrente.}
\switchcolumn
\EN
\RefLine{Lessons I–III: from the Scripture of the day.}
\switchcolumn*

\LA
\LessonHead{Lectio IV}
\switchcolumn
\EN
\LessonHead{Lesson IV}
\switchcolumn*

\LA
\noindent Nereus et Achilleus fratres eunuchi Flaviæ Domitillæ, a beato Petro una cum ipsa ejusque matre Plautilla baptizati, cum Domitillæ persuasissent, ut virginitátem suam Deo consecraret, ab ejus sponso Aureliáno tamquam Christiáni accusati, ob præcláram fidei confessiónem in Pontiam insulam relegántur. Ubi ad quæstiónem íterum vocáti et verbéribus cæsi, mox Tarracinam perducti, a Minucio Rufo equuleo et flammis cruciati, cum constanter negarent se, a sancto Petro Apóstolo baptizatos, ullis tormentis cogi posse ut idolis immolarent, secúri percússi sunt. Quorum córpora ab Auspicio, eórum discipulo et Domitillæ educatore, Romam delata, via Ardeatina sepúlta sunt.\par
\switchcolumn
\EN
\noindent Nereus and Achilles were brethren, eunuchs belonging to Flavia Domitilla, who were baptized by blessed Peter, along with her and her mother Plautilla. They had advised Domitilla to consecrate her virginity to God, and on this account Aurelian, to whom she was betrothed, accused them of being Christians. They nobly confessed the faith, and were banished to the island of Ponza. Then they were again put to the torture, and after being scourged, were taken to Terracina. At Terracina, Minutius Rufus tormented them with the rack and with fire, but as they constantly affirmed that having once been baptized by the blessed Apostle Peter, no torture could ever make them sacrifice to idols, they were beheaded. Auspicius, their own disciple and the tutor of Domitilla, took their bodies to Rome, where they were buried on the road to Ardea.\par
\switchcolumn*

\LA
\begin{Resp}\Rsign Lux perpétua lucébit Sanctis tuis, Dómine, \Star{} Et ætérnitas témporum, allelúja, allelúja.\end{Resp}
\begin{Resp}\Vsign Lætítia sempitérna erit super cápita eórum: gáudium et exsultatiónem obtinébunt.\end{Resp}
\begin{Resp}\Rsign Et ætérnitas témporum, allelúja, allelúja.\end{Resp}
\switchcolumn
\EN
\begin{Resp}\Rsign Eternal light will shine over your saints, O Lord, \Star{} And the eternity of the times, alleluia, alleluia.\end{Resp}
\begin{Resp}\Vsign Everlasting joy shall be upon their heads, they shall obtain joy and gladness\end{Resp}
\begin{Resp}\Rsign And the eternity of the times, alleluia, alleluia.\end{Resp}
\switchcolumn*

\LA
\LessonHead{Lectio V}
\switchcolumn
\EN
\LessonHead{Lesson V}
\switchcolumn*

\LA
\noindent Flavia Domitilla virgo Romana, Titi et Domitiáni imperatórum neptis, cum sacrum virginitátis velámen a beato Cleménte Papa accepísset, ab Aureliáno sponso, Titi Aurelii consulis fílio, deláta quod Christiana esset, a Domitiáno imperatóre in insulam Pontiam est deportata, ubi in carcere longum martyrium duxit. Demum Tarracinam deducta, íterum Christum confessa, cum semper constantior apparéret, sub Trajano imperatóre, júdicis jussu incénso ejus cubiculo, una cum Theodora et Euphrosyna virgínibus et collactaneis suis, gloriosi martyrii cursum confecit, Nonis Maji: quarum corpora, integra invénta, a Cæsario diacono sepúlta sunt. Hac vero die duórum fratrum ac Domitillæ corpora, ex diaconía sancti Hadriáni simul translata, in ipsórum Mártyrum basilicam, tituli Fascíolæ, restituta sunt.\par
\switchcolumn
\EN
\noindent The Virgin Flavia Domitilla was a Roman, the niece of the Emperors Titus and Domitian, and was veiled by the blessed Pope Clement. Aurelian, son of the Consul Titus Aurelius, to whom she was betrothed, accused her of being a Christian, and the Emperor Domitian banished her into the island of Ponza, where she long suffered and testified in prison. At length she was taken to Terracina, where she again confessed Christ, and as she seemed ever to grow firmer, the judge, under the Emperor Trajan, caused her chamber to be set on fire, and there Domitilla, with her foster-sisters the maidens Theodora and Euphrosyne, finished the race of faith by grasping the crown of glory, on the 7th day of May. Their bodies were found whole, and were buried by the Deacon Caesarius. This, the twelfth day of May, is that whereon the bodies of Nereus and Achilles, and that of Domitilla, were carried from the Deaconry of St. Hadrian, and laid in the Church which is properly called by the name of these holy martyrs, but formerly by that of “St. Peter's Bandage.”\par
\switchcolumn*

\LA
\begin{Resp}\Rsign In servis suis, allelúja, \Star{} Consolábitur Deus, allelúja.\end{Resp}
\begin{Resp}\Vsign Judicábit Dóminus pópulum suum, et in servis suis.\end{Resp}
\begin{Resp}\Rsign Consolábitur Deus, allelúja.\end{Resp}
\switchcolumn
\EN
\begin{Resp}\Rsign His servants, alleluia \Star{} God will comfort, alleluia, alleluia.\end{Resp}
\begin{Resp}\Vsign The Lord will judge His people and, His servants,\end{Resp}
\begin{Resp}\Rsign God will comfort, alleluia, alleluia.\end{Resp}
\switchcolumn*

\LA
\LessonHead{Lectio VI}
\switchcolumn
\EN
\LessonHead{Lesson VI}
\switchcolumn*

\LA
\noindent Pancratius, in Phrygia nobili genere natus, puer quatuordecim annórum Romam venit Diocletiáno et Maximiáno imperatóribus. Ubi a Pontifice Romano baptizátus et in fide christiana eruditus, ob eamdem paulo post comprehénsus; cum diis sacrificáre constanter renuísset, virili fortitúdine datis cervicibus, illustrem martyrii corónam consecutus est. Cujus corpus Octavilla matrona noctu sústulit, et unguentis delibutum via Aurelia sepelívit.\par
\switchcolumn
\EN
\noindent Pancras was the son of a noble family of Phrygia. He came to Rome in the reign of the Emperors Diocletian and Maximian, being then a boy of fourteen years of age. There he was baptized by the Bishop of Rome, and Drought up in the Christian faith. On this account he was soon after taken, and having constantly refused to sacrifice to the gods, he offered his neck to the -executioner with manly courage, and won a glorious crown of martyrdom. The Lady Octavilla took his body by night, embalmed it with precious ointments, and buried it on the Aurelian Way.\par
\switchcolumn*

\LA
\begin{Resp}\Rsign Fíliæ Jerúsalem, veníte et vidéte Mártyres cum corónis, quibus coronávit eos Dóminus \Star{} In die solemnitátis et lætítiæ, allelúja.\end{Resp}
\begin{Resp}\Vsign Quóniam confortávit seras portárum tuárum, benedíxit fílios tuos in te.\end{Resp}
\begin{Resp}\Rsign In die solemnitátis et lætítiæ, allelúja.\end{Resp}
\begin{Resp}\Vsign Glória Patri, et Fílio, \Star{} et Spirítui Sancto.\end{Resp}
\begin{Resp}\Rsign In die solemnitátis et lætítiæ, allelúja.\end{Resp}
\switchcolumn
\EN
\begin{Resp}\Rsign Come forth, O ye daughters of Jerusalem, and behold the Martyrs with the crowns wherewith the Lord crowned them; \Star{} In the day of His feasting, and of His gladness. Alleluia.\end{Resp}
\begin{Resp}\Vsign For He hath strengthened the bars of thy gates He hath blessed thy children within thee.\end{Resp}
\begin{Resp}\Rsign In the day of His feasting, and of His gladness. Alleluia.\end{Resp}
\begin{Resp}\Vsign Glory be to the Father, and to the Son, \Star{} and to the Holy Ghost.\end{Resp}
\begin{Resp}\Rsign In the day of His feasting, and of His gladness. Alleluia.\end{Resp}
\switchcolumn*

\LA
\LessonHead{Lectio VII}
\switchcolumn
\EN
\LessonHead{Lesson VII}
\switchcolumn*

\LA
\LessonTitle{Léctio sancti Evangélii secúndum Joánnem}
\switchcolumn
\EN
\LessonTitle{From the Holy Gospel according to John}
\switchcolumn*

\LA

\switchcolumn
\EN
\Cite{John 4:46-53}
\switchcolumn*

\LA
\Source{Joannes 46-53}
\switchcolumn
\EN

\switchcolumn*

\LA
\noindent In illo témpore: Erat quidam regulus, cujus filius infirmabatur Capharnaum. Et réliqua.\par
\switchcolumn
\EN
\noindent At that time There was a certain nobleman, whose son was sick at Capernaum. And so on.\par
\switchcolumn*

\LA
\LessonTitle{Homilía sancti Gregórii Papæ}
\switchcolumn
\EN
\LessonTitle{Homily by Pope St. Gregory the Great.}
\switchcolumn*

\LA
\Source{Ex Homil. 28. habita in Basilica horum Ss. Mártyrum, in die natalis eorum}
\switchcolumn
\EN
\Source{28th Sermon, viz. that preached on the Birth-day and in the Church of these Holy Martyrs.}
\switchcolumn*

\LA
\noindent Quid est quod regulus Dominum rogat ut ad ejus filium veniat, et tamen ire corporaliter recusat: ad servum vero centurionis non invitatur, et tamen se corporaliter ire pollicetur? Reguli filio per corporalem præsentiam non dignatur adesse, centurionis servo non dedignatur occurrere. Quid est hoc, nisi quod superbia nostra retunditur, qui in hominibus non naturam, qua ad imaginem Dei facti sunt, sed honores et divitias veneramur? Redemptor vero noster ut ostenderet, quia quæ alta sunt hominum, despicienda sunt, et quæ despecta sunt hominum, despicienda non sunt: ad filium reguli ire noluit, ad servum centurionis ire paratus fuit.\par
\switchcolumn
\EN
\noindent Wherefore was it that when this nobleman besought the Lord to come down where his child died, the Lord (albeit He healed him) would not come, and yet, when the Centurion prayed Him to heal his servant, albeit not asked to come down, He went with them He deemed not that the nobleman's son was worthy of His bodily presence, but He refused not to go to help the Centurion's servant. What is this but a rebuke to earthly pride, which maketh us to respect in men their honours and riches rather than that Divine image wherein they are created It was not so with our Redeemer, who would not go to the son of the nobleman, but was ready to come down for the Centurion's servant, to show that to Him the things which are great among men are but of little moment, and the things which are little esteemed among men are not beneath His notice.\par
\switchcolumn*

\LA
\begin{Resp}\Rsign Ego sum vitis vera, et vos pálmites: \Star{} Qui manet in me, et ego in eo, hic fert fructum multum, allelúja, allelúja.\end{Resp}
\begin{Resp}\Vsign Sicut diléxit me Pater, et ego diléxi vos.\end{Resp}
\begin{Resp}\Rsign Qui manet in me, et ego in eo, hic fert fructum multum, allelúja, allelúja.\end{Resp}
\switchcolumn
\EN
\begin{Resp}\Rsign I am the vine: you the branches. \Star{} He that abideth in me, and I in him, the same beareth much fruit, alleluia, alleluia.\end{Resp}
\begin{Resp}\Vsign As the Father hath loved me, I also have loved you.\end{Resp}
\begin{Resp}\Rsign He that abideth in me, and I in him, the same beareth much fruit, alleluia, alleluia.\end{Resp}
\switchcolumn*

\LA
\LessonHead{Lectio VIII}
\switchcolumn
\EN
\LessonHead{Lesson VIII}
\switchcolumn*

\LA
\noindent Increpata est ergo superbia nostra, quæ nescit pensare homines propter homines. Sola, ut diximus, quæ circumstant hominibus, pensat, naturam non aspicit, honorem Dei in hominibus non agnoscit. Ecce ire non vult Filius Dei ad filium reguli: et tamen venire paratus est ad salutem servi. Certe si nos cujuspiam servus rogaret, ut ad eum ire deberemus, protinus nobis nostra superbia in cogitatione tacita responderet, dicens: Non eas, quia temetipsum degeneras, honor tuus despicitur, locus vilescit. Ecce de cælo venit, qui servo in terra occurrere non despicit: et tamen humiliari in terra contemnimus, qui de terra sumus.\par
\switchcolumn
\EN
\noindent Our pride then standeth rebuked, that pride which maketh us forget for the sake of one man that another man is a man at all. This pride, as we have said, looketh only at the surroundings of men, not at their nature, and seeth not that God is to be honoured in a man because he is a man. Lo how the Son of God will not go unto the nobleman's son, but is ready to go and heal the servant. Of myself I know that if any one's servant were to ask me to go to him, I have a sort of pride which would say to me silently inside my heart Go not thou wilt lower thyself; the Papal dignity will be lightly esteemed thy exalted station will be degraded. Behold how He Which came down from heaven, doth not deem it below Him to go to help a servant, and yet I who am of the earth earthly, shrink from being trodden on.\par
\switchcolumn*

\LA
\begin{Resp}\Rsign Cándidi facti sunt Nazarǽi ejus, allelúja: splendórem Deo dedérunt, allelúja: \Star{} Et sicut lac coaguláti sunt, allelúja, allelúja.\end{Resp}
\begin{Resp}\Vsign Candidióres nive, nitidióres lacte, rubicundióres ébore antíquo, sapphíro pulchrióres.\end{Resp}
\begin{Resp}\Rsign Et sicut lac coaguláti sunt, allelúja, allelúja.\end{Resp}
\begin{Resp}\Vsign Glória Patri, et Fílio, \Star{} et Spirítui Sancto.\end{Resp}
\begin{Resp}\Rsign Et sicut lac coaguláti sunt, allelúja, allelúja.\end{Resp}
\switchcolumn
\EN
\begin{Resp}\Rsign Her Nazarites are become pure, Alleluia: they reflect the glory of God, Alleluia. \Star{} They are whiter than milk. Alleluia, Alleluia.\end{Resp}
\begin{Resp}\Vsign They are purer than snow, they are whiter than milk, they are more ruddy in body than coral, their polishing is of sapphire.\end{Resp}
\begin{Resp}\Rsign They are whiter than milk. Alleluia, Alleluia.\end{Resp}
\begin{Resp}\Vsign Glory be to the Father, and to the Son, \Star{} and to the Holy Ghost.\end{Resp}
\begin{Resp}\Rsign They are whiter than milk. Alleluia, Alleluia.\end{Resp}
\switchcolumn*

\LA
\LessonHead{Lectio IX}
\switchcolumn
\EN
\LessonHead{Lesson IX}
\switchcolumn*

\LA
\noindent Nolite ergo intra vosmetipsos pensare quod habetis, sed quid estis. Ecce mundus qui diligitur, fugit. Sancti isti, ad quorum tumbam consistimus, florentem mundum mentis despectu calcaverunt. Erat tunc vita longa, salus continua, opulentia in rebus, fœcunditas in propagine, tranquillitas in diuturna pace: et tamen cum in seipso floreret, jam in eorum cordibus mundus aruerat. Ecce jam mundus in seipso aruit, et adhuc in cordibus nostris floret. Ubique mors, ubique luctus, ubique desolatio: undique percutimur, undique amaritudinibus replemur: et tamen cæca mente carnalis concupiscentiæ ipsas ejus amaritudines amamus, fugientem sequimur, labenti inhæremus.\par
\switchcolumn
\EN
\noindent Think not therefore within yourselves what ye have, but what ye are. Behold, the world which I love, is a world which passeth away. Those holy servants of God, by whose grave I am standing, ennobled themselves mentally above the world at its fairest. To them was offered length of days, robust health, plenty in possessions, fruitfulness in offspring, comfort under perpetual peace: and yet while the spring-tide of life was unfolding before them, their hearts had already condemned it to an arid winter. Behold, winter in their hearts, spring in mine Death, and pain, and barrenness occur all around me, I am attacked on all sides, and I feel very bitter, and yet the sting of fleshly lust so blindeth me, that I love the bitter feelings, I hunt after that which flees from me, and cling to that which would leave me.\par
\switchcolumn*

\LA
\TeDeum{Te Deum}
\switchcolumn
\EN
\TeDeum{Te Deum}
\switchcolumn*

\LA
\LessonHead{Lectio IX (pro commemoratione)}
\switchcolumn
\EN
\LessonHead{Lesson IX (when commemorated)}
\switchcolumn*

\LA
\LessonTitle{Commemoratio: Ss. Nerei, Achillei et Domitillæ Virg. atque Pancratii Martyrum}
\switchcolumn
\EN
\LessonTitle{Commemoration of: Ss. Nereus, Achilleus and Domitilla the Virgin and Pancras, Martyrs}
\switchcolumn*

\LA
\noindent Néreus et Achílleus fratres, eunúchi Fláviæ Domitíllæ, a beáto Petro una cum ipsa ejúsque matre Plautílla baptizáti, cum Domitíllæ persuasíssent, ut virginitátem suam Deo consecráret, ab ejus sponso Aureliáno, quod Christiáni essent accusáti, in Póntiam ínsulam relegántur. Mox verbéribus cæsi, ut idólis immolárent, et Tarracínam perdúcti, equúlei et flammárum cruciátibus superátis, secúri percússi sunt; quorum córpora ab Auspício eórum discípulo Romam deláta, via Ardeatína sunt sepúlta. Flávia Domitílla vero, quæ sacrum virginitátis velámen a beáto Cleménte Papa accéperat, et ipsa in ínsulam Póntiam deportáta, et post diutúrna víncula Tarracínam dedúcta, júdicis jussu incénso ejus cubículo, una cum virgínibus Theodóra et Euphrósyna, collactáneis suis, gloriósam mortem oppétiit nonis maji, Trajáno imperatóre; quarum córpora Cæsárius diáconus sepelívit. Pancrátius, nóbili génere in Phrýgia natus, puer quatuórdecim annórum Romæ baptizátus, Diocletiáno et Maximiáno imperatóribus, comprehénditur, et, cum diis sacrificáre constánter renuísset, datis cervícibus, illústrem martýrii corónam consecútus est; cujus corpus ab Octavílla matróna clam via Aurélia sepúltum est.\par
\switchcolumn
\EN
\noindent The brothers Nereus and Achilleus were eunuchs of Flavia Domitilla and were baptized by St. Peter at the same time as she herself and her mother Plautilla. Because they persuaded Domitilla to consecrate her virginity to God, they were accused of being Christians by Aurelian, who had been betrothed to her, and were sent to the island of Ponza. Soon afterwards, they were scourged in an effort to make them sacrifice to idols, and were taken to Terracina, where, after they had overcome the torture of the rack and flaming torches, they were beheaded. Their bodies were taken to Rome by their disciple Auspicius and buried on the Ardeatine Way. As for Flavia Domitilla, who had received the sacred veil of a virgin from Pope St. Clement, she also was deported to the island of Ponza, and after a long imprisonment was taken to Terracina. There, by the judge's orders, her dwelling was set on fire, and she won a glorious death, along with the virgins Theodora and Euphrosyna, her foster-sisters, on May 7, under Emperor Trajan. Their bodies were buried by the Deacon Caesarius. Pancras, born of a noble Phrygian family, was baptized in Rome at the age of fourteen. Under the Emperors Diocletian and Maximian, he was arrested; and when he firmly refused to sacrifice to the gods, he was beheaded and so won the glorious crown of martyrdom. His body was buried secretly on the Via Aurelia by the matron Octavilla.\par
\switchcolumn*

\LA
\TeDeum{Te Deum}
\switchcolumn
\EN
\TeDeum{Te Deum}
\switchcolumn*

\end{paracol}

\par\needspace{10\baselineskip}\addvspace{1.2em}{\centering\small\textsc{Die 13 Maii} · \textsc{13 May}\par}
\Office{S. Roberti Bellarmino Episcopi Confessoris et Ecclesiæ Doctoris}{St. Robert Bellarmine, Bishop, Confessor and Doctor of the Church}{Duplex}
\addcontentsline{toc}{subsection}{Die 13 Maii · S. Roberti Bellarmino Episcopi Confessoris et Ecclesiæ Doctoris\enspace\textit{St. Robert Bellarmine, Bishop, Confessor and Doctor of the Church}}
\begin{paracol}{2}
\LA
\LessonHead{Lectio I}
\switchcolumn
\EN
\LessonHead{Lesson I}
\switchcolumn*

\LA
\LessonTitle{De libro Ecclesiástici}
\switchcolumn
\EN
\LessonTitle{Lesson from the book of Ecclesiasticus}
\switchcolumn*

\LA
\Cite{Sir 39:1-5}
\switchcolumn
\EN
\Cite{Sir 39:1-5}
\switchcolumn*

\LA
\noindent Sapiéntiam ómnium antiquórum exquíret sápiens, et in prophétis vacábit. Narratiónem virórum nominatórum conservábit, et in versútias parabolárum simul introíbit. Occúlta proverbiórum exquíret, et in abscónditis parabolárum conversábitur. In médio magnatórum ministrábit, et in conspéctu prǽsidis apparébit. In terram alienigenárum géntium pertránsiet; bona enim et mala in homínibus tentábit.\par
\switchcolumn
\EN
\noindent The wise men will seek out the wisdom of all the ancients, and will be occupied in the prophets. He will keep the sayings of renowned men, and will enter withal into the subtilties of parables. He will search out the hidden meanings of proverbs, and will be conversant in the secrets of parables. He shall serve among great men, and: appear before the governor. He shall pass into strange countries: for he shall try good and evil among men.\par
\switchcolumn*

\LA
\begin{Resp}\Rsign Euge, serve bone et fidélis, quia in pauca fuísti fidélis, supra multa te constítuam: \Star{} Intra in gáudium Dómini tui, allelúja.\end{Resp}
\begin{Resp}\Vsign Dómine, quinque talénta tradidísti mihi, ecce ália quinque superlucrátus sum.\end{Resp}
\begin{Resp}\Rsign Intra in gáudium Dómini tui, allelúja.\end{Resp}
\switchcolumn
\EN
\begin{Resp}\Rsign Well done, thou good and faithful servant, thou hast been faithful over a few things. I will make thee ruler over many things: \Star{} Enter thou into the joy of thy Lord, alleluia.\end{Resp}
\begin{Resp}\Vsign Lord, thou deliveredst unto me five talents; behold, I have gained beside them five talents more.\end{Resp}
\begin{Resp}\Rsign Enter thou into the joy of thy Lord, alleluia.\end{Resp}
\switchcolumn*

\LA
\LessonHead{Lectio II}
\switchcolumn
\EN
\LessonHead{Lesson II}
\switchcolumn*

\LA
\Cite{Sir 39:6-10}
\switchcolumn
\EN
\Cite{Sir 39:6-10}
\switchcolumn*

\LA
\noindent Cor suum tradet ad vigilándum dilúculo ad Dóminum, qui fecit illum, et in conspéctu Altíssimi deprecábitur. Apériet os suum in oratióne, et pro delíctis suis deprecábitur. Si enim Dóminus magnus volúerit, spíritu intellegéntiæ replébit illum: et ipse tamquam imbres mittet elóquia sapiéntiæ suæ, et in oratióne confitébitur Dómino: et ipse díriget consílium ejus, et disciplínam, et in abscónditis suis consiliábitur.\par
\switchcolumn
\EN
\noindent He will give his heart to resort early to the Lord that made him, and he will pray in the sight of the most High. He will open his mouth in prayer, and will make supplication for his sins. For if it shall please the great Lord, he will fill him with the spirit of understanding: And he will pour forth the words of his wisdom as showers, and in his prayer he will confess to the Lord. And he shall direct his counsel, and his knowledge, and in his secrets shall he meditate.\par
\switchcolumn*

\LA
\begin{Resp}\Rsign Ecce sacérdos magnus, qui in diébus suis plácuit Deo: \Star{} Ideo jurejurándo fecit illum Dóminus créscere in plebem suam, allelúja.\end{Resp}
\begin{Resp}\Vsign Benedictiónem ómnium géntium dedit illi, et testaméntum suum confirmávit super caput ejus.\end{Resp}
\begin{Resp}\Rsign Ideo jurejurándo fecit illum Dóminus créscere in plebem suam, allelúja.\end{Resp}
\switchcolumn
\EN
\begin{Resp}\Rsign Behold an high priest, who in his days pleased God \Star{} Therefore the Lord assured him by an oath that He would multiply his seed among His people, alleluia.\end{Resp}
\begin{Resp}\Vsign He hath made him a blessing unto all nations, and hath established His covenant upon his head.\end{Resp}
\begin{Resp}\Rsign Therefore the Lord assured him by an oath that He would multiply his seed among His people, alleluia.\end{Resp}
\switchcolumn*

\LA
\LessonHead{Lectio III}
\switchcolumn
\EN
\LessonHead{Lesson III}
\switchcolumn*

\LA
\Cite{Sir 39:11-14}
\switchcolumn
\EN
\Cite{Sir 39:11-14}
\switchcolumn*

\LA
\noindent Ipse palam fáciet disciplínam doctrínæ suæ, et in lege testaménti Dómini gloriábitur. Collaudábunt multi sapiéntiam ejus, et usque in sǽculum non delébitur. Non recédet memória ejus, et nomen ejus requirétur a generatióne in generatiónem. Sapiéntiam ejus enarrábunt gentes, et laudem ejus enuntiábit Ecclésia.\par
\switchcolumn
\EN
\noindent He shall show forth the discipline he hath learned, and shall glory in the law of the covenant of the Lord. Many shall praise his wisdom, and it shall never be forgotten. The memory of him shall not depart away, and his name shall be in request from generation to generation. Nations shall declare his wisdom, and the church shall show forth his praise.\par
\switchcolumn*

\LA
\begin{Resp}\Rsign Jurávit Dóminus, et non pœnitébit eum: \Star{} Tu es sacérdos in ætérnum secúndum órdinem Melchísedech, allelúja.\end{Resp}
\begin{Resp}\Vsign Dixit Dóminus Dómino meo, Sede a dextris meis.\end{Resp}
\begin{Resp}\Rsign Tu es sacérdos in ætérnum secúndum órdinem Melchísedech, allelúja.\end{Resp}
\begin{Resp}\Vsign Glória Patri, et Fílio, \Star{} et Spirítui Sancto.\end{Resp}
\begin{Resp}\Rsign Tu es sacérdos in ætérnum secúndum órdinem Melchísedech, allelúja.\end{Resp}
\switchcolumn
\EN
\begin{Resp}\Rsign The Lord hath sworn and will not repent \Star{} Thou art a Priest for ever after the order of Melchisedek, alleluia.\end{Resp}
\begin{Resp}\Vsign The Lord said unto my Lord Sit Thou at My right hand. R.* Thou art a Priest for ever after the order of Melchisedek, alleluia.\end{Resp}
\begin{Resp}\Vsign Glory be to the Father, and to the Son, \Star{} and to the Holy Ghost.\end{Resp}
\begin{Resp}\Rsign Thou art a Priest for ever after the order of Melchisedek, alleluia.\end{Resp}
\switchcolumn*

\LA
\LessonHead{Lectio IV}
\switchcolumn
\EN
\LessonHead{Lesson IV}
\switchcolumn*

\LA
\noindent Robertus, Politianus, e patrícia Bellarminórum gente, matrem pientíssimam hábuit Cynthiam Cervini, Marcelli Papæ secúndi sorórem. Eximia pietáte et castíssimis móribus quamprimum enituit, id unum exoptans, ut Deo soli placeret et ánimas Christo lucrifáceret. Patrium Societátis Jesu collegium summa cum ingenii et modestiæ laude frequentávit; ac duodeviginti annos natus, Romæ eamdem Societátem ingréssus, religiosárum virtútum ómnibus exemplo fuit. Emenso in Romano Collegio philosophíæ curriculo, missus est primum Floréntiam, tum Montem Regalem; deínde Patavium ad sacram theologíam addiscéndam, ac póstea Lovanium, ubi concionatoris munere, nondum sacerdos, mirifice functus est. Lovanii præterea, sacerdotio auctus, theologíam ita docuit, ut plurimos hæreticos ad Ecclésiæ unitátem reduxerit, ac theologus per Europam claríssimus haberétur, eumque Sanctus Cárolus Mediolanénsis epíscopus aliique veheménter sibi expéterent.\par
\switchcolumn
\EN
\noindent Robert, a native of Montepulciano and of the noble family of Bellarmine, had for his mother the most pious Cynthia Cervini, sister of Pope Marcellus II. From the first he was conspicuous for exemplary piety and most chaste manners, earnestly desiring this one thing, to please God alone and to win souls to Christ. He attended the college of the Society of Jesus in his native town where he was highly commended for his intelligence and modesty. At the age of eighteen he entered the same Society at Rome, and was a model of all religious virtues. Having passed through the course of philosophy at the Roman College, he was sent first to Florence, then to Monreale, later to Padua to teach sacred theology, and afterwards to Louvain where, not yet a priest, he ably discharged the office of preacher. After ordination at Louvain, he taught theology with such success that he brought back many heretics to the unity of the Church, and was regarded throughout Europe as a most brilliant theologian; and St. Charles, Bishop of Milan, and others keenly sought after him.\par
\switchcolumn*

\LA
\begin{Resp}\Rsign Invéni David servum meum, óleo sancto meo unxi eum: \Star{} Manus enim mea auxiliábitur ei, allelúja.\end{Resp}
\begin{Resp}\Vsign Nihil profíciet inimícus in eo, et fílius iniquitátis non nocébit ei.\end{Resp}
\begin{Resp}\Rsign Manus enim mea auxiliábitur ei, allelúja.\end{Resp}
\switchcolumn
\EN
\begin{Resp}\Rsign I have found David My servant, with My holy oil have I anointed him \Star{} For My hand shall help him, alleluia.\end{Resp}
\begin{Resp}\Vsign The enemy shall prevail nothing against him, nor the son of wickedness afflict him.\end{Resp}
\begin{Resp}\Rsign For My hand shall help him, alleluia.\end{Resp}
\switchcolumn*

\LA
\LessonHead{Lectio V}
\switchcolumn
\EN
\LessonHead{Lesson V}
\switchcolumn*

\LA
\noindent Romam ex desidério Gregórii Papæ décimi tertii revocatus, theologicam controversiárum disciplínam trádidit in Collegio Romano: ibique vitæ spiritualis magister constitútus, angelicum juvenem Aloísium per sanctitátis sémitas moderátus est. Ipse Collegium Romanum ac deínde Neapolitanam Societátis Jesu Provinciam ad Sancti Ignatii mentem gubernávit. In Urbem íterum accersítus, a Cleménte octavo ad summa Ecclésiæ negotia, máximo cum christianæ rei emolumento, est adhibitus: tum invitus et frustra reluctans, in Cardinalium númerum cooptátus quia, ut palam asseruit ipse Pontifex, tunc non habebat parem Ecclésia Dei quod ad doctrinam. Ab eodem Pontifice consecrátus Epíscopus, Capuanam Archidiœcésim triennium sanctíssime administrávit: quo munere deposito, Romæ ad mortem usque degit, integerrimus ac fidelíssimus Summi Pontificis consiliarius. Multa præcláre scripsit, illud meritum adeptus in primis quod, sanctum Thomam ducem et magistrum secutus, de suórum necessitate témporum próvide cónscius, invicto doctrinæ róbore et amplíssima testimoniórum copia e Sacris Litteris et e Sanctórum Patrum ditíssimo fonte apte deprompta novos erróres debellávit, traditiónis catholicæ et Romani Pontificátus jurium strenuus præprimis adsertor. Complúribus étiam ad pietátem fovéndam libellis exstat insígnis ac præsertim aureo catechismo, quem licet aliis gravíssimis negotiis distentus, tum Capuæ tum Romæ púeros ac rudes docére non prætermittebat. Robertum æquævus Cardinalis a Deo missum judicávit, qui catholicos erudíret, pios cóleret, hæreticos profligaret; Sanctus Franciscus Salesius doctrinæ fontem hábuit; Summus Pontifex Benedíctus décimus quartus hæreticórum malleum dixit, ac Benedíctus décimus quintus catholicam religiónem propagántibus et tuéntibus exemplar indicávit.\par
\switchcolumn
\EN
\noindent Recalled to Rome at the wish of Pope Gregory XIII, he taught the science of controversial theology at the Roman College, and there, as spiritual director he guided the angelic youth Aloysius in the paths of holiness. He governed the Roman College and then the Neapolitan province of the Society of Jesus in accordance with the spirit of St. Ignatius. Again summoned to Rome, he was employed by Clement VIII in the most important affairs of the Church, with the greatest advantage to the Christian state; then against his will and in spite of opposition, he was admitted among the number of the cardinals, because, as the Pontiff publicly declared, he did not have his equal among theologians in the Church of God at the time. He was consecrated bishop by the same Pope, and administered the archdiocese of Capua in a most saintly manner for three years: having resigned this office, he lived in Rome until his death, as a most impartial and trusty counsellor to the Supreme Pontiff. He wrote much, and in an admirable manner. His principal merit lieth in his complete victory in the struggle against the new errors, during which he distinguished himself as a strenuous and outstanding vindicator of Catholic tradition and the rights of the Roman See. He gained this victory by following St. Thomas as his guide and teacher, by a prudent consideration of the needs of his times, by his irrefragable teaching, and by a most abundant wealth of testimony well-chosen from the sacred writings and from the very rich fountain of the Fathers of the Church. He is eminently noted for very numerous short works for fostering piety, and especially for that golden Catechism, which he never failed to explain to the young and ignorant both at Capua and at Rome, although preoccupied with other very important affairs. A contemporary cardinal declared that Robert was sent by God the instruction of Catholics, for the guidance of the good, and for the confusion of heretics; St. Francis de Sales regarded him as a fountain of learning; the Supreme Pontiff Benedict XIV called him the hammer of heretics; and Benedict XV proclaimed him the model of promoters and defenders of the Catholic religion.\par
\switchcolumn*

\LA
\begin{Resp}\Rsign Pósui adjutórium super poténtem, et exaltávi eléctum de plebe mea: \Star{} Manus enim mea auxiliábitur ei, allelúja.\end{Resp}
\begin{Resp}\Vsign Invéni David servum meum, óleo sancto meo unxi eum.\end{Resp}
\begin{Resp}\Rsign Manus enim mea auxiliábitur ei, allelúja.\end{Resp}
\switchcolumn
\EN
\begin{Resp}\Rsign I have laid help upon one that is mighty, and have exalted one chosen out of My people \Star{} For My hand shall help him, alleluia.\end{Resp}
\begin{Resp}\Vsign I have found David My servant, with My holy oil have I anointed him.\end{Resp}
\begin{Resp}\Rsign For My hand shall help him, alleluia.\end{Resp}
\switchcolumn*

\LA
\LessonHead{Lectio VI}
\switchcolumn
\EN
\LessonHead{Lesson VI}
\switchcolumn*

\LA
\noindent Vitæ religiosæ studiosíssimus, eam, inter purpuratos Patres adlectus, in exemplum servávit. Opes ultra necessarias nóluit; módico famulatu, ténui cultu habituque contentus: suórum non studuit opuléntiæ, ac vix adduci pótuit ut inópiam idéntidem levaret. De se humillime sensit, et mira fuit animi simplicitate. Deíparam diléxit unice: plures horas quotídie oratióni tribuebat. Parcíssime víctitans, ter in hebdomada jejunábat: in se constanter austérus, caritate in próximum flagrávit, vocátus sæpenúmero Pater páuperum. E baptismate innocéntiam ne vel levi quidem culpa macularet strenue contendit. Prope octogenarius, ad Sancti Andréæ in Colle Quirinali, extrémum in morbum íncidit, quem sólito virtútum fulgóre illustrávit. Moribundo Gregórius Papa décimus quintus et plures Cardinales adstitérunt, tantum Ecclésiæ cólumen éripi complorantes. Die sacrórum Stígmatum sancti Francisci, quorum memóriæ ubíque celebrandæ auctor fuerat, obdormívit in Dómino, anno millesimo sexcentésimo vigesimo primo. Mortuo tota cívitas parentávit, Sanctum uno ore conclámans. Eum vero Pius undecimus, Pontifex maximus, Beatórum primum ac deínde Sanctórum número adscripsit, et paulo post, ex Sacrórum Rituum Congregatiónis consulto, universalis Ecclésiæ Doctórem declarávit. Ejus corpus Romæ in templo Sancti Ignatii, apud sepúlcrum Sancti Aloísii, ut ipse optarat, pia veneratióne cólitur.\par
\switchcolumn
\EN
\noindent He was most zealous in the religious life and he maintained that manner of life after having been chosen as one of the empurpled cardinals. He did not want to any wealth beyond what was necessary; he was satisfied with a moderate household, and scanty fare and clothing. He did not strive to enrich his relatives, and he could scarcely be induced to relieve their poverty even occasionally. He had the lowest opinion of himself, and was of wonderful simplicity of soul. He had an extraordinary love for the Mother of God; he spent many hours daily in prayer. He ate very sparingly, and fasted three times a week. Uniformly austere with himself, he burned with charity towards his neighbour, and was often called the father of the poor. He earnestly strove that he might not stain his baptismal innocence to even the slightest fault. Almost eighty years old, he fell into his last illness at St. Andrew's on the Quirinal hill, and in it he shewed his usual radiant virtue. Pope Gregory XV and many cardinals visited him on his deathbed, lamenting the loss of such a great pillar of the Church. He fell asleep in the Lord in the year 1621, on the day of the sacred Stigmata of St. Francis, the memory of which he had been instrumental in having celebrated everywhere. The whole city mourned his death, unanimously proclaiming him a Saint. The Supreme Pontiff Pius XI inscribed his name, first, in the number of the Blessed, and then in that of the Saints, and shortly afterwards, by a decree of the Sacred Congregation of Rites, he declared him a Doctor of the universal Church. His body is honoured with pious veneration at Rome in the church of St. Ignatius, near the tomb of St. Aloysius, as he himself had desired.\par
\switchcolumn*

\LA
\begin{Resp}\Rsign Iste est, qui ante Deum magnas virtútes operátus est, et omnis terra doctrína ejus repléta est: \Star{} Ipse intercédat pro peccátis ómnium populórum, allelúja.\end{Resp}
\begin{Resp}\Vsign Iste est, qui contémpsit vitam mundi, et pervénit ad cæléstia regna.\end{Resp}
\begin{Resp}\Rsign Ipse intercédat pro peccátis ómnium populórum, allelúja.\end{Resp}
\begin{Resp}\Vsign Glória Patri, et Fílio, \Star{} et Spirítui Sancto.\end{Resp}
\begin{Resp}\Rsign Ipse intercédat pro peccátis ómnium populórum, allelúja.\end{Resp}
\switchcolumn
\EN
\begin{Resp}\Rsign This is he which wrought great wonders before God, and the whole earth is full of his teaching \Star{} May he pray for all people, that their sins may be forgiven unto them, alleluia.\end{Resp}
\begin{Resp}\Vsign This is he which loved not his life in this world, and hath attained unto the kingdom of heaven.\end{Resp}
\begin{Resp}\Rsign May he pray for all people, that their sins may be forgiven unto them, alleluia.\end{Resp}
\begin{Resp}\Vsign Glory be to the Father, and to the Son, \Star{} and to the Holy Ghost.\end{Resp}
\begin{Resp}\Rsign May he pray for all people, that their sins may be forgiven unto them, alleluia.\end{Resp}
\switchcolumn*

\LA
\LessonHead{Lectio VII}
\switchcolumn
\EN
\LessonHead{Lesson VII}
\switchcolumn*

\LA
\LessonTitle{Léctio sancti Evangélii secúndum Matthǽum}
\switchcolumn
\EN
\LessonTitle{From the Holy Gospel according to Matthew}
\switchcolumn*

\LA
\Cite{Matt 5:13-19}
\switchcolumn
\EN
\Cite{Matt 5:13-19}
\switchcolumn*

\LA
\noindent In illo témpore: Dixit Jesus discípulis suis: Vos estis sal terræ. Quod si sal evanúerit, in quo saliétur? Et réliqua.\par
\switchcolumn
\EN
\noindent At that time: Jesus said unto his disciples: Ye are the salt of the earth: But if the salt have lost his savour, wherewith shall it be salted? And so forth.\par
\switchcolumn*

\LA
\LessonTitle{Homilía sancti Roberti Bellarmino Epíscopi}
\switchcolumn
\EN
\LessonTitle{A Homily by St. Robert Bellarmine the Bishop}
\switchcolumn*

\LA
\Source{Concio 9, de probitate Doctórum Ecclésiæ; ab initio}
\switchcolumn
\EN
\Source{Concio 9, de probitate Doctorum Ecclesiae; ab initio}
\switchcolumn*

\LA
\noindent Quemádmodum in Deo, quem unum in Trinitáte et Trinum in unitate veneramur, tria quædam singuláriter eminent, poténtia, sapiéntia, bonitas; ita quoque, auditores, singulares amicos et fílios suos, Patres ac Doctores nostros, Deus, ut sibi quam simillimos et géntibus ómnibus suspiciéndos atque admirábiles redderet, potentíssimos, sapientíssimos, optimos sanctissimósque esse vóluit. Primum ea poténtia eos armávit, qua multa præter sólitum cursum ordinemque natúræ in elementis, in arbóribus, in brutis animántibus, in ipsis homínibus plane admirabília et singularia fácerent. Deínde sapiéntia ita mentes eórum instruxit, ut non solum præséntia et prætérita cernerent, sed étiam futura multo ante prævidérent atque prædícerent. Postremo dilatávit corda eórum summa atque ardentíssima caritate, tum ut ipsi magno animo opus aggrederéntur, tum ut ii qui converténdi per eos erant, non solum verbis et miraculis, sed étiam exemplis et vitæ probitate moverentur.\par
\switchcolumn
\EN
\noindent Just as in God, whom we venerate as one in the Trinity and three in the Unity, there are three things in particular which are especially clear: power, wisdom, and goodness; so also God, beloved listeners, that he might make his special friends and children, our fathers and teachers, very like unto himself and to be esteemed and admired by all nations, wished them to be in the highest degree powerful, wise, excellent, and holy. First, he furnished them with that power, by which they might do many evidently wonderful and extraordinary things, out of the usual course and order of nature, in regard to the elements, trees, brute beasts, and even to mankind. Then, he gave them such wisdom, that they saw not only the past and present, but they even foresaw the future long before, and predicted it. Finally, he enlarged their hearts with very great and burning charity, enabling them not only to enter whole-heartedly on their labours, but also to influence those whom they were about to convert, as well by their example and holy life, as by their preaching and miracles.\par
\switchcolumn*

\LA
\begin{Resp}\Rsign Amávit eum Dóminus, et ornávit eum: stolam glóriæ índuit eum, \Star{} Et ad portas paradísi coronávit eum, allelúja.\end{Resp}
\begin{Resp}\Vsign Induit eum Dóminus lorícam fídei, et ornávit eum.\end{Resp}
\begin{Resp}\Rsign Et ad portas paradísi coronávit eum, allelúja.\end{Resp}
\switchcolumn
\EN
\begin{Resp}\Rsign The Lord loved him and beautified him He clothed him with a robe of glory \Star{} And crowned him at the gates of Paradise, alleluia.\end{Resp}
\begin{Resp}\Vsign The Lord hath put on him the breast-plate of faith, and hath adorned him.\end{Resp}
\begin{Resp}\Rsign And crowned him at the gates of Paradise, alleluia.\end{Resp}
\switchcolumn*

\LA
\LessonHead{Lectio VIII}
\switchcolumn
\EN
\LessonHead{Lesson VIII}
\switchcolumn*

\LA
\noindent Prædicatores ígitur nostræ legis, tam ii qui primi ad nos fidem detulérunt et Evangélium, quam ii quos deínde singulis sæculis Deus excitávit ad fidem eamdem confírmandam vel propagandam, quales fúerint, quam pii, quam justi, quam religiosi totus mundus novit. Aspicite primum Apóstolos. Quid sublimius atque excellentius móribus Apostolicis? Aspicite deínde sanctos illos hómines, quos Patres et Doctores vocamus, lúmina illa claríssima, quæ Deus in firmaménto Ecclésiæ lucére vóluit, ut iis omnes hæreticórum ténebræ, dissiparéntur, ut Irenæum, Cyprianum, Hilárium, Athanasium, Basilíum, duos Gregórios, Ambrósium, Hierónymum, Augustínum, Chrysóstomum, Cyríllum. Vita et mores eórum nonne in iis monuméntis, quæ nobis reliquérunt, quasi in speculis quibusdam elucent? Nam ex abundántia cordis os lóquitur.\par
\switchcolumn
\EN
\noindent And so, the whole world knew how pious, how just, how religious were the preachers of our law, both those who first brought to us the faith and the Gospel, and those whom God thereafter stirred up in every age to confirm or propagate that same faith. And first, consider the Apostles. What could be better or more sublime than the Apostles' way of life? Next, consider those holy men whom we call Fathers and Doctors, those most shining lights which God hath ordained should shine in the firmament of the Church, that all the darkness of heresy might be dispersed, such as Irenaeus, Cyprian, Hilary, Athanasius, Basil, the two Gregories, Ambrose, Jerome, Augustine, Chrysostom, and Cyril. Do not their lives and conduct shine forth in the records, which they have left us, as in a special kind of mirror? For out of the fulness of the heart the mouth speaketh.\par
\switchcolumn*

\LA
\begin{Resp}\Rsign In médio Ecclésiæ apéruit os ejus, \Star{} Et implévit eum Dóminus spíritu sapiéntiæ et intelléctus, allelúja.\end{Resp}
\begin{Resp}\Vsign Jucunditátem et exsultatiónem thesaurizávit super eum.\end{Resp}
\begin{Resp}\Rsign Et implévit eum Dóminus spíritu sapiéntiæ et intelléctus, allelúja.\end{Resp}
\begin{Resp}\Vsign Glória Patri, et Fílio, \Star{} et Spirítui Sancto.\end{Resp}
\begin{Resp}\Rsign Et implévit eum Dóminus spíritu sapiéntiæ et intelléctus, allelúja.\end{Resp}
\switchcolumn
\EN
\begin{Resp}\Rsign In the midst of the congregation did the Lord open his mouth. \Star{} And filled him with the spirit of wisdom and understanding, alleluia.\end{Resp}
\begin{Resp}\Vsign He made him rich with joy and gladness.\end{Resp}
\begin{Resp}\Rsign And filled him with the spirit of wisdom and understanding, alleluia.\end{Resp}
\begin{Resp}\Vsign Glory be to the Father, and to the Son, \Star{} and to the Holy Ghost.\end{Resp}
\begin{Resp}\Rsign And filled him with the spirit of wisdom and understanding, alleluia.\end{Resp}
\switchcolumn*

\LA
\LessonHead{Lectio IX}
\switchcolumn
\EN
\LessonHead{Lesson IX}
\switchcolumn*

\LA
\noindent Quanta, óbsecro, apparet in libris Sanctórum Patrum cum summa eruditióne conjuncta humilitas? Quanta sobrietas? Nihil ibi obscenum, nihil turpe, nihil súbdolum, nihil arrogans, nihil inflátum. Quam multis modis Spíritus Sanctus, qui eórum péctora inhabitabat, in páginis eórum se prodit? Quis legere potest atténte Cyprianum, qui non statim ardeat amóre martyrii? Quis in Augustino diligenter versátus est, qui non profundíssimam didicerit humilitátem? Quis Hieronymum sæpe evolvit, qui non virginitátem et jejunium adamare incipiat? Spirant scripta Sanctórum religiónem, castitátem, integritátem, caritátem. Isti sunt ígitur epíscopi et pastores (ut verbis utar divi Augustíni) docti, graves, sancti, veritátis acerrimi defensores, qui catholicam fidem in lacte suxérunt, in cibo sumpsérunt: cujus lac et cibum parvis magnisque ministravérunt. Tálibus post Apóstolos sancta Ecclésia plantatóribus, rigatóribus, ædificatóribus, pastóribus, nutritóribus crevit.\par
\switchcolumn
\EN
\noindent Consider, I ask you, the humility, together with the most profound learning, which appeareth in the books of the holy fathers. What moderation! Nothing offensive there, nothing unseemly, no cunning, nothing assuming, nothing pompous. How the manifold working of the Holy Spirit, who dwelt in their hearts, revealeth itself in their pages! Who can read Cyprian attentively without immediately longing for martyrdom? Who can assiduously turn over the pages of Augustine without learning the most profound humility? Who can open Jerome frequently without beginning to love virginity and fasting? The writings of the saints breathe forth religion, chastity, integrity, and charity. Such then are our bishops and pastors (to use the words of the heavenly Augustine), learned men, eminent, holy, intelligent, defenders of the truth, who have taken in the Catholic faith as their milk, and have consumed it as food: and this milk and food they have ministered to great and small. Since the Apostles, holy Church hath flourished by such planters, waterers, builders, shepherds, and nurses.\par
\switchcolumn*

\LA
\TeDeum{Te Deum}
\switchcolumn
\EN
\TeDeum{Te Deum}
\switchcolumn*

\LA
\LessonHead{Lectio IX (pro commemoratione)}
\switchcolumn
\EN
\LessonHead{Lesson IX (when commemorated)}
\switchcolumn*

\LA
\LessonTitle{Commemoratio: S. Roberti Bellarmino Episcopi Confessoris et Ecclesiæ Doctoris}
\switchcolumn
\EN
\LessonTitle{Commemoration of: St. Robert Bellarmine, Bishop, Confessor and Doctor of the Church}
\switchcolumn*

\LA
\noindent Robértus, Politiánus, e patrícia Bellarminórum gente, matrem pientíssimam hábuit Cýnthiam Cervíni, Marcélli Papæ secúndi sorórem. Exímia pietáte et castíssimis móribus ornátus, duodevigínti annórum adoléscens Societátem Jesu Romæ ingréssus est, in eáque, ad mortem usque religiosárum virtútum ómnibus exémplo fuit. Post philosophíæ currículum Floréntiam primum missus, tum Montem Regálem, Patávium et Lovánium, magístri et concionatóris múnere, nondum sacérdos, mirífice functus est. Lovánii prætérea, sacerdótio auctus, theologíam ita dócuit, ut theólogus per Európam claríssimus jam tum haberétur. Romam revocátus, theológicam controversiárum disciplínam in Collégio Románo trádidit, ubi étiam, vitæ spirituális magíster constitútus, angélicum júvenem Aloísium per sanctitátis sémitas moderátus est. A Cleménte Papa octávo, frustra relúctans, in Patrum Cardinálium númerum cooptátus, et paulo post consecrátus Epíscopus, Capuánam Archidiœcésim triénnium sanctíssime rexit; quo múnere depósito, integérrimus ac fidelíssimus Summi Pontíficis consiliárius in Urbe degit, usque dum, prope octogenárius, die décima séptima Septémbris, anno millésimo sexcentésimo vigésimo primo pie in Dómino quiévit. Præter Controversiárum volúmina multa ália præcláre scripsit, inter quæ áureus catechésis libéllus exstat insígnis. Fortíssimum hunc cathólicæ veritátis propugnatórem Pius undécimus Póntifex Máximus in Sanctórum númerum rétulit atque universális Ecclésiæ Doctórem declarávit.\par
\switchcolumn
\EN
\noindent Robert was born at Montepulciano of the patrician family of Bellarmine, and had a most devout mother, Cynthia Cervini, the sister of Pope Marcellus II. Outstanding for his devotion and chastity, the young man entered the Society of Jesus at the age of eighteen, and until his death he was a model of the religious virtues for all his brethren. After his course in philosophy, he was first sent to Florence, then to Monreale, Padua, and Louvain; and he filled the offices of teacher and preacher in an admirable way, even though he was not yet a priest. Later he was ordained to the priesthood at Louvain; and he taught theology in such a way as soon to be considered the most famous theologian in all Europe. Recalled to Rome, he taught apologetics in the Roman College, where he was also appointed spiritual director, and led the angelic young Aloysius along the paths of holiness. Over his protests, he was made Cardinal by Pope Clement VIII, and soon after was consecrated bishop, ruling the archdiocese of Capua for three years in a most holy way. He resigned this office to become a counsellor of the highest integrity and loyalty to the Supreme Pontiff in Rome until, when he was almost eighty years old, on the 17th day of September, 1621, he died a holy death in the Lord. Besides his polemical works he wrote many other noteworthy books, among which his little golden Catechism remaineth famous. Pope Pius XI added the name of this strong defender of the Catholic faith to the number of the Saints and declared him a Doctor of the universal Church.\par
\switchcolumn*

\LA
\TeDeum{Te Deum}
\switchcolumn
\EN
\TeDeum{Te Deum}
\switchcolumn*

\end{paracol}

\par\needspace{10\baselineskip}\addvspace{1.2em}{\centering\small\textsc{Die 14 Maii} · \textsc{14 May}\par}
\Office{S. Bonifatii Martyris}{St. Boniface, Martyr}{Simplex}
\addcontentsline{toc}{subsection}{Die 14 Maii · S. Bonifatii Martyris\enspace\textit{St. Boniface, Martyr}}
\begin{paracol}{2}
\LA
\RefLine{Lectiones I–II: de Scriptura occurrente.}
\switchcolumn
\EN
\RefLine{Lessons I–II: from the Scripture of the day.}
\switchcolumn*

\LA
\LessonHead{Lectio III (pro commemoratione)}
\switchcolumn
\EN
\LessonHead{Lesson III (when commemorated)}
\switchcolumn*

\LA
\noindent Bonifatius, civis Romanus, quod cum Aglaë nobili matrona impudíce versátus esset, tanto illíus intemperantiæ dolóre captus est, ut pœniténtiæ causa se ad conquirenda et sepelienda Mártyrum córpora contulerit. Itaque, relictis peregrinatiónis sociis, cum Tarsi multos propter christianæ fidei professiónem variis tormentis cruciatos vidísset; illórum víncula osculátus, eos veheménter hortabátur, ut constanter supplícia perferrent, quod brevem labórem sempiterna requies consecutura sit. Comprehénsus ígitur, férreis ungulis excarnificátus est; cui étiam inter mánuum ungues et carnem acúti cálami sunt infixi, plumbumque liquefactum in os ejus infusum. Quibus in cruciátibus ea vox tantum Bonifatii audiebátur: Gratias tibi ago, Dómine Jesu Christe, Fili Dei. Mox in ollam ferventis picis demisso cápite conjéctus est; unde cum inviolátus exísset, ira incénsus judex eum secúri pércuti jubet. Quo témpore magnus terræmotus factus est, ita ut multi infidéles ad Christi Dómini fidem converteréntur. Eum sequénti die quæréntes socii, cum martyrio afféctum cognovíssent, quingentis solidis ejus corpus redemérunt, et condítum unguentis linteisque involutum Romam portándum curarunt. Quod factum cum ab Angelo Aglaë matrona, quæ et ipsa pœnitens se piis opéribus addixerat, cognovísset, prodiens obviam sancto corpori, ecclésiam ejus nómine ædificávit; in qua corpus sepúltum est Nonis Junii, cum ejus ánima pridie Idus Maji apud Tarsum Ciliciæ urbem migrasset in cælum, Diocletiáno et Maximiáno imperatóribus.\par
\switchcolumn
\EN
\noindent Boniface was a Roman citizen who had lived in sin with the noble lady Aglae. The memory of this transgression overwhelmed him with exceeding sorrow, so that for penance he gave himself up to look for and bury the bodies of the martyrs. While he was at Tarsus, and apart from his fellow-travellers, he saw a great many persons being diverse ways tormented, because they confessed to believing in Christ. He kissed their chains, and vehemently exhorted them bravely to bear their sufferings, seeing that the same their affliction which was but for a moment, was working for them an exceeding, even an eternal weight of glory. For this cause Boniface also was taken, and his flesh torn off him with iron claws. Sharp reeds also were driven between his finger-nails and the quick, and molten lead poured into his mouth. In his agony he was only heard to say “I thank thee, O Lord Jesus Christ, Son of God.” Afterward he was dipped head foremost into a vessel of boiling pitch, and as he was drawn out unharmed, the judge in fury commanded him to be beheaded. At the time it was done there was a great earthquake, whereby many unbelievers were turned to believe in the Lord Christ. The fellow-travellers of Boniface sought him the next day, and when they knew that he had undergone martyrdom, they bought his body for fifty shillings, and after that they had embalmed it with spices, and wrapped it in linen, they carried it to Rome. The Lady Aglae, who had herself with great contrition given up her life to godly works, was told by an angel what had come to pass. She therefore went forth to meet the holy body, and built a Church in the name of Boniface, wherein his said body was buried upon the fifth day of June next after that fourteenth of May whereon in the city of Tarsus in Cilicia, under the Emperors Diocletian and Maximian, he had passed away to heaven.\par
\switchcolumn*

\LA
\TeDeum{Te Deum}
\switchcolumn
\EN
\TeDeum{Te Deum}
\switchcolumn*

\end{paracol}

\par\needspace{10\baselineskip}\addvspace{1.2em}{\centering\small\textsc{Die 15 Maii} · \textsc{15 May}\par}
\Office{S. Joannis Baptistæ de la Salle Confessoris}{St. John Baptist de la Salle, Confessor}{Duplex}
\addcontentsline{toc}{subsection}{Die 15 Maii · S. Joannis Baptistæ de la Salle Confessoris\enspace\textit{St. John Baptist de la Salle, Confessor}}
\begin{paracol}{2}
\LA
\RefLine{Lectiones I–III: de Scriptura occurrente.}
\switchcolumn
\EN
\RefLine{Lessons I–III: from the Scripture of the day.}
\switchcolumn*

\LA
\LessonHead{Lectio IV}
\switchcolumn
\EN
\LessonHead{Lesson IV}
\switchcolumn*

\LA
\noindent Joánnes Baptista de la Salle, Rhemis claro genere ortus, puer adhuc móribus et factis in sortem Dómini se vocándum et sanctimoniæ laude honestándum portendit. Adoléscens in Rheménsi academía litteras ac philosophicas disciplínas didicit: quo témpore, etsi ob animi virtútes et álacre ingenium ac suave ómnibus carus esset, ab æqualium tamen societáte abhorrebat, ut solitudini addictus facílius Deo vacaret. In clericalem milítiam jampridem cooptatus, sexto décimo ætátis anno inter Rhemenses canonicos adscriptus est, Lutétiam Parisiórum, theologíæ in Sorbónica universitáte datúrus operam, contendit, atque in Sulpitianum seminárium adscítus est. At brevi paréntibus orbatus, domum régredi coactus, fratres educándos suscépit; quod, scientiárum interim sacrárum stúdia non intermittens, optimo cum fructu præstitit, uti éxitus comprobávit.\par
\switchcolumn
\EN
\noindent John Baptist de La Salle, born of an honourable family at Reims, when still a boy shewed by his manners and actions that he was called by destiny to the Lord, and was to be adorned with the excellence of holiness. As a youth he studied literature and the philosophical sciences at the academy at Rheims. During this time, although his mental powers and his lively and pleasant disposition endeared him to all, he nevertheless shrank from the company of his fellows, so that, being inclined to solitude, he might the more easily find time for God. Already having been for some time enlisted in the ranks of the clergy, he was enrolled among the canons of Reims at the age of sixteen years. He went to Paris to study theology at the university of the Sorbonne, and was admitted to the Sulpician seminary. But he was soon forced to return home because of the death of his parents, and undertook the education of his brothers, which he carried on, without meanwhile interrupting his sacred studies, and with the greatest success, as was proved by subsequent events.\par
\switchcolumn*

\LA
\begin{Resp}\Rsign Honéstum fecit illum Dóminus, et custodívit eum ab inimícis, et a seductóribus tutávit illum: \Star{} Et dedit illi claritátem ætérnam, allelúja.\end{Resp}
\begin{Resp}\Vsign Justum dedúxit Dóminus per vias rectas, et osténdit illi regnum Dei.\end{Resp}
\begin{Resp}\Rsign Et dedit illi claritátem ætérnam, allelúja.\end{Resp}
\switchcolumn
\EN
\begin{Resp}\Rsign The Lord made him honourable, and defended him from his enemies, and kept him safe from those that lay in wait for him; \Star{} And gave him perpetual glory, alleluia.\end{Resp}
\begin{Resp}\Vsign He went down with him into the pit, and left him not in bonds.\end{Resp}
\begin{Resp}\Rsign And gave him perpetual glory, alleluia.\end{Resp}
\switchcolumn*

\LA
\LessonHead{Lectio V}
\switchcolumn
\EN
\LessonHead{Lesson V}
\switchcolumn*

\LA
\noindent Sacerdotio demum auctus, qua præstanti fide animique ardore primum ad aram fecit, eisdem toto vitæ témpore Sacris est operatus. Interea, salútis animárum studio incénsus, totum in earumdem utilitátem sese impendit. Sorórum a Jesu infante, puellis educandis institutárum, regimen suscépit; easque non modo prudentíssime est moderátus, sed ab excidio vindicávit. Hinc porro ánimum advértit ad púeros de plebe religióne bonisque móribus informándos. Atque in hoc quidem illum suscitaverat Deus, ut scilicet, nova in Ecclésia sua religiosórum hóminum família cóndita, puerórum, præsertim páuperum, scholis perenni efficacíque ratióne consuleret. Demandátum vero a Dei providéntia munus, per contradictiónes plurimas magnasque ærúmnas, feliciter implévit, fundáta fratrum sodalitate, quam a Scholis christiánis nuncupávit.\par
\switchcolumn
\EN
\noindent He was finally ordained priest, and said his first Mass with the intense faith and ardour of the soul which, throughout his whole life, he brought to those holy Mysteries. Meanwhile, burning with zeal for the salvation of souls, he devoted himself wholly to their service. He undertook the direction of the Sisters of the Infant Jesus, founded for the education of girls; and not only managed them most prudently, but saved their institute from dissolution. From this time onwards, he turned his attention to the education of poor boys in religion and good morals. And God had raised him up for this very end, namely, that he should found in his Church a new family of religious men, and should look after boys' schools, especially of poor boys, with unceasing and efficient care. And, indeed, this duty, entrusted to him by Divine providence, was successfully accomplished, in spite of very much opposition and great hardships, by the foundation of an institute of brothers which he named the Christian Schools.\par
\switchcolumn*

\LA
\begin{Resp}\Rsign Amávit eum Dóminus, et ornávit eum: stolam glóriæ índuit eum, \Star{} Et ad portas paradísi coronávit eum, allelúja.\end{Resp}
\begin{Resp}\Vsign Índuit eum Dóminus lorícam fídei, et ornávit eum.\end{Resp}
\begin{Resp}\Rsign Et ad portas paradísi coronávit eum, allelúja.\end{Resp}
\switchcolumn
\EN
\begin{Resp}\Rsign The Lord loved him and beautified him; He clothed him with a robe of glory, \Star{} And crowned him at the gates of Paradise, alleluia.\end{Resp}
\begin{Resp}\Vsign The Lord hath put on him the breast-plate of faith, and hath adorned him.\end{Resp}
\begin{Resp}\Rsign And crowned him at the gates of Paradise, alleluia.\end{Resp}
\switchcolumn*

\LA
\LessonHead{Lectio VI}
\switchcolumn
\EN
\LessonHead{Lesson VI}
\switchcolumn*

\LA
\noindent Adjunctos ígitur sibi hómines in gravi opere et arduo, apud se primum suscépit; tum aptiori in sede constitutos disciplína sua optime imbuit iis légibus sapientibúsque institutis, quæ póstea a Benedicto décimo tertio sunt confírmata. Ex demissióne animi ac paupertátis amóre primum canonicatu se abdicávit, omniáque sua bona in páuperes erogávit; quin étiam serius, quod frustra sæpius tentaverat, fundáti a se institúti regimen sponte depósuit. Nihil tamen interim de fratrum sollicitúdine remíttens deque scholis ab eo, plúribus jam locis, apertis, impensius Deo vacare cœpit. Assidue jejuniis, flagellis aliisque asperitátibus in se ipsum sæviens, noctes orando ducebat. Donec, virtútibus ómnibus conspicuus, præsertim obediéntia, studio divinæ voluntátis implendæ, amóre ac devotióne in apostolicam Sedem; meritis onustus, sacramentis rite susceptis, obdormívit in Dómino annos natus duo de septuagínta. Eum Leo décimus tertius, Pontifex maximus, Beatórum catálogo inseruit; novisque fulgentem signis, anno jubilæi millesimo nongentésimo, Sanctórum honóribus decorávit. Pius vero duodecimus ómnium magistrórum púeris adulescentibúsque instituéndis præcipuum apud Deum cælestem Patronum constítuit.\par
\switchcolumn
\EN
\noindent His male associates in this great and arduous work he at first received into his own house; and then, establishing them in a more suitable dwelling, thoroughly inspired them with his method and with those wise laws and regulations which were afterwards confirmed by Benedict XIII. Because of humility and love of poverty, he first resigned his canonry and distributed all his property among the poor; and later also, after many unsuccessful attempts to do so, he of his own will resigned the government of the institute which he had founded. But meanwhile his solicitude for the brothers and for the schools which he opened in different places did not lessen, though he began to give himself more diligently to God. Shewing his hatred for self in constant fastings, in the use of the discipline and in other austerities, he spent his nights in prayer. At length, conspicuous for every kind of virtue, especially obedience, and zeal for fulfilling the divine will, and love and devotion to the Apostolic See, full of merit, and having devoutly received the sacraments, he fell asleep in the Lord in the sixty-eighth year of his age. The Supreme Pontiff Leo XIII placed him in the list of the Blessed; and, illustrious by new miracles, he was adorned with honours of the Saints in the year of jubilee, 1900. Pius XII appointed him the special heavenly patron of all teachers of boys and young men.\par
\switchcolumn*

\LA
\begin{Resp}\Rsign Iste homo perfécit ómnia quæ locútus est ei Deus, et dixit ad eum: Ingrédere in réquiem meam: \Star{} Quia te vidi justum coram me ex ómnibus géntibus, allelúja.\end{Resp}
\begin{Resp}\Vsign Iste est, qui contémpsit vitam mundi, et pervénit ad cæléstia regna.\end{Resp}
\begin{Resp}\Rsign Quia te vidi justum coram me ex ómnibus géntibus, allelúja.\end{Resp}
\begin{Resp}\Vsign Glória Patri, et Fílio, \Star{} et Spirítui Sancto.\end{Resp}
\begin{Resp}\Rsign Quia te vidi justum coram me ex ómnibus géntibus, allelúja.\end{Resp}
\switchcolumn
\EN
\begin{Resp}\Rsign This is he which did according unto all that God commanded him; and God said unto him: Enter thou into My rest, \Star{} For thee have I seen righteous before Me among all people, alleluia.\end{Resp}
\begin{Resp}\Vsign This is he which loved not his life in this world, and is come unto an everlasting kingdom.\end{Resp}
\begin{Resp}\Rsign For thee have I seen righteous before Me among all people, alleluia.\end{Resp}
\begin{Resp}\Vsign Glory be to the Father, and to the Son, \Star{} and to the Holy Ghost.\end{Resp}
\begin{Resp}\Rsign For thee have I seen righteous before Me among all people, alleluia.\end{Resp}
\switchcolumn*

\LA
\LessonHead{Lectio VII}
\switchcolumn
\EN
\LessonHead{Lesson VII}
\switchcolumn*

\LA
\LessonTitle{Léctio sancti Evangélii secúndum Matthǽum}
\switchcolumn
\EN
\LessonTitle{The Lesson is taken from the Holy Gospel according to Matthew}
\switchcolumn*

\LA
\Cite{Matt 18:1-5}
\switchcolumn
\EN
\Cite{Matt 18:1-5}
\switchcolumn*

\LA
\noindent In illo témpore: Accessérunt discípuli ad Jesum, dicéntes: Quis, putas, major est in regno cælórum? Et réliqua.\par
\switchcolumn
\EN
\noindent At that time: Came the disciples unto Jesus, saying: Who is the greatest in the kingdom of heaven? And so on, and that which followeth.\par
\switchcolumn*

\LA
\LessonTitle{Homilía sancti Joánnis Chrysostomi}
\switchcolumn
\EN
\LessonTitle{A Homily by St. John Chrysostom}
\switchcolumn*

\LA
\Source{In Cap. 18 Matth. Hom. 60}
\switchcolumn
\EN
\Source{In Cap. 18 Matth. Hom. 60}
\switchcolumn*

\LA
\noindent Vidéte ne aliquem istórum contempséritis parvulórum, quia eórum Angeli Patris mei fáciem semper aspíciunt, et quia ego propter eos veni, et hæc Patris mei volúntas est. Ad tuéndos conservándosque pusillos diligentiores nos reddit. Pérspicis quam ingéntia in tutelam tenúium mœnia erexerit, et quantum studium curamque habeat, ne perdántur; tum quia supremas despiciéntibus eos pœnas státuit, tum quia summam pollicétur mercedem his qui curam eórum suscipiunt, idque tam suo quam Patris exemplo corróborat.\par
\switchcolumn
\EN
\noindent Take heed that ye despise not one of these little ones, saith the Lord, for their Angels do always behold the face of my Father which is in heaven. It was as though he had said: It was for them that I came, because this is the will of my Father. Thus doth he make it our duty to be full of thought and care for the protection and safety of these little ones. Ye see what vast ramparts he hath built for the protection of the little ones, and how much care and trouble he hath taken that they shall not be lost; yea, he pronounceth extreme penalties on them that despise them; and for them that undertake to care wholeheartedly for them he promiseth most high rewards. And all these things of his teaching, he doth further enforce by his own example and by the example of the eternal Father himself.\par
\switchcolumn*

\LA
\begin{Resp}\Rsign Iste est qui ante Deum magnas virtútes operátus est, et de omni corde suo laudávit Dóminum: \Star{} Ipse intercédat pro peccátis ómnium populórum, allelúja.\end{Resp}
\begin{Resp}\Vsign Ecce homo sine queréla, verus Dei cultor, ábstinens se ab omni ópere malo, et pérmanens in innocéntia sua.\end{Resp}
\begin{Resp}\Rsign Ipse intercédat pro peccátis ómnium populórum, allelúja.\end{Resp}
\switchcolumn
\EN
\begin{Resp}\Rsign This is he which wrought great wonders before God, and praised the Lord with all his heart. \Star{} May he pray for all people, that their sins may be forgiven unto them, alleluia.\end{Resp}
\begin{Resp}\Vsign Behold a man without blame, a worshipper of God in truth, keeping himself clean from every evil work, and abiding still in his innocency.\end{Resp}
\begin{Resp}\Rsign May he pray for all people, that their sins may be forgiven unto them, alleluia.\end{Resp}
\switchcolumn*

\LA
\LessonHead{Lectio VIII}
\switchcolumn
\EN
\LessonHead{Lesson VIII}
\switchcolumn*

\LA
\noindent Dóminum ígitur étiam nos imitemur, et nihil pro frátribus omittamus, étiam eórum quæ humília viliáque nimium vidéntur. Sed, si administratióne nostra étiam opus fúerit, quamvis ténuis atque abjéctus quidem, cui administrándum sit, fúerit, quamvis ardua nobis res atque labóris plena esse videátur; ómnia hæc pro fratris salúte tolerabilióra facilióraque, oro, videántur. Tanto enim studio tantáque cura Deus dignam esse ánimam osténdit, ut neque Fílio suo pepercerit.\par
\switchcolumn
\EN
\noindent Let us therefore take to heart what the Lord saith, and imitate his example. Let us neglect nothing which it is in our power to do for these our little brothers and sisters. Let us be ready to undertake any work on their behalf, even the most homely and mean (that is, homely and mean in the eyes of men). And if there should be some further need of our assistance, even to the point of self-denying and laborious effort on our part, let us render it graciously; and let us do these things the more so when our help is required for one that is tiny and unloved and unwanted. And let us practice ourselves in these things until they become tolerable to us, and even easy, because we do them for the sake of one who is our little brother or sister in Christ. For God hath made evident that every soul is worthy of so much diligent care that he spared not his own Son for the sake of such.\par
\switchcolumn*

\LA
\begin{Resp}\Rsign Sint lumbi vestri præcíncti, et lucérnæ ardéntes in mánibus vestris: \Star{} Et vos símiles homínibus exspectántibus dóminum suum, quando revertátur a núptiis, allelúja.\end{Resp}
\begin{Resp}\Vsign Vigiláte ergo, quia nescítis qua hora Dóminus vester ventúrus sit.\end{Resp}
\begin{Resp}\Rsign Et vos símiles homínibus exspectántibus dóminum suum, quando revertátur a núptiis, allelúja.\end{Resp}
\begin{Resp}\Vsign Glória Patri, et Fílio, \Star{} et Spirítui Sancto.\end{Resp}
\begin{Resp}\Rsign Et vos símiles homínibus exspectántibus dóminum suum, quando revertátur a núptiis, allelúja.\end{Resp}
\switchcolumn
\EN
\begin{Resp}\Rsign Let your loins be girded about, and your lights burning; \Star{} And ye yourselves like unto men that wait for their lord, when he will return from the wedding, alleluia.\end{Resp}
\begin{Resp}\Vsign Watch therefore, for ye know not what hour your Lord doth come.\end{Resp}
\begin{Resp}\Rsign And ye yourselves like unto men that wait for their lord, when he will return from the wedding, alleluia.\end{Resp}
\begin{Resp}\Vsign Glory be to the Father, and to the Son, \Star{} and to the Holy Ghost.\end{Resp}
\begin{Resp}\Rsign And ye yourselves like unto men that wait for their lord, when he will return from the wedding, alleluia.\end{Resp}
\switchcolumn*

\LA
\LessonHead{Lectio IX}
\switchcolumn
\EN
\LessonHead{Lesson IX}
\switchcolumn*

\LA
\noindent Si non est nobis satis ad salútem, quod virtuose ipse vivamus, sed oportet aliórum salútem re ipsa desiderare; cum neque nos recte vivámus neque alios hortemur, quid respondébimus? quæ nobis spes salútis réliqua erit? Quid majus quam ánimis moderari, quam adolescentulórum fingere mores? Omni certe pictore, omni certe statuario, ceterisque hujusmodi ómnibus excellentiórem hunc duco, qui juvenum animos fingere non ignoret.\par
\switchcolumn
\EN
\noindent If it be not enough for our salvation that we ourselves live virtuously, but we should desire in very truth the salvation of others; what shall we answer, if we do not live right nor teach others? What hope of salvation will remain unto us? What is nobler than to rule minds or to mould the character of the young? I consider that he who knoweth how to form the youthful mind is truly greater than all painters, sculptors, and all others of that sort.\par
\switchcolumn*

\LA
\TeDeum{Te Deum}
\switchcolumn
\EN
\TeDeum{Te Deum}
\switchcolumn*

\LA
\LessonHead{Lectio IX (pro commemoratione)}
\switchcolumn
\EN
\LessonHead{Lesson IX (when commemorated)}
\switchcolumn*

\LA
\LessonTitle{Commemoratio: S. Joannis Baptistæ de la Salle Confessoris}
\switchcolumn
\EN
\LessonTitle{Commemoration of: St. John Baptist de la Salle, Confessor}
\switchcolumn*

\LA
\noindent Joánnes Baptísta de la Salle, Rhemis claro génere ortus, adoléscens in Rheménsi academía lítteras ac philosophíam dídicit. Clericáli milítiæ adscríptus, sextodécimo ætátis anno inter canónicos Rheménses adscítus fuit, et póstea Parísiis in Sulpitiánum seminárium recéptus. Sacerdótio auctus, Sorórum a Jesu infánte, quæ puéllis educándis incúmbunt, régimen suscépit, quas prudentíssime moderátus est ac deféndit. Púeris de plebe religióne bonísque móribus informándis, post plúrimas contradictiónes, fundávit fratrum sodalitátem, quam a Scholis christiánis nuncupávit, a Benedícto décimo tértio deínde confirmátam. Abdicáto canonicátu, suísque bonis in páuperes erogátis, et fundáti a se institúti regímine ex humilitáte demísso, virtútibus et méritis onústus, obdormívit in Dómino, annos natus duo de septuagínta. Eum Leo Papa décimus tértius primo Beatórum catálogo inséruit, dein in album Sanctórum rétulit. Pius vero duodécimus ómnium magistrórum púeris adulescentibúsque instituéndis præcípuum apud Deum cæléstem Patrónum constítuit.\par
\switchcolumn
\EN
\noindent John Baptist de La Salle was born of a noble family of Reims and as a young man studied literature and philosophy at the academy in that city. Having become a cleric, he was made a canon of Reims at the age of sixteen, and later he was received into the Sulpician seminary at Paris. After his ordination to the priesthood, he undertook the direction of the Sisters of the Child Jesus, whose work was the education of girls, and he guided and defended them most prudently. Turning his attention to the formation of boys from ordinary families in religious and good habits, after many difficulties he founded the Institute of Brothers of the Christian Schools. This institute was confirmed by Benedict XIII. John Baptist resigned his canonry, gave away his goods to the poor, and in humility also resigned the governance of the institute he had founded. Full of virtues and merits, he fell asleep in the Lord in the sixty-eighth year of his age. Pope Leo XIII first included him in the list of the Blessed and then in that of the Saints, and Pius XII appointed him the special heavenly patron of all teachers of boys and young men.\par
\switchcolumn*

\LA
\TeDeum{Te Deum}
\switchcolumn
\EN
\TeDeum{Te Deum}
\switchcolumn*

\end{paracol}

\par\needspace{10\baselineskip}\addvspace{1.2em}{\centering\small\textsc{Die 16 Maii} · \textsc{16 May}\par}
\Office{S. Ubaldi Episcopi et Confessoris}{St. Ubald, Bishop and Confessor}{Semiduplex}
\addcontentsline{toc}{subsection}{Die 16 Maii · S. Ubaldi Episcopi et Confessoris\enspace\textit{St. Ubald, Bishop and Confessor}}
\begin{paracol}{2}
\LA
\RefLine{Lectiones I–III: de Scriptura occurrente.}
\switchcolumn
\EN
\RefLine{Lessons I–III: from the Scripture of the day.}
\switchcolumn*

\LA
\RefLine{Lectiones VII–IX: ex Commúni unius Confessoris Pontificis (C4).}
\switchcolumn
\EN
\RefLine{Lessons VII–IX: from the Common of a Confessor Bishop (C4).}
\switchcolumn*

\LA
\LessonHead{Lectio IV}
\switchcolumn
\EN
\LessonHead{Lesson IV}
\switchcolumn*

\LA
\noindent Ubaldus Eugubii in Umbria nobili genere natus, a primis annis pietáte et litteris egregie est institútus; jamque adoléscens, ut uxórem duceret, sæpe tentatus, numquam tamen a proposito servandæ virginitátis recéssit. Sacerdos effectus, patrimónium suum paupéribus et ecclésiis distríbuit, et canonicórum regulárium ordinis sancti Augustíni institutum suscípiens, illud in pátriam tránstulit, atque in eo aliquámdiu sanctíssime vixit. Cujus sanctitátis opinióne evulgata, ab Honorio secundo summo Pontifice ecclésiæ Eugubinæ invítus præfícitur, et episcopalis consecratiónis munere decorátur.\par
\switchcolumn
\EN
\noindent Ubald was born of a noble family at Gubbio in Umbria, and well established in godliness and learning from his earliest years. When he was a young man, it was often proposed to him to marry, but he never abandoned his determination to preserve his virginity. After that he was ordained Priest he divided his inheritance among the poor and Churches, and embraced the Institute of Canons Regular of St. Augustine. This Institute he brought to Gubbio, and for some time led therein a most holy life. When the fame of his saintliness had got noised abroad, Pope Honorius II set him, contrary to his own wishes, over the Church of Gubbio, and he was honoured with consecration as Bishop by the hands of the said Pope himself, (in the year of our Lord 1129.)\par
\switchcolumn*

\LA
\begin{Resp}\Rsign Invéni David servum meum, óleo sancto meo unxi eum: \Star{} Manus enim mea auxiliábitur ei, allelúja.\end{Resp}
\begin{Resp}\Vsign Nihil profíciet inimícus in eo, et fílius iniquitátis non nocébit ei.\end{Resp}
\begin{Resp}\Rsign Manus enim mea auxiliábitur ei, allelúja.\end{Resp}
\switchcolumn
\EN
\begin{Resp}\Rsign I have found David My servant, with My holy oil have I anointed him \Star{} For My hand shall help him, alleluia.\end{Resp}
\begin{Resp}\Vsign The enemy shall prevail nothing against him, nor the son of wickedness afflict him.\end{Resp}
\begin{Resp}\Rsign For My hand shall help him, alleluia.\end{Resp}
\switchcolumn*

\LA
\LessonHead{Lectio V}
\switchcolumn
\EN
\LessonHead{Lesson V}
\switchcolumn*

\LA
\noindent Ad suam ítaque revertens ecclésiam, cum de consueta vivéndi ratióne nihil ádmodum immutasset, in omni tamen virtútum genere eo magis eminére cœpit, quo efficacius aliórum étiam salútem verbo et exemplo procuraret, factus forma gregis ex animo. Nam victu parco, vestítu moderáto, léctulo áspero et paupérrimo, crucis mortificatiónem jugiter in suo corpore circumferebat, dum inexplébili oratiónis studio spíritum quotídie recrearet. Hinc admirábilem illam mansuetúdinem est adeptus, qua gravíssimas injurias et contumelias non modo æquanimiter tulit, verum étiam mirifico dilectiónis afféctu persecutores suos omni benignitátis testimonio complectebátur.\par
\switchcolumn
\EN
\noindent When Ubald came to live as Bishop in Gubbio, he changed his way of life in no wise from that which he had led before, but his virtues began to be more eminent because his word and example were now more able to benefit his neighbours, to whom the shepherd of their souls was a pattern, not by outward showing only, but from his heart. He ate little, dressed simply, and slept upon a hard and very poor bed. He “always bore in the body the dying of the Lord Jesus,” (2 Cor. iv. 10,) while he daily fed his soul in unceasing and earnest prayer. Hence he acquired such wonderful meekness, that when he was most grievously wronged and insulted he not only took it patiently, but, by a strange impulse of love for them, embraced his persecutors with every proof of affection.\par
\switchcolumn*

\LA
\begin{Resp}\Rsign Pósui adjutórium super poténtem, et exaltávi eléctum de plebe mea: \Star{} Manus enim mea auxiliábitur ei, allelúja.\end{Resp}
\begin{Resp}\Vsign Invéni David servum meum, óleo sancto meo unxi eum.\end{Resp}
\begin{Resp}\Rsign Manus enim mea auxiliábitur ei, allelúja.\end{Resp}
\switchcolumn
\EN
\begin{Resp}\Rsign I have laid help upon one that is mighty, and have exalted one chosen out of My people \Star{} For My hand shall help him, alleluia.\end{Resp}
\begin{Resp}\Vsign I have found David My servant, with My holy oil have I anointed him.\end{Resp}
\begin{Resp}\Rsign For My hand shall help him, alleluia.\end{Resp}
\switchcolumn*

\LA
\LessonHead{Lectio VI}
\switchcolumn
\EN
\LessonHead{Lesson VI}
\switchcolumn*

\LA
\noindent Biennio, antequam ex hac vita migraret, cum diutinis afflicetarétur infirmitátibus, inter acerbíssimos corporis cruciátus velut aurum in fornace purgátum, Deo grátias indesinenter agebat. Adveniénte autem sacro Pentecostes die, cum multis annis ecclésiam sibi commissam cumma cum laude gubernasset, sanctis opéribus ac miraculis clarus quiévit in pace. Quem Cælestinus Papa tertius in Sanctórum númerum retulit. Ejus virtus præcipue in effugandis spirítibus immundis elucet. Corpus vero, per tot sæcula incorruptum, magna fidelium veneratióne in pátria colitur, quam non semel a præsénti discrimine liberávit.\par
\switchcolumn
\EN
\noindent Just the space of two years before Ubald passed away from this present life, he was tried as gold in the furnace, by grievous bodily weakness, and, day after day, amid the sharpest sufferings, he never ceased patiently to give God thanks. He rested in peace on the sacred day of Pentecost, (in the year 1160,) having for many years governed with great praise the Church which had been entrusted to him, and glorious for good works and miracles. Pope Celestine III. numbered him with the Saints. His strength is most chiefly shown in the casting out of evil spirits. His body hath remained without corruption for all these ages, and is reverenced greatly in his native town by Christ's faithful people. To them he hath more than once shown himself good at need.\par
\switchcolumn*

\LA
\begin{Resp}\Rsign Iste est, qui ante Deum magnas virtútes operátus est, et omnis terra doctrína ejus repléta est: \Star{} Ipse intercédat pro peccátis ómnium populórum, allelúja.\end{Resp}
\begin{Resp}\Vsign Iste est, qui contémpsit vitam mundi, et pervénit ad cæléstia regna.\end{Resp}
\begin{Resp}\Rsign Ipse intercédat pro peccátis ómnium populórum, allelúja.\end{Resp}
\begin{Resp}\Vsign Glória Patri, et Fílio, \Star{} et Spirítui Sancto.\end{Resp}
\begin{Resp}\Rsign Ipse intercédat pro peccátis ómnium populórum, allelúja.\end{Resp}
\switchcolumn
\EN
\begin{Resp}\Rsign This is he which wrought great wonders before God, and the whole earth is full of his teaching \Star{} May he pray for all people, that their sins may be forgiven unto them, alleluia.\end{Resp}
\begin{Resp}\Vsign This is he which loved not his life in this world, and hath attained unto the kingdom of heaven.\end{Resp}
\begin{Resp}\Rsign May he pray for all people, that their sins may be forgiven unto them, alleluia.\end{Resp}
\begin{Resp}\Vsign Glory be to the Father, and to the Son, \Star{} and to the Holy Ghost.\end{Resp}
\begin{Resp}\Rsign May he pray for all people, that their sins may be forgiven unto them, alleluia.\end{Resp}
\switchcolumn*

\LA
\LessonHead{Lectio IX (pro commemoratione)}
\switchcolumn
\EN
\LessonHead{Lesson IX (when commemorated)}
\switchcolumn*

\LA
\LessonTitle{Commemoratio: S. Ubaldi Episcopi et Confessoris}
\switchcolumn
\EN
\LessonTitle{Commemoration of: St. Ubald, Bishop and Confessor}
\switchcolumn*

\LA
\noindent Ubáldus, Eugúbii in Umbria nóbili génere natus, a primis annis pietáte et lítteris egrégie est institútus; jamque adoléscens, ut uxórem dúceret, sæpe tentátus, nunquam tamen a propósito servándæ virginitátis recéssit. Sacérdos efféctus, patrimónium suum paupéribus et ecclésiis distríbuit. Canonicórum regulárium órdinis sancti Augustíni institútum suscípiens, illud in pátriam tránstulit. Ab Honório secúndo summo Pontífice ecclésiæ Eugubínæ invítus præfícitur, et episcopális consecratiónis múnere decorátur. Factus forma gregis ex ánimo, de consuéta vivéndi ratióne nihil ádmodum immutávit, et in omni virtútum génere enítuit. Diútius infirmitátibus afflíctus, Deo indesinénter grátias agébat. Cum multis annis ecclésiam sibi commíssam summa cum laude gubernásset, sanctis opéribus et miráculis clarus, quiévit in pace.\par
\switchcolumn
\EN
\noindent St. Ubald was born of a noble family at Gubbio in Umbria, and was well trained, from his earliest years, in devotion and in letters. As a young man, he was often urged to marry, but he never gave up his determination to preserve his virginity. When he was made a priest, he distributed his patrimony to the poor and to churches. He entered the Institute of the Canons Regular of St. Augustine and brought it to his own country. Against his will, he was put over the Church of Gubbio by Pope Honorius II and given episcopal consecration. Making himself with his whole heart a model for his flock, he changed nothing in his way of life and became illustrious for every kind of virtue. Although he suffered from various infirmities for a long time, he never ceased to give thanks to God. When he had governed the Church entrusted to him for many years in a way deserving of the highest praise, he fell asleep in peace, famous for his works and miracles.\par
\switchcolumn*

\LA
\TeDeum{Te Deum}
\switchcolumn
\EN
\TeDeum{Te Deum}
\switchcolumn*

\end{paracol}

\par\needspace{10\baselineskip}\addvspace{1.2em}{\centering\small\textsc{Die 17 Maii} · \textsc{17 May}\par}
\Office{S. Paschalis Baylon Confessoris}{St. Pascal Baylon Confessor}{Duplex}
\addcontentsline{toc}{subsection}{Die 17 Maii · S. Paschalis Baylon Confessoris\enspace\textit{St. Pascal Baylon Confessor}}
\begin{paracol}{2}
\LA
\RefLine{Lectiones I–III: de Scriptura occurrente.}
\switchcolumn
\EN
\RefLine{Lessons I–III: from the Scripture of the day.}
\switchcolumn*

\LA
\RefLine{Lectiones VII–IX: ex Commúni Confessoris non pontificis (C5).}
\switchcolumn
\EN
\RefLine{Lessons VII–IX: from the Common of a Confessor not a Bishop (C5).}
\switchcolumn*

\LA
\LessonHead{Lectio IV}
\switchcolumn
\EN
\LessonHead{Lesson IV}
\switchcolumn*

\LA
\noindent Paschalis Baylon, paupéribus piisque paréntibus in oppido Turris Formosæ, Seguntinæ diœcesis, in Aragónia natus, a téneris annis plura dedit futuræ sanctitátis indicia. Sortítus ánimam bonam ac rerum cæléstium appríme studiosam, puerítiam atque adolescéntiam in gregis custódia transegit; quam ille vivéndi ratiónem ideo præcipue diligebat, quod humilitáti fovendæ ac innocéntiæ conservandæ in primis utilem atque opportunam judicáret. Erat in victu modicus, in oratióne assiduus, tantáque apud coævos et socios florebat auctoritate et grátia, ut eórum lites compónens, erróres corrigens, ignorantiam erudiens ac desidiam éxcitans, velut ómnium parens et magister máximo studio colerétur ac amarétur: beátus étiam tum a plerisque appellatus.\par
\switchcolumn
\EN
\noindent Paschal Baylon was the son of poor and godly parents, in the town of Torre Hermosa, and Diocese of Sagunta in Aragon, (in the year of our Lord 1540.) From his childhood he gave indications of a holy life. He was naturally of a good disposition, and very wishful to learn about heavenly things. His boyhood and youth he passed in the occupation of a shepherd. This way of life pleased him well, because he thought it one useful and fitted to nourish lowliness and keep innocency. He ate little, and was instant in prayer. He had great weight and favour with his fellows and neighbours, whose quarrels he healed, corrected their mistakes, enlightened their ignorance, and roused them from idleness. They all greatly honoured and loved him, as though he were their father and teacher, and even then many called him “Beato,” that is “the Blessed.”\par
\switchcolumn*

\LA
\begin{Resp}\Rsign Honéstum fecit illum Dóminus, et custodívit eum ab inimícis, et a seductóribus tutávit illum: \Star{} Et dedit illi claritátem ætérnam, allelúja.\end{Resp}
\begin{Resp}\Vsign Justum dedúxit Dóminus per vias rectas, et osténdit illi regnum Dei.\end{Resp}
\begin{Resp}\Rsign Et dedit illi claritátem ætérnam, allelúja.\end{Resp}
\switchcolumn
\EN
\begin{Resp}\Rsign The Lord made him honourable, and defended him from his enemies, and kept him safe from those that lay in wait for him; \Star{} And gave him perpetual glory, alleluia.\end{Resp}
\begin{Resp}\Vsign He went down with him into the pit, and left him not in bonds.\end{Resp}
\begin{Resp}\Rsign And gave him perpetual glory, alleluia.\end{Resp}
\switchcolumn*

\LA
\LessonHead{Lectio V}
\switchcolumn
\EN
\LessonHead{Lesson V}
\switchcolumn*

\LA
\noindent Qui vero in sæculo, terra nempe desérta et inaquósa, adeo feliciter adoléverat, flos convallium, plantátus in domo Dómini, mirum ubíque sparsit sanctitátis odórem. Igitur Paschalis arrepto vitæ severioris instituto, atque in ordine Minórum strictioris observantiæ Discalceatórum cooptatus, exsultávit ut gigas ad curréndam viam suam; totumque se Dómino excoléndum tradens, dies noctesque cogitabat, qua se ratióne magis ei magisque conformaret. Ita factum est brevi, ut eum, tamquam seraphicæ perfectiónis exemplar, ipsi quoque provectiores imitándum sibi propónerent. Ipse autem in humili serviéntium gradu constitútus, se velut ómnium peripséma réputans, ardua quæque et abjecta domus ministeria, véluti jure quodam peculiari sibi débita, summa cum hilaritate suscipiebat, et exercebat humilitate ac patiéntia pari. Carnem spiritui quandoque reluctari nitentem jugi maceratióne afflictávit, atque in servitútem redégit; spíritum vero assidua sui abnegatióne ferventiórem in dies ad anterióra extendebat.\par
\switchcolumn
\EN
\noindent In a world which was to him “a dry land, where no water is” the vallies, “planted in the House of the Lord” (Ps. xci. 14,) whose strange sweetness spread all around. When he took upon him an harder life, by entering the Institute of barefooted Grey Friars, of the strict Observance, “he rejoiced as a strong man to run a race” (Ps. xviii. 6,) and gave himself up altogether to serve the Lord, thinking by day and by night only how he might attain more and more to have that mind in him which was also in Christ Jesus (Phil. ii. 5.) And so it came to pass in a little while, that his very elders set him before them for their model, as a pattern of a man seeking to be perfect in the path of the Seraphic Order. Paschal himself held the lowly place of a lay brother, and deemed himself “the off-scouring of all things” (1. Cor. iv. 13.) He took most cheerfully, and discharged with the greatest humility and patience, the hardest and meanest work of the house, as though such were his peculiar right. His flesh would sometimes rebel against his spirit, but he broke it under the yoke of mortification, and brought it into subjection. Day by day the spirit of self-denial waxed stronger in him, and “forgetting those things which were behind, he reached forth unto those things which were before” (Phil. iii. 13.)\par
\switchcolumn*

\LA
\begin{Resp}\Rsign Amávit eum Dóminus, et ornávit eum: stolam glóriæ índuit eum, \Star{} Et ad portas paradísi coronávit eum, allelúja.\end{Resp}
\begin{Resp}\Vsign Índuit eum Dóminus lorícam fídei, et ornávit eum.\end{Resp}
\begin{Resp}\Rsign Et ad portas paradísi coronávit eum, allelúja.\end{Resp}
\switchcolumn
\EN
\begin{Resp}\Rsign The Lord loved him and beautified him; He clothed him with a robe of glory, \Star{} And crowned him at the gates of Paradise, alleluia.\end{Resp}
\begin{Resp}\Vsign The Lord hath put on him the breast-plate of faith, and hath adorned him.\end{Resp}
\begin{Resp}\Rsign And crowned him at the gates of Paradise, alleluia.\end{Resp}
\switchcolumn*

\LA
\LessonHead{Lectio VI}
\switchcolumn
\EN
\LessonHead{Lesson VI}
\switchcolumn*

\LA
\noindent Deíparam Vírginem, cujus clientelæ se ab ineunte ætate dicáverat, tamquam matrem quotidiánis colebat obsequiis atque filiali exorábat fiducia. Porro erga sanctíssimum Eucharistiæ sacraméntum difficile dictu est quam ardénti tenerétur devotiónis afféctu; quem defunctus étiam in cadávere retinére visus est, dum, jacens in féretro, ad sacræ Hostiæ elevatiónem bis óculos reserávit et clausit, magna ómnium, qui áderant, admiratióne. Ejusdem veritátem inter hæreticos publice palamque proféssus, multa et gravia ob eam causam perpéssus est; crebro étiam ad necem petítus, sed singulari Dei providéntia ab impiórum mánibus ereptus. Sæpe inter orándum ómnibus destituebátur sensibus, dulcique languebat amoris deliquio: quo témpore cælestem illam sciéntiam hausisse créditus est, qua, homo rudis et illiteratus, de mystériis fidei difficillimis respondére, atque áliquot étiam libros conscríbere pótuit. Denique meritis plenus, eádem qua prædixerat hora, feliciter migrávit ad Dóminum, anno salútis millesimo quingentésimo nonagesimo secundo, sexto décimo Kalendas Junii, eodem, quo natus fuerat, festo Pentecostes recurrente, annum agens secúndum supra quinquagesimum. Quibus aliisque virtútibus insignem, ac miraculis tam in vita quam post mortem clarum, Paulus quintus Pontifex maximus illum Beátum appellávit; Alexander autem octavus Sanctórum catálogo adscripsit; tandem Leo décimus tertius peculiarem cœtuum eucharisticórum, item societátum ómnium a sanctíssima Eucharistia, sive quæ in posterum futuræ sunt, patronum cælestem declarávit et constítuit.\par
\switchcolumn
\EN
\noindent The Virgin Mother of God he had vowed himself when he was but a little lad, and he paid her every day the services of a son, and trusted her as a mother. It is hard to tell how intense was the love which bound him to the Most Holy Sacrament of the Eucharist, a love which seemed literally stronger than death, for when his dead body was found lying on the bier, its eyes opened and shut twice when the Sacred Host was lifted up, to the amazement of all that were there. When he was among heretics, he suffered much and grievously at their hands for plainly and openly telling the truth touching this Sacrament they often sought after him to murder him, but by the singular Providence of God he was delivered from those wicked men. When he was at prayer he often became utterly insensible, and his soul fainted away with the love of God. During these trances it was believed that he received directly from heaven that knowledge which he had, and which enabled him, although a man altogether rough and unlettered, to answer the hardest questions upon the mysteries of the faith, and even to write some books. At last, full of good works, he joyfully passed away to be ever with the Lord, at the hour foretold by himself, on the Feast of Pentecost, the 17th day of May, in the year of salvation 1592, on which day also he had been born fifty -two years before. Illustrious for the graces above mentioned, and for the miracles which he worked both during his life and after his death, he was named Blessed by Pope Paul V., and Alexander VIII. enrolled him among the Saints.\par
\switchcolumn*

\LA
\begin{Resp}\Rsign Iste homo perfécit ómnia quæ locútus est ei Deus, et dixit ad eum: Ingrédere in réquiem meam: \Star{} Quia te vidi justum coram me ex ómnibus géntibus, allelúja.\end{Resp}
\begin{Resp}\Vsign Iste est, qui contémpsit vitam mundi, et pervénit ad cæléstia regna.\end{Resp}
\begin{Resp}\Rsign Quia te vidi justum coram me ex ómnibus géntibus, allelúja.\end{Resp}
\begin{Resp}\Vsign Glória Patri, et Fílio, \Star{} et Spirítui Sancto.\end{Resp}
\begin{Resp}\Rsign Quia te vidi justum coram me ex ómnibus géntibus, allelúja.\end{Resp}
\switchcolumn
\EN
\begin{Resp}\Rsign This is he which did according unto all that God commanded him; and God said unto him: Enter thou into My rest, \Star{} For thee have I seen righteous before Me among all people, alleluia.\end{Resp}
\begin{Resp}\Vsign This is he which loved not his life in this world, and is come unto an everlasting kingdom.\end{Resp}
\begin{Resp}\Rsign For thee have I seen righteous before Me among all people, alleluia.\end{Resp}
\begin{Resp}\Vsign Glory be to the Father, and to the Son, \Star{} and to the Holy Ghost.\end{Resp}
\begin{Resp}\Rsign For thee have I seen righteous before Me among all people, alleluia.\end{Resp}
\switchcolumn*

\LA
\LessonHead{Lectio IX (pro commemoratione)}
\switchcolumn
\EN
\LessonHead{Lesson IX (when commemorated)}
\switchcolumn*

\LA
\LessonTitle{Commemoratio: S. Paschalis Baylon Confessoris}
\switchcolumn
\EN
\LessonTitle{Commemoration of: St. Pascal Baylon Confessor}
\switchcolumn*

\LA
\noindent Paschális Baylon, paupéribus piísque paréntibus in óppido Turris Formósæ in Aragónia natus puerítiam atque adolescéntiam in gregis custódia transégit. Póstea severióris vitæ institútum ampléxus, et órdini fratrum Minórum adscríptus, júgiter cogitábat, qua se ratióne magis magísque Christo crucifíxo conformáret. Deíparam Vírginem, cujus clientélæ se ab ineúnte ætáte dicáverat, tamquam matrem filiálibus et cotidiánis obséquiis colébat. Erga Eucharístiam ténero et assíduo flagrávit devotiónis afféctu, quem defúnctus étiam retinére visus est, dum jacens in féretro, ad sacræ Hóstiæ elevatiónem bis óculos reserávit et clausit, magna ómnium, qui áderant, admiratióne. Méritis plenus migrávit ad Dóminum, anno millésimo quingentésimo nonagésimo secúndo. Eum Leo décimus tértius peculiárem cœ́tuum eucharisticórum, societatúmque ómnium a sanctíssima Eucharístia patrónum cæléstem declarávit et constítuit.\par
\switchcolumn
\EN
\noindent Paschal Baylon was born of poor and devout parents in the town of Torre Hermosa in Aragon, and spent his boyhood and youth in herding sheep. When he had embraced an austere way of life by entering the Friars Minor, he thought ceaselessly of how he might conform ever more closely to Christ crucified. With daily filial devotion he honoured as mother the Virgin Mother of God, to whom he had consecrated himself from his earliest years. He burned with a tender and steady love for the Eucharist, a love which he manifested even in death, when, lying on his bier, he opened and closed his eyes twice at the elevation of the sacred Host, to the astonishment of all present. Full of merits, he went to the Lord in the year 1592. Leo XIII declared and appointed him the special heavenly patron of all Eucharistic congresses and of all associations in honour of the most holy Eucharist.\par
\switchcolumn*

\LA
\TeDeum{Te Deum}
\switchcolumn
\EN
\TeDeum{Te Deum}
\switchcolumn*

\end{paracol}

\par\needspace{10\baselineskip}\addvspace{1.2em}{\centering\small\textsc{Die 18 Maii} · \textsc{18 May}\par}
\Office{S. Venantii Martyris}{St. Venantius, Martyr}{Duplex}
\addcontentsline{toc}{subsection}{Die 18 Maii · S. Venantii Martyris\enspace\textit{St. Venantius, Martyr}}
\begin{paracol}{2}
\LA
\RefLine{Lectiones I–III: de Scriptura occurrente.}
\switchcolumn
\EN
\RefLine{Lessons I–III: from the Scripture of the day.}
\switchcolumn*

\LA
\RefLine{Lectiones VII–IX: ex Commúni unius Martyris Tempore Paschali (C2ap).}
\switchcolumn
\EN
\RefLine{Lessons VII–IX: from the Common of a Single Martyr in Paschaltide (C2ap).}
\switchcolumn*

\LA
\LessonHead{Lectio IV}
\switchcolumn
\EN
\LessonHead{Lesson IV}
\switchcolumn*

\LA
\noindent Venantius Camers, quindecim annos natus, cum christianæ religiónis accusarétur apud Antíochum, qui sub Decio imperatóre Camerino præerat, in porta civitátis præsidi se obtulit. Quem ille pollicitatiónibus ac terróribus diu tentátum, flagris cædi et vinculis astringi jussit; sed is mirabíliter ab Angelo solutus, lampadibus póstea aduritur, atque inverso ore, fumo supposito, suspénditur. Ejus constantiam in tormentis demirátus Anastasius Cornicularius, et quod eum ab Angelo íterum solutum cándida veste supra fumum ambulántem vidísset, in Christum credidit; et a beato Porphyrio presbytero cum família baptizatus, paulo post martyrii palmam cum eodem promeruit.\par
\switchcolumn
\EN
\noindent Venantius was a lad of Camerino (in the neighbourhood of Ancona,) who at fifteen years of age was accused of Christianity before Antiochus, Praefect of Camerino under the Emperor Decius. Venantius therefore appeared before Antiochus at the gate of the city, and when the Praefect had striven with him for a long while, by promises and threats, he commanded him to be scourged and thrown into irons, but an Angel loosed his bonds. He was afterwards scarified with lamps, and hung head downwards in smoke. Anastasius the trumpeter was amazed at his hardiness under suffering, and when it appeared to him that the Martyr was a second time loosed by an Angel, and was walking in white raiment on the smoke, he believed in Christ, and was baptized, with all his house, by the blessed Priest Porphyry, and a little while after they both together earned the palm of martyrdom.\par
\switchcolumn*

\LA
\begin{Resp}\Rsign Lux perpétua lucébit Sanctis tuis, Dómine, \Star{} Et ætérnitas témporum, allelúja, allelúja.\end{Resp}
\begin{Resp}\Vsign Lætítia sempitérna erit super cápita eórum: gáudium et exsultatiónem obtinébunt.\end{Resp}
\begin{Resp}\Rsign Et ætérnitas témporum, allelúja, allelúja.\end{Resp}
\switchcolumn
\EN
\begin{Resp}\Rsign Eternal light will shine over your saints, O Lord, \Star{} And the eternity of the times, alleluia, alleluia.\end{Resp}
\begin{Resp}\Vsign Everlasting joy shall be upon their heads, they shall obtain joy and gladness\end{Resp}
\begin{Resp}\Rsign And the eternity of the times, alleluia, alleluia.\end{Resp}
\switchcolumn*

\LA
\LessonHead{Lectio V}
\switchcolumn
\EN
\LessonHead{Lesson V}
\switchcolumn*

\LA
\noindent At Venantius præsidi sistitur, et ab eo íterum frustra tentátus ut Christi fidem deséreret, in carcerem conjícitur; quo Attalus præco mittitur, qui ei dicat se quoque christianum fuísse et ei nómini proptérea renuntiásse, quod cognovísset inane esse fidei comméntum, quo Christiáni præséntibus se ábdicant ob vanam futurórum spem. Verum nobilis Christi athleta, callidi hostis insidias non ignorans, diaboli ministrum a se pénitus rejécit. Quare ad præsidem íterum adducto omnes contusi sunt dentes maxillæque confractæ; atque, ita cæsus, in sterquilinium dejícitur. Sed inde ab Angelo quoque ereptus, rursus stetit ante judicem; qui, Venantio adhuc loquente, e tribunali cécidit, et in ea voce, Verus est Venantii Deus, nostros deos destrúite, exclámans exspirávit.\par
\switchcolumn
\EN
\noindent Now Venantius stood before the Praefect, and when he had again vainly tempted him to give up his faith in Christ, he cast him into prison, and sent unto him Attalus the crier. Attalus told him how that he also had been a Christian, but had denied that name, seeing it was a foolish faith which made Christians to throw away things present for a groundless hope of things to come. But Christ's brave champion, well knowing the wiles of our subtle enemy, drove the devil's servant from his presence. When he appeared again before the Praefect, his teeth and jaws were broken, and so mangled he was cast out upon a dunghill. But thence also an Angel delivered him, and he stood again before the judge. And there while Venantius was yet speaking, the judge fell from off the judgment-seat, and when he had cried with a loud voice, “Venantius, his God is true, take away our gods,” he died.\par
\switchcolumn*

\LA
\begin{Resp}\Rsign In servis suis, allelúja, \Star{} Consolábitur Deus, allelúja.\end{Resp}
\begin{Resp}\Vsign Judicábit Dóminus pópulum suum, et in servis suis.\end{Resp}
\begin{Resp}\Rsign Consolábitur Deus, allelúja.\end{Resp}
\switchcolumn
\EN
\begin{Resp}\Rsign His servants, alleluia \Star{} God will comfort, alleluia, alleluia.\end{Resp}
\begin{Resp}\Vsign The Lord will judge His people and, His servants,\end{Resp}
\begin{Resp}\Rsign God will comfort, alleluia, alleluia.\end{Resp}
\switchcolumn*

\LA
\LessonHead{Lectio VI}
\switchcolumn
\EN
\LessonHead{Lesson VI}
\switchcolumn*

\LA
\noindent Quod cum præsidi nuntiátum esset, extemplo Venántium leónibus óbjici jussit; qui, naturali feritáte omissa, ad ejus se pedes abjecérunt. Interim ille pópulum Christi fidem edocébat. Quare inde amotus, íterum in carcerem traditur. Cumque postridie præsidi referret Porphyrius, se per visum noctu pópulos, quos Venantius aqua tingebat, claríssima luce fulgéntes, ipsum vero præsidem obscuríssima caligine operátum vidisse; præses, ira incénsus, eum illico cápite plecti imperat, deínde Venántium per loca vépribus et cárduis cónsita trahi usque ad vésperam. Is, cum semiánimis relíctus esset, mane se íterum præsidi præsentávit; cujus jussu statim e rupe præcipitátur. Sed inde étiam divinitus ereptus, denuo per loca aspera ad mille passus trahitur; ubi, milítibus siti æstuántibus, in próxima convalle, ex lápide, in quo et genuum formam relíquit, sicut étiam nunc in ejus ecclésia vidére licet, crucis signo a Venantio facto, aquæ manarunt. Eo miraculo plures permoti in Christum credidérunt, quos omnes præses eo loci una cum Venantio cápite feriri jussit. Fulgura et terræmotus eo témpore ita magni fuére, ut præses aufugeret; qui paucis tamen post diébus, divinam haud valens effugere justítiam, turpíssimam mortem oppétiit. Christiáni interim Venantii et aliórum córpora honorifico loco sepeliérunt: quæ Camerini, in ecclésia Venantio dicata, cóndita adhuc sunt.\par
\switchcolumn
\EN
\noindent Then they told the President of it, he commanded Venantius to be straightway thrown to the lions. But the beasts were not wild to him, and lay down at his feet. And meanwhile he taught the Christian faith to the people. So they took him away from thence and cast him once more into prison. The next day Porphyry came to the President, and told him how that he had seen in a vision of the night Venantius sprinkling certain ones with water, and they that were sprinkled shone with a marvellous light, and the President himself hidden in deep darkness. Then the President was moved to great anger and commanded forthwith to behead Porphyry. As for Venantius, he bade them drag him about in rough places, full of briars and thistles, until the evening. When it was over, he was left half dead, but in the morning he stood for the last time before the President, who commanded to cast him down from a steep rock. It pleased God that this should not kill him, and he was haled again through rough places for about a mile. There the soldiers were athirst, and Venantius, by the sign of the Cross, made waters to flow from a stone in a gulley hard by. This is that stone whereon also he left the imprint of his knees, and which can be seen to this day in his Church. By this wonder many were moved to believe in Christ and the President commanded them all, and Venantius with them, to be beheaded in the same place where they were. When it was done there were great lightnings and earthquakes, so that the President fled, but he could not fly from the judgment of God, and but a few days thereafter he died a most shameful death. Meanwhile the Christians took the bodies of Venantius and the others, and buried them in an honourable place, wherein they lie to this day, under the Church at Camerino which is dedicated to Venantius.\par
\switchcolumn*

\LA
\begin{Resp}\Rsign Fíliæ Jerúsalem, veníte et vidéte Mártyres cum corónis, quibus coronávit eos Dóminus \Star{} In die solemnitátis et lætítiæ, allelúja.\end{Resp}
\begin{Resp}\Vsign Quóniam confortávit seras portárum tuárum, benedíxit fílios tuos in te.\end{Resp}
\begin{Resp}\Rsign In die solemnitátis et lætítiæ, allelúja.\end{Resp}
\begin{Resp}\Vsign Glória Patri, et Fílio, \Star{} et Spirítui Sancto.\end{Resp}
\begin{Resp}\Rsign In die solemnitátis et lætítiæ, allelúja.\end{Resp}
\switchcolumn
\EN
\begin{Resp}\Rsign Come forth, O ye daughters of Jerusalem, and behold the Martyrs with the crowns wherewith the Lord crowned them; \Star{} In the day of His feasting, and of His gladness. Alleluia.\end{Resp}
\begin{Resp}\Vsign For He hath strengthened the bars of thy gates He hath blessed thy children within thee.\end{Resp}
\begin{Resp}\Rsign In the day of His feasting, and of His gladness. Alleluia.\end{Resp}
\begin{Resp}\Vsign Glory be to the Father, and to the Son, \Star{} and to the Holy Ghost.\end{Resp}
\begin{Resp}\Rsign In the day of His feasting, and of His gladness. Alleluia.\end{Resp}
\switchcolumn*

\LA
\LessonHead{Lectio IX (pro commemoratione)}
\switchcolumn
\EN
\LessonHead{Lesson IX (when commemorated)}
\switchcolumn*

\LA
\LessonTitle{Commemoratio: S. Venantii Martyris}
\switchcolumn
\EN
\LessonTitle{Commemoration of: St. Venantius, Martyr}
\switchcolumn*

\LA
\noindent Venántius Camers, quíndecim annos natus, christiánæ religiónis accusátus apud Antíochum, qui sub Décio imperatóre Cameríno prǽerat, in porta civitátis prǽsidi se óbtulit. Quem ille, pollicitatiónibus ac terróribus diu tentátum, flagris cædi et vínculis astríngi jussit; e quibus mirabíliter ab Angelo solútus, lampádibus póstea adúritur, atque invérso ore, fumo suppósito, suspénditur. Eídem, ad prǽsidem íterum addúcto, omnes contúsi sunt dentes maxillǽque confráctæ, atque ita cæsus in sterquilínium deícitur. Sed inde ab Angelo eréptus, rursus stetit ante júdicem; qui, Venántio adhuc loquénte, e tribunáli cécidit et in ea voce, Verus est Venántii Deus, nostros deos destrúite, exclámans exspirávit. Demum, post nova et exquisíta torménta, una cum áliis decem gloriósi certáminis cursum cervícibus abscíssis implévit. Quorum córpora christiáni honorífico loco sepeliérunt, quæ Cameríni, in ecclésia Venántio dicáta, cóndita sunt.\par
\switchcolumn
\EN
\noindent At the age of fifteen, Venantius of Camerino was denounced for his Christian religion to Antiochus, who was Prefect of Camerino under the Emperor Decius. Venantius presented himself to the prefect at the city gates. For a long time the prefect tempted him by means of promises and threats, then commanded that he be beaten and chained. Miraculously freed by an angel, Venantius was then burned with torches and suspended face down over a smoking fire. Led back again to the governor, he had all his teeth and jaws broken and, thus mutilated, he was thrown into a pit of dung. Rescued from this pit by an Angel, he stood once more before the judge, who, even as Venantius was speaking to him, fell from his tribunal, crying out, “Venantius, his God is true, take away our gods!” and expired. At length, after new and exquisite torments, Venantius was beheaded, along with ten others, and so finished the course of his glorious struggle. The Christians gave honourable burial to the bodies of these martyrs, who now rest in Camerino, in the church dedicated to Venantius.\par
\switchcolumn*

\LA
\TeDeum{Te Deum}
\switchcolumn
\EN
\TeDeum{Te Deum}
\switchcolumn*

\end{paracol}

\par\needspace{10\baselineskip}\addvspace{1.2em}{\centering\small\textsc{Die 19 Maii} · \textsc{19 May}\par}
\Office{S. Petri Celestini Papæ et Confessoris}{St. Peter Celestine, Pope and Confessor}{Duplex}
\addcontentsline{toc}{subsection}{Die 19 Maii · S. Petri Celestini Papæ et Confessoris\enspace\textit{St. Peter Celestine, Pope and Confessor}}
\begin{paracol}{2}
\LA
\RefLine{Lectiones I–III: de Scriptura occurrente.}
\switchcolumn
\EN
\RefLine{Lessons I–III: from the Scripture of the day.}
\switchcolumn*

\LA
\RefLine{Lectiones VII–VIII: ex Commúni Summorum Pontificum Confessorum (C4b).}
\switchcolumn
\EN
\RefLine{Lessons VII–VIII: from the Common of Popes (C4b).}
\switchcolumn*

\LA
\RefLine{Lectio IX: de S. Pudentianæ Virginis (05-19o).}
\switchcolumn
\EN
\RefLine{Lesson IX: of St. Pudentiana, Virgin (05-19o).}
\switchcolumn*

\LA
\LessonHead{Lectio IV}
\switchcolumn
\EN
\LessonHead{Lesson IV}
\switchcolumn*

\LA
\noindent Petrus, a nomine quo Pontifex est appellatus, Cælestinus dictus, honestis catholicisque parentibus Aeserniæ in Samnitibus natus, adolescentiam vix ingressus, ut ánimum a mundi illecebris custodiret, in solitudinem secessit. Ibi contemplationibus mentem mitriens, corpus in servitutem redigens, ferream catenam ad nudam carnem adhibebat. Congregationem, quæ postea Cælestinorum dicta est, sub regula sancti Benedicti instituit. Hinc quasi lucerna supra candelabrum posita, cum abscondi nequiret, (Romana Ecclesia diu viduata pastore) in Petri cathedram ignorans et absens ascitus, magna novitatis admiratione non minus quam repentino gaudio cunctos affecit. Cum autem in pontificatus sublimitate collocatus, variis distentus curis, assuitis incumbere meditationibus vix posse cognosceret, oneri pariter et honori voluntarie cessit. Indeque priscam vitæ rationem repetens, obdormivit in Domino, ejusque pretiosam mortem crux præfulgens in aëre ante cubiculi ostium reddidit amplius gloriosam. Miraculis multis tam vivens, quam post obitum claruit; quibus rite examinatis, Clemens quintus anno postquam decessit undecimo, sanctorum numero adscripsit.\par
\switchcolumn
\EN
\noindent This Peter, who is called Peter Celestine, because when he became Pope he did so under the title of Celestine V, was the son of respectable Catholic parents, and was born at Isernia in Apulia, (about the year of grace 1221.) He was hardly entered on boyhood, when he withdrew into a desert, in order to keep his soul safe from the snares of the world. In solitude he fed his mind with heavenly meditation, and brought his body into subjection, even by wearing an iron chain next to his bare flesh. He founded, under the Rule of St. Benedict, that congregation which was afterwards known as the Celestine. His light, as of a candle set upon a candlestick, could not be kept hidden, and after the Church of Rome had for a long while been widowed of a shepherd, he was chosen without his knowledge and in his absence, to fill the chair of Peter. The news of his election filled himself with as great amazement, as it did all others with sudden joy. When, however, he was seated in the exalted place of the Papal dignity, he found that the many cares by which he was beset made it well-nigh impossible for him to give himself to his accustomed meditations (after four months,) of his own free will he resigned the burden and the honour together (on the 13th day of December, 1294); and, while he sought to return to his old way of life, (on the 19th day of May, 1296,) he fell asleep in the Lord. How precious his death was in His sight was gloriously manifested by a Cross which appeared shining in the air before the door of the cell. He was illustrious for miracles both during his life and after his death, and when these had been duly investigated, Clement V., in the eleventh year after his departure hence, enrolled his name among those of the Saints.\par
\switchcolumn*

\LA
\begin{Resp}\Rsign Invéni David servum meum, óleo sancto meo unxi eum: \Star{} Manus enim mea auxiliábitur ei, allelúja.\end{Resp}
\begin{Resp}\Vsign Nihil profíciet inimícus in eo, et fílius iniquitátis non nocébit ei.\end{Resp}
\begin{Resp}\Rsign Manus enim mea auxiliábitur ei, allelúja.\end{Resp}
\switchcolumn
\EN
\begin{Resp}\Rsign I have found David My servant, with My holy oil have I anointed him \Star{} For My hand shall help him, alleluia.\end{Resp}
\begin{Resp}\Vsign The enemy shall prevail nothing against him, nor the son of wickedness afflict him.\end{Resp}
\begin{Resp}\Rsign For My hand shall help him, alleluia.\end{Resp}
\switchcolumn*

\LA
\LessonHead{Lectio V}
\switchcolumn
\EN
\LessonHead{Lesson V}
\switchcolumn*

\LA
\LessonTitle{Ex libro Moralium sancti Gregorii Papæ.}
\switchcolumn
\EN
\LessonTitle{From the Book of Moral Reflections upon Job, written by Pope St. Gregory the Great}
\switchcolumn*

\LA
\Source{Lib. 10. cap. 16. in cap. 12. Job.}
\switchcolumn
\EN
\Source{(Bk. x Chap. xvi on Job xii)}
\switchcolumn*

\LA
\noindent Deridétur justi simplícitas. Hujus mundi sapiéntia est, cor machinatiónibus tégere, sensum verbis veláre: quæ falsa sunt, vera osténdere; quæ vera sunt, falsa demonstráre. Hæc nimírum prudéntia usu a juvénibus scitur hæc a púeris prétio díscitur; hanc qui sciunt, céteros despiciéndo supérbiunt: hanc qui nésciunt, subjécti et tímidi in áliis mirántur: quia ab eis hæc éadem duplicitátis iníquitas nómine palliáta dilígitur, dum mentis pervérsitas urbánitas vocátur. Hæc sibi obsequéntibus præcipit honórum cúlmina quærere, adépta temporális glóriæ vanitáte gaudére, irrogáta ab áliis mala multiplícius réddere: cum vires súppetunt, nullis resisténtibus cédere: cum virtútis possibílitas deest, quidquid explére per malítiam non valent, hoc in pacífica bonitáte simuláre.\par
\switchcolumn
\EN
\noindent The simplicity of the righteous is made a subject of derision. The wisdom of this world hideth our true feelings by artifice, and useth language to conceal our thoughts this is the wisdom which demonstrated the truth of falsehood, and showeth the falsehood of the truth. This kind of shrewdness the young acquire by practice, and children pay for the learning it. Those who are good at this look down upon their neighbours those who are bad at it are humble and timid, and wonder at it in others they regard this astuteness too, wrong though it be, with wistful admiration, under softened epithets. Un-straightforwardness is called good breeding. The principles of the world teach those who entertain them, to try and rise to distinction, and when they have attained the bubble of glory which is so soon to pass away, to feel it sweet to have at their feet them on whom they may wreak rich revenge. These principles teach a man, as long as he is strong enough, to give way to nobody else, and, if he hath no chance by force, to try and attain his object by diplomacy.\par
\switchcolumn*

\LA
\begin{Resp}\Rsign Pósui adjutórium super poténtem, et exaltávi eléctum de plebe mea: \Star{} Manus enim mea auxiliábitur ei, allelúja.\end{Resp}
\begin{Resp}\Vsign Invéni David servum meum, óleo sancto meo unxi eum.\end{Resp}
\begin{Resp}\Rsign Manus enim mea auxiliábitur ei, allelúja.\end{Resp}
\switchcolumn
\EN
\begin{Resp}\Rsign I have laid help upon one that is mighty, and have exalted one chosen out of My people \Star{} For My hand shall help him, alleluia.\end{Resp}
\begin{Resp}\Vsign I have found David My servant, with My holy oil have I anointed him.\end{Resp}
\begin{Resp}\Rsign For My hand shall help him, alleluia.\end{Resp}
\switchcolumn*

\LA
\LessonHead{Lectio VI}
\switchcolumn
\EN
\LessonHead{Lesson VI}
\switchcolumn*

\LA
\noindent At contra, sapiéntia justórum est: nil per ostensiónem fíngere, sensum verbis aperíre, vera ut sunt dilígere, falsa devitáre; bona gratis exhibére, mala libéntius toleráre quam fácere; nullam injúriæ ultiónem quærere, pro veritáte contuméliam lucrum putáre. Sed hæc justórum simplícitas deridétur; quia ab hujus mundi sapiéntibus puritátis virtus fatúitas créditur. Omne enim quod innocénter ágitur, ab eis proculdúbio stultum putátur; et quidquid in ópere véritas ápprobat, carnáli sapiéntiæ fátuum sonat. Quid namque stúltius vidétur mundo, quam mentem verbis osténdere, nil cállida machinatióne simuláre, nullas injúriis contumélias réddere, pro maledicéntibus oráre, paupertátem quærere, posséssa relínquere, rapiénti non resístere, percutiénti álteram maxíllam præbére?\par
\switchcolumn
\EN
\noindent The wisdom of the righteous is the contrary of all this. They seek to avoid deception, to give their thoughts a clear expression in their words, to love the truth because it is the truth, to avoid falsehood, and rather to suffer than to inflict evil. Such are they who seek not to avenge themselves for wrong, and deem it gain to be despised for the truth's sake. This their simplicity is made a subject of derision, for such as are wise in this world believe the purity of their virtue to be simple foolery. Whatsoever is done innocently, they consider without doubt stupid. Such works as the truth approveth are idiotic, when tried by carnal standards of wisdom. After all, what stupider thing is there in this world than to express our real thoughts in our words, to keep nothing quiet by skilful tact, to repay no injuries, to pray for them which curse us, to seek poverty, to give up property, to strive not with such as take from us, to turn the other cheek to the smiter\par
\switchcolumn*

\LA
\begin{Resp}\Rsign Iste est, qui ante Deum magnas virtútes operátus est, et omnis terra doctrína ejus repléta est: \Star{} Ipse intercédat pro peccátis ómnium populórum, allelúja.\end{Resp}
\begin{Resp}\Vsign Iste est, qui contémpsit vitam mundi, et pervénit ad cæléstia regna.\end{Resp}
\begin{Resp}\Rsign Ipse intercédat pro peccátis ómnium populórum, allelúja.\end{Resp}
\begin{Resp}\Vsign Glória Patri, et Fílio, \Star{} et Spirítui Sancto.\end{Resp}
\begin{Resp}\Rsign Ipse intercédat pro peccátis ómnium populórum, allelúja.\end{Resp}
\switchcolumn
\EN
\begin{Resp}\Rsign This is he which wrought great wonders before God, and the whole earth is full of his teaching \Star{} May he pray for all people, that their sins may be forgiven unto them, alleluia.\end{Resp}
\begin{Resp}\Vsign This is he which loved not his life in this world, and hath attained unto the kingdom of heaven.\end{Resp}
\begin{Resp}\Rsign May he pray for all people, that their sins may be forgiven unto them, alleluia.\end{Resp}
\begin{Resp}\Vsign Glory be to the Father, and to the Son, \Star{} and to the Holy Ghost.\end{Resp}
\begin{Resp}\Rsign May he pray for all people, that their sins may be forgiven unto them, alleluia.\end{Resp}
\switchcolumn*

\LA
\LessonHead{Lectio IX (pro commemoratione)}
\switchcolumn
\EN
\LessonHead{Lesson IX (when commemorated)}
\switchcolumn*

\LA
\LessonTitle{Commemoratio: S. Petri Celestini Papæ et Confessoris}
\switchcolumn
\EN
\LessonTitle{Commemoration of: St. Peter Celestine, Pope and Confessor}
\switchcolumn*

\LA
\noindent Petrus, a nómine quo Póntifex est appellátus, Cælestínus dictus, honéstis catholicísque paréntibus Æsérniæ in Samnítibus natus, adolescéntiam vix ingréssus, ut ánimum a mundi illécebris custodíret, in solitúdinem secéssit. Ibi contemplatiónibus mentem nútriens, corpus in servitútem rédigens, férream caténam ad nudam carnem adhibébat. Congregatiónem, quæ póstea, Cælestinórum dicta est, sub régula sancti Benedícti instítuit. Hinc, quasi lucérna supra candelábrum pósita, cum abscóndi nequíret, (Romána Ecclésia diu viduáta pastóre) in Petri Cáthedram ignórans et absens adscítus, magna novitátis admiratióne non minus quam repentíno gáudio cunctos affécit. Cum autem in pontificátus sublimitáte collocátus, váriis disténtus curis, assuétis incúmbere meditatiónibus vix posse cognósceret, óneri páriter et honóri voluntárie cessit. Indéque priscam vitæ ratiónem répetens, obdormívit in Dómino, ejúsque pretiósam mortem crux præfúlgens in áëre ante cubículi óstium réddidit ámplius gloriósam. Miráculis multis tam vivens quam post óbitum cláruit; quibus rite examinátis, Clemens quintus anno postquam decéssit undécimo, eum Sanctórum número adscrípsit.\par
\switchcolumn
\EN
\noindent Peter, who is called Peter Celestine, because when he became Pope he did so under the title of Celestine V, was the son of respectable Catholic parents, and was born at Isernia in Apulia. He was hardly entered on boyhood, when he withdrew into a desert, in order to keep his soul safe from the snares of the world. In solitude he fed his mind with heavenly meditation, and brought his body into subjection, even by wearing an iron chain next to his bare flesh. He founded, under the Rule of St. Benedict, that congregation which was afterwards known as the Celestines. His light, as of a candle set upon a candlestick, could not be kept hidden, and after the Church of Rome had for a long while been widowed of a shepherd, he was chosen without his knowledge and in his absence, to fill the chair of Peter. The news of his election filled himself with as great amazement, as it did all others with sudden joy. When, however, he was seated in the exalted place of the Papal dignity, he found that the many cares by which he was beset made it well-nigh impossible for him to give himself to his accustomed meditations; of his own free will, he resigned the burden and the honor together; and, while he sought to return to his old way of life, he fell asleep in the Lord. How precious his death was in his sight was gloriously manifested by a Cross which appeared shining in the air before the door of the cell. He was illustrious for miracles both during his life and after his death, and when these had been duly investigated, Clement V, in the eleventh year after his departure hence, enrolled his name among those of the Saints.\par
\switchcolumn*

\LA
\TeDeum{Te Deum}
\switchcolumn
\EN
\TeDeum{Te Deum}
\switchcolumn*

\end{paracol}

\par\needspace{10\baselineskip}\addvspace{1.2em}{\centering\small\textsc{Die 19 Maii} · \textsc{19 May}\par}
\Office{S. Pudentianæ Virginis}{St. Pudentiana, Virgin}{Simplex}
\addcontentsline{toc}{subsection}{Die 19 Maii · S. Pudentianæ Virginis\enspace\textit{St. Pudentiana, Virgin}}
\begin{paracol}{2}
\LA
\LessonHead{Lectio IX (pro commemoratione)}
\switchcolumn
\EN
\LessonHead{Lesson IX (when commemorated)}
\switchcolumn*

\LA
\LessonTitle{Commemoratio: S. Pudentianæ Virginis}
\switchcolumn
\EN
\LessonTitle{Commemoration of: St. Pudentiana, Virgin}
\switchcolumn*

\LA
\noindent Pudentiana virgo, Pudentis Romani filia, parentibus orbata, cum admirabili pietate christianam religionem coleret, una cum sorore Praxede pecuniam ex vendito patrimonio redactam pauperibus distribuit, seque jejuniis et orationibus dedit. Cujus étiam opera, tota ejus familia, in qua erant nonaginta sex homines, a Pio Pontifice baptizata est. Quod autem ab Antonino imperatore sancitum erat, ne Christiani publice sacrificia facerent, Pius Pontifex in sedibus Pudentianæ cum Christianis sacra celebrabat. Quibus illa benigne acceptis, quæ ad vitam necessaria essent, suppeditabat Itaque in his christianæ pietatis officiis migravit e vita, et in sepulcro patris ad cemeterium Priscillæ via Salaria sepulta est, decimo quarto Kalendas Junii.\par
\switchcolumn
\EN
\noindent The maiden Pudentiana was the orphan daughter of Pudens the Roman Senator. She was a Christian of eminent godliness. She with her sister Praxedes distributed to the poor the money which they obtained by the sale of their inheritance. She gave herself continually to fasting and prayer. By her care the whole of the household, being ninety-six persons, were baptized by Pope Pius. Whereas the Emperor Antonine had forbidden the Christians to offer sacrifice in public, Pope Pius used to meet with them in Pudentiana's house, to celebrate the holy rites. She was a gracious hostess to them, and ministered to them in such things as are needful for the body. She thus busied herself in works of Christian godliness until she passed from this present life to a better. She was buried in her father's sepulchre in the cemetery of Priscilla on the Salarian Way upon the 19th day of May.\par
\switchcolumn*

\LA
\TeDeum{Te Deum}
\switchcolumn
\EN
\TeDeum{Te Deum}
\switchcolumn*

\end{paracol}

\par\needspace{10\baselineskip}\addvspace{1.2em}{\centering\small\textsc{Die 20 Maii} · \textsc{20 May}\par}
\Office{S. Bernardini Senensis Confessoris}{St. Bernardine of Siena, Confessor}{Semiduplex}
\addcontentsline{toc}{subsection}{Die 20 Maii · S. Bernardini Senensis Confessoris\enspace\textit{St. Bernardine of Siena, Confessor}}
\begin{paracol}{2}
\LA
\RefLine{Lectiones I–III: de Scriptura occurrente.}
\switchcolumn
\EN
\RefLine{Lessons I–III: from the Scripture of the day.}
\switchcolumn*

\LA
\RefLine{Lectiones VII–IX: ex Commúni Confessoris non pontificis (C5).}
\switchcolumn
\EN
\RefLine{Lessons VII–IX: from the Common of a Confessor not a Bishop (C5).}
\switchcolumn*

\LA
\LessonHead{Lectio IV}
\switchcolumn
\EN
\LessonHead{Lesson IV}
\switchcolumn*

\LA
\noindent Bernardinus Albizesca, nobili Senénsi família ortus, ab ineunte ætate non obscura sanctitátis dedit indícia; nam a piis paréntibus honeste educatus, neglectis puerílibus ludis, inter prima grammaticæ stúdia pietátis opéribus ánimum inténdit, jejuniis, oratióni, et beatíssimæ Vírginis cultui præcipue addictus. Misericórdia vero in páuperes fuit insígnis. Quæ quidem ómnia procedénte témpore quo mélius posset excólere, eórum número adscribi vóluit, qui Senis in hospitali domo beátæ Maríæ de Scala Deo inserviunt, unde complures sanctitáte celebres viri prodiérunt. Ibi corporis afflictatióne et ægrotántium cura, dum atrox pestiléntia grassarétur, incredibili caritate sese exercuit. Inter ceteras autem virtútes castitátem, egregia forma repugnante, sanctíssime custodívit adeo ut eo præsénte nemo umquam, ne impudentíssimus quidem, verbum minus honestum proferre auderet.\par
\switchcolumn
\EN
\noindent This Bernardine was born of the noble family of the Albizeschi, in the Republic of Siena, (on the 8th of September, in the year 1380.) His saintliness began to manifest itself from his earliest years. He was well brought up by a godly father and mother, and even when he was being taught the first rudiments of worldly learning, he used to give up his play-time to occupy himself with devout works, being much drawn to fasting, prayer, and the devotion to the most Blessed Virgin. He abounded likewise in tenderness for the poor. As time went on, that he might the more entirely do these things, it was his will to enroll himself among those who work in the Hospital of Blessed Mary, called “of the Ladder,” at Sienna. There, during the raging of an horrible distemper, he laboured with marvellous charity and great bodily suffering, in serving the sick. In bodily presence he was a very goodly person, but, with all his other virtues, he kept ever so holy a guard over his purity, that it soon came to pass that no one, however shameless, dared to say an unseemly word in his presence.\par
\switchcolumn*

\LA
\begin{Resp}\Rsign Honéstum fecit illum Dóminus, et custodívit eum ab inimícis, et a seductóribus tutávit illum: \Star{} Et dedit illi claritátem ætérnam, allelúja.\end{Resp}
\begin{Resp}\Vsign Justum dedúxit Dóminus per vias rectas, et osténdit illi regnum Dei.\end{Resp}
\begin{Resp}\Rsign Et dedit illi claritátem ætérnam, allelúja.\end{Resp}
\switchcolumn
\EN
\begin{Resp}\Rsign The Lord made him honourable, and defended him from his enemies, and kept him safe from those that lay in wait for him; \Star{} And gave him perpetual glory, alleluia.\end{Resp}
\begin{Resp}\Vsign He went down with him into the pit, and left him not in bonds.\end{Resp}
\begin{Resp}\Rsign And gave him perpetual glory, alleluia.\end{Resp}
\switchcolumn*

\LA
\LessonHead{Lectio V}
\switchcolumn
\EN
\LessonHead{Lesson V}
\switchcolumn*

\LA
\noindent Gravi morbo tentatus, eoque ad quatuor menses patientíssime toleráto, demum incolumis, de religiosæ vitæ instituto capesséndo deliberare cœpit. Quo ut sibi viam muniret, ædiculam in extrema urbe conduxit; in quam cum sese abdidísset, asperrimam omni ex parte vitam trahebat, Deum assidue orans, ut, quid sibi sequéndum esset, osténderet. Quare divinitus factus est, ut beáti Francisci ordinem præ ceteris optaret, in quo humilitate, patiéntia aliisque religiosi hóminis virtútibus excelluit. Id cum cœnobii rector animadverteret, jamque antea Bernardini doctrinam et sacrárum litterárum perítiam perspectam haberet, prædicándi onus eidem impósuit: quo humillime suscepto, cum se minus idoneum agnosceret ob vocis exilitátem ac raucitátem, Dei ope implorata, non sine miraculo ejusmodi impediménto liberátus est.\par
\switchcolumn
\EN
\noindent He suffered a severe sickness, and when, after bearing it with the utmost patience, he recovered his health, he began to think of embracing some institute of the religious life. To make his way sure, he built a little hut in the outskirts of the city, where he hid himself and led a life of hardships of all kinds, continuing instant in prayer to God that He would be pleased to make clear to him what path he should follow. And so it came to pass by God's will that he chose the Order of Blessed Francis. In that Order he shone a bright instance of lowliness, long-suffering, and every other grace of a religious man. When the superior of his convent saw this, and had already considered what his teaching and knowledge of sacred learning were, he laid on Bernardine the duty of preaching. This the Saint humbly accepted, and finding that his usefulness was much impaired by his having a shrill, harsh voice, he betook him to implore the help of God, Who was pleased, not without a miracle, to free him from this drawback.\par
\switchcolumn*

\LA
\begin{Resp}\Rsign Amávit eum Dóminus, et ornávit eum: stolam glóriæ índuit eum, \Star{} Et ad portas paradísi coronávit eum, allelúja.\end{Resp}
\begin{Resp}\Vsign Índuit eum Dóminus lorícam fídei, et ornávit eum.\end{Resp}
\begin{Resp}\Rsign Et ad portas paradísi coronávit eum, allelúja.\end{Resp}
\switchcolumn
\EN
\begin{Resp}\Rsign The Lord loved him and beautified him; He clothed him with a robe of glory, \Star{} And crowned him at the gates of Paradise, alleluia.\end{Resp}
\begin{Resp}\Vsign The Lord hath put on him the breast-plate of faith, and hath adorned him.\end{Resp}
\begin{Resp}\Rsign And crowned him at the gates of Paradise, alleluia.\end{Resp}
\switchcolumn*

\LA
\LessonHead{Lectio VI}
\switchcolumn
\EN
\LessonHead{Lesson VI}
\switchcolumn*

\LA
\noindent Cumque ea témpora vitiis criminibúsque redundarent, et cruentis factiónibus in Italia, divina humanaque ómnia permixta essent, Bernardinus urbes atque oppida concursans, in nómine Jesu, quem semper in ore et in péctore gerebat, collapsam pietátem moresque verbo et exemplo magna ex parte restituit. Quo factum est, ut præcláræ civitátes eum sibi episcopum a summo Pontifice postularent; quod ille munus invicta humilitate constantíssime rejécit. Denique vir Dei imménsis labóribus exhaustus, multis magnisque editis miraculis, libris étiam pie docteque conscriptis, cum vixísset annos sex ac sexagínta, in urbe Aquila in Vestinis beato fine quiévit. Quem novis in dies coruscántem signis, anno post obitum sexto, Nicolaus quintus Pontifex maximus in Sanctórum númerum retulit.\par
\switchcolumn
\EN
\noindent Those were times fruitful in vices and crimes and the bloody civil wars which raged in Italy confounded all things Divine and human. Bernardine went through the cities and towns, and, in the Name of Jesus, that Name which he ever bore upon his lips and in his heart, he prevailed in great measure by his word and example, in setting up falling godliness and morality. Illustrious cities demanded him from the Pope as their Bishop, but this was an honour which his unconquerable humility caused him always steadily to refuse. At last the man of God, after untold labours, the working of many and great miracles, and the writing of godly and learned books, in the 67th year of his age, at Aquila in the Abruzzi, rested in a blessed death, (upon the 20th day of May 1444.) As the fame of new signs and wonders increased day by day, Pope Nicholas V., in the sixth year after his death, added his name to the roll of the Saints.\par
\switchcolumn*

\LA
\begin{Resp}\Rsign Iste homo perfécit ómnia quæ locútus est ei Deus, et dixit ad eum: Ingrédere in réquiem meam: \Star{} Quia te vidi justum coram me ex ómnibus géntibus, allelúja.\end{Resp}
\begin{Resp}\Vsign Iste est, qui contémpsit vitam mundi, et pervénit ad cæléstia regna.\end{Resp}
\begin{Resp}\Rsign Quia te vidi justum coram me ex ómnibus géntibus, allelúja.\end{Resp}
\begin{Resp}\Vsign Glória Patri, et Fílio, \Star{} et Spirítui Sancto.\end{Resp}
\begin{Resp}\Rsign Quia te vidi justum coram me ex ómnibus géntibus, allelúja.\end{Resp}
\switchcolumn
\EN
\begin{Resp}\Rsign This is he which did according unto all that God commanded him; and God said unto him: Enter thou into My rest, \Star{} For thee have I seen righteous before Me among all people, alleluia.\end{Resp}
\begin{Resp}\Vsign This is he which loved not his life in this world, and is come unto an everlasting kingdom.\end{Resp}
\begin{Resp}\Rsign For thee have I seen righteous before Me among all people, alleluia.\end{Resp}
\begin{Resp}\Vsign Glory be to the Father, and to the Son, \Star{} and to the Holy Ghost.\end{Resp}
\begin{Resp}\Rsign For thee have I seen righteous before Me among all people, alleluia.\end{Resp}
\switchcolumn*

\LA
\LessonHead{Lectio IX (pro commemoratione)}
\switchcolumn
\EN
\LessonHead{Lesson IX (when commemorated)}
\switchcolumn*

\LA
\LessonTitle{Commemoratio: S. Bernardini Senensis Confessoris}
\switchcolumn
\EN
\LessonTitle{Commemoration of: St. Bernardine of Siena, Confessor}
\switchcolumn*

\LA
\noindent Bernardínus Albizésca, nóbili Senénsi família ortus, inter prima grammáticæ stúdia, negléctis puerílibus ludis, pietátis opéribus ánimum inténdit, beátæ Vírginis cúltui præcípue addíctus. Caritáte et misericórdia in páuperes insígnis, eórum servítio Senis in hospitáli beátæ Maríæ de Scala se mancipávit. De capesséndo religiósæ vitæ institúto delíberans, Deo sic disponénte, beáti Francísci órdinem præ céteris optávit, in quo humilitáte, patiéntia, aliísque religiósi hóminis virtútibus excélluit. Prædicándi ónere a superióribus suscépto, cum se minus idóneum agnósceret ob vocis exilitátem et raucitátem, Dei ope imploráta, prodigióse ab hoc impediménto liberátus est. Urbes atque óppida concúrsans, in nómine Jesu, quem semper in ore et péctore gerébat, cívium ubíque discórdias exstínxit, et collápsam pietátem morésque verbo et exémplo magna ex parte restítuit. Libros pie doctéque conscrípsit. Plenus méritis et miráculis clarus, annos natus sex ac sexagínta, in urbe Aquila in Vestínis, beáto fine quiévit.\par
\switchcolumn
\EN
\noindent St. Bernardine Albizeschi was born of a noble family of Siena. Even in his first years in school, he turned away from children's games and applied himself to exercises of devotion, especially those honouring the Blessed Virgin. Outstanding for his charity and mercy to the poor, he gave himself over to serving them at the hospital Santa Maria della Scala in Siena. When he came to think of entering the religious life, divine Providence led him to choose, in preference to others, the Franciscan Order, where he excelled in humility, patience, and the other virtues of a religious. His superiors imposed on him the duty of preaching; and, although he knew that his voice was too weak and too hoarse for a preacher's, he accepted the charge and implored the help of God. As a result, he was marvellously freed of his handicap. Going about the towns and villages in the Name of Jesus, which was always carried on his lips and in his heart, he everywhere put an end to the dissensions of the citizens, and by his word and example did much to restore their piety and morality, which had fallen to a low level. He also wrote devout and learned books. At the age of sixty-six, rich in merit and famous for his miracles, he died a happy death at the town of L'Aquila in the Abruzzi.\par
\switchcolumn*

\LA
\TeDeum{Te Deum}
\switchcolumn
\EN
\TeDeum{Te Deum}
\switchcolumn*

\end{paracol}

\par\needspace{10\baselineskip}\addvspace{1.2em}{\centering\small\textsc{Die 25 Maii} · \textsc{25 May}\par}
\Office{S. Gregorii VII Papæ et Confessoris}{St. Gregory VII, Pope and Confessor}{Duplex}
\addcontentsline{toc}{subsection}{Die 25 Maii · S. Gregorii VII Papæ et Confessoris\enspace\textit{St. Gregory VII, Pope and Confessor}}
\begin{paracol}{2}
\LA
\RefLine{Lectiones I–III: de Scriptura occurrente.}
\switchcolumn
\EN
\RefLine{Lessons I–III: from the Scripture of the day.}
\switchcolumn*

\LA
\RefLine{Lectiones VII–VIII: ex Commúni Summorum Pontificum Confessorum (C4b).}
\switchcolumn
\EN
\RefLine{Lessons VII–VIII: from the Common of Popes (C4b).}
\switchcolumn*

\LA
\RefLine{Lectio IX: de S. Urbanis Papæ et Martyris (05-25o).}
\switchcolumn
\EN
\RefLine{Lesson IX: of St. Urban, Pope and Martyr (05-25o).}
\switchcolumn*

\LA
\LessonHead{Lectio IV}
\switchcolumn
\EN
\LessonHead{Lesson IV}
\switchcolumn*

\LA
\noindent Gregórius Papa septimus, antea Hildebrandus, Suánæ in Etruria natus, doctrina, sanctitáte, omnique virtútum genere cum primis nobilis, mirifice universam Dei illustrávit Ecclésiam. Cum párvulus ad fabri ligna edolántis pedes, jam litterárum ínscius, lúderet, ex rejectis tamen segmentis illa Davidici eleménta oraculi, Dominábitur a mari usque ad mare, casu formasse narrátur: manum púeri ductante Númine, quo significarétur ejus fore amplíssimam in mundo auctoritátem. Romam deínde profectus, sub protectióne sancti Petri educátus est. Juvenis, Ecclésiæ libertátem a laicis oppressam ac depravatos ecclesiasticórum mores vehementius dolens, in Cluniacénsi monasterio, ubi sub regula sancti Benedicti austerioris vitæ observántia eo témpore maxime vigebat, monachi habitum induens, tanto pietátis ardore divinæ majestáti deserviébat, ut a sanctis ejusdem cœnobii pátribus prior sit eléctus. Sed, divina providéntia majora de eo disponente, in salútem plurimórum Cluníaco edúctus Hildebrandus, abbas primum monasterii sancti Pauli extra muros Urbis eléctus, ac postmodum Romanæ Ecclésiæ cardinalis creatus, sub Summis Pontificibus Leone nono, Victore secundo, Stephano nono, Nicolao secundo et Alexandro secundo, præcipuis munéribus et legatiónibus perfunctus est ; sanctíssimi et puríssimi consílii vir a beato Petro Damiáno nuncupatus. A Victore Papa secundo legátus a látere in Gálliam missus, Lugduni episcopum, simoníaca labe infectum, ad sui críminis confessiónem miraculo adegit. Berengárium in concílio Turonénsi ad iteratam hæresis abjuratiónem cómpulit. Cadalói quoque schisma sua virtúte compressit.\par
\switchcolumn
\EN
\noindent Hildebrand, who reigned as Pope under the name of Gregory VII, was born at Saona in Tuscany. By his teaching, by his holiness, and by his graces of all kinds, he was a noble light of the Church, whose brightness hath shone throughout all lands. There is a story to the effect that when he was a little child without any schooling, he was playing at the feet of a carpenter who was planing wood, and that God guided his hand to arrange the shavings which fell into the form of letters, making the inspired words of David, He shall have dominion from sea to sea, (Ps. lxxi. 8,) a fore-shadowing, as it were, of that wide lordship over the earth which was afterwards his. He was taken to Rome, and brought up under the shelter of St. Peter. As a young man he bitterly sorrowed over the oppression of the freedom of the Church by the laity, and over the corruption of the clergy themselves. He took the habit of a monk in the Abbey of Cluny, which was then in all the glory of the severest observance of the Rule of St. Benedict. There he served God's majesty with such warmth of earnestness that the saintly fathers of the convent chose him to be their Prior. But the Providence of God had greater things in store for him, whereby to make him a source of health to many, and he was brought away from Cluny. He was first elected Abbot of the monastery of St. Paul-without-the-walls at Rome, and afterwards created a Cardinal of the Roman Church. Under the Popes Leo IX, Victor II, Stephen IX, Nicholas II, and Alexander II, he discharged great offices of trust, and the duties of a Legate, and Blessed Peter Damian, speaking of him at this time, calleth him a man of most holy and honest thoughts. When Pope Victor II. sent him as his Legate into France, he, by a miracle, forced the Bishop of Lyons, who was befouled by the pollution of simony, to acknowledge his sin in the Council of Tours he wrung from Berenger a second abjuration of his heresy and he prevailed against the schism of Cadolaus, and strangled it.\par
\switchcolumn*

\LA
\begin{Resp}\Rsign Invéni David servum meum, óleo sancto meo unxi eum: \Star{} Manus enim mea auxiliábitur ei, allelúja.\end{Resp}
\begin{Resp}\Vsign Nihil profíciet inimícus in eo, et fílius iniquitátis non nocébit ei.\end{Resp}
\begin{Resp}\Rsign Manus enim mea auxiliábitur ei, allelúja.\end{Resp}
\switchcolumn
\EN
\begin{Resp}\Rsign I have found David My servant, with My holy oil have I anointed him \Star{} For My hand shall help him, alleluia.\end{Resp}
\begin{Resp}\Vsign The enemy shall prevail nothing against him, nor the son of wickedness afflict him.\end{Resp}
\begin{Resp}\Rsign For My hand shall help him, alleluia.\end{Resp}
\switchcolumn*

\LA
\LessonHead{Lectio V}
\switchcolumn
\EN
\LessonHead{Lesson V}
\switchcolumn*

\LA
\noindent Mortuo Alexandro secundo, invitus et mærens, unánimi ómnium consénsu, décimo Kalendas Maji anno Christi millesimo septuagesimo tertio, Summus Pontifex eléctus, sicut sol effulsit in domo Dei. Nam, potens opere et sermóne, ecclesiásticæ disciplinæ reparandæ, fidei propagandæ, libertáti Ecclésiæ restituéndæ, exstirpandis erróribus et corruptélis tanto studio incúbuit, ut ex Apostolórum ætate nullus Pontificum fuísse tradátur, qui majores pro Ecclésia Dei labóres molestiasque pertulerit, aut qui pro ejus libertáte acrius pugnaverit. Aliquot provincias a simoníaca labe expurgávit. Contra Henrici imperatóris ímpios conatus, fortis per ómnia athleta, impávidus permansit, seque pro muro dómui Israël pónere non tímuit ; ac eumdem Henricum, in profúndum malórum prolapsum, fidelium communióne regnosque privávit, atque súbditos pópulos fide ei data liberávit.\par
\switchcolumn
\EN
\noindent After the death of Alexander II, Hildebrand, against his own will and to his own grief, was, on the 22nd day of April, in the year of Christ 1073, chosen Pope by one common consent of all. Reigning as Gregory VII., he was as the sun shining upon the Temple of the Most High. (Ecclus. 1. 7.) Mighty both in word and deed, he toiled for the restoration of Ecclesiastical discipline, for the spread of the Faith, for the defence of the freedom of the Church, for the suppression of error and corruption, so that since the time of the Apostles there is said never to have been a Pope who bore more labour and trouble for the sake of God's Church, or contended more manfully for her liberties. He purged diverse provinces of the pollution of simony. Like a brave soldier he withstood without dread the unrighteous contendings of the Emperor Henry IV, against whom he shrank not from setting himself as a wall of defence for the house of Israel. And when the said Henry fell into the depths of sin he cut him off from the communion of the faithful, and from his kingdom, and loosed the nations that were subject to him from their sworn allegiance.\par
\switchcolumn*

\LA
\begin{Resp}\Rsign Pósui adjutórium super poténtem, et exaltávi eléctum de plebe mea: \Star{} Manus enim mea auxiliábitur ei, allelúja.\end{Resp}
\begin{Resp}\Vsign Invéni David servum meum, óleo sancto meo unxi eum.\end{Resp}
\begin{Resp}\Rsign Manus enim mea auxiliábitur ei, allelúja.\end{Resp}
\switchcolumn
\EN
\begin{Resp}\Rsign I have laid help upon one that is mighty, and have exalted one chosen out of My people \Star{} For My hand shall help him, alleluia.\end{Resp}
\begin{Resp}\Vsign I have found David My servant, with My holy oil have I anointed him.\end{Resp}
\begin{Resp}\Rsign For My hand shall help him, alleluia.\end{Resp}
\switchcolumn*

\LA
\LessonHead{Lectio VI}
\switchcolumn
\EN
\LessonHead{Lesson VI}
\switchcolumn*

\LA
\noindent Dum Missárum solemnia perágeret, visa est viris piis columba e cælo delapsa, humero ejus dextro ínsidens, alis exténsis caput ejus velare quo significátum est, Spíritus Sancti afflatu, non humanæ prudéntiæ ratiónibus ipsum duci in Ecclésiæ regímine. Cum ab iníqui Henrici exercitu Romæ gravi obsidióne premerétur, excitátum ab hostibus incendium signo crucis exstinxit. De ejus manu tandem a Roberto Guiscardo duce Northmanno ereptus, Cassinum se cóntulit ; atque inde Salernum ad dedicandam ecclésiam sancti Matthæi Apóstoli contendit. Cum aliquándo in ea civitáte sermónem habuísset ad pópulum, ærumnis confectus in morbum íncidit, quo se interitúrum præscívit. Postrema moriéntis Gregórii verba fuére: Diléxi justítiam et odívi iniquitátem, proptérea morior in exsílio. Innumerabília sunt, quæ vel fortiter sustinuit, vel multis coactis in Urbe Synodis sapiénter constítuit ; vir vere sanctus, criminum vindex, et acerrimus Ecclésiæ defensor. Exactis ítaque in pontificatu annis duodecim, migrávit in cælum anno salútis millesimo octogesimo quinto, plúribus in vita et post mortem miraculis clarus ; ejusque sacrum corpus in cathedrali basilica Salernitana est honorifice cónditum.\par
\switchcolumn
\EN
\noindent While he was celebrating solemn Mass, godly men saw a dove descend from heaven* perch upon his right shoulder, and spread out its wings so as to veil his head, a testimony that it was not by reasonings of man's wisdom, but by the teachings of the Holy Ghost, that he was guided in his rule over the Church. When the armies of the infamous Henry encompassed Rome, and hedged her in on every side, a great fire which the enemy had raised became extinct, when Gregory made the sign of the Cross towards it. The Norman Duke, Robert Guiscard, at length delivered Gregory from the hand of Henry, and he departed from Rome, first to the Abbey of Monte Cassino, and thence onward to Salerno, to dedicate the Church of St. Matthew the Apostle at that place. While he was preaching to the people there, on a certain day he was smitten with grievous pains, and fell into a sickness whereof he foresaw that he should never be healed. As he lay on his death-bed, Gregory's last words were I have loved righteousness and hated iniquity, and therefore I am dying in exile. He was a man really holy, a visitor of sin, and a most leal soldier of the Church. It is past reckoning how many sufferings he manfully bore, and how much he wisely ordained in many Councils, which he gathered together in Rome. He had been Pope twelve years, when, (on the 25th day of May,) in the year of salvation 1085, he went hence to be ever with the Lord. Both during his life and after his death he was marked by signs and wonders not a few. His holy body was honourably buried in the Cathedral Church of Salerno.\par
\switchcolumn*

\LA
\begin{Resp}\Rsign Iste est, qui ante Deum magnas virtútes operátus est, et omnis terra doctrína ejus repléta est: \Star{} Ipse intercédat pro peccátis ómnium populórum, allelúja.\end{Resp}
\begin{Resp}\Vsign Iste est, qui contémpsit vitam mundi, et pervénit ad cæléstia regna.\end{Resp}
\begin{Resp}\Rsign Ipse intercédat pro peccátis ómnium populórum, allelúja.\end{Resp}
\begin{Resp}\Vsign Glória Patri, et Fílio, \Star{} et Spirítui Sancto.\end{Resp}
\begin{Resp}\Rsign Ipse intercédat pro peccátis ómnium populórum, allelúja.\end{Resp}
\switchcolumn
\EN
\begin{Resp}\Rsign This is he which wrought great wonders before God, and the whole earth is full of his teaching \Star{} May he pray for all people, that their sins may be forgiven unto them, alleluia.\end{Resp}
\begin{Resp}\Vsign This is he which loved not his life in this world, and hath attained unto the kingdom of heaven.\end{Resp}
\begin{Resp}\Rsign May he pray for all people, that their sins may be forgiven unto them, alleluia.\end{Resp}
\begin{Resp}\Vsign Glory be to the Father, and to the Son, \Star{} and to the Holy Ghost.\end{Resp}
\begin{Resp}\Rsign May he pray for all people, that their sins may be forgiven unto them, alleluia.\end{Resp}
\switchcolumn*

\LA
\LessonHead{Lectio IX (pro commemoratione)}
\switchcolumn
\EN
\LessonHead{Lesson IX (when commemorated)}
\switchcolumn*

\LA
\LessonTitle{Commemoratio: S. Gregorii VII Papæ et Confessoris}
\switchcolumn
\EN
\LessonTitle{Commemoration of: St. Gregory VII, Pope and Confessor}
\switchcolumn*

\LA
\noindent Gregórius Papa séptimus, ántea Hildebrándus, Suánæ in Etrúria natus, doctrína, sanctitáte, omníque virtútum génere cum primis nóbilis, mirífice univérsam Dei Ecclésiam illustrávit. Adoléscens religiósi hábitum índuit in Cluniacénsi monastério, tantóque pietátis ardóre Deo desérviit, ut a sanctis ejúsdem cœnóbii pátribus prior fúerit eléctus. Póstea factus abbas monastérii sancti Pauli extra muros Urbis, ac deínde creátus Románæ Ecclésiæ cardinális, sub Summis Pontifícibus Leóne nono, Victóre secúndo, Stéphano nono, Nicoláo secúndo et Alexándro secúndo, præcípuis munéribus et legatiónibus perfúnctus est. Mórtuo Alexándro secúndo, unanímiter Summus Póntifex eléctus, ecclesiásticæ libertátis propugnátor ac defénsor acérrimus éxstitit; quaprópter multa passus, et Roma discédere coáctus est. Postréma ejus moriéntis verba fuére: Diléxi justítiam et odívi iniquitátem, proptérea mórior in exsílio. Migrávit in cælum, anno salútis millésimo octogésimo quinto, ejúsque corpus in cathedráli basílica Salernitána honorífice cónditum est.\par
\switchcolumn
\EN
\noindent Pope Gregory VII, the former Hildebrand, was born near Soana in Tuscany. As noble as any of the nobility in learning, in holiness and in every kind of virtue, he was a shining light to the whole Church of God. As a young man, he donned the religious habit at the monastery of Cluny, and served God with such zeal and devotion that he was chosen Prior by the holy religious of that monastery. Later, he was made Abbot of the monastery of St. Paul-outside-the-Walls, and then Cardinal of the Roman Church, performing noteworthy services and missions under Popes Leo IX, Victor II, Stephen IX, Nicholas II and Alexander II. At the death of Alexander, he was unanimously elected Pope, and stood out as a most zealous promoter and defender of the freedom of the Church, for which he suffered many things, even having to leave Rome. His last words, as he lay dying, were: “I have loved righteousness and hated iniquity, and therefore I am dying in exile.” He went to heaven in year of salvation 1085, and his body was buried with honour in the Cathedral of Salerno.\par
\switchcolumn*

\LA
\TeDeum{Te Deum}
\switchcolumn
\EN
\TeDeum{Te Deum}
\switchcolumn*

\end{paracol}

\par\needspace{10\baselineskip}\addvspace{1.2em}{\centering\small\textsc{Die 25 Maii} · \textsc{25 May}\par}
\Office{S. Urbanis Papæ et Martyris}{St. Urban, Pope and Martyr}{Simplex}
\addcontentsline{toc}{subsection}{Die 25 Maii · S. Urbanis Papæ et Martyris\enspace\textit{St. Urban, Pope and Martyr}}
\begin{paracol}{2}
\LA
\LessonHead{Lectio IX (pro commemoratione)}
\switchcolumn
\EN
\LessonHead{Lesson IX (when commemorated)}
\switchcolumn*

\LA
\LessonTitle{Commemoratio: S. Urbanis Papæ et Martyris}
\switchcolumn
\EN

\switchcolumn*

\LA

\switchcolumn
\EN
\Source{For St. Urban, Pope and Martyr}
\switchcolumn*

\LA
\noindent Urbanus Romanus, Alexandro Sevéro imperatóre, doctrina et vitæ sanctitáte multos ad Christi fidem convértit; in illis Valerianum, beátæ Cæcíliæ sponsum, et Tiburtium, Valeriáni fratrem, qui póstea martyrium forti animo subiérunt. Hic de bonis Ecclésiæ attributis scripsit his verbis: Ipsæ res fidelium, quæ Dómino offerúntur, non debent in alios usus quam ecclesiásticos et christianórum fratrum, vel indigéntium, converti; quia vota sunt fidelium, et pretia peccatórum, ac patrimónia páuperum. Sedit annos sex, menses septem, dies quatuor: ac martyrio coronátus, sepúltus est in cœmeterio Prætextati, octavo Kalendas Junii. Ordinatiónibus quinque habitis mense Decembri, creávit presbyteros novem, diaconos quinque, episcopos per diversa loca octo.\par
\switchcolumn
\EN
\noindent This Urban was a Roman, who, in the reign of the Emperor Alexander Severus, by his teaching and holy life, brought many to believe in Christ. Among others was Valerian, the husband of the blessed Cecily, and Tiburtius, the brother of Valerian, both of whom afterwards bravely underwent martyrdom. It was Urban l who wrote the following words\par
concerning the property of the Church: Those things which His faithful ones\par
make offering of unto the Lord, must never be turned to any other use than those\par
of the Church, or of our Christian brethren, or of the poor. They are the\par
free-will offerings of faithful believers, the trespass offerings of sinners,\par
and the inheritance of the poor. He sat in the chair of Peter six years, seven\par
months, and four days, and being crowned with martyrdom, was buried in the\par
cemetery of Praetextatus, on the 25th day of May. He held five ordinations in\par
December, wherein he ordained nine Priests, five Deacons, and eight Bishops for\par
diverse places.\par
\switchcolumn*

\LA
\TeDeum{Te Deum}
\switchcolumn
\EN
\TeDeum{Te Deum}
\switchcolumn*

\end{paracol}

\par\needspace{10\baselineskip}\addvspace{1.2em}{\centering\small\textsc{Die 26 Maii} · \textsc{26 May}\par}
\Office{S. Philippi Neri Confessoris}{St. Philip Neri, Confessor}{Duplex}
\addcontentsline{toc}{subsection}{Die 26 Maii · S. Philippi Neri Confessoris\enspace\textit{St. Philip Neri, Confessor}}
\begin{paracol}{2}
\LA
\RefLine{Lectiones I–III: de Scriptura occurrente.}
\switchcolumn
\EN
\RefLine{Lessons I–III: from the Scripture of the day.}
\switchcolumn*

\LA
\RefLine{Lectiones VII–VIII: ex Commúni Confessoris non pontificis (C5).}
\switchcolumn
\EN
\RefLine{Lessons VII–VIII: from the Common of a Confessor not a Bishop (C5).}
\switchcolumn*

\LA
\RefLine{Lectio IX: de S. Eleutherii Papæ et Martyris (05-26o).}
\switchcolumn
\EN
\RefLine{Lesson IX: of St. Eleutherius, Pope and Martyr (05-26o).}
\switchcolumn*

\LA
\LessonHead{Lectio IV}
\switchcolumn
\EN
\LessonHead{Lesson IV}
\switchcolumn*

\LA
\noindent Philippus Nerius, piis honestisque paréntibus Floréntiæ natus, ab ipsa ineunte ætate non obscura dedit futuræ sanctitátis indicia. Adoléscens, ampla patrui hereditáte dimissa, Romam se contulit; ubi philosophía ac sacris litteris eruditus, totum se Christo dicávit. Ea fuit abstinéntia, ut sæpe jejunus triduum permanserit. Vigiliis et oratiónibus intentus, septem Urbis Ecclésias frequenter visitans, apud cœmetérium Callísti in cæléstium rerum contemplatióne pernoctare consuevit. Sacerdos ex obediéntia factus, in animárum salúte procuranda totus fuit; et in confessiónibus audiéndis ad extremum usque diem perseverans, innumeros penne fílios Christo péperit; quos verbi Dei quotidiáno pabulo, sacramentórum frequéntia, oratiónis assiduitate, aliisque piis exercitatiónibus enutriri cupiens, Oratórii congregatiónem instituit.\par
\switchcolumn
\EN
\noindent Philip Neri was born of godly and respectable parents at Florence, (on the 23rd day of July, in the year of grace 1515) From his earliest childhood he gave signs of the holiness of life which he afterwards attained. As a young man he gave up the rich inheritance which would have come to him from his uncle, and went to Rome, where, in the study of philosophy and theology, he gave himself altogether to Christ. His self-control was such that he sometimes fasted from all food for three days at a time. He was instant in watching and prayer, and during the frequent pilgrimages which he made to the seven churches of Rome, it was his custom to remain all night in prayer to God in the Catacomb of Kallistus. (On the 23rd of May, 1551) he became a Priest in obedience to the advice of his Confessor, and afterwards made the salvation of souls the one object of his existence; he heard confessions with unwearied tenderness until his dying day, and became the spiritual father in Christ of many sons, whom it was his beloved work to feed day by day upon the Word of God, upon the often receiving of the Sacraments, upon instant prayer, and upon other godly works: to the which end he founded the Congregation of the Oratory.\par
\switchcolumn*

\LA
\begin{Resp}\Rsign Honéstum fecit illum Dóminus, et custodívit eum ab inimícis, et a seductóribus tutávit illum: \Star{} Et dedit illi claritátem ætérnam, allelúja.\end{Resp}
\begin{Resp}\Vsign Justum dedúxit Dóminus per vias rectas, et osténdit illi regnum Dei.\end{Resp}
\begin{Resp}\Rsign Et dedit illi claritátem ætérnam, allelúja.\end{Resp}
\switchcolumn
\EN
\begin{Resp}\Rsign The Lord made him honourable, and defended him from his enemies, and kept him safe from those that lay in wait for him; \Star{} And gave him perpetual glory, alleluia.\end{Resp}
\begin{Resp}\Vsign He went down with him into the pit, and left him not in bonds.\end{Resp}
\begin{Resp}\Rsign And gave him perpetual glory, alleluia.\end{Resp}
\switchcolumn*

\LA
\LessonHead{Lectio V}
\switchcolumn
\EN
\LessonHead{Lesson V}
\switchcolumn*

\LA
\noindent Caritate Dei vulnerátus languebat jugiter, tantoque cor ejus æstuábat ardore, ut, cum intra fines suos contineri non posset, illíus sinum, confractis atque elátis duabus cóstulis, mirabíliter Dóminus ampliaverit. Sacrum vero fáciens aut ferventius orans, in áëra quandoque sublatus, mira undique luce fulgere visus fuit. Egenos et páuperes omni caritátis officio prosequebátur: dignus, qui et Angelo in spécie páuperis eleemosynam erogaret; et, dum egéntibus noctu panem deferret, in fóveam lapsus, inde páriter ab Angelo incolumis eriperétur. Humilitáti addictus, ab honóribus semper abhorruit, atque ecclesiásticas dignitates, étiam primárias, non semel ultro delatas, constantíssime recusávit.\par
\switchcolumn
\EN
\noindent He was full of the love of God, and his heart was so hot therewith, that it became straitened in its place, and the Lord was pleased to ease him by(the gristle which joined) the fourth and fifth ribs (on his left side) breaking, and so allowing more play to the internal organs. Sometimes, when he was saying Mass, or more intent than usual in prayer, he was seen to be raised off the ground, and become, in a strange manner, all shining. He was ever ready to succour the poor and needy with kindly services, in which works God was pleased to make him meet once to give alms to an Angel, and again, when once by night he was carrying bread to the hungry, and was fallen into a pit, an Angel drew him out unhurt. He longed to be lowly, and always shrank from honours, and from dignities in the Church, whereof several of the highest were diverse times offered to him, but he always firmly refused them.\par
\switchcolumn*

\LA
\begin{Resp}\Rsign Amávit eum Dóminus, et ornávit eum: stolam glóriæ índuit eum, \Star{} Et ad portas paradísi coronávit eum, allelúja.\end{Resp}
\begin{Resp}\Vsign Índuit eum Dóminus lorícam fídei, et ornávit eum.\end{Resp}
\begin{Resp}\Rsign Et ad portas paradísi coronávit eum, allelúja.\end{Resp}
\switchcolumn
\EN
\begin{Resp}\Rsign The Lord loved him and beautified him; He clothed him with a robe of glory, \Star{} And crowned him at the gates of Paradise, alleluia.\end{Resp}
\begin{Resp}\Vsign The Lord hath put on him the breast-plate of faith, and hath adorned him.\end{Resp}
\begin{Resp}\Rsign And crowned him at the gates of Paradise, alleluia.\end{Resp}
\switchcolumn*

\LA
\LessonHead{Lectio VI}
\switchcolumn
\EN
\LessonHead{Lesson VI}
\switchcolumn*

\LA
\noindent Prophetíæ dono fuit illustris, et in animórum sensibus penetrandis mirifice enituit. Virginitátem perpetuo illibatam servávit; idque assecutus est, ut eos qui puritátem cólerent, ex odore, qui vero secus, ex fœtóre dignosceret. Abséntibus interdum appáruit, iisque periclitántibus opem tulit. Ægrotos plurimos et morti próximos sanitáti restituit. Mortuum quoque ad vitam revocávit. Cælestium spirituum et ipsíus Deiparæ Vírginis frequenter fuit apparitióne dignátus, ac plurimórum ánimas splendore circumfusas in cælum conscéndere vidit. Denique, anno salútis millesimo quingentésimo nonagesimo quinto, octavo Kalendas Junias, in quem diem inciderat festum Corporis Christi, Sacro maxima spíritus exsultatióne peracto, ceterisque functiónibus expletis, post mediam noctem, qua prædixerat hora, octogenarius obdormívit in Dómino. Quem Gregórius décimus quintus, miraculis clarum, in Sanctórum númerum rétulit.\par
\switchcolumn
\EN
\noindent He was illustrious for the gift of prophecy, and had a marked and very wonderful power of reading the thoughts of men's souls. He ever kept his own virginity undefiled, and could distinguish those that were pureminded by a sort of sweet savour, and the unclean, on the contrary, by a kind of stench. He sometimes appeared in double to persons at a distance, and brought them help when they were in peril. He healed many that were sick and dying. He also raised one dead man to life. He was honoured by seeing several times heavenly spirits, and likewise the Virgin Mother of God herself. He saw the souls of diverse persons, radiant with glory, ascend to heaven. In the year of salvation 1595 the Feast of the Body of Christ fell upon the 25th day of May. Philip, on that day, said Mass with extraordinary gladness of spirit, and performed the other religious works of the day, and after the hour of midnight, at the time he had himself foretold, he fell asleep in the Lord, in the eighty-second year of his life. Gregory XV., finding that God had glorified him by many miracles, enrolled his name among those of the saints.\par
\switchcolumn*

\LA
\begin{Resp}\Rsign Iste homo perfécit ómnia quæ locútus est ei Deus, et dixit ad eum: Ingrédere in réquiem meam: \Star{} Quia te vidi justum coram me ex ómnibus géntibus, allelúja.\end{Resp}
\begin{Resp}\Vsign Iste est, qui contémpsit vitam mundi, et pervénit ad cæléstia regna.\end{Resp}
\begin{Resp}\Rsign Quia te vidi justum coram me ex ómnibus géntibus, allelúja.\end{Resp}
\begin{Resp}\Vsign Glória Patri, et Fílio, \Star{} et Spirítui Sancto.\end{Resp}
\begin{Resp}\Rsign Quia te vidi justum coram me ex ómnibus géntibus, allelúja.\end{Resp}
\switchcolumn
\EN
\begin{Resp}\Rsign This is he which did according unto all that God commanded him; and God said unto him: Enter thou into My rest, \Star{} For thee have I seen righteous before Me among all people, alleluia.\end{Resp}
\begin{Resp}\Vsign This is he which loved not his life in this world, and is come unto an everlasting kingdom.\end{Resp}
\begin{Resp}\Rsign For thee have I seen righteous before Me among all people, alleluia.\end{Resp}
\begin{Resp}\Vsign Glory be to the Father, and to the Son, \Star{} and to the Holy Ghost.\end{Resp}
\begin{Resp}\Rsign For thee have I seen righteous before Me among all people, alleluia.\end{Resp}
\switchcolumn*

\LA
\LessonHead{Lectio IX (pro commemoratione)}
\switchcolumn
\EN
\LessonHead{Lesson IX (when commemorated)}
\switchcolumn*

\LA
\LessonTitle{Commemoratio: S. Philippi Neri Confessoris}
\switchcolumn
\EN
\LessonTitle{Commemoration of: St. Philip Neri, Confessor}
\switchcolumn*

\LA
\noindent Philíppus Nérius, piis honestísque paréntibus Floréntiæ natus, ampla pátrui hereditáte dimíssa, Romam se cóntulit; ubi, philosophía ac sacris lítteris erudítus, totum se Christo dicávit. Sacérdos ex obediéntia factus, in animárum salúte procuránda totus fuit, et in confessiónibus audiéndis ad extrémum usque diem persevérans, innúmeros pene fílios Christo péperit; quos verbi Dei cotidiáno pábulo, sacramentórum frequéntia, oratiónis assiduitáte aliísque piis exercitatiónibus enutríri cúpiens, Oratórii congregatiónem instítuit. Caritáte Dei vulnerátum tanto cor ejus æstuábat ardóre, ut, cum intra fines suos continéri non posset, sinum, confráctis atque elátis duábus cóstulis, mirabíliter Dóminus ampliáverit. Prophetíæ dono fuit illústris, et in animórum sénsibus penetrándis mirífice enítuit. Virginitátem perpétuo illibátam servávit; idque assecútus est, ut eos, qui puritátem cólerent, ex odóre, qui vero secus, ex fœtóre dignósceret. Anno salútis millésimo quingentésimo nonagésimo quinto, octogenárius obdormívit in Dómino.\par
\switchcolumn
\EN
\noindent Philip Neri was born in Florence of good and devout parents. Giving up a large inheritance from his uncle, he went to Rome, where he studied philosophy and the sacred sciences and dedicated himself wholly to Christ. He became a priest out of obedience and gave himself up completely to the saving of souls. Through hearing confessions, in which he persevered to the last day of his life, he brought forth innumerable sons for Christ. Desiring to nourish them with the daily bread of God's word, with frequent reception of the sacraments, with constant prayer, and with other exercises of piety, he founded the Congregation of the Oratory. His heart was wounded by the love of God, burning with such ardour that it could only be contained within his breast because the Lord miraculously enlarged the breast by breaking two of his ribs, and forming an arch over the heart. Philip was famed for the gift of prophecy and for his wonderful penetration of the thoughts of men's hearts. He kept his virginity always intact; and he had the gift of distinguishing those who cultivated purity by a good odour, and those who did not by a stench. At the age of eighty, in the year of salvation 1595, he fell asleep in the Lord.\par
\switchcolumn*

\LA
\TeDeum{Te Deum}
\switchcolumn
\EN
\TeDeum{Te Deum}
\switchcolumn*

\end{paracol}

\par\needspace{10\baselineskip}\addvspace{1.2em}{\centering\small\textsc{Die 26 Maii} · \textsc{26 May}\par}
\Office{S. Eleutherii Papæ et Martyris}{St. Eleutherius, Pope and Martyr}{Simplex}
\addcontentsline{toc}{subsection}{Die 26 Maii · S. Eleutherii Papæ et Martyris\enspace\textit{St. Eleutherius, Pope and Martyr}}
\begin{paracol}{2}
\LA
\LessonHead{Lectio IX (pro commemoratione)}
\switchcolumn
\EN
\LessonHead{Lesson IX (when commemorated)}
\switchcolumn*

\LA
\LessonTitle{Commemoratio: S. Eleutherii Papæ et Martyris}
\switchcolumn
\EN
\LessonTitle{Commemoration of: St. Eleutherius, Pope and Martyr}
\switchcolumn*

\LA
\noindent Eleutherius, Nicopoli in Græcia natus, Aniceti Pontificis diaconus, Commodo imperatóre, præfuit Ecclésiæ. Huic, inítio pontificatus, supplices litteræ venérunta Lucio Britannórum rege, ut se ac suos in Christianórum númerum reciperet. Quam ob rem Fugatium et Damianum, doctos et pios viros, misit in Britanniam, per quos rex et réliqui fidem susciperent. Hoc Pontifice, Irenæus Polycarpi discipulus, Romam veniens, ab eo benígne acceptus est. Quo témpore summa pace et quiete fruebátur Ecclésia Dei; ac per totum orbem terrárum, maxime Romæ, fides propagabátur. Vixit Eleutherius in pontificatu annos quindecim, dies viginti tres. Fecit ordinatiónes tres mense Decembri, quibus creávit presbyteros duodecim, diaconos octo, episcopos per diversa loca quindecim: sepultúsque est in Vaticano prope corpus sancti Petri.\par
\switchcolumn
\EN
\noindent Eleutherius was a Greek by race, and was born at Nicopolis, a city of Epirus. His father's name was Abundius. He was a Priest of the holy Roman Church. In the year of our Lord 179, during the reign of the Emperor Marcus Aurelius Augustus, after the death of Soter, he was chosen Bishop of Rome by the votes of all the clergy. He discharged the duties of this office excellently, and with all praise, for fifteen years and twenty-three days. He held three ordinations in the month of December, wherein he ordained twelve Priests, eight Deacons, and fifteen Bishops for diverse places. He was consulted by the church of Lyons by letter concerning certain questions. The bearer of these letters was that most learned Irenaeus. Him he hospitably welcomed, and to him he opened the traditions of the Apostles which the Church of Rome had kept pure. He condemned the superstitious dry-meat system of the Montanists. He laid down excellent rules as to the right form of church sentences. When Marcion and Valentine had oftentimes relapsed he cast them out of the Church. In his days the Church enjoyed the utmost peace, and he brought many even of the chiefest of Rome to believe in Christ. He received letters by messengers from Lleurwg, Prince of the Britons (of Morganwg,) praying him for ministers of the Word of God, and he sent unto him Fagan and Dyfan, Priests of the Roman Church, through whose hands the Prince himself, with his whole household and nearly all his subjects, were born again in the sacred washing of regeneration. At length, when he had done all these things and others for the worship of God, Eleutherius died a holy death upon the 28th day of May, (in the year of our Lord 192,) in the reign of the Emperor Commodus, and was buried upon the Vatican Mount.\par
\switchcolumn*

\LA
\TeDeum{Te Deum}
\switchcolumn
\EN
\TeDeum{Te Deum}
\switchcolumn*

\end{paracol}

\par\needspace{10\baselineskip}\addvspace{1.2em}{\centering\small\textsc{Die 27 Maii} · \textsc{27 May}\par}
\Office{S. Bedæ Venerabilis Confessoris et Ecclesiæ Doctoris}{St. Bede the Venerable, Confessor and Doctor of the Church}{Duplex}
\addcontentsline{toc}{subsection}{Die 27 Maii · S. Bedæ Venerabilis Confessoris et Ecclesiæ Doctoris\enspace\textit{St. Bede the Venerable, Confessor and Doctor of the Church}}
\begin{paracol}{2}
\LA
\RefLine{Lectiones I–III: de Scriptura occurrente.}
\switchcolumn
\EN
\RefLine{Lessons I–III: from the Scripture of the day.}
\switchcolumn*

\LA
\RefLine{Lectio IX: de S. Joannis Papæ et Martyris (05-27o).}
\switchcolumn
\EN
\RefLine{Lesson IX: of St. John I, Pope and Martyr (05-27o).}
\switchcolumn*

\LA
\LessonHead{Lectio IV}
\switchcolumn
\EN
\LessonHead{Lesson IV}
\switchcolumn*

\LA
\noindent Beda presbyter, Girvi in Britanniæ et Scotiæ finibus ortus, septennis sancto Benedicto Biscopio abbati Wiremuthensi educandus traditur. Monachus deínde factus, vitam sic instituit, ut, dum se artium et doctrinarum studiis totum impenderet, nihil umquam de regulari disciplina remitteret. Nullum fuit doctrinæ genus, in quo non esset diligentissime versatus; sed præcipua illi cura divinarum Scripturarum meditatio, quarum sententiam ut plenius assequeretur, Græci Hebraicique sermonis notitiam est adeptus. Trigesimo ætatis anno abbatis sui jussu sacerdos initiatus, statim, suasore Acca Hagulstadensi episcopo sacros explanare libros aggressus est: in quo sanctorum Patrum doctrinis adeo inhæsit, ut nihil proferret nisi illorum judicio comprobatum, eorumdem étiam fere verbis usus. Otium perosus semper, ex lectione ad orationem transibat ac vicissim ex oratione ad lectionem: in qua adeo animo inflammabatur, ut sæpe inter legendum et docendum lacrimis perfunderetur. Ne autem rerum fluxarum curis distraheretur, delatum abbatis munus constantissime detrectavit.\par
\switchcolumn
\EN
\noindent Bede, a priest, was born at Jarrow, on the borders of England and Scotland. At the age of seven years he was placed under the care of holy Benedict Biscop, Abbot of Wearmouth, to be educated. Thereafter he became a monk, and so ordered his life that, whilst he should devote himself wholly to the study of the sciences and of doctrine, he might in nothing relax the discipline of his Order. There was no branch of learning in which he was not most thoroughly versed, but his chief care was the study of Holy Scriptures; and that he might the better understand them he acquired a knowledge of the Greek and Hebrew tongues. When he was thirty years of age he was ordained priest at the command of his Abbot, and immediately, on the advice of Acca, Bishop of Hexham, undertook the work of expounding the Sacred Books. In his interpretations he so strictly adhered to the teaching of the holy Fathers that he would advance nothing which was not approved by their judgment, nay, had the warrant of their very words. He ever hated sloth, and by habitually passing from reading to prayer, and in turn from prayer to reading, he so inflamed his soul that often amid his reading and teaching he was bathed in tears. Lest also his mind should be distracted by the cares of transitory things, he never would take the office of Abbot when it was offered to him.\par
\switchcolumn*

\LA
\begin{Resp}\Rsign Honéstum fecit illum Dóminus, et custodívit eum ab inimícis, et a seductóribus tutávit illum: \Star{} Et dedit illi claritátem ætérnam, allelúja.\end{Resp}
\begin{Resp}\Vsign Justum dedúxit Dóminus per vias rectas, et osténdit illi regnum Dei.\end{Resp}
\begin{Resp}\Rsign Et dedit illi claritátem ætérnam, allelúja.\end{Resp}
\switchcolumn
\EN
\begin{Resp}\Rsign The Lord made him honourable, and defended him from his enemies, and kept him safe from those that lay in wait for him; \Star{} And gave him perpetual glory, alleluia.\end{Resp}
\begin{Resp}\Vsign He went down with him into the pit, and left him not in bonds.\end{Resp}
\begin{Resp}\Rsign And gave him perpetual glory, alleluia.\end{Resp}
\switchcolumn*

\LA
\LessonHead{Lectio V}
\switchcolumn
\EN
\LessonHead{Lesson V}
\switchcolumn*

\LA
\noindent Scientiæ ac pietatis laude Bedæ nomen sic brevi claruit, ut sanctus Sergius Papa de eo Romam arcessendo cogitaverit; quo difficillimis scilicet, quæ de rebus sacris exortæ erant, quæstionibus definiendis conferret operam. Emendandis fidelium moribus, fidei vindicandæ atque asserendæ libros plures conscripsit, quibus tantam sui apud omnes opinionem fecit, ut illum sanctus Bonifatius episcopus et martyr, Ecclesiæ lumen prædicaverit; Lanfrancus, Anglorum doctorem; concilium Aquisgranense, doctorem admirabilem dixerit. Quin ejus scripta eo adhuc vivente, publice in ecclesiis legebantur. Quod cum fieret, quoniam ipsum sanctum minime appellare liceret, Venerabilis titulo efferebant; qui deínde veluti proprius secutis étiam temporibus semper habitus est. Ejus autem doctrinæ eo vis efficacior erat, quod vitæ sanctimonia religiosisque virtutibus confirmabatur. Quam ob rem discipulos, quos multos et egregios imbuendos habuit, studio et exemplo non litteris modo atque scientiis, sed étiam sanctitate fecit insignes.\par
\switchcolumn
\EN
\noindent The name of Bede soon became so famous for learning and piety that St. Sergius the Pope thought of calling him to Rome, where, certainly, he might have helped to solve the very difficult questions which had then arisen concerning sacred things. He wrote many books for the bettering of the lives of the faithful, and defending and extending of the faith. By those he gained everywhere such a reputation that the holy martyr Bishop Boniface styled him a Light of the Church; Lanfranc called him The Teacher of the English, and the Council of Aix-la-Chapelle The Admirable Doctor. But as his writings were publicly read in the churches during his life, and as it was not allowable to call him already a saint, they named him The Venerable, a title which in all times after has remained peculiarly his. The power of his teaching was the greater also, in that it was attested by a holy life and the graces of religious observance. In this way, by his earnestness and example, his disciples, who were many and distinguished, were made eminent, not only in letters and the sciences, but in personal holiness.\par
\switchcolumn*

\LA
\begin{Resp}\Rsign Amávit eum Dóminus, et ornávit eum: stolam glóriæ índuit eum, \Star{} Et ad portas paradísi coronávit eum, allelúja.\end{Resp}
\begin{Resp}\Vsign Índuit eum Dóminus lorícam fídei, et ornávit eum.\end{Resp}
\begin{Resp}\Rsign Et ad portas paradísi coronávit eum, allelúja.\end{Resp}
\switchcolumn
\EN
\begin{Resp}\Rsign The Lord loved him and beautified him; He clothed him with a robe of glory, \Star{} And crowned him at the gates of Paradise, alleluia.\end{Resp}
\begin{Resp}\Vsign The Lord hath put on him the breast-plate of faith, and hath adorned him.\end{Resp}
\begin{Resp}\Rsign And crowned him at the gates of Paradise, alleluia.\end{Resp}
\switchcolumn*

\LA
\LessonHead{Lectio VI}
\switchcolumn
\EN
\LessonHead{Lesson VI}
\switchcolumn*

\LA
\noindent Aetate demum et laboribus fractus, gravi morbo correptus est. Quo cum amplius quinquaginta dies detentus esset, consuetum orandi morem Scripturasque interpretandi non intercepit; eo namque tempore Evangelium Joannis in popularium suorum usum Anglice vertit. Cum autem in Ascensionis præludio instare sibi mortem persentiret, supremis Ecclesiæ sacramentis muniri voluit; tum, sodales amplexatus, atque humi super cilicio stratus, cum illa verba ingeminaret, Gloria Patri, et Filio, et Spiritui Sancto, obdormivit in Domino. Ejus corpus suavissimum uti fertur, spirans odorem sepultum est in monasterio Girvensi, ac postea Dunclinum cum sancti Cuthberti, reliquiis translatum. Eum tamquam Doctorem a Benedictinis aliisque religiosis familiis ac diœcesibus cultum, Leo decimus tertius Pontifex maximus ex sacrorum Rituum Congregationis consulto, universalis Ecclesiæ Doctorem declaravit, et festo ipsius die Missam et Officium de Doctoribus ab omnibus recitari decrevit.\par
\switchcolumn
\EN
\noindent Broken at length by age and labour, he was seized by a grievous illness. Though he suffered under it for more than seven weeks, he ceased not from his prayers and his interpreting of the Scriptures; for at that time he was turning the Gospel of John into English for the use of his people. But when, on the Eve of the Ascension, he perceived that death was coming upon him, he desired to be fortified with the last sacraments of the Church: then, after he had embraced his companions, and was laid on a piece of sackcloth on the ground, he repeated the words, Glory be to the Father, and to the Son, and to the Holy Ghost, and fell asleep in the Lord. His body, very sweet, as it is related, breathing sweet odour, was buried in the monastery of Jarrow, and afterwards was translated to Durham with the relics of St. Cuthbert. Bede, who was already a Doctor among the Benedictines, and in other religious Orders, and venerated in certain dioceses, was declared by Pope Leo XIII., after consulting with the Congregation of Sacred Rites, to be a Doctor of the universal Church; and the Mass and Office for Doctors was ordered to be recited by all on his feast-day.\par
\switchcolumn*

\LA
\begin{Resp}\Rsign Iste homo perfécit ómnia quæ locútus est ei Deus, et dixit ad eum: Ingrédere in réquiem meam: \Star{} Quia te vidi justum coram me ex ómnibus géntibus, allelúja.\end{Resp}
\begin{Resp}\Vsign Iste est, qui contémpsit vitam mundi, et pervénit ad cæléstia regna.\end{Resp}
\begin{Resp}\Rsign Quia te vidi justum coram me ex ómnibus géntibus, allelúja.\end{Resp}
\begin{Resp}\Vsign Glória Patri, et Fílio, \Star{} et Spirítui Sancto.\end{Resp}
\begin{Resp}\Rsign Quia te vidi justum coram me ex ómnibus géntibus, allelúja.\end{Resp}
\switchcolumn
\EN
\begin{Resp}\Rsign This is he which did according unto all that God commanded him; and God said unto him: Enter thou into My rest, \Star{} For thee have I seen righteous before Me among all people, alleluia.\end{Resp}
\begin{Resp}\Vsign This is he which loved not his life in this world, and is come unto an everlasting kingdom.\end{Resp}
\begin{Resp}\Rsign For thee have I seen righteous before Me among all people, alleluia.\end{Resp}
\begin{Resp}\Vsign Glory be to the Father, and to the Son, \Star{} and to the Holy Ghost.\end{Resp}
\begin{Resp}\Rsign For thee have I seen righteous before Me among all people, alleluia.\end{Resp}
\switchcolumn*

\LA
\LessonHead{Lectio VII}
\switchcolumn
\EN
\LessonHead{Lesson VII}
\switchcolumn*

\LA
\LessonTitle{Léctio sancti Evangélii secúndum Matthǽum}
\switchcolumn
\EN
\LessonTitle{From the Holy Gospel according to Matthew}
\switchcolumn*

\LA
\Cite{Matt 5:13-19}
\switchcolumn
\EN
\Cite{Matt 5:13-19}
\switchcolumn*

\LA
\noindent In illo témpore: Dixit Jesus discipulis suis: Vos estis sal terræ. Quod si sal evanuerit, in quo salietur? Et réliqua.\par
\switchcolumn
\EN
\noindent AT that time Jesus said to His disciples: Ye are the salt of the earth. But if the salt have lost its savour, wherewith shall it be salted? And so on.\par
\switchcolumn*

\LA
\LessonTitle{Homilía sancti Bedæ Venerábilis Presbýteri}
\switchcolumn
\EN
\LessonTitle{Homily of the Venerable Bede, Priest.}
\switchcolumn*

\LA
\Source{In Evang. Vos estis sal terra}
\switchcolumn
\EN
\Source{On the Gospel}
\switchcolumn*

\LA
\noindent In terra, humana natura; in sale, sapientia verbis significatur. Salis enim natura, terra efficitur infructuosa; unde quasdam urbes legimus, victorum ira, sale seminatas. Et hoc convenit apostolicæ doctrinæ, ut sale sapientiæ compescat in terra humanæ carnis luxum sæculi aut fœditatem vitiorum germinare. Quod si sal evanuerit, in quo salietur? Id est, si vos, per quos condiendi sunt populi, propter metum persecutionum, aut terrorem, amiseritis regna cælorum, extra Ecclesiam positi, inimicorum opprobria sustinetis non dubium. Vos estis lux mundi: id est, vos, quia vera luce illuminati estis, lux eis qui in mundo sunt, esse debetis. Non potest civitas abscondi supra montem posita: id est, apostolica doctrina super Christum fundata, sive Ecclesia super Christum ex multis gentibus fidei unitate constructa et caritatis bitumine conglutinata; quæ sit tuta intrantibus; et laboriosa adeuntibus, habitatores custodit, et omnes inimicos secludit.\par
\switchcolumn
\EN
\noindent The Gospel saith, Ye are the salt of the earth. In these words the earth signifies human nature, and the salt signifies wisdom. Salt, verily, by its nature renders the earth unfruitful. Hence we read of cities, which in the anger of their victors were sown with salt. And hereto agreeth the teaching of the Apostle that by the salt of wisdom the lust of this world is restrained in the earth of human flesh, lest the foulness of vice should sprout up. But what if the salt shall have lost its savour? That is to say If you, by whom the people are to be seasoned, are, on account of fear of persecution, or terror, you should lose the kingdom of heaven, placed outside the Church, there is no doubt that you will incur the taunts of the enemy. Ye are the light of the world that is to say You, because ye are enlightened by the true light, ought to be the light of them who are in the world. A city set on an hill cannot be hid; that is to say The Apostles' teaching, founded upon Christ; in other words, the Church built upon Christ, out of many nations, in the unity of the faith, and bound together with the cement of love; to those who enter it, a place of safety; to those who go up to it, toilsome; the guardian of those who dwell in it, and excluding every enemy.\par
\switchcolumn*

\LA
\begin{Resp}\Rsign Iste est qui ante Deum magnas virtútes operátus est, et de omni corde suo laudávit Dóminum: \Star{} Ipse intercédat pro peccátis ómnium populórum, allelúja.\end{Resp}
\begin{Resp}\Vsign Ecce homo sine queréla, verus Dei cultor, ábstinens se ab omni ópere malo, et pérmanens in innocéntia sua.\end{Resp}
\begin{Resp}\Rsign Ipse intercédat pro peccátis ómnium populórum, allelúja.\end{Resp}
\switchcolumn
\EN
\begin{Resp}\Rsign This is he which wrought great wonders before God, and praised the Lord with all his heart. \Star{} May he pray for all people, that their sins may be forgiven unto them, alleluia.\end{Resp}
\begin{Resp}\Vsign Behold a man without blame, a worshipper of God in truth, keeping himself clean from every evil work, and abiding still in his innocency.\end{Resp}
\begin{Resp}\Rsign May he pray for all people, that their sins may be forgiven unto them, alleluia.\end{Resp}
\switchcolumn*

\LA
\LessonHead{Lectio VIII}
\switchcolumn
\EN
\LessonHead{Lesson VIII}
\switchcolumn*

\LA
\noindent Neque accendunt lucernam et ponunt eam sub modio, sed super candelabrum. Sub modio ergo lucernam ponit quisquis lucem doctrinæ commodis temporalibus obscurat et tegit; super candelabrum vero, qui se ita ministerio Dei subjicit, ut superior sit doctrina veritatis quam servitus corporis. Aliter Salvator accendit lucernam, qui humanæ testam naturæ flamma suæ divinitatis implevit; et hanc super candelabrum, id est Ecclesiam, posuit, quod in frontibus nostris fidem suæ incarnationis fixit. Quæ lucerna non potuit sub modio poni, id est sub mensura legis includi; nec in sola Judǽa, sed in universo illuxit orbe.\par
\switchcolumn
\EN
\noindent It either doth any man light a candle and put it under a bushel; but upon a candlestick. So he who puts the light under the bushel is he who for his own temporal ends would hide and tamper with the light of doctrine; but upon the candlestick he places it who follows the ministry of God in order that the teaching of the truth may be accounted a greater thing than the service of the body. In another aspect, the Saviour lighted the candle when He filled our mortal body with the flame of the God-head; and He placed it on a candlestick, that is the Church; for He fixed the faith of His incarnation upon our foreheads. Which light cannot be placed under a bushel; that is to say, it cannot be included within the measures of the law, nor in Judea alone, but has lightened the whole earth.\par
\switchcolumn*

\LA
\begin{Resp}\Rsign In médio Ecclésiæ apéruit os ejus, \Star{} Et implévit eum Dóminus spíritu sapiéntiæ et intelléctus, allelúja.\end{Resp}
\begin{Resp}\Vsign Jucunditátem et exsultatiónem thesaurizávit super eum.\end{Resp}
\begin{Resp}\Rsign Et implévit eum Dóminus spíritu sapiéntiæ et intelléctus, allelúja.\end{Resp}
\begin{Resp}\Vsign Glória Patri, et Fílio, \Star{} et Spirítui Sancto.\end{Resp}
\begin{Resp}\Rsign Et implévit eum Dóminus spíritu sapiéntiæ et intelléctus, allelúja.\end{Resp}
\switchcolumn
\EN
\begin{Resp}\Rsign In the midst of the congregation did the Lord open his mouth. \Star{} And filled him with the spirit of wisdom and understanding, alleluia.\end{Resp}
\begin{Resp}\Vsign He made him rich with joy and gladness.\end{Resp}
\begin{Resp}\Rsign And filled him with the spirit of wisdom and understanding, alleluia.\end{Resp}
\begin{Resp}\Vsign Glory be to the Father, and to the Son, \Star{} and to the Holy Ghost.\end{Resp}
\begin{Resp}\Rsign And filled him with the spirit of wisdom and understanding, alleluia.\end{Resp}
\switchcolumn*

\LA
\LessonHead{Lectio IX (pro commemoratione)}
\switchcolumn
\EN
\LessonHead{Lesson IX (when commemorated)}
\switchcolumn*

\LA
\LessonTitle{Commemoratio: S. Bedæ Venerabilis Confessoris et Ecclesiæ Doctoris}
\switchcolumn
\EN
\LessonTitle{Commemoration of: St. Bede the Venerable, Confessor and Doctor of the Church}
\switchcolumn*

\LA
\noindent Beda présbyter, Girvi in Británniæ et Scótiæ fínibus ortus est. Mónachus factus, vitam sic instítuit, ut, dum se ártium et doctrinárum stúdiis totum impénderet, nihil umquam de regulári disciplína remítteret. Nullum fuit doctrínæ genus, in quo non esset diligentíssime versátus; sed præcípua illi cura fuit divinárum Scripturárum meditátio, ita ut, sacerdótio initiátus, sacros explanáre libros aggréssus sit; in quo sanctórum Patrum doctrínis ádeo inhǽsit, ut nihil proférret nisi illórum judício comprobátum, eorúmdem étiam fere verbis usus. Otium perósus semper, ex lectióne ad oratiónem transíbat, ac vicíssim ex oratióne ad lectiónem. Emendándis fidélium móribus, fídei vindicándæ atque asseréndæ libros plures conscrípsit, quibus tantam sui apud omnes opiniónem fecit, ut ejus scripta, eo adhuc vivénte públice in ecclésiis legeréntur. Ætáte demum et labóribus fractus, pie obdormívit in Dómino. Eum Leo décimus tértius universális Ecclésiæ Doctórem declarávit.\par
\switchcolumn
\EN
\noindent Bede the priest was born at Jarrow, on the borders of England and Scotland. When a monk, he so arranged his life as to devote himself completely to the study of the liberal arts and sacred doctrine, without in any way relaxing the discipline of the Rule. There was no kind of learning in which he was not thoroughly versed; but his special interest was the study of the Scriptures; and when he was made a priest, he undertook the task of explaining the holy books. In doing so, he adhered to the teaching of the holy Fathers so closely that he would say nothing not already approved by their judgment, and he even made use of their very words. Abhorring laziness, he would go straight from reading to prayer and from prayer to reading. To raise the level of morality among Christians and to defend and spread the faith, he wrote many books, which gained him such a reputation with everyone that his writings were publicly read in churches during his own lifetime. At length, worn out with age and labours, he fell asleep peacefully in the Lord. Leo XIII declared him a Doctor of the universal Church.\par
\switchcolumn*

\LA
\TeDeum{Te Deum}
\switchcolumn
\EN
\TeDeum{Te Deum}
\switchcolumn*

\end{paracol}

\par\needspace{10\baselineskip}\addvspace{1.2em}{\centering\small\textsc{Die 27 Maii} · \textsc{27 May}\par}
\Office{S. Joannis Papæ et Martyris}{St. John I, Pope and Martyr}{Simplex}
\addcontentsline{toc}{subsection}{Die 27 Maii · S. Joannis Papæ et Martyris\enspace\textit{St. John I, Pope and Martyr}}
\begin{paracol}{2}
\LA
\LessonHead{Lectio IX (pro commemoratione)}
\switchcolumn
\EN
\LessonHead{Lesson IX (when commemorated)}
\switchcolumn*

\LA
\LessonTitle{Commemoratio: S. Joannis Papæ et Martyris}
\switchcolumn
\EN
\LessonTitle{Commemoration of: St. John I, Pope and Martyr}
\switchcolumn*

\LA
\noindent Joannes Etruscus, Justino seniore imperatore, rexit Ecclesiam: ad quem profectus est Constantinopolim auxilii causa, quod Theodoricus rex hæreticus divexabat Italiam. Cujus étiam iter Deus miraculis illustravit, Nam, cum ei nobilis vir ad Corinthum, equum, quo ejus uxor mansueto utebatur, itineris causa commodasset, factum est, ut, domino postea remissus equus ita ferox evaderet, ut fremitu et totius corporis agitatione semper deinceps dominam expulerit: tamquam indignaretur mulierem recipere, ex quo sedisset in eo Jesu Christi Vicarius. Quam ob rem illi equum Pontifici donaverunt. Sed illud majus miraculum, quod Constantinopoli, in aditu portæ Aureæ inspectante frequentissimo populo, qui una cum imperatore Pontifici honoris causa occurrerat, cæco lumen restituit. Ad cujus pedes prostratus étiam imperator eum veneratus est. Rebus cum imperatore compositis, in Italiam rediit, statimque epistolam scripsit ad omnes Italiæ episcopos, jubens eos Arianorum ecclesias ad catholicum ritum consecrare, illud subjungens: Quia et nos, quando fuimus Constantinopoli tam pro religione catholica quam pro regis Theodorici causa, quascumque illis in partibus eorum ecclesias reperire potuimus, catholicas eas consecravimus. Quod iniquissimo animo ferens Theodoricus, dolo accersitum Pontificem Ravennam in carcerem conjecit; ubi, squalore inediaque afflictus, paucis diebus cessit e vita, cum sedisset annos duos, menses novem, dies quatuordecim, ordinatis eo tempore episcopis quindecim. Paulo post moritur Theodoricus: quem quidam eremita, ut scribit sanctus Gregorius, vidit inter Joánnem Pontificem et Symmachum patricium, quem idem occiderat, demergi in ignem Liparitanum; ut videlicet illi, quibus mortem attulerat, tamquam judices essent ejus interitus. Joannis corpus Ravenna Romam portatum est, et in basilica sancti Petri sepultum.\par
\switchcolumn
\EN
\noindent Pope John I was a Tuscan, who ruled the Church during the reign of the Emperor Justinian. He went to Constantinople to get help from Justinian in the troubles which the heretic King Theodoric was then causing in Italy. It pleased the Lord to mark this journey with wonders. A certain nobleman at Corinth lent to the Pope for his journey a very quiet horse on which his own wife was used to ride. But when the horse was returned to his owner he was found become so vicious, that by his restiveness and plunging he was always throwing off his mistress, as though he were not content to carry the lady after having carried the Vicar of Jesus Christ. When the nobleman and his wife found the beast to be thus worthless, they gave him for a present to the Pope. But a thing much more marvellous was that when the Pope, accompanied by the Emperor, and under the gaze of an immense multitude of people, who had come forth with Justinian to do him honour, was at the entering in of the Golden Gate of Constantinople, he gave sight to a blind man. Even the Emperor fell at his feet to show him respect. When he had arranged his business with Justinian he returned into Italy, and forthwith sent out a letter to all the Bishops of Italy, bidding them hallow for Catholic worship the churches of the Arians, and adding these words: We Ourselves when We were at Constantinople on some matters pertaining to the Catholic Religion and others pertaining to the King Theodoric, hallowed as Catholic all their Churches which We were able to find in those parts. Theodoric took this rule very ill, and, having enticed John by fraud to come to Ravenna, he cast him into prison, wherein, in a few days, he died of filth and hunger. He had sat in the chair of Peter two years, nine months, and fourteen days, within which time he had ordained fifteen Bishops. A little while afterward Theodoric also died. St. Gregory writeth that a certain hermit saw him between Pope John and Symmachus the Patrician, whom he had likewise slain, going down into the fiery crater of Lipari, as though they who had been his victims were become the judges of his punishment. The body of John was carried from Ravenna to Rome, and there buried in the Church of St. Peter.\par
\switchcolumn*

\LA
\TeDeum{Te Deum}
\switchcolumn
\EN
\TeDeum{Te Deum}
\switchcolumn*

\end{paracol}

\par\needspace{10\baselineskip}\addvspace{1.2em}{\centering\small\textsc{Die 28 Maii} · \textsc{28 May}\par}
\Office{S. Augustini Episcopi et Confessoris}{St. Augustine of Canterbury, Bishop and Confessor}{Duplex}
\addcontentsline{toc}{subsection}{Die 28 Maii · S. Augustini Episcopi et Confessoris\enspace\textit{St. Augustine of Canterbury, Bishop and Confessor}}
\begin{paracol}{2}
\LA
\RefLine{Lectiones I–III: de Scriptura occurrente.}
\switchcolumn
\EN
\RefLine{Lessons I–III: from the Scripture of the day.}
\switchcolumn*

\LA
\LessonHead{Lectio IV}
\switchcolumn
\EN
\LessonHead{Lesson IV}
\switchcolumn*

\LA
\noindent Augustinus Romæ in Lateranensi cœnobio monachus, a Gregorio Magno cum secus monachis fere quadraginta in Angliam missus est anno quingentesimo nonagesimo septimo, ut gentes illas ad Christum converteret. Erat eo tempore rex Ethelbertus in Cantio potentissimus, qui audita adventus Augustíni causa, eum cum sociis Cantuariam, sui regni metropolim, invitavit; ibique manendi et Christum prædicandi facultatem eidem liberaliter concessit. Quare sanctus vir prope Cantuariam oratorium exstruxit, ubi ipse aliquamdiu consedit, atque apostolicam vivendi rationem cum suis æmulatus est.\par
\switchcolumn
\EN
\noindent Augustine, the first Archbishop of Canterbury, and the Apostle of the English, was sent into England by blessed Gregory, and came thither in the year 597. At that time there was in Kent a most mighty king named Ethelbert, whose power reached even to the Humber. When this King had heard wherefore the holy man was come, he received him kindly, and bade him and his companions, who were all monks, to come to his own capital city of Canterbury being struck with astonishment at the perfect blamelessness of their lives, and the power of the heavenly doctrine which they preached, and which God confirmed with signs following.\par
\switchcolumn*

\LA
\begin{Resp}\Rsign Invéni David servum meum, óleo sancto meo unxi eum: \Star{} Manus enim mea auxiliábitur ei, allelúja.\end{Resp}
\begin{Resp}\Vsign Nihil profíciet inimícus in eo, et fílius iniquitátis non nocébit ei.\end{Resp}
\begin{Resp}\Rsign Manus enim mea auxiliábitur ei, allelúja.\end{Resp}
\switchcolumn
\EN
\begin{Resp}\Rsign I have found David My servant, with My holy oil have I anointed him \Star{} For My hand shall help him, alleluia.\end{Resp}
\begin{Resp}\Vsign The enemy shall prevail nothing against him, nor the son of wickedness afflict him.\end{Resp}
\begin{Resp}\Rsign For My hand shall help him, alleluia.\end{Resp}
\switchcolumn*

\LA
\LessonHead{Lectio V}
\switchcolumn
\EN
\LessonHead{Lesson V}
\switchcolumn*

\LA
\noindent Cælestis doctrina prædicatione plurimis firmata miraculis, ac vitai exemplo sic insulanos illos demulsit, ut eorum plerosque ad christianam fidem perduxerit, ac demum regem ipsum, quem cum innumero suorum comitatu sacro fonte lustravit, summa cum lætitia Berthæ regiæ uxoris, quæ christiana erat. Olim in Natali Domini, cum decem millibus et amplius baptismum in alveo fluminis Eboraci contulisset, quotquot ex iis morbo aliquo affecti erant, cum animæ salute, corporis quoque sanitatem recepisse memoriæ proditum est. Jussu Gregorii ordinatus episcopus, sedem Cantuariæ instituit in ecclesia Salvatoris a se erecta, in qua monachos operis sui subsidiarios collocavit; et sancti Petri monasterium, quod postea et a suo nomine dictum est, in suburbanis construxit. Idem Gregorius usum pallii cum facultate ecclesiasticæ hierarchiæ in Anglia instituendas ei concessit, quo novam étiam operariorum manum misit, nempe Mellitum, Justum, Paulinum, et Rufinianum.\par
\switchcolumn
\EN
\noindent They drew nigh to the city in solemn procession, singing the Litany, and bearing before them for their standard a silver cross and a picture of the Lord our Saviour painted on a panel. Hard by the city, upon the east side, there was a Church builded of old time in honour of St. Martin, and wherein the Queen, who was a Christian, was used to pray. There they first began to meet together, to sing, to pray, to celebrate Masses, to preach, and to baptize, until the King was turned to the faith, and the most part of his people were led by his example, (but not his authority,) to take the name of Christian, for he had learnt from his teachers and his own soul's physicians, that men are to be drawn, and not driven to heaven. And now Augustine, being ordained Archbishop of the English and of Britain, lest he should leave untravailed any part of the Lord's vineyard, asked from the Apostolic See a new band of labourers, Mellitus, Justus, Paulinus, and Rufinian.\par
\switchcolumn*

\LA
\begin{Resp}\Rsign Pósui adjutórium super poténtem, et exaltávi eléctum de plebe mea: \Star{} Manus enim mea auxiliábitur ei, allelúja.\end{Resp}
\begin{Resp}\Vsign Invéni David servum meum, óleo sancto meo unxi eum.\end{Resp}
\begin{Resp}\Rsign Manus enim mea auxiliábitur ei, allelúja.\end{Resp}
\switchcolumn
\EN
\begin{Resp}\Rsign I have laid help upon one that is mighty, and have exalted one chosen out of My people \Star{} For My hand shall help him, alleluia.\end{Resp}
\begin{Resp}\Vsign I have found David My servant, with My holy oil have I anointed him.\end{Resp}
\begin{Resp}\Rsign For My hand shall help him, alleluia.\end{Resp}
\switchcolumn*

\LA
\LessonHead{Lectio VI}
\switchcolumn
\EN
\LessonHead{Lesson VI}
\switchcolumn*

\LA
\noindent Dispositis ejus ecclesiæ rebus, synodum habuit Augustinus cum episcopis atque doctoribus veterum Britonum, qui in Paschæ celebratione aliusque ritibus ab Ecclesia Romana jamdudum dissidebant. Sed cum eos neque Apostolicæ Sedis auctoritate, neque miraculis movere posset ut dissidio cessarent, prophetico spiritu eis excidium prænuntiavit. Denique maximis pro Christo exantlutis laboribus, miraculis clarus, cum Mellitum Londinensi ecclesiæ præfecisset, Justum Roffensi, suæ Laurentium, in cælum migravit septimo Kalendas Junias, Ethelberto regnante, ac sepultus est in monasterio sancti Petri, quod exinde Cantuariensium Antistitum et aliquot regum conditorium fuit. Ejus cultum ferventi studio prosecutæ sunt Anglorum gentes, ac Leo decimus tertius Pontifex Maximus ejus Oificium et Missam ad universam extendit Ecclesiam.\par
\switchcolumn
\EN
\noindent Having arranged the affairs of his church, Augustine held a synod with the bishops and doctors of the ancient Britons, who had long been at variance with the Roman Church in the celebration of Easter and other rites. But since he could not move them, either by the authority of the apostolic see or by miracles, to put an end to these variations, in a prophetic spirit he foretold their ruin. At length, after having endured many difficulties for Christ, and having become noted for miracles, when he had placed Mellitus in charge of the church of London, Justus of that of Rochester, and Laurence in charge of his own church, he passed to heaven on the 26th day of May, in the reign of Ethelbert, and was buried in the monastery of St. Peter, which thereafter became the burying-place of the bishops of Canterbury and of some kings. The English people honoured his memory with fervent zeal; and the Supreme Pontiff Leo XIII extended his Office and Mass to the universal Church.\par
\switchcolumn*

\LA
\begin{Resp}\Rsign Iste est, qui ante Deum magnas virtútes operátus est, et omnis terra doctrína ejus repléta est: \Star{} Ipse intercédat pro peccátis ómnium populórum, allelúja.\end{Resp}
\begin{Resp}\Vsign Iste est, qui contémpsit vitam mundi, et pervénit ad cæléstia regna.\end{Resp}
\begin{Resp}\Rsign Ipse intercédat pro peccátis ómnium populórum, allelúja.\end{Resp}
\begin{Resp}\Vsign Glória Patri, et Fílio, \Star{} et Spirítui Sancto.\end{Resp}
\begin{Resp}\Rsign Ipse intercédat pro peccátis ómnium populórum, allelúja.\end{Resp}
\switchcolumn
\EN
\begin{Resp}\Rsign This is he which wrought great wonders before God, and the whole earth is full of his teaching \Star{} May he pray for all people, that their sins may be forgiven unto them, alleluia.\end{Resp}
\begin{Resp}\Vsign This is he which loved not his life in this world, and hath attained unto the kingdom of heaven.\end{Resp}
\begin{Resp}\Rsign May he pray for all people, that their sins may be forgiven unto them, alleluia.\end{Resp}
\begin{Resp}\Vsign Glory be to the Father, and to the Son, \Star{} and to the Holy Ghost.\end{Resp}
\begin{Resp}\Rsign May he pray for all people, that their sins may be forgiven unto them, alleluia.\end{Resp}
\switchcolumn*

\LA
\LessonHead{Lectio VII}
\switchcolumn
\EN
\LessonHead{Lesson VII}
\switchcolumn*

\LA
\LessonTitle{Léctio sancti Evangélii secúndum Lucam.}
\switchcolumn
\EN
\LessonTitle{From the Holy Gospel according to Luke}
\switchcolumn*

\LA
\Cite{Luc 10:1-9}
\switchcolumn
\EN
\Cite{Luke 10:1-9}
\switchcolumn*

\LA
\noindent In illo témpore: Designavit Dominus et alios septuaginta duos: et misit illos binos ante faciem suam, in omnem civitatem et locum, quo erat ipse venturus. Et reliqua.\par
\switchcolumn
\EN
\noindent At that time, the Lord appointed other seventy-two also, and sent them two and two before His face into every city and place, whither He Himself would come. And so on.\par
\switchcolumn*

\LA
\LessonTitle{Homilia sancti Gregorii Papæ.}
\switchcolumn
\EN
\LessonTitle{Homily by Pope St. Gregory (the Great)}
\switchcolumn*

\LA
\Source{Homilia 17. in Evang.}
\switchcolumn
\EN
\Source{17th Homily on the Gospels}
\switchcolumn*

\LA
\noindent Dominus et Salvator noster, fratres carissimi, aliquando nos sermonibus, aliquando vero operibus admonet. Ipsa etenim facta ejus præcepta sunt: quia dum aliquid tacitus facit, quid agere debeamus innotescit. Ecce enim binos in prædicationem discipulos mittit: quia duo sunt præcepta caritatis, Dei videlicet amor, et proximi: et minus quam inter duos caritas haberi non potest. Nemo enim proprie ad semetipsum habere caritatem dicitur: sed dilectio in alterum tendit, ut caritas esse possit.\par
\switchcolumn
\EN
\noindent Dearly beloved brethren, our Lord and Saviour doth sometimes admonish us by words, and sometimes by works. Yea, His very works do themselves teach us for that which He doth silently His example still moveth us to copy. Behold how He sendeth forth His disciples to preach by two and two since there are two commandments to love, that is, a commandment to love God, and a commandment to love our neighbour and where there are not two, the one, being alone, hath not whereon to do the Lord's commandment. And no man can properly be said to love himself: for love tendeth outward toward our neighbour, if it be thelove whereto the Gospel doth oblige us.\par
\switchcolumn*

\LA
\begin{Resp}\Rsign Amávit eum Dóminus, et ornávit eum: stolam glóriæ índuit eum, \Star{} Et ad portas paradísi coronávit eum, allelúja.\end{Resp}
\begin{Resp}\Vsign Induit eum Dóminus lorícam fídei, et ornávit eum.\end{Resp}
\begin{Resp}\Rsign Et ad portas paradísi coronávit eum, allelúja.\end{Resp}
\switchcolumn
\EN
\begin{Resp}\Rsign The Lord loved him and beautified him He clothed him with a robe of glory \Star{} And crowned him at the gates of Paradise, alleluia.\end{Resp}
\begin{Resp}\Vsign The Lord hath put on him the breast-plate of faith, and hath adorned him.\end{Resp}
\begin{Resp}\Rsign And crowned him at the gates of Paradise, alleluia.\end{Resp}
\switchcolumn*

\LA
\LessonHead{Lectio VIII}
\switchcolumn
\EN
\LessonHead{Lesson VIII}
\switchcolumn*

\LA
\noindent Ecce enim binos ad prædicandum discipulos Dominus mittit: quatenus hoc nobis tacitus innuat, quia qui caritatem erga alterum non habet, prædicationis officium suscipere nullatenus debet. Bene autem dicitur, quia misit eos ante faciem suam in omnem civitatem et locum, quo erat ipse venturus. Prædicatores enim suos Dominus sequitur: quia prædicatio prævenit, et tunc ad mentis nostræ habitaculum Dominus venit, quando verba exhortationis præcurrunt: atque per hoc veritas in mente suscipitur.\par
\switchcolumn
\EN
\noindent Behold, the Lord sendeth forth His disciples to preach by two and two and thus doing, He doth silently teach us that whosoever loveth not his neighbour, such an one it behoveth not to take upon him the office of a preacher. Well also is it said that He sent them before His face into every city and place whither He Himself would come. The Lord followeth His preachers first cometh preaching, and then the Lord Himself cometh to the house of our mind, whither the word of exhortation hath come before and so cometh the truth into our mind.\par
\switchcolumn*

\LA
\begin{Resp}\Rsign Sint lumbi vestri præcíncti, et lucérnæ ardéntes in mánibus vestris: \Star{} Et vos símiles homínibus exspectántibus dóminum suum, quando revertátur a núptiis, allelúja.\end{Resp}
\begin{Resp}\Vsign Vigiláte ergo, quia nescítis qua hora Dóminus vester ventúrus sit.\end{Resp}
\begin{Resp}\Rsign Et vos símiles homínibus exspectántibus dóminum suum, quando revertátur a núptiis, allelúja.\end{Resp}
\begin{Resp}\Vsign Glória Patri, et Fílio, \Star{} et Spirítui Sancto.\end{Resp}
\begin{Resp}\Rsign Et vos símiles homínibus exspectántibus dóminum suum, quando revertátur a núptiis, allelúja.\end{Resp}
\switchcolumn
\EN
\begin{Resp}\Rsign Let your loins be girded about, and your lights burning; \Star{} And ye yourselves like unto men that wait for their lord, when he will return from the wedding, alleluia.\end{Resp}
\begin{Resp}\Vsign Watch therefore, for ye know not what hour your Lord doth come.\end{Resp}
\begin{Resp}\Rsign And ye yourselves like unto men that wait for their lord, when he will return from the wedding, alleluia.\end{Resp}
\begin{Resp}\Vsign Glory be to the Father, and to the Son, \Star{} and to the Holy Ghost.\end{Resp}
\begin{Resp}\Rsign And ye yourselves like unto men that wait for their lord, when he will return from the wedding, alleluia.\end{Resp}
\switchcolumn*

\LA
\LessonHead{Lectio IX}
\switchcolumn
\EN
\LessonHead{Lesson IX}
\switchcolumn*

\LA
\noindent Hinc namque eisdem prædicatoribus Isaias dicit: Parate viam Domini, rectas facite semitas Dei nostri. Hinc filiis Psalmista ait: Iter facite ei, qui ascendit super occasum. Super occasum namque Dominus ascendit: quia unde in passione occubuit, inde majorem suam gloriam resurgendo manifestavit. Super occasum videlicet ascendit; quia mortem quam pertulit, resurgendo calcavit. Ei ergo qui ascendit super occasum, iter facimus, cum nos ejus gloriam vestris mentibus prædicamus, ut eas et ipse post veniens, per amoris sui præsentiam illustret.\par
\switchcolumn
\EN
\noindent Therefore to preachers saith Isaiah: “Prepare ye the way of the Lord, make straight an highway for our God.” (xl. 3.) And again the Psalmist saith: “Spread a path before Him That rideth upon the West.” (lxvii. 4.) The Lord rideth upon the West above that from which in death He veiled His glory hath He royally exalted that glory that excelleth, even the glory of His rising again. He rideth upon the West, Who, being risen again from the dead, is throned high above the death to which He bowed. Before Him, therefore, That rideth upon the West, we spread a path, when we set forth His glory before the eyes of your mind, to the end that He Himself may come after, and Himself enlighten the same your minds by His presence and His love.\par
\switchcolumn*

\LA
\TeDeum{Te Deum}
\switchcolumn
\EN
\TeDeum{Te Deum}
\switchcolumn*

\LA
\LessonHead{Lectio IX (pro commemoratione)}
\switchcolumn
\EN
\LessonHead{Lesson IX (when commemorated)}
\switchcolumn*

\LA
\LessonTitle{Commemoratio: S. Augustini Episcopi et Confessoris}
\switchcolumn
\EN
\LessonTitle{Commemoration of: St. Augustine of Canterbury, Bishop and Confessor}
\switchcolumn*

\LA
\noindent Augustínus, Romæ in Lateranénsi cœnóbio mónachus, a Gregório Magno cum sóciis mónachis fere quadragínta in Angliam missus est, anno quingentésimo nonagésimo séptimo. A rege Ethelbérto Cantuáriam, ejus regni metrópolim, invitátus cum sóciis, prope eam oratórium exstrúxit. Cæléstis doctrínæ prædicatióne plerósque insulános ac regem ipsum ad christiánam fidem perdúxit, summa cum lætítia Berthæ, régiæ uxóris, quæ Christiána erat. Jussu Gregórii ordinátus epíscopus, Sedem Cantuariénsem instítuit, et ab eódem Pontífice usum pállii cum facultáte hierarchíæ in Anglia instituéndæ obtínuit. Máximis demum pro Christo exantlátis labóribus, cum Mellítum Londinénsi ecclésiæ præfecísset, Justum Roffénsi, suæ Lauréntium, in cælum migrávit séptimo kaléndas júnias, et sepúltus est in monastério sancti Petri, quod exínde Cantuariénsium antístitum et áliquot regum conditórium fuit.\par
\switchcolumn
\EN
\noindent Augustine, a monk of the Lateran monastery in Rome, was sent by Gregory the Great in 597 to England with about forty monks as his companions. They were invited by King Ethelbert to Canterbury, the chief city of the kingdom, and they built an oratory nearby. Through preaching the doctrine of heaven, Augustine brought many of the islanders and the king himself to the Christian faith, to the great joy of the king's wife, Bertha, who was a Christian. By order of Pope Gregory, Augustine was ordained bishop and founded the see of Canterbury; by the same Pontiff he was granted the use of the pallium and the right to organize the hierarchy of England. At length, after suffering great hardships for Christ, having set Mellitus over the Church of London, Justus over that of Rochester, and Lawrence over his own Church, he made his journey to heaven on the 26th day of May. He was buried in the monastery of St. Peter, which then became the burial place of bishops of Canterbury and of several kings.\par
\switchcolumn*

\LA
\TeDeum{Te Deum}
\switchcolumn
\EN
\TeDeum{Te Deum}
\switchcolumn*

\end{paracol}

\par\needspace{10\baselineskip}\addvspace{1.2em}{\centering\small\textsc{Die 29 Maii} · \textsc{29 May}\par}
\Office{S. Mariæ Magdalenæ de Pazzis Virginis}{St. Mary Magdalene de Pazzi, Virgin}{Semiduplex}
\addcontentsline{toc}{subsection}{Die 29 Maii · S. Mariæ Magdalenæ de Pazzis Virginis\enspace\textit{St. Mary Magdalene de Pazzi, Virgin}}
\begin{paracol}{2}
\LA
\RefLine{Lectiones I–III: de Scriptura occurrente.}
\switchcolumn
\EN
\RefLine{Lessons I–III: from the Scripture of the day.}
\switchcolumn*

\LA
\RefLine{Lectiones VII–IX: ex Commúni Virginum tantum (C6a).}
\switchcolumn
\EN
\RefLine{Lessons VII–IX: from the Common of a Virgin not a Martyr (C6a).}
\switchcolumn*

\LA
\LessonHead{Lectio IV}
\switchcolumn
\EN
\LessonHead{Lesson IV}
\switchcolumn*

\LA
\noindent Maria Magdalena, illustrióri Pazziórum genere Floréntiæ nata, fere ab incunabulis iter perfectiónis arrípuit. Decennis perpetuam virginitátem vovit, susceptoque habitu in monastério sanctæ Maríæ Angelórum, ordinis Carmelitárum, se ómnium virtútum exemplar exhíbuit. Adeo casta fuit, ut quidquid puritátem lǽdere potest, pénitus ignoraverit. Quinquénnium, Deo jubénte, solo pane et aqua transegit, exceptis diébus Domínicis, quibus cibis quadragesimálibus vescebátur. Corpus suum cilício, flagellis, frígore, inedia, vigiliis, nuditate, atque omni pœnárum genere cruciábat.\par
\switchcolumn
\EN
\noindent This Mary Magdalen was born of the noble Florentine family of the Pazzi, (on the 2nd day of April, in the year of Christ 1566.) She was hardly out of her cradle when she set her feet in the path of perfection. At ten years of age she made a vow of perpetual virginity, and (at fifteen) took the habit of the Order of Mount Carmel, in the convent of Saint Mary of the Angels. In that sisterhood she was in all ways a pattern to all. She was pure to that degree, that she did not even know of the existence of anything which can hurt modesty. For the space of five years, by the command of God, she lived upon nothing but bread and water, the Lord's Day only excepted, in which she used the food which is taken in Lent. She chastised her body with hair-cloth, scourging, cold, hunger, watching, nakedness, and all manner of hardships.\par
\switchcolumn*

\LA
\begin{Resp}\Rsign Propter veritátem, et mansuetúdinem, et justítiam: \Star{} Et dedúcet te mirabíliter déxtera tua, allelúja.\end{Resp}
\begin{Resp}\Vsign Spécie tua et pulchritúdine tua inténde, próspere procéde, et regna.\end{Resp}
\begin{Resp}\Rsign Et dedúcet te mirabíliter déxtera tua, allelúja.\end{Resp}
\switchcolumn
\EN
\begin{Resp}\Rsign Because of truth, and meekness, and righteousness; \Star{} And thy right hand shall lead thee wonderfully, alleluia.\end{Resp}
\begin{Resp}\Vsign In thy comeliness, and thy beauty, go forward, fare prosperously, and reign.\end{Resp}
\begin{Resp}\Rsign And thy right hand shall lead thee wonderfully, alleluia.\end{Resp}
\switchcolumn*

\LA
\LessonHead{Lectio V}
\switchcolumn
\EN
\LessonHead{Lesson V}
\switchcolumn*

\LA
\noindent Tanto igne divini amoris æstuábat, ut, ei feréndo impar, ingesta aqua pectus refrigerare cogerétur. Extra sensus frequenter rapta, diuturnas et admirábiles éxtases passa est, in quibus et arcana cæléstia penetrávit, et exímiis a Deo grátiis illustráta fuit. His autem muníta longum certamen a princípibus tenebrárum sustinuit, árida, desoláta, ab ómnibus derelícta, variisque tentatiónibus vexáta; Deo sic permittente, ut invíctæ patiéntiæ ac profundíssimæ humilitátis exémplar præbéret.\par
\switchcolumn
\EN
\noindent The love of God was so hot within her, that she was sometimes fain to bathe her breast with cold water to allay the agitation. She was oftentimes rapt in the spirit, and that most marvellously, for whole days at a time, during which trances she saw things hidden and heavenly, and was enlightened of God with great gifts. But after all these things she had a stern tussling with the prince of the darkness of this world, while God allowed her spirit to remain dry, deserted, abandoned by all, and tormented with diverse temptations. And all that while she remained an example of unconquered patience and the deepest lowly-mindedness.\par
\switchcolumn*

\LA
\begin{Resp}\Rsign Dilexísti justítiam, et odísti iniquitátem: \Star{} Proptérea unxit te Deus, Deus tuus, óleo lætítiæ, allelúja.\end{Resp}
\begin{Resp}\Vsign Propter veritátem, et mansuetúdinem, et justítiam.\end{Resp}
\begin{Resp}\Rsign Proptérea unxit te Deus, Deus tuus, óleo lætítiæ, allelúja.\end{Resp}
\switchcolumn
\EN
\begin{Resp}\Rsign Thou hast loved righteousness, and hated iniquity; \Star{} Therefore God, thy God, hath anointed thee with the oil of gladness, alleluia.\end{Resp}
\begin{Resp}\Vsign Because of truth, and meekness, and righteousness.\end{Resp}
\begin{Resp}\Rsign Therefore God, thy God, hath anointed thee with the oil of gladness, alleluia.\end{Resp}
\switchcolumn*

\LA
\LessonHead{Lectio VI}
\switchcolumn
\EN
\LessonHead{Lesson VI}
\switchcolumn*

\LA
\noindent Caritáte erga próximum singuláriter enítuit; nam sæpe noctes ducebat insomnes, vel obeúndis sorórum ministériis, vel inserviéndo infirmis occupata, quarum aliquándo úlcera lambens sanávit. Infidélium et peccatórum perditiónem amare deflens, se ad quælibet pro illórum salúte torménta parátam offerébat. Multis ante obitum annis, univérsis cæli deliciis, quibus copiose affluébat, heróica virtúte renuntians, illud frequénter in ore habébat: Pati, non mori. Tandem longa et gravíssima infirmitáte exháusta, transívit ad Sponsum die vigésima quinta Maji, anno millésimo sexcentésimo septimo, expleto anno quadragésimo primo ætátis suæ. Eam, multis in vita et post mortem miráculis claram, Clemens nonus sanctárum Vírginum número adscrípsit: cujus corpus in præséntem diem incorrúptum conservátur.\par
\switchcolumn
\EN
\noindent She was very remarkable for her tender love toward her neighbours. Sometimes she went whole nights without sleep, while she was working for the service of the sisters, Or waiting upon the sick. She sometimes healed sores even by licking them. That there should be unbelievers and sinners perishing caused her bitter weeping, and she offered herself to God to suffer for their conversion whatsoever He chose. For many years, therefore, before her death, her mighty charity towards Others, made her freely to give up that heavenly joy of spirit, wherewith she had once so overflowed. She had often in her mouth the words: To suffer, not to die. At length, in the forty-second year of her age, on the 25th day of May, in the year 1607, after a long and grievous sickness, the Bridegroom came, and she entered with Him into the marriage-chamber. Clement IX, finding that God had glorified her by many miracles, both during her life and after her death, enrolled her name among those of the Holy Virgins. Her body, up to the present day, has never shown the least sign of corruption.\par
\switchcolumn*

\LA
\begin{Resp}\Rsign Afferéntur Regi vírgines post eam, próximæ ejus; \Star{} Afferéntur tibi in lætítia et exsultatióne, allelúja.\end{Resp}
\begin{Resp}\Vsign Spécie tua et pulchritúdine tua; inténde, próspere procéde, et regna.\end{Resp}
\begin{Resp}\Rsign Afferéntur tibi in lætítia et exsultatióne, allelúja.\end{Resp}
\begin{Resp}\Vsign Glória Patri, et Fílio, \Star{} et Spirítui Sancto.\end{Resp}
\begin{Resp}\Rsign Afferéntur tibi in lætítia et exsultatióne, allelúja.\end{Resp}
\switchcolumn
\EN
\begin{Resp}\Rsign After her shall virgins be brought unto the King, her fellows \Star{} They shall be brought unto thee with gladness and rejoicing, alleluia.\end{Resp}
\begin{Resp}\Vsign In thy comeliness and thy beauty, go forward, fare prosperously, and reign.\end{Resp}
\begin{Resp}\Rsign They shall be brought unto thee with gladness and rejoicing, alleluia.\end{Resp}
\begin{Resp}\Vsign Glory be to the Father, and to the Son, \Star{} and to the Holy Ghost.\end{Resp}
\begin{Resp}\Rsign They shall be brought unto thee with gladness and rejoicing, alleluia.\end{Resp}
\switchcolumn*

\LA
\LessonHead{Lectio IX (pro commemoratione)}
\switchcolumn
\EN
\LessonHead{Lesson IX (when commemorated)}
\switchcolumn*

\LA
\LessonTitle{Commemoratio: S. Mariæ Magdalenæ de Pazzis Virginis}
\switchcolumn
\EN
\LessonTitle{Commemoration of: St. Mary Magdalene de Pazzi, Virgin}
\switchcolumn*

\LA
\noindent María Magdaléna Floréntiæ illústri Pazziórum génere nata, fere ab incunábulis iter perfectiónis arrípuit. Decénnis perpétuam virginitátem vovit, susceptóque hábitu in monastério sorórum Carmelitárum, se ómnium virtútum exémplar prǽbuit. Adeo casta fuit, ut, quidquid puritátem lǽdere posset, pénitus ignoráverit. Tanto igne divíni amóris æstuábat, ut, ei feréndo impar, ingésta aqua pectus refrigeráre cogerétur. Caritáte erga próximum excélluit; nam sæpe noctes ducébat insómnes, vel obeúndis sorórum ministériis, vel inserviéndo infírmis occupáta, quarum aliquándo úlcera lambens sanávit. Illud frequénter in ore habébat: Pati, non mori. Tandem, diútino et gravi morbo exháusta, transívit ad Sponsum, anno millésimo sexcentésimo séptimo, expléto ætátis suæ quadragésimo primo. Eam Clemens nonus sanctárum Vírginum número adscrípsit.\par
\switchcolumn
\EN
\noindent Mary Magdalene was born of the famous family of the Pazzi in Florence, and almost from her cradle entered on the way of perfection. At ten years old, she made a vow of perpetual virginity, and when she had taken the habit in the monastery of the Carmelite sisters, she showed herself to be exemplary in all virtues. She was so chaste that she did not even know of anything that could harm purity. She burned with such a fire of divine love that she was unable to bear it, and had to cool her breast by pouring water over it. She was notable for her love of neighbor, often passing sleepless nights either in doing the work of the sisters, or in serving the sick, whose ulcers she sometimes cured by licking them. She frequently repeated this saying: “To suffer, not to die.” At length, worn out by a long and serious illness, she went to the Bridegroom in the year 1607, having completed her forty-first year. Clement IX numbered her among the Holy Virgins.\par
\switchcolumn*

\LA
\TeDeum{Te Deum}
\switchcolumn
\EN
\TeDeum{Te Deum}
\switchcolumn*

\end{paracol}

\par\needspace{10\baselineskip}\addvspace{1.2em}{\centering\small\textsc{Die 30 Maii} · \textsc{30 May}\par}
\Office{S. Felicis I Papæ et Martyris}{St. Felix, Pope and Martyr}{Simplex}
\addcontentsline{toc}{subsection}{Die 30 Maii · S. Felicis I Papæ et Martyris\enspace\textit{St. Felix, Pope and Martyr}}
\begin{paracol}{2}
\LA
\RefLine{Lectiones I–II: de Scriptura occurrente.}
\switchcolumn
\EN
\RefLine{Lessons I–II: from the Scripture of the day.}
\switchcolumn*

\LA
\LessonHead{Lectio III (pro commemoratione)}
\switchcolumn
\EN
\LessonHead{Lesson III (when commemorated)}
\switchcolumn*

\LA
\noindent Felix Románus, patre Constántio, Aureliáno imperatóre præfuit Ecclesiæ. Constítuit ut Missa supra memórias et sepúlcra Mártyrum celebrarétur. Qui cum mense Decémbri habuísset ordinatiónes duas, et creásset presbýteros novem, diáconos quinque, epíscopos per divérsa loca quinque, martyrio coronátus, via Aurelia sepelítur in basílica quam a se ædificátam dedicárat. Vixit in pontificátu annos duos, menses quátuor, dies vigínti novem.\par
\switchcolumn
\EN
\noindent Pope Felix I was a Roman who ruled the Church in the days of the Emperor Aurelian. His father's name was Constantius. His is the ordinance which commands that Mass should be celebrated on the monuments and graves of martyrs. He held two December ordinations, wherein he ordained nine Priests, five Deacons, and five Bishops for diverse places. Having finished his testimony he was buried upon the Aurelian Way, in the Church which he had himself built and dedicated. He lived as Pope two years, four months, and twenty-nine days.\par
\switchcolumn*

\LA
\TeDeum{Te Deum}
\switchcolumn
\EN
\TeDeum{Te Deum}
\switchcolumn*

\end{paracol}

\par\needspace{10\baselineskip}\addvspace{1.2em}{\centering\small\textsc{Die 31 Maii} · \textsc{31 May}\par}
\Office{Beatæ Mariæ Virginis Reginæ}{Beatæ Mariæ Virginis Reginæ}{Duplex II classis}
\addcontentsline{toc}{subsection}{Die 31 Maii · Beatæ Mariæ Virginis Reginæ\enspace\textit{Beatæ Mariæ Virginis Reginæ}}
\begin{paracol}{2}
\LA
\RefLine{Lectio IX: de Scriptura occurrente.}
\switchcolumn
\EN
\RefLine{Lesson IX: from the Scripture of the day.}
\switchcolumn*

\LA
\LessonHead{Lectio I}
\switchcolumn
\EN
\LessonHead{Lesson I}
\switchcolumn*

\LA
\LessonTitle{De libro Ecclesiástici}
\switchcolumn
\EN
\LessonTitle{Lesson from the book of Ecclesiasticus}
\switchcolumn*

\LA
\Cite{Sir 24:5-11}
\switchcolumn
\EN
\Cite{Sir 24:5-11}
\switchcolumn*

\LA
\noindent Ego ex ore Altíssimi prodívi, primogénita ante omnem creatúram: ego feci in cælis ut orirétur lumen indefíciens, et sicut nébula texi omnem terram: ego in altíssimis habitávi et thronus meus in colúmna nubis. Gyrum cæli circuívi sola, et profúndum abýssi penetrávi, in flúctibus maris ambulávi, et in omni terra steti: et in omni pópulo, et in omni gente primátum hábui: et ómnium excelléntium et humílium corda virtúte calcávi: et in his ómnibus réquiem quæsívi, et in hereditáte Dómini morábor.\par
\switchcolumn
\EN
\noindent I came out of the mouth of the most High, the firstborn before all creatures: I made that in the heavens there should rise light that never faileth, and as a cloud I covered all the earth: I dwelt in the highest places, and my throne is in a pillar of a cloud. I alone have compassed the circuit of heaven, and have penetrated into the bottom of the deep, and have walked in the waves of the sea, And have stood in all the earth: and in every people, And in every nation I have had the chief rule: And by my power I have trodden under my feet the hearts of all the high and low: and in all these I sought rest, and I shall abide in the inheritance of the Lord.\par
\switchcolumn*

\LA
\begin{Resp}\Rsign Beáta es, María, quæ credidísti Dómino: perfécta sunt in te quæ dicta sunt tibi. \Star{} Ecce exaltáta es super choros Angelórum ad cæléstia regna.\end{Resp}
\begin{Resp}\Vsign Ave Maria, grátia plena, Dominus tecum.\end{Resp}
\begin{Resp}\Rsign Ecce exaltáta es super choros Angelorum ad cæléstia regna.\end{Resp}
\switchcolumn
\EN
\begin{Resp}\Rsign Blessed art thou, Mary, because thou hast believed the Lord: \Star{} The things promised thee have been accomplished in thee.\end{Resp}
\begin{Resp}\Vsign Hail, Mary, full of grace, the Lord is with thee.\end{Resp}
\begin{Resp}\Rsign The things promised thee have been accomplished in thee.\end{Resp}
\switchcolumn*

\LA
\LessonHead{Lectio II}
\switchcolumn
\EN
\LessonHead{Lesson II}
\switchcolumn*

\LA
\Cite{Sir 24:14-16}
\switchcolumn
\EN
\Cite{Sir 24:14-16}
\switchcolumn*

\LA
\noindent Ab inítio, et ante sǽcula creáta sum, et usque ad futúrum sǽculum non désinam, et in habitatióne sancta coram ipso ministrávi. Et sic in Sion firmáta sum, et in civitáte sanctificáta simíliter requiévi, et in Jerúsalem potéstas mea. Et radicávi in pópulo honorificáto, et in parte Dei mei heréditas illíus, et in plenitúdine sanctórum deténtio mea.\par
\switchcolumn
\EN
\noindent From the beginning, and before the world, was I created, and unto the world to come I shall not cease to be, and in the holy dwelling place I have ministered before him. And so was I established in Sion, and in the holy city likewise I rested, and my power was in Jerusalem. And I took root in an honourable people, and in the portion of my God his inheritance, and my abode is in the full assembly of saints.\par
\switchcolumn*

\LA
\begin{Resp}\Rsign Regálem dignitátem Vírginis Maríæ recolámus. \Star{} Quia cum Christo regnat in ætérnum. (Allelúja.)\end{Resp}
\begin{Resp}\Vsign Glóriam Regínæ nostræ celebrémus.\end{Resp}
\begin{Resp}\Rsign Quia cum Christo regnat in ætérnum. (Allelúja.)\end{Resp}
\switchcolumn
\EN
\begin{Resp}\Rsign Let us reflect upon the royal dignity of the Virgin Mary, \Star{} For with Christ she reigneth forever.\end{Resp}
\begin{Resp}\Vsign Let us proclaim the glory of our Queen.\end{Resp}
\begin{Resp}\Rsign For with Christ she reigneth forever.\end{Resp}
\switchcolumn*

\LA
\LessonHead{Lectio III}
\switchcolumn
\EN
\LessonHead{Lesson III}
\switchcolumn*

\LA
\Cite{Sir 24:24-30}
\switchcolumn
\EN
\Cite{Sir 24:24-30}
\switchcolumn*

\LA
\noindent Ego mater pulchræ dilectiónis, et timóris, et agnitiónis, et sanctæ spei. In me grátia omnis viæ et veritátis: in me omnis spes vitæ et virtútis. Transíte ad me, omnes qui concupíscitis me, et a generatiónibus meis implémini; spíritus enim meus super mel dulcis, et heréditas mea super mel et favum. Memória mea in generatiónes sæculórum. Qui edunt me, adhuc esúrient, et qui bibunt me, adhuc sítient. qui audit me non confundétur, et qui operántur in me non peccábunt: qui elúcidant me, vitam ætérnam habébunt.\par
\switchcolumn
\EN
\noindent I am the mother of fair love, and of fear, and of knowledge, and of holy hope. In me is all grace of the way and of the truth, in me is all hope of life and of virtue. Come over to me, all ye that desire me, and be filled with my fruits. For my spirit is sweet above honey, and my inheritance above honey and the honeycomb. My memory is unto everlasting generations. They that eat me, shall yet hunger: and they that drink me, shall yet thirst. He that hearkeneth to me, shall not be confounded: and they that work by me, shall not sin.\par
\switchcolumn*

\LA
\begin{Resp}\Rsign Elégit eam Deus, et præelégit eam; \Star{} Corónam glóriæ pósuit super caput ejus. (Allelúja.)\end{Resp}
\begin{Resp}\Vsign Et in tabernáculo suo habitáre fecit eam.\end{Resp}
\begin{Resp}\Rsign Corónam glóriæ pósuit super caput ejus. (Allelúja.)\end{Resp}
\begin{Resp}\Vsign Glória Patri, et Fílio, \Star{} et Spirítui Sancto.\end{Resp}
\begin{Resp}\Rsign Corónam glóriæ pósuit super caput ejus. (Allelúja.)\end{Resp}
\switchcolumn
\EN
\begin{Resp}\Rsign God hath chosen her, and preferred her; \Star{} He hath placed a crown of glory on her head.\end{Resp}
\begin{Resp}\Vsign And he maketh her dwell in his tabernacle.\end{Resp}
\begin{Resp}\Rsign He hath placed a crown of glory on her head.\end{Resp}
\begin{Resp}\Vsign Glory be to the Father, and to the Son, \Star{} and to the Holy Ghost.\end{Resp}
\begin{Resp}\Rsign He hath placed a crown of glory on her head.\end{Resp}
\switchcolumn*

\LA
\LessonHead{Lectio IV}
\switchcolumn
\EN
\LessonHead{Lesson IV}
\switchcolumn*

\LA
\LessonTitle{Sermo sancti Petri Canísii Presbýteri}
\switchcolumn
\EN
\LessonTitle{Sermon of St. Peter Canisius, Priest}
\switchcolumn*

\LA
\Source{De Maria Deipara Virgine incomparabili, lib. 15, c. 13}
\switchcolumn
\EN
\Source{On the Incomparable Virgin Mary, Mother of God}
\switchcolumn*

\LA
\noindent Cur beatíssimam Vírginem Maríam Regínæ nómine, Damascénum, Athanásium aliósque secúti, non compellémus, cujus et pater David rex ínclitus, et fílius Rex regum Dominúsque dominántium sine fine ímperans, laudem in Scriptúris præstantíssimam tenent? Regína est ínsuper, si cum illis conferátur, quibus, véluti régibus, cæléste regnum cum Christo, Rege summo, cóntigit, útpote illíus coherédibus, et in eódem véluti throno, ut Scriptúra lóquitur, cum illo collocátis. Regína est étiam nulli electórum secúnda, sed simul Angelis et homínibus tanto præláta dígnius, quo nihil illa sublímius ac sánctius esse potest, quæ sola cum Deo Patre Fílium habet commúnem, et quæ supra se Deum et Christum tantum, infra se vero réliqua videt ómnia.\par
\switchcolumn
\EN
\noindent If we follow St. John Damascene, St. Athanasius and others, are we not forced to call Mary Queen, since her father David receiveth the highest praise in Scripture as a renowned king, and her son as King of kings and Lord of lords, reigning forever? She is Queen, moreover, when compared with the Saints who reign like kings in the heavenly kingdom, co-heirs with Christ, the great King, placed on the same throne with him, as the Scripture saith. And as Queen she is second to none of the elect, but in dignity is raised so high above both Angels and men that nothing can be higher or holier than she, who alone hath the same Son as God the Father, and who seeth above her only God and Christ, and below her creatures other than herself.\par
\switchcolumn*

\LA
\begin{Resp}\Rsign Súscipe verbum, Virgo María, quod tibi a Dómino transmíssum est: \Star{} Ecce concípies et páries Deum páriter et hóminem. (Allelúja.)\end{Resp}
\begin{Resp}\Vsign Et Regína vocáberis super omnes gentes.\end{Resp}
\begin{Resp}\Rsign Ecce concípies et páries Deum páriter et hóminem. (Allelúja.)\end{Resp}
\switchcolumn
\EN
\begin{Resp}\Rsign Receive the word, O Virgin Mary, sent to thee by the Lord: \Star{} Behold, thou wilt conceive and bring forth God as well as man.\end{Resp}
\begin{Resp}\Vsign And thou wilt be called Queen over all the nations.\end{Resp}
\begin{Resp}\Rsign Behold, thou wilt conceive and bring forth God as well as man.\end{Resp}
\switchcolumn*

\LA
\LessonHead{Lectio V}
\switchcolumn
\EN
\LessonHead{Lesson V}
\switchcolumn*

\LA
\noindent Magnus Athanásius perspícue dixit: Maríam non modo Deíparam, sed étiam Regínam et Dóminam próprie veréque censéri, quandóquidem Christus ex ipsa matre Vírgine natus, Deus et Dominus idémque Rex máneat. De hac ígitur Regína dictum illud Psalmógraphi interpretátur: Adstitit Regína a dextris tuis in vestítu deauráto. Non ergo solum cæli, sed et cælórum Regína recte dicitur María, útpote mater Regis Angelórum, Regísque cælórum et amíca et sponsa. A te vero, augustíssima Regína, eadémque fidíssima Mater María, quam sine fructu pie nullus implórat, et cui mortáles omnes beneficiórum memória sempitérna devíncti sunt, étiam atque étiam reverénter oro et óbsecro, ut hoc qualecúmque observántiæ erga te meæ testimónium ratum et gratum habére velis, utque obláti múneris exiguitátem studiósa offeréntis voluntáte metíri ac tuo præpoténti Fílio comprobáre dignéris.\par
\switchcolumn
\EN
\noindent The great Athanasius said clearly: Mary is not only the Mother of God, but also can be properly and truly called Queen and Lady, since in fact the Christ who was born of the Virgin Mother is God and Lord and also King. It is to this Queen, therefore, that the Psalmist's words are applied: The Queen taketh her place at thy right hand in garments of gold. Thus Mary is rightly called Queen, not only of heaven, but also of the heavens, as the Mother of the King of Angels, and as the Bride and beloved of the King of the heavens. O Mary, most august Queen and most faithful Mother, to whom no one prayeth in vain who prayeth devoutly, and to whom all mortal men are bound by the enduring memory of so many benefits, again and again reverently I beseech thee to accept and be pleased with every evidence of my devotion to thee, to value the poor gift I offer according to the zeal with which it is offered, and to recommend it to thine all-powerful Son.\par
\switchcolumn*

\LA
\begin{Resp}\Rsign Ecce pósitus est hic in ruínam et resurrectiónem multórum, \Star{} Et tuam ipsíus ánimam pertransíbit gládius. (Allelúja.)\end{Resp}
\begin{Resp}\Vsign Ave, Christi Mater, passiónis sócia, totíus mundi Regína.\end{Resp}
\begin{Resp}\Rsign Et tuam ipsíus ánimam pertransíbit gládius. (Allelúja.)\end{Resp}
\switchcolumn
\EN
\begin{Resp}\Rsign Behold, this Child is destined for the fall and the rise of many, \Star{} And thine own soul a sword shall pierce.\end{Resp}
\begin{Resp}\Vsign Hail, Mother of Christ, sharing in his passion, Queen of all the world.\end{Resp}
\begin{Resp}\Rsign And thine own soul a sword shall pierce.\end{Resp}
\switchcolumn*

\LA
\LessonHead{Lectio VI}
\switchcolumn
\EN
\LessonHead{Lesson VI}
\switchcolumn*

\LA
\LessonTitle{Ex Lítteris Encýclicis Pii Papæ duodécimi}
\switchcolumn
\EN
\LessonTitle{From the Encyclical Letter of Pope Pius XII}
\switchcolumn*

\LA
\Source{Encyclica « Ad cæli Regínam », diei 11 Octobris 1954}
\switchcolumn
\EN
\Source{Ad cæli Reginam, diei 11 Octobris 1954}
\switchcolumn*

\LA
\noindent E Christiánæ vetustátis monuméntis, e litúrgicis précibus, ex índito christiáno pópulo religiónis sensu, ex opéribus arte conféctis, úndique collégimus voces quæ ásserunt Deíparam Vírginem regáli dignitáte præstáre; ratiónes étiam, quas sancta theología ex divínæ fídei thesáuro deducéndo ástruit, eándem veritátem prorsus confirmáre argúimus. Tot ex allátis testimóniis quasi latíssime résonans concéntus effícitur, qui extóllit régii honóris præcélsum fastígium Dei hominúmque Matris, cui cuncta creáta subsunt, quæ est: « exaltáta super choros Angelórum ad cæléstia regna ». Cum vero, matúro ponderatóque consílio, persuásum Nobis habeámus magna oritúra esse Ecclésiæ emoluménta, si quasi in suo candelábro rutilántior lucérna pósita, illa sólide probáta véritas maniféstior ómnibus refúlgeat, Apostólica Nostra Potestáte decérnimus et institúimus festum Maríæ Regínæ, quod toto terrárum orbe quotánnis die trigésimo primo mensis maji est celebrándum.\par
\switchcolumn
\EN
\noindent From the documents of ancient Christianity, from the prayers of the liturgy, from the innate religious sense of the Christian people, from works of art, from all sides we gather witnesses which assert that the Virgin Mother of God excelleth in queenly dignity. And we have set forth the reasons which sacred theology deduces from the treasury of divine faith to confirm the same truth. All these witnesses form a sort of chorus, proclaiming far and wide the supreme queenly honour granted to the Mother of God and man, who is above all created things and exalted over the choirs of Angels to reign in heaven. Thus it is that after mature and thoughtful consideration we have been persuaded that great benefits would flow to the Church if, like a light that illumines more brightly when placed in its stand, this solidly proved truth were to shine out more clearly to all, and so, by Our Apostolic Authority, we decree and institute the feast of Mary, Queen, which is to be celebrated every year on the thirty-first day of May throughout the world.\par
\switchcolumn*

\LA
\begin{Resp}\Rsign Signum magnum appáruit in cælo: \Star{} Múlier amícta sole, et luna sub pédibus ejus, et in cápite ejus coróna stellárum duódecim. (Allelúja.)\end{Resp}
\begin{Resp}\Vsign Cujus Fílius regnat in ætérnum.\end{Resp}
\begin{Resp}\Rsign Múlier amícta sole, et luna sub pédibus ejus, et in cápite ejus coróna stellárum duódecim. (Allelúja.)\end{Resp}
\begin{Resp}\Vsign Glória Patri, et Fílio, \Star{} et Spirítui Sancto.\end{Resp}
\begin{Resp}\Rsign Múlier amícta sole, et luna sub pédibus ejus, et in cápite ejus coróna stellárum duódecim. (Allelúja.)\end{Resp}
\switchcolumn
\EN
\begin{Resp}\Rsign A great sign appeared in heaven: \Star{} A woman clothed with the sun, and the moon was under her feet, and upon her head a crown of twelve stars.\end{Resp}
\begin{Resp}\Vsign Whose Son reigneth forever.\end{Resp}
\begin{Resp}\Rsign A woman clothed with the sun, and the moon was under her feet, and upon her head a crown of twelve stars.\end{Resp}
\begin{Resp}\Vsign Glory be to the Father, and to the Son, \Star{} and to the Holy Ghost.\end{Resp}
\begin{Resp}\Rsign A woman clothed with the sun, and the moon was under her feet, and upon her head a crown of twelve stars.\end{Resp}
\switchcolumn*

\LA
\LessonHead{Lectio VII}
\switchcolumn
\EN
\LessonHead{Lesson VII}
\switchcolumn*

\LA
\LessonTitle{Léctio sancti Evangélii secúndum Lucam}
\switchcolumn
\EN
\LessonTitle{From the Holy Gospel according to Luke}
\switchcolumn*

\LA
\Cite{Luc 1:26-33}
\switchcolumn
\EN
\Cite{Luke 1:26-33}
\switchcolumn*

\LA
\noindent In illo témpore: Missus est Angelus Gábriel a Deo in civitátem Galilǽæ, cui nomen Názareth, ad Vírginem desponsátam viro, cui nomen erat Joseph de domo David, et nomen Vírginis María. Et réliqua.\par
\switchcolumn
\EN
\noindent In that time: In the sixth month, the angel Gabriel was sent from God into a city of Galilee, called Nazareth, to a virgin espoused to a man whose name was Joseph, of the house of David; and the virgin's name was Mary. And so on.\par
\switchcolumn*

\LA
\LessonTitle{Homilía sancti Bonaventúræ Epíscopi}
\switchcolumn
\EN
\LessonTitle{Homily of St. Bonaventure, Bishop}
\switchcolumn*

\LA
\Source{Sermo de regia dignitate Beatæ Mariæ Virginis}
\switchcolumn
\EN
\Source{Sermon on the Royal Dignity of the Blessed Virgin Mary}
\switchcolumn*

\LA
\noindent Beáta María Virgo, summi Regis mater est per generósam conceptiónem, secúndum quod dícitur in verbo sibi dicto ab Angelo: Ecce inquit, concípies et páries fílium; et póstea: Dabit ei Dóminus sedem David patris ejus, et regnábit in domo Jacob in ætérnum, et regni ejus non erit finis. Ac si apérte dicat: Concípies et páries fílium Regem, in regáli sólio æternáliter residéntem, ac per hoc tamquam Mater Regis regnábis, et ut Regína in regáli sólio residébis. Si enim decet filium honórem matri dare, decet ut ei commúnicet thronum regálem; unde Virgo María, quia concépit eum, qui habet in fémore scriptum: Rex regum et Dóminus dominántium, statim ex quo concépit Fílium Dei, Regína fuit, non solum terræ, sed étiam cæli, quod designátum est in Apocalýpsi ubi dícitur: Signum magnum appáruit in cælo: múlier amícta sole, et luna sub pédibus ejus, et in cápite ejus coróna stellárum duódecim.\par
\switchcolumn
\EN
\noindent The Blessed Virgin Mary is the Mother of the great King by reason of a noble kind of conception, according to the message given her by the Angel. Behold, he said, thou shalt conceive and shalt bring forth a son; and again, the Lord God will give him the throne of David his father, and he shall be king over the house of Jacob forever, and of his kingdom there shall be no end. This is as if to say in so many words, Thou wilt conceive and bear a son who is King, eternally reigning on the royal throne, and as Queen thou wilt be seated on the royal throne. For if it becometh a son to give honour to his mother, it is also fitting that he share his royal throne with her; and so the Virgin Mary, because she conceived him on whose thigh was written, King of kings and Lord of lords, was Queen not only of earth but also of heaven as soon as she conceived the Son of God. This is indicated in the Apocalypse where it saith: A great sign appeared in heaven: a woman clothed with the sun, and the moon was under her feet, and upon her head a crown of twelve stars.\par
\switchcolumn*

\LA
\begin{Resp}\Rsign Ecce Virgo concípiet et páriet Fílium, \Star{} Et vocábitur nomen ejus Admirábilis, Deus, Fortis. (Allelúja.)\end{Resp}
\begin{Resp}\Vsign Super sólium David et super regnum ejus sedébit in ætérnum.\end{Resp}
\begin{Resp}\Rsign Et vocábitur nomen ejus Admirábilis, Deus, Fortis. (Allelúja.)\end{Resp}
\switchcolumn
\EN
\begin{Resp}\Rsign Behold a Virgin shall conceive and bear a Son, \Star{} And his name shall be called Wonderful, God, the Mighty One.\end{Resp}
\begin{Resp}\Vsign He shall sit upon the throne of David, and reign over his kingdom forever.\end{Resp}
\begin{Resp}\Rsign And his name shall be called Wonderful, God, the Mighty One.\end{Resp}
\switchcolumn*

\LA
\LessonHead{Lectio VIII}
\switchcolumn
\EN
\LessonHead{Lesson VIII}
\switchcolumn*

\LA
\noindent María Regína est præclaríssima quantum ad glóriam, quod bene desígnat Prophéta in psalmo qui speciáliter est de Christo et Vírgine María, ubi primo dícitur de Christo: Sedes tua, Deus, in sǽculum sǽculi, et paulo post de Vírgine: Adstitit regína a dextris tuis, hoc est in potióribus bonis, quod quidem dictum est quantum ad glóriam mentis. Séquitur: In vestítu deauráto, in quo exprímitur vestítus gloriósæ immortalitátis, quem décuit habére Vírginem in sua Assumptióne. Absit enim ut vestiméntum illud quo opértus est Christus, quod étiam fuit perfécte sanctificátum in terra per Verbum incarnátum, fiat cibus vérmium. Sicut décuit Christum dare grátiam Matri suæ pleníssimam in sua Conceptióne, sic décuit tribúere pleníssimam glóriam in ipsíus Matris Assumptióne. Et ídeo tenéndum est quod Virgo, gloriósa in ánima et córpore, sédeat juxta Fílium.\par
\switchcolumn
\EN
\noindent Mary the Queen outshineth all others in glory, as the Prophet clearly showeth in the Psalm which particularly concerns Christ and the Virgin Mary. It first saith of Christ: Thy throne, O God, standeth forever and ever, and shortly thereafter of the Virgin: The Queen taketh her place at thy right hand, that is, in the position of highest blessedness, for it referreth to glory of soul. It continueth: In garments of gold, by which is meant the clothing of glorious immortality which was proper to the Virgin in her Assumption. For it could not be that the garment which clothed Christ, the garment completely sanctified on earth by the incarnate Word, should be the food of worms. As it was fitting for Christ to grant the fullness of grace to his Mother at her Conception, so was it fitting that he grant her the fullness of glory at her Assumption. And so we are to hold that the Virgin, glorious in soul and body, is enthroned next to her Son.\par
\switchcolumn*

\LA
\begin{Resp}\Rsign Exaltáta es, sancta Dei Génetrix, \Star{} Super choros Angelórum ad cæléstia regna. (Allelúja.)\end{Resp}
\begin{Resp}\Vsign Intercéde pro nobis ad Dóminum Jesum Christum.\end{Resp}
\begin{Resp}\Rsign Super choros Angelórum ad cæléstia regna. (Allelúja.)\end{Resp}
\begin{Resp}\Vsign Glória Patri, et Fílio, \Star{} et Spirítui Sancto.\end{Resp}
\begin{Resp}\Rsign Super choros Angelórum ad cæléstia regna. (Allelúja.)\end{Resp}
\switchcolumn
\EN
\begin{Resp}\Rsign Thou hast been exalted, holy Mother of God, \Star{} Over the choirs of Angels, to reign in heaven.\end{Resp}
\begin{Resp}\Vsign Intercede for us with the Lord Jesus Christ.\end{Resp}
\begin{Resp}\Rsign Over the choirs of Angels, to reign in heaven.\end{Resp}
\begin{Resp}\Vsign Glory be to the Father, and to the Son, \Star{} and to the Holy Ghost.\end{Resp}
\begin{Resp}\Rsign Over the choirs of Angels, to reign in heaven.\end{Resp}
\switchcolumn*

\LA
\LessonHead{Lectio IX}
\switchcolumn
\EN
\LessonHead{Lesson IX}
\switchcolumn*

\LA
\noindent María Regína est et dispensátrix grátiæ, quod designátum fuit in libro Esther, ubi dícitur: Parvus fons qui crevit in flúvium, et in lucem et solem convérsus est. Virgo Maria sub figura Esther comparátur diffusióni fontis et lúminis, propter diffusiónem grátiæ quantum ad dúplicem usum, actiónis licet et contemplatiónis. Grátia enim Dei, quæ est medicatíva géneris humáni, per ipsam ad nos descéndit quasi per aquædúctum, quia ad ipsam Vírginem pértinet dispensátio grátiæ non per modum princípii, sed per modum mériti. Mérito ígitur Virgo María est excellentíssima Regína in comparatiónem ad plebem, cum ímpetrat véniam, súperat pugnam, et distríbuit grátiam, et consequénter perdúcit ad glóriam.\par
\switchcolumn
\EN
\noindent Mary the Queen is also the distributrix of grace. This is indicated in the book of Esther, where it is said: The little spring which grew into a river and was turned into a light and into the sun. The Virgin Mary, under the type of Esther, is compared to the outpouring of a spring and of light, because of the diffusion of grace for two uses, that is, for action and for contemplation. For the grace of God, which is a healing for the human race, descendeth to us through her as if through an aqueduct, since the dispensing of grace is attributed to the Virgin not as to its beginning, but because of her position through merit. By position the Virgin Mary is a most excellent Queen towards her people: she obtaineth forgiveness, overcometh strife, distributeth grace, and thereby she leadeth them to glory.\par
\switchcolumn*

\LA
\TeDeum{Te Deum}
\switchcolumn
\EN
\TeDeum{Te Deum}
\switchcolumn*

\end{paracol}

\SeasonTitle{Mensis Junius}{June}

\addcontentsline{toc}{section}{Mensis Junius\enspace\textit{June}}

\par\needspace{10\baselineskip}\addvspace{1.2em}{\centering\small\textsc{Die 1 Junii} · \textsc{1 June}\par}
\Office{S. Angelæ Mericiæ Virginis}{St. Angela Merici, Virgin}{Duplex}
\addcontentsline{toc}{subsection}{Die 1 Junii · S. Angelæ Mericiæ Virginis\enspace\textit{St. Angela Merici, Virgin}}
\begin{paracol}{2}
\LA
\RefLine{Lectiones I–III: de Scriptura occurrente.}
\switchcolumn
\EN
\RefLine{Lessons I–III: from the Scripture of the day.}
\switchcolumn*

\LA
\RefLine{Lectiones VII–IX: ex Commúni Virginum tantum (C6a).}
\switchcolumn
\EN
\RefLine{Lessons VII–IX: from the Common of a Virgin not a Martyr (C6a).}
\switchcolumn*

\LA
\LessonHead{Lectio IV}
\switchcolumn
\EN
\LessonHead{Lesson IV}
\switchcolumn*

\LA
\noindent Angela Merícia, Decentiáni, Veronénsis diœcésis óppido ad lacum Benácum in dicióne Venéta, piis orta paréntibus, a prima ætáte virginitátis lílium, quod perpétuo serváre statúerat, sédula sepsit. Ab omni mulíebri ornátu abhórrens, egrégiam vultus formam pulchrámque cæsáriem studióse fœdávit, ut cælésti dumtáxat animárum Sponso placéret. In ipso autem adolescéntiæ flore paréntibus orbáta, austerióris vitæ desidério in desértum locum aufúgere tentávit; sed, ab avúnculo prohíbita, novit præstáre domi, quod in solitúdine non lícuit. Cilício ac flagéllis frequénter usa, carnem nónnisi infírma valetúdine, vinum in Nativitátis et Resurrectiónis Domínicæ tantum celebritáte; complúres vero dies, nihil omníno degustávit. Oratióni dédita, brevíssimum humi carpébat somnum. Dæmónem vero sub lucéntis ángeli forma sibi illúdere conántem agnóvit prótinus, et conjécit in fugam. Tandem patérnis bonis abdicátis, et hábitum ac régulam tértii órdinis sancti Francísci ampléxa, evangélicam paupertátem virginitátis laudi conjúnxit.\par
\switchcolumn
\EN
\noindent Angela Merici was born of godly parents at Decenzano on the western shore of the Lake of Garda, in the diocese of Verona and territory of Venice, (on the 21st day of March, about the year of grace 1474.) From her earliest years she carefully guarded the lily of her virginity, with the intention of keeping it for ever unbroken. She had no taste for women's finery, and purposely marred the exceeding comeliness of her body and her sightly hair, as seeking to appear beautiful only in the eyes of Him Who is the Lover of souls. At ten years of age she lost both her father and mother, and thereafter, being fain to take upon her a life of greater hardness, she essayed to retire into a desert place apart, but this her uncle forbade her to do, and she learnt how to practice at home what she was not allowed to attempt in the wilderness. She often used hair-cloth and scourging never ate flesh-meat, except when she was sick drank wine only on the Feast-days of Christmas and Easter; and many a day took nothing at all. She was instant in prayer. What little sleep she took, she took lying on the ground. The devil strove to beguile her, appearing under the form of an angel of light, but she quickly detected him and put him to flight. At length she added to the glory of virginity that poverty which is commended in the Gospel she gave up all that she had, and adopted the dress and rule of the Third Order of St Francis.\par
\switchcolumn*

\LA
\begin{Resp}\Rsign Propter veritátem, et mansuetúdinem, et justítiam: \Star{} Et dedúcet te mirabíliter déxtera tua, allelúja.\end{Resp}
\begin{Resp}\Vsign Spécie tua et pulchritúdine tua inténde, próspere procéde, et regna.\end{Resp}
\begin{Resp}\Rsign Et dedúcet te mirabíliter déxtera tua, allelúja.\end{Resp}
\switchcolumn
\EN
\begin{Resp}\Rsign Because of truth, and meekness, and righteousness; \Star{} And thy right hand shall lead thee wonderfully, alleluia.\end{Resp}
\begin{Resp}\Vsign In thy comeliness, and thy beauty, go forward, fare prosperously, and reign.\end{Resp}
\begin{Resp}\Rsign And thy right hand shall lead thee wonderfully, alleluia.\end{Resp}
\switchcolumn*

\LA
\LessonHead{Lectio V}
\switchcolumn
\EN
\LessonHead{Lesson V}
\switchcolumn*

\LA
\noindent Nullum pietátis offícium erga próximos omíttens, paupéribus, quidquid sibi ex mendicáto victu superésset, largiebátur; libénter ministrábat ægrótis, plúraque cum magna sanctitátis fama peragrávit loca, ut vel solátio esset afflíctis, vel reis véniam impetráret, vel infénsos ínvicem reconciliáret ánimos, vel e vitiórum cœno sceléstos revocáret. Angelórum pane, quem únice esuriébat, frequentíssime refécta, tanta caritátis vi ferebátur in Deum, ut sæpius extra sensus raperétur. Sacra Palæstínæ loca, summa cum religióne, obívit; quo in itínere et visum, quem ad Cydónias appúlsa oras amíserat, eódem regréssa recuperávit, et barbarórum captivitátem ac naufrágium ímminens divínitus evásit. Romam dénique, firmam Ecclésiæ petram veneratúra et amplíssimæ jubilæi véniæ percúpida, sedénte Cleménte séptimo, accéssit; quam summus Póntifex allocútus, ejúsdem sanctimóniam suspéxit et commendávit summópere; nec ab Urbe ipsam abíre ante permísit, quam álio cælitus vocátam agnóvit.\par
\switchcolumn
\EN
\noindent She left undone no service of kindness which she was able to do to her neighbors. If there remained anything over of the food which was given in alms to herself, she gave that to the poor. She cheerfully waited upon the sick. She journeyed about, with a great reputation for holiness, comforting the afflicted, asking forgiveness for the guilty, reconciling the angry, and recalling the wicked from evil. Her only hunger was for the bread of Angels, and she took the Same right often, and then arose in her vehemence of love bearing her towards God, which oftentimes made her beside herself. She made a pilgrimage, with intense feeling, to the Holy Places in Palestine, during which journey she lost her sight at Canea in Crete on her way out, and recovered it at the same place on her way home. In this journey also, God saved her from being made prisoner by the unbelievers and from shipwreck. She went to Rome, (in 1525,) at once to pray at the immovable Rock of the Church, and to gain the abundant pardons of the Jubilee. Pope Clement VII conversed with her, was edified by her holiness, and highly commended her neither would he let her leave Rome, until he knew that God was calling her elsewhere.\par
\switchcolumn*

\LA
\begin{Resp}\Rsign Dilexísti justítiam, et odísti iniquitátem: \Star{} Proptérea unxit te Deus, Deus tuus, óleo lætítiæ, allelúja.\end{Resp}
\begin{Resp}\Vsign Propter veritátem, et mansuetúdinem, et justítiam.\end{Resp}
\begin{Resp}\Rsign Proptérea unxit te Deus, Deus tuus, óleo lætítiæ, allelúja.\end{Resp}
\switchcolumn
\EN
\begin{Resp}\Rsign Thou hast loved righteousness, and hated iniquity; \Star{} Therefore God, thy God, hath anointed thee with the oil of gladness, alleluia.\end{Resp}
\begin{Resp}\Vsign Because of truth, and meekness, and righteousness.\end{Resp}
\begin{Resp}\Rsign Therefore God, thy God, hath anointed thee with the oil of gladness, alleluia.\end{Resp}
\switchcolumn*

\LA
\LessonHead{Lectio VI}
\switchcolumn
\EN
\LessonHead{Lesson VI}
\switchcolumn*

\LA
\noindent Bríxiam ítaque, ubi domum ad sanctæ Afræ templum condúxit, revérsa, novam ibi vírginum societátem, sicut cælésti voce ac visióne mandátum sibi fúerat, sub certa disciplína sanctísque vivéndi régulis constítuit, quam sanctæ Ursulæ invíctæ vírginum ducis patrocínio ac nómine insignívit: eam vero perénnem futúram morti próxima prædíxit. Tandem prope septuagenária, dives méritis evolávit in cælum sexto Kaléndas Februárii anni millésimi quingentésimi et quadragésimi. Cujus cadáver per ipsos trigínta dies inhumátum, flexíbile ac vivo simíllimum perseverávit. Demum in sanctæ Afræ templo inter céteras, quibus illud abúndat, Sanctórum relíquias repósito, plúrima ad ejus sepúlcrum agi statim cœpére mirácula; quorum fama late diffúsa, non Bríxiæ modo et Decentiáni, sed álibi étiam vulgo cœpit nuncupári beáta, ejúsque imágo aris impóni: immo sanctus ipse Cárolus Borromæus non multis post annis dignam, quæ ab apostólica Sede in sanctárum Vírginum album referrétur, Bríxiæ palam asséruit. Cultum vero illi jámdiu a pópulis exhíbitum, et et tum locórum Ordináriis probátum, tum plúribus étiam summórum Pontíficum indúltis munítum, Clemens Papa décimus tértius solémni decréto ratum hábuit ac confirmávit. Eam tandem, novis miráculis rite probátis insígnem, Pius Papa séptimus, solémni canonizatióne in Vaticána basílica perácta, die vigésima quarta Maji anno millésimo octingentésimo séptimo, sanctárum Vírginum catálogo adscrípsit.\par
\switchcolumn
\EN
\noindent She went back to Brescia, and there hired an house near the Church of St. Afra, in which house, in obedience to a vision and command from heaven, she founded a new Order of religious women, constituted under certain rules and holy regulations of life. This Order she put under the name and patronage of St Ursula, the fearless leader of maidens. When Angela was near to death, she foretold that this Order will never cease. She was well-nigh three score and ten years of age, and full of good works, when, in the (night between the) 27th (and the 28th days) of January, in the year 1540, she winged her flight heavenward. Her dead body lay unburied thirty days, supple and life-like. It was laid at last in the Church of St. Afra, where sleep so many more of God's holy children. Diverse miracles forthwith began to be worked at her grave. The fame of these being noised about, she began to be commonly called Blessed, and that not only at Brescia and Decenzano and pictures of her were put over Altars. Not many years afterward, holy Charles Borromeo said openly at Brescia, that she was one whose name the Apostolic See might well enroll among those of holy virgins. The reverence which had of a long time been shown to her memory was approved by the local Ordinaries, confirmed by diverse Papal Indults, and solemnly ratified and established by decree of Pope Clement XIII. As she continued famous for new and proved miracles, Pope Pius VII, at the solemn canonization held in the Vatican Basilica, upon the 24th day of May, in the year 1807, added her name to the list of holy maids.\par
\switchcolumn*

\LA
\begin{Resp}\Rsign Afferéntur Regi vírgines post eam, próximæ ejus; \Star{} Afferéntur tibi in lætítia et exsultatióne, allelúja.\end{Resp}
\begin{Resp}\Vsign Spécie tua et pulchritúdine tua; inténde, próspere procéde, et regna.\end{Resp}
\begin{Resp}\Rsign Afferéntur tibi in lætítia et exsultatióne, allelúja.\end{Resp}
\begin{Resp}\Vsign Glória Patri, et Fílio, \Star{} et Spirítui Sancto.\end{Resp}
\begin{Resp}\Rsign Afferéntur tibi in lætítia et exsultatióne, allelúja.\end{Resp}
\switchcolumn
\EN
\begin{Resp}\Rsign After her shall virgins be brought unto the King, her fellows \Star{} They shall be brought unto thee with gladness and rejoicing, alleluia.\end{Resp}
\begin{Resp}\Vsign In thy comeliness and thy beauty, go forward, fare prosperously, and reign.\end{Resp}
\begin{Resp}\Rsign They shall be brought unto thee with gladness and rejoicing, alleluia.\end{Resp}
\begin{Resp}\Vsign Glory be to the Father, and to the Son, \Star{} and to the Holy Ghost.\end{Resp}
\begin{Resp}\Rsign They shall be brought unto thee with gladness and rejoicing, alleluia.\end{Resp}
\switchcolumn*

\LA
\LessonHead{Lectio IX (pro commemoratione)}
\switchcolumn
\EN
\LessonHead{Lesson IX (when commemorated)}
\switchcolumn*

\LA
\LessonTitle{Commemoratio: S. Angelæ Mericiæ Virginis}
\switchcolumn
\EN
\LessonTitle{Commemoration of: St. Angela Merici, Virgin}
\switchcolumn*

\LA
\noindent Angela Merícia, piis orta paréntibus, a prima ætáte magna virtútum specímina dedit, cilício ac flagéllis frequénter usa, et oratióni indesinénter dédita. Patérnis bonis abdicátis, ac régulam tértii órdinis sancti Francísci ampléxa, evangélicam paupertátem virginitátis laudi conjúnxit, nullúmque pietátis offícium erga próximos omísit. Sacra Eucharístia frequentíssime refécta, tanta caritátis vi ferebátur in Deum, ut sǽpius extra sensus raperétur. Bríxiæ novam vírginum societátem sub certa disciplína sanctísque vivéndi régulis constítuit, quam sanctæ Ursulæ patrocínio ac nómine insignívit. Tandem prope septuagenária evolávit in cælum, anno Dómini millésimo quingentésimo quadragésimo, sexto kaléndas februárii. Cultum illi jámdiu exhíbitum Clemens Papa décimus tértius solémni decréto ratum hábuit et confirmávit. Pius vero Papa séptimus sanctárum Vírginum catálogo eam adscrípsit.\par
\switchcolumn
\EN
\noindent Angela Merici was born of devout parents and from her earliest years gave evidence of great virtue, frequently using sackcloth and scourges, and praying unceasingly. She renounced her inheritance and embraced the rule of the Third Order of St. Francis. Uniting evangelical poverty with the glory of virginity, she withheld no service of kindness to her neighbour. She was frequently refreshed by the Holy Eucharist, and was carried up to God with such force of love as often to be rapt out of her senses. At Brescia, she founded a new society of Virgins, under a fixed discipline and a rule of life conducive to holiness, which she put under the patronage and name of St. Ursula. At last, when she was almost seventy years old, she went to heaven, in the year of the Lord 1540, on January 27th. Pope Clement XIII by a solemn decree ratified and confirmed the veneration which had already been offered her for a long time; and Pope Pius VII enrolled her in the list of holy Virgins.\par
\switchcolumn*

\LA
\TeDeum{Te Deum}
\switchcolumn
\EN
\TeDeum{Te Deum}
\switchcolumn*

\end{paracol}

\par\needspace{10\baselineskip}\addvspace{1.2em}{\centering\small\textsc{Die 2 Junii} · \textsc{2 June}\par}
\Office{Ss. Marcellini, Petri, atque Erasmi, Episcopi, Martyrum}{Sts. Marcellinus, Peter, and Erasmus, Martyrs}{Simplex}
\addcontentsline{toc}{subsection}{Die 2 Junii · Ss. Marcellini, Petri, atque Erasmi, Episcopi, Martyrum\enspace\textit{Sts. Marcellinus, Peter, and Erasmus, Martyrs}}
\begin{paracol}{2}
\LA
\RefLine{Lectiones I–II: de Scriptura occurrente.}
\switchcolumn
\EN
\RefLine{Lessons I–II: from the Scripture of the day.}
\switchcolumn*

\LA
\LessonHead{Lectio III (pro commemoratione)}
\switchcolumn
\EN
\LessonHead{Lesson III (when commemorated)}
\switchcolumn*

\LA
\noindent Petrus exorcísta, Diocletiáno imperatóre, Romæ a Seréno júdice propter christiánæ fídei confessiónem missus in cárcerem, Paulínam Artémii, qui cárceri præerat, filiam a dæmone agitátam liberávit. Quo facto et paréntes puéllæ cum tota família, et vicínos qui ad rei novitátem concúrrerant, Jesu Christo conciliátos ad Marcellínum presbyterum addúxit, a quo omnes baptizáti sunt. Quod ubi rescívit Serénus, Petrum et Marcellínum ad se vocátos aspérius objúrgat, et ad verbórum acerbitátem minas ac terróres adjúngit, nisi Christo renúntient. Cui cum Marcellínus christiána libertáte respondéret, pugnis contúsum, et a Petro sejúnctum, nudum inclúdit in cárcerem stratum vitri fragméntis, sine cibo ac sine lúmine. Petrum item constríngi ímperat arctíssimis vínculis. Sed cum utríque ex torméntis fides et ánimus crésceret, constánti confessióne, et abscísso cápite, illústre testimónium Jesu Christo dedérunt. Erásmus epíscopus, imperatóribus Diocletiáno et Maximiáno, in Campánia plumbátis et fústibus cæsus, resína quoque, súlphure, plumbo liquefácto, et fervénti pice, cera, oleóque perfúsus, inde tamen ínteger et inviolátus evásit. Quo miráculo multi se ad Christi fidem convertérunt. Verum is, íterum detrúsus in cárcerem, constríctus férreis gravissimísque vínculis, inde ab Angelo mirabíliter eréptus est. Deínde Fórmiis a Maximiáno váriis afféctus supplíciis, tunicáque ærea candénti indútus, illa étiam torménta, divína virtúte superávit. Dénique, plúrimis et in fide confirmátis, et ad fidem convérsis, insígnem martyrii palmam adéptus est.\par
\switchcolumn
\EN
\noindent This Peter was an exorcist, whom, in the reign of the Emperor Diocletian, Serenus the Judge cast into prison at Rome because he confessed the Christian faith. He there set free Paulina, the daughter of Artemius, the keeper of the prison, from an evil spirit which tormented her. Upon this, Artemius and his wife and all their house, with their neighbours who had run together to see the strange thing, would fain be made friends with Jesus Christ. Peter therefore brought them to Marcellinus, the Priest, who baptized them all. When Serenus heard of it, he called Peter and Marcellinus before him, and sharply rebuked them, adding to his bitter words threats and terrors, unless they would deny Christ. Marcellinus answered him with Christian boldness, whereupon he caused him to be buffeted, separated him from Peter and shut him up naked in a prison strewn with broken glass, without either food or light. Peter also he straitly confined. But when both of them were found to wax faithfuller and braver in their bonds, they were beheaded, unshaken in their testimony, and confessing Jesus Christ gloriously by their blood. Elmo was a Bishop in Campania who, (in the year 303,) in the reign of the Emperors Diocletian and Maximian was beaten with clubs and whips loaded with lead, and afterwards anointed with melted pitch, sulphur, and lead, and boiling resin, wax, and oil. From all this he came forth whole and sound which wonder turned many to believe in Christ. He was remanded again to prison, and straitly bound in heavy iron fetters. But from these he was wondrously delivered by an angel. At last, at Formi, Maximian caused him to be subjected to diverse torments, and in the end being clad in a coat of redhot brass the power of God made him to be more than conqueror in this thing also, and to grasp the palm-branch of a glorious testimony, whereby he strengthened many in the faith and turned many to it.\par
\switchcolumn*

\LA
\TeDeum{Te Deum}
\switchcolumn
\EN
\TeDeum{Te Deum}
\switchcolumn*

\end{paracol}

\par\needspace{10\baselineskip}\addvspace{1.2em}{\centering\small\textsc{Die 4 Junii} · \textsc{4 June}\par}
\Office{S. Francisci Caracciolo Confessoris}{St. Francis Caracciolo, Confessor}{Duplex}
\addcontentsline{toc}{subsection}{Die 4 Junii · S. Francisci Caracciolo Confessoris\enspace\textit{St. Francis Caracciolo, Confessor}}
\begin{paracol}{2}
\LA
\RefLine{Lectiones I–III: de Scriptura occurrente.}
\switchcolumn
\EN
\RefLine{Lessons I–III: from the Scripture of the day.}
\switchcolumn*

\LA
\RefLine{Lectiones VII–IX: ex Commúni Confessoris non pontificis (C5).}
\switchcolumn
\EN
\RefLine{Lessons VII–IX: from the Common of a Confessor not a Bishop (C5).}
\switchcolumn*

\LA
\LessonHead{Lectio IV}
\switchcolumn
\EN
\LessonHead{Lesson IV}
\switchcolumn*

\LA
\noindent Francíscus, dictus ántea Ascánius, ex nóbili família Carácciolo in óppido sanctæ Maríæ de Villa in Aprútio ortus, a primis annis exímio enítuit pietátis cultu. Adoléscens gráviter ægrótans státuit sese prorsus Dei proximíque mancipáre servítio. Neápolim proféctus, sacerdótio initiátus sacróque adscríptus sodalítio, contemplatióni lucrandísque animábus se totum devóvit, ac extrémo supplício damnátis hortatórem se præbuit assíduum. Cóntigit autem, ut epístolum álteri destinátum ei per errórem redderétur, quo a piíssimis viris Joánne Augustíno Adórno et Fabrício Carácciolo ad novi religiósi institúti fundatiónem vocabátur. Rei novitáte captus et divínæ voluntátis demirátus consília, álacri ánimo sese illis adjúnxit. Cónditis autem in Camaldulénsium erémo, quo secésserant, novi órdinis légibus, inde Romam simul profécti confirmatiónem a Xysto quinto impetrárunt, qui eósdem cléricos reguláres Minóres appellári vóluit, áddito ad tria consuéta áltero de non ambiéndis dignitátibus voto.\par
\switchcolumn
\EN
\noindent This Francis, whose worldly name was Ascanius, was one of the noble family of Caracciolo. He was born in the town of Santa Maria della Villa, in the Abruzzi, \textasciitilde{}(on the 13th day of October, in the year of grace 1563.) From his earliest years he showed great marks of godliness. When he was a young man he had a severe illness, and on his recovery determined to serve God only, and bade farewell to the world. He betook himself to Naples, where he was ordained Priest, enrolled himself in a devout guild, and gave himself up altogether to seek after God, and to gain souls for Him, in which work he showed himself an unwearied comforter to such prisoners as were condemned to death. It came to pass that those two great servants of God, John Augustine Adorno and Fabricius Caracciolo, wrote a letter to a certain person, wherein they exhorted him to found a new religious Institute. This letter came by a mistake to be delivered to Francis Caracciolo. The newness of the idea and the strange ways of God's Providence took possession of his mind, and he joyfully added himself to their company. They withdrew themselves to the wilderness of the Camaldolese hermits (near Naples,) and there concerted the Rule of the New Order. Thence they went together to Rome, and obtained the confirmation of their work from Sixtus V, who was pleased that they should be called The Lesser Clerks Regular, since they add to the three accustomed vows (of Poverty, Chastity, and Obedience,) a fourth, binding themselves not to seek preferment in the Church.\par
\switchcolumn*

\LA
\begin{Resp}\Rsign Honéstum fecit illum Dóminus, et custodívit eum ab inimícis, et a seductóribus tutávit illum: \Star{} Et dedit illi claritátem ætérnam, allelúja.\end{Resp}
\begin{Resp}\Vsign Justum dedúxit Dóminus per vias rectas, et osténdit illi regnum Dei.\end{Resp}
\begin{Resp}\Rsign Et dedit illi claritátem ætérnam, allelúja.\end{Resp}
\switchcolumn
\EN
\begin{Resp}\Rsign The Lord made him honourable, and defended him from his enemies, and kept him safe from those that lay in wait for him; \Star{} And gave him perpetual glory, alleluia.\end{Resp}
\begin{Resp}\Vsign He went down with him into the pit, and left him not in bonds.\end{Resp}
\begin{Resp}\Rsign And gave him perpetual glory, alleluia.\end{Resp}
\switchcolumn*

\LA
\LessonHead{Lectio V}
\switchcolumn
\EN
\LessonHead{Lesson V}
\switchcolumn*

\LA
\noindent Solémni emíssa professióne, ob singulárem ejus in divum Francíscum Assisinátem cultum, Francísci nomen assúmpsit. Adórno biénnio post vita functo, ipse toti religióni, quamquam invítus, præfícitur; quo in múnere virtútum ómnium præclára præbuit exémpla. Institúti amplificándi studiosíssimus, id assíduis oratiónibus, lácrimis, et jugi córporis maceratióne eníxe a Deo postulábat. Quam ob rem tértio in Hispániam se cóntulit, peregríni hábitu indútus, victúmque ostiátim mendícans. In itínere aspérrima quæque perpéssus, Omnipoténtis auxílium mirum in modum expértus, navim quam conscénderat, ab imminénti naufrágio oratiónis præsídio servávit incólumem. Ut in regnis illis voti compos fíeret, plúrimum laborávit; sed, ejus sanctitátis fama prælucénte, amplissimáque catholicórum regum Philíppi secúndi et Philíppi tértii munificéntia, adversariórum conátibus singulári ánimi fortitúdine superátis, plura sui órdinis domicília fundávit: quod pari evéntu per Itáliam præstitit.\par
\switchcolumn
\EN
\noindent Ascanius Caracciolo, moved by a special love and devotion he had to the holy Francis of Assisi, took, when he made his solemn profession, the name of Francis. After two years, John Adorno departed this life, and Francis, against his own will, was made Head of the Order. In this office he shone a burning light of grace. Devoted to the prosperity of the Institute, he earnestly sought the blessing of God upon it, by constant prayer, by tears, and by stern treatment of his own body. In this work, he thrice travelled into Spain in the guise of a pilgrim, and begging his bread from door to door. In these his journeys he suffered very great hardships, and was most wonderfully helped by the Almighty, especially one while when he was on shipboard and the ship nigh to perish, but for the work of his prayers. He toiled hard in those countries to attain ' his wishes, but through the widespread fame of his holy life, and the noble generosity of the Most Catholic Kings Philip II. and Philip III., he overcame with his brave perseverance the opposition of all that withstood him, and founded several houses of his Order. This he was able to do in Italy also.\par
\switchcolumn*

\LA
\begin{Resp}\Rsign Amávit eum Dóminus, et ornávit eum: stolam glóriæ índuit eum, \Star{} Et ad portas paradísi coronávit eum, allelúja.\end{Resp}
\begin{Resp}\Vsign Índuit eum Dóminus lorícam fídei, et ornávit eum.\end{Resp}
\begin{Resp}\Rsign Et ad portas paradísi coronávit eum, allelúja.\end{Resp}
\switchcolumn
\EN
\begin{Resp}\Rsign The Lord loved him and beautified him; He clothed him with a robe of glory, \Star{} And crowned him at the gates of Paradise, alleluia.\end{Resp}
\begin{Resp}\Vsign The Lord hath put on him the breast-plate of faith, and hath adorned him.\end{Resp}
\begin{Resp}\Rsign And crowned him at the gates of Paradise, alleluia.\end{Resp}
\switchcolumn*

\LA
\LessonHead{Lectio VI}
\switchcolumn
\EN
\LessonHead{Lesson VI}
\switchcolumn*

\LA
\noindent Humilitáte ádeo excélluit, ut Romam véniens, in páuperum hospítio recéptus, se lepróso sociáverit, et ecclesiásticas dignitátes a Paulo quinto sibi oblátas constantíssime recusáverit. Illibátam perpétuo servávit virginitátem, effrontésque mulíeres, ejus castimóniæ insidiántes, Christo lucrifécit. Erga diviníssimum Eucharístiæ mystérium ardénti æstuans amóre, noctes pene íntegras in ejus adoratióne insómnes ducébat: quod pium exercítium, véluti sui órdinis tésseram, in eo perpétuo servándum constítuit. Deíparæ Vírginis cultum impénse fovit. In próximum exímia exársit caritáte. Prophetíæ dono et córdium scrutatióne ditátus fuit. Quadragésimum quartum ætátis suæ annum agens, dum in sacra Lauretána, æde in oratióne persísteret, sibi vitæ finem imminére cognóvit. Aprútium statim defléxit, et in óppido Agnóni apud alúmnos sancti Philíppi Nerii letháli febre corréptus, sacraméntis Ecclésiæ devotíssime suscéptis, prídie Nonas Júnii anni millésimi sexcentésimi octávi, in pervigílio festi Córporis Christi, placidíssime obdormívit in Dómino. Sacrum ejus corpus, Neápolim delátum, in ecclésia sanctæ Maríæ Majóris, ubi prima sui órdinis jécerat fundaménta, honorífice cónditum fuit. Eum póstea, miráculis clarum, Clemens decimus quartus Póntifex máximus solémni ritu inter Beátos; Pius vero séptimus Póntifex máximus, novis fulgéntem signis, anno millésimo octingentésimo séptimo, Sanctórum albo adscrípsit.\par
\switchcolumn
\EN
\noindent There was a great pattern of lowliness, so that when he came to Rome he betook himself to an almshouse, and chose a leper for his familiar friend. Paul V. offered him diverse honours in the Church, but he firmly refused them all. He preserved his purity unspotted, and when certain shameless women set themselves to attack his chastity, he took the occasion to gain over their souls for Christ. Toward God's great mystery of the Eucharist he was drawn with passionate tenderness, and would pass almost whole nights without sleep, simply adoring It. This godly custom he established in his Order, to be kept up therein for ever, the peculiar mark thereof. He was a great encourager of the worship of the Maiden Mother of God. He was hot with strong love for his neighbour. He was gifted with prophecy, and the discerning of spirits. In the forty-fourth year of his age he was continuing long in prayer in the Holy House of Loretto, when it was made known to him that the end of his earthly life was at hand. He straightway took his way to the Abruzzi, and was there seized with illness while he was with the disciples of St. Philip Neri, in the town of Agnone. He received with great devotion the Sacraments of the Church, and then, upon the 4th day of June, being the Eve of the Feast of the Body of Christ, in the year 1608, he very peacefully fell asleep in the Lord. His sacred body was carried to Naples, and there honourably buried in the Church of St. Mary the Greater, where he had laid the first foundations of his Order. As he became distinguished for miracles Pope Clement XIV enrolled his name, with solemn pomp, among those of the Blessed, and Pope Pius VII, in the year 1807, finding his mighty works continue, added it to the list of the Saints.\par
\switchcolumn*

\LA
\begin{Resp}\Rsign Iste homo perfécit ómnia quæ locútus est ei Deus, et dixit ad eum: Ingrédere in réquiem meam: \Star{} Quia te vidi justum coram me ex ómnibus géntibus, allelúja.\end{Resp}
\begin{Resp}\Vsign Iste est, qui contémpsit vitam mundi, et pervénit ad cæléstia regna.\end{Resp}
\begin{Resp}\Rsign Quia te vidi justum coram me ex ómnibus géntibus, allelúja.\end{Resp}
\begin{Resp}\Vsign Glória Patri, et Fílio, \Star{} et Spirítui Sancto.\end{Resp}
\begin{Resp}\Rsign Quia te vidi justum coram me ex ómnibus géntibus, allelúja.\end{Resp}
\switchcolumn
\EN
\begin{Resp}\Rsign This is he which did according unto all that God commanded him; and God said unto him: Enter thou into My rest, \Star{} For thee have I seen righteous before Me among all people, alleluia.\end{Resp}
\begin{Resp}\Vsign This is he which loved not his life in this world, and is come unto an everlasting kingdom.\end{Resp}
\begin{Resp}\Rsign For thee have I seen righteous before Me among all people, alleluia.\end{Resp}
\begin{Resp}\Vsign Glory be to the Father, and to the Son, \Star{} and to the Holy Ghost.\end{Resp}
\begin{Resp}\Rsign For thee have I seen righteous before Me among all people, alleluia.\end{Resp}
\switchcolumn*

\LA
\LessonHead{Lectio IX (pro commemoratione)}
\switchcolumn
\EN
\LessonHead{Lesson IX (when commemorated)}
\switchcolumn*

\LA
\LessonTitle{Commemoratio: S. Francisci Caracciolo Confessoris}
\switchcolumn
\EN
\LessonTitle{Commemoration of: St. Francis Caracciolo, Confessor}
\switchcolumn*

\LA
\noindent Francíscus, dictus ántea Ascánius, ex nóbili família Carácciolo, in óppido sanctæ Maríæ de Villa in Aprútio ortus est. Adoléscens gráviter ægrótans státuit sese prorsus Dei proximíque mancipáre servítio. Neápolim proféctus et sacerdótio initiátus, contemplatióni lucrandísque animábus se totum devóvit, ac extrémo supplício damnátis hortatórem se prǽbuit assíduum. Joánni Augustíno Adórno et Fabrício Carácciolo, mira Dei dispositióne, adjúnctus, Clericórum regulárium Minórum órdinem instítuit, áddito ad tria consuéta áltero de non ambiéndis dignitátibus voto; quem, post óbitum Adórni, sanctíssime rexit, et summo stúdio per Hispániam et Itáliam propagávit. Erga sanctíssimæ Eucharístiæ sacraméntum tanto æstuábat afféctu, ut noctes pene íntegras in ejus adoratióne impénderet; quod pium exercítium, véluti sui órdinis tésseram, perpétuo in eo servándum constítuit. Tandem prophetíæ dono et córdium scrutatióne ditátus, quadragésimum quartum annum agens, in óppido Agnóni in Aprútio letháli febre corréptus, in Dómino obdormívit prídie nonas júnii, anno millésimo sexcentésimo octávo. Sacrum ejus corpus, Neápolim delátum, in ecclésia sui órdinis cónditum est.\par
\switchcolumn
\EN
\noindent Francis, originally called Ascanio, came from the noble family of the Caracciolo, in the city of Santa Maria della Villa in the Abruzzi. During a severe illness in his youth, he resolved to give himself up completely to the service of God and neighbour. He went to Naples and was ordained priest, after which he devoted himself wholly to contemplation and the winning of souls; and he was always on hand to comfort prisoners condemned to death. By a wonderful disposition of God, he joined with John Augustine Adorno and Fabrizio Caracciolo in founding the Order of Clerks Regular Minor, adding to the three usual vows another: not to seek dignities. After the death of Adorno, Francis ruled the Order in a most holy way, and with the greatest zeal promoted it throughout Spain and Italy. He burned with such love for the Most Blessed Sacrament that he would spend almost the whole night in adoring it; and he established adoration of the Blessed Sacrament as a practice to be kept for ever in his Order, the peculiar mark thereof. He was enriched with the gifts of prophecy and the reading of hearts. At length, in his forty-fourth year, he was taken with a deadly fever in the town of Agnone, in the Abruzzi, and fell asleep in the Lord on the 4th day of June, 1608. His holy body was taken to Naples and buried in the church of his Order.\par
\switchcolumn*

\LA
\TeDeum{Te Deum}
\switchcolumn
\EN
\TeDeum{Te Deum}
\switchcolumn*

\end{paracol}

\par\needspace{10\baselineskip}\addvspace{1.2em}{\centering\small\textsc{Die 5 Junii} · \textsc{5 June}\par}
\Office{S. Bonifatii Episcopi et Martyris}{St. Boniface, Bishop and Martyr}{Duplex}
\addcontentsline{toc}{subsection}{Die 5 Junii · S. Bonifatii Episcopi et Martyris\enspace\textit{St. Boniface, Bishop and Martyr}}
\begin{paracol}{2}
\LA
\RefLine{Lectiones I–III: de Scriptura occurrente.}
\switchcolumn
\EN
\RefLine{Lessons I–III: from the Scripture of the day.}
\switchcolumn*

\LA
\LessonHead{Lectio IV}
\switchcolumn
\EN
\LessonHead{Lesson IV}
\switchcolumn*

\LA
\noindent Bonifátius, ántea Winfrídus appellátus, apud Anglos natus est exeúnte sæculo séptimo, et ab ipsa infántia mundum aversátus, vitam monásticam in votis hábuit. Cum ejus pater ánimum sæculi illécebris permutáre frustra tentásset, monastérium ingréditu, et sub beáti Wolphárdi disciplína ómnium virtútum ac scientiárum génere imbúitur. Annum agens trigésimum sacerdótio insignítur, ac verbi divíni prædicátor assíduus, magno animárum lucro hoc in múnere versátur. Attamen, regnum Christi adaugére desíderans, contínuo flebat ingéntem multitúdinem barbarórum, qui ignorántiæ ténebris immérsi dæmoni famulabántur. Qui quidem animárum zelus cum in dies inexstinguíbili ardóre accrésceret, divíno númine per lácrimas et oratiónes exploráto, facultátem a monastérii præpósito obtínuit ad Germánicas oras proficiscéndi.\par
\switchcolumn
\EN
\noindent Winfried, afterwards called Boniface, was an Englishman, and born in England, towards the end of the seventh century. From his very childhood, he turned away from the world, and set his heart upon becoming a monk. His father tried in vain to turn him from his wishes by the beguilements of the world, and he entered a Monastery, where the Blessed Wolphard instructed him in all godliness and diverse kinds of learning. At the age of twenty-nine years he was ordained Priest, and became an unwearied preacher of the Word of God, wherein he had a gift which he used with great gain of souls. Nevertheless, his great desire was to spread the kingdom of Christ, and he continually bewailed the vast number of savages who were plunged in the darkness of ignorance and were the servants of the devil. This zealous love of souls increased in him in intensity day by day, till nothing would serve him, but, having implored the blessing of God by tears and prayers, and obtained authority from the head of his monastery, to set forth for the coast of Germany.\par
\switchcolumn*

\LA
\begin{Resp}\Rsign Lux perpétua lucébit Sanctis tuis, Dómine, \Star{} Et ætérnitas témporum, allelúja, allelúja.\end{Resp}
\begin{Resp}\Vsign Lætítia sempitérna erit super cápita eórum: gáudium et exsultatiónem obtinébunt.\end{Resp}
\begin{Resp}\Rsign Et ætérnitas témporum, allelúja, allelúja.\end{Resp}
\switchcolumn
\EN
\begin{Resp}\Rsign Eternal light will shine over your saints, O Lord, \Star{} And the eternity of the times, alleluia, alleluia.\end{Resp}
\begin{Resp}\Vsign Everlasting joy shall be upon their heads, they shall obtain joy and gladness\end{Resp}
\begin{Resp}\Rsign And the eternity of the times, alleluia, alleluia.\end{Resp}
\switchcolumn*

\LA
\LessonHead{Lectio V}
\switchcolumn
\EN
\LessonHead{Lesson V}
\switchcolumn*

\LA
\noindent Ex Anglia duóbus cum sóciis navem solvens, Dorestádium in Frísiæ óppidum venit. Cum autem bellum gravíssimum inter Frísonum regem Radbódum et Cárolum Martéllum exarsísset, sine fructu Evangélium prædicávit. Quaprópter in Angliam revérsus, ad suum redívit monastérium, cui invítus præfícitur. Post elápsum biénnium, ex consénsu epíscopi Vintoniénsis, munus abdicávit, et Romam proféctus est, ut apostólica auctoritáte ad Gentílium conversiónem delegarétur. Cum ad Urbem pervenísset, a Gregório secúndo benígne excípitur, pro Winfrído Bonifátius a Pontífice nominátur. In Germániam diréctus, Thuríngiæ Saxoniæque pópulis Christum annuntiávit. Cum intérea Radbódus, Frísiæ rex ac infestíssimus christiáni nóminis hostis, occubuísset, Bonifátius ad Frísones rédiit, ubi sancti Willibrórdi sócius per triénnium tanto cum fructu Evangélium prædicávit, ut, destrúctis idolórum simulácris, innúmeræ vero Deo ecclésiæ excitaréntur.\par
\switchcolumn
\EN
\noindent He set sail from England with two companions (in the year 716) and reached the town of Dorestadt in Friesland. A great war being then raging between Radbod, King of the Frieslanders, and Charles Martel, Winfrid preached the Gospel in vain. He went back to England, and betook himself again to his Monastery, whereof he was, against his own will, chosen to be the head. After two years he obtained the consent of the Bishop of Winchester to resign his office, and (in 719) went to Rome, to seek an Apostolic commission to preach to the heathen. When he arrived at the city he was courteously welcomed by Gregory II., who changed his name from Winfrid to Boniface. He departed thence to Germany, and preached Christ to the tribes in Thuringia and Saxony. Radbod, King of Friesland, who bitterly hated the Christian name, being dead, Boniface went a second time among the Frieslanders, and there, with his comrade St. Willibrord, preached the Gospel for three years with so much fruit, that the idols were hewn down, and countless churches arose to the true God.\par
\switchcolumn*

\LA
\begin{Resp}\Rsign In servis suis, allelúja, \Star{} Consolábitur Deus, allelúja.\end{Resp}
\begin{Resp}\Vsign Judicábit Dóminus pópulum suum, et in servis suis.\end{Resp}
\begin{Resp}\Rsign Consolábitur Deus, allelúja.\end{Resp}
\switchcolumn
\EN
\begin{Resp}\Rsign His servants, alleluia \Star{} God will comfort, alleluia, alleluia.\end{Resp}
\begin{Resp}\Vsign The Lord will judge His people and, His servants,\end{Resp}
\begin{Resp}\Rsign God will comfort, alleluia, alleluia.\end{Resp}
\switchcolumn*

\LA
\LessonHead{Lectio VI}
\switchcolumn
\EN
\LessonHead{Lesson VI}
\switchcolumn*

\LA
\noindent A sancto Willibrórdo ad episcopále munus expetítus, illud detrectávit ut prómpius infidélium salúti instáret. In Germániam proféctus, plura Hassórum míllia a dæmonis superstitióne avocávit. A Gregório Pontífice Romam evocátus, post insígnem fídei professiónem epíscopus consecrátur. Exínde ad Germános redux, Hássiam et Thuríngiam ab idololatríæ relíquiis pénitus expurgávit. Tanta propter mérita Bonifátius a Gregório tértio ad dignitátem archiepiscopálem evéhitur, et tértio Romam proféctus a summo Pontífice Sedis apostólicæ legátus constitúitur. Qua insignítus auctoritáte quátuor episcopátus instítuit, et várias synodos celebrávit, inter quas concílium Leptinénse memorábile est, apud Belgas in Cameracénsi diœcési celebrátum, quo quidem témpore ad fidem in Bélgio adaugéndam egrégie cóntulit. A Zacharía Papa creátus Mogúntius archiepíscopus, ipso Pontífice jubénte, Pipínum in regem Francórum unxit. Post mortem sancti Willibrórdi Ultrajecténsem ecclésiam gubernándum suscépit, primo per Eóbanum, deínde per seípsum, dum ab ecclésia Moguntína absolútus Ultrajécti resédit. Frisónibus ad idololatríam relápsis, Evangélium prædicáre rursus aggréditur, cumque offício pastoráli occuparétur, a bárbaris et ímpiis homínibus juxta Bornam flúvium cum Eóbano coëpíscopo multísque áliis cruénta cæde perémptus martyrii palma condecorátur. Corpus sancti Bonifátii Mogúntiam translátum, et, ut ipse vivens petíerat, in Fuldénsi monastério, quod exstrúxerat, recónditum fuit, ubi multis miráculis incláruit. Pius autem nonus Póntifex máximus, ejus Offícium et Missam ad univérsam Ecclésiam exténdit.\par
\switchcolumn
\EN
\noindent Willibrord urged upon him to take the office of a Bishop, but he deferred to seek it, that he might the more instantly toil for the salvation of the unbelievers. Advancing into Germany, he reclaimed thousands of the Hessians from devil-worship. Pope Gregory sent for him to Rome, (whither he came in 723,) and after hearing a noble profession of his faith, consecrated him a Bishop. He again returned to Germany, and thoroughly purged Hesse and Thuringia from all remains of idolatry. On account of such great works Gregory III. advanced Boniface to the dignity of an Archbishop, and on the occasion of a third journey to Rome, (in 738,) he was invested by the Sovereign Pontiff with the powers of Legate of the Apostolic See. As such, he founded (the) four Bishoprics (of Erfurt, Paderborn, Wurtzburg, and Eichstadt,) and held diverse Synods, among which is especially to be remembered that of Lessines, held in Belgium, in the diocese of Cambrai, wherein he made his strongest endeavours to spread the Faith among the Belgians. By Pope Zacharias, he was named Archbishop of Maintz, and by command of the same Pope, he anointed Pepin to be King of the Franks. After the death of St. Willibrord, he undertook the government of the Church of Utrecht, at first through Eoban but he afterwards was released from the care of the Church of Maintz and established his see at Utrecht. The Frieslanders having again fallen back into idolatry, he once more betook himself to preach the Gospel among them, and while he was busied in this duty, he grasped the crown of martyrdom, being murdered by some ungodly savages, along with his fellow-Bishop Eoban, and many others, in a bloody massacre near the River Born, (on the th day of June, in the year of our Lord 755, and of his own age the 75th.) In accordance with the wish expressed by himself during life the body of St. Boniface was carried to Maintz, and buried in the monastery of Fulda, of which he had been the founder, and where God has gloriously honoured it by the working of many signs and wonders. Pope Pius IX. ordered the Office and Mass in bis memory to be used throughout the whole Church.\par
\switchcolumn*

\LA
\begin{Resp}\Rsign Fíliæ Jerúsalem, veníte et vidéte Mártyres cum corónis, quibus coronávit eos Dóminus \Star{} In die solemnitátis et lætítiæ, allelúja.\end{Resp}
\begin{Resp}\Vsign Quóniam confortávit seras portárum tuárum, benedíxit fílios tuos in te.\end{Resp}
\begin{Resp}\Rsign In die solemnitátis et lætítiæ, allelúja.\end{Resp}
\begin{Resp}\Vsign Glória Patri, et Fílio, \Star{} et Spirítui Sancto.\end{Resp}
\begin{Resp}\Rsign In die solemnitátis et lætítiæ, allelúja.\end{Resp}
\switchcolumn
\EN
\begin{Resp}\Rsign Come forth, O ye daughters of Jerusalem, and behold the Martyrs with the crowns wherewith the Lord crowned them; \Star{} In the day of His feasting, and of His gladness. Alleluia.\end{Resp}
\begin{Resp}\Vsign For He hath strengthened the bars of thy gates He hath blessed thy children within thee.\end{Resp}
\begin{Resp}\Rsign In the day of His feasting, and of His gladness. Alleluia.\end{Resp}
\begin{Resp}\Vsign Glory be to the Father, and to the Son, \Star{} and to the Holy Ghost.\end{Resp}
\begin{Resp}\Rsign In the day of His feasting, and of His gladness. Alleluia.\end{Resp}
\switchcolumn*

\LA
\LessonHead{Lectio VII}
\switchcolumn
\EN
\LessonHead{Lesson VII}
\switchcolumn*

\LA
\LessonTitle{Léctio sancti Evangélii secúndum Matthǽum.}
\switchcolumn
\EN
\LessonTitle{From the Holy Gospel according to Matthew}
\switchcolumn*

\LA
\Cite{Matt 5:1-12}
\switchcolumn
\EN
\Cite{Matt 5:1-12}
\switchcolumn*

\LA
\noindent In illo témpore: Videns Jesus turbas, ascéndit in montem, et cum sedísset, accessérunt ad eum discípuli ejus. Et réliqua.\par
\switchcolumn
\EN
\noindent At that time Jesus, seeing the multitudes, went up into a mountain, and, when He was set, His disciples came unto Him. And so on.\par
\switchcolumn*

\LA
\LessonTitle{Homilia sancti Augustíni Epíscopi.}
\switchcolumn
\EN
\LessonTitle{Homily by St. Augustine, Bishop (of Hippo.)}
\switchcolumn*

\LA
\Source{Liber 1 de sermone Dómini in monte, cap. 2. et 3.}
\switchcolumn
\EN
\Source{Bk. i. on the Lord's Sermon.}
\switchcolumn*

\LA
\noindent Beáti mundo corde; quóniam ipsi Deum vidébunt. Quam ergo stulti sunt, qui Deum istis exterióribus óculis quærunt, cum corde videátur, sicut álibi scriptum est: Et in simplicitáte cordis quærite illum. Hoc est enim mundum cor, quod est simplex cor. Et quemádmodum lumen hoc vidéri non potest, nisi óculis mundis; ita nec Deus vidétur, nisi mundum sit illud quo vidéri potest. Beáti pacífici; quóniam ipsi fílii Dei vocabúntur. In pace perféctio est, ubi nihil repúgnat; et ídeo fílii Dei pacífici, quóniam nihil in his resístit Deo, et útique fílii similitúdinem patris habére debent.\par
\switchcolumn
\EN
\noindent Blessed are the pure in heart, for they shall see God. What fools then be they that seek God with their outward eyes, since it is in the heart that He is seen, as it is written elsewhere In simplicity of heart seek Him. (Wisd. i. i.) A simple heart is a pure heart. And even as we cannot see this earthly light, unless the eyes be open, so cannot God be seen, unless that be open which alone can perceive Him. Blessed are the peacemakers, for they shall be called the children of God. The perfection of peace is the absence of contrariety, and the peacemakers are called the children of God because they offer no contrariety against the will of God. As beseemeth children, they have their Father's likeness.\par
\switchcolumn*

\LA
\begin{Resp}\Rsign Ego sum vitis vera, et vos pálmites: \Star{} Qui manet in me, et ego in eo, hic fert fructum multum, allelúja, allelúja.\end{Resp}
\begin{Resp}\Vsign Sicut diléxit me Pater, et ego diléxi vos.\end{Resp}
\begin{Resp}\Rsign Qui manet in me, et ego in eo, hic fert fructum multum, allelúja, allelúja.\end{Resp}
\switchcolumn
\EN
\begin{Resp}\Rsign I am the vine: you the branches. \Star{} He that abideth in me, and I in him, the same beareth much fruit, alleluia, alleluia.\end{Resp}
\begin{Resp}\Vsign As the Father hath loved me, I also have loved you.\end{Resp}
\begin{Resp}\Rsign He that abideth in me, and I in him, the same beareth much fruit, alleluia, alleluia.\end{Resp}
\switchcolumn*

\LA
\LessonHead{Lectio VIII}
\switchcolumn
\EN
\LessonHead{Lesson VIII}
\switchcolumn*

\LA
\noindent Pacífici autem in semetípsis sunt, qui omnes ánimi sui motus componéntes, et subjiciéntes ratióni, id est menti et spirítui, carnalésque concupiscéntias habéntes edómitas, fiunt regnum Dei. In quo ita sunt ordináta ómnia, ut id quod est in hómine præcípuum et excéllens, hoc ímperet, céteris non reluctántibus, quæ sunt nobis bestiísque commúnia; atque idípsum quod excélluit in hómine, id est mens et rátio, subjiciátur potióri, quod est ipsa véritas, unigénitus Fílius Dei. Neque enim imperáre inferióribus potest, nisi superióri se ipse subjíciat. Et hæc est pax, quæ datur in terra homínibus bonæ voluntátis; hæc vita consummáti perfectíque sapiéntis.\par
\switchcolumn
\EN
\noindent They are peacemakers in themselves, who order all the movements of their own mind in obedience to reason, that is, to their intellect and soul, and so doing, and taming the lusts of the flesh, become a kingdom for God. In such kingdom all things are so ordered, that the chiefest and noblest part of man ruleth without contention over those lower things which we have in common with beasts. And just in the same way, must that nobler part of man, that is to say, intellect and reason, needs be put in subjection to what is above it, namely, Truth, the Only begotten Son of God. He only can rule well who hath learnt to obey. And this ordering is that peace which is given on earth to men of good will this is the life of whomsoever is thoroughly and perfectly wise.\par
\switchcolumn*

\LA
\begin{Resp}\Rsign Cándidi facti sunt Nazarǽi ejus, allelúja: splendórem Deo dedérunt, allelúja: \Star{} Et sicut lac coaguláti sunt, allelúja, allelúja.\end{Resp}
\begin{Resp}\Vsign Candidióres nive, nitidióres lacte, rubicundióres ébore antíquo, sapphíro pulchrióres.\end{Resp}
\begin{Resp}\Rsign Et sicut lac coaguláti sunt, allelúja, allelúja.\end{Resp}
\begin{Resp}\Vsign Glória Patri, et Fílio, \Star{} et Spirítui Sancto.\end{Resp}
\begin{Resp}\Rsign Et sicut lac coaguláti sunt, allelúja, allelúja.\end{Resp}
\switchcolumn
\EN
\begin{Resp}\Rsign Her Nazarites are become pure, Alleluia: they reflect the glory of God, Alleluia. \Star{} They are whiter than milk. Alleluia, Alleluia.\end{Resp}
\begin{Resp}\Vsign They are purer than snow, they are whiter than milk, they are more ruddy in body than coral, their polishing is of sapphire.\end{Resp}
\begin{Resp}\Rsign They are whiter than milk. Alleluia, Alleluia.\end{Resp}
\begin{Resp}\Vsign Glory be to the Father, and to the Son, \Star{} and to the Holy Ghost.\end{Resp}
\begin{Resp}\Rsign They are whiter than milk. Alleluia, Alleluia.\end{Resp}
\switchcolumn*

\LA
\LessonHead{Lectio IX}
\switchcolumn
\EN
\LessonHead{Lesson IX}
\switchcolumn*

\LA
\noindent De hujúsmodi regno pacatíssimo et ordinatíssimo missus est foras princeps hujus sæculi, qui pervérsis inordinatísque dominátur. Hac pace intrínsecus constitúta atque firmáta, quascúmque persecutiónes ille, qui foras missus est, forínsecus concitáverit, auget glóriam, quæ secúndum Deum est; non áliquid in illo ædifício labefáctans, sed deficiéntibus máchinis suis innotéscere fáciens, quanta fírmitas intus exstrúcta sit. Ideo séquitur: Beáti, qui persecutiónem patiúntur propter justítiam; quóniam ipsórum est regnum cælórum.\par
\switchcolumn
\EN
\noindent From this most peaceful and most orderly kingdom is cast forth the prince of this world, whose rule is over the contentious and disorderly. When once this peace hath been proclaimed and established within, whatsoever wars he that is without can raise, can but heap more glory upon that glory which is according to God, for nothing of the castle will yield before him, but the yielding of his own engines will witness how strong be its ramparts. And therefore cometh next Blessed are they which are persecuted for righteousness' sake, for theirs is the kingdom of heaven.\par
\switchcolumn*

\LA
\TeDeum{Te Deum}
\switchcolumn
\EN
\TeDeum{Te Deum}
\switchcolumn*

\LA
\LessonHead{Lectio IX (pro commemoratione)}
\switchcolumn
\EN
\LessonHead{Lesson IX (when commemorated)}
\switchcolumn*

\LA
\LessonTitle{Commemoratio: S. Bonifatii Episcopi et Martyris}
\switchcolumn
\EN
\LessonTitle{Commemoration of: St. Boniface, Bishop and Martyr}
\switchcolumn*

\LA
\noindent Bonifátius, ántea Winfrídus appellátus, apud Anglos natus est exeúnte sǽculo séptimo. Monastérium ingréssus et sacerdótio auctus, magno animárum lucro in prædicatóris múnere est versátus. Zelo augéndæ fídei accénsus, apud Frísones Evangélium prædicávit. In Angliam revérsus, cum per biénnium monastério sanctíssime præfuísset, superióris múnere abdicáto, Romam se cóntulit, ubi a Gregório secúndo Bonifátii nomen accépit, et in Germániam missus, Thuríngiæ Saxoniǽque pópulis Christum annuntiávit. Ad Frísones revérsus, cum sancto Willibrórdo, magno fructu Evangélium prædicávit. Mox Romam accersítus, episcopáli dignitáte insignítur, et in Germániam íterum proféctus, Hássiam et Thuríngiam ab idololatríæ relíquiis pénitus expurgávit. Sedis apostólicæ legátus creátus et Moguntínus archiepíscopus, plures eréxit et, per se vel per discípulos, administrávit ecclésias. Frisónibus demum ad idololatríam relápsis Evangélium prædicáre rursus aggréssus, cum Eóbano coëpíscopo multísque áliis, juxta Bornam flúvium cruénta cæde perémptus, martýrii palmam accépit. Ejus corpus in Fuldénsi monastério cónditum est.\par
\switchcolumn
\EN
\noindent Boniface, originally called Winfrid, was born in England towards the end of the seventh century. After entering a monastery and becoming a priest, he showed great skill in winning souls through preaching. Burning with zeal to spread the faith, he preached the Gospel among the Frisians. Then he returned to England, where he ruled his monastery for two years in a most holy manner. Having resigned the office of Superior, he went to Rome, where he received from Gregory II the name of Boniface and the commission to proclaim Christ to the peoples of Thuringia and Saxony. With holy Willibrord, he returned to the Frisians and preached the Gospel with great fruit. Soon he was summoned to Rome and invested with the episcopal dignity; after which, he set out once more for Germany. There he rid Hesse and Thuringia of almost the last vestiges of idolatry. He was made apostolic delegate and Archbishop of Mainz, and he built many churches, and administered them either personally or through his disciples. At length, he went back once again to the Frisians, who had lapsed into idolatry, to preach the Gospel to them. There, with Eobanus his fellow bishop and many others, he was killed in a bloody massacre near the River Born and received the crown of martyrdom. His body lieth in the monastery of Fulda.\par
\switchcolumn*

\LA
\TeDeum{Te Deum}
\switchcolumn
\EN
\TeDeum{Te Deum}
\switchcolumn*

\end{paracol}

\par\needspace{10\baselineskip}\addvspace{1.2em}{\centering\small\textsc{Die 6 Junii} · \textsc{6 June}\par}
\Office{S. Norberti Episcopi et Confessoris}{St. Norbert, Bishop and Confessor}{Duplex}
\addcontentsline{toc}{subsection}{Die 6 Junii · S. Norberti Episcopi et Confessoris\enspace\textit{St. Norbert, Bishop and Confessor}}
\begin{paracol}{2}
\LA
\RefLine{Lectiones I–III: de Scriptura occurrente.}
\switchcolumn
\EN
\RefLine{Lessons I–III: from the Scripture of the day.}
\switchcolumn*

\LA
\RefLine{Lectiones VII–IX: ex Commúni unius Confessoris Pontificis (C4).}
\switchcolumn
\EN
\RefLine{Lessons VII–IX: from the Common of a Confessor Bishop (C4).}
\switchcolumn*

\LA
\LessonHead{Lectio IV}
\switchcolumn
\EN
\LessonHead{Lesson IV}
\switchcolumn*

\LA
\noindent Norbértus, nobilíssimis paréntibus natus, adoléscens liberálibus disciplínis erudítus, in ipsa póstea imperatóris aula, spretis mundi illécebris, ecclesiásticæ milítiæ adscríbi vóluit. Sacris initiátus, rejéctis móllibus ac spléndidis véstibus, pellícea melóte indútus, prædicatióni verbi Dei se totum dedit. Abdicátis ecclesiásticis provéntibus satis amplis, et património in páuperes erogáto, semel in die sub vésperam solo cibo quadragesimáli utens, nudísque pédibus et lácera veste sub brumáli rigóre incédens, miræ austeritátis vitam est aggréssus. Potens ígitur ópere et sermóne, innúmeros hæréticos ad fidem, peccatóres ad pœniténtiam, dissidéntes ad pacem et concórdiam revocávit.\par
\switchcolumn
\EN
\noindent Norbert, born in the year 1080 of parents of the highest rank, thoroughly educated in his youth in worldly knowledge, and a member of the Imperial court, turned his back upon the glory of the world, and chose rather to enlist himself as a soldier of the Church. Being ordained Priest, he laid aside all soft and showy raiment, clad himself in a coat of skins, and made the preaching of the Word of God the one object of his life. He had the right to rich revenues of the Church but these he renounced and to an ample fortune from his father; but this he gave to the poor. He ate only once a day, and that in the evening, and then his meal was of the fare of Lent. His life was one of singular hardness, and he was used even in the depth of winter to go out with bare feet and ragged garments. Hence came that mighty power of his words and deeds, whereby he was enabled to turn countless heretics to the true faith, sinners to repentance, and enemies to peace and brotherly love.\par
\switchcolumn*

\LA
\begin{Resp}\Rsign Invéni David servum meum, óleo sancto meo unxi eum: \Star{} Manus enim mea auxiliábitur ei.\end{Resp}
\begin{Resp}\Vsign Nihil profíciet inimícus in eo, et fílius iniquitátis non nocébit ei.\end{Resp}
\begin{Resp}\Rsign Manus enim mea auxiliábitur ei.\end{Resp}
\switchcolumn
\EN
\begin{Resp}\Rsign I have found David My servant, with My holy oil have I anointed him \Star{} For My hand shall help him.\end{Resp}
\begin{Resp}\Vsign The enemy shall prevail nothing against him, nor the son of wickedness afflict him.\end{Resp}
\begin{Resp}\Rsign For My hand shall help him.\end{Resp}
\switchcolumn*

\LA
\LessonHead{Lectio V}
\switchcolumn
\EN
\LessonHead{Lesson V}
\switchcolumn*

\LA
\noindent Cum Laudúni esset, ab epíscopo rogátus ne a sua diœcési discéderet, desértum in ea locum, qui Præmonstrátus dicebátur, sibi delégit; ibíque, trédecim sóciis aggregátis, Præmonstraténsem órdinem instítuit, divínitus accépta per visum régula a sancto Augustíno. Cum vero ejus fama sanctitátis in dies magis augerétur, ac plúrimi ad eum quotídie discípuli convenírent, idem ordo ab Honório secúndo aliísque summis Pontifícibus confirmátus, ac plúribus ab eo monastériis ædificátis, mirífice propagátus est.\par
\switchcolumn
\EN
\noindent Being one while at Laon, the Bishop besought him not to leave his diocese, and he therefore made choice of a wilderness at the place called Prémontré, whither he withdrew himself with thirteen disciples, and thus founded the Order of the Praemonstratensian Canons, whereof he, by the will of God, received the Rule, in a vision, from St. Augustine. When, however, the fame of his holy life became every day more and more noised abroad, and great numbers sought to become his disciples, and the Order had been approved by Honorius II, and other Popes, many more monasteries were built by him, and the Institute wonderfully extended.\par
\switchcolumn*

\LA
\begin{Resp}\Rsign Pósui adjutórium super poténtem, et exaltávi eléctum de plebe mea: \Star{} Manus enim mea auxiliábitur ei.\end{Resp}
\begin{Resp}\Vsign Invéni David servum meum, óleo sancto meo unxi eum.\end{Resp}
\begin{Resp}\Rsign Manus enim mea auxiliábitur ei.\end{Resp}
\switchcolumn
\EN
\begin{Resp}\Rsign I have laid help upon one that is mighty, and have exalted one chosen out of My people \Star{} For My hand shall help him.\end{Resp}
\begin{Resp}\Vsign I have found David My servant, with My holy oil have I anointed him.\end{Resp}
\begin{Resp}\Rsign For My hand shall help him.\end{Resp}
\switchcolumn*

\LA
\LessonHead{Lectio VI}
\switchcolumn
\EN
\LessonHead{Lesson VI}
\switchcolumn*

\LA
\noindent Antvérpiam accersítus, in ea urbe Tanchelíni nefáriam hæresim profligávit. Prophético spíritu et miráculis cláruit. Archiepíscopus tandem (licet relúctans) Magdeburgénsis creátus, ecclesiásticam disciplínam, præsértim cælibátum, constánter propugnávit. Rhemis in concílio Innocéntium secúndum egrégie adjúvit, et Romam cum áliis epíscopis proféctus, schisma Petri Leónis compréssit. Postrémo vir Dei, méritis et Spíritu Sancto plenus, Magdebúrgi obdormívit in Dómino, anno salútis millésimo centésimo trigésimo quarto, die sexta Júnii.\par
\switchcolumn
\EN
\noindent Being called to Antwerpen, he there gave the death-blow to the shameful heresy of Tanchelm. He was remarkable for the spirit of prophecy and for the gift of miracles. He was created (albeit he would rather not have had it so) Archbishop of Magdeburg, and as such he was a strong upholder of the discipline of the Church, especially contending against the marriage of the clergy. At a Council held at Reims he was a great help to Innocent II, and went with some Other Bishops to Rome, where they stamped out the schism of Peter Leoni. It was at last at Magdeburg that this man of God, full of good works and of the Holy Ghost, fell asleep in the Lord, on the 6th day of June, in the year of salvation 1134.\par
\switchcolumn*

\LA
\begin{Resp}\Rsign Iste est, qui ante Deum magnas virtútes operátus est, et omnis terra doctrína ejus repléta est: \Star{} Ipse intercédat pro peccátis ómnium populórum.\end{Resp}
\begin{Resp}\Vsign Iste est, qui contémpsit vitam mundi, et pervénit ad cæléstia regna.\end{Resp}
\begin{Resp}\Rsign Ipse intercédat pro peccátis ómnium populórum.\end{Resp}
\begin{Resp}\Vsign Glória Patri, et Fílio, \Star{} et Spirítui Sancto.\end{Resp}
\begin{Resp}\Rsign Ipse intercédat pro peccátis ómnium populórum.\end{Resp}
\switchcolumn
\EN
\begin{Resp}\Rsign This is he which wrought great wonders before God, and the whole earth is full of his teaching \Star{} May he pray for all people, that their sins may be forgiven unto them\end{Resp}
\begin{Resp}\Vsign This is he which loved not his life in this world, and hath attained unto the kingdom of heaven.\end{Resp}
\begin{Resp}\Rsign May he pray for all people, that their sins may be forgiven unto them\end{Resp}
\begin{Resp}\Vsign Glory be to the Father, and to the Son, \Star{} and to the Holy Ghost.\end{Resp}
\begin{Resp}\Rsign May he pray for all people, that their sins may be forgiven unto them\end{Resp}
\switchcolumn*

\LA
\LessonHead{Lectio IX (pro commemoratione)}
\switchcolumn
\EN
\LessonHead{Lesson IX (when commemorated)}
\switchcolumn*

\LA
\LessonTitle{Commemoratio: S. Norberti Episcopi et Confessoris}
\switchcolumn
\EN
\LessonTitle{Commemoration of: St. Norbert, Bishop and Confessor}
\switchcolumn*

\LA
\noindent Norbértus, nobilíssimis paréntibus natus, adoléscens liberálibus disciplínis erudítus, in ipsa póstea imperatóris aula, spretis mundi illécebris, ecclesiásticæ milítiæ adscríbi vóluit. Sacris initiátus, prædicatióni verbi Dei se totum dedit. Innúmeros hæréticos ad fidem, peccatóres ad pœniténtiam, dissidéntes ad pacem et concórdiam revocávit. Desértum locum, qui Præmonstrátus dicebátur, in Laudunénsi diœcési sibi delégit; ibíque, trédecim sóciis aggregátis, Præmonstraténsem órdinem instítuit, qui mirífice propagátus est. Archiepíscopus Magdeburgénsis, licet relúctans, creátus, ecclesiásticam disciplínam, cælibátum præsértim, constánter propugnávit. Rhemis in concílio Innocéntium secúndum egrégie adjúvit, et Romam cum áliis epíscopis proféctus, schisma Petri Leónis compréssit. Magdebúrgi obdormívit in Dómino, anno salútis millésimo centésimo trigésimo quarto, die sexta júnii.\par
\switchcolumn
\EN
\noindent Norbert was born of most noble parents and, as a young man, studied the liberal arts. Then, while serving in the court of the Emperor himself, he spurned the seductions of the world and decided to enroll among the soldiers of the Church. After receiving holy orders he devoted himself entirely to preaching the word of God. He brought back innumerable heretics to the faith, sinners to penance, quarrellers to peace and concord. He retired to a desert place called Prémontré in the diocese of Laon; and there, with thirteen companions, he founded the Praemonstratensian Order, which spread in a marvelous way. Against his will, he was made Archbishop of Magdeburg, and constantly defended ecclesiastical discipline and especially celibacy. At the Council of Reims, he was a strong champion of Innocent II; and going to Rome with other bishops, he put an end to the schism of Pierleone. He fell asleep in the Lord at Magdeburg on the 6th day of June in the year of salvation 1134.\par
\switchcolumn*

\LA
\TeDeum{Te Deum}
\switchcolumn
\EN
\TeDeum{Te Deum}
\switchcolumn*

\end{paracol}

\par\needspace{10\baselineskip}\addvspace{1.2em}{\centering\small\textsc{Die 9 Junii} · \textsc{9 June}\par}
\Office{Ss. Primi et Feliciani Martyrum}{Sts. Primus and Felician, Martyrs}{Simplex}
\addcontentsline{toc}{subsection}{Die 9 Junii · Ss. Primi et Feliciani Martyrum\enspace\textit{Sts. Primus and Felician, Martyrs}}
\begin{paracol}{2}
\LA
\RefLine{Lectiones I–II: de Scriptura occurrente.}
\switchcolumn
\EN
\RefLine{Lessons I–II: from the Scripture of the day.}
\switchcolumn*

\LA
\LessonHead{Lectio III (pro commemoratione)}
\switchcolumn
\EN
\LessonHead{Lesson III (when commemorated)}
\switchcolumn*

\LA
\noindent Primus et Felicianus fratres, in persecutióne Diocletiáni et Maximiáni accusáti christianæ religiónis, in víncula conjiciúntur; quibus soluti, inde eripiúntur ab Angelo. Mox ad prætorem adducti, cum christianam fidem acerrime tueréntur, alter ab altero distracti sunt; ac primum varie tentáta est constántia Feliciáni. Sed, cum suasores impietátis se posse quidquam verbis proficere desperárent, affixis stipiti mánibus ejus et pédibus, ipsum sine cibo et potu inde triduum pendentem reliquérunt. Postridie ejus diéi, prætor vocátum ad se Primum sic affátur: Vides quanto sit prudentior, quam tu, frater tuus, qui obsecutus imperatóribus, apud ipsos est honoratus? Quem si tu quoque imitari volueris, particeps eris ejus honoris et gratiæ. Cui Primus: Qui factum sit fratri meo, cognóvi ex Angelo. Utinam, quemádmodum sum cum eo voluntáte conjunctíssimus, sic ab eodem ne martyrio disjungar. Quo dicto excanduit prætor, et ad ceteros cruciatus, quibus Primum affecit, præsénte jam Feliciáno, liquátum igne plumbum in os ejus jussit infundi. Mox utrumque perduci imperat in theatrum, in eosque immitti duos leónis; qui, prostráti ad eórum genua, cápite et cauda ipsis blandiebántur. Ad id spectaculum cum ámplius duódecim míllia hóminum convenissent, quingénti cum suis familiis christianam religiónem suscepérunt. Quibus rebus permotus prætor, eos secúri pércuti jussit.\par
\switchcolumn
\EN
\noindent Primus and Felician were two brothers who were accused of Christianity during the persecution by Diocletian and Maximian, and thrown into irons, which an angel broke, and so freed their limbs. In the presence of the Praetor they most earnestly clave to the profession of their faith, and were immediately parted one from the other. Felician's was the steadfastness which was first tried in diverse ways. They, however, that strove to argue him into sin, when they found that words availed nothing, fastened his hands and feet to a post, and left him to hang there three days without food or drink. On the fourth day the Praetor called Primus before him, and said to him Seest thou how much thy brother is wiser than thou He hath obeyed the Emperors, and they have made him honourable. Thou hast only to follow his example to be made partaker of his honours and favours. Primus answered him: What hath befallen my brother I know, for an angel hath told me. God grant that, seeing I have the same will that he hath, I may not be divided from him in uplifting of testimony. These words raised the wrath of the Praetor, and to the torments which he had already inflicted on Primus, he added this also, that he had boiling lead put into his mouth, compelling his brother Felician to be present and see it done. After that, he had them led into the theatre and two lions let loose upon them, in the presence of about twelve thousand people who were gathered together to see the show. The lions only fawned upon the knees of the Saints, making friends with them with motions of their heads and tails. This exhibition turned five hundred persons and their households to Christ. The Praetor, then, moved beyond all endurance by what had passed, caused Primus and Felician to be beheaded.\par
\switchcolumn*

\LA
\TeDeum{Te Deum}
\switchcolumn
\EN
\TeDeum{Te Deum}
\switchcolumn*

\end{paracol}

\par\needspace{10\baselineskip}\addvspace{1.2em}{\centering\small\textsc{Die 10 Junii} · \textsc{10 June}\par}
\Office{S. Margaritæ Reginæ Viduæ}{St. Margaret, Queen and Widow}{Semiduplex}
\addcontentsline{toc}{subsection}{Die 10 Junii · S. Margaritæ Reginæ Viduæ\enspace\textit{St. Margaret, Queen and Widow}}
\begin{paracol}{2}
\LA
\RefLine{Lectiones I–III: de Scriptura occurrente.}
\switchcolumn
\EN
\RefLine{Lessons I–III: from the Scripture of the day.}
\switchcolumn*

\LA
\RefLine{Lectiones VII–IX: ex Commúni non Virginum non Martyrum (C7a).}
\switchcolumn
\EN
\RefLine{Lessons VII–IX: from the Common of Widows (C7a).}
\switchcolumn*

\LA
\LessonHead{Lectio IV}
\switchcolumn
\EN
\LessonHead{Lesson IV}
\switchcolumn*

\LA
\noindent Margarita, Scotórum regína, paterno Angliæ regum, materno Cæsárum sánguine claríssima, illustrior adhuc fuit christiana virtúte. Hæc in Hungária nata, ubi pater tunc témporis exsulabat, post exactam summa cum pietáte puerílem ætátem, una cum genitore, qui a sancto Eduardo patruo, Anglórum rege, ad paterni regni fastigium vocabátur, in Angliam venit. Mox, alternante paréntum fortuna, ex Angliæ littore solvens, vi tempestátis expulsa, seu verius divinæ providéntiæ consílio deducta est in oram marítimam Scotiæ. Ibi cum ex matris imperio Malchólmo tertio Scotórum regi, egregiis ejus dótibus capto, nupsísset, sanctimoniæ ac pietátis opéribus, trigínta quibus regnávit annis, toti regno mirifice prófuit.\par
\switchcolumn
\EN
\noindent Margaret, Queen of Scots, was most noble by birth, uniting in herself, from her father the blood of the Kings of England and from her mother the blood of the Caesars, but her greatest nobleness was in her brave Christian life. She was born in Hungary, where her father was then an exile, (in the year 1046,) and had passed a religious childhood, when her uncle Edward, the holy King of England, recalled him to his own royal home, and she came to England with him (in 1054.) A few years after, upon the ruin of her family, she was escaping from England by sea, when the violence of the weather, or, to speak more truly, the Providence of God, caused that the ship should take refuge upon the coast of Scotland. There her extraordinary graces of mind and body so attracted King Malcolm III., that by the advice of his mother, he took her to wife (in 1070,) and of Scotland she deserved exceedingly well for the thirty years of her reign, by the holiness of her life and the abundance of her works of mercy.\par
\switchcolumn*

\LA
\begin{Resp}\Rsign Propter veritátem, et mansuetúdinem, et justítiam: \Star{} Et dedúcet te mirabíliter déxtera tua.\end{Resp}
\begin{Resp}\Vsign Spécie tua et pulchritúdine tua inténde, próspere procéde, et regna.\end{Resp}
\begin{Resp}\Rsign Et dedúcet te mirabíliter déxtera tua.\end{Resp}
\switchcolumn
\EN
\begin{Resp}\Rsign Because of truth, and meekness, and righteousness; \Star{} And thy right hand shall lead thee wonderfully.\end{Resp}
\begin{Resp}\Vsign In thy comeliness, and thy beauty, go forward, fare prosperously, and reign.\end{Resp}
\begin{Resp}\Rsign And thy right hand shall lead thee wonderfully.\end{Resp}
\switchcolumn*

\LA
\LessonHead{Lectio V}
\switchcolumn
\EN
\LessonHead{Lesson V}
\switchcolumn*

\LA
\noindent Inter regales delicias corpus afflictatiónibus ac vigiliis mácerans, magnam noctis partem piis precatiónibus extrahebat. Præter alia jejunia, quæ idéntidem usurpabat, íntegros quadragínta dies ante Natalícia festa tanta cum severitáte jejunare consuevit, ut ne in gravíssimis quidem dolóribus intermiserit. Divino cultui addictíssima, templa pluria et cœnobia partim ex íntegro excitávit, partim resarcívit, et sacra supelléctili ac largo censu ditávit. Regem cónjugem ad meliórem frugem et ad simília suis exercitatiónibus ópera saluberrimo exemplo traduxit, liberosque omnes tam sancte et feliciter educávit, ut eórum plerique, quemádmodum et Agatha mater, et Christina soror, sanctíssimum vitæ genus ampléxi sint. Universi demum regni felicitáti cónsulens, a vitiis ómnibus, quæ furtim irrepserant, pópulos expurgávit, eisque mores christiana pietáte dignos restituit.\par
\switchcolumn
\EN
\noindent In the midst of kingly dainties, she afflicted her body with hardships and watching, using to spend great part of the night in earnest prayer. Besides other fasts which she imposed upon herself, it was her custom to observe one of forty days before Christmas, concerning which fast she was so rigid, that she would not relax it even under sharp suffering. She took great delight in the public worship of God, and founded or renewed a great number of Churches and convents, which she enriched at great cost with sacred furniture. Her healthy example drew the King her husband to habits of sobriety, and to imitate her in her good works. To all her children she had the happiness of giving a godly education, and several of them, like her mother Agatha and her sister Christina, led notable holy lives. The happiness of the whole kingdom was the object for which she constantly strove, and she successfully rooted out all the vices which had stealthily crept in, and established among the people a standard of living worthy of Christians.\par
\switchcolumn*

\LA
\begin{Resp}\Rsign Dilexísti justítiam, et odísti iniquitátem: \Star{} Proptérea unxit te Deus, Deus tuus, óleo lætítiæ.\end{Resp}
\begin{Resp}\Vsign Propter veritátem, et mansuetúdinem, et justítiam.\end{Resp}
\begin{Resp}\Rsign Proptérea unxit te Deus, Deus tuus, óleo lætítiæ.\end{Resp}
\switchcolumn
\EN
\begin{Resp}\Rsign Thou hast loved righteousness, and hated iniquity; \Star{} Therefore God, thy God, hath anointed thee with the oil of gladness.\end{Resp}
\begin{Resp}\Vsign Because of truth, and meekness, and righteousness.\end{Resp}
\begin{Resp}\Rsign Therefore God, thy God, hath anointed thee with the oil of gladness.\end{Resp}
\switchcolumn*

\LA
\LessonHead{Lectio VI}
\switchcolumn
\EN
\LessonHead{Lesson VI}
\switchcolumn*

\LA
\noindent Nihil tamen æque in illa mirábile fuit ac flagrantíssima caritas erga próximos, præsertim egénos, quorum numerosis grégibus non modo stipem áffatim suppeditare, verum étiam trecentis quotídie materna benignitate dapes præbére, flexis genibus in morem ancíllæ ministrare, regiis mánibus pedes ablúere, et pressis étiam ósculis úlcera fovére, solemne hábuit. His porro aliisque piis sumptibus non regias tantum vestes et preiosa monília distraxit, sed ipsum non semel exhausit ærárium. Tolerátis demum ad patiéntiæ miraculum acerbíssimis dolóribus, ánimam semestri corporis ægrotatióne purgatam Auctori suo, sextodecimo Kalendas Decembris, réddidit. Quo témporis momento fácies ejus, diuturni morbi macie ac pallore fœdata, insólita quadam venustáte reflóruit. Miris étiam post mortem prodigiis clara, et Clementis décimi auctoritate in Scotiæ patronam accepta, ubíque terrárum religiosíssime cólitur.\par
\switchcolumn
\EN
\noindent The most remarkable feature of her life was the tenderness of her charity toward her neighbour, especially the needy. Of these she would not only order whole flocks to be relieved, but was accustomed to give dinner to three hundred of them every day, treating them with the tenderness of a mother, and waiting upon them on her knees like a maidservant. She held it one of the privileges of her rank to wash their feet with her own Royal hands, and to dress their sores, which latter she would even kiss. To meet the expenses of her charities she sold not only her queenly raiment and her precious jewels, but more than once exhausted her funds entirely. Purified by grievous suffering, which she bore with marvelous patience during an illness of six months, she resigned her soul into the hands of Him Who had created it, upon the 11th day of June, (1093.) At the moment of death, the bystanders saw her poor worn face, pale and disfigured by continual suffering, flush again with a beauty to which it had long been unused. After her death she became illustrious on account of great signs and wonders. With the approval of Clement X., she was chosen Patroness of Scotland, and her memory is held in profound reverence throughout the whole earth.\par
\switchcolumn*

\LA
\begin{Resp}\Rsign Fallax grátia, et vana est pulchritúdo: \Star{} Múlier timens Dóminum, ipsa laudábitur.\end{Resp}
\begin{Resp}\Vsign Date ei de fructu mánuum suárum, et laudent eam in portis ópera ejus.\end{Resp}
\begin{Resp}\Rsign Múlier timens Dóminum, ipsa laudábitur.\end{Resp}
\begin{Resp}\Vsign Glória Patri, et Fílio, \Star{} et Spirítui Sancto.\end{Resp}
\begin{Resp}\Rsign Múlier timens Dóminum, ipsa laudábitur.\end{Resp}
\switchcolumn
\EN
\begin{Resp}\Rsign Favour is deceitful, and beauty is vain \Star{} A woman that feareth God she shall be praised.\end{Resp}
\begin{Resp}\Vsign Give her of the fruit of her hands, and let her own works praise her in the gates.\end{Resp}
\begin{Resp}\Rsign A woman that feareth God she shall be praised.\end{Resp}
\begin{Resp}\Vsign Glory be to the Father, and to the Son, \Star{} and to the Holy Ghost.\end{Resp}
\begin{Resp}\Rsign A woman that feareth God she shall be praised.\end{Resp}
\switchcolumn*

\LA
\LessonHead{Lectio IX (pro commemoratione)}
\switchcolumn
\EN
\LessonHead{Lesson IX (when commemorated)}
\switchcolumn*

\LA
\LessonTitle{Commemoratio: S. Margaritæ Reginæ Viduæ}
\switchcolumn
\EN
\LessonTitle{Commemoration of: St. Margaret, Queen and Widow}
\switchcolumn*

\LA
\noindent Margaríta, ex régia Anglórum stirpe in Hungária nata, post exáctam summa cum pietáte puerítiam, una cum genitóre, qui a sancto Eduárdo pátruo, Anglórum rege, ad patérni regni fastígium vocabátur, in Angliam, dein in Scótiam venit. Ibi, cum ex matris império Scotórum regi Malchólmo tértio nupsísset, sanctimóniæ et pietátis opéribus annis trigínta toti regno prófuit. Máxima erat in ea vitæ austéritas et flagrantíssimum erga próximos caritátis stúdium, præsértim in egénos; pro quibus aléndis non semel exháusit ærárium. Demum acerbis dolóribus et diutúrno morbo patientíssime tolerátis, ánimam Deo réddidit sextodécimo caléndas decémbris. Quo témporis moménto fácies ejus, mácie ac pallóre fœdáta, insólita quadam venustáte reflóruit. Cleméntis décimi auctoritáte in Scótiæ patrónam accépta, ubíque terrárum religiosíssime cólitur.\par
\switchcolumn
\EN
\noindent Margaret, of the royal house of England, was born in Hungary and spent her childhood there as an unusually devout and pious girl. When her father was called to high office in his own country by his uncle, St. Edward, King of England, she went to England and then to Scotland. There, upon instructions from her mother, she married King Malcolm III. The country was blessed by her holy life and by her deeds of charity for the next thirty years. The austerity of her life was exceedingly great, and her charity towards her neighbour most ardent and zealous, especially for those in need, for whom she not infrequently exhausted the treasury. At length, having most patiently endured bitter sorrows and long illness, she rendered her soul to God on the 16th day of November. At the moment of her death her features, emaciated and pale, bloomed again with unusual beauty. On the authority of Clement X she was chosen Patroness of Scotland, and is honoured with great devotion throughout the world.\par
\switchcolumn*

\LA
\TeDeum{Te Deum}
\switchcolumn
\EN
\TeDeum{Te Deum}
\switchcolumn*

\end{paracol}

\par\needspace{10\baselineskip}\addvspace{1.2em}{\centering\small\textsc{Die 11 Junii} · \textsc{11 June}\par}
\Office{S. Barnabæ Apostoli}{St. Barnabas the Apostle}{Duplex majus}
\addcontentsline{toc}{subsection}{Die 11 Junii · S. Barnabæ Apostoli\enspace\textit{St. Barnabas the Apostle}}
\begin{paracol}{2}
\LA
\LessonHead{Lectio I}
\switchcolumn
\EN
\LessonHead{Lesson I}
\switchcolumn*

\LA
\LessonTitle{De Actibus Apostolórum}
\switchcolumn
\EN
\LessonTitle{From the Acts of the Apostles}
\switchcolumn*

\LA
\Cite{Acts 13:43-47}
\switchcolumn
\EN
\Cite{Acts 13:43-47}
\switchcolumn*

\LA
\noindent Cum dimíssa esset synagóga, secúti sunt multi Judæórum et coléntium advenárum, Paulum et Bárnabam; qui loquéntes suadébant eis ut permanérent in grátia Dei. Sequénti vero sábbato pene univérsa cívitas convénit audíre verbum Dei. Vidéntes autem turbas Judǽi, repléti sunt zelo, et contradicébant his, quæ a Paulo dicebántur, blasphemántes. Tunc constánter Paulus et Bárnabas dixérunt: Vobis oportébat primum loqui verbum Dei; sed, quóniam repéllitis illud, et indígnos vos judicátis ætérnæ vitæ, ecce convértimur ad gentes; sic enim præcépit nobis Dóminus: Pósui te in lucem Géntium, ut sis in salútem usque ad extrémum terræ.\par
\switchcolumn
\EN
\noindent And when the synagogue was broken up, many of the Jews, and of the strangers who served God, followed Paul and Barnabas: who speaking to them, persuaded them to continue in the grace of God. But the next sabbath day, the whole city almost came together, to hear the word of God. And the Jews seeing the multitudes, were filled with envy, and contradicted those things which were said by Paul, blaspheming. Then Paul and Barnabas said boldly: To you it behoved us first to speak the word of God: but because you reject it, and judge yourselves unworthy of eternal life, behold we turn to the Gentiles. For so the Lord hath commanded us: I have set thee to be the light of the Gentiles; that thou mayest be for salvation unto the utmost part of the earth.\par
\switchcolumn*

\LA
\begin{Resp}\Rsign Ecce ego mitto vos sicut oves in médio lupórum, dicit Dóminus: \Star{} Estóte ergo prudéntes sicut serpéntes, et símplices sicut colúmbæ.\end{Resp}
\begin{Resp}\Vsign Dum lucem habétis, crédite in lucem, ut fílii lucis sitis.\end{Resp}
\begin{Resp}\Rsign Estóte ergo prudéntes sicut serpéntes, et símplices sicut colúmbæ.\end{Resp}
\switchcolumn
\EN
\begin{Resp}\Rsign Behold, I send you as sheep in the midst of wolves, said the Lord: \Star{} Be ye, therefore, wise as serpents and simple as doves.\end{Resp}
\begin{Resp}\Vsign Whilst you have the light, believe in the light, that you may be the children of light.\end{Resp}
\begin{Resp}\Rsign Be ye therefore wise as serpents and simple as doves.\end{Resp}
\switchcolumn*

\LA
\LessonHead{Lectio II}
\switchcolumn
\EN
\LessonHead{Lesson II}
\switchcolumn*

\LA
\Cite{Act 13:48-52}
\switchcolumn
\EN
\Cite{Acts 13:48-52}
\switchcolumn*

\LA
\noindent Audiéntes autem gentes, gavísæ sunt, et glorificábant verbum Dómini: et credidérunt quotquot erant præordináti ad vitam ætérnam. Disseminabátur autem verbum Dómini per univérsam regiónem. Judǽi autem concitavérunt mulíeres religiósas et honéstas, et primos civitátis, et excitavérunt persecutiónem in Paulum et Bárnabam: et ejecérunt eos de fínibus suis. At illi, excússo púlvere pedum in eos, venérunt Icónium. Discípuli quoque replebántur gáudio, et Spíritu Sancto.\par
\switchcolumn
\EN
\noindent And the Gentiles hearing it, were glad, and glorified the word of the Lord: and as many as were ordained to life everlasting, believed. And the word of the Lord was published throughout the whole country. But the Jews stirred up religious and honourable women, and the chief men of the city, and raised persecution against Paul and Barnabas: and cast them out of their coasts. But they, shaking off the dust of their feet against them, came to Iconium. And the disciples were filled with joy and with the Holy Ghost.\par
\switchcolumn*

\LA
\begin{Resp}\Rsign Tóllite jugum meum super vos, dicit Dóminus, et díscite a me, quia mitis sum et húmilis corde: \Star{} Jugum enim meum suáve est, et onus meum leve.\end{Resp}
\begin{Resp}\Vsign Et inveniétis réquiem animábus vestris.\end{Resp}
\begin{Resp}\Rsign Jugum enim meum suáve est, et onus meum leve.\end{Resp}
\switchcolumn
\EN
\begin{Resp}\Rsign Take up my yoke upon you, said the Lord, and learn of me, because I am meek, and humble of heart. \Star{} For my yoke is sweet and my burden light.\end{Resp}
\begin{Resp}\Vsign And you shall find rest to your souls.\end{Resp}
\begin{Resp}\Rsign For my yoke is sweet and my burden light.\end{Resp}
\switchcolumn*

\LA
\LessonHead{Lectio III}
\switchcolumn
\EN
\LessonHead{Lesson III}
\switchcolumn*

\LA
\Cite{Act 14:1-3}
\switchcolumn
\EN
\Cite{Acts 14:1-3}
\switchcolumn*

\LA
\noindent Factum est autem Icónii, ut simul introírent in synagógam Judæórum, et loqueréntur, ita ut créderet Judæórum et Græcórum copiósa multitúdo. Qui vero incréduli fuérunt Judǽi, suscitavérunt et ad iracúndiam concitavérunt ánimas Géntium advérsus fratres. Multo ígitur témpore demoráti sunt, fiduciáliter agéntes in Dómino, testimónium perhibénte verbo grátiæ suæ, dante signa et prodígia fíeri per manus eórum.\par
\switchcolumn
\EN
\noindent And it came to pass in Iconium, that they entered together into the synagogue of the Jews, and so spoke that a very great multitude both of the Jews and of the Greeks did believe. But the unbelieving Jews stirred up and incensed the minds of the Gentiles against the brethren. A long time therefore they abode there, dealing confidently in the Lord, who gave testimony to the word of his grace, granting signs and wonders to be done by their hands.\par
\switchcolumn*

\LA
\begin{Resp}\Rsign Dum stetéritis ante reges et prǽsides, nolíte cogitáre quómodo aut quid loquámini: \Star{} Dábitur enim vobis in illa hora, quid loquámini.\end{Resp}
\begin{Resp}\Vsign Non enim vos estis qui loquímini: sed Spíritus Patris vestri, qui lóquitur in vobis.\end{Resp}
\begin{Resp}\Rsign Dábitur enim vobis in illa hora, quid loquámini.\end{Resp}
\begin{Resp}\Vsign Glória Patri, et Fílio, \Star{} et Spirítui Sancto.\end{Resp}
\begin{Resp}\Rsign Dábitur enim vobis in illa hora, quid loquámini.\end{Resp}
\switchcolumn
\EN
\begin{Resp}\Rsign But when they shall deliver you up to the judges, take no thought how or what to speak: \Star{} for it shall be given you in that hour what to speak:\end{Resp}
\begin{Resp}\Vsign For it is not you that speak, but the spirit of your Father that speaketh in you.\end{Resp}
\begin{Resp}\Rsign For it shall be given you in that hour what to speak.\end{Resp}
\begin{Resp}\Vsign Glory be to the Father, and to the Son, \Star{} and to the Holy Ghost.\end{Resp}
\begin{Resp}\Rsign For it shall be given you in that hour what to speak.\end{Resp}
\switchcolumn*

\LA
\LessonHead{Lectio IV}
\switchcolumn
\EN
\LessonHead{Lesson IV}
\switchcolumn*

\LA
\noindent Bárnabas Levítes, Cýprius génere, qui et Joseph, cum Paulo Géntium Apóstolus ordinátus est ad prædicándum Jesu Christi Evangélium. Is, agro véndito quem habébat, redáctam ex eo pecúniam áttulit Apóstolis. Missus autem Antiochíam prædicatiónis causa, cum ibi multos ad Christi Dómini fidem convérsos esse comperísset, incredibíliter lætátus, eos hortabátur ut in Christi fide permanérent. Qua cohortatióne multum proficiébat, quod ab ómnibus vir bonus et Spíritu Sancto plenus habebátur.\par
\switchcolumn
\EN
\noindent Joseph, who by the Apostles was surnamed Barnabas, (which is, being interpreted, the Son of Consolation,) a Levite and of the country of Cyprus, having land, sold it, and brought the money, and laid it at the Apostles' feet. (Acts iv. 36, 37.) When Paul, after his conversion, was come to Jerusalem, the disciples were all afraid of him, but Barnabas took him, and brought him to the Apostles, (ix. 26, 27,) When tidings that a great number believed and turned unto the Lord at Antioch came unto the ears of the Church which was at Jerusalem, they sent forth Barnabas that he should go as far as Antioch. Who, when he came, and had seen the grace of God, was glad, and exhorted them all that with purpose of heart they would cleave unto the Lord. For he was a good man, and full of the Holy Ghost, and of faith, and much people was added unto the Lord. (xi. 21-24.)\par
\switchcolumn*

\LA
\begin{Resp}\Rsign Vidi conjúnctos viros, habéntes spléndidas vestes, et Ángelus Dómini locútus est ad me, dicens: \Star{} Isti sunt viri sancti facti amíci Dei.\end{Resp}
\begin{Resp}\Vsign Vidi Ángelum Dei fortem, volántem per médium cælum, voce magna clamántem et dicéntem.\end{Resp}
\begin{Resp}\Rsign Isti sunt viri sancti facti amíci Dei.\end{Resp}
\switchcolumn
\EN
\begin{Resp}\Rsign I saw men standing together, clad in shining raiment, and the Angel of the Lord spoke unto me, saying: \Star{} These men are holy, for they are the friends of God.\end{Resp}
\begin{Resp}\Vsign I saw a strong Angel of God fly into the midst of heaven, saying with a loud voice:\end{Resp}
\begin{Resp}\Rsign These men are holy, for they are the friends of God.\end{Resp}
\switchcolumn*

\LA
\LessonHead{Lectio V}
\switchcolumn
\EN
\LessonHead{Lesson V}
\switchcolumn*

\LA
\noindent Proféctus inde Tarsum ut quæreret Paulum, cum eo Antiochíam venit. In ejus urbis ecclésia annum commoráti, christiánæ fídei et vitæ illis homínibus præcépta dedérunt: ubi étiam Jesu Christi cultóres primum Christiáni sunt appelláti. Discípuli autem Pauli et Bárnabæ, suis facultátibus Christiános, qui in Judǽa erant, sustentábant, eo mitténtes pecúniam per Paulum et Bárnabam. Qui perfúncti illo caritátis offício, adhíbito Joánne, cui cognómen erat Marcus, rediérunt Antiochíam.\par
\switchcolumn
\EN
\noindent Then departed Barnabas to Tarsus for to seek Paul, and, when he had found him, he brought him unto Antioch. And it came to pass that a whole year they assembled themselves with the Church, and taught much people. And the disciples were called Christians first in Antioch. And in these days came prophets from Jerusalem unto Antioch. And there stood up one of them, named Agabus, and signified, by the Spirit, that there should be great dearth throughout all the world which came to pass in the days of Claudius Caesar. Then the disciples, every man according to his ability, determined to send relief unto the brethren which dwelt in Judea, which also they did, and sent it to the elders by the hands of Barnabas and Paul. (xi. 25-30.) And Barnabas and Paul returned from Jerusalem, when they had fulfilled their ministry, and took with them John, whose surname was Mark. (xii. 25.)\par
\switchcolumn*

\LA
\begin{Resp}\Rsign Beáti estis, cum maledíxerint vobis hómines, et persecúti vos fúerint, et díxerint omne malum advérsum vos, mentiéntes, propter me: \Star{} Gaudéte et exsultáte, quóniam merces vestra copiósa est in cælis.\end{Resp}
\begin{Resp}\Vsign Cum vos óderint hómines, et cum separáverint vos, et exprobráverint, et ejécerint nomen vestrum tamquam malum propter Fílium hóminis.\end{Resp}
\begin{Resp}\Rsign Gaudéte et exsultáte, quóniam merces vestra copiósa est in cælis.\end{Resp}
\switchcolumn
\EN
\begin{Resp}\Rsign Blessed are ye when men shall revile you, and persecute you, and shall say all manner of evil against you falsely, for My sake; \Star{} Rejoice, and be exceeding glad, for great is your reward in heaven.\end{Resp}
\begin{Resp}\Vsign When men shall hate you, and when they shall separate you from their company, and shall reproach you, and cast out your name as evil, for the Son of Man's sake.\end{Resp}
\begin{Resp}\Rsign Rejoice, and be exceeding glad, for great is your reward in heaven.\end{Resp}
\switchcolumn*

\LA
\LessonHead{Lectio VI}
\switchcolumn
\EN
\LessonHead{Lesson VI}
\switchcolumn*

\LA
\noindent Cum autem Antiochíæ in Ecclésia, cum céteris prophétis et doctóribus, Paulus et Bárnabas in jejúnio et oratióne Dómino deservírent, dixit Spíritus Sanctus: Segregáte mihi Saulum et Bárnabam in opus, ad quod assúmpsi eos. Tunc jejunántes et orántes, imponentésque eis manus, dimisérunt illos. Itaque Seleucíam venérunt, inde in Cyprum; ac multas prætérea urbes regionésque, prædicántes Evangélium summa cum audéntium utilitáte, peragrárunt. Postrémo Bárnabas digréssus a Paulo, una cum Joánne, qui cognominátus est Marcus, navigávit in Cyprum; ibíque círciter séptimum Nerónis annum, tértio Idus Júnii ad apostólici múneris laudem, martyrii corónam adjúnxit. Ejus corpus, Zenóne imperatóre, repértum est in ínsula Cypro; ad cujus pectus erat Evangélium Matthæi, Bárnabæ manu conscríptum.\par
\switchcolumn
\EN
\noindent Now there were in the Church that was at Antioch, certain Prophets and teachers and, as Paul and Barnabas, together with them, ministered to the Lord and fasted, the Holy Ghost said Separate Me Barnabas and Saul for the work whereunto I have called them. And when they had fasted and prayed, and laid their hands on them, they sent them away. So they, being sent forth by the Holy Ghost, departed unto Seleucia and from thence they sailed to Cyprus (xiii. 1-4) in the which island, and in many other cities and countries, they journeyed about, preaching the Gospel with great gain to them that heard them. Nevertheless, at last, Paul and Barnabas departed asunder one from the other. And so Barnabas took Mark and sailed unto Cyprus, (xv. 39,) once more. And there it was that upon a certain nth of June, in or about the seventh year of the reign of Nero, Barnabas crowned the dignity of the Apostolate with the glory of martyrdom. During the reign of the Emperor Zeno, his body was found in its grave in Cyprus on his breast lay a copy of the Gospel according to Matthew, written by the hand of Barnabas himself.\par
\switchcolumn*

\LA
\begin{Resp}\Rsign Isti sunt triumphatóres et amíci Dei, qui contemnéntes jussa príncipum, meruérunt prǽmia ætérna: \Star{} Modo coronántur, et accípiunt palmam.\end{Resp}
\begin{Resp}\Vsign Isti sunt qui venérunt ex magna tribulatióne, et lavérunt stolas suas in sánguine Agni.\end{Resp}
\begin{Resp}\Rsign Modo coronántur, et accípiunt palmam.\end{Resp}
\begin{Resp}\Vsign Glória Patri, et Fílio, \Star{} et Spirítui Sancto.\end{Resp}
\begin{Resp}\Rsign Modo coronántur, et accípiunt palmam.\end{Resp}
\switchcolumn
\EN
\begin{Resp}\Rsign These are they which have conquered, and are become the friends of God, who recked not of the commandments of princes, and earned the everlasting reward. \Star{} And now have they crowns on their heads, and palms in their hands.\end{Resp}
\begin{Resp}\Vsign These are they which came out of great tribulation, and have washed their robes in the blood of the Lamb.\end{Resp}
\begin{Resp}\Rsign And now have they crowns on their heads, and palms in their hands.\end{Resp}
\begin{Resp}\Vsign Glory be to the Father, and to the Son, \Star{} and to the Holy Ghost.\end{Resp}
\begin{Resp}\Rsign And now have they crowns on their heads, and palms in their hands.\end{Resp}
\switchcolumn*

\LA
\LessonHead{Lectio VII}
\switchcolumn
\EN
\LessonHead{Lesson VII}
\switchcolumn*

\LA
\LessonTitle{Léctio sancti Evangélii secúndum Matthǽum.}
\switchcolumn
\EN
\LessonTitle{From the Holy Gospel according to Matthew}
\switchcolumn*

\LA
\Cite{Matt 10:16-22}
\switchcolumn
\EN
\Cite{Matt 10:6-22}
\switchcolumn*

\LA
\noindent In illo témpore: Dixit Jesus discípulis suis: Ecce ego mitto vos sicut oves in médio lupórum. Et réliqua.\par
\switchcolumn
\EN
\noindent At that time: Jesus said unto His disciples Behold, I send you forth as sheep in the midst of wolves. And so on.\par
\switchcolumn*

\LA
\LessonTitle{Homilía sancti Joánnis Chrysóstomi.}
\switchcolumn
\EN
\LessonTitle{Homily by St. John Chrysostom, Patriarch of Constantinople}
\switchcolumn*

\LA
\Source{Homilia 34 in Matthæum, post initium}
\switchcolumn
\EN
\Source{34/A on Matthew.}
\switchcolumn*

\LA
\noindent Cum Dóminus omnem sollicitúdinem a discipulórum córdibus ejécerit, et ostensióne signórum armáverit, atque, ab ómnibus negótiis sæculáribus alienátos et ab omni temporálium rerum cura liberátos, férreos quodámmodo atque adamántinos fécerit, tum dénique eventúra illis advérsa prædicit. Multa enim ex hac prædicatióne futurárum rerum cómmoda consequebántur. Primum, ut ejus præsciéntiæ vim edíscerent. Deínde, ut nemo suspicarétur, ex Magístri infirmitáte tam grávia mala descéndere. Prætérea, ne, qui ea passúri erant, súbito ac inopináto rerum evéntu perturbaréntur. Dénique, ne, cum ista sub ipsum passiónis suæ tempus audírent, nímium commoveréntur.\par
\switchcolumn
\EN
\noindent When the Lord had cleared the minds of His disciples of all care, and had armed them by showing forth His mighty works, had estranged them from all business of this world, and freed them from all anxiety touching the things of time, moulding them into a frame of iron-like, nay, diamondlike, hardness, then at length He told them of the contendings against the which they were afterward to wrestle. By this foretelling of things to come they were much holpen. First, they learnt the power of His fore-knowledge. Then, they were guarded against all suspicion that these great sorrows flowed from faultiness in their Master. Again, the future sufferers were made safe from all trouble of being taken unawares. Lastly, seeing that they heard these things at a time nigh to His own suffering, they were not over troubled.\par
\switchcolumn*

\LA
\begin{Resp}\Rsign Isti sunt qui vivéntes in carne, plantavérunt Ecclésiam sánguine suo: \Star{} Cálicem Dómini bibérunt, et amíci Dei facti sunt.\end{Resp}
\begin{Resp}\Vsign In omnem terram exívit sonus eórum, et in fines orbis terræ verba eórum.\end{Resp}
\begin{Resp}\Rsign Cálicem Dómini bibérunt, et amíci Dei facti sunt.\end{Resp}
\switchcolumn
\EN
\begin{Resp}\Rsign These are they who while yet they lived in the flesh, planted the Church in their own blood; \Star{} They drank of the Lord's cup, and became the friends of God.\end{Resp}
\begin{Resp}\Vsign Their sound is gone out through all the earth, and their words to the ends of the world.\end{Resp}
\begin{Resp}\Rsign They drank of the Lord's cup, and became the friends of God.\end{Resp}
\switchcolumn*

\LA
\LessonHead{Lectio VIII}
\switchcolumn
\EN
\LessonHead{Lesson VIII}
\switchcolumn*

\LA
\noindent Jam vero, ut intélligant novum hoc esse belli genus et insólitum præliándi morem, cum illos nudos mítteret, una indútos túnica, sine cálceis, absque virga et absque zona et pera, et ab excipiéntibus ali jubéret; non fecit hic dicéndi finem, sed inexplicábilem virtútem suam próferens, Etiam sic eúntes, inquit, mansuetúdinem tamen óvium osténdite, quamvis ad lupos itúri, nec simplíciter ad lupos, sed étiam in médio lupórum: (neque vero óvium tantum mansuetúdinem habére jubet, sed étiam colúmbæ simplicitátem); sic enim virtútem meam máxime osténdam, cum ab óvibus lupi superabúntur; et quamvis illæ sint in médio lupórum, et innúmeris móribus laceréntur, non modo non consúmptæ fúerint, verum étiam illos in sui natúram transmutáverint.\par
\switchcolumn
\EN
\noindent And now, that they may understand how that this is a new kind of warfare, and an unaccustomed manner of contending, when He sendeth them forth unarmed, providing neither gold, nor silver, nor brass in their purses nor scrip for their journey, neither two coats, neither shoes, nor yet staves, (x. 9, 10,) left to the hospitality of whosoever would receive them, He maketh not here an end to His discourse, but, in manifestation of His unspeakable power, He biddeth them, so going, to show forth the meekness of sheep, seeing they were about going unto wolves neither simply unto wolves, but in the very midst of wolves. Neither is it only the meekness of sheep which He biddeth them have, but also the harmlessness of doves, that He might so much the more gloriously display His power, when the sheep overcame the wolves. These are the sheep which albeit they abide in the midst of wolves, and are mangled by many a bite, not only are not destroyed, but do gradually make the wolves change their nature, and become sheep themselves.\par
\switchcolumn*

\LA
\begin{Resp}\Rsign Isti sunt viri sancti, quos elégit Dóminus in caritáte non ficta, et dedit illis glóriam sempitérnam: \Star{} Quorum doctrína fulget Ecclésia, ut sole luna.\end{Resp}
\begin{Resp}\Vsign Sancti per fidem vicérunt regna: operáti sunt justítiam.\end{Resp}
\begin{Resp}\Rsign Quorum doctrína fulget Ecclésia, ut sole luna.\end{Resp}
\begin{Resp}\Vsign Glória Patri, et Fílio, \Star{} et Spirítui Sancto.\end{Resp}
\begin{Resp}\Rsign Quorum doctrína fulget Ecclésia, ut sole luna.\end{Resp}
\switchcolumn
\EN
\begin{Resp}\Rsign These men are saints, whom the Lord hath chosen in love unfeigned, and hath given them glory everlasting. These are they \Star{} By the light of whose teaching the Church is glorified, even as the moon is glorified by the light of the sun.\end{Resp}
\begin{Resp}\Vsign The saints through faith subdued kingdoms, wrought righteousness.\end{Resp}
\begin{Resp}\Rsign By the light of whose teaching the Church is glorified, even as the moon is glorified by the light of the sun.\end{Resp}
\begin{Resp}\Vsign Glory be to the Father, and to the Son, \Star{} and to the Holy Ghost.\end{Resp}
\begin{Resp}\Rsign By the light of whose teaching the Church is glorified, even as the moon is glorified by the light of the sun.\end{Resp}
\switchcolumn*

\LA
\LessonHead{Lectio IX}
\switchcolumn
\EN
\LessonHead{Lesson IX}
\switchcolumn*

\LA
\noindent Majus certe atque admirabílius est mentem adversariórum commutáre, et ánimum in divérsum transférre, quam illos occídere; præsértim cum duódecim tantum essent, et lupis plenus esset orbis univérsus. Erubescámus ígitur, qui, longe divérsa faciéntes, tamquam lupi in adversários rúimus. Nam, quámdiu oves fuérimus, víncimur; tunc enim a nobis pastóris auxílium recédit, qui non lupos, sed oves pascit.\par
\switchcolumn
\EN
\noindent Beyond all doubt it is a greater and more marvellous thing to change the minds of enemies, and to turn their thoughts round, than to kill them more especially when the work is to be done by only twelve sheep, and the whole world is full of the wolves. Shame then upon us, whose deeds are so contrary, and who rather run like wolves upon our enemies. For so long as we are sheep we conquer, yea, though a thousand wolves be gathered round about us, we overcome, and are the conquerors but if we become wolves ourselves, then are we conquered. For then doth the Shepherd's help forsake us, Who feedeth not wolves but sheep.\par
\switchcolumn*

\LA
\TeDeum{Te Deum}
\switchcolumn
\EN
\TeDeum{Te Deum}
\switchcolumn*

\LA
\LessonHead{Lectio IX (pro commemoratione)}
\switchcolumn
\EN
\LessonHead{Lesson IX (when commemorated)}
\switchcolumn*

\LA
\LessonTitle{Commemoratio: S. Barnabæ Apostoli}
\switchcolumn
\EN
\LessonTitle{Commemoration of: St. Barnabas the Apostle}
\switchcolumn*

\LA
\noindent Bárnabas Levítes, Cýprius génere, cum Paulo Apóstolo ordinátus est ad prædicándum Jesu Christi Evangélium. Is, agro véndito quem habébat, redáctam ex eo pecúniam áttulit Apóstolis. Missus Antiochíam prædicatiónis causa, multos ibi ad Christi fidem convérsos suis hortatiónibus confirmávit. Inde proféctus cum eódem Paulo, multas urbes regionésque summa cum audiéntium utilitáte peragrávit. Postrémo, digréssus a Paulo, una cum Joánne, qui cognominátus est Marcus, navigávit in Cyprum, ibíque, séptimo Nerónis anno, ad apostólici múneris laudem martýrii corónam adjúnxit.\par
\switchcolumn
\EN
\noindent Barnabas the Levite, of the people of Cyprus, was sent with Paul the Apostle to preach the Gospel of Jesus Christ. He sold a field he owned and brought the money from it to the Apostles. Sent to Antioch to preach, he gave new strength by his exhortations to many who had been converted to belief in Christ. Thence he set forth with St. Paul and travelled through many cities and countries, to the great benefit of his hearers. Finally he left Paul, and with John, surnamed Mark, he took ship to Cyprus. There, in the seventh year of Nero's reign, he added to the dignity of the apostolic office the crown of martyrdom.\par
\switchcolumn*

\LA
\TeDeum{Te Deum}
\switchcolumn
\EN
\TeDeum{Te Deum}
\switchcolumn*

\end{paracol}

\par\needspace{10\baselineskip}\addvspace{1.2em}{\centering\small\textsc{Die 12 Junii} · \textsc{12 June}\par}
\Office{S. Joannis a S. Facundo Confessoris}{St. John of St. Facundus, Confessor}{Duplex}
\addcontentsline{toc}{subsection}{Die 12 Junii · S. Joannis a S. Facundo Confessoris\enspace\textit{St. John of St. Facundus, Confessor}}
\begin{paracol}{2}
\LA
\RefLine{Lectiones I–III: de Scriptura occurrente.}
\switchcolumn
\EN
\RefLine{Lessons I–III: from the Scripture of the day.}
\switchcolumn*

\LA
\RefLine{Lectiones VII–VIII: ex Commúni Confessoris non pontificis (C5).}
\switchcolumn
\EN
\RefLine{Lessons VII–VIII: from the Common of a Confessor not a Bishop (C5).}
\switchcolumn*

\LA
\RefLine{Lectio IX: de Ss. Basilidis, Cyrini, Naboris and Nazarii Martyrum (06-12o).}
\switchcolumn
\EN
\RefLine{Lesson IX: of Sts. Basilides, Cyrinus, Nabor and Nazarius, Martyrs (06-12o).}
\switchcolumn*

\LA
\LessonHead{Lectio IV}
\switchcolumn
\EN
\LessonHead{Lesson IV}
\switchcolumn*

\LA
\noindent Joánnem Sahaguni in Hispania nobili genere natum, parentes cum diu prole caruissent, piis operibus, et orationibus a Deo impetrarunt. Ab ineunte ætate egregium futuræ sanctitatis specimen dedit: nam e loco superiori ad ceteros pueros crebro verba faciebat, quibus eos ad virtutem, et Dei cultum hortabatur, eorumque dissidia componebat. In patria monachis sancti Facundi ordinis sancti Benedicti primis litterarum rudimentis imbuendus traditur. Dum iis operam daret, curavit pater, ut parochus ecclesiam administraret: quod munus juvenis nullis rationibus adduci potuit, ut retineret. Inter familiares episcopi Burgensis adscriptus, ob spectatam ipsius probitatem intimus ei fuit, ab eoque presbyter, et canonicus factus multis beneficiis auctus est. Sed relicta aula episcopi ut Deo quietius serviret, omnibus ecclesiæ proventibus abdicatis se cuidam sacello addixit, ubi Sacrum quotidie faciebat, ac de rebus divinis magna cum auditorum ædificatione frequenter concionabatur.\par
\switchcolumn
\EN
\noindent John Gonzalez was born, the offspring of a noble race, at San Fagondez in Spain, (on Midsummer Day in the year of grace 1430.) His father and mother after long childlessness, obtained him from God by prayers and good works. From his earliest years he gave clear signs of his after holiness of life. He was used to climb up upon an high place to preach to the other little boys, and to exhort them to be good and to worship God, and he made it his work to reconcile their quarrels. While he was still at home he was given in charge to the monks of the Order of Saint Benedict, at the village of San Fagondez, to teach him his first lessons. While he was thus busied, his father obtained for him the benefice of the Parish, but no persuasions could induce him to keep this preferment. He became one of the household of the Bishop of Burgos, and that Prelate, seeing his uprightness, took him into his counsels, ordained him Priest, and made him a Canon, heaping upon him many kindnesses. However, that he might serve God the more quietly, he left the Bishop's Palace, resigned all his Church income, and betook him to a certain Chapel wherein he celebrated the Holy Liturgy every day, and oftentimes preached concerning the things of God, with great profit to all that heard him.\par
\switchcolumn*

\LA
\begin{Resp}\Rsign Honéstum fecit illum Dóminus, et custodívit eum ab inimícis, et a seductóribus tutávit illum: \Star{} Et dedit illi claritátem ætérnam.\end{Resp}
\begin{Resp}\Vsign Justum dedúxit Dóminus per vias rectas, et osténdit illi regnum Dei.\end{Resp}
\begin{Resp}\Rsign Et dedit illi claritátem ætérnam.\end{Resp}
\switchcolumn
\EN
\begin{Resp}\Rsign The Lord made him honourable, and defended him from his enemies, and kept him safe from those that lay in wait for him; \Star{} And gave him perpetual glory.\end{Resp}
\begin{Resp}\Vsign He went down with him into the pit, and left him not in bonds.\end{Resp}
\begin{Resp}\Rsign And gave him perpetual glory.\end{Resp}
\switchcolumn*

\LA
\LessonHead{Lectio V}
\switchcolumn
\EN
\LessonHead{Lesson V}
\switchcolumn*

\LA
\noindent Postea studiorum causa Salmanticam profectus, in celebre collegium divi Bartholomæi cooptatus, sacerdotis munus ita exercuit, ut simul optatis studiis incumberet, et in sacris étiam concionibus assidue versaretur. Cum vero in gravissimum morbum incidisset, arctioris disciplinæ voto se obstrinxit, quod ut redderet, cum prius cuidam pauperi pene nudo ex duabus, quas tantum habebat vestes, meliorem dedisset, ad cœnobium sancti Augustíni severiori disciplina tum maxime florens se contulit: in quo admissus, obedientia, animi demissione, vigiliis ac oratione provectiores anteibat. Triclinii cura cum ipsi demandata esset, vini doliolum ipso attingente omnibus monachis per annum abunde suffecit. Exacto tirocinii anno, præfecti jussu munus concionandi suscepit. Salmanticæ id temporis adeo cruentis factionibus divina, humanaque omnia permixta erant, ut singulis propemodum horis cædes fierent, et omnium ordinum, ac præsertim nobilium sanguine non viæ solum, et fora, sed templa étiam redundarent.\par
\switchcolumn
\EN
\noindent He went later to Salamanca to study, and there being taken into the celebrated College of St. Bartholomew, he did his priestly office, so that he was at once constant to the studies he desired and busy with sermons. Here he had a severe illness, and vowed to take up a sterner way of living. In fulfillment of this vow, he gave to an half-naked beggar the better of the two garments which were all that he had, and then went to a Convent of the friars of St. Augustine, which was then in the richest bloom of rigid discipline. Being admitted therein, he surpassed the most advanced in obedience, lowliness, watchings, and prayer. At the time that he had charge of the table, one keg of wine abundantly sufficed in his hands for all the friars, throughout an whole year. After his year of novitiate, he undertook the duty of preacher at the command of his Superior. At that time, owing to bloody feuds, all things human and divine at Salamanca were in such utter confusion, that murders were committed almost every hour, and the streets and squares, and the very churches, flowed with the blood of all classes, especially of the nobility.\par
\switchcolumn*

\LA
\begin{Resp}\Rsign Amávit eum Dóminus, et ornávit eum: stolam glóriæ índuit eum, \Star{} Et ad portas paradísi coronávit eum.\end{Resp}
\begin{Resp}\Vsign Índuit eum Dóminus lorícam fídei, et ornávit eum.\end{Resp}
\begin{Resp}\Rsign Et ad portas paradísi coronávit eum.\end{Resp}
\switchcolumn
\EN
\begin{Resp}\Rsign The Lord loved him and beautified him; He clothed him with a robe of glory, \Star{} And crowned him at the gates of Paradise.\end{Resp}
\begin{Resp}\Vsign The Lord hath put on him the breast-plate of faith, and hath adorned him.\end{Resp}
\begin{Resp}\Rsign And crowned him at the gates of Paradise.\end{Resp}
\switchcolumn*

\LA
\LessonHead{Lectio VI}
\switchcolumn
\EN
\LessonHead{Lesson VI}
\switchcolumn*

\LA
\noindent At Joannes, tum concionibus, tum privatis colloquiis civium animos demulcens, ad tranquillitatem urbem reduxit. Virum principem graviter offendit, quod illius in subditos sævitiam increpasset. Qua de causa equites duos immisit, qui eum in itinere confoderent, jamque ad ipsum propinquaverant: cum, stupore divinitus immisso, simul cum equis immobiles steterunt, donec ad pedes sancti viri provoluti, sceleris veniam precarentur. Ipse quoque princeps repentino terrore perculsus, jam de salute desperaverat, cum revocato Joanne, facti pœnitens incolumitati redditus est. Factiosi étiam homines, cum eum fustibus peterent, brachiis diriguere, nec ante redditæ vires quam delicti veniam precarentur. Christum Dominum, dum Sacrum faceret, præsentem contueri, atque ex ipso divinitatis fonte cælestia mysteria haurire solitus. Abdita cordis inspicere, ac futura raro eventu præsagire frequens illi fuit, fratrisque filiam septennem mortuam excitavit. Denique mortis die prænuntiato, et Ecclesiæ sacramentis devotissime susceptis, extremum diem clausit, multis ante et post obitum miraculis gloriosus. Quibus rite probatis Alexander octavus Sanctorum numero eum adscripsit.\par
\switchcolumn
\EN
\noindent It was John, who by public preaching and private conversations, softened the hearts of the citizens so that the town was restored to peace. He grievously offended one of the nobles by rebuking him for his cruelty toward his vassals. This man sent two knights to murder him on the road. They had already come nigh him when God sent a terror upon them, so that they and their horses stood still, until they cast themselves down before the feet of the Saint, imploring his forgiveness for their sin. The Prince himself, also, smitten with a sudden dread, despaired of his salvation, till he had sent for John, who, finding him repent of his deed, restored him to soundness. Some quarrelsome men, likewise, who were fain to give him a cudgelling, found their arms stiffen, nor would their strength come back till they had asked his pardon for their wickedness. Oftentimes when he was celebrating the Holy Liturgy, the Presence of the Lord Christ became sensibly manifest to him, and he drank in things heavenly from their Divine Well-head Himself. Oftentimes also he could see the secrets of men's hearts, and foretell strange things to come. He raised from the dead his own niece, aged seven years. He foretold the day of his own death, and prepared himself by receiving most devoutly the Sacraments of the Church, (and then fell asleep in the Lord, upon the 11th day of June, in the year 1475.) God glorified him by many miracles, both before and after his death. These being duly proved, Alexander VIII. numbered him among the Saints. the Christians took them, and buried them honourably.\par
\switchcolumn*

\LA
\begin{Resp}\Rsign Iste homo perfécit ómnia quæ locútus est ei Deus, et dixit ad eum: Ingrédere in réquiem meam: \Star{} Quia te vidi justum coram me ex ómnibus géntibus.\end{Resp}
\begin{Resp}\Vsign Iste est, qui contémpsit vitam mundi, et pervénit ad cæléstia regna.\end{Resp}
\begin{Resp}\Rsign Quia te vidi justum coram me ex ómnibus géntibus.\end{Resp}
\begin{Resp}\Vsign Glória Patri, et Fílio, \Star{} et Spirítui Sancto.\end{Resp}
\begin{Resp}\Rsign Quia te vidi justum coram me ex ómnibus géntibus.\end{Resp}
\switchcolumn
\EN
\begin{Resp}\Rsign This is he which did according unto all that God commanded him; and God said unto him: Enter thou into My rest, \Star{} For thee have I seen righteous before Me among all people.\end{Resp}
\begin{Resp}\Vsign This is he which loved not his life in this world, and is come unto an everlasting kingdom.\end{Resp}
\begin{Resp}\Rsign For thee have I seen righteous before Me among all people.\end{Resp}
\begin{Resp}\Vsign Glory be to the Father, and to the Son, \Star{} and to the Holy Ghost.\end{Resp}
\begin{Resp}\Rsign For thee have I seen righteous before Me among all people.\end{Resp}
\switchcolumn*

\LA
\LessonHead{Lectio IX (pro commemoratione)}
\switchcolumn
\EN
\LessonHead{Lesson IX (when commemorated)}
\switchcolumn*

\LA
\LessonTitle{Commemoratio: S. Joannis a S. Facundo Confessoris}
\switchcolumn
\EN
\LessonTitle{Commemoration of: St. John of St. Facundus, Confessor}
\switchcolumn*

\LA
\noindent Joánnem, Sahagúni in Hispánia, nóbili génere natum, paréntes cum diu prole caruíssent, piis opéribus et oratiónibus a Deo impetrárunt. Ab ineúnte ætáte futúræ sanctitátis indícia prǽbuit. Présbyter ordinátus, ut Deo quiétius servíret, omnes ecclesiásticos provéntus, quibus mérito auctus fúerat, sponte dimísit. Salmánticæ, cum in gravíssimum morbum incidísset, arctióris disciplínæ voto se obstrínxit, quod ut rédderet, ad cœnóbium sancti Augustíni, severióri disciplína tum máxime florens, se cóntulit; in quo admíssus, virtútibus omnis provectióres anteíbat. Salmanticénses cives, cruéntis factiónibus exagitátos, tum conciónibus tum privátis collóquiis ac vitæ sanctitáte, ad tranquillitátem redúxit, non semel a præsénti discrímine divínitus liberátus. Christum Dóminum, dum Sacrum fáceret, præséntem contuéri, ábdita cordis inspícere, ac futúra præsagíre frequens illi fuit. Dénique, mortis die prænuntiáto, sanctíssime ex hac vita migrávit, multis ante et post óbitum miráculis gloriósus. Quibus rite probátis, Alexánder octávus Sanctórum número eum adscrípsit.\par
\switchcolumn
\EN
\noindent John, born of a noble family at Sahagun [St. Facundus] in Spain, was granted by God to his parents, who had long been childless, in answer to their good works and prayers. From his earliest years, he gave signs of his future holiness. When ordained priest, he renounced, of his own accord, all the ecclesiastical benefits which had been given him, that he might serve God in greater tranquillity. When he incurred a serious illness in Salamanca, he bound himself by vow to observe a severer discipline. To do so, he went to the monastery of St. Augustine, where there flourished the greatest severity of discipline. As a religious, he excelled the most advanced monks in all virtues. Through his public talks and private conversations, as well as the holiness of his life, he brought back to peaceful living the citizens of Salamanca, who had been disturbed by bloody factions. In the course of this work, he was not infrequently saved from imminent death by divine power. The Lord Christ often appeared to him while he was celebrating Mass. Often also he could divine the secrets of hearts and foretell the future. At length, having predicted the day of his death, he departed this life in a most holy way, glorified by many miracles both before and after his death. These miracles were duly proved, and Alexander VIII enrolled him in the number of the Saints.\par
\switchcolumn*

\LA
\TeDeum{Te Deum}
\switchcolumn
\EN
\TeDeum{Te Deum}
\switchcolumn*

\end{paracol}

\par\needspace{10\baselineskip}\addvspace{1.2em}{\centering\small\textsc{Die 12 Junii} · \textsc{12 June}\par}
\Office{Ss. Basilidis, Cyrini, Naboris and Nazarii Martyrum}{Sts. Basilides, Cyrinus, Nabor and Nazarius, Martyrs}{Simplex}
\addcontentsline{toc}{subsection}{Die 12 Junii · Ss. Basilidis, Cyrini, Naboris and Nazarii Martyrum\enspace\textit{Sts. Basilides, Cyrinus, Nabor and Nazarius, Martyrs}}
\begin{paracol}{2}
\LA
\LessonHead{Lectio IX (pro commemoratione)}
\switchcolumn
\EN
\LessonHead{Lesson IX (when commemorated)}
\switchcolumn*

\LA
\LessonTitle{Commemoratio: Ss. Basilidis, Cyrini, Naboris and Nazarii Martyrum}
\switchcolumn
\EN
\LessonTitle{Commemoration of: Sts. Basilides, Cyrinus, Nabor and Nazarius, Martyrs}
\switchcolumn*

\LA
\noindent Basilides, Cyrinus, Nabor et Nazarius, Romani milites, nobiles genere, et virtute illustres, christiana religione suscepta, cum Christum Dei Filium, Diocletiano imperatore, prædicarent, ab Aurelio præfecto Urbis comprehensi, et ut diis sacra facerent admoniti, ejus jussa contemnentes, missi sunt in carcerem. Quibus orantibus, cum subito clarissima lux oborta omnium oculis, qui ibidem essent, carcerem collustrasset: illo cælesti splendore commotus Marcellus custodiæ præpositus, multique alii, Christo Domino crediderunt. Verum postea e carcere emissi, ab imperatore Maximiano, cum ejus étiam neglecto imperio, unum Christum Deum et Dominum in ore haberent, scorpionibus cruciati iterum conjiciuntur in vincula: unde septimo die educti, et ante pedes imperatoris constituti, perstiterunt in irrisione inanium deorum, Jesum Christum Deum constantissime confitentes. Quam ob rem damnati, securi feriuntur. Quorum corpora feris objecta, nec ab illis tacta a Christianis honorifice sepulta sunt.\par
\switchcolumn
\EN
\noindent Basilides, Cyrinus, Nabor, and Nazarius were Roman soldiers, of illustrious birth, and distinguished gallantry. Having embraced the Christian Religion, and being found publishing that Christ was the Son of God, they were arrested by Aurelius, Praefect of Rome under the Emperor Diocletian. As they despised his orders to sacrifice to the gods, they were committed to prison. While they were at prayer there, a brilliant light broke forth before the eyes of all that were there, and shone in all the prison. Marcellinus the keeper of the prison and many others were moved by this heavenly glory to believe in the Lord ChristBasilides, Cyrinus, Nabor, and Nazarius were afterwards discharged out of the prison. However, in the reign of the Emperor Maximian, when they set light by his commands also, and had ever in their mouth that there is but one Christ, one God, and one Lord, they were tormented with whips loaded with metal, and again cast into chains. Thence, on the seventh day, they were brought out, and set before the Emperor, and there still persisted in mocking at the foolish idols, and declaring that Jesus Christ is God. They were accordingly condemned to death and beheaded. Their bodies were given to wild beasts to eat, but, as the creatures would not touch them,\par
\switchcolumn*

\LA
\TeDeum{Te Deum}
\switchcolumn
\EN
\TeDeum{Te Deum}
\switchcolumn*

\end{paracol}

\par\needspace{10\baselineskip}\addvspace{1.2em}{\centering\small\textsc{Die 13 Junii} · \textsc{13 June}\par}
\Office{S. Antonii de Padua Confessoris et Ecclesiæ Doctoris}{St. Anthony of Padua, Confessor}{Duplex}
\addcontentsline{toc}{subsection}{Die 13 Junii · S. Antonii de Padua Confessoris et Ecclesiæ Doctoris\enspace\textit{St. Anthony of Padua, Confessor}}
\begin{paracol}{2}
\LA
\RefLine{Lectiones I–III: de Scriptura occurrente.}
\switchcolumn
\EN
\RefLine{Lessons I–III: from the Scripture of the day.}
\switchcolumn*

\LA
\RefLine{Lectiones VII–IX: ex Commúni Doctoris non Pontificis (C5a).}
\switchcolumn
\EN
\RefLine{Lessons VII–IX: from the Common of a Doctor not a Bishop (C5a).}
\switchcolumn*

\LA
\LessonHead{Lectio IV}
\switchcolumn
\EN
\LessonHead{Lesson IV}
\switchcolumn*

\LA
\noindent Antónius, Ulyssipóne in Lusitánia honéstis ortus paréntibus, et ab iis pie educátus, adoléscens, institútum canonicórum regulárium suscépit. Sed, cum córpora beatórum quinque Mártyrum fratrum Minórum Conímbriam transferéntur, qui paulo ante apud Marróchium pro Christi fide passi erant, martýrii desidério incénsus ad Franciscánum órdinem transívit. Mox eódem ardóre impúlsus, ad Saracénos ire perréxit; sed, advérsa valetúdine afflíctus et redíre coáctus, cum navi ad Hispániæ líttora ténderet, ventórum vi in Sicíliam delátus est.\par
\switchcolumn
\EN
\noindent Ferdinand de Bullones, afterwards called Anthony, was born of decent parents at Lisbon in Portugal, (on the Feast of the Assumption, in the year of grace 1195.) They gave him a godly training, and while he was still a young man, he joined an Institute of Canons Regular. However, when the bodies of the five holy martyred Friars Minor, who had just suffered in Morocco for Christ's sake, were brought to Coimbra, the desire to be himself a martyr took a strong hold upon him, and in 1220 he left the Canons Regular and became a Franciscan. The same yearning led him to attempt to go among the Saracens, but he fell sick on the way, and, being obliged to turn back, the ship in which he had embarked for Spain was driven by stress of weather to Sicily.\par
\switchcolumn*

\LA
\begin{Resp}\Rsign Honéstum fecit illum Dóminus, et custodívit eum ab inimícis, et a seductóribus tutávit illum: \Star{} Et dedit illi claritátem ætérnam.\end{Resp}
\begin{Resp}\Vsign Justum dedúxit Dóminus per vias rectas, et osténdit illi regnum Dei.\end{Resp}
\begin{Resp}\Rsign Et dedit illi claritátem ætérnam.\end{Resp}
\switchcolumn
\EN
\begin{Resp}\Rsign The Lord made him honourable, and defended him from his enemies, and kept him safe from those that lay in wait for him; \Star{} And gave him perpetual glory.\end{Resp}
\begin{Resp}\Vsign He went down with him into the pit, and left him not in bonds.\end{Resp}
\begin{Resp}\Rsign And gave him perpetual glory.\end{Resp}
\switchcolumn*

\LA
\LessonHead{Lectio V}
\switchcolumn
\EN
\LessonHead{Lesson V}
\switchcolumn*

\LA
\noindent Assísium e Sicília ad capítulum generále venit: inde in erémum montis Pauli in Æmília secéssit, ubi divínis contemplatiónibus, jejúniis et vigíliis diu vacávit. Póstea, sacris ordínibus initiátus et ad prædicándum Evangélium missus, dicéndi sapiéntia et cópia tantum profécit, tantámque sui admiratiónem commóvit, ut eum summus Póntifex aliquándo concionántem áudiens, arcam Testaménti appellárit. In primis vero hǽreses summa vi profligávit, ideóque perpétuus hæreticórum málleus est vocátus.\par
\switchcolumn
\EN
\noindent From Sicily he came to Assisi to attend the General Chapter of his Order, and thence withdrew himself to the Hermitage of Monte Paolo near Bologna, where he gave himself up for a long while to consideration of the things of God, to fastings, and to watchings. Being afterwards ordained Priest and sent to preach the Gospel, his wisdom and fluency were very marked, and drew on him such admiration of men, that the Pope, once hearing him preach, called him The Ark of the Covenant. One of his chief points was to expend all his strength in attacking heresies, whence he gained the name of the Heretics' everlasting Hammer.\par
\switchcolumn*

\LA
\begin{Resp}\Rsign Amávit eum Dóminus, et ornávit eum: stolam glóriæ índuit eum, \Star{} Et ad portas paradísi coronávit eum.\end{Resp}
\begin{Resp}\Vsign Índuit eum Dóminus lorícam fídei, et ornávit eum.\end{Resp}
\begin{Resp}\Rsign Et ad portas paradísi coronávit eum.\end{Resp}
\switchcolumn
\EN
\begin{Resp}\Rsign The Lord loved him and beautified him; He clothed him with a robe of glory, \Star{} And crowned him at the gates of Paradise.\end{Resp}
\begin{Resp}\Vsign The Lord hath put on him the breast-plate of faith, and hath adorned him.\end{Resp}
\begin{Resp}\Rsign And crowned him at the gates of Paradise.\end{Resp}
\switchcolumn*

\LA
\LessonHead{Lectio VI}
\switchcolumn
\EN
\LessonHead{Lesson VI}
\switchcolumn*

\LA
\noindent Primus ex suo ordine, ob doctrínæ præstántiam, Bonóniæ et álibi sacras lítteras est interpretátus, fratrúmque suórum stúdiis præfuit. Multis vero peragrátis provínciis, anno ante óbitum Patávium venit, ubi illústria sanctitátis suæ monuménta relíquit. Dénique magnis labóribus pro glória Dei perfúnctus, méritis et miráculis clarus obdormívit in Dómino Idibus Júnii, anno salútis millésimo ducentésimo trigésimo primo. Quem Gregórius nonus, Póntifex máximus, sanctórum Confessórum número adscrípsit, et Pius duodécimus, ex Sacrórum Rítuum Congregatiónis consúlto, universális Ecclésiæ Doctórem declarávit.\par
\switchcolumn
\EN
\noindent He was the first of his Order who, on account of his excellent gift of teaching, publicly lectured at Bologna on the interpretation of Holy Scripture, and directed the studies of his brethren. He traveled through many provinces. The year before his death he came to Padua, where he left some remarkable records of his holy life. After having undergone much toil for the glory of God, full of good works and miracles, he fell asleep in the Lord upon the 13th day of June, in the year of salvation 1231. Pope Gregory IX. enrolled his name among those of the Holy Confessors.\par
\switchcolumn*

\LA
\begin{Resp}\Rsign Iste homo perfécit ómnia quæ locútus est ei Deus, et dixit ad eum: Ingrédere in réquiem meam: \Star{} Quia te vidi justum coram me ex ómnibus géntibus.\end{Resp}
\begin{Resp}\Vsign Iste est, qui contémpsit vitam mundi, et pervénit ad cæléstia regna.\end{Resp}
\begin{Resp}\Rsign Quia te vidi justum coram me ex ómnibus géntibus.\end{Resp}
\begin{Resp}\Vsign Glória Patri, et Fílio, \Star{} et Spirítui Sancto.\end{Resp}
\begin{Resp}\Rsign Quia te vidi justum coram me ex ómnibus géntibus.\end{Resp}
\switchcolumn
\EN
\begin{Resp}\Rsign This is he which did according unto all that God commanded him; and God said unto him: Enter thou into My rest, \Star{} For thee have I seen righteous before Me among all people.\end{Resp}
\begin{Resp}\Vsign This is he which loved not his life in this world, and is come unto an everlasting kingdom.\end{Resp}
\begin{Resp}\Rsign For thee have I seen righteous before Me among all people.\end{Resp}
\begin{Resp}\Vsign Glory be to the Father, and to the Son, \Star{} and to the Holy Ghost.\end{Resp}
\begin{Resp}\Rsign For thee have I seen righteous before Me among all people.\end{Resp}
\switchcolumn*

\LA
\LessonHead{Lectio IX (pro commemoratione)}
\switchcolumn
\EN
\LessonHead{Lesson IX (when commemorated)}
\switchcolumn*

\LA
\LessonTitle{Commemoratio: S. Antonii de Padua Confessoris et Ecclesiæ Doctoris}
\switchcolumn
\EN
\LessonTitle{Commemoration of: St. Anthony of Padua, Confessor}
\switchcolumn*

\LA
\noindent Antónius, Ulyssipóne in Lusitánia honéstis piísque ortus paréntibus, adoléscens institútum canonicórum regulárium suscépit; sed, martýrii desidério incénsus, ad Franciscánum órdinem transívit. Ad Saracénos missus, et advérsa valetúdine redíre coáctus, vi ventórum in Sicíliam delátus est. Mox sacris ordínibus initiátus et prædicatóris múnere fungens, tantam sui admiratiónem commóvit, ut ad sacras lítteras interpretándas Bonóniæ et álibi vocátus sit, fratrum suórum stúdiis fúerit præféctus, et arca Testaménti atque hǽresum málleus merúerit appellári. Multis autem peragrátis provínciis, anno ante óbitum Patávium venit, ubi illústria sanctitátis suæ monuménta relíquit. Méritis et miráculis clarus obdormívit in Dómino idibus júnii, anno salútis millésimo ducentésimo trigésimo primo, ætátis suæ trigésimo sexto, et a summo Pontífice Pio duodécimo universális Ecclésiæ Doctor fuit declarátus.\par
\switchcolumn
\EN
\noindent Anthony came of good and devout parents at Lisbon in Portugal. As a young man, he entered the Institute of Canons Regular, but, fired with the desire for martyrdom, he changed to the Franciscan Order. He was sent to the Saracens, but was forced to return because of ill-health. On the return voyage, strong winds blew the ship off course to Sicily. Anthony was soon given holy orders and took up the work of preaching, in which he aroused so much admiration that he was called to interpret Sacred Scriptures at Bologna and other places. He was put in charge of this brethren's studies, and was deservedly called the Ark of the Covenant. and Hammer of Heretics. After many travels, he came, a year before his death, to Padua, where left shining memories of his holiness. Famous for both merits and miracles, he fell asleep in the Lord on the 13th day of June in the year of salvation 1231, at the age of thirty-six. He was declared a Doctor of the universal Church by Pope Pius XII.\par
\switchcolumn*

\LA
\TeDeum{Te Deum}
\switchcolumn
\EN
\TeDeum{Te Deum}
\switchcolumn*

\end{paracol}

\par\needspace{10\baselineskip}\addvspace{1.2em}{\centering\small\textsc{Die 14 Junii} · \textsc{14 June}\par}
\Office{S. Basilii Magni, Episcopis Confessoris et Ecclesiæ Doctoris}{St. Basil the Great, Confessor and Doctor of the Church}{Duplex}
\addcontentsline{toc}{subsection}{Die 14 Junii · S. Basilii Magni, Episcopis Confessoris et Ecclesiæ Doctoris\enspace\textit{St. Basil the Great, Confessor and Doctor of the Church}}
\begin{paracol}{2}
\LA
\RefLine{Lectiones I–III: de Scriptura occurrente.}
\switchcolumn
\EN
\RefLine{Lessons I–III: from the Scripture of the day.}
\switchcolumn*

\LA
\LessonHead{Lectio IV}
\switchcolumn
\EN
\LessonHead{Lesson IV}
\switchcolumn*

\LA
\noindent Basilíus, nobilis Cáppadox, Athénis una cum Gregório Nazianzéno, ejus amicíssimo, sæculáribus lítteris, deínde in monastério sacris mirabíliter erudítus, eum brevi cursum fecit ad omnem doctrínæ et morum excelléntiam, ut inde Magni cognómen invénerit. Is ad prædicándum Jesu Christi Evangélium in Pontum accersítus, eam provínciam, a christiánis institútis aberrántem, ad viam salútis revocávit. Mox ab Eusébio Cæsaréæ epíscopo ad erudiéndam eam civitátem adjútor adhibétur: in cujus locum póstea succéssit. Is Fílium Patri consubstantiálem esse in primis deféndit, ac Valéntem imperatórem, sibi irátum, miráculis ádeo flexit, ut, incumbéntem ad voluntátem ejiciéndi ipsum in exsílium, a senténtia discédere coégerit.\par
\switchcolumn
\EN
\noindent This Basil was a noble Cappadocian v who studied earthly learning at Athens, in company with Gregory of Nazianzus, to whom he was united in a warm and tender friendship. He afterwards studied things sacred in a monastery, where he quickly attained an eminent degree of excellence in doctrine and life, whereby he gained to himself the surname of the Great. He was called to Pontus to preach the Gospel of Christ Jesus, and brought back into the way of salvation that country which before had been wandering astray from the rules of Christian discipline. He was shortly united as coadjutor to Eusebius, Bishop of Cassarea, for the edification of that city, and afterwards became his successor in the see. One of his greatest labours was to maintain that the Son is of one Substance with the Father, and when the Emperor Valens, moved to wrath against him, was willing to send him into exile, he so bent him by dint of the miracles which he worked that he forced him to forego his intention.\par
\switchcolumn*

\LA
\begin{Resp}\Rsign Invéni David servum meum, óleo sancto meo unxi eum: \Star{} Manus enim mea auxiliábitur ei.\end{Resp}
\begin{Resp}\Vsign Nihil profíciet inimícus in eo, et fílius iniquitátis non nocébit ei.\end{Resp}
\begin{Resp}\Rsign Manus enim mea auxiliábitur ei.\end{Resp}
\switchcolumn
\EN
\begin{Resp}\Rsign I have found David My servant, with My holy oil have I anointed him \Star{} For My hand shall help him.\end{Resp}
\begin{Resp}\Vsign The enemy shall prevail nothing against him, nor the son of wickedness afflict him.\end{Resp}
\begin{Resp}\Rsign For My hand shall help him.\end{Resp}
\switchcolumn*

\LA
\LessonHead{Lectio V}
\switchcolumn
\EN
\LessonHead{Lesson V}
\switchcolumn*

\LA
\noindent Nam et Valéntis sella, in qua factúrus decrétum de ejiciéndo e civitáte Basilío, sedére volébat, confrácta est. Et tribus ab eo cálamis adhíbitis ad scribéndam exsílii legem, nullus eórum réddidit atraméntum; et, cum nihilóminus in propósito scribéndi ímpium decrétum persísteret, ipsíus déxtera, dissolútis nervis, tota contrémuit. His commótus Valens chartam utráque manu conscídit. Ea autem nocte, quæ ad deliberándum Basilío data est, Valéntis uxor íntimis est cruciáta dolóribus, et únicus fílius in gravem morbum íncidit. Quibus ille pertérritus, iniquitátem suam recognóscens, Basilíum accérsit, quo præsénte, puer cœpit convaléscere; verum, vocátis a Valénte ad viséndum púerum hæréticis, paulo post móritur.\par
\switchcolumn
\EN
\noindent The chair upon which Valens sat down, in order to sign the decree of Basil's ejectment from the city, broke down under him, and three pens which he took one after the other to sign the edict of banishment, all would not write and when nevertheless he remained firm to write the ungodly order, his right hand shook. Valens was so frightened at these omens, that he tore the paper in two. During the night which was allowed to Basil to make up his mind, Valens' wife had a severe stomach-ache, and their only son was taken seriously ill. These things alarmed Valens so much that he acknowledged his wickedness, and sent for Basil, during whose visit the child began to get better. However, when Valens sent for some heretics to see it, it presently died.\par
\switchcolumn*

\LA
\begin{Resp}\Rsign Pósui adjutórium super poténtem, et exaltávi eléctum de plebe mea: \Star{} Manus enim mea auxiliábitur ei.\end{Resp}
\begin{Resp}\Vsign Invéni David servum meum, óleo sancto meo unxi eum.\end{Resp}
\begin{Resp}\Rsign Manus enim mea auxiliábitur ei.\end{Resp}
\switchcolumn
\EN
\begin{Resp}\Rsign I have laid help upon one that is mighty, and have exalted one chosen out of My people \Star{} For My hand shall help him.\end{Resp}
\begin{Resp}\Vsign I have found David My servant, with My holy oil have I anointed him.\end{Resp}
\begin{Resp}\Rsign For My hand shall help him.\end{Resp}
\switchcolumn*

\LA
\LessonHead{Lectio VI}
\switchcolumn
\EN
\LessonHead{Lesson VI}
\switchcolumn*

\LA
\noindent Abstinéntia et continéntia fuit admirábili; una túnica conténtus erat: in jejúnio servándo diligentíssimus, in oratióne assíduus, in qua sæpe totam noctem consumébat. Virginitátem perpétuo cóluit. Monastériis exstrúctis, ita monachórum institútum temperávit, ut solitáriæ atque actuósæ vitæ utilitátes præcláre simul conjúngeret. Multa erudíte scripsit; ac nemo, teste Gregório Nazianzéno, sacræ Scriptúræ libros vérius aut ubérius explicávit. Obiit Kaléndis Januárii, cum tantum spíritu vivens, præter ossa et pellem, nulla prætérea córporis parte constáre viderétur.\par
\switchcolumn
\EN
\noindent The abstinence and self-control of Basil were truly wonderful. He was content to wear nothing but one single garment. In observance of fasting he was most earnest, and so instant in prayer, that he would oftentimes pass the whole night therein. His virginity he kept always unsullied. He built monasteries, wherein he so adapted the institution of monasticism, that he exquisitely united for the inmates the advantages of the contemplative and of the active life. He was the author of many learned writings, and, according to the witness of Gregory of Nazianzus, no one hath ever composed more faithful and edifying explanations of the books of the Holy Scripture. He died upon the 1st day of January, (in the year of our Lord 379,) at which time so essentially spiritual was his life, that his body showed nothing but skin and bones.\par
\switchcolumn*

\LA
\begin{Resp}\Rsign Iste est, qui ante Deum magnas virtútes operátus est, et omnis terra doctrína ejus repléta est: \Star{} Ipse intercédat pro peccátis ómnium populórum.\end{Resp}
\begin{Resp}\Vsign Iste est, qui contémpsit vitam mundi, et pervénit ad cæléstia regna.\end{Resp}
\begin{Resp}\Rsign Ipse intercédat pro peccátis ómnium populórum.\end{Resp}
\begin{Resp}\Vsign Glória Patri, et Fílio, \Star{} et Spirítui Sancto.\end{Resp}
\begin{Resp}\Rsign Ipse intercédat pro peccátis ómnium populórum.\end{Resp}
\switchcolumn
\EN
\begin{Resp}\Rsign This is he which wrought great wonders before God, and the whole earth is full of his teaching \Star{} May he pray for all people, that their sins may be forgiven unto them\end{Resp}
\begin{Resp}\Vsign This is he which loved not his life in this world, and hath attained unto the kingdom of heaven.\end{Resp}
\begin{Resp}\Rsign May he pray for all people, that their sins may be forgiven unto them\end{Resp}
\begin{Resp}\Vsign Glory be to the Father, and to the Son, \Star{} and to the Holy Ghost.\end{Resp}
\begin{Resp}\Rsign May he pray for all people, that their sins may be forgiven unto them\end{Resp}
\switchcolumn*

\LA
\LessonHead{Lectio VII}
\switchcolumn
\EN
\LessonHead{Lesson VII}
\switchcolumn*

\LA
\LessonTitle{Léctio sancti Evangélii secúndum Lucam}
\switchcolumn
\EN
\LessonTitle{From the Holy Gospel according to Luke}
\switchcolumn*

\LA
\Cite{Luc 14:26-35}
\switchcolumn
\EN
\Cite{Luke 14:26-35}
\switchcolumn*

\LA
\noindent In illo témpore : Dixit Jesus turbis: Si quis venit ad me, et non odit patrem suum, et matrem, et uxórem, et fílios, et fratres, et soróres, adhuc autem et ánimam suam, non potest meus esse discípulus. Et réliqua.\par
\switchcolumn
\EN
\noindent At that time Jesus said unto the multitudes If any man come to Me, and hate not his father, and mother, and wife, and children, and brethren, and sisters, yea, and his own life also, he cannot be My disciple. And so on.\par
\switchcolumn*

\LA
\LessonTitle{Homilía sancti Basilíi Epíscopi}
\switchcolumn
\EN
\LessonTitle{Homily by St. Basil, Archbishop of Caesarea}
\switchcolumn*

\LA
\Source{Lib. Regul. fusius explic. ad interrog. 8}
\switchcolumn
\EN
\Source{Paraphrase of his Rule.}
\switchcolumn*

\LA
\noindent Perfécta quidem renuntiátio in eo consístit, ut id assequámur, ne ad ipsíus étiam vitæ affectiónem propénsi simus, et respónsum mortis habeámus, ut non simus fidéntes in nobis ipsis. Hujúsmodi autem renuntiátio inítium sumit ab alienatióne rerum externárum, véluti a possessiónibus, ab ináni glória, a vivéndi consuetúdine, a rerum inutílium amóre; quemádmodum étiam suo exémplo nobis ostendérunt sancti Dómini nostri discípuli, Jacóbus quidem et Joánnes, relícto patre Zebedǽo et ipsa quoque navícula, de qua omnis illórum victus rátio pendébat; Matthǽus vero, cum ab ipso telónio surréxit ac Dóminum secútus est.\par
\switchcolumn
\EN
\noindent This is perfect self-renunciation, when we attain to indifference as regards our own lives, and wring from death himself the confession that our trust is not in our own strength. The first step toward this crown of abnegation is to estrange ourselves from outward things, such as property, public reputation, habits of life, and affection for things unnecessary, whereof the immediate disciples of our Holy Lord have left us a fine example James, for instance, and John, who left their father Zebedee and the boats which were their only means of getting their daily bread or Matthew, who got up from the receipt of custom, and straightway followed the Lord.\par
\switchcolumn*

\LA
\begin{Resp}\Rsign Amávit eum Dóminus, et ornávit eum: stolam glóriæ índuit eum, \Star{} Et ad portas paradísi coronávit eum.\end{Resp}
\begin{Resp}\Vsign Induit eum Dóminus lorícam fídei, et ornávit eum.\end{Resp}
\begin{Resp}\Rsign Et ad portas paradísi coronávit eum.\end{Resp}
\switchcolumn
\EN
\begin{Resp}\Rsign The Lord loved him and beautified him He clothed him with a robe of glory \Star{} And crowned him at the gates of Paradise.\end{Resp}
\begin{Resp}\Vsign The Lord hath put on him the breast-plate of faith, and hath adorned him.\end{Resp}
\begin{Resp}\Rsign And crowned him at the gates of Paradise.\end{Resp}
\switchcolumn*

\LA
\LessonHead{Lectio VIII}
\switchcolumn
\EN
\LessonHead{Lesson VIII}
\switchcolumn*

\LA
\noindent Sed quid opus est nostris ratiónibus aut sanctórum virórum exémplis id quod dícimus confirmáre, cum ipsa Dómini verba in médium líceat afférre, iísque ipsis religiósam ac Deum timéntem ánimam commovére, quibus ille perspícue et sine controvérsia protestátur, dicens: Sic ígitur quicúmque ex vobis non renuntiáverit ómnibus quæ póssidet, non potest meus esse discípulus? Et álio in loco, cum prius dixísset: Si vis perféctus esse, vade, et vende ómnia quæ habes, et da paupéribus; póstea subjúnxit: Veni, séquere me.\par
\switchcolumn
\EN
\noindent But what need have we for our own arguments, or for the examples of holy men to confirm what we say, when we are able to cite the very words of the Lord Himself, and by them to move any earnest soul that loveth God Such were they unto whom He plainly and unhesitatingly declared Whosoever he be of you that forsaketh not all that he hath, he cannot be My disciple. And again, in another place, when He had said If thou wilt be perfect, go, and sell that thou hast, and give to the poor the completion of the sentence was, And come and follow Me. \textasciitilde{}(Matth. xix. 19.)\par
\switchcolumn*

\LA
\begin{Resp}\Rsign In médio Ecclésiæ apéruit os ejus, \Star{} Et implévit eum Dóminus spíritu sapiéntiæ et intelléctus.\end{Resp}
\begin{Resp}\Vsign Jucunditátem et exsultatiónem thesaurizávit super eum.\end{Resp}
\begin{Resp}\Rsign Et implévit eum Dóminus spíritu sapiéntiæ et intelléctus.\end{Resp}
\begin{Resp}\Vsign Glória Patri, et Fílio, \Star{} et Spirítui Sancto.\end{Resp}
\begin{Resp}\Rsign Et implévit eum Dóminus spíritu sapiéntiæ et intelléctus.\end{Resp}
\switchcolumn
\EN
\begin{Resp}\Rsign In the midst of the congregation did the Lord open his mouth. \Star{} And filled him with the spirit of wisdom and understanding.\end{Resp}
\begin{Resp}\Vsign He made him rich with joy and gladness.\end{Resp}
\begin{Resp}\Rsign And filled him with the spirit of wisdom and understanding.\end{Resp}
\begin{Resp}\Vsign Glory be to the Father, and to the Son, \Star{} and to the Holy Ghost.\end{Resp}
\begin{Resp}\Rsign And filled him with the spirit of wisdom and understanding.\end{Resp}
\switchcolumn*

\LA
\LessonHead{Lectio IX}
\switchcolumn
\EN
\LessonHead{Lesson IX}
\switchcolumn*

\LA
\noindent Est ígitur renuntiátio, quemádmodum docúimus, vinculórum terrénæ hujus ac temporális vitæ solútio, atque ab humánis negótiis liberátio, per quam ad ineúndam viam, qua ad Deum pervenítur, aptióres et promptióres effícimur; et expedíta rátio ad acquisitiónem usúmque rerum, quæ super aurum et lápidem pretiósum multum longe sunt pretiosióres. Et in summa, cordis humáni ad cæléstem conversatiónem translátio, ita ut dícere líceat: Nostra conversátio in cælis est; et (quod máximum est) inítium unde ad Christi similitúdinem evádimus, qui cum dives esset, propter nos pauper est factus.\par
\switchcolumn
\EN
\noindent This then, as we have taught, is self-renunciation to unlock the chains of this earthly life, which passeth away, and to set oneself free from the business of men, and so to make ourselves lither and meeter to enter on that path which leadeth to God, and let our reason be more unhampered to gain and to use those things which are far more precious than gold or precious stones. (Ps. xviii. 11.) In short, it is to have our heart in heaven and not on earth, so as to be able to say Our conversation is in heaven. (Phil. iii. 20.) And (which is the great thing) this is the first step towards the attaining to be like Christ, who, though He was rich, yet for our sakes He became poor. (2 Cor. viii. 9.)\par
\switchcolumn*

\LA
\TeDeum{Te Deum}
\switchcolumn
\EN
\TeDeum{Te Deum}
\switchcolumn*

\LA
\LessonHead{Lectio IX (pro commemoratione)}
\switchcolumn
\EN
\LessonHead{Lesson IX (when commemorated)}
\switchcolumn*

\LA
\LessonTitle{Commemoratio: S. Basilii Magni, Episcopis Confessoris et Ecclesiæ Doctoris}
\switchcolumn
\EN
\LessonTitle{Commemoration of: St. Basil the Great, Confessor and Doctor of the Church}
\switchcolumn*

\LA
\noindent Basilíus, nóbilis Cáppadox, Athénis una cum Gregório Nazianzéno, ejus amicíssimo, sæculáribus lítteris, deínde in monastério sacris mirabíliter erudítus, eum brevi cursum fecit ad omnem doctrínæ et morum excelléntiam, ut inde Magni cognómen invénerit. Ad prædicándum Jesu Christi Evangélium in Pontum accersítus, eam provínciam ad viam salútis revocávit; mox ab Eusébio Cæsaréæ epíscopo ad erudiéndam eam civitátem adjútor adhibétur, in cujus locum póstea succéssit. Is Fílium Patri consubstantiálem esse in primis deféndit, ac Valéntem imperatórem, sibi irátum et exsílium minitántem, miráculis ádeo flexit, ut a senténtia discédere coégerit. Abstinéntia et continéntia fuit admirábili; in oratióne assíduus, in ea sæpe totam noctem consumébat. Monastériis exstrúctis, ita monachórum institútum temperávit, ut solitáriæ atque actuósæ vitæ utilitátis præcláre simul conjúngeret. Multa erudíte scripsit; ac nemo, teste Gregório Nazianzéno, sacræ Scriptúræ libros vérius aut ubérius explicávit. Obiit kaléndis januárii.\par
\switchcolumn
\EN
\noindent Basil, a Cappodocian nobleman, studied profane letters at Athens together with his close friend, Gregory of Nazianzus, and took his sacred studies in a monastery. Becoming marvelously proficient in both, he soon attained such excellence in learning and in his way of life that from then on he was given the name of The Great. Summoned to preach the Gospel of Christ Jesus in Pontus, he called that province back to the way of salvation. Soon he was asked by Eusebius, Bishop of Caesarea, to aid him in teaching; and he succeeded Eusebius as bishop. Basil was among the first to defend the consubstantiality of the Son with the Father; and by his miracles he caused Emperor Valens, who was angry with him and threatening him with exile, to give up any such intentions. Basil's abstinence and continence were marvelled at; and he was constant in prayer, often spending the whole night in it. He built monasteries, ordering the monastic life so that it would best combine the advantages of the solitary life with those of the active life. He wrote many learned books; and, as Gregory of Nazianzus testifieth, no one hath explained the books of Holy Scripture more truly and fruitfully. He died on the 1st day of January.\par
\switchcolumn*

\LA
\TeDeum{Te Deum}
\switchcolumn
\EN
\TeDeum{Te Deum}
\switchcolumn*

\end{paracol}

\par\needspace{10\baselineskip}\addvspace{1.2em}{\centering\small\textsc{Die 15 Junii} · \textsc{15 June}\par}
\Office{Ss. Viti, Modesti atque Crescentiæ Martyrum}{Sts. Vitus, Modestus and Crescentia, Martyrs}{Simplex}
\addcontentsline{toc}{subsection}{Die 15 Junii · Ss. Viti, Modesti atque Crescentiæ Martyrum\enspace\textit{Sts. Vitus, Modestus and Crescentia, Martyrs}}
\begin{paracol}{2}
\LA
\RefLine{Lectiones I–II: de Scriptura occurrente.}
\switchcolumn
\EN
\RefLine{Lessons I–II: from the Scripture of the day.}
\switchcolumn*

\LA
\LessonHead{Lectio III (pro commemoratione)}
\switchcolumn
\EN
\LessonHead{Lesson III (when commemorated)}
\switchcolumn*

\LA
\noindent Vitus admodum puer inscio patre baptizatus est: quod cum ille rescivisset, nihil prætermisit, quo filium a christiana religione removeret. Qua in voluntate permanentem, Valeriano judici verberibus castigandum tradidit. Sed nihilominus in sententia persistens, patri redditus est. Sed dum eum pater gravius punire cogitat, Vitus Angeli monitu, comitibus Modesto et Crescentia, ejus educatoribus, migrat in alienas terras: ibique eam sanctitatis laudem adeptus est, ut ejus fama ad Diocletianum perlata, ipsum imperator accerseret, ut filium suum a dæmone vexatum liberaret: Quo liberato, cum ei amplissimis præmiis ingratus imperator, ut deos coleret, persuadere non potuisset, una cum Modesto et Crescentia vinculis constrictum mittit in carcerem. Quos ubi constantiores esse comperit, demitti jubet in ingens vas liquato plumbo, ferventi resina ac pice plenum: in quo cum trium Hebræorum puerorum more divinos hymnos canerent, inde erepti leoni objiciuntur: qui prostérnens se, eorum pedes lambebat. Quare inflammatus ira imperator, quod multitudinem videbat miraculo commoveri, eos in catasta sterni jubet et ita cædi eorum membra, atque ossa divelli. Quo tempore tonitrua, fulgura, magnique terræmotus fuere, quibus templa deorum corruerunt, et multi oppressi sunt. Eorum reliquias Florentia nobilis femina, unguentis conditas, honorifice sepelivit.\par
\switchcolumn
\EN
\noindent This Vitus was a child who was baptized without his father's knowledge. When his father had found it out, he used his best endeavours to dissuade his son from the Christian religion, but as he found him persistent in it, he handed him over to Valerian the judge to be whipped. But as he still remained as unshaken as before, he was given back to his father. But while his father was turning over in his mind to what severe discipline to subject him, Vitus, being warned by an Angel, fled out of the country, in company with his foster-parents Modestus and Crescentia. In his new home he gained great praise for holiness, so that the fame of it came to Diocletian, which Emperor sent for him to deliver his own child which was vexed with a devil. Him Vitus delivered, but when the Emperor found that with all his great gifts he could not bring him to worship the gods, he had the ingratitude to cast him and Modestus and Crescentia into prison, binding them in fetters. But when they were found in their prison more faithful than ever to their confession, the Emperor commanded them to be thrown into a great vessel full of melted lead, resin, and pitch. Therein these three, like the three Holy Children in the burning fiery furnace, sang praise to God and upon that they were haled forth and cast to a lion, but he lay down before them, and licked their feet. Then the Emperor, being filled with fury, more especially because he saw that the multitude that looked on were stirred up at the miracle, commanded Vitus, Modestus, and Crescentia to be stretched upon a block and their limbs crushed, and their bones rent one from the other. While as they were dying there came great thunderings, and lightnings, and earthquakes, so that temples of the gods fell down, and many men were killed. As for that which remained of the Martyrs, the noble lady Florence took it, and embalmed it with spices, and honourably buried it.\par
\switchcolumn*

\LA
\TeDeum{Te Deum}
\switchcolumn
\EN
\TeDeum{Te Deum}
\switchcolumn*

\end{paracol}

\par\needspace{10\baselineskip}\addvspace{1.2em}{\centering\small\textsc{Die 18 Junii} · \textsc{18 June}\par}
\Office{S. Ephræm Syri Confessoris et Ecclesiæ Doctoris}{St. Ephraem of Syria, Confessor and Doctor of the Church}{Duplex}
\addcontentsline{toc}{subsection}{Die 18 Junii · S. Ephræm Syri Confessoris et Ecclesiæ Doctoris\enspace\textit{St. Ephraem of Syria, Confessor and Doctor of the Church}}
\begin{paracol}{2}
\LA
\RefLine{Lectiones I–III: de Scriptura occurrente.}
\switchcolumn
\EN
\RefLine{Lessons I–III: from the Scripture of the day.}
\switchcolumn*

\LA
\RefLine{Lectio IX: de SS. Marci et Marcelliani Martyrum (06-18t).}
\switchcolumn
\EN
\RefLine{Lesson IX: of Sts. Mark and Marcellian, Martyrs (06-18t).}
\switchcolumn*

\LA
\LessonHead{Lectio IV}
\switchcolumn
\EN
\LessonHead{Lesson IV}
\switchcolumn*

\LA
\noindent Ephræm genere Syrus, Nisibeno patre natus est. Adhuc juvenis ad sanctum Jacobum episcopum se contulit, a quo baptizatus, brevi ita sanctitate et doctrina profecit, ut in schola Nisibi, Mesopotamiæ urbe, florente magister fuerit constitutus. Post Jacobi episcopi mortem, Nisibi a Persis capta, Edessam profectus est; ubi primum in monte inter monachos consedit, deínde, ut plurimos ad se confluentes homines vitaret, vitam duxit eremiticam. Edessenæ ecclesiæ diaconus ordinatus, et ob humilitatem sacerdotium recusans, omnium virtutum splendore enituit, et pietatem et religionem vera sapientiæ professione sibi comparare sategit. Spem omnem in solo Deo defixam habens, quævis humana ac transitoria contemnens, divina ac sempiterna assidue concupiscebat.\par
\switchcolumn
\EN
\noindent Ephraem was of Syrian descent, and son of a citizen of Nisibis. While yet a young man he went to the holy bishop James, by whom he was baptized, and he soon made such progress in holiness and learning as to be appointed master of a flourishing school at Nisibis, a city of Mesopotamia. After the death of the bishop James, Nisibis was captured by the Persians, and Ephraem went to Edessa. Here he settled first on the mountain among the monks, and then, that he might avoid the great numbers of men who flocked to him, he adopted the eremitical life. He was ordained deacon of the Church of Edessa, but refused the priesthood out of humility. He was conspicuous with the splendour of every virtue and strove to acquire piety and religion by professing true wisdom. He placed all his hope in God alone, despised all human and transitory things, and always longed for the divine and eternal.\par
\switchcolumn*

\LA
\begin{Resp}\Rsign Honéstum fecit illum Dóminus, et custodívit eum ab inimícis, et a seductóribus tutávit illum: \Star{} Et dedit illi claritátem ætérnam.\end{Resp}
\begin{Resp}\Vsign Justum dedúxit Dóminus per vias rectas, et osténdit illi regnum Dei.\end{Resp}
\begin{Resp}\Rsign Et dedit illi claritátem ætérnam.\end{Resp}
\switchcolumn
\EN
\begin{Resp}\Rsign The Lord made him honourable, and defended him from his enemies, and kept him safe from those that lay in wait for him; \Star{} And gave him perpetual glory.\end{Resp}
\begin{Resp}\Vsign He went down with him into the pit, and left him not in bonds.\end{Resp}
\begin{Resp}\Rsign And gave him perpetual glory.\end{Resp}
\switchcolumn*

\LA
\LessonHead{Lectio V}
\switchcolumn
\EN
\LessonHead{Lesson V}
\switchcolumn*

\LA
\noindent Cæsaream Cappadociæ, divino ductus Spiritu, cum petiisset, ipsum ibi os Ecclesiæ Basilium vidit, et uterque mutua consuetudine opportunum in modum usus est. Ad innumeros errores refellendos, qui, tunc temporis grassantes, Ecclesiam Dei divexabant, atque ad mysteria Domini nostri Jesu Christi sedulo illustranda, plurimas edidit lucubrationes, Syro sermone compositas, et fere omnes in linguam Græcam versas; atque, teste sancto Hieronymo, ipse ad tantum venit claritudinem, ut, post lectionem Scripturarum, publice in quibusdam ecclesiis ejus scripta recitarentur.\par
\switchcolumn
\EN
\noindent When, led by the Spirit of God, he went to Caesarea in Cappadocia, there he saw Basil, that mouthpiece of the Church, and both enjoyed mutual companionship in a suitable manner. In order to refute the countless errors which were rife at that time, and which were troubling the Church of God, and in order to expound zealously the divine mysteries of our Lord Jesus Christ, he wrote many studies in Syrian, almost all of which have been translated into Greek. St. Jerome beareth witness that he attained such fame, that his writings were read publicly in certain churches after the reading from the Scriptures.\par
\switchcolumn*

\LA
\begin{Resp}\Rsign Amávit eum Dóminus, et ornávit eum: stolam glóriæ índuit eum, \Star{} Et ad portas paradísi coronávit eum.\end{Resp}
\begin{Resp}\Vsign Índuit eum Dóminus lorícam fídei, et ornávit eum.\end{Resp}
\begin{Resp}\Rsign Et ad portas paradísi coronávit eum.\end{Resp}
\switchcolumn
\EN
\begin{Resp}\Rsign The Lord loved him and beautified him; He clothed him with a robe of glory, \Star{} And crowned him at the gates of Paradise.\end{Resp}
\begin{Resp}\Vsign The Lord hath put on him the breast-plate of faith, and hath adorned him.\end{Resp}
\begin{Resp}\Rsign And crowned him at the gates of Paradise.\end{Resp}
\switchcolumn*

\LA
\LessonHead{Lectio VI}
\switchcolumn
\EN
\LessonHead{Lesson VI}
\switchcolumn*

\LA
\noindent Universa illius opera, tam splendido doctrinæ lumine referta effecerunt, ut idem Sanctus, adhuc vivens, tamquam Ecclesiæ Doctor, magno honore habitus fuerit. Metrica quoque cantica composuit in laudem beatissimæ Virginis Mariæ ac Sanctorum: quam ob causam a Syris Spiritus Sancti cithara merito fuit appellatus. In mirifica ac pia devotione erga eamdem Virginem immaculatam primum excelluit. Meritis plenus, Edessæ, in Mesopotamia, decimo quarto Kalendas Julii, decessit sub Valente principe: eumque, instántibus plúribus Sanctæ Románæ Ecclésiæ Cardinálibus, Patriárchis, Archiepíscopis, Epíscopis, Abbátibus, et religiósis famíliis, Benedíctus Papa décimus quintus, ex Sacrórum Rítuum Congregatiónis consúlto, universális Ecclésiæ Doctórem declarávit.\par
\switchcolumn
\EN
\noindent His works taken as a whole, so infused with the bright light of learning, brought it about that this holy man, while yet alive, was held in great honour, and was even considered a Doctor of the Church. He also composed songs in verse, in honour of the Most Blessed Virgin Mary, and of the Saints, and for this reason he was appropriately named by the Syrians the Harp of the Holy Ghost. He was noted for his great and tender devotion towards the Immaculate Virgin. He died, rich in merits, at Edessa in Mesopotamia on the 18th day of June in the reign of Valens. Pope Benedict XV, at the instance of many Cardinals of the Holy Roman Church, Patriarchs, Archbishops, Bishops, Abbots, and religious communities, declared him by a decree of the Congregation of Sacred Rites to be a Doctor of the universal Church.\par
\switchcolumn*

\LA
\begin{Resp}\Rsign Iste homo perfécit ómnia quæ locútus est ei Deus, et dixit ad eum: Ingrédere in réquiem meam: \Star{} Quia te vidi justum coram me ex ómnibus géntibus.\end{Resp}
\begin{Resp}\Vsign Iste est, qui contémpsit vitam mundi, et pervénit ad cæléstia regna.\end{Resp}
\begin{Resp}\Rsign Quia te vidi justum coram me ex ómnibus géntibus.\end{Resp}
\begin{Resp}\Vsign Glória Patri, et Fílio, \Star{} et Spirítui Sancto.\end{Resp}
\begin{Resp}\Rsign Quia te vidi justum coram me ex ómnibus géntibus.\end{Resp}
\switchcolumn
\EN
\begin{Resp}\Rsign This is he which did according unto all that God commanded him; and God said unto him: Enter thou into My rest, \Star{} For thee have I seen righteous before Me among all people.\end{Resp}
\begin{Resp}\Vsign This is he which loved not his life in this world, and is come unto an everlasting kingdom.\end{Resp}
\begin{Resp}\Rsign For thee have I seen righteous before Me among all people.\end{Resp}
\begin{Resp}\Vsign Glory be to the Father, and to the Son, \Star{} and to the Holy Ghost.\end{Resp}
\begin{Resp}\Rsign For thee have I seen righteous before Me among all people.\end{Resp}
\switchcolumn*

\LA
\LessonHead{Lectio VII}
\switchcolumn
\EN
\LessonHead{Lesson VII}
\switchcolumn*

\LA
\LessonTitle{Léctio sancti Evangélii secúndum Matthǽum.}
\switchcolumn
\EN
\LessonTitle{From the Holy Gospel according to Matthew}
\switchcolumn*

\LA
\Cite{Matt 5:13-19}
\switchcolumn
\EN
\Cite{Matt 5:13-19}
\switchcolumn*

\LA
\noindent In illo témpore: Dixit Jesus discípulis suis: Vos estis sal terræ. Quod si sal evanúerit, in quo saliétur? Et réliqua.\par
\switchcolumn
\EN
\noindent At that time: Jesus said unto His disciples: Ye are the salt of the earth. But if the salt have lost his savour, wherewith shall it be salted. And so on.\par
\switchcolumn*

\LA
\LessonTitle{Homilia sancti Ephræm Syri, Diaconi}
\switchcolumn
\EN
\LessonTitle{Homily of St. Ephrem of Syria Deacon}
\switchcolumn*

\LA
\Source{Sermo de vita et exercitatione monastica}
\switchcolumn
\EN
\Source{Sermo de vita et exercitatione monastica}
\switchcolumn*

\LA
\noindent Præclarum est bonum inchoare atque perficere, et gratum Deo esse et utilem proximo, ipsique summo ac dulcissimo rectori nostro Christo Jesu placere, qui ait: Vos estis sal terræ, et columna cælorum. Labor afflictionis tuæ, dilectissime, tamquam somnus est; porro laboris requies inenarrabilis atque inæstimabilis. Attende ergo tibi ipsi sollicite, ne utrumque pariter amittas, dum neutrum plene persequeris, præsentem scilicet sempiternamque lætitiam. Stude potius perfectam virtutem consequi, ornatam atque insignitam omnibus quæ diligit Deus. Hanc si assequaris, numquam irritabis Deum, neque proximum tuum violabis.\par
\switchcolumn
\EN
\noindent Clearly it is a good thing to begin and to accomplish, to be acceptable to God and useful to one's neighbour, to please our most high and gracious ruler, Christ Jesus, who saith: Ye are the salt of the earth and the pillars of the heavens. The burden of your affliction, dearly beloved, is as a sleep; then there followeth unspeakable and priceless rest from labour. Therefore, watch thyself carefully, so that while thou followest after neither, wholeheartedly, thou shouldst not lose both the present and the eternal joy. Study rather to attain to the perfect virtue, adorned and stamped by all that God loveth. If thou dost strive after this, thou wilt never either anger God nor injure thy neighbour.\par
\switchcolumn*

\LA
\begin{Resp}\Rsign Iste est qui ante Deum magnas virtútes operátus est, et de omni corde suo laudávit Dóminum: \Star{} Ipse intercédat pro peccátis ómnium populórum.\end{Resp}
\begin{Resp}\Vsign Ecce homo sine queréla, verus Dei cultor, ábstinens se ab omni ópere malo, et pérmanens in innocéntia sua.\end{Resp}
\begin{Resp}\Rsign Ipse intercédat pro peccátis ómnium populórum.\end{Resp}
\switchcolumn
\EN
\begin{Resp}\Rsign This is he which wrought great wonders before God, and praised the Lord with all his heart. \Star{} May he pray for all people, that their sins may be forgiven unto them.\end{Resp}
\begin{Resp}\Vsign Behold a man without blame, a worshipper of God in truth, keeping himself clean from every evil work, and abiding still in his innocency.\end{Resp}
\begin{Resp}\Rsign May he pray for all people, that their sins may be forgiven unto them.\end{Resp}
\switchcolumn*

\LA
\LessonHead{Lectio VIII}
\switchcolumn
\EN
\LessonHead{Lesson VIII}
\switchcolumn*

\LA
\noindent Porro virtus ista, unica uniusque speciei dicitur, variarum virtutum in se ipsa habens pulchritudinem. Diadema regium absque pretiosis lapidibus candentibusque margaritis connecti texique non potest; ita et hæc unica virtus sine variarum fulgore virtutum constare nequit. Est enim profecto simillima diademati regio. Nam, ut illi, si lapis unus aut margarita defuerit, in regio capite lucere pleniter nequit; ita et hæc unica virtus, nisi virtutum ceterarum honore conseritur, perfecta virtus non appellatur. Similis item est pretiosissimis epulis, exquisitissimis condimentis præparatis, sed sale carentibus. Sicut enim pretiosi illi cibi sine sale comedi nequeunt; ita et ista virtus uniformis, si variarum virtutum gloria et honore decoretur, absit autem Dei proximique dilectio, vilis prorsus atque contemptibilis est.\par
\switchcolumn
\EN
\noindent Moreover this virtue is called special and singular, having within itself the beauty of different virtues. We cannot have a royal diadem without precious stones and gleaming pearls arranged and fitted together; so likewise this single virtue cannot remain without the splendour of other different virtues. For it is, indeed, most like a kingly crown. For as in the latter case, if one stone or one pearl be missing it cannot shine perfectly upon the royal head; so, then, this special virtue cannot be called a perfect virtue, unless it is worthily connected with other virtues. Again, it is like unto very rich food, furnished with exquisite seasonings, but lacking salt. For as those rich dishes cannot be eaten without salt, so this simple virtue may be adorned with the glory and honour of different virtues, but if a man lack the love of God and of his neighbour, he is wholly worthless and contemptible.\par
\switchcolumn*

\LA
\begin{Resp}\Rsign In médio Ecclésiæ apéruit os ejus, \Star{} Et implévit eum Dóminus spíritu sapiéntiæ et intelléctus.\end{Resp}
\begin{Resp}\Vsign Jucunditátem et exsultatiónem thesaurizávit super eum.\end{Resp}
\begin{Resp}\Rsign Et implévit eum Dóminus spíritu sapiéntiæ et intelléctus.\end{Resp}
\begin{Resp}\Vsign Glória Patri, et Fílio, \Star{} et Spirítui Sancto.\end{Resp}
\begin{Resp}\Rsign Et implévit eum Dóminus spíritu sapiéntiæ et intelléctus.\end{Resp}
\switchcolumn
\EN
\begin{Resp}\Rsign In the midst of the congregation did the Lord open his mouth. \Star{} And filled him with the spirit of wisdom and understanding.\end{Resp}
\begin{Resp}\Vsign He made him rich with joy and gladness.\end{Resp}
\begin{Resp}\Rsign And filled him with the spirit of wisdom and understanding.\end{Resp}
\begin{Resp}\Vsign Glory be to the Father, and to the Son, \Star{} and to the Holy Ghost.\end{Resp}
\begin{Resp}\Rsign And filled him with the spirit of wisdom and understanding.\end{Resp}
\switchcolumn*

\LA
\LessonHead{Lectio IX (pro commemoratione)}
\switchcolumn
\EN
\LessonHead{Lesson IX (when commemorated)}
\switchcolumn*

\LA
\LessonTitle{Commemoratio: S. Ephræm Syri Confessoris et Ecclesiæ Doctoris}
\switchcolumn
\EN
\LessonTitle{Commemoration of: St. Ephraem of Syria, Confessor and Doctor of the Church}
\switchcolumn*

\LA
\noindent Ephræm, génere Syrus, Nisibéno patre natus est. Adhuc júvenis ad sanctum Jacóbum epíscopum se cóntulit, a quo baptizátus, brevi ita sanctitáte et doctrína profécit, ut in schola Nísibi, Mesopotámiæ urbe, florénte magíster fúerit constitútus. Edessénæ ecclésiæ diáconus ordinátus et ob humilitátem sacerdótium recúsans, ómnium virtútum splendóre enítuit, et pietátem et religiónem vera sapiéntiæ professióne sibi comparáre satégit. Univérsa illíus ópera, tam spléndido doctrínæ lúmine reférta, effecérunt, ut idem Sanctus, adhuc vivens, tamquam Ecclésiæ Doctor, magno honóre hábitus fúerit. In mirífica ac pia devotióne erga Vírginem Immaculátam primum excélluit. Méritis plenus, Edéssæ in Mesopotámia, décimo quarto kaléndas júlii, decéssit sub Valénte príncipe: eúmque Benedíctus Papa décimus quintus, ex Sacrórum Rítuum Congregatiónis consúlto, universális Ecclésiæ Doctórem declarávit.\par
\switchcolumn
\EN
\noindent Ephraem was of Syrian stock, his father being a citizen of Nisibis. While he was still young, he went to the holy bishop James to be baptized. In a short time he advanced so much in holiness and learning that he was made master of a flourishing school at Nisibis, a city in Mesopotamia. Ordained deacon of the Church of Edessa, and refusing the priesthood out of humility, he shone with the splendour of all virtues, and sought to acquire devotion and religion by the profession of true wisdom. All his works, illuminated with the bright light of learning, caused this Saint to be treated with great honour as a Doctor of the Church even in his lifetime. He excelled, above all, in a wonderful and loving devotion to the Immaculate Virgin. Rich in merits, he died on the 18th day of June at Edessa in Mesopotamia, under the Emperor Valens. Pope Benedict XV, after consulting the Congregation of Sacred Rites, declared him a Doctor of the Universal Church.\par
\switchcolumn*

\LA
\TeDeum{Te Deum}
\switchcolumn
\EN
\TeDeum{Te Deum}
\switchcolumn*

\end{paracol}

\par\needspace{10\baselineskip}\addvspace{1.2em}{\centering\small\textsc{Die 18 Junii} · \textsc{18 June}\par}
\Office{SS. Marci et Marcelliani Martyrum}{Sts. Mark and Marcellian, Martyrs}{Simplex}
\addcontentsline{toc}{subsection}{Die 18 Junii · SS. Marci et Marcelliani Martyrum\enspace\textit{Sts. Mark and Marcellian, Martyrs}}
\begin{paracol}{2}
\LA
\LessonHead{Lectio IX (pro commemoratione)}
\switchcolumn
\EN
\LessonHead{Lesson IX (when commemorated)}
\switchcolumn*

\LA
\LessonTitle{Commemoratio: SS. Marci et Marcelliani Martyrum}
\switchcolumn
\EN
\LessonTitle{Commemoration of: Sts. Mark and Marcellian, Martyrs}
\switchcolumn*

\LA
\noindent Marcus et Marcellianus fratres Romani, propter christianam fidem a Fabiano duce comprehensi, ad stipitem alligati sunt, pedibus clavis confixis. Ad quos cum ita loqueretur judex: Resipiscite miseri, et vos ipsos ab his cruciatibus eripite: responderunt: Numquam tam jucunde epulati sumus, quam hæc libenter Jesu Christi causa perferimus, in cujus amóre nunc fixi esse cœpimus: utinam tamdiu nos hæc pati sinat, quamdiu hoc corruptibili corpore vestiti erimus. Qui diem noctemque in tormentis divinas laudes canentes, denique telis transfixi, ad martyrii gloriam pervenerunt. Quorum corpora via Ardeatina sepulta sunt.\par
\switchcolumn
\EN
\noindent Mark and Marcellian were two brothers, Romans, who were arrested by Duke Fabian for believing in Christ, and fastened to a beam, to which their feet were nailed. The Judge said to them Wretched creatures, do think for a moment, and free yourselves from such suffering. But they answered him We have never enjoyed any dinner so much as we do what we are now undergoing here for Jesus Christ's sake. We have got ourselves a little fast to His love now. Would that He would let us suffer this as long as we are clad in this corruptible body. Still suffering, they for a day and a night sang the praises of God continually, and in the end were thrust through with darts, and so attained the glory of Martyrdom. Their bodies are buried upon the Way to Ardea.\par
\switchcolumn*

\LA
\TeDeum{Te Deum}
\switchcolumn
\EN
\TeDeum{Te Deum}
\switchcolumn*

\end{paracol}

\par\needspace{10\baselineskip}\addvspace{1.2em}{\centering\small\textsc{Die 19 Junii} · \textsc{19 June}\par}
\Office{S. Julianæ de Falconeriis Virginis}{St. Juliana Falconeri, Virgin}{Duplex}
\addcontentsline{toc}{subsection}{Die 19 Junii · S. Julianæ de Falconeriis Virginis\enspace\textit{St. Juliana Falconeri, Virgin}}
\begin{paracol}{2}
\LA
\RefLine{Lectiones I–III: de Scriptura occurrente.}
\switchcolumn
\EN
\RefLine{Lessons I–III: from the Scripture of the day.}
\switchcolumn*

\LA
\RefLine{Lectiones VII–VIII: ex Commúni Virginum tantum (C6a).}
\switchcolumn
\EN
\RefLine{Lessons VII–VIII: from the Common of a Virgin not a Martyr (C6a).}
\switchcolumn*

\LA
\RefLine{Lectio IX: de Ss. Gervasii et Protasii Martyrum (06-19o).}
\switchcolumn
\EN
\RefLine{Lesson IX: of Ss. Gervase and Protase, Martyrs (06-19o).}
\switchcolumn*

\LA
\LessonHead{Lectio IV}
\switchcolumn
\EN
\LessonHead{Lesson IV}
\switchcolumn*

\LA
\noindent Juliana, ex nobili Falconéria família, Claríssimo patre, qui templum Deiparæ ab Angelo salutátæ ære suo magnifice a fundamentis Floréntiæ, ut nunc visitur, eréxit, matre Reguardata, ambóbus jam senescéntibus ac ad id tempus sterílibus, nata est anno millesimo ducentésimo septuagesimo. Ab incunabulis non éxiguum futuræ sanctitátis specimen dedit ; vagiéntibus quippe labris, suavíssima Jesu et Maríæ nómina ultro proferre audíta est. Puerítiam postmodum ingressa totam se christiánis virtútibus mancipávit, in quibus adeo excelluit, ut beátus Aléxius patruus, cujus institútis ac exemplis instruebátur, matri dicere non dubitaverit, ipsam non feminam peperisse, sed angelum. Nam ita modesto vultu animoque ab omni vel levíssima erroris macula pura fuit, ut óculos vel levíssima, erroris macula pura fuit, ut óculos numquam in toto vitæ cursu ad hóminis fáciem intuéndam erexerit, auditoque peccáti vocabulo contremúerit. Expleto nondum décimo quinto ætátis suæ anno, re familiari, licet opulenta, terrenisque posthabitis nuptiis, Deo virginitátem in mánibus divi Philippi Benitii solemniter vovit, ab eoque ómnium prima religiósum Mantellatárum habitum, ut dicunt, sumpsit.\par
\switchcolumn
\EN
\noindent Juliana was a daughter of the noble family of the Falconieri, and was born in the year 1270. Her father was the same who at his own costs so splendidly built from the foundations the Church of Our Lady of the Annunciation as it now standeth at Florence. Her mother's name was Reguardata. They were both well stricken in years, and, until the birth of Juliana, had been childless. From her very cradle she gave tokens of the holiness of life to which she afterwards attained. And from the murmuring of her baby lips was caught the sweet sound of the names of Jesus and Mary. As she entered on her girlhood, she delivered herself up entirely to the pursuit of Christian godliness, and so excellently shone therein, that her uncle, the Blessed Alexius, scrupled not to tell her mother that she had given birth to an Angel rather than to a woman. So modest was her carriage, and so clean her soul from the lightest speck of indiscretion, that she never in her whole life stared a man in the face, and that the very mention of sin made her shiver, and when the story of a grievous crime was told her, she dropped down nearly fainting. Before she had finished her fifteenth year, she renounced her inheritance, although a rich one, and all prospect of an earthly marriage, and made to God a vow of virginity, before holy Philip Benizi, from whom she was the first to receive the religious habit of what are called the mantled nuns.\par
\switchcolumn*

\LA
\begin{Resp}\Rsign Propter veritátem, et mansuetúdinem, et justítiam: \Star{} Et dedúcet te mirabíliter déxtera tua.\end{Resp}
\begin{Resp}\Vsign Spécie tua et pulchritúdine tua inténde, próspere procéde, et regna.\end{Resp}
\begin{Resp}\Rsign Et dedúcet te mirabíliter déxtera tua.\end{Resp}
\switchcolumn
\EN
\begin{Resp}\Rsign Because of truth, and meekness, and righteousness; \Star{} And thy right hand shall lead thee wonderfully.\end{Resp}
\begin{Resp}\Vsign In thy comeliness, and thy beauty, go forward, fare prosperously, and reign.\end{Resp}
\begin{Resp}\Rsign And thy right hand shall lead thee wonderfully.\end{Resp}
\switchcolumn*

\LA
\LessonHead{Lectio V}
\switchcolumn
\EN
\LessonHead{Lesson V}
\switchcolumn*

\LA
\noindent Julianæ exemplum secutæ sunt plurimæ ex nobilióribus familiis feminæ, ac mater ipsa fíliæ sese religiose instituéndam dedit, ita ut, aucto paulatim numero, ordinem Mantellatárum institúerit, ac illi pie vivéndi leges summa prudéntia ac sanctitáte tradíderit. Ejus virtútes cum optime perspectas divus Benitius haberet, morti próximus nulli mélius quam Julianæ, non feminas tantum, sed et totum Servórum ordinem, cujus propagator et moderátor exstiterat, commendátum vóluit. Verum ipsa demisse semper de se cogitábat ; et, cum ceterárum esset magistra, in re quaque domestica, licet vili, soróribus famulabátur. Assiduitate orándi íntegras insumebat dies in éxtasim sæpíssime rapta, et, si reliquum, in sedandis civium dissidiis, criminosis a via iniquitátis retrahendis, ac inserviéndis impendebat ægrotis ; quorum quandoque saniem ex ulcéribus manántem, admoto ore lambens, eos sanitáti restituébat. Corpus suum flagris, nodosis funículis, férreis cingulis, vigiliis, humi nudæ cubando, térere sólita fuit. Parcíssimo cibo, et hoc vili, quatuor hebdomadæ diebus, et reliquis duobus solo Angelórum pane contenta, excepto die sábbati, quo pane solo et aqua nutriebátur.\par
\switchcolumn
\EN
\noindent She put herself for instruction under her daughter. Thus in a little while their number increased, and she became the foundress of the Order of Mantled nuns, to whom she gave a rule of life full of wisdom and godliness. Holy Philip Benizi, having thorough knowledge of her excellence, chose her above all living to whom at his death to leave the care not of the women only but of the whole of the Order of the Servants of the Blessed Virgin Mary, of which he had been the propagator and director. Juliana, who deemed ever lowly of herself, even when she was the mistress of the others, ministered to her sisters in the meanest offices of the work of the house. She passed whole days in incessant prayer, and was often rapt in spirit, and the remainder of her time she toiled to make peace among the citizens, who were at variance together, to recall transgressors from the ways of iniquity, and to nurse the sick, to cure whom she would sometimes even use her tongue to remove the matter that ran from their sores. It was her custom to afflict her own body with whips, knotted cords, iron girdles, watching, and sleeping upon the ground. Upon Mondays, Tuesdays, Wednesdays, and Thursdays, she ate very sparingly some unpalatable food, upon Fridays she took nothing except the Bread of Angels, and upon Saturdays, besides the Holy Communion, only bread and water.\par
\switchcolumn*

\LA
\begin{Resp}\Rsign Dilexísti justítiam, et odísti iniquitátem: \Star{} Proptérea unxit te Deus, Deus tuus, óleo lætítiæ.\end{Resp}
\begin{Resp}\Vsign Propter veritátem, et mansuetúdinem, et justítiam.\end{Resp}
\begin{Resp}\Rsign Proptérea unxit te Deus, Deus tuus, óleo lætítiæ.\end{Resp}
\switchcolumn
\EN
\begin{Resp}\Rsign Thou hast loved righteousness, and hated iniquity; \Star{} Therefore God, thy God, hath anointed thee with the oil of gladness.\end{Resp}
\begin{Resp}\Vsign Because of truth, and meekness, and righteousness.\end{Resp}
\begin{Resp}\Rsign Therefore God, thy God, hath anointed thee with the oil of gladness.\end{Resp}
\switchcolumn*

\LA
\LessonHead{Lectio VI}
\switchcolumn
\EN
\LessonHead{Lesson VI}
\switchcolumn*

\LA
\noindent Dura hujusmodi vivéndi ratióne in stómachi morbum íncidit, quo ingravescente, cum septuagesimum ætátis annum ágeret, ad extremum vitæ spatium redacta est. Diuturnæ valetúdinis incómmoda hílari vultu constantique animo pértulit ; de uno tantum cónqueri audíta est, quod, cum cibum cápere ac retinére nullo modo posset, ab Eucharistica mensa ob sacraménti reveréntiam arcerétur. Verum, his in angustiis constituta, sacerdotem rogávit ut allátum divinum Panem, quem ore sumere nequíbat, péctori saltem exterius admovéret. Precibus illíus morem gessit sacerdos ; et mirum! eodem témporis momento divinus Panis dispáruit, et Juliana sereno ac ridénti vultu exspirávit. Res supra fidem tamdiu fuit, donec virgineum de more curarétur corpus ; invénta enim est circa sinistrum péctoris latus carni véluti sigillo impressa forma hostiæ, quæ Christi crucifixi effigiem repræsentábat. Hujus prodigii fama ceterorúmque miraculórum non Floréntiæ tantum, sed totius christiáni orbis veneratiónem illi conciliávit, ac per quatuor prope íntegra sæcula adeo aucta est, ut tandem Benedíctus Papa décimus tertius in ejus celebritate Offícium proprium recitari ab universo ordine beátæ Maríæ Vírginis Servórum jússerit. Clemens vero duodecimus, munificentíssimus ejusdem ordinis protéctor, novis in dies miraculis coruscántem, sanctárum Vírginum catálogo adscripsit.\par
\switchcolumn
\EN
\noindent The self-inflicted hardships of her life brought upon her a disease of the stomach, whereby, when she was seventy years of age, she was brought to the point of death. She bore the daily sufferings of her illness with a smiling face and a brave heart. The only thing of which she was heard to complain was that, her stomach being so weak that she could not keep down any food, she was withheld by reverence for the Sacrament from drawing near to the Lord's Table. Finding herself in these straits she begged the Priest to bring the Bread of God, and, as she dared not take It into her mouth, to put It as near as possible to her heart. The Priest did as she wished, and, to the amazement of all present, the Divine Bread at once disappeared from sight, and at the same instant a smile of joyous peace crossed the face of Juliana, and she gave up the ghost. All were confounded until the virgin body was being laid out after death in the accustomed manner. Then there was found upon the left side of the bosom a mark like the stamp of a seal, reproducing the form of the Sacred Host, the mould of which was one of those that bear a figure of Christ crucified. The noise of this and other wonders got for Juliana a reverence not only from Florence, but from all parts of the Christian world, which so increased through the course of four hundred years, that Pope Benedict XIII. commanded an office in her honour to be said by the whole Order of Servants of the Blessed Virgin Mary, and Clement XII., a munificent Protector of the same Order, finding new signs and wonders shedding lustre upon her memory every day, numbered her among Holy Virgins.\par
\switchcolumn*

\LA
\begin{Resp}\Rsign Afferéntur Regi vírgines post eam, próximæ ejus; \Star{} Afferéntur tibi in lætítia et exsultatióne.\end{Resp}
\begin{Resp}\Vsign Spécie tua et pulchritúdine tua; inténde, próspere procéde, et regna.\end{Resp}
\begin{Resp}\Rsign Afferéntur tibi in lætítia et exsultatióne.\end{Resp}
\begin{Resp}\Vsign Glória Patri, et Fílio, \Star{} et Spirítui Sancto.\end{Resp}
\begin{Resp}\Rsign Afferéntur tibi in lætítia et exsultatióne.\end{Resp}
\switchcolumn
\EN
\begin{Resp}\Rsign After her shall virgins be brought unto the King, her fellows \Star{} They shall be brought unto thee with gladness and rejoicing.\end{Resp}
\begin{Resp}\Vsign In thy comeliness and thy beauty, go forward, fare prosperously, and reign.\end{Resp}
\begin{Resp}\Rsign They shall be brought unto thee with gladness and rejoicing.\end{Resp}
\begin{Resp}\Vsign Glory be to the Father, and to the Son, \Star{} and to the Holy Ghost.\end{Resp}
\begin{Resp}\Rsign They shall be brought unto thee with gladness and rejoicing.\end{Resp}
\switchcolumn*

\LA
\LessonHead{Lectio IX (pro commemoratione)}
\switchcolumn
\EN
\LessonHead{Lesson IX (when commemorated)}
\switchcolumn*

\LA
\LessonTitle{Commemoratio: S. Julianæ de Falconeriis Virginis}
\switchcolumn
\EN
\LessonTitle{Commemoration of: St. Juliana Falconeri, Virgin}
\switchcolumn*

\LA
\noindent Juliána, ex nóbili Falconéria família, ab incunábulis, vagiéntibus lábiis suavíssima Jesu et Maríæ nómina ultro proférre audíta est. Expléto nondum décimo quinto ætátis anno, re familiári, licet opulénta, terrenísque núptiis posthábitis, Deo virginitátem in mánibus divi Philíppi Benítii solémniter vovit, et ab eo ómnium prima religiósum Mantellatárum, quas vocant, hábitum sumpsit. Cum vero ejus exémplum plúrimæ ex nobilióribus féminis sequeréntur, et mater fíliæ sese instituéndam dedísset, Juliána órdinem Mantellatárum instítuit. Mira humilitáte, assíduo oratiónis stúdio, singulári abstinéntia excélluit. Cum ob advérsam valetúdinem cibum cápere ac retinére nullo modo posset, ideóque ab Eucharística mensa arcerétur, sacerdótem rogávit, ut allátum divínum Panem, quem ore súmere nequíbat, péctori saltem extérius admovéret. Quod cum sacérdos præstitísset, íllico divínus Panis dispáruit, et Juliána ridénti vultu exspirávit.\par
\switchcolumn
\EN
\noindent When Juliana, of the noble family of the Falconieri, was still in her cradle, her baby lips were heard to utter, without any prompting, the sweet names of Jesus and Mary. Before she was fifteen years old, she renounced a rich inheritance and an earthly wedding and took a solemn vow of virginity in the presence of St. Philip Benizi. She was the first to receive from him the habit of the religious called the Mantellates. When many noble ladies followed her example, and even her mother gave herself over to her daughter to be instructed in the religious life, Juliana founded the Order of the Mantellate Nuns. She excelled in a wonderful humility, a constant zeal for prayer and an amazing abstinence. When her health failed so that she could take and retain no food at all, and was therefore kept from the Eucharistic table, she asked the Priest to place the divine Bread on her breast, since she could not receive it with her mouth. When he did so, the holy Bread disappeared at once, and Juliana, smiling, departed this life.\par
\switchcolumn*

\LA
\TeDeum{Te Deum}
\switchcolumn
\EN
\TeDeum{Te Deum}
\switchcolumn*

\end{paracol}

\par\needspace{10\baselineskip}\addvspace{1.2em}{\centering\small\textsc{Die 19 Junii} · \textsc{19 June}\par}
\Office{Ss. Gervasii et Protasii Martyrum}{Ss. Gervase and Protase, Martyrs}{Simplex}
\addcontentsline{toc}{subsection}{Die 19 Junii · Ss. Gervasii et Protasii Martyrum\enspace\textit{Ss. Gervase and Protase, Martyrs}}
\begin{paracol}{2}
\LA
\LessonHead{Lectio IX (pro commemoratione)}
\switchcolumn
\EN
\LessonHead{Lesson IX (when commemorated)}
\switchcolumn*

\LA
\LessonTitle{Commemoratio: Ss. Gervasii et Protasii Martyrum}
\switchcolumn
\EN
\LessonTitle{Commemoration of: Ss. Gervase and Protase, Martyrs}
\switchcolumn*

\LA
\noindent Gervasius et Protasius, Vitalis et Valériæ fílii, quorum pater Ravennæ, mater Mediolani pro Christi Dómini fide martyrium subiérunt, distributo paupéribus patrimonio, domesticos servos libertáte donarunt. Quo facto Gentílium sacerdótes immane in illos conceptum ódium habebant. Quare, cum Astasius comes in bellum proficisci vellet, hanc occasiónem perdéndi pios fratres se nactos esse putavérunt. Itaque Astasio persuadent se a diis admonitos esse, nullo modo eum in bello victórem futurum, nisi Gervasio et Protasio coactis Christum negare, eosdem ad sacra diis facienda compélleret. Quod cum illi detestaréntur, Astasius imperávit Gervasium tamdiu cædi, dum inter verbera exspiraret; Protasium fustibus contusum secúri pércuti jubet. Quorum córpora Philippus Christi servus clam sústulit et in suis ædibus sepelívit: quæ póstea sanctus Ambrosius, Dei mónitu invénta, in loco sacro et insígni collocanda curávit. Passi sunt Mediolani décimo tertio Kalendas Julii.\par
\switchcolumn
\EN
\noindent Gervase and Protase were the sons of Vitalis and Valeria, both of whom testified even unto death for the Lord Christ's sake, the father at Ravenna, and the mother at Milan. After the victory of their parents, Gervase and Protase gave all the inheritance to the poor, and set free their slaves. This act of theirs stirred up against them a savage hatred on the part of the heathen priests, and when the Count Astasius was about setting forth to war, they believed they had got a good occasion for the destruction of the two godly brethren. They persuaded Astasius that their gods had revealed to them that he had no chance of conquering in the war, unless he had first made Gervase and Protase to deny Christ and to offer sacrifice to the gods. Being commanded so to do, they flatly refused, and Astasius then ordered Gervase to be lashed until he died between the stripes, and Protase to be cudgelled and beheaded. A servant of Christ named Philip took away their dead bodies by stealth, and buried them in his own house, and, in after times, St. Ambrose, being warned of God, found them, and bestowed them in an hallowed and honourable place. They suffered at Milan upon the 19th day of June.\par
\switchcolumn*

\LA
\TeDeum{Te Deum}
\switchcolumn
\EN
\TeDeum{Te Deum}
\switchcolumn*

\end{paracol}

\PartTitle{Commune Sanctorum}{Common of the Saints}

\addcontentsline{toc}{chapter}{Commune Sanctorum\enspace\textit{Common of the Saints}}

\Office{Commune Evangelistarum tempore Paschali}{Commune Evangelistarum tempore Paschali}{}
\addcontentsline{toc}{subsection}{Commune Evangelistarum tempore Paschali\enspace\textit{Commune Evangelistarum tempore Paschali}}
\begin{paracol}{2}
\LA
\LessonHead{Lectio I}
\switchcolumn
\EN
\LessonHead{Lesson I}
\switchcolumn*

\LA
\LessonTitle{Incipit liber Ezechiélis Prophétæ}
\switchcolumn
\EN
\LessonTitle{Lesson from the book of Ezekiel}
\switchcolumn*

\LA
\Cite{Ezech 1:1-4}
\switchcolumn
\EN
\Cite{Ezek 1:1-4}
\switchcolumn*

\LA
\noindent Et factum est in trigésimo anno, in quarto, in quinta mensis, cum essem in médio captivórum juxta flúvium Chobar, apérti sunt cæli, et vidi visiónes Dei. In quinta mensis, ipse est annus quintus transmigratiónis regis Jóachim, factum est verbum Dómini ad Ezechiélem fílium Buzi sacerdótem, in terra Chaldæórum, secus flumen Chobar: et facta est super eum ibi manus Dómini. Et vidi, et ecce ventus túrbinis veniébat ab Aquilóne, et nubes magna, et ignis invólvens, et splendor in circúitu ejus: et de médio ejus, quasi spécies eléctri, id est, de médio ignis:\par
\switchcolumn
\EN
\noindent Now it came to pass in the thirtieth year, in the fourth month, on the fifth day of the month, when I was in the midst of the captives by the river Chobar, the heavens were opened, and I saw the visions of God. On the fifth day of the month, the same was the fifth year of the captivity of king Joachin, The word of the Lord came to Ezechiel the priest the son of Buzi in the land of the Chaldeans, by the river Chobar: and the hand of the Lord was there upon him. And I saw, and behold a whirlwind came out of the north: and a great cloud, and a fire infolding it, and brightness was about it: and out of the midst thereof, that is, out of the midst of the fire, as it were the resemblance of amber.\par
\switchcolumn*

\LA
\begin{Resp}\Rsign Beátus vir qui métuit Dóminum, allelúja: \Star{} In mandátis ejus cupit nimis, allelúja, allelúja, allelúja.\end{Resp}
\begin{Resp}\Vsign Glória et divítiæ in domo ejus, et justítia ejus manet in sǽculum sǽculi.\end{Resp}
\begin{Resp}\Rsign In mandátis ejus cupit nimis, allelúja, allelúja, allelúja.\end{Resp}
\switchcolumn
\EN
\begin{Resp}\Rsign Blessed is the man that feareth the Lord, alleluia. \Star{} He shall delight exceedingly in his commandments, alleluia, alleluia, alleluia.\end{Resp}
\begin{Resp}\Vsign Glory and wealth shall be in his house: and his justice remaineth for ever and ever.\end{Resp}
\begin{Resp}\Rsign He shall delight exceedingly in his commandments, alleluia, alleluia, alleluia.\end{Resp}
\switchcolumn*

\LA
\LessonHead{Lectio II}
\switchcolumn
\EN
\LessonHead{Lesson II}
\switchcolumn*

\LA
\Cite{Ezech 1:5-9}
\switchcolumn
\EN
\Cite{Ezek 1:5-9}
\switchcolumn*

\LA
\noindent Et in médio ejus similitúdo quátuor animálium, et hic aspéctus eórum, similitúdo hóminis in eis. Quátuor fácies uni, et quátuor pennæ uni. Pedes eórum, pedes recti, et planta pedis eórum quasi planta pedis vítuli: et scintíllæ quasi aspéctus æris candéntis. Et manus hóminis sub pennis eórum, in quátuor pártibus: et fácies et pennas per quátuor partes habébant, junctǽque erant pennæ eórum altérius ad álterum. Non revertebántur cum incéderent, sed unumquódque ante fáciem suam gradiebátur.\par
\switchcolumn
\EN
\noindent And in the midst thereof the likeness of four living creatures: and this was their appearance: there was the likeness of a man in them. Every one had four faces, and every one four wings. Their feet were straight feet, and the sole of their foot was like the sole of a calf's foot, and they sparkled like the appearance of glowing brass. And they had the hands of a man under their wings on their four sides: and they had faces, and wings on the four sides, And the wings of one were joined to the wings of another. They turned not when they went: but every one went straight forward.\par
\switchcolumn*

\LA
\begin{Resp}\Rsign Tristítia vestra, allelúja, \Star{} Convertétur in gáudium, allelúja, allelúja.\end{Resp}
\begin{Resp}\Vsign Mundus autem gaudébit, vos vero contristabímini, sed tristítia vestra.\end{Resp}
\begin{Resp}\Rsign Convertétur in gáudium, allelúja, allelúja.\end{Resp}
\switchcolumn
\EN
\begin{Resp}\Rsign Your sorrow, alleluia. \Star{} Shall be turned into joy, alleluia.\end{Resp}
\begin{Resp}\Vsign The world shall rejoice: and you shall be made sorrowful, but your sorrow.\end{Resp}
\begin{Resp}\Rsign Shall be turned into joy, alleluia.\end{Resp}
\switchcolumn*

\LA
\LessonHead{Lectio III}
\switchcolumn
\EN
\LessonHead{Lesson III}
\switchcolumn*

\LA
\Cite{Ezech 1:10-12}
\switchcolumn
\EN
\Cite{Ezek 1:10-12}
\switchcolumn*

\LA
\noindent Similitúdo autem vultus eórum, fácies hóminis et fácies leónis a dextris ipsórum quátuor, fácies autem bovis a sinístris ipsórum quátuor, et fácies áquilæ désuper ipsórum quátuor. Fácies eórum et pennæ eórum exténtæ désuper: duæ pennæ singulórum jungebántur, et duæ tegébant córpora eórum: et unumquódque eórum coram fácie sua ambulábat: ubi erat ímpetus spíritus, illuc gradiebántur, nec revertebántur cum ambulárent.\par
\switchcolumn
\EN
\noindent And as for the likeness of their faces: there was the face of a man, and the face of a lion on the right side of all the four: and the face of an ox, on the left side of all the four: and the face of an eagle over all the four. And their faces, and their wings were stretched upward: two wings of every one were joined, and two covered their bodies: And every one of them went straight forward: whither the impulse of the spirit was to go, thither they went: and they turned not when they went.\par
\switchcolumn*

\LA
\begin{Resp}\Rsign Pretiósa in conspéctu Dómini, allelúja, \Star{} Mors Sanctórum ejus, allelúja.\end{Resp}
\begin{Resp}\Vsign Custódit Dóminus ómnia ossa eórum, unum ex his non conterétur.\end{Resp}
\begin{Resp}\Rsign Mors Sanctórum ejus, allelúja.\end{Resp}
\begin{Resp}\Vsign Glória Patri, et Fílio, \Star{} et Spirítui Sancto.\end{Resp}
\begin{Resp}\Rsign Mors Sanctórum ejus, allelúja.\end{Resp}
\switchcolumn
\EN
\begin{Resp}\Rsign Precious in the sight of the Lord, alleluia. \Star{} Is the death of his saints, alleluia.\end{Resp}
\begin{Resp}\Vsign The Lord is nigh unto them that are of a contrite heart: and he will save the humble of spirit.\end{Resp}
\begin{Resp}\Rsign Is the death of his saints, alleluia.\end{Resp}
\begin{Resp}\Vsign Glory be to the Father, and to the Son, \Star{} and to the Holy Ghost.\end{Resp}
\begin{Resp}\Rsign Is the death of his saints, alleluia.\end{Resp}
\switchcolumn*

\LA
\LessonHead{Lectio VII}
\switchcolumn
\EN
\LessonHead{Lesson VII}
\switchcolumn*

\LA
\LessonTitle{Léctio sancti Evangélii secúndum Lucam}
\switchcolumn
\EN
\LessonTitle{From the Holy Gospel according to Luke}
\switchcolumn*

\LA
\Cite{Luc 10:1-9}
\switchcolumn
\EN
\Cite{Luke 10:1-9}
\switchcolumn*

\LA
\noindent In illo témpore: Designávit Dóminus et álios septuagínta duos: et misit illos binos ante fáciem suam, in omnem civitátem et locum, quo erat ipse ventúrus. Et réliqua.\par
\switchcolumn
\EN
\noindent At that time, the Lord appointed other seventy-two also, and sent them two and two before His face into every city and place, whither He Himself would come. And so on.\par
\switchcolumn*

\LA
\LessonTitle{Homilía sancti Gregórii Papæ}
\switchcolumn
\EN
\LessonTitle{Homily by Pope St. Gregory (the Great.)}
\switchcolumn*

\LA
\Source{Homilia 17 in Evangelia}
\switchcolumn
\EN
\Source{17th Homily on the Gospels}
\switchcolumn*

\LA
\noindent Dóminus et Salvátor noster, fratres caríssimi, aliquándo nos sermónibus, aliquándo vero opéribus ádmonet. Ipsa étenim facta ejus præcépta sunt: quia dum áliquid tácitus facit, quid ágere debeámus innotéscit. Ecce enim binos in prædicatiónem discípulos mittit: quia duo sunt præcépta caritátis, Dei vidélicet amor, et próximi: et minus quam inter duos cáritas habéri non potest. Nemo enim próprie ad semetípsum habére caritátem dícitur: sed diléctio in álterum tendit, ut cáritas esse possit.\par
\switchcolumn
\EN
\noindent Dearly beloved brethren, our Lord and Saviour doth sometimes admonish us by words, and sometimes by works. Yea, His very works do themselves teach us for that which He doth silently His example still moveth us to copy. Behold how He sendeth forth His disciples to preach by two and two since there are two commandments to love, that is, a commandment to love God, and a commandment to love our neighbour and where there are not two, the one, being alone, hath not whereon to do the Lord's commandment. And no man can properly be said to love himself: for love tendeth outward toward our neighbour, if it be the love whereto the Gospel doth oblige us.\par
\switchcolumn*

\LA
\begin{Resp}\Rsign Ego sum vitis vera, et vos pálmites: \Star{} Qui manet in me, et ego in eo, hic fert fructum multum, allelúja, allelúja.\end{Resp}
\begin{Resp}\Vsign Sicut diléxit me Pater, et ego diléxi vos.\end{Resp}
\begin{Resp}\Rsign Qui manet in me, et ego in eo, hic fert fructum multum, allelúja, allelúja.\end{Resp}
\switchcolumn
\EN
\begin{Resp}\Rsign I am the vine: you the branches. \Star{} He that abideth in me, and I in him, the same beareth much fruit, alleluia, alleluia.\end{Resp}
\begin{Resp}\Vsign As the Father hath loved me, I also have loved you.\end{Resp}
\begin{Resp}\Rsign He that abideth in me, and I in him, the same beareth much fruit, alleluia, alleluia.\end{Resp}
\begin{Resp}\Vsign Glory be to the Father, and to the Son, \Star{} and to the Holy Ghost.\end{Resp}
\begin{Resp}\Rsign He that abideth in me, and I in him, the same beareth much fruit, alleluia, alleluia.\end{Resp}
\switchcolumn*

\LA
\LessonHead{Lectio VIII}
\switchcolumn
\EN
\LessonHead{Lesson VIII}
\switchcolumn*

\LA
\noindent Ecce enim binos ad prædicándum discípulos Dóminus mittit: quátenus hoc nobis tácitus ínnuat, quia qui caritátem erga álterum non habet, prædicatiónis offícium suscípere nullátenus debet. Bene autem dícitur, quia misit eos ante fáciem suam in omnem civitátem et locum, quo erat ipse ventúrus. Prædicatóres enim suos Dóminus séquitur: quia prædicátio prǽvenit, et tunc ad mentis nostræ habitáculum Dóminus venit, quando verba exhortatiónis præcúrrunt: atque per hoc véritas in mente suscípitur.\par
\switchcolumn
\EN
\noindent Behold, the Lord sendeth forth His disciples to preach by two and two and thus doing, He doth silently teach us that whosoever loveth not his neighbour, such a one it behoveth not to take upon him the office of a preacher. Well also is it said that He sent them before His face into every city and place whither He Himself would come. The Lord followeth His preachers first cometh preaching, and then the Lord Himself cometh to the house of our mind, whither the word of exhortation hath come before and so cometh the truth into our mind.\par
\switchcolumn*

\LA
\begin{Resp}\Rsign Cándidi facti sunt Nazarǽi ejus, allelúja: splendórem Deo dedérunt, allelúja: \Star{} Et sicut lac coaguláti sunt, allelúja, allelúja.\end{Resp}
\begin{Resp}\Vsign Candidióres nive, nitidióres lacte, rubicundióres ébore antíquo, sapphíro pulchrióres.\end{Resp}
\begin{Resp}\Rsign Et sicut lac coaguláti sunt, allelúja, allelúja.\end{Resp}
\begin{Resp}\Vsign Glória Patri, et Fílio, \Star{} et Spirítui Sancto.\end{Resp}
\begin{Resp}\Rsign Et sicut lac coaguláti sunt, allelúja, allelúja.\end{Resp}
\switchcolumn
\EN
\begin{Resp}\Rsign Her Nazarites are become pure, Alleluia: they reflect the glory of God, Alleluia. \Star{} They are whiter than milk. Alleluia, Alleluia.\end{Resp}
\begin{Resp}\Vsign They are purer than snow, they are whiter than milk, they are more ruddy in body than coral, their polishing is of sapphire.\end{Resp}
\begin{Resp}\Rsign They are whiter than milk. Alleluia, Alleluia.\end{Resp}
\begin{Resp}\Vsign Glory be to the Father, and to the Son, \Star{} and to the Holy Ghost.\end{Resp}
\begin{Resp}\Rsign They are whiter than milk. Alleluia, Alleluia.\end{Resp}
\switchcolumn*

\LA
\LessonHead{Lectio IX}
\switchcolumn
\EN
\LessonHead{Lesson IX}
\switchcolumn*

\LA
\noindent Hinc namque eísdem prædicatóribus Isaías dicit: Paráte viam Dómini, rectas fácite sémitas Dei nostri. Hinc fíliis Psalmísta ait: Iter fácite ei, qui ascéndit super occásum. Super occásum namque Dóminus ascéndit: quia unde in passióne occúbuit, inde majórem suam glóriam resurgéndo manifestávit. Super occásum vidélicet ascéndit; quia mortem quam pértulit, resurgéndo calcávit. Ei ergo qui ascéndit super occásum, iter fácimus, cum nos ejus glóriam vestris méntibus prædicámus, ut eas et ipse post véniens, per amóris sui præséntiam illústret.\par
\switchcolumn
\EN
\noindent Therefore, to preachers saith Isaiah: “Prepare ye the way of the Lord, make straight an highway for our God.” (xl. 3.) And again the Psalmist saith: “Spread a path before him that rideth upon the West.” (lxvii. 4.) The Lord rideth upon the West above that from which in death He veiled His glory, hath He royally exalted that glory that excelleth, even the glory of His rising again. He rideth upon the West, Who, being risen again from the dead, is throned high above the death to which He bowed. Before Him, therefore, That rideth upon the West, we spread a path, when we set forth His glory before the eyes of your mind, to the end that He Himself may come after, and Himself enlighten the same your minds by His presence and His love.\par
\switchcolumn*

\LA
\TeDeum{Te Deum}
\switchcolumn
\EN
\TeDeum{Te Deum}
\switchcolumn*

\end{paracol}

\Office{Commune unius Martyris Tempore Paschali}{Common of a Single Martyr in Paschaltide}{}
\addcontentsline{toc}{subsection}{Commune unius Martyris Tempore Paschali\enspace\textit{Common of a Single Martyr in Paschaltide}}
\begin{paracol}{2}
\LA
\LessonHead{Lectio VII}
\switchcolumn
\EN
\LessonHead{Lesson VII}
\switchcolumn*

\LA
\LessonTitle{Léctio sancti Evangélii secúndum Joánnem}
\switchcolumn
\EN
\LessonTitle{From the holy Gospel according to John}
\switchcolumn*

\LA
\Cite{Joannes 15:1-7}
\switchcolumn
\EN
\Cite{John 15:1-7}
\switchcolumn*

\LA
\noindent In illo témpore: Dixit Jesus discípulis suis: Ego sum vitis vera: et Pater meus agrícola est. Et réliqua.\par
\switchcolumn
\EN
\noindent In that time Jesus said to his disciples: I an the true vine; and my Father is the husbandman. And so on.\par
\switchcolumn*

\LA
\LessonTitle{Homilía sancti Augustíni Epíscopi}
\switchcolumn
\EN
\LessonTitle{Homily by St. Augustine, Bishop (of Hippo.)}
\switchcolumn*

\LA
\Source{Tract. 80 in Joannem}
\switchcolumn
\EN
\Source{Tract 80 on John}
\switchcolumn*

\LA
\noindent Iste locus evangélicus, fratres, ubi se dicit Dóminus vitem, et discípulos suos pálmites, secúndum hoc dicit, quod est caput Ecclésiæ, nosque membra ejus, mediátor Dei et hóminum, homo Christus Jesus. Uníus quippe natúræ sunt vitis et pálmites. Propter quod cum esset Deus, cujus natúræ non sumus, factus est homo, ut in illo esset vitis humána natúra, cujus et nos hómines pálmites esse possémus.\par
\switchcolumn
\EN
\noindent Dearly beloved brethren, this passage of the Gospel, wherein the Lord saith that He is the vine, and that His disciples are the branches, is to be taken in that sense wherein it is also said, that He is the Head of the Church, (Eph. v. 23,) and that we are the members of Him (30) Who is the Mediator between God and men, the man Christ Jesus (1 Tim. ii. 5.) The vine and his branches are of one and the same nature. Therefore, seeing that He was God, of which nature we are not, He was made man, to the end that He might have in Himself this vine, that is, the manhood, whereof we men can be made branches.\par
\switchcolumn*

\LA
\begin{Resp}\Rsign Ego sum vitis vera, et vos pálmites: \Star{} Qui manet in me, et ego in eo, hic fert fructum multum, allelúja, allelúja.\end{Resp}
\begin{Resp}\Vsign Sicut diléxit me Pater, et ego diléxi vos.\end{Resp}
\begin{Resp}\Rsign Qui manet in me, et ego in eo, hic fert fructum multum, allelúja, allelúja.\end{Resp}
\switchcolumn
\EN
\begin{Resp}\Rsign I am the vine: you the branches. \Star{} He that abideth in me, and I in him, the same beareth much fruit, alleluia, alleluia.\end{Resp}
\begin{Resp}\Vsign As the Father hath loved me, I also have loved you.\end{Resp}
\begin{Resp}\Rsign He that abideth in me, and I in him, the same beareth much fruit, alleluia, alleluia.\end{Resp}
\begin{Resp}\Vsign Glory be to the Father, and to the Son, \Star{} and to the Holy Ghost.\end{Resp}
\begin{Resp}\Rsign He that abideth in me, and I in him, the same beareth much fruit, alleluia, alleluia.\end{Resp}
\switchcolumn*

\LA
\LessonHead{Lectio VIII}
\switchcolumn
\EN
\LessonHead{Lesson VIII}
\switchcolumn*

\LA
\noindent Quid ergo est, Ego sum vitis vera? Numquid ut ádderet, vera, hoc ad eam vitem rétulit, unde ista similitúdo transláta est? Sic enim dícitur vitis per similitúdinem, non per proprietátem: quemádmodum dícitur ovis, agnus, leo, petra, lapis anguláris, et cétera hujúsmodi, quæ magis ipsa sunt vera, ex quibus ducúntur istæ similitúdines, non proprietátes. Sed cum dicit, Ego sum vitis vera: ab illa se útique discérnit, cui dícitur: Quómodo convérsa es in amaritúdinem vitis aliéna? Nam quo pacto est vitis vera, quæ exspectáta est ut fáceret uvam, fecit autem spinas?\par
\switchcolumn
\EN
\noindent Why saith He: I am the true vine? As touching this word true, hath He not here regard to that other parable of a vine, the like figure whereto He doth here apply to Himself? (Jer. ii. 21.) Here is He called a vine, not plainly, but in parable, as also He is called elsewhere a sheep, (Isa. liii. 7, Acts viii. 32,) a lamb, (John i. 36,) a lion, (Apoc. v. 5,) a rock, (1 Cor. x. 4,) a corner-stone, \textasciitilde{}(Eph. ii. 20,) and other things of the like kind. But these things are in themselves that which they seem to be, albeit He is called by their names, not plainly, but in a parable, and herein are they different from that vine, whereof in this place He taketh on Him the name. For when He saith: I am the true vine, doth He not make distinction between Himself, and that which indeed seemed to be a vine, but to which it is said: How art thou turned into the degenerate plant of a strange vine unto Me? (Jer. ii. 21.) For by what title shall that plant be called other than a false vine, whereto they looked that she should bring forth grapes, and she brought forth thorns?\par
\switchcolumn*

\LA
\begin{Resp}\Rsign Cándidi facti sunt Nazarǽi ejus, allelúja: splendórem Deo dedérunt, allelúja: \Star{} Et sicut lac coaguláti sunt, allelúja, allelúja.\end{Resp}
\begin{Resp}\Vsign Candidióres nive, nitidióres lacte, rubicundióres ébore antíquo, sapphíro pulchrióres.\end{Resp}
\begin{Resp}\Rsign Et sicut lac coaguláti sunt, allelúja, allelúja.\end{Resp}
\begin{Resp}\Vsign Glória Patri, et Fílio, \Star{} et Spirítui Sancto.\end{Resp}
\begin{Resp}\Rsign Et sicut lac coaguláti sunt, allelúja, allelúja.\end{Resp}
\switchcolumn
\EN
\begin{Resp}\Rsign Her Nazarites are become pure, Alleluia: they reflect the glory of God, Alleluia. \Star{} They are whiter than milk. Alleluia, Alleluia.\end{Resp}
\begin{Resp}\Vsign They are purer than snow, they are whiter than milk, they are more ruddy in body than coral, their polishing is of sapphire.\end{Resp}
\begin{Resp}\Rsign They are whiter than milk. Alleluia, Alleluia.\end{Resp}
\begin{Resp}\Vsign Glory be to the Father, and to the Son, \Star{} and to the Holy Ghost.\end{Resp}
\begin{Resp}\Rsign They are whiter than milk. Alleluia, Alleluia.\end{Resp}
\switchcolumn*

\LA
\LessonHead{Lectio IX}
\switchcolumn
\EN
\LessonHead{Lesson IX}
\switchcolumn*

\LA
\noindent Ego sum, inquit, vitis vera: et Pater meus agrícola est. Numquid unum sunt agrícola et vitis? Secúndum hoc ergo vitis Christus, secúndum quod ait: Pater major me est. Secúndum autem id, quod ait: Ego, et Pater unum sumus; et ipse agrícola est: nec talis, quales sunt, qui extrínsecus operándo éxhibent ministérium: sed talis, ut det étiam intrínsecus increméntum. Nam neque qui plantat est áliquid, neque qui rigat: sed, qui increméntum dat, Deus. Sed útique Deus est Christus, quia Deus erat Verbum: unde ipse, et Pater unum sunt. Et, si Verbum caro factum est, quod non erat, manet quod erat.\par
\switchcolumn
\EN
\noindent He saith: I am the true vine, and My Father is the husbandman. Is the vine one with the husbandman? These words then are to be taken in that sense wherein He also saith: My Father is greater than I. (John xiv. 28.) In this sense is He the vine, and the Father is the husbandman. But again, in regard to those words: I and the Father are one, and again and: My Father is the husbandman, we understand that They are not the vine and the husbandman, after the manner of a vine, and the husbandman that from without doth care for and keep it, but after the manner of a vine and Him That from within doth make it to bring forth fruit. For neither is he that planteth anything, neither he that watereth, but God that giveth the increase. (1 Cor. iii. 7.) But Christ is God, for the Word was God. (John i. 1.) Therefore He and the Father are one: and, albeit the Word was made flesh, (John i. 14,) which, before, He was not, He ceased not to be still That Which He was.\par
\switchcolumn*

\LA
\TeDeum{Te Deum}
\switchcolumn
\EN
\TeDeum{Te Deum}
\switchcolumn*

\LA
\LessonHead{Lectio IV (in 2 loco)}
\switchcolumn
\EN
\LessonHead{Lesson IV (set 2)}
\switchcolumn*

\LA
\LessonTitle{De Epístola sancti Cypriáni Epíscopi et Mártyris ad Mártyres et Confessóres.}
\switchcolumn
\EN
\LessonTitle{From the Epistle of St Cyprian, Bishop [of Carthage,] and himself Martyr, to the Martyrs and Confessors}
\switchcolumn*

\LA
\Source{Lib. 2. Epist. 6.}
\switchcolumn
\EN
\Source{Bk. ii, ep. 6}
\switchcolumn*

\LA
\noindent Quibus ego vos láudibus prǽdicem, fortíssimi Mártyres? robur péctoris vestri, et perseverántiam fídei, quo præcónio vocis exórnem? Tolerástis usque ad consummatiónem glóriæ duríssimam quæstiónem: nec cessístis supplíciis, sed vobis pótius supplícia cessérunt. Finem dolóribus, quem torménta non dabant, corónæ dedérunt. Laniéna grávior ad hoc diu perseverávit, non ut stantem fidem dejíceret, sed ut hómines Dei ad Deum velócius mítteret.\par
\switchcolumn
\EN
\noindent How shall I praise you, Martyrs so brilliantly victorious? Can the voice of man's praise add anything to the glory of your manful heart and unshaken faithfulness? Ye have borne all the hardness of the torment, and have attained unto the excellent height of glory: the tormentors have not worn you out, nay, ye rather have worn out the tormentors. When they that kill the body would give you no rest from suffering, ye suffered until ye gained the crown. And the torment waxing still more dread, waxed not to the casting down of your strong faith, but to the sooner sending God's men home to God.\par
\switchcolumn*

\LA
\begin{Resp}\Rsign Lux perpétua lucébit Sanctis tuis, Dómine, \Star{} Et ætérnitas témporum, allelúja, allelúja.\end{Resp}
\begin{Resp}\Vsign Lætítia sempitérna erit super cápita eórum: gáudium et exsultatiónem obtinébunt.\end{Resp}
\begin{Resp}\Rsign Et ætérnitas témporum, allelúja, allelúja.\end{Resp}
\switchcolumn
\EN
\begin{Resp}\Rsign Eternal light will shine over your saints, O Lord, \Star{} And the eternity of the times, alleluia, alleluia.\end{Resp}
\begin{Resp}\Vsign Everlasting joy shall be upon their heads, they shall obtain joy and gladness\end{Resp}
\begin{Resp}\Rsign And the eternity of the times, alleluia, alleluia.\end{Resp}
\switchcolumn*

\LA
\LessonHead{Lectio V (in 2 loco)}
\switchcolumn
\EN
\LessonHead{Lesson V (set 2)}
\switchcolumn*

\LA
\noindent Vidit admírans præséntium multitúdo cæléste certamen, certámen Dei, certámen spiritále, prǽlium Christi: stetísse servos ejus voce líbera, mente incorrúpta, virtúte divína, telis quidem sæculáribus nudos, sed armis fídei ardéntis armátos. Stetérunt torti torquéntibus fortióres: et pulsántes ac laniántes úngulas, pulsáta ac laniáta membra vicérunt. Inexpugnábilem fidem superáre non pótuit sǽviens diu plaga repetíta, quamvis rupta compáge víscerum torqueréntur in servis Dei jam non membra, sed vúlnera. Fluébat sanguis, qui incéndium persecutiónis exstíngueret, qui flammas et ignes gehénnæ glorióso cruóre sopíret.\par
\switchcolumn
\EN
\noindent They that stood by looked in wonder at your heavenly conflict, that battle of God, that wrestling of spirit, that combat of Christ. There they saw His servants standing with voice unshaken, with spirit unbroken, strong in God's strength, naked indeed, as to the arms of this world, but clothed on with the armor of God, and equipped with the fiery weapons of faith. There they that were tormented stood braver than they that tormented them. Their bruised and mangled bodies overcame the instruments of cruelty that bruised and mangled them. The bloody stripes, so often laid on, could not beat down the impregnable castle of their faith, even when the covering of their bowels was broken, and that which was tormented in God's servants was no longer limbs but wounds. The blood that ran down, ran down to quench the rage of persecution, noble blood, that can put out the flames and fire of hell.\par
\switchcolumn*

\LA
\begin{Resp}\Rsign In servis suis, allelúja, \Star{} Consolábitur Deus, allelúja.\end{Resp}
\begin{Resp}\Vsign Judicábit Dóminus pópulum suum, et in servis suis.\end{Resp}
\begin{Resp}\Rsign Consolábitur Deus, allelúja.\end{Resp}
\switchcolumn
\EN
\begin{Resp}\Rsign His servants, alleluia \Star{} God will comfort, alleluia, alleluia.\end{Resp}
\begin{Resp}\Vsign The Lord will judge His people and, His servants,\end{Resp}
\begin{Resp}\Rsign God will comfort, alleluia, alleluia.\end{Resp}
\switchcolumn*

\LA
\LessonHead{Lectio VI (in 2 loco)}
\switchcolumn
\EN
\LessonHead{Lesson VI (set 2)}
\switchcolumn*

\LA
\noindent O quale illud fuit spectáculum Dómino, quam sublíme, quam magnum, quam Dei óculis sacraménto ac devotióne mílitis ejus accéptum: sicut scriptum est in Psalmis, Spíritu Sancto loquénte ad nos páriter et monénte: Pretiósa est in conspéctu Dómini mors justórum ejus. Pretiósa mors hæc est, quæ emit immortalitátem prétio sui sánguinis; quæ accépit corónam de consummatióne virtútis. Quam lætus illic Christus fuit, quam libens in tálibus servis suis et pugnávit, et vicit protéctor fídei, dans credéntibus tantum, quantum se credit cápere qui sumit! Certámini suo ádfuit, præliatóres atque assertóres sui nóminis eréxit, corroborávit, animávit. Et qui pro nobis mortem semel vicit, semper vincit in nobis.\par
\switchcolumn
\EN
\noindent O what a spectacle was that in the eyes of the Lord! O how noble! O how mighty! O how precious in the sight of God were His soldiers' loyalty and faithfulness! Even as it is written in the Psalms, the Holy Ghost therein at once speaking to us and warning us, “Precious in the sight of the Lord is the death of His Saints.” O what a precious death is his, who maketh purchase of life that can never die, at the price of his own blood, and seizeth on the crown, when courage hath no more left to meet! O how joyful was Christ! How gladly fought He in such servants as these, in these how gladly did He triumph, the Keeper of their faith, and, in the end, to them how gladly did He give that reward which no man knoweth saving he that receiveth it! (Apoc. ii. 17.) He it was, Who was there when they fought, He it was, Who raised them up to be the champions and defenders of His holy Name, He, who gave them the strength, He, Who nerved them. He, That by death hath once conquered for us, liveth now for ever to conquer in us.\par
\switchcolumn*

\LA
\begin{Resp}\Rsign Fíliæ Jerúsalem, veníte et vidéte Mártyres cum corónis, quibus coronávit eos Dóminus \Star{} In die solemnitátis et lætítiæ, allelúja.\end{Resp}
\begin{Resp}\Vsign Quóniam confortávit seras portárum tuárum, benedíxit fílios tuos in te.\end{Resp}
\begin{Resp}\Rsign In die solemnitátis et lætítiæ, allelúja.\end{Resp}
\begin{Resp}\Vsign Glória Patri, et Fílio, \Star{} et Spirítui Sancto.\end{Resp}
\begin{Resp}\Rsign In die solemnitátis et lætítiæ, allelúja.\end{Resp}
\switchcolumn
\EN
\begin{Resp}\Rsign Come forth, O ye daughters of Jerusalem, and behold the Martyrs with the crowns wherewith the Lord crowned them; \Star{} In the day of His feasting, and of His gladness. Alleluia.\end{Resp}
\begin{Resp}\Vsign For He hath strengthened the bars of thy gates He hath blessed thy children within thee.\end{Resp}
\begin{Resp}\Rsign In the day of His feasting, and of His gladness. Alleluia.\end{Resp}
\begin{Resp}\Vsign Glory be to the Father, and to the Son, \Star{} and to the Holy Ghost.\end{Resp}
\begin{Resp}\Rsign In the day of His feasting, and of His gladness. Alleluia.\end{Resp}
\switchcolumn*

\end{paracol}

\Office{Commune unius Martyris Tempore Paschali}{Common of a Single Martyr in Paschaltide}{}
\addcontentsline{toc}{subsection}{Commune unius Martyris Tempore Paschali\enspace\textit{Common of a Single Martyr in Paschaltide}}
\begin{paracol}{2}
\LA
\LessonHead{Lectio VII}
\switchcolumn
\EN
\LessonHead{Lesson VII}
\switchcolumn*

\LA
\LessonTitle{Léctio sancti Evangélii secúndum Joánnem}
\switchcolumn
\EN
\LessonTitle{From the holy Gospel according to John}
\switchcolumn*

\LA
\Cite{Joannes 15:1-7}
\switchcolumn
\EN
\Cite{John 15:1-7}
\switchcolumn*

\LA
\noindent In illo témpore: Dixit Jesus discípulis suis: Ego sum vitis vera: et Pater meus agrícola est. Et réliqua.\par
\switchcolumn
\EN
\noindent In that time Jesus said to his disciples: I an the true vine; and my Father is the husbandman. And so on.\par
\switchcolumn*

\LA
\LessonTitle{Homilía sancti Augustíni Epíscopi}
\switchcolumn
\EN
\LessonTitle{Homily by St. Augustine, Bishop (of Hippo.)}
\switchcolumn*

\LA
\Source{Tract. 80 in Joannem}
\switchcolumn
\EN
\Source{Tract 80 on John}
\switchcolumn*

\LA
\noindent Iste locus evangélicus, fratres, ubi se dicit Dóminus vitem, et discípulos suos pálmites, secúndum hoc dicit, quod est caput Ecclésiæ, nosque membra ejus, mediátor Dei et hóminum, homo Christus Jesus. Uníus quippe natúræ sunt vitis et pálmites. Propter quod cum esset Deus, cujus natúræ non sumus, factus est homo, ut in illo esset vitis humána natúra, cujus et nos hómines pálmites esse possémus.\par
\switchcolumn
\EN
\noindent Dearly beloved brethren, this passage of the Gospel, wherein the Lord saith that He is the vine, and that His disciples are the branches, is to be taken in that sense wherein it is also said, that He is the Head of the Church, (Eph. v. 23,) and that we are the members of Him (30) Who is the Mediator between God and men, the man Christ Jesus (1 Tim. ii. 5.) The vine and his branches are of one and the same nature. Therefore, seeing that He was God, of which nature we are not, He was made man, to the end that He might have in Himself this vine, that is, the manhood, whereof we men can be made branches.\par
\switchcolumn*

\LA
\begin{Resp}\Rsign Ego sum vitis vera, et vos pálmites: \Star{} Qui manet in me, et ego in eo, hic fert fructum multum, allelúja, allelúja.\end{Resp}
\begin{Resp}\Vsign Sicut diléxit me Pater, et ego diléxi vos.\end{Resp}
\begin{Resp}\Rsign Qui manet in me, et ego in eo, hic fert fructum multum, allelúja, allelúja.\end{Resp}
\switchcolumn
\EN
\begin{Resp}\Rsign I am the vine: you the branches. \Star{} He that abideth in me, and I in him, the same beareth much fruit, alleluia, alleluia.\end{Resp}
\begin{Resp}\Vsign As the Father hath loved me, I also have loved you.\end{Resp}
\begin{Resp}\Rsign He that abideth in me, and I in him, the same beareth much fruit, alleluia, alleluia.\end{Resp}
\begin{Resp}\Vsign Glory be to the Father, and to the Son, \Star{} and to the Holy Ghost.\end{Resp}
\begin{Resp}\Rsign He that abideth in me, and I in him, the same beareth much fruit, alleluia, alleluia.\end{Resp}
\switchcolumn*

\LA
\LessonHead{Lectio VIII}
\switchcolumn
\EN
\LessonHead{Lesson VIII}
\switchcolumn*

\LA
\noindent Quid ergo est, Ego sum vitis vera? Numquid ut ádderet, vera, hoc ad eam vitem rétulit, unde ista similitúdo transláta est? Sic enim dícitur vitis per similitúdinem, non per proprietátem: quemádmodum dícitur ovis, agnus, leo, petra, lapis anguláris, et cétera hujúsmodi, quæ magis ipsa sunt vera, ex quibus ducúntur istæ similitúdines, non proprietátes. Sed cum dicit, Ego sum vitis vera: ab illa se útique discérnit, cui dícitur: Quómodo convérsa es in amaritúdinem vitis aliéna? Nam quo pacto est vitis vera, quæ exspectáta est ut fáceret uvam, fecit autem spinas?\par
\switchcolumn
\EN
\noindent Why saith He: I am the true vine? As touching this word true, hath He not here regard to that other parable of a vine, the like figure whereto He doth here apply to Himself? (Jer. ii. 21.) Here is He called a vine, not plainly, but in parable, as also He is called elsewhere a sheep, (Isa. liii. 7, Acts viii. 32,) a lamb, (John i. 36,) a lion, (Apoc. v. 5,) a rock, (1 Cor. x. 4,) a corner-stone, \textasciitilde{}(Eph. ii. 20,) and other things of the like kind. But these things are in themselves that which they seem to be, albeit He is called by their names, not plainly, but in a parable, and herein are they different from that vine, whereof in this place He taketh on Him the name. For when He saith: I am the true vine, doth He not make distinction between Himself, and that which indeed seemed to be a vine, but to which it is said: How art thou turned into the degenerate plant of a strange vine unto Me? (Jer. ii. 21.) For by what title shall that plant be called other than a false vine, whereto they looked that she should bring forth grapes, and she brought forth thorns?\par
\switchcolumn*

\LA
\begin{Resp}\Rsign Cándidi facti sunt Nazarǽi ejus, allelúja: splendórem Deo dedérunt, allelúja: \Star{} Et sicut lac coaguláti sunt, allelúja, allelúja.\end{Resp}
\begin{Resp}\Vsign Candidióres nive, nitidióres lacte, rubicundióres ébore antíquo, sapphíro pulchrióres.\end{Resp}
\begin{Resp}\Rsign Et sicut lac coaguláti sunt, allelúja, allelúja.\end{Resp}
\begin{Resp}\Vsign Glória Patri, et Fílio, \Star{} et Spirítui Sancto.\end{Resp}
\begin{Resp}\Rsign Et sicut lac coaguláti sunt, allelúja, allelúja.\end{Resp}
\switchcolumn
\EN
\begin{Resp}\Rsign Her Nazarites are become pure, Alleluia: they reflect the glory of God, Alleluia. \Star{} They are whiter than milk. Alleluia, Alleluia.\end{Resp}
\begin{Resp}\Vsign They are purer than snow, they are whiter than milk, they are more ruddy in body than coral, their polishing is of sapphire.\end{Resp}
\begin{Resp}\Rsign They are whiter than milk. Alleluia, Alleluia.\end{Resp}
\begin{Resp}\Vsign Glory be to the Father, and to the Son, \Star{} and to the Holy Ghost.\end{Resp}
\begin{Resp}\Rsign They are whiter than milk. Alleluia, Alleluia.\end{Resp}
\switchcolumn*

\LA
\LessonHead{Lectio IX}
\switchcolumn
\EN
\LessonHead{Lesson IX}
\switchcolumn*

\LA
\noindent Ego sum, inquit, vitis vera: et Pater meus agrícola est. Numquid unum sunt agrícola et vitis? Secúndum hoc ergo vitis Christus, secúndum quod ait: Pater major me est. Secúndum autem id, quod ait: Ego, et Pater unum sumus; et ipse agrícola est: nec talis, quales sunt, qui extrínsecus operándo éxhibent ministérium: sed talis, ut det étiam intrínsecus increméntum. Nam neque qui plantat est áliquid, neque qui rigat: sed, qui increméntum dat, Deus. Sed útique Deus est Christus, quia Deus erat Verbum: unde ipse, et Pater unum sunt. Et, si Verbum caro factum est, quod non erat, manet quod erat.\par
\switchcolumn
\EN
\noindent He saith: I am the true vine, and My Father is the husbandman. Is the vine one with the husbandman? These words then are to be taken in that sense wherein He also saith: My Father is greater than I. (John xiv. 28.) In this sense is He the vine, and the Father is the husbandman. But again, in regard to those words: I and the Father are one, and again and: My Father is the husbandman, we understand that They are not the vine and the husbandman, after the manner of a vine, and the husbandman that from without doth care for and keep it, but after the manner of a vine and Him That from within doth make it to bring forth fruit. For neither is he that planteth anything, neither he that watereth, but God that giveth the increase. (1 Cor. iii. 7.) But Christ is God, for the Word was God. (John i. 1.) Therefore He and the Father are one: and, albeit the Word was made flesh, (John i. 14,) which, before, He was not, He ceased not to be still That Which He was.\par
\switchcolumn*

\LA
\TeDeum{Te Deum}
\switchcolumn
\EN
\TeDeum{Te Deum}
\switchcolumn*

\end{paracol}

\Office{Commune Plurimorum Martyrum Pontificum}{Commune Plurimorum Martyrum Pontificum}{}
\addcontentsline{toc}{subsection}{Commune Plurimorum Martyrum Pontificum\enspace\textit{Commune Plurimorum Martyrum Pontificum}}
\begin{paracol}{2}
\LA
\LessonHead{Lectio I}
\switchcolumn
\EN
\LessonHead{Lesson I}
\switchcolumn*

\LA
\LessonTitle{De Epístola beáti Pauli Apóstoli ad Romános}
\switchcolumn
\EN
\LessonTitle{Lesson from the letter of St Paul the Apostle to the Romans}
\switchcolumn*

\LA
\Cite{Rom 8:12-19}
\switchcolumn
\EN
\Cite{Rom 8:12-19}
\switchcolumn*

\LA
\noindent Fratres: Debitóres sumus non carni, ut secúndum carnem vivámus. Si enim secúndum carnem vixéritis, moriémini: si autem spíritu facta carnis mortificavéritis, vivétis. Quicúmque enim spíritu Dei agúntur, ii sunt fílii Dei. Non enim accepístis spíritum servitútis íterum in timóre, sed accepístis spíritum adoptiónis filiórum, in quo clamámus: Abba (Pater). Ipse enim Spíritus testimónium reddit spirítui nostro quod sumus fílii Dei. Si autem fílii, et herédes: herédes, quidem Dei, coherédes autem Christi: si tamen compátimur ut et conglorificémur. Exístimo enim quod non sunt condígnæ passiónes hujus témporis ad futúram glóriam, quæ revelábitur in nobis. Nam exspectátio creatúræ revelatiónem filiórum Dei exspéctat.\par
\switchcolumn
\EN
\noindent Therefore, brethren, we are debtors, not to the flesh, to live according to the flesh. For if you live according to the flesh, you shall die: but if by the Spirit you mortify the deeds of the flesh, you shall live. For whosoever are led by the Spirit of God, they are the sons of God. For you have not received the spirit of bondage again in fear; but you have received the spirit of adoption of sons, whereby we cry: Abba, Father. For the Spirit himself giveth testimony to our spirit, that we are the sons of God. And if sons, heirs also; heirs indeed of God, and joint heirs with Christ: yet so, if we suffer with him, that we may be also glorified with him. For I reckon that the sufferings of this time are not worthy to be compared with the glory to come, that shall be revealed in us. For the expectation of the creature waiteth for the revelation of the sons of God\par
\switchcolumn*

\LA
\begin{Resp}\Rsign Abstérget Deus omnem lácrimam ab óculis Sanctórum: et jam non erit ámplius neque luctus, neque clamor, sed nec ullus dolor: \Star{} Quóniam prióra transiérunt.\end{Resp}
\begin{Resp}\Vsign Non esúrient, neque sítient ámplius, neque cadet super illos sol, neque ullus æstus.\end{Resp}
\begin{Resp}\Rsign Quóniam prióra transiérunt.\end{Resp}
\switchcolumn
\EN
\begin{Resp}\Rsign God shall wipe away all tears from the eyes of His Saints, and there shall be no more sorrow, there shall be no more death, neither sorrow, nor crying, neither shall there be any more pain; \Star{} For the former things are passed away.\end{Resp}
\begin{Resp}\Vsign They shall hunger no more, neither thirst any more, neither shall the sun light on them, nor any heat.\end{Resp}
\begin{Resp}\Rsign For the former things are passed away.\end{Resp}
\switchcolumn*

\LA
\LessonHead{Lectio II}
\switchcolumn
\EN
\LessonHead{Lesson II}
\switchcolumn*

\LA
\Cite{Rom 8:28-34}
\switchcolumn
\EN
\Cite{Rom 8:28-34}
\switchcolumn*

\LA
\noindent Scimus autem quóniam diligéntibus Deum ómnia cooperántur in bonum iis qui secúndum propósitum vocáti sunt sancti. Nam quos præscívit et prædestinávit confórmes fíeri imáginis Fílii sui ut sit ipse primogénitus in multis frátribus. Quos autem prædestinávit, hos et vocávit: et quos vocávit, hos et justificávit: quos autem justificávit, illos et glorificávit. Quid ergo dicémus ad hæc? Si Deus pro nobis, quis contra nos? Qui étiam próprio Fílio suo non pepércit, sed pro nobis ómnibus trádidit illum, quómodo non étiam cum illo ómnia nobis donávit? Quis accusábit advérsus eléctos Dei? Deus qui justíficat, quis est qui condémnet? Christus Jesus, qui mórtuus est, immo qui et resurréxit, qui est ad déxteram Dei, qui étiam interpéllat pro nobis.\par
\switchcolumn
\EN
\noindent And we know that to them that love God, all things work together unto good, to such as, according to his purpose, are called to be saints. For whom he foreknew, he also predestinated to be made conformable to the image of his Son; that he might be the firstborn amongst many brethren. And whom he predestinated, them he also called. And whom he called, them he also justified. And whom he justified, them he also glorified. What shall we then say to these things? If God be for us, who is against us? He that spared not even his own Son, but delivered him up for us all, how hath he not also, with him, given us all things? Who shall accuse against the elect of God? God that justifieth. Who is he that shall condemn? Christ Jesus that died, yea that is risen also again; who is at the right hand of God, who also maketh intercession for us.\par
\switchcolumn*

\LA
\begin{Resp}\Rsign Viri sancti gloriósum sánguinem fudérunt pro Dómino, amavérunt Christum in vita sua, imitáti sunt eum in morte sua: \Star{} Et ídeo corónas triumpháles meruérunt.\end{Resp}
\begin{Resp}\Vsign Unus spíritus, et una fides erat in eis.\end{Resp}
\begin{Resp}\Rsign Et ídeo corónas triumpháles meruérunt.\end{Resp}
\switchcolumn
\EN
\begin{Resp}\Rsign These men are holy, who have gloriously shed their blood for the Lord's sake, yea, who loved Christ in their lives, and were made like unto Him in their flesh, \Star{} And therefore they have earned crowns of victory.\end{Resp}
\begin{Resp}\Vsign One spirit, and one faith was in them.\end{Resp}
\begin{Resp}\Rsign And therefore they have earned crowns of victory.\end{Resp}
\switchcolumn*

\LA
\LessonHead{Lectio III}
\switchcolumn
\EN
\LessonHead{Lesson III}
\switchcolumn*

\LA
\Cite{Rom 8:35-39}
\switchcolumn
\EN
\Cite{Rom 8:35-39}
\switchcolumn*

\LA
\noindent Quis ergo nos separábit a caritáte Christi? tribulátio? an angústia? an fames? an núditas? an perículum? an persecútio? an gládius? (sicut scriptum est: Quia propter te mortificámur tota die: æstimáti sumus sicut oves occisiónis.) Sed in his ómnibus superámus propter eum qui diléxit nos. Certus sum enim quia neque mors, neque vita, neque Ángeli, neque Principátus, neque Virtútes, neque instántia, neque futúra, neque fortitúdo, neque altitúdo, neque profúndum, neque creatúra ália póterit nos separáre a caritáte Dei, quæ est in Christo Jesu Dómino nostro.\par
\switchcolumn
\EN
\noindent Who then shall separate us from the love of Christ? Shall tribulation? or distress? or famine? or nakedness? or danger? or persecution? or the sword? As it is written: For thy sake we are put to death all the day long. We are accounted as sheep for the slaughter. But in all these things we overcome, because of him that hath loved us. For I am sure that neither death, nor life, nor angels, nor principalities, nor powers, nor things present, nor things to come, nor might, Nor height, nor depth, nor any other creature, shall be able to separate us from the love of God, which is in Christ Jesus our Lord.\par
\switchcolumn*

\LA
\begin{Resp}\Rsign Tradidérunt córpora sua propter Deum ad supplícia: \Star{} Et meruérunt habére corónas perpétuas.\end{Resp}
\begin{Resp}\Vsign Isti sunt qui venérunt ex magna tribulatióne, et lavérunt stolas suas in sánguine Agni.\end{Resp}
\begin{Resp}\Rsign Et meruérunt habére corónas perpétuas.\end{Resp}
\begin{Resp}\Vsign Glória Patri, et Fílio, \Star{} et Spirítui Sancto.\end{Resp}
\begin{Resp}\Rsign Et meruérunt habére corónas perpétuas.\end{Resp}
\switchcolumn
\EN
\begin{Resp}\Rsign They gave their bodies for God’s sake to death \Star{} And gained the everlasting crown.\end{Resp}
\begin{Resp}\Vsign These are they which came out of great tribulation, and have washed their robes in the Blood of the Lamb.\end{Resp}
\begin{Resp}\Rsign And gained the everlasting crown.\end{Resp}
\begin{Resp}\Vsign Glory be to the Father, and to the Son, \Star{} and to the Holy Ghost.\end{Resp}
\begin{Resp}\Rsign And gained the everlasting crown.\end{Resp}
\switchcolumn*

\end{paracol}

\Office{Commune plurium Summorum Pontificum Martyrum Tempore Paschali}{Commune plurium Summorum Pontificum Martyrum Tempore Paschali}{}
\addcontentsline{toc}{subsection}{Commune plurium Summorum Pontificum Martyrum Tempore Paschali\enspace\textit{Commune plurium Summorum Pontificum Martyrum Tempore Paschali}}
\begin{paracol}{2}
\LA
\LessonHead{Lectio VII}
\switchcolumn
\EN
\LessonHead{Lesson VII}
\switchcolumn*

\LA
\LessonTitle{Léctio sancti Evangélii secúndum Matthǽum}
\switchcolumn
\EN
\LessonTitle{From the Holy Gospel according to Matthew}
\switchcolumn*

\LA
\Cite{Matt 16:13-19}
\switchcolumn
\EN
\Cite{Matt 16:13-19}
\switchcolumn*

\LA
\noindent In illo témpore: Venit Jesus in partes Cæsaréæ Philíppi, et interrogábat discípulos suos, dicens: Quem dicunt hómines esse Fílium hóminis? Et réliqua.\par
\switchcolumn
\EN
\noindent At that time: When Jesus came into the coasts of Caesarea Philippi, he asked his disciples, saying, Whom do men say that I the Son of man am? And so on.\par
\switchcolumn*

\LA
\LessonTitle{Homilía sancti Leónis Papæ}
\switchcolumn
\EN
\LessonTitle{A Homily by St. Leo the Pope}
\switchcolumn*

\LA
\Source{Sermo 2 in anniversario assumpt. suæ, ante medium}
\switchcolumn
\EN
\Source{Sermo 2 in anniversario assumpt. suæ ante medium}
\switchcolumn*

\LA
\noindent Cum, sicut evangélica lectióne reserátum est, interrogásset Dóminus discípulos, quem ipsum (multis divérsa opinántibus) créderent; respondissétque beátus Petrus, dicens: Tu es Christus Fílius Dei vivi; Dóminus ait: Beátus es, Simon Bar-Jona, quia caro et sanguis non revelávit tibi, sed Pater meus, qui in cælis est: et ego dico tibi, quia tu es Petrus, et super hanc petram ædificábo Ecclésiam meam, et portæ ínferi non prævalébunt advérsus eam. Et tibi dabo claves regni cælórum: et quodcúmque ligáveris super terram, erit ligátum et in cælis: et quodcúmque sólveris super terram, erit solútum et in cælis. Manet ergo disposítio veritátis, et beátus Petrus, in accépta fortitúdine petræ persevérans, suscépta Ecclésiæ gubernácula non relíquit.\par
\switchcolumn
\EN
\noindent When the Lord, as we read in the Gospel, asked his disciples who did men, amid their diverse speculations, believe him the Son of Man to be, blessed Peter answered and said: Thou art the Christ, the Son of the living God. And the Lord answered and said unto him: Blessed art thou, Simon Bar-Jona: for flesh and blood hath not revealed it unto thee, but my Father, which is in heaven: and I say also unto thee: That thou art Peter, and upon this rock I will build my Church, and the gates of hell shall not prevail against it; and I will give unto thee the keys of the kingdom of heaven; and whatsoever thou shalt bind on earth shall be bound in heaven; and whatsoever thou shalt loose on earth shall be loosed in heaven. But the dispensation of truth perdures, and blessed Peter, persevering in the strength of the rock which he hath received, hath not relinquished the position he assumed at the helm of the Church.\par
\switchcolumn*

\LA
\begin{Resp}\Rsign Ego sum vitis vera, et vos pálmites: \Star{} Qui manet in me, et ego in eo, hic fert fructum multum, allelúja, allelúja.\end{Resp}
\begin{Resp}\Vsign Sicut diléxit me Pater, et ego diléxi vos.\end{Resp}
\begin{Resp}\Rsign Qui manet in me, et ego in eo, hic fert fructum multum, allelúja, allelúja.\end{Resp}
\switchcolumn
\EN
\begin{Resp}\Rsign I am the vine: you the branches. \Star{} He that abideth in me, and I in him, the same beareth much fruit, alleluia, alleluia.\end{Resp}
\begin{Resp}\Vsign As the Father hath loved me, I also have loved you.\end{Resp}
\begin{Resp}\Rsign He that abideth in me, and I in him, the same beareth much fruit, alleluia, alleluia.\end{Resp}
\begin{Resp}\Vsign Glory be to the Father, and to the Son, \Star{} and to the Holy Ghost.\end{Resp}
\begin{Resp}\Rsign He that abideth in me, and I in him, the same beareth much fruit, alleluia, alleluia.\end{Resp}
\switchcolumn*

\LA
\LessonHead{Lectio VIII}
\switchcolumn
\EN
\LessonHead{Lesson VIII}
\switchcolumn*

\LA
\noindent In univérsa namque Ecclésia, Tu es Christus Fílius Dei vivi, quotídie Petrus dicit; et omnis lingua, quæ confitétur Dóminum, magistério hujus vocis imbúitur. Hæc fides diábolum vincit et captivórum ejus víncula dissólvit. Hæc érutos mundo, ínserit cælo, et portæ ínferi advérsus eam prævalére non possunt. Tanta enim divínitus soliditáte muníta est, ut eam neque hærética umquam corrúmpere právitas, nec pagána potúerit superáre perfídia. His ítaque modis, dilectíssimi, rationábile obséquio celebrétur hodiérna festívitas: ut in persóna humilitátis meæ ille intelligátur, ille honorétur, in quo et ómnium pastórum sollicitúdo, cum commendatárum sibi óvium custódia persevérat, et cujus étiam dígnitas in indígno heréde non déficit.\par
\switchcolumn
\EN
\noindent In the universal Church it is as if Peter were still saying every day: Thou art the Christ, the Son of the living God. For every tongue which confesseth the Lord is taught that confession by the teaching of Peter. This is the Faith that overcometh the devil and looseth the bonds of his prisoners. This is the Faith which maketh men free of the world and bringeth them to heaven, and the gates of hell are impotent to prevail against it. This is the rock which God hath fortified with such ramparts of salvation, that the contagion of heresy will never be able to infect it, nor idolatry and unbelief to overcome it. And therefore, dearly beloved, we celebrate today’s festival with reasonable obedience, that in my humble person he may be acknowledged and honoured who doth continue to care for all the shepherds as well as sheep entrusted unto him, and who doth lose none of his dignity even in an unworthy successor.\par
\switchcolumn*

\LA
\begin{Resp}\Rsign Cándidi facti sunt Nazarǽi ejus, allelúja: splendórem Deo dedérunt, allelúja: \Star{} Et sicut lac coaguláti sunt, allelúja, allelúja.\end{Resp}
\begin{Resp}\Vsign Candidióres nive, nitidióres lacte, rubicundióres ébore antíquo, sapphíro pulchrióres.\end{Resp}
\begin{Resp}\Rsign Et sicut lac coaguláti sunt, allelúja, allelúja.\end{Resp}
\begin{Resp}\Vsign Glória Patri, et Fílio, \Star{} et Spirítui Sancto.\end{Resp}
\begin{Resp}\Rsign Et sicut lac coaguláti sunt, allelúja, allelúja.\end{Resp}
\switchcolumn
\EN
\begin{Resp}\Rsign Her Nazarites are become pure, Alleluia: they reflect the glory of God, Alleluia. \Star{} They are whiter than milk. Alleluia, Alleluia.\end{Resp}
\begin{Resp}\Vsign They are purer than snow, they are whiter than milk, they are more ruddy in body than coral, their polishing is of sapphire.\end{Resp}
\begin{Resp}\Rsign They are whiter than milk. Alleluia, Alleluia.\end{Resp}
\begin{Resp}\Vsign Glory be to the Father, and to the Son, \Star{} and to the Holy Ghost.\end{Resp}
\begin{Resp}\Rsign They are whiter than milk. Alleluia, Alleluia.\end{Resp}
\switchcolumn*

\LA
\LessonHead{Lectio IX}
\switchcolumn
\EN
\LessonHead{Lesson IX}
\switchcolumn*

\LA
\noindent Cum ergo cohortatiónes nostras áuribus vestræ sanctitátis adhibémus, ipsum vobis, cujus vice fúngimur, loqui crédite: quia et illíus vos afféctu monémus, et non áliud vobis, quam quod dócuit, prædicámus; obsecrántes, ut succíncti lumbos mentis vestræ, castam et sóbriam vitam in Dei timóre ducátis. Coróna mea, sicut Apóstolus ait, et gáudium vos estis, si fides vestra, quæ ab inítio Evangélii in univérso mundo prædicáta est, in dilectióne et sanctitáte permánserit. Nam licet omnem Ecclésiam, quæ in toto est orbe terrárum, cunctis opórteat florére virtútibus; vos tamen præcípue inter céteros pópulos decet méritis pietátis excéllere, quos in ipsa apostólicæ petræ arce fundátos, et Dóminus noster Jesus Christus cum ómnibus redémit, et beátus Apóstolus Petrus præ ómnibus erudívit.\par
\switchcolumn
\EN
\noindent When, therefore, we address our exhortations to your godly ears, believe ye that ye are hearing him speak whose office we are discharging. Yea, it is with his love for you that we warn you. And we preach unto you no other thing than that which he taught, entreating you as did he: Gird up the loins of your mind; be sober; be ye holy in all manner of living; pass the time of your sojourning here in the fear of God. My disciples, dearly beloved, ye are to me as the disciples of the Apostle Paul were to him, namely: My crown and joy; if so be that your faith, abide, still in all lowliness and holiness, like unto the first times of the Gospel. For although the whole Church, which is in all the world, should indeed abound in all the virtues, it becometh especially you among all others to excel in acts of piety, founded as ye be on the very citadel of the Apostolic Rock ye who have not only been redeemed with the rest of men by our Lord Jesus Christ, but who have been instructed by the blessed Apostle Peter far beyond all others.\par
\switchcolumn*

\LA
\TeDeum{Te Deum}
\switchcolumn
\EN
\TeDeum{Te Deum}
\switchcolumn*

\end{paracol}

\Office{Commune unius Confessoris Pontificis}{Common of a Confessor Bishop}{}
\addcontentsline{toc}{subsection}{Commune unius Confessoris Pontificis\enspace\textit{Common of a Confessor Bishop}}
\begin{paracol}{2}
\LA
\LessonHead{Lectio I}
\switchcolumn
\EN
\LessonHead{Lesson I}
\switchcolumn*

\LA
\LessonTitle{De Epístola prima beáti Pauli Apóstoli ad Timótheum}
\switchcolumn
\EN
\LessonTitle{Lesson from the first letter of St. Paul the Apostle to Timothy}
\switchcolumn*

\LA
\Cite{1 Tim 3:1-7}
\switchcolumn
\EN
\Cite{1 Tim 3:1-7}
\switchcolumn*

\LA
\noindent Fidélis sermo: si quis episcopátum desíderat, bonum opus desíderat. Opórtet ergo epíscopum irreprehensíbilem esse, uníus uxóris virum, sóbrium, prudéntem, ornátum, pudícum, hospitálem, doctórem, non vinoléntum, non percussórem, sed modéstum; non litigiósum, non cúpidum, sed suæ dómui bene præpósitum, fílios habéntem súbditos cum omni castitáte. Si quis autem dómui suæ præésse nescit, quómodo Ecclésiæ Dei diligéntiam habébit? Non neóphytum, ne in supérbiam elátus, in judícium íncidat diáboli. Opórtet autem illum et testimónium habére bonum ab iis qui foris sunt, ut non in oppróbrium íncidat, et in láqueum diáboli.\par
\switchcolumn
\EN
\noindent A faithful saying: if a man desire the office of a bishop, he desireth a good work. It behoveth therefore a bishop to be blameless, the husband of one wife, sober, prudent, of good behaviour, chaste, given to hospitality, a teacher, Not given to wine, no striker, but modest, not quarrelsome, not covetous, but One that ruleth well his own house, having his children in subjection with all chastity. But if a man know not how to rule his own house, how shall he take care of the church of God? Not a neophyte: lest being puffed up with pride, he fall into the judgment of the devil. Moreover he must have a good testimony of them who are without: lest he fall into reproach and the snare of the devil.\par
\switchcolumn*

\LA
\begin{Resp}\Rsign Euge, serve bone et fidélis, quia in pauca fuísti fidélis, supra multa te constítuam: \Star{} Intra in gáudium Dómini tui.\end{Resp}
\begin{Resp}\Vsign Dómine, quinque talénta tradidísti mihi, ecce ália quinque superlucrátus sum.\end{Resp}
\begin{Resp}\Rsign Intra in gáudium Dómini tui.\end{Resp}
\switchcolumn
\EN
\begin{Resp}\Rsign Well done, thou good and faithful servant, thou hast been faithful over a few things. I will make thee ruler over many things: \Star{} Enter thou into the joy of thy Lord.\end{Resp}
\begin{Resp}\Vsign Lord, thou deliveredst unto me five talents; behold, I have gained beside them five talents more.\end{Resp}
\begin{Resp}\Rsign Enter thou into the joy of thy Lord.\end{Resp}
\switchcolumn*

\LA
\LessonHead{Lectio II}
\switchcolumn
\EN
\LessonHead{Lesson II}
\switchcolumn*

\LA
\LessonTitle{De Epístola ad Titum}
\switchcolumn
\EN
\LessonTitle{From the letter of St. Paul the Apostle to Titus}
\switchcolumn*

\LA
\Cite{Titus 1:7-11}
\switchcolumn
\EN
\Cite{Titus 1:7-11}
\switchcolumn*

\LA
\noindent Opórtet enim epíscopum sine crímine esse, sicut Dei dispensatórem: non supérbum, non iracúndum, non vinoléntum, non percussórem, non turpis lucri cúpidum: sed hospitálem, benígnum, sóbrium, justum, sanctum, continéntem, amplecténtem eum, qui secúndum doctrínam est, fidélem sermónem: ut potens sit exhortári in doctrína sana, et eos qui contradícunt, argúere. Sunt enim multi étiam inobediéntes, vaníloqui, et seductóres: máxime qui de circumcisióne sunt: quos opórtet redárgui: qui univérsas domos subvértunt, docéntes quæ non opórtet, turpis lucri grátia.\par
\switchcolumn
\EN
\noindent For a bishop must be without crime, as the steward of God: not proud, not subject to anger, not given to wine, no striker, not greedy of filthy lucre: But given to hospitality, gentle, sober, just, holy, continent: Embracing that faithful word which is according to doctrine, that he may be able to exhort in sound doctrine, and to convince the gainsayers. For there are also many disobedient, vain talkers, and seducers: especially they who are of the circumcision: Who must be reproved, who subvert whole houses, teaching things which they ought not, for filthy lucre’s sake.\par
\switchcolumn*

\LA
\begin{Resp}\Rsign Ecce sacérdos magnus, qui in diébus suis plácuit Deo: \Star{} Ideo jurejurándo fecit illum Dóminus créscere in plebem suam.\end{Resp}
\begin{Resp}\Vsign Benedictiónem ómnium géntium dedit illi, et testaméntum suum confirmávit super caput ejus.\end{Resp}
\begin{Resp}\Rsign Ideo jurejurándo fecit illum Dóminus créscere in plebem suam.\end{Resp}
\switchcolumn
\EN
\begin{Resp}\Rsign Behold an high priest, who in his days pleased God \Star{} Therefore the Lord assured him by an oath that He would multiply his seed among His people.\end{Resp}
\begin{Resp}\Vsign He hath made him a blessing unto all nations, and hath established His covenant upon his head.\end{Resp}
\begin{Resp}\Rsign Therefore the Lord assured him by an oath that He would multiply his seed among His people\end{Resp}
\switchcolumn*

\LA
\LessonHead{Lectio III}
\switchcolumn
\EN
\LessonHead{Lesson III}
\switchcolumn*

\LA
\Cite{Titus 2:1-8}
\switchcolumn
\EN
\Cite{Titus 2:1-8}
\switchcolumn*

\LA
\noindent Tu autem lóquere quæ decent sanam doctrínam: senes ut sóbrii sint, pudíci, prudéntes, sani in fide, in dilectióne, in patiéntia: anus simíliter in hábitu sancto, non criminatríces, non multo vino serviéntes, bene docéntes: ut prudéntiam dóceant adolescéntulas, ut viros suos ament, fílios suos díligant, prudéntes, castas, sóbrias, domus curam habéntes, benígnas, súbditas viris suis, ut non blasphemétur verbum Dei. Júvenes simíliter hortáre ut sóbrii sint. In ómnibus teípsum præbe exémplum bonórum óperum, in doctrína, in integritáte, in gravitáte, verbum sanum, irreprehensíbile: ut is qui ex advérso est, vereátur, nihil habens malum dícere de nobis.\par
\switchcolumn
\EN
\noindent But speak thou the things that become sound doctrine: That the aged men be sober, chaste, prudent, sound in faith, in love, in patience. The aged women, in like manner, in holy attire, not false accusers, not given to much wine, teaching well: That they may teach the young women to be wise, to love their husbands, to love their children, To be discreet, chaste, sober, having a care of the house, gentle, obedient to their husbands, that the word of God be not blasphemed. Young men, in like manner, exhort that they be sober. In all things shew thyself an example of good works, in doctrine, in integrity, in gravity, The sound word that can not be blamed: that he, who is on the contrary part, may be afraid, having no evil to say of us.\par
\switchcolumn*

\LA
\begin{Resp}\Rsign Jurávit Dóminus, et non pœnitébit eum: \Star{} Tu es sacérdos in ætérnum secúndum órdinem Melchísedech.\end{Resp}
\begin{Resp}\Vsign Dixit Dóminus Dómino meo, Sede a dextris meis.\end{Resp}
\begin{Resp}\Rsign Tu es sacérdos in ætérnum secúndum órdinem Melchísedech.\end{Resp}
\begin{Resp}\Vsign Glória Patri, et Fílio, \Star{} et Spirítui Sancto.\end{Resp}
\begin{Resp}\Rsign Tu es sacérdos in ætérnum secúndum órdinem Melchísedech.\end{Resp}
\switchcolumn
\EN
\begin{Resp}\Rsign The Lord hath sworn and will not repent \Star{} Thou art a Priest for ever after the order of Melchisedek.\end{Resp}
\begin{Resp}\Vsign The Lord said unto my Lord Sit Thou at My right hand. R.* Thou art a Priest for ever after the order of Melchisedek.\end{Resp}
\begin{Resp}\Vsign Glory be to the Father, and to the Son, \Star{} and to the Holy Ghost.\end{Resp}
\begin{Resp}\Rsign Thou art a Priest for ever after the order of Melchisedek.\end{Resp}
\switchcolumn*

\LA
\LessonHead{Lectio VII}
\switchcolumn
\EN
\LessonHead{Lesson VII}
\switchcolumn*

\LA
\LessonTitle{Léctio sancti Evangélii secúndum Matthǽum}
\switchcolumn
\EN
\LessonTitle{From the Holy Gospel according to Matthew}
\switchcolumn*

\LA
\Cite{Matt 25:14-23}
\switchcolumn
\EN
\Cite{Matt 25:14-23}
\switchcolumn*

\LA
\noindent In illo témpore: Dixit Jesus discípulis suis parábolam hanc: Homo péregre proficíscens, vocávit servos suos, et trádidit illis bona sua. Et réliqua.\par
\switchcolumn
\EN
\noindent At that time Jesus spake unto His disciples this parable A man, travelling into a far country, called his own servants, and delivered unto them his goods. And so on.\par
\switchcolumn*

\LA
\LessonTitle{Homilía sancti Gregórii Papæ}
\switchcolumn
\EN
\LessonTitle{Homily by Pope St Gregory (the Great.)}
\switchcolumn*

\LA
\Source{Homilia 9 in Evangelia}
\switchcolumn
\EN
\Source{9th on the Gospels.}
\switchcolumn*

\LA
\noindent Léctio sancti Evangélii, fratres caríssimi, solícite consideráre nos ádmonet, ne nos, qui plus céteris in hoc mundo accepísse áliquid cérnimur, ab Auctóre mundi grávius inde judicémur. Cum enim augéntur dona, ratiónes étiam crescunt donórum. Tanto ergo esse humílior atque ad serviéndum Deo prómptior quisque debet ex múnere, quanto se obligatiórem esse cónspicit in reddénda ratióne. Ecce homo, qui péregre proficíscitur, servos suos vocat, eisque ad negótium talénta partítur. Post multum vero témporis positúrus ratiónem revértitur: bene operántes pro apportáto lucro remúnerat, servum vero a bono ópere torpéntem damnat.\par
\switchcolumn
\EN
\noindent Dearly beloved brethren, this Lesson from the Holy Gospel moveth us to take good heed lest we, who are seen in this world to have received more than others, should thereby bring ourselves into greater condemnation from the Maker of this world. To whom much is given, of the same is much required. Therefore, let him that receiveth much, strive to be all the more lowly, and all the more ready to do God service, for his very gifts’ sake, knowing that he will be obliged to give account thereof. Behold, a man, travelling into a far country, calleth his own servants, and delivereth unto them talents, to the end that they may trade therewith. After a long time, the lord of those servants cometh, and reckoneth with them, and to them that have done well He rendereth a reward of their labours, but that servant which was careless of his master’s work He condemneth.\par
\switchcolumn*

\LA
\begin{Resp}\Rsign Amávit eum Dóminus, et ornávit eum: stolam glóriæ índuit eum, \Star{} Et ad portas paradísi coronávit eum.\end{Resp}
\begin{Resp}\Vsign Induit eum Dóminus lorícam fídei, et ornávit eum.\end{Resp}
\begin{Resp}\Rsign Et ad portas paradísi coronávit eum.\end{Resp}
\switchcolumn
\EN
\begin{Resp}\Rsign The Lord loved him and beautified him He clothed him with a robe of glory \Star{} And crowned him at the gates of Paradise.\end{Resp}
\begin{Resp}\Vsign The Lord hath put on him the breast-plate of faith, and hath adorned him.\end{Resp}
\begin{Resp}\Rsign And crowned him at the gates of Paradise.\end{Resp}
\switchcolumn*

\LA
\LessonHead{Lectio VIII}
\switchcolumn
\EN
\LessonHead{Lesson VIII}
\switchcolumn*

\LA
\noindent Quis ítaque iste homo est qui péregre proficíscitur, nisi Redémptor noster, qui, in ea carne quam assúmpserat, ábiit in cælum? Carnis enim locus próprius terra est; quæ quasi ad peregrína dúcitur, dum per Redemptórem nostrum in cælo collocátur. Sed homo iste, péregre proficíscens, servis suis bona sua trádidit; quia fidélibus suis spirituália dona concéssit. Et uni quidem quinque talénta, álii duo, álii vero commísit unum. Quinque étenim sunt córporis sensus, vidélicet: visus, audítus, gustus, odorátus et tactus. Quinque ergo taléntis, donum quinque sénsuum, id est, exteriórum sciéntia, exprímitur. Duóbus vero, intelléctus et operátio designátur. Uníus autem talénti nómine, intelléctus tantúmmodo designátur.\par
\switchcolumn
\EN
\noindent What other, then, is that man travelling into a far country but our Redeemer, Who is gone up from us into heaven in that Flesh Which He had taken into Himself? For the earth is the home of the Flesh, Which travelleth into a far country when our Redeemer giveth It a place in heaven. But that man travelling into a far country delivered unto his servants his goods and so doth our Redeemer give spiritual gifts unto His faithful people. “And unto one he gave five talents, to another two, and to another one.” There are five bodily senses that is, sight, hearing, taste, smell, and touch. By the five talents therefore are signified the five senses, that is, outward knowledge. By the two, wit and work. And by the figure of the one talent, understanding, which is alone.\par
\switchcolumn*

\LA
\begin{Resp}\Rsign Sint lumbi vestri præcíncti, et lucérnæ ardéntes in mánibus vestris: \Star{} Et vos símiles homínibus exspectántibus dóminum suum, quando revertátur a núptiis.\end{Resp}
\begin{Resp}\Vsign Vigiláte ergo, quia nescítis qua hora Dóminus vester ventúrus sit.\end{Resp}
\begin{Resp}\Rsign Et vos símiles homínibus exspectántibus dóminum suum, quando revertátur a núptiis.\end{Resp}
\begin{Resp}\Vsign Glória Patri, et Fílio, \Star{} et Spirítui Sancto.\end{Resp}
\begin{Resp}\Rsign Et vos símiles homínibus exspectántibus dóminum suum, quando revertátur a núptiis.\end{Resp}
\switchcolumn
\EN
\begin{Resp}\Rsign Let your loins be girded about, and your lights burning; \Star{} And ye yourselves like unto men that wait for their lord, when he will return from the wedding.\end{Resp}
\begin{Resp}\Vsign Watch therefore, for ye know not what hour your Lord doth come.\end{Resp}
\begin{Resp}\Rsign And ye yourselves like unto men that wait for their lord, when he will return from the wedding.\end{Resp}
\begin{Resp}\Vsign Glory be to the Father, and to the Son, \Star{} and to the Holy Ghost.\end{Resp}
\begin{Resp}\Rsign And ye yourselves like unto men that wait for their lord, when he will return from the wedding.\end{Resp}
\switchcolumn*

\LA
\LessonHead{Lectio IX}
\switchcolumn
\EN
\LessonHead{Lesson IX}
\switchcolumn*

\LA
\noindent Sed is, qui quinque talénta accéperat, ália quinque lucrátus est: quia sunt nonnúlli, qui, etsi intérna ac mýstica penetráre nésciunt, pro intentióne tamen supérnæ pátriæ docent recta quos possunt; de ipsis exterióribus, quæ accepérunt, duplum taléntum portant; dumque se a carnis petulántia et a terrenárum rerum ámbitu atque a visibílium voluptáte custódiunt, ab his étiam álios admonéndo compéscunt. Et sunt nonnúlli, qui, quasi duóbus taléntis ditáti, intelléctum atque operatiónem percípiunt, subtília de intérnis intélligunt, mira in exterióribus operántur; cumque et intelligéndo et operándo áliis prǽdicant, quasi duplicátum de negótio lucrum repórtant.\par
\switchcolumn
\EN
\noindent And so he that had received five talents, gained other five talents for some there be who, while yet they are not able to go on unto things inward and mystic, do yet so desire our Fatherland which is above, that they teach well all whom they can, and of those very outward things which they have received make gain double. These are they which keep themselves clean from the unruly motions of the flesh, and from the lust of the world, and from the delight of things which are seen, and, by their preaching, keep other men also clean from all these things. And some there are who receive, as their two talents, the power to think and the power to work. These are they which inwardly understand dark things, and outwardly work wonders. And these, since they preach unto others, both through their understanding and their works, gain, as it were, double, for the talents which they have received.\par
\switchcolumn*

\LA
\TeDeum{Te Deum}
\switchcolumn
\EN
\TeDeum{Te Deum}
\switchcolumn*

\end{paracol}

\Office{Commune Doctoris Pontificis}{Common of a Doctor Bishop}{}
\addcontentsline{toc}{subsection}{Commune Doctoris Pontificis\enspace\textit{Common of a Doctor Bishop}}
\begin{paracol}{2}
\LA
\LessonHead{Lectio I}
\switchcolumn
\EN
\LessonHead{Lesson I}
\switchcolumn*

\LA
\LessonTitle{De libro Ecclesiástici}
\switchcolumn
\EN
\LessonTitle{Lesson from the book of Ecclesiasticus}
\switchcolumn*

\LA
\Cite{Sir 39:1-5}
\switchcolumn
\EN
\Cite{Sir 39:1-5}
\switchcolumn*

\LA
\noindent Sapiéntiam ómnium antiquórum exquíret sápiens, et in prophétis vacábit. Narratiónem virórum nominatórum conservábit, et in versútias parabolárum simul introíbit. Occúlta proverbiórum exquíret, et in abscónditis parabolárum conversábitur. In médio magnatórum ministrábit, et in conspéctu prǽsidis apparébit. In terram alienigenárum géntium pertránsiet; bona enim et mala in homínibus tentábit.\par
\switchcolumn
\EN
\noindent The wise men will seek out the wisdom of all the ancients, and will be occupied in the prophets. He will keep the sayings of renowned men, and will enter withal into the subtilties of parables. He will search out the hidden meanings of proverbs, and will be conversant in the secrets of parables. He shall serve among great men, and: appear before the governor. He shall pass into strange countries: for he shall try good and evil among men.\par
\switchcolumn*

\LA
\begin{Resp}\Rsign Euge, serve bone et fidélis, quia in pauca fuísti fidélis, supra multa te constítuam: \Star{} Intra in gáudium Dómini tui.\end{Resp}
\begin{Resp}\Vsign Dómine, quinque talénta tradidísti mihi, ecce ália quinque superlucrátus sum.\end{Resp}
\begin{Resp}\Rsign Intra in gáudium Dómini tui.\end{Resp}
\switchcolumn
\EN
\begin{Resp}\Rsign Well done, thou good and faithful servant, thou hast been faithful over a few things. I will make thee ruler over many things: \Star{} Enter thou into the joy of thy Lord.\end{Resp}
\begin{Resp}\Vsign Lord, thou deliveredst unto me five talents; behold, I have gained beside them five talents more.\end{Resp}
\begin{Resp}\Rsign Enter thou into the joy of thy Lord.\end{Resp}
\switchcolumn*

\LA
\LessonHead{Lectio II}
\switchcolumn
\EN
\LessonHead{Lesson II}
\switchcolumn*

\LA
\Cite{Sir 39:6-10}
\switchcolumn
\EN
\Cite{Sir 39:6-10}
\switchcolumn*

\LA
\noindent Cor suum tradet ad vigilándum dilúculo ad Dóminum, qui fecit illum, et in conspéctu Altíssimi deprecábitur. Apériet os suum in oratióne, et pro delíctis suis deprecábitur. Si enim Dóminus magnus volúerit, spíritu intellegéntiæ replébit illum: et ipse tamquam imbres mittet elóquia sapiéntiæ suæ, et in oratióne confitébitur Dómino: et ipse díriget consílium ejus, et disciplínam, et in abscónditis suis consiliábitur.\par
\switchcolumn
\EN
\noindent He will give his heart to resort early to the Lord that made him, and he will pray in the sight of the most High. He will open his mouth in prayer, and will make supplication for his sins. For if it shall please the great Lord, he will fill him with the spirit of understanding: And he will pour forth the words of his wisdom as showers, and in his prayer he will confess to the Lord. And he shall direct his counsel, and his knowledge, and in his secrets shall he meditate.\par
\switchcolumn*

\LA
\begin{Resp}\Rsign Ecce sacérdos magnus, qui in diébus suis plácuit Deo: \Star{} Ideo jurejurándo fecit illum Dóminus créscere in plebem suam.\end{Resp}
\begin{Resp}\Vsign Benedictiónem ómnium géntium dedit illi, et testaméntum suum confirmávit super caput ejus.\end{Resp}
\begin{Resp}\Rsign Ideo jurejurándo fecit illum Dóminus créscere in plebem suam.\end{Resp}
\switchcolumn
\EN
\begin{Resp}\Rsign Behold an high priest, who in his days pleased God \Star{} Therefore the Lord assured him by an oath that He would multiply his seed among His people.\end{Resp}
\begin{Resp}\Vsign He hath made him a blessing unto all nations, and hath established His covenant upon his head.\end{Resp}
\begin{Resp}\Rsign Therefore the Lord assured him by an oath that He would multiply his seed among His people\end{Resp}
\switchcolumn*

\LA
\LessonHead{Lectio III}
\switchcolumn
\EN
\LessonHead{Lesson III}
\switchcolumn*

\LA
\Cite{Sir 39:11-14}
\switchcolumn
\EN
\Cite{Sir 39:11-14}
\switchcolumn*

\LA
\noindent Ipse palam fáciet disciplínam doctrínæ suæ, et in lege testaménti Dómini gloriábitur. Collaudábunt multi sapiéntiam ejus, et usque in sǽculum non delébitur. Non recédet memória ejus, et nomen ejus requirétur a generatióne in generatiónem. Sapiéntiam ejus enarrábunt gentes, et laudem ejus enuntiábit Ecclésia.\par
\switchcolumn
\EN
\noindent He shall show forth the discipline he hath learned, and shall glory in the law of the covenant of the Lord. Many shall praise his wisdom, and it shall never be forgotten. The memory of him shall not depart away, and his name shall be in request from generation to generation. Nations shall declare his wisdom, and the church shall show forth his praise.\par
\switchcolumn*

\LA
\begin{Resp}\Rsign Jurávit Dóminus, et non pœnitébit eum: \Star{} Tu es sacérdos in ætérnum secúndum órdinem Melchísedech.\end{Resp}
\begin{Resp}\Vsign Dixit Dóminus Dómino meo, Sede a dextris meis.\end{Resp}
\begin{Resp}\Rsign Tu es sacérdos in ætérnum secúndum órdinem Melchísedech.\end{Resp}
\begin{Resp}\Vsign Glória Patri, et Fílio, \Star{} et Spirítui Sancto.\end{Resp}
\begin{Resp}\Rsign Tu es sacérdos in ætérnum secúndum órdinem Melchísedech.\end{Resp}
\switchcolumn
\EN
\begin{Resp}\Rsign The Lord hath sworn and will not repent \Star{} Thou art a Priest for ever after the order of Melchisedek.\end{Resp}
\begin{Resp}\Vsign The Lord said unto my Lord Sit Thou at My right hand. R.* Thou art a Priest for ever after the order of Melchisedek.\end{Resp}
\begin{Resp}\Vsign Glory be to the Father, and to the Son, \Star{} and to the Holy Ghost.\end{Resp}
\begin{Resp}\Rsign Thou art a Priest for ever after the order of Melchisedek.\end{Resp}
\switchcolumn*

\LA
\LessonHead{Lectio VII}
\switchcolumn
\EN
\LessonHead{Lesson VII}
\switchcolumn*

\LA
\LessonTitle{Léctio sancti Evangélii secúndum Matthǽum}
\switchcolumn
\EN
\LessonTitle{From the Holy Gospel according to Matthew}
\switchcolumn*

\LA
\Cite{Matt 5:13-19}
\switchcolumn
\EN
\Cite{Matt 5:13-19}
\switchcolumn*

\LA
\noindent In illo témpore: Dixit Jesus discípulis suis: Vos estis sal terræ. Quod si sal evanúerit, in quo saliétur? Et réliqua.\par
\switchcolumn
\EN
\noindent At that time: Jesus said unto his disciples: Ye are the salt of the earth: But if the salt have lost his savour, wherewith shall it be salted? And so on, and that which followeth.\par
\switchcolumn*

\LA
\LessonTitle{Homilía sancti Augustíni Epíscopi}
\switchcolumn
\EN
\LessonTitle{A Homily by St. Augustine the Bishop}
\switchcolumn*

\LA
\Source{Lib. 1 de Sermone Domini in monte, cap. 6}
\switchcolumn
\EN
\Source{Lib. 1 de Sermóne Dómini in monte, cap. 6}
\switchcolumn*

\LA
\noindent Osténdit Dóminus fátuos esse judicándos, qui temporálium bonórum vel cópiam sectántes vel inópiam metuéntes, amíttunt ætérna, quæ nec dari possunt ab homínibus nec auférri. Itaque, si sal infatuátum fúerit, in quo saliétur? Id est, si vos, per quos condiéndi sunt quodámmodo pópuli, metu persecutiónum temporálium amiséritis regna cælórum; qui erunt hómines, per quos a vobis error auferátur, cum vos elégerit Deus, per quos errórem áuferat ceterórum?\par
\switchcolumn
\EN
\noindent The Lord would have us understand how that men do lose their power of savouring others with righteousness when they are willing to place their eternal welfare in jeopardy for the sake of any témporal advantage, like as attainment of ease or luxury, or escape from suffering or toil. For that which is eternal, unlike things of this world, can neither be bestowed by men, nor by them taken away. Hence, when he asketh: If the salt have lost his savour, wherewith shall it be salted? he would have us understand the question to be: If ye, by whom mankind is preserved from corruption, be willing to lose the kingdom of heaven so as to escape trials or persecutions in this world, who is there to preserve you from corruption, seeing ye are they that God hath chosen to preserve all others from corruption?\par
\switchcolumn*

\LA
\begin{Resp}\Rsign Amávit eum Dóminus, et ornávit eum: stolam glóriæ índuit eum, \Star{} Et ad portas paradísi coronávit eum.\end{Resp}
\begin{Resp}\Vsign Induit eum Dóminus lorícam fídei, et ornávit eum.\end{Resp}
\begin{Resp}\Rsign Et ad portas paradísi coronávit eum.\end{Resp}
\switchcolumn
\EN
\begin{Resp}\Rsign The Lord loved him and beautified him He clothed him with a robe of glory \Star{} And crowned him at the gates of Paradise.\end{Resp}
\begin{Resp}\Vsign The Lord hath put on him the breast-plate of faith, and hath adorned him.\end{Resp}
\begin{Resp}\Rsign And crowned him at the gates of Paradise.\end{Resp}
\switchcolumn*

\LA
\LessonHead{Lectio VIII}
\switchcolumn
\EN
\LessonHead{Lesson VIII}
\switchcolumn*

\LA
\noindent Ergo ad níhilum valet sal infatuátum, nisi ut mittátur foras et calcétur ab homínibus. Non ítaque calcátur ab homínibus qui pátitur persecutiónem; sed qui, persecutiónem timéndo, infatuátur. Calcári enim non potest nisi inférior; sed inférior non est, qui, quamvis córpore multa in terra sustíneat, corde tamen fixus in cælo est.\par
\switchcolumn
\EN
\noindent Those that should be the salt of the earth, but have lost their savour, are thenceforth good for nothing, but to be cast out, and to be trodden under foot of men. But no one that suffereth persecution is truly said to be trodden under foot of men. Rather, that one is truly trodden under foot of men who through fear of persecution hath lost the savour of righteousness. For no one can be trodden upon, unless he be beneath him which treadeth upon him. And certainly no one who hath his heart in heaven, no matter how grievously he doth suffer in his body on earth, is rightly said to be beneath anyone who misuseth him.\par
\switchcolumn*

\LA
\begin{Resp}\Rsign In médio Ecclésiæ apéruit os ejus, \Star{} Et implévit eum Dóminus spíritu sapiéntiæ et intelléctus.\end{Resp}
\begin{Resp}\Vsign Jucunditátem et exsultatiónem thesaurizávit super eum.\end{Resp}
\begin{Resp}\Rsign Et implévit eum Dóminus spíritu sapiéntiæ et intelléctus.\end{Resp}
\begin{Resp}\Vsign Glória Patri, et Fílio, \Star{} et Spirítui Sancto.\end{Resp}
\begin{Resp}\Rsign Et implévit eum Dóminus spíritu sapiéntiæ et intelléctus.\end{Resp}
\switchcolumn
\EN
\begin{Resp}\Rsign In the midst of the congregation did the Lord open his mouth. \Star{} And filled him with the spirit of wisdom and understanding.\end{Resp}
\begin{Resp}\Vsign He made him rich with joy and gladness.\end{Resp}
\begin{Resp}\Rsign And filled him with the spirit of wisdom and understanding.\end{Resp}
\begin{Resp}\Vsign Glory be to the Father, and to the Son, \Star{} and to the Holy Ghost.\end{Resp}
\begin{Resp}\Rsign And filled him with the spirit of wisdom and understanding.\end{Resp}
\switchcolumn*

\LA
\LessonHead{Lectio IX}
\switchcolumn
\EN
\LessonHead{Lesson IX}
\switchcolumn*

\LA
\noindent Vos estis lumen mundi. Quómodo dixit supérius sal terræ, sic nunc dicit lumen mundi. Nam, neque supérius ista terra accipiénda est, quam pédibus corpóreis calcámus; sed hómines, qui in terra hábitant, vel étiam peccatóres, quorum condiéndis et exstinguéndis putóribus apostólicum salem Dóminus misit. Et hic mundum non cælum et terram, sed hómines qui sunt in mundo vel díligunt mundum, opórtet intélligi; quibus illuminándis Apóstoli missi sunt. Non potest cívitas abscóndi super montem pósita; id est, fundáta super insígnem magnámque justítiam, quam signíficat étiam ipse mons, in quo dísputat Dóminus.\par
\switchcolumn
\EN
\noindent Ye are the light of the world. And we are to understand the word World in the same sense as the word Earth when he spoke above of the salt of the earth, that is, not that earth whereupon we walk with our bodily feet, but the men which dwell upon the earth; in other words, sinners, for the sweetening and correction of whose corruption, the Lord hath sent his Apostles, as it were, as so much salt. And so by the world we are to understand, not the heaven and the earth, but the men who are in the world and love the world, for the enlightening of whom the Apostles have been sent. A city that is set on a hill cannot be hid: that is, what is founded upon the heights of righteousness, whereof the mountain upon which the Lord gave this discourse was itself a figure, is magnificent in the eyes of all men.\par
\switchcolumn*

\LA
\TeDeum{Te Deum}
\switchcolumn
\EN
\TeDeum{Te Deum}
\switchcolumn*

\end{paracol}

\Office{Commune Summorum Pontificum Confessorum}{Common of Popes}{}
\addcontentsline{toc}{subsection}{Commune Summorum Pontificum Confessorum\enspace\textit{Common of Popes}}
\begin{paracol}{2}
\LA
\LessonHead{Lectio VII}
\switchcolumn
\EN
\LessonHead{Lesson VII}
\switchcolumn*

\LA
\LessonTitle{Léctio sancti Evangélii secúndum Matthǽum}
\switchcolumn
\EN
\LessonTitle{From the Holy Gospel according to Matthew}
\switchcolumn*

\LA
\Cite{Matt 16:13-19}
\switchcolumn
\EN
\Cite{Matt 16:13-19}
\switchcolumn*

\LA
\noindent In illo témpore: Venit Jesus in partes Cæsaréæ Philíppi, et interrogábat discípulos suos, dicens: Quem dicunt hómines esse Fílium hóminis? Et réliqua.\par
\switchcolumn
\EN
\noindent At that time: When Jesus came into the coasts of Caesarea Philippi, he asked his disciples, saying, Whom do men say that I the Son of man am? And so on.\par
\switchcolumn*

\LA
\LessonTitle{Homilía sancti Leónis Papæ}
\switchcolumn
\EN
\LessonTitle{A Homily by St. Leo the Pope}
\switchcolumn*

\LA
\Source{Sermo 2 in anniversario assumpt. suæ, ante medium}
\switchcolumn
\EN
\Source{Sermo 2 in anniversario assumpt. suæ ante medium}
\switchcolumn*

\LA
\noindent Cum, sicut evangélica lectióne reserátum est, interrogásset Dóminus discípulos, quem ipsum (multis divérsa opinántibus) créderent; respondissétque beátus Petrus, dicens: Tu es Christus Fílius Dei vivi; Dóminus ait: Beátus es, Simon Bar-Jona, quia caro et sanguis non revelávit tibi, sed Pater meus, qui in cælis est: et ego dico tibi, quia tu es Petrus, et super hanc petram ædificábo Ecclésiam meam, et portæ ínferi non prævalébunt advérsus eam. Et tibi dabo claves regni cælórum: et quodcúmque ligáveris super terram, erit ligátum et in cælis: et quodcúmque sólveris super terram, erit solútum et in cælis. Manet ergo disposítio veritátis, et beátus Petrus, in accépta fortitúdine petræ persevérans, suscépta Ecclésiæ gubernácula non relíquit.\par
\switchcolumn
\EN
\noindent When the Lord, as we read in the Gospel, asked his disciples who did men, amid their diverse speculations, believe him the Son of Man to be, blessed Peter answered and said: Thou art the Christ, the Son of the living God. And the Lord answered and said unto him: Blessed art thou, Simon Bar-Jona: for flesh and blood hath not revealed it unto thee, but my Father, which is in heaven: and I say also unto thee: That thou art Peter, and upon this rock I will build my Church, and the gates of hell shall not prevail against it; and I will give unto thee the keys of the kingdom of heaven; and whatsoever thou shalt bind on earth shall be bound in heaven; and whatsoever thou shalt loose on earth shall be loosed in heaven. But the dispensation of truth perdures, and blessed Peter, persevering in the strength of the rock which he hath received, hath not relinquished the position he assumed at the helm of the Church.\par
\switchcolumn*

\LA
\begin{Resp}\Rsign Amávit eum Dóminus, et ornávit eum: stolam glóriæ índuit eum, \Star{} Et ad portas paradísi coronávit eum.\end{Resp}
\begin{Resp}\Vsign Induit eum Dóminus lorícam fídei, et ornávit eum.\end{Resp}
\begin{Resp}\Rsign Et ad portas paradísi coronávit eum.\end{Resp}
\switchcolumn
\EN
\begin{Resp}\Rsign The Lord loved him and beautified him He clothed him with a robe of glory \Star{} And crowned him at the gates of Paradise.\end{Resp}
\begin{Resp}\Vsign The Lord hath put on him the breast-plate of faith, and hath adorned him.\end{Resp}
\begin{Resp}\Rsign And crowned him at the gates of Paradise.\end{Resp}
\switchcolumn*

\LA
\LessonHead{Lectio VIII}
\switchcolumn
\EN
\LessonHead{Lesson VIII}
\switchcolumn*

\LA
\noindent In univérsa namque Ecclésia, Tu es Christus Fílius Dei vivi, quotídie Petrus dicit; et omnis lingua, quæ confitétur Dóminum, magistério hujus vocis imbúitur. Hæc fides diábolum vincit et captivórum ejus víncula dissólvit. Hæc érutos mundo, ínserit cælo, et portæ ínferi advérsus eam prævalére non possunt. Tanta enim divínitus soliditáte muníta est, ut eam neque hærética umquam corrúmpere právitas, nec pagána potúerit superáre perfídia. His ítaque modis, dilectíssimi, rationábile obséquio celebrétur hodiérna festívitas: ut in persóna humilitátis meæ ille intelligátur, ille honorétur, in quo et ómnium pastórum sollicitúdo, cum commendatárum sibi óvium custódia persevérat, et cujus étiam dígnitas in indígno heréde non déficit.\par
\switchcolumn
\EN
\noindent In the universal Church it is as if Peter were still saying every day: Thou art the Christ, the Son of the living God. For every tongue which confesseth the Lord is taught that confession by the teaching of Peter. This is the Faith that overcometh the devil and looseth the bonds of his prisoners. This is the Faith which maketh men free of the world and bringeth them to heaven, and the gates of hell are impotent to prevail against it. This is the rock which God hath fortified with such ramparts of salvation, that the contagion of heresy will never be able to infect it, nor idolatry and unbelief to overcome it. And therefore, dearly beloved, we celebrate today’s festival with reasonable obedience, that in my humble person he may be acknowledged and honoured who doth continue to care for all the shepherds as well as sheep entrusted unto him, and who doth lose none of his dignity even in an unworthy successor.\par
\switchcolumn*

\LA
\begin{Resp}\Rsign Sint lumbi vestri præcíncti, et lucérnæ ardéntes in mánibus vestris: \Star{} Et vos símiles homínibus exspectántibus dóminum suum, quando revertátur a núptiis.\end{Resp}
\begin{Resp}\Vsign Vigiláte ergo, quia nescítis qua hora Dóminus vester ventúrus sit.\end{Resp}
\begin{Resp}\Rsign Et vos símiles homínibus exspectántibus dóminum suum, quando revertátur a núptiis.\end{Resp}
\begin{Resp}\Vsign Glória Patri, et Fílio, \Star{} et Spirítui Sancto.\end{Resp}
\begin{Resp}\Rsign Et vos símiles homínibus exspectántibus dóminum suum, quando revertátur a núptiis.\end{Resp}
\switchcolumn
\EN
\begin{Resp}\Rsign Let your loins be girded about, and your lights burning; \Star{} And ye yourselves like unto men that wait for their lord, when he will return from the wedding.\end{Resp}
\begin{Resp}\Vsign Watch therefore, for ye know not what hour your Lord doth come.\end{Resp}
\begin{Resp}\Rsign And ye yourselves like unto men that wait for their lord, when he will return from the wedding.\end{Resp}
\begin{Resp}\Vsign Glory be to the Father, and to the Son, \Star{} and to the Holy Ghost.\end{Resp}
\begin{Resp}\Rsign And ye yourselves like unto men that wait for their lord, when he will return from the wedding.\end{Resp}
\switchcolumn*

\LA
\LessonHead{Lectio IX}
\switchcolumn
\EN
\LessonHead{Lesson IX}
\switchcolumn*

\LA
\noindent Cum ergo cohortatiónes nostras áuribus vestræ sanctitátis adhibémus, ipsum vobis, cujus vice fúngimur, loqui crédite: quia et illíus vos afféctu monémus, et non áliud vobis, quam quod dócuit, prædicámus; obsecrántes, ut succíncti lumbos mentis vestræ, castam et sóbriam vitam in Dei timóre ducátis. Coróna mea, sicut Apóstolus ait, et gáudium vos estis, si fides vestra, quæ ab inítio Evangélii in univérso mundo prædicáta est, in dilectióne et sanctitáte permánserit. Nam licet omnem Ecclésiam, quæ in toto est orbe terrárum, cunctis opórteat florére virtútibus; vos tamen præcípue inter céteros pópulos decet méritis pietátis excéllere, quos in ipsa apostólicæ petræ arce fundátos, et Dóminus noster Jesus Christus cum ómnibus redémit, et beátus Apóstolus Petrus præ ómnibus erudívit.\par
\switchcolumn
\EN
\noindent When, therefore, we address our exhortations to your godly ears, believe ye that ye are hearing him speak whose office we are discharging. Yea, it is with his love for you that we warn you. And we preach unto you no other thing than that which he taught, entreating you as did he: Gird up the loins of your mind; be sober; be ye holy in all manner of living; pass the time of your sojourning here in the fear of God. My disciples, dearly beloved, ye are to me as the disciples of the Apostle Paul were to him, namely: My crown and joy; if so be that your faith, abide, still in all lowliness and holiness, like unto the first times of the Gospel. For although the whole Church, which is in all the world, should indeed abound in all the virtues, it becometh especially you among all others to excel in acts of piety, founded as ye be on the very citadel of the Apostolic Rock ye who have not only been redeemed with the rest of men by our Lord Jesus Christ, but who have been instructed by the blessed Apostle Peter far beyond all others.\par
\switchcolumn*

\LA
\TeDeum{Te Deum}
\switchcolumn
\EN
\TeDeum{Te Deum}
\switchcolumn*

\end{paracol}

\Office{Commune Confessoris non pontificis}{Common of a Confessor not a Bishop}{}
\addcontentsline{toc}{subsection}{Commune Confessoris non pontificis\enspace\textit{Common of a Confessor not a Bishop}}
\begin{paracol}{2}
\LA
\LessonHead{Lectio I}
\switchcolumn
\EN
\LessonHead{Lesson I}
\switchcolumn*

\LA
\LessonTitle{De libro Ecclesiástici}
\switchcolumn
\EN
\LessonTitle{Lesson from the book of Ecclesiasticus}
\switchcolumn*

\LA
\Cite{Sir 31:8-11}
\switchcolumn
\EN
\Cite{Sir 31:8-11}
\switchcolumn*

\LA
\noindent Beátus vir, qui invéntus est sine mácula, et qui post aurum non ábiit, nec sperávit in pecúnia et thesáuris. Quis est hic? et laudábimus eum: fecit enim mirabília in vita sua. Qui probátus est in illo, et perféctus est, erit illi glória ætérna: Qui pótuit tránsgredi, et non est transgréssus; fácere mala, et non fecit: ídeo stabilíta sunt bona illíus in Dómino, et eleemósynas illíus enarrábit omnis ecclésia sanctórum.\par
\switchcolumn
\EN
\noindent Blessed is the rich man that is found without blemish: and that hath not gone after gold, nor put his trust in money nor in treasures. Who is he, and we will praise him? for he hath done wonderful things in his life. Who hath been tried thereby, and made perfect, he shall have glory everlasting. He that could have transgressed, and hath not transgressed: and could do evil things, and hath not done them: Therefore are his goods established in the Lord, and all the church of the saints shall declare his alms\par
\switchcolumn*

\LA
\begin{Resp}\Rsign Euge, serve bone et fidélis, quia in pauca fuísti fidélis, supra multa te constítuam: \Star{} Intra in gáudium Dómini tui.\end{Resp}
\begin{Resp}\Vsign Dómine, quinque talénta tradidísti mihi, ecce ália quinque superlucrátus sum.\end{Resp}
\begin{Resp}\Rsign Intra in gáudium Dómini tui.\end{Resp}
\switchcolumn
\EN
\begin{Resp}\Rsign Well done, thou good and faithful servant, thou hast been faithful over a few things. I will make thee ruler over many things: \Star{} Enter thou into the joy of thy Lord.\end{Resp}
\begin{Resp}\Vsign Lord, thou deliveredst unto me five talents; behold, I have gained beside them five talents more.\end{Resp}
\begin{Resp}\Rsign Enter thou into the joy of thy Lord.\end{Resp}
\switchcolumn*

\LA
\LessonHead{Lectio II}
\switchcolumn
\EN
\LessonHead{Lesson II}
\switchcolumn*

\LA
\Cite{Sir 32:18-20; 32:28; 33:1-3}
\switchcolumn
\EN
\Cite{Sir 32:18-20; 32:28; 33:1-3}
\switchcolumn*

\LA
\noindent Qui timet Dóminum excípiet doctrínam ejus: et qui vigiláverint ad illum invénient benedictiónem. Qui quærit legem, replébitur ab ea: et qui insidióse agit, scandalizábitur in ea. Qui timent Dóminum invénient judícium justum, et justítias quasi lumen accéndent. Qui credit Deo atténdit mandátis: et qui confídit in illo non minorábitur. Timénti Dóminum non occúrrent mala: sed in tentatióne Deus illum conservábit, et liberábit a malis. Sápiens non odit mandáta et justítias, et non illidétur quasi in procélla navis. Homo sensátus credit legi Dei, et lex illi fidélis.\par
\switchcolumn
\EN
\noindent He that feareth the Lord, will receive his discipline: and they that will seek him early, shall find a blessing. He that seeketh the law, shall be filled with it: and he that dealeth deceitfully, shall meet with a stumblingblock therein. They that fear the Lord, shall find just judgment, and shall kindle justice as a light. He that believeth God, taketh heed to the commandments: and he that trusteth in him, shall fare never the worse. No evils shall happen to him that feareth the Lord, but in temptation God will keep him, and deliver him from evils. A wise man hateth not the commandments and justices, and he shall not be dashed in pieces as a ship in a storm. A man of understanding is faithful to the law of God, and the law is faithful to him.\par
\switchcolumn*

\LA
\begin{Resp}\Rsign Justus germinábit sicut lílium: \Star{} Et florébit in ætérnum ante Dóminum.\end{Resp}
\begin{Resp}\Vsign Plantátus in domo Dómini, in átriis domus Dei nostri.\end{Resp}
\begin{Resp}\Rsign Et florébit in ætérnum ante Dóminum.\end{Resp}
\switchcolumn
\EN
\begin{Resp}\Rsign The righteous shall grow as the lily; \Star{} Yea, he shall flourish in the presence of the Lord for ever.\end{Resp}
\begin{Resp}\Vsign Those that be planted in the house of the Lord, shall flourish in the courts of the house of our God.\end{Resp}
\begin{Resp}\Rsign Yea, he shall flourish in the presence of the Lord for ever.\end{Resp}
\switchcolumn*

\LA
\LessonHead{Lectio III}
\switchcolumn
\EN
\LessonHead{Lesson III}
\switchcolumn*

\LA
\Cite{Sir 34:14-20}
\switchcolumn
\EN
\Cite{Sir 34:14-20}
\switchcolumn*

\LA
\noindent Spíritus timéntium Deum quǽritur, et in respéctu illíus benedicétur. Spes enim illórum in salvántem illos, et óculi Dei in diligéntes se. Qui timet Dóminum nihil trepidábit: et non pavébit, quóniam ipse est spes ejus. Timéntis Dóminum, beáta est ánima ejus. Ad quem réspicit, et quis est fortitúdo ejus? Óculi Dómini super timéntes eum: protéctor poténtiæ, firmaméntum virtútis, tégimen ardóris, et umbráculum meridiáni: deprecátio offensiónis, et adjutórium casus: exáltans ánimam, et illúminans óculos, dans sanitátem, et vitam, et benedictiónem.\par
\switchcolumn
\EN
\noindent The spirit of those that fear God; is sought after, and by his regard shall be blessed. For their hope is on him that saveth them, and the eyes of God are upon them that love him. He that feareth the Lord shall tremble at nothing, and shall not be afraid for he is his hope. The soul of him that feareth the Lord is blessed. To whom doth he look, and who in his strength? The eyes of the Lord are upon them that fear him, he is their powerful protector, and strong stay, a defence from the heat, and a cover from the sun at noon, A preservation from stumbling, and a help from falling; he raiseth up the soul, and enlighteneth the eyes, and giveth health, and life, and blessing.\par
\switchcolumn*

\LA
\begin{Resp}\Rsign Iste cognóvit justítiam, et vidit mirabília magna, et exorávit Altíssimum: \Star{} Et invéntus est in número Sanctórum.\end{Resp}
\begin{Resp}\Vsign Iste est qui contémpsit vitam mundi, et pervénit ad cæléstia regna.\end{Resp}
\begin{Resp}\Rsign Et invéntus est in número Sanctórum.\end{Resp}
\begin{Resp}\Vsign Glória Patri, et Fílio, \Star{} et Spirítui Sancto.\end{Resp}
\begin{Resp}\Rsign Et invéntus est in número Sanctórum.\end{Resp}
\switchcolumn
\EN
\begin{Resp}\Rsign This is he which knew righteousness, and saw great wonders, and made his prayer unto the Most High; \Star{} And he is numbered among the Saints.\end{Resp}
\begin{Resp}\Vsign This is he which loved not his life in this world, and is come unto an everlasting kingdom.\end{Resp}
\begin{Resp}\Rsign And he is numbered among the Saints.\end{Resp}
\begin{Resp}\Vsign Glory be to the Father, and to the Son, \Star{} and to the Holy Ghost.\end{Resp}
\begin{Resp}\Rsign And he is numbered among the Saints.\end{Resp}
\switchcolumn*

\LA
\LessonHead{Lectio VII}
\switchcolumn
\EN
\LessonHead{Lesson VII}
\switchcolumn*

\LA
\LessonTitle{Léctio sancti Evangélii secúndum Lucam}
\switchcolumn
\EN
\LessonTitle{From the Holy Gospel according to Luke}
\switchcolumn*

\LA
\Cite{Luc 12:35-40}
\switchcolumn
\EN
\Cite{Luke 12:35-40}
\switchcolumn*

\LA
\noindent In illo témpore: Dixit Jesus discípulis suis: Sint lumbi vestri præcincti, et lucernæ ardentes in mánibus vestris. Et réliqua.\par
\switchcolumn
\EN
\noindent At that time, Jesus said unto His disciples: Let your loins be girded about, and your lights burning. And so on.\par
\switchcolumn*

\LA
\LessonTitle{Homilía sancti Gregórii Papæ}
\switchcolumn
\EN
\LessonTitle{Homily by Pope St. Gregory (the Great.)}
\switchcolumn*

\LA
\Source{Homilia 13 in Evang.}
\switchcolumn
\EN
\Source{13th on the Gospels.}
\switchcolumn*

\LA
\noindent Sancti Evangélii, fratres caríssimi, apérta vobis est léctio recitata. Sed ne alíquibus ipsa ejus planíties alta fortásse videátur, eam sub brevitáte transcúrrimus, quátenus ejus exposítio ita nesciéntibus fiat cógnita, ut tamen sciéntibus non sit onerósa. Dóminus dicit: Sint lumbi vestri præcíncti. Lumbos enim præcíngimus, cum carnis luxúriam per continéntiam coarctámus. Sed quia minus est, mala non ágere, nisi étiam quisque stúdeat, et bonis opéribus insudáre, prótinus ádditur: Et lucérnæ ardéntes in mánibus vestris. Lucérnas quippe ardéntes in mánibus tenémus, cum per bona ópera próximis nostris lucis exémpla monstrámus. De quibus profécto opéribus Dóminus dicit: Lúceat lux vestra coram homínibus, ut vídeant ópera vestra bona, et gloríficent Patrem vestrum qui in cælis est.\par
\switchcolumn
\EN
\noindent Dearly beloved brethren, the words of the Holy Gospel, which have just been read, lie open before you, and, lest their very plainness should make them seem to some to be hard, we will go through them with such shortness as that neither may they which understand not remain unenlightened, nor they which understand be wearied. The Lord saith Let your loins be girded about. Now, we gird our loins about, when by continency we master the lustful inclination of the flesh. But, forasmuch as it sufficeth not for a man to abstain from evil deeds, if he strive not to join thereto the earnest doing of good works, it is immediately added And your lights burning. Our lights burn when, by good works, we give bright example to our neighbour; concerning which works the Lord saith Let your light so shine before men, that they may see your good works, and glorify your Father Which is in heaven. (Matth. v. 16.)\par
\switchcolumn*

\LA
\begin{Resp}\Rsign Iste est qui ante Deum magnas virtútes operátus est, et de omni corde suo laudávit Dóminum: \Star{} Ipse intercédat pro peccátis ómnium populórum.\end{Resp}
\begin{Resp}\Vsign Ecce homo sine queréla, verus Dei cultor, ábstinens se ab omni ópere malo, et pérmanens in innocéntia sua.\end{Resp}
\begin{Resp}\Rsign Ipse intercédat pro peccátis ómnium populórum.\end{Resp}
\switchcolumn
\EN
\begin{Resp}\Rsign This is he which wrought great wonders before God, and praised the Lord with all his heart. \Star{} May he pray for all people, that their sins may be forgiven unto them.\end{Resp}
\begin{Resp}\Vsign Behold a man without blame, a worshipper of God in truth, keeping himself clean from every evil work, and abiding still in his innocency.\end{Resp}
\begin{Resp}\Rsign May he pray for all people, that their sins may be forgiven unto them.\end{Resp}
\switchcolumn*

\LA
\LessonHead{Lectio VIII}
\switchcolumn
\EN
\LessonHead{Lesson VIII}
\switchcolumn*

\LA
\noindent Duo autem sunt, quæ jubéntur, et lumbos restríngere, et lucérnas tenére: ut et mundítia sit castitátis in córpore, et lumen veritátis in operatióne. Redemptóri étenim nostro unum sine áltero placére nequáquam potest: si aut is qui bona agit, adhuc luxúriæ inquinaménta non déserit: aut is qui castitáte præéminet, necdum se per bona ópera exércet. Nec cástitas ergo magna est sine bono ópere, nec opus bonum est áliquod sine castitáte. Sed et si utrúmque ágitur, restat, ut quisquis ille est, spe ad supérnam pátriam tendat, et nequáquam se a vítiis pro mundi hujus honestáte contíneat.\par
\switchcolumn
\EN
\noindent Here, then, are two commandments, to gird our loins about, and to keep our lights burning the cleanness of purity in our body, and the light of the truth in our works. Whoso hath the one and not the other, pleaseth not thereby our Redeemer; that is, he pleaseth Him not which doth good works, but bridleth not himself from the pollutions of lust, neither he which is eminent in chastity, but exerciseth not himself in good works. Neither is chastity a great thing without good works, nor good works anything without chastity. And if any man do both, it remaineth that he must look by hope toward our Fatherland above, and not have for his reason wherethrough he turneth himself away from vice, the love of honour in this present world.\par
\switchcolumn*

\LA
\begin{Resp}\Rsign Sint lumbi vestri præcíncti, et lucérnæ ardéntes in mánibus vestris: \Star{} Et vos símiles homínibus exspectántibus dóminum suum, quando revertátur a núptiis.\end{Resp}
\begin{Resp}\Vsign Vigiláte ergo, quia nescítis qua hora Dóminus vester ventúrus sit.\end{Resp}
\begin{Resp}\Rsign Et vos símiles homínibus exspectántibus dóminum suum, quando revertátur a núptiis.\end{Resp}
\begin{Resp}\Vsign Glória Patri, et Fílio, \Star{} et Spirítui Sancto.\end{Resp}
\begin{Resp}\Rsign Et vos símiles homínibus exspectántibus dóminum suum, quando revertátur a núptiis.\end{Resp}
\switchcolumn
\EN
\begin{Resp}\Rsign Let your loins be girded about, and your lights burning; \Star{} And ye yourselves like unto men that wait for their lord, when he will return from the wedding.\end{Resp}
\begin{Resp}\Vsign Watch therefore, for ye know not what hour your Lord doth come.\end{Resp}
\begin{Resp}\Rsign And ye yourselves like unto men that wait for their lord, when he will return from the wedding.\end{Resp}
\begin{Resp}\Vsign Glory be to the Father, and to the Son, \Star{} and to the Holy Ghost.\end{Resp}
\begin{Resp}\Rsign And ye yourselves like unto men that wait for their lord, when he will return from the wedding.\end{Resp}
\switchcolumn*

\LA
\LessonHead{Lectio IX}
\switchcolumn
\EN
\LessonHead{Lesson IX}
\switchcolumn*

\LA
\noindent Et vos símiles homínibus exspectántibus dóminum suum, quando revertátur a núptiis: ut cum vénerit et pulsáverit, conféstim apériant ei. Venit quippe Dóminus, cum ad judícium próperat: pulsat vero, cum jam per ægritúdinis moléstias esse mortem vicínam desígnat. Cui conféstim aperímus, si hunc cum amóre suscípimus. Aperíre enim júdici pulsánti non vult, qui exíre de córpore trépidat: et vidére eum, quem contempsísse se méminit, júdicem formídat. Qui autem de sua spe et operatióne secúrus est, pulsanti conféstim áperit, quia lætus júdicem sústinet; et, cum tempus propínquæ mortis advénerit, de glória retributiónis hilaréscit.\par
\switchcolumn
\EN
\noindent And ye yourselves like unto men that wait for their lord, when he will return from the wedding that, when he cometh and knocketh, they may open unto him immediately. The Lord cometh at the hour of judgment He knocketh when, by the pains of sickness, He biddeth us know that death is nigh. To Him open we immediately, if we receive Him in love. Whoso feareth to leave this body, will not open to the Judge when He knocketh, for he dreadeth to see that Judge, Whom he knoweth that he hath despised. But whosoever knoweth that his hope and works are built upon a good foundation, when he heareth the Judge knock, openeth to Him immediately, for to such an one that coming is blessed, yea, when the hour of death is at hand, such an one haileth with gladness a glorious reward.\par
\switchcolumn*

\LA
\TeDeum{Te Deum}
\switchcolumn
\EN
\TeDeum{Te Deum}
\switchcolumn*

\LA
\LessonHead{Lectio I (in 2 loco)}
\switchcolumn
\EN
\LessonHead{Lesson I (set 2)}
\switchcolumn*

\LA
\LessonTitle{De libro Sapiéntiæ}
\switchcolumn
\EN
\LessonTitle{Lesson from the book of Wisdom}
\switchcolumn*

\LA
\Cite{Sap 4:7-14}
\switchcolumn
\EN
\Cite{Wis 4:7-14}
\switchcolumn*

\LA
\noindent Justus si morte præoccupátus fúerit, in refrigério erit; Senéctus enim venerábilis est non diutúrna, neque annórum número computáta: cani autem sunt sensus hóminis, et ætas senectútis vita immaculáta. Placens Deo factus est diléctus, et vivens inter peccatóres translátus est. Raptus est, ne malítia mutáret intelléctum ejus, aut ne fíctio decíperet ánimam illíus. Fascinátio enim nugacitátis obscúrat bona, et inconstántia concupiscéntiæ transvértit sensum sine malítia. Consummátus in brevi, explévit témpora multa; plácita enim erat Deo ánima illíus: propter hoc properávit edúcere illum de médio iniquitátum.\par
\switchcolumn
\EN
\noindent But the just man, if he be prevented with death, shall be in rest. For venerable old age is not that of long time, nor counted by the number of years: but the understanding of a man is grey hairs. And a spotless life is old age. He pleased God and was beloved, and living among sinners he was translated. He was taken away lest wickedness should alter his understanding, or deceit beguile his soul. For the bewitching of vanity obscureth good things, and the wandering of concupiscence overturneth the innocent mind. Being made perfect in a short space, he fulfilled a long time: For his soul pleased God: therefore he hastened to bring him out of the midst of iniquities:\par
\switchcolumn*

\LA
\begin{Resp}\Rsign Euge, serve bone et fidélis, quia in pauca fuísti fidélis, supra multa te constítuam: \Star{} Intra in gáudium Dómini tui.\end{Resp}
\begin{Resp}\Vsign Dómine, quinque talénta tradidísti mihi, ecce ália quinque superlucrátus sum.\end{Resp}
\begin{Resp}\Rsign Intra in gáudium Dómini tui.\end{Resp}
\switchcolumn
\EN
\begin{Resp}\Rsign Well done, thou good and faithful servant, thou hast been faithful over a few things. I will make thee ruler over many things: \Star{} Enter thou into the joy of thy Lord.\end{Resp}
\begin{Resp}\Vsign Lord, thou deliveredst unto me five talents; behold, I have gained beside them five talents more.\end{Resp}
\begin{Resp}\Rsign Enter thou into the joy of thy Lord.\end{Resp}
\switchcolumn*

\LA
\LessonHead{Lectio II (in 2 loco)}
\switchcolumn
\EN
\LessonHead{Lesson II (set 2)}
\switchcolumn*

\LA
\Cite{Sap 4:14-19}
\switchcolumn
\EN
\Cite{Wis 4:14-19}
\switchcolumn*

\LA
\noindent Pópuli autem vidéntes, et non intelligéntes, nec ponéntes in præcórdiis tália, quóniam grátia Dei et misericórdia est in sanctos ejus, et respéctus in eléctos illíus. Condémnat autem justus mórtuus vivos ímpios, et juvéntus celérius consummáta longam vitam injústi. Vidébunt enim finem sapiéntis, et non intélligent quid cogitáverit de illo Deus, et quare muníerit illum Dóminus. Vidébunt, et contémnent eum; illos autem Dóminus irridébit. Et erunt post hæc decidéntes sine honóre, et in contumélia inter mórtuos in perpétuum: quóniam disrúmpet illos inflátos sine voce, et commovébit illos a fundaméntis, et usque ad suprémum desolabúntur.\par
\switchcolumn
\EN
\noindent but the people see this, and understand not, nor lay up such things in their hearts: That the grace of God, and his mercy is with his saints, and that he hath respect to his chosen. But the just that is dead, condemneth the wicked that are living, and youth soon ended, the long life of the unjust. For they shall see the end of the wise man, and shall not understand what God hath designed for him, and why the Lord hath set him in safety. They shall see him, and shall despise him: but the Lord shall laugh them to scorn. And they shall fall after this without honour, and be a reproach among the dead for ever: for he shall burst them puffed up and speechless, and shall shake them from the foundations, and they shall be utterly laid waste:\par
\switchcolumn*

\LA
\begin{Resp}\Rsign Justus germinábit sicut lílium: \Star{} Et florébit in ætérnum ante Dóminum.\end{Resp}
\begin{Resp}\Vsign Plantátus in domo Dómini, in átriis domus Dei nostri.\end{Resp}
\begin{Resp}\Rsign Et florébit in ætérnum ante Dóminum.\end{Resp}
\switchcolumn
\EN
\begin{Resp}\Rsign The righteous shall grow as the lily; \Star{} Yea, he shall flourish in the presence of the Lord for ever.\end{Resp}
\begin{Resp}\Vsign Those that be planted in the house of the Lord, shall flourish in the courts of the house of our God.\end{Resp}
\begin{Resp}\Rsign Yea, he shall flourish in the presence of the Lord for ever.\end{Resp}
\switchcolumn*

\LA
\LessonHead{Lectio III (in 2 loco)}
\switchcolumn
\EN
\LessonHead{Lesson III (set 2)}
\switchcolumn*

\LA
\Cite{Sap 4:19-20; 5:1-5}
\switchcolumn
\EN
\Cite{Wis 4:19-20; 5:1-5}
\switchcolumn*

\LA
\noindent Et erunt geméntes, et memória illórum períbit. Vénient in cogitatióne peccatórum suórum tímidi, et tradúcent illos ex advérso iniquitátes ipsórum. Tunc stabunt justi in magna constántia advérsus eos qui se angustiavérunt, et qui abstulérunt labóres eórum. Vidéntes turbabúntur timóre horríbili, et mirabúntur in subitatióne insperátæ salútis; dicéntes intra se, pœniténtiam agéntes, et præ angústia spíritus geméntes: Hi sunt quos habúimus aliquándo in derísum, et in similitúdinem impropérii. Nos insensáti, vitam illórum æstimabámus insániam, et finem illórum sine honóre; ecce quómodo computáti sunt inter fílios Dei, et inter sanctos sors illórum est.\par
\switchcolumn
\EN
\noindent they shall be in sorrow, and their memory shall perish. They shall come with fear at the thought of their sins, and their iniquities shall stand against them to convict them. Then shall the just stand with great constancy against those that have afflicted them, and taken away their labours. These seeing it, shall be troubled with terrible fear, and shall be amazed at the suddenness of their unexpected salvation. Saying within themselves, repenting, and groaning for anguish of spirit: These are they, whom we had some time in derision, and for a parable of reproach. We fools esteemed their life madness, and their end without honour. Behold how they are numbered among the children of God, and their lot is among the saints.\par
\switchcolumn*

\LA
\begin{Resp}\Rsign Iste cognóvit justítiam, et vidit mirabília magna, et exorávit Altíssimum: \Star{} Et invéntus est in número Sanctórum.\end{Resp}
\begin{Resp}\Vsign Iste est qui contémpsit vitam mundi, et pervénit ad cæléstia regna.\end{Resp}
\begin{Resp}\Rsign Et invéntus est in número Sanctórum.\end{Resp}
\begin{Resp}\Vsign Glória Patri, et Fílio, \Star{} et Spirítui Sancto.\end{Resp}
\begin{Resp}\Rsign Et invéntus est in número Sanctórum.\end{Resp}
\switchcolumn
\EN
\begin{Resp}\Rsign This is he which knew righteousness, and saw great wonders, and made his prayer unto the Most High; \Star{} And he is numbered among the Saints.\end{Resp}
\begin{Resp}\Vsign This is he which loved not his life in this world, and is come unto an everlasting kingdom.\end{Resp}
\begin{Resp}\Rsign And he is numbered among the Saints.\end{Resp}
\begin{Resp}\Vsign Glory be to the Father, and to the Son, \Star{} and to the Holy Ghost.\end{Resp}
\begin{Resp}\Rsign And he is numbered among the Saints.\end{Resp}
\switchcolumn*

\end{paracol}

\Office{Commune Doctoris non Pontificis}{Common of a Doctor not a Bishop}{}
\addcontentsline{toc}{subsection}{Commune Doctoris non Pontificis\enspace\textit{Common of a Doctor not a Bishop}}
\begin{paracol}{2}
\LA
\LessonHead{Lectio I}
\switchcolumn
\EN
\LessonHead{Lesson I}
\switchcolumn*

\LA
\LessonTitle{De libro Ecclesiástici}
\switchcolumn
\EN
\LessonTitle{Lesson from the book of Ecclesiasticus}
\switchcolumn*

\LA
\Cite{Sir 39:1-5}
\switchcolumn
\EN
\Cite{Sir 39:1-5}
\switchcolumn*

\LA
\noindent Sapiéntiam ómnium antiquórum exquíret sápiens, et in prophétis vacábit. Narratiónem virórum nominatórum conservábit, et in versútias parabolárum simul introíbit. Occúlta proverbiórum exquíret, et in abscónditis parabolárum conversábitur. In médio magnatórum ministrábit, et in conspéctu prǽsidis apparébit. In terram alienigenárum géntium pertránsiet; bona enim et mala in homínibus tentábit.\par
\switchcolumn
\EN
\noindent The wise men will seek out the wisdom of all the ancients, and will be occupied in the prophets. He will keep the sayings of renowned men, and will enter withal into the subtilties of parables. He will search out the hidden meanings of proverbs, and will be conversant in the secrets of parables. He shall serve among great men, and: appear before the governor. He shall pass into strange countries: for he shall try good and evil among men.\par
\switchcolumn*

\LA
\begin{Resp}\Rsign Euge, serve bone et fidélis, quia in pauca fuísti fidélis, supra multa te constítuam: \Star{} Intra in gáudium Dómini tui.\end{Resp}
\begin{Resp}\Vsign Dómine, quinque talénta tradidísti mihi, ecce ália quinque superlucrátus sum.\end{Resp}
\begin{Resp}\Rsign Intra in gáudium Dómini tui.\end{Resp}
\switchcolumn
\EN
\begin{Resp}\Rsign Well done, thou good and faithful servant, thou hast been faithful over a few things. I will make thee ruler over many things: \Star{} Enter thou into the joy of thy Lord.\end{Resp}
\begin{Resp}\Vsign Lord, thou deliveredst unto me five talents; behold, I have gained beside them five talents more.\end{Resp}
\begin{Resp}\Rsign Enter thou into the joy of thy Lord.\end{Resp}
\switchcolumn*

\LA
\LessonHead{Lectio II}
\switchcolumn
\EN
\LessonHead{Lesson II}
\switchcolumn*

\LA
\Cite{Sir 39:6-10}
\switchcolumn
\EN
\Cite{Sir 39:6-10}
\switchcolumn*

\LA
\noindent Cor suum tradet ad vigilándum dilúculo ad Dóminum, qui fecit illum, et in conspéctu Altíssimi deprecábitur. Apériet os suum in oratióne, et pro delíctis suis deprecábitur. Si enim Dóminus magnus volúerit, spíritu intellegéntiæ replébit illum: et ipse tamquam imbres mittet elóquia sapiéntiæ suæ, et in oratióne confitébitur Dómino: et ipse díriget consílium ejus, et disciplínam, et in abscónditis suis consiliábitur.\par
\switchcolumn
\EN
\noindent He will give his heart to resort early to the Lord that made him, and he will pray in the sight of the most High. He will open his mouth in prayer, and will make supplication for his sins. For if it shall please the great Lord, he will fill him with the spirit of understanding: And he will pour forth the words of his wisdom as showers, and in his prayer he will confess to the Lord. And he shall direct his counsel, and his knowledge, and in his secrets shall he meditate.\par
\switchcolumn*

\LA
\begin{Resp}\Rsign Justus germinábit sicut lílium: \Star{} Et florébit in ætérnum ante Dóminum.\end{Resp}
\begin{Resp}\Vsign Plantátus in domo Dómini, in átriis domus Dei nostri.\end{Resp}
\begin{Resp}\Rsign Et florébit in ætérnum ante Dóminum.\end{Resp}
\switchcolumn
\EN
\begin{Resp}\Rsign The righteous shall grow as the lily; \Star{} Yea, he shall flourish in the presence of the Lord for ever.\end{Resp}
\begin{Resp}\Vsign Those that be planted in the house of the Lord, shall flourish in the courts of the house of our God.\end{Resp}
\begin{Resp}\Rsign Yea, he shall flourish in the presence of the Lord for ever.\end{Resp}
\switchcolumn*

\LA
\LessonHead{Lectio III}
\switchcolumn
\EN
\LessonHead{Lesson III}
\switchcolumn*

\LA
\Cite{Sir 39:11-14}
\switchcolumn
\EN
\Cite{Sir 39:11-14}
\switchcolumn*

\LA
\noindent Ipse palam fáciet disciplínam doctrínæ suæ, et in lege testaménti Dómini gloriábitur. Collaudábunt multi sapiéntiam ejus, et usque in sǽculum non delébitur. Non recédet memória ejus, et nomen ejus requirétur a generatióne in generatiónem. Sapiéntiam ejus enarrábunt gentes, et laudem ejus enuntiábit Ecclésia.\par
\switchcolumn
\EN
\noindent He shall show forth the discipline he hath learned, and shall glory in the law of the covenant of the Lord. Many shall praise his wisdom, and it shall never be forgotten. The memory of him shall not depart away, and his name shall be in request from generation to generation. Nations shall declare his wisdom, and the church shall show forth his praise.\par
\switchcolumn*

\LA
\begin{Resp}\Rsign Iste cognóvit justítiam, et vidit mirabília magna, et exorávit Altíssimum: \Star{} Et invéntus est in número Sanctórum.\end{Resp}
\begin{Resp}\Vsign Iste est qui contémpsit vitam mundi, et pervénit ad cæléstia regna.\end{Resp}
\begin{Resp}\Rsign Et invéntus est in número Sanctórum.\end{Resp}
\begin{Resp}\Vsign Glória Patri, et Fílio, \Star{} et Spirítui Sancto.\end{Resp}
\begin{Resp}\Rsign Et invéntus est in número Sanctórum.\end{Resp}
\switchcolumn
\EN
\begin{Resp}\Rsign This is he which knew righteousness, and saw great wonders, and made his prayer unto the Most High; \Star{} And he is numbered among the Saints.\end{Resp}
\begin{Resp}\Vsign This is he which loved not his life in this world, and is come unto an everlasting kingdom.\end{Resp}
\begin{Resp}\Rsign And he is numbered among the Saints.\end{Resp}
\begin{Resp}\Vsign Glory be to the Father, and to the Son, \Star{} and to the Holy Ghost.\end{Resp}
\begin{Resp}\Rsign And he is numbered among the Saints.\end{Resp}
\switchcolumn*

\LA
\LessonHead{Lectio VII}
\switchcolumn
\EN
\LessonHead{Lesson VII}
\switchcolumn*

\LA
\LessonTitle{Léctio sancti Evangélii secúndum Matthǽum}
\switchcolumn
\EN
\LessonTitle{From the Holy Gospel according to Matthew}
\switchcolumn*

\LA
\Cite{Matt 5:13-19}
\switchcolumn
\EN
\Cite{Matt 5:13-19}
\switchcolumn*

\LA
\noindent In illo témpore: Dixit Jesus discípulis suis: Vos estis sal terræ. Quod si sal evanúerit, in quo saliétur? Et réliqua.\par
\switchcolumn
\EN
\noindent At that time: Jesus said unto his disciples: Ye are the salt of the earth: But if the salt have lost his savour, wherewith shall it be salted? And so on, and that which followeth.\par
\switchcolumn*

\LA
\LessonTitle{Homilía sancti Augustíni Epíscopi}
\switchcolumn
\EN
\LessonTitle{A Homily by St. Augustine the Bishop}
\switchcolumn*

\LA
\Source{Lib. 1 de Sermone Domini in monte, cap. 6}
\switchcolumn
\EN
\Source{Lib. 1 de Sermóne Dómini in monte, cap. 6}
\switchcolumn*

\LA
\noindent Osténdit Dóminus fátuos esse judicándos, qui temporálium bonórum vel cópiam sectántes vel inópiam metuéntes, amíttunt ætérna, quæ nec dari possunt ab homínibus nec auférri. Itaque, si sal infatuátum fúerit, in quo saliétur? Id est, si vos, per quos condiéndi sunt quodámmodo pópuli, metu persecutiónum temporálium amiséritis regna cælórum; qui erunt hómines, per quos a vobis error auferátur, cum vos elégerit Deus, per quos errórem áuferat ceterórum?\par
\switchcolumn
\EN
\noindent The Lord would have us understand how that men do lose their power of savouring others with righteousness when they are willing to place their eternal welfare in jeopardy for the sake of any témporal advantage, like as attainment of ease or luxury, or escape from suffering or toil. For that which is eternal, unlike things of this world, can neither be bestowed by men, nor by them taken away. Hence, when he asketh: If the salt have lost his savour, wherewith shall it be salted? he would have us understand the question to be: If ye, by whom mankind is preserved from corruption, be willing to lose the kingdom of heaven so as to escape trials or persecutions in this world, who is there to preserve you from corruption, seeing ye are they that God hath chosen to preserve all others from corruption?\par
\switchcolumn*

\LA
\begin{Resp}\Rsign Iste est qui ante Deum magnas virtútes operátus est, et de omni corde suo laudávit Dóminum: \Star{} Ipse intercédat pro peccátis ómnium populórum.\end{Resp}
\begin{Resp}\Vsign Ecce homo sine queréla, verus Dei cultor, ábstinens se ab omni ópere malo, et pérmanens in innocéntia sua.\end{Resp}
\begin{Resp}\Rsign Ipse intercédat pro peccátis ómnium populórum.\end{Resp}
\switchcolumn
\EN
\begin{Resp}\Rsign This is he which wrought great wonders before God, and praised the Lord with all his heart. \Star{} May he pray for all people, that their sins may be forgiven unto them.\end{Resp}
\begin{Resp}\Vsign Behold a man without blame, a worshipper of God in truth, keeping himself clean from every evil work, and abiding still in his innocency.\end{Resp}
\begin{Resp}\Rsign May he pray for all people, that their sins may be forgiven unto them.\end{Resp}
\switchcolumn*

\LA
\LessonHead{Lectio VIII}
\switchcolumn
\EN
\LessonHead{Lesson VIII}
\switchcolumn*

\LA
\noindent Ergo ad níhilum valet sal infatuátum, nisi ut mittátur foras et calcétur ab homínibus. Non ítaque calcátur ab homínibus qui pátitur persecutiónem; sed qui, persecutiónem timéndo, infatuátur. Calcári enim non potest nisi inférior; sed inférior non est, qui, quamvis córpore multa in terra sustíneat, corde tamen fixus in cælo est.\par
\switchcolumn
\EN
\noindent Those that should be the salt of the earth, but have lost their savour, are thenceforth good for nothing, but to be cast out, and to be trodden under foot of men. But no one that suffereth persecution is truly said to be trodden under foot of men. Rather, that one is truly trodden under foot of men who through fear of persecution hath lost the savour of righteousness. For no one can be trodden upon, unless he be beneath him which treadeth upon him. And certainly no one who hath his heart in heaven, no matter how grievously he doth suffer in his body on earth, is rightly said to be beneath anyone who misuseth him.\par
\switchcolumn*

\LA
\begin{Resp}\Rsign In médio Ecclésiæ apéruit os ejus, \Star{} Et implévit eum Dóminus spíritu sapiéntiæ et intelléctus.\end{Resp}
\begin{Resp}\Vsign Jucunditátem et exsultatiónem thesaurizávit super eum.\end{Resp}
\begin{Resp}\Rsign Et implévit eum Dóminus spíritu sapiéntiæ et intelléctus.\end{Resp}
\begin{Resp}\Vsign Glória Patri, et Fílio, \Star{} et Spirítui Sancto.\end{Resp}
\begin{Resp}\Rsign Et implévit eum Dóminus spíritu sapiéntiæ et intelléctus.\end{Resp}
\switchcolumn
\EN
\begin{Resp}\Rsign In the midst of the congregation did the Lord open his mouth. \Star{} And filled him with the spirit of wisdom and understanding.\end{Resp}
\begin{Resp}\Vsign He made him rich with joy and gladness.\end{Resp}
\begin{Resp}\Rsign And filled him with the spirit of wisdom and understanding.\end{Resp}
\begin{Resp}\Vsign Glory be to the Father, and to the Son, \Star{} and to the Holy Ghost.\end{Resp}
\begin{Resp}\Rsign And filled him with the spirit of wisdom and understanding.\end{Resp}
\switchcolumn*

\LA
\LessonHead{Lectio IX}
\switchcolumn
\EN
\LessonHead{Lesson IX}
\switchcolumn*

\LA
\noindent Vos estis lumen mundi. Quómodo dixit supérius sal terræ, sic nunc dicit lumen mundi. Nam, neque supérius ista terra accipiénda est, quam pédibus corpóreis calcámus; sed hómines, qui in terra hábitant, vel étiam peccatóres, quorum condiéndis et exstinguéndis putóribus apostólicum salem Dóminus misit. Et hic mundum non cælum et terram, sed hómines qui sunt in mundo vel díligunt mundum, opórtet intélligi; quibus illuminándis Apóstoli missi sunt. Non potest cívitas abscóndi super montem pósita; id est, fundáta super insígnem magnámque justítiam, quam signíficat étiam ipse mons, in quo dísputat Dóminus.\par
\switchcolumn
\EN
\noindent Ye are the light of the world. And we are to understand the word World in the same sense as the word Earth when he spoke above of the salt of the earth, that is, not that earth whereupon we walk with our bodily feet, but the men which dwell upon the earth; in other words, sinners, for the sweetening and correction of whose corruption, the Lord hath sent his Apostles, as it were, as so much salt. And so by the world we are to understand, not the heaven and the earth, but the men who are in the world and love the world, for the enlightening of whom the Apostles have been sent. A city that is set on a hill cannot be hid: that is, what is founded upon the heights of righteousness, whereof the mountain upon which the Lord gave this discourse was itself a figure, is magnificent in the eyes of all men.\par
\switchcolumn*

\LA
\TeDeum{Te Deum}
\switchcolumn
\EN
\TeDeum{Te Deum}
\switchcolumn*

\end{paracol}

\Office{Commune Virginum tantum}{Common of a Virgin not a Martyr}{}
\addcontentsline{toc}{subsection}{Commune Virginum tantum\enspace\textit{Common of a Virgin not a Martyr}}
\begin{paracol}{2}
\LA
\LessonHead{Lectio VII}
\switchcolumn
\EN
\LessonHead{Lesson VII}
\switchcolumn*

\LA
\LessonTitle{Léctio sancti Evangélii secúndum Matthǽum}
\switchcolumn
\EN
\LessonTitle{From the Holy Gospel according to St. Matthew}
\switchcolumn*

\LA
\Cite{Matt 25:1-13}
\switchcolumn
\EN
\Cite{Matt 25:1}
\switchcolumn*

\LA
\noindent In illo témpore: Dixit Jesus discípulis suis parábolam hanc: Símile erit regnum cælórum decem virgínibus, quæ, accipiéntes lámpades suas, exiérunt óbviam sponso et sponsæ. Et réliqua.\par
\switchcolumn
\EN
\noindent At that time, Jesus said to His disciples: The Kingdom of heaven shall be likened unto ten virgins, which took their lamps, and went forth to meet the Bridegroom and the Bride. And so on.\par
\switchcolumn*

\LA
\LessonTitle{Homilía sancti Gregórii Papæ}
\switchcolumn
\EN
\LessonTitle{Homily by Pope St. Gregory the Great}
\switchcolumn*

\LA
\Source{Homilía 12 in Evangelia}
\switchcolumn
\EN
\Source{12th on the Gospels}
\switchcolumn*

\LA
\noindent Sæpe vos, fratres caríssimi, admóneo prava ópera fúgere, mundi hujus inquinaménta devitáre, sed hodiérna sancti Evangélii lectióne compéllor dícere, ut, et bona, quæ ágitis, cum magna cautéla teneátis; ne per hoc, quod a vobis rectum géritur, favor aut grátia humána requirátur; ne appetítus laudis subrépat, et, quod foris osténditur, intus a mercéde vacuétur. Ecce enim Redemptóris voce decem vírgines, et omnes dicúntur vírgines, et tamen intra beatitúdinis jánuam non omnes sunt recéptæ; quia eárum quædam, dum de virginitáte sua glóriam foris éxpetunt, in vasis suis óleum habére noluérunt.\par
\switchcolumn
\EN
\noindent Dearly beloved brethren oftentimes do I warn you to fly corrupt conversation, and to keep yourselves unspotted from the world. But the portion which is this day read from the Holy Gospel doth oblige me to say that even to these good things which ye do, ye must needs take all careful heed. Look ye well to it, that, when ye work righteousness, ye do it not as seeking the praise and admiration of men, for if the lust of praise do once creep in, that which seemeth so fair without, loseth its reward within. Behold how the Redeemer speaketh of these ten virgins. He calleth them all virgins, yet entered not all of them into the door of blessedness, for there were some of them who sought outwardly the honour of virginity, but would take no oil within their vessels with their lamps.\par
\switchcolumn*

\LA
\begin{Resp}\Rsign Hæc est Virgo sápiens, quam Dóminus vigilántem invénit, quæ accéptis lampádibus sumpsit secum óleum: \Star{} Et veniénte Dómino, introívit cum eo ad núptias.\end{Resp}
\begin{Resp}\Vsign Média nocte clamor factus est: Ecce sponsus venit, exíte óbviam ei.\end{Resp}
\begin{Resp}\Rsign Et veniénte Dómino, introívit cum eo ad núptias.\end{Resp}
\switchcolumn
\EN
\begin{Resp}\Rsign This is one of those wise virgins, whom the Lord found watching, for when she took her lamp, she took oil with her. \Star{} And when the Lord came, she went in with him to the marriage.\end{Resp}
\begin{Resp}\Vsign At midnight there was a cry made: Behold, the Bridegroom cometh, go ye out to meet him.\end{Resp}
\begin{Resp}\Rsign And when the Lord came, she went in with him to the marriage.\end{Resp}
\switchcolumn*

\LA
\LessonHead{Lectio VIII}
\switchcolumn
\EN
\LessonHead{Lesson VIII}
\switchcolumn*

\LA
\noindent Sed prius quæréndum nobis est quid sit regnum cælórum, aut cur decem virgínibus comparétur, quæ étiam vírgines prudéntes et fátuæ dicántur. Dum enim cælórum regnum constat quia reprobórum nullus ingréditur, étiam fátuis virgínibus cur símile esse perhibétur? Sed sciéndum nobis est quod sæpe in sacro elóquio regnum cælórum præséntis témporis Ecclésia dícitur. De quo álio in loco Dóminus dicit: Mittet Fílius hóminis Ángelos suos, et cólligent de regno ejus ómnia scándala. Neque enim in illo regno beatitúdinis, in quo pax summa est, inveníri scándala póterunt, quæ colligántur.\par
\switchcolumn
\EN
\noindent First of all, it is for us to ask What is the kingdom of Heaven? And wherefore shall the same be likened unto ten virgins, whereof, albeit five were wise, yet five were foolish For if the kingdom of heaven be such that there shall in no wise enter into it anything that defileth, neither whatsoever worketh abomination, or maketh a lie, (Apoc. xxi. 27,) how can it be like unto five virgins which were foolish? But we must know that, in the word of God, the kingdom of heaven doth oftentimes signify the Church as she now is, touching the which the Lord saith in another place “The Son of Man shall send forth His Angels, and they shall gather out of His kingdom all things that offend.” (Matth. xiii. 41.) In that kingdom of Blessedness, wherein peace shall have her perfect reign, there shall be nothing found that offendeth for the angels to gather out.\par
\switchcolumn*

\LA
\begin{Resp}\Rsign Média nocte clamor factus est: \Star{} Ecce sponsus venit, exíte óbviam ei.\end{Resp}
\begin{Resp}\Vsign Prudéntes vírgines, aptáte vestras lámpades.\end{Resp}
\begin{Resp}\Rsign Ecce sponsus venit, exíte óbviam ei.\end{Resp}
\begin{Resp}\Vsign Glória Patri, et Fílio, \Star{} et Spirítui Sancto.\end{Resp}
\begin{Resp}\Rsign Ecce sponsus venit, exíte óbviam ei.\end{Resp}
\switchcolumn
\EN
\begin{Resp}\Rsign At midnight there was a cry made: \Star{} Behold, the Bridegroom cometh, go ye out to meet him.\end{Resp}
\begin{Resp}\Vsign Trim your lamps, O ye wise virgins.\end{Resp}
\begin{Resp}\Rsign Behold, the Bridegroom cometh, go ye out to meet him.\end{Resp}
\begin{Resp}\Vsign Glory be to the Father, and to the Son, \Star{} and to the Holy Ghost.\end{Resp}
\begin{Resp}\Rsign Behold, the Bridegroom cometh, go ye out to meet him.\end{Resp}
\switchcolumn*

\LA
\LessonHead{Lectio IX}
\switchcolumn
\EN
\LessonHead{Lesson IX}
\switchcolumn*

\LA
\noindent In quinque autem córporis sénsibus unusquísque subsístit; geminátus autem quinárius denárium pérficit. Et, quia ex utróque sexu fidélium multitúdo collígitur, sancta Ecclésia decem virgínibus símilis esse denuntiátur. In qua quia mali cum bonis et réprobi cum eléctis admíxti sunt, recte símilis virgínibus prudéntibus et fátuis esse perhibétur. Sunt namque pleríque continéntes, qui ab appetítu se exterióri custódiunt et spe ad interióra rapiúntur, carnem mácerant, et toto desidério ad supérnam pátriam anhélant, ætérna prǽmia éxpetunt, pro labóribus suis recípere laudes humánas nolunt. Hi nimírum glóriam suam non in ore hóminum ponunt, sed intra consciéntiam cóntegunt. Et sunt pleríque, qui corpus per abstinéntiam afflígunt, sed de ipsa sua abstinéntia humános favóres éxpetunt.\par
\switchcolumn
\EN
\noindent The body of every man doth consist of five senses, and five being doubled, is ten. Forasmuch, therefore, as the whole body of the faithful doth consist of two sexes, the Holy Church is likened unto ten virgins. And forasmuch as in the Church the good are for the present mingled with the bad, and the reprobate with the elect, it is rightly said that, of the ten virgins, five are wise and five are foolish. There are many who have self-control, which do keep themselves from lusting after things outward, whose hope beareth them to things inward, who chastise the flesh, who long with intense home-sickness for their Fatherland which is in heaven, who seek an eternal reward, and who will not to receive for their labours the praise of men. These are they who reckon their glory, not in the mouths of men, but in the testimony of their own conscience. And many there be likewise who afflict the body by self-control, and yet who seek for their self-control applause from men.\par
\switchcolumn*

\LA
\TeDeum{Te Deum}
\switchcolumn
\EN
\TeDeum{Te Deum}
\switchcolumn*

\end{paracol}

\Office{Commune non Virginum non Martyrum}{Common of Widows}{}
\addcontentsline{toc}{subsection}{Commune non Virginum non Martyrum\enspace\textit{Common of Widows}}
\begin{paracol}{2}
\LA
\LessonHead{Lectio I}
\switchcolumn
\EN
\LessonHead{Lesson I}
\switchcolumn*

\LA
\LessonTitle{De Parábolis Salomónis}
\switchcolumn
\EN
\LessonTitle{Lesson from the book of Proverbs}
\switchcolumn*

\LA
\Cite{Prov 31:10-17}
\switchcolumn
\EN
\Cite{Prov 31:10-17}
\switchcolumn*

\LA
\noindent Mulíerem fortem quis invéniet? procul et de últimis fínibus prétium ejus. Confídit in ea cor viri sui, et spóliis non indigébit. Reddet ei bonum, et non malum, ómnibus diébus vitæ suæ. Quæsívit lanam et linum, et operáta est consília mánuum suárum. Facta est quasi navis institóris, de longe portans panem suum. Et de nocte surréxit, dedítque prædam domésticis suis, et cibária ancíllis suis. Considerávit agrum, et emit eum; de fructu mánuum suárum plantávit víneam. Accínxit fortitúdine lumbos suos, et roborávit bráchium suum.\par
\switchcolumn
\EN
\noindent Who shall find a valiant woman? Far and from the uttermost coasts is the price of her. The heart of her husband trusteth in her, and he shall have no need of spoils. She will render him good, and not evil, all the days of her life. She hath sought wool and flax, and hath wrought by the counsel of her hands. She is like the merchant's ship, she bringeth her bread from afar. And she hath risen in the night, and given a prey to her household, and victuals to her maidens. She hath considered a field, and bought it: with the fruit of her hands she hath planted a vineyard. She hath girded her loins with strength, and hath strengthened her arm\par
\switchcolumn*

\LA
\begin{Resp}\Rsign Veni, elécta mea, et ponam in te thronum meum \Star{} Quia concupívit Rex spéciem tuam.\end{Resp}
\begin{Resp}\Vsign Spécie tua et pulchritúdine tua inténde, próspere procéde, et regna.\end{Resp}
\begin{Resp}\Rsign Quia concupívit Rex spéciem tuam.\end{Resp}
\switchcolumn
\EN
\begin{Resp}\Rsign Come, O My chosen one, and I will establish My throne in thee; \Star{} For the King hath greatly desired thy beauty.\end{Resp}
\begin{Resp}\Vsign In thy comeliness and thy beauty, go forward, fare prosperously, and reign.\end{Resp}
\begin{Resp}\Rsign For the King hath greatly desired thy beauty.\end{Resp}
\switchcolumn*

\LA
\LessonHead{Lectio II}
\switchcolumn
\EN
\LessonHead{Lesson II}
\switchcolumn*

\LA
\Cite{Prov 31:18-24}
\switchcolumn
\EN
\Cite{Prov 31:18-24}
\switchcolumn*

\LA
\noindent Gustávit, et vidit quia bona est negotiátio ejus; non extinguétur in nocte lucérna ejus. Manum suam misit ad fórtia, et dígiti ejus apprehendérunt fusum. Manum suam apéruit ínopi, et palmas suas exténdit ad páuperem. Non timébit dómui suæ a frigóribus nivis; omnes enim doméstici ejus vestíti sunt duplícibus. Stragulátam vestem fecit sibi; byssus et púrpura induméntum ejus. Nóbilis in portis vir ejus, quando séderit cum senatóribus terræ. Síndonem fecit, et véndidit, et cíngulum trádidit Chananǽo.\par
\switchcolumn
\EN
\noindent She hath tasted and seen that her traffic is good: her lamp shall not be put out in the night. She hath put out her hand to strong things, and her fingers have taken hold of the spindle. She hath opened her hand to the needy, and stretched out her hands to the poor. She shall not fear for her house in the cold of snow: for all her domestics are clothed with double garments. She hath made for herself clothing of tapestry: fine linen, and purple is her covering. Her husband is honourable in the gates, when he sitteth among the senators of the land. She made fine linen, and sold it, and delivered a girdle to the Chanaanite.\par
\switchcolumn*

\LA
\begin{Resp}\Rsign Diffúsa est grátia in lábiis tuis, \Star{} Proptérea benedíxit te Deus in ætérnum.\end{Resp}
\begin{Resp}\Vsign Spécie tua et pulchritúdine tua inténde, próspere procéde, et regna.\end{Resp}
\begin{Resp}\Rsign Proptérea benedíxit te Deus in ætérnum.\end{Resp}
\switchcolumn
\EN
\begin{Resp}\Rsign Grace is poured into thy lips, therefore; \Star{} God hath blessed thee for ever.\end{Resp}
\begin{Resp}\Vsign In thy comeliness and thy beauty, go forward, fare prosperously, and reign.\end{Resp}
\begin{Resp}\Rsign God hath blessed thee for ever.\end{Resp}
\switchcolumn*

\LA
\LessonHead{Lectio III}
\switchcolumn
\EN
\LessonHead{Lesson III}
\switchcolumn*

\LA
\Cite{Prov 31:25-31}
\switchcolumn
\EN
\Cite{Prov 31:25-31}
\switchcolumn*

\LA
\noindent Fortitúdo et decor induméntum ejus, et ridébit in die novíssimo. Os suum apéruit sapiéntiæ, et lex cleméntiæ in lingua ejus. Considerávit sémitas domus suæ, et panem otiósa non comédit. Surrexérunt fílii ejus, et beatíssimam prædicavérunt; vir ejus, et laudávit eam. Multæ fíliæ congregavérunt divítias; tu supergréssa es univérsas. Fallax grátia, et vana est pulchritúdo: múlier timens Dóminum, ipsa laudábitur. Date ei de fructu mánuum suárum, et laudent eam in portis ópera ejus.\par
\switchcolumn
\EN
\noindent Strength and beauty are her clothing, and she shall laugh in the latter day. She hath opened her mouth to wisdom, and the law of clemency is on her tongue. She hath looked well to the paths of her house, and hath not eaten her bread idle. Her children rose up, and called her blessed: her husband, and he praised her. Many daughters have gathered together riches: thou hast surpassed them all. Favour is deceitful, and beauty is vain: the woman that feareth the Lord, she shall be praised. Give her of the fruit of her hands: and let her works praise her in the gates.\par
\switchcolumn*

\LA
\begin{Resp}\Rsign Spécie tua et pulchritúdine tua \Star{} Inténde, próspere procéde, et regna.\end{Resp}
\begin{Resp}\Vsign Diffúsa est grátia in lábiis tuis, proptérea benedíxit te Deus in ætérnum.\end{Resp}
\begin{Resp}\Rsign Inténde, próspere procéde, et regna.\end{Resp}
\begin{Resp}\Vsign Glória Patri, et Fílio, \Star{} et Spirítui Sancto.\end{Resp}
\begin{Resp}\Rsign Inténde, próspere procéde, et regna.\end{Resp}
\switchcolumn
\EN
\begin{Resp}\Rsign In thy comeliness and thy beauty; \Star{} Go forward, fare prosperously, and reign.\end{Resp}
\begin{Resp}\Vsign Grace is poured into thy lips, therefore God hath blessed thee for ever.\end{Resp}
\begin{Resp}\Rsign Go forward, fare prosperously, and reign.\end{Resp}
\begin{Resp}\Vsign Glory be to the Father, and to the Son, \Star{} and to the Holy Ghost.\end{Resp}
\begin{Resp}\Rsign Go forward, fare prosperously, and reign.\end{Resp}
\switchcolumn*

\LA
\LessonHead{Lectio VII}
\switchcolumn
\EN
\LessonHead{Lesson VII}
\switchcolumn*

\LA
\LessonTitle{Léctio sancti Evangélii secúndum Matthǽum}
\switchcolumn
\EN
\LessonTitle{From the Holy Gospel according to Matthew}
\switchcolumn*

\LA
\Cite{Matt 13:44-52}
\switchcolumn
\EN
\Cite{Matt 13:44}
\switchcolumn*

\LA
\noindent In illo témpore: Dixit Jesus discípulis suis parábolam hanc: Simile est regnum cælórum thesáuro abscóndito in agro. Et réliqua.\par
\switchcolumn
\EN
\noindent At that time Jesus spake unto His disciples this parable The kingdom of heaven is like unto treasure hid in a field. And so on.\par
\switchcolumn*

\LA
\LessonTitle{Homilía sancti Gregórii Papæ}
\switchcolumn
\EN
\LessonTitle{Homily by Pope St Gregory the Great}
\switchcolumn*

\LA
\Source{Homilia 11 in Evangelia}
\switchcolumn
\EN
\Source{11th on the Gospels}
\switchcolumn*

\LA
\noindent Cælórum regnum, fratres caríssimi, idcírco terrénis rebus símile dícitur, ut, ex his quæ ánimus novit, surgat ad incógnita, quæ non novit: quátenus exémplo visibílium se ad invisibília rápiat, et, per ea quæ usu dídicit quasi confricátus incaléscat; ut per hoc, quod scit notum dilígere, discat et incógnita amáre. Ecce enim cælórum regnum thesáuro abscóndito in agro comparátur; quem, qui invénit homo, abscóndit, et præ gáudio illíus vadit, et vendit univérsa quæ habet, et emit agrum illum.\par
\switchcolumn
\EN
\noindent Dearly beloved brethren, the kingdom of heaven is likened unto the things of earth, to the end that by the mean of things which we know, our mind may rise to the contemplation of the things which we know not by the example of things which are seen, may fix her gaze on things which are not seen by the touch of things which she useth, may be warmed towards the things which she useth not; by things which she knoweth and loveth, to love also the things which she knoweth not. For, behold, “the kingdom of heaven is likened unto treasure hid in a field, the which when a man hath found, he hideth, and, for joy thereof, goeth and selleth all that he hath and buyeth that field.”\par
\switchcolumn*

\LA
\begin{Resp}\Rsign Os suum apéruit sapiéntiæ, et lex cleméntiæ in lingua ejus: considerávit sémitas domus suæ, \Star{} Et panem otiósa non comédit.\end{Resp}
\begin{Resp}\Vsign Gustávit et vidit quia bona est negotiátio ejus: non exstinguétur in nocte lucérna ejus.\end{Resp}
\begin{Resp}\Rsign Et panem otiósa non comédit.\end{Resp}
\switchcolumn
\EN
\begin{Resp}\Rsign She openeth her mouth with wisdom, and in her tongue is the law of kindness. She looketh well to the ways of her household \Star{} And eateth not the bread of idleness.\end{Resp}
\begin{Resp}\Vsign She tasteth and perceiveth that her merchandise is good. Her candle goeth not out by night.\end{Resp}
\begin{Resp}\Rsign And she eateth not the bread of idleness.\end{Resp}
\switchcolumn*

\LA
\LessonHead{Lectio VIII}
\switchcolumn
\EN
\LessonHead{Lesson VIII}
\switchcolumn*

\LA
\noindent Qua in re hoc quoque notándum est, quod invéntus thesáurus abscónditur, ut servétur: quia stúdium cæléstis desidérii a malígnis spirítibus custodíre non súfficit, qui hoc ab humánis láudibus non abscóndit. In præsénti étenim vita, quasi in via sumus, qua ad pátriam pérgimus. Malígni autem spíritus iter nostrum quasi quidam latrúnculi óbsident. Deprædári ergo desíderat, qui thesáurum públice portat in via. Hoc autem dico, non ut próximi ópera nostra bona non vídeant, cum scriptum sit: Vídeant ópera vestra bona, et gloríficent Patrem vestrum qui in cælis est; sed, ut per hoc quod ágimus, laudes extérius non quærámus. Sic autem sit opus in público, quátenus inténtio máneat in occúlto; ut, et de bono ópere próximis præbeámus exémplum, et tamen per intentiónem, qua Deo soli placére quǽrimus, semper optémus secrétum.\par
\switchcolumn
\EN
\noindent And herein we must remark that the treasure, when once it hath been found, is hidden to keep it safe. He whose intimate yearnings after God are not hidden from the praise of men, is open thereby to the attacks of evil spirits. In this life we are, as it were, journeying homewards on a road beset by evil spirits who are like highwaymen. He therefore inviteth robbery who carrieth his treasure ostentatiously. Doubtless our neighbour should be able to see our good works, as it is written: “Let your light so shine before men that they may see your good works, and glorify your Father which is in heaven.” But this is not to be understood to mean that we are to seek the praise of men by what we do. Rather, let us in such wise work in the open that the inner intention of devotion is not advertised. So we shall give an example to our neighbour, and yet keep hidden, except from the sight of God, our purpose of pleasing him.\par
\switchcolumn*

\LA
\begin{Resp}\Rsign Regnum mundi et omnem ornátum sǽculi contémpsi, propter amórem Dómini mei Jesu Christi: \Star{} Quem vidi, quem amávi, in quem crédidi, quem diléxi.\end{Resp}
\begin{Resp}\Vsign Eructávit cor meum verbum bonum: dico ego ópera mea Regi.\end{Resp}
\begin{Resp}\Rsign Quem vidi, quem amávi, in quem crédidi, quem diléxi.\end{Resp}
\begin{Resp}\Vsign Glória Patri, et Fílio, \Star{} et Spirítui Sancto.\end{Resp}
\begin{Resp}\Rsign Quem vidi, quem amávi, in quem crédidi, quem diléxi.\end{Resp}
\switchcolumn
\EN
\begin{Resp}\Rsign The kingdom of this world and all the beauty of life I have esteemed as nothing, for the excellency of the love of Jesus Christ my Lord \Star{} Whom, having seen, I loved Whom, having believed, I longed after.\end{Resp}
\begin{Resp}\Vsign My heart is overflowing with a good matter I speak of my works unto the King.\end{Resp}
\begin{Resp}\Rsign Whom, having seen, I loved Whom, having believed, I longed after.\end{Resp}
\begin{Resp}\Vsign Glory be to the Father, and to the Son, \Star{} and to the Holy Ghost.\end{Resp}
\begin{Resp}\Rsign Whom, having seen, I loved Whom, having believed, I longed after.\end{Resp}
\switchcolumn*

\LA
\LessonHead{Lectio IX}
\switchcolumn
\EN
\LessonHead{Lesson IX}
\switchcolumn*

\LA
\noindent Thesáurus autem cæléste est desidérium; ager vero, in quo thesáurus abscónditur, disciplína stúdii cæléstis. Quem profécto agrum, vénditis ómnibus, cómparat, qui, voluptátibus carnis renúntians, cuncta sua terréna desidéria per disciplínæ cæléstis custódiam calcat: ut nihil jam quod caro blandítur, líbeat; nihil, quod carnálem vitam trucídat, spíritus perhorréscat.\par
\switchcolumn
\EN
\noindent The treasure is the desire for heaven, the field wherein it is hidden is the earnest observance wherewith this desire is surrounded. Whosoever turneth his back upon the enjoyments of the flesh, and by earnest striving heavenward, putteth all earthly lusts under the feet of discipline, so that he smileth back no more when the flesh smileth at him, and shuddereth no more at anything that can only kill the body whosoever doth thus, hath sold all that he had, and bought that field.\par
\switchcolumn*

\LA
\TeDeum{Te Deum}
\switchcolumn
\EN
\TeDeum{Te Deum}
\switchcolumn*

\end{paracol}

\Office{Commune plurimarum non Virginum Martyrum}{Common of multiple widowed Martyrs}{}
\addcontentsline{toc}{subsection}{Commune plurimarum non Virginum Martyrum\enspace\textit{Common of multiple widowed Martyrs}}
\begin{paracol}{2}
\LA
\LessonHead{Lectio I}
\switchcolumn
\EN
\LessonHead{Lesson I}
\switchcolumn*

\LA
\noindent Confitébor tibi, Dómine, Rex, et collaudábo te Deum Salvatórem meum. Confitébor nómini tuo: quóniam adjútor et protéctor factus es mihi, et liberásti corpus meum a perditióne, a láqueo linguæ iníquæ et a lábiis operántium mendácium, et in conspéctu astántium factus es mihi adjútor. Et liberásti me secúndum multitúdinem misericórdiæ nóminis tui a rugiéntibus, præparátis ad escam: de mánibus quæréntium ánimam meam, et de portis tribulatiónum, quæ circumdedérunt me; a pressúra flammæ, quæ circúmdedit me, et in médio ignis non sum æstuáta; de altitúdine ventris ínferi, et a lingua coinquináta, et a verbo mendácii, a rege iníquo, et a lingua injústa.\par
\switchcolumn
\EN
\noindent A prayer of Jesus the son of Sirach. I will give glory to thee, O Lord, O King, and I will praise thee, O God my Saviour. I will give glory to thy name: for thou hast been a helper and protector to me. And hast preserved my body from destruction, from the snare of an unjust tongue, and from the lips of them that forge lies, and in the sight of them that stood by, thou hast been my helper. And thou hast delivered me, according to the multitude of the mercy of thy name, from them that did roar, prepared to devour. Out of the hands of them that sought my life, and from the gates of afflictions, which compassed me about: From the oppression of the flame which surrounded me, and in the midst of the fire I was not burnt. From the depth of the belly of hell, and from an unclean tongue, and from lying words, from an unjust king, and from a slanderous tongue:\par
\switchcolumn*

\LA
\begin{Resp}\Rsign Veni Sponsa Christi, áccipe corónam, quam tibi Dóminus præparávit in ætérnum; pro cujus amóre sánguinem tuum fudísti, \Star{} Et cum Ángelis in paradísum introísti.\end{Resp}
\begin{Resp}\Vsign Veni, elécta mea, et ponam in te thronum meum; quia concupívit Rex spéciem tuam.\end{Resp}
\begin{Resp}\Rsign Et cum Ángelis in paradísum introísti.\end{Resp}
\switchcolumn
\EN

\switchcolumn*

\LA
\LessonHead{Lectio II}
\switchcolumn
\EN
\LessonHead{Lesson II}
\switchcolumn*

\LA
\Cite{Sir 51:8-12}
\switchcolumn
\EN
\Cite{Sir 51:8-12}
\switchcolumn*

\LA
\noindent Laudábit usque ad mortem ánima mea Dóminum, et vita mea appropínquans erat in inférno deórsum. Circumdedérunt me úndique, et non erat qui adjuváret. Respíciens eram ad adjutórium hóminum, et non erat. Memoráta sum misericórdiæ tuæ, Dómine, et operatiónis tuæ, quæ a sǽculo sunt: quóniam éruis sustinéntes te, Dómine, et líberas eos de mánibus géntium.\par
\switchcolumn
\EN
\noindent My soul shall praise the Lord even to death. And my life was drawing near to hell beneath. They compassed me on every side, and there was no one that would help me. I looked for the succour of men, and there was none. I remembered thy mercy, O Lord, and thy works, which are from the beginning of the world. How thou deliverest them that wait for thee, O Lord, and savest them out of the hands of the nations.\par
\switchcolumn*

\LA
\begin{Resp}\Rsign Diffúsa est grátia in lábiis tuis, \Star{} Proptérea benedíxit te Deus in ætérnum.\end{Resp}
\begin{Resp}\Vsign Spécie tua et pulchritúdine tua inténde, próspere procéde, et regna.\end{Resp}
\begin{Resp}\Rsign Proptérea benedíxit te Deus in ætérnum.\end{Resp}
\switchcolumn
\EN
\begin{Resp}\Rsign Grace is poured into thy lips, therefore; \Star{} God hath blessed thee for ever.\end{Resp}
\begin{Resp}\Vsign In thy comeliness and thy beauty, go forward, fare prosperously, and reign.\end{Resp}
\begin{Resp}\Rsign God hath blessed thee for ever.\end{Resp}
\switchcolumn*

\LA
\LessonHead{Lectio III}
\switchcolumn
\EN
\LessonHead{Lesson III}
\switchcolumn*

\LA
\Cite{Sir 51:13-17}
\switchcolumn
\EN
\Cite{Sir 51:13-17}
\switchcolumn*

\LA
\noindent Exaltásti super terram habitatiónem meam, et pro morte defluénte deprecáta sum. Invocávi Dóminum, Patrem Dómini mei, ut non derelínquat me in die tribulatiónis meæ, et in témpore superbórum, sine adjutório. Laudábo nomen tuum assídue, et collaudábo illud in confessióne, et exaudíta est orátio mea. Et liberásti me de perditióne, et eripuísti me de témpore iníquo. Proptérea confitébor, et laudem dicam tibi, et benedícam nómini Dómini.\par
\switchcolumn
\EN
\noindent Thou hast exalted my dwelling place upon the earth and I have prayed for death to pass away. I called upon the Lord, the father of my Lord, that he would not leave me in the day of my trouble, and in the time of the proud without help. I will praise thy name continually, and will praise it with thanksgiving, and my prayer was heard. And thou hast saved me from destruction, and hast delivered me from the evil time. Therefore I will give thanks, and praise thee, and bless the name of the Lord.\par
\switchcolumn*

\LA
\begin{Resp}\Rsign Spécie tua et pulchritúdine tua \Star{} Inténde, próspere procéde, et regna.\end{Resp}
\begin{Resp}\Vsign Diffúsa est grátia in lábiis tuis, proptérea benedíxit te Deus in ætérnum.\end{Resp}
\begin{Resp}\Rsign Inténde, próspere procéde, et regna.\end{Resp}
\begin{Resp}\Vsign Glória Patri, et Fílio, \Star{} et Spirítui Sancto.\end{Resp}
\begin{Resp}\Rsign Inténde, próspere procéde, et regna.\end{Resp}
\switchcolumn
\EN
\begin{Resp}\Rsign In thy comeliness and thy beauty; \Star{} Go forward, fare prosperously, and reign.\end{Resp}
\begin{Resp}\Vsign Grace is poured into thy lips, therefore God hath blessed thee for ever.\end{Resp}
\begin{Resp}\Rsign Go forward, fare prosperously, and reign.\end{Resp}
\begin{Resp}\Vsign Glory be to the Father, and to the Son, \Star{} and to the Holy Ghost.\end{Resp}
\begin{Resp}\Rsign Go forward, fare prosperously, and reign.\end{Resp}
\switchcolumn*

\LA
\LessonHead{Lectio VII}
\switchcolumn
\EN
\LessonHead{Lesson VII}
\switchcolumn*

\LA
\LessonTitle{Léctio sancti Evangélii secúndum Matthǽum}
\switchcolumn
\EN
\LessonTitle{From the Holy Gospel according to Matthew}
\switchcolumn*

\LA
\Cite{Matt 13:44-52}
\switchcolumn
\EN
\Cite{Matt 13:44}
\switchcolumn*

\LA
\noindent In illo témpore: Dixit Jesus discípulis suis parábolam hanc: Simile est regnum cælórum thesáuro abscóndito in agro. Et réliqua.\par
\switchcolumn
\EN
\noindent At that time Jesus spake unto His disciples this parable The kingdom of heaven is like unto treasure hid in a field. And so on.\par
\switchcolumn*

\LA
\LessonTitle{Homilía sancti Gregórii Papæ}
\switchcolumn
\EN
\LessonTitle{Homily by Pope St Gregory the Great}
\switchcolumn*

\LA
\Source{Homilia 11 in Evangelia}
\switchcolumn
\EN
\Source{11th on the Gospels}
\switchcolumn*

\LA
\noindent Cælórum regnum, fratres caríssimi, idcírco terrénis rebus símile dícitur, ut, ex his quæ ánimus novit, surgat ad incógnita, quæ non novit: quátenus exémplo visibílium se ad invisibília rápiat, et, per ea quæ usu dídicit quasi confricátus incaléscat; ut per hoc, quod scit notum dilígere, discat et incógnita amáre. Ecce enim cælórum regnum thesáuro abscóndito in agro comparátur; quem, qui invénit homo, abscóndit, et præ gáudio illíus vadit, et vendit univérsa quæ habet, et emit agrum illum.\par
\switchcolumn
\EN
\noindent Dearly beloved brethren, the kingdom of heaven is likened unto the things of earth, to the end that by the mean of things which we know, our mind may rise to the contemplation of the things which we know not by the example of things which are seen, may fix her gaze on things which are not seen by the touch of things which she useth, may be warmed towards the things which she useth not; by things which she knoweth and loveth, to love also the things which she knoweth not. For, behold, “the kingdom of heaven is likened unto treasure hid in a field, the which when a man hath found, he hideth, and, for joy thereof, goeth and selleth all that he hath and buyeth that field.”\par
\switchcolumn*

\LA
\begin{Resp}\Rsign Os suum apéruit sapiéntiæ, et lex cleméntiæ in lingua ejus: considerávit sémitas domus suæ, \Star{} Et panem otiósa non comédit.\end{Resp}
\begin{Resp}\Vsign Gustávit et vidit quia bona est negotiátio ejus: non exstinguétur in nocte lucérna ejus.\end{Resp}
\begin{Resp}\Rsign Et panem otiósa non comédit.\end{Resp}
\switchcolumn
\EN
\begin{Resp}\Rsign She openeth her mouth with wisdom, and in her tongue is the law of kindness. She looketh well to the ways of her household \Star{} And eateth not the bread of idleness.\end{Resp}
\begin{Resp}\Vsign She tasteth and perceiveth that her merchandise is good. Her candle goeth not out by night.\end{Resp}
\begin{Resp}\Rsign And she eateth not the bread of idleness.\end{Resp}
\switchcolumn*

\LA
\LessonHead{Lectio VIII}
\switchcolumn
\EN
\LessonHead{Lesson VIII}
\switchcolumn*

\LA
\noindent Qua in re hoc quoque notándum est, quod invéntus thesáurus abscónditur, ut servétur: quia stúdium cæléstis desidérii a malígnis spirítibus custodíre non súfficit, qui hoc ab humánis láudibus non abscóndit. In præsénti étenim vita, quasi in via sumus, qua ad pátriam pérgimus. Malígni autem spíritus iter nostrum quasi quidam latrúnculi óbsident. Deprædári ergo desíderat, qui thesáurum públice portat in via. Hoc autem dico, non ut próximi ópera nostra bona non vídeant, cum scriptum sit: Vídeant ópera vestra bona, et gloríficent Patrem vestrum qui in cælis est; sed, ut per hoc quod ágimus, laudes extérius non quærámus. Sic autem sit opus in público, quátenus inténtio máneat in occúlto; ut, et de bono ópere próximis præbeámus exémplum, et tamen per intentiónem, qua Deo soli placére quǽrimus, semper optémus secrétum.\par
\switchcolumn
\EN
\noindent And herein we must remark that the treasure, when once it hath been found, is hidden to keep it safe. He whose intimate yearnings after God are not hidden from the praise of men, is open thereby to the attacks of evil spirits. In this life we are, as it were, journeying homewards on a road beset by evil spirits who are like highwaymen. He therefore inviteth robbery who carrieth his treasure ostentatiously. Doubtless our neighbour should be able to see our good works, as it is written: “Let your light so shine before men that they may see your good works, and glorify your Father which is in heaven.” But this is not to be understood to mean that we are to seek the praise of men by what we do. Rather, let us in such wise work in the open that the inner intention of devotion is not advertised. So we shall give an example to our neighbour, and yet keep hidden, except from the sight of God, our purpose of pleasing him.\par
\switchcolumn*

\LA
\begin{Resp}\Rsign Regnum mundi et omnem ornátum sǽculi contémpsi, propter amórem Dómini mei Jesu Christi: \Star{} Quem vidi, quem amávi, in quem crédidi, quem diléxi.\end{Resp}
\begin{Resp}\Vsign Eructávit cor meum verbum bonum: dico ego ópera mea Regi.\end{Resp}
\begin{Resp}\Rsign Quem vidi, quem amávi, in quem crédidi, quem diléxi.\end{Resp}
\begin{Resp}\Vsign Glória Patri, et Fílio, \Star{} et Spirítui Sancto.\end{Resp}
\begin{Resp}\Rsign Quem vidi, quem amávi, in quem crédidi, quem diléxi.\end{Resp}
\switchcolumn
\EN
\begin{Resp}\Rsign The kingdom of this world and all the beauty of life I have esteemed as nothing, for the excellency of the love of Jesus Christ my Lord \Star{} Whom, having seen, I loved Whom, having believed, I longed after.\end{Resp}
\begin{Resp}\Vsign My heart is overflowing with a good matter I speak of my works unto the King.\end{Resp}
\begin{Resp}\Rsign Whom, having seen, I loved Whom, having believed, I longed after.\end{Resp}
\begin{Resp}\Vsign Glory be to the Father, and to the Son, \Star{} and to the Holy Ghost.\end{Resp}
\begin{Resp}\Rsign Whom, having seen, I loved Whom, having believed, I longed after.\end{Resp}
\switchcolumn*

\end{paracol}

\Office{Sanctæ Mariæ Sabbato (tempore Paschali)}{Our Lady's Saturday (in Paschaltide)}{Simplex}
\addcontentsline{toc}{subsection}{Sanctæ Mariæ Sabbato (tempore Paschali)\enspace\textit{Our Lady's Saturday (in Paschaltide)}}
\begin{paracol}{2}
\LA
\RefLine{Lectiones I–II: de Scriptura occurrente.}
\switchcolumn
\EN
\RefLine{Lessons I–II: from the Scripture of the day.}
\switchcolumn*

\LA
\LessonHead{Lectio III (mense Aprili)}
\switchcolumn
\EN
\LessonHead{Lesson III (in April)}
\switchcolumn*

\LA
\LessonTitle{De Expositióne sancti Hierónymi Presbýteri in Ezechiélem Prophétam}
\switchcolumn
\EN
\LessonTitle{From the Commentary of St. Jerome, Priest, on the prophet Ezechiel}
\switchcolumn*

\LA
\Source{Lib. 13 in cap. 44}
\switchcolumn
\EN
\Source{Book 13 on ch. 44}
\switchcolumn*

\LA
\noindent Porta hæc clausa erit, et non aperiétur. Pulchre quidam portam clausam, per quam solus Dóminus Deus Israël ingréditur, et dux cui porta clausa est, Maríam Vírginem intéllegunt, quæ et ante partum et post partum virgo permánsit. Etenim témpore, quo Angelus loquebátur: Spíritus Sanctus véniet super te, et virtus Altíssimi obumbrábit te: quod autem nascétur ex te Sanctum, vocábitur Fílius Dei; et quando natus est, virgo permánsit ætérna; ad confundéndos eos, qui arbitrántur eam post nativitátem Salvatóris habuísse de Joseph fílios, ex occasióne fratrum ejus, qui vocántur in Evangélio.\par
\switchcolumn
\EN
\noindent “This gate is to remain closed it is not to be opened”. In a beautiful figure, some persons understand the closed which only the Lord, the God of Israel enters the leader on whose account it has been - as a type of the Virgin Mary who remained a virgin both before and after she gave birth - she remained a virgin, while the Angel was speaking to her: “The Holy Spirit shall come upon you and the power of the Most High shall overshadow you; and therefore the Holy One is be born shall be called the Son of God.” And when He was born, she remained a virgin, her perpetual virginity confounding those who think, because of the mention in the Gospel of the Saviour's brethren,” that after His birth she had children by Joseph.\par
\switchcolumn*

\LA
\TeDeum{Te Deum}
\switchcolumn
\EN
\TeDeum{Te Deum}
\switchcolumn*

\LA
\LessonHead{Lectio III (mense Majo)}
\switchcolumn
\EN
\LessonHead{Lesson III (in May)}
\switchcolumn*

\LA
\LessonTitle{Ex Tractátu S. Augustíni Epíscopi de Sýmbolo ad Catechúmenos}
\switchcolumn
\EN
\LessonTitle{From the Treatise of St. Augustine, Bishop, on the Creed, to the Catechumens}
\switchcolumn*

\LA
\Source{Lib. 3 cap. 4 in fine}
\switchcolumn
\EN
\Source{Book 3, ch. 4, at the end}
\switchcolumn*

\LA
\noindent Per féminam mors, per féminam vita: per Hevam intéritus, per Maríam salus. Illa, corrúpta, secúta est seductórem: hæc, íntegra, péperit Salvatórem. Illa póculum a serpénte propinátum libénter accépit et viro trádidit, ex quo simul mereréntur occídi: hæc, grátia cælésti désuper infúsa, vitam prótulit, per quam caro mórtua possit resuscitári. Quis est qui hæc operátus est, nisi Vírginis Fílius et vírginum Sponsus, qui áttulit Matri fecunditátem, sed non ábstulit integritátem?\par
\switchcolumn
\EN
\noindent Through a woman came death; through a woman, life: through Eve, ruin; through Mary, salvation. The former, corrupted, followed the deceiver; the latter, uncorrupted, gave birth to the Saviour. Eve willingly accepted the drink offered by the serpent and handed it on to her husband; and by their action both deserved the penalty of death. Mary, filled with heavenly grace from above, brought forth life, by which mankind, already dead, can be revived. Who has worked this miracle, if not the Son of the Virgin and the Spouse of virgin who brought fruitfulness to His mother without taking away her integrity?\par
\switchcolumn*

\LA
\TeDeum{Te Deum}
\switchcolumn
\EN
\TeDeum{Te Deum}
\switchcolumn*

\end{paracol}
\backmatter
\IndexOfOffices
\end{document}
